% Generated mechanically from frozen input inputs/p02/ja/tex/more-algebra.tex.
% Do not edit this generated file; edit the bound locale source instead.

% OK, start here.
%

\setcounter{chapter}{14}
\chapter{代数学のさらなる話題}



\phantomsection
\label{more-algebra-section-phantom}




\section{序論}
\label{more-algebra-section-introduction}

\noindent
本章では、可換代数に関する結果のうち、最初の可換代数の章にあるものより
高度なものをいくつか証明する。Algebra, Section
\ref{algebra-section-introduction} を参照せよ。
参考文献として \cite{MatCA} がある。





\section{読者への助言}
\label{more-algebra-section-advice}

\noindent
可換代数の章以上に、本章の各節はそれぞれ独立している。
Section \ref{more-algebra-section-derived-modules} 以降では、環上の加群の（非有界）導来圏と、
それに伴うすべての機構を自由に用いる。




\section{安定自由加群}
\label{more-algebra-section-stably-free}

\noindent
以下が一般に受け入れられている定義と思われる。

\begin{definition}
\label{more-algebra-definition-stably-free}
$R$ を環とする。
\begin{enumerate}
\item $R$ 上の二つの加群 $M$, $N$ が {it 安定同型} であるとは、ある
$n, m \geq 0$ が存在して、$R$-加群として
$M \oplus R^{\oplus m} \cong N \oplus R^{\oplus n}$
となることをいう。
\item 加群 $M$ が {it 安定自由} であるとは、自由加群と安定同型であることをいう。
\end{enumerate}
\end{definition}

\noindent
安定自由加群は射影加群であることに注意せよ。

\begin{lemma}
\label{more-algebra-lemma-exact-category-stably-free}
$R$ を環とする。有限生成射影 $R$-加群の短完全列
$0 \to P' \to P \to P'' \to 0$ をとる。この $3$ 個の加群のうち $2$ 個が
安定自由ならば、残る一つも安定自由である。
\end{lemma}

\begin{proof}
これらの加群は射影的なので、列は分裂する。従って同型
$P = P' \oplus P''$ を選べる。$P' \oplus R^{\oplus n}$ と
$P'' \oplus R^{\oplus m}$ が自由ならば、
$P \oplus R^{\oplus n + m}$ が自由であることが分かる。$P'$ と $P$ が
安定自由であると仮定し、$P \oplus R^{\oplus n}$ が自由であり、かつ
$P' \oplus R^{\oplus m}$ が自由であるとする。このとき
$$
P'' \oplus (P' \oplus R^{\oplus m}) \oplus R^{\oplus n} =
(P'' \oplus P') \oplus R^{\oplus m} \oplus R^{\oplus n} =
(P \oplus R^{\oplus n}) \oplus R^{\oplus m}
$$
は自由である。従って $P''$ は安定自由である。対称性により、残る場合も従う。
\end{proof}










\begin{lemma}
\label{more-algebra-lemma-lift-stably-free}
$R$ を環とし、$I \subset R$ をイデアルとする。$1 + I$ のすべての元が単元で
あると仮定する（言い換えれば、$I$ は $R$ の Jacobson 根基に含まれる）。
任意の有限生成安定自由 $R/I$-加群 $E$ に対し、$M/IM \cong E$ を満たす
有限生成安定自由 $R$-加群 $M$ が存在する。
\end{lemma}

\begin{proof}
整数 $n$, $m$ と同型
$E \oplus (R/I)^{\oplus n} \cong (R/I)^{\oplus m}$ を選ぶ。
射影 $(R/I)^{\oplus m} \to (R/I)^{\oplus n}$ および単射
$(R/I)^{\oplus n} \to (R/I)^{\oplus m}$ の持ち上げとなる $R$-線型写像
$\varphi : R^{\oplus m} \to R^{\oplus n}$ と
$\psi : R^{\oplus n} \to R^{\oplus m}$ を選ぶ。
このとき $\varphi \circ \psi : R^{\oplus n} \to R^{\oplus n}$ は $I$ を法として
恒等写像に帰着する。従って、この写像の行列式は $I$ に関する仮定により可逆である。
ゆえに $P = \Ker(\varphi)$ は安定自由であり、$E$ を持ち上げる。
\end{proof}

\begin{lemma}
\label{more-algebra-lemma-lift-projective}
$R$ を環とし、$I \subset R$ をイデアルとする。
$1 + I$ のすべての元が単元であると仮定する
（言い換えれば、$I$ は $R$ の Jacobson 根基に含まれる）。
$M$ を有限生成平坦 $R$-加群とし、$M/IM$ が射影 $R/I$-加群であるとする。
このとき $M$ は有限生成射影 $R$-加群である。
\end{lemma}

\begin{proof}
Algebra, Lemma \ref{algebra-lemma-finite-flat-local} により、すべての素イデアル
$\mathfrak p \subset R$ に対して $M_\mathfrak p$ は有限階数自由である。
Algebra, Lemma \ref{algebra-lemma-finite-projective} により、関数
$\rho_M : \mathfrak p \mapsto
\dim_{\kappa(\mathfrak p)} M \otimes_R \kappa(\mathfrak p)$
が $\Spec(R)$ 上局所定値であることを示せば十分である。$M/IM$ は有限生成射影なので、
これは $V(I) \subset \Spec(R)$ 上で成り立つ。$\Spec(R)$ のすべての閉点は
$V(I)$ に属し、素イデアル $\mathfrak p \subset \mathfrak q \subset R$ に対して
$\rho_M(\mathfrak p) = \rho_M(\mathfrak q)$ であるから、省略する初等的な位相空間の
議論によって結論を得る。
\end{proof}

\noindent
Lemma \ref{more-algebra-lemma-lift-stably-free} の持ち上げは、次の補題により同型を除いて一意である。

\begin{lemma}
\label{more-algebra-lemma-isomorphic-finite-projective-lifts}
$R$ を環とし、$I \subset R$ をイデアルとする。$1 + I$ のすべての元が単元であると
仮定する（言い換えれば、$I$ は $R$ の Jacobson 根基に含まれる）。$P$ と $P'$ を
有限生成射影 $R$-加群とする。このとき
\begin{enumerate}
\item $\varphi : P \to P'$ が同型
$\overline{\varphi} : P/IP \to P'/IP'$ を誘導する $R$-加群写像ならば、
$\varphi$ は同型である。
\item $P/IP \cong P'/IP'$ ならば、$P \cong P'$ である。
\end{enumerate}
\end{lemma}

\begin{proof}
(1) の証明。$P'$ は $R$-加群として射影的なので、写像
$P' \to P'/IP' \xrightarrow{\overline{\varphi}^{-1}} P/IP$ の持ち上げ
$\psi : P' \to P$ を選べる。中山の補題
（Algebra, Lemma \ref{algebra-lemma-NAK}）により、$\psi \circ \varphi$ と
$\varphi \circ \psi$ は全射である。従って、これらの写像は同型である
（Algebra, Lemma \ref{algebra-lemma-fun}）。ゆえに $\varphi$ は同型である。

\medskip\noindent
(2) の証明。同型 $P/IP \cong P'/IP'$ を選ぶ。$P$ は射影的なので、写像
$P \to P/IP \to P'/IP'$ の持ち上げ $\varphi : P \to P'$ を選べる。
(1) により $\varphi$ は同型である。
\end{proof}





\section{Artin--Rees 性質についての注意}
\label{more-algebra-section-artin-rees}

\noindent
この節の内容の一部は \cite{conrad-dejong} による。さらに参考文献を含む一般的な
議論は \cite[Section 1]{Eis} にある。

\medskip\noindent
$A$ を Noether 環とし、$I \subset A$ をイデアルとする。有限生成 $A$-加群の準同型
$f : M \to N$ が与えられたとき、すべての $n \geq c$ に対して
$$
f(M) \cap I^nN \subset f(I^{n - c}M)
$$
を満たす $c \geq 0$ が存在する。Algebra, Lemma
\ref{algebra-lemma-map-AR} を参照せよ。この状況で、{it $c$ は
Artin--Rees 補題において $f$ に対して有効である} という。

\begin{lemma}
\label{more-algebra-lemma-approximate-complex}
$A$ を Noether 環とする。$I \subset A$ を $A$ の Jacobson 根基に含まれる
イデアルとする。有限生成 $A$-加群の二つの複体
$$
S : L \xrightarrow{f} M \xrightarrow{g} N
\quad\text{and}\quad
S' : L \xrightarrow{f'} M \xrightarrow{g'} N
$$
を上のようにとる。次を仮定する。
\begin{enumerate}
\item $c$ は Artin--Rees 補題において $f$ と $g$ に対して有効である。
\item 複体 $S$ は完全である。
\item $f' = f \bmod I^{c + 1}M$ かつ $g' = g \bmod I^{c + 1}N$ である。
\end{enumerate}
このとき、$c$ は Artin--Rees 補題において $g'$ に対して有効であり、複体 $S'$ は
完全である。
\end{lemma}

\begin{proof}
まず、$n \geq c$ に対して
$g'(M) \cap I^nN \subset g'(I^{n - c}M)$ を示す。
$g'(a) \in I^nN$ を満たす $M$ の元 $a$ をとる。$g'(a)$ を変えずに $a$ を
$f'(L)$ の元で調整し、$a \in I^{n-c}M$ としたい。
$r < n - c$ かつ $a \in I^rM$ であると仮定する。このとき
$$
g(a) = g'(a) + (g - g')(a) \in
I^n N + I^{r + c + 1}N = I^{r + c + 1}N.
$$
$g$ に対する Artin--Rees により、$g(a) \in g(I^{r + 1}M)$ である。
$a_1 \in I^{r + 1}M$ を用いて $g(a) = g(a_1)$ と書く。列 $S$ は完全なので、
$a - a_1 \in f(L)$ である。従って、ある $b \in L$ によって
$a = f(b) + a_1$ と書ける。すると $f(b) = a - a_1 \in I^rM$ である。
$f$ に対する Artin--Rees により、$r \geq c$ ならば $b$ を
$I^{r - c}L$ の元に取り替えられる。従っていずれの場合にも
$a = f'(b) + a_2$ と書ける。ここで
$a_2 = (f - f')(b) + a_1 \in I^{r + 1}M$ である。（実際、$c \geq r$ ならば
仮定により $(f - f')(b) \in I^{r + 1}M$ であり、$c < r$ ならば
$b \in I^{r - c}$ なので、やはり
$(f - f')(b) \in I^{c + 1} I^{r - c} M = I^{r + 1}M$ である。）
ゆえに $f'(b) \in f'(L)$ によって $a$ を調整し、$r$ を $1$ 増やせる。

\medskip\noindent
実際、上の議論はすべての $n \geq c$ に対して
$(g')^{-1}(I^nN) \subset f'(L) + I^{n - c}M$ であることを示す。
従って $I$ は $A$ の Jacobson 根基に含まれるので、Algebra, Lemma
\ref{algebra-lemma-intersection-powers-ideal-module} により
$$
(g')^{-1}(0) = (g')^{-1}(\bigcap I^nN) \subset
\bigcap f'(L) + I^{n - c}M = f'(L)
$$
であり、$S'$ は完全である。
\end{proof}

\noindent
環 $A$ のイデアル $I \subset A$ と $A$-加群 $M$ が与えられたとき、
$$
\text{Gr}_I(M) = \bigoplus I^nM/I^{n + 1}M.
$$
とおく。これを次数付き $\text{Gr}_I(A)$-加群とみなす。

\begin{lemma}
\label{more-algebra-lemma-approximate-complex-graded}
Lemma \ref{more-algebra-lemma-approximate-complex} と同じ仮定をおく。
$Q = \Coker(g)$ および $Q' = \Coker(g')$ とする。このとき、次数付き
$\text{Gr}_I(A)$-加群として
$\text{Gr}_I(Q) \cong \text{Gr}_I(Q')$ である。
\end{lemma}

\begin{proof}
次数 $n$ において
$\text{Gr}_I(Q)_n = I^nN/(I^{n + 1}N + g(M) \cap I^nN)$
であり、$Q'$ についても同様である。次を主張する。
$$
g(M) \cap I^nN \subset I^{n + 1}N + g'(M) \cap I^nN.
$$
対称性により（この主張の証明では $c$ が $g$ に対して有効であることしか用いず、
補題により同じことが $g'$ に対しても成り立つ）、これから
$$
I^{n + 1}N + g(M) \cap I^nN = I^{n + 1}N + g'(M) \cap I^nN
$$
が従う。従って $\text{Gr}_I(Q)_n$ と $\text{Gr}_I(Q')_n$ は $N$ の部分商として
一致し、補題が従う。$g = g' \bmod I^{c + 1}N$ なので、$n \leq c$ の場合に
主張は明らかである。$n > c$ とし、$b \in g(M) \cap I^nN$ と仮定する。
$a \in I^{n - c}M$ を用いて $b = g(a)$ と書く。$b' = g'(a)$ とおく。このとき
$b - b' = (g - g')(a) \in I^{n + 1}N$ であり、主張を得る。
\end{proof}

\begin{lemma}
\label{more-algebra-lemma-works-flat-extension}
$A \to B$ を Noether 環の平坦な環準同型とする。$I \subset A$ をイデアルとし、
$f : M \to N$ を有限生成 $A$-加群の準同型とする。$c$ が Artin--Rees 補題に
おいて $f$ に対して有効であると仮定する。このとき、$c$ はイデアル $IB$ に関する
Artin--Rees 補題において
$f \otimes 1 : M \otimes_A B \to N \otimes_A B$ に対して有効である。
\end{lemma}

\begin{proof}
次に注意する。
$$
(f \otimes 1)(M) \cap I^n N \otimes_A B
= (f \otimes 1)\left((f \otimes 1)^{-1}(I^n N \otimes_A B)\right)
$$
一方、
\begin{align*}
(f \otimes 1)^{-1}(I^n N \otimes_A B) &
= \Ker(M \otimes_A B \to N \otimes_A B/(I^n N \otimes_A B)) \\
& =
\Ker(M \otimes_A B \to (N/I^nN) \otimes_A B)
\end{align*}
$A \to B$ は平坦なので、核および余核をとる操作は $B$ とのテンソル積と可換する。
従って、これは $f^{-1}(I^nN) \otimes_A B$ に等しい。仮定により
$f^{-1}(I^nN)$ は $\Ker(f) + I^{n - c}M$ に含まれる。ゆえに補題が成り立つ。
\end{proof}
\section{環のファイバー積 I}
\label{more-algebra-section-fibre-products-rings}

\noindent
環のファイバー積はスキームの押し出しと関係する。スキームの押し出しのいくつかの
場合は More on Morphisms, Section \ref{more-morphisms-section-pushouts} で論じる。

\begin{lemma}
\label{more-algebra-lemma-fibre-product-finite-type}
$R$ を環とする。$A \to B$ と $C \to B$ を $R$-代数準同型とする。次を仮定する。
\begin{enumerate}
\item $R$ は Noether 環である。
\item $A$, $B$, $C$ は $R$ 上有限型である。
\item $A \to B$ は全射である。
\item $B$ は $C$ 上有限生成である。
\end{enumerate}
このとき $A \times_B C$ は $R$ 上有限型である。
\end{lemma}

\begin{proof}
$D = A \times_B C$ とおく。完全な行をもつ可換図式
$$
\xymatrix{
0 &
B \ar[l] &
A \ar[l] &
I \ar[l] &
0 \ar[l] \\
0 &
C \ar[l] \ar[u] &
D \ar[l] \ar[u] &
I \ar[l] \ar[u] &
0 \ar[l]
}
$$
がある。$C$-加群としての $B$ の生成元 $y_1, \ldots, y_n \in B$ を選び、
$y_i$ に写る $x_i \in A$ を選ぶ。このとき $1, x_1, \ldots, x_n$ は
$D$-加群として $A$ を生成する。写像 $D \to A \times C$ は単射であり、環
$A \times C$ は $D$-加群として有限生成である（実際、有限生成 $D$-加群
$A$ と $C$ の直和である）。従って Artin--Tate の補題
（Algebra, Lemma \ref{algebra-lemma-Artin-Tate}）から結論が従う。
\end{proof}

\begin{lemma}
\label{more-algebra-lemma-formal-consequence}
$R$ を Noether 環とし、$I$ を有限集合とする。デカルト図式
$$
\xymatrix{
\prod B_i &
\prod A_i \ar[l]^{\prod \varphi_i} \\
Q \ar[u]^{\prod \psi_i} &
P \ar[u] \ar[l]
}
$$
が与えられているとする。ここで $\psi_i$ と $\varphi_i$ は全射であり、
$Q$, $A_i$, $B_i$ は $R$ 上有限型である。このとき $P$ は $R$ 上有限型である。
\end{lemma}

\begin{proof}
Lemma \ref{more-algebra-lemma-fibre-product-finite-type} と $I$ の濃度に関する帰納法から従う。
実際、$I = I' \amalg \{i_0\}$ とする。補題の図式で $I'$ を用いて定義される環を
$P'$ とする。帰納法の仮定により $P'$ は有限型である。最後に $P$ はファイバー積図式
$$
\xymatrix{
B_{i_0} &
A_{i_0} \ar[l] \\
P' \ar[u] &
P \ar[u] \ar[l] &
}
$$
に入るので、前の補題を適用できる。
\end{proof}

\begin{lemma}
\label{more-algebra-lemma-diagram-localize}
環のデカルト図式
$$
\xymatrix{
R &
R' \ar[l]^t \\
B \ar[u]_s &
B'\ar[u] \ar[l]
}
$$
が与えられている、すなわち $B' = B \times_R R'$ であるとする。
$h \in B'$ が $g \in B$ および $f \in R'$ に対応し、$s(g) = t(f)$ を満たすならば、
図式
$$
\xymatrix{
R_{s(g)} = R_{t(f)} &
(R')_f \ar[l]^-t \\
B_g \ar[u]_s &
(B')_h \ar[u] \ar[l]
}
$$
もデカルトである。
\end{lemma}

\begin{proof}
等式 $B' = B \times_R R'$ は、$B'$-加群の列
$$
0 \to B' \to B \oplus R' \xrightarrow{s, -t} R
$$
が完全であることを意味する。$B'$-加群として $B_g = B_h$,
$R'_f = R'_h$, および $R_{s(g)} = R_{t(f)} = R_h$ である。
局所化の完全性（Algebra, Proposition
\ref{algebra-proposition-localization-exact}）により、列
$$
0 \to B'_h \to B_g \oplus R'_f \xrightarrow{s, -t} R_{s(g)} = R_{t(f)}
$$
は完全である。これで補題が証明された。
\end{proof}

\noindent
環の可換図式
$$
\xymatrix{
R &
R' \ar[l] \\
B \ar[u] &
B' \ar[u] \ar[l]
}
$$
を考える。次の関手を考える（ここで圏のファイバー積は Categories, Example
\ref{categories-example-2-fibre-product-categories} で構成したものである）。
\begin{equation}
\label{more-algebra-equation-modules}
\text{Mod}_{B'} \longrightarrow
\text{Mod}_B \times_{\text{Mod}_R} \text{Mod}_{R'},\quad
L' \longmapsto (L' \otimes_{B'} B, L' \otimes_{B'} R', can)
\end{equation}
ここで $can$ は標準的な同一視
$L' \otimes_{B'} B \otimes_B R = L' \otimes_{B'} R' \otimes_{R'} R$ である。
以下では右辺の対象を $(N, M', \varphi)$ と書く。すなわち、$N$ は $B$-加群、
$M'$ は $R'$-加群であり、
$\varphi : N \otimes_B R \to M' \otimes_{R'} R$ は同型である。

\begin{lemma}
\label{more-algebra-lemma-modules}
環の可換図式
$$
\xymatrix{
R &
R' \ar[l] \\
B \ar[u] &
B' \ar[u] \ar[l]
}
$$
が与えられているとき、関手 (\ref{more-algebra-equation-modules}) は右随伴をもつ。
その右随伴は関手
$$
F : (N, M', \varphi) \longmapsto N \times_\varphi M'
$$
である（説明は証明を参照せよ）。
\end{lemma}

\begin{proof}
圏 $\text{Mod}_B \times_{\text{Mod}_R} \text{Mod}_{R'}$ の対象
$(N, M', \varphi)$ に対し、
$$
N \times_\varphi M' = \{(n, m') \in N \times M' \mid
\varphi(n \otimes 1) = m' \otimes 1\text{ in }M' \otimes_{R'} R\}
$$
とおき、これを $B'$-加群とみなす。随伴性の主張は、$B'$-加群 $L'$ と三つ組
$(N, M', \varphi)$ に対して
$$
\Hom_{B'}(L', N \times_\varphi M') =
\Hom_B(L' \otimes_{B'} B, N)
\times_{\Hom_R(L' \otimes_{B'} R, M' \otimes_{R'} R)}
\Hom_{R'}(L' \otimes_{B'} R', M')
$$
となることである。Algebra, Lemma
\ref{algebra-lemma-adjoint-tensor-restrict} により、右辺は
$$
\Hom_{B'}(L', N) \times_{\Hom_{B'}(L', M' \otimes_{R'} R)} \Hom_{B'}(L', M')
$$
に等しい。従って、このファイバー積の元の対 $(g, f')$ から $B'$-線型写像
% SOURCE-ERRATUM-CANDIDATE P02-E0743: frozen source says “an B'-linear map”; Japanese has no indefinite article.
$L' \to N \times_\varphi M'$, $l' \mapsto (g(l'), f'(l'))$ が得られることは明らかである。
逆に、$B'$-線型写像 $g' : L' \to N \times_\varphi M'$ が与えられたとき、
$g$ を合成 $L' \to N \times_\varphi M' \to N$ とし、$f'$ を合成
$L' \to N \times_\varphi M' \to M'$ とする。これらの構成は互いに逆であり、
所望の同型を定める。
\end{proof}

\section{環のファイバー積 II}
\label{more-algebra-section-fibre-products-rings-II}

\noindent
本節では、次の状況におけるファイバー積を論じる。

\begin{situation}
\label{more-algebra-situation-module-over-fibre-product}
以下では環準同型
$$
\xymatrix{
B \ar[r] & A & A' \ar[l]
}
$$
を考え、$A' \to A$ は全射で核が $I$ であると仮定する。この状況で
$B' = B \times_A A'$ とおくと、デカルト正方形
$$
\xymatrix{
A & A' \ar[l] \\
B \ar[u] & B' \ar[l] \ar[u]
}
$$
を得る。
\end{situation}

\begin{lemma}
\label{more-algebra-lemma-points-of-fibre-product}
Situation \ref{more-algebra-situation-module-over-fibre-product} において、位相空間として
$$
\Spec(B') = \Spec(B) \amalg_{\Spec(A)} \Spec(A')
$$
である。
\end{lemma}

\begin{proof}
$B' = B \times_A A'$ なので、スペクトルの可換正方形を得る。その始域は
位相空間の圏における押し出し（Topology, Section
\ref{topology-section-colimits} により存在する）であるから、これは連続写像
$$
can : \Spec(B) \amalg_{\Spec(A)} \Spec(A') \longrightarrow \Spec(B')
$$
を誘導する。

\medskip\noindent
写像 $can$ が全射であることを示すため、素イデアル $\mathfrak q' \subset B'$ をとる。
$I \subset \mathfrak q'$ ならば（ここでも以下でも、$I$ を $A'$ のイデアルであると
同時に $B'$ のイデアルともみなす）、$\mathfrak q'$ は $B$ の素イデアルに対応し、
像に入る。そうでなければ、$h \in I$, $h \not \in \mathfrak q'$ を選ぶ。この場合
$B_h = A_h = 0$ であり、Lemma \ref{more-algebra-lemma-diagram-localize} により環準同型
$B'_h \to A'_h$ は同型である。従って $\mathfrak q'$ は $I$ を含まない一意的な
素イデアル $\mathfrak p' \subset A'$ に対応する。

\medskip\noindent
$B' \to B$ は全射なので、$can$ は成分 $\Spec(B)$ 上で単射である。上で、
$\Spec(A') \to \Spec(B')$ は $V(I) \subset \Spec(A')$ の補集合上で単射であることを
見た。$V(I) \subset \Spec(A')$ はちょうど $\Spec(A) \to \Spec(A')$ の像なので、
自明な集合論的議論により $can$ は単射である。

\medskip\noindent
証明を終えるには $can$ が開写像であることを示せばよい。このため、押し出しの開集合は
$V \amalg U'$ の形であることに注意する。ここで $V \subset \Spec(B)$ と
$U' \subset \Spec(A')$ は、その $\Spec(A)$ への逆像が一致する開集合である。
$v \in V$ をとる。$v \in D(g) \subset V$ を満たす $g \in B$ を選べる。その像を
$f \in A$ とし、$f$ に写る $f' \in A'$ を選ぶ。このとき
$D(f') \cap U' \cap V(I) = D(f') \cap V(I)$ である。従って
$V(I) \cap D(f')$ と $D(f') \cap (U')^c$ は
$D(f') = \Spec(A'_{f'})$ の互いに素な閉部分集合である。あるイデアル
$J \subset A'$ によって $(U')^c = V(J)$ と書く。今示した非交叉性により
$$
A'_{f'} \to A'_{f'}/IA'_{f'} \times A'_{f'}/JA'_{f'}
$$
は全射なので、$A'_{f'}/IA'_{f'}$ で $1$ に写り、
$A'_{f'}/JA'_{f'}$ で零に写る $a'' \in A'_{f'}$ を選べる。分母を払うと、$A$ で
$f^n$ に写る元 $a' \in J$ を得る。このとき $D(a'f') \subset U'$ である。
$h' = (g^{n + 1}, a'f') \in B'$ とおく。先に引用した補題により
$B'_{h'} = B_{g^{n + 1}} \times_{A_{f^{n + 1}}} A'_{a'f'}$ なので、$D(h')$ は
押し出しにおける $v$ の開近傍に引き戻される。すなわち $V \amalg U'$ の像は
$v$ の像の開近傍を含む。$u' \in U'$ かつ $u' \not \in V(I)$ の場合にも同じことが
成り立つことの（より容易な）証明は省略する。
\end{proof}

\begin{lemma}
\label{more-algebra-lemma-fibre-product-integral}
Situation \ref{more-algebra-situation-module-over-fibre-product} において、$B \to A$ が整ならば
$B' \to A'$ も整である。
\end{lemma}

\begin{proof}
$a' \in A'$ とし、その像を $a \in A$ とする。$a$ が満たす $B$ 係数のモニック多項式
% SOURCE-ERRATUM-CANDIDATE P02-E0744: frozen source says “monical polynomial”; Japanese uses the intended standard term.
$x^d + b_1 x^{d - 1} + \ldots + b_d$ をとる。$b_i \in B$ に写る
$b'_i \in B'$ を選ぶ（これは可能である）。すると
$(a')^d + b'_1 (a')^{d - 1} + \ldots + b'_d$ は $A' \to A$ の核に属する。
$\Ker(B' \to B) = \Ker(A' \to A)$ なので、$b'_d$ の選び方を変更して、所望の
$$
(a')^d + b'_1 (a')^{d - 1} + \ldots + b'_d = 0
$$
を得る。
\end{proof}

\noindent
Situation \ref{more-algebra-situation-module-over-fibre-product} において、$B'$-加群を
$A'$, $A$, $B$ 上の加群によって理解したい。このため、次の関手を考える
（圏のファイバー積は Categories, Example
\ref{categories-example-2-fibre-product-categories} で構成したものである）。
\begin{equation}
\label{more-algebra-equation-functor}
\text{Mod}_{B'} \longrightarrow
\text{Mod}_B \times_{\text{Mod}_A} \text{Mod}_{A'},\quad
L' \longmapsto (L' \otimes_{B'} B, L' \otimes_{B'} A', can)
\end{equation}
ここで $can$ は標準的な同一視
$L' \otimes_{B'} B \otimes_B A = L' \otimes_{B'} A' \otimes_{A'} A$ である。
以下では右辺の対象を $(N, M', \varphi)$ と書く。すなわち、$N$ は $B$-加群、
$M'$ は $A'$-加群であり、
$\varphi : N \otimes_B A \to M' \otimes_{A'} A$ は同型である。しかし、多くの場合
% SOURCE-ERRATUM-CANDIDATE P02-E0745: frozen source omits “to” in “more convenient to think”.
$\varphi$ を、同型 $N \otimes_B A \to M' \otimes_{A'} A = M'/IM'$ を誘導する
$B$-線型写像 $\varphi : N \to M'/IM'$ とみなす方が便利である。

\begin{lemma}
\label{more-algebra-lemma-module-over-fibre-product}
Situation \ref{more-algebra-situation-module-over-fibre-product} において、関手
(\ref{more-algebra-equation-functor}) は右随伴をもつ。その右随伴は
$$
F : (N, M', \varphi) \longmapsto N \times_{\varphi, M} M'
$$
である。ここで $M = M'/IM'$ とする。さらに、$F$ と
(\ref{more-algebra-equation-functor}) の合成は
$\text{Mod}_B \times_{\text{Mod}_A} \text{Mod}_{A'}$ 上の恒等関手である。
言い換えれば、$N' = N \times_{\varphi, M} M'$ とおくと、
$N' \otimes_{B'} B = N$ かつ $N' \otimes_{B'} A' = M'$ である。
\end{lemma}

\begin{proof}
随伴性の主張は、より一般的な Lemma \ref{more-algebra-lemma-modules} から従う。
最後の主張を証明するため、$B' = B \times_A A'$ および
$N' = N \times_{\varphi, M} M'$ を思い出し、これらの等式を
$$
\vcenter{
\xymatrix{
A & A' \ar[l] & I \ar[l] \\
B \ar[u] & B' \ar[l] \ar[u] & J \ar[l] \ar[u]
}
}
\quad\text{and}\quad
\vcenter{
\xymatrix{
M & M' \ar[l] & K \ar[l] \\
N \ar[u]_\varphi & N' \ar[l] \ar[u] & L \ar[l] \ar[u]
}
}
$$
に拡張する。ここで $I, J, K, L$ は元の図式の水平写像の核である。
証明を一連の観察として述べる。
\begin{enumerate}
\item $K = IM'$ である（補題の主張を参照せよ）。
\item $B' \to B$ は核 $J$ をもつ全射であり、$J \to I$ は全単射である。
\item $N' \to N$ は核 $L$ をもつ全射であり、$L \to K$ は全単射である。
\item $JN' \subset L$ である。
\item $\Im(N \to M)$ は $A$-加群として $M$ を生成する
（$N \otimes_B A = M$ だからである）。
\item $\Im(N' \to M')$ は $A'$-加群として $M'$ を生成する
（これは $K$ を法として成り立ち、$L$ は $K$ に同型に写るからである）。
\item $JN' = L$ である（前項により、$L \cong K = I M'$ は
$x \in I$ と $n' \in N'$ に対する元 $x n'$ の像で生成されるからである）。
\item $N' \otimes_{B'} B = N$ である（$N = N'/L$, $B = B'/J$ および前項による）。
\item 写像 $\gamma : N' \otimes_{B'} A' \to M'$ がある。
\item $\gamma$ は全射である（上を参照せよ）。
\item 合成 $N' \otimes_{B'} A' \to M' \to M$ の核は、
% SOURCE-ERRATUM-CANDIDATE P02-E0746: frozen source says l in K here; the tensor-source kernel is L.
$l \in K$, $n' \in N'$, $x \in I$ に対する元 $l \otimes 1$ と
$n' \otimes x$ で生成される（仮定により $M = N \otimes_B A$ であり、
$N' \to N$ と $A' \to A$ はそれぞれ核 $L$ と $I$ をもつ全射だからである）。
\item $l \in L$, $n' \in N'$, $x \in I$ に対する元 $l \otimes 1$ と
$n' \otimes x$ で生成される $N' \otimes_{B'} A'$ の部分加群の任意の元は、
ある $l \in L$ によって $l \otimes 1$ と書ける（$J$ は $I$ に同型に写るので、
$N' \otimes_{B'} A'$ において $n' \otimes x = n'x \otimes 1$ である。同様に、
$n' \in N'$, $x \in J$, $a' \in A'$ のとき、$N' \otimes_{B'} A'$ において
$x n' \otimes a' = n' \otimes xa' = n'(xa') \otimes 1$ である。
既に $JN' = L$ を見たので、主張が従う）。
\item $\gamma$ の核は零である（(10) と (11) により、核の任意の元は
$l \in L$ による $l \otimes 1$ の形であり、$\gamma$ によって
$l \in K \subset M'$ に写るからである）。
\end{enumerate}
これで証明が完了する。
\end{proof}

\begin{lemma}
\label{more-algebra-lemma-module-over-fibre-product-bis}
Lemma \ref{more-algebra-lemma-module-over-fibre-product} の状況で、$B'$-加群 $L'$ に対する随伴写像
$$
L' \longrightarrow
(L' \otimes_{B'} B) \times_{(L' \otimes_{B'} A)} (L' \otimes_{B'} A')
$$
は全射であるが、一般には単射ではない。
\end{lemma}

\begin{proof}
Lemma \ref{more-algebra-lemma-module-over-fibre-product} の証明と同様に、$J \subset B'$ を
写像 $B' \to B$ の核とする。このとき $L' \otimes_{B'} B = L'/JL'$ である。
従って全射性を示すには、ファイバー積の $(0, z)$ の形の元が補題の写像の像に
入ることを示せば十分である。写像
$L' \otimes_{B'} A' \to L' \otimes_{B'} A$ の核は
$L' \otimes_{B'} I \to L' \otimes_{B'} A'$ の像である。$B' \to A'$ が誘導する
写像 $J \to I$ は同型なので、合成
$$
L' \otimes_{B'} J \to L' \to
(L' \otimes_{B'} B) \times_{(L' \otimes_{B'} A)} (L' \otimes_{B'} A')
$$
は $L' \otimes_{B'} J$ から $(0, z)$ の形の元の集合への全射を誘導する。
写像が一般には単射でないことを見るため、簡単な例を挙げる。体 $k$ をとり、
$B' = k[x, y]/(xy)$, $A' = B'/(x)$, $B = B'/(y)$, $A = B'/(x, y)$,
$L' = B'/(x - y)$ とおく。この場合、$L'$ における $x$ の類は零でないが、
表示した矢印によって零に写る。
\end{proof}

\begin{lemma}
\label{more-algebra-lemma-surjection-module-over-fibre-product}
Situation \ref{more-algebra-situation-module-over-fibre-product} において、
$\text{Mod}_B \times_{\text{Mod}_A} \text{Mod}_{A'}$ の射
$(N_1, M'_1, \varphi_1) \to (N_2, M'_2, \varphi_2)$ をとり、
$N_1 \to N_2$ と $M'_1 \to M'_2$ は全射であるとする。このとき、
$M_1 = M'_1/IM'_1$ および $M_2 = M'_2/IM'_2$ とおけば、
$$
N_1 \times_{\varphi_1, M_1} M'_1 \to N_2 \times_{\varphi_2, M_2} M'_2
$$
は全射である。
\end{lemma}

\begin{proof}
$(x_2, y_2) \in N_2 \times_{\varphi_2, M_2} M'_2$ をとる。$x_2$ に写る
$x_1 \in N_1$ を選ぶ。$M'_1 \to M_1$ は全射なので、
$\varphi_1(x_1)$ に写る $y_1 \in M'_1$ を選べる。このとき
$(x_1, y_1)$ は $N_2 \times_{\varphi_2, M_2} M'_2$ における
$(x_2, y'_2)$ に写る。従って、$(0, y_2)$ の形の元が写像の像に入ることを
示せば十分である。このとき $y_2 \in IM'_2$ である。
$t_i \in I$ を用いて $y_2 = \sum t_i y_{2, i}$ と書く。
$y_{2, i}$ に写る $y_{1, i} \in M'_1$ を選ぶ。このとき
$y_1 = \sum t_iy_{1, i} \in IM'_1$ であり、元 $(0, y_1)$ が所望のものである。
\end{proof}

\begin{lemma}
\label{more-algebra-lemma-finite-module-over-fibre-product}
Lemma \ref{more-algebra-lemma-module-over-fibre-product} と同じ
$A, A', B, B', I, M, M', N, \varphi$ をとる。
% SOURCE-ERRATUM-CANDIDATE P02-E0747: frozen source omits “is” in “If N is finite over B”.
$N$ が $B$ 上有限生成で、$M'$ が $A'$ 上有限生成ならば、
$N' = N \times_{\varphi, M} M'$ は $B'$ 上有限生成である。
\end{lemma}

\begin{proof}
以下では Lemma \ref{more-algebra-lemma-module-over-fibre-product} の結果を断りなく用いる。
$B$ 上の $N$ の生成元 $y_1, \ldots, y_r$ と、$A'$ 上の $M'$ の生成元
$x_1, \ldots, x_s$ を選ぶ。$N = N' \otimes_{B'} B$ であり
$B' \to B$ は全射なので、$N$ で $y_1, \ldots, y_r$ に写る
$u_1, \ldots, u_r \in N'$ を選べる。$M' = N' \otimes_{B'} A'$ なので、
ある $a'_{ij} \in A'$ によって $x_i = \sum v_j \otimes a'_{ij}$ となる
$v_1, \ldots, v_t \in N'$ を選べる。特に、$v_j$ の像
$\overline{v}_j \in M'$ は $A'$ 上 $M'$ を生成する。
$u_1, \ldots, u_r, v_1, \ldots, v_t$ が $B'$-加群として $N'$ を生成すると主張する。
実際、$\xi \in N'$ をとる。まず、$\xi$ と $\sum b'_i u_i$ が $N$ の同じ元に写る
ように $b'_1, \ldots, b'_r \in B'$ を選ぶ。これは $B' \to B$ が全射で、
$y_1, \ldots, y_r$ が $B$ 上 $N$ を生成するため可能である。
差 $\xi - \sum b'_i u_i$ は、ある $\theta$ によって $(0, \theta)$ の形である。
この元は $IM'$ に属する。$t_j \in I$ を用いて
$\theta$ を $\sum t_j\overline{v}_j$ と書く。
$J = \Ker(B' \to B)$ は $I$ に同型に写るので、$t_j$ に写る
$s_j \in J \subset B'$ を選べる。$N' = N \times_{\varphi, M} M'$ なので、
所望の $\xi = \sum b'_i u_i + \sum s_j v_j$ が従う。
\end{proof}

\begin{lemma}
\label{more-algebra-lemma-flat-module-over-fibre-product}
Situation \ref{more-algebra-situation-module-over-fibre-product} と同じ
$A, A', B, B', I$ をとる。
\begin{enumerate}
\item $(N, M', \varphi)$ を
$\text{Mod}_B \times_{\text{Mod}_A} \text{Mod}_{A'}$ の対象とする。
$M'$ が $A'$ 上平坦で、$N$ が $B$ 上平坦ならば、
$N' = N \times_{\varphi, M} M'$ は $B'$ 上平坦である。
\item $L'$ が平坦な $B'$-加群ならば、
% SOURCE-ERRATUM-CANDIDATE P02-E0748: frozen formula uses undefined L rather than the stated L'; preserved.
$L' = (L \otimes_{B'} B) \times_{(L \otimes_{B'} A)} (L \otimes_{B'} A')$ である。
\item 平坦 $B'$-加群の圏は、
$N$ が $B$ 上平坦で $M'$ が $A'$ 上平坦である三つ組 $(N, M', \varphi)$ からなる
$\text{Mod}_B \times_{\text{Mod}_A} \text{Mod}_{A'}$ の充満部分圏と同値である。
\end{enumerate}
\end{lemma}

\begin{proof}
証明では Lemma \ref{more-algebra-lemma-module-over-fibre-product} を断りなく用いる。

\medskip\noindent
(1) の証明。$J = \Ker(B' \to B)$ とおく。これは $B'$ のイデアルであり、
$I = \Ker(A' \to A)$ に同型に写る。$\mathfrak b' \subset B'$ をイデアルとする。
Algebra, Lemma \ref{algebra-lemma-flat} により、
$\mathfrak b' \otimes_{B'} N' \to N'$ が単射であることを示せばよい。
$N$ は $B$ 上平坦なので、
$$
\mathfrak b'/(\mathfrak b' \cap J) \otimes_{B'} N' =
\mathfrak b'/(\mathfrak b' \cap J) \otimes_B N \to N
$$
は単射である。列
$\mathfrak b' \cap J \to \mathfrak b' \to
\mathfrak b'/(\mathfrak b' \cap J) \to 0$ は完全なので、
$(\mathfrak b' \cap J) \otimes_{B'} N' \to N'$ が単射であることを示せば十分である。
従って $\mathfrak b' \subset J$ と仮定してよい。次に、$J \to I$ は同型なので
$$
J \otimes_{B'} N' =
I \otimes_{A'} A' \otimes_{B'} N' =
I \otimes_{A'} M'
$$
であり、$M'$ は平坦な $A'$-加群なので、これは $M'$ に単射で写る。
従って $J \otimes_{B'} N' \to N'$ は単射であり、Algebra, Remark
\ref{algebra-remark-Tor-ring-mod-ideal} により
$\text{Tor}_1^{B'}(B'/J, N') = 0$ と結論する。ゆえに、$B'$ 上の $N'$ と
イデアル $J$ に Algebra, Lemma \ref{algebra-lemma-what-does-it-mean} を適用できる。
イデアル $\mathfrak b' \subset J$ に戻り、その $I$ における像が $I$ の
$A'$-部分加群となるような最小のイデアル
$\mathfrak b' \subset \mathfrak b'' \subset J$ をとる。言い換えれば、
$J = I$ を $A'$-加群とみなすと $\mathfrak b'' = A' \mathfrak b'$ である。
このとき $\mathfrak b''/\mathfrak b'$ は $J$ で零化され、補題の適用による
$\text{Tor}_1^{B'}(\mathfrak b''/\mathfrak b', N')$ の消滅から短完全列
$$
0 \to \mathfrak b' \otimes_{B'} N' \to
\mathfrak b'' \otimes_{B'} N' \to
\mathfrak b''/\mathfrak b' \otimes_{B'} N' \to 0
$$
を得る。従って $\mathfrak b'$ を $\mathfrak b''$ に置き換えてよい。
特に、$\mathfrak b'$ が $A'$-加群であり、$A'$ のイデアルに写ると仮定してよい。
このとき
$$
\mathfrak b' \otimes_{B'} N' =
\mathfrak b' \otimes_{A'} A' \otimes_{B'} N' =
\mathfrak b' \otimes_{A'} M'
$$
である。このテンソル積は、$M'$ が $A'$ 上平坦であるという仮定により
$M'$ に単射で写る。従って
$\mathfrak b' \otimes_{B'} N' \to N' \to M'$ は単射であり、所望のように
最初の写像も単射である。

\medskip\noindent
(2) の証明。短完全列 $0 \to B' \to B \oplus A' \to A \to 0$ と
$L'$ との $B'$ 上のテンソル積をとれば従う。

\medskip\noindent
(3) の証明。(1) と (2) の直ちの帰結である。
\end{proof}

\begin{lemma}
\label{more-algebra-lemma-finitely-presented-module-over-fibre-product}
Situation \ref{more-algebra-situation-module-over-fibre-product} と同じ
$A, A', B, B', I$ をとる。有限生成射影 $B'$-加群の圏は、
$N$ が $B$ 上有限生成射影で $M'$ が $A'$ 上有限生成射影である三つ組
$(N, M', \varphi)$ からなる
$\text{Mod}_B \times_{\text{Mod}_A} \text{Mod}_{A'}$ の充満部分圏と同値である。
\end{lemma}

\begin{proof}
加群が有限生成射影であることと、有限表示かつ平坦であることは同値である。
Algebra, Lemma \ref{algebra-lemma-finite-projective} を参照せよ。
Lemmas \ref{more-algebra-lemma-flat-module-over-fibre-product} および
\ref{more-algebra-lemma-finite-module-over-fibre-product} を用いると、$N$ が $B$ 上有限生成射影で、
$M'$ が $A'$ 上有限生成射影であるとき、
$N' = N \times_{\varphi, M} M'$ が有限表示 $B'$-加群であることを示せばよい。

\medskip\noindent
Lemma \ref{more-algebra-lemma-finite-module-over-fibre-product} により、$N'$ は $B'$ 上有限生成である。
核を $K'$ とする全射 $(B')^{\oplus n} \to N'$ を選ぶ。基底変換により、核がそれぞれ
$K_B$, $K_{A'}$, $K_A$ である写像
$B^{\oplus n} \to N$, $(A')^{\oplus n} \to M'$, $A^{\oplus n} \to M$ を得る。
標準写像
$$
K' \longrightarrow K_B \times_{K_A} K_{A'}
$$
がある。一方、$N' = N \times_{\varphi, M} M'$ および
$B' = B \times_A A'$ なので、表示した矢印の逆となる標準写像
$K_B \times_{K_A} K_{A'} \to K'$ もある。従って表示した写像は同型である。
Algebra, Lemma \ref{algebra-lemma-extension} により、$K_B$ と $K_{A'}$ は有限生成である。
Lemma \ref{more-algebra-lemma-finite-module-over-fibre-product} により、
$K_B \to K_A$ と $K_{A'} \to K_A$ が同型
$K_B \otimes_B A = K_A = K_{A'} \otimes_{A'} A$ を誘導するならば、
$K'$ は有限生成 $B'$-加群であると結論する。
% SOURCE-ERRATUM-CANDIDATE P02-E0749: frozen source has plural subject with “implies”.
これは、平坦性の仮定により、テンソル積をとっても列
$$
0 \to K_B \to B^{\oplus n} \to N \to 0
\quad\text{and}\quad
0 \to K_{A'} \to (A')^{\oplus n} \to M' \to 0
$$
が完全に保たれるため成り立つ。Algebra, Lemma
\ref{algebra-lemma-flat-tor-zero} を参照せよ。
\end{proof}

\section{環のファイバー積 III}
\label{more-algebra-section-fibre-products-rings-III}

\noindent
本節では、次の状況におけるファイバー積を論じる。

\begin{situation}
\label{more-algebra-situation-relative-module-over-fibre-product}
$A, A', B, B', I$ を Situation
\ref{more-algebra-situation-module-over-fibre-product} と同じものとする。
環準同型 $B' \to D'$ をとり、$D = D' \otimes_{B'} B$,
$C' = D' \otimes_{B'} A'$, $C = D' \otimes_{B'} A$ とおく。これにより環の大きな
可換図式
$$
\xymatrix{
C & & & C' \ar[lll] \\
& A \ar[ul] & A' \ar[l] \ar[ru] \\
& B \ar[u] \ar[ld] & B' \ar[l] \ar[u] \ar[rd] \\
D \ar[uuu] & & & D' \ar[lll] \ar[uuu]
}
$$
を得る。写像 $D' \to D \times_C C'$ が同型であるとは {\bf 仮定しない} ことに
注意せよ\footnote{ただし、Lemma \ref{more-algebra-lemma-module-over-fibre-product-bis} により
$D' \to D \times_C C'$ は全射である。}。
この状況で、(\ref{more-algebra-equation-functor}) と同様の関手
\begin{equation}
\label{more-algebra-equation-relative-functor}
\text{Mod}_{D'} \longrightarrow
\text{Mod}_D \times_{\text{Mod}_C} \text{Mod}_{C'},\quad
L' \longmapsto (L' \otimes_{D'} D, L' \otimes_{D'} C', can)
\end{equation}
を得る。次に注意せよ。
% SOURCE-ERRATUM-CANDIDATE P02-E0750: frozen identities use undefined L on the right; preserved.
$L' \otimes_{D'} D = L \otimes_{D'} (D' \otimes_{B'} B) = L \otimes_{B'} B$
であり、同様に
$L' \otimes_{D'} C' = L \otimes_{D'} (D' \otimes_{B'} A') = L \otimes_{B'} A'$
であるから、図式
$$
\xymatrix{
\text{Mod}_{D'} \ar[r] \ar[d] &
\text{Mod}_D \times_{\text{Mod}_C} \text{Mod}_{C'} \ar[d] \\
\text{Mod}_{B'} \ar[r] &
\text{Mod}_B \times_{\text{Mod}_A} \text{Mod}_{A'}
}
$$
は可換である。以下では
$\text{Mod}_D \times_{\text{Mod}_C} \text{Mod}_{C'}$ の対象を
$(N, M', \varphi)$ と書く。すなわち、$N$ は $D$-加群、$M'$ は $C'$-加群であり、
% SOURCE-ERRATUM-CANDIDATE P02-E0751: frozen English uses “an C'-module”.
$\varphi : N \otimes_B A \to M' \otimes_{A'} A$ は $C$-加群の同型である。
しかし、多くの場合、$\varphi$ を、同型
% SOURCE-ERRATUM-CANDIDATE P02-E0752: frozen source omits “to” in “more convenient to think”.
$N \otimes_B A \to M' \otimes_{A'} A = M'/IM'$ を誘導する $D$-線型写像
$\varphi : N \to M'/IM'$ とみなす方が便利である。
\end{situation}

\begin{lemma}
\label{more-algebra-lemma-relative-module-over-fibre-product}
Situation \ref{more-algebra-situation-relative-module-over-fibre-product} において、関手
(\ref{more-algebra-equation-relative-functor}) は右随伴をもつ。その右随伴は関手
$$
F : (N, M', \varphi) \longmapsto N \times_{\varphi, M} M'
$$
である。ここで $M = M'/IM'$ とする。さらに、$F$ と
(\ref{more-algebra-equation-relative-functor}) の合成は
$\text{Mod}_D \times_{\text{Mod}_C} \text{Mod}_{C'}$ 上の恒等関手である。
言い換えれば、$N' = N \times_{\varphi, M} M'$ とおくと、
$N' \otimes_{D'} D = N$ かつ $N' \otimes_{D'} C' = M'$ である。
\end{lemma}

\begin{proof}
随伴性の主張は、より一般的な Lemma \ref{more-algebra-lemma-modules} から従う。
最後の主張は Lemma \ref{more-algebra-lemma-module-over-fibre-product} の対応する主張から従う。
実際、
$N' \otimes_{D'} D = N' \otimes_{D'} D' \otimes_{B'} B = N' \otimes_{B'} B$
および
$N' \otimes_{D'} C' = N' \otimes_{D'} D' \otimes_{B'} A' = N' \otimes_{B'} A'$
だからである。
\end{proof}

\begin{lemma}
\label{more-algebra-lemma-relative-surjection-ideals}
Situation \ref{more-algebra-situation-relative-module-over-fibre-product} において、
$J = \Ker(B' \to B)$ とおくと、写像 $JD' \to IC'$ は全射である。
\end{lemma}

\begin{proof}
$C' = D' \otimes_{B'} A'$ なので、$IC'$ は写像
$D' \otimes_{B'} I = C' \otimes_{A'} I \to C'$ の像である。
環準同型 $B' \to A'$ は同型 $J \to I$ を誘導するので、補題が従う。
\end{proof}

\begin{lemma}
\label{more-algebra-lemma-relative-finite-module-over-fibre-product}
Lemma \ref{more-algebra-lemma-relative-module-over-fibre-product} と同じ
$A, A', B, B', C, C', D, D', I, M', M, N, \varphi$ をとる。
% SOURCE-ERRATUM-CANDIDATE P02-E0753: frozen source omits “is” in “If N is finite over D”.
$N$ が $D$ 上有限生成で、$M'$ が $C'$ 上有限生成ならば、
$N' = N \times_{\varphi, M} M'$ は $D'$ 上有限生成である。
\end{lemma}

\begin{proof}
Lemma \ref{more-algebra-lemma-module-over-fibre-product-bis} により
$D' \to D \times_C C'$ は全射であることを思い出そう。
$N' = N \times_{\varphi, M} M'$ は $D \times_C C'$ 上の加群である。
データ $C, C', D, D', IC', M', M, N, \varphi$ に
Lemma \ref{more-algebra-lemma-finite-module-over-fibre-product} を適用すると、
$N' = N \times_{\varphi, M} M'$ は $D \times_C C'$ 上有限生成である。
従って $D'$ 上有限生成である。
\end{proof}

\begin{lemma}
\label{more-algebra-lemma-relative-flat-module-over-fibre-product}
Situation \ref{more-algebra-situation-relative-module-over-fibre-product} と同じ
$A, A', B, B', C, C', D, D', I$ をとる。
\begin{enumerate}
\item $(N, M', \varphi)$ を
$\text{Mod}_D \times_{\text{Mod}_C} \text{Mod}_{C'}$ の対象とする。
$M'$ が $A'$ 上平坦で、$N$ が $B$ 上平坦ならば、
$N' = N \times_{\varphi, M} M'$ は $B'$ 上平坦である。
\item $L'$ が $B'$ 上平坦な $D'$-加群ならば、
% SOURCE-ERRATUM-CANDIDATE P02-E0754: frozen formula uses undefined L rather than the stated L'; preserved.
$L' = (L \otimes_{D'} D) \times_{(L \otimes_{D'} C)} (L \otimes_{D'} C')$ である。
\item $B'$ 上平坦な $D'$-加群の圏は、
$N$ が $B$ 上平坦で $M'$ が $A'$ 上平坦である
$\text{Mod}_D \times_{\text{Mod}_C} \text{Mod}_{C'}$ の対象
$(N, M', \varphi)$ の圏と同値である。
\end{enumerate}
\end{lemma}

\begin{proof}
(1) は Lemma \ref{more-algebra-lemma-flat-module-over-fibre-product} の (1) から従う。

\medskip\noindent
(2) は Lemma \ref{more-algebra-lemma-flat-module-over-fibre-product} の (2) と、
$L' \otimes_{D'} D = L' \otimes_{B'} B$,
$L' \otimes_{D'} C' = L' \otimes_{B'} A'$, および
$L' \otimes_{D'} C = L' \otimes_{B'} A$ から従う。
Situation \ref{more-algebra-situation-relative-module-over-fibre-product} の議論を参照せよ。

\medskip\noindent
(3) は (1) と (2) の直ちの帰結である。
\end{proof}

\noindent
次の補題は、絶対的な場合の対応物
（Lemma \ref{more-algebra-lemma-finitely-presented-module-over-fibre-product}）よりはるかに
興味深いが、証明は本質的に同じである。

\begin{lemma}
\label{more-algebra-lemma-relative-finitely-presented-module-over-fibre-product}
Lemma \ref{more-algebra-lemma-relative-module-over-fibre-product} と同じ
$A, A', B, B', C, C', D, D', I, M', M, N, \varphi$ をとる。次を仮定する。
\begin{enumerate}
\item $N$ は $D$ 上有限表示であり、$B$ 上平坦である。
\item $M'$ は $C'$ 上有限表示であり、$A'$ 上平坦である。
% SOURCE-ERRATUM-CANDIDATE P02-E0755: frozen condition (2) omits “is”.
\item 環準同型 $B' \to D'$ は $B' \to D'' \to D'$ と分解し、
$B' \to D''$ は平坦、$D'' \to D'$ は有限表示である。
\end{enumerate}
このとき $N' = N \times_M M'$ は $D'$ 上有限表示である。
\end{lemma}

\begin{proof}
有限生成な核 $J$ をもつ全射
$D''' = D''[x_1, \ldots, x_n] \to D'$ を選ぶ。Algebra, Lemma
\ref{algebra-lemma-finite-finitely-presented-extension} により、$N'$ が
$D'''$-加群として有限表示であることを示せば十分である。さらに、
$D''' \otimes_{B'} B \to D' \otimes_{B'} B = D$ および
$D''' \otimes_{B'} A' \to D' \otimes_{B'} A' = C'$ は、その核が $J$ の像で
生成される全射である。従って、再び Algebra, Lemma
\ref{algebra-lemma-finite-finitely-presented-extension} により、$N$ は
有限表示 $D''' \otimes_{B'} B$-加群であり、$M'$ は有限表示
$D''' \otimes_{B'} A'$-加群である。ゆえに $D'$ を $D'''$ に、
$D$ を $D''' \otimes_{B'} B$ に置き換え、以下同様にしてよい。
$D'''$ は $B'$ 上平坦なので、$B' \to D'$ が平坦であると仮定してよい。

\medskip\noindent
$B' \to D'$ が平坦であると仮定する。
Lemma \ref{more-algebra-lemma-relative-finite-module-over-fibre-product} により、
$N'$ は $D'$ 上有限生成である。核を $K'$ とする全射
$(D')^{\oplus n} \to N'$ を選ぶ。基底変換により、核がそれぞれ
$K_D$, $K_{C'}$, $K_C$ である写像
$D^{\oplus n} \to N$, $(C')^{\oplus n} \to M'$, $C^{\oplus n} \to M$ を得る。
標準写像
$$
K' \longrightarrow K_D \times_{K_C} K_{C'}
$$
がある。一方、$N' = N \times_M M'$ であり、
$D' = D \times_C C'$ であるから（平坦 $B'$-加群 $D'$ に
Lemma \ref{more-algebra-lemma-flat-module-over-fibre-product} を適用せよ）、表示した矢印の逆となる
標準写像 $K_D \times_{K_C} K_{C'} \to K'$ もある。従って表示した写像は同型である。
Algebra, Lemma \ref{algebra-lemma-extension} により、$K_D$ と $K_{C'}$ は有限生成である。
Lemma \ref{more-algebra-lemma-relative-finite-module-over-fibre-product} により、
$K_D \to K_C$ と $K_{C'} \to K_C$ が同型
$K_D \otimes_B A = K_C = K_{C'} \otimes_{A'} A$ を誘導するならば、
$K'$ は有限生成 $D'$-加群であると結論する。
% SOURCE-ERRATUM-CANDIDATE P02-E0756: frozen source has plural subject with “implies”.
これは、平坦性の仮定により、テンソル積をとっても列
$$
0 \to K_D \to D^{\oplus n} \to N \to 0
\quad\text{and}\quad
0 \to K_{C'} \to (C')^{\oplus n} \to M' \to 0
$$
が完全に保たれるため成り立つ。Algebra, Lemma
\ref{algebra-lemma-flat-tor-zero} を参照せよ。
\end{proof}

\begin{lemma}
\label{more-algebra-lemma-properties-algebras-over-fibre-product}
Situation \ref{more-algebra-situation-module-over-fibre-product} と同じ
$A, A', B, B', I$ をとる。$B$-代数 $D$、$A'$-代数 $C'$、および同型
% SOURCE-ERRATUM-CANDIDATE P02-E0757: frozen English reverses the articles before B and A'.
$D \otimes_B A \to C'/IC' = C$ からなる系 $(D, C', \varphi)$ をとる。
Lemma \ref{more-algebra-lemma-module-over-fibre-product} と同様に
$D' = D \times_C C'$ とおく。このとき
\begin{enumerate}
\item $B' \to D'$ が有限型であることと、$B \to D$ および $A' \to C'$ が
有限型であることは同値である。
% SOURCE-ERRATUM-CANDIDATE P02-E0758: frozen item (1) omits “of” in the finite-type predicates.
\item $B' \to D'$ が平坦であることと、$B \to D$ および $A' \to C'$ が平坦で
あることは同値である。
\item $B' \to D'$ が平坦かつ有限表示であることと、$B \to D$ および
$A' \to C'$ が平坦かつ有限表示であることは同値である。
\item $B' \to D'$ が滑らかであることと、$B \to D$ および $A' \to C'$ が滑らかで
あることは同値である。
\item $B' \to D'$ が \'etale であることと、$B \to D$ および $A' \to C'$ が
\'etale であることは同値である。
\end{enumerate}
さらに、$D'$ が平坦な $B'$-代数ならば、
$D' \to (D' \otimes_{B'} B) \times_{(D' \otimes_{B'} A)} (D' \otimes_{B'} A')$
は同型である。このようにして、平坦 $B'$-代数の圏は、上のような系
$(D, C', \varphi)$ のうち $D$ が $B$ 上平坦で $C'$ が $A'$ 上平坦であるものの圏と
同値である。
\end{lemma}

\begin{proof}
含意「$\Rightarrow$」は、Algebra, Lemmas
\ref{algebra-lemma-base-change-finiteness},
\ref{algebra-lemma-flat-base-change},
\ref{algebra-lemma-base-change-smooth}, \ref{algebra-lemma-etale} から従う。
実際、Lemma \ref{more-algebra-lemma-module-over-fibre-product} により
$D' \otimes_{B'} B = D$ および $D' \otimes_{B'} A' = C'$ である。
従って、逆向きの含意を証明すれば十分である。

\medskip\noindent
(1) について。$D$ が $B$ 上有限型で、$C'$ が $A'$ 上有限型であると仮定する。
% SOURCE-ERRATUM-CANDIDATE P02-E0759: frozen proof omits “is” in both finite-type clauses.
以下では Lemma \ref{more-algebra-lemma-module-over-fibre-product} の結果を断りなく用いる。
$B$ 上の $D$ の生成元 $x_1, \ldots, x_r$ と、$A'$ 上の $C'$ の生成元
$y_1, \ldots, y_s$ を選ぶ。$D = D' \otimes_{B'} B$ であり
$B' \to B$ は全射なので、$D$ で $x_1, \ldots, x_r$ に写る
$u_1, \ldots, u_r \in D'$ を選べる。$C' = D' \otimes_{B'} A'$ なので、
ある $a'_{ij} \in A'$ によって $y_i = \sum v_j \otimes a'_{ij}$ となる
$v_1, \ldots, v_t \in D'$ を選べる。特に $v_j$ の $C'$ における像は
$A'$-代数として $C'$ を生成する。$N = r + t$ とおき、環の立方体
$$
\xymatrix{
A[x_1, \ldots, x_N] & & A'[x_1, \ldots, x_N] \ar[ll] \\
& A \ar[lu] & & A' \ar[ll] \ar[lu] \\
B[x_1, \ldots, x_N] \ar[uu] & & B'[x_1, \ldots, x_N] \ar[uu] \ar[ll] \\
& B \ar[uu] \ar[lu] & & B' \ar[ll] \ar[uu] \ar[lu]
}
$$
を考える。背面の正方形もデカルトであることに注意せよ。環準同型
$$
B'[x_1, \ldots, x_N] \to D',\quad
x_i \mapsto u_i \quad\text{and}\quad x_{r + j} \mapsto v_j.
$$
を考える。このとき誘導される写像 $B[x_1, \ldots, x_N] \to D$ および
$A'[x_1, \ldots, x_N] \to C'$ は全射、特に有限である。
Lemma \ref{more-algebra-lemma-relative-finite-module-over-fibre-product} により
$B'[x_1, \ldots, x_N] \to D'$ は有限である。従って、例えば Algebra, Lemma
\ref{algebra-lemma-compose-finite-type} により、$D'$ は $B'$ 上有限型である。

\medskip\noindent
(2) について。含意「$\Leftarrow$」は
Lemma \ref{more-algebra-lemma-relative-flat-module-over-fibre-product} から従う。
さらに、最後の主張は Lemma
\ref{more-algebra-lemma-relative-flat-module-over-fibre-product} の最後の主張から従う。

\medskip\noindent
(3) について。$B \to D$ と $A' \to C'$ が平坦かつ有限表示であると仮定する。
$B' \to D'$ の平坦性は (2) で見た。(1) により $B' \to D'$ は有限型である。
全射 $B'[x_1, \ldots, x_N] \to D'$ を選ぶ。Algebra, Lemma
\ref{algebra-lemma-finite-presentation-independent} により、環 $D$ は
$B[x_1, \ldots, x_N]$-加群として有限表示であり、環 $C'$ は
$A'[x_1, \ldots, x_N]$-加群として有限表示である。
Lemma \ref{more-algebra-lemma-relative-finitely-presented-module-over-fibre-product} により、
$D'$ は $B'[x_1, \ldots, x_N]$-加群として有限表示である。すなわち
$B' \to D'$ は有限表示である。

\medskip\noindent
(4) について。$B \to D$ と $A' \to C'$ が滑らかであると仮定する。
(3) により $B' \to D'$ は平坦かつ有限表示である。Algebra, Lemma
\ref{algebra-lemma-flat-fibre-smooth} により、$B'$ 上の任意の体 $k$ に対して
$D' \otimes_{B'} k$ が滑らかであることを確認すれば十分である。
合成 $J \to B' \to k$ が零ならば、$B' \to k$ は
$B' \to B \to k$ と分解し、
$$
D' \otimes_{B'} k = D' \otimes_{B'} B \otimes_B k
= D \otimes_B k
$$
を得る。$B \to D$ は滑らかなので、これは滑らかである。
合成 $J \to B' \to k$ が零でなければ、$k$ で零に写らない $h \in J$ が存在する。
このとき $B' \to k$ は $B' \to B'_h \to k$ と分解する。$h$ は $B$ で零に写るので
$B_h = 0$ である。従って Lemma \ref{more-algebra-lemma-diagram-localize} により
$B'_h = A'_h$ であり、
$$
D' \otimes_{B'} k = D' \otimes_{B'} B'_h \otimes_{B'_h} k
= C'_h \otimes_{A'_h} k
$$
を得る。$A' \to C'$ は滑らかなので、これは滑らかである。

\medskip\noindent
(5) について。$B \to D$ と $A' \to C'$ が \'etale であると仮定する。
(4) により $B' \to D'$ は滑らかである。滑らかな写像が \'etale であるか否かは
ファイバーの次元から読み取れるので、(5) が従う
（ファイバーを同定するには (4) の証明と同様に論じればよい。詳細の一部は省略する）。
\end{proof}

\begin{remark}
\label{more-algebra-remark-relative-modules-over-fibre-product}
Situation \ref{more-algebra-situation-relative-module-over-fibre-product} において、
$B' \to D'$ は有限表示であり、$D'$-加群 $L'$ が与えられていると仮定する。
次の二つの間に全単射な対応があると主張する。
\begin{enumerate}
\item $Q'$ が $D'$ 上有限表示かつ $B'$ 上平坦であるような $D'$-加群の全射
$L' \to Q'$。
\item 加群の全射の対
$(L' \otimes_{D'} D \to Q_1, L' \otimes_{D'} C' \to Q_2)$ で、次を満たすもの。
\begin{enumerate}
\item $Q_1$ は $D$ 上有限表示であり、$B$ 上平坦である。
\item $Q_2$ は $C'$ 上有限表示であり、$A'$ 上平坦である。
\item $L' \otimes_{D'} C$ の商として
$Q_1 \otimes_D C = Q_2 \otimes_{C'} C$ である。
\end{enumerate}
\end{enumerate}
% SOURCE-ERRATUM-CANDIDATE P02-E0760: frozen source switches from Q' to undefined Q and omits “is”; preserved.
これらの対応は $Q \mapsto (Q_1, Q_2)$ により与えられる。ここで
$Q_1 = Q \otimes_{D'} D$ および $Q_2 = Q \otimes_{D'} C'$ である。
逆向きには $Q = Q_1 \times_{Q_{12}} Q_2$ を用いる。ここで $Q_{12}$ は
$L' \otimes_{D'} C$ の共通の商
$Q_1 \otimes_D C = Q_2 \otimes_{C'} C$ である。商写像として
$$
L' \longrightarrow
(L' \otimes_{D'} D) \times_{(L' \otimes_{D'} C)} (L' \otimes_{D'} C')
\longrightarrow Q_1 \times_{Q_{12}} Q_2 = Q
$$
を用いる。最初の矢印は Lemma \ref{more-algebra-lemma-module-over-fibre-product-bis} により全射で、
第二の矢印は Lemma \ref{more-algebra-lemma-surjection-module-over-fibre-product} により全射である。
主張は Lemmas \ref{more-algebra-lemma-relative-flat-module-over-fibre-product} および
\ref{more-algebra-lemma-relative-finitely-presented-module-over-fibre-product} から従う。
\end{remark}














\section{Fitting イデアル}
\label{more-algebra-section-fitting-ideals}

\noindent
有限加群の Fitting イデアルとは、Lemma \ref{more-algebra-lemma-fitting-ideal}
の構成によって定まるイデアルである。

\begin{lemma}
\label{more-algebra-lemma-ideals-generated-by-minors}
$R$ を環とする。$A$ を成分が $R$ に属する $n \times m$ 行列とする。
$I_r(A)$ を $A$ の $r \times r$ 小行列式が生成するイデアルとし、
$I_0(A) = R$、また $r > \min(n, m)$ のとき $I_r(A) = 0$ と約束する。
このとき
\begin{enumerate}
\item $I_0(A) \supset I_1(A) \supset I_2(A) \supset \ldots$,
\item $B$ が $(n + n') \times m$ 行列で、$A$ が $B$ の最初の
$n$ 行からなるならば、$I_{r + n'}(B) \subset I_r(A)$,
\item $C$ が $n \times n$ 行列ならば、$I_r(CA) \subset I_r(A)$.
\item $A$ が分割行列
$$
\left(
\begin{matrix}
A_1 & 0 \\
0 & A_2 
\end{matrix}
\right)
$$
ならば、$I_r(A) = \sum_{r_1 + r_2 = r} I_{r_1}(A_1) I_{r_2}(A_2)$.
\item ここにさらに追加する。
\end{enumerate}
\end{lemma}

\begin{proof}
省略する。（ヒント：行列式は一つの列または行に沿って展開することで計算できる。）
\end{proof}

\begin{lemma}
\label{more-algebra-lemma-fitting-ideal}
$R$ を環とし、$M$ を有限 $R$-加群とする。$M$ の表示
$$
\bigoplus\nolimits_{j \in J} R \longrightarrow R^{\oplus n}
\longrightarrow M \longrightarrow 0.
$$
を選ぶ。$A = (a_{ij})_{i = 1, \ldots, n, j \in J}$ を写像
$\bigoplus_{j \in J} R \to R^{\oplus n}$ の行列とする。
$A$ の $(n - k) \times (n - k)$ 小行列式が生成するイデアル
$\text{Fit}_k(M)$ は、表示の選び方によらない。
\end{lemma}

\begin{proof}
$K \subset R^{\oplus n}$ を全射 $R^{\oplus n} \to M$ の核とする。
$z_1, \ldots, z_{n - k} \in K$ を選び、
$z_j = (z_{1j}, \ldots, z_{nj})$ と書く。
イデアル $\text{Fit}_k(M)$ は、こうして得られるすべての行列
$(z_{ij})$ の $(n - k) \times (n - k)$ 小行列式によって生成される
イデアルとしても記述できる。

\medskip\noindent
この全射を、核が $K'$ である全射 $R^{\oplus n + n'} \to M$ に取り替えるとする。
ここで $R^{\oplus n + n'}$ の最初の $n$ 個の標準基底ベクトル上では元の写像を用い、
最後の $n'$ 個の基底ベクトル上では $0$ とする。このとき、対応するイデアルは等しい。
実際、上のように $z_1, \ldots, z_{n - k} \in K$ をとり、
$j = 1, \ldots, n - k$ に対して
$z'_j = (z_{1j}, \ldots, z_{nj}, 0, \ldots, 0) \in K'$ と置き、さらに
$z'_{n + j'} = (0, \ldots, 0, 1, 0, \ldots, 0) \in K'$ と置く。
すると、$(z_{ij})$ の $(n - k) \times (n - k)$ 小行列式が生成するイデアルは、
$(z'_{ij})$ の $(n + n' - k) \times (n + n' - k)$ 小行列式が生成する
イデアルと一致する。これで一方の包含が得られる。
逆に、$K'$ の元 $z'_1, \ldots, z'_{n + n' - k}$ が与えられたとする。
これらを $R^{\oplus n}$ に射影すると、$K$ の元
$z_1, \ldots, z_{n + n' - k}$ が得られる。Lemma
\ref{more-algebra-lemma-ideals-generated-by-minors} により、$(z'_{ij})$ の
$(n + n' - k) \times (n + n' - k)$ 小行列式が生成するイデアルは、
$(z_{ij})$ の $(n - k) \times (n - k)$ 小行列式が生成するイデアルに含まれる。
これでもう一方の包含が得られる。

\medskip\noindent
$R^{\oplus m} \to M$ を、核が $L$ である別の全射とする。
Schanuel の補題（Algebra, Lemma \ref{algebra-lemma-Schanuel}）と前段の結果により、
$m = n$ であり、かつ $M$ への全射と可換する同型
$R^{\oplus n} \to R^{\oplus m}$ があると仮定してよい。
$C = (c_{li})$ をこの写像の（可逆な）行列とする
（$n = m$ なので、これは正方行列である）。すると、上のように
$z'_1, \ldots, z'_{n - k} \in L$ が与えられたとき、
$z_1, \ldots, z_{n - k} \in K$ で
$z_1' = Cz_1, \ldots, z'_{n - k} = Cz_{n - k}$ を満たすものがとれる。
Lemma \ref{more-algebra-lemma-ideals-generated-by-minors} により一方の包含が得られ、
対称性によりもう一方が得られる。
\end{proof}

\begin{definition}
\label{more-algebra-definition-fitting-ideal}
$R$ を環、$M$ を有限 $R$-加群、$k \geq 0$ とする。
$M$ の {\it 第 $k$ Fitting イデアル}とは、Lemma
\ref{more-algebra-lemma-fitting-ideal} で構成したイデアル $\text{Fit}_k(M)$ のことである。
$\text{Fit}_{-1}(M) = 0$ と置く。
\end{definition}

\noindent
Fitting イデアルは一つの大きな行列の小行列式イデアルであり
（Lemma \ref{more-algebra-lemma-ideals-generated-by-minors} における順序とは逆向きに番号が付いている）、
ある $t \gg 0$ に対して
$$
0 = \text{Fit}_{-1}(M) \subset \text{Fit}_0(M) \subset \text{Fit}_1(M)
\subset \ldots \subset \text{Fit}_t(M) = R
$$
となることが分かる。以下に Fitting イデアルの基本的な性質を挙げる。

\begin{lemma}
\label{more-algebra-lemma-fitting-ideal-basics}
$R$ を環とし、$M$ を有限 $R$-加群とする。
\begin{enumerate}
\item $M$ が $n$ 個の元で生成できるならば、$\text{Fit}_n(M) = R$.
\item 二つ目の有限 $R$-加群 $M'$ が与えられたとき、
$\text{Fit}_0(M \oplus M') = \text{Fit}_0(M)\text{Fit}_0(M')$ であり、
より一般に
$$
\text{Fit}_l(M \oplus M') =
\sum\nolimits_{k + k' = l} \text{Fit}_k(M)\text{Fit}_{k'}(M')
$$
\item $R \to R'$ が環写像ならば、$\text{Fit}_k(M \otimes_R R')$ は
$\text{Fit}_k(M)$ の像が $R'$ において生成するイデアルである。
\item $M$ が有限表示ならば、$\text{Fit}_k(M)$ は有限生成イデアルである。
\item $M \to M'$ が全射ならば、
$\text{Fit}_k(M) \subset \text{Fit}_k(M')$.
\item $\text{Fit}_0(M) \subset \text{Ann}_R(M)$ である。
\item $V(\text{Fit}_0(M)) = \text{Supp}(M)$ である。
\item $I$ が $R$ のイデアルならば、$\text{Fit}_0(R/I) = I$.
\item ここにさらに追加する。
\end{enumerate}
\end{lemma}

\begin{proof}
(1) は Lemma \ref{more-algebra-lemma-ideals-generated-by-minors} における
$I_0(A) = R$ という事実から従う。

\medskip\noindent
% SOURCE-ERRATUM-CANDIDATE P02-E0761: frozen source has “follows form”; Japanese supplies the intended postposition.
(2) は Lemma \ref{more-algebra-lemma-ideals-generated-by-minors} の対応する主張から従う。

\medskip\noindent
(3) は $\otimes_R R'$ が右完全であることから従う。従って、$M$ の表示を
基底変換して得られる列は $M \otimes_R R'$ の表示である。

\medskip\noindent
(4) の証明。表示
$R^{\oplus m} \xrightarrow{A} R^{\oplus n} \to M \to 0$
をとる。このとき $\text{Fit}_k(M)$ は、行列 $A$ の
$(n - k) \times (n - k)$ 小行列式が生成するイデアルである。

\medskip\noindent
(5) は定義から直ちに従う。

\medskip\noindent
(6) の証明。Lemma \ref{more-algebra-lemma-fitting-ideal} のように、行列 $A$ をもつ
$M$ の表示を選ぶ。$J' \subset J$ を濃度 $n$ の部分集合とする。
$f = \det(a_{ij})_{i = 1, \ldots, n, j \in J'}$ が $M$ を零化することを
示せば十分である。これは明らかである。というのも、
$$
R^{\oplus n} \xrightarrow{A' = (a_{ij})_{i = 1, \ldots, n, j \in J'}}
R^{\oplus n} \to M \to 0
$$
の余核は $f$ によって零化されるからである。実際、
$A' B = f1_{n \times n}$ を満たす行列 $B$ が存在する。

\medskip\noindent
(7) の証明。Lemma \ref{more-algebra-lemma-fitting-ideal} のように、行列 $A$ をもつ
$M$ の表示を選ぶ。Nakayama の補題（Algebra, Lemma
\ref{algebra-lemma-NAK}）により
$$
M_\mathfrak p \not = 0
\Leftrightarrow
M \otimes_R \kappa(\mathfrak p) \not = 0
\Leftrightarrow
\text{rank}(\text{image }A\text{ in }\kappa(\mathfrak p)) < n
$$
である。$\text{Fit}_0(M)$ が、この性質をもつ素イデアルの集合をちょうど
切り出すことは明らかである。
\end{proof}

\begin{example}
\label{more-algebra-example-fitting-free}
$R$ を環とする。有限自由加群 $M = R^{\oplus n}$ の Fitting イデアルは、
$k < n$ のとき $\text{Fit}_k(M) = 0$ であり、$k \geq n$ のとき
$\text{Fit}_k(M) = R$ である。
\end{example}


\begin{example}
\label{more-algebra-example-laurent}
$R$ を環とし、$I, J\subset R$ をイデアルとする。このとき
$\text{Fit}_0(R/I) = I$, $\text{Fit}_0(R/I \oplus R/J) = IJ$,
$\text{Fit}_1(R/I\oplus R/J) = I+J$.
\end{example}

\begin{lemma}
\label{more-algebra-lemma-fitting-ideal-generate-locally}
$R$ を環、$M$ を有限 $R$-加群、$k \geq 0$ とする。
$\mathfrak p \subset R$ を素イデアルとする。次の条件は同値である。
\begin{enumerate}
\item $\text{Fit}_k(M) \not \subset \mathfrak p$,
\item $\dim_{\kappa(\mathfrak p)} M \otimes_R \kappa(\mathfrak p) \leq k$,
\item $M_\mathfrak p$ は $R_\mathfrak p$ 上 $k$ 個の元で生成できる、かつ
\item ある $f \in R$, $f \not \in \mathfrak p$ に対して、$M_f$ は
$R_f$ 上 $k$ 個の元で生成できる。
\end{enumerate}
\end{lemma}

\begin{proof}
Nakayama の補題（Algebra, Lemma \ref{algebra-lemma-NAK}）により、
$M \otimes_R \kappa(\mathfrak p)$ が $k$ 個の元で生成できるならば、
ある $f \in R$, $f \not \in \mathfrak p$ に対して $M_f$ は
$R_f$ 上 $k$ 個の元で生成できる。従って (2), (3), (4) は同値である。
Lemma \ref{more-algebra-lemma-fitting-ideal-basics} の (3) を用いると、問題は
$R$ が体で $\mathfrak p = (0)$ である場合に帰着する。この場合の結果は
Example \ref{more-algebra-example-fitting-free} から従う。
\end{proof}

\begin{lemma}
\label{more-algebra-lemma-fitting-ideal-finite-locally-free}
$R$ を環、$M$ を有限 $R$-加群、$r \geq 0$ とする。
次の条件は同値である。
\begin{enumerate}
\item $M$ は rank $r$ の有限局所自由加群である
（Algebra, Definition \ref{algebra-definition-locally-free}）、
\item $\text{Fit}_{r - 1}(M) = 0$ かつ $\text{Fit}_r(M) = R$、かつ
\item $k < r$ に対して $\text{Fit}_k(M) = 0$ であり、$k \geq r$ に対して
$\text{Fit}_k(M) = R$ である。
\end{enumerate}
\end{lemma}

\begin{proof}
Fitting イデアルは増大するイデアル列をなすので、(2) と (3) が同値であることは
直ちに分かる。$\text{Fit}_k(M)$ の形成は基底変換と可換するので
（Lemma \ref{more-algebra-lemma-fitting-ideal-basics}）、Example
\ref{more-algebra-example-fitting-free} と貼り合わせに関する結果
（Algebra, Section \ref{algebra-section-more-glueing}）から (1) は (2) を含意する。
逆に (2) を仮定する。Lemma \ref{more-algebra-lemma-fitting-ideal-generate-locally} により、
$M$ は $r$ 個の元で生成されると仮定してよい。
% SOURCE-ERRATUM-CANDIDATE P02-E0762: frozen source has the sentence fragment “Thus a presentation”; Japanese supplies the missing predicate.
従って表示 $\bigoplus_{j \in J} R \to R^{\oplus r} \to M \to 0$ がとれる。
ここで $\text{Fit}_{r - 1}(M) = 0$ という仮定から、写像
$\bigoplus_{j \in J} R \to R^{\oplus r}$ の行列のすべての成分は零である。
従って $M$ は自由である。
\end{proof}

\begin{lemma}
\label{more-algebra-lemma-principal-fitting-ideal}
$R$ を局所環、$M$ を有限 $R$-加群、$k \geq 0$ とする。
ある $f \in R$ に対して $\text{Fit}_k(M) = (f)$ と仮定する。
$M'$ を $M$ の $\{x \in M \mid fx = 0\}$ による商とする。
このとき $M'$ は $k$ 個の元で生成できる。
\end{lemma}

\begin{proof}
全射 $R^{\oplus n} \to M$ に対応する生成元 $x_1, \ldots, x_n \in M$ を選ぶ。
$R$ は局所環なので、元の集合 $E \subset (f)$ が $(f)$ を生成するならば、
ある $e \in E$ が $(f)$ を生成する。Algebra, Lemma
\ref{algebra-lemma-NAK} を参照せよ。従って、$R^{\oplus n} \to M$ の核から
$z_1, \ldots, z_{n - k}$ を、$n \times (n - k)$ 行列
$A = (z_{ij})$ のある $(n - k) \times (n - k)$ 小行列式が $(f)$ を
生成するように選べる。$x_i$ の番号を付け替えることにより、最初の小行列式
$\det(z_{ij})_{1 \leq i, j \leq n - k}$ が $(f)$ を生成すると仮定してよい。
すなわち、ある単元 $u \in R$ に対して
$\det(z_{ij})_{1 \leq i, j \leq n - k} = uf$ である。
他のすべての小行列式は $f$ の倍数である。Algebra, Lemma
\ref{algebra-lemma-matrix-right-inverse} により、ある
$n - k \times n - k$ 行列 $B$ が存在して
$$
AB = f
\left(
\begin{matrix}
u 1_{n - k \times n - k} \\
C
\end{matrix}
\right)
$$
となる。ここで $C$ は成分が $R$ に属するある行列である。従って、すべての
$i \leq n - k$ に対して、元 $y_i = ux_i + \sum_j c_{ji}x_j$ は $f$ により
零化される。$M/\sum Ry_i$ は $x_{n - k + 1}, \ldots, x_n$ の像によって
生成されるので、主張が従う。
\end{proof}

\begin{lemma}
\label{more-algebra-lemma-fitting-ideals-and-pd1}
$R$ を環、$M$ を有限表示 $R$-加群、$k \geq 0$ とする。
ある非零因子 $f \in R$ に対して $\text{Fit}_k(M) = (f)$ であり、
$\text{Fit}_{k - 1}(M) = 0$ であると仮定する。このとき
\begin{enumerate}
\item $M$ の射影次元は $\leq 1$ である、
\item $M' = \Ker(f : M \to M)$ は $M$ の $f$-冪ねじれ部分加群である、
\item $M'$ の射影次元は $\leq 1$ である、
\item $M/M'$ は rank $k$ の有限局所自由加群である、かつ
\item $M \cong M/M' \oplus M'$.
\end{enumerate}
\end{lemma}

\begin{proof}
成分が $R$ に属するある行列 $A$ による表示
$$
R^{\oplus m} \xrightarrow{A} R^{\oplus n} \to M \to 0
$$
を選ぶ。

\medskip\noindent
まず $R$ が局所環である場合に補題を証明する。主張におけるように
$M' = \{x \in M \mid fx = 0\}$ と置く。Lemma
\ref{more-algebra-lemma-principal-fitting-ideal} により、$M/M'$ を生成する
$x_1, \ldots, x_k \in M$ を選べる。すると $x_1, \ldots, x_k$ は
$M_f = (M/M')_f$ を生成する。従って、$M$ において関係
$\sum a_ix_i = 0$ があるならば、$a_1, \ldots, a_k$ の $R_f$ における像は
零である。そうでなければ
$\text{Fit}_{k - 1}(M) R_f = \text{Fit}_{k - 1}(M_f)$ が非零となるからである。
$f$ は非零因子なので、$a_1 = \ldots = a_k = 0$ と結論できる。
従って $M \cong R^{\oplus k} \oplus M'$ である。上の表示の基底を取り替えることで、
$R^{\oplus n}$ の最初の $n - k$ 個の基底ベクトルは $M$ の直和因子 $M'$ に写り、
% SOURCE-ERRATUM-CANDIDATE P02-E0763: frozen source hyphenates “last k-basis vectors”; Japanese expresses the intended noun phrase.
$R^{\oplus n}$ の最後の $k$ 個の基底ベクトルは $M$ の直和因子
$R^{\oplus k}$ の基底元に写ると
仮定してよい。こうすると、行列 $A$ の最後の $k$ 行は零になる。
このようにして、$M$ を $M'$ に、$k$ を $0$ に、$n$ を $n - k$ に、
$A$ を最後の $k$ 行を削除した小行列にそれぞれ取り替えることにより、
次段で論じる場合に帰着する。

\medskip\noindent
$R$ が局所環、$k = 0$ であり、$M$ が $f$ により零化されると仮定する。
このとき $M$ の第 $0$ Fitting イデアルは $(f)$ であり、大きさ
$n \times m$ の行列 $A$ の $n \times n$ 小行列式によって生成される
（特に、これは $m \geq n$ を含意する）。$R$ は局所環なので、$A$ のある
$n \times n$ 小行列式は、単元 $u \in R$ に対して $uf$ である。
番号を付け替えて、これが最初の小行列式であると仮定してよい。
さらに、$A$ の他のすべての $n \times n$ 小行列式は $f$ で割り切れる。
$A = (A_1 A_2)$ と分割して書く。ここで $A_1$ は $n \times n$ 行列、
$A_2$ は $n \times (m - n)$ 行列である。Algebra, Lemma
\ref{algebra-lemma-matrix-right-inverse} を $A$ の転置に適用すると（！）、
ある $n \times n$ 行列 $B$ が存在して
$$
BA = B(A_1 A_2) = f
\left(
\begin{matrix}
u 1_{n \times n} & C
\end{matrix}
\right)
$$
となる。ここで $C$ は成分が $R$ に属する $n \times (m - n)$ 行列である。
まず $BA_1 = fu 1_{n \times n}$ が従う。従って
$$
BA_2 = fC = u^{-1}fuC = u^{-1}BA_1C
$$
である。$B$ の行列式は非零因子なので、$A_2 = u^{-1}A_1C$ と結論できる。
従って $A$ の像は $A_1$ の像に等しく、後者は $A_1$ の行列式が非零因子であるため
$R^{\oplus n}$ と同型である。よって $M$ の射影次元は $\leq 1$ である。

\medskip\noindent
一般の環 $R$ の場合に戻る。局所的な場合から、$M/M'$ は rank $k$ の
有限局所自由加群であることが分かる。Algebra, Lemma
\ref{algebra-lemma-finite-projective} を参照せよ。従って完全列
$0 \to M' \to M \to M/M' \to 0$ は分裂する。これより $M'$ は有限表示加群である。
短完全列 $0 \to K \to R^{\oplus a} \to M' \to 0$ を選ぶ。
このとき $K$ は有限 $R$-加群である。Algebra, Lemma
\ref{algebra-lemma-extension} を参照せよ。局所的な場合から、すべての素イデアルに対して
$K_\mathfrak p \cong R_\mathfrak p^{\oplus a}$ である。従って再び Algebra, Lemma
\ref{algebra-lemma-finite-projective} により、$K$ は rank $a$ の有限局所自由加群である。
これより $M'$ の射影次元は $\leq 1$ であり、補題が証明された。
\end{proof}

\section{持ち上げ}
\label{more-algebra-section-lifting}

\noindent
% SOURCE-ERRATUM-CANDIDATE P02-E0764: frozen source says “we collection some lemmas”; Japanese supplies the intended verb.
この節では、次の型の持ち上げに関する主張を扱う補題をいくつか集める。
$A$ が環、$I \subset A$ がイデアルで、$\overline{\xi}$ が $A/I$ 上のある種の構造ならば、
$\overline{\xi}$ を $A$ 上、または $A$ のある \'etale 拡大上の同種の構造 $\xi$ に
持ち上げることができる。すでにいくつかの結果を証明した構造の種類を以下に挙げる。
\begin{enumerate}
\item 冪等元。Algebra, Lemmas \ref{algebra-lemma-lift-idempotents} および
\ref{algebra-lemma-lift-idempotents-noncommutative} を参照せよ、
\item 射影加群。Algebra, Lemmas \ref{algebra-lemma-lift-projective-module} および
\ref{algebra-lemma-lift-finite-projective-module} を参照せよ、
\item 有限安定自由加群。Lemma \ref{more-algebra-lemma-lift-stably-free} を参照せよ、
\item 基底元。Algebra, Lemmas \ref{algebra-lemma-local-artinian-basis-when-flat} および
\ref{algebra-lemma-lift-basis} を参照せよ、
\item 環写像、すなわち、ある代数が形式的に滑らかであることの証明。Algebra, Lemma
\ref{algebra-lemma-polynomial-ring-formally-smooth}, Proposition
\ref{algebra-proposition-smooth-formally-smooth}, および Lemma
\ref{algebra-lemma-smooth-strong-lift} を参照せよ、
\item syntomic 環写像。Algebra, Lemma \ref{algebra-lemma-lift-syntomic} を参照せよ、
\item 滑らかな環写像。Algebra, Lemma \ref{algebra-lemma-lift-smooth} を参照せよ、
\item \'etale 環写像。Algebra, Lemma \ref{algebra-lemma-lift-etale} を参照せよ、
\item 多項式の因数分解。Algebra, Lemma
\ref{algebra-lemma-factor-mod-lift-etale} を参照せよ、かつ
\item Algebra, Section \ref{algebra-section-henselian} では Hensel 的局所環を論じる。
\end{enumerate}
関心のある読者は、特に Smoothing Ring Maps, Section
\ref{smoothing-section-presentations} および Smoothing Ring Maps, Proposition
\ref{smoothing-proposition-lift-smooth} に、この種の結果をさらに見出すだろう。

\medskip\noindent
$A$ を環、$I \subset A$ をイデアルとし、$\overline{\xi}$ を $A/I$ 上の
ある種の構造とする。以下の補題では、同型 $A/I \to A'/IA'$ を誘導する
\'etale 環写像 $A \to A'$ と、$\overline{\xi}$ を持ち上げる $A'$ 上の対象
$\xi'$ を探す。一般的な注意として、$A/I \cong A'/IA'$ および
$A'/IA' \cong A''/IA''$ を満たす \'etale 環写像
$A \to A' \to A''$ が与えられると、合成 $A \to A''$ も \'etale であり
（Algebra, Lemma \ref{algebra-lemma-etale}）、さらに
$A/I \cong A''/IA''$ を満たす。以下の補題では、このことを断らずに頻繁に用いる。
まず、求めている型の結果の自明な例を挙げる。

\begin{lemma}
\label{more-algebra-lemma-lift-invertible-element}
$A$ を環、$I \subset A$ をイデアルとし、$\overline{u} \in A/I$ を可逆元とする。
同型 $A/I \to A'/IA'$ を誘導する \'etale 環写像 $A \to A'$ と、
$\overline{u}$ を持ち上げる可逆元 $u' \in A'$ が存在する。
\end{lemma}

\begin{proof}
% SOURCE-ERRATUM-CANDIDATE P02-E0765: frozen source defines u although the statement and intended lift use u'; Japanese preserves that variable switch.
$\overline{u}$ の任意の持ち上げ $f \in A$ を選び、$A' = A_f$ と置き、
$u$ を $A'$ における $f$ の像とする。
\end{proof}

\begin{lemma}
\label{more-algebra-lemma-lift-idempotent}
$A$ を環、$I \subset A$ をイデアルとし、$\overline{e} \in A/I$ を冪等元とする。
同型 $A/I \to A'/IA'$ を誘導する \'etale 環写像 $A \to A'$ と、
$\overline{e}$ を持ち上げる冪等元 $e' \in A'$ が存在する。
\end{lemma}

\begin{proof}
$\overline{e}$ の任意の持ち上げ $x \in A$ を選ぶ。次のように置く。
$$
A' = A[t]/(t^2 - t)\left[\frac{1}{t - 1 + x}\right].
$$
$(2t - 1)\text{d}t = 0$ かつ $(2t - 1)(2t - 1) = 1$ は可逆なので、
環写像 $A \to A'$ は \'etale である。さらに
$A'/IA' = A/I[t]/(t^2 - t)[\frac{1}{t - 1 + \overline{e}}] \cong A/I$
である。最後の写像は $t$ を $\overline{e}$ に送り、$\overline{e}$ は
$t^2 - t$ の根なので、これはうまくいく。また、$e'$ を $A'$ における
$t$ の類とすればよいことも分かる。
\end{proof}

\begin{lemma}
\label{more-algebra-lemma-lift-open-covering}
$A$ を環、$I \subset A$ をイデアルとする。
$\Spec(A/I) = \coprod_{j \in J} \overline{U}_j$ を有限の互いに素な開被覆とする。
このとき、同型 $A/I \to A'/IA'$ を誘導する \'etale 環写像 $A \to A'$ と、
与えられた被覆を持ち上げる有限の互いに素な開被覆
$\Spec(A') = \coprod_{j \in J} U'_j$ が存在する。
\end{lemma}

\begin{proof}
これは Lemma \ref{more-algebra-lemma-lift-idempotent} と、スペクトルの開かつ閉な部分集合が
冪等元に対応するという事実から従う。Algebra, Lemma
\ref{algebra-lemma-disjoint-decomposition} を参照せよ。
\end{proof}

\begin{lemma}
\label{more-algebra-lemma-localize-upstairs}
$A \to B$ を環写像、$J \subset B$ をイデアルとする。
$A \to B$ が $V(J)$ のすべての素イデアルで \'etale ならば、
$B/J$ で可逆元に写る $g \in B$ で、$A' = B_g$ が $A$ 上 \'etale となるものが存在する。
\end{lemma}

\begin{proof}
$A \to B$ が \'etale でない $\Spec(B)$ の点の集合は $\Spec(B)$ の閉集合である。
Algebra, Definition \ref{algebra-definition-etale} を参照せよ。
これをあるイデアル $J' \subset B$ に対する $V(J')$ と書く。このとき
$V(J') \cap V(J) = \emptyset$ なので、Algebra, Lemma
\ref{algebra-lemma-Zariski-topology} により $J + J' = B$ である。
$f \in J$, $g \in J'$ によって $1 = f + g$ と書けば、この $g$ が求めるものである。
\end{proof}

\noindent
次の三つの補題は、多項式の因数分解を持ち上げられることを述べる。

\begin{lemma}
\label{more-algebra-lemma-lift-factorization-monic}
$A$ を環、$I \subset A$ をイデアルとする。$f \in A[x]$ をモニック多項式とする。
$\overline{f} = \overline{g} \overline{h}$ を $A/I[x]$ における $f$ の因数分解とし、
$\overline{g}$ と $\overline{h}$ はモニックで $A/I[x]$ の単位イデアルを生成すると仮定する。
このとき、同型 $A/I \to A'/IA'$ を誘導する \'etale 環写像 $A \to A'$ と、
与えられた $A/I$ 上の因数分解を持ち上げる、$g'$, $h'$ がモニックである
$A'[x]$ における因数分解 $f = g' h'$ が存在する。
\end{lemma}

\begin{proof}
これを、先に証明した普遍因数分解に関する結果から導く。ただし、この技巧を使わない
独自の証明を見つけることを読者に勧める。
$\deg(\overline{g}) = n$, $\deg(\overline{h}) = m$ とする。従って
$\deg(f) = n + m$ である。ある $\alpha_1, \ldots, \alpha_{n + m} \in A$ によって
$f = x^{n + m} + \sum \alpha_i x^{n + m - i}$ と書く。Algebra, Example
\ref{algebra-example-factor-polynomials-etale} の環写像
$$
R = \mathbf{Z}[a_1, \ldots, a_{n + m}]
\longrightarrow
S = \mathbf{Z}[b_1, \ldots, b_n, c_1, \ldots, c_m]
$$
を考える。$a_i$ を $\alpha_i$ に送る環写像を $R \to A$ とする。次のように置く。
$$
B = A \otimes_R S
$$
構成により、$B[x]$ における $f$ の像 $f_B$ は因数分解する。
$f_B = g_B h_B$ と書く。ここで
$g_B = x^n + \sum (1 \otimes b_i) x^{n - i}$ であり、$h_B$ についても同様である。
$\overline{g} = x^n + \sum \overline{\beta}_i x^{n - i}$ および
$\overline{h} = x^m + \sum \overline{\gamma}_i x^{m - i}$ と書く。
$A$-代数写像
$$
B \longrightarrow A/I, \quad
1 \otimes b_i \mapsto \overline{\beta}_i, \quad
1 \otimes c_i \mapsto \overline{\gamma}_i
$$
は、$g_B$ と $h_B$ をそれぞれ $A/I[x]$ の $\overline{g}$ と $\overline{h}$ に送る。
この表示された写像は全射である。その核を $J \subset B$ と書く。
Algebra, Example \ref{algebra-example-factor-polynomials-etale} の議論から、
% SOURCE-ERRATUM-CANDIDATE P02-E0766: frozen source drops the TeX accent in “etale”; Japanese uses the established spelling.
$A \to B$ が $V(J) \subset \Spec(B)$ のすべての点で \'etale であることは明らかである。
Lemma \ref{more-algebra-lemma-localize-upstairs} のように $g \in B$ を選び、$A$-代数 $B_g$ を考える。
$g$ は $B/J = A/I$ において単元に写るので、$A/I$-代数の写像
$B_g/I B_g \to A/I$ も得られる。$A/I \to B_g/I B_g$ は \'etale なので、
$B_g/IB_g \to A/I$ も \'etale である
（Algebra, Lemma \ref{algebra-lemma-map-between-etale}）。従って、
$A/I = (B_g/I B_g)_e$ となる冪等元 $e \in B_g/I B_g$ が存在する
（Algebra, Lemma \ref{algebra-lemma-surjective-flat-finitely-presented}）。
$e$ の持ち上げ $h \in B_g$ を選ぶ。このとき $A \to A' = (B_g)_h$ と、
$A'$ における因数分解 $f_B = g_B h_B$ の像が与える因数分解とが、
補題の問題の解である。
\end{proof}

\begin{lemma}
\label{more-algebra-lemma-lift-factorization-easy}
$A$ を環、$I \subset A$ をイデアルとする。$f \in A[x]$ をモニック多項式とする。
$\overline{f} = \overline{g} \overline{h}$ を $A/I[x]$ における $f$ の因数分解とし、
次を仮定する。
\begin{enumerate}
\item $\overline{g}$ の最高次係数は $A/I$ の可逆元である、かつ
\item $\overline{g}$, $\overline{h}$ は $A/I[x]$ の単位イデアルを生成する。
\end{enumerate}
このとき、同型 $A/I \to A'/IA'$ を誘導する \'etale 環写像 $A \to A'$ と、
与えられた $A/I$ 上の因数分解を持ち上げる、$A'[x]$ における因数分解
$f = g' h'$ が存在する。
\end{lemma}

\begin{proof}
Lemma \ref{more-algebra-lemma-lift-invertible-element} を適用することで、$\overline{g}$ の
最高次係数は可逆元 $u \in A$ の剰余であると仮定してよい。そこで
$\overline{g}$ を $\overline{u}^{-1}\overline{g}$ に、$\overline{h}$ を
$\overline{u}\overline{h}$ に取り替えられる。従って $\overline{g}$ はモニックと
仮定してよい。$f$ はモニックなので、$\overline{h}$ もモニックである。
この場合、結果は Lemma \ref{more-algebra-lemma-lift-factorization-monic} から従う。
\end{proof}

\begin{example}
\label{more-algebra-example-cannot-lift-factorization}
この例は、Lemma \ref{more-algebra-lemma-lift-factorization-easy} における最高次係数に関する条件を
取り除けないことを示す。実際、$A = \mathbf{Z}$, $I = 4\mathbf{Z}$, $f = 1$ とし、
$A/I[x]$ において $\bar{g} = \bar{h} = 2x^2 + 2x + 1$ とする。
$A \to A'$ が \'etale で $A'/I A' = A/I$ ならば、$A'$ の $2$-進完備化は
$\mathbf{Z}_2$ と同型である。ところが、$4$ を法として $2x^2 + 2x + 1$ と合同な
$\mathbf{Z}_2[x]$ の任意の多項式 $g'$ は、$\mathbf{Z}_2[x]$ において多項式の
乗法逆元をもたない。従って補題の結論は成立しない。
\end{example}

\begin{lemma}
\label{more-algebra-lemma-separate-image-closed-from-closed}
$R \to S$ を環写像とする。$I \subset R$ を $R$ のイデアル、$J \subset S$ を
$S$ のイデアルとする。$\Spec(R)$ における $V(J)$ の像の閉包が $V(I)$ と
交わらないならば、$R/I$ で $1$ に写り、$S$ では $J$ の元に写る元
$f \in R$ が存在する。
\end{lemma}

\begin{proof}
$V(I')$ が $V(J)$ の像の閉包となるイデアル $I' \subset R$ をとる。
仮定により $V(I) \cap V(I') = \emptyset$ なので、Algebra, Lemma
\ref{algebra-lemma-Zariski-topology} により $I + I' = R$ である。
$g \in I$, $f \in I'$ によって $1 = g + f$ と書く。
$f'$ を $S$ における $f$ の像とすれば、$V(f') \supset V(J)$ である。
従って、ある $n$ に対して $(f')^n \in J$ である。Algebra, Lemma
\ref{algebra-lemma-Zariski-topology} を参照せよ。$f$ を $f^n$ に取り替えればよい。
\end{proof}

\begin{lemma}
\label{more-algebra-lemma-helper-integral}
$I$ を環 $A$ のイデアルとし、$A \to B$ を整な環写像とする。
$b \in B$ は $B/IB$ で冪等元に写るとする。このとき、ある $d \geq 1$ に対して
$f(b) = 0$ かつ $f \bmod I = x^d(x - 1)^d$ を満たすモニック多項式
$f \in A[x]$ が存在する。
\end{lemma}

\begin{proof}
$z = b^2 - b$ は $IB$ の元であることに注意する。Algebra, Lemma
\ref{algebra-lemma-integral-integral-over-ideal} により、次数 $d$ のモニック多項式
$g(x) = x^d + \sum a_j x^j$ で、$a_j \in I$ かつ $B$ において $g(z) = 0$ を
満たすものが存在する。従って $f(x) = g(x^2 - x) \in A[x]$ はモニックで、
$f(x) \equiv x^d(x - 1)^d \bmod I$ かつ $B$ において $f(b) = 0$ を満たす。
\end{proof}

\begin{lemma}
\label{more-algebra-lemma-lift-idempotent-upstairs}
$A$ を環、$I \subset A$ をイデアルとする。$A \to B$ を整な環写像とし、
$\overline{e} \in B/IB$ を冪等元とする。このとき、同型
$A/I \to A'/IA'$ を誘導する \'etale 環写像 $A \to A'$ と、
$\overline{e}$ を持ち上げる冪等元 $e' \in B \otimes_A A'$ が存在する。
\end{lemma}

\begin{proof}
$\overline{e}$ を持ち上げる元 $y \in B$ を選ぶ。$y$ に対して Lemma
\ref{more-algebra-lemma-helper-integral} のような $f \in A[x]$ を選ぶ。Lemma
\ref{more-algebra-lemma-lift-factorization-easy} により、同型 $A/I \to A'/IA'$ を誘導する
\'etale 環写像 $A \to A'$ であって、
% SOURCE-ERRATUM-CANDIDATE P02-E0767: frozen source places the lifted factorization in A[x] although it is obtained after A -> A'; preserved.
$A[x]$ において $f = gh$ となり、$g(x) = x^d \bmod IA'$ かつ
$h(x) = (x - 1)^d \bmod IA'$ となるものを見つけられる。
$A$ を $A'$ に取り替えることで、この因数分解は $A$ 上定義されていると仮定してよい。
この場合、$b_1 = g(y) \in B$ は $\overline{e}^d = \overline{e}$ の持ち上げであり、
$b_2 = h(y) \in B$ は
$(\overline{e} - 1)^d = (-1)^d (1 - \overline{e})^d = (-1)^d(1 - \overline{e})$
の持ち上げである。さらに $b_1b_2 = 0$ である。従って
$(b_1, b_2)B/IB = B/IB$ であり、$V(b_1, b_2) \subset \Spec(B)$ は
$V(IB)$ と交わらない。$\Spec(B) \to \Spec(A)$ は閉写像なので
（Algebra, Lemmas \ref{algebra-lemma-integral-going-up} および
\ref{algebra-lemma-going-up-closed} を参照）、$A/I$ で可逆元に写り、$B$ における像が
$(b_1, b_2)$ に属する元 $a \in A$ を見つけられる。Lemma
\ref{more-algebra-lemma-separate-image-closed-from-closed} を参照せよ。
$A$ を局所化 $A_a$ に取り替えると、$(b_1, b_2) = B$ となる。このとき
$\Spec(B) = D(b_1) \amalg D(b_2)$ である。これは $b_1b_2 = 0$ なので非交和であり、
$(b_1, b_2) = B$ なので $\Spec(B)$ を被覆する。開かつ閉な部分集合 $D(b_1)$ に
対応する冪等元を $e \in B$ とする。Algebra, Lemma
\ref{algebra-lemma-disjoint-decomposition} を参照せよ。
$b_1$ は $\overline{e}$ の持ち上げで、$b_2$ は $\pm (1 - \overline{e})$ の
持ち上げなので、Algebra, Lemma \ref{algebra-lemma-disjoint-decomposition} の
一意性の主張から、$e$ は $\overline{e}$ の持ち上げであると結論できる。
\end{proof}

\begin{lemma}
\label{more-algebra-lemma-lift-projective-module}
$A$ を環、$I \subset A$ をイデアルとする。$\overline{P}$ を有限生成射影
$A/I$-加群とする。このとき、同型 $A/I \to A'/IA'$ を誘導する
\'etale 環写像 $A \to A'$ と、$\overline{P}$ を持ち上げる有限生成射影
$A'$-加群 $P'$ が存在する。
\end{lemma}

\begin{proof}
ある整数 $n$ と、ある $R/I$-加群 $\overline{K}$ による直和分解
% SOURCE-ERRATUM-CANDIDATE P02-E0768: frozen source calls Kbar an R/I-module although only A and I are defined; preserved.
$(A/I)^{\oplus n} = \overline{P} \oplus \overline{K}$ を選べる。
直和因子 $\overline{P}$ に対応する射影子 $\overline{p}$ の持ち上げ
$\varphi : A^{\oplus n} \to A^{\oplus n}$ を選ぶ。
$f \in A[x]$ を $\varphi$ の特性多項式とし、$B = A[x]/(f)$ と置く。
Cayley--Hamilton により（Algebra, Lemma \ref{algebra-lemma-charpoly}）、
$x$ を $\varphi$ に送る写像 $B \to \text{End}_A(A^{\oplus n})$ がある。
各素イデアル $\mathfrak p \supset I$ に対して、$\kappa(\mathfrak p)$ における
$f$ の像は $(x - 1)^rx^{n - r}$ である。ここで $r$ は
$\overline{P} \otimes_{A/I} \kappa(\mathfrak p)$ の次元である。
従って、すべての $\mathfrak p \supset I$ に対して $(x - 1)^nx^n$ は
$B \otimes_A \kappa(\mathfrak p)$ で零に写る。よって $x(1 - x)$ は
$B/IB$ のすべての素イデアルに含まれる。従って、ある $N \geq 1$ に対して
$x^N(1 - x)^N$ は $IB$ に含まれる。
これより $x^N + (1 - x)^N$ は $B/IB$ の単元であり、
$$
\overline{e} = \text{image of }\frac{x^N}{x^N + (1 - x)^N}\text{ in }B/IB
$$
は冪等元である。この二つの主張はいずれも
$\mathbf{Z}[x]/(x^N(x - 1)^N)$ において成り立つからである。
$\text{End}_{A/I}((A/I)^{\oplus n})$ における $\overline{e}$ の像は
$$
\frac{\overline{p}^N}{\overline{p}^N + (1 - \overline{p})^N} = \overline{p}
$$
である。$\overline{p}$ は冪等元だからである。補題にあるような \'etale 拡大
$A'$ で $A$ を取り替えることで、$B/IB$ で $\overline{e}$ に写る冪等元
$e \in B$ が存在すると仮定してよい。Lemma
\ref{more-algebra-lemma-lift-idempotent-upstairs} を参照せよ。このとき写像
$$
B = A[x]/(f) \longrightarrow \text{End}_A(A^{\oplus n}).
$$
による $e$ の像は、$\overline{p}$ を持ち上げる冪等元 $p$ である。
$P = \Im(p)$ と置けばよい。
\end{proof}

\begin{lemma}
\label{more-algebra-lemma-cotangent-complex-symmetric-algebra}
$A$ を環とする。
% SOURCE-ERRATUM-CANDIDATE P02-E0769: frozen statement calls the displayed 0-to-0 complex merely a sequence although exactness is used throughout; preserved as “列”.
$0 \to K \to A^{\oplus m} \to M \to 0$ を $A$-加群の列とする。全射
$A^{\oplus m} \to M$ から得られる表示
$\alpha : A[y_1, \ldots, y_m] \to C$ を備えた $A$-代数
$C = \text{Sym}^*_A(M)$ を考える。このとき
$$
\NL(\alpha) =
(K \otimes_A C \to \bigoplus\nolimits_{j = 1, \ldots, m} C \text{d}y_j)
$$
であり（Algebra, Section \ref{algebra-section-netherlander} を参照）、特に
$\Omega_{C/A} = M \otimes_A C$ である。
\end{lemma}

\begin{proof}
$J = \Ker(\alpha)$ と置く。補題の主張は
$J/J^2 \cong K \otimes_A C$ である。$\alpha$ は次数付き代数の準同型であることに
注意する。次数 $d$ において
$(J/J^2)_d = K \otimes_A C_{d - 1}$ であることを証明する。まず
$$
J_d = \Ker(\text{Sym}^d_A(A^{\oplus m}) \to \text{Sym}^d_A(M))
= \Im(K \otimes_A \text{Sym}^{d - 1}_A(A^{\oplus m})
\to \text{Sym}^d_A(A^{\oplus m})),
$$
である。Algebra, Lemma \ref{algebra-lemma-presentation-sym-exterior} を参照せよ。
従って $(J^2)_d = \sum_{a + b = d} J_a \cdot J_b$ は
% SOURCE-ERRATUM-CANDIDATE P02-E0770: frozen source uses A^{\otimes m} at four loci where the surrounding symmetric-power construction requires A^{\oplus m}; formulas preserved.
$$
K \otimes_A K \otimes_A \text{Sym}^{d - 2}_A(A^{\otimes m})
\to \text{Sym}^d_A(A^{\oplus m}).
$$
の像である。上で参照した補題により、写像
$K \otimes_A \text{Sym}^{d - 2}_A(A^{\otimes m}) \to
\text{Sym}^{d - 1}_A(A^{\oplus m})$ の余核は
$\text{Sym}^{d - 1}_A(M)$ である。従って
$(J/J^2)_d = J_d/(J^2)_d$ が次に等しいことは明らかである。
\begin{align*}
\Coker(
K \otimes_A K \otimes_A \text{Sym}^{d - 2}_A(A^{\otimes m})
\to K \otimes_A \text{Sym}^{d - 1}_A(A^{\otimes m}))
& = K \otimes_A \text{Sym}^{d - 1}_A(M) \\
& = K \otimes_A C_{d -1}
\end{align*}
これが求める主張である。
\end{proof}

\begin{lemma}
\label{more-algebra-lemma-symmetric-algebra-smooth}
$A$ を環、$M$ を $A$-加群とする。このとき $C = \text{Sym}_A^*(M)$ が
$A$ 上滑らかであるための必要十分条件は、$M$ が有限生成射影 $A$-加群であることである。
\end{lemma}

\begin{proof}
$\sigma : C \to A$ を $C$ の次数 $0$ 部分への射影とする。このとき
$J = \Ker(\sigma)$ は次数 $> 0$ の部分であり、$A$-加群として $J/J^2 = M$ である。
従って $A \to C$ が滑らかならば、Algebra, Lemma
\ref{algebra-lemma-section-smooth} により $M$ は有限生成射影 $A$-加群である。

\medskip\noindent
逆に、$M$ が有限生成射影であると仮定し、核が $K$ である全射
$A^{\oplus n} \to M$ を選ぶ。もちろん $M$ は射影的なので、列
$0 \to K \to A^{\oplus n} \to M \to 0$ は分裂する。特に $K$ は有限
$A$-加群であり、従って $C$ は $A$ 上有限表示である。実際、$C$ は
$A[x_1, \ldots, x_n]$ を $K \subset \bigoplus Ax_i$ が生成するイデアルで割った
商である。Lemma \ref{more-algebra-lemma-cotangent-complex-symmetric-algebra} の計算は、
% SOURCE-ERRATUM-CANDIDATE P02-E0771: frozen source writes (K -> M) although the cotangent presentation uses K -> A^{\oplus n}; preserved.
$\NL_{C/A}$ が $(K \to M) \otimes_A C$ とホモトピー同値であることを示す。
従って $\NL_{C/A}$ は次数 $0$ に置かれた $C \otimes_A M$ と擬同型である。
これは、Algebra, Definition \ref{algebra-definition-smooth} により、$C$ が
$A$ 上滑らかであることを意味する。
\end{proof}

\begin{lemma}
\label{more-algebra-lemma-lift-section-smooth-morphism}
$A$ を環、$I \subset A$ をイデアルとする。可換図式
$$
\xymatrix{
B \ar[rd] \\
A \ar[u] \ar[r] & A/I
}
$$
を考える。ここで $B$ は滑らかな $A$-代数である。このとき、同型
$A/I \to A'/IA'$ を誘導する \'etale 環写像 $A \to A'$ と、環写像
$B \to A/I$ を持ち上げる $A$-代数写像 $B \to A'$ が存在する。
\end{lemma}

\begin{proof}
$J \subset B$ を $B \to A/I$ の核とする。従って $B/J = A/I$ である。
Algebra, Lemma \ref{algebra-lemma-application-NL-smooth} により列
$$
0 \to I/I^2 \to J/J^2 \to \Omega_{B/A} \otimes_B B/J \to 0
$$
は分裂完全である。従って
$\overline{P} = J/(J^2 + IB) = \Omega_{B/A} \otimes_B B/J$ は有限生成射影
$A/I$-加群である。ある整数 $n$ と直和分解
% SOURCE-ERRATUM-CANDIDATE P02-E0772: frozen source omits parentheses in A/I^{\oplus n}; exact formal span preserved.
$A/I^{\oplus n} = \overline{P} \oplus \overline{K}$ を選ぶ。
Lemma \ref{more-algebra-lemma-lift-projective-module} により、同型 $A/I \to A'/IA'$ を誘導する
\'etale 環写像 $A \to A'$ と、$\overline{K}$ を持ち上げる有限生成射影
% SOURCE-ERRATUM-CANDIDATE P02-E0773: frozen source says A-module after passing to A'; preserved before the stated replacement.
$A$-加群 $K$ を見つけられる。$A$ を $A'$ に取り替えることにする。
$B' = B \otimes_A \text{Sym}_A^*(K)$ と置く。Lemma
\ref{more-algebra-lemma-symmetric-algebra-smooth} により $A \to \text{Sym}_A^*(K)$ は滑らかなので、
$B \to B'$ は滑らかであり、従って $A \to B'$ も滑らかである
（Algebra, Lemmas \ref{algebra-lemma-base-change-smooth} および
\ref{algebra-lemma-locally-smooth} を参照）。さらに、切断
$\text{Sym}^*_A(K) \to A$ は切断 $B' \to B$ を定め、合成
$B' \to B \to A/I$ を $B' \to A/I$ とする。$J' \subset B'$ を
$B' \to A/I$ の核とする。
$JB' \subset J'$ および $B \otimes_A K \subset J'$ である。これらの写像は合わさって
同型
$$
(A/I)^{\oplus n} \cong J/(J^2 + IB) \oplus \overline{K}
\longrightarrow
J'/((J')^2 + IB')
$$
を与える。従って $B$ を $B'$ に取り替えることで、
$J/(J^2 + IB) = \Omega_{B/A} \otimes_B B/J$ は rank $n$ の自由
$A/I$-加群であると仮定してよい。

\medskip\noindent
この場合、$J/(J^2 + IB)$ の基底に写る $f_1, \ldots, f_n \in J$ を選ぶ。
有限表示 $A$-代数 $C = B/(f_1, \ldots, f_n)$ を考える。完全列
$$
0 \to H_1(L_{C/A}) \to (f_1, \ldots, f_n)/(f_1, \ldots, f_n)^2
\to \Omega_{B/A} \otimes_B C \to \Omega_{C/A} \to 0
$$
があることに注意する。Algebra, Lemma
\ref{algebra-lemma-exact-sequence-NL} を参照せよ
（$H_1(L_{B/A}) = 0$ であり、$\Omega_{B/A}$ は有限生成射影、特に平坦なので、
Tor 群が消えることに注意せよ）。$B$ の任意の素イデアル
$\mathfrak q \supset J$ に対して、加群 $\Omega_{B/A, \mathfrak q}$ は rank $n$ の
自由加群である。これは $\Omega_{B/A}$ が有限生成射影であり、かつ
$\Omega_{B/A} \otimes_B B/J$ が rank $n$ の自由加群だからである
（Algebra, Lemma \ref{algebra-lemma-finite-projective} を参照）。
$f_1, \ldots, f_n$ の選び方により、写像
$$
\left((f_1, \ldots, f_n)/(f_1, \ldots, f_n)^2\right)_{\mathfrak q}
\to
\Omega_{B/A, \mathfrak q}
$$
は $J$ を法として全射である。従って、局所環 $C_{\mathfrak q}$ 上のこの加群写像は
同型でなければならないことが分かる
（Algebra, Lemma \ref{algebra-lemma-NAK} により全射であり、例えば Algebra, Lemma
\ref{algebra-lemma-fun} により、
$((f_1, \ldots, f_n)/(f_1, \ldots, f_n)^2)_{\mathfrak q}$ は $n$ 個の元で
生成されるので単射である）。従って
$H_1(L_{C/A})_{\mathfrak q} = 0$ かつ $\Omega_{C/A, \mathfrak q} = 0$ である。
Algebra, Lemma \ref{algebra-lemma-smooth-at-point} により、$A \to C$ は
$\mathfrak q$ に対応する $C$ の素イデアル $\overline{\mathfrak q}$ で滑らかである。
$\Omega_{C/A, \mathfrak q} = 0$ なので、実際には $\overline{\mathfrak q}$ で
\'etale である。従って $A \to C$ は、$JC$ を含む $C$ のすべての素イデアルで
\'etale である。Lemma \ref{more-algebra-lemma-localize-upstairs} により、$C/JC$ で可逆元に写り、
$A \to C_f$ が \'etale となる $f \in C$ を見つけられる。$f$ の選び方により、
依然として $C_f/JC_f = A/I$ である。写像 $C_f/IC_f \to A/I$ は全射であり、
Algebra, Lemma \ref{algebra-lemma-map-between-etale} により \'etale である。
従って $A/I$ は $C_f/IC_f$ をある元 $g \in C$ で局所化したものと同型である。
Algebra, Lemma \ref{algebra-lemma-surjective-flat-finitely-presented} を参照せよ。
$A' = C_{fg}$ と置けば証明が完了する。
\end{proof}

\section{Zariski 対}
\label{more-algebra-section-zariski-pairs}

\noindent
この節と次節において、{it 対}とは、$A$ が環で $I \subset A$ がイデアルであるような
対 $(A, I)$ のことである。{it 対の射} $(A, I) \to (B, J)$ とは、
$\varphi(I) \subset J$ を満たす環写像 $\varphi : A \to B$ のことである。

\begin{definition}
\label{more-algebra-definition-zariski-pair}
{it Zariski 対}とは、$I$ が $A$ の Jacobson 根基に含まれるような対 $(A, I)$ のことである。
\end{definition}

\begin{lemma}
\label{more-algebra-lemma-idempotents-determined-modulo-radical}
$(A, I)$ を Zariski 対とする。$A$ の冪等元から $A/I$ の冪等元への写像は単射である。
\end{lemma}

\begin{proof}
局所環の冪等元は $0$ または $1$ である。従って、Algebra, Lemma
\ref{algebra-lemma-characterize-zero-local} により、冪等元はそれが消える極大イデアルの
集合によって決まる。
\end{proof}

\begin{lemma}
\label{more-algebra-lemma-check-isomorphism-zariski}
$(A, I)$ を Zariski 対とする。$A \to B$ を平坦かつ整で有限表示な環写像とし、
$A/I \to B/IB$ は同型であるとする。このとき $A \to B$ は同型である。
\end{lemma}

\begin{proof}
Algebra, Lemma \ref{algebra-lemma-characterize-finite-in-terms-of-integral} により、
環写像 $A \to B$ は有限である。従って Algebra, Lemma
\ref{algebra-lemma-finite-finitely-presented-extension} により、$B$ は
$A$-加群として有限表示である。従って Algebra, Lemma
\ref{algebra-lemma-finite-projective} により、$B$ は有限局所自由 $A$-加群である。
加群 $B$ は $V(I)$ に沿って rank $1$ であり
（Algebra, Lemma \ref{algebra-lemma-finite-projective} で記述される rank 関数を参照）、
$(A, I)$ は Zariski 対なので、rank は至る所 $1$ であると結論できる。
これより $A \to B$ は同型である。実際、$S$ が rank $1$ の局所自由
$R$-加群となる環写像 $R \to S$ は同型であることを示すのは、よい演習である
（ヒント：局所環を考えよ）。
\end{proof}

\begin{lemma}
\label{more-algebra-lemma-helper-finite}
$(A, I)$ を Zariski 対とし、$A \to B$ を有限環写像とする。次を仮定する。
\begin{enumerate}
% SOURCE-ERRATUM-CANDIDATE P02-E0774: frozen source omits terminal punctuation after item (1); Japanese supplies it.
\item $B/IB = B_1 \times B_2$ は $A/I$-代数の積である、
\item $A/I \to B_1/IB_1$ は全射である、
\item $b \in B$ はこの積において $(1, 0)$ に写る。
\end{enumerate}
このとき、ある $d \geq 1$ に対して $f(b) = 0$ かつ
$f \bmod I = (x - 1)x^d$ を満たすモニック多項式 $f \in A[x]$ が存在する。
\end{lemma}

\begin{proof}
Lemma \ref{more-algebra-lemma-lift-idempotent-upstairs} により、同型
$A/I \to A'/IA'$ を誘導する \'etale 環写像 $A \to A'$ で、
$B' = B \otimes_A A'$ が $B'/IB'$ における $b$ の像を持ち上げる冪等元
$e'$ を含むものを見つけられる。対応する $A'$-代数の分解
$$
B' = B'_1 \times B'_2
$$
を考える。これは $I$ を法として簡約すると、補題で与えられた分解と両立する。
写像 $A' \to B'_1$ は $IA'$ を法として全射である。Nakayama の補題
（Algebra, Lemma \ref{algebra-lemma-NAK}）により、ある $i \in IA'$ を見つけ、
$A'$ を $A'_{1 + i}$ に取り替えた後で $A' \to B'_1$ が全射となるようにできる。
第1因子における $b$ の像 $b'_1 \in B'_1$ は
$b'_1 - 1 \in IB'_1$ を満たすことに注意する。従って、$b'_1 - 1$ に写る
$a' \in IA'$ を選べる。一方、第2因子における $b$ の像 $b'_2 \in B'_2$ は
$IB'_2$ に属する。Algebra, Lemma
\ref{algebra-lemma-integral-integral-over-ideal} により、次数 $d$ のモニック多項式
$g(x) = x^d + \sum a'_j x^j$ で、$a'_j \in IA'$ かつ $B'_2$ において
$g(b'_2) = 0$ を満たすものが存在する。従って積における $b$ の像
$b' = (b'_1, b'_2) \in B'$ は多項式 $(x - 1 - a')g(x)$ の根である。
これより
$$
(b' - 1)(b')^d \in \sum\nolimits_{j = 0, \ldots, d} IA' \cdot (b')^j
$$
と結論する。我々は、これから $B$ において
$$
(b - 1)b^d \in \sum\nolimits_{j = 0, \ldots, d} I \cdot b^j
$$
が従うと主張する。そのためには、環写像 $A \to A'$ が忠実平坦であることを
示せば十分である。というのも、この条件は
$B/\sum_{j = 0, \ldots, d} Ib^j$ において $(b - 1)b^d$ の像が零であるという
条件だからである（Algebra, Lemma
\ref{algebra-lemma-faithfully-flat-universally-injective} を用いよ）。
% SOURCE-ERRATUM-CANDIDATE P02-E0775: frozen source omits “is” in “The map A to A' flat”; Japanese supplies the copula.
写像 $A \to A'$ は \'etale なので平坦である
（Algebra, Lemma \ref{algebra-lemma-etale}）。一方、誘導されるスペクトル上の写像は
開写像であり（Algebra, Proposition \ref{algebra-proposition-fppf-open} と先に参照した
補題を用いよ）、その像は $V(I)$ を含む。$I$ は $A$ の Jacobson 根基に含まれるので、
結論が従う。
\end{proof}

\begin{lemma}
\label{more-algebra-lemma-noetherian-zariski-jacobson-complement}
$(A, I)$ を Zariski 対とし、$A$ は Noether 環であるとする。
$f \in I$ とする。このとき $A_f$ は Jacobson 環である。
\end{lemma}

\begin{proof}
Algebra, Lemma \ref{algebra-lemma-noetherian-dim-1-Jacobson} の判定法を用いる。
$\mathfrak p_f = \mathfrak p A_f$ が素イデアルで極大でないような素イデアル
$\mathfrak p \subset A$ をとる。$A_f/\mathfrak p_f = (A/\mathfrak p)_f$ が無限個の
素イデアルをもつことを示す必要がある。$A$ を $A/\mathfrak p$ に取り替えることで、
$A$ は整域、$\dim A_f > 0$ であり、目標は $\Spec(A_f)$ が無限であることだと
仮定してよい。$\dim A_f > 0$ なので、$f$ を含まない非零素イデアル
$\mathfrak q \subset A$ を見つけられる。$\mathfrak q$ を含む極大イデアル
$\mathfrak m \subset A$ を選ぶ。$(A, I)$ は Zariski 対なので
$I \subset \mathfrak m$ である。従って $\mathfrak m \not = \mathfrak q$ かつ
$\dim(A_\mathfrak m) > 1$ である。Algebra, Lemma
\ref{algebra-lemma-Noetherian-local-domain-dim-2-infinite-opens} により
$\Spec((A_\mathfrak m)_f) \subset \Spec(A_f)$ は無限であり、主張が従う。
\end{proof}

\section{Hensel 対}
\label{more-algebra-section-henselian-pairs}

\noindent
Section \ref{more-algebra-section-lifting} の結果のいくつかは、Hensel 対に関する結果と
みなすことができる。この節において、{it 対}とは、$A$ が環で
$I \subset A$ がイデアルであるような対 $(A, I)$ のことである。
{it 対の射} $(A, I) \to (B, J)$ とは、$\varphi(I) \subset J$ を満たす
環写像 $\varphi : A \to B$ のことである。Section \ref{more-algebra-section-lifting} と同様に、
$A$ 上の対象 $\xi$ が与えられたとき、$\xi$ を $A/I$ 上の対象へ「基底変換」したものを
$\overline{\xi}$ で表す（これが意味をもつ場合に限る）。

\begin{definition}
\label{more-algebra-definition-henselian-pair}
{it Hensel 対}とは、次を満たす対 $(A, I)$ のことである。
\begin{enumerate}
\item $I$ は $A$ の Jacobson 根基に含まれる、かつ
\item 任意のモニック多項式 $f \in A[T]$ と、$A/I[T]$ の単位イデアルを生成する
モニック多項式 $g_0, h_0 \in A/I[T]$ による因数分解
$\overline{f} = g_0h_0$ に対して、$g, h$ がモニックで
$g_0 = \overline{g}$, $h_0 = \overline{h}$ を満たす $A[T]$ における因数分解
$f = gh$ が存在する。
\end{enumerate}
\end{definition}

\noindent
$A$ が局所環で $I = \mathfrak m$ が極大イデアルならば、$(A, I)$ が Hensel 対である
ための必要十分条件は、$A$ が Hensel 局所環であることである。Algebra, Lemma
\ref{algebra-lemma-characterize-henselian} を参照せよ。Lemma
\ref{more-algebra-lemma-characterize-henselian-pair} では Hensel 対のいくつかの同値な特徴付けを
与える（そして今後さらに追加する）。

\begin{lemma}
\label{more-algebra-lemma-locally-nilpotent-henselian}
$(A, I)$ を、$I$ が局所冪零である対とする。関手 $B \mapsto B/IB$ は、
$A$ 上の \'etale 代数の圏と $A/I$ 上の \'etale 代数の圏との間の同値を誘導する。
さらに、この対は Hensel 対である。
\end{lemma}

\begin{proof}
本質的全射性は Algebra, Lemma \ref{algebra-lemma-lift-etale} により成立する。
$B$, $B'$ が $A$ 上 \'etale で、$B/IB \to B'/IB'$ が $A/I$-代数の射ならば、
Algebra, Lemma \ref{algebra-lemma-smooth-strong-lift} によりこれを持ち上げられる。
最後に、$f, g : B \to B'$ が $f \bmod I = g \bmod I$ を満たす二つの
$A$-代数写像であるとする。乗法写像 $B \otimes_A B \to B$ の核を生成する冪等元
$e \in B \otimes_A B$ を選ぶ。Algebra, Lemmas
\ref{algebra-lemma-diagonal-unramified} および \ref{algebra-lemma-unramified}
（\'etale が不分岐であることを見るため）を参照せよ。このとき
$(f \otimes g)(e) \in IB'$ である。$IB'$ は局所冪零なので
（Algebra, Lemma \ref{algebra-lemma-locally-nilpotent}）、Algebra, Lemma
\ref{algebra-lemma-lift-idempotents} により $(f \otimes g)(e) = 0$ が従う。
従って $f = g$ である。

\medskip\noindent
$I$ が $A$ の Jacobson 根基に含まれることは明らかである。
$f \in A[T]$ をモニック多項式とし、$g_0, h_0 \in A/I[T]$ がモニックで
$A/I[T]$ の単位イデアルを生成するような
$\overline{f} = f \bmod I$ の因数分解 $\overline{f} = g_0h_0$ をとる。
Lemma \ref{more-algebra-lemma-lift-factorization-monic} により、同型
$A/I \to A'/IA'$ を誘導する \'etale 環写像 $A \to A'$ で、因数分解が
$A'$ 上のモニック多項式への因数分解に持ち上がるものが存在する。
上で示したことにより $A = A'$ であり、この因数分解は $A$ 上にある。
\end{proof}

\begin{lemma}
\label{more-algebra-lemma-limit-henselian}
$A = \lim A_n$ とする。ここで $(A_n)$ は環の逆系であり、その遷移写像は全射で
局所冪零な核をもつ。$I_n = \Ker(A \to A_n)$ と置くと、$(A, I_n)$ は Hensel 対である。
\end{lemma}

\begin{proof}
$n$ を固定する。$A_n$ で $1$ に写る元 $a \in A$ をとる。Algebra, Lemma
\ref{algebra-lemma-locally-nilpotent-unit} により、$a$ はすべての $m \geq n$ に対して
$A_m$ の単元に写る。従って $a$ は $A$ の単元である。よって Algebra, Lemma
\ref{algebra-lemma-contained-in-radical} により、イデアル $I_n$ は $A$ の Jacobson
根基に含まれる。$f \in A[T]$ をモニック多項式とし、$g_n, h_n \in A_n[T]$ が
モニックで $A_n[T]$ の単位イデアルを生成するような
$\overline{f} = f \bmod I_n$ の因数分解 $\overline{f} = g_nh_n$ をとる。
Lemma \ref{more-algebra-lemma-locally-nilpotent-henselian} により、この因数分解を逐次的に、
すべての $m \geq n$ に対する $A_m[T]$ のモニック多項式 $g_m, h_m$ による
$f \bmod I_m = g_m h_m$ に持ち上げられる。各段階で、持ち上げ
$g_m, h_m$ が $A_n[T]$ の単位イデアルを生成することを確認しなければならない。
% SOURCE-ERRATUM-CANDIDATE P02-E0776: frozen source says the lifts g_m,h_m generate in A_n[T], although they live in A_m[T]; preserved.
これは $g_n, h_n$ に対する対応する事実と、$A_m \to A_n$ の核が局所冪零なので
$\Spec(A_n[T]) = \Spec(A_m[T])$ であることから従う。
$A = \lim A_m$ なので証明が完了する。
\end{proof}

\begin{lemma}
\label{more-algebra-lemma-complete-henselian}
$(A, I)$ を対とする。$A$ が $I$-進完備ならば、この対は Hensel 対である。
\end{lemma}

\begin{proof}
Algebra, Lemma \ref{algebra-lemma-radical-completion} により、イデアル $I$ は
$A$ の Jacobson 根基に含まれる。$f \in A[T]$ をモニック多項式とし、
$g_0, h_0 \in A/I[T]$ がモニックで $A/I[T]$ の単位イデアルを生成するような
$\overline{f} = f \bmod I$ の因数分解 $\overline{f} = g_0h_0$ をとる。
Lemma \ref{more-algebra-lemma-locally-nilpotent-henselian} により、この因数分解を逐次的に、
すべての $n \geq 1$ に対する $A/I^n[T]$ のモニック多項式 $g_n, h_n$ による
$f \bmod I^n = g_n h_n$ に持ち上げられる。$A = \lim A/I^n$ なので証明が完了する。
\end{proof}

\begin{lemma}
\label{more-algebra-lemma-helper-finite-type}
$(A, I)$ を対とする。$A \to B$ を有限型環写像とし、
$B/IB = C_1 \times C_2$、かつ $A/I \to C_1$ は有限であるとする。
$B'$ を $B$ における $A$ の整閉包とする。このとき、写像
$B'/IB' \to B/IB$ が積分解を保つように $B'/IB' = C_1 \times C'_2$ と書くことができ、
さらに $C_1 \times C'_2$ で $(1, 0)$ に写る $g \in B'$ で
$B'_g \to B_g$ が同型となるものが存在する。
\end{lemma}

\begin{proof}
$A \to B$ は閉部分集合 $T = \Spec(C_1) \subset \Spec(B)$ のすべての素イデアルで
準有限であることに注意する（ファイバー環を見れば従う。Algebra, Definition
\ref{algebra-definition-quasi-finite} を参照）。位相空間の図式
$$
\xymatrix{
\Spec(B) \ar[rr]_\phi \ar[rd]_\psi & & \Spec(B') \ar[ld]^{\psi'} \\
& \Spec(A)
}
$$
を考える。Algebra, Theorem \ref{algebra-theorem-main-theorem} により、
各 $\mathfrak p \in T$ に対して $h_\mathfrak p \in B'$ で、
$h_\mathfrak p \not \in \mathfrak p$ かつ
% SOURCE-ERRATUM-CANDIDATE P02-E0777: frozen source switches from h_p to undefined h in the localization map; preserved.
$B'_h \to B_h$ が同型となるものが存在する。和集合
$U = \bigcup D(h_\mathfrak p)$ は、$\phi^{-1}(U) \to U$ が同相で
$T \subset \phi^{-1}(U)$ となる開集合 $U \subset \Spec(B')$ を与える。
$T$ は $\psi^{-1}(V(I))$ において開なので、$\phi(T)$ は
$U \cap (\psi')^{-1}(V(I))$ において開である。従って $\phi(T)$ は
$(\psi')^{-1}(V(I))$ において開である。一方、$C_1$ は $A/I$ 上有限なので
$B'$ 上有限である。従って Algebra, Lemmas \ref{algebra-lemma-going-up-closed} および
\ref{algebra-lemma-integral-going-up} により、$\phi(T)$ は $\Spec(B')$ の閉部分集合である。
従って $\Spec(B'/IB') \supset \phi(T)$ であり、この部分集合は開かつ閉である。Algebra, Lemma
\ref{algebra-lemma-disjoint-implies-product} により、対応する積分解
$B'/IB' = C'_1 \times C'_2$ を得る。スペクトル上で何が起こるかを見ることで、
写像 $B'/IB' \to B/IB$ は $C'_1$ を $C_1$ に、$C'_2$ を $C_2$ に送ることが分かる
（ヒント：$\phi(T)$ の逆像はちょうど $T$ である。詳細の一部は省略する）。
$C'_1 \times C'_2$ において $(1, 0)$ に写り、$D(g) \subset U$ となる
$g \in B'$ を選ぶ。これは $\Spec(C'_1)$ と $\Spec(C'_2)$ が $\Spec(B')$ の
互いに素な閉部分集合で、$\Spec(C'_1)$ が $U$ に含まれるので可能である。
すると $B'_g \to B_g$ はスペクトル上の同相と局所環上の同型を定める
（上の $U$ の選び方による）。従って、例えば Algebra, Lemma
\ref{algebra-lemma-characterize-zero-local} により、これは同型である。
最後に $C'_1 = C_1$ が従い、証明が完了する。
\end{proof}

\begin{lemma}
\label{more-algebra-lemma-characterize-henselian-pair}
\begin{reference}
\cite[Chapter XI]{Henselian} and \cite[Proposition 1]{Gabber-henselian}
\end{reference}
$(A, I)$ を対とする。次の条件は同値である。
\begin{enumerate}
\item $(A, I)$ は Hensel 対である、
\item \'etale 環写像 $A \to A'$ と $A$-代数写像
$\sigma : A' \to A/I$ が与えられると、$\sigma$ を持ち上げる $A$-代数写像
$A' \to A$ が存在する、
\item 任意の有限 $A$-代数 $B$ に対して、写像 $B \to B/IB$ は冪等元上の
全単射を誘導する、
\item 任意の整 $A$-代数 $B$ に対して、写像 $B \to B/IB$ は冪等元上の
全単射を誘導する、かつ
\item (Gabber) $I$ は $A$ の Jacobson 根基に含まれ、$a_n, \ldots, a_0 \in I$,
$n \ge 1$ を満たす次の形の任意のモニック多項式 $f(T) \in A[T]$ は、
根 $\alpha \in 1 + I$ をもつ。
$$
f(T) = T^n(T - 1) + a_n T^n + \ldots + a_1 T + a_0
$$
\end{enumerate}
さらに、(5) における根は一意である。
\end{lemma}

\begin{proof}
(2) が成り立つと仮定する。このとき $I$ は $A$ の Jacobson 根基に含まれる。
そうでなければ、$I$ を法として $1$ と合同な非単元 $f \in A$ が存在し、写像
$A \to A_f$ が (2) に矛盾するからである。従って $B$ が $A$ 上整ならば、
$IB \subset B$ は $B$ の Jacobson 根基に含まれる。実際、Algebra, Lemmas
\ref{algebra-lemma-going-up-closed} および \ref{algebra-lemma-integral-going-up}
により $\Spec(B) \to \Spec(A)$ は閉写像である。よって Lemma
\ref{more-algebra-lemma-idempotents-determined-modulo-radical} により、$B$ の冪等元から
$B/IB$ の冪等元への写像は単射である。一方、(2) が成り立つので、Lemma
\ref{more-algebra-lemma-lift-idempotent-upstairs} により $B/IB$ の各冪等元は $B$ の冪等元に
持ち上がる。これで (2) が (4) を含意することが分かる。

\medskip\noindent
含意 (4) $\Rightarrow$ (3) は自明である。

\medskip\noindent
(3) を仮定する。$\mathfrak m$ を極大イデアルとし、有限写像
$A \to B = A/(I \cap \mathfrak m)$ を考える。$B \to B/IB$ が冪等元上の
全単射を誘導するという条件から、$I \subset \mathfrak m$ が従う
（そうでなければ、$B = A/I \times A/\mathfrak m$ かつ $B/IB = A/I$ となる）。
従って $I$ は $A$ の Jacobson 根基に含まれる。
$f \in A[T]$ をモニックとし、$A/I[T]$ の単位イデアルを生成するモニック多項式
$g_0, h_0 \in A/I[T]$ による因数分解 $\overline{f} = g_0h_0$ が与えられたとする。
$B = A[T]/(f)$ と置く。$A$-代数の分解
$$
B/IB = A/I[T]/(g_0) \times A/I[T]/(h_0)
$$
に対応する $B/IB$ の冪等元を $\overline{e}$ とする。(3) の仮定により存在する
$\overline{e}$ の持ち上げを $e \in B$ とする。これから積分解
$$
B = eB \times (1 - e)B
$$
が得られる。$B$ は $A$-加群として rank $\deg(f)$ の自由加群であることに注意する。
従って $eB$ と $(1 - e)B$ は有限局所自由 $A$-加群である。しかも、$A/I$ 上で
$eB$ と $(1 - e)B$ はそれぞれ一定の rank $\deg(g_0)$ と $\deg(h_0)$ をもつので、
$\Spec(A)$ 上でも同じことが成り立つ。従って
\begin{align*}
f & = \text{CharPol}_A(T : B \to B) \\
& =
\text{CharPol}_A(T : eB \to eB)
\text{CharPol}_A(T : (1 - e)B \to (1 - e)B)
\end{align*}
は、$I$ を法として与えられた因数分解に簡約するモニック多項式への因数分解である。
ここで $\text{CharPol}_A$ は、有限局所自由 $A$-加群の自己準同型の特性多項式を表す。
加群が自由ならば、$\text{CharPol}_A$ は対応する行列の特性多項式として定義され、
一般には Algebra, Lemma \ref{algebra-lemma-standard-covering} を用いて貼り合わせる。
詳細は省略する。従って (3) は (1) を含意する。

\medskip\noindent
(1) を仮定する。$f$ を (5) のような多項式とする。$f \bmod I$ の
$T^n$ と $T - 1$ への因数分解は、Definition \ref{more-algebra-definition-henselian-pair} により、
$g$ と $h$ がモニックである因数分解 $f = gh$ に持ち上がる。このとき $h$ の次数は
$1$ でなければならず、$f$ は $I$ を法として $1$ に簡約する根をもつことが分かる。
最後に、Hensel 対の定義により $I$ は Jacobson 根基に含まれる。
従って (1) は (5) を含意する。

\medskip\noindent
最後の段階の証明を与える前に、(5) の根 $\alpha$ が存在するならば一意であることを
示そう。実際、$f(T)$ の明示的な形により $f'(\alpha) \in 1 + I$ である。
ここで $f'$ は $T$ に関する $f$ の導関数である。初等的な議論により
$$
f(T) = f(\alpha + T - \alpha) =
f(\alpha) + f'(\alpha) \cdot (T - \alpha) \bmod
(T - \alpha)^2 A[T]
$$
となる。従って $f(T)$ の別の根 $\alpha' \in 1 + I$ は、ある $i \in I$ に対して
$0 = f(\alpha') - f(\alpha) = (\alpha' - \alpha)(1 + i)$ を満たす。
$1 + i$ は $A$ の単元なので、$\alpha = \alpha'$ である。

\medskip\noindent
(5) を仮定する。(2) が成り立つこと、すなわち、任意の \'etale 写像
$A \to A'$ に対して、$I$ を法とする任意の切断 $\sigma : A' \to A/I$ が
切断 $A' \to A$ に持ち上がることを示す。$A \to A'$ は \'etale なので、切断
$\sigma$ は $A/I$-代数の分解
\begin{equation}
\label{more-algebra-equation-GCHP}
A'/IA' \cong A/I \times C
\end{equation}
を定める。実際、全射環写像 $A'/IA' \to A/I$ は Algebra, Lemma
\ref{algebra-lemma-map-between-etale} により \'etale であり、Algebra, Lemma
\ref{algebra-lemma-surjective-flat-finitely-presented} により求める冪等元が得られる。
この分解が、$A'_1$ が $A$ 上整であるような $A$-代数の分解
\begin{equation}
\label{more-algebra-equation-GCHP-want}
A' \cong A'_1 \times A'_2
\end{equation}
に持ち上がることを示す。すると $A \to A'_1$ は整かつ \'etale で、
$A/I \to A'_1/IA'_1$ は同型である。従って Lemma
\ref{more-algebra-lemma-check-isomorphism-zariski} により $A \to A'_1$ は同型である
（ここでは \'etale 環写像が平坦かつ有限表示であることも用いる。Algebra, Lemma
\ref{algebra-lemma-etale} を参照せよ）。

\medskip\noindent
$B'$ を $A'$ における $A$ の整閉包とする。Lemma
\ref{more-algebra-lemma-helper-finite-type} により、(\ref{more-algebra-equation-GCHP}) と両立するように
$A/I$-代数として
\begin{equation}
\label{more-algebra-equation-dec-mod-I}
B'/IB' \cong A/I \times C'
\end{equation}
と分解でき、$(1, 0)$ を持ち上げる $b \in B'$ で
$B'_b \to A'_b$ が同型となるものを見つけられる。分解
(\ref{more-algebra-equation-dec-mod-I}) が $A$-代数の分解
\begin{equation}
\label{more-algebra-equation-want-2}
B' \cong B'_1 \times B'_2
\end{equation}
に持ち上がるならば、誘導される分解 $A' = A'_1 \times A'_2$ は、求める
(\ref{more-algebra-equation-GCHP-want}) を与える。実際、$b$ は $B'_1$ の単元なので
（詳細は省略する）、$B'_1 \cong A'_1$ となり、$A'_1$ は $A$ 上整となる。

\medskip\noindent
$b$ を含む有限 $A$-部分代数 $B'' \subset B'$ を選ぶ
（$B'$ の任意の有限生成 $A$-部分代数は $A$ 上有限であることに注意せよ）。
$B''$ を拡大することで、$b$ は $B''/IB''$ において冪等元に写り、分解
\begin{equation}
\label{more-algebra-equation-again-dec-mod-I}
B''/IB'' \cong C''_1 \times C''_2
\end{equation}
を生じると仮定してよい。$B'_b \cong A'_b$ なので、$B'_b$ は $A$ 上有限型である。
$B'_b$ は $A$ 上 $b_1/b^n, \ldots, b_t/b^n$ によって生成されるとし、
$b_1, \ldots, b_t \in B''$ となるように $B''$ を拡大する。
すると $B''_b \to B'_b$ は全射かつ単射、従って同型である。特に
$C''_1 = A/I$ である！ 従って $A/I \to C''_1$ は同型、特に全射である。
Lemma \ref{more-algebra-lemma-helper-finite} により、$a_n, \ldots, a_0 \in I$,
$n \ge 1$ で、$f(b) = 0$ を満たす形
$$
f(T) = T^n(T - 1) + a_n T^n + \ldots + a_1 T + a_0
$$
の $f(T) \in A[T]$ を見つけられる。特に $B'$ は
% SOURCE-ERRATUM-CANDIDATE P02-E0778: frozen source says “a A[T]/(f)-algebra”; Japanese has no article.
$A[T]/(f)$-代数である。(5) により、$f$ の根 $a \in 1 + I$ が存在する。
これは、$B'/IB'$ の分解 (\ref{more-algebra-equation-dec-mod-I}) と両立する積分解
$A[T]/(f) = A[T]/(T - a) \times D$ を生じる。そこから誘導される $B'$ の分解が
求める (\ref{more-algebra-equation-want-2}) である。
\end{proof}

\begin{lemma}
\label{more-algebra-lemma-change-ideal-henselian-pair}
$A$ を環とし、$I, J \subset A$ を $V(I) = V(J)$ を満たすイデアルとする。
このとき $(A, I)$ が Hensel 対であるための必要十分条件は、$(A, J)$ が
Hensel 対であることである。
\end{lemma}

\begin{proof}
任意の整環写像 $A \to B$ に対して $V(IB) = V(JB)$ である。従って $B/IB$ の
冪等元と $B/JB$ の冪等元は全単射に対応する
（Algebra, Lemma \ref{algebra-lemma-disjoint-decomposition}）。これより、
$B \to B/IB$ が冪等元の集合上の全単射を誘導するための必要十分条件は、
$B \to B/JB$ が冪等元の集合上の全単射を誘導することである。
従って Lemma \ref{more-algebra-lemma-characterize-henselian-pair} により結論が従う。
\end{proof}

\begin{lemma}
\label{more-algebra-lemma-integral-over-henselian-pair}
$(A, I)$ を Hensel 対とし、$A \to B$ を整環写像とする。このとき
$(B, IB)$ は Hensel 対である。
\end{lemma}

\begin{proof}
Lemma \ref{more-algebra-lemma-characterize-henselian-pair} における Hensel 対の第4の特徴付けと、
整環写像の合成が整であるという事実から直ちに従う。
\end{proof}

\begin{lemma}
\label{more-algebra-lemma-henselian-henselian-pair}
$I \subset J \subset A$ を環 $A$ のイデアルとする。次の条件は同値である。
\begin{enumerate}
\item $(A, I)$ と $(A/I, J/I)$ は Hensel 対である、かつ
% SOURCE-ERRATUM-CANDIDATE P02-E0779: frozen source uses “an henselian pair”; Japanese has no indefinite article.
\item $(A, J)$ は Hensel 対である。
\end{enumerate}
\end{lemma}

\begin{proof}
(1) を仮定する。$B$ を整 $A$-代数とする。環写像
$$
B \to B/IB \to B/JB
$$
を考える。Lemma \ref{more-algebra-lemma-characterize-henselian-pair} により、両方の矢印が
冪等元上の全単射を誘導することが分かる。従って合成もそうである。
よって Lemma \ref{more-algebra-lemma-characterize-henselian-pair} により $(A, J)$ は Hensel 対である。

\medskip\noindent
逆に (2) が成り立つと仮定する。Lemma \ref{more-algebra-lemma-integral-over-henselian-pair} により
$(A/I, J/I)$ は Hensel 対である。$B$ を整 $A$-代数とし、環写像
$$
B \to B/IB \to B/JB
$$
を考える。Lemma \ref{more-algebra-lemma-characterize-henselian-pair} により、合成と第2の矢印は
冪等元上の全単射を誘導する。従って第1の矢印もそうである。よって再び同じ補題により
$(A, I)$ は Hensel 対である。
\end{proof}

\begin{lemma}
\label{more-algebra-lemma-sum-henselian}
$A$ を環とし、$(A, I)$ と $(A, I')$ を Hensel 対とする。このとき
% SOURCE-ERRATUM-CANDIDATE P02-E0780: frozen source uses “an henselian pair”; Japanese has no indefinite article.
$(A, I + I')$ は Hensel 対である。
\end{lemma}

\begin{proof}
Lemma \ref{more-algebra-lemma-integral-over-henselian-pair} により、対
$(A/I, (I' + I)/I)$ は Hensel 対である。従って Lemma
\ref{more-algebra-lemma-henselian-henselian-pair} から結論を得る。
\end{proof}

\begin{lemma}
\label{more-algebra-lemma-product-henselian-pairs}
$J$ を集合、$\{ (A_j, I_j)\}_{j \in J}$ を対の族とする。このとき
$(\prod_{j \in J} A_j, \prod_{j\in J} I_j)$ が Hensel 対であるための必要十分条件は、
各 $(A_j, I_j)$ が Hensel 対であることである。
\end{lemma}
 
\begin{proof}
各 $j \in J$ に対して、射影 $\prod_{j \in J} A_j \rightarrow A_j$ は整環写像なので、
$(\prod_{j \in J} A_j, \prod_{j\in J} I_j)$ が Hensel 対ならば、Lemma
\ref{more-algebra-lemma-integral-over-henselian-pair} により各 $(A_j, I_j)$ も Hensel 対である。

\medskip\noindent
逆に、各 $(A_j, I_j)$ が Hensel 対であると仮定する。このとき
$x \in \prod_{j \in J} I_j$ に対する各 $1 + x$ は、成分ごとに単元なので
$\prod_{j \in J} A_j$ の単元である。ここで Algebra, Lemma
\ref{algebra-lemma-contained-in-radical} と Definition
\ref{more-algebra-definition-henselian-pair} を用いた。従って再び Algebra, Lemma
\ref{algebra-lemma-contained-in-radical} により、$\prod_{j \in J} I_j$ は
$\prod_{j \in J} A_j$ の Jacobson 根基に含まれる。引き続き成分ごとに考えると、
任意のモニック $f \in (\prod_{j \in J} A_j)[T]$ と、
$\overline{f} = g_0h_0$ という因数分解で、
$g_0, h_0 \in (\prod_{j \in J} A_j / \prod_{j \in J} I_j)[T] =
(\prod_{j \in J} A_j/I_j)[T]$ がモニックで
$(\prod_{j \in J} A_j / \prod_{j \in J} I_j)[T]$ の単位イデアルを
生成するものに対して、$(\prod_{j \in J} A_j)[T]$ における
$g$, $h$ がモニックで $g_0$, $h_0$ に簡約する因数分解 $f = gh$ が存在する。
結論として Definition \ref{more-algebra-definition-henselian-pair} により
$(\prod_{j \in J} A_j, \prod_{j\in J} I_j)$ は Hensel 対である。
\end{proof}

\begin{lemma}
\label{more-algebra-lemma-limits-henselian}
Hensel 対であるという性質は、対の極限をとることで保たれる。より正確には、
$J$ を前順序集合、$(A_j, I_j)$ を $J$ 上の Hensel 対の逆系とする。このとき、
イデアル $I = \lim I_j$ を備えた $A = \lim A_j$ は Hensel 対 $(A, I)$ である。
\end{lemma}

\begin{proof}
Categories, Lemma \ref{categories-lemma-limits-products-equalizers} により、積と等化子だけを
考えればよい。積については、主張は Lemma \ref{more-algebra-lemma-product-henselian-pairs} から従う。
そこで、等化子図式
$$
\xymatrix{
(A, I) \ar[r] & (A', I')
\ar@<1ex>[r]^{\varphi} \ar@<-1ex>[r]_{\psi}
&
(A'', I'')  
}
$$
を考える。ここで対 $(A', I')$ と $(A'', I'')$ は Hensel 対である。
対 $(A, I)$ も Hensel 対であることを確認するため、Lemma
\ref{more-algebra-lemma-characterize-henselian-pair} の Gabber の判定法を用いる。
$1 + I$ の各元は $A$ の単元である。実際、単元の逆元の一意性により、これは
$(A', I')$ において確認できる。従って Algebra, Lemma
\ref{algebra-lemma-contained-in-radical} により $I$ は $A$ の Jacobson 根基に含まれる。
そこで、$a_{N - 1}, \dotsc, a_0 \in I$, $N \ge 1$ に対する多項式
% SOURCE-ERRATUM-CANDIDATE P02-E0781: frozen source reindexes Gabber's exponent as N-1 but retains N >= 1, whereas the cited criterion requires exponent at least 1; formula preserved.
$$
f(T) = T^{N - 1}(T - 1) + a_{N - 1} T^{N - 1} + \dotsb + a_1 T + a_0
$$
を $A[T]$ においてとる。$A'[T]$ における $f(T)$ の像は一意な根
$\alpha' \in 1 + I'$ をもち、$A''[T]$ におけるそのさらなる像についても同様である。
従って一意性により $\varphi(\alpha') = \psi(\alpha')$ であり、これは
$\alpha'$ が求めるように $1 + I$ における $f(T)$ の根を定めることを意味する。
\end{proof}

\begin{lemma}
\label{more-algebra-lemma-filtered-colimits-henselian}
Hensel 対であるという性質は、対のフィルター付き余極限をとることで保たれる。
より正確には、$J$ を有向集合、$(A_j, I_j)$ を $J$ 上の Hensel 対の系とする。
このとき、イデアル $I = \colim I_j$ を備えた $A = \colim A_j$ は Hensel 対 $(A, I)$ である。
\end{lemma}

\begin{proof}
$u \in 1 + I$ ならば、ある $j \in J$ に対して $u$ はある
$u_j \in 1 + I_j$ の像である。$I_j$ は $A_j$ の Jacobson 根基に含まれるという仮定と
Algebra, Lemma \ref{algebra-lemma-contained-in-radical} により、$u_j$ は $A_j$ で可逆である。
従って $u$ は $A$ で可逆である。よって同じ補題により、$I$ は $A$ の Jacobson
根基に含まれる。

\medskip\noindent
$f \in A[T]$ をモニック多項式とし、$A/I[T]$ の単位イデアルを生成するモニック多項式
$g_0, h_0 \in A/I[T]$ による因数分解 $\overline{f} = g_0 h_0$ をとる。
ある $g'_0, h'_0 \in A/I[T]$ によって $1 = g_0 g'_0 + h_0 h'_0$ と書く。
$A = \colim A_j$ と $A/I = \colim A_j/I_j$ はフィルター付き余極限なので、
ある $j \in J$ と
% SOURCE-ERRATUM-CANDIDATE P02-E0782: frozen source declares polynomial f_j as an element of A_j rather than A_j[T]; preserved.
$f_j \in A_j$、およびモニック多項式 $g_{j, 0}, h_{j, 0} \in A_j/I_j[T]$ による
因数分解 $\overline{f}_j = g_{j, 0} h_{j, 0}$ で、ある
$g'_{j, 0}, h'_{j, 0} \in A_j/I_j[T]$ に対して
$1 = g_{j, 0} g'_{j, 0} + h_{j, 0} h'_{j, 0}$ を満たし、
$f_j, g_{j, 0}, h_{j, 0}, g'_{j, 0}, h'_{j, 0}$ がそれぞれ
$f, g_0, h_0, g'_0, h'_0$ に写るものを見つけられる。
$(A_j, I_j)$ は Hensel 対なので、$\overline{f}_j = g_{j, 0} h_{j, 0}$ を
$A_j$ 上の因数分解に持ち上げられる。その像を $A$ にとると、対応する
$A$ における因数分解が得られる。従って $(A, I)$ は Hensel 対である。
\end{proof}

\begin{example}[Moret-Bailly]
\label{more-algebra-example-moret-bailly}
Lemma \ref{more-algebra-lemma-filtered-colimits-henselian} は、余極限がフィルター付きでなければ
正しくない。例えば Hensel 対 $(\mathbf{Z}_p, (p))$ と
$(\mathbf{Z}_p, (p))$ の余積をとると、$A = \mathbf{Z}_p \otimes_\mathbf{Z}
\mathbf{Z}_p$ に対する $(A, pA)$ が得られる。これは Hensel 対ではない。
% SOURCE-ERRATUM-CANDIDATE P02-E0784: frozen source says “if (A,pA) where henselian”; Japanese supplies “were”.
$A/pA = \mathbf{F}_p$ なので、もし $(A, pA)$ が Hensel 対ならば、$A$ は局所環で
なければならない。しかし $\Spec(A)$ は非連結である。例えば奇素数 $p$ に対して、
非自明な冪等元
$$
(1/2 \otimes 1)
\left(1 \otimes 1 - (1 + p)^{-1}u \otimes u\right)
$$
がある。ここで $u \in \mathbf{Z}_p$ は $1 + p$ の平方根である。詳細の一部は省略する。
\end{example}

\begin{lemma}
\label{more-algebra-lemma-largest-ideal-henselian}
$A$ を環とする。$(A, I)$ が Hensel 対となる最大のイデアル $I \subset A$ が存在する。
\end{lemma}

\begin{proof}
Lemmas \ref{more-algebra-lemma-henselian-henselian-pair}, \ref{more-algebra-lemma-sum-henselian}, および
\ref{more-algebra-lemma-filtered-colimits-henselian} を組み合わせればよい。
\end{proof}

\begin{lemma}
\label{more-algebra-lemma-irreducible-henselian-pair-connected}
$(A, I)$ を Hensel 対とし、$\mathfrak p \subset A$ を素イデアルとする。
このとき $V(\mathfrak p + I)$ は連結である。
\end{lemma}

\begin{proof}
Lemma \ref{more-algebra-lemma-integral-over-henselian-pair} により、
% SOURCE-ERRATUM-CANDIDATE P02-E0783: frozen source writes I + p/p instead of the ideal (I+p)/p in A/p; preserved.
$(A/\mathfrak p, I + \mathfrak p/\mathfrak p)$ は Hensel 対である。
従って、$(A, I)$ が Hensel 対で $A$ が整域ならば
$\Spec(A/I) = V(I)$ は連結である、という主張を証明すれば十分である。
そうでなければ、Algebra, Lemma \ref{algebra-lemma-characterize-spec-connected} により
$A/I$ は非自明な冪等元をもつ。Lemma \ref{more-algebra-lemma-characterize-henselian-pair} により、
これは $A$ が非自明な冪等元をもつことを含意する。これは矛盾である。
\end{proof}

\section{対の Hensel 化}
\label{more-algebra-section-henselization}

\noindent
Section \ref{more-algebra-section-henselian-pairs} で始めた議論を続ける。

\begin{lemma}
\label{more-algebra-lemma-henselization}
包含関手
$$
\text{category of henselian pairs}
\longrightarrow
\text{category of pairs}
$$
は左随伴 $(A, I) \mapsto (A^h, I^h)$ をもつ。
\end{lemma}

\begin{proof}
$(A, I)$ を対とする。同型 $A/I \to B/IB$ を誘導する \'etale 環写像
$A \to B$ からなる圏を $\mathcal{C}$ とする。圏 $\mathcal{C}$ が有向であり、
イデアル $I^h = IA^h$ を備えた $A^h = \colim_{B \in \mathcal{C}} B$ が
求める随伴を与えることを示す。

\medskip\noindent
まず $\mathcal{C}$ が有向であることを示す
（Categories, Definition \ref{categories-definition-directed}）。
$\text{id} : A \to A$ は対象なので、この圏は空でない。$B$ と $B'$ が
$\mathcal{C}$ の二つの対象ならば、$B'' = B \otimes_A B'$ は $\mathcal{C}$ の対象であり
（Algebra, Lemma \ref{algebra-lemma-etale} を用いよ）、射 $B \to B''$ と
$B' \to B''$ がある。$f, g : B \to B'$ が $\mathcal{C}$ の対象間の二つの写像ならば、
余等化子は
$$
(B' \otimes_{f, B, g} B') \otimes_{(B' \otimes_A B')} B'
$$
である。これは Algebra, Lemmas \ref{algebra-lemma-etale} および
\ref{algebra-lemma-map-between-etale} により $A$ 上 \'etale である。
従って圏 $\mathcal{C}$ は有向である。

\medskip\noindent
$\mathcal{C}$ のすべての対象 $B$ に対して $B/IB = A/I$ なので、
$A^h/I^h = A^h/IA^h = \colim B/IB = \colim A/I = A/I$ である。

\medskip\noindent
次に、$I^h = IA^h$ を備えた $A^h = \colim_{B \in \mathcal{C}} B$ が Hensel 対で
あることを示す。そのため Lemma \ref{more-algebra-lemma-characterize-henselian-pair} の条件 (2) を
確認する。すなわち、\'etale 環写像 $A^h \to A'$ と $A^h$-代数写像
$\sigma : A' \to A^h/I^h$ が与えられたとする。このとき、ある
$B \in \mathcal{C}$ と \'etale 環写像 $B \to B'$ が存在して
$A' = B' \otimes_B A^h$ となる。Algebra, Lemma \ref{algebra-lemma-etale} を参照せよ。
$A^h/I^h = A/I$ なので、写像 $\sigma$ は $A$-代数写像 $s : B' \to A/I$ を誘導する。
このとき $A/I$-代数として $B'/IB' = A/I \times C$ である。ここで $C$ は
% SOURCE-ERRATUM-CANDIDATE P02-E0785: frozen source calls C the kernel of a unital projection although the kernel is the nonunital ideal 0×C; preserved.
$s$ が誘導する写像 $B'/IB' \to A/I$ の核である。
$(1, 0) \in A/I \times C$ に写る $g \in B'$ をとる。このとき
$B \to B'_g$ は \'etale で、$A/I \to B'_g/IB'_g$ は同型である。すなわち
$B'_g$ は $\mathcal{C}$ の対象である。従って標準写像 $B'_g \to A^h$ が得られ、
$$
\vcenter{
\xymatrix{
B'_g \ar[r] & A^h \\
B \ar[u] \ar[ur]
}
}
\quad\text{and}\quad
\vcenter{
\xymatrix{
B' \ar[r] \ar[rrd]_s & B'_g \ar[r] & A^h \ar[d] \\
& & A/I
}
}
$$
は可換である。これは求めるように $\sigma$ と両立する写像
$A' = B' \otimes_B A^h \to A^h$ を誘導する。

\medskip\noindent
$(A, I) \to (A', I')$ を対の射とし、$(A', I')$ は Hensel 対であるとする。
一意な因数分解 $A \to A^h \to A'$ が存在することを示す。これで証明が完了する。
実際、$\mathcal{C}$ の各 $A \to B$ に対して、環写像
$A' \to B' = A' \otimes_A B$ は \'etale で、同型
$A'/I' \to B'/I'B'$ を誘導する。従って Lemma
\ref{more-algebra-lemma-characterize-henselian-pair} により切断 $\sigma_B : B' \to A'$ が存在する。
$\mathcal{C}$ の射 $B_1 \to B_2$ が与えられたとき、図式
$$
\xymatrix{
B'_1 \ar[rr] \ar[rd]_{\sigma_{B_1}} & &
B'_2 \ar[ld]^{\sigma_{B_2}} \\
& A'
}
$$
が可換であると主張する。これは、$\mathcal{C}$ の各 $B$ に対して切断
$\sigma_B$ が一意な $A'$-代数写像 $B' \to A'$ であることを証明すれば従う。
ある環 $R$ に対して $B' \otimes_{A'} B' = B' \times R$ である。
Algebra, Lemma \ref{algebra-lemma-diagonal-unramified} を参照せよ。この場合
$B'/I'B' = A'/I'$ なので $R/I'R = 0$ である。従って二つの $A'$-代数写像
$\sigma_B, \sigma_B' : B' \to A'$ が与えられると、
$e = (\sigma_B \otimes \sigma_B')(0, 1) \in A'$ は $I'$ に含まれる冪等元である。
Lemma \ref{more-algebra-lemma-idempotents-determined-modulo-radical} により $e = 0$ と結論する。
従って求めるように $\sigma_B = \sigma_B'$ である。この可換性を用いて
$$
A^h = \colim_{B \in \mathcal{C}} B \to
\colim_{B \in \mathcal{C}} A' \otimes_A B \xrightarrow{\colim \sigma_B} A'
$$
を得る。写像 $\sigma_B$ の一意性は、この写像も一意であることを保証する。
従って $(A, I) \mapsto (A^h, I^h)$ が求める随伴である。
\end{proof}

\begin{lemma}
\label{more-algebra-lemma-henselization-flat}
$(A, I)$ を対とし、$(A^h, I^h)$ を Lemma \ref{more-algebra-lemma-henselization} のものとする。
このとき $A \to A^h$ は平坦、$I^h = IA^h$ であり、すべての $n$ に対して
$A/I^n \to A^h/I^nA^h$ は同型である。
\end{lemma}

\begin{proof}
Lemma \ref{more-algebra-lemma-henselization} の証明で、$A^h$ は、同型
$A/I \to B/IB$ を誘導する \'etale $A$-代数 $B$ のフィルター付き余極限であり、
$I^h = IA^h$ であることを見た。\'etale 環写像は平坦なので
（Algebra, Lemma \ref{algebra-lemma-etale}）、Algebra, Lemma
\ref{algebra-lemma-colimit-flat} により $A \to A^h$ は平坦である。
各 $A \to B$ は平坦なので、写像 $A/I^n \to B/I^nB$ も同型である
（例えば Algebra, Lemma \ref{algebra-lemma-lift-basis} による）。
余極限をとると、求める $A/I^n = A^h/I^nA^h$ が得られる。
\end{proof}

\begin{lemma}
\label{more-algebra-lemma-henselization-local-ring}
\begin{slogan}
% SOURCE-ERRATUM-CANDIDATE P02-E0786: frozen slogan omits “of” in “Compatibility henselization”; Japanese supplies the relation.
対の Hensel 化と局所環の Hensel 化の両立性。
\end{slogan}
Lemma \ref{more-algebra-lemma-henselization} の関手は、局所環 $(A, \mathfrak m)$ にその
Hensel 化を対応させる。
\end{lemma}

\begin{proof}
$(A^h, \mathfrak m^h)$ を、Lemma \ref{more-algebra-lemma-henselization} で構成した対
$(A, \mathfrak m)$ の Hensel 化とする。Lemma \ref{more-algebra-lemma-henselization-flat} により
$\mathfrak m^h = \mathfrak m A^h$ は極大イデアルであり、これは Jacobson 根基に
含まれるので、$A^h$ は極大イデアル $\mathfrak m^h$ をもつ局所環である。
ここから証明を終える方法は二つある。

\medskip\noindent
第1の証明：Algebra, Lemma \ref{algebra-lemma-henselization} の証明における
余極限としての構成は、Lemma \ref{more-algebra-lemma-henselization} で $A^h$ を構成するために
用いた余極限と同じであることに注意する。
第2の証明：Hensel 化 $A \to S$ と Lemma \ref{more-algebra-lemma-henselization} の
$A \to A^h$ はともに局所環準同型であり、$S$ と $A^h$ はともに \'etale
$A$-代数のフィルター付き余極限であり、$S$ と $A^h$ はともに Hensel 局所環であり、
$S$ と $A^h$ の剰余体はともに $\kappa(\mathfrak m)$ に等しい
（後者については Lemma \ref{more-algebra-lemma-henselization-flat} による）。従って Algebra, Lemma
\ref{algebra-lemma-uniqueness-henselian} により、これらは標準的に同型である。
\end{proof}

\begin{lemma}
\label{more-algebra-lemma-henselization-Noetherian-pair}
\begin{slogan}
Noether 対の Hensel 化は、同じ完備化をもつ Noether 対である。
\end{slogan}
$(A, I)$ を $A$ が Noether 環である対とし、$(A^h, I^h)$ を Lemma
\ref{more-algebra-lemma-henselization} のものとする。このとき $I$-進完備化の写像
$$
A^\wedge \to (A^h)^\wedge
$$
は同型である。さらに $A^h$ は Noether 環であり、写像
$A \to A^h \to A^\wedge$ は平坦で、$A^h \to A^\wedge$ は忠実平坦である。
\end{lemma}

\begin{proof}
最初の主張は Lemma \ref{more-algebra-lemma-henselization-flat} の直ちに従う帰結であり、実際
$A$ が Noether 環であるという仮定なしに成り立つ。Lemma
\ref{more-algebra-lemma-henselization} の証明で、$A^h$ は、同型 $A/I \to B/IB$ を誘導する
\'etale $A$-代数 $B$ のフィルター付き余極限であることを見た。このような各
$A \to B$ に対して、誘導される写像 $A^\wedge \to B^\wedge$ は同型である
（Lemma \ref{more-algebra-lemma-henselization-flat} の証明を参照）。Algebra, Lemma
\ref{algebra-lemma-completion-flat} により、環写像
$B \to A^\wedge = B^\wedge = (A^h)^\wedge$ は各 $B$ に対して平坦である。
従って Algebra, Lemma \ref{algebra-lemma-colimit-rings-flat} により
$A^h \to A^\wedge = (A^h)^\wedge$ は平坦である。$I^h = IA^h$ は $A^h$ の
Jacobson 根基に含まれ、$A^h \to A^\wedge$ は同型 $A^h/I^h \to A/I$ を誘導するので、
Algebra, Lemma \ref{algebra-lemma-ff} により $A^h \to A^\wedge$ は忠実平坦である。
Algebra, Lemma \ref{algebra-lemma-completion-Noetherian-Noetherian} により
環 $A^\wedge$ は Noether 環である。従って Algebra, Lemma
\ref{algebra-lemma-descent-Noetherian} により $A^h$ は Noether 環であると結論する。
\end{proof}

\begin{lemma}
\label{more-algebra-lemma-henselization-colimit}
$(A, I) = \colim (A_i, I_i)$ を対のフィルター付き余極限とする。Lemma
\ref{more-algebra-lemma-henselization} の関手は $A^h = \colim A_i^h$ および
$I^h = \colim I_i^h$ を与える。
\end{lemma}

\noindent
この補題はフィルター付きでない余極限に対しては正しくない。Example
\ref{more-algebra-example-moret-bailly} を参照せよ。

\begin{proof}
Categories, Lemma \ref{categories-lemma-adjoint-exact} により、$(A^h, I^h)$ は
Hensel 対の圏における系 $(A_i^h, I_i^h)$ の余極限である。従って Hensel 対
$(B, J)$ に対して
$$
\Mor((A^h, I^h), (B, J)) =
\lim \Mor((A_i^h, I_i^h), (B, J)) =
\Mor(\colim (A_i^h, I_i^h), (B, J))
$$
である。ここで余極限は対の圏におけるものである。余極限はフィルター付きなので、
対の圏において
$\colim (A_i^h, I_i^h) = (\colim A_i^h, \colim I_i^h)$ を得る。詳細は省略する。
% SOURCE-ERRATUM-CANDIDATE P02-E0787: frozen source omits “that” in “Again using the colimit is filtered”; Japanese supplies the connective.
余極限がフィルター付きであることを再び用いると、これは Hensel 対である
（Lemma \ref{more-algebra-lemma-filtered-colimits-henselian}）。従って Yoneda の補題により
$(A^h, I^h) = (\colim A_i^h, \colim I_i^h)$ を得る。
\end{proof}

\begin{lemma}
\label{more-algebra-lemma-henselization-change-ideal}
\begin{slogan}
対の Hensel 化はイデアルの根基のみに依存する。
\end{slogan}
$A$ をイデアル $I$, $J$ をもつ環とする。$V(I) = V(J)$ ならば、Lemma
\ref{more-algebra-lemma-henselization} の関手は対 $(A, I)$ と対 $(A, J)$ に同じ環を与える。
\end{lemma}

\begin{proof}
対 $(A, I)$ から Lemma \ref{more-algebra-lemma-henselization} によって得られる対を
$(A', IA')$ とする。Lemma \ref{more-algebra-lemma-henselization-flat} を参照せよ。
対 $(A, J)$ から Lemma \ref{more-algebra-lemma-henselization} によって得られる対を
$(A'', JA'')$ とする。Lemma \ref{more-algebra-lemma-change-ideal-henselian-pair} により、
$(A', JA')$ は Hensel 対であり、$(A'', IA'')$ も Hensel 対である。
構成の普遍性により、一意な $A$-代数写像 $A'' \to A'$ と $A' \to A''$ を得る。
一意性から、これらは互いに逆である。
\end{proof}

\begin{lemma}
\label{more-algebra-lemma-henselization-integral}
\begin{slogan}
Hensel 化は整基底変換と可換である。
\end{slogan}
$(A, I) \to (B, J)$ を $V(J) = V(IB)$ を満たす対の写像とする。
$(A^h , I^h) \to (B^h, J^h)$ を Hensel 化上の誘導写像とする
（Lemma \ref{more-algebra-lemma-henselization}）。$A \to B$ が整ならば、誘導写像
$A^h \otimes_A B \to B^h$ は同型である。
\end{lemma}

\begin{proof}
Lemma \ref{more-algebra-lemma-henselization-change-ideal} により $J = IB$ と仮定してよい。
Lemma \ref{more-algebra-lemma-integral-over-henselian-pair} により、対
$(A^h \otimes_A B, I^h(A^h \otimes_A B))$ は Hensel 対である。
$(B^h, IB^h)$ の普遍性により、写像 $B^h \to A^h \otimes_A B$ を得る。
この写像が補題の写像の逆であることの証明は省略する。
\end{proof}


\begin{lemma}
\label{more-algebra-lemma-CRT-henselization}
$I_1, I_2, \dots, I_n \subset A$ を環 $A$ のイデアルとする。
$i \not = j$, $1 \leq i, j \leq n$ に対して $I_i + I_j = A$ と仮定する。
$A^h$ を対 $(A, I_1 \cap \ldots \cap I_n)$ の Hensel 化、$A_i^h$ を対
$(A, I_i)$ の Hensel 化とする。このとき自然な射
$$
A^h \to \prod\nolimits_{i = 1, \ldots, n} A^h_i
$$
は同型である。
\end{lemma}

\begin{proof}
この写像は構成の関手性から得られる。$n = 1$ ならば結果は直ちに従う。
$n > 1$ ならば、イデアルに関する仮定から
$(I_1 \cap \ldots \cap I_{n - 1}) + I_n = A$ であることに注意する。
従って $n = 2$ に対して補題を証明すれば、すべての $n$ に対して帰納法で従う。

\medskip\noindent
$n = 2$ の場合。$f_1 + f_2 = 1$ を満たす $f_i \in I_i$ を選ぶ。
$A^h$ は、$A \to B$ が \'etale で、
% SOURCE-ERRATUM-CANDIDATE P02-E0788: frozen source omits parentheses in A/I_1 \cap I_2; exact formal span preserved.
$A/I_1 \cap I_2 \to B/(I_1 \cap I_2)B$ が同型となる $A$-代数 $B$ の余極限として
構成されることを思い出す。Lemma \ref{more-algebra-lemma-henselization} の証明を参照せよ。
このような $A$-代数 $B$ の圏を $Hens$ と書く。同様に、$i = 1, 2$ に対して、
$A/I_i \to B/I_iB$ が同型となる \'etale $A$-代数 $B$ の圏を $Hens_i$ と書く。

\medskip\noindent
$f_2$ が可逆となるような $B$ からなる部分圏
$Hens_1(f_2) \subset Hens_1$ を考える。これは共終部分圏である。実際、
$B$ が $Hens_1$ に属すれば $B_{f_2}$ も $Hens_1$ に属する
（ヒント：$f_2$ は $I_1$ を法として単元なので
$B_{f_2}/I_1B_{f_2} = (B/I_1B)_{f_2} = B/I_1B$）。同様に部分圏
$Hens_2(f_1) \subset Hens_2$ も共終である。従って $A_1^h$ と $A_2^h$ の計算には
% SOURCE-ERRATUM-CANDIDATE P02-E0789: frozen source says “this smaller categories”; Japanese supplies plural agreement.
これらのより小さい圏を用いてよい。Categories, Lemma
\ref{categories-lemma-cofinal} を参照せよ。

\medskip\noindent
関手
$$
Hens_1(f_2) \times Hens_2(f_1) \to Hens,\quad (B_1, B_2) \mapsto B_1 \times B_2
$$
を考える。次の計算により、この関手は意味をもつ。
\begin{align*}
(B_1 \times B_2)/(I_1 \cap I_2)(B_1 \times B_2)
& =
B_1/(I_1 \cap I_2)B_1 \times B_2/(I_1 \cap I_2)B_2 \\
& =
B_1/I_1B_1 \times B_2/I_2B_2 \\
& = A/I_1 \times A/I_2 \\
& = A/(I_1 \cap I_2)
\end{align*}
第2の等式は、$f_2 \in I_2$ が $B_1$ において可逆なので
$(I_1 \cap I_2)B_1 = I_1B_1 \cap I_2B_1 = I_1B_1$ であることから従う
（Algebra, Lemma \ref{algebra-lemma-flat-intersect-ideals} も用いた）。
最後の等式は Algebra, Lemma \ref{algebra-lemma-chinese-remainder} による。
証明を終えるには、表示された関手が共終であることを示せば十分である
（上で引用した補題を参照）。Categories, Definition
\ref{categories-definition-cofinal} の第1の条件は、$Hens$ の対象 $B$ が
関手の像に属する $B_{f_2} \times B_{f_1}$ に写るので成り立つ。
第2の条件も常に成り立つ。実際、$B \to B_1 \times B_2$ と
$B \to B'_1 \times B'_2$ が与えられ、$B_1, B'_1 \in Hens_1(f_2)$ および
$B_2, B'_2 \in Hens_2(f_1)$ ならば、$Hens_1(f_2)$ に属する
$B''_1 = B_1 \otimes_B B'_1$ と、同様に $Hens_2(f_1)$ に属する
% SOURCE-ERRATUM-CANDIDATE P02-E0790: frozen source defines B''_2 self-referentially using B''_2 instead of B'_2; preserved.
$B''_2 = B_2 \otimes_B B''_2$ を置くことで、可換図式に入る写像
$B \to B''_1 \times B''_2$ を見つけられる。
$$
\xymatrix{
& B \ar[ld] \ar[d] \ar[rd] \\
B_1 \times B_2 \ar[r] & B_1'' \times B_2'' & B_1' \times B_2' \ar[l]
}
$$
これで証明が完了する。
\end{proof}

\section{持ち上げと Hensel 対}
\label{more-algebra-section-lifting-henselian-pairs}

\noindent
この節では主に Sections \ref{more-algebra-section-lifting} と
\ref{more-algebra-section-henselian-pairs} の結果を組み合わせる。

\begin{lemma}
\label{more-algebra-lemma-lift-finite-projective-module}
$(R, I)$ を Hensel 対とする。写像
$$
P \longrightarrow P/IP
$$
は、有限生成射影 $R$-加群の同型類の集合と有限生成射影 $R/I$-加群の同型類の
集合との間の全単射を誘導する。特に、任意の有限生成射影 $R/I$-加群は、
ある有限生成射影 $R$-加群 $P$ に対する $P/IP$ と同型である。
\end{lemma}

\begin{proof}
まず最後の主張を証明する。$\overline{P}$ を有限生成射影 $R/I$-加群とする。
$R/I = R'/IR'$ を満たす、$R$ 上 \'etale なある $R'$ 上の有限生成射影加群
$P'$ で、$P'/IP'$ が $\overline{P}$ と同型となるものを見つけられる。
Lemma \ref{more-algebra-lemma-lift-projective-module} を参照せよ。$(R, I)$ は Hensel 対なので、
\'etale 環写像 $R \to R'$ は切断 $\tau : R' \to R$ をもつ
（Lemma \ref{more-algebra-lemma-characterize-henselian-pair}）。
$P = P' \otimes_{R', \tau} R$ と置くと、$P/IP$ は $\overline{P}$ と同型である。
もちろん、これは補題の主張にある写像が全射であることを示す。

\medskip\noindent
単射性。$P_1$ と $P_2$ を有限生成射影 $R$-加群とし、$R/I$-加群として
$P_1/IP_1 \cong P_2/IP_2$ であるとする。$P_1$ は射影的なので、与えられた同型を
持ち上げる $R$-加群写像 $u : P_1 \to P_2$ を見つけられる。
すると Nakayama の補題（Algebra, Lemma \ref{algebra-lemma-NAK}）により $u$ は全射である。
同様に全射 $v : P_2 \to P_1$ を見つけられる。Algebra, Lemma
\ref{algebra-lemma-fun} により写像 $v \circ u$ は同型であり、$u$ も同型と結論する。
\end{proof}

\begin{lemma}
\label{more-algebra-lemma-finite-etale-equivalence}
$(A, I)$ を Hensel 対とする。関手 $B \to B/IB$ は、有限 \'etale $A$-代数と
有限 \'etale $A/I$-代数との間の同値を定める。
\end{lemma}

\begin{proof}
$B, B'$ を $A$ 上有限 \'etale な二つの $A$-代数とする。このとき
$B' \to B'' = B \otimes_A B'$ も有限 \'etale である
（Algebra, Lemmas \ref{algebra-lemma-etale} および
\ref{algebra-lemma-base-change-integral}）。次の間に $1$ 対 $1$ の対応がある。
\begin{enumerate}
\item $A$-代数写像 $B \to B'$,
\item $B' \to B''$ の切断、かつ
\item $B' \to B'' \to eB''$ が同型となる $B''$ の冪等元 $e$.
\end{enumerate}
(2) と (3) の間の全単射は、$\sigma : B'' \to B'$ を、$(1 - e)$ が
$\sigma$ の核を生成する冪等元となるような $e$ に送る。この冪等元は Algebra, Lemmas
\ref{algebra-lemma-map-between-etale} および
\ref{algebra-lemma-surjective-flat-finitely-presented} により存在する。
同様に、$A/I$-代数写像 $B/IB \to B'/IB'$ と、
$B'/IB' \to B''/IB'' \to \overline{e}(B''/IB'')$ が同型となる
$B''/IB''$ の冪等元 $\overline{e}$ との間にも対応がある。しかし
$B''/IB''$ の各冪等元 $\overline{e}$ は $B''$ の冪等元 $e$ に一意に持ち上がる
（Lemma \ref{more-algebra-lemma-characterize-henselian-pair}）。さらに、
$B'/IB' \to \overline{e}(B''/IB'')$ が同型ならば、Nakayama の補題
（Algebra, Lemma \ref{algebra-lemma-NAK}）により $B' \to eB''$ も同型である。
こうしてこの関手が充満忠実であることが分かる。

\medskip\noindent
本質的全射性。$A/I \to C$ を有限 \'etale 写像とする。Algebra, Lemma
\ref{algebra-lemma-lift-etale} により、$B/IB \cong C$ となる \'etale 写像
$A \to B$ が存在する。$B'$ を $B$ における $A$ の整閉包とする。
Lemma \ref{more-algebra-lemma-helper-finite-type} により、ある環 $C'$ に対して
$B'/IB' = C \times C'$ であり、$(1, 0) \in C \times C'$ に写るある
$g \in B'$ に対して $B'_g \cong B_g$ である。冪等元は持ち上がるので
（Lemma \ref{more-algebra-lemma-characterize-henselian-pair}）、
$C = B'_1/IB'_1$ および $C' = B'_2/IB'_2$ を満たす
$B' = B'_1 \times B'_2$ を得る。$B'_1$ における $g$ の像は可逆である。
% SOURCE-ERRATUM-CANDIDATE P02-E0791: frozen source uses undefined B_2 instead of B'_2 in the localized decomposition; preserved.
すると $B_g = B'_g = B'_1 \times (B_2)_g$ であり、これは $A \to B'_1$ が
\'etale であることを含意する。従って $B'_1$ は $A$ 上有限 \'etale である
（例えば Algebra, Lemma
\ref{algebra-lemma-characterize-finite-in-terms-of-integral} により、整かつ \'etale ならば
有限 \'etale である）。これで証明が完了する。
\end{proof}

\begin{lemma}
\label{more-algebra-lemma-lift-smooth-henselian}
$R$ を環、$S$ を滑らかな $R$-代数とする。$A$ は $R$-代数で、$(A,I)$ は
Hensel 対であると仮定する。このとき任意の $R$-代数写像 $S \to A/I$ は
$R$-代数写像 $S \to A$ に持ち上げられる。
\end{lemma}

\begin{proof}
$\tau : S \to A/I$ を $R$-代数写像とする。Algebra, Lemma
\ref{algebra-lemma-base-change-smooth} により、$S \otimes_R A$ は滑らかな
$A$-代数である。従って Lemma \ref{more-algebra-lemma-lift-section-smooth-morphism} により、
誘導写像 $S \otimes_R A \to A/I$ を $A$-代数準同型
% SOURCE-ERRATUM-CANDIDATE P02-E0792: frozen source misspells “homomorphism” as “homorphism”; Japanese supplies the correct term.
$S \otimes_R A \to A'$ に持ち上げられる。ここで $A \to A'$ は \'etale で、
同型 $A/I \to A'/IA'$ を誘導する。$(A, I)$ は Hensel 対なので、Lemma
\ref{more-algebra-lemma-characterize-henselian-pair} により $A$-代数写像 $A' \to A$ がある。
合成 $S \to S \otimes_R A \to A' \to A$ が求める持ち上げである。
\end{proof}

\begin{lemma}
\label{more-algebra-lemma-lim-finite-projective-gives-finite-projective}
$A = \lim A_n$ を環の逆系 $(A_n)$ の極限とする。$A_n$-加群 $M_n$ と
$A_{n + 1}$-加群写像 $M_{n + 1} \to M_n$ が与えられているとし、次を仮定する。
\begin{enumerate}
\item 遷移写像 $A_{n + 1} \to A_n$ は全射で、その核は局所冪零である、
\item $M_1$ は有限生成射影 $A_1$-加群である、
\item $M_n$ は有限平坦 $A_n$-加群である、かつ
\item 写像は同型 $M_{n + 1} \otimes_{A_{n + 1}} A_n \to M_n$ を誘導する。
\end{enumerate}
このとき $M = \lim M_n$ は有限生成射影 $A$-加群であり、すべての $n$ に対して
$M \otimes_A A_n \to M_n$ は同型である。
\end{lemma}

\begin{proof}
Lemma \ref{more-algebra-lemma-limit-henselian} により、対 $(A, \Ker(A \to A_1))$ は Hensel 対である。
Lemma \ref{more-algebra-lemma-lift-finite-projective-module} により、有限生成射影 $A$-加群 $P$ と
同型 $P \otimes_A A_1 \to M_1$ を選べる。$P$ は射影的なので、$A$-加群写像
$P \to M_1$ を $A$-加群写像 $P \to M_2$, $P \to M_3$ などへ逐次的に
持ち上げられる。従って写像
$$
P \longrightarrow M
$$
を得る。$P$ は有限生成射影なので、ある $m \geq 0$ と $A$-加群 $Q$ に対して
$A^{\oplus m} = P \oplus Q$ と書ける。$A = \lim A_n$ なので
$P = \lim P \otimes_A A_n$ と結論する。従って表示された $A$-加群写像が同型で
あることを示すには、写像 $P \otimes_A A_n \to M_n$ が同型であることを示せば十分である。
Lemma \ref{more-algebra-lemma-lift-projective} により、$M_n$ は有限生成射影加群である。
Lemma \ref{more-algebra-lemma-isomorphic-finite-projective-lifts} により、写像
$P \otimes_A A_n \to M_n$ は同型である。
\end{proof}

\section{絶対整閉包}
\label{more-algebra-section-absolute-integral-closure}

\noindent
ここで定義を与える。

\begin{definition}
\label{more-algebra-definition-absolutely-integrally-closed}
任意のモニック多項式 $f \in A[T]$ が一次因子の積であるとき、環 $A$ は
{\it 絶対整閉}であるという。
\end{definition}

\noindent
注意が必要である。$f$ は一次因子の積として多くの異なる仕方で書けることがある。

\begin{lemma}
\label{more-algebra-lemma-absolutely-integrally-closed}
$A$ を環とする。次は同値である。
\begin{enumerate}
\item $A$ は絶対整閉である、かつ
\item 任意のモニック多項式 $f \in A[T]$ は $A$ に根をもつ。
\end{enumerate}
\end{lemma}

\begin{proof}
省略する。
\end{proof}








\begin{lemma}
\label{more-algebra-lemma-absolutely-integrally-closed-quotient-localization}
$A$ を絶対整閉環とする。
\begin{enumerate}
\item $A$ の任意の商環 $A/I$ は絶対整閉である。
\item 任意の局所化 $S^{-1}A$ は絶対整閉である。
\end{enumerate}
\end{lemma}

\begin{proof}
省略する。
\end{proof}

\begin{lemma}
\label{more-algebra-lemma-integrally-closed-in-absolutely-integrally-closed}
$A$ を環とする。$S \subset A$ を非零因子からなる乗法的部分集合とする。
$S^{-1}A$ が絶対整閉であり、$A \subset S^{-1}A$ が $S^{-1}A$ の中で整閉ならば、
$A$ は絶対整閉である。
\end{lemma}

\begin{proof}
省略する。
\end{proof}

\begin{lemma}
\label{more-algebra-lemma-normal-domain-absolutely-integrally-closed}
$A$ を正規整域とする。このとき $A$ が絶対整閉であることと、その分数体が
代数閉であることは同値である。
\end{lemma}

\begin{proof}
体が代数閉であることと、それが環として絶対整閉であることは同値である。
従って補題は Lemmas
\ref{more-algebra-lemma-absolutely-integrally-closed-quotient-localization} および
\ref{more-algebra-lemma-integrally-closed-in-absolutely-integrally-closed} から従う。
\end{proof}

\begin{lemma}
\label{more-algebra-lemma-construct-absolute-integral-closure}
任意の環 $A$ に対して、次を満たす拡大 $A \subset B$ が存在する。
\begin{enumerate}
\item $B$ は有限自由 $A$-代数のフィルター付き余極限である、
\item $B$ は $A$-加群として自由である、かつ
\item $B$ は絶対整閉である。
\end{enumerate}
\end{lemma}

\begin{proof}
$I$ を $A$ 上のモニック多項式全体の集合とする。$i \in I$ に対し、
% SOURCE-ERRATUM-CANDIDATE P02-E0793: frozen source says “denote x_i a variable”, omitting “by/as”; Japanese supplies the grammatical relation.
$x_i$ を変数とし、$P_i$ を変数 $x_i$ における対応するモニック多項式とする。
そして
$$
F(A) = A[x_i; i \in I]/(P_i; i \in I)
$$
とおく。記号が示唆するように、$F$ は環の圏からそれ自身への関手である。
$A \subset F(A)$ であること、$F(A)$ が $A$-加群として自由であること、および
$F(A)$ が有限自由 $A$-代数のフィルター付き余極限であることに注意せよ。次に
$$
B = \colim F^n(A)
$$
とおく。ここで遷移写像は包含
$F^n(A) \subset F(F^n(A)) = F^{n + 1}(A)$
である。係数が $B$ に属する任意のモニック多項式は、実際にはある $n$ に対して
係数が $F^n(A)$ に属し、従って構成により $F^{n + 1}(A)$ に解をもつ。
Lemma \ref{more-algebra-lemma-absolutely-integrally-closed} により、これは $B$ が絶対整閉で
あることを含意する。他の性質の証明は省略する。
\end{proof}

\begin{lemma}
\label{more-algebra-lemma-absolutely-integrally-closed-strictly-henselian}
$A$ を絶対整閉環とする。$\mathfrak p \subset A$ を素イデアルとする。
このとき局所環 $A_\mathfrak p$ は強 Hensel 環である。
\end{lemma}

\begin{proof}
Lemma \ref{more-algebra-lemma-absolutely-integrally-closed-quotient-localization} により、
$A$ が局所環で $\mathfrak p$ がその極大イデアルであると仮定してよい。
剰余体は Lemma
\ref{more-algebra-lemma-absolutely-integrally-closed-quotient-localization} により代数閉である。
任意のモニック多項式は一次因子に完全分解するので、Algebra, Definition
\ref{algebra-definition-henselian} が直接適用される。
\end{proof}

\begin{lemma}
\label{more-algebra-lemma-absolutely-integrally-closed-henselian-pair}
$A$ を絶対整閉環とする。$I \subset A$ をイデアルとする。このとき
$(A, I)$ が Hensel 対であるための必要十分条件は、次が成り立つことである。
\begin{enumerate}
\item $I$ は $A$ の Jacobson 根基に含まれる、
\item $A \to A/I$ は冪等元の集合の間の全単射を誘導する。
\end{enumerate}
\end{lemma}

\begin{proof}
$f \in A[T]$ をモニック多項式とし、
$f \bmod I = g_0 h_0$ を $A/I$ 上の因数分解で、$g_0$, $h_0$ はモニックであり、
さらに $g_0$ と $h_0$ が $A/I[T]$ の単位イデアルを生成するものとする。これは
$$
A/I[T]/(f) = A/I[T]/(g_0) \times A/I[T]/(h_0)
$$
を意味する。
% SOURCE-ERRATUM-CANDIDATE P02-E0794: frozen source misspells “corresponding” as “correspoing”; Japanese supplies the intended word.
$e \in A/I[T]/(f)$ を、右辺の環の冪等元 $(1, 0)$ に対応する元とする。
$a_i \in A$ を用いて $f = (T - a_1) \ldots (T - a_d)$ と書く。
各 $i \in \{1, \ldots, d\}$ に対し、$A$-代数写像
$\varphi_i : A[T]/(f) \to A$, $T \mapsto a_i$ を得て、これは同様の
$A/I$-代数写像 $\overline{\varphi}_i : A/I[T]/(f) \to A/I$ を誘導する。
$e_i = \overline{\varphi}_i(e) \in A/I$ とおく。これらは冪等元である。
仮定 (2) により、$e_i$ を $A$ の冪等元に持ち上げられる。これは
$A = \prod A_j$ を環の有限積として書き、各 $A_j/IA_j$ において各 $e_i$ が
$0$ または $1$ となるようにできることを意味する。いくつかの詳細は省略する。
$A_j$ は $A$ の因子環として絶対整閉であることに注意せよ。
$f$ の $A_j/IA_j$ 上の因数分解を $A_j$ 上へ持ち上げれば十分である。
これにより次段落で論じる状況へ帰着する。

\medskip\noindent
$i = 1, \ldots, r$ に対して $e_i = 1$、$i = r + 1, \ldots, d$ に対して
$e_i = 0$ と仮定する。$(g_0, h_0) = A/I[T]$ から、
$g_0 k_0 + h_0 l_0 = 1$ を満たす $k_0, l_0 \in A/I[T]$ が存在する。
$e = h_0 l_0$ および $e_i = h_0(a_i) l_0(a_i)$ である。
従って $i = 1, \dots ,r$ に対して $h_0(a_i)$ は単元である。
$f(a_i) = 0$ なので $0 = h_0(a_i)g_0(a_i)$ であり、
$i = 1, \ldots, r$ に対して $g_0(a_i) = 0$ と結論する。
従って $(T - a_1)$ は $A/I[T]$ において $g_0$ を割り、
$g_0 = (T - a_1) g_0'$ と書ける。
$f' = (T - a_2) \ldots (T - a_d)$ および $h'_0 = h_0$ とおく。
% SOURCE-ERRATUM-CANDIDATE P02-E0795: frozen source repeats “over” in “over over A”; Japanese supplies the intended relation once.
$d$ に関する帰納法により、因数分解 $f' \bmod I = g'_0 h'_0$ を
$A$ 上の因数分解 $f' = g' h'$ に持ち上げられ、これは所望のように
因数分解 $f \bmod I = g_0 h_0$ を持ち上げる因数分解
$f = (T - a_1) g' h'$ を与える。
\end{proof}

\section{自己随伴環}
\label{more-algebra-section-auto-ass}

\noindent
この内容の一部は \cite{Autour} にある。

\begin{definition}
\label{more-algebra-definition-auto-ass}
環 $R$ が局所環であり、その極大イデアル $\mathfrak m$ が $R$ に弱随伴するとき、
$R$ は {\it 自己随伴}であるという。
\end{definition}

\begin{lemma}
\label{more-algebra-lemma-auto-ass-implies-P}
自己随伴環 $R$ は次の性質をもつ。(P) 任意の真の有限生成イデアル
$I \subset R$ は非零の零化イデアルをもつ。
\end{lemma}

\begin{proof}
仮定により、任意の $f \in \mathfrak m$ に対して
% SOURCE-ERRATUM-CANDIDATE P02-E0796: frozen source uses f^n without quantifying n; weak association requires an exponent depending on f.
ある指数について $f^n x = 0$ となるような非零元 $x \in R$ が存在する。
$I = (f_1, \ldots, f_r)$ と書く。このとき $x$ は
$R \to \bigoplus R_{f_i}$ の核に属する。従って Algebra, Lemma
\ref{algebra-lemma-when-injective-covering} により、すべての $i$ に対して
$f_i y = 0$ となる非零元 $y \in R$ が存在する。
$y \in \text{Ann}_R(I)$ なので主張が従う。
\end{proof}

\begin{lemma}
\label{more-algebra-lemma-P-universally-injective}
$R$ を Lemma \ref{more-algebra-lemma-auto-ass-implies-P} の性質 (P) をもつ環とする。
$u : N \to M$ を射影 $R$-加群の準同型とする。このとき $u$ が普遍的単射であることと
$u$ が単射であることは同値である。
\end{lemma}

\begin{proof}
$u$ が単射であると仮定する。目標は $u$ が普遍的単射であることを示すことである。
まず $N \oplus Q$ が自由となるような加群 $Q$ を選ぶ。
写像 $N \oplus Q \to M \oplus Q$ を考えることにより、$N$ が自由な場合に
補題を証明すれば十分である。この場合 $N$ は有限自由 $R$-加群の有向余極限である。
従って $N$ が有限自由 $R$-加群である場合に帰着する。
$N = R^{\oplus n}$ とおき、$n$ に関する帰納法で補題を証明する。
$n = 0$ の場合は自明である。

\medskip\noindent
$u : R^{\oplus n} \to M$ を単射な加群写像とし、$M$ は射影的であるとする。
$M \oplus Q$ が自由となるような $R$-加群 $Q$ を選ぶ。$u$ を合成
$R^{\oplus n} \to M \to M \oplus Q$ で置き換えると、$M$ は自由であると
仮定してよい。このとき $u(R^{\oplus n}) \subset R^{\oplus m}$ を満たす直和因子
$R^{\oplus m} \subset M$ を見つけられる。従って $M = R^{\oplus m}$ と仮定してよい。
この場合 $u$ は行列 $A = (a_{ij})$ により
$u(x_1, \ldots, x_n) = (\sum x_i a_{i1}, \ldots, \sum x_i a_{im})$
と与えられる。$u$ は単射なので、特に $x \not = 0$ ならば
$u(x, 0, \ldots, 0) = (xa_{11}, xa_{12}, \ldots, xa_{1m}) \not = 0$ である。
$R$ は性質 (P) をもつので
$a_{11}R + a_{12}R + \ldots + a_{1m}R = R$ である。従って
% SOURCE-ERRATUM-CANDIDATE P02-E0797: frozen source says “Hence see that”, omitting “we”.
$R(a_{11}, \ldots, a_{1m}) \subset R^{\oplus m}$ は $R^{\oplus m}$ の直和因子であり、
特に $R^{\oplus m}/R(a_{11}, \ldots, a_{1m})$ は射影 $R$-加群である。
可換図式
$$
\xymatrix{
0 \ar[r] &
R \ar[rr] \ar[d]^1 & & R^{\oplus n} \ar[r] \ar[d]^u &
R^{\oplus n - 1} \ar[r] \ar[d] & 0 \\
0 \ar[r] & R \ar[rr]^{(a_{11}, \ldots, a_{1m})} & &
R^{\oplus m} \ar[r] & R^{\oplus m}/R(a_{11}, \ldots, a_{1m}) \ar[r] & 0
}
$$
を得る。各行は分裂完全である。従って右の縦の矢印は単射であり、帰納法の仮定を
適用して右の縦の矢印が普遍的単射であると結論できる。ゆえに中央の縦の矢印は
普遍的単射である。
\end{proof}

\begin{lemma}
\label{more-algebra-lemma-P-fPD-zero}
$R$ を環とする。次は同値である。
\begin{enumerate}
\item $R$ は Lemma \ref{more-algebra-lemma-auto-ass-implies-P} の性質 (P) をもつ、
\item 射影 $R$-加群の任意の単射は普遍的単射である、
\item $u : N \to M$ が単射で、$N$, $M$ が有限射影 $R$-加群ならば、
$\Coker(u)$ は有限射影 $R$-加群である、
\item $N \subset M$ であり、$N$, $M$ が $R$-加群として有限射影ならば、
$N$ は $M$ の直和因子である、かつ
\item 任意の単射 $R \to R^{\oplus n}$ は分裂単射である。
\end{enumerate}
\end{lemma}

\begin{proof}
(1) $\Rightarrow$ (2) は Lemma \ref{more-algebra-lemma-P-universally-injective} である。
(3) と (4) が同値であることは明らかである。Algebra, Lemma
\ref{algebra-lemma-universally-exact-split} により
(2) $\Rightarrow$ (3), (4) である。(5) は (4) の特殊な場合である。
(5) を仮定する。$I = (a_1, \ldots, a_n)$ を $R$ の真の有限生成イデアルとする。
$I \not = R$ なので
$R \to R^{\oplus n}$, $x \mapsto (xa_1, \ldots, xa_n)$
は分裂単射でない。従って非零の核をもち、
$\text{Ann}_R(I) \not = 0$ と結論する。ゆえに (1) が成り立つ。
\end{proof}

\begin{example}
\label{more-algebra-example-auto-ass-weird-flat}
Lemma \ref{more-algebra-lemma-P-fPD-zero} の同値な条件が成り立っても、自由 $R$-加群の任意の
単射が常に分裂単射になるとは限らない。例えば
$R = k[x_1, x_2, x_3, \ldots]/(x_i^2)$ と仮定する。これは自己随伴環である。
自由 $R$-加群の写像
$$
u :
\bigoplus\nolimits_{i \geq 1} Re_i
\longrightarrow
\bigoplus\nolimits_{i \geq 1} Rf_i, \quad
e_i \longmapsto f_i - x_if_{i + 1}.
$$
を考える。任意の整数 $n$ に対して、$u$ の
$\bigoplus_{i = 1, \ldots, n} Re_i$ への制限は単射である。実際、
像 $u(e_1), \ldots, u(e_n)$ は $R$ 上線形独立である。
従って $u$ は単射であり、補題により普遍的単射である。
$u \otimes \text{id}_k$ は全単射なので、もし $u$ が分裂単射ならば
$u$ は全射である。しかし $u$ は全射でない。実際、$f_1$ の逆像は元
$$
\sum\nolimits_{i \geq 0} x_1 \ldots x_ie_{i + 1} =
e_1 + x_1e_2 + x_1x_2e_3 + \ldots
$$
でなければならないが、これは直和の元ではない。付言すると、
$\Coker(u)$ は平坦で（$u$ が普遍的単射だからである）、可算生成な
$R$-加群であるが、射影的でなく（$u$ が分裂しないからである）、
従って Mittag-Leffler でない（Algebra, Lemma
\ref{algebra-lemma-countgen-projective} を参照）。
\end{example}

\noindent
次の補題は、局所環が Noether 環である場合の Algebra, Proposition
\ref{algebra-proposition-what-exact} の特殊な場合である。

\begin{lemma}
\label{more-algebra-lemma-exact-length-1}
$(R, \mathfrak m)$ を局所環とする。
$\varphi : R^m \to R^n$ を有限自由加群の写像とする。次は同値である。
\begin{enumerate}
\item $\varphi$ は単射である、
\item $\varphi$ の階数は $m$ であり、$R$ における
$I(\varphi)$ の零化イデアルは零である。
\end{enumerate}
$R$ が Noether 環ならば、これらはさらに次と同値である。
\begin{enumerate}
\item[(3)] $\varphi$ の階数は $m$ であり、
$I(\varphi) = R$ であるか、またはそれは非零因子を含む。
\end{enumerate}
ここで $\varphi$ の階数および $I(\varphi)$ は Algebra, Definition
\ref{algebra-definition-rank} と同様に定義される。
\end{lemma}

\begin{proof}
$\varphi$ のいずれかの行列成分が $\mathfrak m$ に属さないならば、Algebra, Lemma
\ref{algebra-lemma-add-trivial-complex} を適用し、$\varphi$ を
$1 : R \to R$ と写像 $\varphi' : R^{m-1} \to R^{n-1}$ の和として書く。
$\varphi'$ に対する補題が $\varphi$ に対する補題を含意することは容易に分かる。
従って初めから $\varphi$ のすべての行列成分が $\mathfrak m$ に属すると
仮定してよい。

\medskip\noindent
$\varphi$ が単射であると仮定する。$m > 0$ と仮定してよい。
$\mathfrak q \in \text{WeakAss}(R)$ とし、従って
$R_\mathfrak q$ は自己随伴環である。このとき $\varphi$ は単射
$R_\mathfrak q^m \to R_\mathfrak q^n$ を誘導し、これは Lemmas
\ref{more-algebra-lemma-auto-ass-implies-P} および \ref{more-algebra-lemma-P-universally-injective} により
普遍的単射である。従って
$\varphi : \kappa(\mathfrak q)^m \to \kappa(\mathfrak q)^n$
は単射である。ゆえに $\varphi \bmod \mathfrak q$ の階数は $m$ であり、
$I(\varphi \otimes \kappa(\mathfrak q))$ は零イデアルでない。
$m$ は $\varphi$ がもち得る最大の階数なので、$\varphi$ の階数も $m$ であると
結論する（行列の階数は基底変換後には下がることしかない）。従って
$I(\varphi) \cdot \kappa(\mathfrak q) =
I(\varphi \otimes \kappa(\mathfrak q))$ は非零である。
ゆえに $I(\varphi)$ は $\mathfrak q$ に含まれない。従って $R$ の弱随伴素イデアルは
どれも $R$-加群 $\text{Ann}_R I(\varphi)$ の弱随伴素イデアルでない。
従って $\text{Ann}_R I(\varphi)$ は弱随伴素イデアルをもたない
（Algebra, Lemma \ref{algebra-lemma-weakly-ass} を参照）。
Algebra, Lemma \ref{algebra-lemma-weakly-ass-zero} により
$\text{Ann}_R I(\varphi)$ は零である。

\medskip\noindent
逆に (2) を仮定する。階数が $m$ であることは $n \geq m$ を含意する。
$I(\varphi)$ は有限生成なので $I(\varphi) = (f_1, \ldots, f_r)$ と書く。
Algebra, Lemma \ref{algebra-lemma-matrix-left-inverse} により、
% SOURCE-ERRATUM-CANDIDATE P02-E0798: frozen source defines psi_i but the following equation uses undefined psi; preserved.
$\psi \circ \varphi = f_i \text{id}_{R^m}$ を満たす写像
$\psi_i : R^n \to R^m$ を見つけられる。従って
$\varphi(x) = 0$ は $i = 1, \ldots, r$ に対して $f_i x = 0$ を含意する。
これは $x = 0$ を含意するので、$\varphi$ は単射である。

\medskip\noindent
Noether 局所環の場合の (1) と (3) の同値性については、Algebra, Proposition
\ref{algebra-proposition-what-exact} を参照する。
環 $R$ が Noether 環だが局所環でないならば、読者は局所環の場合からこれを
導ける。詳細は省略する。別の方法として、随伴素イデアルを用いて上の議論を
やり直すこともできる。その際には、それらが有限個しかないこと、素イデアル回避、
および Noether 環の非零因子はどの随伴素イデアルにも含まれない元として特徴づけ
られることを用いる。
\end{proof}

\begin{lemma}
\label{more-algebra-lemma-coker-injective-free}
$R$ を環とする。
% SOURCE-ERRATUM-CANDIDATE P02-E0799: frozen source combines “Suppose that” with “be”; Japanese supplies one grammatical construction.
$\varphi : R^n \to R^n$ が同じ階数の有限自由加群の単射であると仮定する。
このとき $\Hom_R(\Coker(\varphi), R) = 0$ である。
\end{lemma}

\begin{proof}
$\varphi^t : R^n \to R^n$ を $\varphi$ の転置とする。
補題の主張は $\varphi^t$ が単射であるということである。
Lemma \ref{more-algebra-lemma-exact-length-1} と同じ記号を用いると、$\varphi^t$ の階数は
$n$ であり、$I(\varphi) = I(\varphi^t)$ であることが分かる。
従って補題の (1) と (2) の同値性から結論を得る。
\end{proof}

\section{平坦化分層}
\label{more-algebra-section-flattening}

\noindent
$R \to S$ を環写像、$M$ を $S$-加群とする。任意の $R$-代数 $R'$ に対して、
基底変換 $S' = S \otimes_R R'$ および $M' = M \otimes_R R'$ を考えられる。
加群 $M'$ が $R'$ 上平坦であるとき、$R \to R'$ は $M$ を
{\it 平坦化する}という。$M$ を平坦化する環写像 $R \to R'$ 全体の構造を
理解したい。特に、{\it $M$ の普遍平坦化 $R \to R_{univ}$}、すなわち
$M$ を平坦化し、$M$ を平坦化する任意の環写像 $R \to R'$ が
$R \to R_{univ}$ を経由して分解するという性質をもつ環写像
$R \to R_{univ}$ が存在するかを知りたい。そのような普遍解は通常存在しない
ことが分かる。

\medskip\noindent
More on Flatness, Section \ref{flat-section-flattening} において、
スキーム論的な設定 $\mathcal{F}/X/S$ で {\it 普遍平坦化}と
{\it 平坦化分層}を論じる。普遍平坦化 $R \to R_{univ}$ が存在するならば、
スキームの射 $\Spec(R_{univ}) \to \Spec(R)$ は
$\Spec(S)$ 上の準連接加群 $\widetilde{M}$ の普遍平坦化である。

\medskip\noindent
本節および続くいくつかの節では、これに関係する基本的な代数の事実を証明する。
最も基本的な結果はおそらく次である。

\begin{lemma}
\label{more-algebra-lemma-intersection-flat}
$R$ を環、$M$ を $R$-加群、$I_1$, $I_2$ を $R$ のイデアルとする。
$M/I_1M$ が $R/I_1$ 上平坦で、$M/I_2M$ が $R/I_2$ 上平坦ならば、
$M/(I_1 \cap I_2)M$ は $R/(I_1 \cap I_2)$ 上平坦である。
\end{lemma}

\begin{proof}
$R$ を $R/(I_1 \cap I_2)$ で、$M$ を $M/(I_1 \cap I_2)M$ で置き換えることにより、
$I_1 \cap I_2 = 0$ と仮定してよい。$J \subset R$ をイデアルとする。
$M$ が $R$ 上平坦であることを証明するには、
$J \otimes_R M \to M$ が単射であることを示さなければならない
（Algebra, Lemma \ref{algebra-lemma-flat} を参照）。
$M/I_1M$ は $R/I_1$ 上平坦なので、写像
$$
J/(J \cap I_1) \otimes_R M =
(J + I_1)/I_1 \otimes_{R/I_1} M/I_1M
\longrightarrow M/I_1M
$$
は単射である。$0 \to (J \cap I_1) \to J \to J/(J \cap I_1) \to 0$
は完全なので、図式
$$
\xymatrix{
(J \cap I_1) \otimes_R M \ar[r] \ar[d] &
J \otimes_R M \ar[r] \ar[d] &
J/(J \cap I_1) \otimes_R M \ar[r] \ar[d] & 0 \\
M \ar@{=}[r] &
M \ar[r] &
M/I_1M
}
$$
を得る。従って $(J \cap I_1) \otimes_R M \to M$ が単射であることを
示せば十分である。$I_1 \cap I_2 = 0$ なので、イデアル $J \cap I_1$ は
イデアル $J' \subset R/I_2$ に同型に写り、
$(J \cap I_1) \otimes_R M = J' \otimes_{R/I_2} M/I_2M$ である。
$M/I_2M$ は $R/I_2$ 上平坦なので、写像
$J' \otimes_{R/I_2} M/I_2M \to M/I_2M$ は単射であり、これは明らかに
$(J \cap I_1) \otimes_R M \to M$ が単射であることを含意する。
\end{proof}

\section{Artin 環上の平坦化}
\label{more-algebra-section-flattening-artinian}

\noindent
基礎環が Artin 局所環であるとき、普遍平坦化が存在する。
これは任意の加群に対して存在する。従って後で見るように、基礎スキームが
% SOURCE-ERRATUM-CANDIDATE P02-E0800: frozen source says “flatting stratification”; Japanese supplies “flattening”.
Artin 局所環のスペクトルであるとき、平坦化分層が存在する。

\begin{lemma}
\label{more-algebra-lemma-flattening-artinian}
$R$ を Artin 環、$M$ を $R$-加群とする。このとき
$M/IM$ が $R/I$ 上平坦となるような最小のイデアル $I \subset R$ が存在する。
\end{lemma}

\begin{proof}
これは Lemma \ref{more-algebra-lemma-intersection-flat} と Artin 性から直接従う。
\end{proof}

\noindent
このイデアルは次の普遍性をもつ。

\begin{lemma}
\label{more-algebra-lemma-flattening-artinian-universal-property}
$R$ を Artin 環、$M$ を $R$-加群とする。$I \subset R$ を、
$M/IM$ が $R/I$ 上平坦となるような最小のイデアル $I \subset R$ とする。
このとき $I$ は次の普遍性をもつ。任意の環写像
$\varphi : R \to R'$ に対して
$$
R' \otimes_R M\text{ is flat over }R'
\Leftrightarrow
\text{we have }\varphi(I) = 0.
$$
である。
\end{lemma}

\begin{proof}
$I$ は Lemma \ref{more-algebra-lemma-flattening-artinian} により存在することに注意せよ。
% SOURCE-ERRATUM-CANDIDATE P02-E0801: frozen source attributes the forward implication to flat base change, although that lemma proves the reverse implication; the following argument proves the forward direction.
含意 $\Rightarrow$ は Algebra, Lemma \ref{algebra-lemma-flat-base-change} から従う。
$M \otimes_R R'$ が $R'$ 上平坦となるような
$\varphi : R \to R'$ をとる。$J = \Ker(\varphi)$ とおく。Algebra, Lemma
\ref{algebra-lemma-descent-flatness-injective-map-artinian-rings} と、
$R' \otimes_R M = R' \otimes_{R/J} M/JM$ が $R'$ 上平坦であることから、
$M/JM$ は $R/J$ 上平坦であると結論する。従って所望のように $I \subset J$ である。
\end{proof}

\section{基底の閉部分集合上の平坦化}
\label{more-algebra-section-flattening-local-base}

\noindent
$R \to S$ を環写像、$I \subset R$ をイデアル、$M$ を $S$-加群とする。
以下では次の条件を考える。
\begin{equation}
\label{more-algebra-equation-flat-at-primes-over}
\forall \mathfrak q \in V(IS) \subset \Spec(S) :
M_{\mathfrak q}\text{ is flat over }R.
\end{equation}
幾何学的には、これは $M$ が $\Spec(S)$ における $V(I)$ の逆像に沿って
$R$ 上平坦であることを意味する。$R$ と $S$ が Noether 環で、$M$ が有限
$S$-加群ならば、(\ref{more-algebra-equation-flat-at-primes-over}) は、すべての
$n \geq 1$ に対して $M/I^nM$ が $R/I^n$ 上平坦であるという条件と同値である
（Algebra, Lemma \ref{algebra-lemma-flat-module-powers} を参照）。

\begin{lemma}
\label{more-algebra-lemma-base-change-flat-at-primes-over}
$R \to S$ を環写像、$I \subset R$ をイデアル、$M$ を $S$-加群とする。
$R \to R'$ を環写像とし、$IR' \subset I' \subset R'$ をイデアルとする。
$(R \to S, I, M)$ に対して (\ref{more-algebra-equation-flat-at-primes-over}) が成り立つならば、
$(R' \to S \otimes_R R', I', M \otimes_R R')$ に対しても
(\ref{more-algebra-equation-flat-at-primes-over}) が成り立つ。
\end{lemma}

\begin{proof}
$(R \to S, I \subset R, M)$ に対して
(\ref{more-algebra-equation-flat-at-primes-over}) が成り立つと仮定する。
$I'(S \otimes_R R') \subset \mathfrak q'$ を $S \otimes_R R'$ の素イデアルとし、
$\mathfrak q \subset S$ を対応する $S$ の素イデアルとする。
このとき $IS \subset \mathfrak q$ である。
$(M \otimes_R R')_{\mathfrak q'}$ は基底変換
$M_{\mathfrak q} \otimes_R R'$ の局所化であることに注意せよ。
従って $(M \otimes_R R')_{\mathfrak q'}$ は平坦加群の局所化として
$R'$ 上平坦である（Algebra, Lemmas
\ref{algebra-lemma-flat-base-change} および
\ref{algebra-lemma-flat-localization} を参照）。
\end{proof}

\begin{lemma}
\label{more-algebra-lemma-flat-descent-flat-at-primes-over}
$R \to S$ を環写像、$I \subset R$ をイデアル、$M$ を $S$-加群とする。
$R \to R'$ を環写像とし、$IR' \subset I' \subset R'$ を次を満たす
イデアルとする。
\begin{enumerate}
\item $\Spec(R') \to \Spec(R)$ が誘導する写像
$V(I') \to V(I)$ は全射である、かつ
\item すべての素イデアル $\mathfrak p' \in V(I')$ に対して
$R'_{\mathfrak p'}$ は $R$ 上平坦である。
\end{enumerate}
$(R' \to S \otimes_R R', I', M \otimes_R R')$ に対して
(\ref{more-algebra-equation-flat-at-primes-over}) が成り立つならば、
$(R \to S, I, M)$ に対して (\ref{more-algebra-equation-flat-at-primes-over}) が成り立つ。
\end{lemma}

\begin{proof}
% SOURCE-ERRATUM-CANDIDATE P02-E0802: frozen proof assumes the condition for IR' although the lemma assumes it for I'; the proof uses primes over I'.
$(R' \to S \otimes_R R', IR', M \otimes_R R')$ に対して
(\ref{more-algebra-equation-flat-at-primes-over}) が成り立つと仮定する。
素イデアル $IS \subset \mathfrak q \subset S$ をとり、
$I \subset \mathfrak p \subset R$ を対応する $R$ の素イデアルとする。
仮定により、$\mathfrak p$ の上にある $R'$ の素イデアル
$\mathfrak p' \in V(I')$ で $R_{\mathfrak p} \to R'_{\mathfrak p'}$ が
平坦となるものが存在する。
$\overline{\mathfrak q}' \subset
\kappa(\mathfrak q) \otimes_{\kappa(\mathfrak p)} \kappa(\mathfrak p')$
を素イデアルとして選ぶ。これは $\mathfrak q$ と $\mathfrak p'$ の上にある
素イデアル $\mathfrak q' \subset S \otimes_R R'$ に対応する。
$(S \otimes_R R')_{\mathfrak q'}$ は
$S_{\mathfrak q} \otimes_{R_{\mathfrak p}} R'_{\mathfrak p'}$
の局所化であることに注意せよ。仮定により加群
$(M \otimes_R R')_{\mathfrak q'}$ は $R'_{\mathfrak p'}$ 上平坦である。
従って Algebra, Lemma \ref{algebra-lemma-base-change-flat-up-down} は
$M_{\mathfrak q}$ が $R_{\mathfrak p}$ 上平坦であることを含意し、
これが示すべきことであった。
\end{proof}








\begin{lemma}
\label{more-algebra-lemma-limit-preserving-flat-at-primes-over}
$R \to S$ を有限表示な環写像、$M$ を有限表示 $S$-加群とする。
$R' = \colim_{\lambda \in \Lambda} R_\lambda$
を $R$-代数の有向余極限とする。$\mu \geq \lambda$ のとき
$I_\lambda R_\mu \subset I_\mu$ を満たすイデアル
$I_\lambda \subset R_\lambda$ をとり、$I' = \colim_\lambda I_\lambda$ とおく。
$(R' \to S \otimes_R R', I', M \otimes_R R')$ に対して
(\ref{more-algebra-equation-flat-at-primes-over}) が成り立つならば、ある
$\lambda \in \Lambda$ が存在して
$(R_\lambda \to S \otimes_R R_\lambda, I_\lambda,
M \otimes_R R_\lambda)$ に対して
(\ref{more-algebra-equation-flat-at-primes-over}) が成り立つ。
\end{lemma}

\begin{proof}
$S_\lambda = S \otimes_R R_\lambda$,
$S' = S \otimes_R R'$, $M_\lambda = M \otimes_R R_\lambda$ および
$M' = M \otimes_R R'$ と書く。
基底変換 $S'$ は $R'$ 上有限表示で、$M'$ は $S'$ 上有限表示であり、
添字 $\lambda$ をもつ版についても同様である
（Algebra, Lemma \ref{algebra-lemma-base-change-finiteness} を参照）。
Algebra, Theorem \ref{algebra-theorem-openness-flatness} により、集合
$$
U' = \{\mathfrak q' \in \Spec(S') \mid
M'_{\mathfrak q'}\text{ is flat over }R'\}
$$
は $\Spec(S')$ の開集合である。$V(I'S')$ は仮定により $U'$ に含まれる
準コンパクト空間であることに注意せよ。従って有限個の
$g'_j \in S'$, $j = 1, \ldots, m$ が存在し、$D(g'_j) \subset U'$ かつ
$V(I'S') \subset \bigcup D(g'_j)$ となる。特に
$(M')_{g'_j}$ は $R'$ 上平坦な加群である。

\medskip\noindent
以下、$\lambda \in \Lambda$ を順次十分大きくとる。まず、
$g'_j$ に写る $g_{j, \lambda} \in S_{\lambda}$ を見つけられるほど大きくとる。
包含 $V(I'S') \subset \bigcup D(g'_j)$ は
$I'S' + (g'_1, \ldots, g'_m) = S'$ を意味し、これは、ある
$z_s \in I'$, $h_s, f_j \in S'$ による
$1 = \sum z_sh_s + \sum f_jg'_j$ と表せる。
$\lambda$ を大きくした後、このような等式が $S_\lambda$ で成り立つと仮定してよい。
従って $V(I_\lambda S_\lambda) \subset \bigcup D(g_{j, \lambda})$
と仮定してよい。Algebra, Lemma
\ref{algebra-lemma-flat-finite-presentation-limit-flat} により、十分大きいある
$\lambda$ に対して加群 $(M_\lambda)_{g_{j, \lambda}}$ は
$R_\lambda$ 上平坦である。特に加群 $M_\lambda$ は、イデアル
$I_\lambda$ の上にあるすべての素イデアルにおいて $R_\lambda$ 上平坦である。
\end{proof}

% SOURCE-ERRATUM-CANDIDATE P02-E0803: frozen section title says “a closed subsets”, with singular article and plural noun; Japanese supplies grammatical number.
\section{始域と基底の閉部分集合上の平坦化}
\label{more-algebra-section-flattening-local-source-base}

\noindent
本節では Section \ref{more-algebra-section-flattening-local-base} の議論を少し一般化する。
読者にはまず同節を読み、理解することを強く勧める。

\begin{situation}
\label{more-algebra-situation-flattening-general}
$R \to S$ を環写像、$J \subset S$ をイデアル、$M$ を $S$-加群とする。
\end{situation}

\noindent
この状況で、$R$-代数 $R'$ とイデアル $I' \subset R'$ が与えられたとき、
$S' = S \otimes_R R'$ および $M' = M \otimes_R R'$ とおく。
次の条件を考える。
\begin{equation}
\label{more-algebra-equation-flat-at-primes}
\forall \mathfrak q' \in V(I'S' + JS') \subset \Spec(S') :
M'_{\mathfrak q'}\text{ is flat over }R'.
\end{equation}
幾何学的には、これは $M'$ が、$V(I')$ の逆像と $V(J)$ の逆像との交わりに沿って
$R'$ 上平坦であることを意味する。$(R \to S, J, M)$ は固定されているので、
条件 (\ref{more-algebra-equation-flat-at-primes}) は、$R'$ を $R$-代数とみたときの対
$(R', I')$ のみに依存する。

\begin{lemma}
\label{more-algebra-lemma-base-change-flat-at-primes}
Situation \ref{more-algebra-situation-flattening-general} において、
$R' \to R''$ を $R$-代数写像とする。$I' \subset R'$ および
$I'R'' \subset I'' \subset R''$ をイデアルとする。
$(R', I')$ に対して (\ref{more-algebra-equation-flat-at-primes}) が成り立つならば、
$(R'', I'')$ に対しても (\ref{more-algebra-equation-flat-at-primes}) が成り立つ。
\end{lemma}

\begin{proof}
$(R', I')$ に対して (\ref{more-algebra-equation-flat-at-primes}) が成り立つと仮定する。
$I''S'' + JS'' \subset \mathfrak q''$ を $S''$ の素イデアルとし、
$\mathfrak q' \subset S'$ を対応する $S'$ の素イデアルとする。
$\mathfrak q''$ に対して対応する条件が成り立つので、
$I'S' \subset \mathfrak q'$ と $JS' \subset \mathfrak q'$ の両方が成り立つ。
$(M'')_{\mathfrak q''}$ は基底変換
% SOURCE-ERRATUM-CANDIDATE P02-E0804: frozen source tensors M' over R rather than over R'; preserved.
$M'_{\mathfrak q'} \otimes_R R''$ の局所化であることに注意せよ。
従って $(M'')_{\mathfrak q''}$ は平坦加群の局所化として $R''$ 上平坦である
（Algebra, Lemmas \ref{algebra-lemma-flat-base-change} および
\ref{algebra-lemma-flat-localization} を参照）。
\end{proof}

\begin{lemma}
\label{more-algebra-lemma-flat-descent-flat-at-primes}
Situation \ref{more-algebra-situation-flattening-general} において、
$R' \to R''$ を $R$-代数写像とする。$I' \subset R'$ および
$I'R'' \subset I'' \subset R''$ をイデアルとする。次を仮定する。
\begin{enumerate}
\item $\Spec(R'') \to \Spec(R')$ が誘導する写像
$V(I'') \to V(I')$ は全射である、かつ
\item すべての素イデアル $\mathfrak p'' \in V(I'')$ に対して
$R''_{\mathfrak p''}$ は $R'$ 上平坦である。
\end{enumerate}
$(R'', I'')$ に対して (\ref{more-algebra-equation-flat-at-primes}) が成り立つならば、
$(R', I')$ に対して (\ref{more-algebra-equation-flat-at-primes}) が成り立つ。
\end{lemma}

\begin{proof}
$(R'', I'')$ に対して (\ref{more-algebra-equation-flat-at-primes}) が成り立つと仮定する。
素イデアル $I'S' + JS' \subset \mathfrak q' \subset S'$ をとり、
$I' \subset \mathfrak p' \subset R'$ を対応する $R'$ の素イデアルとする。
仮定により、$\mathfrak p'$ の上にある $R''$ の素イデアル
$\mathfrak p'' \in V(I'')$ で、
$R'_{\mathfrak p'} \to R''_{\mathfrak p''}$ が平坦となるものが存在する。
素イデアル
$\overline{\mathfrak q}'' \subset
\kappa(\mathfrak q') \otimes_{\kappa(\mathfrak p')} \kappa(\mathfrak p'')$
を選ぶ。これは $\mathfrak q'$ と $\mathfrak p''$ の上にある素イデアル
$\mathfrak q'' \subset S'' = S' \otimes_{R'} R''$ に対応する。
特に $I''S'' \subset \mathfrak q''$ および $JS'' \subset \mathfrak q''$ である。
$(S' \otimes_{R'} R'')_{\mathfrak q''}$ は
$S'_{\mathfrak q'} \otimes_{R'_{\mathfrak p'}} R''_{\mathfrak p''}$
の局所化であることに注意せよ。仮定により加群
$(M' \otimes_{R'} R'')_{\mathfrak q''}$ は $R''_{\mathfrak p''}$ 上平坦である。
従って Algebra, Lemma \ref{algebra-lemma-base-change-flat-up-down} は
$M'_{\mathfrak q'}$ が $R'_{\mathfrak p'}$ 上平坦であることを含意し、
これが示すべきことであった。
\end{proof}

\begin{lemma}
\label{more-algebra-lemma-limit-preserving-flat-at-primes}
Situation \ref{more-algebra-situation-flattening-general} において、$R \to S$ は本質的有限表示で、
$M$ は有限表示 $S$-加群であると仮定する。
$R' = \colim_{\lambda \in \Lambda} R_\lambda$
を $R$-代数の有向余極限とする。すべての $\mu \geq \lambda$ に対して
$I_\lambda R_\mu \subset I_\mu$ を満たすイデアル
$I_\lambda \subset R_\lambda$ をとり、$I' = \colim_\lambda I_\lambda$ とおく。
$(R', I')$ に対して (\ref{more-algebra-equation-flat-at-primes}) が成り立つならば、
ある $\lambda \in \Lambda$ が存在し、$(R_\lambda, I_\lambda)$ に対して
(\ref{more-algebra-equation-flat-at-primes}) が成り立つ。
\end{lemma}

\begin{proof}
まず $R \to S$ が有限表示である場合に補題を証明し、その後で一般の場合に
必要な変更を説明する。$S_\lambda = S \otimes_R R_\lambda$,
$S' = S \otimes_R R'$, $M_\lambda = M \otimes_R R_\lambda$ および
$M' = M \otimes_R R'$ と書く。
基底変換 $S'$ は $R'$ 上有限表示で、$M'$ は $S'$ 上有限表示であり、
添字 $\lambda$ をもつ版についても同様である
（Algebra, Lemma \ref{algebra-lemma-base-change-finiteness} を参照）。
Algebra, Theorem \ref{algebra-theorem-openness-flatness} により、集合
$$
U' = \{\mathfrak q' \in \Spec(S') \mid
M'_{\mathfrak q'}\text{ is flat over }R'\}
$$
は $\Spec(S')$ の開集合である。$V(I'S' + JS')$ は仮定により $U'$ に含まれる
準コンパクト空間であることに注意せよ。従って有限個の
$g'_j \in S'$, $j = 1, \ldots, m$ が存在し、$D(g'_j) \subset U'$ かつ
$V(I'S' + JS') \subset \bigcup D(g'_j)$ となる。特に
$(M')_{g'_j}$ は $R'$ 上平坦な加群である。

\medskip\noindent
以下、$\lambda \in \Lambda$ を順次十分大きくとる。まず、
$g'_j$ に写る $g_{j, \lambda} \in S_{\lambda}$ を見つけられるほど大きくとる。
包含 $V(I'S' + JS') \subset \bigcup D(g'_j)$ は
$I'S' + JS' + (g'_1, \ldots, g'_m) = S'$ を意味し、これは、ある
$z_s \in I'$, $y_t \in J$, $k_t, h_s, f_j \in S'$ に対する
$$
1 = \sum y_tk_t + \sum z_sh_s + \sum f_jg'_j
$$
と表せる。$\lambda$ を大きくした後、このような等式が $S_\lambda$ で
成り立つと仮定してよい。従って
$V(I_\lambda S_\lambda + J S_\lambda) \subset \bigcup D(g_{j, \lambda})$
と仮定してよい。Algebra, Lemma
\ref{algebra-lemma-flat-finite-presentation-limit-flat} により、十分大きいある
$\lambda$ に対して加群 $(M_\lambda)_{g_{j, \lambda}}$ は
$R_\lambda$ 上平坦である。特に加群 $M_\lambda$ は、
$V(I_\lambda S_\lambda + J S_\lambda)$ の点に対応するすべての素イデアルにおいて
$R_\lambda$ 上平坦である。

\medskip\noindent
$S$ が本質的有限表示である場合、$R \to C$ が有限表示で
$\Sigma \subset C$ が乗法的部分集合となるように
$S = \Sigma^{-1}C$ と書ける。また、有限表示 $C$-加群 $N$ を用いて
$M = \Sigma^{-1}N$ と書ける
（Algebra, Lemma \ref{algebra-lemma-construct-fp-module} を参照）。
ここで $C_\lambda$, $C'$, $N_\lambda$, $N'$ を導入する。
すると上の議論で、$N'$ が $R'$ 上平坦となる開集合
$U' \subset \Spec(C')$ を得る。(\ref{more-algebra-equation-flat-at-primes}) が成り立つという仮定は
$V(I'S' + JS')$ が $U'$ に写ることを意味する。実際、
素イデアル $\mathfrak q' \subset S'$ に対応する
素イデアル $\mathfrak r' \subset C'$ に対して
$M'_{\mathfrak q'} = N'_{\mathfrak r'}$ である。
従って $\bigcup D(g'_j)$ が $V(I'S' + JS')$ の像を含むような
$g'_j \in C'$ を見つけられる。証明の残りは前とまったく同じである。
\end{proof}

\begin{lemma}
\label{more-algebra-lemma-flat-module-powers-variant}
Situation \ref{more-algebra-situation-flattening-general} におく。
$I \subset R$ をイデアルとする。次を仮定する。
\begin{enumerate}
\item $R$ は Noether 環である、
\item $S$ は Noether 環である、
\item $M$ は有限 $S$-加群である、かつ
\item 各 $n \geq 1$ と任意の素イデアル
$\mathfrak q \in V(J + IS)$ に対して、加群
$(M/I^n M)_{\mathfrak q}$ は $R/I^n$ 上平坦である。
\end{enumerate}
このとき $(R, I)$ に対して (\ref{more-algebra-equation-flat-at-primes}) が成り立つ。
すなわち、任意の素イデアル $\mathfrak q \in V(J + IS)$ に対して
局所化 $M_{\mathfrak q}$ は $R$ 上平坦である。
\end{lemma}

\begin{proof}
$\mathfrak q \in V(J + IS)$ とする。
Algebra, Lemma \ref{algebra-lemma-flat-module-powers} を
$R \to S_{\mathfrak q}$ および $M_{\mathfrak q}$ に適用すると、
$M_{\mathfrak q}$ は $R$ 上平坦であることが従う。
\end{proof}

\section{完備 Noether 局所環上の平坦化}
\label{more-algebra-section-flattening-Noetherian-complete-local}

\noindent
基底が完備 Noether 局所環 $R$ で、与えられた加群が有限型 $R$-代数 $S$ 上
有限であるとき、次の三つの補題は「局所」平坦化分層の存在に対する
完全に代数的な証明を与える。

\begin{lemma}
\label{more-algebra-lemma-flattening-complete-local-noetherian}
$R \to S$ を環写像、$M$ を $S$-加群とする。次を仮定する。
\begin{enumerate}
\item $(R, \mathfrak m)$ は完備 Noether 局所環である、
\item $S$ は Noether 環である、かつ
\item $M$ は $S$ 上有限である。
\end{enumerate}
このとき、次を満たすイデアル $I \subset \mathfrak m$ が存在する。
\begin{enumerate}
\item $\mathfrak m$ の上にある $S/IS$ のすべての素イデアル
$\mathfrak q$ に対して、$(M/IM)_{\mathfrak q}$ は $R/I$ 上平坦である、かつ
\item $J \subset R$ が、$\mathfrak m$ の上にあるすべての素イデアル
$\mathfrak q$ に対して $(M/JM)_{\mathfrak q}$ が $R/J$ 上平坦となる
イデアルならば、$I \subset J$ である。
\end{enumerate}
言い換えれば、$I$ は
$(\overline{R} \to \overline{S}, \overline{\mathfrak m}, \overline{M})$
に対して (\ref{more-algebra-equation-flat-at-primes-over}) が成り立つような $R$ の最小の
イデアルである。ここで $\overline{R} = R/I$, $\overline{S} = S/IS$,
$\overline{\mathfrak m} = \mathfrak m/I$ および $\overline{M} = M/IM$ である。
\end{lemma}

\begin{proof}
$J \subset R$ をイデアルとする。Algebra, Lemma
\ref{algebra-lemma-flat-module-powers} を環 $R/J$ 上の加群 $M/JM$ に適用する。
このとき
% SOURCE-ERRATUM-CANDIDATE P02-E0805: frozen source says all primes of S/JS, but the cited powers lemma yields primes over the maximal ideal; preserved.
$(M/JM)_{\mathfrak q}$ が $S/JS$ のすべての素イデアル $\mathfrak q$ に対して
$R/J$ 上平坦であることと、すべての $n \geq 1$ に対して
$M/(J + \mathfrak m^n)M$ が $R/(J + \mathfrak m^n)$ 上平坦であることは同値である。
以下ではこの注意を用いる。

\medskip\noindent
すべての $n \geq 1$ に対して、局所環 $R/\mathfrak m^n$ は Artin 環である。
従って Lemma \ref{more-algebra-lemma-flattening-artinian} により、
$M/I_nM$ が $R/I_n$ 上平坦となるような最小のイデアル
$I_n \supset \mathfrak m^n$ が存在する。
% SOURCE-ERRATUM-CANDIDATE P02-E0806: frozen source says “is contains”; Japanese supplies the intended verb once.
$I_{n + 1} + \mathfrak m^n$ が $I_n$ を含むことは明らかであり、
Lemma \ref{more-algebra-lemma-intersection-flat} を適用すると
$I_n = I_{n + 1} + \mathfrak m^n$ が分かる。
$R = \lim_n\ R/\mathfrak m^n$ なので
$I = \lim_n\ I_n/\mathfrak m^n$ は $R$ のイデアルであり、すべての
$n \geq 1$ に対して $I_n = I + \mathfrak m^n$ である。
証明冒頭の注意により、$I$ は (1) と (2) を満たす。いくつかの詳細は省略する。
\end{proof}

\begin{lemma}
\label{more-algebra-lemma-flattening-complete-local-noetherian-property-by-finite-type}
Lemma \ref{more-algebra-lemma-flattening-complete-local-noetherian} と同じ記号
$R \to S$, $M$, $I$ および仮定を用いる。
$R'$ が Noether 環となるような局所環の局所準同型
$\varphi : (R, \mathfrak m) \to (R', \mathfrak m')$ を考える。
このとき次は同値である。
\begin{enumerate}
\item $(R' \to S \otimes_R R', \mathfrak m', M \otimes_R R')$ に対して
条件 (\ref{more-algebra-equation-flat-at-primes-over}) が成り立つ、かつ
\item $\varphi(I) = 0$.
\end{enumerate}
\end{lemma}

\begin{proof}
含意 (2) $\Rightarrow$ (1) は Lemma
\ref{more-algebra-lemma-base-change-flat-at-primes-over} から従う。
$\varphi : R \to R'$ を補題の通りで (1) を満たすものとする。
$\varphi(I) = 0$ を示さなければならない。これは条件
$\varphi(I)R' = 0$ と同値である。Noether 局所環 $R'$ における
Artin--Rees により（Algebra, Lemma
\ref{algebra-lemma-intersect-powers-ideal-module-zero} を参照）、これは
すべての $n > 0$ に対して
$\varphi(I)R' + (\mathfrak m')^n = (\mathfrak m')^n$ であるという条件と同値である。
従ってこれは、各 $n$ に対して合成
$\varphi_n : R \to R' \to R'/(\mathfrak m')^n$ が $I$ を零化するという条件と
同値である。$\varphi$ に対する仮定 (1) は、Lemma
\ref{more-algebra-lemma-base-change-flat-at-primes-over} により $\varphi_n$ に対する
仮定 (1) を含意する。これにより $R'$ が Artin 局所環である場合に帰着する。

\medskip\noindent
% SOURCE-ERRATUM-CANDIDATE P02-E0807: frozen source says “Assume R' Artinian”, omitting “is”.
$R'$ が Artin 環であると仮定する。$J = \Ker(\varphi)$ とおく。
$I \subset J$ を示さなければならない。Lemma
\ref{more-algebra-lemma-flattening-complete-local-noetherian} における $I$ の構成により、
$\mathfrak m$ の上にある $S/JS$ の各素イデアル $\mathfrak q$ に対して
$(M/JM)_{\mathfrak q}$ が $R/J$ 上平坦であることを示せば十分である。
$R'$ は Artin 環なので、条件 (1) は $M \otimes_R R'$ が $R'$ 上平坦であることを
意味する。$R'$ は Artin 環であり、$R/J \to R'$ は単射な局所環写像なので、
$R/J$ も Artin 環である。従って
$M \otimes_R R' = M/JM \otimes_{R/J} R'$ が $R'$ 上平坦であることは、
Algebra, Lemma
\ref{algebra-lemma-descent-flatness-injective-map-artinian-rings} により
$M/JM$ が $R/J$ 上平坦であることを含意する。これで証明が完了する。
\end{proof}

\begin{lemma}
\label{more-algebra-lemma-flattening-complete-local-universal-property}
Lemma \ref{more-algebra-lemma-flattening-complete-local-noetherian} と同じ記号
$R \to S$, $M$, $I$ および仮定を用いる。さらに $R \to S$ は有限型であると仮定する。
このとき任意の局所環の局所準同型
$\varphi : (R, \mathfrak m) \to (R', \mathfrak m')$ に対して次は同値である。
\begin{enumerate}
\item $(R' \to S \otimes_R R', \mathfrak m', M \otimes_R R')$ に対して
条件 (\ref{more-algebra-equation-flat-at-primes-over}) が成り立つ、かつ
\item $\varphi(I) = 0$.
\end{enumerate}
\end{lemma}

\begin{proof}
含意 (2) $\Rightarrow$ (1) は Lemma
\ref{more-algebra-lemma-base-change-flat-at-primes-over} から従う。
$\varphi : R \to R'$ を補題の通りで (1) を満たすものとする。
$R$ は Noether 環なので、$R \to S$ は有限表示で、$M$ は有限表示 $S$-加群である。
この環を、極大イデアル
$\mathfrak m_\lambda = R_\lambda \cap \mathfrak m'$ をもち、各 $R_\lambda$ が
$R$ 上本質的有限型であるような局所 $R$-部分代数
$R_\lambda \subset R'$ の有向余極限
$R' = \colim_\lambda R_\lambda$ と書く。Lemma
\ref{more-algebra-lemma-limit-preserving-flat-at-primes-over} により、ある $\lambda$ に対して
$(R_\lambda \to S \otimes_R R_\lambda, \mathfrak m_\lambda,
M \otimes_R R_\lambda)$ に条件 (\ref{more-algebra-equation-flat-at-primes-over}) が成り立つ。
従って Lemma
\ref{more-algebra-lemma-flattening-complete-local-noetherian-property-by-finite-type} を
環写像 $R \to R_\lambda$ に適用でき、$I$ は $R_\lambda$ で零に写ることが分かる。
従って、まして $R'$ で零に写る。
\end{proof}

\section{整写像に沿う平坦性の降下}
\label{more-algebra-section-descent-flatness-integral}

\noindent
まずいくつかの簡単な補題を示す。

\begin{lemma}
\label{more-algebra-lemma-have-one-root}
$R$ を環とし、$P(T)$ を係数が $R$ に属するモニック多項式とする。
$P(\alpha) = 0$ を満たす $\alpha \in R$ をとる。このとき、あるモニック多項式
$Q(T) \in R[T]$ に対して $P(T) = (T - \alpha)Q(T)$ である。
\end{lemma}

\begin{proof}
$P$ の次数に関する帰納法による。$\deg(P) = 1$ ならば
$P(T) = T - \alpha$ なので結果は正しい。$\deg(P) > 1$ ならば、
長除法により、次数が $< \deg(P)$ である多項式 $Q \in R[T]$ と
$r \in R$ を用いて $P(T) = (T - \alpha)Q(T) + r$ と書ける。
仮定により $0 = P(\alpha) = (\alpha - \alpha)Q(\alpha) + r = r$ なので、
所望のように $r = 0$ と結論する。
\end{proof}

\begin{lemma}
\label{more-algebra-lemma-adjoin-one-root}
$R$ を環とし、$P(T)$ を係数が $R$ に属するモニック多項式とする。
ある $\alpha \in R'$ とモニック多項式 $Q(T) \in R'[T]$ に対して
$P(T) = (T - \alpha)Q(T)$ となるような有限自由な環写像
$R \to R'$ が存在する。
\end{lemma}

\begin{proof}
$P(T) = T^d + a_1T^{d - 1} + \ldots + a_0$ と書き、
$R' = R[x]/(x^d + a_1x^{d - 1} + \ldots + a_0)$ とおく。
$\alpha$ を $x$ の剰余類とする。このとき $P(\alpha) = 0$ は明らかである。
従って Lemma \ref{more-algebra-lemma-have-one-root} により主張が従う。
\end{proof}

\begin{lemma}
\label{more-algebra-lemma-finite-split}
$R \to S$ を有限な環写像とする。$S \otimes_R R'$ が
$$
R'[T_1, \ldots, T_n]/(P_1(T_1), \ldots, P_n(T_n))
$$
の形の環の商となるような有限自由環拡大 $R \subset R'$ が存在する。
ここで、ある $\alpha_{ij} \in R'$ に対して
$P_i(T) = \prod_{j = 1, \ldots, d_i} (T - \alpha_{ij})$ である。
\end{lemma}

\begin{proof}
$x_1, \ldots, x_n \in S$ を $R$ 上の $S$ の生成元とする。
各 $i$ に対し、$S$ において $P_i(x_i) = 0$ となるモニック多項式
$P_i(T) \in R[T]$ を選べる
（Algebra, Lemma \ref{algebra-lemma-finite-is-integral} を参照）。
$\deg(P_i) = d_i$ とおく。Lemma \ref{more-algebra-lemma-adjoin-one-root} を
$\sum d_i$ 回適用することにより、各 $P_i$ が
$$
P_i(T) = \prod\nolimits_{j = 1, \ldots, d_i} (T - \alpha_{ij})
$$
と完全分解するような有限自由環拡大 $R \subset R'$ が存在する。
% SOURCE-ERRATUM-CANDIDATE P02-E0808: frozen source changes the second index from j to undefined k in alpha_{ik}; preserved.
ここで $\alpha_{ik} \in R'$ である。
$T_i$ を $x_i \otimes 1$ に送る $R'$-代数写像
$R'[T_1, \ldots, T_n] \to S \otimes_R R'$ を考える。
これは $P_i(T_i)$ を零に送るので、所望の全射を誘導する。
\end{proof}

\begin{lemma}
\label{more-algebra-lemma-split-image}
$R$ を環とし、$S = R[T_1, \ldots, T_n]/J$ とする。
ある $\alpha_{ij} \in R$ に対して
$P_i(T) = \prod_{j = 1, \ldots, d_i} (T - \alpha_{ij})$
となる形の元 $P_i(T_i)$ を $J$ が含むと仮定する。
$1 \leq k_i \leq d_i$ を満たす
$\underline{k} = (k_1, \ldots, k_n)$ に対し、環写像
$$
\Phi_{\underline{k}} : R[T_1, \ldots, T_n] \to R,
\quad
T_i \longmapsto \alpha_{ik_i}
$$
を考える。$J_{\underline{k}} = \Phi_{\underline{k}}(J)$ とおく。
このとき $\Spec(S) \to \Spec(R)$ の像は
$V(\bigcap J_{\underline{k}})$ に等しい。
\end{lemma}

\begin{proof}
この補題はほとんど自明である。ヒント：
$V(\bigcap J_{\underline{k}}) = \bigcup V(J_{\underline{k}})$.
\end{proof}

\noindent
次の結果は Ferrand による。文献 \cite{Ferrand} を参照せよ。

\begin{lemma}
\label{more-algebra-lemma-descent-flatness-injective-finite-Noetherian-rings}
$R \to S$ を Noether 環の有限単射準同型とし、$M$ を $R$-加群とする。
$M \otimes_R S$ が平坦 $S$-加群ならば、$M$ は平坦 $R$-加群である。
\end{lemma}

\begin{proof}
$M \otimes_R S$ が $S$ 上平坦となる $R$-加群 $M$ をとる。
Algebra, Lemma \ref{algebra-lemma-flatness-descends} により、$M$ の平坦性を
証明する際には $R$ を任意の忠実平坦環拡大で置き換えてよい。
Lemma \ref{more-algebra-lemma-finite-split} により、あるイデアル
$J \subset R'[T_1, \ldots, T_n]$ に対して
$S' = S \otimes_R R' = R'[T_1, \ldots, T_n]/J$ となるような有限局所自由環拡大
$R \subset R'$ を見つけられる。ここでこのイデアルは、ある
$\alpha_{ij} \in R'$ に対して
$P_i(T) = \prod_{j = 1, \ldots, d_i} (T - \alpha_{ij})$
となる形の元 $P_i(T_i)$ を含む。
% SOURCE-ERRATUM-CANDIDATE P02-E0809: frozen source has a doubled space in “the  elements” here and in the parallel projectivity proof.
$R'$ は Noether 環で、$R' \subset S'$ は有限環拡大であることに注意せよ。
従って $R$ を $R'$ で置き換え、$S$ が Lemma
\ref{more-algebra-lemma-split-image} のような表示をもつと仮定してよい。
$\Spec(S) \to \Spec(R)$ は全射であることに注意せよ
（Algebra, Lemma \ref{algebra-lemma-integral-overring-surjective} を参照）。
従って Lemma \ref{more-algebra-lemma-split-image} を用いると、
$I = \bigcap J_{\underline{k}}$ は $V(I) = \Spec(R)$ を満たすイデアルであると
結論する。これは $I \subset \sqrt{(0)}$ を意味し、$R$ は Noether 環なので
$I$ は冪零である。写像 $\Phi_{\underline{k}}$ は可換図式
$$
\xymatrix{
S \ar[rr] & & R/J_{\underline{k}} \\
& R \ar[lu] \ar[ru]
}
$$
を誘導する。これから $M/J_{\underline{k}}M$ は
$R/J_{\underline{k}}$ 上平坦であると結論する。
Lemma \ref{more-algebra-lemma-intersection-flat} により $M/IM$ は $R/I$ 上平坦である。
最後に Algebra, Lemma \ref{algebra-lemma-lift-flatness} を適用して、
$M$ は $R$ 上平坦であると結論する。
\end{proof}

\begin{lemma}
\label{more-algebra-lemma-descent-flatness-injective-integral}
$R \to S$ を単射な整環写像とする。$M$ を
$R[x_1, \ldots, x_n]$ 上の有限表示加群とする。
$M \otimes_R S$ が $S$ 上平坦ならば、$M$ は $R$ 上平坦である。
\end{lemma}

\begin{proof}
表示
$$
R[x_1, \ldots, x_n]^{\oplus t} \to R[x_1, \ldots, x_n]^{\oplus r} \to M \to 0.
$$
を選ぶ。最初の写像が、$f_{ij} \in R[x_1, \ldots, x_n]$ を成分とする
$r \times t$-行列 $T = (f_{ij})$ で与えられるとする。
$f_{ij, I} \in R$ を用いて
$f_{ij} = \sum f_{ij, I} x^I$ と書く（多重添字記法）。
図式
$$
\xymatrix{
R \ar[r] & S \\
R_\lambda \ar[u] \ar[r] & S_\lambda \ar[u]
}
$$
を考える。ここで $R_\lambda$ はすべての $f_{ij, I}$ を含む
$R$ の有限生成 $\mathbf{Z}$-部分代数で、$S_\lambda$ は $S$ の有限
$R_\lambda$-部分代数である。同じ表示と同じ行列 $T$ を
$R_\lambda[x_1, \ldots, x_n]$ 上の行列とみなして定まる有限
$R_\lambda[x_1, \ldots, x_n]$-加群を $M_\lambda$ とする。
$S$ は $S_\lambda$ の有向余極限であることに注意せよ（詳細は省略する）。
Algebra, Lemma \ref{algebra-lemma-flat-finite-presentation-limit-flat} により、
ある $\lambda$ に対して加群
$M_\lambda \otimes_{R_\lambda} S_\lambda$ は $S_\lambda$ 上平坦である。
Lemma \ref{more-algebra-lemma-descent-flatness-injective-finite-Noetherian-rings} により、
$M_\lambda$ は $R_\lambda$ 上平坦であると結論する。
$M = M_\lambda \otimes_{R_\lambda} R$ なので、Algebra, Lemma
\ref{algebra-lemma-flat-base-change} により主張が従う。
\end{proof}

\begin{lemma}
\label{more-algebra-lemma-descent-projective-injective-finite-Noetherian-rings}
$R \to S$ を Noether 環の有限単射準同型とし、$P$ を $R$-加群とする。
$P \otimes_R S$ が射影 $S$-加群ならば、$P$ は射影 $R$-加群である。
\end{lemma}

\begin{proof}
$P \otimes_R S$ が $S$ 上射影的となる $R$-加群 $P$ をとる。
Algebra, Theorem \ref{algebra-theorem-ffdescent-projectivity} により、$P$ の
射影性を証明する際には $R$ を任意の忠実平坦環拡大で置き換えてよい。
Lemma \ref{more-algebra-lemma-finite-split} により、あるイデアル
$J \subset R'[T_1, \ldots, T_n]$ に対して
$S' = S \otimes_R R' = R'[T_1, \ldots, T_n]/J$ となるような有限局所自由環拡大
$R \subset R'$ を見つけられる。ここでこのイデアルは、ある
$\alpha_{ij} \in R'$ に対して
$P_i(T) = \prod_{j = 1, \ldots, d_i} (T - \alpha_{ij})$
となる形の元 $P_i(T_i)$ を含む。
% SOURCE-ERRATUM-CANDIDATE P02-E0809: second frozen occurrence of doubled spacing in “the  elements”.
$R'$ は Noether 環で、$R' \subset S'$ は有限環拡大であることに注意せよ。
従って $R$ を $R'$ で置き換え、$S$ が Lemma
\ref{more-algebra-lemma-split-image} のような表示をもつと仮定してよい。
$\Spec(S) \to \Spec(R)$ は全射であることに注意せよ
（Algebra, Lemma \ref{algebra-lemma-integral-overring-surjective} を参照）。
従って Lemma \ref{more-algebra-lemma-split-image} を用いると、
$I = \bigcap J_{\underline{k}}$ は $V(I) = \Spec(R)$ を満たすイデアルであると
結論する。これは $I \subset \sqrt{(0)}$ を意味し、$R$ は Noether 環なので
$I$ は冪零である。写像 $\Phi_{\underline{k}}$ は可換図式
$$
\xymatrix{
S \ar[rr] & & R/J_{\underline{k}} \\
& R \ar[lu] \ar[ru]
}
$$
を誘導する。これから $P/J_{\underline{k}}P$ は
$R/J_{\underline{k}}$ 上射影的であると結論する。
Algebra, Lemma \ref{algebra-lemma-projective-modulo-two-ideals} により、
$P/IP$ は $R/I$ 上射影的である。Lemma
\ref{more-algebra-lemma-descent-flatness-injective-finite-Noetherian-rings} により
$P$ は $R$ 上平坦なので、Algebra, Lemma
\ref{algebra-lemma-lift-projective} を適用して $P$ は $R$ 上射影的であると
結論できる。
\end{proof}

\section{ねじれのない加群}
\label{more-algebra-section-torsion-flat}

\noindent
本節では、ねじれのない加群と平坦性との関係（特に次元 1 の環上での関係）を論じる。

\begin{definition}
\label{more-algebra-definition-torsion}
$R$ を整域、$M$ を $R$-加群とする。
\begin{enumerate}
\item 非零元 $f \in R$ で $fx = 0$ となるものが存在するとき、
元 $x \in M$ は {\it ねじれ元}であるという。
\item $M$ のねじれ元が $0$ のみであるとき、$M$ は
{\it ねじれがない}という。
\item $M$ のすべての元がねじれ元であるとき、$M$ は
{\it ねじれ加群}であるという。
\end{enumerate}
\end{definition}

\noindent
$R$ を整域とし、$S = R \setminus \{0\}$ を $R$ の非零元からなる乗法的集合とする。
$R$-加群 $M$ がねじれをもたないことと、$M \to S^{-1}M$ が単射であることは
同値である。言い換えれば、$K = S^{-1}R$ を $R$ の分数体とするとき、写像
$M \to M \otimes_R K$ が単射であることと同値である。

\begin{lemma}
\label{more-algebra-lemma-torsion}
$R$ を整域、$M$ を $R$-加群とする。$M$ のねじれ元全体の集合
$M_{tors}$ は写像 $M \to M \otimes_R K$ の核である。
従って $M_{tors}$ は $M$ の $R$-部分加群である。
商加群 $M/M_{tors}$ はねじれをもたない。
\end{lemma}

\begin{proof}
上の議論を参照せよ。
\end{proof}

\begin{lemma}
\label{more-algebra-lemma-localize-torsion}
$R$ を整域、$M$ をねじれのない $R$-加群とする。任意の乗法的集合
$S \subset R$ に対して、加群 $S^{-1}M$ はねじれのない
$S^{-1}R$-加群である。
\end{lemma}

\begin{proof}
省略する。
\end{proof}

\begin{lemma}
\label{more-algebra-lemma-flat-pullback-torsion}
$R \to R'$ を整域の平坦準同型とする。$M$ がねじれのない
$R$-加群ならば、$M \otimes_R R'$ はねじれのない $R'$-加群である。
\end{lemma}

\begin{proof}
$M$ がねじれをもたないならば、$K$ を $R$ の分数体として
$M \subset M \otimes_R K$ は単射である。$R'$ は $R$ 上平坦なので、
$M \otimes_R R' \to (M \otimes_R K) \otimes_R R'$ は単射である。
$M \otimes_R K$ は $K$ のコピーの直和と同型なので、
$K \otimes_R R'$ がねじれをもたないことを示せば十分である。
これは $R'$ の局所化なので正しい。
\end{proof}

\begin{lemma}
\label{more-algebra-lemma-extension-torsion-free}
$R$ を整域とし、$0 \to M \to M' \to M'' \to 0$ を
$R$-加群の短完全列とする。$M$ と $M''$ がねじれをもたないならば、
$M'$ もねじれをもたない。
\end{lemma}

\begin{proof}
省略する。
\end{proof}

\begin{lemma}
\label{more-algebra-lemma-check-torsion}
$R$ を整域、$M$ を $R$-加群とする。このとき $M$ がねじれをもたないことと、
$R$ のすべての極大イデアル $\mathfrak m$ に対して
$M_\mathfrak m$ がねじれのない $R_\mathfrak m$-加群であることは同値である。
\end{lemma}

\begin{proof}
省略する。ヒント：Lemma \ref{more-algebra-lemma-localize-torsion} および Algebra, Lemma
\ref{algebra-lemma-characterize-zero-local} を用いよ。
\end{proof}

\begin{lemma}
\label{more-algebra-lemma-finite-torsion-free-submodule-free}
$R$ を整域、$M$ を有限 $R$-加群とする。このとき $M$ がねじれをもたないことと、
$M$ が有限自由加群の部分加群であることは同値である。
\end{lemma}

\begin{proof}
$M$ が $R^{\oplus n}$ の部分加群ならば、$M$ はねじれをもたない。
逆に $M$ がねじれをもたないと仮定する。$K$ を $R$ の分数体とする。
このとき $M \otimes_R K$ は有限次元 $K$-ベクトル空間である。
このベクトル空間の基底 $e_1, \ldots, e_r$ を選び、
$x_1, \ldots, x_n$ を $M$ の生成元とする。
$b_{ij} \not = 0$ を満たすある $a_{ij}, b_{ij} \in R$ により
$x_i = \sum (a_{ij}/b_{ij}) e_j$ と書く。$b = \prod_{i, j} b_{ij}$ とおく。
$M$ はねじれをもたないので、写像 $M \to M \otimes_R K$ は単射であり、
その像は $R^{\oplus r} = R e_1/b \oplus \ldots \oplus Re_r/b$ に含まれる。
\end{proof}

\begin{lemma}
\label{more-algebra-lemma-torsion-free-finite-noetherian-domain}
$R$ を Noether 整域、$M$ を非零有限 $R$-加群とする。次は同値である。
\begin{enumerate}
\item $M$ はねじれをもたない、
\item $M$ は有限自由加群の部分加群である、
\item $(0)$ は $M$ の唯一の随伴素イデアルである、
\item $(0)$ は $M$ の台に属し、$M$ は性質 $(S_1)$ をもつ、かつ
\item $(0)$ は $M$ の台に属し、$M$ は埋込み随伴素イデアルをもたない。
\end{enumerate}
\end{lemma}

\begin{proof}
(1) と (2) の同値性は Lemma
\ref{more-algebra-lemma-finite-torsion-free-submodule-free} で見た。
(4) と (5) の同値性は Algebra, Lemma
\ref{algebra-lemma-criterion-no-embedded-primes} で見た。
(3) と (5) の同値性は定義から直ちに従う。
ねじれのない加群の局所化はねじれをもたない
（Lemma \ref{more-algebra-lemma-localize-torsion}）。従って
% SOURCE-ERRATUM-CANDIDATE P02-E0810: frozen source says “a M has”, with a spurious article; Japanese supplies the intended subject.
$M$ が $(0)$ と異なる随伴素イデアルをもたないことは明らかである。
従って (1) は (5) を含意する。逆に (5) を仮定する。
$M$ がねじれをもつならば、$R$ のある非零イデアル $I$ に対して
埋め込み $R/I \subset M$ が存在する。従って $M$ は $(0)$ と異なる
随伴素イデアルをもつ
（Algebra, Lemmas \ref{algebra-lemma-ass} および
\ref{algebra-lemma-ass-zero} を参照）。
これは仮定に反する埋込み随伴素イデアルである。
\end{proof}

\begin{lemma}
\label{more-algebra-lemma-flat-torsion-free}
$R$ を整域とする。任意の平坦 $R$-加群はねじれをもたない。
\end{lemma}

\begin{proof}
$x \in R$ が非零ならば $x : R \to R$ は単射である。従って $M$ が
$R$ 上平坦ならば $x : M \to M$ は単射である。ゆえに $M$ が
$R$ 上平坦ならば $M$ はねじれをもたない。
\end{proof}

\begin{lemma}
\label{more-algebra-lemma-valuation-ring-torsion-free-flat}
$A$ を付値環とする。$A$-加群 $M$ が $A$ 上平坦であることと、
$M$ がねじれをもたないことは同値である。
\end{lemma}

\begin{proof}
含意「flat $\Rightarrow$ torsion free」は Lemma
\ref{more-algebra-lemma-flat-torsion-free} である。逆に $M$ がねじれをもたないと仮定する。
平坦性の等式判定法により（Algebra, Lemma
\ref{algebra-lemma-flat-eq} を参照）、$M$ の任意の関係が自明であることを
示さなければならない。そのため、$x_i \in M$ と $a_i \in A$ に対して
$\sum_{i = 1, \ldots, n} a_i x_i = 0$ であると仮定する。
番号を付け直すことにより、すべての $i$ に対して
$v(a_1) \leq v(a_i)$ と仮定してよい。従って、ある $a'_i \in A$ に対して
$a_i = a'_i a_1$ と書ける。$a'_1 = 1$ であることに注意せよ。
$M$ はねじれをもたないので
$x_1 = - \sum_{i \geq 2} a'_i x_i$ である。従って
$y_i = x_i$, $i = 2, \ldots, n$ と選べば
$$
x_1 = \sum\nolimits_{j \geq 2} -a'_j y_j, \quad
x_i = y_i, (i \geq 2)\quad
0 = a_1 \cdot (-a'_j) + a_j \cdot 1 (j \geq 2)
$$
は関係が自明であったことを示す（明示的には、元 $a_{ij}$ は
$a_{11} = 0$, $a_{1j} = -a'_j$ ($j > 1$) および
$a_{ij} = \delta_{ij}$ ($i, j \geq 2$) とおくことで定義される）。
\end{proof}

\begin{lemma}
\label{more-algebra-lemma-dedekind-torsion-free-flat}
$A$ を Dedekind 整域（例えば離散付値環、より一般には PID）とする。
\begin{enumerate}
\item $A$-加群が平坦であることと、ねじれをもたないことは同値である。
\item 有限でねじれのない $A$-加群は有限局所自由である。
\item $A$ が PID ならば、有限でねじれのない $A$-加群は有限自由である。
\end{enumerate}
\end{lemma}

\begin{proof}
（補題の括弧内の注意については、Algebra, Lemma
\ref{algebra-lemma-PID-dedekind} を参照。）(1) の証明。
Lemma \ref{more-algebra-lemma-check-torsion} および Algebra, Lemma
\ref{algebra-lemma-flat-localization} により、$\mathfrak m \subset A$ が
極大であるとき $A_\mathfrak m$ 上で主張を確認すれば十分である。
$A_\mathfrak m$ は離散付値環なので
（Algebra, Lemma \ref{algebra-lemma-characterize-Dedekind}）、
Lemma \ref{more-algebra-lemma-valuation-ring-torsion-free-flat} により主張が従う。

\medskip\noindent
(2) の証明。Algebra, Lemma \ref{algebra-lemma-finite-projective} と (1) から従う。

\medskip\noindent
(3) の証明。$A$ を PID とし、$M$ を有限でねじれのない加群とする。
Lemma \ref{more-algebra-lemma-finite-torsion-free-submodule-free} により、ある $n$ に対して
$M \subset A^{\oplus n}$ である。$n$ に関する帰納法で $M$ が自由であることを示す。
$n = 1$ の場合は、$A$ が PID であるという事実そのものである。
$n > 1$ ならば、$M' \subset A^{\oplus n - 1}$ を
$A^{\oplus n}$ の最後の $n - 1$ 個の直和因子への射影の像とする。
このとき短完全列 $0 \to I \to M \to M' \to 0$ を得る。ここで $I$ は
$M$ と $A^{\oplus n}$ の最初の直和因子 $A$ との交わりである。
帰納法により $M$ は有限自由 $A$-加群の拡大であり、従って有限自由である。
\end{proof}

\begin{lemma}
\label{more-algebra-lemma-hom-into-torsion-free}
$R$ を整域、$M$, $N$ を $R$-加群とする。$N$ がねじれをもたないならば、
$\Hom_R(M, N)$ もねじれをもたない。
\end{lemma}

\begin{proof}
全射 $\bigoplus_{i \in I} R \to M$ を選ぶ。
このとき $\Hom_R(M, N) \subset \prod_{i \in I} N$ である。
\end{proof}

\section{加群の階数}
\label{more-algebra-section-rank}

\noindent
ここで定義を与える。

\begin{definition}
\label{more-algebra-definition-rank}
$R$ を分数体 $K$ をもつ整域とする。$R$-加群 $M$ の {\it 階数}とは
$\dim_K(M \otimes_R K)$ のことである。
\end{definition}

\noindent
両方の定義が適用できる場合、この定義は Algebra, Definition
\ref{algebra-definition-locally-free} における階数 $r$ の局所自由加群という
概念と矛盾しない。

\begin{lemma}
\label{more-algebra-lemma-torsion-does-not-matter-rank}
$R$ を整域とする。$M \to M'$ が、核と余核がねじれ加群である
$R$-加群写像ならば、$M$ の階数は $M'$ の階数に等しい。
\end{lemma}

\begin{proof}
省略する。ヒント：$K$ が $R$ の分数体ならば、誘導写像
$M \otimes_R K \to M' \otimes_R K$ は同型である。
\end{proof}

\begin{lemma}
\label{more-algebra-lemma-rank-additive}
$R$ を整域とし、$0 \to M \to M' \to M'' \to 0$ を
$R$-加群の短完全列とする。このとき $M'$ の階数は $M$ と $M''$ の階数の和である。
\end{lemma}

\begin{proof}
省略する。
\end{proof}

\begin{lemma}
\label{more-algebra-lemma-rank-tensor}
$R$ を整域とし、$M$ と $N$ を $R$-加群とする。
\begin{enumerate}
\item $M \otimes_R N$ の階数は $M$ と $N$ の階数の積である。
\item $M$ が有限表示 $R$-加群ならば、$\Hom_R(M, N)$ の階数は
$M$ と $N$ の階数の積である。
\end{enumerate}
\end{lemma}

\begin{proof}
(1) はベクトル空間に対する同じ事実から従う。
$M$ が有限表示 $R$-加群であると仮定する。このとき Algebra, Lemma
\ref{algebra-lemma-hom-from-finitely-presented} により、
$K$-ベクトル空間として $\Hom_R(M, N) \otimes_R K$ は
$\Hom_K(M \otimes_R K, N \otimes_R K)$ と同型である。
$M \otimes_R K$ も有限次元なので、ベクトル空間に対する対応する事実から
結論が従う。
\end{proof}

\begin{lemma}
\label{more-algebra-lemma-dominant-pullback-rank}
$R \subset R'$ を整域の拡大とする。$M$ が $R$-加群ならば、
$R$ 上の $M$ の階数は $R'$ 上の $M \otimes_R R'$ の階数に等しい。
\end{lemma}

\begin{proof}
これはベクトル空間の次元が基礎体の拡大で不変だからである。
\end{proof}

\begin{lemma}
\label{more-algebra-lemma-finite-module-rank}
$R$ を整域、$M$ を有限 $R$-加群とする。このとき
\begin{enumerate}
\item $M$ は有限階数 $r \geq 0$ をもつ、
\item 核と余核がねじれ加群である写像 $M \to R^{\oplus r}$ が存在する、
\item $M_f$ が階数 $r$ の自由加群となるような
$f \in R$, $f \not = 0$ が存在する、かつ
\item 余核がねじれ加群である単射 $R^{\oplus r} \to M$ が存在する。
\end{enumerate}
\end{lemma}

\begin{proof}
$M$ は有限なので、$M \otimes_R K$ は $R$ の分数体 $K$ 上の有限次元
ベクトル空間である。従って階数 $r$ は有限である。

\medskip\noindent
$M$ の生成元 $x_1, \ldots, x_n$ と、$M \otimes_R K$ の基底
$e_1, \ldots, e_r$ を選ぶ。$b_{ij} \not = 0$ を満たすある
% SOURCE-ERRATUM-CANDIDATE P02-E0811: frozen source has doubled spacing before “some” and a space before the comma in a_{ij} , b_{ij}; Japanese normalizes typography.
$a_{ij} , b_{ij} \in R$ によって
$x_i \otimes 1 = \sum (a_{ij}/b_{ij}) e_j$ と書ける。
$b = \prod b_{ij}$ とおき、$a_{ij}$ を $a_{ij}b/b_{ij}$ で置き換えることにより、
$x_i \otimes 1 = \sum (a_{ij}/b) e_j$ と仮定してよい。
$M$ に関係式 $\sum r_i x_i = 0$ があるならば、すべての
$j = 1, \ldots, r$ に対して $R$ で $\sum r_ia_{ij} = 0$ であることに注意せよ。

\medskip\noindent
(2) の証明。上の最後の注意により、$x_i$ を $\sum a_{ij} e_j$ に送る写像
% SOURCE-ERRATUM-CANDIDATE P02-E0812: frozen source places a period before “is well defined,” splitting the subject from its predicate.
$M \to \bigoplus Re_j$ が定まり、これは良定義である。
$e_1, \ldots, e_r$ が $M \otimes_R K$ の基底をなすことから、
この写像は $K$ をテンソルした後に同型となる。従ってこれは (2) の写像である。

\medskip\noindent
(3) の証明。(2) で構成した写像 $M \to R^{\oplus r}$ の余核は有限かつ
ねじれ加群である。従ってこの余核を零化する
$g \in R$, $g \not = 0$ を見つけられる。言い換えれば、写像
% SOURCE-ERRATUM-CANDIDATE P02-E0813: frozen source begins a sentence fragment “In other words, such that the map …”.
$M_g \to R_g^{\oplus r}$ は全射である。従って分裂
$M_g = N \oplus R_g^{\oplus r}$ を選べる。このとき $N$ は
$N \otimes_{R_g} K = 0$ を満たす有限 $R_g$-加群である。
言い換えれば、$N$ は $R_g$ 上有限なねじれ加群である。
従って $N$ を零化する $g' \in R$ を見つけられる。
$f= gg'$ とおくと (3) が成り立つ。

\medskip\noindent
(4) を証明するため、(3) の $f$ に対する $M_f$ の基底
$y_1, \ldots, y_r$ を選ぶ。ある $m_j \in M$ と $t_j \geq 0$ により
$y_j = m_j/f^{t_j}$ と書く。第 $j$ 基底ベクトルを $m_j$ に送る写像
$R^{\oplus r} \to M$ が (4) の写像となる。詳細は省略する。
\end{proof}

\section{反射的加群}
\label{more-algebra-section-reflexive}

\noindent
ここで定義を与える。

\begin{definition}
\label{more-algebra-definition-reflexive}
$R$ を整域とする。$m \in M$ を、$\varphi \in \Hom_R(M, R)$ を
$\varphi(m) \in R$ に送る写像に対応させる標準写像
$$
j : M \longrightarrow \Hom_R(\Hom_R(M, R), R)
$$
が同型であるとき、$R$-加群 $M$ は {\it 反射的}であるという。
\end{definition}

\noindent
この定義はより一般の環に対しても行えるが、既に上の定義には難点がある。
Noether 整域と有限でねじれのない加群に制限し、さらに（おそらく）
$R$ に何らかの正則性条件（例えば $R$ が正規であること）を課すのが賢明であろう。

\begin{lemma}
\label{more-algebra-lemma-reflexive-torsion-free}
$R$ を整域、$M$ を $R$-加群とする。
\begin{enumerate}
\item $M$ が反射的ならば、$M$ はねじれをもたない。
\item $M$ が有限ならば、写像
$j : M \to \Hom_R(\Hom_R(M, R), R)$ の核と余核はねじれ加群である。
\item $M$ が有限ならば、$j$ が単射であることと $M$ がねじれをもたないことは
同値である。
\end{enumerate}
\end{lemma}

\begin{proof}
(1) を見るには、Lemma \ref{more-algebra-lemma-hom-into-torsion-free} により
$M$ の二重双対がねじれをもたないことを用いる。
$M$ が有限であると仮定し、$R$ 上の $M$ の生成元
$x_1, \ldots, x_n$ を選ぶ。$K$ を $R$ の分数体とする。
番号を付け直すことにより、ある $0 \leq r \leq n$ に対して
$x_1, \ldots, x_r$ の像が $K$ 上 $M \otimes_R K$ の基底になると仮定してよい。
$r < i \leq n$ に対して、ある $k_{ij} \in K$ により
$x_i \otimes 1 = \sum_{1 \leq j \leq r} x_j \otimes k_{ij}$ と書く。
$a_{ij} = ck_{ij} \in R$ となり、かつ $M$ において
% SOURCE-ERRATUM-CANDIDATE P02-E0814: frozen source contains doubled spacing after the equality sign; the formal span remains exact.
$cx_i =  \sum_{1 \leq j \leq r} a_{ij} x_j$ となるような $c \in R$ を選ぶ
（細部を一つ省略する）。すると写像
$$
R^{\oplus r} \xrightarrow{\alpha} M \xrightarrow{\beta} R^{\oplus r}
\xrightarrow{\alpha} M
$$
を得て、$\alpha \circ \beta$ と $\beta \circ \alpha$ はともに $c$ 倍で与えられる。
ここで $\alpha$ は $x_1, \ldots, x_r$ を用い、$\beta$ は
$1 \leq i \leq r$ に対して $x_i$ を第 $i$ 基底ベクトルの $c$ 倍に送り、
$r < i \leq n$ に対してベクトル $(a_{i1}, \ldots, a_{ir})$ に送る。
可換図式
$$
\xymatrix{
R^{\oplus r} \ar[r] \ar[d]^{\text{id}} &
M \ar[r] \ar[d]^j &
R^{\oplus r} \ar[d]^{\text{id}} \ar[r] &
M \ar[d]^j \\
R^{\oplus r} \ar[r] &
\Hom_R(\Hom_R(M, R), R) \ar[r] &
R^{\oplus r} \ar[r] &
\Hom_R(\Hom_R(M, R), R)
}
$$
を考える。もちろん、下段の連続する矢印の合成は $c$ 倍に等しい。
この図式は $j$ の核と余核が $c$ で零化されることを示し、(2) が従う。
(3) については、$M$ がねじれをもたないならば $\beta$ は単射であり、
これは $j$ が単射であることを含意する。逆は明らかである。実際、
$M$ の二重双対がねじれをもたないことは既に見た。
\end{proof}

\begin{lemma}
\label{more-algebra-lemma-cokernel-map-double-dual-dvr}
$R$ を離散付値環、$M$ を有限 $R$-加群とする。このとき写像
$j : M \to \Hom_R(\Hom_R(M, R), R)$ は全射である。
\end{lemma}

\begin{proof}
$M_{tors} \subset M$ をねじれ部分加群とする。このとき
$$
\Hom_R(M, R) = \Hom_R(M/M_{tors}, R)
$$
である（これは任意の整域上で成り立つ）。
従って $M$ がねじれをもたないと仮定してよい。このとき Lemma
\ref{more-algebra-lemma-dedekind-torsion-free-flat} により $M$ は自由であり、補題は明らかである。
\end{proof}

\begin{lemma}
\label{more-algebra-lemma-check-reflexive}
$R$ を Noether 整域、$M$ を有限 $R$-加群とする。次は同値である。
\begin{enumerate}
\item $M$ は反射的である、
\item すべての素イデアル $\mathfrak p \subset R$ に対して
$M_\mathfrak p$ は反射的 $R_\mathfrak p$-加群である、かつ
\item $R$ のすべての極大イデアル $\mathfrak m$ に対して
$M_\mathfrak m$ は反射的 $R_\mathfrak m$-加群である。
\end{enumerate}
\end{lemma}

\begin{proof}
写像 $j : M \to \Hom_R(\Hom_R(M, R), R)$ の素イデアル
$\mathfrak p$ における局所化は、Noether 局所整域 $R_\mathfrak p$ 上の加群
$M_\mathfrak p$ に対する対応する写像である。Algebra, Lemma
\ref{algebra-lemma-hom-from-finitely-presented} を参照せよ。
従って補題は Algebra, Lemma
\ref{algebra-lemma-characterize-zero-local} により成り立つ。
\end{proof}

\begin{lemma}
\label{more-algebra-lemma-sequence-reflexive}
$R$ を Noether 整域とする。
% SOURCE-ERRATUM-CANDIDATE P02-E0815: frozen source says “Let 0→M→M'→M'' an exact sequence”, omitting “be”.
$0 \to M \to M' \to M''$ を有限 $R$-加群の完全列とする。
$M'$ が反射的で $M''$ がねじれをもたないならば、$M$ は反射的である。
\end{lemma}

\begin{proof}
有限 $R$-加群 $N$, $N'$ に対して $\Hom_R(N, N')$ は有限
$R$-加群であることを、以下断りなく用いる。Algebra, Lemma
\ref{algebra-lemma-ext-noetherian} を参照せよ。
双対をとると列
$$
\Hom_R(M, R) \leftarrow \Hom_R(M', R) \leftarrow \Hom_R(M'', R)
$$
を得る。もう一度双対をとると可換図式
{\small
$$
\xymatrix{
\Hom_R(\Hom_R(M, R), R) \ar[r]_j &
\Hom_R(\Hom_R(M', R), R) \ar[r] &
\Hom_R(\Hom_R(M'', R), R) \\
M \ar[u] \ar[r] & M' \ar[u] \ar[r] & M'' \ar[u]
}
$$
}
を得る。上段が完全であるかは分からない。しかし仮定により中央の縦の矢印は
同型で、右の縦の矢印は単射である
（Lemma \ref{more-algebra-lemma-reflexive-torsion-free}）。$j$ が単射であると主張する。
この主張を仮定すれば、図式追跡により左の縦の矢印が同型であること、すなわち
$M$ が反射的であることが分かる。

\medskip\noindent
主張の証明。$Q$ を定義する完全列
$\Hom_R(M',R)\to \Hom_R(M,R)\to Q \to 0$ を考える。
Algebra, Lemma \ref{algebra-lemma-hom-from-finitely-presented} を適用すると
$$
\Hom_K(M'\otimes_R K,K) \to
\Hom_K(M\otimes_R K,K) \to
Q\otimes_R K\to 0
$$
を得る。しかし $M \otimes_R K \to M' \otimes_R K$ はベクトル空間の単射で、
従って分裂単射であるから、$Q \otimes_R K = 0$、すなわち $Q$ はねじれ加群である。
すると完全列
$$
0 \to \Hom_R(Q,R) \to
\Hom_R(\Hom_R(M,R),R) \to
\Hom_R(\Hom_R(M',R),R)
$$
を得て、$Q$ はねじれ加群なので $\Hom_R(Q,R)=0$ である。
\end{proof}

\begin{lemma}
\label{more-algebra-lemma-characterize-reflexive}
$R$ を Noether 整域、$M$ を有限 $R$-加群とする。次は同値である。
\begin{enumerate}
\item $M$ は反射的である、
\item $F$ が有限自由で $N$ がねじれをもたない短完全列
$0 \to M \to F \to N \to 0$ が存在する。
\end{enumerate}
\end{lemma}

\begin{proof}
有限自由加群が反射的であることに注意せよ。
Lemma \ref{more-algebra-lemma-sequence-reflexive} により (2) は (1) を含意する。
$M$ が反射的であると仮定し、表示
$R^{\oplus m} \to R^{\oplus n} \to \Hom_R(M, R) \to 0$ を選ぶ。
双対をとると完全列
$$
0 \to \Hom_R(\Hom_R(M, R), R) \to R^{\oplus n} \to N \to 0
$$
を得る。ここで
$N = \Im(R^{\oplus n} \to R^{\oplus m})$ はねじれのない加群である。
$M = \Hom_R(\Hom_R(M, R), R)$ なので、(2) のような完全列を得る。
\end{proof}

\begin{lemma}
\label{more-algebra-lemma-flat-pullback-reflexive}
$R \to R'$ を Noether 整域の平坦準同型とする。$M$ が有限反射的
$R$-加群ならば、$M \otimes_R R'$ は有限反射的 $R'$-加群である。
\end{lemma}

\begin{proof}
$F$ が有限自由で $N$ がねじれをもたない短完全列
$0 \to M \to F \to N \to 0$ を選ぶ。Lemma
\ref{more-algebra-lemma-characterize-reflexive} を参照せよ。$R \to R'$ は平坦なので、
$F \otimes_R R'$ が有限自由で $N \otimes_R R'$ がねじれをもたない
短完全列
$0 \to M \otimes_R R' \to F \otimes_R R' \to N \otimes_R R' \to 0$
を得る（Lemma \ref{more-algebra-lemma-flat-pullback-torsion}）。
従って Lemma \ref{more-algebra-lemma-characterize-reflexive} により
$M \otimes_R R'$ は反射的である。
\end{proof}

\begin{lemma}
\label{more-algebra-lemma-dual-reflexive}
$R$ を Noether 整域、$M$ を有限 $R$-加群、$N$ を有限反射的
$R$-加群とする。このとき $\Hom_R(M, N)$ は反射的である。
\end{lemma}

\begin{proof}
表示 $R^{\oplus m} \to R^{\oplus n} \to M \to 0$ を選ぶ。
すると
$$
0 \to \Hom_R(M, N) \to N^{\oplus n} \to N' \to 0
$$
を得る。ここで
$N' = \Im(N^{\oplus n} \to N^{\oplus m})$ はねじれをもたない。
Lemma \ref{more-algebra-lemma-sequence-reflexive} により結論する。
\end{proof}

\begin{definition}
\label{more-algebra-definition-reflexive-hull}
$R$ を Noether 整域、$M$ を有限 $R$-加群とする。加群
$$
M^{**} = \Hom_R(\Hom_R(M, R), R)
$$
を $M$ の {\it 反射包}という。
\end{definition}

\noindent
反射包は Lemma \ref{more-algebra-lemma-dual-reflexive} により反射的なので、この定義には意味がある。
対応 $M \mapsto M^{**}$ は関手である。
$\varphi : M \to N$ が反射的 $R$-加群 $N$ への $R$-加群写像ならば、
$\varphi$ は $M$ の反射包を経由して $M \to M^{**} \to N$ と分解する。
別の言い方をすれば、反射包をとる関手は、Noether 整域 $R$ 上の包含関手
$$
\text{finite reflexive modules} \subset
\text{finite modules}
$$
の左随伴である。

\begin{lemma}
\label{more-algebra-lemma-hom-into-depth}
$R$ を Noether 局所環とし、$M$, $N$ を有限 $R$-加群とする。
\begin{enumerate}
\item $N$ の深さが $\geq 1$ ならば、$\Hom_R(M, N)$ の深さも $\geq 1$ である。
\item $N$ の深さが $\geq 2$ ならば、$\Hom_R(M, N)$ の深さも $\geq 2$ である。
\end{enumerate}
\end{lemma}

\begin{proof}
表示 $R^{\oplus m} \to R^{\oplus n} \to M \to 0$ を選ぶ。
双対をとると完全列
$$
0 \to \Hom_R(M, N) \to N^{\oplus n} \to N' \to 0
$$
を得る。ここで $N' = \Im(N^{\oplus n} \to N^{\oplus m})$ である。
深さ $\geq 1$ の加群の部分加群は深さ $\geq 1$ であり、これは定義から
直ちに従う。従って (1) は明らかである。(2) については、ここでは仮定と
先の注意から $N'$ の深さが $\geq 1$ であることに注意せよ。
% SOURCE-ERRATUM-CANDIDATE P02-E0816: frozen source has subject--verb disagreement in “the assumption and the previous remark implies”.
加群 $N^{\oplus n}$ の深さは $\geq 2$ である。
Algebra, Lemma \ref{algebra-lemma-depth-in-ses} から
$\Hom_R(M, N)$ の深さが $\geq 2$ であると結論する。
\end{proof}

\begin{lemma}
\label{more-algebra-lemma-hom-into-S2}
$R$ を Noether 環とし、$M$, $N$ を有限 $R$-加群とする。
\begin{enumerate}
\item $N$ が性質 $(S_1)$ をもつならば、$\Hom_R(M, N)$ も性質 $(S_1)$ をもつ。
\item $N$ が性質 $(S_2)$ をもつならば、$\Hom_R(M, N)$ も性質 $(S_2)$ をもつ。
\item $R$ が整域で、$N$ がねじれをもたず $(S_2)$ ならば、
$\Hom_R(M, N)$ はねじれをもたず、性質 $(S_2)$ をもつ。
\end{enumerate}
\end{lemma}

\begin{proof}
有限 $R$-加群については素イデアルでの局所化と $\Hom_R$ をとることが可換なので
（Algebra, Lemma \ref{algebra-lemma-ext-noetherian}）、(1) と (2) は
Lemma \ref{more-algebra-lemma-hom-into-depth} から直ちに従う。(3) は (2) と
Lemma \ref{more-algebra-lemma-hom-into-torsion-free} から従う。
\end{proof}

\begin{lemma}
\label{more-algebra-lemma-check-injective-on-ass}
$R$ を Noether 環とし、$\varphi : M \to N$ を $R$-加群の写像とする。
$R$ の任意の素イデアル $\mathfrak p$ に対して、次の少なくとも一方が
成り立つと仮定する。
\begin{enumerate}
\item $M_\mathfrak p \to N_\mathfrak p$ は単射である、または
\item $\mathfrak p \not \in \text{Ass}(M)$ である。
\end{enumerate}
このとき $\varphi$ は単射である。
\end{lemma}

\begin{proof}
$\mathfrak p$ を $\Ker(\varphi)$ の随伴素イデアルとする。このとき
$M_\mathfrak p \to N_\mathfrak p$ の核に属し、その零化イデアルが
$\mathfrak pR_\mathfrak p$ である元 $x \in M_\mathfrak p$ が存在する
（Algebra, Lemma \ref{algebra-lemma-associated-primes-localize}）。
これはいずれの場合にも不可能である。従って
$\text{Ass}(\Ker(\varphi)) = \emptyset$ であり、Algebra, Lemma
\ref{algebra-lemma-ass-zero} により $\Ker(\varphi) = 0$ と結論する。
\end{proof}

\begin{lemma}
\label{more-algebra-lemma-check-isomorphism-via-depth-and-ass}
$R$ を Noether 環とし、$\varphi : M \to N$ を $R$-加群の写像とする。
$M$ は有限であり、$R$ の任意の素イデアル $\mathfrak p$ に対して
次のいずれかが成り立つと仮定する。
\begin{enumerate}
\item $M_\mathfrak p \to N_\mathfrak p$ は同型である、または
\item $\text{depth}(M_\mathfrak p) \geq 2$ かつ
$\mathfrak p \not \in \text{Ass}(N)$ である。
\end{enumerate}
このとき $\varphi$ は同型である。
\end{lemma}

\begin{proof}
Lemma \ref{more-algebra-lemma-check-injective-on-ass} により $\varphi$ は単射である。
$M$ の像を含む有限生成 $R$-部分加群 $N' \subset N$ をとる。
% SOURCE-ERRATUM-CANDIDATE P02-E0817: frozen source says “an finitely generated R-module”; Japanese uses the grammatical determiner-free construction.
このとき $\text{Ass}(N_\mathfrak p) = \emptyset$ ならば
$\text{Ass}(N'_\mathfrak p) = \emptyset$ である。
% SOURCE-ERRATUM-CANDIDATE P02-E0818: frozen proof invokes emptiness of Ass(N_p), whereas the hypothesis is p notin Ass(N); the needed inherited condition is p notin Ass(N').
従ってこの補題の仮定は $M \to N'$ に対しても成り立つ。
$\varphi$ が同型であることを示すには、そのようなすべての
$N' \subset N$ について同じことを示せば十分である。従って $N$ を有限
$R$-加群としてよい。この場合、
$\mathfrak p \not \in \text{Ass}(N) \Rightarrow
\text{depth}(N_\mathfrak p) \geq 1$ である。Algebra, Lemma
\ref{algebra-lemma-ideal-nonzerodivisor} を参照せよ。$Q$ を定める短完全列
$$
0 \to M \to N \to Q \to 0
$$
を考える。条件を見ると、場合 (1) には $Q_\mathfrak p = 0$ であり、
場合 (2) には Algebra, Lemma \ref{algebra-lemma-depth-in-ses} により
$\text{depth}(Q_\mathfrak p) \geq 1$ である。従って $Q$ は随伴素イデアルを
一つももたず、Algebra, Lemma \ref{algebra-lemma-ass-zero} により $Q = 0$ である。
\end{proof}

\begin{lemma}
\label{more-algebra-lemma-isom-depth-2-torsion-free}
$R$ を Noether 整域とし、$\varphi : M \to N$ を $R$-加群の写像とする。
$M$ は有限で $N$ はねじれをもたず、$R$ の任意の素イデアル
$\mathfrak p$ に対して次のいずれかが成り立つと仮定する。
\begin{enumerate}
\item $M_\mathfrak p \to N_\mathfrak p$ は同型である、または
\item $\text{depth}(M_\mathfrak p) \geq 2$ である。
\end{enumerate}
このとき $\varphi$ は同型である。
\end{lemma}

\begin{proof}
これは Lemma \ref{more-algebra-lemma-check-isomorphism-via-depth-and-ass} の特別な場合である。
\end{proof}

\begin{lemma}
\label{more-algebra-lemma-reflexive-depth-2}
$R$ を Noether 整域とし、$M$ を有限 $R$-加群とする。次は同値である。
\begin{enumerate}
\item $M$ は反射的である、
\item $R$ の任意の素イデアル $\mathfrak p$ に対して次のいずれかが成り立つ。
\begin{enumerate}
\item $M_\mathfrak p$ は反射的 $R_\mathfrak p$-加群である、または
\item $\text{depth}(M_\mathfrak p) \geq 2$ である。
\end{enumerate}
\end{enumerate}
\end{lemma}

\begin{proof}
(1) が成り立つならば、Lemma \ref{more-algebra-lemma-check-reflexive} により
$M_\mathfrak p$ はすべての素イデアル $\mathfrak p$ に対して反射的である。
% SOURCE-ERRATUM-CANDIDATE P02-E0819: frozen source says “for all primes of p”; Japanese supplies “for all primes p”.
従って (1) $\Rightarrow$ (2) である。(2) を仮定する。
$N = \Hom_R(\Hom_R(M, R), R)$ とおく。このとき $R$ の任意の素イデアル
$\mathfrak p$ に対して
$$
N_\mathfrak p =
\Hom_{R_\mathfrak p}(
\Hom_{R_\mathfrak p}(M_\mathfrak p, R_\mathfrak p), R_\mathfrak p)
$$
である。Algebra, Lemma \ref{algebra-lemma-hom-from-finitely-presented} を
参照せよ。写像 $j : M \to N$ に Lemma
\ref{more-algebra-lemma-isom-depth-2-torsion-free} を適用する。$M$ は有限であり、
Lemma \ref{more-algebra-lemma-hom-into-torsion-free} により $N$ はねじれをもたないので、
これは可能である。場合 (2)(a) には写像
$M_\mathfrak p \to N_\mathfrak p$ は同型であり、場合 (2)(b) には
$\text{depth}(M_\mathfrak p) \geq 2$ である。
\end{proof}

\begin{lemma}
\label{more-algebra-lemma-reflexive-S2}
$R$ を Noether 整域、$M$ を有限反射的 $R$-加群、
$\mathfrak p \subset R$ を素イデアルとする。
\begin{enumerate}
\item $\text{depth}(R_\mathfrak p) \geq 2$ ならば、
$\text{depth}(M_\mathfrak p) \geq 2$ である。
\item $R$ が $(S_2)$ ならば、$M$ も $(S_2)$ である。
\end{enumerate}
\end{lemma}

\begin{proof}
反射包 $\Hom_R(\Hom_R(M, R), R)$ の形成は局所化と可換なので
（Algebra, Lemma \ref{algebra-lemma-hom-from-finitely-presented}）、
(1) は Lemma \ref{more-algebra-lemma-hom-into-depth} から従う。
(2) は Lemma \ref{more-algebra-lemma-hom-into-S2} から直ちに従う。
\end{proof}

\begin{example}
\label{more-algebra-example-ring-not-S2}
上と下の結果は反射性が条件 $(S_2)$ と関係することを示唆するが、
過度に楽観的な予想を防ぐ例を挙げる。$k$ を体とする。
$R$ を $1, y, x^2, xy, x^3$ で生成される $k[x, y]$ の
$k$-部分代数とする。このとき $R$ は $(S_2)$ ではない。従って
$R$ を $R$-加群とみれば、$(S_2)$ でない反射的 $R$-加群の例となる。
$M = k[x, y]$ を $R$-加群とみなす。このとき
$$
\Hom_R(M, R) = \mathfrak m = (y, x^2, xy, x^3)
\quad\text{and}\quad
\Hom_R(\mathfrak m, R) = M
$$
であるから $M$ は反射的 $R$-加群である。また $M$ は $R$-加群として
$(S_2)$ である（計算は省略する）。従って $R$ は反射的な $(S_2)$ 加群をもつ
Noether 整域であるが、$R$ 自身は $(S_2)$ ではない。
\end{example}

\begin{lemma}
\label{more-algebra-lemma-reflexive-over-normal}
$R$ を分数体 $K$ をもつ Noether 正規整域とし、$M$ を有限
$R$-加群とする。次は同値である。
\begin{enumerate}
\item $M$ は反射的である、
\item $M$ はねじれをもたず、性質 $(S_2)$ をもつ、
\item $M$ はねじれをもたず、
$M = \bigcap_{\text{height}(\mathfrak p) = 1} M_{\mathfrak p}$
である。ここで共通部分は $M_K = M \otimes_R K$ の中でとる。
\end{enumerate}
\end{lemma}

\begin{proof}
Algebra, Lemma \ref{algebra-lemma-criterion-normal} により、
$R$ は $(R_1)$ と $(S_2)$ を満たす。

\medskip\noindent
(1) を仮定する。このとき Lemma \ref{more-algebra-lemma-reflexive-torsion-free} により
$M$ はねじれをもたず、Lemma \ref{more-algebra-lemma-reflexive-S2} により $(S_2)$ を
満たす。従って (2) が成り立つ。

\medskip\noindent
(2) を仮定する。定義により
$M' = \bigcap_{\text{height}(\mathfrak p) = 1} M_{\mathfrak p}$ は写像
$$
M_K
\longrightarrow
\bigoplus\nolimits_{\text{height}(\mathfrak p) = 1} M_K/M_\mathfrak p
\subset
\prod\nolimits_{\text{height}(\mathfrak p) = 1} M_K/M_\mathfrak p
$$
の核である。示したとおり、この写像は実際に直和を経由することに注意せよ。
実際、$a/b \in K$ を与えると、$b \in \mathfrak p$ となる高さ $1$ の
素イデアル $\mathfrak p$ は高々有限個である。$\mathfrak p_0$ を高さ $1$ の
素イデアルとする。このとき $\mathfrak p = \mathfrak p_0$ でない限り
$(M_K/M_\mathfrak p)_{\mathfrak p_0} = 0$ であり、等しい場合には
$(M_K/M_\mathfrak p)_{\mathfrak p_0} = M_K/M_{\mathfrak p_0}$ を得る。
従って局所化の完全性と局所化が直和と可換であることから、
$M'_{\mathfrak p_0} = M_{\mathfrak p_0}$ である。$M$ は高さ $> 1$ の
素イデアルにおいて深さ $\geq 2$ をもつので、Lemma
\ref{more-algebra-lemma-isom-depth-2-torsion-free} により $M \to M'$ は同型である。
従って (3) が成り立つ。

\medskip\noindent
(3) を仮定する。$\mathfrak p$ を高さ $1$ の素イデアルとする。
$(R_1)$ により $R_\mathfrak p$ は離散付値環である。Lemma
\ref{more-algebra-lemma-dedekind-torsion-free-flat} により $M_\mathfrak p$ は有限自由であり、
特に反射的である。従って写像 $M \to M^{**}$ は高さ $1$ のすべての
素イデアル $\mathfrak p$ で同型を誘導する。ゆえに条件
$M = \bigcap_{\text{height}(\mathfrak p) = 1} M_{\mathfrak p}$ から
$M = M^{**}$ が従い、(1) が成り立つ。
\end{proof}

\begin{lemma}
\label{more-algebra-lemma-describe-reflexive-hull}
$R$ を Noether 正規整域とし、$M$ を有限 $R$-加群とする。このとき
$M$ の反射包は
$$
M^{**} =
\bigcap\nolimits_{\text{height}(\mathfrak p) = 1}
M_{\mathfrak p}/(M_\mathfrak p)_{tors} =
\bigcap\nolimits_{\text{height}(\mathfrak p) = 1}
(M/M_{tors})_\mathfrak p
$$
という共通部分であり、これは $M \otimes_R K$ の中でとる。
\end{lemma}

\begin{proof}
$\mathfrak p$ を高さ $1$ の素イデアルとする。
$M_\mathfrak p \to M \otimes_R K$ の核は $M_\mathfrak p$ のねじれ部分加群
$(M_\mathfrak p)_{tors}$ である。さらに
$(M/M_{tors})_\mathfrak p = M_\mathfrak p/(M_\mathfrak p)_{tors}$
であり、これは離散付値環 $R_\mathfrak p$ 上の有限自由加群である
（Lemma \ref{more-algebra-lemma-dedekind-torsion-free-flat}）。従って
$M_\mathfrak p/(M_\mathfrak p)_{tors} \to
(M_\mathfrak p)^{**} = (M^{**})_\mathfrak p$ は同型であり、この補題は
Lemma \ref{more-algebra-lemma-reflexive-over-normal} から従う。
\end{proof}

\begin{lemma}
\label{more-algebra-lemma-integral-closure-reflexive}
$A$ を分数体 $K$ をもつ Noether 正規整域とし、$L$ を $K$ の有限拡大とする。
$L$ における $A$ の整閉包 $B$ が $A$ 上有限ならば、$B$ は
$A$-加群として反射的である。
\end{lemma}

\begin{proof}
共通部分を高さ $1$ の素イデアル $\mathfrak p \subset A$ にわたってとるとき
$B = \bigcap B_\mathfrak p$ であることを示せば十分である。Lemma
\ref{more-algebra-lemma-reflexive-over-normal} を参照せよ。
$b \in \bigcap B_\mathfrak p$ とし、
$x^d + a_1x^{d - 1} + \ldots + a_d$ を $K$ 上の $b$ の最小多項式とする。
$a_i \in A$ を示したい。Algebra, Lemma
\ref{algebra-lemma-minimal-polynomial-normal-domain} により、すべての $i$ と
すべての高さ一の素イデアル $\mathfrak p$ に対して
$a_i \in A_\mathfrak p$ である。従って Algebra, Lemma
\ref{algebra-lemma-normal-domain-intersection-localizations-height-1} から
所望の結論を得る（あるいは $A$ がそれ自身上の反射的加群であることから、
すでに引用した補題を用いてもよい）。
\end{proof}

\section{内容イデアル}
\label{more-algebra-section-content}

\noindent
この定義は予想しているものとは異なるかもしれない。

\begin{definition}
\label{more-algebra-definition-content-ideal}
$A$ を環、$M$ を平坦 $A$-加群、$x \in M$ とする。
$x \in IM$ を満たす $A$ のイデアル $I$ 全体の集合が最小元をもつならば、
それを $x$ の {\it 内容イデアル}という。
\end{definition}

\noindent
$M$ は $A$ 上平坦なので、$A$ のイデアルの対 $I, I'$ に対して
$IM \cap I'M = (I \cap I')M$ であることに注意せよ。これは完全列
$0 \to I \cap I' \to I \oplus I' \to I + I' \to 0$ を $M$ と
テンソルすることから分かる。

\begin{lemma}
\label{more-algebra-lemma-content-finitely-generated}
$A$ を環、$M$ を平坦 $A$-加群、$x \in M$ とする。
$x$ の内容イデアルが存在するならば、それは有限生成である。
\end{lemma}

\begin{proof}
$x \in IM$ とする。このとき $f_i \in I$ および $x_i \in M$ を用いて
$x = \sum_{i = 1, \ldots, n} f_i x_i$ と書ける。従って
$I' = (f_1, \ldots, f_n)$ とおけば $x \in I'M$ である。
\end{proof}

\begin{lemma}
\label{more-algebra-lemma-equal-content}
$(A, \mathfrak m)$ を局所環とし、$u : M \to N$ を平坦 $A$-加群の写像で、
$\overline{u} : M/\mathfrak m M \to N/\mathfrak m N$ が単射であるものとする。
$x \in M$ が内容イデアル $I$ をもつならば、$u(x)$ も内容イデアル $I$ をもつ。
\end{lemma}

\begin{proof}
$u(x) \in IN$ は明らかである。$u(x) \in I'N$ ならば
$u(x) \in (I' \cap I)N$ である。Definition
\ref{more-algebra-definition-content-ideal} に続く議論を参照せよ。従って次を示せば十分である：
% SOURCE-ERRATUM-CANDIDATE P02-E0820: frozen source assumes x∈I'N even though x belongs to M, then concludes u(x)∉I'N; the typed premise is extraneous and contradictory to the intended reduction.
$x \in I'N$ かつ $I' \subset I$, $I' \not =  I$ ならば
% SOURCE-ERRATUM-CANDIDATE P02-E0821: frozen source contains doubled spacing after the relation in I' \not =  I; the formal span remains exact.
$u(x) \not \in I'N$ である。$I/I'$ は零でない有限 $A$-加群なので
（Lemma \ref{more-algebra-lemma-content-finitely-generated}）、Nakayama の補題
（Algebra, Lemma \ref{algebra-lemma-NAK}）により、零でない $A$-加群写像
$\chi : I/I' \to A/\mathfrak m$ が存在する。$I$ は $x$ の内容イデアルなので、
$I'' = \Ker(\chi)$ とおくと $x \not \in I''M$ である。
従って $x$ は写像
$$
IM = I \otimes_A M \xrightarrow{\chi \otimes 1}
A/\mathfrak m \otimes M \cong M/\mathfrak m M
$$
の核に属さない。$\overline{u}$ に関する仮定を適用すると、$u(x)$ は写像
$$
IN = I \otimes_A N \xrightarrow{\chi \otimes 1}
A/\mathfrak m \otimes N \cong N/\mathfrak m N
$$
によって零に写らないと結論でき、証明が完了する。
\end{proof}

\begin{lemma}
\label{more-algebra-lemma-content-exists-flat-Mittag-Leffler}
$A$ を環とし、$M$ を平坦 Mittag--Leffler 加群とする。
このとき $M$ のすべての元は内容イデアルをもつ。
\end{lemma}

\begin{proof}
これは Algebra, Lemma \ref{algebra-lemma-flat-ML-criterion} の特別な場合である。
\end{proof}

\section{平坦性と有限性条件}
\label{more-algebra-section-flat-finite}

\noindent
この節では「平坦 $+$ 有限型 $\Rightarrow$ 有限表示」という型のいくつかの
含意を論じる。この結果には平坦性を扱う章で再び戻る。More on Flatness,
Section \ref{flat-section-introduction} を参照せよ。この型の最初の結果は
Algebra, Lemma \ref{algebra-lemma-finite-flat-module-finitely-presented}
で証明された。

\begin{lemma}
\label{more-algebra-lemma-flat-finite-type-finite-presentation-local-module}
$R$ を環とし、$S = R[x_1, \ldots, x_n]$ を $R$ 上の多項式環、
$M$ を $S$-加群とする。次を仮定する。
\begin{enumerate}
\item $R \to \prod R_{\mathfrak p_j}$ が単射となる $R$ の有限個の
素イデアル $\mathfrak p_1, \ldots, \mathfrak p_m$ が存在する、
\item $M$ は有限 $S$-加群である、
\item $M$ は $R$ 上平坦である、かつ
% SOURCE-ERRATUM-CANDIDATE P02-E0822: frozen source says “M flat over R,” omitting “is”; Japanese supplies the predicate.
\item $R$ の任意の素イデアル $\mathfrak p$ に対して、加群
$M_{\mathfrak p}$ は $S_{\mathfrak p}$ 上有限表示である。
\end{enumerate}
このとき $M$ は $S$ 上有限表示である。
\end{lemma}

\begin{proof}
$S$-加群としての $M$ の表示
$$
0 \to K \to S^{\oplus r} \to M \to 0
$$
を選ぶ。$\mathfrak q$ を $R$ の素イデアル $\mathfrak p$ 上にある
$S$ の素イデアルとする。仮定により、$K' = \sum Sk_j \subset K$ とおけば
$K'_{\mathfrak p} = K_{\mathfrak p}$ かつ
$j = 1, \ldots, m$ に対して
$K'_{\mathfrak p_j} = K_{\mathfrak p_j}$ となるような有限個の元
$k_1, \ldots, k_t \in K$ が存在する。$M' = S^{\oplus r}/K'$ とおけば、
特に $M'_{\mathfrak q} = M_{\mathfrak q}$ である。平坦性の開性
（Algebra, Theorem \ref{algebra-theorem-openness-flatness}）により、
$M'_g$ が $R$ 上平坦となる $g \in S$, $g \not \in \mathfrak q$ が存在する。
従って $M'_g \to M_g$ は平坦 $R$-加群の全射である。可換図式
$$
\xymatrix{
M'_g \ar[r] \ar[d] & M_g \ar[d] \\
\prod (M'_g)_{\mathfrak p_j} \ar[r] & \prod (M_g)_{\mathfrak p_j}
}
$$
を考える。下の矢印は $k_1, \ldots, k_t$ の選び方により同型である。
$R \to \prod R_{\mathfrak p_j}$ は単射で $M'_g$ は $R$ 上平坦なので、
左の垂直矢印は単射である。従って上の水平矢印は単射、ゆえに同型である。
これにより $M_g$ は $S_g$ 上有限表示である。Algebra, Lemma
\ref{algebra-lemma-cover} を適用して結論する。
\end{proof}

\begin{lemma}
\label{more-algebra-lemma-flat-finite-type-finite-presentation-local}
$R \to S$ を環準同型とする。次を仮定する。
\begin{enumerate}
\item $R \to \prod R_{\mathfrak p_j}$ が単射となる $R$ の有限個の
素イデアル $\mathfrak p_1, \ldots, \mathfrak p_m$ が存在する、
\item $R \to S$ は有限型である、
\item $S$ は $R$ 上平坦である、かつ
% SOURCE-ERRATUM-CANDIDATE P02-E0823: frozen source says “S flat over R,” omitting “is”; Japanese supplies the predicate.
\item $R$ の任意の素イデアル $\mathfrak p$ に対して、環
$S_{\mathfrak p}$ は $R_{\mathfrak p}$ 上有限表示である。
\end{enumerate}
このとき $S$ は $R$ 上有限表示である。
\end{lemma}

\begin{proof}
仮定により $S$ は $R$ 上の多項式環の商である。従って結果は
Lemma \ref{more-algebra-lemma-flat-finite-type-finite-presentation-local-module}
から直接従う。
\end{proof}

\begin{lemma}
\label{more-algebra-lemma-flat-graded-finite-type-finite-presentation-module}
$R$ を環とする。$S = R[x_1, \ldots, x_n]$ を $R$ 上の次数付き
多項式代数、すなわち $\deg(x_i) > 0$ だが必ずしも $1$ に等しくないものとする。
$M$ を次数付き $S$-加群とする。次を仮定する。
\begin{enumerate}
\item $R$ は局所環である、
\item $M$ は有限 $S$-加群である、かつ
\item $M$ は $R$ 上平坦である。
\end{enumerate}
このとき $M$ は $S$-加群として有限表示である。
\end{lemma}

\begin{proof}
$M = \bigoplus M_d$ を $M$ の次数付けとする。$M$ の斉次生成元
$m_1, \ldots, m_r \in M$ を選び、$\deg(m_i) = d_i \in \mathbf{Z}$ とする。
これにより表示
$$
0 \to K \to \bigoplus\nolimits_{i = 1, \ldots, r} S(-d_i) \to M \to 0
$$
を得て、各次数 $d$ において短完全列
$$
0 \to K_d \to \bigoplus\nolimits_{i = 1, \ldots, r} S_{d - d_i} \to
M_d \to 0.
$$
を得る。仮定により各 $M_d$ は有限平坦 $R$-加群である。Algebra,
Lemma \ref{algebra-lemma-finite-flat-local} により、各 $M_d$ は有限自由
$R$-加群である。従って各 $K_d$ は有限 $R$-加群である。また Algebra,
Lemma \ref{algebra-lemma-flat-ses} により各 $K_d$ は $R$ 上平坦でもある。
ゆえに再び Algebra, Lemma \ref{algebra-lemma-finite-flat-local} により、
各 $K_d$ は有限自由である。

\medskip\noindent
$\mathfrak m$ を $R$ の極大イデアルとする。$M$ は $R$ 上平坦なので、上の
短完全列は $\kappa = \kappa(\mathfrak m)$ とテンソルした後も短完全である。
環 $S \otimes_R \kappa$ は Noether なので、その像 $\overline{k}_j$ が
$K \otimes_R \kappa$ を $S \otimes_R \kappa$ 上生成する斉次元
$k_1, \ldots, k_t \in K$ が存在する。$\deg(k_j) = e_j$ とする。
従って任意の $d$ に対する写像
$$
\bigoplus\nolimits_{j = 1, \ldots, t} S_{d - e_j}
\longrightarrow
K_d
$$
は $\kappa$ とテンソルした後に全射となる。Nakayama の補題
（Algebra, Lemma \ref{algebra-lemma-NAK}）により、この写像は $R$ 上
全射である。従って $K$ は $k_1, \ldots, k_t$ によって $S$ 上生成され、
証明が完了する。
\end{proof}


\begin{lemma}
\label{more-algebra-lemma-flat-graded-finite-type-finite-presentation}
$R$ を環とし、$S = \bigoplus_{n \geq 0} S_n$ を次数付き $R$-代数、
$M = \bigoplus_{d \in \mathbf{Z}} M_d$ を次数付き $S$-加群とする。
$S$ は $R$-代数として有限生成であり、$S_0$ は有限 $R$-代数であり、
$R \to \prod R_{\mathfrak p_j}$ が単射となる有限個の素イデアル
$\mathfrak p_j$, $i = 1, \ldots, m$ が存在すると仮定する。
% SOURCE-ERRATUM-CANDIDATE P02-E0824: frozen source indexes the primes by j but quantifies i=1,…,m; the exact mismatched span is preserved.
\begin{enumerate}
\item $S$ が $R$ 上平坦ならば、$S$ は有限表示 $R$-代数である。
\item $M$ が $R$-加群として平坦で $S$-加群として有限ならば、
$M$ は $S$-加群として有限表示である。
\end{enumerate}
\end{lemma}

\begin{proof}
$S$ は $R$-代数として有限生成なので、$S_0$-代数としても有限生成である。
次数 $d_1, \ldots, d_n > 0$ の斉次元 $t_1, \ldots, t_n \in S$ で生成されるとする。
$P = R[x_1, \ldots, x_n]$ とおき、$\deg(x_i) = d_i$ とする。
$S_0$ は有限 $R$-加群なので、環写像 $P \to S$, $x_i \to t_i$ は有限である。
(1) を証明するには $S$ が有限表示 $P$-加群であることを示せば十分である。
(2) を証明するには $M$ が有限表示 $P$-加群であることを示せば十分である。
% SOURCE-ERRATUM-CANDIDATE P02-E0825: frozen source has doubled intersentence spacing before “To prove (2)”; Japanese normalizes typography.
従って、$S = P$ が次数付き多項式環で、$M$ が $R$ 上平坦な有限
$S$-加群ならば $M$ は有限表示 $S$-加群であることを示せば十分である。
Lemma \ref{more-algebra-lemma-flat-graded-finite-type-finite-presentation-module} により、
$R$ の任意の素イデアル $\mathfrak p$ に対して
$M_{\mathfrak p}$ は有限表示 $S_{\mathfrak p}$-加群である。従って結果は
Lemma \ref{more-algebra-lemma-flat-finite-type-finite-presentation-local-module} から従う。
\end{proof}

\begin{remark}
\label{more-algebra-remark-when-does-condition-hold}
$R$ を環とする。Lemmas
\ref{more-algebra-lemma-flat-finite-type-finite-presentation-local-module},
\ref{more-algebra-lemma-flat-finite-type-finite-presentation-local}, および
\ref{more-algebra-lemma-flat-graded-finite-type-finite-presentation} で述べた条件を
$R$ が満たすのはいつか。次の各場合には満たされる。
\begin{enumerate}
\item $R$ が局所環である、
\item $R$ が Noether である、
\item $R$ が整域である、
\item $R$ が有限個の極小素イデアルをもつ被約環である、または
\item $R$ が有限個の弱随伴素イデアルをもつ。Algebra, Lemma
\ref{algebra-lemma-zero-at-weakly-ass-zero} を参照せよ。
\end{enumerate}
従ってこれらの補題は上に挙げたすべての場合に成り立つ。
\end{remark}

\noindent
次の補題は More on Flatness, Proposition
\ref{flat-proposition-flat-finite-type-finite-presentation-domain}
で強化される。

\begin{lemma}
\label{more-algebra-lemma-flat-finite-type-valuation-ring-finite-presentation}
\begin{reference}
\cite[Theorem 3]{Nagata-Finitely}
\end{reference}
$A$ を付値環、$A \to B$ を有限型の環写像、$M$ を有限 $B$-加群とする。
\begin{enumerate}
\item $B$ が $A$ 上平坦ならば、$B$ は有限表示 $A$-代数である。
\item $M$ が $A$-加群として平坦ならば、$M$ は $B$-加群として有限表示である。
\end{enumerate}
\end{lemma}

\begin{proof}
$A$-加群が平坦であることとねじれをもたないことが同値であるという事実を用いる。
Lemma \ref{more-algebra-lemma-valuation-ring-torsion-free-flat} を参照せよ。
Algebra, Lemma \ref{algebra-lemma-homogenize} により、$S_0 = A$ で
次数 $1$ の有限個の元によって生成される次数付き $A$-代数 $S$、元
$f \in S_1$、および有限次数付き $S$-加群 $N$ で、
$B \cong S_{(f)}$ および $M \cong N_{(f)}$ を満たすものを見つけられる。
$M$ がねじれをもたないならば、$N$ を $N/N_{tors}$ で置き換えることにより、
$N$ もねじれをもたないとしてよい。Lemma \ref{more-algebra-lemma-torsion} を参照せよ。
同様に、$B$ がねじれをもたないならば、$S$ を $S/S_{tors}$ で置き換える
ことにより、ねじれをもたないものとしてよい。従って場合 (1) には
Lemma \ref{more-algebra-lemma-flat-graded-finite-type-finite-presentation} を適用して、
$S$ が有限表示 $A$-代数であることが分かり、それゆえ
$B = S_{(f)}$ も有限表示 $A$-代数である。(2) を示すには、まず $S$ を
次数付き多項式環で置き換え、次に Lemma
\ref{more-algebra-lemma-flat-graded-finite-type-finite-presentation-module} を適用すればよい。
\end{proof}

\begin{lemma}
\label{more-algebra-lemma-valuation-ring-flat-essentially-finite-type}
$A$ を付値環、$A \to B$ を本質的有限型の局所準同型、
$M$ を有限 $B$-加群とする。
\begin{enumerate}
\item $B$ が $A$ 上平坦ならば、$B$ は $A$ 上本質的有限表示である。
\item $M$ が $A$-加群として平坦ならば、$M$ は $B$-加群として有限表示である。
\end{enumerate}
\end{lemma}

\begin{proof}
仮定により、$B$ は多項式代数 $P = A[x_1, \ldots, x_n]$ を素イデアル
$\mathfrak q$ で局所化した環の商として書ける。場合 (1) には $M = B$ を
$P_\mathfrak q$ 上の有限加群とみなし、場合 (2) には $M$ を
$P_\mathfrak q$ 上の有限加群とみなす。いずれの場合にも、これが有限表示
$P_\mathfrak q$-加群であることを示さなければならない。場合 (2) については
Algebra, Lemma \ref{algebra-lemma-finitely-presented-over-subring}
を参照せよ。

\medskip\noindent
$M$ は $P_\mathfrak q$ 上有限なので、表示
$0 \to K \to P_\mathfrak q^{\oplus r} \to M \to 0$ を選ぶ。
$L = P^{\oplus r} \cap K$ とおく。Algebra, Lemma
\ref{algebra-lemma-submodule-localization} により $K = L_\mathfrak q$ である。
このとき $N = P^{\oplus r}/L$ は $M$ の部分加群であり、従って Lemma
\ref{more-algebra-lemma-valuation-ring-torsion-free-flat} により平坦である。さらに $N$ は
有限 $P$-加群なので、Lemma
\ref{more-algebra-lemma-flat-finite-type-valuation-ring-finite-presentation} により
$P$-加群として有限表示である。局所化は完全なので
（Algebra, Proposition \ref{algebra-proposition-localization-exact}）、
$N_\mathfrak q = M$ であり、結論を得る。
\end{proof}

\section{ブローアップと平坦性}
\label{more-algebra-section-blowup-flat}

\noindent
この節では「ブローアップの後に狭義変換が平坦になる」という形の結果の議論を
始める。この型のさらなる結果は Divisors, Section
\ref{divisors-section-blowup-flat} および More on Flatness, Section
\ref{flat-section-blowup-flat} にある。

\begin{definition}
\label{more-algebra-definition-strict-transform}
$R$ を環、$I \subset R$ をイデアル、$a \in I$ とする。
$R[\frac{I}{a}]$ をアフィン・ブローアップ代数とする。Algebra, Definition
\ref{algebra-definition-blow-up} を参照せよ。$M$ を $R$-加群とする。
{\it $R \to R[\frac{I}{a}]$ に沿う $M$ の狭義変換}とは、
$R[\frac{I}{a}]$-加群
$$
M' = \left(M \otimes_R R[\textstyle{\frac{I}{a}}]\right)/a\text{-power torsion}
$$
のことである。
\end{definition}

\noindent
次はブローアップによる平坦化の非常に弱い形だが、それでも有用な場合がある。

\begin{lemma}
\label{more-algebra-lemma-flatten-on-affine-blowup}
$(R, \mathfrak m)$ を分数体 $K$ をもつ局所整域とする。
$S$ を有限型 $R$-代数、$M$ を有限 $S$-加群とする。
$R$ を支配する任意の付値環 $A \subset K$ に対して、イデアル
$I \subset \mathfrak m$ と零でない元 $a \in I$ で、次を満たすものが存在する。
\begin{enumerate}
\item $I$ は有限生成である、
\item $A$ は $R[\frac{I}{a}]$ 上に中心をもつ、
\item $\mathfrak m$ における $R \to R[\frac{I}{a}]$ のファイバー環は
零でない、かつ
\item $R \to R[\frac{I}{a}]$ に沿う $S$ の狭義変換 $S_{I, a}$ は
$R$ 上平坦かつ有限表示であり、$R \to R[\frac{I}{a}]$ に沿う
$M$ の狭義変換 $M_{I, a}$ は $R$ 上平坦かつ $S_{I, a}$ 上有限表示である。
\end{enumerate}
\end{lemma}

\begin{proof}
$S = R[x_1, \ldots, x_n]/J$ と書き、$N = S \oplus M$ を
$P = R[x_1, \ldots, x_n]$ 上の加群とみなす。$S$ が $R$ 上の多項式代数である
場合に補題を証明できるならば、(1), (2), (3) を満たす $I, a$ で、
$R \to R[\frac{I}{a}]$ に沿う $N$ の狭義変換 $N_{I, a}$ が $R$ 上平坦で、
$P$ の狭義変換 $P_{I, a}$ 上有限表示となるものを見つけられる。
$P_{I, a} = R[\frac{I}{a}][x_1, \ldots, x_n]$ なので（細部を一つ省略する）、
直和因子 $S_{I, a} \subset N_{I, a}$ は $R$ 上平坦であり、
$R[\frac{I}{a}][x_1, \ldots, x_n]$ 上有限表示である。従って
$S_{I, a}$ は有限表示 $R[\frac{I}{a}]$-代数である。さらに直和因子
$M_{I, a} \subset N_{I, a}$ は $R$ 上平坦で、$P_{I, a}$ 上有限表示、
従って $S_{I, a}$ 上でも有限表示である。Algebra, Lemma
\ref{algebra-lemma-finitely-presented-over-subring} を参照せよ。
これにより次の段落で論じる場合に帰着する。

\medskip\noindent
$S = R[x_1, \ldots, x_n]$ と仮定し、表示
$$
0 \to K \to S^{\oplus r} \to M \to 0.
$$
を選ぶ。$M_A$ を $M \otimes_R A$ のねじれ部分加群による商とする。
Lemma \ref{more-algebra-lemma-torsion} を参照せよ。このとき $M_A$ は
$S_A = A[x_1, \ldots, x_n]$ 上の有限加群である。Lemma
\ref{more-algebra-lemma-valuation-ring-torsion-free-flat} により $M_A$ は $A$ 上平坦であり、
Lemma \ref{more-algebra-lemma-flat-finite-type-valuation-ring-finite-presentation} により
$M_A$ は有限表示である。従って、表示 $S_A^{\oplus r} \to M_A$ の核を
$S_A$-加群として生成する有限個の元
$k_1, \ldots, k_t \in S_A^{\oplus r}$ が存在する。(1), (2), (3) を満たす
任意の選択 $a \in I \subset \mathfrak m$ に対し、
$R \to R[\frac{I}{a}]$ に沿う $M$ の狭義変換を $M_{I, a}$ と表す。
これは $S_{I, a} = R[\frac{I}{a}][x_1, \ldots, x_n]$ 上の有限加群である。
Algebra, Lemma \ref{algebra-lemma-valuation-ring-colimit-affine-blowups} により
$A = \colim_{I, a} R[\frac{I}{a}]$ である。従って
$S_A = \colim S_{I, a}$ であり、
$$
\colim M \otimes_R R[\textstyle{\frac{I}{a}}] = M \otimes_R A
$$
である。$I, a$ と持ち上げ $k_1, \ldots, k_t \in S_{I, a}^{\oplus r}$ を選ぶ。
$M_A$ は $M \otimes_R A$ のねじれによる商なので、
$M \otimes_R A$ における $k_1, \ldots, k_t$ の像は、ある零でない元
$\alpha \in A$ によって零化される。$I, a$ を別の対で置き換えれば
（余極限がフィルター付きであることを想起せよ）、
ある零でない $x \in I^n$ に対して $\alpha = x/a^n$ と仮定してよい。
すると $x k_1, \ldots, x k_t$ は $M \otimes_R A$ で零に写る。
さらに $I, a$ を別の対で置き換えると、ある零でない $x \in R$ に対して
$x k_1, \ldots, x k_t$ が
$M \otimes_R R[\frac{I}{a}]$ で零に写ると仮定してよい。
最後に $I, a$ を $xI, xa$ で置き換えることで、
$k_1, \ldots, k_t$ は $M \otimes_R R[\frac{I}{a}]$ の
$a$-冪ねじれ元に写ると仮定してよい。そのような各対 $(I, a)$ に対して
$$
M'_{I, a} = S_{I, a}^{\oplus r}/ \sum S_{I, a}k_j.
$$
とおく。$M_A = S_A^{\oplus r}/ \sum S_Ak_j$ なので、
$M_A = \colim_{I, a} M'_{I, a}$ である。ここでようやく Algebra, Lemma
\ref{algebra-lemma-flat-finite-presentation-limit-flat} (3) を適用し、
上のようなある対 $(I, a)$ に対して $M'_{I, a}$ が平坦であると結論する。
狭義変換の系
$$
M_{I, a} = (M \otimes_R R[\textstyle{\frac{I}{a}}])/a\text{-power torsion}
$$
には、遷移写像が補題の仮定を満たすとは限らないため、この補題を先験的には
適用できない。しかし平坦性はねじれがないことを含意し
（Lemma \ref{more-algebra-lemma-flat-torsion-free}）、さらに $M_{I, a}$ は
$M'_{I, a}$ を $a$-冪ねじれ部分加群で割った商である
（$k_1, \ldots, k_t$ が $M_{I, a}$ で零に写るようにしておいた）ので、
そのような対に対して $M'_{I, a} = M_{I, a}$ と結論でき、証明が完了する。
\end{proof}

\begin{lemma}
\label{more-algebra-lemma-blowup-fitting-ideal}
$R$ を環、$M$ を有限 $R$-加群とする。$k \geq 0$ とし、
$I = \text{Fit}_k(M)$ とおく。任意の $a \in I$ に対し
$R' = R[\frac{I}{a}]$ とおけば、狭義変換
$$
M' = (M \otimes_R R')/a\text{-power torsion}
$$
は $\text{Fit}_k(M') = R'$ を満たす。
\end{lemma}

\begin{proof}
まず $\text{Fit}_k(M \otimes_R R') = IR' = aR'$ であることに注意せよ。
最初の等式は Lemma \ref{more-algebra-lemma-fitting-ideal-basics} の (3) により、
二つ目は Algebra, Lemma \ref{algebra-lemma-affine-blowup} による。
Lemma \ref{more-algebra-lemma-principal-fitting-ideal} と局所化の完全性から、
$R'$ の任意の素イデアル $\mathfrak p'$ に対して
$M'_{\mathfrak p'}$ は $\leq k$ 個の元で生成できる。従って、例えば
Lemma \ref{more-algebra-lemma-fitting-ideal-generate-locally} により
$\text{Fit}_k(M') = R'$ である。
\end{proof}

\begin{lemma}
\label{more-algebra-lemma-blowup-fitting-ideal-locally-free}
$R$ を環、$M$ を有限 $R$-加群とする。$k \geq 0$ とし、
$I = \text{Fit}_k(M)$ とおく。任意の
$\mathfrak p \not \in V(I)$ に対して $M_\mathfrak p$ は階数 $k$ の自由加群である
と仮定する。このとき任意の $a \in I$ に対して
$R' = R[\frac{I}{a}]$ とおけば、狭義変換
$$
M' = (M \otimes_R R')/a\text{-power torsion}
$$
は階数 $k$ の局所自由加群である。
\end{lemma}

\begin{proof}
Lemma \ref{more-algebra-lemma-blowup-fitting-ideal} により
$\text{Fit}_k(M') = R'$ である。Lemma
\ref{more-algebra-lemma-fitting-ideal-finite-locally-free} により、
$\text{Fit}_{k - 1}(M') = 0$ を示せば十分である。Algebra, Lemma
\ref{algebra-lemma-affine-blowup} により
$R' \subset R'_a = R_a$ であることを想起せよ。従って
$\text{Fit}_{k - 1}(M')$ が $R'_a = R_a$ で零に写ることを示せば十分である。
明らかに $(M')_a = M_a$ なので、Fitting イデアルの形成は基底変換と可換する
（Lemma \ref{more-algebra-lemma-fitting-ideal-basics} の (3)）ことから、
$\text{Fit}_{k - 1}(M_a) = 0$ を示すことに帰着する。これは仮定により
$M_a$ が階数 $k$ の有限局所自由加群であること
（Algebra, Lemma \ref{algebra-lemma-finite-projective}）と、すでに引用した
Lemma \ref{more-algebra-lemma-fitting-ideal-finite-locally-free} から従う。
\end{proof}

\begin{lemma}
\label{more-algebra-lemma-blowup-module}
$R$ を環、$M$ を有限 $R$-加群とする。$f \in R$ を、
$M_f$ が階数 $r$ の有限局所自由加群となる元とする。このとき
$V(f) = V(I)$ を満たす有限生成イデアル $I \subset R$ で、
任意の $a \in I$ に対し $R' = R[\frac{I}{a}]$ とおけば、狭義変換
$$
M' = (M \otimes_R R')/a\text{-power torsion}
$$
が階数 $r$ の局所自由加群となるものが存在する。
\end{lemma}

\begin{proof}
全射 $R^{\oplus n} \to M$ を選ぶ。$f$ を逆にすると
$R^{\oplus n}/K \to M$ が同型となるような有限部分加群
$K \subset \Ker(R^{\oplus n} \to M)$ を選ぶ。これは、例えば Algebra,
Lemma \ref{algebra-lemma-finite-projective} により $M_f$ が有限表示なので可能である。
$M_1 = R^{\oplus n}/K$ とおき、$M_1$ に対して補題を証明できると仮定する。
$I \subset R$ を対応するイデアルとする。このとき $a \in I$ に対して写像
$$
M_1' = (M_1 \otimes_R R')/a\text{-power torsion}
\longrightarrow
M' = (M \otimes_R R')/a\text{-power torsion}
$$
は全射である。また、Algebra, Lemma
\ref{algebra-lemma-blowup-in-principal} により
$R'$ で $a$ を逆にすると $R'_a = R_f$ なので、この写像は同型となる。
% SOURCE-ERRATUM-CANDIDATE P02-E0826: frozen source twice states R'_a=R_f, whereas the cited lemma gives R'_f=R'_a=R_a; the exact formula is preserved at both loci.
しかし $a$ は $M'_1$ 上の非零因子なので、表示された写像は同型である。
従って $M$ が有限表示 $R$-加群である場合に補題を証明すれば十分である。

\medskip\noindent
$M$ を有限表示 $R$-加群と仮定する。このとき
$J = \text{Fit}_r(M) \subset S$ は有限生成イデアルである。
% SOURCE-ERRATUM-CANDIDATE P02-E0827: frozen source places J inside undefined S; the ambient ring in this lemma is R and the exact source span is preserved.
$I = fJ$ が使えると主張する。

\medskip\noindent
まず $V(f) = V(I)$ を確認する。包含 $V(f) \subset V(I)$ は明らかである。
逆に $f \not \in \mathfrak p$ ならば、Lemma
\ref{more-algebra-lemma-fitting-ideal-generate-locally} により
$\mathfrak p$ は $V(J)$ の元ではない。従って
$\mathfrak p \not \in V(fJ) = V(I)$ である。

\medskip\noindent
$a \in I$ とし、$R' = R[\frac{I}{a}]$ とおく。ある $b \in J$ によって
$a = fb$ と書ける。Algebra, Lemmas
\ref{algebra-lemma-affine-blowup} および
\ref{algebra-lemma-blowup-add-principal} により
$J R' = b R'$ であり、$b$ は $R'$ の非零因子である。
$\mathfrak p' \subset R' = R[\frac{I}{a}]$ を素イデアルとする。
このとき $JR'_{\mathfrak p'}$ は $b$ で生成される。Lemma
\ref{more-algebra-lemma-principal-fitting-ideal} により、
$M'_{\mathfrak p'}$ は $r$ 個の元で生成できる。$M'$ は有限なので、
対応する写像 $(R')^{\oplus r} \to M'$ が $g$ を逆にした後に全射となる
$m_1, \ldots, m_r \in M'$ と
$g \in R'$, $g \not \in \mathfrak p'$ が存在する。

\medskip\noindent
最後にイデアル $J' = \text{Fit}_{k - 1}(M')$ を考える。
% SOURCE-ERRATUM-CANDIDATE P02-E0828: frozen source uses undefined k in Fit_{k-1}; the module rank in this lemma is r and the exact source formula is preserved.
$J' R'_g$ は $m_1, \ldots, m_r$ の間の関係式の係数で生成されることに
注意せよ（Fitting イデアルと基底変換との両立性）。従って
Lemma \ref{more-algebra-lemma-fitting-ideal-finite-locally-free} により
$J' = 0$ を示せば十分である。Algebra, Lemma
\ref{algebra-lemma-blowup-in-principal} により $R'_a = R_f$ であり、
% SOURCE-ERRATUM-CANDIDATE P02-E0826: second frozen occurrence of the same localization mismatch.
$M'_a = M_f$ は階数 $r$ の自由加群なので、$J'_a = 0$ である。
$a$ は $R'$ の非零因子なので、$J' = 0$ と結論でき、証明が完了する。
\end{proof}

\section{完備化と平坦性}
\label{more-algebra-section-completion-flat}

\noindent
この節では「大きな」平坦加群の完備化がいつ平坦になるかを論じる。

\begin{lemma}
\label{more-algebra-lemma-ui-completion-direct-sum-into-product}
$R$ を環、$I \subset R$ をイデアル、$A$ を集合とする。
$R$ は Noether で $I$ に関して完備であると仮定する。直和の $I$-進完備化から
直積への標準写像
$$
\left(\bigoplus\nolimits_{\alpha \in A} R\right)^\wedge
\longrightarrow
\prod\nolimits_{\alpha \in A} R
$$
が存在し、これは普遍単射である。
\end{lemma}

\begin{proof}
定義により、左辺の元 $x$ は $x = (x_n)$ という形であり、ここで
$x_n = (x_{n, \alpha}) \in \bigoplus\nolimits_{\alpha \in A} R/I^n$ かつ
$x_{n, \alpha} = x_{n + 1, \alpha} \bmod I^n$ である。
$R = R^\wedge$ なので、任意の $\alpha$ に対して
$x_{n, \alpha} = y_\alpha \bmod I^n$ を満たす $y_\alpha \in R$ が存在する。
各 $n$ に対して、元 $x_{n, \alpha}$ が零でない $\alpha$ は有限個しかない
ことに注意せよ。逆に $(y_\alpha) \in \prod_\alpha R$ が与えられ、各 $n$ に
対して $y_{\alpha} \bmod I^n$ が零でない $\alpha$ が有限個しかなければ、
これは左辺の元を定める。従って左辺の元は、$y_\alpha \in R$ で、各 $n$ に
対して $I^n$ を法として零でない $y_\alpha$ が有限個しかないような無限の
「収束和」$\sum_\alpha y_\alpha$ と考えられる。表示された写像はこの元を
直積の元 $(y_\alpha)$ に写す。
% SOURCE-ERRATUM-CANDIDATE P02-E0829: frozen source says “maps this element to the element to (y_alpha),” with a duplicated “to”; Japanese supplies one destination particle.
特にこの写像は単射である。

\medskip\noindent
$Q$ を有限 $R$-加群とする。Algebra, Theorem
\ref{algebra-theorem-universally-exact-criteria} により、写像
$$
Q \otimes_R \left(\bigoplus\nolimits_{\alpha \in A} R\right)^\wedge
\longrightarrow
Q \otimes_R \left(\prod\nolimits_{\alpha \in A} R\right)
$$
が単射であることを示さなければならない。表示
$R^{\oplus k} \to R^{\oplus m} \to Q \to 0$ を選び、対応する $Q$ の生成元を
$q_1, \ldots, q_m \in Q$ と表す。Artin--Rees
（Algebra, Lemma \ref{algebra-lemma-Artin-Rees}）により、定数 $c$ で
$$
\Im(R^{\oplus k} \to R^{\oplus m}) \cap (I^N)^{\oplus m}
\subset \Im((I^{N - c})^{\oplus k} \to R^{\oplus m})
$$
を満たすものが存在する。図式
$$
\xymatrix{
\bigoplus_{l = 1}^k \left(\bigoplus\nolimits_{\alpha \in A} R\right)^\wedge
\ar[r] \ar[d] &
\bigoplus_{j = 1}^m \left(\bigoplus\nolimits_{\alpha \in A} R\right)^\wedge
\ar[r] \ar[d] &
Q \otimes_R \left(\bigoplus\nolimits_{\alpha \in A} R\right)^\wedge
\ar[r] \ar[d] &
0 \\
\bigoplus_{l = 1}^k \left(\prod\nolimits_{\alpha \in A} R\right)
\ar[r] &
\bigoplus_{j = 1}^m \left(\prod\nolimits_{\alpha \in A} R\right)
\ar[r] &
Q \otimes_R \left(\prod\nolimits_{\alpha \in A} R\right)
\ar[r] &
0
}
$$
を考える。その行は完全である。
$\bigoplus_{j = 1, \ldots, m}
\left(\bigoplus\nolimits_{\alpha \in A} R\right)^\wedge$ の元
$\sum_j \sum_\alpha y_{j, \alpha}$ を選ぶ。この元が
$Q \otimes_R \left(\prod\nolimits_{\alpha \in A} R\right)$ で零に写るならば、
特に各 $\alpha$ に対して
$\sum_j q_j \otimes y_{j, \alpha} = 0$ が $Q$ で成り立つ。
従って
$(y_{1, \alpha}, \ldots, y_{m, \alpha}) \in
\bigoplus_{j = 1, \ldots, m} R$ に写る元
$(z_{1, \alpha}, \ldots, z_{k, \alpha}) \in
\bigoplus_{l = 1, \ldots, k} R$ を見つけられる。さらに
$j = 1, \ldots, m$ に対して $y_{j, \alpha} \in I^{N_\alpha}$ ならば、
$l = 1, \ldots, k$ に対して
$z_{l, \alpha} \in I^{N_\alpha - c}$ と仮定してよい。従って和
$\sum_l \sum_\alpha z_{l, \alpha}$ は「収束」し、
$$
\bigoplus_{l = 1, \ldots, k}
\left(\bigoplus\nolimits_{\alpha \in A} R\right)^\wedge
$$
の元を定め、最初の元 $\sum_j \sum_\alpha y_{j, \alpha}$ に写る。
従って右の垂直矢印は単射であり、証明が完了する。
\end{proof}

\noindent
次の補題は下の Lemma \ref{more-algebra-lemma-limit-flat} から導くこともできる。

\begin{lemma}
\label{more-algebra-lemma-completed-direct-sum-flat}
$R$ を環、$I \subset R$ をイデアル、$A$ を集合とする。
$R$ は Noether であると仮定する。完備化
$(\bigoplus\nolimits_{\alpha \in A} R)^\wedge$ は平坦 $R$-加群である。
\end{lemma}

\begin{proof}
$R^\wedge$ を $I$ に関する $R$ の完備化とする。Algebra, Lemma
\ref{algebra-lemma-completion-flat} により $R \to R^\wedge$ は平坦なので、
$(\bigoplus\nolimits_{\alpha \in A} R)^\wedge$ が平坦
$R^\wedge$-加群であることを示せば十分である（Algebra, Lemma
\ref{algebra-lemma-composition-flat} を用いよ）。
$$
(\bigoplus\nolimits_{\alpha \in A} R)^\wedge
=
(\bigoplus\nolimits_{\alpha \in A} R^\wedge)^\wedge
$$
なので、$R$ を $R^\wedge$ で置き換え、$R$ は $I$ に関して完備であると
仮定してよい（Algebra, Lemma
\ref{algebra-lemma-completion-complete} を参照せよ）。この場合 Lemma
\ref{more-algebra-lemma-ui-completion-direct-sum-into-product} により写像
$(\bigoplus\nolimits_{\alpha \in A} R)^\wedge \to \prod_{\alpha \in A} R$
は普遍単射である。従って Algebra, Lemma
\ref{algebra-lemma-ui-flat-domain} により、
$\prod_{\alpha \in A} R$ が平坦であることを示せば十分である。
Algebra, Proposition \ref{algebra-proposition-characterize-coherent}
および Algebra, Lemma \ref{algebra-lemma-Noetherian-coherent} により
$\prod_{\alpha \in A} R$ は平坦である。
\end{proof}

\begin{lemma}
\label{more-algebra-lemma-tor-strictly-pro-zero}
\begin{reference}
これは \cite[Lemma 9.9]{quillenhomology} である。ただし、明らかに意図されて
いたにもかかわらず、著者は主張で「strict」という語を落としている。
\end{reference}
$A$ を Noether 環、$I$ を $A$ のイデアル、$M$ を有限 $A$-加群とする。
任意の $p > 0$ に対して $c > 0$ が存在し、すべての $n \geq c$ に対して
写像 $\text{Tor}_p^A(M, A/I^n) \to \text{Tor}_p^A(M, A/I^{n - c})$
は零である。
\end{lemma}

\begin{proof}
$p = 1$ の場合を証明する。短完全列
$0 \to K \to A^{\oplus t} \to M \to 0$ を選ぶ。このとき
$\text{Tor}_1^A(M, A/I^n) = K \cap (I^n)^{\oplus t}/I^nK$ である。
Artin--Rees（Algebra, Lemma \ref{algebra-lemma-Artin-Rees}）により、
$n \geq c$ に対して
$K \cap (I^n)^{\oplus t} \subset I^{n - c}K$ となる定数
$c \geq 0$ が存在する。従って $p = 1$ の場合の結果を得る。
% SOURCE-ERRATUM-CANDIDATE P02-E0830: frozen source has the sentence fragment “Thus the result for p = 1”; Japanese supplies the predicate.
$p > 1$ に対して
$\text{Tor}_p^A(M, A/I^n) = \text{Tor}^A_{p - 1}(K, A/I^n)$ である。
従って補題は帰納法から従う。
\end{proof}

\begin{lemma}
\label{more-algebra-lemma-limit-flat}
$A$ を Noether 環、$I$ を $A$ のイデアルとする。$(M_n)$ を次を満たす
$A$-加群の逆系とする。
\begin{enumerate}
\item $M_n$ は平坦 $A/I^n$-加群である、
\item $M_{n + 1} \to M_n$ は全射である。
\end{enumerate}
このとき $M = \lim M_n$ は平坦 $A$-加群であり、任意の有限
$A$-加群 $Q$ に対して
$Q \otimes_A M = \lim Q \otimes_A M_n$ である。
\end{lemma}

\begin{proof}
まず、任意の有限 $A$-加群 $Q$ に対して
$Q \otimes_A M = \lim Q \otimes_A M_n$ であることを示す。
有限自由 $A$-加群 $F_i$ による分解
$F_2 \to F_1 \to F_0 \to Q \to 0$ を選ぶ。このとき
$$
F_2 \otimes_A M_n \to F_1 \otimes_A M_n \to F_0 \otimes_A M_n
$$
は鎖複体で、その次数 $0$ のホモロジーは $Q \otimes_A M_n$、次数 $1$ の
ホモロジーは
$$
\text{Tor}_1^A(Q, M_n) = \text{Tor}_1^A(Q, A/I^n) \otimes_{A/I^n} M_n
$$
である。ここで $M_n$ は $A/I^n$ 上平坦である。Lemma
\ref{more-algebra-lemma-tor-strictly-pro-zero} により、この系は本質的に値 $0$ の定数系である。
Homology, Lemma \ref{homology-lemma-apply-Mittag-Leffler-again} から
$\lim Q \otimes_A M_n =
\Coker(\lim F_1 \otimes_A M_n \to \lim F_0 \otimes_A M_n)$ が従う。
$F_i$ は有限自由なので、これは
$\Coker(F_1 \otimes_A M \to F_0 \otimes_A M) = Q \otimes_A M$ に等しい。

\medskip\noindent
次に $Q \to Q'$ を有限 $A$-加群の単射とする。Algebra, Lemma
\ref{algebra-lemma-flat} により、
$Q \otimes_A M \to Q' \otimes_A M$ が単射であることを示さなければならない。
上の結果から
$$
\Ker(Q \otimes_A M \to Q' \otimes_A M) =
\Ker(\lim Q \otimes_A M_n \to \lim Q' \otimes_A M_n).
$$
である。各 $n$ に対して完全列
$$
\text{Tor}_1^A(Q', M_n) \to \text{Tor}_1^A(Q'', M_n) \to
Q \otimes_A M_n \to Q' \otimes_A M_n
$$
がある。ここで $Q'' = \Coker(Q \to Q')$ である。上で Tor の逆系が
本質的に値 $0$ の定数系であることを見た。Homology, Lemma
\ref{homology-lemma-apply-Mittag-Leffler-again} により、最も右の写像の
逆極限は単射である。
\end{proof}

\begin{lemma}
\label{more-algebra-lemma-flat-after-completion}
$R$ を環、$I \subset R$ をイデアル、$M$ を $R$-加群とする。次を仮定する。
\begin{enumerate}
\item $I$ は有限生成である、
\item $R/I$ は Noether である、
\item $M/IM$ は $R/I$ 上平坦である、
\item $\text{Tor}_1^R(M, R/I) = 0$ である。
\end{enumerate}
このとき $I$-進完備化 $R^\wedge$ は Noether 環であり、
$M^\wedge$ は $R^\wedge$ 上平坦である。
\end{lemma}

\begin{proof}
Algebra, Lemma \ref{algebra-lemma-what-does-it-mean} により、すべての
$n$ に対して加群 $M/I^nM$ は $R/I^n$ 上平坦である。
Algebra, Lemma \ref{algebra-lemma-hathat-finitely-generated} により、
(a) $R^\wedge$ と $M^\wedge$ は $I$-進完備であり、(b) すべての $n$ に対して
$R/I^n = R^\wedge/I^nR^\wedge$ である。Algebra, Lemma
\ref{algebra-lemma-completion-Noetherian} により環 $R^\wedge$ は Noether である。
Lemma \ref{more-algebra-lemma-limit-flat} を適用すると、
$M^\wedge = \lim M/I^nM$ は $R^\wedge$-加群として平坦であると結論する。
\end{proof}

\section{Koszul 複体}
\label{more-algebra-section-koszul}

\noindent
Koszul 複体を次のように定義する。

\begin{definition}
\label{more-algebra-definition-koszul}
$R$ を環、$\varphi : E \to R$ を $R$-加群写像とする。
$\varphi$ に付随する {\it Koszul 複体} $K_\bullet(\varphi)$ とは、
次のように定義される可換微分次数付き代数である。
\begin{enumerate}
\item その台となる次数付き代数は外積代数
$K_\bullet(\varphi) = \wedge(E)$ である、
\item 微分 $d : K_\bullet(\varphi) \to K_\bullet(\varphi)$ は、
すべての $e \in E = K_1(\varphi)$ に対して $d(e) = \varphi(e)$ を満たす
一意な導分である。
\end{enumerate}
\end{definition}

\noindent
明示的には、$e_1 \wedge \ldots \wedge e_n$ が $K_\bullet(\varphi)$ の
次数 $n$ の生成元の一つならば、
$$
d(e_1 \wedge \ldots \wedge e_n) =
\sum\nolimits_{i = 1, \ldots, n} (-1)^{i + 1}
\varphi(e_i)e_1 \wedge \ldots \wedge \widehat{e_i} \wedge \ldots \wedge e_n.
$$
である。これがテンソル代数上の良定義な導分を与え、
$e \otimes e$ を零化するので外積代数を経由することは直ちに分かる。

\medskip\noindent
しばしば $E$ は有限自由加群、例えば $E = R^{\oplus n}$ であると仮定する。
この場合、写像 $\varphi$ は元の列 $f_1, \ldots, f_n \in R$ によって与えられる。

\begin{definition}
\label{more-algebra-definition-koszul-complex}
$R$ を環とし、$f_1, \ldots, f_r \in R$ とする。
{\it $f_1, \ldots, f_r$ 上の Koszul 複体}とは、写像
$(f_1, \ldots, f_r) : R^{\oplus r} \to R$ に付随する Koszul 複体である。
記号は $K_\bullet(f_\bullet)$, $K_\bullet(f_1, \ldots, f_r)$,
$K_\bullet(R, f_1, \ldots, f_r)$, または $K_\bullet(R, f_\bullet)$ とする。
\end{definition}

\noindent
もちろん $E$ が有限局所自由ならば、$K_\bullet(\varphi)$ は
$\Spec(R)$ 上局所的に Koszul 複体
$K_\bullet(f_1, \ldots, f_r)$ と同型である。この複体は多くの興味深い
形式的性質をもつ。

\begin{lemma}
\label{more-algebra-lemma-functorial}
$\varphi : E \to R$ と $\varphi' : E' \to R$ を $R$-加群写像とする。
$\psi : E \to E'$ を $\varphi' \circ \psi = \varphi$ を満たす
$R$-加群写像とする。このとき $\psi$ は微分次数付き代数の準同型
$K_\bullet(\varphi) \to K_\bullet(\varphi')$ を誘導する。
\end{lemma}

\begin{proof}
これは定義から直ちに従う。
\end{proof}

\begin{lemma}
\label{more-algebra-lemma-change-basis}
$f_1, \ldots, f_r \in R$ を列とし、$(x_{ij})$ を $R$ に係数をもつ
可逆な $r \times r$-行列とする。このとき複体 $K_\bullet(f_\bullet)$ と
$$
K_\bullet(\sum x_{1j}f_j, \sum x_{2j}f_j, \ldots, \sum x_{rj}f_j)
$$
は同型である。
\end{lemma}

\begin{proof}
$g_i = \sum x_{ij}f_j$ とおく。行列 $(x_{ji})$ は同型
$x : R^{\oplus r} \to R^{\oplus r}$ を与え、
$(g_1, \ldots, g_r) = (f_1, \ldots, f_r) \circ x$ を満たす。
従って Lemma \ref{more-algebra-lemma-functorial} で述べた Koszul 複体の関手性から従う。
\end{proof}

\begin{lemma}
\label{more-algebra-lemma-homotopy-koszul-abstract}
$R$ を環、$\varphi : E \to R$ を $R$-加群写像とする。
$e \in E$ の $R$ における像を $f = \varphi(e)$ とする。このとき
$$
f = de + ed
$$
が $K_\bullet(\varphi)$ の自己準同型として成り立つ。
\end{lemma}

\begin{proof}
$d(ea) = d(e)a - ed(a) = fa - ed(a)$ なので成り立つ。
\end{proof}

\begin{lemma}
\label{more-algebra-lemma-homotopy-koszul}
$R$ を環とし、$f_1, \ldots, f_r \in R$ を列とする。
$K_\bullet(f_\bullet)$ 上の $f_i$ 倍写像は零写像にホモトピックであり、
特にコホモロジー加群 $H_i(K_\bullet(f_\bullet))$ はイデアル
$(f_1, \ldots, f_r)$ によって零化される。
% SOURCE-ERRATUM-CANDIDATE P02-E0831: frozen source calls H_i of the chain Koszul complex “cohomology modules”; the homological indexing and context indicate “homology modules,” but Japanese preserves the printed claim.
\end{lemma}

\begin{proof}
Lemma \ref{more-algebra-lemma-homotopy-koszul-abstract} の特別な場合である。
\end{proof}

\noindent
Derived Categories, Section \ref{derived-section-cones} で余鎖複体の射の錐を
定義した。鎖複体の射 $f : A_\bullet \to B_\bullet$ の錐
$C(f)_\bullet$ は複体 $C(f)_\bullet$ であり、
$C(f)_n = B_n \oplus A_{n - 1}$ で微分が
\begin{equation}
\label{more-algebra-equation-differential-cone}
d_{C(f), n} =
\left(
\begin{matrix}
d_{B, n} & f_{n - 1} \\
0 & -d_{A, n - 1}
\end{matrix}
\right)
\end{equation}
によって与えられる複体である。これには、明らかな写像
$B_n \to C(f)_n \to A_{n - 1}$ が誘導する複体の標準射
$i : B_\bullet \to C(f)_\bullet$ と
$p : C(f)_\bullet \to A_\bullet[-1]$ が備わる。

\begin{lemma}
\label{more-algebra-lemma-cone-koszul-abstract}
$R$ を環、$\varphi : E \to R$ を $R$-加群写像とする。$f \in R$ とする。
$E' = E \oplus R$ とおき、$\varphi' : E' \to R$ を $E$ 上では
$\varphi$、$R$ 上では $f$ 倍写像として定める。複体
$K_\bullet(\varphi')$ は複体の写像
$$
f :
K_\bullet(\varphi)
\longrightarrow
K_\bullet(\varphi).
$$
の錐と同型である。
\end{lemma}

\begin{proof}
$e_0 \in E'$ を $1 \in R \subset R \oplus E$ という元と表す。
% SOURCE-ERRATUM-CANDIDATE P02-E0832: frozen source reverses the declared direct-sum order E'=E⊕R when it writes R⊂R⊕E; Japanese preserves the exact inclusion span.
上の錐の定義により
$$
C(f)_n = K_n(\varphi) \oplus K_{n - 1}(\varphi) =
\wedge^n(E) \oplus \wedge^{n - 1}(E) = \wedge^n(E')
$$
である。最後の $=$ において
$(0, e_1 \wedge \ldots \wedge e_{n - 1})$ を
$\wedge^n(E')$ の
$e_0 \wedge e_1 \wedge \ldots \wedge e_{n - 1}$ に写す。
計算により、この同型が微分と両立することが分かる。実際、第一直和因子の
元については、$\varphi'|_E = \varphi$ であり、第一直和因子に制限した
$d_{C(f)}$ はちょうど $d_{K_\bullet(\varphi)}$ なので明らかである。
一方、$e_1 \wedge \ldots \wedge e_{n - 1}$ が第二直和因子に属するならば、
$$
d_{C(f)}(0, e_1 \wedge \ldots \wedge e_{n - 1}) =
fe_1 \wedge \ldots \wedge e_{n - 1}
- d_{K_\bullet(\varphi)}(e_1 \wedge \ldots \wedge e_{n - 1})
$$
であり、他方
% SOURCE-ERRATUM-CANDIDATE P02-E0833: frozen source applies d_{K(phi')} to the pair-like argument (0,e0∧...), although K(phi') contains the wedge element rather than a cone pair; the formula is preserved exactly.
\begin{align*}
& d_{K_\bullet(\varphi')}(0, e_0 \wedge e_1 \wedge \ldots \wedge e_{n - 1}) \\
& =
\sum\nolimits_{i = 0, \ldots, n - 1}
(-1)^i \varphi'(e_i)e_0 \wedge \ldots \wedge \widehat{e_i}
\wedge \ldots \wedge e_{n - 1} \\
& =
fe_1 \wedge \ldots \wedge e_{n - 1} +
\sum\nolimits_{i = 1, \ldots, n - 1}
(-1)^i \varphi(e_i)e_0 \wedge \ldots \wedge \widehat{e_i}
\wedge \ldots \wedge e_{n - 1} \\
& =
fe_1 \wedge \ldots \wedge e_{n - 1} -
e_0 \left(\sum\nolimits_{i = 1, \ldots, n - 1}
(-1)^{i + 1} \varphi(e_i)e_1 \wedge \ldots \wedge \widehat{e_i}
\wedge \ldots \wedge e_{n - 1}\right)
\end{align*}
であり、これは先の計算結果の像である。
\end{proof}

\begin{lemma}
\label{more-algebra-lemma-cone-koszul}
$R$ を環とし、$f_1, \ldots, f_r$ を $R$ の元の列とする。複体
$K_\bullet(f_1, \ldots, f_r)$ は複体の写像
$$
f_r :
K_\bullet(f_1, \ldots, f_{r - 1})
\longrightarrow
K_\bullet(f_1, \ldots, f_{r - 1}).
$$
の錐と同型である。
\end{lemma}

\begin{proof}
Lemma \ref{more-algebra-lemma-cone-koszul-abstract} の特別な場合である。
\end{proof}

\begin{lemma}
\label{more-algebra-lemma-cone-squared}
$R$ を環、$A_\bullet$ を $R$-加群の複体、$f, g \in R$ とする。
$C(f)_\bullet$ を $f : A_\bullet \to A_\bullet$ の錐とし、
$C(g)_\bullet$ と $C(fg)_\bullet$ も同様に定義する。このとき
$C(fg)_\bullet$ は写像
$$
C(f)_\bullet[1] \longrightarrow C(g)_\bullet
$$
の錐とホモトピー同値である。
\end{lemma}

\begin{proof}
まず $A_\bullet$ が次数 $0$ に置かれた $R$ からなる複体である場合を証明する。
この場合、複体 $C(f)_\bullet$ は
$$
\ldots \to 0 \to R \xrightarrow{f} R \to 0 \to \ldots
$$
であり、$R$ は（ホモロジー）次数 $1$ と $0$ に置かれる。用いる複体の写像は
$$
\xymatrix{
0 \ar[r] \ar[d] &
0 \ar[r] \ar[d] &
R \ar[r]^f \ar[d]^1 &
R \ar[r] \ar[d] & 0 \ar[d] \\
0 \ar[r] & R \ar[r]^g & R \ar[r] & 0 \ar[r] & 0
}
$$
である。この錐は次数 $1$ と $0$ に置かれた $R^{\oplus 2}$ からなる鎖複体で、
その微分（\ref{more-algebra-equation-differential-cone}）は
$$
\left(
\begin{matrix}
g & 1 \\
0 & -f
\end{matrix}
\right) :
R^{\oplus 2} \longrightarrow R^{\oplus 2}
$$
である。この鎖複体が $C(fg)_\bullet$、すなわち
$R \xrightarrow{fg} R$ とホモトピー同値であることを見るため、複体の写像
$$
\xymatrix{
R \ar[d]_{(1, -g)} \ar[r]_{fg} & R \ar[d]^{(0, 1)} \\
R^{\oplus 2} \ar[r] & R^{\oplus 2}
}
\quad\quad
\xymatrix{
R^{\oplus 2} \ar[d]_{(1, 0)} \ar[r] & R^{\oplus 2} \ar[d]^{(f, 1)} \\
R \ar[r]^{fg} & R
}
$$
を考える。記法の意味は明らかである。一方の向きの二写像の合成は
$C(fg)_\bullet$ 上の恒等写像だが、他方の向きでは恒等写像でない。
必要なホモトピーを書き下すことは省略する。

\medskip\noindent
一般の場合には、複体の圏上の関手
$K_\bullet \mapsto \text{Tot}(A_\bullet \otimes_R K_\bullet)$ が
ホモトピーおよび錐と両立することを用いる。次に
$C(f)_\bullet$ と $C(g)_\bullet$ を二重複体
$$
(R \xrightarrow{f} R) \otimes_R A_\bullet
\quad\text{and}\quad
(R \xrightarrow{g} R) \otimes_R A_\bullet
$$
の全複体として書く。このようにして上で論じた特別な場合から結果を導く。
細部をいくつか省略する。
\end{proof}

\begin{lemma}
\label{more-algebra-lemma-koszul-mult-abstract}
$R$ を環、$\varphi : E \to R$ を $R$-加群写像、$f, g \in R$ とする。
$E' = E \oplus R$ とおき、$\varphi'_f, \varphi'_g, \varphi'_{fg} : E' \to R$
を $E$ 上では $\varphi$、$R$ 上ではそれぞれ $f, g, fg$ 倍写像として定める。
複体 $K_\bullet(\varphi'_{fg})$ は複体のある写像の錐とホモトピー同値である：
$$
K_\bullet(\varphi'_f)[1]
\longrightarrow
K_\bullet(\varphi'_g).
$$
\end{lemma}

\begin{proof}
Lemma \ref{more-algebra-lemma-cone-koszul-abstract} により、
複体 $K_\bullet(\varphi'_f)$ は $K_\bullet(\varphi)$ 上の $f$ 倍写像の錐と
同型であり、他の二つの場合も同様である。従って補題は
Lemma \ref{more-algebra-lemma-cone-squared} から従う。
\end{proof}

\begin{lemma}
\label{more-algebra-lemma-koszul-mult}
$R$ を環、$f_1, \ldots, f_{r - 1}$ を $R$ の元の列とし、
$f, g \in R$ とする。複体
$K_\bullet(f_1, \ldots, f_{r - 1}, fg)$ は複体のある写像の錐と
ホモトピー同値である：
$$
K_\bullet(f_1, \ldots, f_{r - 1}, f)[1]
\longrightarrow
K_\bullet(f_1, \ldots, f_{r - 1}, g)
$$
\end{lemma}

\begin{proof}
Lemma \ref{more-algebra-lemma-koszul-mult-abstract} の特別な場合である。
\end{proof}

\begin{lemma}
\label{more-algebra-lemma-join-sequences-koszul-complex}
$R$ を環とし、$f_1, \ldots, f_r$, $g_1, \ldots, g_s$ を $R$ の元とする。
このとき Koszul 複体の同型
$$
K_\bullet(R, f_1, \ldots, f_r, g_1, \ldots, g_s) =
\text{Tot}(K_\bullet(R, f_1, \ldots, f_r) \otimes_R
K_\bullet(R, g_1, \ldots, g_s)).
$$
が存在する。
\end{lemma}

\begin{proof}
省略する。ヒント：$K_\bullet(R, f_1, \ldots, f_r)$ が微分次数付き代数として
$\text{d}(x_i) = f_i$ を満たす $x_1, \ldots, x_r$ によって生成され、
$K_\bullet(R, g_1, \ldots, g_s)$ が微分次数付き代数として
$\text{d}(y_j) = g_j$ を満たす $y_1, \ldots, y_s$ によって生成されるならば、
$K_\bullet(R, f_1, \ldots, f_r, g_1, \ldots, g_s)$ は
$\text{d}(x_i) = f_i$ と $\text{d}(y_j) = g_j$ を満たす元の列
$x_1, \ldots, x_r, y_1, \ldots, y_s$ によって生成される微分次数付き代数と
考えられる。
\end{proof}

\section{拡張交代 {\v C}ech 複体}
\label{more-algebra-section-alternating-cech}

\noindent
$R$ を環とし、$f_1, \ldots, f_r \in R$ とする。$R$ の拡張交代
{\v C}ech 複体とは余鎖複体
$$
R \to \bigoplus\nolimits_{i_0} R_{f_{i_0}} \to
\bigoplus\nolimits_{i_0 < i_1} R_{f_{i_0}f_{i_1}} \to
\ldots \to R_{f_1\ldots f_r}
$$
のことである。ここで $R$ は次数 $0$ にあり、項
$\bigoplus_{i_0} R_{f_{i_0}}$ は次数 $1$ にあり、以下同様である。
% SOURCE-ERRATUM-CANDIDATE P02-E0834: frozen source misspells “degree” as “degre”; Japanese supplies the standard term.
写像は次のように定義する。
\begin{enumerate}
\item 写像 $R \to \bigoplus\nolimits_{i_0} R_{f_{i_0}}$ は標準写像
$R \to R_{f_{i_0}}$ によって与えられる。
\item $1 \leq i_0 < \ldots < i_{p + 1} \leq r$ および
$0 \leq j \leq p + 1$ が与えられれば、標準局所化写像
$$
R_{f_{i_0} \ldots \hat f_{i_j} \ldots f_{i_{p + 1}}} \to
R_{f_{i_0} \ldots f_{i_{p + 1}}}
$$
がある。
\item 微分は (2) の標準写像を符号 $(-1)^j$ とともに用いる。
\end{enumerate}
$M$ を任意の $R$-加群とする。$M$ の拡張交代 {\v C}ech 複体とは、
同様に構成される余鎖複体
$$
M \to \bigoplus\nolimits_{i_0} M_{f_{i_0}} \to
\bigoplus\nolimits_{i_0 < i_1} M_{f_{i_0}f_{i_1}} \to
\ldots \to M_{f_1\ldots f_r}
$$
のことであり、前と同様に $M$ は次数 $0$ にある。

\begin{lemma}
\label{more-algebra-lemma-extended-alternating-is-complex}
上で定義した拡張交代 {\v C}ech 複体は $R$-加群の複体である。
\end{lemma}

\begin{proof}
省略する。
\end{proof}

\begin{lemma}
\label{more-algebra-lemma-extended-alternating-form-module}
$R$ を環、$f_1, \ldots, f_r \in R$、$M$ を $R$-加群とする。
$M$ の拡張交代 {\v C}ech 複体は、$M$ と $R$ の拡張交代
{\v C}ech 複体との $R$ 上のテンソル積である。
\end{lemma}

\begin{proof}
省略する。
\end{proof}

\begin{lemma}
\label{more-algebra-lemma-extended-alternating-base-change}
$R$ を環、$f_1, \ldots, f_r \in R$、$M$ を $R$-加群とする。
$R \to S$ を環写像とし、$f_1, \ldots, f_r$ の像を
$g_1, \ldots, g_r \in S$ と表し、$N = M \otimes_R S$ とおく。
$S$, $g_1, \ldots, g_r$, $N$ を用いて構成した拡張交代
{\v C}ech 複体は、$M$ の拡張交代 {\v C}ech 複体と $S$ との
$R$ 上のテンソル積である。
\end{lemma}

\begin{proof}
省略する。
\end{proof}

\begin{lemma}
\label{more-algebra-lemma-extended-alternating-homotopy-zero}
$R$ を環、$f_1, \ldots, f_r \in R$、$M$ を $R$-加群とする。
$f_i$ が単元となる $i \in \{1, \ldots, r\}$ が存在するならば、
$M$ の拡張交代 {\v C}ech 複体は $0$ とホモトピー同値である。
\end{lemma}

\begin{proof}
次の記法を用いる。$M$ の拡張交代 {\v C}ech 複体の次数 $p + 1$ の余鎖
$x$ を $x = (x_{i_0 \ldots i_p})$ と書く。ここで
$x_{i_0 \ldots i_p}$ は $M_{f_{i_0} \ldots f_{i_p}}$ に属する。
この記法では
$$
d(x)_{i_0 \ldots i_{p + 1}} =
\sum\nolimits_j (-1)^j x_{i_0 \ldots \hat i_j \ldots i_{p + 1}}
$$
である。ホモトピーとして、規則
$$
h : \text{cochains of degree }p + 2 \to \text{cochains of degree }p + 1
$$
および
$$
h(x)_{i_0 \ldots i_p} = 0 \text{ if } i \in \{i_0, \ldots, i_p\}
\text{ and }
h(x)_{i_0 \ldots i_p} = 
(-1)^j x_{i_0 \ldots i_j i i_{j + 1} \ldots i_p} \text{ if not}
$$
で与えられる写像を用いる。第二の場合、$j$ は
$i_j < i < i_{j + 1}$ を満たす一意な添字である。また $f_i$ は単元なので
$$
M_{f_{i_0} \ldots f_{i_p}} =
M_{f_{i_0} \ldots f_{i_j} f_i f_{i_{j + 1}} \ldots f_{i_p}}
$$
であり、この等式により
$(-1)^j x_{i_0 \ldots i_j i i_{j + 1} \ldots i_p}$ を
$M_{f_{i_0} \ldots f_{i_p}}$ の元と考えられる。
% SOURCE-ERRATUM-CANDIDATE P02-E0835: frozen source calls j unique without specifying endpoint conventions when i<i_0 or i>i_p; Japanese preserves the printed rule.
計算により $d h + h d$ が恒等写像の負に等しいことを示せば証明が完了する。
そのため次数 $p + 1$ の余鎖 $x$ を固定し、
$1 \leq i_0 < \ldots < i_p \leq r$ とする。

\medskip\noindent
場合 I：$i \in \{i_0, \ldots, i_p\}$。$i = i_t$ とする。このとき
$h(d(x))_{i_0 \ldots i_p} = 0$ である。他方、
$$
d(h(x))_{i_0 \ldots i_p} =
\sum (-1)^j h(x)_{i_0 \ldots \hat i_j \ldots i_p} =
(-1)^t h(x)_{i_0 \ldots \hat i \ldots i_p} =
(-1)^t (-1)^{t - 1} x_{i_0 \ldots i_p}
$$
である。従って所望のとおり
$(dh + hd)(x)_{i_0 \ldots i_p} = -x_{i_0 \ldots i_p}$ である。

\medskip\noindent
場合 II：$i \not \in \{i_0, \ldots, i_p\}$。
$i_j < i < i_{j + 1}$ となる $j$ をとる。このとき
\begin{align*}
h(d(x))_{i_0 \ldots i_p}
& =
(-1)^j d(x)_{i_0 \ldots i_j i i_{j + 1} \ldots i_p} \\
& =
\sum\nolimits_{j' \leq j} (-1)^{j + j'}
x_{i_0 \ldots \hat i_{j'} \ldots i_j i i_{j + 1} \ldots i_p} -
x_{i_0 \ldots i_p} \\
&
+ \sum\nolimits_{j' > j} (-1)^{j + j' + 1}
x_{i_0 \ldots i_j i i_{j + 1} \ldots \hat i_{j'} \ldots i_p}
\end{align*}
である。他方、
\begin{align*}
d(h(x))_{i_0 \ldots i_p}
& =
\sum\nolimits_{j'} (-1)^{j'} h(x)_{i_0 \ldots \hat i_{j'} \ldots i_p} \\
& =
\sum\nolimits_{j' \leq j} (-1)^{j' + j - 1}
x_{i_0 \ldots \hat i_{j'} \ldots i_j i i_{j + 1} \ldots i_p} \\
& +
\sum\nolimits_{j' > j} (-1)^{j' + j}
x_{i_0 \ldots i_j i i_{j + 1} \ldots \hat i_{j'} \ldots i_p}
\end{align*}
である。これらを加えると、所望のとおり
$(dh + hd)(x)_{i_0 \ldots i_p} = - x_{i_0 \ldots i_p}$ を得る。
\end{proof}

\begin{lemma}
\label{more-algebra-lemma-extended-alternating-torsion}
$R$ を環、$f_1, \ldots, f_r \in R$、$M$ を $R$-加群とする。
$H^q$ を $M$ の拡張交代 {\v C}ech 複体の第 $q$ コホモロジー加群とする。
% SOURCE-ERRATUM-CANDIDATE P02-E0836: frozen source says “extended alternation Cech complex” instead of “extended alternating”; Japanese uses the established construction name.
このとき
\begin{enumerate}
\item $q \not \in [0, r]$ ならば $H^q = 0$ である、
\item $x \in H^i$ に対して $n \geq 1$ が存在し、
$i = 1, \ldots, r$ に対して $f_i^n x = 0$ である、
% SOURCE-ERRATUM-CANDIDATE P02-E0837: frozen item reuses i both as the cohomological degree in H^i and as the quantified generator index; H^q is intended, but the exact spans remain.
\item $H^q$ の台は $V(f_1, \ldots, f_r)$ に含まれる、
\item $f \in (f_1, \ldots, f_r)$ で $M$ 上可逆に作用するものが存在すれば、
$H^q = 0$ である。
\end{enumerate}
\end{lemma}

\begin{proof}
(1) は、拡張交代 {\v C}ech 複体が次数 $< 0$ と $> r$ で零であることから従う。
(2) を示すには、各 $i$ に対して $f_i^n x = 0$ となる $n \geq 1$ が存在する
ことを示せば十分である。そのためには $(H^q)_{f_i} = 0$ を示せばよい。
局所化は完全なので、$(H^q)_{f_i}$ は $M$ の拡張交代複体を $f_i$ で
局所化した複体の第 $q$ コホモロジー加群である。Lemma
\ref{more-algebra-lemma-extended-alternating-base-change} により、この局所化は
$R_{f_i}$ 上の $M_{f_i}$ の、$R_{f_i}$ における
$f_1, \ldots, f_r$ の像に関する拡張交代 {\v C}ech 複体である。従って $f_i$ が可逆ならば
$H^q$ は零であることに帰着し、これは Lemma
\ref{more-algebra-lemma-extended-alternating-homotopy-zero} から従う。(3) は、今証明した
すべての $i$ に対する $(H^q)_{f_i} = 0$ から従う。(4) については、この場合
$f$ は $H^q$ 上可逆に作用する一方、(3) により $H^q$ は $V(f)$ 上に
台をもつことに注意せよ。従って $H^q$ は零である（細部を一つ省略する）。
\end{proof}

\begin{lemma}
\label{more-algebra-lemma-extended-alternating-Cech-is-colimit-koszul}
$R$ を環とし、$f_1, \ldots, f_r \in R$ とする。拡張交代
{\v C}ech 複体
$$
R \to \bigoplus\nolimits_{i_0} R_{f_{i_0}} \to
\bigoplus\nolimits_{i_0 < i_1} R_{f_{i_0}f_{i_1}} \to
\ldots \to R_{f_1\ldots f_r}
$$
は Koszul 複体 $K(R, f_1^n, \ldots, f_r^n)$ の余極限である。
正確な主張については証明を参照せよ。
\end{lemma}

\begin{proof}
読者自身でこれを証明することを強く勧める。Definition
\ref{more-algebra-definition-koszul-complex} の Koszul 複体
$K(R, f_1^n, \ldots, f_r^n)$ を次数 $0, \ldots, r$ にある余鎖複体とみなす。
従って
$$
K(R, f_1^n, \ldots, f_r^n) :
0 \to \wedge^r(R^{\oplus r}) \to
\wedge^{r - 1}(R^{\oplus r}) \to \ldots \to
R^{\oplus r} \to R \to 0
$$
であり、項 $\wedge^r(R^{\oplus r})$ は次数 $0$ にある。
$e^n_1, \ldots, e^n_r$ を $R^{\oplus r}$ の標準基底とする。
$1 \leq j_1 < \ldots < j_{r - p} \leq r$ に対する元
$e^n_{j_1} \wedge \ldots \wedge e^n_{j_{r - p}}$ は Koszul 複体の
次数 $p$ の項の基底をなす。さらに Section \ref{more-algebra-section-koszul} における
Koszul 複体の構成により
$$
d(e^n_{j_1} \wedge \ldots \wedge e^n_{j_{r - p}}) =
\sum (-1)^{a + 1} f_{j_a}^n e^n_{j_1} \wedge \ldots \wedge \hat e^n_{j_a} \wedge
\ldots \wedge e^n_{j_{r - p}}
$$
である。系の遷移写像
$$
K(R, f_1^n, \ldots, f_r^n) \to K(R, f_1^{n + 1}, \ldots, f_r^{n + 1})
$$
は規則
$$
e^n_{j_1} \wedge \ldots \wedge e^n_{j_{r - p}}
\longmapsto
f_{i_0} \ldots f_{i_{p - 1}}
e^{n + 1}_{j_1} \wedge \ldots \wedge e^{n + 1}_{j_{r - p}}
$$
によって与えられる。ここで添字
$1 \leq i_0 < \ldots < i_{p - 1} \leq r$ は
$$
\{1, \ldots r\} =
\{i_0, \ldots, i_{p - 1}\} \amalg \{j_1, \ldots, j_{r - p}\}
$$
を満たす。
% SOURCE-ERRATUM-CANDIDATE P02-E0838: frozen source omits the comma in “1, \ldots r” at this locus and again below; exact formal spans are preserved.
これが微分と両立することを示す短い計算は省略する。遷移写像は次数 $0$ では
常に $1$ で、次数 $r$ では $f_1 \ldots f_r$ に等しいことに注意せよ。

\medskip\noindent
$K^p(R, f_1^n, \ldots, f_r^n)$ で Koszul 複体の次数 $p$ の項を表す。
任意の $f \in R$ に対して
$$
R_f = \colim (R \xrightarrow{f} R \xrightarrow{f} R \to \ldots )
$$
であることに注意せよ。従って次数 $p$ では
$$
\colim K^p(R, f_1^n, \ldots f_r^n) =
\bigoplus\nolimits_{1 \leq i_0 < \ldots < i_{p - 1} \leq r}
R_{f_{i_0} \ldots f_{i_{p - 1}}}
$$
を得る。
% SOURCE-ERRATUM-CANDIDATE P02-E0839: frozen source omits the comma after \ldots in the Koszul argument list; the exact formal span is preserved.
上の Koszul 複体の元
$e^n_{j_1} \wedge \ldots \wedge e^n_{j_{r - p}}$ は余極限において、
$$
\{1, \ldots r\} = \{i_0, \ldots, i_{p - 1}\}
\amalg \{j_1, \ldots, j_{r - p}\}
$$
となるように添字を選ぶとき、直和因子
$R_{f_{i_0} \ldots f_{i_{p - 1}}}$ の元
$(f_{i_0} \ldots f_{i_{p - 1}})^{-n}$ に写る。
% SOURCE-ERRATUM-CANDIDATE P02-E0838: second frozen occurrence of “1, \ldots r” without the comma.
従ってこの複体の微分は
$$
d(1\text{ in }R_{f_{i_0} \ldots f_{i_{p - 1}}}) =
\sum\nolimits_{i \not \in \{i_0, \ldots, i_{p - 1}\}}
(-1)^{i - t}\text{ in }
R_{f_{i_0} \ldots f_{i_t} f_i f_{i_{t + 1}} \ldots f_{i_{p - 1}}}
$$
によって与えられる。従って次数 $p$ において写像
$$
\bigoplus\nolimits_{1 \leq i_0 < \ldots < i_{p - 1} \leq r}
R_{f_{i_0} \ldots f_{i_{p - 1}}}
\longrightarrow
\bigoplus\nolimits_{1 \leq i_0 < \ldots < i_{p - 1} \leq r}
R_{f_{i_0} \ldots f_{i_{p - 1}}}
$$
で、規則
$$
1\text{ in }R_{f_{i_0} \ldots f_{i_{p - 1}}}
\longmapsto
(-1)^{i_0 + \ldots + i_{p - 1} + p}\text{ in }R_{f_{i_0} \ldots f_{i_{p - 1}}}
$$
によって定まるものを考えると、
$\colim K(R, f_1^n, \ldots, f_r^n)$ からこの節で定義した拡張交代
{\v C}ech 複体への複体の同型を得る。符号が整合することの確認は省略する。
\end{proof}

\section{Koszul 正則列}
\label{more-algebra-section-koszul-regular}

\noindent
この節を読む前に、Algebra, Sections
\ref{algebra-section-regular-sequences},
\ref{algebra-section-quasi-regular}, および
\ref{algebra-section-depth} を参照されたい。

\begin{definition}
\label{more-algebra-definition-koszul-regular-sequence}
$R$ を環とする。$r \geq 0$ とし、$f_1, \ldots, f_r \in R$ を元の列、
$M$ を $R$-加群とする。列 $f_1, \ldots, f_r$ を次のように呼ぶ。
\begin{enumerate}
\item すべての $i \not = 0$ に対して
$H_i(K_\bullet(f_1, \ldots, f_r) \otimes_R M) = 0$ であるとき、
{\it $M$-Koszul-正則}であるという。
\item $H_1(K_\bullet(f_1, \ldots, f_r) \otimes_R M) = 0$ であるとき、
{\it $M$-$H_1$-正則}であるという。
\item すべての $i \not = 0$ に対して
$H_i(K_\bullet(f_1, \ldots, f_r)) = 0$ であるとき、
{\it Koszul-正則}であるという。
\item $H_1(K_\bullet(f_1, \ldots, f_r)) = 0$ であるとき、
{\it $H_1$-正則}であるという。
\end{enumerate}
\end{definition}

\noindent
環 $R$ の元 $f_1, \ldots, f_r$ について、Lemmas
\ref{more-algebra-lemma-regular-koszul-regular},
\ref{more-algebra-lemma-koszul-regular-H1-regular}, および
\ref{more-algebra-lemma-H1-regular-quasi-regular} により、次の含意が成り立つことが分かる。
\begin{align*}
f_1, \ldots, f_r\text{ is a regular sequence}
& \Rightarrow f_1, \ldots, f_r\text{ is a Koszul-regular sequence} \\
& \Rightarrow f_1, \ldots, f_r\text{ is an }H_1\text{-regular sequence} \\
& \Rightarrow f_1, \ldots, f_r\text{ is a quasi-regular sequence.}
\end{align*}
一般にはこれらの含意のいずれも逆にできない。しかし $R$ が Noether 局所環で
$f_1, \ldots, f_r \in \mathfrak m_R$ ならば、四条件はすべて同値である
（Lemma \ref{more-algebra-lemma-noetherian-finite-all-equivalent}）。
$f = f_1 \in R$ が長さ $1$ の列で、$f$ が $R$ の単元でないとき、
次の条件がすべて同値であることも明らかである。
\begin{enumerate}
\item $f$ は長さ一の正則列である。
\item $f$ は長さ一の Koszul-正則列である。
\item $f$ は長さ一の $H_1$-正則列である。
\end{enumerate}
これらが $f$ は長さ一の準正則列であることを含意するのも明らかである。
しかし、正則列でない長さ $1$ の準正則列は実際に存在する。すなわち
$$
R = k[x, y_0, y_1, \ldots]/(xy_0, xy_1 - y_0, xy_2 - y_1, \ldots)
$$
とし、$f$ を $R$ における $x$ の像とする。このとき $f$ は零因子だが、
$\bigoplus_{n \geq 0} (f^n)/(f^{n + 1}) \cong k[x]$ は多項式環である。

\begin{lemma}
\label{more-algebra-lemma-regular-koszul-regular}
$R$ を環、$M$ を $R$-加群、$f_1, \ldots, f_r \in R$ とし、
$i = 1, \ldots, r$ に対して $f_i$ 倍写像が
$M/(f_1, \ldots, f_{i - 1})M$ 上単射であるとする。このとき
$f_1, \ldots, f_r$ は $M$-Koszul 正則である。特に、$M$-正則列は
$M$-Koszul-正則であり、任意の正則列は Koszul-正則である。
\end{lemma}

\begin{proof}
$R$, $M$, $f_1, \ldots, f_r$ を補題の第一文のとおりとする。
$r = 1$ ならば、$f_1$ が $M$-Koszul-正則であることは直ちに分かる。
$r > 1$ と仮定する。$f_1$ は $M$ 上の非零因子なので、複体の短完全列
$$
0 \to K_\bullet(f_2, \ldots, f_r) \otimes M
\xrightarrow{f_1}
K_\bullet(f_2, \ldots, f_r) \otimes M \to
K_\bullet(\overline{f}_2, \ldots, \overline{f}_r) \otimes M/f_1M \to 0
$$
を得る。ここで $\overline{f}_i$ は $R/(f_1)$ における $f_i$ の像である。
Lemma \ref{more-algebra-lemma-cone-koszul} により、複体
$K_\bullet(f_1, \ldots, f_r)$ は
$K_\bullet(f_2, \ldots, f_r)$ 上の $f_1$ 倍写像の錐と同型である。従って
$K_\bullet(R, f_1, \ldots, f_r) \otimes M$ は第一の写像の錐と同型である。
ゆえに
$K_\bullet(\overline{f}_2, \ldots, \overline{f}_r) \otimes M/f_1M$
は $K_\bullet(f_1, \ldots, f_r) \otimes M$ と擬同型である。
$R/(f_1)$, $M/f_1M$, $\overline{f}_2, \ldots, \overline{f}_r$ は
補題の条件を満たすので、帰納法によりこの複体は正次数で非輪状である。
% SOURCE-ERRATUM-CANDIDATE P02-E0840: frozen source misspells “positive” as “postive”; Japanese supplies the standard term.
これで第一の主張が証明された。第二の主張は第一の主張から直ちに従う。
\end{proof}

\begin{lemma}
\label{more-algebra-lemma-koszul-regular-H1-regular}
$M$-Koszul-正則列は $M$-$H_1$-正則である。
Koszul-正則列は $H_1$-正則である。
\end{lemma}

\begin{proof}
定義から直ちに従う。
\end{proof}

\begin{lemma}
\label{more-algebra-lemma-mult-koszul-regular}
$f_1, \ldots, f_{r - 1} \in R$ を列とし、$f, g \in R$ とする。
$M$ を $R$-加群とする。
\begin{enumerate}
\item $f_1, \ldots, f_{r - 1}, f$ と $f_1, \ldots, f_{r - 1}, g$ が
$M$-$H_1$-正則ならば、$f_1, \ldots, f_{r - 1}, fg$ も
$M$-$H_1$-正則である。
\item $f_1, \ldots, f_{r - 1}, f$ と $f_1, \ldots, f_{r - 1}, g$ が
$M$-Koszul-正則ならば、$f_1, \ldots, f_{r - 1}, fg$ も
$M$-Koszul-正則である。
\end{enumerate}
\end{lemma}

\begin{proof}
Lemma \ref{more-algebra-lemma-koszul-mult} により、すべての $i$ に対して完全列
{\small
$$
H_i(K_\bullet(f_1, \ldots, f_{r - 1}, f) \otimes M) \to
H_i(K_\bullet(f_1, \ldots, f_{r - 1}, fg) \otimes M) \to
H_i(K_\bullet(f_1, \ldots, f_{r - 1}, g) \otimes M)
$$
}
が存在する。
\end{proof}

\begin{lemma}
\label{more-algebra-lemma-koszul-regular-flat-base-change}
$\varphi : R \to S$ を平坦な環写像とし、$f_1, \ldots, f_r \in R$ とする。
$M$ を $R$-加群とし、$N = M \otimes_R S$ とおく。
\begin{enumerate}
\item $R$ における $f_1, \ldots, f_r$ が $M$-$H_1$-正則列ならば、
$\varphi(f_1), \ldots, \varphi(f_r)$ は $S$ における
$N$-$H_1$-正則列である。
\item $R$ における $f_1, \ldots, f_r$ が $M$-Koszul-正則列ならば、
$\varphi(f_1), \ldots, \varphi(f_r)$ は $S$ における
$N$-Koszul-正則列である。
\end{enumerate}
\end{lemma}

\begin{proof}
これは
$$
K_\bullet(f_1, \ldots, f_r) \otimes_R S =
K_\bullet(\varphi(f_1), \ldots, \varphi(f_r))
$$
であり、従って
$$
(K_\bullet(f_1, \ldots, f_r) \otimes_R M) \otimes_R S =
K_\bullet(\varphi(f_1), \ldots, \varphi(f_r)) \otimes_S N
$$
であることから従う。
\end{proof}

\begin{lemma}
\label{more-algebra-lemma-H1-regular-quasi-regular}
$M$-$H_1$-正則列は $M$-準正則である。
\end{lemma}

\begin{proof}
$R$ を環、$M$ を $R$-加群とする。
$f_1, \ldots, f_r$ を $M$-$H_1$-正則列とする。
$J = (f_1, \ldots, f_r)$ とおく。仮定は完全列
$$
\wedge^2(R^r) \otimes M \to R^{\oplus r} \otimes M \to JM \to 0
$$
が存在することを意味する。ここで第一の矢印は
$e_i \wedge e_j \otimes m \mapsto (f_ie_j - f_je_i) \otimes m$
によって与えられる。この列を $R/J$ とテンソルすると、
$$
JM/J^2M = (R/J)^{\oplus r} \otimes_R M = (M/JM)^{\oplus r}
$$
が有限自由加群であることが分かる。
% SOURCE-ERRATUM-CANDIDATE P02-E0841: for arbitrary M, (M/JM)^{\oplus r} need not be finite free; the frozen claim is preserved.
証明を終えるには、すべての $n \geq 2$ に対して次を示さなければならない：
$$
\xi = \sum\nolimits_{|I| = n, I = (i_1, \ldots, i_r)}
m_I f_1^{i_1} \ldots f_r^{i_r} \in J^{n + 1}M
$$
ならば、すべての $I$ に対して $m_I \in JM$ である。
次の段落では $I = (0, \ldots, 0, n)$ に対する $m_I \in JM$ を証明し、
最後の段落でこの特別な場合から一般の場合を導く。

\medskip\noindent
$I = (0, \ldots, 0, n)$ とする。$\xi$ を上のとおりとする。
これを
$\xi = m_1 f_1 + \ldots + m_{r - 1}f_{r - 1} + m_I f_r^n$
と書ける。$\xi \in J^{n + 1}M$ と仮定したので、さらに
$\xi = \sum_{1 \leq i \leq j \leq r - 1} m_{ij}f_if_j +
\sum_{1 \leq i \leq r - 1}m'_i f_if_r^n + m'' f_r^{n + 1}$
と書ける。従って
$$
\begin{matrix}
(m_1 - m_{11}f_1 - m'_1f_r^n)f_1 + \\
(m_2 - m_{12}f_1 - m_{22}f_2 - m'_2f_r^n)f_2 + \\
\ldots + \\
(m_{r - 1} - m_{1 r - 1}f_1 - \ldots - m_{r - 1 r - 1}f_{r - 1}
- m'_{r - 1}f_r^n)f_{r - 1} + \\
(m_I - m'' f_r)f_r^n = 0
\end{matrix}
$$
である。Lemma \ref{more-algebra-lemma-mult-koszul-regular} により
$f_1, \ldots, f_{r - 1}, f_r^n$ は $M$-$H_1$-正則だから、
$m_I - m'' f_r$ は部分加群
$f_1M + \ldots + f_{r - 1}M + f_r^nM$ に属する。
従って $m_I \in f_1M + \ldots + f_rM$ である。

\medskip\noindent
$S = R[x_1, x_2, \ldots, x_r, 1/x_r]$ とする。環写像 $R \to S$ は
忠実平坦なので、Lemma \ref{more-algebra-lemma-koszul-regular-flat-base-change} により
$f_1, \ldots, f_r$ は $S$ における $M$-$H_1$-正則列である。
Lemma \ref{more-algebra-lemma-change-basis} により
$$
g_1 = f_1 - \frac{x_1}{x_r} f_r,
\ \ldots,
\ g_{r - 1} = f_{r - 1} - \frac{x_{r - 1}}{x_r} f_r,
\ g_r = \frac{1}{x_r}f_r
$$
は $S$ における $M$-$H_1$-正則列である。最後に、元 $\xi$ は
$$
\xi = \sum\nolimits_{|I| = n, I = (i_1, \ldots, i_r)}
m_I (g_1 + x_1 g_r)^{i_1} \ldots (g_{r - 1} + x_{r - 1} g_r)^{i_{r - 1}}
(x_rg_r)^{i_r}
$$
と書き直せ、この式における $g_r^n$ の係数は
$$
\sum m_I x_1^{i_1} \ldots x_r^{i_r}
$$
である。前段落で論じた場合により、この和は $J(M \otimes_R S)$ に属する。
単項式 $x_1^{i_1} \ldots x_r^{i_r}$ は $S$ の $R$ 上のある $R$-基の一部を
なすので、所望のとおりすべての $I$ に対して $m_I \in J$ と結論する。
% SOURCE-ERRATUM-CANDIDATE P02-E0842: m_I is an element of M, so the frozen conclusion m_I \in J is ill-typed; JM is intended.
\end{proof}

\noindent
Noether 局所環上の非零有限加群について、これまでに導入した各種の正則列は
すべて同値である。

\begin{lemma}
\label{more-algebra-lemma-noetherian-finite-all-equivalent}
$(R, \mathfrak m)$ を Noether 局所環、$M$ を非零有限 $R$-加群とし、
$f_1, \ldots, f_r \in \mathfrak m$ とする。次の条件は同値である。
\begin{enumerate}
\item $f_1, \ldots, f_r$ は $M$-正則列である。
\item $f_1, \ldots, f_r$ は $M$-Koszul-正則列である。
\item $f_1, \ldots, f_r$ は $M$-$H_1$-正則列である。
\item $f_1, \ldots, f_r$ は $M$-準正則列である。
\end{enumerate}
特に、列 $f_1, \ldots, f_r$ が $R$ における正則列であること、
Koszul 正則列であること、$H_1$-正則列であること、準正則列であることは
すべて同値である。
\end{lemma}

\begin{proof}
含意 (1) $\Rightarrow$ (2) は
Lemma \ref{more-algebra-lemma-regular-koszul-regular} である。
含意 (2) $\Rightarrow$ (3) は
Lemma \ref{more-algebra-lemma-koszul-regular-H1-regular} である。
含意 (3) $\Rightarrow$ (4) は
Lemma \ref{more-algebra-lemma-H1-regular-quasi-regular} である。
含意 (4) $\Rightarrow$ (1) は
Algebra, Lemma \ref{algebra-lemma-quasi-regular-regular} である。
\end{proof}

\begin{lemma}
\label{more-algebra-lemma-H1-regular-in-quotient}
$A$ を環、$I \subset A$ をイデアルとする。
$g_1, \ldots, g_m$ を $A$ の列とし、その $A/I$ における像が
$H_1$-正則であるとする。このとき
$I \cap (g_1, \ldots, g_m) = I(g_1, \ldots, g_m)$ である。
\end{lemma}

\begin{proof}
複体の完全列
$$
0 \to I \otimes_A K_\bullet(A, g_1, \ldots, g_m)
\to K_\bullet(A, g_1, \ldots, g_m) \to
K_\bullet(A/I, g_1, \ldots, g_m) \to 0
$$
を考える。仮定により右の複体は $H_1 = 0$ をもつので、
$$
\Coker(I^{\oplus m} \to I)
\longrightarrow
\Coker(A^{\oplus m} \to A)
$$
が単射であることが分かる。これは補題の主張と同値である。
\end{proof}

\begin{lemma}
\label{more-algebra-lemma-conormal-sequence-H1-regular}
$A$ を環とし、$I \subset J \subset A$ をイデアルとする。
$J/I \subset A/I$ が $H_1$-正則列によって生成されると仮定する。
このとき $I \cap J^2 = IJ$ である。
\end{lemma}

\begin{proof}
これを示すため、$g_1, \ldots, g_m \in J$ を、その $A/I$ における像が
$J/I$ を生成する $H_1$-正則列をなすように選ぶ。
特に $J = I + (g_1, \ldots, g_m)$ である。
$x \in I \cap J^2$ とする。$x \in J^2$ なので
$$
x =
\sum a_{ij} g_ig_j +
\sum a_j g_j +
a
$$
と書ける。ここで $a_{ij} \in A$, $a_j \in I$, $a \in I^2$ である。
% SOURCE-ERRATUM-CANDIDATE P02-E0843: frozen source omits the subject in “Because x in J^2 can write”; Japanese supplies the necessary grammar.
すると $\sum a_{ij}g_ig_j \in I \cap (g_1, \ldots, g_m)$ だから、
Lemma \ref{more-algebra-lemma-H1-regular-in-quotient} により
$\sum a_{ij}g_ig_j \in I(g_1, \ldots, g_m)$ である。
従って所望のとおり $x \in IJ$ である。
\end{proof}

\begin{lemma}
\label{more-algebra-lemma-join-quasi-regular-H1-regular}
$A$ を環とする。$I$ を $A$ における準正則列
$f_1, \ldots, f_n$ によって生成されるイデアルとする。
$g_1, \ldots, g_m \in A$ を、その像
$\overline{g}_1, \ldots, \overline{g}_m$ が $A/I$ において
$H_1$-正則列をなす元とする。このとき
$f_1, \ldots, f_n, g_1, \ldots, g_m$ は $A$ における準正則列である。
\end{lemma}

\begin{proof}
すべての $d$ に対して $g_1, \ldots, g_m$ が $A/I^d$ における
$H_1$-正則列をなすことを主張する。帰納的に、このことが
$A/I^{d - 1}$ において成り立つと仮定する。複体の短完全列
$$
0 \to K_\bullet(A, g_\bullet) \otimes_A I^{d - 1}/I^d
\to K_\bullet(A/I^d, g_\bullet) \to
K_\bullet(A/I^{d - 1}, g_\bullet) \to 0
$$
がある。$f_1, \ldots, f_n$ は準正則なので、第一の複体は
$K_\bullet(A/I, g_1, \ldots, g_m)$ のコピーの直和であり、
従って次数 $1$ で非輪状である。帰納法の仮定により最後の複体は
次数 $1$ で非輪状であり、従って中間の複体もそうである。
特に、Lemma \ref{more-algebra-lemma-H1-regular-quasi-regular} により、すべての
$d \geq 1$ に対して列 $g_1, \ldots, g_m$ は $A/I^d$ において準正則である。
ここで $f_1, \ldots, f_n, g_1, \ldots, g_m$ が $A$ における
準正則列であることを証明できる。
実際、$J = (f_1, \ldots, f_n, g_1, \ldots, g_m)$ とおき、
（多重指数記法で）
$$
\sum\nolimits_{|N| + |M| = d} a_{N, M} f^N g^M \in J^{d + 1}
$$
がある $a_{N, M} \in A$ に対して成り立つとする。すべての $N, M$ に対して
$a_{N, M} \in J$ であることを示さなければならない。
$e \in \{0, 1, \ldots, d\}$ とする。このとき
$$
\sum\nolimits_{|N| = d - e, \ |M| = e} a_{N, M} f^N g^M \in
(g_1, \ldots, g_m)^{e + 1} + I^{d - e + 1}
$$
である。$g_1, \ldots, g_m$ は $A/I^{d - e + 1}$ における
準正則列なので、$|M| = e$ を満たす各 $M$ に対して
$$
\sum\nolimits_{|N| = d - e} a_{N, M} f^N \in
(g_1, \ldots, g_m) + I^{d - e + 1}
$$
を得る。環 $A/I^{d - e + 1}$ において
$I^{d - e}/I^{d - e + 1}$ に
Lemma \ref{more-algebra-lemma-H1-regular-in-quotient} を適用すると、これは
$\sum_{|N| = d - e} a_{N, M} f^N \in I^{d - e}(g_1, \ldots, g_m)$
を含意する。$f_1, \ldots, f_n$ は $A$ において準正則なので、
$|N| = d - e$ と $|M| = e$ を満たす各 $N, M$ に対して
$a_{N, M} \in J$ であることが従う。これで補題が証明された。
\end{proof}

\begin{lemma}
\label{more-algebra-lemma-join-H1-regular-sequences}
$A$ を環とする。$I$ を $A$ における $H_1$-正則列
$f_1, \ldots, f_n$ によって生成されるイデアルとする。
$g_1, \ldots, g_m \in A$ を、その像
$\overline{g}_1, \ldots, \overline{g}_m$ が $A/I$ において
$H_1$-正則列をなす元とする。このとき
$f_1, \ldots, f_n, g_1, \ldots, g_m$ は $A$ における
$H_1$-正則列である。
\end{lemma}

\begin{proof}
$H_1(A, f_1, \ldots, f_n, g_1, \ldots, g_m) = 0$ を示さなければならない。
そのため交換図
$$
\xymatrix{
\wedge^2(A^{\oplus n + m}) \ar[r] \ar[d] &
A^{\oplus n + m} \ar[r] \ar[d] &
A \ar[r] \ar[d] & 0 \\
\wedge^2(A/I^{\oplus m}) \ar[r] &
A/I^{\oplus m} \ar[r] &
A/I \ar[r] & 0
}
$$
を考える。
% SOURCE-ERRATUM-CANDIDATE P02-E0844: A/I^{\oplus m} is formally ambiguous or ill-typed here and below; (A/I)^{\oplus m} is intended, but the frozen spans remain.
$A$ において零に写る元
$(a_1, \ldots, a_{n + m}) \in A^{\oplus n + m}$ を考える。
$\overline{g}_1, \ldots, \overline{g}_m$ は $A/I$ における
$H_1$-正則列をなすので、
$(\overline{a}_{n + 1}, \ldots, \overline{a}_{n + m})$ は
$\wedge^2(A/I^{\oplus m})$ のある元 $\overline{\alpha}$ の像である。
$\overline{\alpha}$ を元
$\alpha \in \wedge^2(A^{\oplus n + m})$ に持ち上げ、その
$A^{\oplus n + m}$ における像を元 $(a_1, \ldots, a_{n + m})$ から引ける。
従って $a_{n + 1}, \ldots, a_{n + m} \in I$ と仮定してよい。
$I = (f_1, \ldots, f_n)$ なので、元 $(a_1, \ldots, a_{n + m})$ を
左上の水平矢印の像に属する元
$$
(0, \ldots, g_j, 0, \ldots, 0, f_i, 0, \ldots, 0)
$$
の線形結合によって変更し、
$a_{n + 1}, \ldots, a_{n + m}$ が零である場合に帰着できる。
% SOURCE-ERRATUM-CANDIDATE P02-E0845: the displayed vector has two positive entries and is not a Koszul relation in general; one entry needs a minus sign.
この場合 $(a_1, \ldots, a_n, 0, \ldots, 0)$ は
$H_1(A, f_1, \ldots, f_n)$ の元を定めるが、これは零であると仮定した。
\end{proof}

\begin{lemma}
\label{more-algebra-lemma-truncate-H1-regular}
$A$ を環とする。$f_1, \ldots, f_n, g_1, \ldots, g_m \in A$ を
$H_1$-正則列とする。このとき $A/(f_1, \ldots, f_n)$ における像
$\overline{g}_1, \ldots, \overline{g}_m$ は $H_1$-正則列をなす。
\end{lemma}

\begin{proof}
$I = (f_1, \ldots, f_n)$ とおく。$A/I$ における任意の関係
$\sum_{j = 1, \ldots, m} \overline{a}_j \overline{g}_j$ が
自明な関係の線形結合であることを示さなければならない。
$I = (f_1, \ldots, f_n)$ なので、この関係を $A$ における関係
$$
\sum\nolimits_{j = 1, \ldots, m} a_j g_j +
\sum\nolimits_{i = 1, \ldots, n} b_if_i = 0
$$
に持ち上げられる。
% SOURCE-ERRATUM-CANDIDATE P02-E0846: the three sums use malformed subscripts “j = 1, \ldots, m” or “i = 1, \ldots, n” instead of bounded summation notation; exact spans remain.
仮定により、この $A$ における関係は自明な関係の線形結合である。
$A/I$ における像をとれば所望の結論を得る。
\end{proof}

\begin{lemma}
\label{more-algebra-lemma-join-koszul-regular-sequences}
$A$ を環とする。$I$ を $A$ における Koszul-正則列
$f_1, \ldots, f_n$ によって生成されるイデアルとする。
$g_1, \ldots, g_m \in A$ を、その像
$\overline{g}_1, \ldots, \overline{g}_m$ が $A/I$ における
Koszul-正則列をなす元とする。このとき
$f_1, \ldots, f_n, g_1, \ldots, g_m$ は $A$ における
Koszul-正則列である。
\end{lemma}

\begin{proof}
仮定により $K_\bullet(A, f_1, \ldots, f_n)$ は $A/I$ の有限自由分解であり、
$K_\bullet(A/I, \overline{g}_1, \ldots, \overline{g}_m)$ は
$A/(f_i, g_j)$ の $A/I$ 上の有限自由分解である。従って
\begin{align*}
K_\bullet(A, f_1, \ldots, f_n, g_1, \ldots, g_m)
& = \text{Tot}(K_\bullet(A, f_1, \ldots, f_n) \otimes_A
K_\bullet(A, g_1, \ldots, g_m)) \\
& \cong A/I \otimes_A K_\bullet(A, g_1, \ldots, g_m) \\
& = K_\bullet(A/I, \overline{g}_1, \ldots, \overline{g}_m) \\
& \cong A/(f_i, g_j)
\end{align*}
である。第一の等式は
Lemma \ref{more-algebra-lemma-join-sequences-koszul-complex} による。
第一の擬同型 $\cong$ は Homology, Lemma
\ref{homology-lemma-double-complex-gives-resolution} の双対による。
実際、二重複体
$K_\bullet(A, f_1, \ldots, f_n) \otimes_A K_\bullet(A, g_1, \ldots, g_m)$
の第 $q$ 行は
$A/I \otimes_A K_q(A, g_1, \ldots, g_m)$ の分解である。
第二の等式は明らかであり、最後の擬同型は仮定による。従って主張が従う。
\end{proof}

\noindent
次の補題の結論を得るには、$f_1, \ldots, f_n$ と
$f_1, \ldots, f_n, g_1, \ldots, g_m$ の両方が Koszul-正則であると
仮定する必要がある。$f_1, \ldots, f_n$ が Koszul-正則であるという仮定を
取り除いた場合の反例は Examples, Lemma
\ref{examples-lemma-strange-regular-sequence} にある。

\begin{lemma}
\label{more-algebra-lemma-truncate-koszul-regular}
$A$ を環とし、
$$
f_1, \ldots, f_n, g_1, \ldots, g_m \in A
$$
とする。列
$$
f_1, \ldots, f_n
$$
と
$$
f_1, \ldots, f_n, g_1, \ldots, g_m
$$
の両方が
$A$ における Koszul-正則列ならば、
$A/(f_1, \ldots, f_n)$ における
$\overline{g}_1, \ldots, \overline{g}_m$ は Koszul-正則列をなす。
\end{lemma}

\begin{proof}
$I = (f_1, \ldots, f_n)$ とおく。仮定により
$K_\bullet(A, f_1, \ldots, f_n)$ は $A/I$ の有限自由分解であり、
$K_\bullet(A, f_1, \ldots, f_n, g_1, \ldots, g_m)$ は
$A/(f_i, g_j)$ の $A$ 上の有限自由分解である。従って
\begin{align*}
A/(f_i, g_j) & \cong K_\bullet(A, f_1, \ldots, f_n, g_1, \ldots, g_m) \\
& = \text{Tot}(K_\bullet(A, f_1, \ldots, f_n) \otimes_A
K_\bullet(A, g_1, \ldots, g_m)) \\
& \cong A/I \otimes_A K_\bullet(A, g_1, \ldots, g_m) \\
& = K_\bullet(A/I, \overline{g}_1, \ldots, \overline{g}_m)
\end{align*}
である。第一の擬同型 $\cong$ は仮定による。第一の等式は
Lemma \ref{more-algebra-lemma-join-sequences-koszul-complex} による。
第二の擬同型は Homology, Lemma
\ref{homology-lemma-double-complex-gives-resolution} の双対による。
実際、二重複体
$K_\bullet(A, f_1, \ldots, f_n) \otimes_A K_\bullet(A, g_1, \ldots, g_m)$
の第 $q$ 行は
$A/I \otimes_A K_q(A, g_1, \ldots, g_m)$ の分解である。
第二の等式は明らかである。従って主張が従う。
% SOURCE-ERRATUM-CANDIDATE P02-E0847: both parallel source proofs contain sentence fragments (“The first equality by…”, “The first quasi-isomorphism by…”), and one uses the informal “Hence we win”; Japanese supplies complete mathematical prose.
\end{proof}

\begin{lemma}
\label{more-algebra-lemma-independence-of-generators}
$R$ を環とする。$I$ を $f_1, \ldots, f_r \in R$ によって
生成されるイデアルとする。
\begin{enumerate}
\item $I$ が長さ $r$ の準正則列によって生成できるならば、
$f_1, \ldots, f_r$ は準正則列である。
\item $I$ が長さ $r$ の $H_1$-正則列によって生成できるならば、
$f_1, \ldots, f_r$ は $H_1$-正則列である。
\item $I$ が長さ $r$ の Koszul-正則列によって生成できるならば、
$f_1, \ldots, f_r$ は Koszul-正則列である。
\end{enumerate}
\end{lemma}

\begin{proof}
$I$ が長さ $r$ の準正則列によって生成できるならば、
$I/I^2$ は $R/I$ 上階数 $r$ の自由加群である。
$f_1, \ldots, f_r$ は仮定により生成するので、その像
$\overline{f}_i$ は $R/I$ 上の $I/I^2$ の基底をなす。
従って $f_1, \ldots, f_r$ は準正則列である。実際、
$I/I^2$ の自由性を除けば、これはまさに写像
$\text{Sym}^n_{R/I}(I/I^2) \to I^n/I^{n + 1}$
が同型であることを意味する。

\medskip\noindent
引き続き $I$ が、$g_1, \ldots, g_r$ とするある準正則列によって
生成できると仮定する。$g_j = \sum a_{ij}f_i$ と書く。
前段落により $f_1, \ldots, f_r$ は準正則なので、
$\det(a_{ij})$ は $I$ を法として可逆である。行列 $a_{ij}$ は写像
$R^{\oplus r} \to R^{\oplus r}$ を与え、これは Koszul 複体の写像
$\alpha : K_\bullet(R, f_1, \ldots, f_r) \to K_\bullet(R, g_1, \ldots, g_r)$
を誘導する。Lemma \ref{more-algebra-lemma-functorial} を参照せよ。
この写像は $\det(a_{ij})$ を可逆化すると同型になる。
$K_\bullet(R, f_1, \ldots, f_r)$ と
$K_\bullet(R, g_1, \ldots, g_r)$ のコホモロジー加群はいずれも
$I$ によって零化されるので（Lemma \ref{more-algebra-lemma-homotopy-koszul}）、
$\alpha$ は擬同型である。
% SOURCE-ERRATUM-CANDIDATE P02-E0848: these chain Koszul complexes have homology modules, not cohomology modules; Japanese preserves the frozen claim.

\medskip\noindent
ここで $g_1, \ldots, g_r$ が $I$ を生成する $H_1$-正則列であると仮定する。
Lemma \ref{more-algebra-lemma-H1-regular-quasi-regular} により
$g_1, \ldots, g_r$ は準正則列である。前段落により
$f_1, \ldots, f_r$ は $H_1$-正則列であると結論する。
Koszul-正則列についても同様である。
\end{proof}

\begin{lemma}
\label{more-algebra-lemma-make-nonzero-divisor}
\begin{reference}
これは \cite[Corollary]{McCoy} の特別な場合である。
\end{reference}
$R$ を環とする。$a_1, \ldots, a_n \in R$ を、
$R \to R^{\oplus n}$, $x \mapsto (xa_1, \ldots, xa_n)$ が単射となる元とする。
このとき多項式環 $R[t_1, \ldots, t_n]$ の元
$\sum a_i t_i$ は非零因子である。
\end{lemma}

\begin{proof}
$a_i$ の一つが単元ならば、これは $R$ 上の多項式環において
$t_1 + a_2 t_2 + \ldots + a_n t_n$ という形の任意の元が
非零因子であるという主張にほかならない。

\medskip\noindent
場合 I：$R$ は Noether 環である。$\mathfrak q_j$, $j = 1, \ldots, m$
を $R$ の付随素イデアルとする。写像
$$
\sum a_i t_i :
\text{Sym}^d(R^{\oplus n})
\longrightarrow
\text{Sym}^{d + 1}(R^{\oplus n})
$$
の各々が単射であることを示さなければならない。
$\text{Sym}^d(R^{\oplus n})$ は自由 $R$-加群なので、その付随素イデアルは
$\mathfrak q_j$, $j = 1, \ldots, m$ である。各 $j$ に対し、
$a_i \not \in \mathfrak q_j$ を満たす $i = i(j)$ が存在する。
実際、$\mathfrak q_jx = 0$ を満たす $x \in R$ が存在する一方、
仮定によりある $i$ に対して $a_i x \not = 0$ である。
従って $a_i$ は $R_{\mathfrak q_j}$ における単元であり、
写像は $\mathfrak q_j$ で局所化すると単射である。ゆえに写像は単射である。
Algebra, Lemma \ref{algebra-lemma-zero-at-ass-zero} を参照せよ。

\medskip\noindent
場合 II：一般の $R$。$R$ を、$a_1, \ldots, a_n \in R_\lambda$ を満たす
Noether 環 $R_\lambda$ の合併として書ける。各 $R_\lambda$ に対して
結果が成り立つので、$R$ に対しても結果が成り立つ。
\end{proof}

\begin{lemma}
\label{more-algebra-lemma-Koszul-regular-flat-locally-regular}
$R$ を環とする。$f_1, \ldots, f_n$ を
$(f_1, \ldots, f_n) \not = R$ を満たす $R$ における Koszul-正則列とする。
忠実平坦かつ滑らかな環写像
$$
R \longrightarrow
S = R[\{t_{ij}\}_{i \leq j}, t_{11}^{-1}, t_{22}^{-1}, \ldots, t_{nn}^{-1}]
$$
を考える。$1 \leq i \leq n$ に対し
$$
g_i = \sum\nolimits_{i \leq j} t_{ij} f_j \in S.
$$
とおく。このとき $g_1, \ldots, g_n$ は $S$ における正則列であり、
$(f_1, \ldots, f_n)S = (g_1, \ldots, g_n)$ である。
\end{lemma}

\begin{proof}
イデアルの等式は、行列
$$
\left(
\begin{matrix}
t_{11} & t_{12} & t_{13} & \ldots \\
0 & t_{22} & t_{23} & \ldots \\
0 & 0 & t_{33} & \ldots \\
\ldots & \ldots & \ldots & \ldots
\end{matrix}
\right)
$$
が $S$ において可逆であることから明らかである。
$f_1, \ldots, f_n$ は Koszul-正則列なので、写像
$R \to R^{\oplus n}$, $x \mapsto (xf_1, \ldots, xf_n)$ の核は零である
（これは $f_1, \ldots, f_n$ に関する $R$ の第 $n$ Koszul ホモロジーを計算する）。
% SOURCE-ERRATUM-CANDIDATE P02-E0849: frozen source prints “n the” after the math span instead of the ordinal “nth”; Japanese supplies the ordinal construction.
従って Lemma \ref{more-algebra-lemma-make-nonzero-divisor} により
$g_1 = f_1 t_{11} + \ldots + f_n t_{1n}$ は
$S' = R[t_{11}, t_{12}, \ldots, t_{1n}, t_{11}^{-1}]$ における非零因子である。
Lemmas \ref{more-algebra-lemma-koszul-regular-flat-base-change} および
\ref{more-algebra-lemma-independence-of-generators} により、
$g_1, f_2, \ldots, f_n$ は $S'$ における Koszul-列である。
% SOURCE-ERRATUM-CANDIDATE P02-E0850: frozen source says the undefined “Koszul-sequence”; “Koszul-regular sequence” is intended, but Japanese preserves the missing modifier.
Lemma \ref{more-algebra-lemma-truncate-koszul-regular} により
$\overline{f}_2, \ldots, \overline{f}_n$ は $S'/(g_1)$ における
Koszul-正則列であると結論する。従って $n$ に関する帰納法により、
$S'/(g_1)[\{t_{ij}\}_{2 \leq i \leq j}, t_{22}^{-1}, \ldots, t_{nn}^{-1}]$
における $g_2, \ldots, g_n$ の像
$\overline{g}_2, \ldots, \overline{g}_n$ は正則列をなす。
これはさらに $g_1, \ldots, g_n$ が $S$ における正則列をなすことを意味する。
\end{proof}

\section{Koszul 正則列についての補足}
\label{more-algebra-section-more-koszul-regular}

\noindent
Section \ref{more-algebra-section-koszul-regular} の議論を続ける。

\begin{lemma}
\label{more-algebra-lemma-vanishing-extended-alternating-koszul}
$R$ を環とし、$f_1, \ldots, f_r \in R$ を Koszul-正則列とする。
このとき Section \ref{more-algebra-section-alternating-cech} の拡張交代 {\v C}ech 複体
$R \to \bigoplus\nolimits_{i_0} R_{f_{i_0}} \to
\bigoplus\nolimits_{i_0 < i_1} R_{f_{i_0}f_{i_1}} \to
\ldots \to R_{f_1\ldots f_r}$ は次数 $r$ にしかコホモロジーをもたない。
\end{lemma}

\begin{proof}
Lemma \ref{more-algebra-lemma-mult-koszul-regular} と帰納法により、すべての
$n \geq 1$ に対して列 $f_1, \ldots, f_{r - 1}, f_r^n$ は Koszul 正則である。
Lemma \ref{more-algebra-lemma-change-basis} により、Koszul 正則列の任意の置換は
Koszul 正則列である。従って、$f_i$ のいずれか（またはすべて）を
その $n$ 乗で置き換えても Koszul 正則列が得られる。ゆえに
$K_\bullet(R, f_1^n, \ldots, f_r^n)$ はホモロジー次数 $0$ にしか
非零コホモロジーをもたない。これは
Lemma \ref{more-algebra-lemma-extended-alternating-Cech-is-colimit-koszul} により
所望の結論を含意する。
% SOURCE-ERRATUM-CANDIDATE P02-E0851: the chain Koszul complex has homology, not cohomology, in homological degree 0; Japanese preserves the frozen claim.
\end{proof}

\begin{lemma}
\label{more-algebra-lemma-blowup-regular-sequence}
$a, a_2, \ldots, a_r$ を環 $R$ における $H_1$-正則列とする
（例えば Koszul 正則列または正則列である。Lemmas
\ref{more-algebra-lemma-regular-koszul-regular} および
\ref{more-algebra-lemma-koszul-regular-H1-regular} を参照せよ）。
$I = (a, a_2, \ldots, a_r)$ とおくと、ブローアップ代数
$R' = R[\frac{I}{a}]$ は
$R'' = R[y_2, \ldots, y_r]/(a y_i - a_i)$ と同型である。
\end{lemma}

\begin{proof}
Algebra, Lemma \ref{algebra-lemma-affine-blowup-quotient-description} により、
$R''$ が $a$-ねじれなしであることを示せば十分である。

\medskip\noindent
$a, ay_2 - a_2, \ldots, ay_n - a_r$ は
$R[y_2, \ldots, y_r]$ における $H_1$-正則列であると主張する。
実際、$a, ay_2 - a_2, \ldots, ay_n - a_r$ 上の Koszul 複体を定義するために
用いる写像
$$
(a, ay_2 - a_2, \ldots, ay_n - a_r) :
R[y_2, \ldots, y_r]^{\oplus r}
\longrightarrow
R[y_2, \ldots, y_r]
$$
は、$a, a_2, \ldots, a_r$ 上の Koszul 複体を定義するために用いる写像
$$
(a, a_2, \ldots, a_r) :
R[y_2, \ldots, y_r]^{\oplus r} \longrightarrow
R[y_2, \ldots, y_r]
$$
と、同型
$$
R[y_2, \ldots, y_r]^{\oplus r}
\longrightarrow
R[y_2, \ldots, y_r]^{\oplus r}
$$
を介して同型である。この同型は $(b_1, \ldots, b_r)$ を
$(b_1 - b_2y_2 \ldots - b_ry_r, -b_2, \ldots, - b_r)$ に送る。
% SOURCE-ERRATUM-CANDIDATE P02-E0852: four occurrences use ay_n-a_r although the indexed variables run through y_r; ay_r-a_r is intended, but all frozen spans remain.
% SOURCE-ERRATUM-CANDIDATE P02-E0853: frozen source says “used to the define”; Japanese supplies complete grammar.
% SOURCE-ERRATUM-CANDIDATE P02-E0854: the frozen change-of-basis formula omits the operator before \ldots; b_1-b_2y_2-\ldots-b_ry_r is intended.
Lemma \ref{more-algebra-lemma-functorial} により、これらの Koszul 複体は同型である。
平坦な環写像 $R \to R[y_2, \ldots, y_r]$ に
Lemma \ref{more-algebra-lemma-koszul-regular-flat-base-change} を適用すると、
主張が正しいと結論する。Lemma \ref{more-algebra-lemma-cone-koszul} により、
$a, ay_2 - a_2, \ldots, ay_n - a_r$ 上の Koszul 複体 $K$ は
$a : L \to L$ の錐である。ここで $L$ は
$ay_2 - a_2, \ldots, ay_n - a_r$ 上の Koszul 複体である。
主張により $H_1(K) = 0$ だから、$a : H_0(L) \to H_0(L)$ は単射である。
言い換えれば、所望のとおり
$R'' = R[y_2, \ldots, y_r]/(a y_i - a_i)$ は非零の
$a$-ねじれ元をもたない。
\end{proof}

\begin{lemma}
\label{more-algebra-lemma-base-change-H1-regular}
$A \to B$ を環写像とする。$f_1, \ldots, f_r$ を、
$B/(f_1, \ldots, f_r)$ が $A$-平坦となる $B$ の列とする。
$A \to A'$ を環写像とする。このとき標準写像
$$
H_1(K_\bullet(B, f_1, \ldots, f_r)) \otimes_A A'
\longrightarrow
H_1(K_\bullet(B', f'_1, \ldots, f'_r))
$$
は全射である。ここで $B' = B \otimes_A A'$ であり、
$f_i' \in B'$ は $f_i$ の像である。
\end{lemma}

\begin{proof}
列
$$
\wedge^2(B^{\oplus r}) \to B^{\oplus r} \to B \to B/J \to 0
$$
は $A$-加群の複体で、$B/J$ は $A$ 上平坦であり、
$B^{\oplus r}$ の位置にコホモロジー群
$H_1 = H_1(K_\bullet(B, f_1, \ldots, f_r))$ をもつ。
% SOURCE-ERRATUM-CANDIDATE P02-E0855: H_1 of the displayed chain complex is a homology group, not a cohomology group; Japanese preserves the frozen claim.
これを $A'$ とテンソルすると、複体
$$
\wedge^2((B')^{\oplus r}) \to (B')^{\oplus r} \to B' \to B'/J' \to 0
$$
を得る。これは $B'$ と $B'/J'$ において完全である。
$(B')^{\oplus r}$ におけるコホモロジー群
$H'_1 = H_1(K_\bullet(B', f'_1, \ldots, f'_r))$ を計算するため、
上の第一の列を完全列
$0 \to J \to B \to B/J \to 0$,
$0 \to K \to B^{\oplus r} \to J \to 0$, および
$\wedge^2(B^{\oplus r}) \to K \to H_1 \to 0$ に分割する。
$A$ 上 $A'$ とテンソルすると、完全列
$$
\begin{matrix}
0 \to J \otimes_A A' \to B \otimes_A A' \to (B/J) \otimes_A A' \to 0 \\
K \otimes_A A' \to B^{\oplus r} \otimes_A A' \to J \otimes_A A' \to 0 \\
\wedge^2(B^{\oplus r}) \otimes_A A' \to K \otimes_A A' \to H_1 \otimes_A A'
\to 0
\end{matrix}
$$
を得る。第一の列が完全なのは、$B/J$ が $A$ 上平坦だからである。
Algebra, Lemma \ref{algebra-lemma-flat-tor-zero} を参照せよ。
従って $J' = J \otimes_A A' \subset B'$ であり、
$K \otimes_A A' \to \Ker((B')^{\oplus r} \to B')$ は全射である。
ゆえに
\begin{align*}
H_1 \otimes_A A'
& =
\Coker\left(\wedge^2(B^{\oplus r}) \otimes_A A' \to K \otimes_A A'\right) \\
& \to
\Coker\left(
\wedge^2((B')^{\oplus r})  \to \Ker((B')^{\oplus r} \to B')
\right) = H'_1
\end{align*}
も全射である。
\end{proof}

\begin{lemma}
\label{more-algebra-lemma-relative-regular-immersion-algebra}
$A \to B$ と $A \to A'$ を環写像とし、$B' = B \otimes_A A'$ とおく。
$f_1, \ldots, f_r \in B$ とする。
$B/(f_1, \ldots, f_r)B$ が $A$ 上平坦であると仮定する。
% SOURCE-ERRATUM-CANDIDATE P02-E0856: frozen source omits terminal punctuation before the enumeration; Japanese supplies it.
\begin{enumerate}
\item $f_1, \ldots, f_r$ が準正則列ならば、$B'$ における像も準正則列である。
\item $f_1, \ldots, f_r$ が $H_1$-正則列ならば、
$B'$ における像も $H_1$-正則列である。
\end{enumerate}
\end{lemma}

\begin{proof}
$f_1, \ldots, f_r$ が準正則であると仮定し、
$J = (f_1, \ldots, f_r)$ とおく。仮定により
$J^n/J^{n + 1}$ は $B/J$ のコピーの直和と同型であり、
従って $A$ 上平坦である。帰納法と Algebra, Lemma
\ref{algebra-lemma-flat-ses} により、$B/J^n$ は $A$ 上平坦である。
イデアル $(J')^n$ は $J^n \otimes_A A'$ に等しい。
Algebra, Lemma \ref{algebra-lemma-flat-tor-zero} を参照せよ。
従って
$(J')^n/(J')^{n + 1} = J^n/J^{n + 1} \otimes_A A'$ であり、
これは明らかに $f_1, \ldots, f_r$ が $B'$ における準正則列であることを含意する。

\medskip\noindent
$f_1, \ldots, f_r$ が $H_1$-正則であると仮定する。
Lemma \ref{more-algebra-lemma-base-change-H1-regular} により、Koszul ホモロジー群
$H_1(K_\bullet(B, f_1, \ldots, f_r))$ が消滅することは
$H_1(K_\bullet(B', f'_1, \ldots, f'_r))$ の消滅を含意し、結論を得る。
\end{proof}

\begin{lemma}
\label{more-algebra-lemma-cut-by-koszul-flat}
$A' \to B'$ を環写像とし、$I \subset A'$ をイデアルとする。
$A = A'/I$ および $B = B'/IB'$ とおく。
$f'_1, \ldots, f'_r \in B'$ とする。次を仮定する。
\begin{enumerate}
\item $A' \to B'$ は平坦かつ有限表示である。
\item $I$ は局所冪零である。
\item 像 $f_1, \ldots, f_r \in B$ は準正則列をなす。
\item $B/(f_1, \ldots, f_r)$ は $A$ 上平坦である。
\end{enumerate}
このとき $B'/(f'_1, \ldots, f'_r)$ は $A'$ 上平坦である。
\end{lemma}

\begin{proof}
$C' = B'/(f'_1, \ldots, f'_r)$ とおく。
$A' \to C'$ が平坦であることを示さなければならない。
$\mathfrak p' \subset A'$ の上にある素イデアル
$\mathfrak r' \subset C'$ をとる。$\mathfrak q' \subset B'$ を
$\mathfrak r'$ の逆像とする。Algebra, Lemma
\ref{algebra-lemma-flat-localization} により、
$A'_{\mathfrak p'} \to C'_{\mathfrak q'}$ が平坦であることを示せば十分である。
Algebra, Lemma
\ref{algebra-lemma-grothendieck-regular-sequence-general} によれば、
$f'_1, \ldots, f'_r$ が、明らかな記法のもとで
$$
B'_{\mathfrak q'}/\mathfrak p'B'_{\mathfrak q'} =
B_\mathfrak q/\mathfrak p B_\mathfrak q =
(B \otimes_A \kappa(\mathfrak p))_\mathfrak q
$$
における正則列に写ることを示せば十分である。
分かっているのは、$f_1, \ldots, f_r$ が $B$ における準正則列であり、
$B/(f_1, \ldots, f_r)$ が $A$ 上平坦であることである。
Lemma \ref{more-algebra-lemma-relative-regular-immersion-algebra} により、
$B \otimes_A \kappa(\mathfrak p)$ における
$f'_1, \ldots, f'_r$ の像
$\overline{f}_1, \ldots, \overline{f}_r$ は準正則列をなす。
$(B \otimes_A \kappa(\mathfrak p))_\mathfrak q$ は Noether 局所環なので、
Lemma \ref{more-algebra-lemma-noetherian-finite-all-equivalent} により結論する。
\end{proof}

\begin{lemma}
\label{more-algebra-lemma-cut-by-koszul}
$A' \to B'$ を環写像とし、$I \subset A'$ をイデアルとする。
$A = A'/I$ および $B = B'/IB'$ とおく。
$f'_1, \ldots, f'_r \in B'$ とする。次を仮定する。
\begin{enumerate}
\item $A' \to B'$ は平坦かつ有限表示である（例えば滑らかである）。
\item $I$ は局所冪零である。
\item 像 $f_1, \ldots, f_r \in B$ は準正則列をなす。
\item $B/(f_1, \ldots, f_r)$ は $A$ 上滑らかである。
\end{enumerate}
このとき $B'/(f'_1, \ldots, f'_r)$ は $A'$ 上滑らかである。
\end{lemma}

\begin{proof}
$C' = B'/(f'_1, \ldots, f'_r)$ および
$C = B/(f_1, \ldots, f_r)$ とおく。$A' \to C'$ は有限表示である。
Lemma \ref{more-algebra-lemma-cut-by-koszul-flat} により $A' \to C'$ は平坦である。
$A' \to C'$ のファイバー環は $A \to C$ のファイバー環に等しく、
従って仮定 (4) により滑らかである。ゆえに Algebra, Lemma
\ref{algebra-lemma-flat-fibre-smooth} により $A' \to C'$ は滑らかである。
\end{proof}

\section{正則イデアル}
\label{more-algebra-section-ideals}

\noindent
正則イデアル層の概念は Divisors, Section
\ref{divisors-section-regular-ideal-sheaves} で非常に一般的に論じる。
ここではアフィンの場合、すなわち環のイデアルの場合に対応する概念を定義する。

\begin{definition}
\label{more-algebra-definition-regular-ideal}
$R$ を環、$I \subset R$ をイデアルとする。
\begin{enumerate}
\item すべての $\mathfrak p \in V(I)$ に対して
$g \in R$, $g \not \in \mathfrak p$ と正則列
$f_1, \ldots, f_r \in R_g$ が存在し、$I_g$ が
$f_1, \ldots, f_r$ によって生成されるとき、$I$ を
{\it 正則イデアル}という。
\item すべての $\mathfrak p \in V(I)$ に対して
$g \in R$, $g \not \in \mathfrak p$ と Koszul-正則列
$f_1, \ldots, f_r \in R_g$ が存在し、$I_g$ が
$f_1, \ldots, f_r$ によって生成されるとき、$I$ を
{\it Koszul-正則イデアル}という。
\item すべての $\mathfrak p \in V(I)$ に対して
$g \in R$, $g \not \in \mathfrak p$ と $H_1$-正則列
$f_1, \ldots, f_r \in R_g$ が存在し、$I_g$ が
$f_1, \ldots, f_r$ によって生成されるとき、$I$ を
{\it $H_1$-正則イデアル}という。
\item すべての $\mathfrak p \in V(I)$ に対して
$g \in R$, $g \not \in \mathfrak p$ と準正則列
$f_1, \ldots, f_r \in R_g$ が存在し、$I_g$ が
$f_1, \ldots, f_r$ によって生成されるとき、$I$ を
{\it 準正則イデアル}という。
\end{enumerate}
\end{definition}

\noindent
$I \subset R$ が与えられたとき、含意
\begin{align*}
I\text{ is a regular ideal}
& \Rightarrow
I\text{ is a Koszul-regular ideal} \\
& \Rightarrow
I\text{ is an }H_1\text{-regular ideal} \\
& \Rightarrow
I\text{ is a quasi-regular ideal}
\end{align*}
が成り立つことは明らかである。Lemmas
\ref{more-algebra-lemma-regular-koszul-regular},
\ref{more-algebra-lemma-koszul-regular-H1-regular}, および
\ref{more-algebra-lemma-H1-regular-quasi-regular} を参照せよ。
このようなイデアルは常に有限生成である。

\begin{lemma}
\label{more-algebra-lemma-quasi-regular-ideal-finite}
準正則イデアルは有限生成である。
\end{lemma}

\begin{proof}
$I \subset R$ を準正則イデアルとする。$V(I)$ は準コンパクトなので、
$g_1, \ldots, g_m \in R$ で
$V(I) \subset D(g_1) \cup \ldots \cup D(g_m)$ を満たし、かつ
$I_{g_j}$ が準正則列
$g_{j1}, \ldots, g_{jr_j} \in R_{g_j}$ によって生成されるものが存在する。
ある $g'_{ij} \in I$ に対して
$g_{ji} = g'_{ji}/g_j^{e_{ij}}$ と書く。
% SOURCE-ERRATUM-CANDIDATE P02-E0857: frozen source swaps the indices on g' between g'_{ji} and “some g'_{ij}”; Japanese preserves both exact spans.
$V(I) \subset D(g_1) \cup \ldots \cup D(g_m)$ なので可能であるように、
ある $x \in I$ に対して $1 + x = \sum g_j h_j$ と書く。
$\Spec(R) = D(g_1) \cup \ldots \cup D(g_m) \bigcup D(x)$ である。
% SOURCE-ERRATUM-CANDIDATE P02-E0858: frozen source omits terminal punctuation after the displayed affine cover; Japanese supplies it.
この被覆の各部分上で元 $g'_{ij}$ と $x$ が生成するので、
これらは $I$ を生成する。Algebra, Lemma
\ref{algebra-lemma-cover} を参照せよ。
\end{proof}

\begin{lemma}
\label{more-algebra-lemma-quasi-regular-ideal-finite-projective}
$I \subset R$ を環の準正則イデアルとする。
このとき $I/I^2$ は有限射影 $R/I$-加群である。
\end{lemma}

\begin{proof}
Algebra, Lemma \ref{algebra-lemma-finite-projective} と定義から従う。
\end{proof}

\noindent
Koszul-正則、$H_1$-正則、および準正則イデアルに対する平坦降下を証明する。

\begin{lemma}
\label{more-algebra-lemma-flat-descent-regular-ideal}
$A \to B$ を忠実平坦な環写像、$I \subset A$ をイデアルとする。
$IB$ が $B$ における Koszul-正則
（resp.\ $H_1$-正則、resp.\ 準正則）イデアルならば、
$I$ は $A$ における Koszul-正則
（resp.\ $H_1$-正則、resp.\ 準正則）イデアルである。
\end{lemma}

\begin{proof}
証明を通じて素イデアル $\mathfrak p \supset I$ を固定する。
$IB$ が準正則であると仮定する。
Lemma \ref{more-algebra-lemma-quasi-regular-ideal-finite} により
$IB$ は有限加群だから、Algebra, Lemma
\ref{algebra-lemma-descend-properties-modules} により
$I$ は有限 $A$-加群である。$A \to B$ は平坦なので
$$
I/I^2 \otimes_{A/I} B/IB = I/I^2 \otimes_A B = IB/(IB)^2.
$$
である。$IB$ は準正則だから、$B/IB$-加群 $IB/(IB)^2$ は
有限局所自由である。従って Algebra, Proposition
\ref{algebra-proposition-ffdescent-finite-projectivity} により
$I/I^2$ は有限射影である。特に、ある
$f \in A$, $f \not \in \mathfrak p$ に対して $A$ を $A_f$ で置き換えれば、
$I/I^2$ は階数 $r$ の自由加群であると仮定してよい。
$I/I^2$ の基底を与える $f_1, \ldots, f_r \in I$ を選ぶ。
中山の補題（Algebra, Lemma \ref{algebra-lemma-NAK}）により、
上と同様にもう一度 $A \leadsto A_f$ と置き換えれば、
$I$ は $f_1, \ldots, f_r$ によって生成される。

\medskip\noindent
「準正則」の場合の証明。上で、$I/I^2$ が $r$ 個の生成元
$f_1, \ldots, f_r$ 上自由であることを見た。この場合の証明を終えるには、
各 $d \geq 2$ に対して写像
$\text{Sym}^d(I/I^2) \to I^d/I^{d + 1}$ が同型であることを示せばよい。
これは忠実平坦な基底変換
$\text{Sym}^d(IB/(IB)^2) \to (IB)^d/(IB)^{d + 1}$ が仮定により
$B$ 上局所的に同型であることから明らかである。細部は省略する。

\medskip\noindent
「$H_1$-正則」および「Koszul-正則」の場合の証明。
上で構成した $I$ を生成する元の列 $f_1, \ldots, f_r$ を考える。
Lemma \ref{more-algebra-lemma-independence-of-generators} により、$IB$ が
$H_1$-正則列または Koszul-正則列によって生成されるような任意の $g \in B$
について、$f_1, \ldots, f_r$ は $B_g$ における
$H_1$-正則列または Koszul-正則列に写る。従って
$K_\bullet(A, f_1, \ldots, f_r) \otimes_A B_g$ の
$H_1$、または $H_i$, $i > 0$ は消滅する。
$K_\bullet(B, f_1, \ldots, f_r) = K_\bullet(A, f_1, \ldots, f_r) \otimes_A B$
のホモロジーは $IB$ によって零化され
（Lemma \ref{more-algebra-lemma-homotopy-koszul}）、かつ
$V(IB) \subset \bigcup_{g\text{ as above}} D(g)$ なので、
$K_\bullet(A, f_1, \ldots, f_r) \otimes_A B$ は次数 $1$、
またはすべての正次数で消滅するホモロジーをもつ。
$A \to B$ は忠実平坦なので、
$K_\bullet(A, f_1, \ldots, f_r)$ についても同じことが成り立つ。
\end{proof}

\begin{lemma}
\label{more-algebra-lemma-conormal-sequence-H1-regular-ideal}
$A$ を環、$I \subset J \subset A$ をイデアルとする。
$J/I \subset A/I$ が $H_1$-正則イデアルであると仮定する。
このとき $I \cap J^2 = IJ$ である。
\end{lemma}

\begin{proof}
局所化して Lemma \ref{more-algebra-lemma-conormal-sequence-H1-regular} を適用すれば
直ちに従う。
\end{proof}

\section{局所完全交差写像}
\label{more-algebra-section-lci}

\noindent
上の結果を用いると、（有限）多項式代数による表示を使って
環の間の局所完全交差写像を定義できる。

\begin{lemma}
\label{more-algebra-lemma-koszul-independence-presentation}
$A \to B$ を有限型環写像とする。ある表示
$\alpha : A[x_1, \ldots, x_n] \to B$ の核 $I$ が Koszul-正則イデアルならば、
任意の表示 $\beta : A[y_1, \ldots, y_m] \to B$ の核 $J$ も
Koszul-正則イデアルである。
\end{lemma}

\begin{proof}
$\alpha(f_j) = \beta(y_j)$ を満たす
$f_j \in A[x_1, \ldots, x_n]$ と、
$\beta(g_i) = \alpha(x_i)$ を満たす
$g_i \in A[y_1, \ldots, y_m]$ を選ぶ。すると交換図
$$
\xymatrix{
A[x_1, \ldots, x_n, y_1, \ldots, y_m]
\ar[d]^{x_i \mapsto g_i} \ar[rr]_-{y_j \mapsto f_j} & &
A[x_1, \ldots, x_n] \ar[d] \\
A[y_1, \ldots, y_m] \ar[rr] & & B
}
$$
を得る。$A[x_i, y_j] \to B$ の核 $K$ は
$K = (I, y_j - f_j) = (J, x_i - g_i)$ に等しいことに注意せよ。
特に、Lemma \ref{more-algebra-lemma-quasi-regular-ideal-finite} により $I$ は有限生成だから、
$J = K/(x_i - g_i)$ も有限生成である。

\medskip\noindent
素イデアル $\mathfrak q \subset B$ を選ぶ。
$I/I^2 \oplus B^{\oplus m} = J/J^2 \oplus B^{\oplus n}$
（Algebra, Lemma \ref{algebra-lemma-conormal-module}）なので、
$$
\dim J/J^2 \otimes_B \kappa(\mathfrak q) + n =
\dim I/I^2 \otimes_B \kappa(\mathfrak q) + m.
$$
である。
$I/I^2 \otimes \kappa(\mathfrak q) =
I \otimes_{A[x_i]} \kappa(\mathfrak q)$ の基底に写る
$p_1, \ldots, p_t \in I$ を選ぶ。
$J/J^2 \otimes \kappa(\mathfrak q) =
J \otimes_{A[y_j]} \kappa(\mathfrak q)$ の基底に写る
$q_1, \ldots, q_s \in J$ を選ぶ。従って $s + n = t + m$ である。
中山の補題により、いずれも $\kappa(\mathfrak q)$ の非零元に写り、かつ
$A[x_i, 1/h]$ において $I_h = (p_1, \ldots, p_t)$、
$A[y_j, 1/h']$ において $J_{h'} = (q_1, \ldots, q_s)$ となる
$h \in A[x_i]$ と $h' \in A[y_j]$ が存在する。
$I$ は Koszul-正則だから、$I_h$ が Koszul 正則列によって生成されるとも
仮定してよい。この列の長さは必ず
$t = \dim I/I^2 \otimes_B \kappa(\mathfrak q)$ なので、
Lemma \ref{more-algebra-lemma-independence-of-generators} により
$p_1, \ldots, p_t$ は Koszul-正則列である。
$y_1 - f_1, \ldots, y_m - f_m$ も正則列なので、
Lemma \ref{more-algebra-lemma-join-koszul-regular-sequences} により
$$
y_1 - f_1, \ldots, y_m - f_m, p_1, \ldots, p_t
$$
は $A[x_i, y_j, 1/h]$ における Koszul-正則列である。
この列はイデアル $K_h$ を生成する。従ってイデアル $K_{hh'}$ は
長さ $m + t = n + s$ の Koszul-正則列によって生成される。
一方、これは同じ長さの列
$$
x_1 - g_1, \ldots, x_n - g_n, q_1, \ldots, q_s
$$
によっても生成されるので、Lemma \ref{more-algebra-lemma-independence-of-generators} により
この列も Koszul-正則である。最後に Lemma
\ref{more-algebra-lemma-truncate-koszul-regular} により、
$$
A[x_i, y_j, 1/hh']/(x_1 - g_1, \ldots, x_n - g_n)
\cong A[y_j, 1/h'']
$$
における $q_1, \ldots, q_s$ の像は $J_{h''}$ を生成する
Koszul-正則列をなす。$h''$ は $hh'$ の像なので
$\kappa(\mathfrak q)$ において零に写らず、結論を得る。
\end{proof}

\noindent
この補題により次の定義を置ける。

\begin{definition}
\label{more-algebra-definition-local-complete-intersection}
環写像 $A \to B$ が有限型であり、ある（同値に任意の）表示
$B = A[x_1, \ldots, x_n]/I$ においてイデアル $I$ が Koszul-正則であるとき、
この写像を {\it 局所完全交差}という。
\end{definition}

\noindent
この概念は局所的である。

\begin{lemma}
\label{more-algebra-lemma-lci-local}
$R \to S$ を環写像とする。$g_1, \ldots, g_m \in S$ が単位イデアルを
生成するとする。各 $R \to S_{g_j}$ が局所完全交差ならば、
$R \to S$ もそうである。
\end{lemma}

\begin{proof}
$S = R[x_1, \ldots, x_n]/I$ を表示とする。
$S$ において $g_j$ に写る $h_j \in R[x_1, \ldots, x_n]$ を選ぶ。
すると
$R[x_1, \ldots, x_n, x_{n + 1}]/(I, x_{n + 1}h_j - 1)$
は $S_{g_j}$ の表示である。従って
$I_j = (I, x_{n + 1}h_j - 1)$ は
$R[x_1, \ldots, x_n, x_{n + 1}]$ における Koszul-正則イデアルである。
素イデアル $I \subset \mathfrak q \subset R[x_1, \ldots, x_n]$ を選ぶ。
このときある $j$ に対して $h_j \not \in \mathfrak q$ であり、
$\mathfrak q_j = (\mathfrak q, x_{n + 1}h_j - 1)$ は
$\mathfrak q$ の上にある $V(I_j)$ の素イデアルである。
$I/I^2 \otimes \kappa(\mathfrak q)$ の基底に写る
$f_1, \ldots, f_r \in I$ を選ぶ。このとき
$x_{n + 1}h_j - 1, f_1, \ldots, f_r$ は $I_j$ の元の列で、
$I_j \otimes \kappa(\mathfrak q_j)$ の基底に写る。
Algebra, Lemma \ref{algebra-lemma-principal-localization-NL} を参照せよ。
中山の補題により、$(I_j)_h$ が
$x_{n + 1}h_j - 1, f_1, \ldots, f_r$ によって生成されるような
$h \in R[x_1, \ldots, x_n, x_{n + 1}]$ が存在する。
さらに $(I_j)_h$ が長さ $e$ のある Koszul 正則列によって生成されると
仮定してよい。$I_j \otimes \kappa(\mathfrak q_j)$ の次元を見ると
$e = r + 1$ である。従って Lemma
\ref{more-algebra-lemma-independence-of-generators} により
$x_{n + 1}h_j - 1, f_1, \ldots, f_r$ は、ある
$h \in R[x_1, \ldots, x_n, x_{n + 1}]$,
$h \not \in \mathfrak q_j$ に対して $(I_j)_h$ を生成する
Koszul-正則列である。Lemma \ref{more-algebra-lemma-truncate-koszul-regular} により、
ある $h' \in R[x_1, \ldots, x_n]$, $h' \not \in \mathfrak q$ に対して
$I_{h'}$ は Koszul-正則列によって生成され、所望の結論を得る。
\end{proof}

\begin{lemma}
\label{more-algebra-lemma-relative-global-complete-intersection-koszul}
$R$ を環とする。$R[x_1, \ldots, x_n]/(f_1, \ldots, f_c)$ が
相対大域完全交差ならば、$f_1, \ldots, f_c$ は Koszul 正則列である。
\end{lemma}

\begin{proof}
ホモロジー群 $H_i(K_\bullet(f_\bullet))$ はイデアル
$(f_1, \ldots, f_c)$ によって零化されることを思い出そう。
従って $(f_1, \ldots, f_c)$ を含むすべての素イデアル
$\mathfrak q \subset R[x_1, \ldots, x_n]$ に対して
$H_i(K_\bullet(f_\bullet))_\mathfrak q$ が零であることを示せば十分である。
これは Algebra, Lemma
\ref{algebra-lemma-relative-global-complete-intersection-conormal} と、
正則列が Koszul 正則であること
（Lemma \ref{more-algebra-lemma-regular-koszul-regular}）から従う。
\end{proof}

\begin{lemma}
\label{more-algebra-lemma-syntomic-lci}
$R \to S$ を環写像とする。次は同値である。
\begin{enumerate}
\item $R \to S$ は syntomic である
（Algebra, Definition \ref{algebra-definition-lci}）。
\item $R \to S$ は平坦かつ局所完全交差である。
\end{enumerate}
\end{lemma}

\begin{proof}
(1) を仮定する。このとき定義により $R \to S$ は平坦である。
Algebra, Lemma \ref{algebra-lemma-syntomic} と
Lemma \ref{more-algebra-lemma-lci-local} により、相対大域完全交差が
局所完全交差準同型であることを示せば十分であり、これは
Lemma \ref{more-algebra-lemma-relative-global-complete-intersection-koszul} である。

\medskip\noindent
(2) を仮定する。Koszul-正則イデアルは有限生成なので、
局所完全交差は有限表示である。$R \to k$ を体への写像とする。
Algebra, Definition \ref{algebra-definition-lci-field} により、
$S' = S \otimes_R k$ が $k$ 上の局所完全交差であることを示せば十分である。
素イデアル $\mathfrak q' \subset S'$ を選ぶ。
$S = R[x_1, \ldots, x_n]/I$ と書く。このとき
$S' = k[x_1, \ldots, x_n]/I'$ であり、
$I' \subset k[x_1, \ldots, x_n]$ は $I$ の像である。
対応する素イデアルを
$\mathfrak p' \subset k[x_1, \ldots, x_n]$,
$\mathfrak q \subset S$, および
$\mathfrak p \subset R[x_1, \ldots, x_n]$ とする。
Definition \ref{more-algebra-definition-regular-ideal} により、
$g \in R[x_1, \ldots, x_n]$, $g \not \in \mathfrak p$ と
$I_g$ を生成する Koszul-正則列をなす
$f_1, \ldots, f_r \in R[x_1, \ldots, x_n]_g$ が存在する。
% SOURCE-ERRATUM-CANDIDATE P02-E0859: frozen source says “By Definition … exists an g,” omitting “there” and using the wrong article; Japanese supplies complete grammar.
$S$、従って $S_g$ は $R$ 上平坦なので、
Lemma \ref{more-algebra-lemma-relative-regular-immersion-algebra} により
$k[x_1, \ldots, x_n]_g$ における像
$f'_1, \ldots, f'_r$ は $I'_g$ を生成する $H_1$-正則列をなす。
従って Lemma \ref{more-algebra-lemma-noetherian-finite-all-equivalent} により
$f'_1, \ldots, f'_r$ は
$k[x_1, \ldots, x_n]_{\mathfrak p'}$ において
$I'_{\mathfrak p'}$ を生成する正則列に写る。
Algebra, Lemma \ref{algebra-lemma-lci} を適用すると、ある
$g' \in S$, $g' \not \in \mathfrak q'$ に対する $S'_{gg'}$ は
所望のとおり $k$ 上の大域完全交差である。
\end{proof}

\noindent
局所完全交差 $R \to S$ に対して、$n \geq 2$ ならば
$H_n(L_{S/R}) = 0$ である。まだ完全な余接複体を定義していないので
これを述べて証明することはできないが、その帰結の一つを導ける。

\begin{lemma}
\label{more-algebra-lemma-transitive-lci-at-end}
$A \to B \to C$ を環写像とし、$B \to C$ が局所完全交差準同型であると仮定する。
核を $I$ とする表示 $\alpha : A[x_s, s \in S] \to B$ を選ぶ。
核を $J$ とする表示 $\beta : B[y_1, \ldots, y_m] \to C$ を選ぶ。
$\gamma : A[x_s, y_t] \to C$ を、核を $K$ とする $C$ の誘導された表示とする。
このとき完全な行をもつ標準的な交換図
$$
\xymatrix{
0 \ar[r] &
\Omega_{A[x_s]/A} \otimes C \ar[r] &
\Omega_{A[x_s, y_t]/A} \otimes C \ar[r] &
\Omega_{B[y_t]/B} \otimes C \ar[r] &
0 \\
0 \ar[r] &
I/I^2 \otimes C \ar[r] \ar[u] &
K/K^2 \ar[r] \ar[u] &
J/J^2 \ar[r] \ar[u] &
0
}
$$
を得る。特に Algebra, Lemma \ref{algebra-lemma-exact-sequence-NL} の
六項完全列は左側に零を加えて完成できる。すなわち列
$$
0 \to H_1(\NL_{B/A} \otimes_B C) \to
H_1(L_{C/A}) \to
H_1(L_{C/B}) \to
\Omega_{B/A} \otimes_B C \to
\Omega_{C/A} \to
\Omega_{C/B} \to 0
$$
は完全である。
\end{lemma}

\begin{proof}
証明すべきことは写像 $I/I^2 \otimes C \to K/K^2$ の単射性だけである。
仮定によりイデアル $J$ は Koszul-正則である。
従って Lemma \ref{more-algebra-lemma-conormal-sequence-H1-regular-ideal} により
$IA[x_s, y_j] \cap K^2 = IK$ である。これは
$K/K^2 \to J/J^2$ の核が $IA[x_s, y_j]/IK$ と同型であることを意味する。
テンソル積の右完全性により
$I/I^2 \otimes_A C = IA[x_s, y_j]/IK$ なので、
所望の単射 $I/I^2 \otimes_A C \to K/K^2$ を得る。
% SOURCE-ERRATUM-CANDIDATE P02-E0860: I/I^2 is naturally a B-module and the diagram requires tensoring with C over B, not over A; Japanese preserves both frozen _A spans.
\end{proof}

\begin{lemma}
\label{more-algebra-lemma-transitive-colimit-lci-at-end}
$A \to B \to C$ を環写像とする。$B \to C$ が局所完全交差準同型の
フィルター余極限ならば、Lemma \ref{more-algebra-lemma-transitive-lci-at-end} の結論は
引き続き成り立つ。
\end{lemma}

\begin{proof}
Lemma \ref{more-algebra-lemma-transitive-lci-at-end} と
Algebra, Lemma \ref{algebra-lemma-colimits-NL} から従う。
\end{proof}

\begin{lemma}
\label{more-algebra-lemma-henselization-NL}
$A \to B$ を局所環の局所準同型とする。
$A^h \to B^h$, resp.\ $A^{sh} \to B^{sh}$ を Hensel 化、
resp.\ 強 Hensel 化の上に誘導される写像とする
（Algebra, Lemma \ref{algebra-lemma-henselian-functorial},
resp.\ Lemma \ref{algebra-lemma-strictly-henselian-functorial}）。
このとき
$\NL_{B/A} \otimes_B B^h \to \NL_{B^h/A^h}$ および
$\NL_{B/A} \otimes_B B^{sh} \to \NL_{B^{sh}/A^{sh}}$
はコホモロジー群上の同型を誘導する。
\end{lemma}

\begin{proof}
$A^h$ は $A$ 上の \'etale 代数のフィルター余極限なので、
Algebra, Lemma \ref{algebra-lemma-colimits-NL} と
Algebra, Definition \ref{algebra-definition-etale} により
$\NL_{A^h/A}$ は非輪状複体である。同じことが $B^h/B$ にも成り立つ。
$A \to A^h \to B^h$ に対する Jacobi--Zariski 列
（Algebra, Lemma \ref{algebra-lemma-exact-sequence-NL}）を用いると、
$\NL_{B^h/A} \to \NL_{B^h/A^h}$ はコホモロジー群上の同型を誘導する。
さらに、\'etale 環写像は大域完全交差でさえあるので局所完全交差である。
Algebra, Lemma \ref{algebra-lemma-etale-standard-smooth} を参照せよ。
Lemma \ref{more-algebra-lemma-transitive-colimit-lci-at-end} により
$A \to B \to B^h$ に付随する六項完全 Jacobi--Zariski 列を得る。
これは $\NL_{B/A} \otimes_B B^h \to \NL_{B^h/A}$ が
コホモロジー群上の同型を誘導することを証明する。
これで Hensel 化上の写像の場合の証明が完了する。
強 Hensel 化の場合もまったく同様に証明される。
\end{proof}

\section{Cartier の等式と幾何学的正則性}
\label{more-algebra-section-cartier-equality}

\noindent
この節と次節の参考文献は \cite[Section 39]{MatCA} である。
この節を無理なく読むには、素朴な余接複体とその性質に習熟していることが望ましい。
Algebra, Section \ref{algebra-section-netherlander} を参照せよ。

\begin{lemma}[Cartier equality]
\label{more-algebra-lemma-cartier-equality}
$K/k$ を有限生成体拡大とする。このとき $\Omega_{K/k}$ と
$H_1(L_{K/k})$ は有限次元であり、
$\text{trdeg}_k(K) = \dim_K \Omega_{K/k} - \dim_K H_1(L_{K/k})$ である。
\end{lemma}

\begin{proof}
$k$ 上の大域完全交差
$A = k[x_1, \ldots, x_n]/(f_1, \ldots, f_c)$ で、
$K$ が $A$ の分数体と同型になるものを見いだせる。
Algebra, Lemma \ref{algebra-lemma-colimit-syntomic} とその証明を参照せよ。
この場合 Algebra, Lemmas \ref{algebra-lemma-NL-homotopy} および
\ref{algebra-lemma-localize-NL} により、$\NL_{K/k}$ は複体
$$
\bigoplus\nolimits_{j = 1, \ldots, c} K \longrightarrow
\bigoplus\nolimits_{i = 1, \ldots, n} K\text{d}x_i
$$
とホモトピー同値である。
% SOURCE-ERRATUM-CANDIDATE P02-E0861: both direct-sum subscripts use malformed “j = 1, \ldots, c” / “i = 1, \ldots, n” notation rather than bounded indices; exact spans remain.
$K$ の $k$ 上の超越次数は $A$ の次元
（Algebra, Lemma \ref{algebra-lemma-dimension-prime-polynomial-ring}）、
すなわち $n - c$ であり、結論を得る。
\end{proof}

\begin{lemma}
\label{more-algebra-lemma-transitivity-gamma}
$M/L/K$ を体拡大とする。このとき Jacobi--Zariski 列
$$
0 \to H_1(L_{L/K}) \otimes_L M \to
H_1(L_{M/K}) \to
H_1(L_{M/L}) \to
\Omega_{L/K} \otimes_L M \to
\Omega_{M/K} \to
\Omega_{M/L} \to 0
$$
は完全である。
\end{lemma}

\begin{proof}
Lemma \ref{more-algebra-lemma-transitive-colimit-lci-at-end} と
Algebra, Lemma \ref{algebra-lemma-colimit-syntomic} を組み合わせる。
\end{proof}

\begin{lemma}
\label{more-algebra-lemma-gamma-commutative-diagram}
体の交換図
$$
\xymatrix{
K \ar[r] & K' \\
k \ar[u] \ar[r] & k' \ar[u]
}
$$
が与えられ、$k'/k$ と $K'/K$ が有限生成体拡大であるとする。
このとき写像
$$
\alpha : \Omega_{K/k} \otimes_K K' \to \Omega_{K'/k'}
\quad\text{and}\quad
\beta : H_1(L_{K/k}) \otimes_K K' \to H_1(L_{K'/k'})
$$
の核と余核は有限次元であり、
$$
\dim \Ker(\alpha) - \dim \Coker(\alpha)
-\dim \Ker(\beta) + \dim \Coker(\beta)
=
\text{trdeg}_k(k') - \text{trdeg}_K(K')
$$
である。
\end{lemma}

\begin{proof}
$k \subset k' \subset K'$ および $k \subset K \subset K'$ に対する
Jacobi--Zariski 列はそれぞれ
$$
0 \to H_1(L_{k'/k}) \otimes K'  \to
H_1(L_{K'/k}) \to
H_1(L_{K'/k'}) \to
\Omega_{k'/k} \otimes K' \to
\Omega_{K'/k} \to
\Omega_{K'/k'} \to 0
$$
および
$$
0 \to H_1(L_{K/k}) \otimes K' \to
H_1(L_{K'/k}) \to
H_1(L_{K'/K}) \to
\Omega_{K/k} \otimes K' \to
\Omega_{K'/k} \to
\Omega_{K'/K} \to 0
$$
である。Lemma \ref{more-algebra-lemma-cartier-equality} により、ベクトル空間
$\Omega_{k'/k}$, $\Omega_{K'/K}$, $H_1(L_{K'/K})$, および
$H_1(L_{k'/k})$ は有限次元であり、それらの次元の交代和は
$\text{trdeg}_k(k') - \text{trdeg}_K(K')$ である。補題が従う。
\end{proof}

\section{幾何学的正則性}
\label{more-algebra-section-geometrically-regular}

\noindent
$k$ を体とし、$(A, \mathfrak m, K)$ を Noether 局所 $k$-代数とする。
Jacobi--Zariski 列（Algebra, Lemma
\ref{algebra-lemma-exact-sequence-NL}）は標準的な完全列
$$
H_1(L_{K/k}) \to \mathfrak m/\mathfrak m^2 \to \Omega_{A/k} \otimes_A K \to
\Omega_{K/k} \to 0
$$
である。実際、Algebra, Lemma \ref{algebra-lemma-NL-surjection} により
$H_1(L_{K/A}) = \mathfrak m/\mathfrak m^2$ である。
この列が左端でも完全であることが、正則局所環 $A$ が $k$ 上
幾何学的正則であるか否かを特徴づけることを示す。
Section \ref{more-algebra-section-regular-fs} ではこれを形式的滑らかさの概念と結びつける。

\begin{proposition}
\label{more-algebra-proposition-characterization-geometrically-regular}
$k$ を標数 $p > 0$ の体とし、$(A, \mathfrak m, K)$ を
Noether 局所 $k$-代数とする。次は同値である。
\begin{enumerate}
\item $A$ は $k$ 上幾何学的正則である。
\item $k$ 上有限なすべての $k \subset k' \subset k^{1/p}$ に対して、
環 $A \otimes_k k'$ は正則である。
\item $A$ は正則で、標準写像
$H_1(L_{K/k}) \to \mathfrak m/\mathfrak m^2$ は単射である。
\item $A$ は正則で、写像
$\Omega_{k/\mathbf{F}_p} \otimes_k K \to \Omega_{A/\mathbf{F}_p} \otimes_A K$
は単射である。
\end{enumerate}
\end{proposition}

\begin{proof}
(3) $\Rightarrow$ (1) の証明。
(3) を仮定する。$k'/k$ を有限純非分離拡大とする。
$A' = A \otimes_k k'$ とおく。これは極大イデアル
$\mathfrak m'$ をもつ局所環である。$K' = A'/\mathfrak m'$ とおく。
完全な行をもつ交換図
{\small
$$
\xymatrix{
0 \ar[r] &
H_1(L_{K/k}) \otimes K' \ar[r] \ar[d]_\beta &
\mathfrak m/\mathfrak m^2 \otimes K' \ar[r] \ar[d] &
\Omega_{A/k} \otimes_A K' \ar[r] \ar[d]_{\cong} &
\Omega_{K/k} \otimes K' \ar[r] \ar[d]_\alpha &
0
\\
 &
H_1(L_{K'/k'}) \ar[r] &
\mathfrak m'/(\mathfrak m')^2 \ar[r] &
\Omega_{A'/k'} \otimes_{A'} K' \ar[r] &
\Omega_{K'/k'} \ar[r] &
0
}
$$
}
を得る。第三の垂直矢印は微分加群の基底変換
（Algebra, Lemma \ref{algebra-lemma-differentials-base-change}）
により同型である。従って $\alpha$ は全射である。
Lemma \ref{more-algebra-lemma-gamma-commutative-diagram} により
$$
\dim \Ker(\alpha) - \dim \Ker(\beta) + \dim \Coker(\beta) = 0
$$
である（これらの次元はすべて有限である）。
図式追跡により
$\dim \mathfrak m'/(\mathfrak m')^2 \leq \dim \mathfrak m/\mathfrak m^2$
が分かる。一方、$A \to A'$ は有限平坦なので、
Algebra, Lemma \ref{algebra-lemma-dimension-base-fibre-total} により
$\dim(A) = \dim(A')$ である。従って定義により $A'$ は正則である。

\medskip\noindent
(3) と (4) の同値性。交換図
$$
\xymatrix{
\mathbf{F}_p \ar[r] & A \ar[r] & K \\
\mathbf{F}_p \ar[r] \ar[u] & k \ar[r] \ar[u] & K \ar[u]
}
$$
の各行に対する Jacobi--Zariski 列を考えると、完全な行をもつ交換図
$$
\xymatrix{
0 \ar[r] &
\mathfrak m/\mathfrak m^2 \ar[r] &
\Omega_{A/\mathbf{F}_p} \otimes_A K \ar[r] &
\Omega_{K/\mathbf{F}_p} \ar[r] & 0 & \\
0 \ar[r] &
H_1(L_{K/k}) \ar[r] \ar[u] &
\Omega_{k/\mathbf{F}_p} \otimes_k K \ar[r] \ar[u] &
\Omega_{K/\mathbf{F}_p} \ar[r] \ar[u] &
\Omega_{K/k} \ar[r] \ar[u] &
0
}
$$
を得る。ここで $H_1(L_{K/A}) = \mathfrak m/\mathfrak m^2$ であることと、
$K/\mathbf{F}_p$ が分離的なので $H_1(L_{K/\mathbf{F}_p}) = 0$ であることを
用いた。Algebra, Proposition
\ref{algebra-proposition-characterize-separable-field-extensions} を参照せよ。
従って
$H_1(L_{K/k}) \to \mathfrak m/\mathfrak m^2$ と
$\Omega_{k/\mathbf{F}_p} \otimes_k K \to
\Omega_{A/\mathbf{F}_p} \otimes_A K$
の核が同じ次元をもつことは明らかである。

\medskip\noindent
Faltings に従う (2) $\Rightarrow$ (4) の証明。
\cite{Faltings-einfacher} を参照せよ。
$\text{d}a_1, \ldots, \text{d}a_n$ が
$\Omega_{k/\mathbf{F}_p}$ において線形独立となるように
$a_1, \ldots, a_n \in k$ をとる。体拡大
$k' = k(a_1^{1/p}, \ldots, a_n^{1/p})$ を考える。
Algebra, Lemma \ref{algebra-lemma-size-extension-pth-roots} により
$k' = k[x_1, \ldots, x_n]/(x_1^p - a_1, \ldots, x_n^p - a_n)$ である。
特に $k'/k$ の素朴な余接複体は、零微分をもつ複体
$\bigoplus_{j = 1, \ldots, n} k' \rightarrow
\bigoplus_{i = 1, \ldots, n} k'$ とホモトピー同値である。
実際、$\Omega_{k[x_1, \ldots, x_n]/k}$ において
$\text{d}(x_j^p - a_j) = 0$ である。
% SOURCE-ERRATUM-CANDIDATE P02-E0862: both direct-sum subscripts here and in the base-changed complex use malformed bounded-index notation; exact spans remain.
上と同様に $A' = A \otimes_k k'$ および
$K' = A'/\mathfrak m'$ とおく。Algebra, Lemma
\ref{algebra-lemma-change-base-NL} により、$\NL_{A'/A}$ は零微分をもつ複体
$\bigoplus_{j = 1, \ldots, n} A' \rightarrow
\bigoplus_{i = 1, \ldots, n} A'$ とホモトピー同値である。
すなわち $H_1(L_{A'/A})$ と $\Omega_{A'/A}$ は階数 $n$ の自由加群である。
$\mathbf{F}_p \to A \to A'$ に対する Jacobi--Zariski 列は
$$
H_1(L_{A'/A}) \to \Omega_{A/\mathbf{F}_p} \otimes_A A'
\to \Omega_{A'/\mathbf{F}_p} \to \Omega_{A'/A} \to 0
$$
である。核を $(x_j^p - a_j)$ とする表示
$A[x_1, \ldots, x_n] \to A'$ を用い、Algebra, Lemma
\ref{algebra-lemma-exact-sequence-NL} の写像を書き下すと、
$H_1(L_{A'/A})$ の第 $j$ 基底ベクトルは
$\Omega_{A/\mathbf{F}_p} \otimes A'$ における
$\text{d}a_j \otimes 1$ に写る。
$\Omega_{A'/A}$ は自由、従って平坦なので、$K'$ とテンソルすると完全列
$$
K'^{\oplus n} \to \Omega_{A/\mathbf{F}_p} \otimes_A K'
\xrightarrow{\beta} \Omega_{A'/\mathbf{F}_p} \otimes_{A'} K' \to
K'^{\oplus n} \to 0
$$
を得る。元 $\text{d}a_j \otimes 1$ は $\Ker(\beta)$ を生成するので、
それらが線形独立であること、すなわち
$\dim(\Ker(\beta)) = n$ を示さなければならない。
% SOURCE-ERRATUM-CANDIDATE P02-E0863: frozen source says “show that are linearly independent,” omitting “they”; Japanese supplies the subject.
次の大きな図を考える。
$$
\xymatrix{
0 \ar[r] &
\mathfrak m'/(\mathfrak m')^2 \ar[r] &
\Omega_{A'/\mathbf{F}_p} \otimes K' \ar[r] &
\Omega_{K'/\mathbf{F}_p} \ar[r] & 0 \\
0 \ar[r] &
\mathfrak m/\mathfrak m^2 \otimes K' \ar[r] \ar[u]^\alpha &
\Omega_{A/\mathbf{F}_p} \otimes K' \ar[r] \ar[u]^\beta &
\Omega_{K/\mathbf{F}_p} \otimes K' \ar[r] \ar[u]^\gamma & 0
}
$$
Lemma \ref{more-algebra-lemma-cartier-equality} と
$\mathbf{F}_p \to K \to K'$ に対する Jacobi--Zariski 列により、
$\gamma$ の核と余核は同じ有限次元をもつ。
仮定により $A'$ は正則であり（また上で見たように $A$ と同じ次元をもつ）、
従って $\alpha$ の核と余核も同じ次元をもつ。
ゆえに $\beta$ の核と余核は同じ次元をもち、これが示すべきことであった。

\medskip\noindent
含意 (1) $\Rightarrow$ (2) は自明である。これで命題の証明が完了する。
\end{proof}

\begin{lemma}
\label{more-algebra-lemma-geometrically-regular-over-field}
$k$ を標数 $p > 0$ の体とし、$(A, \mathfrak m, K)$ を
Noether 局所 $k$-代数とする。$A$ は $k$ 上幾何学的正則であると仮定する。
$K/F/k$ を有限生成部分拡大とする。
$\varphi : k[y_1, \ldots, y_m] \to A$ を $k$-代数写像で、
$y_i$ が $K$ において $F$ の元に写り、かつ
$\text{d}y_1, \ldots, \text{d}y_m$ が
$\Omega_{F/k}$ の基底に写るものとする。
$\mathfrak p = \varphi^{-1}(\mathfrak m)$ とおく。このとき
$$
k[y_1, \ldots, y_m]_\mathfrak p \to A
$$
は平坦で、$A/\mathfrak pA$ は正則である。
\end{lemma}

\begin{proof}
$A_0 = k[y_1, \ldots, y_m]_\mathfrak p$ とおき、その極大イデアルを
$\mathfrak m_0$、剰余体を $K_0$ とする。
$\Omega_{A_0/k}$ は階数 $m$ の自由加群であり、
$\Omega_{A_0/k} \otimes K_0 \to \Omega_{K_0/k}$ は同型であることに注意せよ。
$A_0$ が $k$ 上幾何学的正則であることは明らかである。
従って Proposition
\ref{more-algebra-proposition-characterization-geometrically-regular} により
$H_1(L_{K_0/k}) \to \mathfrak m_0/\mathfrak m_0^2$ は同型である。
ここで
$$
\xymatrix{
H_1(L_{K_0/k}) \otimes K \ar[d] \ar[r] &
\mathfrak m_0/\mathfrak m_0^2 \otimes K \ar[d] \\
H_1(L_{K/k}) \ar[r] & \mathfrak m/\mathfrak m^2
}
$$
を考える。左の垂直矢印は
Lemma \ref{more-algebra-lemma-transitivity-gamma} により単射であり、
下の水平矢印も Proposition
\ref{more-algebra-proposition-characterization-geometrically-regular} により単射なので、
右の垂直矢印も単射であると結論する。
% SOURCE-ERRATUM-CANDIDATE P02-E0864: frozen source omits the predicate “is injective” after “the lower horizontal”; Japanese supplies it.
従って $A_0$ の正則パラメータ系は $A$ の正則パラメータ系の一部に写る。
Algebra, Lemmas \ref{algebra-lemma-flat-over-regular} および
\ref{algebra-lemma-regular-ring-CM} により結論を得る。
\end{proof}

\section{位相環と位相加群}
\label{more-algebra-section-topological-ring}

\noindent
位相アーベル群のいくつかの性質を手短に論じよう。
アーベル群 $M$ に、加法 $M \times M \to M$、
$(x, y) \mapsto x + y$ および逆元写像 $M \to M$、$x \mapsto -x$ が
連続となるような位相が与えられているとき、$M$ を
{\it 位相アーベル群}という。{\it 位相アーベル群の準同型}とは、
連続なアーベル群準同型にほかならない。
可換位相群の圏は加法圏であり核と余核をもつが、
公理 $\Im = \Coim$ が成り立たないためアーベル圏ではない。
$N \subset M$ が部分群ならば、$N$ と $M/N$ にも次のように位相を入れて、
これらを位相群とみなす。すなわち、$N$ には誘導位相を、$M/N$ には
商位相（すなわち $M \to M/N$ が商写像となるような位相）を入れる。
$N \subset M$ が開部分群ならば、
$M/N$ の位相は離散位相であることに注意せよ。

\medskip\noindent
$0$ の部分群からなる基本近傍系が存在するとき、$M$ の位相を
{\it 線形}であるという。このとき、これらの部分群は開でもある。
次が一例である。$I$ を有向集合とし、$G_i$ を $I$ 上の（離散）
アーベル群の逆系とする。このとき
$$
G = \lim_{i \in I} G_i
$$
は、逆極限位相に関して線形位相をもち、$0$ の基本近傍系は
$\Ker(G \to G_i)$ によって与えられる。逆に、$M$ を線形位相をもつ
アーベル群とする。開部分群からなる任意の基本系
$U_i \subset M$、$i \in I$ を選ぶ（すなわち $U_i$ は開近傍の
基本系をなし、各 $U_i$ は $M$ の部分群である）。
$i \geq i' \Leftrightarrow U_i \subset U_{i'}$ と定めると、$I$ は
有向集合である。線形位相アーベル群の準同型
$$
c : M \longrightarrow \lim_{i \in I} M/U_i.
$$
を得る。$M$ が（位相空間として）{\it 分離的}であることと
$c$ が単射であることは明らかに同値である。$c$ が同型であるとき、
$M$ を{\it 完備}であるという\footnote{完備群では極限が一意であって
ほしいので、分離的であることを完備性の一部に含める。任意の位相群に
対する完備性の定義もあり、分離性の問題を除けば、ここでの特別な場合の
定義と一致する。}。
% SOURCE-ERRATUM-CANDIDATE P02-E0865: frozen source says “a unique limits,” with incompatible article and number; Japanese expresses the intended unique-limit statement.
この条件が、上で選んだ開部分群の基本系
$\{U_i\}_{i \in I}$ の選び方に依存しないことの確認は読者に委ねる。
実際、位相アーベル群 $M^\wedge = \lim_{i \in I} M/U_i$ はこの選択に
依存せず、$M$ の{\it 完備化}と呼ばれることがある。
上のような任意の $G = \lim G_i$ は完備であり、とりわけ完備化
$M^\wedge$ は常に完備である。

\begin{definition}[位相環]
\label{more-algebra-definition-topological-ring}
\begin{reference}
\cite[Chapter 0, Sections 7.1 and 7.2]{EGA1}
\end{reference}
$R$ を環、$M$ を $R$-加群とする。
\begin{enumerate}
\item $R$ に位相が与えられ、加法と乗法がともに
$R \times R \to R$ として連続であるとき、$R$ を{\it 位相環}という。
ここで $R \times R$ には積位相を入れる。この場合、$M$ に、加法
$M \times M \to M$ とスカラー乗法 $R \times M \to M$ が連続となる
ような位相が与えられているとき、$M$ を{\it 位相加群}という。
\item {\it 位相加群の準同型}とは、連続な $R$-加群写像にほかならない。
{\it 位相環の準同型}とは、与えられた位相に関して連続な環準同型である。
\item $0$ が部分加群からなる基本近傍系をもつとき、$M$ は
{\it 線形位相をもつ}という。$0$ がイデアルからなる基本近傍系をもつ
とき、$R$ は{\it 線形位相をもつ}という。
\item $R$ が線形位相をもつとする。$I \subset R$ が開であり、かつ
$0$ の任意の近傍が、ある $n$ に対する $I^n$ を含むとき、$I$ を
{\it 定義イデアル}という。
\item $R$ が線形位相をもつとする。$R$ が定義イデアルをもつとき、
$R$ は{\it pre-admissible（前許容）}であるという。
\item $R$ が線形位相をもつとする。pre-admissible で、かつ
完備\footnote{本章の規約では、これは分離的であることを含む。}ならば、
$R$ は{\it admissible（許容）}であるという。
\item $R$ が線形位相をもつとする。ある定義イデアル $I$ が存在し、
$\{I^n\}_{n \geq 0}$ が $0$ の基本近傍系をなすとき、$R$ は
{\it pre-adic} であるという。
\item $R$ が線形位相をもつとする。$R$ が pre-adic かつ完備ならば、
$R$ は{\it adic} であるという。
\end{enumerate}
(pre)adic 位相環とは、すべての $n \geq 1$ に対して $I^n$ が開となる
定義イデアル $I$ をもつ (pre)admissible 位相環にほかならないことに
注意せよ。
\end{definition}

\noindent
$R$ を環、$M$ を $R$-加群とし、$I \subset R$ をイデアルとする。
$\{I^n\}_{n \geq 0}$ を $0$ の基本近傍系とする $R$ 上の線形位相を
考えることができる。この位相を $I$-進位相という。$I$-進位相に関して
$R$ は pre-adic 位相環である\footnote{従って、$I$-進位相を
$I$-pre-adic 位相と呼ぶこともある。}。さらに、
$\{I^nM\}_{n \geq 0}$ を $0$ の開近傍の基本系とする $M$ 上の線形位相に
よって、$M$ は位相 $R$-加群になる。これを $M$ 上の{\it $I$-進位相}
という。$M$ が $I$-進完備であること（Algebra, Definition
\ref{algebra-definition-complete} の意味で）と、$M$ が $I$-進位相に
関して完備であることは同値である\footnote{$I$-進完備化
$M^\wedge$ が $I$-進完備でないこともあり得るが、$M^\wedge$ は極限位相に
関して常に完備である。$I$ が有限生成ならば、$M^\wedge$ 上の $I$-進位相と
極限位相は一致する。Algebra, Lemma
\ref{algebra-lemma-hathat-finitely-generated} とその証明を参照せよ。}。
とりわけ、$R$ が $I$-進完備であることと、$R$ が $I$-進位相に関して
adic 位相環であることは同値である。

\medskip\noindent
特別な場合として、離散位相は $0$-進位相であり、離散位相を入れた任意の
環は adic であることに注意せよ。

\begin{lemma}
\label{more-algebra-lemma-continuous}
$\varphi : R \to S$ を環写像とする。$I \subset R$ および $J \subset S$ を
イデアルとし、$R$ に $I$-進位相、$S$ に $J$-進位相を入れる。このとき、
$\varphi$ が位相環の準同型であるための必要十分条件は、ある
$n \geq 1$ に対して $\varphi(I^n) \subset J$ となることである。
\end{lemma}

\begin{proof}
省略する。
\end{proof}

\begin{lemma}[Baire のカテゴリー定理]
\label{more-algebra-lemma-baire-category-complete-module}
$M$ を位相アーベル群とする。$M$ は線形位相をもち、完備で、かつ
$0$ の可算基本近傍系をもつと仮定する。$U_n \subset M$、$n \geq 1$ が
開稠密部分集合ならば、$\bigcap_{n \geq 1} U_n$ は稠密である。
\end{lemma}

\begin{proof}
$U_n$ を補題の主張のとおりとする。$U_n$ を
$U_1 \cap \ldots \cap U_n$ で置き換えることにより、
$U_1 \supset U_2 \supset \ldots$ と仮定してよい。
$M_n$、$n \in \mathbf{N}$ を $0$ の基本近傍系とする。
$M_{n + 1} \subset M_n$ と仮定してよい。$x \in M$ を選ぶ。
任意の $k \geq 1$ に対して、$x - y \in M_k$ を満たす
$y \in \bigcap_{n \geq 1} U_n$ が存在することを示す。

\medskip\noindent
$y$ を次のように構成する。まず、$y_1 \in x + M_k$ を満たす
$y_1 \in U_1$ を選ぶ。$U_1$ は稠密で $x + M_k$ は開なので、これは可能で
ある。次に $y_1 + M_{k_1} \subset U_1$ を満たす $k_1 > k$ を選ぶ。
$U_1$ は開なので、これは可能である。続いて
$y_2 \in y_1 + M_{k_1}$ を満たす $y_2 \in U_2$ を選ぶ。
$U_2$ は稠密で $y_2 + M_{k_1}$ は開なので、これは可能である。
% SOURCE-ERRATUM-CANDIDATE P02-E0866: frozen source cites y_2 + M_{k_1} as the open set justifying the choice of y_2; the relevant previously known open set is y_1 + M_{k_1}. Japanese preserves the frozen claim.
次に $y_2 + M_{k_2} \subset U_2$ を満たす $k_2 > k_1$ を選ぶ。
$U_2$ は開なので、これは可能である。

\medskip\noindent
このように続けると、$M$ の元の収束列 $y_i$ とその極限 $y$ を得る。
構成により $x - y \in M_k$ である。また
$$
y - y_i = (y_{i + 1} - y_i) + (y_{i + 2} - y_{i + 1}) + \ldots
$$
は $M_{k_i}$ に属するので、望みどおりすべての $i$ に対して
$y \in y_i + M_{k_i} \subset U_i$ である。
\end{proof}

\begin{lemma}
\label{more-algebra-lemma-consequence-baire-complete-module}
Lemma \ref{more-algebra-lemma-baire-category-complete-module} と同じ仮定のもとで、
$M = \bigcup_{n \geq 1} N_n$ が閉部分群 $N_n$ による表示ならば、
ある $n$ に対して $N_n$ は開である。
% SOURCE-ERRATUM-CANDIDATE P02-E0867: frozen source says “With same assumptions,” omitting “the”; Japanese supplies natural grammar.
\end{lemma}

\begin{proof}
そうでなければ、すべての $n$ に対して $U_n = M \setminus N_n$ は稠密で、
Lemma \ref{more-algebra-lemma-baire-category-complete-module} に矛盾する。
\end{proof}

\begin{lemma}[開写像補題]
\label{more-algebra-lemma-open-mapping}
$u : N \to M$ を、位相づけられたアーベル群の連続写像とする。
$M$ は分離的であり、$N$ は完備で線形位相をもち、かつ $0$ の可算基本
近傍系をもつと仮定する。このとき、次のちょうど一方が成り立つ。
% SOURCE-ERRATUM-CANDIDATE P02-E0868: frozen source says “Assume that M separated,” omitting “is”; Japanese supplies grammatical syntax.
\begin{enumerate}
\item $u$ は開写像である。または
\item ある開部分群 $N' \subset N$ に対して、像 $u(N')$ は $M$ の
至る所非稠密な部分集合である。
\end{enumerate}
\end{lemma}

\begin{proof}
$N_n$、$n \in \mathbf{N}$ を $0$ の基本近傍系とする。
$N_n$ は部分群で、$N_{n + 1} \subset N_n$ と仮定してよい。
(2) が成り立たないならば、$u(N_n)$ の閉包 $M_n$ は
$n = 1, 2, 3, \ldots$ に対して開部分群である。$u$ は連続なので、
$M_n$、$n \in \mathbf{N}$ は $M$ における $0$ の開近傍の基本系で
なければならない。また、$M_n$ は $u(N_n)$ の閉包なので、
$$
u(N_n) + M_{n + 1} = M_n
$$
がすべての $n \geq 1$ に対して成り立つ。$x_1 \in M_1$ を選ぶ。
このとき、帰納的に $y_i \in N_i$ と $x_{i + 1} \in M_{i + 1}$ を
$$
u(y_i) + x_{i + 1} = x_i
$$
となるように選べる。$N$ は完備なので、$N$ の元
$y = y_1 + y_2 + y_3 + \ldots$ が存在する。すると $M$ は分離的なので
$x_1 = u(y)$ である。従って $M_1 = u(N_1)$ である。まったく同様にして、
すべての $i \geq 2$ に対して $M_i = u(N_i)$ となることを読者は示せる。
従って $u$ は開写像である。
\end{proof}

\section{位相環の形式的に滑らかな写像}
\label{more-algebra-section-formally-smooth}

\noindent
位相環の準同型に適用できる形式的滑らかさの概念がある。

\begin{definition}
\label{more-algebra-definition-formally-smooth}
$R \to S$ を位相環の準同型とし、$R$ と $S$ は線形位相をもつとする。
位相環の準同型からなる、実線矢印が可換な任意の図式
$$
\xymatrix{
S \ar[r] \ar@{-->}[rd] & A/J \\
R \ar[r] \ar[u] & A \ar[u]
}
$$
で、$A$ が離散環、$J \subset A$ が平方零イデアルであるものに対し、
図式を可換にする点線矢印が存在するとき、$S$ は $R$ 上
{\it 形式的に滑らか}であるという。
\end{definition}

\noindent
この概念は主として、イデアル $\mathfrak m \subset R$ と
$\mathfrak n \subset S$ が与えられ、$R$ に $\mathfrak m$-進位相、
$S$ に $\mathfrak n$-進位相を入れる場合に用いる。
$\varphi : R \to S$ が連続であるための必要十分条件は、ある
$m \geq 1$ に対して
$\varphi(\mathfrak m^m) \subset  \mathfrak n$ となることである。
Lemma \ref{more-algebra-lemma-continuous} を参照せよ。この場合には、実は $S$ 上の
位相だけが関係する。

\begin{lemma}
\label{more-algebra-lemma-formally-smooth}
$\varphi : R \to S$ を環写像とする。
\begin{enumerate}
\item $R \to S$ が Algebra, Definition
\ref{algebra-definition-formally-smooth} の意味で形式的に滑らかならば、
$R \to S$ を連続にする $R$ 上の任意の線形位相と $S$ 上の任意の
pre-adic 位相に関して、$R \to S$ は形式的に滑らかである。
\item $\mathfrak n \subset S$ と $\mathfrak m \subset R$ をイデアルとし、
$R$ 上の $\mathfrak m$-進位相および $S$ 上の $\mathfrak n$-進位相に
関して $\varphi$ は連続であるとする。このとき、次は同値である。
\begin{enumerate}
\item $\varphi$ は、$R$ 上の $\mathfrak m$-進位相と $S$ 上の
$\mathfrak n$-進位相に関して形式的に滑らかである。
\item $\varphi$ は、$R$ 上の離散位相と $S$ 上の $\mathfrak n$-進位相に
関して形式的に滑らかである。
\end{enumerate}
\end{enumerate}
\end{lemma}

\begin{proof}
$R \to S$ が Algebra, Definition
\ref{algebra-definition-formally-smooth} の意味で形式的に滑らかであると
仮定する。$S$ が pre-adic 位相をもつならば、$S$ が
$\mathfrak n$-進位相をもつようなイデアル $\mathfrak n \subset S$ が
存在する。Definition \ref{more-algebra-definition-formally-smooth} における実線部分が
可換な図式が与えられたとする。$S \to A/J$ の連続性は、ある
$k \geq 1$ に対して $\mathfrak n^k$ が $A/J$ において零に写ることを
意味する。Lemma \ref{more-algebra-lemma-continuous} を参照せよ。$S$ が $R$ 上
形式的に滑らかであるという仮定から、環写像 $\psi : S \to A$ を得る。
このとき $\psi(\mathfrak n^k) \subset J$ であるから、$J^2 = 0$ より
$\psi(\mathfrak n^{2k}) = 0$ である。従って Lemma
\ref{more-algebra-lemma-continuous} により $\psi$ は連続である。これで (1) が
示された。

\medskip\noindent
(2)(b) $\Rightarrow$ (2)(a) の証明は (1) の証明と同じである。
(2)(a) を仮定する。Definition \ref{more-algebra-definition-formally-smooth} において、
$R$ 上に離散位相を用いた実線部分が可換な図式が与えられたとする。
$\varphi$ は連続なので、ある $n \geq 1$ に対して
$\varphi(\mathfrak m^n) \subset \mathfrak n$ である。
$S \to A/J$ は連続なので、ある $k \geq 1$ に対して
$\mathfrak n^k$ は $A/J$ において零に写る。従って写像 $R \to A$ のもとで
$\mathfrak m^{nk}$ は $J$ に写る。ゆえに $\mathfrak m^{2nk}$ は $A$ に
おいて零に写り、$R \to A$ は $\mathfrak m$-進位相に関して連続である。
従って (2)(a) により、望む点線矢印を得る。
\end{proof}

\begin{definition}
\label{more-algebra-definition-formally-smooth-adic}
$R \to S$ を環写像、$\mathfrak n \subset S$ をイデアルとする。
Lemma \ref{more-algebra-lemma-formally-smooth} の同値な条件 (2)(a) と (2)(b) が
成り立つとき、$R \to S$ は{\it $\mathfrak n$-進位相に関して形式的に
滑らか}であるという。
\end{definition}

\noindent
この性質は完備化に遺伝する。

\begin{lemma}
\label{more-algebra-lemma-formally-smooth-completion}
$(R, \mathfrak m)$ と $(S, \mathfrak n)$ を、有限生成イデアルが指定された
環とする。$R$ と $S$ にそれぞれ $\mathfrak m$-進位相と
$\mathfrak n$-進位相を入れる。$R \to S$ を位相環の準同型とする。
次は同値である。
\begin{enumerate}
\item $R \to S$ は $\mathfrak n$-進位相に関して形式的に滑らかである。
\item $R \to S^\wedge$ は $\mathfrak n^\wedge$-進位相に関して
形式的に滑らかである。
\item $R^\wedge \to S^\wedge$ は $\mathfrak n^\wedge$-進位相に関して
形式的に滑らかである。
\end{enumerate}
ここで $R^\wedge$ と $S^\wedge$ は、それぞれ $R$ と $S$ の
$\mathfrak m$-進完備化と $\mathfrak n$-進完備化である。
\end{lemma}

\begin{proof}
$\mathfrak m$ が有限生成であるという仮定から、$R^\wedge$ は
$\mathfrak mR^\wedge$-進完備であり、
$\mathfrak mR^\wedge = \mathfrak m^\wedge$ であり、かつ
$R^\wedge/\mathfrak m^nR^\wedge = R/\mathfrak m^n$ である。
Algebra, Lemma \ref{algebra-lemma-hathat-finitely-generated} とその証明を
参照せよ。$(S, \mathfrak n)$ についても同様である。従って、
Definition \ref{more-algebra-definition-formally-smooth} における (1)、(2)、(3) の
場合の図式が一対一に対応することは明らかである。
\end{proof}

\noindent
adic 環を用いる利点は、より強い持ち上げ性が得られることである。

\begin{lemma}
\label{more-algebra-lemma-lift-continuous}
$R \to S$ を環写像とし、$\mathfrak n$ を $S$ のイデアルとする。
$R \to S$ は $\mathfrak n$-進位相に関して形式的に滑らかであると仮定する。
実線部分が可換な図式
$$
\xymatrix{
S \ar[r]_\psi \ar@{-->}[rd] & A/J \\
R \ar[r] \ar[u] & A \ar[u]
}
$$
を考える。ここで、これは位相環の準同型からなる図式であり、$A$ は
adic、$A/J$ は $A$ の閉イデアル $J \subset A$ による（位相環としての）
商であり、ある $t \geq 1$ に対して $J^t$ は $A$ のある定義イデアルに
含まれるとする。このとき、図式を可換にする点線矢印が位相環の圏に
おいて存在する。
\end{lemma}

\begin{proof}
$I \subset A$ を、ある $t$ に対して $I \supset J^t$ となる定義イデアルと
する。このとき $A = \lim A/I^n$ であり、$J$ は閉であると仮定したので
$A/J = \lim A/J + I^n$ である。離散 $R$-代数
$A_{n, m} = A/J^n + I^m$ の次の図式を考える。
% SOURCE-ERRATUM-CANDIDATE P02-E0869: frozen source repeatedly writes quotient expressions A/J + I^n and A/J^n + I^m without parentheses; the intended objects appear to be A/(J + I^n) and A/(J^n + I^m). Japanese preserves every frozen formal span.
$$
\xymatrix{
A/J^3 + I^3 \ar[r] \ar[d] &
A/J^2 + I^3 \ar[r] \ar[d] &
A/J + I^3 \ar[d] \\
A/J^3 + I^2 \ar[r] \ar[d] &
A/J^2 + I^2 \ar[r] \ar[d] &
A/J + I^2 \ar[d] \\
A/J^3 + I \ar[r] &
A/J^2 + I \ar[r] &
A/J + I
}
$$
各可換正方形が $R$-代数の全射
$$
A_{n + 1, m + 1} \longrightarrow A_{n + 1, m} \times_{A_{n, m}} A_{n, m + 1}
$$
を定め、その核は平方零であることに注意せよ。$R$-代数写像
$\varphi_{n, m} : S \to A_{n, m}$ を帰納的に構成する。
まず、写像 $\varphi_{1, m} = \psi \bmod J + I^m$ がある。
$\psi$ は連続なので、これらの各写像も連続であることに注意せよ。
$\varphi_{1, 1}$ の選択から始め、上の図式の最下行に沿って、$S$ が
$R$ 上形式的に滑らかであることを用いて順次持ち上げることにより、
写像 $\varphi_{n, 1}$ を帰納的に選べる。残りの写像
$\varphi_{n, m}$ は $n + m$ に関する帰納法で構成する。すなわち、
上に表示した全射に沿って組
$(\varphi_{n + 1, m}, \varphi_{n, m + 1})$ を持ち上げることにより
$\varphi_{n + 1, m + 1}$ を選ぶ（ここでも $S$ が $R$ 上形式的に
滑らかであることを用いる）。このようにして、すべての写像
$\varphi_{n, m}$ は系の遷移写像と整合的になる。
$J^t \subset I$ なので、例えば
$\varphi_n = \varphi_{nt, n} \bmod I^n$ は写像 $S \to A/I^n$ を誘導する。
極限 $\varphi = \lim \varphi_n$ をとると、写像
$S \to A = \lim A/I^n$ を得る。$A/J = \lim A/J + I^n$ であることを
すでに見たので、$A/J$ への合成は $\psi$ と一致する。
最後に $\varphi$ が連続であることを示す。$\psi$ が位相環の射であるという
仮定から、ある $r \geq 1$ に対して
$\psi(\mathfrak n^r) \subset J + I/J$ であることが分かる。Lemma
\ref{more-algebra-lemma-continuous} を参照せよ。従って
$\varphi(\mathfrak n^r) \subset J + I$ であり、ゆえに望みどおり
$\varphi(\mathfrak n^{rt}) \subset I$ である。
% SOURCE-ERRATUM-CANDIDATE P02-E0870: frozen source's J + I/J is ill-typed; the quotient target needed for the inference is (J + I)/J. Japanese preserves the exact frozen span.
\end{proof}

\begin{lemma}
\label{more-algebra-lemma-increase-ideal}
$R \to S$ を環写像とし、$\mathfrak n \subset \mathfrak n' \subset S$ を
イデアルとする。$R \to S$ が $\mathfrak n$-進位相に関して形式的に
滑らかならば、$R \to S$ は $\mathfrak n'$-進位相に関して形式的に
滑らかである。
\end{lemma}

\begin{proof}
省略する。
\end{proof}

\begin{lemma}
\label{more-algebra-lemma-compose-formally-smooth}
線形位相環の、形式的に滑らかな連続準同型の合成は形式的に滑らかである。
\end{lemma}

\begin{proof}
省略する。（ヒント：これは完全に形式的であり、適切な図式を考えれば従う。）
\end{proof}

\begin{lemma}
\label{more-algebra-lemma-base-change-fs}
$R$、$S$ を環とし、$\mathfrak n \subset S$ をイデアルとする。
$R \to S$ は $\mathfrak n$-進位相に関して形式的に滑らかであるとする。
$R \to R'$ を任意の環写像とする。このとき
$R' \to S' = S \otimes_R R'$ は
$\mathfrak n' = \mathfrak nS'$-進位相に関して形式的に滑らかである。
\end{lemma}

\begin{proof}
Definition \ref{more-algebra-definition-formally-smooth} における実線部分が可換な図式
$$
\xymatrix{
S \ar[r] \ar@{-->}[rrd] & S' \ar[r] \ar@{-->}[rd] & A/J \\
R  \ar[u] \ar[r] & R' \ar[r] \ar[u] & A \ar[u]
}
$$
が与えられたとする。このとき合成 $S \to S' \to A/J$ は連続である。
仮定により、長い方の点線矢印が存在する。テンソル積の普遍性により、
短い方の点線矢印を得る。
\end{proof}

\noindent
形式的滑らかさの忠実平坦環写像に沿う降下は Algebra, Lemma
\ref{algebra-lemma-descent-formally-smooth} で見た。現在の位相環の状況でも
同様のことが成り立つ。しかしここでは、すでにかなり有用で、証明も
きわめて容易な次の簡単な版だけを示す。

\begin{lemma}
\label{more-algebra-lemma-descent-fs}
$R$、$S$ を環、$\mathfrak n \subset S$ をイデアルとする。
$R \to R'$ を環写像とし、$S' = S \otimes_R R'$ および
$\mathfrak n' = \mathfrak nS$ とおく。次を仮定する。
% SOURCE-ERRATUM-CANDIDATE P02-E0871: frozen source defines n' as nS although n' is subsequently an ideal of S'; the preceding base-change lemma uses the required nS'. Japanese preserves the frozen span.
\begin{enumerate}
\item 写像 $R \to R'$ は、$R$-加群として $R$ を $R'$ の直和因子に
埋め込む。
\item $R' \to S'$ は $\mathfrak n'$-進位相に関して形式的に滑らかである。
\end{enumerate}
このとき $R \to S$ は $\mathfrak n$-進位相に関して形式的に滑らかである。
\end{lemma}

\begin{proof}
Definition \ref{more-algebra-definition-formally-smooth} における実線部分が可換な図式
$$
\xymatrix{
S \ar[r] & A/J \\
R \ar[u] \ar[r] & A \ar[u]
}
$$
が与えられたとする。$A' = A \otimes_R R'$ および
$J' = \Im(J \otimes_R R' \to A')$ とおく。上の図式の基底変換は図式
$$
\xymatrix{
S' \ar[r] \ar@{-->}[rd]^{\psi'} & A'/J' \\
R' \ar[u] \ar[r] & A' \ar[u]
}
$$
であり、その矢印は連続である。条件 (2) により点線矢印
$\psi' : S' \to A'$ を得る。条件 (1) を用いて、$R$-加群としての直和分解
$R' = R \oplus C$ を選ぶ。（注意：$C$ は $R'$ のイデアルではない。）
このとき $A' = A \oplus A \otimes_R C$ である。おく：
$$
J'' = \Im(J \otimes_R C \to A \otimes_R C) \subset J' \subset A'.
$$
このとき $A$-加群として $J' = J \oplus J''$ である。
$\psi'$ と $S \to S'$ の合成 $\psi : S \to A'$ の像は
$A + J' = A \oplus J''$ に含まれる。しかし環
$A + J' = A \oplus J''$ において、$A$-部分加群 $J''$ はイデアルである！
（$J^2 = 0$ を用いよ。）従って合成
$S \to A + J' \to (A + J')/J'' = A$ が求める矢印である。
\end{proof}

\section{局所環の形式的に滑らかな写像}
\label{more-algebra-section-formally-smooth-local}

\noindent
局所環の局所準同型の場合には、持ち上げ性を確認すべき図式を限定できる。
Algebra, Lemma \ref{algebra-lemma-smooth-test-artinian} と比較されたい。

\begin{lemma}
\label{more-algebra-lemma-fs-local}
$(R, \mathfrak m) \to (S, \mathfrak n)$ を局所環の局所準同型とする。
次は同値である。
\begin{enumerate}
\item $R \to S$ は $\mathfrak n$-進位相に関して形式的に滑らかである。
\item 局所環の局所準同型からなる、実線部分が可換な任意の図式
$$
\xymatrix{
S \ar[r] \ar@{-->}[rd] & A/J \\
R \ar[r] \ar[u] & A \ar[u]
}
$$
で、$J \subset A$ が平方零イデアルであり、ある $n > 0$ に対して
$\mathfrak m_A^n = 0$ であり、$S \to A/J$ が剰余体上の同型を誘導する
ものに対し、図式を可換にする点線矢印が存在する。
\end{enumerate}
$S$ が Noether 環ならば、これらの条件はさらに次とも同値である。
\begin{enumerate}
\item[(3)] (2) と同じであるが、さらに $A \to A/J$ が小拡大である図式
だけを考える（Algebra, Definition
\ref{algebra-definition-small-extension}）。
\end{enumerate}
\end{lemma}

\begin{proof}
含意 (1) $\Rightarrow$ (2) は定義から従う。$R$ 上の
$\mathfrak m$-進位相と $S$ 上の $\mathfrak n$-進位相に対する Definition
\ref{more-algebra-definition-formally-smooth} における図式
$$
\xymatrix{
S \ar[r] \ar@{-->}[rd] & A/J \\
R \ar[r] \ar[u] & A \ar[u]
}
$$
を考える。$\mathfrak n^m(A/J) = 0$ となる $m > 0$ を選ぶ
（図式中の写像の連続性により可能である）。$A$ の部分環 $A'$ を、
$A/J$ における $S$ の像の逆像として定める。$J' = J$ とおき、これを
$A'$ のイデアルとみなす。このとき $J'$ は $A'$ の平方零イデアルであり、
$A'/J'$ は $S/\mathfrak n^m$ の商である。従って $A'$ は局所環で、
$\mathfrak m_{A'}^{2m} = 0$ である。こうして (2) における図式
$$
\xymatrix{
S \ar[r] \ar@{-->}[rd] & A'/J' \\
R \ar[r] \ar[u] & A' \ar[u]
}
$$
を得る。この図式で点線矢印を構成できれば、$A' \to A$ と合成して
元の図式の点線矢印を得る。このようにして (2) が (1) を含意することが
分かる。

\medskip\noindent
$S$ は Noether 環であると仮定する。含意 (1) $\Rightarrow$ (3) は直ちに
従う。(3) を仮定し、(2) における図式が与えられたとする。
このとき、ある $n > 0$ に対して $\mathfrak m_A^n J = 0$ である。
写像
$$
A \to A/\mathfrak m_A^{n - 1}J \to \ldots \to A/\mathfrak mJ \to A/J
$$
を考えると、$\mathfrak m_A J = 0$ の場合に持ち上げを構成すれば十分で
あることが分かる。$\mathfrak m_A J = 0$ と仮定し、$A' \subset A$ を
上で構成した環とする。このとき $A'/J'$ は Artin 局所環
$S/\mathfrak n^m$ の商として Artin 環である。
従って、性質 (3) から、$A/J$ が Artin 環で $\mathfrak m_A J = 0$ である
(2) の図式に点線矢印を構成できることを示せば十分である。
$\kappa$ を $A$、$A/J$、$S$ の共通剰余体とする。(3) により、
$J_0 \subset J$ が $\dim_\kappa(J/J_0) = 1$ を満たすイデアルならば、
点線矢印 $S \to A/J_0$ を構成できる。積をとると
$$
S \longrightarrow \prod\nolimits_{J_0 \text{ as above}} A/J_0
$$
を得る。この矢印の像は、小対角
$A/J \subset \prod_{J_0} A/J$ に写る元からなる $R$-部分代数 $A'$ に
含まれることは明らかである。$J' \subset A'$ を $A/J$ に零写される元の
集合とする。このとき $J'$ は平方零イデアルで、$\kappa$-ベクトル空間
として
$$
J' = \prod\nolimits_{J_0 \text{ as above}} J/J_0
$$
に等しい。従って写像 $J \to J'$ は単射である。ベクトル空間の理論により、
分裂 $J' = J \oplus M$ を選べる。すると
$$
A' = A \oplus M
$$
が $R$-代数として成り立つ。従って写像 $S \to A'$ と射影
$A' \to A$ を合成すれば、望む点線矢印が得られ、補題の証明が完了する。
\end{proof}

\noindent
次の補題は Section \ref{more-algebra-section-regular-fs} で改良される。

\begin{lemma}
\label{more-algebra-lemma-fs-implies-regular}
$k$ を体、$(A, \mathfrak m, K)$ を Noether 局所 $k$-代数とする。
$k \to A$ が $\mathfrak m$-進位相に関して形式的に滑らかならば、
$A$ は正則局所環である。
\end{lemma}

\begin{proof}
$k_0 \subset k$ を素体とする。すると $k_0$ は完全体なので、
$k / k_0$ は分離的であり、従って Algebra, Lemma
\ref{algebra-lemma-formally-smooth-extensions-easy} により形式的に滑らかで
ある。Lemmas \ref{more-algebra-lemma-formally-smooth} および
\ref{more-algebra-lemma-compose-formally-smooth} により、$k_0 \to A$ は $A$ 上の
$\mathfrak m$-進位相に関して形式的に滑らかである。従って
$k = \mathbf{Q}$ または $k = \mathbf{F}_p$ と仮定してよい。

\medskip\noindent
Algebra, Lemmas \ref{algebra-lemma-completion-faithfully-flat} および
\ref{algebra-lemma-flat-under-regular} により、完備化 $A^\wedge$ が正則で
あることを示せば十分である。Lemma
\ref{more-algebra-lemma-formally-smooth-completion} により $A$ を $A^\wedge$ で
置き換えてよい。従って $A$ は Noether 完備局所環であると仮定してよい。
Cohen 構造定理（Algebra, Theorem
\ref{algebra-theorem-cohen-structure-theorem}）により、写像 $K \to A$ が
存在する。$k$ は素体なので、$K \to A$ は $k$-代数写像である。
% SOURCE-ERRATUM-CANDIDATE P02-E0872: frozen source says “there exist a map”; Japanese supplies the required singular agreement.

\medskip\noindent
$x_1, \ldots, x_n \in \mathfrak m$ を、その像が
$\mathfrak m/\mathfrak m^2$ の基底をなす元とする。
$T = K[[X_1, \ldots, X_n]]$ とおく。次に注意せよ：
$$
A/\mathfrak m^2 \cong K[x_1, \ldots, x_n]/(x_ix_j)
$$
かつ
$$
T/\mathfrak m_T^2 \cong K[X_1, \ldots, X_n]/(X_iX_j).
$$
$x_i$ の類を $X_i$ の類に写す局所 $K$-代数同型
$A/\mathfrak m^2 \to T/m_T^2$ をとる。
% SOURCE-ERRATUM-CANDIDATE P02-E0873: frozen source writes T/m_T^2 although every adjacent occurrence uses the maximal ideal \mathfrak m_T; Japanese preserves the exact frozen span.
$f_1 : A \to T/\mathfrak m_T^2$ を、この同型と商写像
$A \to A/\mathfrak m^2$ の合成とする。$k \to A$ が $\mathfrak m$-進位相に
関して形式的に滑らかであるという仮定により、$f_1$ を写像
$f_2 : A \to T/\mathfrak{m}_T^3$ に、次いで写像
$f_3 : A \to T/\mathfrak{m}_T^4$ に、さらに同様にすべての
$n \geq 1$ に対して持ち上げられる。注意：写像 $f_n$ は連続な
$k$-代数写像であり、$K$-代数写像とは限らない。局所 $k$-代数の誘導写像
$f : A \to T = \lim T/\mathfrak m_T^n$ を得る。$f_1$ の選び方により、
$f$ は同型 $\mathfrak m/\mathfrak m^2 \to
\mathfrak m_T/\mathfrak m_T^2$ を誘導する。従って各 $f_n$ は全射であり、
$A$ は完備なので $f$ も全射であると結論する。これは
$\dim(A) \geq \dim(T) = n$ を含意する。従って定義により $A$ は正則で
ある。（さらに $f$ が同型であることも従う。）
\end{proof}

\begin{lemma}
\label{more-algebra-lemma-lift-residue-field}
$k$ を体、$(A, \mathfrak m, \kappa)$ を完備局所 $k$-代数とする。
$\kappa/k$ が分離的ならば、$\kappa \to A \to \kappa$ が
$\text{id}_\kappa$ となる $k$-代数写像 $\kappa \to A$ が存在する。
\end{lemma}

\begin{proof}
Algebra, Proposition
\ref{algebra-proposition-characterize-separable-field-extensions} により、
拡大 $\kappa/k$ は形式的に滑らかである。Lemma
\ref{more-algebra-lemma-formally-smooth} により、$k \to \kappa$ は Definition
\ref{more-algebra-definition-formally-smooth} の意味で形式的に滑らかである。
すると Lemma \ref{more-algebra-lemma-lift-continuous} から $\kappa \to A$ を得る。
\end{proof}

\begin{lemma}
\label{more-algebra-lemma-power-series-over-residue-field}
$k$ を体、$(A, \mathfrak m, \kappa)$ を完備局所 $k$-代数とする。
$\kappa/k$ が分離的で $A$ が正則ならば、$k$-代数としての同型
$A \cong \kappa[[t_1, \ldots, t_d]]$ が存在する。
% SOURCE-ERRATUM-CANDIDATE P02-E0874: frozen source says “there exists an isomorphism of A ≅ ...”; Japanese supplies grammatical phrasing without changing the formal span.
\end{lemma}

\begin{proof}
Lemma \ref{more-algebra-lemma-lift-residue-field} のように $\kappa \to A$ を選び、
Algebra, Lemma
\ref{algebra-lemma-regular-complete-containing-coefficient-field} を適用する。
\end{proof}

\noindent
次の結果は Section \ref{more-algebra-section-regular-fs} で改良される。

\begin{lemma}
\label{more-algebra-lemma-regular-implies-fs}
$k$ を体、$(A, \mathfrak m, K)$ を正則局所 $k$-代数とし、$K/k$ は
分離的であるとする。このとき $k \to A$ は $\mathfrak m$-進位相に関して
形式的に滑らかである。
\end{lemma}

\begin{proof}
$A$ の完備化が $k$ 上形式的に滑らかであることを示せば十分である。
Lemma \ref{more-algebra-lemma-formally-smooth-completion} を参照せよ。従って $A$ は
完備正則局所 $k$-代数で、剰余体 $K$ は $k$ 上分離的であると仮定して
よい。Lemma \ref{more-algebra-lemma-power-series-over-residue-field} により
$A = K[[x_1, \ldots, x_n]]$ である。

\medskip\noindent
冪級数環 $K[[x_1, \ldots, x_n]]$ は $k$ 上形式的に滑らかである。
実際、$K$ は $k$ 上形式的に滑らかであり、
$K[x_1, \ldots, x_n]$ は多項式代数として $K$ 上形式的に滑らかである。
従って Algebra, Lemma \ref{algebra-lemma-compose-formally-smooth} により
$K[x_1, \ldots, x_n]$ は $k$ 上形式的に滑らかである。
Lemma \ref{more-algebra-lemma-formally-smooth} により、
$k \to K[x_1, \ldots, x_n]$ は $(x_1, \ldots, x_n)$-進位相に関して
形式的に滑らかである。最後に Lemma
\ref{more-algebra-lemma-formally-smooth-completion} により、
$k \to K[[x_1, \ldots, x_n]]$ は $(x_1, \ldots, x_n)$-進位相に関して
形式的に滑らかである。
\end{proof}

\begin{lemma}
\label{more-algebra-lemma-formally-smooth-finite-type}
$A \to B$ を有限型環写像とし、$A$ は Noether 環であるとする。
$\mathfrak q \subset B$ を $\mathfrak p \subset A$ の上にある素イデアルと
する。次は同値である。
\begin{enumerate}
\item $A \to B$ は $\mathfrak q$ において滑らかである。
\item $A_\mathfrak p \to B_\mathfrak q$ は $\mathfrak q$-進位相に関して
形式的に滑らかである。
\end{enumerate}
\end{lemma}

\begin{proof}
含意 (2) $\Rightarrow$ (1) は Algebra, Lemma
\ref{algebra-lemma-smooth-test-artinian} から従う。逆に、$A \to B$ が
$\mathfrak q$ において滑らかならば、ある $g \in B$、
$g \not \in \mathfrak q$ に対して $A \to B_g$ は滑らかである。
すると Algebra, Proposition
\ref{algebra-proposition-smooth-formally-smooth} により $A \to B_g$ は
形式的に滑らかである。局所化は形式的滑らかさを保つので、
$A_\mathfrak p \to B_\mathfrak q$ は形式的に滑らかである
（例えば Algebra, Proposition
\ref{algebra-proposition-characterize-formally-smooth} の判定法と、
余接複体が局所化に関して良く振る舞うという事実を用いる。Algebra,
Lemmas \ref{algebra-lemma-NL-localize-bottom} および
\ref{algebra-lemma-localize-NL} を参照せよ）。最後に Lemma
\ref{more-algebra-lemma-formally-smooth} により、$A_\mathfrak p \to B_\mathfrak q$ は
$\mathfrak q$-進位相に関して形式的に滑らかである。
\end{proof}

\section{冪級数環に関するいくつかの結果}
\label{more-algebra-section-power-series}

\noindent
Noether 局所環の間の形式的に滑らかな写像に関する問題は、しばしば
冪級数環の間の写像に関する問題へ帰着できる。この節では、この種の
議論を容易にするいくつかの補助補題を証明する。

\begin{lemma}
\label{more-algebra-lemma-power-series-ring-over-Cohen-fs}
$K$ を標数 $0$ の体とし、$A = K[[x_1, \ldots, x_n]]$ とおく。
$L$ を標数 $p > 0$ の体とし、$B = L[[x_1, \ldots, x_n]]$ とおく。
$\Lambda$ を Cohen 環とし、$C = \Lambda[[x_1, \ldots, x_n]]$ とおく。
\begin{enumerate}
\item $\mathbf{Q} \to A$ は $\mathfrak m_A$-進位相に関して形式的に
滑らかである。
\item $\mathbf{F}_p \to B$ は $\mathfrak m_B$-進位相に関して形式的に
滑らかである。
\item $\mathbf{Z} \to C$ は $\mathfrak m_C$-進位相に関して形式的に
滑らかである。
\end{enumerate}
\end{lemma}

\begin{proof}
冪級数環の普遍性により、次を証明すれば十分である。
\begin{enumerate}
\item $\mathbf{Q} \to K$ は形式的に滑らかである。
\item $\mathbf{F}_p \to L$ は形式的に滑らかである。
\item $\mathbf{Z} \to \Lambda$ は $\mathfrak m_\Lambda$-進位相に関して
形式的に滑らかである。
\end{enumerate}
最初の二つは Algebra, Proposition
\ref{algebra-proposition-characterize-separable-field-extensions} である。
第三は Algebra, Lemma \ref{algebra-lemma-cohen-ring-formally-smooth} から
従う。実際、Definition \ref{more-algebra-definition-formally-smooth} における任意の
試験図式では、$p$ のある冪が $A/J$ において零となり、従って $p$ の
ある冪が $A$ において零となる。
\end{proof}

\begin{lemma}
\label{more-algebra-lemma-quotient-power-series-ring-over-Cohen}
$K$ を体とし、$A = K[[x_1, \ldots, x_n]]$ とおく。
$\Lambda$ を Cohen 環とし、$B = \Lambda[[x_1, \ldots, x_n]]$ とおく。
\begin{enumerate}
\item $y_1, \ldots, y_n \in A$ が正則パラメータ系ならば、
$K[[y_1, \ldots, y_n]] \to A$ は同型である。
% SOURCE-ERRATUM-CANDIDATE P02-E0875: frozen source uses singular “is” with the plural subject y_1,...,y_n; Japanese supplies grammatical agreement.
\item $z_1, \ldots, z_r \in A$ が $A$ の正則パラメータ系の一部を
なすならば、$r \leq n$ であり、
$A/(z_1, \ldots, z_r) \cong K[[y_1, \ldots, y_{n - r}]]$ である。
\item $p, y_1, \ldots, y_n \in B$ が正則パラメータ系ならば、
$\Lambda[[y_1, \ldots, y_n]] \to B$ は同型である。
% SOURCE-ERRATUM-CANDIDATE P02-E0876: frozen source again uses singular “is” with the plural list p,y_1,...,y_n; Japanese supplies grammatical agreement.
\item $p, z_1, \ldots, z_r \in B$ が $B$ の正則パラメータ系の一部を
なすならば、$r \leq n$ であり、
$B/(z_1, \ldots, z_r) \cong \Lambda[[y_1, \ldots, y_{n - r}]]$ である。
\end{enumerate}
\end{lemma}

\begin{proof}
(1) の証明。$A' = K[[y_1, \ldots, y_n]]$ とおく。写像 $A' \to A$ が
すべての $n \geq 1$ に対して同型
$A'/\mathfrak m_{A'}^n \to A/\mathfrak m_A^n$ を誘導することは明らかで
ある。$A$ と $A'$ はともに完備なので、$A' \to A$ は同型である。
(2) の証明。$z_1, \ldots, z_r$ を $A$ の正則パラメータ系
$z_1, \ldots, z_r, y_1, \ldots, y_{n - r}$ に延長する。写像
$A' = K[[z_1, \ldots, z_r, y_1, \ldots, y_{n - r}]] \to A$ を考える。
これは (1) により同型である。従って
$A'/(z_1, \ldots, z_r) \cong K[[y_1, \ldots, y_{n - r}]]$ は明らかなので、
(2) が従う。(3) と (4) の証明は、それぞれ (1) と (2) の証明と
まったく同じである。
\end{proof}

\begin{lemma}
\label{more-algebra-lemma-embed-map-Noetherian-complete-local-rings}
$A \to B$ を Noether 完備局所環の局所準同型とする。このとき、次の性質を
もつ可換図式
$$
\xymatrix{
S \ar[r] & B \\
R \ar[u] \ar[r] & A \ar[u]
}
$$
が存在する。
\begin{enumerate}
\item 水平矢印は全射である。
\item $A/\mathfrak m_A$ の標数が零ならば、$S$ と $R$ は体上の冪級数環で
ある。
\item $A/\mathfrak m_A$ の標数が $p > 0$ ならば、$S$ と $R$ は Cohen 環
上の冪級数環である。
\item $R \to S$ は $R$ の正則パラメータ系を $S$ の正則パラメータ系の
一部に写す。
\end{enumerate}
とりわけ $R \to S$ は平坦で（Algebra, Lemma
\ref{algebra-lemma-flat-over-regular} を参照）、そのファイバー
$S/\mathfrak m_R S$ は正則である（Algebra, Lemma
\ref{algebra-lemma-regular-ring-CM} を参照）。
\end{lemma}

\begin{proof}
Cohen 構造定理（Algebra, Theorem
\ref{algebra-theorem-cohen-structure-theorem}）を用いて、補題の主張に
あるような全射 $S \to B$ を選ぶ。この際、剰余標数が $p > 0$ ならば
$S$ を Cohen 環上の冪級数環に、それ以外ならば体上の冪級数環に選ぶ。
$J \subset S$ を $S \to B$ の核とする。次に全射
$R = \Lambda[[x_1, \ldots, x_n]] \to A$ を選ぶ。ここで、$A$ の剰余標数が
$p > 0$ ならば $\Lambda$ を Cohen 環に選び、それ以外ならば
$\Lambda$ を $A$ の剰余体に等しくとる。合成
$\Lambda[[x_1, \ldots, x_n]] \to A \to B$ を写像
$\varphi : R \to S$ に持ち上げることができる。実際、Lemma
\ref{more-algebra-lemma-power-series-ring-over-Cohen-fs} によれば
$\Lambda[[x_1, \ldots, x_n]]$ は $\mathfrak m$-進位相に関して
$\mathbf{Z}$ 上形式的に滑らかなので、Lemma
\ref{more-algebra-lemma-lift-continuous} を適用すればよい。最後に $\varphi$ を写像
$\varphi' : R = \Lambda[[x_1, \ldots, x_n]] \to
S' = S[[y_1, \ldots, y_n]]$ で置き換える。ここで
$\varphi'|_\Lambda = \varphi|_\Lambda$ および
$\varphi'(x_i) = \varphi(x_i) + y_i$ とする。また $S \to B$ を、
$y_i$ を零に写す写像 $S' \to B$ で置き換える。この置き換えの後、
$R$ の正則パラメータ系が $S'$ の正則列の一部に写ることは明らかであり、
結論を得る。
% SOURCE-ERRATUM-CANDIDATE P02-E0877: the lemma requires part of a regular system of parameters, but the proof closes with the weaker phrase “part of a regular sequence”; the construction appears to prove the stated stronger claim via cotangent-space independence.
\end{proof}

\noindent
次の補題には初等的な証明があるはずである。

\begin{lemma}
\label{more-algebra-lemma-dominate-two-surjections}
$S \to R$ と $S' \to R$ を完備 Noether 局所環の全射とする。このとき
$S \times_R S'$ は完備 Noether 局所環である。
\end{lemma}

\begin{proof}
$k$ を $R$ の剰余体とする。$k$ の標数が $p > 0$ ならば、剰余体 $k$ を
もつ Cohen 環 $\Lambda$ をとる（Algebra, Definition
\ref{algebra-definition-cohen-ring} および Algebra, Lemma
\ref{algebra-lemma-cohen-rings-exist}）。$k$ の標数が $0$ ならば
$\Lambda = k$ とおく。
% SOURCE-ERRATUM-CANDIDATE P02-E0878: frozen source says “we denote Lambda a Cohen ring,” missing the preposition normally required by this construction; Japanese supplies natural grammar.
全射 $\Lambda[[x_1, \ldots, x_n]] \to R$ を選び（Cohen 構造定理の場合と
同様。Algebra, Theorem \ref{algebra-theorem-cohen-structure-theorem} を
参照）、Lemmas \ref{more-algebra-lemma-power-series-ring-over-Cohen-fs} および
\ref{more-algebra-lemma-lift-continuous} を用いて、これを写像
$\Lambda[[x_1, \ldots, x_n]] \to S$、
$\varphi : \Lambda[[x_1, \ldots, x_n]] \to S$、および
$\varphi' : \Lambda[[x_1, \ldots, x_n]] \to S'$ に持ち上げる。
% SOURCE-ERRATUM-CANDIDATE P02-E0879: frozen source redundantly lists an unnamed lift to S and then the named lift phi to S; Japanese preserves all three formal spans.
次に、$S \to R$ の核を生成する $f_1, \ldots, f_m \in S$ と、
$S' \to R$ の核を生成する $f'_1, \ldots, f'_{m'} \in S'$ を選ぶ。
このとき写像
$$
\Lambda[[x_1, \ldots, x_n, y_1, \ldots, y_m, z_1, \ldots, z_{m'}]]
\longrightarrow S \times_R S,
$$
% SOURCE-ERRATUM-CANDIDATE P02-E0880: frozen displayed codomain is S times_R S, but the lemma, construction, and next sentence require S times_R S'. Japanese preserves the frozen display.
で、$x_i$ を $(\varphi(x_i), \varphi'(x_i))$ に、$y_j$ を $(f_j, 0)$ に、
$z_{j'}$ を $(0, f'_j)$ に写すものは全射である。
% SOURCE-ERRATUM-CANDIDATE P02-E0881: frozen source maps z_{j'} to f'_j, mismatching the primed index; the intended component is f'_{j'}. Japanese preserves the frozen span.
従って $S \times_R S'$ は完備局所環の商であり、ゆえに完備である。
\end{proof}
\section{幾何学的正則性と形式的滑らかさ}
\label{more-algebra-section-regular-fs}

\noindent
この節では、前節までの結果を組み合わせて、体上の幾何学的正則局所環の
次の特徴付けを証明する。続いて、いくつかの議論を再利用して、Noether
局所環の間の $\mathfrak m$-進位相に関する形式的に滑らかな写像の特徴付けを
証明する。

\begin{theorem}
\label{more-algebra-theorem-regular-fs}
$k$ を体、$(A, \mathfrak m, K)$ を Noether 局所 $k$-代数とする。
$k$ の標数が零ならば、次は同値である。
\begin{enumerate}
\item $A$ は正則局所環である。
\item $k \to A$ は $\mathfrak m$-進位相に関して形式的に滑らかである。
\end{enumerate}
$k$ の標数が $p > 0$ ならば、次は同値である。
\begin{enumerate}
\item $A$ は $k$ 上幾何学的正則である。
\item $k \to A$ は $\mathfrak m$-進位相に関して形式的に滑らかである。
\item $k \subset k' \subset k^{1/p}$ を満たし、$k$ 上有限なすべての
環 $A \otimes_k k'$ は正則である。
\item $A$ は正則であり、標準写像
$H_1(L_{K/k}) \to \mathfrak m/\mathfrak m^2$ は単射である。
\item $A$ は正則であり、写像
$\Omega_{k/\mathbf{F}_p} \otimes_k K \to
\Omega_{A/\mathbf{F}_p} \otimes_A K$ は単射である。
\end{enumerate}
\end{theorem}

\begin{proof}
$k$ の標数が零の場合、(1) と (2) の同値性は Lemmas
\ref{more-algebra-lemma-fs-implies-regular} および \ref{more-algebra-lemma-regular-implies-fs} から
従う。

\medskip\noindent
$k$ の標数が $p > 0$ の場合、Proposition
\ref{more-algebra-proposition-characterization-geometrically-regular} により (1)、(3)、
(4)、(5) は同値である。(2) を仮定する。Lemma
\ref{more-algebra-lemma-base-change-fs} により、$k' \to A' = A \otimes_k k'$ は
$\mathfrak m' = \mathfrak mA'$-進位相に関して形式的に滑らかである。
従って $k \subset k'$ が有限純非分離拡大ならば、Lemma
\ref{more-algebra-lemma-fs-implies-regular} により $A'$ は正則局所環である。
よって (1) が成り立つ。

\medskip\noindent
最後に (5) から (2) を証明する。Definition
\ref{more-algebra-definition-formally-smooth} における実線部分が可換な図式
$$
\xymatrix{
A \ar[r]_{\bar\psi} \ar@{-->}[rd] & B/J \\
k \ar[u]^i \ar[r]^\varphi & B \ar[u]_\pi
}
$$
を選ぶ。$J^2 = 0$ なので、$J$ は標準的な $B/J$-加群構造をもち、
$\bar\psi$ を介して $A$-加群構造をもつ。$\bar\psi$ は
$\mathfrak m$-進位相に関して連続なので、ある $n$ に対して
$\mathfrak m^nJ = 0$ である。従って $J$ を $B/J$-部分加群で
$$
0 \subset J_1 \subset J_2 \subset \ldots \subset J_n = J
$$
とフィルターでき、各商 $J_{t + 1}/J_t$ は $\mathfrak m$ により消える。
環写像の列
$$
B \to B/J_1 \to B/J_2 \to \ldots \to B/J
$$
を考えると、$J$ が $\mathfrak m$ により消える場合、すなわち $J$ が
$K$-ベクトル空間である場合に点線矢印の存在を示せば十分である。

\medskip\noindent
ここで、上のような図式で $J$ が $\mathfrak m$ により消えると仮定する。
Lemma \ref{more-algebra-lemma-regular-implies-fs} により、
$\mathbf{F}_p \to A$ は $\mathfrak m$-進位相に関して形式的に滑らかである。
従って、$\pi \circ \psi = \bar \psi$ を満たす環写像
$\psi : A \to B$ を得ることができる。すると
$\psi \circ i, \varphi : k \to B$ は、$\pi$ との合成が等しい二つの写像で
ある。従って
$$
D = \psi \circ i - \varphi : k \to J
$$
は導分である。Algebra, Lemma \ref{algebra-lemma-universal-omega} により、
ある $k$-線形写像 $\xi : \Omega_{k/\mathbf{F}_p} \to J$ を用いて
$D = \xi \circ \text{d}$ と書ける。$J$ の $K$-ベクトル空間構造を用いて、
$\xi$ を $K$-線形写像
$$
\xi' : \Omega_{k/\mathbf{F}_p} \otimes_k K \to J
$$
へ拡張する。(5) により、$K$-線形写像
$\xi'' : \Omega_{A/\mathbf{F}_p} \otimes_A K$ で、その
$\Omega_{k/\mathbf{F}_p} \otimes_k K$ への制限が $\xi'$ となるものを選べる。
$$
D' : A \xrightarrow{\text{d}} \Omega_{A/\mathbf{F}_p}
\to \Omega_{A/\mathbf{F}_p} \otimes_A K \xrightarrow{\xi''} J.
$$
とおく。最後に $\psi' = \psi - D' : A \to B$ とする。
$\psi'$ が $\pi \circ \psi' = \bar \psi$ および
$\psi' \circ i = \varphi$ を満たす環写像であることは読者が確認できる。
\end{proof}

\begin{example}
\label{more-algebra-example-fs}
$k$ を標数 $p > 0$ の体とする。$a \in k$ が $p$ 乗ではない元であるとする。
剰余体が $k$ 上純非分離な幾何学的正則局所 $k$-代数の標準的な例は環
$$
A = k[x, y]_{(x, y^p - a)}/(y^p - a - x)
$$
である。実際、$A$ は $k$ 上の滑らかな代数の局所化なので、
$k \to A$ は形式的に滑らかであり、従って $k \to A$ は
$\mathfrak m$-進位相に関して形式的に滑らかである。
密接に関連する例は次である。$k = \mathbf{F}_p(s)$、
$K = \mathbf{F}_p(t)^{perf}$ とおく。環写像
$$
k \longrightarrow A = K[[x]],\quad s \longmapsto t + x
$$
は $A$ 上の $(x)$-進位相に関して形式的に滑らかであると主張する。
実際、$\Omega_{k/\mathbf{F}_p}$ は $\text{d}s$ を基底とする $1$ 次元であり、
$\Omega_{A/\mathbf{F}_p}$ において元
$$
\text{d}x + \text{d}t = \text{d}x
$$
へ写る。$\Omega_{A/\mathbf{F}_p}$ が $A$-加群として
$\text{d}x$ 上自由であることは読者に委ねる。従って定理
\ref{more-algebra-theorem-regular-fs} の条件 (5) が成り立ち、$k \to A$ は
$(x)$-進位相に関して形式的に滑らかである。
\end{example}

\noindent
上の結果の応用として、形式的に滑らかな代数の変形には障害がないことを
証明する。





\begin{lemma}
\label{more-algebra-lemma-formally-smooth-flat}
$A \to B$ を Noether 局所環の局所準同型とする。$A \to B$ は
$\mathfrak m_B$-進位相に関して形式的に滑らかであると仮定する。
このとき $A \to B$ は平坦である。
\end{lemma}

\begin{proof}
Lemma \ref{more-algebra-lemma-formally-smooth-completion} および Algebra, Lemma
\ref{algebra-lemma-completion-Noetherian-Noetherian} により、$A$ と $B$ は
% P02-E0882: frozen English source line 9988 says “A and B a Noetherian complete local rings”; canon verification should restore “are”.
Noether 完備局所環であると仮定してよい\footnote{ここではさらに、完備化上の
写像の平坦性が $A \to B$ の平坦性を含意することを示すために Algebra,
Lemmas \ref{algebra-lemma-flatness-descends-more-general} および
\ref{algebra-lemma-completion-faithfully-flat} も用いる。}。
Lemma \ref{more-algebra-lemma-embed-map-Noetherian-complete-local-rings} における交換図
$$
\xymatrix{
S \ar[r] & B \\
R \ar[u] \ar[r] & A \ar[u]
}
$$
を $R \to S$ が平坦となるように選ぶ。$I \subset R$ を $R \to A$ の核とする。
$B$ は $A$ 上形式的に滑らかなので、$A$-代数写像
$$
S/IS \longrightarrow B
$$
は切断をもつ。Lemma \ref{more-algebra-lemma-lift-continuous} を参照せよ。従って
$B$ は平坦な $A$-加群 $S/IS$ の直和因子である（平坦性の基底変換により。
Algebra, Lemma \ref{algebra-lemma-flat-base-change} を参照）。ゆえに平坦で
ある。
\end{proof}

\begin{lemma}
\label{more-algebra-lemma-formally-smooth-JZ}
$A \to B$ を Noether 局所環の局所準同型とする。$A \to B$ は
$\mathfrak m_B$-進位相に関して形式的に滑らかであると仮定する。$K$ を
$B$ の剰余体とする。このとき、$A \to B \to K$ に対する
Jacobi--Zariski 列は完全列
$$
0 \to H_1(\NL_{K/A}) \to \mathfrak m_B/\mathfrak m_B^2
\to \Omega_{B/A} \otimes_B K \to \Omega_{K/A} \to 0
$$
を与える。
\end{lemma}

\begin{proof}
Algebra, Lemma \ref{algebra-lemma-NL-surjection} により、
$\mathfrak m_B/\mathfrak m_B^2 = H_1(\NL_{K/B})$ であることに注意せよ。
従って Algebra, Lemma \ref{algebra-lemma-exact-sequence-NL} により、
$$
H_1(\NL_{K/A}) \to \mathfrak m_B/\mathfrak m_B^2
$$
の単射性を示せばよい。$A$ の剰余体を $k$ とすると、$A \to k \to K$ に
対する Jacobi--Zariski 列から $\Omega_{K/A} = \Omega_{K/k}$ および完全列
$$
\mathfrak m_A/\mathfrak m_A^2 \otimes_k K \to
H_1(\NL_{K/A}) \to
H_1(\NL_{K/k}) \to 0
$$
を得る。$\overline{B} = B \otimes_A k$ とおく。$\overline{B}$ は正則なので、
イデアル $\mathfrak m_{\overline{B}}$ は正則列で生成される。
Lemmas \ref{more-algebra-lemma-conormal-sequence-H1-regular} および
\ref{more-algebra-lemma-noetherian-finite-all-equivalent} を
$\mathfrak m_A B \subset \mathfrak m_B$ に適用すると、
$$
\mathfrak m_A B / (\mathfrak m_AB \cap \mathfrak m_B^2) =
\mathfrak m_A B / \mathfrak m_A \mathfrak m_B
$$
を得る。これは、Lemma \ref{more-algebra-lemma-formally-smooth-flat} により
$A \to B$ が平坦なので、$\mathfrak m_A/\mathfrak m_A^2 \otimes_k K$ に等しい。
従って短完全列
$$
0 \to
\mathfrak m_A/\mathfrak m_A^2 \otimes_k K \to
\mathfrak m_B/\mathfrak m_B^2 \to
\mathfrak m_{\overline{B}}/\mathfrak m_{\overline{B}}^2 \to 0
$$
を得る。Jacobi--Zariski 列の函手性により、可換図式
$$
\xymatrix{
&
\mathfrak m_A/\mathfrak m_A^2 \otimes_k K \ar[d] \ar[r] &
H_1(\NL_{K/A}) \ar[d] \ar[r] &
H_1(\NL_{K/k}) \ar[d] \ar[r] & 0 \\
0 \ar[r] &
\mathfrak m_A/\mathfrak m_A^2 \otimes_k K \ar[r] &
\mathfrak m_B/\mathfrak m_B^2 \ar[r] &
\mathfrak m_{\overline{B}}/\mathfrak m_{\overline{B}}^2 \ar[r] & 0
}
$$
を得る。左の垂直矢印は定理
\ref{more-algebra-theorem-regular-fs} により単射である。実際、Lemma
\ref{more-algebra-lemma-base-change-fs} により $k \to \overline{B}$ は
$\mathfrak m_{\overline{B}}$-進位相に関して形式的に滑らかである。
従って蛇補題により証明が完了する。
\end{proof}

\begin{proposition}
\label{more-algebra-proposition-fs-flat-fibre-fs}
$A \to B$ を Noether 局所環の局所準同型とする。$A$ の剰余体を $k$ とし、
$\overline{B} = B \otimes_A k$ を特殊ファイバーとする。次は同値である。
\begin{enumerate}
\item $A \to B$ は平坦で、$\overline{B}$ は $k$ 上幾何学的正則である。
\item $A \to B$ は平坦で、$k \to \overline{B}$ は
$\mathfrak m_{\overline{B}}$-進位相に関して形式的に滑らかである。
\item $A \to B$ は $\mathfrak m_B$-進位相に関して形式的に滑らかである。
\end{enumerate}
\end{proposition}

\begin{proof}
(1) と (2) の同値性は Theorem \ref{more-algebra-theorem-regular-fs} から従う。

\medskip\noindent
(3) を仮定する。Lemma \ref{more-algebra-lemma-formally-smooth-flat} により、
$A \to B$ は平坦である。Lemma \ref{more-algebra-lemma-base-change-fs} により、
$k \to \overline{B}$ は $\mathfrak m_{\overline{B}}$-進位相に関して
形式的に滑らかである。従って (2) が成り立つ。

\medskip\noindent
(2) を仮定する。Lemma \ref{more-algebra-lemma-formally-smooth-completion} により、
形式的滑らかさは完備化で保たれる。平坦性についても Algebra, Lemma
\ref{algebra-lemma-completion-faithfully-flat} により同じことが成り立つ。
従って $A$ と $B$ をそれぞれの完備化で置き換え、$A$ と $B$ は Noether
完備局所環であると仮定してよい。この場合、Lemma
\ref{more-algebra-lemma-embed-map-Noetherian-complete-local-rings} における図式
$$
\xymatrix{
S \ar[r] & B \\
R \ar[u] \ar[r] & A \ar[u]
}
$$
を選ぶ。以下、この図式のすべての性質を断りなく用いる。
$R$ の正則パラメータ系 $t_1, \ldots, t_d$ を固定し、$k$ の標数が
$p > 0$ の場合には $t_1 = p$ とする。$\overline{S} = S \otimes_R k$ と
おく。短完全列
$$
0 \to J \to S \to B \to 0
$$
を考える。$\overline{B}$ と $\overline{S}$ は正則なので、
$\overline{S} \to \overline{B}$ の核は、$\overline{S}$ の正則パラメータ系の
一部をなす元 $\overline{x}_1, \ldots, \overline{x}_r$ で生成される。
Algebra, Lemma \ref{algebra-lemma-regular-quotient-regular} を参照せよ。
これらの元を $x_1, \ldots, x_r \in J$ に持ち上げる。このとき
$t_1, \ldots, t_d, x_1, \ldots, x_r$ は $S$ の正則パラメータ系の一部で
ある。従って $S/(x_1, \ldots, x_r)$ は、$k$ の標数が零ならば体上の
冪級数環、$k$ の標数が $p > 0$ ならば Cohen 環上の冪級数環である。
Lemma \ref{more-algebra-lemma-quotient-power-series-ring-over-Cohen} を参照せよ。
さらに、$R \to S/(x_1, \ldots, x_r)$ が
$t_1, \ldots, t_d$ を $S/(x_1, \ldots, x_r)$ の正則パラメータ系の一部に
写すという性質も保たれる。従って $S$ を
$S/(x_1, \ldots, x_r)$ で置き換え、さらに
$$
\xymatrix{
S \ar[r] & B \\
R \ar[u] \ar[r] & A \ar[u]
}
$$
が Lemma \ref{more-algebra-lemma-embed-map-Noetherian-complete-local-rings} の図式であり、
しかも $\overline{S} = \overline{B}$ であると仮定してよい。
この場合、写像
$$
S \otimes_R A \longrightarrow B
$$
は全射で、特殊ファイバー上の同型を誘導し、始域と終域がともに $A$ 上平坦
なので同型である（例えば Algebra, Lemma
\ref{algebra-lemma-mod-injective} を用いるか、短完全列
$0 \to I \to S \otimes_R A \to B \to 0$ を $A$ 上 $k$ とテンソルすると
$I \otimes_A k = 0$ となることから中山の補題により $I = 0$ となることを用いる）。
従って Lemma \ref{more-algebra-lemma-base-change-fs} により、$R \to S$ が
$\mathfrak m_S$-進位相に関して形式的に滑らかであることを示せば十分である。
もちろん $\overline{S} = \overline{B}$ なので、$\overline{S}$ は
$k = R/\mathfrak m_R$ 上形式的に滑らかである。

\medskip\noindent
$t_1, \ldots, t_d, y_1, \ldots, y_m$ が $S$ の正則パラメータ系となるような
$y_1, \ldots, y_m \in S$ を選ぶ。$k$ の標数が零ならば
$K \subset S$ なる係数体を選び、$k$ の標数が $p > 0$ ならば剰余体 $K$ を
もつ Cohen 環 $\Lambda \subset S$ を選ぶ。この時点で、
（標数零の場合）
$$
K[[t_1, \ldots, t_d, y_1, \ldots, y_m]] \to S
$$
または（標数 $p > 0$ の場合）
$$
\Lambda[[t_2, \ldots, t_d, y_1, \ldots, y_m]] \to S
$$
は同型である。Lemma \ref{more-algebra-lemma-quotient-power-series-ring-over-Cohen} を
参照せよ。以下、$S$ を上記の冪級数環とみなす。

\medskip\noindent
残りの証明は Theorem \ref{more-algebra-theorem-regular-fs} の証明における議論と
類似している。Definition \ref{more-algebra-definition-formally-smooth} における実線部分が
可換な図式
$$
\xymatrix{
S \ar[r]_{\bar\psi} \ar@{-->}[rd] & N/J \\
R \ar[u]^i \ar[r]^\varphi & N \ar[u]_\pi
}
$$
を選ぶ。$J^2 = 0$ なので、$J$ は標準的な $N/J$-加群構造をもち、
$\bar\psi$ を介して $S$-加群構造をもつ。$\bar\psi$ は
$\mathfrak m_S$-進位相に関して連続なので、ある $n$ に対して
$\mathfrak m_S^nJ = 0$ である。従って $J$ を $N/J$-部分加群で
$$
0 \subset J_1 \subset J_2 \subset \ldots \subset J_n = J
$$
とフィルターでき、各商 $J_{t + 1}/J_t$ は $\mathfrak m_S$ により消える。
環写像の列
$$
N \to N/J_1 \to N/J_2 \to \ldots \to N/J
$$
を考えると、$J$ が $\mathfrak m_S$ により消える場合、すなわち $J$ が
$K$-ベクトル空間である場合に点線矢印の存在を示せば十分である。

\medskip\noindent
上のような図式で $J$ が $\mathfrak m_S$ により消えると仮定する。
$\mathbf{Q} \to S$（標数零の場合）または $\mathbf{Z} \to S$（標数
$p > 0$ の場合）は $\mathfrak m_S$-進位相に関して形式的に滑らかである
（Lemma \ref{more-algebra-lemma-power-series-ring-over-Cohen-fs}）。従って
$\pi \circ \psi = \bar \psi$ を満たす環写像 $\psi : S \to N$ を得る。
$S$ は、部分環上の $t_1, \ldots, t_d$（標数零の場合）または
$t_2, \ldots, t_d$（標数 $p > 0$ の場合）に関する冪級数環なので、
$\psi(t_i)$ が $\varphi(t_i)$ に等しくなるように $\psi$ の選択を変更できる
（標数 $p$ の場合の $t_1 = p$ については自動的である）。
すると $\psi \circ i$ と $\varphi : R \to N$ は $\pi$ との合成が等しく、
かつ $t_1, \ldots, t_d$ 上で一致する二つの写像である。従って
$$
D = \psi \circ i - \varphi : R \to J
$$
は $t_1, \ldots, t_d$ を消す導分である。Algebra, Lemma
\ref{algebra-lemma-universal-omega} により、ある $R$-線形写像
$$
\xi : \Omega_{R/\mathbf{Z}} \to J
$$
を用いて $D = \xi \circ \text{d}$ と書ける。この写像は構成により
$\text{d}t_1, \ldots, \text{d}t_d$ および
$\mathfrak m_R \Omega_{R/\mathbf{Z}}$（$J$ は $\mathfrak m_R$ により消える）を
消す。従って $\xi$ は合成
$$
\Omega_{R/\mathbf{Z}} \to \Omega_{k/\mathbf{Z}} \xrightarrow{\xi'} J
$$
を介して因子化する。ここで $\xi'$ は $k$-線形である。$J$ の $K$-ベクトル
空間構造を用いて、$\xi'$ を $K$-線形写像
$\xi'' : \Omega_{k/\mathbf{Z}} \otimes_k K \longrightarrow J.$
へ拡張する。$\overline{S}/k$ が形式的に滑らかであることから、
Theorem \ref{more-algebra-theorem-regular-fs} により
$$
\Omega_{k/\mathbf{Z}} \otimes_k K \to
\Omega_{\overline{S}/\mathbf{Z}} \otimes_S K
$$
は単射である（標数零の場合にも成り立つ。実際、Algebra, Proposition
\ref{algebra-proposition-characterize-separable-field-extensions} により、
この場合は $\Omega_{k/\mathbf{Z}} \to \Omega_{K/\mathbf{Z}}$ も単射である）。
従って、その $\Omega_{k/\mathbf{Z}} \otimes_k K$ への制限が $\xi''$ となる
$K$-線形写像
$$
\xi''' : \Omega_{\overline{S}/\mathbf{Z}} \otimes_S K \to J
$$
を選べる。$$
D' : S \xrightarrow{\text{d}} \Omega_{S/\mathbf{Z}}
\to \Omega_{\overline{S}/\mathbf{Z}} \to
\Omega_{\overline{S}/\mathbf{Z}} \otimes_S K \xrightarrow{\xi'''} J.
$$
とおく。最後に $\psi' = \psi - D' : S \to N$ とする。
$\psi'$ が $\pi \circ \psi' = \bar \psi$ および
$\psi' \circ i = \varphi$ を満たす環写像であることは読者が確認できる。
\end{proof}

\begin{lemma}
\label{more-algebra-lemma-lift-fs}
$A$ を剰余体 $k$ をもつ Noether 完備局所環とする。$B$ を Noether 完備局所
$k$-代数とし、$k \to B$ は $\mathfrak m_B$-進位相に関して形式的に滑らか
であると仮定する。このとき、Noether 完備局所環 $C$ と、
$\mathfrak m_C$-進位相に関して形式的に滑らかな局所準同型 $A \to C$ で
$C \otimes_A k \cong B$ を満たすものが存在する。
\end{lemma}

\begin{proof}
Lemma \ref{more-algebra-lemma-embed-map-Noetherian-complete-local-rings} における図式
$$
\xymatrix{
S \ar[r] & B \\
R \ar[u] \ar[r] & A \ar[u]
}
$$
を選ぶ。$R$ の正則パラメータ系 $t_1, \ldots, t_d$ を、$k$ の標数が
$p > 0$ の場合には $t_1 = p$ となるように選ぶ。$B$ と
$\overline{S} = S \otimes_R k$ は正則なので、
$\Ker(\overline{S} \to B)$ は、$\overline{S}$ の正則パラメータ系の一部を
なす $\overline{x}_1, \ldots, \overline{x}_r$ で生成される。
Algebra, Lemma \ref{algebra-lemma-regular-quotient-regular} を参照せよ。
これらの元を $x_1, \ldots, x_r \in S$ に持ち上げる。すると
$t_1, \ldots, t_d, x_1, \ldots, x_r$ は $S$ の正則パラメータ系の一部で
ある。従って $S/(x_1, \ldots, x_r)$ は、$k$ の標数が零ならば体上の
冪級数環、$k$ の標数が $p > 0$ ならば Cohen 環上の冪級数環である。
Lemma \ref{more-algebra-lemma-quotient-power-series-ring-over-Cohen} を参照せよ。
さらに $R \to S/(x_1, \ldots, x_r)$ は $t_1, \ldots, t_d$ を
$S/(x_1, \ldots, x_r)$ の正則パラメータ系の一部に写す。従って $S$ を
$S/(x_1, \ldots, x_r)$ で置き換え、
$$
\xymatrix{
S \ar[r] & B \\
R \ar[u] \ar[r] & A \ar[u]
}
$$
が Lemma \ref{more-algebra-lemma-embed-map-Noetherian-complete-local-rings} の図式で、
さらに $\overline{S} = B$ であると仮定してよい。この場合、
Proposition \ref{more-algebra-proposition-fs-flat-fibre-fs} により $R \to S$ は
$\mathfrak m_S$-進位相に関して形式的に滑らかである。従って基底変換
$C = S \otimes_R A$ は、Lemma \ref{more-algebra-lemma-base-change-fs} により $A$ 上
$\mathfrak m_C$-進位相に関して形式的に滑らかである。
\end{proof}

\begin{remark}
\label{more-algebra-remark-what-does-it-mean}
Lemma \ref{more-algebra-lemma-lift-fs} の主張はかなり強い。実際、次のような図式が
与えられたとする。
$$
\xymatrix{
& B \\
A \ar[r] & A' \ar[u]
}
$$
ここで、Noether 完備局所環の局所準同型からなる図式であり、
$A \to A'$ は剰余体
$k = A/\mathfrak m_A = A'/\mathfrak m_{A'}$ 上の同型を誘導し、
$B \otimes_{A'} k$ は $k$ 上形式的に滑らかであるとする。
このとき、次のような可換図式へ拡張できる。
$$
\xymatrix{
C \ar[r] & B \\
A \ar[r] \ar[u] & A' \ar[u]
}
$$
これは Noether 完備局所環の局所準同型からなり、$A \to C$ は
$\mathfrak m_C$-進位相に関して形式的に滑らかで、
$C \otimes_A k \cong B \otimes_{A'} k$ を満たす。
実際、$B \otimes_{A'} k$ を $k$ 上持ち上げる Lemma
\ref{more-algebra-lemma-lift-fs} のような $A \to C$ を選ぶ。形式的滑らかさにより、
Lemma \ref{more-algebra-lemma-lift-continuous} から $C \to B$ の矢印を得る。
$C \otimes_A^\wedge A'$ を、イデアル $C \otimes_A \mathfrak m_{A'}$ に
関する $C \otimes_A A'$ の完備化とする。$C \otimes_A^\wedge A'$ は
Noether 完備局所環であり（Algebra, Lemma
\ref{algebra-lemma-completion-Noetherian}）、$A'$ 上平坦である
（Algebra, Lemma \ref{algebra-lemma-flat-module-powers}）ことに注意せよ。
さらに次が成り立つ。
\begin{enumerate}
\item $C \otimes_A^\wedge A' \to B$ は全射である。
\item $A \to A'$ が全射ならば、$C \to B$ は全射である。
\item $A \to A'$ が有限ならば、$C \to B$ は有限である。
\item $A' \to B$ が平坦ならば、$C \otimes_A^\wedge A' \cong B$ である。
\end{enumerate}
実際、冪零イデアルに対する中山の補題（Algebra, Lemma
\ref{algebra-lemma-NAK}）により、$C \otimes_A k \cong B \otimes_{A'} k$ から
すべての $n$ に対して
$$
C \otimes_A A'/\mathfrak m_{A'}^n \to B/\mathfrak m_{A'}^nB
$$
が全射であることが分かる。これで (1) が示される。(2) と (3) は (1) から
従い、(4) は Algebra, Lemma \ref{algebra-lemma-mod-injective} から従う。
\end{remark}
\section{正則環写像}
\label{more-algebra-section-regular}

\noindent
$k$ を体とする。Noether $k$-代数 $A$ が $k$ 上 {\it 幾何学的に正則}
であるとは、体 $k$ の任意の有限純非分離拡大 $k'$ に対して
$A \otimes_k k'$ が正則であることと同値であることを思い出そう。
Algebra, Definition \ref{algebra-definition-geometrically-regular} を参照せよ。
さらにこの条件が成り立つなら、任意の有限生成体拡大 $k'/k$ に対しても
$A \otimes_k k'$ は正則である。Algebra, Lemma
\ref{algebra-lemma-geometrically-regular} を参照せよ。
この概念を次の定義で用いる。

\begin{definition}
\label{more-algebra-definition-regular}
環写像 $R \to \Lambda$ が {\it 正則} であるとは、それが平坦であり、
任意の素イデアル $\mathfrak p \subset R$ に対してファイバー環
$$
\Lambda \otimes_R \kappa(\mathfrak p) =
\Lambda_\mathfrak p/\mathfrak p\Lambda_\mathfrak p
$$
が Noether で $\kappa(\mathfrak p)$ 上幾何学的に正則であることをいう。
\end{definition}

\noindent
$R \to \Lambda$ が環写像で $\Lambda$ が Noether ならば、
ファイバー環は常に Noether である。

\begin{lemma}[正則性は局所的性質]
\label{more-algebra-lemma-regular-local}
$R \to \Lambda$ を $\Lambda$ が Noether である環写像とする。
次の条件は同値である。
\begin{enumerate}
\item $R \to \Lambda$ は正則である。
\item $\mathfrak p \subset R$ の上にあるすべての素イデアル
$\mathfrak q \subset \Lambda$ に対して、$R_\mathfrak p \to \Lambda_\mathfrak q$ は正則であり、
\item $R$ の $\mathfrak m$ の上にあるすべての極大イデアル
$\mathfrak m' \subset \Lambda$ に対して、
$R_\mathfrak m \to \Lambda_{\mathfrak m'}$ は正則である。
\end{enumerate}
\end{lemma}

\begin{proof}
これは、Noether 環が正則であることとすべての局所環が正則局所環であることが同値であり、
環写像が平坦であることと誘導されるすべての局所環の写像が平坦であることが同値だからである。
Algebra, Definition \ref{algebra-definition-regular} および Algebra, Lemma
\ref{algebra-lemma-flat-localization} を参照せよ。
\end{proof}

\begin{lemma}[正則写像と基底変換]
\label{more-algebra-lemma-regular-base-change}
$R \to \Lambda$ を正則環写像とする。
任意の有限型環写像 $R \to R'$ に対して、基底変換
$R' \to \Lambda \otimes_R R'$ も正則である。
\end{lemma}

\begin{proof}
平坦性は任意の基底変換で保たれる。Algebra, Lemma
\ref{algebra-lemma-flat-base-change} を参照せよ。
$\mathfrak p \subset R$ の上にある素イデアル
$\mathfrak p' \subset R'$ を考える。剰余体拡大
$\kappa(\mathfrak p')/\kappa(\mathfrak p)$ は、$R'$ が $R$ 上有限型なので有限生成である。
したがってファイバー環
$$
(\Lambda \otimes_R R') \otimes_{R'} \kappa(\mathfrak p') =
\Lambda \otimes_R \kappa(\mathfrak p) \otimes_{\kappa(\mathfrak p)} 
\kappa(\mathfrak p')
$$
は、Algebra, Lemma \ref{algebra-lemma-Noetherian-field-extension} と
$R \to \Lambda$ のファイバー環に関する仮定により Noether である。
ファイバーの幾何学的正則性は Algebra, Lemma
\ref{algebra-lemma-geometrically-regular} により保たれる。
\end{proof}

\begin{lemma}[正則写像の合成]
\label{more-algebra-lemma-regular-composition}
$A \to B$ および $B \to C$ を正則環写像とする。
$A \to C$ のファイバー環が Noether ならば、$A \to C$ は正則である。
\end{lemma}

\begin{proof}
$\mathfrak p \subset A$ の素イデアルをとる。
$\kappa(\mathfrak p) \subset k$ を有限純非分離拡大とする。
$C \otimes_A k$ が正則であることを示せばよい。
Lemma \ref{more-algebra-lemma-regular-base-change} により $A = k$ と仮定してよく、
$C$ が正則であることの証明に帰着する。
仮定は $B$ が正則で、$B \to C$ が平坦かつ正則なファイバーをもつことである。
したがって Algebra, Lemma \ref{algebra-lemma-flat-over-regular-with-regular-fibre} により
$C$ は正則である。詳細は省略する。
\end{proof}

\begin{lemma}
\label{more-algebra-lemma-colimit-smooth-regular}
$R$ を環とする。$(A_i, \varphi_{ii'})$ を滑らかな $R$-代数の有向系とする。
$\Lambda = \colim A_i$ とおく。すべての $\mathfrak p \subset R$ に対して
ファイバー環 $\Lambda \otimes_R \kappa(\mathfrak p)$ が Noether ならば、
$R \to \Lambda$ は正則である。
\end{lemma}

\begin{proof}
Algebra, Lemmas \ref{algebra-lemma-colimit-flat} および
\ref{algebra-lemma-smooth-syntomic} により、$\Lambda$ は $R$ 上平坦であることに注意する。
$\kappa(\mathfrak p) \subset k$ を有限純非分離拡大とする。すると
$$
\Lambda \otimes_R \kappa(\mathfrak p) \otimes_{\kappa(\mathfrak p)} k =
\Lambda \otimes_R k = \colim A_i \otimes_R k
$$
は滑らかな $k$-代数の余極限である。Algebra, Lemma
\ref{algebra-lemma-base-change-smooth} を参照せよ。
滑らかな $k$-代数の各局所環は Algebra, Lemma
\ref{algebra-lemma-characterize-smooth-over-field} により正則である。
したがって Algebra, Lemma \ref{algebra-lemma-colimit-regular} により、
$\Lambda \otimes_R k$ のすべての局所環は正則である。
これで補題が示された。
\end{proof}

\noindent
体拡大がいつ正則環写像を定めるかを見よう。

\begin{lemma}
\label{more-algebra-lemma-regular-field-extension}
$K/k$ を体拡大とする。このとき $k \to K$ が正則環写像であることと、
$K$ が $k$ の分離拡大であることは同値である。
\end{lemma}

\begin{proof}
$k \to K$ が正則ならば、$K$ は $k$ 上幾何学的に被約である。
ゆえに Algebra, Proposition
\ref{algebra-proposition-characterize-separable-field-extensions} により、
$K$ は $k$ 上分離的である。
逆に $K/k$ が分離的ならば、$K$ は滑らかな $k$-代数の余極限である。
Algebra, Lemma \ref{algebra-lemma-colimit-syntomic} を参照せよ。
したがって Lemma \ref{more-algebra-lemma-colimit-smooth-regular} により正則である。
\end{proof}

\begin{lemma}
\label{more-algebra-lemma-regular-permanence}
$A \to B \to C$ を環写像とする。$A \to C$ が正則で、
$B \to C$ が平坦かつスペクトル上全射ならば、$A \to B$ は正則である。
\end{lemma}

\begin{proof}
Algebra, Lemma \ref{algebra-lemma-flat-permanence} により $A \to B$ は平坦である。
$\mathfrak p \subset A$ の素イデアルをとる。環写像
$B \otimes_A \kappa(\mathfrak p) \to C \otimes_A \kappa(\mathfrak p)$
は平坦かつスペクトル上全射である。したがって Algebra, Lemma
\ref{algebra-lemma-geometrically-regular-descent} により
$B \otimes_A \kappa(\mathfrak p)$ は幾何学的に正則である。
\end{proof}



\section{正則環写像に沿った性質の上昇}
\label{more-algebra-section-ascending-properties}

\noindent
この節は Algebra, Section \ref{algebra-section-ascending-properties} の類似であるが、
ここでは環写像 $R \to S$ が正則である。
\begin{lemma}
\label{more-algebra-lemma-reduced-goes-up}
環写像 $\varphi : R \to S$ をとる。次を仮定する。
\begin{enumerate}
\item $\varphi$ は正則である。
\item $S$ は Noether である。
\item $R$ は Noether かつ被約である。
\end{enumerate}
このとき $S$ は被約である。
\end{lemma}

\begin{proof}
Noether 環が被約であることは性質 $(S_1)$ および $(R_0)$ をもつことと同値である。
Algebra, Lemma \ref{algebra-lemma-criterion-reduced} を参照せよ。
したがって Algebra, Lemmas \ref{algebra-lemma-Sk-goes-up} および
\ref{algebra-lemma-Rk-goes-up} を適用できる。
\end{proof}

\begin{lemma}
\label{more-algebra-lemma-normal-goes-up}
環写像 $\varphi : R \to S$ をとる。次を仮定する。
\begin{enumerate}
\item $\varphi$ は正則である。
\item $S$ は Noether である。
\item $R$ は Noether かつ正規である。
\end{enumerate}
このとき $S$ は正規である。
\end{lemma}

\begin{proof}
Noether 環が正規であることは性質 $(S_2)$ および $(R_1)$ をもつことと同値である。
Algebra, Lemma \ref{algebra-lemma-criterion-normal} を参照せよ。
したがって Algebra, Lemmas \ref{algebra-lemma-Sk-goes-up} および
\ref{algebra-lemma-Rk-goes-up} を適用できる。
\end{proof}

\begin{lemma}
\label{more-algebra-lemma-regular-goes-up}
環写像 $\varphi : R \to S$ をとる。次を仮定する。
\begin{enumerate}
\item $\varphi$ は正則である。
\item $S$ は Noether である。
\item $R$ は Noether かつ正則である。
\end{enumerate}
このとき $S$ は正則である。
\end{lemma}

\begin{proof}
Noether 環が正則であることはすべての $k$ に対する性質 $(R_k)$ をもつことと同値である。
したがって Algebra, Lemma \ref{algebra-lemma-Rk-goes-up} を適用できる。
\end{proof}

\begin{lemma}
\label{more-algebra-lemma-CM-goes-up}
環写像 $\varphi : R \to S$ をとる。次を仮定する。
\begin{enumerate}
\item $\varphi$ は正則である。
\item $S$ は Noether である。
\item $R$ は Noether かつ Cohen-Macaulay である。
\end{enumerate}
このとき $S$ は Cohen-Macaulay である。
\end{lemma}

\begin{proof}
Noether 環が Cohen-Macaulay であることはすべての $k$ に対する性質 $(S_k)$ をもつことと同値である。
したがって Algebra, Lemma \ref{algebra-lemma-Sk-goes-up} を適用できる。
\end{proof}



\section{完備化の下での性質の保存}
\label{more-algebra-section-permanence-completion}

\noindent
Noether 局所環 $(A, \mathfrak m)$ が与えられたとき、$A^\wedge$ を、$A$ の $\mathfrak m$ に関する
完備化と呼ぶ。以下では、$A^\wedge$ が極大イデアル
$\mathfrak m^\wedge = \mathfrak m A^\wedge$ をもつ Noether 完備局所環であり、
$A \to A^\wedge$ が忠実平坦であることを断りなく用いる。
Algebra, Lemmas \ref{algebra-lemma-completion-Noetherian-Noetherian},
\ref{algebra-lemma-completion-complete}, および
\ref{algebra-lemma-completion-faithfully-flat} を参照せよ。

\begin{lemma}
\label{more-algebra-lemma-completion-dimension}
$A$ を Noether 局所環とする。
このとき $\dim(A) = \dim(A^\wedge)$ である。
\end{lemma}

\begin{proof}
Algebra, Lemma \ref{algebra-lemma-completion-complete} により、写像
$A \to A^\wedge$ は
$A/\mathfrak m^n = A^\wedge/(\mathfrak m^\wedge)^n$（$n \geq 1$）という同型を誘導する。
Algebra, Lemma \ref{algebra-lemma-pushdown-module} により、これは
$$
\text{length}_A(A/\mathfrak m^n) =
\text{length}_{A^\wedge}(A^\wedge/(\mathfrak m^\wedge)^n)
$$
をすべての $n \geq 1$ について意味する。したがって
$d(A) = d(A^\wedge)$ であり、Algebra, Proposition
\ref{algebra-proposition-dimension} により結論を得る。
別証明として Algebra, Lemma
\ref{algebra-lemma-dimension-base-fibre-equals-total} を用いてもよい。
\end{proof}

\begin{lemma}
\label{more-algebra-lemma-completion-depth}
$A$ を Noether 局所環とする。このとき
$\text{depth}(A) = \text{depth}(A^\wedge)$ である。
\end{lemma}

\begin{proof}
Algebra, Lemma \ref{algebra-lemma-apply-grothendieck} を参照せよ。
\end{proof}

\begin{lemma}
\label{more-algebra-lemma-completion-CM}
$A$ を Noether 局所環とする。
$A$ が Cohen-Macaulay であることと $A^\wedge$ がそうであることは同値である。
\end{lemma}

\begin{proof}
局所環 $A$ が Cohen-Macaulay であることと $\dim(A) = \text{depth}(A)$ は同値である。
これら二つの不変量はいずれも完備化で保存される（Lemmas
\ref{more-algebra-lemma-completion-dimension} および \ref{more-algebra-lemma-completion-depth}）ので、主張が従う。
\end{proof}

\begin{lemma}
\label{more-algebra-lemma-completion-regular}
$A$ を Noether 局所環とする。
$A$ が正則であることと $A^\wedge$ が正則であることは同値である。
\end{lemma}

\begin{proof}
$A^\wedge$ が正則ならば、Algebra, Lemma
\ref{algebra-lemma-flat-under-regular} により $A$ は正則である。
$A$ が正則であると仮定する。$A$ の極大イデアルを $\mathfrak m$ とする。このとき
$\dim_{\kappa(\mathfrak m)} \mathfrak m/\mathfrak m^2 =
\dim(A) = \dim(A^\wedge)$（Lemma \ref{more-algebra-lemma-completion-dimension}）である。
一方、$\mathfrak mA^\wedge$ は $A^\wedge$ の極大イデアルであり、したがって
$\mathfrak m_{A^\wedge}$ は高々 $\dim(A^\wedge)$ 個の元で生成される。
ゆえに $A^\wedge$ は正則である。（Algebra, Lemma
\ref{algebra-lemma-flat-over-regular-with-regular-fibre} を用いてもよい。）
\end{proof}

\begin{lemma}
\label{more-algebra-lemma-completion-dvr}
$A$ を Noether 局所環とする。
$A$ が離散付値環であることと $A^\wedge$ がそうであることは同値である。
\end{lemma}

\begin{proof}
これは Lemmas \ref{more-algebra-lemma-completion-dimension} および
\ref{more-algebra-lemma-completion-regular} と Algebra, Lemma
\ref{algebra-lemma-characterize-dvr} から従う。
\end{proof}

\begin{lemma}
\label{more-algebra-lemma-completion-reduced}
$A$ を Noether 局所環とする。
\begin{enumerate}
\item $A^\wedge$ が被約ならば、$A$ も被約である。
\item 一般には、$A$ が被約でも $A^\wedge$ が被約とは限らない。
\item $A$ が Nagata ならば、$A$ が被約であることと $A^\wedge$ が
被約であることは同値である。
\end{enumerate}
\end{lemma}

\begin{proof}
$A \to A^\wedge$ は忠実平坦なので、(1) は Algebra, Lemma
\ref{algebra-lemma-descent-reduced} による。(2) は Algebra, Example
\ref{algebra-example-bad-dvr-char-p} を参照せよ（標数零にも例があり、
Algebra, Remark \ref{algebra-remark-resolution-dim-1} を参照せよ）。
(3) は Algebra, Lemmas
\ref{algebra-lemma-local-nagata-domain-analytically-unramified} および
\ref{algebra-lemma-analytically-unramified-easy} による。
\end{proof}

\begin{lemma}
\label{more-algebra-lemma-completion-normal}
$A$ を Noether 局所環とする。$A^\wedge$ が正規ならば、$A$ も正規である。
\end{lemma}

\begin{proof}
$A \to A^\wedge$ は忠実平坦なので、Algebra, Lemma
\ref{algebra-lemma-descent-normal} により従う。
\end{proof}

\begin{lemma}
\label{more-algebra-lemma-flat-completion}
$A \to B$ を Noether 局所環の局所準同型とする。
誘導される完備化の写像 $A^\wedge \to B^\wedge$ が平坦であることと、
$A \to B$ が平坦であることは同値である。
\end{lemma}

\begin{proof}
可換図式
$$
\xymatrix{
A^\wedge \ar[r] & B^\wedge \\
A \ar[r] \ar[u] & B \ar[u]
}
$$
を考える。鉛直矢印は忠実平坦である。
$A^\wedge \to B^\wedge$ が平坦であると仮定する。このとき $A \to B^\wedge$ は平坦である。
したがって Algebra, Lemma \ref{algebra-lemma-flatness-descends-more-general} により
$B$ は $A$ 上平坦である。

\medskip\noindent
$A \to B$ が平坦であると仮定する。このとき $A \to B^\wedge$ は平坦である。
ゆえに $B^\wedge/\mathfrak m_A^n B^\wedge$ はすべての $n \geq 1$ に対して
$A/\mathfrak m_A^n$ 上平坦である。$\mathfrak m_A^n A^\wedge$ は
$A^\wedge$ の極大イデアル $\mathfrak m_A^\wedge$ の $n$ 乗であり、
$A/\mathfrak m_A^n = A^\wedge/(\mathfrak m_A^\wedge)^n$ であることに注意せよ。
したがって Algebra, Lemma \ref{algebra-lemma-flat-module-powers} を
（$R = A^\wedge$, $I = \mathfrak m_A^\wedge$,
$S = B^\wedge$, $M = S$ として）適用すれば、$B^\wedge$ は $A^\wedge$ 上平坦である。
\end{proof}

\begin{lemma}
\label{more-algebra-lemma-flat-unramified}
$A \to B$ を Noether 局所環の平坦な局所準同型とし、
$\mathfrak m_A B = \mathfrak m_B$ および
$\kappa(\mathfrak m_A) = \kappa(\mathfrak m_B)$ を満たすとする。
このとき $A \to B$ は完備化の同型 $A^\wedge \to B^\wedge$ を誘導する。
\end{lemma}

\begin{proof}
Algebra, Lemma \ref{algebra-lemma-finite-after-completion} により、
$B^\wedge$ は $B$ の $\mathfrak m_A$-進完備化であり、$A^\wedge \to B^\wedge$ は有限である。
$A \to B$ は平坦なので $\text{Tor}_1^A(B, \kappa(\mathfrak m_A)) = 0$ である。
したがって Lemma \ref{more-algebra-lemma-flat-after-completion} により、$B^\wedge$ は $A^\wedge$ 上平坦である。
ゆえに Algebra, Lemma \ref{algebra-lemma-finite-flat-local} により
$B^\wedge$ は $A^\wedge$ 上自由である。
$A^\wedge \to B^\wedge$ は仮定（および Algebra, Lemma
\ref{algebra-lemma-hathat-finitely-generated}）により同型
$\kappa(\mathfrak m_A) = A^\wedge/\mathfrak m_A A^\wedge \to
B^\wedge/\mathfrak m_A B^\wedge = B^\wedge/\mathfrak m_B B^\wedge =
\kappa(\mathfrak m_B)$ を誘導するので、$B^\wedge$ は階数 $1$ の自由加群である。
したがって $A^\wedge \to B^\wedge$ は同型である。
\end{proof}




\section{\'Etale 写像の下での性質の保存}
\label{more-algebra-section-permanence-etale}

\noindent
この節では \'etale 環写像 $\varphi : A \to B$ を考え、
$A$ の性質が $B$ に継承される場合、および $\mathfrak q$ における $B$ の局所環の性質が
$\mathfrak p = \varphi^{-1}(\mathfrak q)$ における $A$ の局所環に継承される場合を調べる。
基本的には、この節は以前の結果を再検討し集めるもので、新しい内容は加えない。

\medskip\noindent
以下では断りなく、\'etale 環写像は平坦であること（Algebra, Lemma
\ref{algebra-lemma-etale}）、および局所環の平坦な局所準同型は忠実平坦であること
（Algebra, Lemma \ref{algebra-lemma-local-flat-ff}）を用いる。
\begin{lemma}
\label{more-algebra-lemma-Noetherian-etale-extension}
$A \to B$ が \'etale 環写像で、$\mathfrak q$ が $\mathfrak p \subset A$ の上にある
$B$ の素イデアルならば、$A_{\mathfrak p}$ が Noether であることと
$B_{\mathfrak q}$ が Noether であることは同値である。
\end{lemma}

\begin{proof}
$A_\mathfrak p \to B_\mathfrak q$ は忠実平坦なので、$B_\mathfrak q$ が Noether なら
$A_\mathfrak p$ も Noether である。Algebra, Lemma
\ref{algebra-lemma-descent-Noetherian} を参照せよ。
逆に $A_\mathfrak p$ が Noether なら、$B_\mathfrak q$ は有限型
$A_\mathfrak p$-代数の局所化なので Noether である。
\end{proof}

\begin{lemma}
\label{more-algebra-lemma-dimension-etale-extension}
$A \to B$ が \'etale 環写像で、$\mathfrak q$ が $\mathfrak p \subset A$ の上にある
$B$ の素イデアルならば、
$\dim(A_{\mathfrak p}) = \dim(B_{\mathfrak q})$ である。
\end{lemma}

\begin{proof}
$A_{\mathfrak p} \to B_{\mathfrak q}$ は平坦なので going down をもつ。
したがって不等式
$\dim(A_{\mathfrak p}) \leq \dim(B_{\mathfrak q})$ が成り立つ。Algebra, Lemma
\ref{algebra-lemma-dimension-going-up} を参照せよ。
一方、
$\mathfrak q_0 \subset \mathfrak q_1 \subset \ldots \subset \mathfrak q_n$
を $B_{\mathfrak q}$ の素イデアルの鎖とする。対応する素イデアルの列
$\mathfrak p_0 \subset \mathfrak p_1 \subset \ldots \subset \mathfrak p_n$
（$\mathfrak p_i = \mathfrak q_i \cap A_{\mathfrak p}$）も鎖である（すなわち列の中に等号はない）。
なぜなら \'etale 環写像は準有限だからである（Algebra, Lemma
\ref{algebra-lemma-etale-quasi-finite}）。また準有限環写像は、定義により、
スペクトルの間に離散ファイバーをもつ写像を誘導する。
したがって所望どおり $\dim(A_{\mathfrak p}) \geq \dim(B_{\mathfrak q})$ である。
\end{proof}

\begin{lemma}
\label{more-algebra-lemma-regular-etale-extension}
$A \to B$ が \'etale 環写像で、$\mathfrak q$ が $\mathfrak p \subset A$ の上にある
$B$ の素イデアルならば、$A_{\mathfrak p}$ が正則であることと
$B_{\mathfrak q}$ が正則であることは同値である。
\end{lemma}

\begin{proof}
同値性を示すため、Lemma \ref{more-algebra-lemma-Noetherian-etale-extension} により
$A_\mathfrak p$ と $B_\mathfrak q$ の双方が Noether であると仮定してよい。
$x_1, \ldots, x_t \in \mathfrak pA_\mathfrak p$ を最小生成系としてとる。
$A_\mathfrak p \to B_\mathfrak q$ は忠実平坦なので、$B_\mathfrak q$ における像
$y_1, \ldots, y_t$ は
$\mathfrak pB_\mathfrak q = \mathfrak q B_\mathfrak q$ の最小生成系をなす。
Algebra, Lemma \ref{algebra-lemma-etale-at-prime} を参照せよ。
$A_\mathfrak p$ の正則性は定義により $t = \dim(A_\mathfrak p)$ を意味し、
$B_\mathfrak q$ についても同様である。したがって
Lemma \ref{more-algebra-lemma-dimension-etale-extension} の
$\dim(A_\mathfrak p) = \dim(B_\mathfrak q)$ から補題が従う。
\end{proof}

\begin{lemma}
\label{more-algebra-lemma-Dedekind-etale-extension}
$A \to B$ が \'etale 環写像で $A$ が Dedekind 整域ならば、$B$ は Dedekind 整域の有限積である。
特に、$\mathfrak q \subset B$ が極大であるときの局所化 $B_\mathfrak q$ は
離散付値環である。
\end{lemma}

\begin{proof}
局所環についての主張は Lemmas \ref{more-algebra-lemma-dimension-etale-extension} および
\ref{more-algebra-lemma-regular-etale-extension} と Algebra, Lemma
\ref{algebra-lemma-characterize-dvr} から従う。
したがって $B$ は次元 $1$ の Noether 正規環である。
Algebra, Lemma \ref{algebra-lemma-characterize-reduced-ring-normal} により、
$B$ は次元 $1$ の正規整域の有限積である。これらは Algebra, Lemma
\ref{algebra-lemma-characterize-Dedekind} により Dedekind 整域である。
\end{proof}



\section{Hensel 化の下での性質の保存}
\label{more-algebra-section-permanence-henselization}

\noindent
局所環 $R$ が与えられたとき、$R$ の Hensel 化を $R^h$、強 Hensel 化を
$R^{sh}$ と書く。Algebra, Definition \ref{algebra-definition-henselization} を参照せよ。
この節で示すように、$R$ の多くの性質は $R^h$ および $R^{sh}$ に反映される。

\begin{lemma}
\label{more-algebra-lemma-dumb-properties-henselization}
$(R, \mathfrak m, \kappa)$ を局所環とする。このとき次が成り立つ。
\begin{enumerate}
\item $R \to R^h \to R^{sh}$ は忠実平坦な環写像である。
\item $\mathfrak m R^h = \mathfrak m^h$ および
$\mathfrak m R^{sh} = \mathfrak m^h R^{sh} = \mathfrak m^{sh}$ である。
\item すべての $n$ に対して $R/\mathfrak m^n = R^h/\mathfrak m^nR^h$ である。
\item $x_i \in R^{sh}$ であって、$R^{sh}/\mathfrak m^nR^{sh}$ が
$x_i \bmod \mathfrak m^nR^{sh}$ を基底とする $R/\mathfrak m^n$-自由加群となるものが存在する。
\end{enumerate}
\end{lemma}

\begin{proof}
構成により $R^h$ は \'etale $R$-代数の余極限である。Algebra, Lemma
\ref{algebra-lemma-henselization} を参照せよ。\'etale 環写像は平坦である
（Algebra, Lemma \ref{algebra-lemma-etale}）から、Algebra, Lemma
\ref{algebra-lemma-colimit-flat} により $R^h$ は $R$ 上平坦である。
平坦な局所環準同型は忠実平坦である（Algebra, Lemma
\ref{algebra-lemma-local-flat-ff}）ので、$R \to R^h$ は忠実平坦である。
$R^h \to R^{sh}$ は有限 \'etale 環写像の余極限である。Algebra, Lemma
\ref{algebra-lemma-strict-henselization} の証明を参照せよ。
したがって上と同じ議論により $R^h \to R^{sh}$ も忠実平坦である。

\medskip\noindent
(2) は Algebra, Lemmas \ref{algebra-lemma-henselization} および
\ref{algebra-lemma-strict-henselization} から従う。(3) は Algebra, Lemma
\ref{algebra-lemma-local-artinian-basis-when-flat} による。実際、
$R/\mathfrak m \to R^h/\mathfrak mR^h$ は同型であり、
$R/\mathfrak m^n \to R^h/\mathfrak m^nR^h$ は平坦な環写像 $R \to R^h$ の
基底変換として平坦である（Algebra, Lemma \ref{algebra-lemma-flat-base-change}）。
$\kappa^{sep}$ を $R^{sh}$ の剰余体とする（これは $\kappa$ の分離代数閉包である）。
$\kappa$-ベクトル空間 $\kappa^{sep}$ の基底の像となるような
$x_i \in R^{sh}$ を選ぶ。(4) は Algebra, Lemma
\ref{algebra-lemma-local-artinian-basis-when-flat} から、上とまったく同様に従う。
\end{proof}

\begin{lemma}
\label{more-algebra-lemma-henselization-formally-smooth}
$(R, \mathfrak m, \kappa)$ を局所環とする。このとき
\begin{enumerate}
\item $R \to R^h$、$R^h \to R^{sh}$、および $R \to R^{sh}$ は形式的 \'etale である。
\item $R \to R^h$、$R^h \to R^{sh}$、resp.\ $R \to R^{sh}$ は、それぞれ
$\mathfrak m^h$、$\mathfrak m^{sh}$、resp.\ $\mathfrak m^{sh}$-位相に関して形式的に滑らかである。
\end{enumerate}
\end{lemma}

\begin{proof}
(1) は、$R^h$ と $R^{sh}$ が（構成により）\'etale 代数の有向余極限であり、
\'etale 代数が形式的 \'etale であること（Algebra, Lemma
\ref{algebra-lemma-formally-etale-etale}）、および形式的 \'etale 代数の余極限が形式的 \'etale
であること（Algebra, Lemma \ref{algebra-lemma-colimit-formally-etale}）から従う。
(2) は、形式的 \'etale 環写像が形式的に滑らかであることと Lemma
\ref{more-algebra-lemma-formally-smooth} から従う。
\end{proof}

\begin{lemma}
\label{more-algebra-lemma-henselization-noetherian}
\begin{reference}
\cite[IV, Theorem 18.6.6 and Proposition 18.8.8]{EGA}
\end{reference}
$R$ を局所環とする。次の条件は同値である。
\begin{enumerate}
\item $R$ は Noether である。
\item $R^h$ は Noether である。
\item $R^{sh}$ は Noether である。
\end{enumerate}
この場合、次が成り立つ。
\begin{enumerate}
\item[(a)] $(R^h)^\wedge$ および $(R^{sh})^\wedge$ は Noether 完備局所環である。
\item[(b)] $R^\wedge \to (R^h)^\wedge$ は同型である。
\item[(c)] $R^h \to (R^h)^\wedge$ および $R^{sh} \to (R^{sh})^\wedge$ は平坦である。
\item[(d)] $R^\wedge \to (R^{sh})^\wedge$ は
$\mathfrak m_{(R^{sh})^\wedge}$-進位相に関して形式的に滑らかである。
\item[(e)] $(R^\wedge)^{sh} = R^\wedge \otimes_{R^h} R^{sh}$ である。
\item[(f)] $((R^\wedge)^{sh})^\wedge = (R^{sh})^\wedge$ である。
\end{enumerate}
\end{lemma}
\begin{proof}
$R \to R^h \to R^{sh}$ は忠実平坦である（Lemma
\ref{more-algebra-lemma-dumb-properties-henselization}）から、$R^h$ または $R^{sh}$ が Noether なら
$R$ も Noether である。Algebra, Lemma \ref{algebra-lemma-descent-Noetherian} を参照せよ。
以下の証明では $R$ が Noether であると仮定する。

\medskip\noindent
$\mathfrak m \subset R$ は有限生成なので、$\mathfrak m^h = \mathfrak m R^h$ および
$\mathfrak m^{sh} = \mathfrak mR^{sh}$ も有限生成である。
Lemma \ref{more-algebra-lemma-dumb-properties-henselization} を参照せよ。
したがって $(R^h)^\wedge$ および $(R^{sh})^\wedge$ は Algebra, Lemma
\ref{algebra-lemma-complete-local-ring-Noetherian} により Noether である。
これで (a) が示された。

\medskip\noindent
(b) は Lemma \ref{more-algebra-lemma-dumb-properties-henselization} から直ちに従う。
特に Algebra, Lemma \ref{algebra-lemma-completion-faithfully-flat} により
$(R^h)^\wedge$ は $R$ 上平坦である。

\medskip\noindent
次に $R^h \to (R^h)^\wedge$ が平坦であることを示す。
$R^h = \colim_i R_i$ を \'etale $R$-代数の局所化からなる有向余極限と書く。
Algebra, Lemma \ref{algebra-lemma-colimit-rings-flat} により、$(R^h)^\wedge$ が各 $R_i$ 上平坦ならば
$R^h \to (R^h)^\wedge$ は平坦である。構成により $R^h = R_i^h$ であることに注意する。
したがって (b) により $R_i^\wedge = (R^h)^\wedge$ は $R_i$ 上平坦であり、所望の結果を得る。
(c) の証明を終えるため、$R^{sh} \to (R^{sh})^\wedge$ が平坦であることを示す。
上と同様の極限論法により、$(R^{sh})^\wedge$ が $R$ 上平坦であることを示せば十分である。
Lemma \ref{more-algebra-lemma-dumb-properties-henselization} から、$(R^{sh})^\wedge$ は
自由 $R$-加群の完備化であることが分かる。Lemma \ref{more-algebra-lemma-completed-direct-sum-flat} により、
これは $R$ 上平坦である。これで (c) の証明が終わる。

\medskip\noindent
ここまでで (c) が成り立ち、$(R^h)^\wedge$ と $(R^{sh})^\wedge$ が Noether であることが分かっている。
Algebra, Lemma \ref{algebra-lemma-descent-Noetherian} により、$R^h$ と $R^{sh}$ は Noether である。

\medskip\noindent
(d) は Lemma \ref{more-algebra-lemma-henselization-formally-smooth} および
Lemma \ref{more-algebra-lemma-formally-smooth-completion} から従う。

\medskip\noindent
(e) は Algebra, Lemma \ref{algebra-lemma-sh-from-h-map} と、
$R^\wedge$ が Hensel 環であること（Algebra, Lemma
\ref{algebra-lemma-complete-henselian}）から従う。

\medskip\noindent
(f) の証明。 (e) により $R^{sh} \to (R^\wedge)^{sh}$ という写像があり、
完備化により $(R^{sh})^\wedge \to ((R^\wedge)^{sh})^\wedge$ を誘導する。
(e) により $R^\wedge \to (R^{sh})^\wedge$ という写像もある。
$(R^{sh})^\wedge$ は（上で見たように）強 Hensel なので、この写像は Algebra, Lemma
\ref{algebra-lemma-strictly-henselian-functorial} により
$(R^\wedge)^{sh} \to (R^{sh})^\wedge$ を誘導する。
完備化して $((R^\wedge)^{sh})^\wedge \to (R^{sh})^\wedge$ を得る。
これら二つの写像が互いに逆であることの検証は省略する。
\end{proof}
\begin{lemma}
\label{more-algebra-lemma-henselization-reduced}
\begin{slogan}
被約性は（強）Hensel 化へ移る。
\end{slogan}
局所環 $R$ をとる。次の条件は同値である。$R$ が被約であること、
$R$ の Hensel 化 $R^h$ が被約であること、および $R$ の強 Hensel 化
$R^{sh}$ が被約であること。
\end{lemma}

\begin{proof}
環写像 $R \to R^h \to R^{sh}$ は忠実平坦である。
したがって含意の一方向は Algebra, Lemma
\ref{algebra-lemma-descent-reduced} から従う。
逆に $R$ が被約であると仮定する。$R^h$ と $R^{sh}$ は、\'etale、したがって滑らかな
$R$-代数のフィルター付き余極限であるから、結果は Algebra, Lemma
\ref{algebra-lemma-reduced-goes-up} から従う。
\end{proof}

\begin{lemma}
\label{more-algebra-lemma-henselization-nil}
局所環 $R$ をとる。$nil(R)$ を $R$ の冪零元からなるイデアルとする。
このとき $nil(R)R^h = nil(R^h)$ および
$nil(R)R^{sh} = nil(R^{sh})$ である。
\end{lemma}

\begin{proof}
$nil(R)$ は、冪零元からなるイデアルのうち最大のものであり、商
$R/nil(R)$ が被約になることに注意する。また Algebra, Lemma
\ref{algebra-lemma-locally-nilpotent} により $nil(R)R^h$ は冪零元からなる。
さらに Algebra, Lemma \ref{algebra-lemma-quotient-henselization} により
$R^h/nil(R) R^h$ は $R/nil(R)$ の Hensel 化であることに注意する。
したがって Lemma \ref{more-algebra-lemma-henselization-reduced} により
$R^h/nil(R)R^h$ は被約である。ゆえに所望どおり
$nil(R) R^h = nil(R^h)$ である。
強 Hensel 化についても同様であるが、Algebra, Lemma
\ref{algebra-lemma-quotient-strict-henselization} を用いる。
\end{proof}

\begin{lemma}
\label{more-algebra-lemma-henselization-normal}
局所環 $R$ をとる。次の条件は同値である。$R$ が正規整域であること、
$R$ の Hensel 化 $R^h$ が正規整域であること、および $R$ の強 Hensel 化
$R^{sh}$ が正規整域であること。
\end{lemma}

\begin{proof}
まず、局所環が正規であることと正規整域であることは同値であることに注意する
（Algebra, Definition \ref{algebra-definition-ring-normal} を参照せよ）。
環写像 $R \to R^h \to R^{sh}$ は忠実平坦である。
したがって含意の一方向は Algebra, Lemma
\ref{algebra-lemma-descent-normal} から従う。
逆に $R$ が正規であると仮定する。$R^h$ と $R^{sh}$ は、\'etale、したがって滑らかな
$R$-代数のフィルター付き余極限であるから、結果は Algebra, Lemmas
\ref{algebra-lemma-normal-goes-up} および
\ref{algebra-lemma-colimit-normal-ring} から従う。
\end{proof}

\begin{lemma}
\label{more-algebra-lemma-henselization-dimension}
任意の局所環 $R$ に対して $\dim(R) = \dim(R^h) = \dim(R^{sh})$ である。
\end{lemma}

\begin{proof}
$R \to R^{sh}$ は忠実平坦である（Lemma \ref{more-algebra-lemma-dumb-properties-henselization}）から、
going down により $\dim(R^{sh}) \geq \dim(R)$ である。Algebra, Lemma
\ref{algebra-lemma-dimension-going-up} を参照せよ。
逆向きの不等式のため、$R^{sh} = \colim R_i$ を、各 $R_i$ が \'etale $R$-代数の局所化である
局所環の有向余極限と書く。$R^{sh}$ の素イデアルの鎖
$\mathfrak q_0 \subset \mathfrak q_1 \subset \ldots \subset \mathfrak q_n$
が与えられたとする。このとき十分大きな $i$ に対して
$$
R_i \cap \mathfrak q_0 \subset
R_i \cap \mathfrak q_1 \subset \ldots \subset
R_i \cap \mathfrak q_n
$$
は $R_i$ の素イデアルの鎖である。したがって
$\dim(R^{sh}) \leq \sup_i \dim(R_i)$ である。
しかし Lemma \ref{more-algebra-lemma-dimension-etale-extension} の結果により、各 $i$ について
$\dim(R_i) = \dim(R)$ である。これで示された。
\end{proof}

\begin{lemma}
\label{more-algebra-lemma-henselization-depth}
Noether 局所環 $R$ に対して
$\text{depth}(R) = \text{depth}(R^h) = \text{depth}(R^{sh})$ である。
\end{lemma}

\begin{proof}
Lemma \ref{more-algebra-lemma-henselization-noetherian} により $R^h$ と $R^{sh}$ は Noether である。
したがって補題は Algebra, Lemma \ref{algebra-lemma-apply-grothendieck} から従う。
\end{proof}

\begin{lemma}
\label{more-algebra-lemma-henselization-CM}
$R$ を Noether 局所環とする。次の条件は同値である。$R$ が Cohen-Macaulay であること、
$R$ の Hensel 化 $R^h$ が Cohen-Macaulay であること、および $R$ の強 Hensel 化
$R^{sh}$ が Cohen-Macaulay であること。
\end{lemma}

\begin{proof}
Lemma \ref{more-algebra-lemma-henselization-noetherian} により $R^h$ と $R^{sh}$ は Noether なので、
補題の主張は意味をもつ。さらに
$\text{depth}(R) = \text{depth}(R^h) = \text{depth}(R^{sh})$
および
$\dim(R) = \dim(R^h) = \dim(R^{sh})$
が Lemmas \ref{more-algebra-lemma-henselization-depth} および
\ref{more-algebra-lemma-henselization-dimension} により成り立つので、結論を得る。
\end{proof}

\begin{lemma}
\label{more-algebra-lemma-henselization-regular}
Noether 局所環 $R$ をとる。次の条件は同値である：
$R$ が正則局所環であること、$R$ の Hensel 化 $R^h$ が正則局所環で
あること、および $R$ の強 Hensel 化 $R^{sh}$ が正則局所環であること。
\end{lemma}

\begin{proof}
Lemma \ref{more-algebra-lemma-henselization-noetherian} により、$R^h$ と $R^{sh}$ は
Noether であるので、この補題は意味をもつ。$R$ の極大イデアルを
$\mathfrak m$ とする。
$\mathfrak m$ の最小生成系 $x_1, \ldots, x_t \in \mathfrak m$ をとる。
すなわち、その $\mathfrak m/\mathfrak m^2$ における像が
$\kappa = R/\mathfrak m$ 上の基底をなすようなものとする。
$R \to R^h$ および $R \to R^{sh}$ は忠実平坦なので、$R^h$ における
$x_1^h, \ldots, x_t^h$ の像、および $R^{sh}$ における
$x_1^{sh}, \ldots, x_t^{sh}$ の像は、それぞれ
$\mathfrak m^h = \mathfrak mR^h$ および
$\mathfrak m^{sh} = \mathfrak mR^{sh}$ の最小生成系である。
定義により $R$ の正則性は $t = \dim(R)$ を意味し、$R^h$ と
$R^{sh}$ についても同様である。したがってこの補題は、
Lemma \ref{more-algebra-lemma-henselization-dimension} の次元の等式
$\dim(R) = \dim(R^h) = \dim(R^{sh})$ から従う。
\end{proof}

\begin{lemma}
\label{more-algebra-lemma-henselization-dvr}
Noether 局所環 $R$ をとる。$R$ が離散付値環であることと、$R^h$ が
離散付値環であることと、$R^{sh}$ が離散付値環であることは同値である。
\end{lemma}

\begin{proof}
これは Lemmas \ref{more-algebra-lemma-henselization-dimension} および
\ref{more-algebra-lemma-henselization-regular} と Algebra, Lemma
\ref{algebra-lemma-characterize-dvr} から従う。
\end{proof}

\begin{lemma}
\label{more-algebra-lemma-filtered-colimit-etale-noetherian-fibres}
環 $A$ をとる。$B$ が \'etale $A$-代数のフィルター付き余極限であるとする。
$\mathfrak p$ を $A$ の素イデアルとする。$B$ が Noether ならば、
$\mathfrak p$ の上にある素イデアル $\mathfrak q_1, \ldots, \mathfrak q_r$
が有限個存在し、
$B \otimes_A \kappa(\mathfrak p) = \prod \kappa(\mathfrak q_i)$
であり、各体拡大
$\kappa(\mathfrak q_i)/\kappa(\mathfrak p)$ は分離代数的である。
\end{lemma}

\begin{proof}
$B$ を、$A \to B_i$ が \'etale となるようなフィルター付き余極限
$B = \colim B_i$ と書く。一方では
$B \otimes_A \kappa(\mathfrak p) = \colim B_i \otimes_A \kappa(\mathfrak p)$
は \'etale $\kappa(\mathfrak p)$-代数のフィルター付き余極限であり、
他方では Noether である。\'etale $\kappa(\mathfrak p)$-代数は有限分離体
拡大の有限積である（Algebra, Lemma
\ref{algebra-lemma-etale-over-field}）。
したがって、各代数
$B_i \otimes_A \kappa(\mathfrak p)$ の（すべて極大かつ極小である）
素イデアルの間には非自明な特殊化がなく、ゆえに
$B \otimes_A \kappa(\mathfrak p)$ の素イデアルの間にも非自明な
特殊化はない。したがって
$B \otimes_A \kappa(\mathfrak p)$ は被約であり、素イデアルは有限個で、
すべて極小である。
ゆえにこれは体の有限積である（Algebra, Lemma
\ref{algebra-lemma-total-ring-fractions-no-embedded-points} または
Algebra, Proposition \ref{algebra-proposition-dimension-zero-ring} を用いる）。
これらの体はそれぞれ有限分離拡大の余極限であるから、補題の最後の主張が従う。
\end{proof}

\begin{lemma}
\label{more-algebra-lemma-fibres-henselization}
Noether 局所環 $R$ をとる。$\mathfrak p \subset R$ を素イデアルとする。
このとき
$$
R^h \otimes_R \kappa(\mathfrak p) =
\prod\nolimits_{i = 1, \ldots, t} \kappa(\mathfrak q_i)
\quad\text{resp.}\quad
R^{sh} \otimes_R \kappa(\mathfrak p) =
\prod\nolimits_{i = 1, \ldots, s} \kappa(\mathfrak r_i)
$$
である。ここで $\mathfrak q_1, \ldots, \mathfrak q_t$、
resp.\ $\mathfrak r_1, \ldots, \mathfrak r_s$ は、それぞれ
$\mathfrak p$ の上にある $R^h$、resp.\ $R^{sh}$ の素イデアルである。
さらに体拡大
$\kappa(\mathfrak q_i)/\kappa(\mathfrak p)$
resp.\ $\kappa(\mathfrak r_i)/\kappa(\mathfrak p)$
は分離代数的である。
\end{lemma}

\begin{proof}
これは、より一般的な Lemma
\ref{more-algebra-lemma-filtered-colimit-etale-noetherian-fibres} から、Hensel 化と
強 Hensel 化が Noether であること（上で見た）を使って導ける。しかし、
ここでは次のように直接証明も与える。
\medskip\noindent
以下では、特に断らない限り Lemmas
\ref{more-algebra-lemma-dumb-properties-henselization} および
\ref{more-algebra-lemma-henselization-noetherian} の結果を用いる。
$R^h/\mathfrak pR^h$、resp.\ $R^{sh}/\mathfrak pR^{sh}$ は、
$R/\mathfrak p$ の Hensel 化、resp.\ 強 Hensel 化であることに注意する。
Algebra, Lemma \ref{algebra-lemma-quotient-henselization}
resp.\ Algebra, Lemma \ref{algebra-lemma-quotient-strict-henselization} を参照せよ。
したがって $R$ を $R/\mathfrak p$ で置き換え、
$R$ が Noether 局所整域で $\mathfrak p = (0)$ であると仮定してよい。
$R^h$、resp.\ $R^{sh}$ は Noether なので、極小素イデアル
$\mathfrak q_1, \ldots, \mathfrak q_t$、resp.\
$\mathfrak r_1, \ldots, \mathfrak r_s$ を有限個もつ。
$R \to R^h$、resp.\ $R \to R^{sh}$ は平坦なので、これらは
（going down により）$\mathfrak p = (0)$ の上にある素イデアルにほかならない。
最後に、$R$ は整域なので、$R^h$、resp.\ $R^{sh}$ は被約であることが
Lemma \ref{more-algebra-lemma-henselization-reduced} から分かる。
したがって $R^h \otimes_R \kappa(\mathfrak p)$、resp.\
$R^{sh} \otimes_R \kappa(\mathfrak p)$ は、有限個の素イデアルをもつ
被約 Noether 環であり、それらはすべて極小（したがって極大）である。
ゆえにこれらの環は Artinian であり、極大イデアルでの局所化の積である。
各局所化は必ず体である（Algebra, Proposition
\ref{algebra-proposition-dimension-zero-ring} および Algebra, Lemma
\ref{algebra-lemma-minimal-prime-reduced-ring} を参照せよ）。
\medskip\noindent
最後の主張は、$R \to R^h$、resp.\ $R \to R^{sh}$ が \'etale 環写像の
余極限であり、したがって誘導される剰余体拡大が有限分離拡大の余極限である
ことから従う。Algebra, Lemma \ref{algebra-lemma-etale-at-prime} を参照せよ。
\end{proof}


\section{体拡大、再訪}
\label{more-algebra-section-p-bases}

\noindent
この節では、標数 $p > 0$ の体拡大に関するいくつかの特異な性質を調べる。

\begin{definition}
\label{more-algebra-definition-p-basis}
素数 $p$ をとる。体の拡大 $k \to K$ をとり、標数 $p$ であるとする。
$K$ の中で $k$ と $K^p$ の合成体を $kK^p$ と記す。
\begin{enumerate}
\item $K$ の部分集合 $\{x_i\} \subset K$ が $k$ 上 {\it p-独立} であるとは、
$0 \leq e_i < p$ を満たす元
$x^E = \prod x_i^{e_i}$ が $kK^p$ 上線型独立であることをいう。
\item $K$ の部分集合 $\{x_i\}$ が $k$ 上の
{\it p-基底} であるとは、元 $x^E$ が $kK^p$ 上の $K$ の基底をなすことをいう。
\end{enumerate}
\end{definition}

\noindent
これは $\mathbf{F}_p$-代数の $p$-基底という概念（後で議論する）と関連する
（将来の参照をここに挿入）。

\begin{lemma}
\label{more-algebra-lemma-p-basis}
$K/k$ を体拡大とする。$k$ が標数 $p > 0$ をもつと仮定する。
$\{x_i\}$ を $K$ の部分集合とする。次は同値である。
\begin{enumerate}
\item 元 $\{x_i\}$ は $k$ 上 $p$-独立であり、
\item 元 $\text{d}x_i$ は $\Omega_{K/k}$ において $K$-線型独立である。
\end{enumerate}
任意の $p$-独立な集合は、$K$ の $k$ 上の $p$-基底に拡張できる。
特に体 $K$ は $k$ 上の $p$-基底をもつ。
さらに次は同値である。
\begin{enumerate}
\item[(a)] $\{x_i\}$ は $K$ の $k$ 上の $p$-基底であり、
\item[(b)] $\text{d}x_i$ は $K$-ベクトル空間 $\Omega_{K/k}$ の基底である。
\end{enumerate}
\end{lemma}

\begin{proof}
(2) を仮定し、$\sum a_E x^E = 0$ が $a_E \in k K^p$ を係数とする
線型関係であるとする。$\theta_i : K \to K$ を、
$\theta_i(x_j) = \delta_{ij}$（Kronecker のデルタ）を満たす $k$-導分とする。
任意の $k$-導分は $kK^p$ を消すことに注意する。与えられた関係に
$\theta_i$ を適用すると、次の新しい関係を得る。
$$
\sum\nolimits_{E, e_i > 0}
e_i a_E x_1^{e_1}\ldots x_i^{e_i - 1} \ldots x_n^{e_n} = 0
$$
したがって、ある $a_E \not = 0$ に対して $\sum a_E x^E$ を全次数
$|E| = \sum e_i$ が最小の関係として選ぶと、矛盾を得る。ゆえに (1) が成り立つ。
\medskip\noindent
$\{x_i\}$ が $K$ の $k$ 上の $p$-基底ならば、
$K \cong kK^p[X_i]/(X_i^p - x_i^p)$ である。したがって
$\text{d}x_i$ は $K$ 上の $\Omega_{K/k}$ の基底をなすことが分かる。
ゆえに (a) は (b) を含意する。
\medskip\noindent
$\{x_i\}$ を $k$ 上の $K$ の $p$-独立部分集合とする。Zorn の補題の適用により、
これを $K$ の極大 $p$-独立部分集合に拡張できる。
$K$ の任意の極大 $p$-独立部分集合 $\{x_i\}$ は $K$ の $k$ 上の
$K$ の $k$ 上の $p$-基底であると主張する。この主張は、(1) が (2) を含意することと
$p$-基底の存在を示す。
主張を証明するため、$kK^p$ と $x_i$ たちが生成する $K$ の部分体を
$L$ とする。$L = K$ を示さなければならない。$x \in K$ だが
$x \not \in L$ ならば、$x^p \in L$ であり、
$L(x) \cong L[z]/(z^p - x)$ である。したがって
$\{x_i\} \cup \{x\}$ は $k$ 上 $p$-独立となり、矛盾である。
\medskip\noindent
最後に、(b) が (a) を含意することを示す。 (1) と (2) の同値性により、
$\{x_i\}$ は $k$ 上の $K$ の極大 $p$-独立部分集合である。したがって
上の主張により、これは $p$-基底である。
\end{proof}


\begin{lemma}
\label{more-algebra-lemma-intersection-subfields-subspace}
$K/k$ を体拡大とする。$\{K_\alpha\}_{\alpha \in A}$ を $K$ の部分体族とし、
次の性質を満たすとする。
\begin{enumerate}
\item すべての $\alpha \in A$ に対して $k \subset K_\alpha$、
\item $k = \bigcap_{\alpha \in A} K_\alpha$、
\item $\alpha, \alpha' \in A$ に対して、ある $\alpha'' \in A$ が存在し、
$K_{\alpha''} \subset K_\alpha \cap K_{\alpha'}$ となる。
\end{enumerate}
このとき、$n \geq 1$ および $V \subset K^{\oplus n}$ が $K$-ベクトル空間ならば、
$V \cap k^{\oplus n} \not = 0$ であることと、すべての $\alpha \in A$ に対して
$V \cap K_\alpha^{\oplus n} \not = 0$ であることは同値である。
\end{lemma}

\begin{proof}
$n$ に関する帰納法による。$n = 1$ の場合は仮定から従う。
$K^{\oplus n - 1}$ の部分空間について結果が証明済みであると仮定する。
$V \subset K^{\oplus n}$ がすべての $\alpha \in A$ に対して
$K_\alpha^{\oplus n}$ と非零の共通部分をもつと仮定する。
$V \cap 0 \oplus k^{\oplus n - 1}$ が非零ならば示された。
したがって、これが成り立たないと仮定してよい。帰納法の仮定により、
ある $\alpha$ について $V \cap 0 \oplus K_\alpha^{\oplus n - 1}$ が零となる。
$v = (x_1, \ldots, x_n) \in V \cap K_\alpha^{\oplus n}$ を非零元とする。
$\alpha$ の選び方により $x_1$ は零でない。
$v$ を $x_1^{-1}v$ に置き換え、$v = (1, x_2, \ldots, x_n)$ とする。
$v' = (x_1', \ldots, x'_n) \in V \cap K_\alpha^{\oplus n}$ ならば、
$\alpha$ の選び方により $v' - x_1'v = 0$ であることに注意する。したがって
$V \cap K_\alpha^{\oplus n} = K_\alpha v$ である。
$K_{\alpha'} \subset K_\alpha$ となる $\alpha'$ を選ぶと、
（$\alpha'$ に同じ議論を適用すれば）必然的に
$v \in V \cap K_{\alpha'}^{\oplus n}$ であることが分かる。したがって
$$
x_2, \ldots, x_n \in
\bigcap\nolimits_{\alpha' \in A, K_{\alpha'} \subset K_\alpha} K_{\alpha'}
$$
これは (2) と (3) により $k$ に等しい。
\end{proof}

\begin{lemma}
\label{more-algebra-lemma-intersection-subfields}
標数 $p$ の体 $K$ をとる。$\{K_\alpha\}_{\alpha \in A}$ を $K$ の部分体族とし、
次の性質を満たすとする。
\begin{enumerate}
\item すべての $\alpha \in A$ に対して $K^p \subset K_\alpha$、
\item $K^p = \bigcap_{\alpha \in A} K_\alpha$、
\item $\alpha, \alpha' \in A$ に対して、ある $\alpha'' \in A$ が存在し、
$K_{\alpha''} \subset K_\alpha \cap K_{\alpha'}$ となる。
\end{enumerate}
このとき
\begin{enumerate}
\item 写像 $\Omega_{K/\mathbf{F}_p} \to \Omega_{K/K_\alpha}$ の核の共通部分は零であり、
\item 任意の有限拡大 $L/K$ に対して $L^p = \bigcap_{\alpha \in A} L^pK_\alpha$ である。
\end{enumerate}
\end{lemma}


\begin{proof}
(1) の証明。
$\mathbf{F}_p$ 上の $K$ の $p$-基底 $\{x_i\}$ を選ぶ。
$\eta = \sum_{i \in I'} y_i \text{d}x_i$ が、すべての $\alpha \in A$ に対して
$\Omega_{K/K_\alpha}$ で零に写ると仮定する。ここで添字集合 $I'$ は有限である。
Lemma \ref{more-algebra-lemma-p-basis} により、これはすべての $\alpha$ に対して関係
$$
\sum\nolimits_E a_{E, \alpha} x^E = 0,\quad a_{E, \alpha} \in K_\alpha
$$
が存在することを意味する。ここで $E$ は、$0 \leq e_i < p$ を満たす
多重添字 $E = (e_i)_{i \in I'}$ を走る。一方、Lemma
\ref{more-algebra-lemma-p-basis} は、$a_E \in K^p$ を係数とする関係
$\sum a_E x^E = 0$ は存在しないことを保証する。
これは Lemma \ref{more-algebra-lemma-intersection-subfields-subspace} による矛盾である。
\medskip\noindent
(2) の証明。有限体拡大の塔 $L/M/K$ があるとする。
$M_\alpha = M^p K_\alpha$ および $L_\alpha = L^p K_\alpha = L^p M_\alpha$
とおく。まず $M^p = \bigcap_{\alpha \in A} M_\alpha$ を証明し、その後で
$L^p = \bigcap_{\alpha \in A} L_\alpha$ を証明できる。したがって、
非自明な部分体をもたない原始体拡大について (2) を証明すれば十分である。
まず、$L = K(\theta)$ が $K$ 上分離的であると仮定する。このとき
$L$ は $K$ 上 $\theta^p$ により生成されるので、$\theta \in L^p$ と
仮定してよい。この場合
$$
L^p = K^p \oplus K^p\theta \oplus \ldots K^p\theta^{d - 1}
\quad\text{and}\quad
L^pK_\alpha =
K_\alpha \oplus K_\alpha \theta \oplus \ldots K_\alpha\theta^{d - 1}
$$
であり、ここで $d = [L : K]$ である。したがってこの場合の結論は明らかである。
もう一つの場合は、$L = K(\theta)$ かつ $\theta^p = t \in K$、
$t \not \in K^p$ である。この場合
$$
L^p = K^p \oplus K^pt \oplus \ldots K^pt^{p - 1}
\quad\text{and}\quad
L^pK_\alpha =
K_\alpha \oplus K_\alpha t \oplus \ldots K_\alpha t^{p - 1}
$$
である。再び結果は明らかである。
\end{proof}

\begin{lemma}
\label{more-algebra-lemma-power-series-ring-subfields}
標数 $p > 0$ の体 $k$ をとる。$\{x_i\}_{i \in I}$ を $k$ の $p$-基底とする。
$n, m \geq 0$ とする。
$K$ を
$A = k[[x_1, \ldots, x_n]][y_1, \ldots, y_m]$ の分数体とする。
$J$ を $I$ の有限部分集合とする。$k/k_J/k^p$ を、$k^p$ と
$i \in I \setminus J$ に対する $x_i$ により生成される部分体とする。
$$
A_J = k_J[[x_1^p, \ldots, x_n^p]][y_1^p, \ldots, y_m^p]
$$
の分数体 $K_J$ は、$K$ の部分体族として Lemma \ref{more-algebra-lemma-intersection-subfields} における
部分体族をなす。さらに、各環拡大 $A_J \subset A$ は有限である。
\end{lemma}

\begin{proof}
$k/k_J$ は有限なので、環拡大
$k_J[[x_1^p, \ldots, x_d^p]] \subset k[[x_1, \ldots, x_d]]$ は
Algebra, Lemma \ref{algebra-lemma-finite-after-completion} により有限である。
これは $A_J \to A$ が有限であることを含意する。
\medskip\noindent
Lemma \ref{more-algebra-lemma-intersection-subfields} の性質 (1)、(2)、(3) を確認する。
(1) の証明。$a \in A$ に対して $a^p \in A_J$ であることが分かる。
したがって $K^p \subset K_J$。(2) の証明。
$f/g^p \in K$（$f, g \in A$、$g \not = 0$）が、あらゆる $J$ の選択に対して
$K_J$ に含まれると仮定する。
いま $J$ を固定する。$f/g^p \in K_J$ なので、
$a \in A_J$ および非零の $b \in A$ を用いて
$f/g^p = a/b^p$ と書ける。したがって $b^p f \in A_J$ である。
任意の $A_J$-導分 $D : A \to A$ に対して
$0 = D(b^pf) = b^p D(f)$ であり、$A$ は整域なので $D(f) = 0$ である。
$D = \partial_{x_i}$ および $D = \partial_{y_j}$ と取ると、
$f \in k[[x_1^p, \ldots, x_n^p]][y_1^p, \ldots, y_d^p]$
であることが分かる。
$k_J$-導分 $\theta : k \to k$ を適用して同様に、
$f$ のすべての係数が $k_J$ にある、すなわち $f \in A_J$ と結論する。
$A^p = \bigcap\nolimits_J A_J$ であることは明らかであり、ここで $J$ は
補題におけるすべての部分体を走るので、所望どおり $f \in A^p$ を得る。
(3) の証明。これは $K_{J \cup J'} \subset K_J \cap K_{J'}$ なので明らかである。
\end{proof}






\section{特異点集合}
\label{more-algebra-section-singular-locus}

\noindent
$R$ を Noether 環とする。$X = \Spec(R)$ の {\it 正則点集合}
$\text{Reg}(X)$ とは、$R_\mathfrak p$ が正則局所環となる素イデアル
$\mathfrak p$ の集合である。$X = \Spec(R)$ の {\it 特異点集合}
$\text{Sing}(X)$ とは、その補集合
$X \setminus \text{Reg}(X)$、すなわち $R_\mathfrak p$ が正則局所環でない
素イデアル $\mathfrak p$ の集合である。Algebra, Definition
\ref{algebra-definition-regular} に先立つ議論により、
$\text{Reg}(X)$ は一般化に関して安定であることが分かる。
この節では、$\text{Reg}(X)$ が開となることを保証する条件を調べる。

\begin{definition}
\label{more-algebra-definition-J}
\begin{reference}
\cite[(32.B)]{MatCA}
\end{reference}
$R$ を Noether 環とする。$X = \Spec(R)$ とおく。
\begin{enumerate}
\item $\text{Reg}(X)$ が非空開集合を含むとき、$R$ は {\it J-0} であるという。
\item $\text{Reg}(X)$ が開であるとき、$R$ は {\it J-1} であるという。
\item 任意の有限型 $R$-代数が J-1 であるとき、$R$ は {\it J-2} であるという。
\end{enumerate}
\end{definition}

\noindent
環 $\mathbf{Q}[x]/(x^2)$ は J-0 を満たさないが、J-1 は満たす。
一方、Noether 整域、より一般には非零被約 Noether 環については、そのような環は
極小素イデアルにおいて正則なので、J-1 は J-0 を含意する。
J-1 性質の特徴づけを次に与える。

\begin{lemma}
\label{more-algebra-lemma-J-1}
$R$ を Noether 環とする。$X = \Spec(R)$ とおく。
$R$ が J-1 であるための必要十分条件は、すべての
$\mathfrak p \in \text{Reg}(X)$ に対して
$V(\mathfrak p) \cap \text{Reg}(X)$ が $V(\mathfrak p)$ の非空開部分集合を
含むことである。
\end{lemma}

\begin{proof}
これは Topology, Lemma
\ref{topology-lemma-characterize-open-Noetherian} と、Algebra, Lemma
\ref{algebra-lemma-localization-of-regular-local-is-regular} により
$\text{Reg}(X)$ が一般化に関して安定であることから従う。
\end{proof}

\begin{lemma}
\label{more-algebra-lemma-intersection-regular-with-closed}
$R$ を Noether 環とする。$X = \Spec(R)$ とおく。すべての素イデアル
$\mathfrak p \subset R$ に対して環 $R/\mathfrak p$ が J-0 であると仮定する。
このとき $R$ は J-1 である。
\end{lemma}

\begin{proof}
Lemma \ref{more-algebra-lemma-J-1} の判定条件が適用できることを示す。
$\mathfrak p \in \text{Reg}(X)$ を高さ $r$ の素イデアルとする。
$\mathfrak pR_\mathfrak p$ の生成元に写る
$f_1, \ldots, f_r \in \mathfrak p$ を選ぶ。
$\mathfrak p \in \text{Reg}(X)$ なので、$f_1, \ldots, f_r$ は
$R_\mathfrak p$ において正則列に写る。Algebra, Lemma
\ref{algebra-lemma-regular-ring-CM} を参照せよ。したがって Algebra, Lemma
\ref{algebra-lemma-regular-sequence-in-neighbourhood} により、
ある $g \in R$、$g \not \in \mathfrak p$ に対して $R$ を $R_g$ で
置き換えた後、列 $f_1, \ldots, f_r$ は $R$ における正則列となる。
さらにもう一度置き換えることで、$f_1, \ldots, f_r$ が
$\mathfrak p$ を生成すると仮定してよい。
次に、$\mathfrak p \subset \mathfrak q$ を、
$(R/\mathfrak p)_\mathfrak q$ が正則局所環となる素イデアルとする。
補題の仮定により、そのような素イデアルからなる $V(\mathfrak p)$ の
非空開部分集合が存在するので、$R_\mathfrak q$ が正則であることを
証明すれば十分である。$f_1, \ldots, f_r$ は
$R_\mathfrak q$ における正則列であり、
$R_\mathfrak q/(f_1, \ldots, f_r)R_\mathfrak q$ は正則であることに注意する。
したがって Algebra, Lemma \ref{algebra-lemma-regular-mod-x} により
$R_\mathfrak q$ は正則である。
\end{proof}

\begin{lemma}
\label{more-algebra-lemma-J-0-goes-down}
$R \to S$ を環写像とする。次を仮定する。
\begin{enumerate}
\item $R$ は Noether 整域であり、
\item $R \to S$ は単射かつ有限型であり、
\item $S$ は整域かつ J-0 である。
\end{enumerate}
このとき $R$ は J-0 である。
\end{lemma}

\begin{proof}
ある非零元 $g \in S$ に対して $S$ を $S_g$ で置き換え、
$S$ が正則環であると仮定してよい。一般平坦性により、
$R \to S$ も忠実平坦であると仮定してよい。Algebra, Lemma
\ref{algebra-lemma-generic-flatness-Noetherian} を参照せよ。
すると Algebra, Lemma \ref{algebra-lemma-descent-regular} により
$R$ は正則である。
\end{proof}

\begin{lemma}
\label{more-algebra-lemma-J-0-goes-up}
$R \to S$ を環写像とする。次を仮定する。
\begin{enumerate}
\item $R$ は Noether 整域かつ J-0 であり、
\item $R \to S$ は単射かつ有限型であり、
\item $S$ は整域であり、
\item 誘導される分数体の拡大は分離的である。
\end{enumerate}
このとき $S$ は J-0 である。
\end{lemma}

\begin{proof}
$R$ を単項局所化で置き換え、$R$ が正則環であると仮定してよい。
Algebra, Lemma \ref{algebra-lemma-smooth-at-generic-point} により、
環写像 $R \to S$ は $(0)$ において滑らかである。
したがって $S$ を単項局所化で置き換えた後、
$S$ が $R$ 上滑らかであると仮定してよい。すると $S$ も正則である。
Algebra, Lemma \ref{algebra-lemma-regular-goes-up} を参照せよ。
\end{proof}

\begin{lemma}
\label{more-algebra-lemma-J-2}
$R$ を Noether 環とする。次は同値である。
\begin{enumerate}
\item $R$ は J-2 である。
\item 整域である任意の有限型 $R$-代数は J-0 である。
\item 任意の有限 $R$-代数は J-1 である。
\item 任意の素イデアル $\mathfrak p$ と任意の有限純非分離拡大
$L/\kappa(\mathfrak p)$ に対して、整域かつ J-0 で、分数体が $L$ となる
有限 $R$-代数 $R'$ が存在する。
\end{enumerate}
\end{lemma}

\begin{proof}
含意 (1) $\Rightarrow$ (2) および (2) $\Rightarrow$ (4) は明らかである。
J-1 である整域は J-0 であることを想起せよ。したがって含意
(1) $\Rightarrow$ (3) および (3) $\Rightarrow$ (4) も成り立つ。
\medskip\noindent
$R \to S$ を有限型環写像とし、$S$ が J-1 であることを示そう。
Lemma \ref{more-algebra-lemma-intersection-regular-with-closed} により、$S$ のすべての
素イデアル $\mathfrak q$ に対して $S/\mathfrak q$ が J-0 であることを
証明すれば十分である。このようにして (2) $\Rightarrow$ (1) が分かる。
\medskip\noindent
(4) を仮定する。(2) が成り立つことを示せば証明が完了する。
$R \to S$ を有限型環写像とし、$S$ が整域であるとする。
$\mathfrak p = \Ker(R \to S)$ とおく。$K$ を $S$ の分数体とする。
次の体の図式が存在する。
$$
\xymatrix{
K \ar[r] & K' \\
\kappa(\mathfrak p) \ar[u] \ar[r] & L \ar[u]
}
$$
ここで水平矢印は有限純非分離体拡大であり、$K'/L$ は分離的である。
Algebra, Lemma \ref{algebra-lemma-make-separably-generated} を参照せよ。
(4) のように $R' \subset L$ を選び、写像
$S \otimes_R R' \to K'$ の像を $S'$ とする。
すると $S'$ は分数体 $K'$ をもつ整域であり、したがって Lemma
\ref{more-algebra-lemma-J-0-goes-up} と $R'$ の選び方により $S'$ は J-0 である。
次に Lemma \ref{more-algebra-lemma-J-0-goes-down} を適用すると、
所望どおり $S$ は J-0 である。
\end{proof}




\section{正則性と導分}
\label{more-algebra-section-regularity-derivations}

\noindent
$R \to S$ を環写像とする。$D : R \to R$ を導分とする。
次の図式を可換にする導分 $D' : S \to S$ が存在するとき、
$D$ は $S$ に {\it 延長される} という。
$$
\xymatrix{
S \ar[r]_{D'} & S \\
R \ar[u] \ar[r]^D & R \ar[u]
}
$$

\begin{lemma}
\label{more-algebra-lemma-derivation-extends}
$R$ を環とする。$D : R \to R$ を導分とする。
\begin{enumerate}
\item 任意のイデアル $I \subset R$ に対して、導分 $D$ は $I$-進完備化上の
導分 $D^\wedge : R^\wedge \to R^\wedge$ に標準的に延長される。
\item 任意の乗法的部分集合 $S \subset R$ に対して、導分 $D$ は
$R$ の局所化 $S^{-1}R$ に一意的に延長される。
\end{enumerate}
$R \subset R'$ が有限型環拡大で、$R'$ における非零因子
$g \in R$ に対して $R_g \cong R'_g$ であるならば、
ある $N \geq 0$ に対して $g^ND$ は $R'$ に延長される。
\end{lemma}

\begin{proof}
(1) の証明。Leibniz 則により、$n \geq 2$ に対して
$D(I^n) \subset I^{n - 1}$ である。したがって $D$ は写像
$D_n : R/I^n \to R/I^{n - 1}$ を誘導する。極限を取って
$D^\wedge$ を得る。$D^\wedge$ が導分であることの確認は省略する。
\medskip\noindent
(2) の証明。$D$ を $S^{-1}R$ に延長するには
$D(r/s) = D(r)/s - rD(s)/s^2$ と定め、諸公理を確認すればよい。
\medskip\noindent
最後の主張の証明。$x_1, \ldots, x_n \in R'$ を $R$ 上の $R'$ の
生成元とする。$g^Nx_i \in R$ となる $N$ を選ぶ。
$g^{N + 1}D$ を考える。(2) により、これは $R_g$ に延長される。
さらに Leibniz 則と上の延長の構成により
$$
g^{N + 1}D(x_i) = g^{N + 1}D(g^{-N} g^Nx_i) = -Ng^Nx_iD(g) +
gD(g^Nx_i)
$$
であり、両項とも $R$ に属する。これは
$$
g^{N + 1}D(x_1^{e_1} \ldots x_n^{e_n}) =
\sum e_i x_1^{e_1} \ldots x_i^{e_i - 1} \ldots x_n^{e_n} g^{N + 1}D(x_i)
$$
が $R'$ の元であることを含意する。したがって $R'$ の任意の元
（$R$ の係数をもつ $x_i$ の単項式の和として書ける）は
$g^{N + 1}D$ により $R'$ の元に写り、証明が終わる。
\end{proof}

\begin{lemma}
\label{more-algebra-lemma-quotient-regular}
\begin{slogan}
超曲面に対する Jacobian 判定法を、正しく述べたもの。
\end{slogan}
$R$ を正則環とする。$f \in R$ とする。$D(f)$ が $R/(f)$ の単元となる
導分 $D : R \to R$ が存在すると仮定する。このとき $R/(f)$ は正則である。
\end{lemma}

\begin{proof}
$R$ が極大イデアル $\mathfrak m$ と剰余体 $\kappa$ をもつ局所環である場合を
証明すれば十分である。この場合、Algebra, Lemma
\ref{algebra-lemma-regular-ring-CM} により
$f \not \in \mathfrak m^2$ を証明すれば十分である。
しかし $f \in \mathfrak m^2$ ならば、Leibniz 則により
$D(f) \in \mathfrak m$ となり、矛盾である。
\end{proof}

\begin{lemma}
\label{more-algebra-lemma-quotient-sequence-regular}
$(R, \mathfrak m, \kappa)$ を正則局所環とする。$m \geq 1$ とする。
$f_1, \ldots, f_m \in \mathfrak m$ とする。導分
$D_1, \ldots, D_m : R \to R$ が存在し、
$\det_{1 \leq i, j \leq m}(D_i(f_j))$ が $R$ の単元であると仮定する。
このとき $R/(f_1, \ldots, f_m)$ は正則であり、
$f_1, \ldots, f_m$ は正則列である。
\end{lemma}

\begin{proof}
Algebra, Lemma \ref{algebra-lemma-regular-ring-CM} により、
$f_1, \ldots, f_m$ が $\mathfrak m/\mathfrak m^2$ において
$\kappa$-線型独立であることを示せば十分である。
しかし非自明な線型関係があれば、すべてが
$a_i \in \mathfrak m$ ではないある $a_i \in R$ に対して
$\sum a_i f_i \in \mathfrak m^2$ を得る。導分の Leibniz 則により
$D_i(\mathfrak m^2) \subset \mathfrak m$ および
$D_i(a_j f_j) \equiv a_j D_i(f_j) \bmod \mathfrak m$
であることに注意する。したがってこれは
$$
\sum a_j D_i(f_j) \in \mathfrak m
$$
を含意し、行列式に関する仮定に矛盾する。
\end{proof}

\begin{lemma}
\label{more-algebra-lemma-degree-p-extension-regular}
$R$ を正則環とする。$f \in R$ とする。$D(f)$ が $R$ の単元となる導分
$D : R \to R$ が存在すると仮定する。このとき任意の整数 $n \geq 1$ に対して
$R[z]/(z^n - f)$ は正則である。より一般に、任意の
$p \in \mathbf{Z}[z]$ に対して $R[z]/(p(z) - f)$ は正則である。
\end{lemma}

\begin{proof}
Algebra, Lemma \ref{algebra-lemma-regular-goes-up} により、
$R[z]$ は正則環である。$z$ を零に写す $R[z]$ への $D$ の延長に
Lemma \ref{more-algebra-lemma-quotient-regular} を適用する。
$D$ は整数係数の任意の多項式を消し、$f$ を単元に写すので、これでよい。
\end{proof}

\begin{lemma}
\label{more-algebra-lemma-find-D}
$p$ を素数とする。$B$ を、$B$ において $p = 0$ が成り立つ整域とする。
$f \in B$ を、$B$ の分数体における $p$ 乗でない元とする。
$B$ が Noether 完備局所環上有限型ならば、
$D(f)$ が零でない導分 $D : B \to B$ が存在する。
\end{lemma}

\begin{proof}
有限型環写像 $R \to B$ が存在するような Noether 完備局所環 $R$ をとる。
もちろん $R$ を $B$ における像で置き換えてよいので、
$R$ は標数 $p > 0$ の整域（かつ Noether 完備局所）であると仮定してよい。
Algebra, Lemma
\ref{algebra-lemma-complete-local-Noetherian-domain-finite-over-regular}
により、ある体 $k$ と整数 $n$ に対して $R$ を
$k[[x_1, \ldots, x_n]]$ の有限拡大として書ける。
したがって $R$ を $k[[x_1, \ldots, x_n]]$ で置き換えてよい。
次に Algebra, Lemma
\ref{algebra-lemma-Noether-normalization-over-a-domain} を用いて
$R \to B$ を
$$
R \subset R[y_1, \ldots, y_d] \subset B' \subset B
$$
と分解する。ここで $B'$ は $R[y_1, \ldots, y_d]$ 上有限であり、
ある非零元 $g \in R$ に対して $B'_g \cong B_g$ である。
十分大きな整数 $N$ に対して $f' = g^{pN} f \in B'$ であることに注意する。
$f'$ が $B'$ の分数体における $p$ 乗でないことは明らかである。
$D'(f') \not = 0$ を満たす導分 $D' : B' \to B'$ を見つけられれば、
Lemma \ref{more-algebra-lemma-derivation-extends} により、ある $M > 0$ に対して
$D = g^MD'$ は $B$ に延長される。そのとき
$D(f) = g^MD'(f) = g^MD'(g^{-pN}f') = g^{M - pN}D'(f')$
は非零である。したがって
$B$ が $A = k[[x_1, \ldots, x_n]][y_1, \ldots, y_m]$ の有限拡大である
場合に補題を証明すれば十分である。
\medskip\noindent
$B$ が
$A = k[[x_1, \ldots, x_n]][y_1, \ldots, y_m]$
の有限拡大であると仮定する。$B$ の分数体を $L$ と記す。
$\text{d}f$ は $\Omega_{L/\mathbf{F}_p}$ において非零であることに注意する。
Algebra, Lemma \ref{algebra-lemma-derivative-zero-pth-power} を参照せよ。
Lemma \ref{more-algebra-lemma-power-series-ring-subfields} を適用すると、有限指数の部分体
$k' \subset k$ であって、
$A' = k'[[x_1^p, \ldots, x_n^p]][y_1^p, \ldots, y_m^p]$
としたとき、$\text{d}f$ が $\Omega_{L/K'}$ において零に写らないものを
見つけられる。ここで $K'$ は $A'$ の分数体である。
したがって $D'(f) \not = 0$ を満たす $K'$-導分
$D' : L \to L$ を選べる。構成により $A' \subset A$ および
$A \subset B$ は有限なので、$A' \subset B$ は有限である。
$B$ を $A'$-加群として生成する $b_1, \ldots, b_t \in B$ を選ぶ。
このとき、ある $f_i, g_i \in B$、$g_i \not = 0$ によって
$D'(b_i) = f_i/g_i$ と書ける。$D = g_1 \ldots g_t D'$ とおけば示された。
\end{proof}

\begin{lemma}
\label{more-algebra-lemma-complete-Noetherian-domain-J-0}
$A$ を Noether 完備局所整域とする。このとき $A$ は J-0 である。
\end{lemma}

\begin{proof}
Algebra, Lemma
\ref{algebra-lemma-complete-local-Noetherian-domain-finite-over-regular}
により、$A$ が $A_0$ 上有限となる正則部分環 $A_0 \subset A$ を見つけられる。
誘導される分数体の拡大 $K/K_0$ は有限である。
$K/K_0$ が分離的ならば、Lemma \ref{more-algebra-lemma-J-0-goes-up} により示される。
そうでなければ $A_0$ と $A$ は標数 $p > 0$ をもつ。
任意の中間拡大 $K/M/K_0$ に対して、分数体が $M$ となる有限中間拡大
$A_0 \subset B \subset A$ が存在する。したがって $[K : K_0]$ に関する
帰納法により、$B$ が J-0 で、$K/M$ が非自明な中間拡大をもたないような
$A_0 \subset B \subset A$ が存在すると仮定してよい。
この場合、$K/M$ が分離的ならば Lemma \ref{more-algebra-lemma-J-0-goes-up} により
$A$ は J-0 である。
そうでなければ、ある $b_1, b_2 \in B$、$b_2 \not = 0$ で、
$b_1/b_2$ が $M$ における $p$ 乗でないものに対して
$K = M[z]/(z^p - b_1/b_2)$ である。
$az \in A$ となる非零元 $a \in A$ を選ぶ。
$z$ を $b_2 a^p z$ で置き換えると、$z \in A$ かつ
$b \in B$ が $M$ における $p$ 乗でないような
$K = M[z]/(z^p - b)$ を得る。
Lemma \ref{more-algebra-lemma-find-D} により、$D(b) \not = 0$ を満たす導分
$D : B \to B$ を見つけられる。
Lemma \ref{more-algebra-lemma-degree-p-extension-regular} を適用すると、$B$ の正則素イデアル
の上にあり、$D(b)$ を含まない $A$ の任意の素イデアル
$\mathfrak p$ に対して $A_\mathfrak p$ は正則である。
$B$ は J-0 なので、$A$ も J-0 であると結論する。
\end{proof}


\begin{proposition}
\label{more-algebra-proposition-ubiquity-J-2}
次の種類の環は J-2 である。
\begin{enumerate}
\item 体、
\item Noether 完備局所環、
\item $\mathbf{Z}$、
\item 次元 $1$ の Noether 局所環、
\item 次元 $1$ の Nagata 環、
\item 標数零の分数体をもつ Dedekind 整域、
\item 以上のいずれかの有限型環拡大。
\end{enumerate}
\end{proposition}

\begin{proof}
場合 (1)、(3)、(5)、(6) については、Lemma \ref{more-algebra-lemma-J-2} の条件 (4) を
確認すれば証明される。ここでは $R$ が次元 $1$ の Nagata 環である場合だけを扱う。
$\mathfrak p \subset R$ を素イデアルとし、
$L/\kappa(\mathfrak p)$ を有限純非分離拡大とする。
$\mathfrak p \subset R$ が極大イデアルならば、$R \to L$ は有限で、
$L$ は正則環なので、条件は確認された。
$\mathfrak p \subset R$ が極小素イデアルならば、Nagata 条件により、
$L$ における $R$ の整閉包 $R' \subset L$ は $R$ 上有限である。
すると $R'$ は次元 $1$ の正規整域であり
（Algebra, Lemma \ref{algebra-lemma-integral-dim-up}）、
したがって正則である
（Algebra, Lemma \ref{algebra-lemma-criterion-normal}）。
ゆえにこの場合にも条件は確認された。
\medskip\noindent
場合 (2) については、Lemma \ref{more-algebra-lemma-J-2} の条件 (3) を用いる。
$R$ を Noether 完備局所環とする。$R \to R'$ が有限ならば、
$R'$ は Noether 完備局所環の積であることに注意する。Algebra, Lemma
\ref{algebra-lemma-quotient-complete-local} を参照せよ。
したがって Lemma \ref{more-algebra-lemma-intersection-regular-with-closed} により、
整域である Noether 完備局所環が J-0 であることを証明すれば十分であり、
これは Lemma \ref{more-algebra-lemma-complete-Noetherian-domain-J-0} である。
\medskip\noindent
場合 (4) についても、Lemma \ref{more-algebra-lemma-J-2} の条件 (3) を用いる。
実際、$R$ が次元 $1$ の Noether 局所環で $R \to R'$ が有限ならば、
$\Spec(R')$ は有限である。正則点集合は一般化に関して安定なので、
$R'$ は J-1 である。
\end{proof}





\section{形式的滑らかさと正則性}
\label{more-algebra-section-fs-regular}

\noindent
この節の表題は Andr\'e の Theorem \ref{more-algebra-theorem-fs-regular} を指している。

\begin{lemma}
\label{more-algebra-lemma-lift-derivation-through-fs}
$A \to B$ を Noether 局所環の局所準同型とする。$D : A \to A$ を導分とする。
$B$ は完備で、$A \to B$ は $\mathfrak m_B$-進位相について形式的に滑らかで
あると仮定する。このとき $D$ の延長 $D' : B \to B$ が存在する。
\end{lemma}

\begin{proof}
$B$ 上の双対数環を $B[\epsilon] = B[x]/(x^2)$ と記す。
環写像 $\psi : A \to B[\epsilon]$、
$a \mapsto a + \epsilon D(a)$ を考える。可換図式
$$
\xymatrix{
B \ar[r]_1 & B \\
A \ar[u] \ar[r]^\psi & B[\epsilon] \ar[u]
}
$$
を考える。Lemma \ref{more-algebra-lemma-lift-continuous} と $B/A$ の形式的滑らかさの仮定により、
この図式に入る写像 $\varphi : B \to B[\epsilon]$ を得る。
$\varphi(b) = b + \epsilon D'(b)$ と書けば、
$D' : B \to B$ が求める延長である。
\end{proof}

\begin{proposition}
\label{more-algebra-proposition-fs-regular}
$A \to B$ を Noether 完備局所環の局所準同型とする。
$k$ を $A$ の剰余体とし、$\overline{B} = B \otimes_A k$ を特殊ファイバーとする。
次は同値である。
\begin{enumerate}
\item $A \to B$ は正則である。
\item $A \to B$ は平坦で、$\overline{B}$ は $k$ 上幾何学的に正則である。
\item $A \to B$ は平坦で、$k \to \overline{B}$ は
$\mathfrak m_{\overline{B}}$-進位相について形式的に滑らかである。
\item $A \to B$ は $\mathfrak m_B$-進位相について形式的に滑らかである。
\end{enumerate}
\end{proposition}

\begin{proof}
(2)、(3)、(4) の同値性は Proposition
\ref{more-algebra-proposition-fs-flat-fibre-fs} で見た。(1) が (2) を含意することは明らかである。
そこで同値な条件 (2)、(3)、(4) が成り立つと仮定し、(1) を証明する。
\medskip\noindent
$\mathfrak p$ を $A$ の素イデアルとする。
$B \otimes_A \kappa(\mathfrak p)$ が $\kappa(\mathfrak p)$ 上
幾何学的に正則であることを示す。
Lemma \ref{more-algebra-lemma-base-change-fs} により、$A$ を $A/\mathfrak p$ で、
$B$ を $B/\mathfrak pB$ で置き換えてよい。
したがって $A$ は整域で、$\mathfrak p = (0)$ であると仮定してよい。
\medskip\noindent
Algebra, Lemma
\ref{algebra-lemma-complete-local-Noetherian-domain-finite-over-regular}
のように $A_0 \subset A$ を選ぶ。以後、同補題に述べられた性質を
特に断らずにすべて用いる。$A_0 \to A$ は剰余体上の同型を誘導し、
$B/\mathfrak m_A B$ は $A/\mathfrak m_A$ 上幾何学的に正則なので、
次の図式を見つけられる。
$$
\xymatrix{
C \ar[r] & B \\
A_0 \ar[r] \ar[u] & A \ar[u]
}
$$
ここで $A_0 \to C$ は $\mathfrak m_C$-進位相について形式的に滑らかで、
$B = C \otimes_{A_0} A$ である。Remark
\ref{more-algebra-remark-what-does-it-mean} を参照せよ。
（$A_0 \to A$ は有限なので、テンソル積における完備化は不要である。
Algebra, Lemma \ref{algebra-lemma-completion-tensor} を参照せよ。）
したがって $C \otimes_{A_0} K_0$ が、$A_0$ の分数体 $K_0$ 上の
幾何学的に正則な代数であることを示せば十分である。
\medskip\noindent
前段落の要点は、$k$ を体として $A = k[[x_1, \ldots, x_n]]$ であるか、
$\Lambda$ を Cohen 環として
$A = \Lambda[[x_1, \ldots, x_n]]$ であると仮定してよいことである。
この場合、Algebra, Lemma
\ref{algebra-lemma-flat-over-regular-with-regular-fibre} により
$B$ は正則環である。したがって $B \otimes_A K$ も正則環である
（ここで $K$ は $A$ の分数体）。ゆえに $K$ の標数が零ならば示された。
\medskip\noindent
残るのは、$A = k[[x_1, \ldots, x_n]]$ で、$k$ が標数 $p > 0$ の体である場合である。
$L/K$ を有限純非分離体拡大とする。$[L : K]$ に関する帰納法により、
$B \otimes_A L$ が正則であることを示す。基底の場合は $L = K$ で、
これは上で見た。$K \subset M \subset L$ を部分体とし、
$L$ が元 $f \in M$ の $p$ 乗根を添加して得られる
$M$ の次数 $p$ の拡大であるとする。
$A'$ を、分数体が $M$ となる $M$ の有限 $A$-部分代数とする。
分母を払うことで、$f \in A'$ と仮定してよい。
$A'' = A'[z]/(z^p -f)$ とおく。$A' \subset A''$ は有限で、
$A''$ の分数体は $L$ であることに注意する。
帰納法により、環 $B \otimes_A M$ は正則である。次が成り立つ。
$$
B \otimes_A L = B \otimes_A M[z]/(z^p - f)
$$
Lemma \ref{more-algebra-lemma-find-D} により、$D(f) \not = 0$ を満たす導分
$D : A' \to A'$ が存在する。Lemma \ref{more-algebra-lemma-descent-fs} により
$A' \to B \otimes_A A'$ は $\mathfrak m$-進位相について形式的に滑らかなので、
Lemma \ref{more-algebra-lemma-lift-derivation-through-fs} を用いて $D$ を導分
$D' : B \otimes_A A' \to B \otimes_A A'$ に延長できる。
$D(f)$ は $A' \subset M$ において非零なので、
$D'(f) = D(f)$ は $B \otimes_A M$ の単元であることに注意する。
したがって Lemma \ref{more-algebra-lemma-degree-p-extension-regular} により
$B \otimes_A L$ は正則であり、証明が終わる。
\end{proof}

\begin{theorem}[Andr\'e]
\label{more-algebra-theorem-fs-regular}
\begin{reference}
含意 (4) $\Rightarrow$ (1) は \cite{Andre-smooth} の主結果である。
\end{reference}
$A \to B$ を Noether 局所環の局所準同型とする。
$k$ を $A$ の剰余体とし、$\overline{B} = B \otimes_A k$ を特殊ファイバーとする。
$A \to A^\wedge$ は正則であると仮定する\footnote{たとえば、$A$ が G-環または
（準）excellent である場合。Sections \ref{more-algebra-section-G-ring} および
\ref{more-algebra-section-excellent} を参照せよ。}。
次は同値である。
\begin{enumerate}
\item $A \to B$ は正則である。
\item $A \to B$ は平坦で、$\overline{B}$ は $k$ 上幾何学的に正則である。
\item $A \to B$ は平坦で、$k \to \overline{B}$ は
$\mathfrak m_{\overline{B}}$-進位相について形式的に滑らかである。
\item $A \to B$ は $\mathfrak m_B$-進位相について形式的に滑らかである。
\end{enumerate}
\end{theorem}

\begin{proof}
(2)、(3)、(4) の同値性は Proposition
\ref{more-algebra-proposition-fs-flat-fibre-fs} で見た。(1) が (2) を含意することは明らかである。
そこで (4) が成り立つと仮定し、(1) を証明する。
\medskip\noindent
Lemma \ref{more-algebra-lemma-formally-smooth-completion} により、
$A^\wedge \to B^\wedge$ は形式的に滑らかである。
Proposition \ref{more-algebra-proposition-fs-regular} により、
$A^\wedge \to B^\wedge$ は正則である。
仮定により $A \to A^\wedge$ は正則なので、Lemma
\ref{more-algebra-lemma-regular-composition} により $A \to B^\wedge$ は正則である。
$B \to B^\wedge$ は忠実平坦なので、Lemma
\ref{more-algebra-lemma-regular-permanence} により $A \to B$ は正則であると結論する。
\end{proof}





\section{G-環}
\label{more-algebra-section-G-ring}

\noindent
$A$ を Noether 局所環とする。Section \ref{more-algebra-section-permanence-completion} では、
$A$ の性質のうち一部は $A$ の完備化 $A^\wedge$ に反映されるが、
すべてがそうではないことを見た。これをさらに調べるため、いくつか用語を導入する。
$A$ の素イデアル $\mathfrak q$ に対し、ファイバー環
% P02-E0883: the frozen source writes A/q; see ERRATA_CANDIDATES.jsonl.
$$
A^\wedge \otimes_A \kappa(\mathfrak q) =
(A^\wedge)_\mathfrak q/\mathfrak q(A^\wedge)_\mathfrak q =
(A/\mathfrak q)^\wedge \otimes_{A/q} \kappa(\mathfrak q)
$$
を $A$ の {\it 形式ファイバー} という。形式ファイバーは
$\kappa(\mathfrak q)$ 上の代数とみなす。したがって、
$A \to A^\wedge$ が正則環準同型であるための必要十分条件は、
すべての形式ファイバーが幾何学的に正則な代数であることである。

\begin{definition}
\label{more-algebra-definition-G-ring}
環 $R$ が {\it G-環} であるとは、$R$ が Noether であり、$R$ のすべての
素イデアル $\mathfrak p$ に対して環写像
$R_\mathfrak p \to (R_\mathfrak p)^\wedge$ が正則であることをいう。
\end{definition}

\noindent
上の議論により、$R$ が G-環であるための必要十分条件は、すべての局所環
$R_\mathfrak p$ が幾何学的に正則な形式ファイバーをもつことである。
$\mathbf{Q} \subset R$ ならば、形式ファイバーが正則であることだけを
確認すれば十分であることに注意する。
G-環条件の別の表し方を次の補題で述べる。

\begin{lemma}
\label{more-algebra-lemma-check-G-ring-easy}
$R$ を Noether 環とする。$R$ が G-環であるための必要十分条件は、すべての
素イデアル対 $\mathfrak q \subset \mathfrak p \subset R$ に対して代数
$$
(R/\mathfrak q)_\mathfrak p^\wedge \otimes_{R/\mathfrak q} \kappa(\mathfrak q)
$$
が $\kappa(\mathfrak q)$ 上幾何学的に正則であることである。
\end{lemma}

\begin{proof}
これは、$\kappa(\mathfrak q)$ 上の代数として
$$
R_\mathfrak p^\wedge \otimes_R \kappa(\mathfrak q) =
(R/\mathfrak q)_\mathfrak p^\wedge \otimes_{R/\mathfrak q} \kappa(\mathfrak q)
$$
であることから従う。
\end{proof}

\begin{lemma}
\label{more-algebra-lemma-G-ring-goes-up-quasi-finite}
$R \to R'$ を Noether 環の有限型写像とし、
$$
\xymatrix{
\mathfrak q' \ar[r] & \mathfrak p' \ar[r] & R' \\
\mathfrak q \ar[r] \ar@{-}[u] &
\mathfrak p \ar[r] \ar@{-}[u] & R \ar[u]
}
$$
を素イデアルとする。$R \to R'$ が $\mathfrak p'$ において準有限であると仮定する。
\begin{enumerate}
\item 形式ファイバー $R_\mathfrak p^\wedge \otimes_R \kappa(\mathfrak q)$ が
$\kappa(\mathfrak q)$ 上幾何学的に正則ならば、形式ファイバー
$(R'_{\mathfrak p'})^\wedge \otimes_{R'} \kappa(\mathfrak q')$ は
$\kappa(\mathfrak q')$ 上幾何学的に正則である。
\item $R_\mathfrak p$ の形式ファイバーが幾何学的に正則ならば、
$R'_{\mathfrak p'}$ の形式ファイバーも幾何学的に正則である。
\item $R \to R'$ が準有限で $R$ が G-環ならば、$R'$ は G-環である。
\end{enumerate}
\end{lemma}

\begin{proof}
(1) $\Rightarrow$ (2) $\Rightarrow$ (3) は明らかである。
$R_\mathfrak p^\wedge \otimes_R \kappa(\mathfrak q)$ が
$\kappa(\mathfrak q)$ 上幾何学的に正則であると仮定する。
Algebra, Lemma \ref{algebra-lemma-completion-at-quasi-finite-prime} により
$$
R_\mathfrak p^\wedge \otimes_R R'
=
(R'_{\mathfrak p'})^\wedge \times B
$$
となる $R_\mathfrak p^\wedge$-代数 $B$ が存在する。
したがって $R'_{\mathfrak p'} \to (R'_{\mathfrak p'})^\wedge$ は、
写像 $R_\mathfrak p \to R_\mathfrak p^\wedge$ の基底変換の直積因子である。
従って $(R'_{\mathfrak p'})^\wedge \otimes_{R'} \kappa(\mathfrak q')$ は
$$
R_\mathfrak p^\wedge \otimes_R R' \otimes_{R'} \kappa(\mathfrak q') =
R_\mathfrak p^\wedge \otimes_R \kappa(\mathfrak q)
\otimes_{\kappa(\mathfrak q)} \kappa(\mathfrak q').
$$
の直積因子である。基礎体の拡大は幾何学的正則性を保つので、結果が従う。
Algebra, Lemma \ref{algebra-lemma-geometrically-regular} を参照せよ。
\end{proof}

\begin{lemma}
\label{more-algebra-lemma-check-G-ring}
$R$ を Noether 環とする。$R$ が G-環であるための必要十分条件は、
任意の有限自由環写像 $R \to S$ に対して $S$ の形式ファイバーが
正則環であることである。
\end{lemma}

\begin{proof}
任意の有限自由環写像 $R \to S$ に対して環 $S$ が正則な形式ファイバーを
もつと仮定する。$\mathfrak q \subset \mathfrak p \subset R$ を素イデアルとし、
$\kappa(\mathfrak q) \subset L$ を有限純非分離拡大とする。
$R$ が G-環であることを示すには、
$$
R_\mathfrak p^\wedge \otimes_R \kappa(\mathfrak q)
\otimes_{\kappa(\mathfrak q)} L
$$
が正則環であることを示せば十分である。
$\mathfrak q' = \mathfrak qR'$ が素イデアルで、
$\kappa(\mathfrak q')$ が $\kappa(\mathfrak q)$ 上 $L$ と同型になるような
有限自由拡大 $R \to R'$ を選ぶ。Algebra, Lemma
\ref{algebra-lemma-finite-free-given-residue-field-extension} を参照せよ。
Algebra, Lemma \ref{algebra-lemma-completion-finite-extension} により
$$
R_\mathfrak p^\wedge \otimes_R R' = \prod (R'_{\mathfrak p_i'})^\wedge
$$
である。ここで $\mathfrak p_i'$ は $\mathfrak p$ の上にある $R'$ の
素イデアルである。したがって
$$
R_\mathfrak p^\wedge \otimes_R \kappa(\mathfrak q)
\otimes_{\kappa(\mathfrak q)} L =
R_\mathfrak p^\wedge \otimes_R R'
\otimes_{R'} \kappa(\mathfrak q')
=
\prod (R'_{\mathfrak p_i'})^\wedge
\otimes_{R'_{\mathfrak p'_i}} \kappa(\mathfrak q')
$$
である。仮定により右辺の環は正則なので、左辺の環も正則である。
したがって $R$ は G-環である。逆は Lemma
\ref{more-algebra-lemma-G-ring-goes-up-quasi-finite} から従う。
\end{proof}


\begin{lemma}
\label{more-algebra-lemma-helper-G-ring}
$k$ を標数 $p$ の体とする。
$A = k[[x_1, \ldots, x_n]][y_1, \ldots, y_m]$ とし、
$A$ の分数体を $K$ と記す。$\mathfrak p \subset A$ を素イデアルとする。
このとき $A_\mathfrak p^\wedge \otimes_A K$ は
$K$ 上幾何学的に正則である。
\end{lemma}

\begin{proof}
$L/K$ を有限純非分離体拡大とする。$[L : K]$ に関する帰納法により、
$A_\mathfrak p^\wedge \otimes L$ が正則であることを示す。
基底の場合は $L = K$ である。$A$ は正則なので、
$A_\mathfrak p^\wedge$ は正則であり
（Lemma \ref{more-algebra-lemma-completion-regular}）、
したがって局所化 $A_\mathfrak p^\wedge \otimes K$ は正則である。
$K \subset M \subset L$ を部分体とし、$L$ が元 $f \in M$ の
$p$ 乗根を添加して得られる $M$ の次数 $p$ の拡大であるとする。
$B$ を、分数体が $M$ となる $M$ の有限 $A$-部分代数とする。
分母を払うことで、$f \in B$ と仮定してよい。
$C = B[z]/(z^p -f)$ とおく。$B \subset C$ は有限で、
$C$ の分数体は $L$ であることに注意する。
$A \subset B \subset C$ は有限で、$L/M/K$ は純非分離なので、
$B$ または $C$ のすべての元は、そのある冪が $A$ に属する。
したがって $\mathfrak p$ の上にある素イデアル
$\mathfrak r \subset B$、resp.\ $\mathfrak q \subset C$ は一意である。
次に注意する。
$$
A_\mathfrak p^\wedge \otimes_A M = B_\mathfrak r^\wedge \otimes_B M
$$
Algebra, Lemma \ref{algebra-lemma-completion-finite-extension} を参照せよ。
帰納法により、この環は正則である。同様に
% P02-E0884: the frozen source indexes C by \mathfrak r; see ERRATA_CANDIDATES.jsonl.
$$
A_\mathfrak p^\wedge \otimes_A L =
C_\mathfrak r^\wedge \otimes_C L =
B_\mathfrak r^\wedge \otimes_B M[z]/(z^p - f)
$$
である。最後の等式は、$C = B[z]/(z^p - f)$ の完備化が
$B_\mathfrak r^\wedge[z]/(z^p -f)$ に等しいことによる。
Lemma \ref{more-algebra-lemma-find-D} により、$D(f) \not = 0$ を満たす導分
$D : B \to B$ が存在する。言い換えれば、$g = D(f)$ は $M$ の単元である！
Lemma \ref{more-algebra-lemma-derivation-extends} により、$D$ は $B_\mathfrak r$、
$B_\mathfrak r^\wedge$、および $B_\mathfrak r^\wedge \otimes_B M$ の導分へ
延長される（局所化、完備化、局所化を順に通して延長する）。
これは延長なので、最終的に $B_\mathfrak r^\wedge \otimes_B M$ の導分で、
$f$ を $g$ に写すものを得る。そして $g$ は環
$B_\mathfrak r^\wedge \otimes_B M$ の単元である。
したがって Lemma \ref{more-algebra-lemma-degree-p-extension-regular} により
$A_\mathfrak p^\wedge \otimes_A L$ は正則であり、証明が終わる。
\end{proof}

\begin{proposition}
\label{more-algebra-proposition-Noetherian-complete-G-ring}
Noether 完備局所環は G-環である。
\end{proposition}

\begin{proof}
$A$ を Noether 完備局所環とする。Lemma \ref{more-algebra-lemma-check-G-ring-easy} により、
$B = A/\mathfrak q$ が $B$ の極小素イデアル $(0)$ の上に
幾何学的に正則な形式ファイバーをもつことを確認すれば十分である。
したがって $A$ は整域であると仮定してよく、$A$ の極小素イデアル
$(0)$ の上の形式ファイバーについて条件を確認すれば十分である。
$A$ の分数体を $K$ とする。
\medskip\noindent
$A$ が $A_0$ 上有限となる正則完備局所部分環 $A_0 \subset A$ を選べる。
Algebra, Lemma
\ref{algebra-lemma-complete-local-Noetherian-domain-finite-over-regular}
を参照せよ。さらに $A_0$ は、体または Cohen 環上の冪級数環であると
仮定してよい。Lemma \ref{more-algebra-lemma-G-ring-goes-up-quasi-finite} により、
$A_0$ について結果を証明すれば十分である。
\medskip\noindent
$A$ が体または Cohen 環上の冪級数環であると仮定する。
$A$ は正則なので、局所化 $A_\mathfrak p$ は正則である
（Algebra, Definition \ref{algebra-definition-regular} とその直前の議論を参照せよ）。
したがって完備化 $A_\mathfrak p^\wedge$ は正則である。
Lemma \ref{more-algebra-lemma-completion-regular} を参照せよ。
ゆえにファイバー $A_{\mathfrak p}^\wedge \otimes_A K$ も、
$A_\mathfrak p^\wedge$ の局所化として正則である。
従って $K$ の標数が $0$ ならば示された。正標数の場合は
$A = k[[x_1, \ldots, x_d]]$ の場合であり、これは Lemma
\ref{more-algebra-lemma-helper-G-ring} の特殊な場合である。
\end{proof}


\begin{lemma}
\label{more-algebra-lemma-check-G-ring-maximal-ideals}
$R$ を Noether 環とする。$R$ が G-環であるための必要十分条件は、
$R$ のすべての極大イデアル $\mathfrak m$ に対して
$R_\mathfrak m$ が幾何学的に正則な形式ファイバーをもつことである。
\end{lemma}

\begin{proof}
すべての極大イデアル $\mathfrak m$ に対して
$R_\mathfrak m \to R_\mathfrak m^\wedge$ が正則であると仮定する。
$\mathfrak p$ を $R$ の素イデアルとし、
$\mathfrak p \subset \mathfrak m$ となる $R$ の極大イデアルを選ぶ。
$R_\mathfrak m \to R_\mathfrak m^\wedge$ は忠実平坦なので、
$\mathfrak pR_\mathfrak m$ の上にある
$R_\mathfrak m^\wedge$ の素イデアル $\mathfrak p'$ を選べる。
可換図式
$$
\xymatrix{
R_\mathfrak m^\wedge \ar[r] &
(R_\mathfrak m^\wedge)_{\mathfrak p'} \ar[r] &
(R_\mathfrak m^\wedge)_{\mathfrak p'}^\wedge
\\
R_\mathfrak m \ar[u] \ar[r] & R_\mathfrak p \ar[u] \ar[r] &
R_\mathfrak p^\wedge \ar[u]
}
$$
を考える。仮定により環写像
$R_\mathfrak m \to R_\mathfrak m^\wedge$ は正則である。
Proposition \ref{more-algebra-proposition-Noetherian-complete-G-ring} により
$(R_\mathfrak m^\wedge)_{\mathfrak p'} \to
(R_\mathfrak m^\wedge)_{\mathfrak p'}^\wedge$ は正則である。
局所化
$R_\mathfrak m^\wedge \to (R_\mathfrak m^\wedge)_{\mathfrak p'}$
は正則である。したがって Lemma \ref{more-algebra-lemma-regular-composition} により
$R_\mathfrak m \to (R_\mathfrak m^\wedge)_{\mathfrak p'}^\wedge$
は正則である。これは局所化 $R_\mathfrak p$ を経由するので、環写像
$R_\mathfrak p \to (R_\mathfrak m^\wedge)_{\mathfrak p'}^\wedge$
も正則である。したがって Lemma \ref{more-algebra-lemma-regular-permanence} を適用すると
$R_\mathfrak p \to R_\mathfrak p^\wedge$ が正則であることが分かる。
\end{proof}

\begin{lemma}
\label{more-algebra-lemma-henselization-G-ring}
$R$ を G-環である Noether 局所環とする。このとき Hensel 化 $R^h$ と
強 Hensel 化 $R^{sh}$ は G-環である。
\end{lemma}

\begin{proof}
Lemma \ref{more-algebra-lemma-check-G-ring-maximal-ideals} の判定条件を用いる。
$\mathfrak q \subset R^h$ を素イデアルとし、
$\mathfrak p = R \cap \mathfrak q$ とおく。
$\mathfrak q_1 = \mathfrak q$ とおき、
$\mathfrak q_2, \ldots, \mathfrak q_t$ を $\mathfrak p$ の上にある
$R^h$ の他の素イデアルとする。このとき Lemma
\ref{more-algebra-lemma-fibres-henselization} により
$R^h \otimes_R \kappa(\mathfrak p) =
\prod\nolimits_{i = 1, \ldots, t} \kappa(\mathfrak q_i)$ である。
$(R^h)^\wedge = R^\wedge$ であること
（Lemma \ref{more-algebra-lemma-henselization-noetherian}）を用いると、
$$
\prod\nolimits_{i = 1, \ldots, t}
(R^h)^\wedge \otimes_{R^h} \kappa(\mathfrak q_i) =
(R^h)^\wedge \otimes_{R^h} (R^h \otimes_R \kappa(\mathfrak p)) =
R^\wedge \otimes_R \kappa(\mathfrak p)
$$
を得る。したがって仮定により
$(R^h)^\wedge \otimes_{R^h} \kappa(\mathfrak q_i)$ は
$\kappa(\mathfrak p)$ 上幾何学的に正則である。
$\kappa(\mathfrak q_i)$ は $\kappa(\mathfrak p)$ 上分離代数的なので、
Algebra, Lemma
\ref{algebra-lemma-geometrically-regular-over-separable-algebraic} により
$(R^h)^\wedge \otimes_{R^h} \kappa(\mathfrak q_i)$ は
$\kappa(\mathfrak q_i)$ 上幾何学的に正則である。
\medskip\noindent
$\mathfrak r \subset R^{sh}$ を素イデアルとし、
$\mathfrak p = R \cap \mathfrak r$ とおく。
$\mathfrak r_1 = \mathfrak r$ とおき、
$\mathfrak r_2, \ldots, \mathfrak r_s$ を $\mathfrak p$ の上にある
$R^{sh}$ の他の素イデアルとする。このとき Lemma
\ref{more-algebra-lemma-fibres-henselization} により
$R^{sh} \otimes_R \kappa(\mathfrak p) =
\prod\nolimits_{i = 1, \ldots, s} \kappa(\mathfrak r_i)$ である。
すると
$$
\prod\nolimits_{i = 1, \ldots, s}
(R^{sh})^\wedge \otimes_{R^{sh}} \kappa(\mathfrak r_i) =
(R^{sh})^\wedge \otimes_{R^{sh}} (R^{sh} \otimes_R \kappa(\mathfrak p)) =
(R^{sh})^\wedge \otimes_R \kappa(\mathfrak p)
$$
であることが分かる。Lemma \ref{more-algebra-lemma-henselization-noetherian} により
$R^\wedge \to (R^{sh})^\wedge$ は
$\mathfrak m_{(R^{sh})^\wedge}$-進位相について形式的に滑らかであることに注意する。
したがって Proposition \ref{more-algebra-proposition-fs-regular} により
$R^\wedge \to (R^{sh})^\wedge$ は正則である。
仮定により $R^\wedge \otimes_R \kappa(\mathfrak p)$ は
$\kappa(\mathfrak p)$ 上正則なので、Lemma
\ref{more-algebra-lemma-regular-composition} により
$(R^{sh})^\wedge \otimes_{R^{sh}} \kappa(\mathfrak r_i)$ は
$\kappa(\mathfrak p)$ 上正則であると結論する。
$\kappa(\mathfrak r_i)$ は $\kappa(\mathfrak p)$ 上分離代数的なので、
Algebra, Lemma
\ref{algebra-lemma-geometrically-regular-over-separable-algebraic} により
$(R^{sh})^\wedge \otimes_{R^{sh}} \kappa(\mathfrak r_i)$ は
$\kappa(\mathfrak r_i)$ 上幾何学的に正則である。
\end{proof}


\begin{lemma}
\label{more-algebra-lemma-another-helper-G-ring}
$p$ を素数とする。$A$ を、分数体 $K$ が標数 $p$ をもつ
Noether 完備局所整域とする。$\mathfrak q \subset A[x]$ を
$A$ の極大イデアルの上にある極大イデアルとし、
$(0) \not = \mathfrak r \subset \mathfrak q$ を
$(0) \subset A$ の上にある素イデアルとする。このとき
$A[x]_\mathfrak q^\wedge \otimes_{A[x]} \kappa(\mathfrak r)$ は
$\kappa(\mathfrak r)$ 上幾何学的に正則である。
\end{lemma}

\begin{proof}
$K \subset \kappa(\mathfrak r)$ は有限であることに注意する。
したがって、有限純非分離拡大 $L/\kappa(\mathfrak r)$ が与えられると、
$\kappa(\mathfrak r) \otimes_A B$ が $L$ の上に全射となるような
Noether 完備局所整域の有限拡大 $A \subset B$ が存在する。
実際、分数体が $L$ となる有限 $A$-部分代数 $B \subset L$ を取ればよい。
写像
$B[x] = A[x] \otimes_A B \to \kappa(\mathfrak r) \otimes_A B \to L$
の核を $\mathfrak r' \subset B[x]$ と記すと、
$\kappa(\mathfrak r') = L$ である。このとき
$$
A[x]_\mathfrak q^\wedge \otimes_{A[x]} L =
A[x]_\mathfrak q^\wedge \otimes_{A[x]} B[x]
\otimes_{B[x]} \kappa(\mathfrak r') =
\prod B[x]_{\mathfrak q_i}^\wedge \otimes_{B[x]} \kappa(\mathfrak r')
$$
である。ここで $\mathfrak q_1, \ldots, \mathfrak q_t$ は
$\mathfrak q$ の上にある $B[x]$ の素イデアルである。
Algebra, Lemma \ref{algebra-lemma-completion-finite-extension} を参照せよ。
したがって、環
$B[x]_{\mathfrak q_i}^\wedge \otimes_{B[x]} \kappa(\mathfrak r')$
が正則であることを証明すれば十分である。
これは $K = \kappa(\mathfrak r)$ という特殊な場合に
$A[x]_\mathfrak q^\wedge \otimes_{A[x]} \kappa(\mathfrak r)$ が
正則であることを示す問題に帰着する。
\medskip\noindent
$K = \kappa(\mathfrak r)$ と仮定する。この場合、ある $f \in K$ に対して
$\mathfrak r K[x]$ は $x - f$ で生成され、次が成り立つ。
$$
A[x]_\mathfrak q^\wedge \otimes_{A[x]} \kappa(\mathfrak r)
=
(A[x]_\mathfrak q^\wedge \otimes_A K)/(x - f)
$$
$A[x]$ の導分 $D = \text{d}/\text{d}x$ は $K[x]$ に延長され、
$x - f$ を $K[x]$ の単元に写す。さらに Lemma
\ref{more-algebra-lemma-derivation-extends} により、$D$ は
$A[x]_\mathfrak q^\wedge \otimes_A K$ に延長される。
$A \to A[x]_\mathfrak q^\wedge$ は形式的に滑らかなので
（Lemmas \ref{more-algebra-lemma-formally-smooth} および
\ref{more-algebra-lemma-formally-smooth-completion} を参照せよ）、
Proposition \ref{more-algebra-proposition-fs-regular} により
$A[x]_\mathfrak q^\wedge \otimes_A K$ は正則である
（この特別な場合、同命題の証明の議論は大幅に簡単になる）。
Lemma \ref{more-algebra-lemma-quotient-regular} により結論を得る。
\end{proof}


\begin{proposition}
\label{more-algebra-proposition-finite-type-over-G-ring}
$R$ を G-環とする。$R \to S$ が本質的有限型ならば、$S$ は G-環である。
\end{proposition}

\begin{proof}
G-環であることは局所環の性質なので、G-環の局所化が G-環であることは明らかである。
逆に、すべての素イデアルにおける局所化が G-環ならば、その環は G-環である。
したがって、任意の有限型 $R$-代数 $S$ と $S$ の任意の素イデアル
$\mathfrak q$ に対して $S_\mathfrak q$ が G-環であることを示せば十分である。
$S$ を $R[x_1, \ldots, x_n]$ の商として書くと、Lemma
\ref{more-algebra-lemma-G-ring-goes-up-quasi-finite} により、
$R[x_1, \ldots, x_n]$ が G-環であることを証明すれば十分である。
$n$ に関する帰納法により、$R[x]$ が G-環であることを証明すれば十分である。
$\mathfrak q \subset R[x]$ を極大イデアルとする。Lemma
\ref{more-algebra-lemma-check-G-ring-maximal-ideals} により、
$$
R[x]_\mathfrak q \longrightarrow R[x]_\mathfrak q^\wedge
$$
が正則であることを示せば十分である。
$\mathfrak q$ が $\mathfrak p \subset R$ の上にあるならば、
$R$ を $R_\mathfrak p$ で置き換えてよい。
したがって $R$ は極大イデアル $\mathfrak m$ をもつ Noether 局所 G-環で、
$\mathfrak q \subset R[x]$ は $\mathfrak m$ の上にあると仮定してよい。
$\mathfrak q$ の上にある素イデアル $\mathfrak q' \subset R^\wedge[x]$
は一意であることに注意する。図式
$$
\xymatrix{
R[x]_\mathfrak q^\wedge \ar[r] &
(R^\wedge[x]_{\mathfrak q'})^\wedge \\
R[x]_\mathfrak q \ar[r] \ar[u] & R^\wedge[x]_{\mathfrak q'} \ar[u]
}
$$
を考える。$R$ は G-環なので、下の水平矢印は正則である
（正則環写像 $R \to R^\wedge$ の基底変換の局所化だからである）。
右の垂直矢印が正則であることを証明できたとする。
すると合成 $R[x]_\mathfrak q \to (R^\wedge[x]_{\mathfrak q'})^\wedge$
が正則であり、従って Lemma \ref{more-algebra-lemma-regular-permanence} により
左の垂直矢印は正則である。
ゆえに $R$ は Noether 完備局所環で、$\mathfrak q$ は $R$ の極大イデアルの
上にある素イデアルであると仮定してよい。
\medskip\noindent
$R$ を Noether 完備局所環とし、$\mathfrak q \subset R[x]$ を
$R$ の極大イデアルの上にある極大イデアルとする。
$\mathfrak r \subset \mathfrak q$ を素イデアルとする。
$R[x]_\mathfrak q^\wedge \otimes_{R[x]} \kappa(\mathfrak r)$ が
$\kappa(\mathfrak r)$ 上幾何学的に正則な代数であることを示したい。
$\mathfrak p = R \cap \mathfrak r$ とおく。Lemma
\ref{more-algebra-lemma-check-G-ring-easy} により、$R$ を $R/\mathfrak p$ で、
$\mathfrak q$ と $\mathfrak r$ を $R/\mathfrak p[x]$ におけるそれらの像で
置き換えられる。したがって $R$ は整域で
$\mathfrak r \cap R = (0)$ であると仮定してよい。
\medskip\noindent
Algebra, Lemma
\ref{algebra-lemma-complete-local-Noetherian-domain-finite-over-regular}
により、$R$ が $R_0$ 上有限となる正則環 $R_0 \subset R$ を見つけられる。
Lemma \ref{more-algebra-lemma-G-ring-goes-up-quasi-finite} を適用すると、上の仮定に加えて
$R$ が正則完備局所環である場合に
$R[x]_\mathfrak q^\wedge \otimes_{R[x]} \kappa(\mathfrak r)$ が
$\kappa(\mathfrak r)$ 上幾何学的に正則であることを証明すれば十分である。
\medskip\noindent
いま $R$ は正則完備局所環で、
$\mathfrak r \subset \mathfrak q \subset R[x]$、
$(0) = R \cap \mathfrak r$ であり、$\mathfrak q$ は $R$ の極大イデアルの
上にある極大イデアルである。
$R$ は正則なので環 $R[x]$ は正則である
（Algebra, Lemma \ref{algebra-lemma-regular-goes-up}）。
したがって局所化 $R[x]_\mathfrak q$ は正則である。
ゆえに完備化 $R[x]_\mathfrak q^\wedge$ は正則である。
Lemma \ref{more-algebra-lemma-completion-regular} を参照せよ。
従ってファイバー
$R[x]_{\mathfrak q}^\wedge \otimes_{R[x]} \kappa(\mathfrak r)$ も、
$R[x]_\mathfrak q^\wedge$ の局所化として正則である。
ゆえに $R$ の分数体の標数が $0$ ならば示された。
\medskip\noindent
$R$ の標数が正ならば、$R = k[[x_1, \ldots, x_n]]$ である。
この場合、議論を二つの場合に分ける。
\begin{enumerate}
\item $\mathfrak r = (0)$ の場合。結果は Lemma
\ref{more-algebra-lemma-helper-G-ring} の直接の帰結である。
\item $\mathfrak r \not = (0)$ の場合。これは Lemma
\ref{more-algebra-lemma-another-helper-G-ring} である。
\end{enumerate}
\end{proof}

\begin{remark}
\label{more-algebra-remark-G-does-not-survive-completion}
$R$ を G-環とし、$I \subset R$ をイデアルとする。
一般に、$I$-進完備化 $R^\wedge$ が G-環であるとは限らない。
Nishimura が \cite{Nishimura} で例を与えた。この例の一般化と、
ある意味での明確化は \cite{Dumitrescu} の最終節にある。
\end{remark}

\begin{proposition}
\label{more-algebra-proposition-ubiquity-G-ring}
次の種類の環は G-環である。
\begin{enumerate}
\item 体、
\item Noether 完備局所環、
\item $\mathbf{Z}$、
\item 分数体の標数が 0 である Dedekind 整域、
\item 上のいずれかの有限型環拡大。
\end{enumerate}
\end{proposition}

\begin{proof}
体、$\mathbf{Z}$、および標数 0 の Dedekind 整域については、
これは定義と、離散付値環の完備化が離散付値環であることから直ちに従う。
Noether 完備局所環は
Proposition \ref{more-algebra-proposition-Noetherian-complete-G-ring} により G-環である。
有限型上位環に関する主張は
Proposition \ref{more-algebra-proposition-finite-type-over-G-ring} である。
\end{proof}

\begin{lemma}
\label{more-algebra-lemma-henselian-local-limit-G-rings}
$(A, \mathfrak m)$ を Hensel 局所環とする。
このとき $A$ は、局所遷移写像を持つ
Hensel 局所 G-環の系のフィルター付き余極限である。
\end{lemma}

\begin{proof}
$A = \colim A_i$ と、有限型 $\mathbf{Z}$-代数のフィルター付き余極限として書く。
$\mathfrak p_i$ を $\mathfrak m$ の下にある $A_i$ の素イデアルとする。
$A_i$ を $\mathfrak p_i$ における $A_i$ の局所化で置き換えてよい。
そのとき $A_i$ は Noether 局所 G-環である
(Proposition \ref{more-algebra-proposition-ubiquity-G-ring})。
Lemma \ref{more-algebra-lemma-henselization-colimit} により
$A = \colim A_i^h$ である。
Lemma \ref{more-algebra-lemma-henselization-G-ring} により、環 $A_i^h$ は G-環である。
\end{proof}

\begin{lemma}
\label{more-algebra-lemma-map-G-ring-to-completion-regular}
\begin{reference}
\cite[Theorem 79]{MatCA}
\end{reference}
$A$ を G-環とする。$I \subset A$ をイデアルとし、
$A^\wedge$ を $I$ に関する $A$ の完備化とする。
このとき $A \to A^\wedge$ は正則である。
\end{lemma}

\begin{proof}
環写像 $A \to A^\wedge$ は
Algebra, Lemma \ref{algebra-lemma-completion-flat} により平坦である。
環 $A^\wedge$ は
Algebra, Lemma \ref{algebra-lemma-completion-Noetherian-Noetherian} により Noether である。
したがって、Lemma \ref{more-algebra-lemma-regular-local} の第三条件を確かめれば十分である。
$\mathfrak m \subset A$ の上にある極大イデアルを
$\mathfrak m' \subset A^\wedge$ とする。
Algebra, Lemma \ref{algebra-lemma-radical-completion} により
$IA^\wedge \subset \mathfrak m'$ である。
$A^\wedge/IA^\wedge = A/I$ であるから、
$I \subset \mathfrak m$、$\mathfrak m/I = \mathfrak m'/IA^\wedge$、および
$A/\mathfrak m = A^\wedge/\mathfrak m'$ を得る。
$A^\wedge/\mathfrak m'$ は体なので、$\mathfrak m$ も極大イデアルであると結論される。
そこで $A_\mathfrak m \to A^\wedge_{\mathfrak m'}$ は、剰余体を同一視し、かつ
$\mathfrak m A^\wedge_{\mathfrak m'} = \mathfrak m'A^\wedge_{\mathfrak m'}$
を満たす Noether 局所環の平坦局所環準同型である。
したがって完備局所環上の同型を誘導する。
Lemma \ref{more-algebra-lemma-flat-unramified} を参照せよ。
$(A_\mathfrak m)^\wedge$ を、$A_\mathfrak m$ の極大イデアルに関する完備化とする。
環写像
$$
(A^\wedge)_{\mathfrak m'} \to
((A^\wedge)_{\mathfrak m'})^\wedge = (A_\mathfrak m)^\wedge
$$
は忠実平坦である
(Algebra, Lemma \ref{algebra-lemma-completion-faithfully-flat})。
ゆえに環写像
$$
A_\mathfrak m \to (A^\wedge)_{\mathfrak m'} \to (A_\mathfrak m)^\wedge
$$
に Lemma \ref{more-algebra-lemma-regular-permanence} を適用できる。
$A$ が G-環であるため
$A_\mathfrak m \to (A_\mathfrak m)^\wedge$ は正則であり、これで結論が従う。
\end{proof}

\begin{lemma}
\label{more-algebra-lemma-henselization-pair-G-ring}
\begin{slogan}
G-環であることは、イデアルに沿う Hensel 化で保たれる
\end{slogan}
\begin{reference}
\cite[Theorem 5.3 i)]{Greco}
\end{reference}
$A$ を G-環とし、$I \subset A$ をイデアルとする。
$(A^h, I^h)$ を対 $(A, I)$ の Hensel 化とする。
Lemma \ref{more-algebra-lemma-henselization} を参照せよ。
このとき $A^h$ は G-環である。
\end{lemma}

\begin{proof}
$\mathfrak m^h \subset A^h$ を極大イデアルとする。
Lemma \ref{more-algebra-lemma-check-G-ring-maximal-ideals} により、
$A^h_{\mathfrak m^h}$ からその完備化への写像が幾何学的に正則なファイバーを持つことを示さなければならない。
$\mathfrak m$ を $A$ における $\mathfrak m^h$ の逆像とする。
$(A^h, I^h)$ は Hensel 対なので、
$I^h \subset \mathfrak m^h$ であり、したがって $I \subset \mathfrak m$ であることに注意せよ。
$A^h$ は Noether で、$I^h = IA^h$ であり、さらに $A \to A^h$ は
$I$-進完備化上の同型を誘導することを想起せよ。
Lemma \ref{more-algebra-lemma-henselization-Noetherian-pair} を参照せよ。
すると Noether 局所環の局所準同型
$$
A_\mathfrak m \to A^h_{\mathfrak m^h}
$$
は、Lemma \ref{more-algebra-lemma-flat-unramified} により極大イデアルにおける完備化上の同型を誘導する
(詳細は省略する)。
$\mathfrak q \subset A_\mathfrak m$ の上にある
$A^h_{\mathfrak m^h}$ の素イデアルを $\mathfrak q^h$ とする。
$\mathfrak q_1 = \mathfrak q^h$ とおき、
$\mathfrak q_2, \ldots, \mathfrak q_t$ を $\mathfrak q$ の上にある $A^h$ の他の素イデアルとすると、
$$
A^h \otimes_A \kappa(\mathfrak q) =
\prod\nolimits_{i = 1, \ldots, t} \kappa(\mathfrak q_i)
$$
である。Lemma \ref{more-algebra-lemma-filtered-colimit-etale-noetherian-fibres} を参照せよ。
上で論じたように
$(A^h)_{\mathfrak m^h}^\wedge = (A_\mathfrak m)^\wedge$ であることを用いると、
$$
\prod\nolimits_{i = 1, \ldots, t}
(A^h_{\mathfrak m^h})^\wedge \otimes_{A^h_{\mathfrak m^h}}
\kappa(\mathfrak q_i) =
(A^h_{\mathfrak m^h})^\wedge \otimes_{A^h_{\mathfrak m^h}}
(A^h_{\mathfrak m^h} \otimes_{A_{\mathfrak m}} \kappa(\mathfrak q)) =
(A_{\mathfrak m})^\wedge \otimes_{A_{\mathfrak m}} \kappa(\mathfrak q)
$$
を得る。したがって、その一つの成分である環
$$
(A^h_{\mathfrak m^h})^\wedge \otimes_{A^h_{\mathfrak m^h}}
\kappa(\mathfrak q^h)
$$
は、$A$ に関する仮定により $\kappa(\mathfrak q)$ 上幾何学的に正則である。
$\kappa(\mathfrak q^h)$ は $\kappa(\mathfrak q)$ 上分離代数的なので、
Algebra, Lemma
\ref{algebra-lemma-geometrically-regular-over-separable-algebraic} から、
$$
(A^h_{\mathfrak m^h})^\wedge \otimes_{A^h_{\mathfrak m^h}}
\kappa(\mathfrak q^h)
$$
は望まれた通り $\kappa(\mathfrak q^h)$ 上幾何学的に正則である。
\end{proof}







\section{形式的ファイバーの性質}
\label{more-algebra-section-properties-formal-fibres}

\noindent
% P02-E0885: frozen English has the malformed phrase "for to be able" at source line 12818.
本節では、Noether 環の形式的ファイバーの性質について適切に論じられるように、
Section \ref{more-algebra-section-G-ring} の議論のいくつかをやり直す。

\medskip\noindent
$P$ を、$k$ が体で $R$ が Noether であるような環写像
$k \to R$ の性質とする。
$B$ が Noether である環準同型 $A \to B$ のファイバーが
$P$ を持つとは、$A$ のすべての素イデアル $\mathfrak q$ に対して
$\kappa(\mathfrak q) \to B \otimes_A \kappa(\mathfrak q)$ が $P$ を持つこととする。
以下では次の主張を用いる。
\begin{enumerate}
\item[(A)] 有限生成体拡大 $k'/k$ に対して
$P(k \to R) \Rightarrow P(k' \to R \otimes_k k')$、
\item[(B)] $P(k \to R_\mathfrak p),\ \forall \mathfrak p \in \Spec(R)
\Leftrightarrow P(k \to R)$、
\item[(C)] Noether 環の平坦写像 $A \to B \to C$ が与えられたとき、
$A \to B$ のファイバーが $P$ を持ち、$B \to C$ が正則ならば、
$A \to C$ のファイバーは $P$ を持つ、
\item[(D)] Noether 環の平坦写像 $A \to B \to C$ が与えられたとき、
$A \to C$ のファイバーが $P$ を持ち、$B \to C$ が忠実平坦ならば、
$A \to B$ のファイバーは $P$ を持つ、
\item[(E)] $k \to k' \to R$ が与えられ、$R$ が Noether、
$k'/k$ が分離代数的で $P(k \to R)$ が成り立つならば、
$P(k' \to R)$ が成り立つ。および
% P02-E0886: frozen source item (F) is an unfinished editorial placeholder.
\item[(F)] ここにさらに追加する。
\end{enumerate}
Noether 局所環 $A$ が与えられたとき、
``{\it $A$ の形式的ファイバーが $P$ を持つ}''
とは、$A \to A^\wedge$ のファイバーが $P$ を持つこととする。
{\it $R$ が $P$-環である}とは、$R$ が Noether であり、
$R$ のすべての素イデアル $\mathfrak p$ に対して
$R_\mathfrak p$ の形式的ファイバーが $P$ を持つこととする。

\begin{lemma}
\label{more-algebra-lemma-check-P-ring-easy}
$R$ を Noether 環とし、$P$ を上のような性質とする。
このとき $R$ が $P$-環であることと、
素イデアルのすべての対 $\mathfrak q \subset \mathfrak p \subset R$ に対して
$\kappa(\mathfrak q)$-代数
$$
(R/\mathfrak q)_\mathfrak p^\wedge \otimes_{R/\mathfrak q} \kappa(\mathfrak q)
$$
が性質 $P$ を持つことは同値である。
\end{lemma}

\begin{proof}
これは、$\kappa(\mathfrak q)$-代数として
$$
R_\mathfrak p^\wedge \otimes_R \kappa(\mathfrak q) =
(R/\mathfrak q)_\mathfrak p^\wedge \otimes_{R/\mathfrak q} \kappa(\mathfrak q)
$$
が成り立つことから従う。
\end{proof}

\begin{lemma}
\label{more-algebra-lemma-P-local}
$R \to \Lambda$ を Noether 環の準同型とする。
$P$ が性質 (B) を持つと仮定する。次は同値である。
\begin{enumerate}
\item $R \to \Lambda$ のファイバーは $P$ を持つ。
\item $\mathfrak p \subset R$ の上にあるすべての $\mathfrak q \subset \Lambda$ に対して、
$R_\mathfrak p \to \Lambda_\mathfrak q$ のファイバーは $P$ を持つ。
\item $R$ の $\mathfrak m$ の上にあるすべての極大イデアル
$\mathfrak m' \subset \Lambda$ に対して、
$R_\mathfrak m \to \Lambda_{\mathfrak m'}$ のファイバーは $P$ を持つ。
\end{enumerate}
\end{lemma}

\begin{proof}
$\mathfrak p \subset R$ を素イデアルとする。
$\mathfrak p$ 上のファイバーは環 $\Lambda \otimes_R \kappa(\mathfrak p)$ であり、
そのスペクトルは $\mathfrak p$ の上にある素イデアル $\mathfrak q$ からなる
$\Spec(\Lambda)$ の部分集合へ全単射的に写る。
Algebra, Remark \ref{algebra-remark-fundamental-diagram} を参照せよ。
そのような素イデアル $\mathfrak q$ に対し、
極大イデアル $\mathfrak q \subset \mathfrak m'$ を選び、
$\mathfrak m = R \cap \mathfrak m'$ とおく。
このとき $\mathfrak p \subset \mathfrak m$ であり、
$$
(\Lambda \otimes_R \kappa(\mathfrak p))_\mathfrak q \cong
(\Lambda_{\mathfrak m'} \otimes_{R_\mathfrak m}
\kappa(\mathfrak p))_\mathfrak q
$$
は $\kappa(\mathfrak q)$-代数の同型である。
したがって (B) により性質 $P$ は局所環上で確かめられるから、
(1), (2), (3) は同値である。
\end{proof}

\begin{lemma}
\label{more-algebra-lemma-P-ring-goes-up-quasi-finite}
$R \to R'$ を Noether 環の有限型写像とし、
$$
\xymatrix{
\mathfrak q' \ar[r] & \mathfrak p' \ar[r] & R' \\
\mathfrak q \ar[r] \ar@{-}[u] &
\mathfrak p \ar[r] \ar@{-}[u] & R \ar[u]
}
$$
を素イデアルとする。
$R \to R'$ は $\mathfrak p'$ で準有限であると仮定する。
$P$ は (A) と (B) を満たすと仮定する。
\begin{enumerate}
\item $\kappa(\mathfrak q) \to
R_\mathfrak p^\wedge \otimes_R \kappa(\mathfrak q)$ が $P$ を持つならば、
$\kappa(\mathfrak q') \to
(R'_{\mathfrak p'})^\wedge \otimes_{R'} \kappa(\mathfrak q')$ は $P$ を持つ。
\item $R_\mathfrak p$ の形式的ファイバーが $P$ を持つならば、
$R'_{\mathfrak p'}$ の形式的ファイバーは $P$ を持つ。
\item $R \to R'$ が準有限で $R$ が $P$-環ならば、$R'$ は $P$-環である。
\end{enumerate}
\end{lemma}

\begin{proof}
(1) $\Rightarrow$ (2) $\Rightarrow$ (3) は明らかである。
$P$ が
$\kappa(\mathfrak q) \to R_\mathfrak p^\wedge \otimes_R \kappa(\mathfrak q)$
に対して成り立つと仮定する。
Algebra, Lemma \ref{algebra-lemma-completion-at-quasi-finite-prime} により、
ある $R_\mathfrak p^\wedge$-代数 $B$ に対して
$$
R_\mathfrak p^\wedge \otimes_R R'
=
(R'_{\mathfrak p'})^\wedge \times B
$$
である。したがって
$R'_{\mathfrak p'} \to (R'_{\mathfrak p'})^\wedge$ は、
写像 $R_\mathfrak p \to R_\mathfrak p^\wedge$ の基底変換の一成分である。
それゆえ $(R'_{\mathfrak p'})^\wedge \otimes_{R'} \kappa(\mathfrak q')$ は
$$
R_\mathfrak p^\wedge \otimes_R R' \otimes_{R'} \kappa(\mathfrak q') =
R_\mathfrak p^\wedge \otimes_R \kappa(\mathfrak q)
\otimes_{\kappa(\mathfrak q)} \kappa(\mathfrak q').
$$
の一成分である。従って結果は $P$ に関する仮定から従う。
\end{proof}

\begin{lemma}
\label{more-algebra-lemma-check-P-ring-maximal-ideals}
$R$ を Noether 環とする。$P$ は (C) と (D) を満たすと仮定する。
このとき $R$ が $P$-環であることと、
$R$ のすべての極大イデアル $\mathfrak m$ に対して
$R_\mathfrak m$ の形式的ファイバーが $P$ を持つことは同値である。
\end{lemma}

\begin{proof}
$R$ のすべての極大イデアル $\mathfrak m$ に対して、
$R_\mathfrak m$ の形式的ファイバーが $P$ を持つと仮定する。
$R$ の素イデアルを $\mathfrak p$ とし、極大イデアル
$\mathfrak p \subset \mathfrak m$ を選ぶ。
$R_\mathfrak m \to R_\mathfrak m^\wedge$ は忠実平坦なので、
$\mathfrak pR_\mathfrak m$ の上にある $R_\mathfrak m^\wedge$ の素イデアル $\mathfrak p'$ を選べる。
可換図式
$$
\xymatrix{
R_\mathfrak m^\wedge \ar[r] &
(R_\mathfrak m^\wedge)_{\mathfrak p'} \ar[r] &
(R_\mathfrak m^\wedge)_{\mathfrak p'}^\wedge
\\
R_\mathfrak m \ar[u] \ar[r] & R_\mathfrak p \ar[u] \ar[r] &
R_\mathfrak p^\wedge \ar[u]
}
$$
を考える。
仮定により、環写像 $R_\mathfrak m \to R_\mathfrak m^\wedge$ のファイバーは $P$ を持つ。
Proposition \ref{more-algebra-proposition-Noetherian-complete-G-ring} により、
$(R_\mathfrak m^\wedge)_{\mathfrak p'} \to
(R_\mathfrak m^\wedge)_{\mathfrak p'}^\wedge$ は正則である。
局所化 $R_\mathfrak m^\wedge \to (R_\mathfrak m^\wedge)_{\mathfrak p'}$ は正則である。
従って Lemma \ref{more-algebra-lemma-regular-composition} により、
$R_\mathfrak m^\wedge \to (R_\mathfrak m^\wedge)_{\mathfrak p'}^\wedge$ は正則である。
それゆえ (C) により、
$R_\mathfrak m \to (R_\mathfrak m^\wedge)_{\mathfrak p'}^\wedge$ のファイバーは $P$ を持つ。
$R_\mathfrak m \to (R_\mathfrak m^\wedge)_{\mathfrak p'}^\wedge$ は局所化
$R_\mathfrak p$ を経由するので、
$R_\mathfrak p \to (R_\mathfrak m^\wedge)_{\mathfrak p'}^\wedge$ のファイバーも $P$ を持つ。
よって (D) を適用して、
$R_\mathfrak p \to R_\mathfrak p^\wedge$ のファイバーが $P$ を持つことが分かる。
\end{proof}

\begin{proposition}
\label{more-algebra-proposition-finite-type-over-P-ring}
$R$ を、(A), (B), (C), (D) を満たす性質 $P$ に対する $P$-環とする。
$R \to S$ が本質的に有限型ならば、$S$ は $P$-環である。
\end{proposition}

\begin{proof}
$P$-環であることは局所環の性質なので、
$P$-環の局所化が $P$-環であることは明らかである。
逆に、すべての素イデアルにおける局所化が $P$-環ならば、もとの環は $P$-環である。
従って、すべての有限型 $R$-代数 $S$ と、その $S$ のすべての素イデアル $\mathfrak q$ に対して
$S_\mathfrak q$ が $P$-環であることを示せば十分である。
$S$ を $R[x_1, \ldots, x_n]$ の商として書くと、
Lemma \ref{more-algebra-lemma-P-ring-goes-up-quasi-finite} により、
$R[x_1, \ldots, x_n]$ が $P$-環であることを示せば十分である。
$n$ に関する帰納法により、$R[x]$ が $P$-環であることを示せば十分である。
$\mathfrak q \subset R[x]$ を極大イデアルとする。
Lemma \ref{more-algebra-lemma-check-P-ring-maximal-ideals} により、
$$
R[x]_\mathfrak q \longrightarrow R[x]_\mathfrak q^\wedge
$$
のファイバーが $P$ を持つことを示せば十分である。
$\mathfrak q$ が $\mathfrak p \subset R$ の上にあるならば、
$R$ を $R_\mathfrak p$ で置き換えてよい。
従って、$R$ は極大イデアル $\mathfrak m$ を持つ Noether 局所 $P$-環であり、
$\mathfrak q \subset R[x]$ は $\mathfrak m$ の上にあると仮定できる。
$\mathfrak q$ の上にある素イデアル
$\mathfrak q' \subset R^\wedge[x]$ は一意であることに注意せよ。
図式
$$
\xymatrix{
R[x]_\mathfrak q^\wedge \ar[r] &
(R^\wedge[x]_{\mathfrak q'})^\wedge \\
R[x]_\mathfrak q \ar[r] \ar[u] & R^\wedge[x]_{\mathfrak q'} \ar[u]
}
$$
を考える。
$R$ は $P$-環なので、$R[x] \to R^\wedge[x]$ のファイバーは $P$ を持つ。
実際、これらは $R \to R^\wedge$ のファイバーの有限生成体拡大による基底変換なので、(A) が適用される。
ゆえに例えば Lemma \ref{more-algebra-lemma-P-local} により、
下の水平矢印のファイバーは $P$ を持つ。
$R^\wedge$ は G-環なので、右の垂直矢印は正則である
(Propositions \ref{more-algebra-proposition-Noetherian-complete-G-ring} and
\ref{more-algebra-proposition-finite-type-over-G-ring})。
したがって (C) により、合成
$R[x]_\mathfrak q \to (R^\wedge[x]_{\mathfrak q'})^\wedge$
のファイバーは $P$ を持つ。
よって (D) により左の垂直矢印のファイバーは $P$ を持ち、証明が完了する。
\end{proof}

\begin{lemma}
\label{more-algebra-lemma-map-P-ring-to-completion-P}
$A$ を $P$-環とし、$P$ は (B) と (D) を満たすとする。
$I \subset A$ をイデアルとし、$A^\wedge$ を $I$ に関する $A$ の完備化とする。
このとき $A \to A^\wedge$ のファイバーは $P$ を持つ。
\end{lemma}

\begin{proof}
環写像 $A \to A^\wedge$ は
Algebra, Lemma \ref{algebra-lemma-completion-flat} により平坦である。
環 $A^\wedge$ は
Algebra, Lemma \ref{algebra-lemma-completion-Noetherian-Noetherian} により Noether である。
従って Lemma \ref{more-algebra-lemma-P-local} の第三条件を確かめれば十分である。
$\mathfrak m \subset A$ の上にある極大イデアルを
$\mathfrak m' \subset A^\wedge$ とする。
Algebra, Lemma \ref{algebra-lemma-radical-completion} により
$IA^\wedge \subset \mathfrak m'$ である。
$A^\wedge/IA^\wedge = A/I$ であるから、
$I \subset \mathfrak m$、$\mathfrak m/I = \mathfrak m'/IA^\wedge$、および
$A/\mathfrak m = A^\wedge/\mathfrak m'$ を得る。
$A^\wedge/\mathfrak m'$ は体なので、$\mathfrak m$ も極大イデアルである。
そこで $A_\mathfrak m \to A^\wedge_{\mathfrak m'}$ は、剰余体を同一視し、かつ
$\mathfrak m A^\wedge_{\mathfrak m'} = \mathfrak m'A^\wedge_{\mathfrak m'}$
を満たす Noether 局所環の平坦局所環準同型である。
したがって完備局所環上の同型を誘導する。
Lemma \ref{more-algebra-lemma-flat-unramified} を参照せよ。
$(A_\mathfrak m)^\wedge$ を $A_\mathfrak m$ の極大イデアルに関する完備化とする。
環写像
$$
(A^\wedge)_{\mathfrak m'} \to
((A^\wedge)_{\mathfrak m'})^\wedge = (A_\mathfrak m)^\wedge
$$
は忠実平坦である
(Algebra, Lemma \ref{algebra-lemma-completion-faithfully-flat})。
従って環写像
$$
A_\mathfrak m \to (A^\wedge)_{\mathfrak m'} \to (A_\mathfrak m)^\wedge
$$
に (D) を適用できる。
$A$ は $P$-環なので $A_\mathfrak m \to (A_\mathfrak m)^\wedge$
のファイバーは $P$ を持ち、これで結論が従う。
\end{proof}

\begin{lemma}
\label{more-algebra-lemma-henselization-pair-P-ring}
\begin{slogan}
環の Hensel 化は、形式的ファイバーの良い性質を受け継ぐ
\end{slogan}
$A$ を $P$-環とし、$P$ は (B), (C), (D), (E) を満たすとする。
$I \subset A$ をイデアルとする。
$(A^h, I^h)$ を対 $(A, I)$ の Hensel 化とする。
Lemma \ref{more-algebra-lemma-henselization} を参照せよ。
このとき $A^h$ は $P$-環である。
\end{lemma}

\begin{proof}
$\mathfrak m^h \subset A^h$ を極大イデアルとする。
Lemma \ref{more-algebra-lemma-check-P-ring-maximal-ideals} により、
$A^h_{\mathfrak m^h} \to (A^h_{\mathfrak m^h})^\wedge$ のファイバーが
$P$ を持つことを示さなければならない。
$\mathfrak m$ を $A$ における $\mathfrak m^h$ の逆像とする。
$(A^h, I^h)$ は Hensel 対なので、
$I^h \subset \mathfrak m^h$ であり、従って $I \subset \mathfrak m$ であることに注意せよ。
$A^h$ は Noether で、$I^h = IA^h$ であり、さらに $A \to A^h$ は
$I$-進完備化上の同型を誘導することを想起せよ。
Lemma \ref{more-algebra-lemma-henselization-Noetherian-pair} を参照せよ。
すると Noether 局所環の局所準同型
$$
A_\mathfrak m \to A^h_{\mathfrak m^h}
$$
は、Lemma \ref{more-algebra-lemma-flat-unramified} により極大イデアルにおける完備化上の同型を誘導する
(詳細は省略する)。
$\mathfrak q \subset A_\mathfrak m$ の上にある $A^h_{\mathfrak m^h}$ の素イデアルを
$\mathfrak q^h$ とする。
$\mathfrak q_1 = \mathfrak q^h$ とおき、
$\mathfrak q_2, \ldots, \mathfrak q_t$ を $\mathfrak q$ の上にある $A^h$ の他の素イデアルとすると、
$A^h \otimes_A \kappa(\mathfrak q) =
\prod\nolimits_{i = 1, \ldots, t} \kappa(\mathfrak q_i)$ である。
Lemma \ref{more-algebra-lemma-filtered-colimit-etale-noetherian-fibres} を参照せよ。
上で論じたように $(A^h)_{\mathfrak m^h}^\wedge = (A_\mathfrak m)^\wedge$ であることを用いると、
$$
\prod\nolimits_{i = 1, \ldots, t}
(A^h_{\mathfrak m^h})^\wedge \otimes_{A^h_{\mathfrak m^h}}
\kappa(\mathfrak q_i) =
(A^h_{\mathfrak m^h})^\wedge \otimes_{A^h_{\mathfrak m^h}}
(A^h_{\mathfrak m^h} \otimes_{A_{\mathfrak m}} \kappa(\mathfrak q)) =
(A_{\mathfrak m})^\wedge \otimes_{A_{\mathfrak m}} \kappa(\mathfrak q)
$$
を得る。したがって局所環を見て (B) を用いると、
$$
\kappa(\mathfrak q) \longrightarrow
(A^h_{\mathfrak m^h})^\wedge \otimes_{A^h_{\mathfrak m^h}}
\kappa(\mathfrak q^h)
$$
は $P$ を持つ。実際、$A$ に関する仮定により
$\kappa(\mathfrak q) \to
(A_\mathfrak m)^\wedge \otimes_{A_\mathfrak m} \kappa(\mathfrak q)$
がこの性質を持つからである。
$\kappa(\mathfrak q^h)/\kappa(\mathfrak q)$ は分離代数的なので、(E) により
$\kappa(\mathfrak q^h) \to
(A^h_{\mathfrak m^h})^\wedge \otimes_{A^h_{\mathfrak m^h}}
\kappa(\mathfrak q^h)$ は望まれた通り $P$ を持つ。
\end{proof}

\begin{lemma}
\label{more-algebra-lemma-henselization-P-ring}
$R$ を $P$-環である Noether 局所環とし、$P$ は
(B), (C), (D), (E) を満たすとする。
このとき Hensel 化 $R^h$ と強 Hensel 化 $R^{sh}$ は $P$-環である。
\end{lemma}

\begin{proof}
Hensel 化については Lemma \ref{more-algebra-lemma-henselization-pair-P-ring} で示した。
強 Hensel 化について示すには、
Lemma \ref{more-algebra-lemma-check-P-ring-maximal-ideals} により、
$R^{sh}$ の形式的ファイバーが $P$ を持つことを示せば十分である。
$\mathfrak r \subset R^{sh}$ を素イデアルとし、
$\mathfrak p = R \cap \mathfrak r$ とおく。
$\mathfrak r_1 = \mathfrak r$ とおき、
$\mathfrak r_2, \ldots, \mathfrak r_s$ を $\mathfrak p$ の上にある $R^{sh}$ の他の素イデアルとすると、
$R^{sh} \otimes_R \kappa(\mathfrak p) =
\prod\nolimits_{i = 1, \ldots, s} \kappa(\mathfrak r_i)$ である。
Lemma \ref{more-algebra-lemma-fibres-henselization} を参照せよ。
すると、
% P02-E0887: frozen display indexes the product to t although the primes are indexed through s.
$$
\prod\nolimits_{i = 1, \ldots, t}
(R^{sh})^\wedge \otimes_{R^{sh}} \kappa(\mathfrak r_i) =
(R^{sh})^\wedge \otimes_{R^{sh}} (R^{sh} \otimes_R \kappa(\mathfrak p)) =
(R^{sh})^\wedge \otimes_R \kappa(\mathfrak p)
$$
を得る。
$R^\wedge \to (R^{sh})^\wedge$ は
$\mathfrak m_{(R^{sh})^\wedge}$-進位相に関して形式的に滑らかであることに注意せよ。
Lemma \ref{more-algebra-lemma-henselization-noetherian} を参照せよ。
従って Proposition \ref{more-algebra-proposition-fs-regular} により
$R^\wedge \to (R^{sh})^\wedge$ は正則である。
それゆえ (C) と $R$ に関する仮定から、性質 $P$ は
$\kappa(\mathfrak p) \to (R^{sh})^\wedge \otimes_R \kappa(\mathfrak p)$
に対して成り立つ。
性質 (B) と上の分解を用い、局所環を見ると、性質 $P$ は
$\kappa(\mathfrak p) \to (R^{sh})^\wedge \otimes_{R^{sh}} \kappa(\mathfrak r)$
に対して成り立つ。
$\kappa(\mathfrak r)/\kappa(\mathfrak p)$ は分離代数的なので、(E) から $P$ は
$\kappa(\mathfrak r) \to (R^{sh})^\wedge \otimes_{R^{sh}} \kappa(\mathfrak r)$
に対して成り立つ。
\end{proof}

\begin{lemma}
\label{more-algebra-lemma-formal-fibres-reduced}
性質 (A), (B), (C), (D), (E) は
$P(k \to R) =$``$R$ は $k$ 上幾何学的に被約である''
に対して成り立つ。
\end{lemma}

\begin{proof}
(A) は幾何学的に被約な代数の定義から従う
(Algebra, Definition \ref{algebra-definition-geometrically-reduced})。
(B) も従う。実際、環が被約であることと、そのすべての局所環が被約であることは同値である。
(C)。これは Lemma \ref{more-algebra-lemma-reduced-goes-up} から従う。
(D)。これは Algebra, Lemma \ref{algebra-lemma-descent-reduced} から従う。
(E)。これは Algebra, Lemma
\ref{algebra-lemma-geometrically-reduced-over-separable-algebraic} から従う。
\end{proof}

\begin{lemma}
\label{more-algebra-lemma-formal-fibres-normal}
性質 (A), (B), (C), (D), (E) は
$P(k \to R) =$``$R$ は $k$ 上幾何学的に正規である''
に対して成り立つ。
\end{lemma}

\begin{proof}
(A) は幾何学的に正規な代数の定義から従う
(Algebra, Definition \ref{algebra-definition-geometrically-normal})。
(B) も従う。実際、環が正規であることと、そのすべての局所環が正規であることは同値である。
(C)。これは Lemma \ref{more-algebra-lemma-normal-goes-up} から従う。
(D)。これは Algebra, Lemma \ref{algebra-lemma-descent-normal} から従う。
(E)。これは Algebra, Lemma
\ref{algebra-lemma-geometrically-normal-over-separable-algebraic} から従う。
\end{proof}

\begin{lemma}
\label{more-algebra-lemma-formal-fibres-Sk}
$n \geq 1$ を固定する。性質 (A), (B), (C), (D), (E) は
$P(k \to R) =$``$R$ は $(S_n)$ を満たす''
に対して成り立つ。
\end{lemma}

\begin{proof}
$k$ を体、$R$ を Noether 環とする環写像を $k \to R$ とする。
$k'/k$ を有限生成体拡大とする。
環写像 $R \to R \otimes_k k'$ のファイバーは
Algebra, Lemma \ref{algebra-lemma-tensor-fields-CM} により Cohen--Macaulay である。
従って環写像 $R \to R \otimes_k k'$ に
Algebra, Lemma \ref{algebra-lemma-Sk-goes-up} を適用でき、
$R$ が $(S_n)$ を満たすならば $R \otimes_k k'$ も満たすことが分かる。
これで (A) が示された。
% P02-E0888: frozen English says "a Noetherian rings has" at source line 13245.
(B) も従う。実際、Noether 環が $(S_n)$ を満たすことと、
そのすべての局所環が $(S_n)$ を満たすことは同値である。
(C)。正則準同型のファイバーは正則で、特に Cohen--Macaulay なので、
これは Algebra, Lemma \ref{algebra-lemma-Sk-goes-up} から従う。
(D)。これは Algebra, Lemma \ref{algebra-lemma-descent-Sk} から従う。
(E)。この条件は基礎体に言及しないので、これは直ちに従う。
\end{proof}

\begin{lemma}
\label{more-algebra-lemma-formal-fibres-CM}
性質 (A), (B), (C), (D), (E) は
$P(k \to R) =$``$R$ は Cohen--Macaulay である''
に対して成り立つ。
\end{lemma}

\begin{proof}
これは Lemma \ref{more-algebra-lemma-formal-fibres-Sk} と、Noether 環が Cohen--Macaulay であることと
すべての $n$ に対して条件 $(S_n)$ を満たすことが同値であることから直ちに従う。
\end{proof}

\begin{lemma}
\label{more-algebra-lemma-formal-fibres-Rk}
$n \geq 0$ を固定する。性質 (A), (B), (C), (D), (E) は
$P(k \to R) =$``すべての有限拡大 $k'/k$ に対して
$R \otimes_k k'$ は $(R_n)$ を満たす''
に対して成り立つ。
\end{lemma}

\begin{proof}
$k$ を体、$R$ を Noether 環とする環写像を $k \to R$ とする。
$P(k \to R)$ が真であると仮定する。
$K/k$ を有限生成体拡大とする。
Algebra, Lemma \ref{algebra-lemma-make-separable} により、図式
$$
\xymatrix{
K \ar[r] & K' \\
k \ar[u] \ar[r] & k' \ar[u]
}
$$
で、$k'/k$, $K'/K$ が有限純非分離体拡大であり、
$K'/k'$ が分離的であるものを取れる。
Algebra, Lemma \ref{algebra-lemma-localization-smooth-separable} により、
滑らかな $k'$-代数 $B$ で、$K'$ がその $B$ の分数体であるものが存在する。
ここで次のように議論できる。
Step 1: $k \to R$ に対して $P$ を仮定したので、
$R \otimes_k k'$ は $(R_n)$ を満たす。
Step 2: $R \otimes_k k' \to R \otimes_k k' \otimes_{k'} B$ は滑らかな環写像である
(Algebra, Lemma \ref{algebra-lemma-base-change-smooth})。
したがって Algebra, Lemma \ref{algebra-lemma-Rk-goes-up} により、
$R \otimes_k k' \otimes_{k'} B$ は $(R_n)$ を満たす
(仮定が満たされることを確かめるため、Algebra, Lemma
\ref{algebra-lemma-characterize-smooth-over-field} も用いる)。
Step 3. $R \otimes_k k' \otimes_{k'} K' = R \otimes_k K'$ は、
$(R_n)$ を満たす環の局所化なので $(R_n)$ を満たす。
Step 4. 最後に、忠実平坦環写像
$K \otimes_k R \to K' \otimes_k R$ に沿う $(R_n)$ の降下により、
$R \otimes_k K$ は $(R_n)$ を満たす
(Algebra, Lemma \ref{algebra-lemma-descent-Rk})。
これで (A) が示された。
(B) も従う。実際、Noether 環が $(R_n)$ を満たすことと、
そのすべての局所環が $(R_n)$ を満たすことは同値である。
(C)。正則準同型のファイバーは正則であるから、これは
Algebra, Lemma \ref{algebra-lemma-Rk-goes-up} から従う
(小さな詳細は省略する)。
(D)。これは Algebra, Lemma \ref{algebra-lemma-descent-Rk} から従う
(小さな詳細は省略する)。

\medskip\noindent
(E)。$l/k$ を体の分離代数拡大とし、
$l \to R$ を $R$ が Noether である環写像とする。
$k \to R$ が $P$ を持つと仮定する。
$l \to R$ が $P$ を持つことを示さなければならない。
$l'/l$ を有限拡大とする。
まず、有限中間拡大 $l/m/k$ と有限拡大 $m'/m$ で、
$l' = l \otimes_m m'$ となるものが存在することに注意する。
すると $R \otimes_l l' = R \otimes_m m'$ である。
ゆえに $m \to R$ が性質 $P$ を持つことを示せば十分である。
すなわち、$l/k$ は有限であると仮定できる。
$l/k$ が有限ならば $l'/k$ は有限であり、
$$
l' \otimes_l R = (l' \otimes_k R) \otimes_{l \otimes_k l} l
$$
は、仮定 $k \to R$ に対する $P$ により性質 $(R_n)$ を持つ Noether 環
$l' \otimes_k R$ の局所化である
(Algebra, Lemma \ref{algebra-lemma-separable-algebraic-diagonal})。
したがって $l' \otimes_l R$ は望まれた通り性質 $(R_n)$ を持つ。
\end{proof}








\section{Excellent 環}
\label{more-algebra-section-excellent}

\noindent
本節では Grothendieck の excellent 環の概念を論じる。
G-環、J-2 環、および普遍的鏖状環の定義については、
Definition \ref{more-algebra-definition-G-ring},
Definition \ref{more-algebra-definition-J}, および
Algebra, Definition \ref{algebra-definition-universally-catenary} を参照せよ。

\begin{definition}
\label{more-algebra-definition-excellent}
$R$ を環とする。
\begin{enumerate}
\item $R$ が Noether、G-環、かつ J-2 であるとき、
$R$ は {\it 準 excellent} であるという。
\item $R$ が準 excellent かつ普遍的鏖状であるとき、
$R$ は {\it excellent} であるという。
\end{enumerate}
\end{definition}

\noindent
従って Noether 環が準 excellent であるとは、
幾何学的に正則な形式的ファイバーを持ち、かつその上の任意の有限型代数が閉じた特異点集合を持つことである。
そのような環が excellent であるためには、さらに（局所的に）良い次元関数が存在することを要請する。
後で、普遍的鏖状であることが、$R$ の極大イデアル $\mathfrak m$ に対する写像
$R_\mathfrak m \to R_\mathfrak m^\wedge$ 上の条件として定式化できることを見る
(Section \ref{more-algebra-section-formally-catenary})。

\begin{lemma}
\label{more-algebra-lemma-finite-type-over-excellent}
(準)excellent 環上の有限型環の任意の局所化は (準)excellent である。
\end{lemma}

\begin{proof}
有限型代数については、J-2 および普遍的鏖状という性質の定義からこれが従う。
G-環については Proposition \ref{more-algebra-proposition-finite-type-over-G-ring} を参照せよ。
局所化が (準)excellence を保つことの証明は省略する。
\end{proof}

\begin{proposition}
\label{more-algebra-proposition-ubiquity-excellent}
次の種類の環は excellent である。
\begin{enumerate}
\item 体、
\item Noether 完備局所環、
\item $\mathbf{Z}$、
\item 分数体の標数が 0 である Dedekind 整域、
\item 上のいずれかの有限型環拡大。
\end{enumerate}
\end{proposition}

\begin{proof}
これらの環が G-環で J-2 を満たすことは、
Propositions \ref{more-algebra-proposition-ubiquity-G-ring} and
\ref{more-algebra-proposition-ubiquity-J-2} を参照せよ。
任意の Cohen--Macaulay 環は普遍的鏖状である。
Algebra, Lemma \ref{algebra-lemma-CM-ring-catenary} を参照せよ。
特に、体、Dedekind 環、さらに一般の正則環は普遍的鏖状である。
Cohen 構造定理を通じて完備局所環が普遍的鏖状であることが分かる。
Algebra, Remark
\ref{algebra-remark-Noetherian-complete-local-ring-universally-catenary} を参照せよ。
\end{proof}

\noindent
上で展開した内容は Nagata 環にいくつかの帰結を持つ。

\begin{lemma}
\label{more-algebra-lemma-Nagata-local-ring}
$(A, \mathfrak m)$ を Noether 局所環とする。次は同値である。
\begin{enumerate}
\item $A$ は Nagata である。
\item $A$ の形式的ファイバーは幾何学的に被約である。
\end{enumerate}
\end{lemma}

\begin{proof}
(2) を仮定する。
Algebra, Lemma \ref{algebra-lemma-local-nagata-and-analytically-unramified} により、
$A \to B$ が有限、$B$ が整域、かつ $\mathfrak m' \subset B$ が極大イデアルならば、
$B_{\mathfrak m'}$ が解析的に不分岐であることを示さなければならない。
Lemmas \ref{more-algebra-lemma-formal-fibres-reduced} and
\ref{more-algebra-lemma-check-P-ring-maximal-ideals} と
Proposition \ref{more-algebra-proposition-finite-type-over-P-ring} を合わせると、
$B_{\mathfrak m'}$ の形式的ファイバーは幾何学的に被約である。
特に、$B$ の分数体を $L$ とすると、
$B_{\mathfrak m'}^\wedge \otimes_B L$ は被約である。
したがって $B_{\mathfrak m'}^\wedge$ は被約である。すなわち、
$B_{\mathfrak m'}$ は解析的に不分岐である。

\medskip\noindent
(1) を仮定する。$\mathfrak q \subset A$ を素イデアルとし、
$K/\kappa(\mathfrak q)$ を有限拡大とする。
$A^\wedge \otimes_A K$ が被約であることを示さなければならない。
$A/\mathfrak q \subset B \subset K$ を、$A$ 上有限で分数体が $K$ である局所部分環とする。
$B$ を構成するため、$\kappa(\mathfrak q)$ 上 $K$ を生成し、かつモニック多項式
$$
P_i(T) = T^{d_i} + a_{i, 1} T^{d_i - 1} + \ldots + a_{i, d_i} = 0
$$
で $a_{i, j} \in \mathfrak m$ を満たす $x_1, \ldots, x_n \in K$ を選ぶ。
そして $B$ を $x_1, \ldots, x_n$ で生成される $K$ の $A$-部分代数とする。
(詳細は Algebra, Lemma
\ref{algebra-lemma-local-nagata-and-analytically-unramified} の証明を参照せよ。)
すると
$$
A^\wedge \otimes_A K =
(A^\wedge \otimes_A B)_\mathfrak q =
B^\wedge_\mathfrak q
$$
である。Algebra, Lemma
\ref{algebra-lemma-local-nagata-and-analytically-unramified} により $B^\wedge$ は被約なので、
証明が完了する。
\end{proof}

\begin{lemma}
\label{more-algebra-lemma-quasi-excellent-nagata}
準 excellent 環は Nagata である。
\end{lemma}

\begin{proof}
$R$ を準 excellent とする。
$R$ 上の有限型代数は準 excellent であること
(Lemma \ref{more-algebra-lemma-finite-type-over-excellent}) を用いると、
任意の準 excellent 整域が N-1 であることを示せば十分である。
Algebra, Lemma \ref{algebra-lemma-check-universally-japanese} を参照せよ。
Algebra, Lemma \ref{algebra-lemma-characterize-N-1} を適用し
(かつ準 excellent 環が J-2 であることを用い)、
準 excellent 局所整域 $R$ が N-1 であることの証明に帰着する。
$R \to R^\wedge$ は正則なので、
Lemma \ref{more-algebra-lemma-reduced-goes-up} により $R^\wedge$ は被約である。
言い換えれば、$R$ は解析的に不分岐である。
従って Algebra, Lemma \ref{algebra-lemma-analytically-unramified-easy} により
$R$ は N-1 である。
\end{proof}

\begin{lemma}
\label{more-algebra-lemma-completion-normal-local-ring}
$(A, \mathfrak m)$ を Noether 局所環とする。
$A$ が正規で、$A$ の形式的ファイバーが正規ならば
(例えば $A$ が excellent または準 excellent ならば)、
$A^\wedge$ は正規である。
\end{lemma}

\begin{proof}
これは Algebra, Lemma \ref{algebra-lemma-normal-goes-up-noetherian} から直ちに従う。
\end{proof}





\section{加群の Abel 圏}
\label{more-algebra-section-abelian-categories-modules}

\noindent
$R$ を環とする。$R$-加群の圏 $\text{Mod}_R$ は Abel 圏である。
$\text{Mod}_R$ の Abel な部分圏の例をいくつか挙げる
(ここでは Homology, Definition \ref{homology-definition-serre-subcategory}
および Homology, Lemmas \ref{homology-lemma-characterize-serre-subcategory} and
\ref{homology-lemma-characterize-weak-serre-subcategory} で導入した用語を用いる)。
\begin{enumerate}
\item 連接 $R$-加群の圏は $\text{Mod}_R$ の弱 Serre 部分圏である。
これは Algebra, Lemma \ref{algebra-lemma-coherent} から従う。
\item $S \subset R$ を乗法的部分集合とする。
$s \in S$ による乗法が $M$ 上の同型であるような
$R$-加群 $M$ からなる充満部分圏は $\text{Mod}_R$ の Serre 部分圏である。
これは Algebra, Lemma \ref{algebra-lemma-localization-and-modules} から従う。
\item $I \subset R$ を有限生成イデアルとする。
$I$-冪捩れ加群の充満部分圏は $\text{Mod}_R$ の Serre 部分圏である。
Lemma \ref{more-algebra-lemma-I-power-torsion} を参照せよ。
\item いくつかの文献では、すべての $x \in M$ に対して
$fx = 0$ を満たす非零因子 $f \in R$ が存在するような加群 $M$ を
{\it 捩れ加群}と定義する。捩れ加群の充満部分圏は
$\text{Mod}_R$ の Serre 部分圏である。
\item $R$ が Noether でないならば、有限生成 $R$-加群の圏
$\text{Mod}^{fg}_R$ は Abel 圏では {\bf ない}。
実際、$I \subset R$ が有限生成でないイデアルならば、
$R \to R/I$ は $\text{Mod}^{fg}_R$ に核を持たない。
\item $R$ が Noether ならば、連接 $R$-加群と有限生成（すなわち有限）
$R$-加群は一致する。Algebra, Lemmas
\ref{algebra-lemma-Noetherian-coherent},
\ref{algebra-lemma-coherent-ring}, and
\ref{algebra-lemma-Noetherian-finite-type-is-finite-presentation} を参照せよ。
従って上の (1) により $\text{Mod}^{fg}_R$ は Abel 圏であるが、
% P02-E0889: frozen English has "in fact,in" without a separating space at source line 13536.
実際この場合、圏 $\text{Mod}_R^{fg}$ は $\text{Mod}_R$ の（強）Serre 部分圏である。
\end{enumerate}




\section{入射 Abel 群}
\label{more-algebra-section-abelian-groups}

\noindent
本節では、Abel 群の圏が十分な入射対象を持つことを示す。
Abel 群 $M$ が {\it 可除}であることと、すべての $x \in M$ と
$n \in \mathbf{N}$ に対して $n y = x$ を満たす $y \in M$ が存在することは同値であることを想起せよ。

\begin{lemma}
\label{more-algebra-lemma-injective-abelian}
Abel 群 $J$ が Abel 群の圏の入射対象であることと、
$J$ が可除であることは同値である。
\end{lemma}

\begin{proof}
$J$ は可除でないと仮定する。すると、
$n y = x$ を満たす $y \in J$ が存在しないような
$x \in J$ と $n \in \mathbf{N}$ が存在する。
このとき準同型 $\mathbf{Z} \to J$、$m \mapsto mx$ は
$\frac{1}{n}\mathbf{Z} \supset \mathbf{Z}$ へ拡張できない。
従って $J$ は入射的でない。

\medskip\noindent
$A \subset B$ を Abel 群とする。$J$ は可除 Abel 群であると仮定する。
$\varphi : A \to J$ を準同型とする。
$A \subset A' \subset B$ かつ $\varphi'|_A = \varphi$ を満たす準同型
$\varphi' : A' \to J$ の集合を考える。
$A' \supset A''$ かつ $\varphi'|_{A''} = \varphi''$ であるとき、およびそのときに限り
$(A', \varphi') \geq (A'', \varphi'')$ と定義する。
$(A_i, \varphi_i)_{i \in I}$ がそのような対の全順序族ならば、
$a \in A_i$ を $\varphi_i(a)$ に写す写像
$\bigcup_{i \in I} A_i \to J$ を得る。
したがって Zorn の補題が適用できる。
結論するには、対 $(A', \varphi')$ が極大ならば $A' = B$ であることを示さなければならない。
言い換えれば、部分群 $A \subset B$、$A \not = B$ と任意の
$\varphi : A \to J$ が与えられたとき、
$A \subset A' \subset B$ である $\varphi' : A' \to J$ で、
(a) 包含 $A \subset A'$ が真であり、かつ
(b) 準同型 $\varphi'$ が $\varphi$ を拡張するものを見つければ十分である。

\medskip\noindent
これを示すため、$x \in B$、$x \not \in A$ を選ぶ。
$nx \in A$ を満たす $n\in \mathbf{N}$ が存在しないならば、
$A \oplus \mathbf{Z} \cong A + \mathbf{Z}x$ である。
したがって、$A$ 上で $\varphi$ を用い、例えば $x$ を零に写すことにより、
$\varphi$ を $A' = A + \mathbf{Z}x$ へ拡張できる。
$nx \in A$ を満たす $n \in \mathbf{N}$ が存在するならば、
$n$ をそのような最小の整数とする。
$nz = \varphi(nx)$ を満たす元 $z \in J$ を取る。
準同型 $\tilde\varphi : A \oplus \mathbf{Z} \to J$ を
$(a, m) \mapsto \varphi(a) + mz$ で定義する。
$z$ の選び方により、$\tilde \varphi$ の核は写像
$A \oplus \mathbf{Z} \to B$、$(a, m) \mapsto a + mx$ の核を含む。
したがって $\tilde \varphi$ は像 $A' = A + \mathbf{Z}x$ を経由し、
これは準同型 $\varphi$ を拡張する。
\end{proof}

\noindent
この補題を用いて、任意の Abel 群が入射 Abel 群に埋め込めることを示せる。
しかしこれは次節の結果の特別な場合である。






\section{入射加群}
\label{more-algebra-section-injectives-modules}

\noindent
入射加群に関するいくつかの補題を述べる。

\begin{definition}
% P02-E0890: frozen label says "projective" although this definition is injectivity.
\label{more-algebra-definition-projective}
$R$ を環とする。$R$-加群 $J$ が {\it 入射的}であることと、関手
$\Hom_R(-, J) : \text{Mod}_R \to \text{Mod}_R$ が完全関手であることは同値である。
\end{definition}

\noindent
任意の $R$-加群 $M$ に対して関手 $\Hom_R(- , M)$ は左完全である。
Algebra, Lemma \ref{algebra-lemma-hom-exact} を参照せよ。
従って $J$ が入射的であるための条件は、
$R$-加群の単射 $M \to M'$ が与えられたとき、写像
$\Hom_R(M', J) \to \Hom_R(M, J)$ が全射であることを意味する。

\medskip\noindent
これを ${Ext}$-加群によって言い換える前に、
$\Ext^1_R(M, N)$ と Homology, Section \ref{homology-section-extensions} の意味での拡大との関係を論じる。

\begin{lemma}
\label{more-algebra-lemma-relation-ext-ext}
$R$ を環とし、$\mathcal{A}$ を $R$-加群の Abel 圏とする。
Algebra, Lemmas \ref{algebra-lemma-long-exact-seq-ext} and
\ref{algebra-lemma-reverse-long-exact-seq-ext} の長完全列、および
Homology, Lemma \ref{homology-lemma-six-term-sequence-ext} の $6$-項完全列と互換な標準同型
$\Ext_\mathcal{A}(M, N) = \Ext^1_R(M, N)$ が存在する。
\end{lemma}

\begin{proof}
省略する。
\end{proof}

\begin{lemma}
\label{more-algebra-lemma-characterize-injective}
$R$ を環とし、$J$ を $R$-加群とする。次は同値である。
\begin{enumerate}
\item $J$ は入射的である。
\item すべての $R$-加群 $M$ に対して $\Ext^1_R(M, J) = 0$ である。
\end{enumerate}
\end{lemma}

\begin{proof}
$0 \to M'' \to M' \to M \to 0$ を $R$-加群の短完全列とする。
Algebra, Lemma \ref{algebra-lemma-reverse-long-exact-seq-ext} の長完全列
$$
\begin{matrix}
0
\to \Hom_R(M, J)
\to \Hom_R(M', J)
\to \Hom_R(M'', J)
\\
\phantom{0\ }
\to \Ext^1_R(M, J)
\to \Ext^1_R(M', J)
\to \Ext^1_R(M'', J)
\to \ldots
\end{matrix}
$$
を考える。すると (2) が (1) を導くことが分かる。
逆に $J$ が入射的ならば、Homology, Lemma
\ref{homology-lemma-characterize-injectives} と Lemma \ref{more-algebra-lemma-relation-ext-ext} により
$\Ext$-群は零である。
\end{proof}

\begin{lemma}
\label{more-algebra-lemma-characterize-injective-bis}
$R$ を環とし、$J$ を $R$-加群とする。次は同値である。
\begin{enumerate}
\item $J$ は入射的である。
\item すべてのイデアル $I \subset R$ に対して $\Ext^1_R(R/I, J) = 0$ である。
\item イデアル $I \subset R$ と加群写像 $I \to J$ に対し、拡張 $R \to J$ が存在する。
\end{enumerate}
\end{lemma}

\begin{proof}
$I \subset R$ がイデアルならば、短完全列
$0 \to I \to R \to R/I \to 0$ は完全列
$$
\Hom_R(R, J) \to
\Hom_R(I, J) \to
\Ext^1_R(R/I, J) \to 0
$$
を与える。これは Algebra, Lemma
\ref{algebra-lemma-reverse-long-exact-seq-ext} と、$R$ が射影的なので
$\Ext^1_R(R, J) = 0$ であることによる
(Algebra, Lemma \ref{algebra-lemma-characterize-projective})。
したがって (2) と (3) は同値である。
この証明では Baer の判定法として知られる
(1) $\Leftrightarrow$ (3) を示す。

\medskip\noindent
(1) を仮定する。(3) のような加群写像 $I \to J$ が与えられると、
定義により写像 $\Hom_R(R, J) \to \Hom_R(I, J)$ が全射なので、
拡張 $R \to J$ を得る。

\medskip\noindent
(3) を仮定する。$M \subset N$ を $R$-加群の包含とする。
$\varphi : M \to J$ を準同型とする。
$\varphi$ が $N$ へ拡張することを示せば、補題の証明が完了する。
$M \subset M' \subset N$ かつ $\varphi'|_M = \varphi$ を満たす準同型
$\varphi' : M' \to J$ の集合を考える。
$M' \supset M''$ かつ $\varphi'|_{M''} = \varphi''$ であるとき、およびそのときに限り
$(M', \varphi') \geq (M'', \varphi'')$ と定義する。
$(M_i, \varphi_i)_{i \in I}$ がそのような対の全順序族ならば、
$a \in M_i$ を $\varphi_i(a)$ に写す写像 $\bigcup_{i \in I} M_i \to J$ を得る。
従って Zorn の補題が適用できる。
結論するには、対 $(M', \varphi')$ が極大ならば $M' = N$ であることを示さなければならない。
言い換えれば、
% P02-E0891: frozen English says "subgroup" although M and N are R-modules.
任意の部分加群 $M \subset N$、$M \not = N$ と任意の
$\varphi : M \to J$ が与えられたとき、
$M \subset M' \subset N$ である $\varphi' : M' \to J$ で、
(a) 包含 $M \subset M'$ が真であり、かつ
(b) 準同型 $\varphi'$ が $\varphi$ を拡張するものを見つければ十分である。

\medskip\noindent
これを示すため、$x \in N$、$x \not \in M$ を選ぶ。
$I = \{f \in R \mid fx \in M\}$ とおく。これは $R$ のイデアルである。
準同型 $\psi : I \to J$ を $f \mapsto \varphi(fx)$ で定義する。
仮定 (3) により可能な写像 $\tilde\psi : R \to J$ へ拡張する。
$I$ の選び方により、写像
$M \oplus R \to J$、$(y, f) \mapsto \varphi(y) + \tilde\psi(f)$ の核は、写像
$M \oplus R \to N$、$(y, f) \mapsto y + fx$ の核を含む。
従ってこの準同型は像 $M' = M + Rx$ を経由し、
望まれた通り与えられた準同型を拡張する。
\end{proof}

\noindent
本節の残りでは、環 $R$ 上に十分な入射加群が存在することを証明する。
まず $\mathbf{Q}/\mathbf{Z}$ が入射 Abel 群であることから始める。
これは Lemma \ref{more-algebra-lemma-injective-abelian} から従う。

\begin{definition}
\label{more-algebra-definition-simple-functors}
$R$ を環とする。
\begin{enumerate}
\item $R$ 上の任意の $R$-加群 $M$ に対し、自然な $R$-加群構造を入れた
$M^\vee = \Hom(M, \mathbf{Q}/\mathbf{Z})$ を表す。
{\it $M \mapsto M^\vee$} を $R$-加群の圏からそれ自身への反変関手とみなす。
\item 任意の $R$-加群 $M$ に対し、
$$
F(M) = \bigoplus\nolimits_{m \in M} R[m]
$$
を、元 $[m]$、$m \in M$ で与えられる基底を持つ {\it 自由加群}として表す。
$R$-加群の自然な全射を
$F(M)\to M$、$\sum f_i [m_i] \mapsto \sum f_i m_i$ とする。
$M \mapsto (F(M) \to M)$ を $R$-加群の圏から
$R$-加群の矢印の圏への関手とみなす。
\end{enumerate}
\end{definition}

\begin{lemma}
\label{more-algebra-lemma-vee-exact}
$R$ を環とする。関手 $M \mapsto M^\vee$ は完全である。
\end{lemma}

\begin{proof}
% P02-E0892: frozen English reads "This because" at source line 13793.
これは Lemma \ref{more-algebra-lemma-injective-abelian} により
$\mathbf{Q}/\mathbf{Z}$ が入射 Abel 群であるからである。
\end{proof}

\noindent
評価によって与えられる標準写像 $ev : M \to (M^\vee)^\vee$ がある。
すなわち、$x \in M$ が与えられたとき、
$ev(x) \in (M^\vee)^\vee = \Hom(M^\vee, \mathbf{Q}/\mathbf{Z})$ を写像
$\varphi \mapsto \varphi(x)$ とする。

\begin{lemma}
\label{more-algebra-lemma-ev-injective}
任意の $R$-加群 $M$ に対し、評価写像
$ev : M \to (M^\vee)^\vee$ は単射である。
\end{lemma}

\begin{proof}
$\mathbf{Q}/\mathbf{Z}$ が入射 Abel 群であることを用いれば確かめられる。
実際、$x \in M$ が零でないならば、$M' \subset M$ をそれが生成する循環群とする。
$x$ を零化しない非零写像 $M' \to \mathbf{Q}/\mathbf{Z}$ が存在する。
これは写像 $\varphi : M \to \mathbf{Q}/\mathbf{Z}$ に拡張し、そのとき
$ev(x)(\varphi) = \varphi(x) \not = 0$ である。
\end{proof}

\noindent
$R$-加群の標準全射 $F(M) \to M$ は、上の構成によって
$R$-加群の標準単射
$$
(M^\vee)^\vee \longrightarrow (F(M^\vee))^\vee.
$$
に移る。$J(M) = (F(M^\vee))^\vee$ とおく。
% P02-E0893: frozen English says "with this the displayed map" at source line 13825.
$ev$ と表示した写像との合成は、$M$ に関して関手的な $M \to J(M)$ を与える。

\begin{lemma}
\label{more-algebra-lemma-JM-injective}
$R$ を環とする。すべての $R$-加群 $M$ に対して、
$R$-加群 $J(M)$ は入射的である。
\end{lemma}

\begin{proof}
$R$-加群として $J(M) \cong \prod_{\varphi \in M^\vee} R^\vee$ であることに注意せよ。
入射加群の積は入射的なので、$R^\vee$ が入射的であることを示せば十分である。
このために、
$$
\Hom_R(N, R^\vee) =
\Hom_R(N, \Hom_{\mathbf{Z}}(R, \mathbf{Q}/\mathbf{Z})) =
N^\vee
$$
と、Lemma \ref{more-algebra-lemma-vee-exact} により $(-)^\vee$ が完全関手であることを用いる。
\end{proof}

\begin{lemma}
\label{more-algebra-lemma-injectives-modules}
$R$ を環とする。上の構成は、$R$-加群の圏から
$R$-加群の矢印の圏への共変関手
$M \mapsto (M \to J(M))$ を定め、すべての加群 $M$ に対する出力
$M \to J(M)$ は、$M$ から入射 $R$-加群 $J(M)$ への単射である。
\end{lemma}

\begin{proof}
上記から従う。
\end{proof}

\noindent
特に、$R$-加群の任意の写像 $M \to N$ に対し、次の図式を可換にする対応する準同型
$J(M) \to J(N)$ がある。
$$
\xymatrix{
M \ar[d] \ar[r] & N \ar[d] \\
J(M) \ar[r] & J(N) }
$$
これが一般に欲しい構成の種類である。
Homology, Section \ref{homology-section-injectives} でこれを表す用語を導入した。
すなわち、これは $R$-加群の圏が関手的入射埋め込みを持つことを意味するという。











\section{加群の導来圏}
\label{more-algebra-section-derived-modules}

\noindent
本節では、環上の加群の導来圏に関する一般的事項をまとめる。

\medskip\noindent
$A$ を環とする。$A$-加群の圏を $\text{Mod}_A$ と書く。
$A$-加群の複体のホモトピー圏を表すのに記号 $K(A)$ を用いる。
すなわち、圏として $K(A) = K(\text{Mod}_A)$ とおく。
Derived Categories, Section \ref{derived-section-homotopy} を参照せよ。
有界版は $K^+(A)$, $K^-(A)$, $K^b(A)$ である。
Derived Categories, Section \ref{derived-section-homotopy-triangulated} のように
$K(A)$ を三角圏とみなす。
$A$ の {\it 導来圏}、$D(A)$ とは、$K(A)$ の擬同型を可逆にして得られる圏である。
すなわち $D(A) = D(\text{Mod}_A)$ とおく。Derived Categories, Section
\ref{derived-section-derived-categories}\footnote{Injectives, Remark \ref{injectives-remark-existence-D} も参照せよ。}
を参照せよ。有界版は $D^+(A)$, $D^-(A)$, $D^b(A)$ である。

\medskip\noindent
$A$ を環とする。$A$-加群の圏は積を持ち、積は完全である。
Lemma \ref{more-algebra-lemma-injectives-modules} により、$A$-加群の圏は十分な入射対象を持つ。
従ってすべての $A$-加群の複体は K-入射複体と擬同型である
(Derived Categories, Lemma \ref{derived-lemma-enough-K-injectives-Ab4-star})。
それゆえ $D(A)$ は可算積を持ち
(Derived Categories, Lemma \ref{derived-lemma-products})、
実際には任意の積を持つ
(Injectives, Lemma \ref{injectives-lemma-derived-products})。
これにより、$D(A)$ の対象の任意の逆系は（同型を除いて定まる）導来極限を持つ。
Derived Categories, Section \ref{derived-section-derived-limit} を参照せよ。

\begin{lemma}
\label{more-algebra-lemma-K-injective-flat}
$R \to S$ を平坦な環写像とする。
$I^\bullet$ が $S$-加群の K-入射複体ならば、
$I^\bullet$ は $R$-加群の複体として K-入射的である。
\end{lemma}

\begin{proof}
これは、Algebra, Lemma \ref{algebra-lemma-adjoint-tensor-restrict} と
$S$ とのテンソル積を取ることが完全であることにより、
$$
\Hom_{K(R)}(M^\bullet, I^\bullet) =
\Hom_{K(S)}(M^\bullet \otimes_R S, I^\bullet)
$$
であるから真である。
\end{proof}

\begin{lemma}
\label{more-algebra-lemma-K-injective-epimorphism}
$R \to S$ を環のエピ射とし、$I^\bullet$ を $S$-加群の複体とする。
$I^\bullet$ が $R$-加群の複体として K-入射的ならば、
$I^\bullet$ は $S$-加群の K-入射複体である。
\end{lemma}

\begin{proof}
これは、任意の $S$-加群の複体 $N^\bullet$ に対して
$$
\Hom_{K(R)}(N^\bullet, I^\bullet) =
\Hom_{K(S)}(N^\bullet, I^\bullet)
$$
であるから真である。Algebra, Lemma
\ref{algebra-lemma-epimorphism-modules} を参照せよ。
\end{proof}

\begin{lemma}
\label{more-algebra-lemma-hom-K-injective}
$A \to B$ を環写像とする。$I^\bullet$ が $A$-加群の K-入射複体ならば、
$\Hom_A(B, I^\bullet)$ は $B$-加群の K-入射複体である。
\end{lemma}

\begin{proof}
これは Algebra, Lemma \ref{algebra-lemma-adjoint-hom-restrict} により
$$
\Hom_{K(B)}(N^\bullet, \Hom_A(B, I^\bullet)) =
\Hom_{K(A)}(N^\bullet, I^\bullet)
$$
であるから真である。
\end{proof}










\section{Tor の計算}
\label{more-algebra-section-computing-tor}

\noindent
$R$ を環とする。$R$-加群の Abel 圏 $\text{Mod}_R$ の導来圏を $D(R)$ と書く。
すべての自由 $R$-加群は射影的なので、$\text{Mod}_R$ は十分な射影対象を持つことに注意せよ。
従って $\text{Mod}_R$ から任意の Abel 圏への任意の加法的関手の左導来関手を定義できる。

\medskip\noindent
特にこれは関手
$ - \otimes_R M : \text{Mod}_R \to \text{Mod}_R$
に適用される。その左導来関手は Tor 関手 $\text{Tor}_i^R(-, M)$ である。
Algebra, Section \ref{algebra-section-tor} を参照せよ。
全左導来関手
\begin{equation}
\label{more-algebra-equation-derived-tensor-module}
-\otimes_R^{\mathbf{L}} M :
D^{-}(R)
\longrightarrow
D^{-}(R)
\end{equation}
もあり、$-\otimes_R^{\mathbf{L}} M$ と書く。
その satellite は Tor 加群である。すなわち、
$$
H^{-p}(N \otimes_R^{\mathbf{L}} M) = \text{Tor}_p^R(N, M).
$$
である。

\medskip\noindent
$R$-代数 $A$ とのテンソル積を考えるときには特別な状況が生じる。
この場合、$- \otimes_R A$ を $\text{Mod}_R$ から $\text{Mod}_A$ への関手とみなす。
% P02-E0894: frozen English calls the tensor construction a total right derived functor at source line 14011.
そこで全左導来関手
\begin{equation}
\label{more-algebra-equation-derived-tensor-algebra}
-\otimes_R^{\mathbf{L}} A :
D^{-}(R)
\longrightarrow
D^{-}(A)
\end{equation}
があり、$-\otimes_R^{\mathbf{L}} A$ と書く。
その satellite は Tor 群である。すなわち、
$$
H^{-p}(N \otimes_R^{\mathbf{L}} A) = \text{Tor}_p^R(N, A).
$$
である。特に、これらの Tor 群は自然な $A$-加群構造を持つ。

\medskip\noindent
次のいくつかの節で、本節の内容を非有界複体に一般化する。






\section{複体のテンソル積}
\label{more-algebra-section-symmetric-monoidal}

\noindent
$R$ を環とする。$R$-加群の複体の圏 $\text{Comp}(R)$ は対称モノイド構造を持つ。
実際、$R$-加群の二つの複体 $L^\bullet$ と $M^\bullet$ があるとする。
Homology, Example \ref{homology-example-double-complex-as-tensor-product-of} と
Homology, Definition \ref{homology-definition-associated-simple-complex} を用いて、第三の $R$-加群の複体
$$
\text{Tot}(L^\bullet \otimes_R M^\bullet)
$$
を得る。この構成が $L^\bullet$ と $M^\bullet$ の両方に関して関手的であることは明らかである。
結合制約は複体の標準同型
$$
\text{Tot}(\text{Tot}(K^\bullet \otimes_R L^\bullet) \otimes_R M^\bullet)
\longrightarrow
\text{Tot}(K^\bullet \otimes_R \text{Tot}(L^\bullet \otimes_R M^\bullet))
$$
とする。これは Homology, Remark \ref{homology-remark-triple-complex} において
三重複体 $K^\bullet \otimes_R L^\bullet \otimes_R M^\bullet$ から構成されたものである。
可換制約は標準同型
$$
\text{Tot}(L^\bullet \otimes_R M^\bullet) \to
\text{Tot}(M^\bullet \otimes_R L^\bullet)
$$
であり、直和成分 $L^p \otimes_R M^q$ 上で符号 $(-1)^{pq}$ を用いる。
これが複体の写像であることを見るため、
$x \in L^p$ と $y \in M^q$ に対して次を計算する。
$$
\text{d}(x \otimes y) =
\text{d}_L(x) \otimes y + (-1)^px \otimes \text{d}_M(y)
$$
我々の規則により、右辺は
$$
(-1)^{(p + 1)q}y \otimes \text{d}_L(x) +
(-1)^{p + p(q + 1)} \text{d}_M(y) \otimes x
$$
に写る。一方、
$$
\text{d}((-1)^{pq}y \otimes x) =
(-1)^{pq} \text{d}_M(y) \otimes x +
(-1)^{pq + q} y \otimes \text{d}_L(x)
$$
である。これら二つの式は望まれた通り見れば一致する。

\begin{lemma}
\label{more-algebra-lemma-symmetric-monoidal-cat-complexes}
$R$ を環とする。$R$-加群の複体の圏 $\text{Comp}(R)$ に、関手
$$
(L^\bullet, M^\bullet) \mapsto \text{Tot}(L^\bullet \otimes_R M^\bullet)
$$
と上記の結合制約および可換制約を入れると、これは対称モノイド圏である。
\end{lemma}

\begin{proof}
省略する。ヒントとして、単位 $\mathbf{1}$ には次数 $0$ で $R$ を持ち、
他の次数で零である複体を取る。すると明らかな同型
$\text{Tot}(\mathbf{1} \otimes_R M^\bullet) = M^\bullet$ および
$\text{Tot}(K^\bullet \otimes_R \mathbf{1}) = K^\bullet$ がある。
% P02-E0895: frozen English starts the next sentence with lowercase "to" at source line 14096.
補題を証明するには、いくつかの図式の可換性を確かめる必要がある。
Categories, Definitions \ref{categories-definition-monoidal-category} and
\ref{categories-definition-symmetric-monoidal-category} を参照せよ。
いずれの場合も検証は直接的である。
\end{proof}

\begin{lemma}
\label{more-algebra-lemma-derived-tor-homotopy}
$R$ を環とし、$P^\bullet$ を $R$-加群の複体とする。
$\alpha, \beta : L^\bullet \to M^\bullet$ をホモトピーな複体写像とする。
このとき $\alpha$ と $\beta$ が誘導する写像
$$
\text{Tot}(\alpha \otimes \text{id}_P),
\text{Tot}(\beta \otimes \text{id}_P) :
\text{Tot}(L^\bullet \otimes_R P^\bullet)
\longrightarrow
\text{Tot}(M^\bullet \otimes_R P^\bullet).
$$
はホモトピーである。特に、構成
$L^\bullet \mapsto \text{Tot}(L^\bullet \otimes_R P^\bullet)$
は複体のホモトピー圏の自己関手を定める。
\end{lemma}

\begin{proof}
$h^n : L^n \to M^{n - 1}$ で定められるあるホモトピー $h$ に対して
$\alpha = \beta + dh + hd$ とする。
$$
H^n = \bigoplus\nolimits_{a + b = n} h^a \otimes \text{id}_{P^b} :
\bigoplus\nolimits_{a + b = n} L^a \otimes_R P^b
\longrightarrow
\bigoplus\nolimits_{a + b = n} M^{a - 1} \otimes_R P^b
$$
とおく。すると直接の計算により、写像
$\text{Tot}(L^\bullet \otimes_R P^\bullet) \to
\text{Tot}(M^\bullet \otimes_R P^\bullet)$ として
$$
\text{Tot}(\alpha \otimes \text{id}_P) =
\text{Tot}(\beta \otimes \text{id}_P) + dH + Hd
$$
である。
\end{proof}

\begin{lemma}
\label{more-algebra-lemma-symmetric-monoidal-homotopy-category}
$R$ を環とする。$R$-加群の複体のホモトピー圏 $K(R)$ に、関手
$$
(L^\bullet, M^\bullet) \mapsto \text{Tot}(L^\bullet \otimes_R M^\bullet)
$$
と上記の結合制約および可換制約を入れると、これは対称モノイド圏である。
\end{lemma}

\begin{proof}
これは Lemmas \ref{more-algebra-lemma-symmetric-monoidal-cat-complexes} and
\ref{more-algebra-lemma-derived-tor-homotopy} から従う。詳細は省略する。
\end{proof}

\begin{lemma}
\label{more-algebra-lemma-derived-tor-exact}
$R$ を環とし、$P^\bullet$ を $R$-加群の複体とする。関手
$$
K(R) \longrightarrow K(R), \quad
L^\bullet \longmapsto \text{Tot}(P^\bullet \otimes_R L^\bullet)
$$
および
$$
K(R) \longrightarrow K(R), \quad
L^\bullet \longmapsto \text{Tot}(L^\bullet \otimes_R P^\bullet)
$$
は三角圏の完全関手である。
\end{lemma}

\begin{proof}
これは Derived Categories, Remark
\ref{derived-remark-double-complex-as-tensor-product-of} から従う。
\end{proof}




\section{導来テンソル積}
\label{more-algebra-section-derived-tensor-product}

\noindent
導来テンソル積はより一般に構成できる。
実際、K-平坦分解を選べば、有界性の仮定は必要ないことが分かる。

\begin{definition}
\label{more-algebra-definition-K-flat}
$R$ を環とする。複体 $K^\bullet$ は、任意の非輪状複体 $M^\bullet$ に対して全複体
$\text{Tot}(M^\bullet \otimes_R K^\bullet)$ が非輪状であるとき {\it K-平坦}と呼ばれる。
\end{definition}

\begin{lemma}
\label{more-algebra-lemma-K-flat-quasi-isomorphism}
$R$ を環とし、$K^\bullet$ を K-平坦複体とする。このとき関手
$$
K(R) \longrightarrow K(R), \quad
L^\bullet \longmapsto \text{Tot}(L^\bullet \otimes_R K^\bullet)
$$
は擬同型を擬同型に写す。
\end{lemma}

\begin{proof}
これは Lemma \ref{more-algebra-lemma-derived-tor-exact} と、$K(R)$ における擬同型が非輪状なコーンを持つことによって特徴づけられることから従う。
\end{proof}

\begin{lemma}
\label{more-algebra-lemma-base-change-K-flat}
$R \to R'$ を環写像とする。$K^\bullet$ が $R$-加群の K-平坦複体ならば、
$K^\bullet \otimes_R R'$ は $R'$-加群の K-平坦複体である。
\end{lemma}

\begin{proof}
これは定義と、任意の $R'$-加群の複体 $L^\bullet$ に対する等式
$$(K^\bullet \otimes_R R') \otimes_{R'} L^\bullet =
K^\bullet \otimes_R L^\bullet$$
から従う。
\end{proof}

\begin{lemma}
\label{more-algebra-lemma-tensor-product-K-flat}
$R$ を環とする。$K^\bullet$, $L^\bullet$ が $R$-加群の K-平坦複体ならば、
$\text{Tot}(K^\bullet \otimes_R L^\bullet)$ は $R$-加群の K-平坦複体である。
\end{lemma}

\begin{proof}
これは同型
$$
\text{Tot}(M^\bullet \otimes_R \text{Tot}(K^\bullet \otimes_R L^\bullet))
=
\text{Tot}(\text{Tot}(M^\bullet \otimes_R K^\bullet) \otimes_R L^\bullet)
$$
と定義から従う。
\end{proof}

\begin{lemma}
\label{more-algebra-lemma-K-flat-two-out-of-three}
$R$ を環とする。$(K_1^\bullet, K_2^\bullet, K_3^\bullet)$ を $K(R)$ の完全三角とする。
$K_i^\bullet$ のうち三つ中二つが K-平坦ならば、残る一つも K-平坦である。
\end{lemma}

\begin{proof}
これは Lemma \ref{more-algebra-lemma-derived-tor-exact} と、$K(R)$ の完全三角では三つ中二つが非輪状ならば残る一つも非輪状であることから従う。
\end{proof}

\begin{lemma}
\label{more-algebra-lemma-K-flat-two-out-of-three-ses}
$R$ を環とする。
$0 \to K_1^\bullet \to K_2^\bullet \to K_3^\bullet \to 0$ を複体の短完全列とする。
すべての $n \in \mathbf{Z}$ に対して $K_3^n$ が平坦であり、
$K_i^\bullet$ のうち三つ中二つが K-平坦ならば、残る一つも K-平坦である。
\end{lemma}

\begin{proof}
$L^\bullet$ を $R$-加群の複体とする。このとき
$$
0 \to
\text{Tot}(L^\bullet \otimes_R K_1^\bullet) \to
\text{Tot}(L^\bullet \otimes_R K_2^\bullet) \to
\text{Tot}(L^\bullet \otimes_R K_3^\bullet) \to 0
$$
は複体の短完全列である。
実際、各 $n, m$ に対して加群の列
$0 \to L^n \otimes_R K_1^m \to
L^n \otimes_R K_2^m \to
L^n \otimes_R K_3^m \to 0$
は Algebra, Lemma \ref{algebra-lemma-flat-tor-zero} により完全であり、
複体の列はこれらの直和である。
従って補題は、これと、複体の短完全列では三つ中二つが非輪状ならば残る一つも非輪状であることから従う。
\end{proof}

\begin{lemma}
\label{more-algebra-lemma-derived-tor-quasi-isomorphism}
$R$ を環とし、$P^\bullet$ を平坦 $R$-加群の上に有界な複体とする。
このとき $P^\bullet$ は K-平坦である。
\end{lemma}

\begin{proof}
$L^\bullet$ を $R$-加群の非輪状複体とする。
$\xi \in H^n(\text{Tot}(L^\bullet \otimes_R P^\bullet))$ とする。
$\xi = 0$ を示さなければならない。
$\text{Tot}^n(L^\bullet \otimes_R P^\bullet)$ は項 $L^a \otimes_R P^b$ を持つ直和なので、
ある $m \in \mathbf{Z}$ に対して $\xi$ は元
$$
H^n(\text{Tot}(\tau_{\leq m}L^\bullet \otimes_R P^\bullet))
$$
から来る。$\tau_{\leq m}L^\bullet$ も非輪状なので、
$L^\bullet$ を $\tau_{\leq m}L^\bullet$ で置き換えてよい。
したがって $L^\bullet$ は上に有界であると仮定できる。
この場合、Homology, Lemma \ref{homology-lemma-first-quadrant-ss} のスペクトル列は
$$
{}'E_1^{p, q} = H^p(L^\bullet \otimes_R P^q)
$$
を持つ。$P^q$ が平坦で $L^\bullet$ が非輪状なので、これは零である。
したがって $H^*(\text{Tot}(L^\bullet \otimes_R P^\bullet)) = 0$ である。
\end{proof}

\noindent
次の補題で複体の系の余極限というとき、項ごとの余極限を意味する。

\begin{lemma}
\label{more-algebra-lemma-colimit-K-flat}
$R$ を環とする。
$K_1^\bullet \to K_2^\bullet \to \ldots$ を K-平坦複体の系とする。
このとき $\colim_i K_i^\bullet$ は K-平坦である。
より一般に、K-平坦複体の任意のフィルター付き余極限は K-平坦である。
\end{lemma}

\begin{proof}
項ごとの余極限を取っているので、Algebra, Lemma
\ref{algebra-lemma-tensor-products-commute-with-limits} により
$$
\colim_i \text{Tot}(M^\bullet \otimes_R K_i^\bullet) =
\text{Tot}(M^\bullet \otimes_R \colim_i K_i^\bullet)
$$
である。従って補題は、フィルター付き余極限が完全であることから従う。
Algebra, Lemma \ref{algebra-lemma-directed-colimit-exact} を参照せよ。
\end{proof}

\begin{lemma}
\label{more-algebra-lemma-universally-acyclic-K-flat}
$R$ を環とし、$K^\bullet$ を $R$-加群の複体とする。
すべての有限表示 $R$-加群 $M$ に対して
$K^\bullet \otimes_R M$ が非輪状ならば、$K^\bullet$ は K-平坦である。
\end{lemma}

\begin{proof}
% P02-E0896: frozen source line 14342 reads "tensor product commute".
テンソル積が余極限と可換であることを繰り返し用いる
(Algebra, Lemma \ref{algebra-lemma-tensor-products-commute-with-limits})。
したがって、任意の $R$-加群は有限表示 $R$-加群 $M$ のフィルター付き余極限であるから、
任意の $R$-加群 $M$ に対して $K^\bullet \otimes_R M$ は非輪状である。
Algebra, Lemma \ref{algebra-lemma-module-colimit-fp} を参照せよ。
$M^\bullet$ を $R$-加群の非輪状複体とする。
$\text{Tot}(M^\bullet \otimes_R K^\bullet)$ が非輪状であることを示さなければならない。
$M^\bullet = \colim \tau_{\leq n} M^\bullet$（項ごとの余極限）なので
$$
\text{Tot}(M^\bullet \otimes_R K^\bullet) =
\colim \text{Tot}(\tau_{\leq n} M^\bullet \otimes_R K^\bullet)
$$
である。ここで切断は Homology, Section \ref{homology-section-truncations} のものとする。
フィルター付き余極限は完全なので
(Algebra, Lemma \ref{algebra-lemma-directed-colimit-exact})、
$M^\bullet$ を $\tau_{\leq n}M^\bullet$ で置き換え、
$M^\bullet$ が上に有界であると仮定してよい。
上に有界な場合には
$M^\bullet = \colim \sigma_{\geq -n} M^\bullet$
と書ける。ここで複体 $\sigma_{\geq -n} M^\bullet$ は有界だが、
もはや非輪状とは限らない。
上と同様に論じれば、$M^\bullet$ が有界複体である場合に帰着する。
最後に、有界複体 $M^a \to \ldots \to M^b$ に対し、
% P02-E0897: frozen source lines 14367--14370 starts induction at b-a=1 although the preceding module case is b-a=0.
複体の長さ $b - a$ に関する帰納法を用いる。
$b - a = 1$ の場合は上で既に示した。
$b - a > 1$ の場合には複体の分裂短完全列
$$
0 \to \sigma_{\geq a + 1}M^\bullet \to M^\bullet \to M^a[-a] \to 0
$$
を考え、Lemma \ref{more-algebra-lemma-derived-tor-exact} を適用して帰納段階を得る。
いくつかの詳細は省略する。
\end{proof}

\begin{lemma}
\label{more-algebra-lemma-K-flat-resolution}
$R$ を環とする。任意の複体 $M^\bullet$ に対し、各項が平坦 $R$-加群である
K-平坦複体 $K^\bullet$ と、項ごとに全射な擬同型
$K^\bullet \to M^\bullet$ が存在する。
\end{lemma}

\begin{proof}
$\mathcal{P} \subset \Ob(\text{Mod}_R)$ を平坦 $R$-加群の類とする。
Derived Categories, Lemma \ref{derived-lemma-special-direct-system} により、系
$K_1^\bullet \to K_2^\bullet \to \ldots$
および図式
$$
\xymatrix{
K_1^\bullet \ar[d] \ar[r] &
K_2^\bullet \ar[d] \ar[r] & \ldots \\
\tau_{\leq 1}M^\bullet \ar[r] &
\tau_{\leq 2}M^\bullet \ar[r] & \ldots
}
$$
であって、同補題に列挙された性質 (1), (2), (3) を持つものが存在する。
これらの性質から、各複体 $K_i^\bullet$ は平坦加群からなる上に有界な複体である。
したがって Lemma \ref{more-algebra-lemma-derived-tor-quasi-isomorphism} により
$K_i^\bullet$ は K-平坦である。
誘導される写像 $\colim_i K_i^\bullet \to M^\bullet$ は、構成により擬同型かつ項ごとに全射である。
複体 $\colim_i K_i^\bullet$ は Lemma \ref{more-algebra-lemma-colimit-K-flat} により K-平坦である。
また、フィルター付き余極限は平坦加群なので、各項 $\colim K_i^n$ は平坦である。
Algebra, Lemma \ref{algebra-lemma-colimit-flat} を参照せよ。
\end{proof}

\begin{remark}
\label{more-algebra-remark-P-resolution}
実際には Lemma \ref{more-algebra-lemma-K-flat-resolution} より強い結果が得られる。
すなわち、$R$-加群の複体 $P^\bullet$ であって、部分複体によるフィルトレーション
$$
0 = F_{-1}P^\bullet \subset F_0P^\bullet \subset
F_1P^\bullet \subset \ldots \subset P^\bullet
$$
を備え、次を満たすものからの擬同型 $P^\bullet \to M^\bullet$ を見つけられる。
\begin{enumerate}
\item $P^\bullet = \bigcup F_pP^\bullet$,
\item 包含 $F_iP^\bullet \to F_{i + 1}P^\bullet$ は項ごとに分裂単射である,
\item 商 $F_{i + 1}P^\bullet/F_iP^\bullet$ はシフト $R[k]$ の直和と
（複体として、従って微分は零である）同型である。
\end{enumerate}
これは Differential Graded Algebra, Lemma \ref{dga-lemma-resolve} で示される
（または Lemma \ref{more-algebra-lemma-resolution-filtered-complex} の証明と同様に論じてもよい）。
さらに、このような複体が与えられると、$D(R)$ における完全三角
$$
\bigoplus F_iP^\bullet \to \bigoplus F_iP^\bullet \to M^\bullet
\to \bigoplus F_iP^\bullet[1]
$$
を得る。これを用いると、一般の複体に関する命題を $R[k]$ に関する命題へ
帰着できることがある（もちろん、命題が直和を取る操作で保たれる場合に限る）。
より正確には、$T$ を $D(R)$ の対象の性質とする。次を仮定する。
\begin{enumerate}
\item $K_i \in D(R)$, $i \in I$ が対象の族で、すべての $i \in I$ に対して
$T(K_i)$ が成り立つならば、$T(\bigoplus K_i)$ が成り立つ,
\item $K \to L \to M \to K[1]$ が完全三角であり、三対象のうち二つで
$T$ が成り立つならば、残る対象でも $T$ が成り立つ,
\item すべての $k$ に対して $T(R[k])$ が成り立つ。
\end{enumerate}
このとき $D(R)$ のすべての対象で $T$ が成り立つ。
\end{remark}

\begin{lemma}
\label{more-algebra-lemma-derived-tor-quasi-isomorphism-other-side}
$R$ を環とする。
$\alpha : P^\bullet \to Q^\bullet$ を K-平坦 $R$-加群複体の擬同型とする。
任意の $R$-加群複体 $L^\bullet$ に対し、誘導される写像
$$
\text{Tot}(\text{id}_L \otimes \alpha) :
\text{Tot}(L^\bullet \otimes_R P^\bullet)
\longrightarrow
\text{Tot}(L^\bullet \otimes_R Q^\bullet)
$$
は擬同型である。
\end{lemma}

\begin{proof}
$K^\bullet$ が K-平坦である擬同型 $K^\bullet \to L^\bullet$ を選ぶ。
Lemma \ref{more-algebra-lemma-K-flat-resolution} を参照せよ。
可換図式
$$
\xymatrix{
\text{Tot}(K^\bullet \otimes_R P^\bullet) \ar[r] \ar[d] &
\text{Tot}(K^\bullet \otimes_R Q^\bullet) \ar[d] \\
\text{Tot}(L^\bullet \otimes_R P^\bullet) \ar[r] &
\text{Tot}(L^\bullet \otimes_R Q^\bullet)
}
$$
を考える。Lemma \ref{more-algebra-lemma-K-flat-quasi-isomorphism} により、
縦の二写像と上の横写像が擬同型なので、結果が従う。
\end{proof}

\noindent
$R$ を環とし、$M^\bullet$ を $D(R)$ の対象とする。
K-平坦分解 $K^\bullet \to M^\bullet$ を選ぶ。
Lemma \ref{more-algebra-lemma-K-flat-resolution} を参照せよ。
Lemmas \ref{more-algebra-lemma-derived-tor-homotopy} and \ref{more-algebra-lemma-derived-tor-exact} により、
三角圏の間の完全関手
$$
K(R) \longrightarrow K(R), \quad
L^\bullet \longmapsto \text{Tot}(L^\bullet \otimes_R K^\bullet)
$$
を得る。Lemma \ref{more-algebra-lemma-K-flat-quasi-isomorphism} により、この関手は
$D(R)$ が擬同型に関する $K(R)$ の局所化であることから、関手 $D(R) \to D(R)$ を誘導する。
$M^\bullet$ の $K$-平坦分解の圏は余フィルター付きであり、
Lemma \ref{more-algebra-lemma-derived-tor-quasi-isomorphism-other-side} が成り立つので、
得られる関手は（同型を除いて）$K^\bullet$ の選択に依存しない。

\begin{definition}
\label{more-algebra-definition-derived-tor}
$R$ を環とし、$M^\bullet$ を $D(R)$ の対象とする。
{it 導来テンソル積}とは、上で述べた三角圏の完全関手
$$
- \otimes_R^{\mathbf{L}} M^\bullet : D(R) \longrightarrow D(R)
$$
のことである。
\end{definition}

\noindent
この関手は関手 (\ref{more-algebra-equation-derived-tensor-module}) を拡張する。
明示的な構成から、$L^\bullet$ と $M^\bullet$ がともに $D(R)$ に属するとき、
同型（符号の選択を伴う。下を参照）
$$
M^\bullet \otimes_R^{\mathbf{L}} L^\bullet
\cong
L^\bullet \otimes_R^{\mathbf{L}} M^\bullet
$$
があることは明らかである。
従って $M^\bullet \otimes_R^{\mathbf{L}} L^\bullet$ と書くとき、
通常はどちらの変数を用いて導来テンソル積を定義しているかを区別しない。

\begin{lemma}
\label{more-algebra-lemma-flip-douoble-tensor-product}
$R$ を環とし、$K^\bullet, L^\bullet$ を $R$-加群の複体とする。
両複体に関して関手的な標準同型
$$
K^\bullet \otimes_R^\mathbf{L} L^\bullet \longrightarrow
L^\bullet \otimes_R^\mathbf{L} K^\bullet
$$
が存在し、写像 $K^p \otimes_R L^q \to L^q \otimes_R K^p$ には
符号 $(-1)^{pq}$ を用いる（説明は証明を参照）。
\end{lemma}

\begin{proof}
複体を K-平坦複体 $K^\bullet$ と $L^\bullet$ で置き換え、
Section \ref{more-algebra-section-symmetric-monoidal} で述べた可換性制約を用いる。
\end{proof}

\begin{lemma}
\label{more-algebra-lemma-triple-tensor-product}
$R$ を環とし、$K^\bullet, L^\bullet, M^\bullet$ を $R$-加群の複体とする。
三複体すべてに関して関手的な標準同型
$$
(K^\bullet \otimes_R^\mathbf{L} L^\bullet) \otimes_R^\mathbf{L} M^\bullet
=
K^\bullet \otimes_R^\mathbf{L} (L^\bullet \otimes_R^\mathbf{L} M^\bullet)
$$
が存在する。
\end{lemma}

\begin{proof}
複体を K-平坦複体で置き換え、
Section \ref{more-algebra-section-symmetric-monoidal} の結合性制約を用いる。
\end{proof}

\begin{lemma}
\label{more-algebra-lemma-factor-through-K-flat}
$R$ を環とし、$a : K^\bullet \to L^\bullet$ を $R$-加群複体の写像とする。
$K^\bullet$ が K-平坦ならば、複体 $N^\bullet$ と複体の写像
$b : K^\bullet \to N^\bullet$, $c : N^\bullet \to L^\bullet$ であって次を満たすものが存在する。
\begin{enumerate}
\item $N^\bullet$ は K-平坦である,
\item $c$ は擬同型である,
\item $a$ は $c \circ b$ とホモトピックである。
\end{enumerate}
$K^\bullet$ の各項が平坦ならば、$N^\bullet$ についても同じことが成り立つように
$N^\bullet$, $b$, $c$ を選べる。
\end{lemma}

\begin{proof}
ホモトピー圏 $K(R)$ が三角圏であることを用いる。
Derived Categories, Proposition
\ref{derived-proposition-homotopy-category-triangulated} を参照せよ。
完全三角 $K^\bullet \to L^\bullet \to C^\bullet \to K^\bullet[1]$ を選ぶ。
各項が平坦な K-平坦複体 $M^\bullet$ からの擬同型
$M^\bullet \to C^\bullet$ を選ぶ。Lemma \ref{more-algebra-lemma-K-flat-resolution} を参照せよ。
三角圏の公理により、合成 $M^\bullet \to C^\bullet \to K^\bullet[1]$ を
完全三角 $K^\bullet \to N^\bullet \to M^\bullet \to K^\bullet[1]$ に組み込める。
Lemma \ref{more-algebra-lemma-K-flat-two-out-of-three} により $N^\bullet$ は K-平坦である。
再び三角圏の公理を用いると、次の完全三角の射に組み込まれる写像
$N^\bullet \to L^\bullet$ を選べる。
$$
\xymatrix{
K^\bullet \ar[r] \ar[d] &
N^\bullet \ar[r] \ar[d] &
M^\bullet \ar[r] \ar[d] &
K^\bullet[1] \ar[d] \\
K^\bullet \ar[r] &
L^\bullet \ar[r] &
C^\bullet \ar[r] &
K^\bullet[1]
}
$$
三写像のうち二つが擬同型なので、これらの完全三角に付随するコホモロジー長完全列により、
第三の写像 $N^\bullet \to L^\bullet$ も擬同型である
（または、この図式の $D(R)$ における像を見て Derived Categories, Lemma
\ref{derived-lemma-third-isomorphism-triangle} を用いてもよい）。
これで (1), (2), (3) の証明が完了する。
最後の主張を示すには、$N^n \cong M^n \oplus K^n$ となるように
$N^\bullet$ を選べばよい。Derived Categories, Lemma
\ref{derived-lemma-improve-distinguished-triangle-homotopy} を参照せよ。
従って $K^\bullet$ の各項が平坦ならば、所望の平坦性が得られる。
\end{proof}

\section{環の導来変更}
\label{more-algebra-section-derived-base-change}

\noindent
$R \to A$ を環写像とし、$N^\bullet$ を $A$-加群の複体とする。
K-平坦分解を用いて関手
$$
- \otimes_R^{\mathbf{L}} N^\bullet : D(R) \to D(A)
$$
を、関手 $K(R) \to K(A)$,
$M^\bullet \mapsto \text{Tot}(M^\bullet \otimes_R N^\bullet)$ の左導来関手として定義することもできる。
特に $N^\bullet = A[0]$ とすると、関手 (\ref{more-algebra-equation-derived-tensor-algebra}) を拡張する
導来基底変換関手
$$
- \otimes_R^{\mathbf{L}} A : D(R) \to D(A)
$$
を得る。すなわち、任意の $R$-加群複体 $M^\bullet$ に対して
K-平坦分解 $K^\bullet \to M^\bullet$ を選び、
$$
M^\bullet \otimes_R^{\mathbf{L}} N^\bullet =
\text{Tot}(K^\bullet \otimes_R N^\bullet).
$$
と置ける。
Lemmas \ref{more-algebra-lemma-K-flat-resolution} and
\ref{more-algebra-lemma-derived-tor-quasi-isomorphism-other-side} を用いれば、
これが良定義であることが分かる。しかし細部まですべて厳密に確認するには、
一般論を用いる方が便利かもしれない。

\begin{lemma}
\label{more-algebra-lemma-derived-base-change}
上の構成は選択に依存せず、三角圏の完全関手
$- \otimes_R^\mathbf{L} N^\bullet : D(R) \to D(A)$ を定める。
$D(R)$ の $E^\bullet$ に対して関手的同型
$$
E^\bullet \otimes_R^\mathbf{L} N^\bullet =
(E^\bullet \otimes_R^\mathbf{L} A) \otimes_A^\mathbf{L} N^\bullet
$$
がある。
\end{lemma}

\begin{proof}
導来関手 $- \otimes_R^\mathbf{L} N^\bullet$ の存在を示すため、
Derived Categories, Section \ref{derived-section-derived-functors} で展開された一般論を用いる。
$\mathcal{D} = K(R)$, $\mathcal{D}' = D(A)$ と置く。
規則 $F(M^\bullet) = \text{Tot}(M^\bullet \otimes_R N^\bullet)$ により定まる
三角圏の完全関手を $F : \mathcal{D} \to \mathcal{D}'$ と書く。
$F$ の述べた性質は
Lemmas \ref{more-algebra-lemma-derived-tor-homotopy} and \ref{more-algebra-lemma-derived-tor-exact} から従う。
$S$ を $\mathcal{D} = K(R)$ における擬同型の集合とする。
これは Derived Categories, Situation \ref{derived-situation-derived-functor} の状況を与えるので、
Derived Categories, Definition
\ref{derived-definition-right-derived-functor-defined} が適用できる。
$LF$ は至る所で定義されると主張する。
これは Derived Categories, Lemma \ref{derived-lemma-find-existence-computes} で
$\mathcal{P} \subset \Ob(\mathcal{D})$ を K-平坦複体の集まりとすれば従う。
実際、(1) は Lemma \ref{more-algebra-lemma-K-flat-resolution} から、
(2) は Lemma \ref{more-algebra-lemma-derived-tor-quasi-isomorphism-other-side} から従う。
従って導来関手
$$
LF : D(R) = S^{-1}\mathcal{D} \longrightarrow \mathcal{D}' = D(A)
$$
を得る。Derived Categories, Equation (\ref{derived-equation-everywhere}) を参照せよ。
最後に Derived Categories, Lemma \ref{derived-lemma-find-existence-computes} により、
$K^\bullet$ が K-平坦なとき
$LF(K^\bullet) = F(K^\bullet) =
\text{Tot}(K^\bullet \otimes_R N^\bullet)$ である。
従って $LF$ は確かに上で述べた方法で計算される。
さらに Lemma \ref{more-algebra-lemma-base-change-K-flat} により、複体
$K^\bullet \otimes_R A$ は K-平坦 $A$-加群複体である。
% P02-E0898: frozen source lines 14719--14723 tensor K over A before K has an A-module structure.
従って
$$
(K^\bullet \otimes_R^\mathbf{L} A) \otimes_A^\mathbf{L} N^\bullet =
\text{Tot}((K^\bullet \otimes_R A) \otimes_A N^\bullet) =
\text{Tot}(K^\bullet \otimes_A N^\bullet) =
K^\bullet \otimes_A^\mathbf{L} N^\bullet
$$
であり、補題の最後の主張が従う。
\end{proof}

\begin{lemma}
\label{more-algebra-lemma-functoriality-derived-base-change}
$R \to A$ を環写像とし、$f : L^\bullet \to N^\bullet$ を
$A$-加群複体の写像とする。このとき $f$ は関手の変換
% P02-E0899: frozen statement tensors D(R) inputs over A although its proof uses tensor over R.
$$
1 \otimes f :
- \otimes_A^\mathbf{L} L^\bullet
\longrightarrow
- \otimes_A^\mathbf{L} N^\bullet
$$
を誘導する。$f$ が擬同型ならば、$1 \otimes f$ は関手の同型である。
\end{lemma}

\begin{proof}
% P02-E0900: frozen source line 14743 reads "functors are computing".
これらの関手は K-平坦複体 $K^\bullet$ 上で評価して計算されるので、関手性
$$
\text{Tot}(K^\bullet \otimes_R L^\bullet) \to
\text{Tot}(K^\bullet \otimes_R N^\bullet)
$$
をそのまま用いて変換を定義できる。
最後の主張は Lemma \ref{more-algebra-lemma-K-flat-quasi-isomorphism} から従う。
\end{proof}

\begin{lemma}
\label{more-algebra-lemma-tensor-hom-adjoint}
$R \to A$ を環写像とする。
Lemma \ref{more-algebra-lemma-derived-base-change} の関手
$D(R) \to D(A)$, $E \mapsto E \otimes_R^\mathbf{L} A$ は、
制限関手 $D(A) \to D(R)$ の左随伴である。
\end{lemma}

\begin{proof}
これは Derived Categories, Lemma
\ref{derived-lemma-pre-derived-adjoint-functors-general} と、
Algebra, Lemma \ref{algebra-lemma-adjoint-tensor-restrict} により
$- \otimes_R A$ と制限が随伴であることから従う。
\end{proof}

\begin{remark}[注意]
\label{more-algebra-remark-warning-compute-base-change}
$R \to A$ を環写像とし、$N$, $N'$ を $A$-加群とする。
$N$, $N'$ の $R$-加群への制限をそれぞれ $N_R$, $N'_R$ と書く。
Algebra, Section \ref{algebra-section-base-change} を参照せよ。
この状況では、$D(A)$ の対象 $N_R \otimes_R^\mathbf{L} N'$ と
$N \otimes_R^\mathbf{L} N'_R$ は一般には同型でない！
言い換えると、二つの側のどちらを用いて $A$-加群構造を与えるかに
十分注意しなければならない。

\medskip\noindent
具体例として $R = k[x, y]$, $A = R/(xy)$, $N = R/(x)$,
$N' = A = R/(xy)$ と置く。分解
$0 \to R \xrightarrow{xy} R \to N'_R \to 0$ から、
$D(A)$ において $N \otimes_R^\mathbf{L} N'_R = N[1] \oplus N$ である。
分解 $0 \to R \xrightarrow{x} R \to N_R \to 0$ から、
$N_R \otimes_R^\mathbf{L} N'$ は複体 $A \xrightarrow{x} A$ で表される。
これら二複体が同型でないことを見るには、後者が $D(A)$ において
そのコホモロジー群の直和と同型でないことを示してもよいし、
前者が $D(A)$ の完全対象でない一方、後者は完全対象であることを示してもよい。
いくつかの詳細は省略する。
\end{remark}

\begin{lemma}
\label{more-algebra-lemma-double-base-change}
$A \to B \to C$ を環写像とし、$N^\bullet$ を $B$-加群の複体、
$K^\bullet$ を $C$-加群の複体とする。
関手の合成
% P02-E0901: frozen source line 14804 uses singular "is" with plural subject "compositions".
$$
D(A) \xrightarrow{- \otimes_A^\mathbf{L} N^\bullet}
D(B) \xrightarrow{- \otimes_B^\mathbf{L} K^\bullet} D(C)
$$
は関手
$- \otimes_A^\mathbf{L} (N^\bullet \otimes_B^\mathbf{L} K^\bullet) :
D(A) \to D(C)$ である。
$M$, $N$, $K$ がそれぞれ $A$, $B$, $C$ 上の加群ならば、$D(C)$ において
$$
(M \otimes_A^\mathbf{L} N) \otimes_B^\mathbf{L} K =
M \otimes_A^\mathbf{L} (N \otimes_B^\mathbf{L} K) =
(M \otimes_A^\mathbf{L} C) \otimes_C^\mathbf{L} (N \otimes_B^\mathbf{L} K)
$$
である。また、符号を用いる標準同型
$$
(M \otimes_A^\mathbf{L} N) \otimes_B^\mathbf{L} K \longrightarrow
(M \otimes_A^\mathbf{L} K) \otimes_C^\mathbf{L} (N \otimes_B^\mathbf{L} C)
$$
もある。
% P02-E0902: frozen source line 14818 reads "Similar results holds".
複体についても同様の結果が成り立つ。
\end{lemma}

\begin{proof}
$B$-加群の K-平坦複体 $P^\bullet$ と擬同型
$P^\bullet \to N^\bullet$ を選ぶ（Lemma \ref{more-algebra-lemma-K-flat-resolution}）。
$M^\bullet$ を $D(A)$ の任意の対象を表す K-平坦 $A$-加群複体とする。
すると、括弧内の部分に Lemma \ref{more-algebra-lemma-functoriality-derived-base-change} を適用することにより
$$
(M^\bullet \otimes_A^\mathbf{L} P^\bullet) \otimes_B^\mathbf{L} K^\bullet
\longrightarrow
(M^\bullet \otimes_A^\mathbf{L} N^\bullet) \otimes_B^\mathbf{L} K^\bullet
$$
は同型である。
Lemmas \ref{more-algebra-lemma-base-change-K-flat} and \ref{more-algebra-lemma-tensor-product-K-flat} により、複体
% P02-E0903: frozen source lines 14836--14838 contain an undefined R and an unmatched opening parenthesis.
$$
\text{Tot}(M^\bullet \otimes_A P^\bullet) =
\text{Tot}((M^\bullet \otimes_R A) \otimes_A P^\bullet
$$
は $B$-加群複体として K-平坦であり、構成により $D(B)$ における導来テンソル積を表す。
従って
$(M^\bullet \otimes_A^\mathbf{L} P^\bullet) \otimes_B^\mathbf{L} K^\bullet$
は $C$-加群複体
$$
\text{Tot}(\text{Tot}(M^\bullet \otimes_A P^\bullet)\otimes_B K^\bullet) =
\text{Tot}(M^\bullet \otimes_A \text{Tot}(P^\bullet \otimes_B K^\bullet))
$$
で表される。等号は Homology, Remark \ref{homology-remark-triple-complex} による。
来た道を戻ると、これは
$$
M^\bullet \otimes_A^\mathbf{L} (P^\bullet \otimes_B^\mathbf{L} K^\bullet)
\longleftarrow
M^\bullet \otimes_A^\mathbf{L} (N^\bullet \otimes_B^\mathbf{L} K^\bullet)
$$
に等しいことが分かる。この矢印は関手
$-\otimes_B^\mathbf{L} K^\bullet$ の定義により同型である。
これらすべての構成は複体 $M^\bullet$ に関して関手的なので、所望の関手の同型を得る。

\medskip\noindent
上から、次の最初の等号が得られる。
$$
(M \otimes_A^\mathbf{L} N) \otimes_B^\mathbf{L} K =
M \otimes_A^\mathbf{L} (N \otimes_B^\mathbf{L} K) =
(M \otimes_A^\mathbf{L} C) \otimes_C^\mathbf{L} (N \otimes_B^\mathbf{L} K)
$$
第二の等号は Lemma \ref{more-algebra-lemma-derived-base-change} の最後の主張から従う。
同じことから
$N \otimes_B^\mathbf{L} K = (N \otimes_B^\mathbf{L} C) \otimes_C^\mathbf{L} K$
と書けるので、これを代入すると
\begin{align*}
(M \otimes_A^\mathbf{L} N) \otimes_B^\mathbf{L} K
& =
(M \otimes_A^\mathbf{L} C) \otimes_C^\mathbf{L}
((N \otimes_B^\mathbf{L} C) \otimes_C^\mathbf{L} K) \\
& =
(M \otimes_A^\mathbf{L} C) \otimes_C^\mathbf{L}
(K \otimes_C^\mathbf{L} (N \otimes_B^\mathbf{L} C)) \\
% P02-E0904/P02-E0905: frozen source lines 14882--14886 have an extra closing parenthesis and tensor M over C although M is only an A-module.
& =
((M \otimes_A^\mathbf{L} C) \otimes_C^\mathbf{L} K)
\otimes_C^\mathbf{L} (N \otimes_B^\mathbf{L} C)) \\
& =
(M \otimes_C^\mathbf{L} K)
\otimes_C^\mathbf{L} (N \otimes_B^\mathbf{L} C)
\end{align*}
である。ここでは Lemmas \ref{more-algebra-lemma-flip-douoble-tensor-product} and
\ref{more-algebra-lemma-triple-tensor-product} および前述の補題を用いた。
\end{proof}

\section{Tor 独立性}
\label{more-algebra-section-tor-independence}

\noindent
環の可換図式
$$
\xymatrix{
A \ar[r] & A' \\
R \ar[r] \ar[u] & R' \ar[u]
}
$$
を考える。$D(A)$ の対象 $K$ が与えられると、その導来基底変換
$K \otimes_A^\mathbf{L} A'$ を $D(A')$ の対象として考えられる。
または、$K$ を $D(R)$ の対象に制限し、これを $D(R')$ の対象へ導来基底変換したもの
$K \otimes_R^\mathbf{L} R'$ を考えられる。
$D(R')$ において関手的な比較写像
\begin{equation}
\label{more-algebra-equation-comparison-map}
K \otimes_R^{\mathbf{L}} R'
\longrightarrow
K \otimes_A^{\mathbf{L}} A'
\end{equation}
があると主張する。この比較写像を構成するため、$K$ を表す
K-平坦 $A$-加群複体 $K^\bullet$ を選ぶ。
次に、$E^\bullet$ が K-平坦 $R$-加群複体である擬同型
$E^\bullet \to K^\bullet$ を選ぶ。上の写像は
$$
K \otimes_R^{\mathbf{L}} R' =
E^\bullet \otimes_R R'
\longrightarrow
K^\bullet \otimes_A A' =
K \otimes_A^{\mathbf{L}} A'
$$
である。一般には、この写像が同型となる見込みはない。

\medskip\noindent
しかし、上の図式が環の「基底変換」図式、すなわち
$A' = A \otimes_R R'$ である状況にはしばしば出会う。
この場合、任意の $A$-加群 $M$ に対して
$M \otimes_A A' = M \otimes_R R'$ である。
従って $- \otimes_R R'$ と $- \otimes_A A'$ は関手
$\text{Mod}_A \to \text{Mod}_{A'}$ として等しい。
一般には、この等式は導来テンソル積へ{\bf 拡張されない}。
言い換えると、比較写像は同型でない。
簡単な例として
$R = k[x]$, $A = R' = A' = k[x]/(x) = k$, $K^\bullet = A[0]$ を取ればよい。
明らかに必要な条件は、すべての $p > 0$ に対して
$\text{Tor}_p^R(A, R') = 0$ となることである。

\begin{definition}
\label{more-algebra-definition-tor-independent}
$R$ を環とし、$A$, $B$ を $R$-代数とする。
すべての $p > 0$ に対して $\text{Tor}_p^R(A, B) = 0$ であるとき、
$A$ と $B$ は $R$ 上 {\it Tor 独立}であるという。
\end{definition}

\begin{lemma}
\label{more-algebra-lemma-base-change-comparison}
$A' = A \otimes_R R'$ であり、$A$ と $R'$ が $R$ 上 Tor 独立ならば、
比較写像 (\ref{more-algebra-equation-comparison-map}) は同型である。
\end{lemma}

\begin{proof}
これを示すため、$R$-加群としての $R'$ の自由分解
$F^\bullet \to R'$ を選ぶ。
$A$ と $R'$ は $R$ 上 Tor 独立なので、
$F^\bullet \otimes_R A$ は $A$ 上の $A'$ の自由 $A$-加群分解である。
上の導来テンソル積の一般構成により
$$
K^\bullet \otimes_A A' \cong
\text{Tot}(K^\bullet \otimes_A (F^\bullet \otimes_R A)) =
\text{Tot}(K^\bullet \otimes_R F^\bullet) \cong
\text{Tot}(E^\bullet \otimes_R F^\bullet) \cong
E^\bullet \otimes_R R'
$$
となり、所望の結果を得る。
\end{proof}

\begin{lemma}
\label{more-algebra-lemma-tor-independent-flat}
環の可換図式
$$
\xymatrix{
A' & R' \ar[r] \ar[l] & B' \\
A \ar[u] & R \ar[l] \ar[u] \ar[r] & B \ar[u]
}
$$
を考える。$R'$ は $R$ 上平坦であり、$A'$ は $A \otimes_R R'$ 上平坦、
$B'$ は $R' \otimes_R B$ 上平坦であると仮定する。このとき
$$
\text{Tor}_i^R(A, B) \otimes_{(A \otimes_R B)} (A' \otimes_{R'} B') =
\text{Tor}_i^{R'}(A', B')
$$
である。
\end{lemma}

\begin{proof}
Algebra, Section \ref{algebra-section-functoriality-tor} により標準写像
$$
\text{Tor}_i^R(A, B) \longrightarrow
\text{Tor}_i^{R'}(A \otimes_R R', B \otimes_R R') \longrightarrow
\text{Tor}_i^{R'}(A', B')
$$
がある。これらは補題の式の左辺から右辺への写像を誘導する。

\medskip\noindent
$R$-加群としての $A$ の自由分解 $F_\bullet \to A$ を取る。
すると $F_\bullet \otimes_R R'$ は $A \otimes_R R'$ の分解である。
従って $\text{Tor}_i^{R'}(A \otimes_R R', B \otimes_R R')$ は
$F_\bullet \otimes_R B \otimes_R R'$ により計算される。
$R'$ は $R$ 上平坦であるという仮定により、これは
$\text{Tor}_i^R(A, B) \otimes_R R'$ を計算する。
従って
$\text{Tor}_i^{R'}(A \otimes_R R', B \otimes_R R') =
\text{Tor}_i^R(A, B) \otimes_R R'$ である
（ここでは $R'$ の $R$ 上の平坦性しか用いていない）。

\medskip\noindent
Lazard の定理 (Algebra, Theorem \ref{algebra-theorem-lazard}) により、
$A'$, resp.\ $B'$ を、有限自由
$A \otimes_R R'$, resp.\ $B \otimes_R R'$-加群のフィルター付き余極限として書ける。
$A' = \colim M_i$, $B' = \colim N_j$ としよう。
上の結果から
$$
\text{Tor}_i^{R'}(M_i, N_j) =
\text{Tor}_i^R(A, B) \otimes_{A \otimes_R B} (M_i \otimes_{R'} N_j)
$$
を得る。これは基底を用いてすべて書き下せば分かる。
余極限を取れば補題の結果を得る。
\end{proof}

\begin{lemma}
\label{more-algebra-lemma-flat-base-change-tor-independent}
$R \to A$, $R \to B$ を環写像とする。$R \to R'$ を環写像とし、
$A' = A \otimes_R R'$, $B' = B \otimes_R R'$ と置く。
% P02-E0906: frozen source line 15029 lowercases "tor" despite using Tor as the named functor throughout.
$A$ と $B$ が $R$ 上 Tor 独立で、$R \to R'$ が平坦ならば、
$A'$ と $B'$ は $R'$ 上 Tor 独立である。
\end{lemma}

\begin{proof}
Lemma \ref{more-algebra-lemma-tor-independent-flat} と
Definition \ref{more-algebra-definition-tor-independent} から直ちに従う。
\end{proof}

\begin{lemma}
\label{more-algebra-lemma-lemma-tor-independent-flat-compare}
Lemma \ref{more-algebra-lemma-tor-independent-flat} と同じ仮定を置く。
$M \in D(A)$ に対して、$A' \otimes_{R'} B'$-加群の標準同型
$$
H^i((M \otimes_A^\mathbf{L} A') \otimes_{R'}^\mathbf{L} B') =
H^i(M \otimes_R^\mathbf{L} B) \otimes_{(A \otimes_R B)} (A' \otimes_{R'} B')
$$
がある。
\end{lemma}

\begin{proof}
等式の両辺を詳しく説明しよう。
左辺では、関手 $D(A) \to D(A') \to D(R') \to D(B')$ と
関手 $H^i : D(B') \to \text{Mod}_{B'}$ の合成を考えている。
$D(B')$ の対象
$(M \otimes_A^\mathbf{L} A') \otimes_{R'}^\mathbf{L} B'$ の自己準同型への
$A'$ からの写像があるので、左辺は確かに $A' \otimes_{R'} B'$-加群である。
同じ議論により、$H^i(M \otimes_R^\mathbf{L} B)$ は
$A \otimes_R B$-加群構造を持つ。

\medskip\noindent
まず $B' = R' \otimes_R B$ の場合に結果を示す。
自由 $R$-加群による分解 $F^\bullet \to B$ を選ぶ。
また、$M$ を表す K-平坦 $A$-加群複体 $M^\bullet$ を選ぶ。
このとき左辺は
\begin{align*}
H^i(\text{Tot}((M^\bullet \otimes_A A') \otimes_{R'} (R' \otimes_R F^\bullet)))
& =
H^i(\text{Tot}(M^\bullet \otimes_A A' \otimes_R F^\bullet)) \\
& =
H^i(\text{Tot}(M^\bullet \otimes_R F^\bullet) \otimes_A A') \\
& =
H^i(M \otimes_R^\mathbf{L} B) \otimes_A A'
\end{align*}
で表される。
% P02-E0907: frozen source line 15076 is the incomplete sentence "The final equality because ...".
最後の等号は $A \to A'$ が平坦であることによる。
この段落では $B' = R' \otimes_R B$ と仮定したので
$A' \otimes_{R'} B' = A' \otimes_R B$ であり、最後の加群は所望の加群である。

\medskip\noindent
一般の場合を考える。$B' \to B''$ を平坦な環写像とする。
このとき
$$
H^i((M \otimes_A^\mathbf{L} A') \otimes_{R'}^\mathbf{L} B'') =
H^i((M \otimes_A^\mathbf{L} A') \otimes_{R'}^\mathbf{L} B')
\otimes_{B'} B''
$$
および
$$
H^i(M \otimes_R^\mathbf{L} B) \otimes_{(A \otimes_R B)} (A' \otimes_{R'} B'')
=
\left(
H^i(M \otimes_R^\mathbf{L} B) \otimes_{(A \otimes_R B)} (A' \otimes_{R'} B')
\right) \otimes_{B'} B''
$$
であることは容易に分かる。
従って $B'$ に対する結果から $B''$ に対する結果が従う。
前段落で $R' \otimes_R B$ に対する結果を示したので、一般の場合も従う。
\end{proof}

\begin{lemma}
\label{more-algebra-lemma-tor-independent}
$R$ を環とし、$A$, $B$ を $R$-代数とする。次は同値である。
\begin{enumerate}
\item $A$ と $B$ は $R$ 上 Tor 独立である,
\item 同じ素イデアル $\mathfrak r \subset R$ 上にある任意の素イデアル対
$\mathfrak p \subset A$, $\mathfrak q \subset B$ に対して、環
$A_\mathfrak p$ と $B_\mathfrak q$ は $R_\mathfrak r$ 上 Tor 独立である,
\item $A \otimes_R B$ の任意の素イデアル $\mathfrak s$ に対して加群
$$
\text{Tor}_i^R(A, B)_\mathfrak s =
\text{Tor}_i^{R_\mathfrak r}(A_\mathfrak p, B_\mathfrak q)_\mathfrak s
$$
は零である。ここで $\mathfrak p = A \cap \mathfrak s$,
$\mathfrak q = B \cap \mathfrak s$, $\mathfrak r = R \cap \mathfrak s$ とする。
\end{enumerate}
\end{lemma}

\begin{proof}
$\mathfrak s$ を (3) のような $A \otimes_R B$ の素イデアルとする。
等式
$$
\text{Tor}_i^R(A, B)_\mathfrak s =
\text{Tor}_i^{R_\mathfrak r}(A_\mathfrak p, B_\mathfrak q)_\mathfrak s
$$
は、$\mathfrak p = A \cap \mathfrak s$,
$\mathfrak q = B \cap \mathfrak s$, $\mathfrak r = R \cap \mathfrak s$ とすれば、
Lemma \ref{more-algebra-lemma-tor-independent-flat} から従う。
従って (2) は (3) を含意する。
加群の消滅は素イデアルで局所化して判定できるので
(Algebra, Lemma \ref{algebra-lemma-characterize-zero-local})、(3) は (1) を含意する。
(1) $\Rightarrow$ (2) には、再び Lemma \ref{more-algebra-lemma-tor-independent-flat} による等式
$$
\text{Tor}_i^{R_\mathfrak r}(A_\mathfrak p, B_\mathfrak q) =
\text{Tor}_i^R(A, B) \otimes_{(A \otimes_R B)}
(A_\mathfrak p \otimes_{R_{\mathfrak r}} B_\mathfrak q)
$$
を用いる。
\end{proof}


\section{Tor のスペクトル列}
\label{more-algebra-section-spectral-sequence-tor}

\noindent
この節では、Tor 関手を考えるときに現れるさまざまなスペクトル列を集める。

\begin{example}
\label{more-algebra-example-cohomology-complex-tensored}
$R$ を環とする。$K_\bullet$ を $K_n = 0$ for $n \ll 0$ を満たす
$R$-加群の鎖複体とする。$M$ を $R$-加群とする。
自由 $R$-加群による $M$ の分解 $P_\bullet \to M$ を選ぶ。
二重鎖複体 $K_\bullet \otimes_R P_\bullet$ を得る。
Homology, Section \ref{homology-section-double-complex} の内容
（特に Homology, Lemma \ref{homology-lemma-first-quadrant-ss}）を
鎖複体の言葉へ移すと、$H_*(K_\bullet \otimes_R^\mathbf{L} M)$ に収束する
二つのスペクトル列を得る。
一方は $E_2$-項
$$
(E_2)_{i, j} = \text{Tor}^R_j(H_i(K_\bullet), M)
\Rightarrow
H_{i + j}(K_\bullet \otimes^{\mathbf{L}}_R M)
$$
を持ち、微分 $d_2$ は写像
$\text{Tor}^R_j(H_i(K_\bullet), M) \to
\text{Tor}^R_{j - 2}(H_{i + 1}(K_\bullet), M)$
で与えられる。もう一方は $E_1$-項
$$
(E_1)_{i, j} = \text{Tor}^R_j(K_i, M)
\Rightarrow
H_{i + j}(K_\bullet \otimes^{\mathbf{L}}_R M)
$$
を持ち、微分 $d_1$ は $K_i \to K_{i - 1}$ が誘導する写像
$\text{Tor}^R_j(K_i, M) \to \text{Tor}^R_j(K_{i - 1}, M)$ で与えられる。
\end{example}

\begin{example}
\label{more-algebra-example-tor-change-rings}
$R \to S$ を環写像とし、$M$ を $R$-加群、$N$ を $S$-加群とする。
このときスペクトル列
$$
\text{Tor}^S_n(\text{Tor}^R_m(M, S), N) \Rightarrow
\text{Tor}^R_{n + m}(M, N).
$$
がある。
% P02-E0908: frozen source line 15199 says "choose a R-free resolution".
これを構成するには、$M$ の $R$-自由分解 $P^\bullet$ を選ぶ。このとき
$$
M \otimes_R^{\mathbf{L}} N = P^\bullet \otimes_R N =
(P^\bullet \otimes_R S) \otimes_S N
$$
であり、Example \ref{more-algebra-example-cohomology-complex-tensored} の第一のスペクトル列を適用すればよい。
\end{example}

\begin{example}
\label{more-algebra-example-tor-base-change}
可換図式
$$
\xymatrix{
B \ar[r] & B' = B \otimes_A A' \\
A \ar[r] \ar[u] & A' \ar[u]
}
$$
と $B$-加群 $M, N$ を考える。
$M' = M \otimes_A A' = M \otimes_B B'$,
$N' = N \otimes_A A' = N \otimes_B B'$ と置く。
{\it $A \to B$ が平坦であり、$M$ と $N$ が $A$-平坦であると仮定する。}
このときスペクトル列
$$
\text{Tor}^A_i(\text{Tor}_j^B(M, N), A')
\Rightarrow
\text{Tor}^{B'}_{i + j}(M', N')
$$
がある。理由は次の通りである。
% P02-E0909: frozen source line 15226 omits the article in "Choose free resolution".
自由 $B$-加群による分解 $F_\bullet \to M$ を選ぶ。
$B$ と $M$ は $A$-平坦なので、$F_\bullet \otimes_A A'$ は
$M'$ の自由 $B'$-分解である。従って群 $\text{Tor}^{B'}_n(M', N')$ は複体
$$
(F_\bullet \otimes_A A') \otimes_{B'} N' =
(F_\bullet \otimes_B N) \otimes_A A' =
(F_\bullet \otimes_B N) \otimes^{\mathbf{L}}_A A'
$$
によって計算される。
% P02-E0910: frozen source line 15236 begins the incomplete clause "the last equality because".
最後の等号は、$N$ が $A$ 上平坦なので
$F_\bullet \otimes_B N$ が平坦 $A$-加群の複体であることによる。
従って Example \ref{more-algebra-example-cohomology-complex-tensored} のスペクトル列を適用して、
所望のスペクトル列を得る。
\end{example}

\begin{example}
\label{more-algebra-example-tor}
$K^\bullet, L^\bullet$ を $D^{-}(R)$ の対象とする。
このとき
$$
E_2^{p, q} = H^p(K^\bullet \otimes_R^{\mathbf{L}} H^q(L^\bullet))
\Rightarrow H^{p + q}(K^\bullet \otimes_R^{\mathbf{L}} L^\bullet)
$$
を持つスペクトル列と、
$$
E_2^{p, q} =
H^p(H^q(K^\bullet) \otimes_R^{\mathbf{L}} L^\bullet)
\Rightarrow H^{p + q}(K^\bullet \otimes_R^{\mathbf{L}} L^\bullet)
$$
を持つ別のスペクトル列がある。
どちらのスペクトル列も
$d_2^{p, q} : E_2^{p, q} \to E_2^{p + 2, q - 1}$ を持つ。
$K^\bullet$ と $L^\bullet$ を射影加群からなる上に有界な複体で置き換えると、
これらは Homology, Section \ref{homology-section-double-complex} で述べた
$\text{Tot}(K^\bullet \otimes L^\bullet)$ のコホモロジーを計算する二つのスペクトル列にほかならない。
\end{example}

\section{積と Tor}
\label{more-algebra-section-products-tor}

\noindent
積写像の最も簡単な例は次の状況から得られる。
$K^\bullet, L^\bullet \in D(R)$ と仮定する。このとき写像
\begin{equation}
\label{more-algebra-equation-simple-tor-product}
H^i(K^\bullet) \otimes_R H^j(L^\bullet)
\longrightarrow
H^{i + j}(K^\bullet \otimes_R^{\mathbf{L}} L^\bullet)
\end{equation}
がある。実際、これらの写像を定義するには、$K^\bullet, L^\bullet$ の一方が
K-平坦 $R$-加群複体であると仮定してよい
（例えば、自由または射影 $R$-加群からなる上に有界な複体）。
この場合 $K^\bullet \otimes_R^{\mathbf{L}} L^\bullet$ は複体
$\text{Tot}(K^\bullet \otimes_R L^\bullet)$ で表される。
Section \ref{more-algebra-section-derived-tensor-product}
（または Section \ref{more-algebra-section-computing-tor}）を参照せよ。
次に $\xi \in H^i(K^\bullet)$, $\zeta \in H^j(L^\bullet)$ とする。
$\xi$, $\zeta$ を表す
$k \in \Ker(K^i \to K^{i + 1})$,
$l \in \Ker(L^j \to L^{j + 1})$ を選ぶ。このとき
$$
\xi \cup \zeta =
\text{class of }k \otimes l\text{ in }
H^{i + j}(\text{Tot}(K^\bullet \otimes_R L^\bullet)).
$$
と置く。
% P02-E0911: frozen source line 15300 reads "This make sense".
これは意味を持つ。実際、全複体の微分 $\text{d}$ の公式
（Homology, Definition \ref{homology-definition-associated-simple-complex} を参照）
から $k \otimes l$ はコサイクルである。
さらに、ある $k'' \in K^{i - 1}$ に対して $k' = d_K(k'')$ ならば、
$l$ がコサイクルなので $k' \otimes l = \text{d}(k'' \otimes l)$ である。
同様に、$\zeta$ を表す $l$ の選択を変えても $k \otimes l$ の類は変わらない。
$\cup$ が双線形であることも同様に明らかなので、
$H^i(K^\bullet) \otimes_R H^j(L^\bullet)$ の一般の元には
$$
\sum \xi_i \otimes \zeta_i \longmapsto \sum \xi_i \cup \zeta_i
$$
を $H^{i + j}(\text{Tot}(K^\bullet \otimes_R L^\bullet))$ の元として対応させる。

\medskip\noindent
$R \to A$ を環写像とし、$K^\bullet, L^\bullet \in D(R)$ とする。
このとき $D(A)$ における標準同一視
\begin{equation}
\label{more-algebra-equation-pullback-derived-tensor-product}
(K^\bullet \otimes_R^{\mathbf{L}} A)
\otimes_A^{\mathbf{L}}
(L^\bullet \otimes_R^{\mathbf{L}} A)
=
(K^\bullet \otimes_R^{\mathbf{L}} L^\bullet) \otimes_R^{\mathbf{L}} A
\end{equation}
がある。これは次のように構成される。
まず $R$ 上の K-平坦分解 $P^\bullet \to K^\bullet$,
$Q^\bullet \to L^\bullet$ を選ぶ。
すると左辺は複体
$\text{Tot}((P^\bullet \otimes_R A) \otimes_A (Q^\bullet \otimes_R A))$ で、
右辺は複体
$\text{Tot}(P^\bullet \otimes_R Q^\bullet) \otimes_R A$ で表される。
これらの複体は標準的に同型である。従って上の構成は、場合によって有用な積
$$
\text{Tor}^R_n(K^\bullet, A) \otimes_A \text{Tor}^R_m(L^\bullet, A)
\longrightarrow
\text{Tor}_{n + m}^R(K^\bullet \otimes_R^\mathbf{L} L^\bullet, A)
$$
を誘導する。

\medskip\noindent
$M$, $N$ を $R$-加群とする。上の一般構成、標準写像
$M \otimes_R^\mathbf{L} N \to M \otimes_R N$、および
$\text{Tor}$ の関手性を用いて標準写像
\begin{equation}
\label{more-algebra-equation-tor-product}
\text{Tor}^R_n(M, A) \otimes_A \text{Tor}^R_m(N, A)
\longrightarrow \text{Tor}_{n + m}^R(M \otimes_R N, A)
\end{equation}
を得る。射影分解を用いる直接的な構成は次の通りである。
まず $R$ 上の射影分解
$$
P_\bullet \to M, \quad Q_\bullet \to N, \quad T_\bullet \to M \otimes_R N
$$
を選ぶ。$\otimes_R$ の右完全性により
$H_0(\text{Tot}(P_\bullet \otimes_R Q_\bullet)) = M \otimes_R N$ である。
従って Derived Categories, Lemmas
\ref{derived-lemma-morphisms-lift-projective} and
\ref{derived-lemma-morphisms-equal-up-to-homotopy-projective} は、
$H_0(\mu) = \text{id}_{M \otimes_R N}$ を満たす複体の写像
$\mu : \text{Tot}(P_\bullet \otimes_R Q_\bullet) \to T_\bullet$ の
存在と一意性を保証する。これは $D(A)$ における標準写像
\begin{align*}
(M \otimes_R^{\mathbf{L}} A) \otimes_A^{\mathbf{L}}
(N \otimes_R^{\mathbf{L}} A)
& =
\text{Tot}((P_\bullet \otimes_R A) \otimes_A (Q_\bullet \otimes_R A)) \\
& =
\text{Tot}(P_\bullet \otimes_R Q_\bullet) \otimes_R A \\
& \to
T_\bullet \otimes_R A \\
& = (M \otimes_R N) \otimes_R^{\mathbf{L}} A
\end{align*}
を誘導する。従って上の積 (\ref{more-algebra-equation-tor-product}) は、
$A$ 上で (\ref{more-algebra-equation-simple-tor-product}) を用いて
$$
\text{Tor}^R_n(M, A) \otimes_A \text{Tor}^R_m(N, A) \to
H^{-n-m}((M \otimes_R^{\mathbf{L}} A) \otimes_A^{\mathbf{L}}
(N \otimes_R^{\mathbf{L}} A))
$$
を構成し、次に上の表示写像と合成して
$\text{Tor}_{n + m}^R(M \otimes_R N, A)$ に到達することにより構成される。

\medskip\noindent
上の興味深い特別の場合として、$B$ が $R$-代数で $M = N = B$ である場合がある。
この場合、写像
$$
\text{Tor}_n^R(B, A) \otimes_A \text{Tor}_m^R(B, A)
\longrightarrow
\text{Tor}_{n + m}^R(B \otimes_R B, A)
\longrightarrow
\text{Tor}_{n + m}^R(B, A)
$$
を得る。第二の矢印は、$\text{Tor}$ の関手性を通じて乗法写像
$B \otimes_R B \to B$ から誘導される。
言い換えると、$\text{Tor}^R_{\star}(B, A)$ 上に $A$-代数構造を得る。
% P02-E0912/P02-E0913/P02-E0914: frozen source lines 15397--15400 mix noun/adjective forms, singular "case", and an editorial placeholder.
この代数構造は多くの興味深い性質
（結合性、次数付き可換性、$B$-代数構造、場合によっては分割冪など）を持つ。
これらについては別の箇所で論じる（ここに将来の参照を挿入する）。

\begin{lemma}
\label{more-algebra-lemma-functoriality-product-tor}
$R$ を環とし、$A, B, C$ を $R$-代数、$B \to C$ を $R$-代数写像とする。
このとき誘導される写像
$$
\text{Tor}^R_{\star}(B, A)
\longrightarrow
\text{Tor}^R_{\star}(C, A)
$$
は $A$-代数準同型である。
\end{lemma}

\begin{proof}
省略する。ヒント：すべての複体を明示的に書き下し、
定義をたどることで証明できる。
\end{proof}

\section{K\"unneth スペクトル列}
\label{more-algebra-section-kunneth-spectral-sequence}

\noindent
$R$ を環とし、$K^\bullet$ と $L^\bullet$ を $R$-加群のフィルター付き複体とする
（Homology, Definition \ref{homology-definition-filtered-complex} を参照。
ここでフィルトレーションは減少的であることに注意せよ）。
このとき複体
$$
T^\bullet = \text{Tot}(K^\bullet \otimes_R L^\bullet)
$$
にも、公式
$$
F^nT^\bullet =
\Im\left(
\bigoplus\nolimits_{i + j = n}
\text{Tot}(F^iK^\bullet \otimes_R F^jL^\bullet) \to
\text{Tot}(K^\bullet \otimes_R L^\bullet)
\right)
$$
で定義される減少フィルトレーションがある。
フィルター付き複体にいくつかの仮定を置き、
Homology, Section \ref{homology-section-filtered-complex} の構成によって
ここから生じるスペクトル列を決定する。

\medskip\noindent
各 $K^n$, $F^iK^n$, $\text{gr}^iK^n$, $L^m$, $F^jL^m$,
$\text{gr}^jL^m$ が平坦 $R$-加群であると仮定する。
この場合、加群
$F^iK^n \otimes_R L^m$, $K^n \otimes_R F^jL^m$,
$F^iK^n \otimes_R F^jL^m$ は $K^n \otimes_R L^m$ の部分加群である。
同様に、加群
$\text{gr}^iK^n \otimes_R \text{gr}^jL^m$ は
$K^n/F^{i + 1}K^n \otimes_R L^m/F^{j + 1}L^m$ の部分加群である。
写像
% P02-E0915: frozen source line 15459 writes F^nT^n although the displayed arrow is a map of complexes.
$$
F^nT^n \longrightarrow
\bigoplus\nolimits_{i + j = n}
\text{Tot}(K^\bullet / F^{i + 1}K^\bullet
\otimes_R L^\bullet / F^{j + 1}L^\bullet)
$$
を考える。$F^nT^\bullet$ の定義により、実際に直積ではなく直和に入ることの確認は読者に委ねる。
$a + b = n$ に対して、表示写像を $F^nT^\bullet$ の部分複体
$\text{Tot}(F^aK^\bullet \otimes_R F^bL^\bullet)$ に制限すると、
$i = a$, $j = b$ の直和因子へ写る。
さらに平坦性の仮定により、その像は
$\text{Tot}(\text{gr}^iK^\bullet \otimes_R \text{gr}^jL^\bullet)$ と同型であり、
写像
$\text{Tot}(F^iK^\bullet \otimes_R F^jL^\bullet) \to
\text{Tot}(\text{gr}^iK^\bullet \otimes_R \text{gr}^jL^\bullet)$ の核は
$\text{Tot}(F^{i + 1}K^\bullet \otimes_R F^jL^\bullet) +
\text{Tot}(F^iK^\bullet \otimes_R F^{j + 1}L^\bullet)$ である。
従って、仮定のもとで複体の短完全列
$$
0 \to F^{n + 1}T^\bullet \to F^nT^\bullet \to
\bigoplus\nolimits_{i + j = n}
\text{Tot}(\text{gr}^iK^\bullet \otimes_R \text{gr}^jL^\bullet) \to 0
$$
がある。言い換えると、$\text{gr}^nT^\bullet$ は、
$i + j = n$ にわたる複体
$\text{Tot}(\text{gr}^iK^\bullet \otimes_R \text{gr}^jL^\bullet)$ の直和である。

\medskip\noindent
さらに $R$-加群の複体
$K^\bullet$, $F^iK^\bullet$, $\text{gr}^iK^\bullet$, $L^\bullet$,
$F^jL^\bullet$, $\text{gr}^jL^\bullet$ が K-平坦であると仮定する。
この場合、フィルター付き複体 $T^\bullet$ に付随する
Homology, Section \ref{homology-section-filtered-complex} のスペクトル列は
$$
E_1^{p, q} =
\bigoplus\nolimits_{i + j = p}
H^{p + q}(\text{gr}^iK^\bullet \otimes_R^\mathbf{L} \text{gr}^jL^\bullet)
$$
という項を持ち、その微分は短完全列
$$
0 \to \text{gr}^{n + 1}T^\bullet \to
F^nT^\bullet/F^{n + 2}T^\bullet \to
\text{gr}^nT^\bullet \to 0
$$
から誘導される。
Homology, Lemma
\ref{homology-lemma-spectral-sequence-filtered-complex-d1} の説明を参照せよ。
特に、$E_1^{p, q}$ の直和因子
$H^{p + q}(\text{gr}^iK^\bullet \otimes_R^\mathbf{L} \text{gr}^jL^\bullet)$ は、
$E_1^{p + 1, q}$ の中の直和
$$
H^{p + q + 1}(\text{gr}^{i + 1}K^\bullet \otimes_R^\mathbf{L}
\text{gr}^jL^\bullet) \oplus
H^{p + q + 1}(\text{gr}^iK^\bullet \otimes_R^\mathbf{L}
\text{gr}^{j + 1}L^\bullet)
$$
へ写ることを読者は確認できる。

\begin{lemma}
\label{more-algebra-lemma-kunneth-converges}
上の仮定に加え、次のいずれかを仮定する。
\begin{enumerate}
\item $K^\bullet$ 上のフィルトレーションと $L^\bullet$ 上のフィルトレーションが有限である, または
\item より一般に、次が成り立つ。
\begin{enumerate}
\item $i \gg 0$ に対して $F^iK^\bullet$ は非輪状である,
\item $i \ll 0$ に対して $F^iK^\bullet \to K^\bullet$ は擬同型である,
\item $j \gg 0$ に対して $F^jL^\bullet$ は非輪状である,
\item $j \ll 0$ に対して $F^jL^\bullet \to L^\bullet$ は擬同型である。
\end{enumerate}
\end{enumerate}
このときスペクトル列は有界であり、各
$H^n(T^\bullet) = H^n(K^\bullet \otimes_R^\mathbf{L} L^\bullet)$ 上の
付随フィルトレーションは有限であり、収束
% P02-E0916: frozen source line 15539 uses H^n as the abutment although the spectral-sequence indices require H^{p+q}.
$$
E_1^{p, q} =
\bigoplus\nolimits_{i + j = p}
H^{p + q}(\text{gr}^iK^\bullet \otimes_R^\mathbf{L} \text{gr}^jL^\bullet)
\Rightarrow
H^n(K^\bullet \otimes_R^\mathbf{L} L^\bullet)
$$
がある。
\end{lemma}

\begin{proof}
(1) の場合、$T^\bullet$ 上のフィルトレーションは有限であり、
補題は Homology, Lemma \ref{homology-lemma-biregular-ss-converges} から直ちに従う。
(2) の場合、$i, j > b$ に対して $F^iK^\bullet$, $F^jL^\bullet$ が非輪状であり、
$i, j < a$ に対して $F^iK^\bullet \to K^\bullet$,
$F^jL^\bullet \to L^\bullet$ が擬同型となるような $a < b$ を選ぶ。
このとき、$n > 2b$ に対して複体 $F^nT^\bullet$ は非輪状であり、
$n < 2a - 1$ に対して $F^nT^\bullet \to T^\bullet$ は擬同型であると主張する。
この主張から Homology, Lemma \ref{homology-lemma-ss-converges-trivial} が
適用できることが分かるので、補題が従う。

\medskip\noindent
主張を証明する。$n$ と整数 $i_1 \leq i_2$ が与えられたとき
$$
S^{n, i_1, i_2} =
\Im\left(
\bigoplus\nolimits_{i_1 \leq i \leq i_2}
\text{Tot}(F^iK^\bullet \otimes_R F^{n - i}L^\bullet)
\to
\text{Tot}(K^\bullet \otimes_R L^\bullet)
\right)
$$
を $\text{Tot}(K^\bullet \otimes_R L^\bullet)$ の部分複体とみなす。
$F^nT^\bullet = \colim S^{n, i_1, i_2}$ はフィルター付き余極限であることに注意せよ。
従って主張を示すには、$n > 2b$ に対して
$S^{n, i_1, i_2}$ が非輪状であり、$n < 2a - 1$ かつ
余終的な組 $i_1, i_2$ に対して
$S^{n, i_1, i_2} \to \text{Tot}(K^\bullet \otimes L^\bullet)$ が
擬同型であることを示せば十分である。

\medskip\noindent
$i_1 = i_2 = i$ かつ $n > 2b$ ならば
$$
S^{n, i, i} =
\text{Tot}(F^iK^\bullet \otimes_R F^{n - i}L^\bullet)
$$
は非輪状である。実際、$i > b$ ならば
$F^{n - i}L^\bullet$ が K-平坦であるという仮定によりこの複体は非輪状であり、
$i \leq b$ ならば $n - i \geq n - b > b$ であるから同じことが成り立つ。
平坦性の仮定を用いると、複体の短完全列
$$
0 \to
S^{n + 1, i_2 + 1, i_2 + 1}
\to
S^{n, i_1, i_2} \oplus S^{n, i_2 + 1, i_2 + 1}
\to
S^{n, i_1, i_2 + 1}
\to
0
$$
があることを読者は確認できる。
従って $i_2 - i_1$ に関する帰納法により、$n > 2b$ に対する
すべての複体 $S^{n, i_1, i_2}$ は非輪状である。

\medskip\noindent
$n < 2a - 1$ と仮定する。写像
$$
S^{n, a - 1, a - 1} =
\text{Tot}(F^{a - 1}K^\bullet \otimes_R F^{n - a + 1}L^\bullet)
\to \text{Tot}(K^\bullet \otimes_R L^\bullet)
$$
は擬同型である。実際、$a - 1 < a$ かつ $n - a + 1 < a$ であり、
すべての複体が K-平坦なので、両辺は $D(R)$ の同じ対象を計算する。
$i_2 \geq a - 1$, $n < 2a - 1$ に対して写像
$$
S^{n + 1, i_2 + 1, i_2 + 1} \to S^{n, i_2 + 1, i_2 + 1}
$$
は擬同型であることに注意せよ。
実際、両辺は
$\text{Tot}(F^{i_2 + 1}K^\bullet \otimes L^\bullet)$ と擬同型である。
従って前段落の複体の短完全列を用いると、
$S^{n, a - 1, a - 1} \to S^{n, a - 1, t}$ は
すべての $t \geq a - 1$ に対して擬同型である。
一方、同様に複体の短完全列
$$
0 \to
S^{n + 1, i_1, i_1}
\to
S^{n, i_1 - 1, i_1 - 1} \oplus
S^{n, i_1, i_2}
\to
S^{n, i_1 - 1, i_2}
\to
0
$$
がある。上と同様に、$i_1 \leq a - 1$ に対して写像
$$
S^{n + 1, i_1, i_1} \to S^{n, i_1 - 1, i_1 - 1}
$$
は擬同型であることに注意する。実際、両辺は
$\text{Tot}(K^\bullet \otimes F^{n - i_1 + 1}L^\bullet)$ と擬同型である。
従って、すべての $s \leq a - 1$ に対して
$S^{n, s, t} \to S^{n, a - 1, t}$ は擬同型である。
以上を合わせると、$s \leq a - 1$, $t \geq a - 1$ に対して
$S^{n, s, t}$ は $\text{Tot}(K^\bullet \otimes_R L^\bullet)$ へ擬同型に写る。
これで補題の証明が完了する。
\end{proof}

\begin{lemma}
\label{more-algebra-lemma-resolution-filtered-complex}
$R$ を環とし、$K^\bullet$ をフィルター付き複体とする。
フィルター付き複体の写像 $f : P^\bullet \to K^\bullet$ であって、
次を満たすものが存在する。
\begin{enumerate}
\item 各 $P^n$, $F^iP^n$, $\text{gr}^iP^n$ は自由 $R$-加群である,
\item $R$-加群の複体 $P^\bullet$, $F^iP^\bullet$, $\text{gr}^iP^\bullet$ は K-平坦である,
\item $f$ は擬同型
$P^\bullet \to K^\bullet$,
$F^iP^\bullet \to F^iK^\bullet$,
$\text{gr}^iP^\bullet \to \text{gr}^iK^\bullet$
を誘導する。
\end{enumerate}
\end{lemma}

\begin{proof}
フィルター付き複体 $L^\bullet$ が、各 $L^n$, $F^iL^n$, $\text{gr}^iL^n$ が
自由 $R$-加群であり、すべての微分が零であるとき、{\it 基本的}であるということにする。

\medskip\noindent
基本的フィルター付き複体 $P_0^\bullet$ とフィルター付き複体の写像
$f_0 : P_0^\bullet \to K^\bullet$ であって、$f_0$ および
$F^if_0$, $i \in \mathbf{Z}$ がコホモロジー上全射となるものが存在する。
実際、
$$
P_0^n = \bigoplus\nolimits_{z \in \Ker(d_K^n)} R \xi_z
\oplus
\bigoplus\nolimits_{j \in \mathbf{Z}}
\bigoplus\nolimits_{z \in \Ker(d_{F^jK}^n)} R \xi_z
$$
と置き、微分を零、写像 $f_0$ を $f_0(\xi_z) = z$ で定める。
フィルトレーションについては
% P02-E0917: frozen source line 15679 uses the malformed chained condition j >= i in Z.
$$
F^iP_0^n = \bigoplus\nolimits_{j \geq i \in \mathbf{Z}}
\bigoplus\nolimits_{z \in \Ker(d_{F^jK}^n)} R \xi_z
$$
と置く。これが主張通りの $f_0 : P_0^\bullet \to K^\bullet$ を与えることの確認は読者に委ねる。

\medskip\noindent
$m \geq 0$ に関する帰納法により、フィルター付き複体の埋め込み
% P02-E0918: frozen source line 15688 omits a subset relation between the ellipsis and P_m.
$$
P_0^\bullet \subset \ldots P_m^\bullet \subset P_{m + 1}^\bullet
$$
と写像 $f_m : P_m^\bullet \to K^\bullet$ で、次の性質を持つものを構成する。
\begin{enumerate}
\item フィルター付き複体 $P_{m + 1}^\bullet / P_m^\bullet$ は基本的である,
\item $H^n(f_m)$ と $H^n(F^if_m)$ の核は
$H^n(P_{m + 1}^\bullet)$ と $H^n(F^iP_{m + 1}^\bullet)$ の中で零へ写る,
\item 写像 $f_{m + 1} : P_{m + 1}^\bullet \to K^\bullet$ は写像 $f_m$ を拡張する。
\end{enumerate}
% P02-E0919: frozen source line 15699 reads "To to this".
これを行うため、
$$
\Omega_{n, m} = \Ker\left(\Ker(d^n_{P_m}) \to H^n(K^\bullet)\right),
\text{ resp. }
\Omega_{n, i, m} = \Ker\left(\Ker(d^n_{F^iP_m}) \to H^n(F^iK^\bullet)\right),
$$
と置く。$\Omega_{n, m}$ は $H^n(f_m)$ の核へ全射し、
% P02-E0920: frozen source line 15706 uses f_n although every object in this induction stage is indexed by m.
$\Omega_{n, i, m}$ は $H^n(F^if_n)$ の核へ全射することに注意せよ。
各 $z \in \Omega_{n, m}$, resp.\ $z \in \Omega_{n, i, m}$ に対して、
$d_K(y_z) = f_m(z)$ を満たす
$y_z \in K^{n - 1}$, resp.\ $y_z \in F^iK^{n - 1}$ を選ぶ。次に
$$
P^n_{m + 1} = P_m^n \oplus
\bigoplus\nolimits_{z \in \Omega_{n + 1, m}} R \eta_z
\oplus
\bigoplus\nolimits_{j \in \mathbf{Z}}
\bigoplus\nolimits_{z \in \Omega_{n + 1, j, m}} R \eta_z
$$
と置く。
% P02-E0921: frozen source line 15717 applies d_{P_m} to the newly adjoined eta_z, which is not an element of P_m.
$d_{P_m}(\eta_z) = z$, $f_{m + 1}(\eta_z) = y_z$ と置く。
最後に
$$
F^iP_{m + 1}^n = F^iP_m^n \oplus
\bigoplus\nolimits_{j \geq i}
\bigoplus\nolimits_{z \in \Omega_{n + 1, j, m}} R \eta_z
$$
と置く。これが主張通りの
$P_m^\bullet \subset P_{m + 1}^\bullet$ と
$f_{m + 1} : P_{m + 1}^\bullet \to K^\bullet$ を与えることの確認は読者に委ねる。

\medskip\noindent
ここで単に
$$
P^\bullet = \bigcup P_m^\bullet
$$
をフィルター付き複体とし、写像 $f : P^\bullet \to K^\bullet$ を
$\bigcup f_m$ で与える。

\medskip\noindent
補題の主張 (1) が成り立つのは、フィルター付き加群としての $P^n$ が、
$P_0^n$ と $m \geq 0$ にわたる $P_{m + 1}^n/P_m^n$ の直和に同型だからである。
細部を一つ省略する。

\medskip\noindent
主張 (3) を示す。$H^n(P^\bullet) = \colim H^n(P_m^\bullet)$ であることに注意する。
写像 $f_0$ はコホモロジー上全射であり、従って各 $f_m$ についても同じことが成り立つ。
各 $m$ に対して、構成により埋め込み
$P_m^\bullet \subset P_{m + 1}^\bullet$ は $H^n(f_m)$ の核を消す。
これらを合わせると $H^n(f)$ が同型であることを読者は容易に結論できる。
% P02-E0922: frozen source line 15749 says H^n(F^i f_n) although the limiting map under discussion is f.
$H^n(F^if_n)$ についても同様である。
すると、短完全列
$0 \to F^{i + 1} \to F^i \to \text{gr}^i \to 0$
（例えばフィルター付き複体の圏上の関手の列）により、
$H^n(\text{gr}^if)$ も同型でなければならない。
細部を一つ省略する。

\medskip\noindent
主張 (2) を示す。$P^\bullet$ が K-平坦であることを見るには、
Lemma \ref{more-algebra-lemma-colimit-K-flat} により各 $P_m^\bullet$ が K-平坦であることを示せば十分である。
Lemma \ref{more-algebra-lemma-K-flat-two-out-of-three-ses} と帰納法により、
微分が零で各項が自由である複体が K-平坦であることを注意すれば十分である。
同じ議論により、すべての $i \in \mathbf{Z}$ に対して $F^iP^\bullet$ は K-平坦である。
最後に Lemma \ref{more-algebra-lemma-K-flat-two-out-of-three-ses} をもう一度適用すると、
$\text{gr}^iP^\bullet$ が K-平坦であることが分かる。
\end{proof}

\begin{proposition}
\label{more-algebra-proposition-kunneth-filtered}
$R$ を環とし、$K^\bullet$, $L^\bullet$ を $R$-加群のフィルター付き複体とする。
$D(R)$ において $K^\bullet \otimes_R^\mathbf{L} L^\bullet$ を表す
フィルター付き複体 $T^\bullet$ であって、付随するスペクトル列の $E_1$-項が
$$
E_1^{p, q} =
\bigoplus\nolimits_{i + j = p}
H^{p + q}(\text{gr}^iK^\bullet \otimes_R^\mathbf{L} \text{gr}^jL^\bullet)
$$
となるものが存在する。もし
\begin{enumerate}
\item $i \gg 0$ に対して $F^iK^\bullet$ は非輪状である,
\item $i \ll 0$ に対して $F^iK^\bullet \to K^\bullet$ は擬同型である,
\item $j \gg 0$ に対して $F^jL^\bullet$ は非輪状である,
\item $j \ll 0$ に対して $F^jL^\bullet \to L^\bullet$ は擬同型である
\end{enumerate}
ならば、スペクトル列は有界であり、各
$H^n(K^\bullet \otimes_R^\mathbf{L} L^\bullet)$ 上の付随フィルトレーションは有限であり、
% P02-E0923: frozen source line 15786 uses the noun "convergences" where the verb "converges" is required.
スペクトル列は収束する。
\end{proposition}

\begin{proof}
Lemma \ref{more-algebra-lemma-resolution-filtered-complex} のような
$P^\bullet \to K^\bullet$, $Q^\bullet \to L^\bullet$ を選ぶ。
次に、この節の本文と Lemma \ref{more-algebra-lemma-kunneth-converges} で述べた
フィルター付き複体
$$
T^\bullet = \text{Tot}(P^\bullet \otimes_R Q^\bullet)
$$
のスペクトル列を用いる。
\end{proof}

\begin{lemma}[K\"unneth スペクトル列]
\label{more-algebra-lemma-kunneth-spectral-sequence}
$R$ を環とし、$K$, $L$ を $D^b(R)$ の対象とする。
$$
E_2^{p, q} = \bigoplus\nolimits_{i + j = q} \text{Tor}^R_{-p}(H^i(K), H^j(L))
$$
を持ち、双次数 $(r, -r + 1)$ の $d_r$ を備え、
$H^{p + q}(K \otimes_R^\mathbf{L} L)$ に収束する
双次数付き有界スペクトル列 $\{E_r\}_{r \geq 2}$ が存在する。
\end{lemma}

\begin{proof}
$K$ を表す $R$-加群複体を $K^\bullet$、$L$ を表す $R$-加群複体を
$L^\bullet$ とする。
$F^iK^\bullet = \tau_{\leq -i}K^\bullet$ と置き、$L$ についても同様にする。
Homology, Section \ref{homology-section-truncations} を参照せよ。
$\text{gr}^iK^\bullet = H^{-i}(K)[i]$ であることに注意して
Proposition \ref{more-algebra-proposition-kunneth-filtered} を適用すると、
$$
(E')^{p, q}_1 =
\bigoplus\nolimits_{i + j = p} \text{Tor}^R_{-2p - q}(H^{-i}(K), H^{-j}(L))
$$
を持ち、$K \otimes_R^\mathbf{L} L$ のコホモロジーに収束する
有界スペクトル列 $\{(E')_r\}_{r \geq 1}$ を得る。
構成により、このスペクトル列は $r \geq 1$ に対して微分
$d_r^{p, q} : (E')_r^{p, q} \to (E')_r^{p + r, q - r + 1}$ を持つ。
補題のスペクトル列を得るには、$r \geq 2$ に対して
$$
E_r^{p, q} = (E')_{r - 1}^{-q, p + 2q}
$$
と置く。これでうまくいくことの確認は読者に委ねる。
\end{proof}

\begin{example}
\label{more-algebra-example-kunneth-dedekind}
$R = \mathbf{Z}$、またはより一般に $R$ が Dedekind 整域ならば、
任意の $R$-加群 $M$, $N$ と $i \not \in \{0, 1\}$ に対して
$\text{Tor}_i^R(M, N) = 0$ である。
従って Lemma \ref{more-algebra-lemma-kunneth-spectral-sequence} のスペクトル列は
$E_2$-項で退化し、すべての $n \in \mathbf{Z}$ に対して短完全列
$$
0 \to \bigoplus_{i + j = n} H^i(K) \otimes_R H^j(L) \to
H^n(K \otimes_R^\mathbf{L} L) \to
\bigoplus_{i + j = n + 1} \text{Tor}^R_1(H^i(K), H^j(L)) \to 0
$$
を得る。
\end{example}

\section{擬連接加群 I}
\label{more-algebra-section-pseudo-coherent}

\noindent
$R$ を環とする。$R$-加群 $M$ が有限型であるとは、
全射 $R^{\oplus a} \to M$ が存在すること、有限表示であるとは、表示
$R^{\oplus a_1} \to R^{\oplus a_0} \to M \to 0$ が存在することであった。
同様に、有限自由 $R$-加群による長さ $n$ の分解
\begin{equation}
\label{more-algebra-equation-pseudo-coherent}
R^{\oplus a_n} \to R^{\oplus a_{n - 1}} \to \ldots \to R^{\oplus a_0} \to
M \to 0
\end{equation}
が存在する $R$-加群を考えられる。
すべての $n$ に対してこのような分解を見つけられる加群を{\it 擬連接}という。
形式的な定義は次の通りである。

\begin{definition}
\label{more-algebra-definition-pseudo-coherent}
$R$ を環とする。
% P02-E0924: frozen source line 15875 omits "by" in "Denote D(R) its derived category".
その導来圏を $D(R)$ と書く。$m \in \mathbf{Z}$ とする。
\begin{enumerate}
\item $D(R)$ の対象 $K^\bullet$ が {\it $m$-擬連接}であるとは、
有限自由 $R$-加群からなる有界複体 $E^\bullet$ と射
$\alpha : E^\bullet \to K^\bullet$ が存在し、
$i > m$ に対して $H^i(\alpha)$ が同型、$H^m(\alpha)$ が全射となることをいう。
\item $D(R)$ の対象 $K^\bullet$ が{\it 擬連接}であるとは、
有限自由 $R$-加群からなる上に有界な複体と擬同型であることをいう。
\item $R$-加群 $M$ が {\it $m$-擬連接}であるとは、
$M[0]$ が $D(R)$ の $m$-擬連接対象であることをいう。
\item $R$-加群 $M$ が
{\it 擬連接}\footnote{\cite{Bourbaki-CA} で擬連接加群と呼ばれるものとは衝突する。}
であるとは、$M[0]$ が $D(R)$ の擬連接対象であることをいう。
\end{enumerate}
\end{definition}

\noindent
通常通り、この用語を $R$-加群の複体にも適用する。
$D(R)$ における任意の射 $E^\bullet \to K^\bullet$ は実際の複体写像で表されるので曖昧さはない。
Derived Categories, Lemma \ref{derived-lemma-morphisms-from-projective-complex} を参照せよ。
$K^\bullet$ が擬連接であることと、すべての $m \in \mathbf{Z}$ に対して
$K^\bullet$ が $m$-擬連接であることは同値である。Lemma \ref{more-algebra-lemma-pseudo-coherent} を参照せよ。
また、環が Noether 環ならば、この条件はコホモロジーの有限生成条件として理解できる。
Lemma \ref{more-algebra-lemma-Noetherian-pseudo-coherent} を参照せよ。
まず、これを上の非形式的な議論と関係づけよう。

\begin{lemma}
\label{more-algebra-lemma-cone-pseudo-coherent}
$R$ を環、$m \in \mathbf{Z}$ とし、
$(K^\bullet, L^\bullet, M^\bullet, f, g, h)$ を $D(R)$ の完全三角とする。
\begin{enumerate}
\item $K^\bullet$ が $(m + 1)$-擬連接で $L^\bullet$ が $m$-擬連接ならば、
$M^\bullet$ は $m$-擬連接である。
\item $K^\bullet, M^\bullet$ が $m$-擬連接ならば、
$L^\bullet$ は $m$-擬連接である。
\item $L^\bullet$ が $(m + 1)$-擬連接で $M^\bullet$ が $m$-擬連接ならば、
$K^\bullet$ は $(m + 1)$-擬連接である。
\end{enumerate}
\end{lemma}

\begin{proof}
(1) を証明する。有限自由加群からなる有界複体 $P^\bullet$ と
$\alpha : P^\bullet \to K^\bullet$ で、$i > m + 1$ に対して
$H^i(\alpha)$ が同型、$i = m + 1$ に対して全射となるものを選ぶ。
$P^\bullet$ を $\sigma_{\geq m + 1}P^\bullet$ で置き換えられるので、
$i < m + 1$ に対して $P^i = 0$ と仮定してよい。
有限自由加群からなる有界複体 $E^\bullet$ と
$\beta : E^\bullet \to L^\bullet$ で、$i > m$ に対して
$H^i(\beta)$ が同型、$i = m$ に対して全射となるものを選ぶ。
Derived Categories, Lemma \ref{derived-lemma-lift-map} により、
$D(R)$ において図式
$$
\xymatrix{
K^\bullet \ar[r] & L^\bullet \\
P^\bullet \ar[u] \ar[r]^\gamma & E^\bullet \ar[u]_\beta
}
$$
が可換となる写像 $\gamma : P^\bullet \to E^\bullet$ を見つけられる。
コーン $C(\gamma)^\bullet$ は有限自由 $R$-加群からなる有界複体であり、
図式の可換性から完全三角の射
$$
(P^\bullet, E^\bullet, C(\gamma)^\bullet)
\longrightarrow
(K^\bullet, L^\bullet, M^\bullet).
$$
が存在する。誘導されるコホモロジー長完全列の写像と
Homology, Lemmas \ref{homology-lemma-four-lemma} and
\ref{homology-lemma-five-lemma} により、
$C(\gamma)^\bullet \to M^\bullet$ は次数 $> m$ のコホモロジー上で同型、
次数 $m$ で全射を誘導する。従って $M^\bullet$ は $m$-擬連接である。

\medskip\noindent
主張 (2), (3) は、完全三角を回転して (1) を適用すれば従う。
\end{proof}

\begin{lemma}
\label{more-algebra-lemma-finite-cohomology}
$R$ を環、$K^\bullet$ を $R$-加群の複体、$m \in \mathbf{Z}$ とする。
\begin{enumerate}
\item $K^\bullet$ が $m$-擬連接で、$i > m$ に対して $H^i(K^\bullet) = 0$ ならば、
$H^m(K^\bullet)$ は有限型 $R$-加群である。
\item $K^\bullet$ が $m$-擬連接で、$i > m + 1$ に対して
$H^i(K^\bullet) = 0$ ならば、$H^{m + 1}(K^\bullet)$ は有限表示 $R$-加群である。
\end{enumerate}
\end{lemma}

\begin{proof}
(1) を証明する。有限射影 $R$-加群からなる有界複体 $E^\bullet$ と写像
$\alpha : E^\bullet \to K^\bullet$ で、次数 $> m$ のコホモロジー上で同型、
次数 $m$ で全射を誘導するものを選ぶ。
$E^\bullet$ に対して結果を証明すれば十分であることは明らかである。
$E^n \not = 0$ を満たす最大の整数を $n$ とする。
$n = m$ ならば結果は明らかである。
$n > m$ ならば、$H^n(E^\bullet) = 0$ なので
$E^{n - 1} \to E^n$ は全射である。
$E^n$ は有限射影なので $E^{n - 1} = E' \oplus E^n$ である。
% P02-E0925: frozen source lines 15987--15988 read "the same ... except has".
従って、次数 $n - 1$ に $E'$、次数 $n$ に $0$ を持ち、
それ以外は $E^\bullet$ と同じ複体 $(E')^\bullet$ に対して結果を示せば十分である。
$n$ に関する帰納法で証明が終わる。

\medskip\noindent
(2) を証明する。有限射影 $R$-加群からなる有界複体 $E^\bullet$ と写像
$\alpha : E^\bullet \to K^\bullet$ で、次数 $> m$ のコホモロジー上で同型、
次数 $m$ で全射を誘導するものを選ぶ。
(1) の証明と同様に、$i > m + 1$ に対して $E^i = 0$ である場合に帰着できる。
すると
$H^{m + 1}(K^\bullet) \cong
H^{m + 1}(E^\bullet) = \Coker(E^m \to E^{m + 1})$
は有限表示である。
\end{proof}

\begin{lemma}
\label{more-algebra-lemma-n-pseudo-module}
$R$ を環、$M$ を $R$-加群とする。このとき
\begin{enumerate}
\item $M$ が $0$-擬連接であることと、$M$ が有限 $R$-加群であることは同値である,
\item $M$ が $(-1)$-擬連接であることと、$M$ が有限表示 $R$-加群であることは同値である,
\item $M$ が $(-d)$-擬連接であることと、長さ $d$ の分解
$$
R^{\oplus a_d} \to R^{\oplus a_{d - 1}} \to \ldots \to R^{\oplus a_0} \to
M \to 0
$$
が存在することは同値である,
\item $M$ が擬連接であることと、有限自由 $R$-加群による無限分解
$$
\ldots \to R^{\oplus a_1} \to R^{\oplus a_0} \to M \to 0
$$
が存在することは同値である。
\end{enumerate}
\end{lemma}

\begin{proof}
$M$ が有限型（resp.\ 有限表示）ならば、
Definition \ref{more-algebra-definition-pseudo-coherent} に先立つ議論から、
$M$ は $0$-擬連接（resp.\ $(-1)$-擬連接）である。
逆に $M$ が $0$-擬連接ならば、
Lemma \ref{more-algebra-lemma-finite-cohomology} により
$M = H^0(M[0])$ は有限型である。
$M$ が $(-1)$-擬連接ならば $0$-擬連接でもあり、従って有限型である。
全射 $R^{\oplus a} \to M$ を選び、
$K = \Ker(R^{\oplus a} \to M)$ と書く。
Lemma \ref{more-algebra-lemma-cone-pseudo-coherent} により $K$ は $0$-擬連接であり、
従って有限型なので、$M$ は有限表示である。

\medskip\noindent
% P02-E0926: frozen source line 16043 uses singular "statement" for the third and fourth assertions.
第三、第四の主張には帰納法と上と同様の議論を用いる（詳細は省略する）。
\end{proof}

\begin{lemma}
\label{more-algebra-lemma-pseudo-coherent}
$R$ を環、$K^\bullet$ を $R$-加群の複体とする。次は同値である。
\begin{enumerate}
\item $K^\bullet$ は擬連接である,
\item すべての $m \in \mathbf{Z}$ に対して $K^\bullet$ は $m$-擬連接である,
\item $K^\bullet$ は有限射影 $R$-加群からなる上に有界な複体と擬同型である。
\end{enumerate}
(1), (2), (3) が成り立ち、$i > b$ に対して $H^i(K^\bullet) = 0$ ならば、
$F^i$ が有限自由 $R$-加群で $i > b$ に対して $F^i = 0$ となる擬同型
$F^\bullet \to K^\bullet$ を見つけられる。
\end{lemma}

\begin{proof}
有限自由加群は有限射影 $R$-加群なので、(1) $\Rightarrow$ (3) である。
逆に $P^\bullet$ を有限射影 $R$-加群からなる上に有界な複体とする。
$i > n_0$ に対して $P^i = 0$ としよう。
% P02-E0927: frozen source line 16066 reads "a direct sum decompositions".
$F^{n_0}$ が有限自由 $R$-加群となる直和分解
$F^{n_0} = P^{n_0} \oplus C^{n_0}$ を選び、$n \leq n_0$ に対して帰納的に
$$
F^{n - 1} = P^{n - 1} \oplus C^n \oplus C^{n - 1}
$$
を選ぶ。
% P02-E0928: frozen source line 16072 repeats F^{n_0}, but the inductive decomposition needs F^{n-1} finite free.
ここで $F^{n_0}$ は有限自由 $R$-加群とする。
複体として $F^\bullet$ は写像 $F^{n - 1} \to F^n$ を持ち、
これは $P^{n - 1} \to P^n$ と一致し、$C^n \to C^n$ 上恒等写像を誘導し、
$C^{n - 1}$ 上零である。
写像 $F^\bullet \to P^\bullet$ は擬同型（実際、ホモトピー同値）なので、
(3) は (1) を含意する。

\medskip\noindent
(1) を仮定する。$E^\bullet$ を有限自由 $R$-加群からなる上に有界な複体とし、
$E^\bullet \to K^\bullet$ を擬同型とする。
$E^\bullet$ の愚直切断から $K^\bullet$ への誘導写像
$\sigma_{\geq m}E^\bullet \to K^\bullet$ は、
$K^\bullet$ が $m$-擬連接であることを示す。
従って (1) は (2) を含意する。

\medskip\noindent
(2) を仮定する。$K^\bullet$ は $0$-擬連接なので、$K^\bullet$ は特に上に有界である。
$i > b$ に対して $H^i(K^\bullet) = 0$ となる整数 $b$ を取る。
$n \in \mathbf{Z}$ に関する降下帰納法により、$i \geq n$ に対する有限自由
$R$-加群 $F^i$、$i \geq n$ に対する微分
$d^i : F^i \to F^{i + 1}$、微分と可換な写像
$\alpha : F^i \to K^i$ であって、(1) $i > n$ に対して
$H^i(\alpha)$ は同型、$i = n$ に対して全射であり、
(2) $i > b$ に対して $F^i = 0$ となるものを構成する。図式は
$$
\xymatrix{
& F^n \ar[r] \ar[d]^\alpha & F^{n + 1} \ar[d]^\alpha \ar[r] & \ldots \\
K^{n - 1} \ar[r] & K^n \ar[r] & K^{n + 1} \ar[r] & \ldots
}
$$
である。基底の場合は $n = b + 1$ であり、すべての $i$ に対して $F^i = 0$ と取れる。
帰納段階を考える。$C^\bullet$ を $\alpha$ のコーンとする
（Derived Categories, Definition \ref{derived-definition-cone}）。
コホモロジー長完全列
$$
0 \to H^{n - 1}(K^\bullet) \to H^{n - 1}(C^\bullet) \to
H^n(F^\bullet) \to H^n(K^\bullet) \to H^n(C^\bullet) \to \ldots
$$
から、$i \geq n$ に対して $H^i(C^\bullet) = 0$ である。
Lemma \ref{more-algebra-lemma-cone-pseudo-coherent} により $C^\bullet$ は
$(n - 1)$-擬連接である。
Lemma \ref{more-algebra-lemma-finite-cohomology} により
$H^{n - 1}(C^\bullet)$ は有限 $R$-加群である。
特に $H^n(F^\bullet) \to H^n(K^\bullet)$ の核は有限 $R$-加群である。
有限自由 $R$-加群 $F^{n - 1}$ と写像
% P02-E0929: frozen source line 16115 takes the kernel of F^n -> F^{n-1}, opposite to the cochain differential.
$F^{n - 1} \to \Ker(F^n \to F^{n - 1})$ で、
$F^{n - 1}$ が $H^n(F^\bullet) \to H^n(K^\bullet)$ の核へ全射するものを選ぶ。
% P02-E0930: frozen source line 16117 reads "We extend our of complexes".
複体の写像を
$$
\xymatrix{
& F^{n - 1} \ar[d]^{\alpha^{n - 1}} \ar[r] &
F^n \ar[r] \ar[d]^\alpha &
F^{n + 1} \ar[d]^\alpha \ar[r] & \ldots \\
\ldots \ar[r] &
K^{n - 1} \ar[r] & K^n \ar[r] & K^{n + 1} \ar[r] & \ldots
}
$$
へ拡張する。
$F^{n - 1} \to F^n \to K^n$ の像は構成により
$K^{n - 1} \to K^n$ の像に含まれるので、$\alpha^{n - 1}$ は存在する。
% P02-E0931: frozen source line 16129 begins "Dnote again".
この拡張された複体写像のコーンを再び $C^\bullet$ と書く。
この時点で完全列
$$
H^{n - 1}(F^\bullet) \to H^{n - 1}(K^\bullet) \to H^{n - 1}(C^\bullet) \to
H^n(F^\bullet) \cong H^n(K^\bullet) \to H^n(C^\bullet) \to \ldots
$$
を得る。言い換えると、
$H^{n - 1}(F^\bullet) \to H^{n - 1}(K^\bullet)$ の余核は有限 $R$-加群であり、
$\xi_1, \ldots, \xi_r \in \Ker(K^{n - 1} \to K^n)$ の類で生成されるとしよう。
そこで $F^{n - 1}$ を $F^{n - 1} \oplus R^{\oplus r}$ で置き換える。
% P02-E0932: frozen source lines 16139--16140 read "the basis elements ... map to zero ... and the to xi_i".
自由直和因子の基底元は $F^n$ の中で零へ、$K^{n - 1}$ の中で $\xi_i$ へ写す。
これで帰納段階の証明が完了する。
\end{proof}

\begin{lemma}
\label{more-algebra-lemma-two-out-of-three-pseudo-coherent}
$R$ を環とし、$(K^\bullet, L^\bullet, M^\bullet, f, g, h)$ を
$D(R)$ の完全三角とする。$K^\bullet, L^\bullet, M^\bullet$ のうち
二つが擬連接ならば、残る一つも擬連接である。
\end{lemma}

\begin{proof}
Lemmas \ref{more-algebra-lemma-cone-pseudo-coherent} and \ref{more-algebra-lemma-pseudo-coherent}
を組み合わせればよい。
\end{proof}

\begin{lemma}
\label{more-algebra-lemma-recognize-pseudo-coherent}
$R$ を環、$K^\bullet$ を $R$-加群の複体、$m \in \mathbf{Z}$ とする。
\begin{enumerate}
\item すべての $i \geq m$ に対して $H^i(K^\bullet) = 0$ ならば、
$K^\bullet$ は $m$-擬連接である。
\item $i > m$ に対して $H^i(K^\bullet) = 0$ であり、
$H^m(K^\bullet)$ が有限 $R$-加群ならば、$K^\bullet$ は $m$-擬連接である。
\item $i > m + 1$ に対して $H^i(K^\bullet) = 0$ であり、
$H^{m + 1}(K^\bullet)$ が有限表示、$H^m(K^\bullet)$ が有限型の加群ならば、
$K^\bullet$ は $m$-擬連接である。
\end{enumerate}
\end{lemma}

\begin{proof}
(3) を証明すれば十分である。$M = H^{m + 1}(K^\bullet)$ と置く。
$\tau_{\geq m + 1}K^\bullet$ は $M[- m - 1]$ と擬同型であることに注意する。
Lemma \ref{more-algebra-lemma-n-pseudo-module} により、$M[- m - 1]$ は $m$-擬連接である。
完全三角
$$
(\tau_{\leq m}K^\bullet, K^\bullet, \tau_{\geq m + 1}K^\bullet)
$$
があるので（Derived Categories, Remark
\ref{derived-remark-truncation-distinguished-triangle}）、
Lemma \ref{more-algebra-lemma-cone-pseudo-coherent} により、
% P02-E0933: frozen source line 16186 asks for pseudo-coherence, while the argument proves only m-pseudo-coherence and that is what the cited cone lemma needs.
$\tau_{\leq m}K^\bullet$ が擬連接であることを示せば十分である。
仮定により $H^m(\tau_{\leq m}K^\bullet)$ は有限型 $R$-加群である。
従って有限自由 $R$-加群 $E$ と写像 $E \to \Ker(d_K^m)$ で、その合成
$E \to \Ker(d_K^m) \to H^m(\tau_{\leq m}K^\bullet)$ が全射となるものを取れる。
このとき $E[-m] \to \tau_{\leq m}K^\bullet$ は、
$\tau_{\leq m}K^\bullet$ が $m$-擬連接であることを示す。
\end{proof}

\begin{lemma}
\label{more-algebra-lemma-summands-pseudo-coherent}
$R$ を環、$m \in \mathbf{Z}$ とする。$K^\bullet \oplus L^\bullet$ が
$m$-擬連接（resp.\ 擬連接）ならば、$K^\bullet$ と $L^\bullet$ もそうである。
\end{lemma}

\begin{proof}
この証明では上付き添字 ${}^\bullet$ を省略する。
$K \oplus L$ が $m$-擬連接であると仮定する。
$K, L \in D^{-}(R)$ であることは明らかである。完全三角
$$
(K \oplus L, K \oplus L, L \oplus L[1]) =
(K, K, 0) \oplus (L, L, L \oplus L[1])
$$
があることに注意する。Derived Categories,
Lemma \ref{derived-lemma-direct-sum-triangles} を参照せよ。
Lemma \ref{more-algebra-lemma-cone-pseudo-coherent} により、
$L \oplus L[1]$ は $m$-擬連接である。
従って $L[1] \oplus L[2]$ も $m$-擬連接である。
帰納法により $L[n] \oplus L[n + 1]$ は $m$-擬連接である。
Lemma \ref{more-algebra-lemma-recognize-pseudo-coherent} により、十分大きな $n$ に対して
$L[n]$ は $m$-擬連接である。そこで完全三角
$$
(L[n], L[n] \oplus L[n - 1], L[n - 1])
$$
を用いて逆向きに進めれば、$L[n], L[n - 1], \ldots, L$ が
$m$-擬連接であることが従う。擬連接の場合は、これと
Lemma \ref{more-algebra-lemma-pseudo-coherent} から従う。
\end{proof}

\begin{lemma}
\label{more-algebra-lemma-complex-pseudo-coherent-modules}
$R$ を環、$m \in \mathbf{Z}$ とする。$K^\bullet$ を上に有界な
$R$-加群複体で、すべての $i$ に対して $K^i$ が
$(m - i)$-擬連接であるものとする。このとき $K^\bullet$ は $m$-擬連接である。
特に、$K^\bullet$ が擬連接 $R$-加群からなる上に有界な複体ならば、
$K^\bullet$ は擬連接である。
\end{lemma}

\begin{proof}
$K^\bullet$ を（例えば）$\sigma_{\geq m - 1}K^\bullet$ で置き換えられるので、
$K^\bullet$ は有界であると仮定してよい。
各 $K^i[-i]$ は $m$-擬連接なので、複体の長さに関する帰納法により
$K^\bullet$ は $m$-擬連接である。Lemma \ref{more-algebra-lemma-cone-pseudo-coherent}
と愚直切断を用いればよい。
最後の主張については、Lemma \ref{more-algebra-lemma-pseudo-coherent} により、
すべての $m \in \mathbf{Z}$ に対して $K^\bullet$ が $m$-擬連接であることを
示せば十分であり、これは前半から従う。
\end{proof}

\begin{lemma}
\label{more-algebra-lemma-cohomology-pseudo-coherent}
$R$ を環、$m \in \mathbf{Z}$ とする。$K^\bullet \in D^{-}(R)$ とし、
すべての $i$ に対して $H^i(K^\bullet)$ が
$(m - i)$-擬連接（resp.\ 擬連接）であるとする。
このとき $K^\bullet$ は $m$-擬連接（resp.\ 擬連接）である。
\end{lemma}

\begin{proof}
$K^\bullet$ は $D^{-}(R)$ の対象であり、各 $H^i(K^\bullet)$ が
$(m - i)$-擬連接であると仮定する。
$H^n(K^\bullet)$ が非零となる最大の整数を $n$ とする。
$n$ に関する帰納法で補題を証明する。
$n < m$ ならば、Lemma \ref{more-algebra-lemma-recognize-pseudo-coherent} により
$K^\bullet$ は $m$-擬連接である。
$n \geq m$ ならば完全三角
$$
(\tau_{\leq n - 1}K^\bullet, K^\bullet, H^n(K^\bullet)[-n])
$$
を持つ（Derived Categories, Remark
\ref{derived-remark-truncation-distinguished-triangle}）。
仮定により $H^n(K^\bullet)[-n]$ は $m$-擬連接なので、
Lemma \ref{more-algebra-lemma-cone-pseudo-coherent} を用いれば、
$\tau_{\leq n - 1}K^\bullet$ が $m$-擬連接であることを示せば十分である。
$n$ に関する帰納法で証明が終わる。（擬連接の場合は、これと
Lemma \ref{more-algebra-lemma-pseudo-coherent} から従う。）
\end{proof}

\begin{lemma}
\label{more-algebra-lemma-finite-push-pseudo-coherent}
$A \to B$ を環準同型とし、$B$ は $A$-加群として擬連接であると仮定する。
$K^\bullet$ を $B$-加群の複体とする。次は同値である。
\begin{enumerate}
\item $K^\bullet$ は $B$-加群の複体として $m$-擬連接である。
\item $K^\bullet$ は $A$-加群の複体として $m$-擬連接である。
\end{enumerate}
擬連接についても同じ同値が成り立つ。
\end{lemma}

\begin{proof}
(1) を仮定する。有限自由 $B$-加群からなる有界複体 $E^\bullet$ と写像
$\alpha : E^\bullet \to K^\bullet$ で、次数 $> m$ のコホモロジー上で同型、
次数 $m$ で全射となるものを選ぶ。完全三角
$(E^\bullet, K^\bullet, C(\alpha)^\bullet)$ を考える。
Lemma \ref{more-algebra-lemma-recognize-pseudo-coherent} により、
$C(\alpha)^\bullet$ は $A$-加群の複体として $m$-擬連接である。
従って、$E^\bullet$ が $A$-加群の複体として擬連接であることを示せば十分であり、
これは Lemma \ref{more-algebra-lemma-complex-pseudo-coherent-modules} から従う。
(1) $\Rightarrow$ (2) の擬連接の場合は、これと
Lemma \ref{more-algebra-lemma-pseudo-coherent} から従う。

\medskip\noindent
(2) を仮定する。$H^n(K^\bullet) \not = 0$ となる最大の整数を $n$ とする。
$n - m$ に関する帰納法により、$K^\bullet$ が $B$-加群の複体として
$m$-擬連接であることを示す。$n < m$ の場合は
Lemma \ref{more-algebra-lemma-recognize-pseudo-coherent} から従う。
有限自由 $A$-加群からなる有界複体 $E^\bullet$ と写像
$\alpha : E^\bullet \to K^\bullet$ で、次数 $> m$ のコホモロジー上で同型、
次数 $m$ で全射となるものを選ぶ。誘導された複体写像
$$
\alpha \otimes 1 : E^\bullet \otimes_A B \to K^\bullet.
$$
を考える。
% P02-E0934: frozen source line 16323 uses H^n(E), although only E^\bullet is defined and the adjacent terms use E^\bullet.
$H^n(E) \to H^n(E^\bullet \otimes_A B) \to H^n(K^\bullet)$ は構成により全射であり、
Example \ref{more-algebra-example-tor} のスペクトル列により $i > n$ に対して
$H^i(E^\bullet \otimes_A B) = 0$ なので、$C(\alpha \otimes 1)^\bullet$ は
次数 $\geq n$ で非輪状である。
一方、$K^\bullet$ と $E^\bullet \otimes_A B$ はともに
$A$-加群の複体として $m$-擬連接なので（Lemma
\ref{more-algebra-lemma-complex-pseudo-coherent-modules} を参照）、
Lemma \ref{more-algebra-lemma-cone-pseudo-coherent} により
$C(\alpha \otimes 1)^\bullet$ もそうである。
従って帰納法により、$C(\alpha \otimes 1)^\bullet$ は
$B$-加群の複体として $m$-擬連接である。最後にもう一度
Lemma \ref{more-algebra-lemma-cone-pseudo-coherent} を適用すれば、
$K^\bullet$ は $B$-加群の複体として $m$-擬連接であることが分かる
（$E^\bullet \otimes_A B$ が $B$-加群の複体として擬連接であることは明らかである）。
(2) $\Rightarrow$ (1) の擬連接の場合は、これと
Lemma \ref{more-algebra-lemma-pseudo-coherent} から従う。
\end{proof}

\begin{lemma}
\label{more-algebra-lemma-pull-pseudo-coherent}
$A \to B$ を環準同型とする。$K^\bullet$ を $m$-擬連接（resp.\ 擬連接）な
$A$-加群複体とする。このとき $K^\bullet \otimes_A^{\mathbf{L}} B$ は
$m$-擬連接（resp.\ 擬連接）な $B$-加群複体である。
\end{lemma}

\begin{proof}
まず、$K^\bullet$ は上に有界であり、従って
$K^\bullet \otimes_A^{\mathbf{L}} B$ は Equation
(\ref{more-algebra-equation-derived-tensor-algebra}) によって定義されるので、
補題の主張には意味がある。有限自由 $A$-加群からなる有界複体 $E^\bullet$ と
$\alpha : E^\bullet \to K^\bullet$ で、$H^i(\alpha)$ が $i > m$ に対して同型、
$i = m$ に対して全射となるものを選ぶ。
このときコーン $C(\alpha)^\bullet$ は次数 $\geq m$ で非輪状である。
$-\otimes_A^{\mathbf{L}} B$ は完全関手なので、$B$-加群複体の完全三角
$$
(E^\bullet \otimes_A^{\mathbf{L}} B, K^\bullet \otimes_A^{\mathbf{L}} B,
C(\alpha)^\bullet \otimes_A^{\mathbf{L}} B)
$$
を得る。Derived Categories,
Lemma \ref{derived-lemma-negative-vanishing} の双対により、
$i \geq m$ に対して $H^i(C(\alpha)^\bullet \otimes_A^{\mathbf{L}} B) = 0$
であることが分かる。
% P02-E0935: frozen source line 16369 calls E^\bullet a complex of projective R-modules, but the lemma defines it over A.
$E^\bullet$ は射影 $R$-加群の複体なので、
$E^\bullet \otimes_A^{\mathbf{L}} B = E^\bullet \otimes_A B$ であり、従って
$$
E^\bullet \otimes_A B
\longrightarrow
K^\bullet \otimes_A^{\mathbf{L}} B
$$
は、$K^\bullet \otimes_A^{\mathbf{L}} B$ が $m$-擬連接であることを示す
$B$-加群複体の射である。擬連接複体の場合は、$m$-擬連接複体の場合と
Lemma \ref{more-algebra-lemma-pseudo-coherent} から従う。
\end{proof}

\begin{lemma}
\label{more-algebra-lemma-flat-base-change-pseudo-coherent}
$A \to B$ を平坦な環準同型とする。$M$ を $m$-擬連接（resp.\ 擬連接）な
$A$-加群とする。このとき $M \otimes_A B$ は $m$-擬連接（resp.\ 擬連接）な
$B$-加群である。
\end{lemma}

\begin{proof}
Lemma \ref{more-algebra-lemma-pull-pseudo-coherent} と、$B$ が $A$ 上平坦なので
$M \otimes_A^{\mathbf{L}} B = M \otimes_A B$ であることから直ちに従う。
\end{proof}

\noindent
次の補題は、より強い Lemma \ref{more-algebra-lemma-flat-descent-pseudo-coherent}
からも従う。

\begin{lemma}
\label{more-algebra-lemma-glue-pseudo-coherent}
$R$ を環とし、$f_1, \ldots, f_r \in R$ を単位イデアルを生成する元とする。
$m \in \mathbf{Z}$ とし、$K^\bullet$ を $R$-加群の複体とする。
各 $i$ に対して複体 $K^\bullet \otimes_R R_{f_i}$ が
$m$-擬連接（resp.\ 擬連接）ならば、$K^\bullet$ は
$m$-擬連接（resp.\ 擬連接）である。
\end{lemma}

\begin{proof}
$- \otimes_R R_{f_i}$ が完全関手であり、従って
$$
H^i(K^\bullet)_{f_i} =
H^i(K^\bullet) \otimes_R R_{f_i} = H^i(K^\bullet \otimes_R R_{f_i}).
$$
であることを、以下断りなく用いる。
$i = 1, \ldots, r$ に対して $K^\bullet \otimes_R R_{f_i}$ が
$m$-擬連接であると仮定する。ある $i$ に対して
$H^n(K^\bullet \otimes_R R_{f_i})$ が非零となる最大の整数
$n \in \mathbf{Z}$ を取る。これは特に、$i > n$ に対して
$H^i(K^\bullet) = 0$ であること（および $H^n(K^\bullet) \not = 0$）を含意する。
Algebra, Lemma \ref{algebra-lemma-cover} を参照せよ。
$n - m$ に関する帰納法で補題を証明する。
$n < m$ ならば、Lemma \ref{more-algebra-lemma-recognize-pseudo-coherent} により補題は真である。
$n \geq m$ ならば、各 $i$ に対して $H^n(K^\bullet)_{f_i}$ は有限
$R_{f_i}$-加群である。Lemma \ref{more-algebra-lemma-finite-cohomology} を参照せよ。
従って $H^n(K^\bullet)$ は有限 $R$-加群である。
Algebra, Lemma \ref{algebra-lemma-cover} を参照せよ。
有限自由 $R$-加群 $E$ と全射 $E \to H^n(K^\bullet)$ を選ぶ。
$E$ は射影なので、これを複体の写像 $\alpha : E[-n] \to K^\bullet$ に持ち上げられる。
このときコーン $C(\alpha)^\bullet$ の次数 $\geq n$ のコホモロジーは消える。
一方、各 $i$ に対して複体 $C(\alpha)^\bullet \otimes_R R_{f_i}$ は
$m$-擬連接である。Lemma \ref{more-algebra-lemma-cone-pseudo-coherent} を参照せよ。
従って帰納法により、$C(\alpha)^\bullet$ は $R$-加群の複体として
$m$-擬連接である。Lemma \ref{more-algebra-lemma-cone-pseudo-coherent} をもう一度
適用すれば結論を得る。
\end{proof}

\begin{lemma}
\label{more-algebra-lemma-flat-descent-pseudo-coherent}
$R$ を環、$m \in \mathbf{Z}$ とし、$K^\bullet$ を $R$-加群の複体とする。
$R \to R'$ を忠実平坦な環準同型とする。複体
$K^\bullet \otimes_R R'$ が $m$-擬連接（resp.\ 擬連接）ならば、
$K^\bullet$ は $m$-擬連接（resp.\ 擬連接）である。
\end{lemma}

\begin{proof}
$- \otimes_R R'$ が完全関手であり、従って
$$
H^i(K^\bullet) \otimes_R R' = H^i(K^\bullet \otimes_R R').
$$
であることを、以下断りなく用いる。
$K^\bullet \otimes_R R'$ が $m$-擬連接であると仮定する。
$H^n(K^\bullet)$ が非零となる最大の整数を $n \in \mathbf{Z}$ とする。
このとき $n$ は、$H^n(K^\bullet \otimes_R R')$ が非零となる最大の整数でもある。
$n - m$ に関する帰納法で補題を証明する。
$n < m$ ならば、Lemma \ref{more-algebra-lemma-recognize-pseudo-coherent} により補題は真である。
$n \geq m$ ならば、$H^n(K^\bullet) \otimes_R R'$ は有限 $R'$-加群である。
Lemma \ref{more-algebra-lemma-finite-cohomology} を参照せよ。
従って $H^n(K^\bullet)$ は有限 $R$-加群である。
Algebra, Lemma \ref{algebra-lemma-descend-properties-modules} を参照せよ。
有限自由 $R$-加群 $E$ と全射 $E \to H^n(K^\bullet)$ を選ぶ。
$E$ は射影なので、これを複体の写像 $\alpha : E[-n] \to K^\bullet$ に持ち上げられる。
このときコーン $C(\alpha)^\bullet$ の次数 $\geq n$ のコホモロジーは消える。
一方、複体 $C(\alpha)^\bullet \otimes_R R'$ は $m$-擬連接である。
Lemma \ref{more-algebra-lemma-cone-pseudo-coherent} を参照せよ。
従って帰納法により、$C(\alpha)^\bullet$ は $R$-加群の複体として
$m$-擬連接である。Lemma \ref{more-algebra-lemma-cone-pseudo-coherent} をもう一度
適用すれば結論を得る。
\end{proof}

\begin{lemma}
\label{more-algebra-lemma-tensor-pseudo-coherent}
$R$ を環とし、$K, L$ を $D(R)$ の対象とする。
\begin{enumerate}
\item $K$ が $n$-擬連接で $i > a$ に対して $H^i(K) = 0$、
$L$ が $m$-擬連接で $j > b$ に対して $H^j(L) = 0$ ならば、
$K \otimes_R^\mathbf{L} L$ は $t = \max(m + a, n + b)$ に対して
$t$-擬連接である。
\item $K$ と $L$ が擬連接ならば、$K \otimes_R^\mathbf{L} L$ は擬連接である。
\end{enumerate}
\end{lemma}

\begin{proof}
(1) を証明する。有限自由 $R$-加群からなる有界複体 $K^\bullet$ と
$L^\bullet$、および写像 $\alpha : K^\bullet \to K$ と
$\beta : L^\bullet \to L$ が存在し、$H^i(\alpha)$ が $i > n$ に対して同型、
$i = n$ に対して全射であり、$H^i(\beta)$ が $i > m$ に対して同型、
% P02-E0936: frozen source lines 16502--16503 say "and isomorphism" twice instead of "an isomorphism".
$i = m$ に対して全射であると仮定してよい。このとき写像
$$
\alpha \otimes^\mathbf{L} \beta :
\text{Tot}(K^\bullet \otimes_R L^\bullet)
\to K \otimes_R^\mathbf{L} L
$$
は $i > t$ に対して次数 $i$ のコホモロジー上で同型を、
$i = t$ に対して全射を誘導する。これは Tor のスペクトル列から従う
（詳細は省略する）。(2) は (1) と Lemma \ref{more-algebra-lemma-pseudo-coherent} から従う。
\end{proof}

\begin{lemma}
\label{more-algebra-lemma-Noetherian-pseudo-coherent}
$R$ を Noether 環とする。このとき
\begin{enumerate}
\item $R$-加群の複体 $K^\bullet$ が $m$-擬連接であることと、
$K^\bullet \in D^{-}(R)$ かつ $i \geq m$ に対して $H^i(K^\bullet)$ が
有限 $R$-加群であることは同値である。
\item $R$-加群の複体 $K^\bullet$ が擬連接であることと、
$K^\bullet \in D^{-}(R)$ かつすべての $i$ に対して $H^i(K^\bullet)$ が
有限 $R$-加群であることは同値である。
\item $R$-加群が擬連接であることと有限であることは同値である。
\end{enumerate}
\end{lemma}

\begin{proof}
Algebra, Lemma \ref{algebra-lemma-resolution-by-finite-free} で、任意の有限
$R$-加群が擬連接であることを見た。一方、擬連接加群は有限である。
Lemma \ref{more-algebra-lemma-n-pseudo-module} を参照せよ。従って (3) が成り立つ。
$K^\bullet$ を $m$-擬連接複体とする。このとき有限自由 $R$-加群からなる
有界複体 $E^\bullet$ が存在し、$i > m$ に対して $H^i(K^\bullet)$ は
$H^i(E^\bullet)$ と同型であり、$H^m(K^\bullet)$ は $H^m(E^\bullet)$ の商である。
従って各 $H^i(K^\bullet)$、$i \geq m$ が有限加群であることは明らかである。
(1) の逆向きの含意は Lemma \ref{more-algebra-lemma-cohomology-pseudo-coherent} と (3) から従う。
(2) は (1) と Lemma \ref{more-algebra-lemma-pseudo-coherent} から従う。
\end{proof}

\begin{lemma}
\label{more-algebra-lemma-coherent-pseudo-coherent}
$R$ をコヒーレント環とする
（Algebra, Definition \ref{algebra-definition-coherent}）。
$K \in D^-(R)$ とする。次は同値である。
\begin{enumerate}
\item $K$ は $m$-擬連接である。
\item $H^m(K)$ は有限 $R$-加群であり、$i > m$ に対して $H^i(K)$ はコヒーレントである。
\item $H^m(K)$ は有限 $R$-加群であり、$i > m$ に対して
$H^i(K)$ は有限表示である。
\end{enumerate}
従って、$K$ が擬連接であることと、すべての $i$ に対して
$H^i(K)$ がコヒーレント加群であることは同値である。
\end{lemma}

\begin{proof}
$R$-加群 $M$ がコヒーレントであることと有限表示であることは同値であった
（Algebra, Lemma \ref{algebra-lemma-coherent-ring}）。
これで (2) と (3) の同値性が説明される。その条件が成り立ち、完全列
$0 \to N \to R^{\oplus m} \to M \to 0$ を選ぶと、
Algebra, Lemma \ref{algebra-lemma-coherent} により $N$ はコヒーレントである。
従ってこの場合、$N$ に対してこの手続きを繰り返すと、有限自由 $R$-加群による分解
$$
\ldots \to R^{\oplus n} \to R^{\oplus m} \to M \to 0
$$
を得る。すなわち $M$ は擬連接である。
(1) と (2) の同値性は、これと Lemmas
\ref{more-algebra-lemma-cohomology-pseudo-coherent} and
\ref{more-algebra-lemma-n-pseudo-module} から従う。
最後の主張は、(1) と (2) の同値性を
Lemma \ref{more-algebra-lemma-pseudo-coherent} と組み合わせれば従う。
\end{proof}









\section{擬連接加群 II}
\label{more-algebra-section-pseudo-coherent-bis}

\noindent
Section \ref{more-algebra-section-pseudo-coherent} で始めた議論を続ける。

\begin{lemma}
\label{more-algebra-lemma-pseudo-coherence-colimit-ext}
$R$ を環とし、$M = \colim M_i$ を $R$-加群のフィルター余極限とする。
$K \in D(R)$ は $m$-擬連接であるとする。このとき $n < -m$ に対して
$\colim \Ext^n_R(K, M_i) = \Ext^n_R(K, M)$ であり、
$\colim \Ext^{-m}_R(K, M_i) \to \Ext^{-m}_R(K, M)$ は単射である。
\end{lemma}

\begin{proof}
定義により、$D(R)$ の完全三角
$$
E \to K \to L \to E[1]
$$
で、$E$ が有限自由 $R$-加群からなる有界複体で表され、
$i \geq m$ に対して $H^i(L) = 0$ となるものを見つけられる。
このとき、任意の $R$-加群 $N$ と $n \leq -m$ に対して
$\Ext^n_R(L, N) = 0$ である。Derived Categories,
Lemma \ref{derived-lemma-negative-exts} を参照せよ。
この完全三角に付随する $\Ext$ の長完全列により、
$\Ext^n_R(K, N) \to \Ext^n_R(E, N)$ は $n < -m$ に対して同型、
$n = -m$ に対して単射である。従って、$E$ が有限自由 $R$-加群からなる
有界複体 $E^\bullet$ で表されるとき、$M \mapsto \Ext_R^n(E, M)$ が
フィルター余極限と可換であることを示せば十分である。
加群 $\Ext^n_R(E, M)$ は複体 $\Hom_R(E^\bullet, M)$ によって計算される。
Derived Categories, Lemma
\ref{derived-lemma-morphisms-from-projective-complex} を参照せよ。
$E^p$ は有限自由なので、関手 $M \mapsto \Hom_R(E^p, M)$ は
フィルター余極限と可換である。従って複体として
$\Hom_R(E^\bullet, M) = \colim \Hom_R(E^\bullet, M_i)$ である。
フィルター余極限は完全なので（Algebra, Lemma
\ref{algebra-lemma-directed-colimit-exact}）、結論を得る。
\end{proof}

\begin{lemma}
\label{more-algebra-lemma-characterize-pseudo-coherent-colimit-ext}
$R$ を環、$K \in D^-(R)$、$m \in \mathbf{Z}$ とする。
$K$ が $m$-擬連接であるための必要十分条件は、$R$-加群の任意の
フィルター余極限 $M = \colim M_i$ に対して、$n < -m$ ならば
$\colim \Ext^n_R(K, M_i) = \Ext^n_R(K, M)$ であり、
$\colim \Ext^{-m}_R(K, M_i) \to \Ext^{-m}_R(K, M)$ が単射となることである。
\end{lemma}

\begin{proof}
一方の含意は Lemma \ref{more-algebra-lemma-pseudo-coherence-colimit-ext} で示した。
$R$-加群の任意のフィルター余極限 $M = \colim M_i$ に対して、
$n < -m$ ならば $\colim \Ext^n_R(K, M_i) = \Ext^n_R(K, M)$ であり、
$\colim \Ext^{-m}_R(K, M_i) \to \Ext^{-m}_R(K, M)$ が単射であると仮定する。
$K$ が $m$-擬連接であることを示す。

\medskip\noindent
$H^t(K)$ が非零となる最大の整数を $t$ とし、$t$ に関する帰納法を用いる。
$t < m$ ならば、Lemma \ref{more-algebra-lemma-recognize-pseudo-coherent} により
$K$ は $m$-擬連接である。$t \geq m$ ならば、
$$
\Hom_R(H^t(K), M) = \Ext^{-t}_R(K, M)
$$
なので、任意のフィルター余極限
$M = \colim M_i$ に対して
$$
\colim \Hom_R(H^t(K), M_i) \to \Hom_R(H^t(K), M)
$$
が単射である。
Algebra, Lemma \ref{algebra-lemma-characterize-finite-module-hom} により、
これは $H^t(K)$ が有限 $R$-加群であることを含意する。
有限自由 $R$-加群 $F$ と全射 $F \to H^t(K)$ を選ぶ。
これを $D(R)$ の射 $F[-t] \to K$ に持ち上げ、$D(R)$ の完全三角
$$
F[-t] \to K \to L \to F[-t + 1]
$$
を選べる。このとき $i \geq t$ に対して $H^i(L) = 0$ である。
さらに、この完全三角に付随する $\Ext$ の長完全列から、
$L$ が $K$ に課した仮定を受け継ぐことが、ここでは省略する短い議論で分かる。
$t$ に関する帰納法により $L$ は $m$-擬連接である。
従って Lemma \ref{more-algebra-lemma-cone-pseudo-coherent} により
$K$ は $m$-擬連接である。
\end{proof}

\begin{lemma}
\label{more-algebra-lemma-pseudo-coherence-and-ext}
$R$ を環とし、$L$, $M$, $N$ を $R$-加群とする。
\begin{enumerate}
\item $M$ が有限表示で $L$ が平坦ならば、標準写像
$\Hom_R(M, N) \otimes_R L \to \Hom_R(M, N \otimes_R L)$ は同型である。
\item $M$ が $(-m)$-擬連接で $L$ が平坦ならば、標準写像
$\Ext^i_R(M, N) \otimes_R L \to \Ext^i_R(M, N \otimes_R L)$ は
$i < m$ に対して同型である。
\end{enumerate}
\end{lemma}

\begin{proof}
各項が自由 $R$-加群である分解 $F_\bullet \to M$ を選ぶ。
Algebra, Lemma \ref{algebra-lemma-resolution-by-finite-free} を参照せよ。
複体 $\Hom_R(F_\bullet, N)$ は $\Ext^i_R(M, N)$ を計算し、複体
$\Hom_R(F_\bullet, N \otimes_R L)$ は $\Ext^i_R(M, N \otimes_R L)$ を計算する。
常に余鎖複体の写像
$$
\Hom_R(F_\bullet, N) \otimes_R L
\longrightarrow
\Hom_R(F_\bullet, N \otimes_R L)
$$
があり、すべての $i \geq 0$ に対して標準写像
$\Ext^i_R(M, N) \otimes_R L \to \Ext^i_R(M, N \otimes_R L)$ を誘導する
（例えば、これらの写像が分解 $F_\bullet$ の選択に依存しないという意味で標準的である）。
$L$ が平坦ならば、$L$ とのテンソル積を取る操作はコホモロジーと可換なので、
複体 $\Hom_R(F_\bullet, N) \otimes_R L$ は
$\Ext^i_R(M, N) \otimes_R L$ を計算する。

\medskip\noindent
以上を踏まえ、$M$ が $(-m)$-擬連接ならば、$i = 0, \ldots, m$ に対して
$F_i$ が有限自由となるように $F_\bullet$ を選べる。
すると上に表示した余鎖複体の写像は次数 $\leq m$ で同型であり、
従って次数 $< m$ のコホモロジー群上で同型である。これで (2) が示された。
$M$ が有限表示ならば、Lemma \ref{more-algebra-lemma-n-pseudo-module} により
$M$ は $(-1)$-擬連接であり、$\Hom = \Ext^0$ なので結果を得る。
\end{proof}

\begin{lemma}
\label{more-algebra-lemma-pseudo-coherence-and-base-change-ext}
$R \to R'$ を平坦な環準同型とし、$M$, $N$ を $R$-加群とする。
\begin{enumerate}
\item $M$ が有限表示 $R$-加群ならば、
$\Hom_R(M, N) \otimes_R R' = \Hom_{R'}(M \otimes_R R', N \otimes_R R')$ である。
\item $M$ が $(-m)$-擬連接ならば、$i < m$ に対して
$\Ext^i_R(M, N) \otimes_R R' = \Ext^i_{R'}(M \otimes_R R', N \otimes_R R')$ である。
\end{enumerate}
特に $R$ が Noether 環で $M$ が有限加群ならば、これはすべての $i$ に対して成り立つ。
\end{lemma}

\begin{proof}
Algebra, Lemma \ref{algebra-lemma-flat-base-change-ext} により
$\Ext^i_{R'}(M \otimes_R R', N \otimes_R R') =
\Ext^i_R(M, N \otimes_R R')$ である。
% P02-E0937: frozen source line 16727 says "(1) and (2) holds" although the subject is plural.
Lemma \ref{more-algebra-lemma-pseudo-coherence-and-ext} と組み合わせれば (1), (2) が従う。
最後の主張は、これと Lemma \ref{more-algebra-lemma-Noetherian-pseudo-coherent} から従う。
\end{proof}

\begin{lemma}
\label{more-algebra-lemma-pseudo-coherent-tensor}
$R$ を環、$K \in D^-(R)$ とする。次は同値である。
\begin{enumerate}
\item $K$ は擬連接である。
\item $R$-加群の任意の族 $(Q_{\alpha})_{\alpha \in A}$ に対して、標準写像
$$
\alpha :
K \otimes_R^\mathbf{L} \left( \prod\nolimits_\alpha Q_{\alpha} \right)
\longrightarrow
\prod\nolimits_\alpha (K \otimes_R^\mathbf{L} Q_{\alpha})
$$
は $D(R)$ における同型である。
\item 任意の $R$-加群 $Q$ と任意の集合 $A$ に対して、標準写像
$$
\beta : K \otimes_R^\mathbf{L} Q^A \longrightarrow (K \otimes_R^\mathbf{L} Q)^A
$$
は $D(R)$ における同型である。
\item 任意の集合 $A$ に対して、標準写像
$$
\gamma : K \otimes_R^\mathbf{L} R^A \longrightarrow K^A
$$
は $D(R)$ における同型である。
\end{enumerate}
$m \in \mathbf{Z}$ を与えると、次は同値である。
\begin{enumerate}
\item[(a)] $K$ は $m$-擬連接である。
\item[(b)] $R$-加群の任意の族 $(Q_{\alpha})_{\alpha \in A}$ に対し、
上の $\alpha$ について、$H^i(\alpha)$ は $i > m$ に対して同型、
$i = m$ に対して全射である。
\item[(c)] 任意の $R$-加群 $Q$ と任意の集合 $A$ に対し、上の $\beta$ について、
$H^i(\beta)$ は $i > m$ に対して同型、$i = m$ に対して全射である。
\item[(d)] 任意の集合 $A$ に対し、上の $\gamma$ について、
$H^i(\gamma)$ は $i > m$ に対して同型、$i = m$ に対して全射である。
\end{enumerate}
\end{lemma}

\begin{proof}
$K$ が擬連接ならば、$K$ は有限自由 $R$-加群からなる上に有界な複体で表せる。
従って導来テンソル積は、この複体とのテンソル積を取ることで計算される。
また、$D(R)$ の積は、代表複体の任意の選択について積を取ることで与えられる。
従って加群についての対応する事実、Algebra, Proposition
\ref{algebra-proposition-fp-tensor} により、(1) は (2), (3), (4) を含意する。

\medskip\noindent
% P02-E0939: frozen source line 16779 says "using the tensor product is right exact", omitting "that" or "the fact that".
同じ仕方で（テンソル積が右完全であることを用いて）、(a) が (b), (c), (d) を
含意することを読者は示せる。

\medskip\noindent
(4) が成り立つと仮定する。$K$ が擬連接であることを示すには、
すべての $m$ に対して $K$ が $m$-擬連接であることを示せば十分である
（Lemma \ref{more-algebra-lemma-pseudo-coherent}）。
% P02-E0938: frozen source line 16785 reads "finish then proof" instead of "finish the proof".
従って証明を終えるには、(d) が (a) を含意することを証明すれば十分である。

\medskip\noindent
(d) を仮定する。$H^i(K)$ が非零となる最大の整数を $i$ とする。
$i < m$ ならば証明は終わる。そうでなければ、(d) と上で述べた
$D(R)$ における積の記述から、$H^i(K) \otimes_R R^A \to H^i(K)^A$ は全射である。
従って Algebra, Proposition \ref{algebra-proposition-fg-tensor} により、
$H^i(K)$ は有限生成 $R$-加群である。
そこで次数 $i$ に置かれた単一の有限自由加群からなる複体 $L$ と複体写像
$L \to K$ で、$H^i(L) \to H^i(K)$ が全射となるものを選べる。
特に $L$ は (1), (2), (3), (4) を満たす。完全三角
$$
L \to K \to M \to L[1]
$$
を選ぶ。このとき $j \geq i$ に対して $H^j(M) = 0$ である。
一方、ここでは省略する短い議論により、$M$ はなお性質 (d) を持つ。
$i$ に関する帰納法により $M$ は $m$-擬連接である。
従って Lemma \ref{more-algebra-lemma-cone-pseudo-coherent} により
$K$ は $m$-擬連接である。
\end{proof}

\begin{lemma}
\label{more-algebra-lemma-detect-cohomology-pseudo-coherent}
$R$ を環とし、$K \in D(R)$ は擬連接であるとする。
$i \in \mathbf{Z}$ とする。このとき有限表示 $R$-加群 $M$ と
$D(R)$ の写像 $K \to M[-i]$ で、単射 $H^i(K) \to M$ を誘導するものが存在する。
\end{lemma}

\begin{proof}
Definition \ref{more-algebra-definition-pseudo-coherent} により、$K$ を有限自由 $R$-加群からなる
複体 $P^\bullet$ で表せる。$M = \Coker(P^{i - 1} \to P^i)$ と置く。
\end{proof}

\begin{lemma}
\label{more-algebra-lemma-detect-cohomology}
$A$ を Noether 環とし、$K \in D(A)$ は擬連接、すなわち
$K \in D^-(A)$ でそのコホモロジー加群は有限であるとする。
$\mathfrak m$ を $A$ の極大イデアルとする。
$H^i(K)/\mathfrak m H^i(K) \not = 0$ ならば、$\mathfrak m$ のある冪で零化される
有限 $A$-加群 $E$ と写像 $K \to E[-i]$ で、$H^i(K)$ 上非零となるものが存在する。
\end{lemma}

\begin{proof}
（主張にある擬連接性の同値な定式化は
Lemma \ref{more-algebra-lemma-Noetherian-pseudo-coherent} である。）
Lemma \ref{more-algebra-lemma-detect-cohomology-pseudo-coherent} のように $K \to M[-i]$ を選ぶ。
Artin--Rees（Algebra, Lemma \ref{algebra-lemma-Artin-Rees}）により、
$H^i(K) \cap \mathfrak m^n M \subset \mathfrak m H^i(K)$ となる $n$ を見つけられる。
$E = M/\mathfrak m^n M$ と取る。
\end{proof}









\section{Tor 次元}
\label{more-algebra-section-tor}

\noindent
射影加群による分解の代わりに、平坦加群による分解を考えられる。
これは次の概念へ導く。

\begin{definition}
\label{more-algebra-definition-tor-amplitude}
$R$ を環とする。
% P02-E0940: frozen source line 16862 repeats the malformed "Denote D(R) its derived category" construction recorded at P02-E0924 for line 15875.
その導来圏を $D(R)$ と書く。$a, b \in \mathbf{Z}$ とする。
\begin{enumerate}
\item $D(R)$ の対象 $K^\bullet$ が {it $[a, b]$ に Tor 振幅を持つ}とは、
すべての $R$-加群 $M$ とすべての $i \not \in [a, b]$ に対して
$H^i(K^\bullet \otimes_R^\mathbf{L} M) = 0$ となることをいう。
\item $D(R)$ の対象 $K^\bullet$ が{it 有限 Tor 次元を持つ}とは、
ある $a, b$ に対して $[a, b]$ に Tor 振幅を持つことをいう。
\item $R$-加群 $M$ が {it Tor 次元 $\leq d$ を持つ}とは、
$D(R)$ の対象としての $M[0]$ が $[-d, 0]$ に Tor 振幅を持つことをいう。
\item $R$-加群 $M$ が{it 有限 Tor 次元を持つ}とは、
$D(R)$ の対象としての $M[0]$ が有限 Tor 次元を持つことをいう。
\end{enumerate}
\end{definition}

\noindent
$K^\bullet$ が有限 Tor 次元を持つならば $K^\bullet \in D^b(R)$ であることに注意する。

\begin{lemma}
\label{more-algebra-lemma-last-one-flat}
$R$ を環とし、$K^\bullet$ を $[a, b]$ に Tor 振幅を持つ平坦
$R$-加群の上に有界な複体とする。このとき
$\Coker(d_K^{a - 1})$ は平坦 $R$-加群である。
\end{lemma}

\begin{proof}
$K^\bullet$ は平坦加群の上に有界な複体なので、
$K^\bullet \otimes_R M = K^\bullet \otimes_R^{\mathbf{L}} M$ である。
従って任意の $R$-加群 $M$ に対して列
$$
K^{a - 2} \otimes_R M \to K^{a - 1} \otimes_R M \to K^a \otimes_R M
$$
は中央で完全である。
$K^{a - 2} \to K^{a - 1} \to K^a \to \Coker(d_K^{a - 1}) \to 0$
は平坦分解なので、これはすべての $R$-加群 $M$ に対して
$\text{Tor}_1^R(\Coker(d_K^{a - 1}), M) = 0$ であることを含意する。
従って $\Coker(d_K^{a - 1})$ は平坦である。
Algebra, Lemma \ref{algebra-lemma-characterize-flat} を参照せよ。
\end{proof}

\begin{lemma}
\label{more-algebra-lemma-tor-amplitude}
$R$ を環、$K^\bullet$ を $D(R)$ の対象とし、$a, b \in \mathbf{Z}$ とする。
次は同値である。
\begin{enumerate}
\item $K^\bullet$ は $[a, b]$ に Tor 振幅を持つ。
\item $K^\bullet$ は、$i \not \in [a, b]$ に対して $E^i = 0$ となる
平坦 $R$-加群の複体 $E^\bullet$ と擬同型である。
\end{enumerate}
\end{lemma}

\begin{proof}
(2) が成り立つならば、
$K^\bullet \otimes_R^\mathbf{L} M = E^\bullet \otimes_R M$ と計算でき、
(1) が成り立つことは明らかである。(1) を仮定する。
$i > b$ に対して $K^i = 0$ となる射影分解で $K^\bullet$ を置き換えてよい。
Derived Categories, Lemma \ref{derived-lemma-projective-resolutions-exist} を参照せよ。
$E^\bullet = \tau_{\geq a}K^\bullet$ と置く。
$E^a$ が平坦であること以外はすべて明らかであり、それは
Lemma \ref{more-algebra-lemma-last-one-flat} と定義から直ちに従う。
\end{proof}

\begin{lemma}
\label{more-algebra-lemma-bounded-below-tor-amplitude}
$R$ を環、$a \in \mathbf{Z}$ とし、$K$ を $D(R)$ の対象とする。
次は同値である。
\begin{enumerate}
\item $K$ は $[a, \infty]$ に Tor 振幅を持つ。
\item $K$ は、各項が平坦 $R$-加群で $i \not \in [a, \infty]$ に対して
$E^i = 0$ となる K-平坦複体 $E^\bullet$ と擬同型である。
\end{enumerate}
\end{lemma}

\begin{proof}
(2) $\Rightarrow$ (1) は直ちに従う。(1) を仮定する。
まず、$K$ を表す平坦な項を持つ K-平坦複体 $K^\bullet$ を選ぶ。
Lemma \ref{more-algebra-lemma-K-flat-resolution} を参照せよ。
任意の $R$-加群 $M$ に対して
$$
K^{n - 1} \otimes_R M \to K^n \otimes_R M \to K^{n + 1} \otimes_R M
$$
のコホモロジーは $H^n(K \otimes_R^\mathbf{L} M)$ を計算する。
これは $n < a$ に対して常に零である。従って複体
$\ldots \to K^{a - 1} \to K^a \to K^{a + 1}$ に
Lemma \ref{more-algebra-lemma-last-one-flat} を適用すれば、
$N = \Coker(K^{a - 1} \to K^a)$ が平坦 $R$-加群であることが分かる。
$$
E^\bullet = \tau_{\geq a}K^\bullet =
(\ldots \to 0 \to N \to K^{a + 1} \to \ldots )
$$
と置く。$K^\bullet \to E^\bullet$ の核 $L^\bullet$ は複体
$$
L^\bullet = (\ldots \to K^{a - 1} \to I \to 0 \to \ldots)
$$
であり、ここで $I \subset K^a$ は $K^{a - 1} \to K^a$ の像である。
短完全列 $0 \to I \to K^a \to N \to 0$ があるので、$I$ は平坦 $R$-加群である。
従って $L^\bullet$ は平坦加群からなる上に有界な複体であり、
Lemma \ref{more-algebra-lemma-derived-tor-quasi-isomorphism} により K-平坦である。
すると Lemma \ref{more-algebra-lemma-K-flat-two-out-of-three-ses} により
$E^\bullet$ は K-平坦である。
\end{proof}

\begin{lemma}
\label{more-algebra-lemma-cone-tor-amplitude}
$R$ を環とし、$(K^\bullet, L^\bullet, M^\bullet, f, g, h)$ を
$D(R)$ の完全三角とする。$a, b \in \mathbf{Z}$ とする。
\begin{enumerate}
\item $K^\bullet$ が $[a + 1, b + 1]$ に Tor 振幅を持ち、
$L^\bullet$ が $[a, b]$ に Tor 振幅を持つならば、
$M^\bullet$ は $[a, b]$ に Tor 振幅を持つ。
\item $K^\bullet, M^\bullet$ が $[a, b]$ に Tor 振幅を持つならば、
$L^\bullet$ は $[a, b]$ に Tor 振幅を持つ。
\item $L^\bullet$ が $[a + 1, b + 1]$ に Tor 振幅を持ち、
$M^\bullet$ が $[a, b]$ に Tor 振幅を持つならば、
$K^\bullet$ は $[a + 1, b + 1]$ に Tor 振幅を持つ。
\end{enumerate}
\end{lemma}

\begin{proof}
省略する。ヒント：完全三角に付随するコホモロジー長完全列と、
$- \otimes_R^{\mathbf{L}} M$ が完全三角を保つことから従うだけである。
最も証明しやすいのは (2) であり、他は移動によってこれから従う。
\end{proof}

\begin{lemma}
\label{more-algebra-lemma-tor-dimension}
$R$ を環、$M$ を $R$-加群とし、$d \geq 0$ とする。次は同値である。
\begin{enumerate}
\item $M$ は Tor 次元 $\leq d$ を持つ。
\item $F_i$ が平坦 $R$-加群である分解
$$
0 \to F_d \to \ldots \to F_1 \to F_0 \to M \to 0
$$
が存在する。
\end{enumerate}
特に、$R$-加群が Tor 次元 $0$ を持つことと平坦 $R$-加群であることは同値である。
\end{lemma}

\begin{proof}
(2) を仮定する。$E^{-i} = F_i$ を持つ複体 $E^\bullet$ は $M$ と擬同型である。
従って Lemma \ref{more-algebra-lemma-tor-amplitude} により $M$ の Tor 次元は高々 $d$ である。
逆に (1) を仮定する。$P^\bullet \to M$ を $M$ の射影分解とする。
Lemma \ref{more-algebra-lemma-last-one-flat} により $\tau_{\geq -d}P^\bullet$ は
$M$ の長さ $d$ の平坦分解である。すなわち (2) が成り立つ。
\end{proof}

\begin{lemma}
\label{more-algebra-lemma-summands-tor-amplitude}
$R$ を環、$a, b \in \mathbf{Z}$ とする。
$K^\bullet \oplus L^\bullet$ が $[a, b]$ に Tor 振幅を持つならば、
$K^\bullet$ と $L^\bullet$ もそうである。
\end{lemma}

\begin{proof}
Tor 関手が加法的であることから明らかである。
\end{proof}

\begin{lemma}
\label{more-algebra-lemma-complex-finite-tor-dimension-modules}
$R$ を環とし、$K^\bullet$ を有界な $R$-加群複体で、すべての $i$ に対して
$K^i$ が $[a - i, b - i]$ に Tor 振幅を持つものとする。
このとき $K^\bullet$ は $[a, b]$ に Tor 振幅を持つ。
特に、$K^\bullet$ が有限 Tor 次元を持つ $R$-加群からなる有限複体ならば、
$K^\bullet$ は有限 Tor 次元を持つ。
\end{lemma}

\begin{proof}
有限複体の長さに関する帰納法から従う。
Lemma \ref{more-algebra-lemma-cone-tor-amplitude} と愚直切断を用いればよい。
\end{proof}

\begin{lemma}
\label{more-algebra-lemma-cohomology-tor-amplitude}
$R$ を環、$a, b \in \mathbf{Z}$ とする。$K^\bullet \in D^b(R)$ とし、
すべての $i$ に対して $H^i(K^\bullet)$ が $[a - i, b - i]$ に
Tor 振幅を持つとする。このとき $K^\bullet$ は $[a, b]$ に Tor 振幅を持つ。
特に、$K^\bullet \in D^b(R)$ でそのすべてのコホモロジー群が有限 Tor 次元を
持つならば、$K^\bullet$ は有限 Tor 次元を持つ。
\end{lemma}

\begin{proof}
有限複体の長さに関する帰納法から従う。
Lemma \ref{more-algebra-lemma-cone-tor-amplitude} と標準切断を用いればよい。
\end{proof}

\begin{lemma}
\label{more-algebra-lemma-push-tor-amplitude}
$A \to B$ を環準同型とし、$K^\bullet$ と $L^\bullet$ を $B$-加群の複体とする。
$a, b, c, d \in \mathbf{Z}$ とする。次を仮定する。
\begin{enumerate}
\item $K^\bullet$ は $B$-加群の複体として $[a, b]$ に Tor 振幅を持つ。
\item $L^\bullet$ は $A$-加群の複体として $[c, d]$ に Tor 振幅を持つ。
\end{enumerate}
このとき $K^\bullet \otimes^\mathbf{L}_B L^\bullet$ は $A$-加群の複体として
$[a + c, b + d]$ に Tor 振幅を持つ。
\end{lemma}

\begin{proof}
Lemma \ref{more-algebra-lemma-tor-amplitude} により、$K^\bullet$ は $i \not \in [a, b]$ に対して
$K^i = 0$ となる平坦 $B$-加群の複体であると仮定してよい。
$M$ を $A$-加群とし、自由分解 $F^\bullet \to M$ を選ぶ。このとき
$$
(K^\bullet \otimes_B^\mathbf{L} L^\bullet) \otimes_A^{\mathbf{L}} M =
\text{Tot}(\text{Tot}(K^\bullet \otimes_B L^\bullet) \otimes_A F^\bullet) =
\text{Tot}(K^\bullet \otimes_B \text{Tot}(L^\bullet \otimes_A F^\bullet))
$$
である。第二の等式については Homology, Remark
\ref{homology-remark-triple-complex} を参照せよ。
仮定 (2) により、複体 $\text{Tot}(L^\bullet \otimes_A F^\bullet)$ の
コホモロジーが非零となり得る次数は $[c, d]$ のみである。
従って二重複体 $K^\bullet \otimes_B \text{Tot}(L^\bullet \otimes_A F^\bullet)$
に対する Homology, Lemma \ref{homology-lemma-ss-double-complex} の
スペクトル列により、$(K^\bullet \otimes_B^\mathbf{L} L^\bullet) \otimes_A^{\mathbf{L}} M$
のコホモロジーが非零となり得る次数は $[a + c, b + d]$ のみである。
\end{proof}

\begin{lemma}
\label{more-algebra-lemma-flat-push-tor-amplitude}
$A \to B$ を環準同型とし、$B$ は $A$-加群として平坦であると仮定する。
$K^\bullet$ を $B$-加群の複体とし、$a, b \in \mathbf{Z}$ とする。
$K^\bullet$ が $B$-加群の複体として $[a, b]$ に Tor 振幅を持つならば、
$K^\bullet$ は $A$-加群の複体としても $[a, b]$ に Tor 振幅を持つ。
\end{lemma}

\begin{proof}
これは Lemma \ref{more-algebra-lemma-push-tor-amplitude} の特別な場合だが、次のように直接にも分かる。
$$K^\bullet \otimes_A^{\mathbf{L}} M =
K^\bullet \otimes_B^{\mathbf{L}} (M \otimes_A B)$$
である。実際、$B$-加群の複体としての $K^\bullet$ の任意の射影分解は、
$A$-加群の複体としての $K^\bullet$ の平坦分解であり、
$K^\bullet \otimes_A^{\mathbf{L}} M$ の計算に用いられる。
\end{proof}

\begin{lemma}
\label{more-algebra-lemma-finite-tor-dimension-push-tor-amplitude}
$A \to B$ を環準同型とし、$B$ は $A$-加群として Tor 次元 $\leq d$ を持つと仮定する。
$K^\bullet$ を $B$-加群の複体とし、$a, b \in \mathbf{Z}$ とする。
$K^\bullet$ が $B$-加群の複体として $[a, b]$ に Tor 振幅を持つならば、
$K^\bullet$ は $A$-加群の複体として $[a - d, b]$ に Tor 振幅を持つ。
\end{lemma}

\begin{proof}
これは Lemma \ref{more-algebra-lemma-push-tor-amplitude} の特別な場合だが、次のように直接にも分かる。
$M$ を $A$-加群とし、自由分解 $F^\bullet \to M$ を選ぶ。このとき
$$
K^\bullet \otimes_A^{\mathbf{L}} M =
\text{Tot}(K^\bullet \otimes_A F^\bullet) =
\text{Tot}(K^\bullet \otimes_B (F^\bullet \otimes_A B)) =
K^\bullet \otimes_B^{\mathbf{L}} (M \otimes_A^{\mathbf{L}} B).
$$
$A$-加群としての $B$ に関する仮定により、
$M \otimes_A^{\mathbf{L}} B$ のコホモロジーが非零となり得る次数は
$-d, -d + 1, \ldots, 0$ のみである。
$K^\bullet$ は $[a, b]$ に Tor 振幅を持つので、
Example \ref{more-algebra-example-tor} のスペクトル列から、
$K^\bullet \otimes_B^{\mathbf{L}} (M \otimes_A^{\mathbf{L}} B)$ の
コホモロジーが非零となり得る次数は $[-d + a, b]$ のみである。
\end{proof}

\begin{lemma}
\label{more-algebra-lemma-pull-tor-amplitude}
$A \to B$ を環準同型とし、$a, b \in \mathbf{Z}$ とする。
$K^\bullet$ を $[a, b]$ に Tor 振幅を持つ $A$-加群の複体とする。
このとき $K^\bullet \otimes_A^{\mathbf{L}} B$ は $B$-加群の複体として
$[a, b]$ に Tor 振幅を持つ。
\end{lemma}

\begin{proof}
Lemma \ref{more-algebra-lemma-tor-amplitude} により、擬同型 $E^\bullet \to K^\bullet$ で、
$E^\bullet$ が平坦 $A$-加群の複体であり、$i \not \in [a, b]$ に対して
$E^i = 0$ となるものを見つけられる。このとき構成により
$E^\bullet \otimes_A B$ は $K^\bullet \otimes_A ^{\mathbf{L}} B$ を計算し、
Algebra, Lemma \ref{algebra-lemma-flat-base-change} により
各 $E^i \otimes_A B$ は平坦 $B$-加群である。
従って Lemma \ref{more-algebra-lemma-tor-amplitude} により結論を得る。
\end{proof}

\begin{lemma}
\label{more-algebra-lemma-flat-base-change-finite-tor-dimension}
$A \to B$ を平坦な環準同型とし、$d \geq 0$ とする。
$M$ を Tor 次元 $\leq d$ を持つ $A$-加群とする。
このとき $M \otimes_A B$ は Tor 次元 $\leq d$ を持つ $B$-加群である。
\end{lemma}

\begin{proof}
Lemma \ref{more-algebra-lemma-pull-tor-amplitude} と、$B$ が $A$ 上平坦なので
$M \otimes_A^{\mathbf{L}} B = M \otimes_A B$ であることから直ちに従う。
\end{proof}

\begin{lemma}
\label{more-algebra-lemma-tor-amplitude-localization}
$A  \to B$ を環準同型とし、$K^\bullet$ を $B$-加群の複体、
$a, b \in \mathbf{Z}$ とする。次は同値である。
\begin{enumerate}
\item $K^\bullet$ は $A$-加群の複体として $[a, b]$ に Tor 振幅を持つ。
\item $\mathfrak p = A \cap \mathfrak q$ を持つすべての素イデアル
$\mathfrak q \subset B$ に対して、$K^\bullet_\mathfrak q$ は
$A_\mathfrak p$-加群の複体として $[a, b]$ に Tor 振幅を持つ。
\item $\mathfrak p = A \cap \mathfrak m$ を持つすべての極大イデアル
$\mathfrak m \subset B$ に対して、$K^\bullet_\mathfrak m$ は
$A_\mathfrak p$-加群の複体として $[a, b]$ に Tor 振幅を持つ。
\end{enumerate}
\end{lemma}

\begin{proof}
(3) を仮定し、$M$ を $A$-加群とする。このとき
$H^i = H^i(K^\bullet \otimes_A^\mathbf{L} M)$ は $B$-加群であり、
$(H^i)_\mathfrak m =
H^i(K^\bullet_\mathfrak m \otimes_{A_\mathfrak p}^\mathbf{L} M_\mathfrak p)$ である。
従って Algebra, Lemma \ref{algebra-lemma-characterize-zero-local} により、
$i \not \in [a, b]$ に対して $H^i = 0$ である。
これで (3) $\Rightarrow$ (1) が示された。
(1) $\Rightarrow$ (2) と (2) $\Rightarrow$ (3) の証明は省略する。
\end{proof}

\begin{lemma}
\label{more-algebra-lemma-glue-tor-amplitude}
$R$ を環とし、$f_1, \ldots, f_r \in R$ を単位イデアルを生成する元とする。
$a, b \in \mathbf{Z}$ とし、$K^\bullet$ を $R$-加群の複体とする。
各 $i$ に対して複体 $K^\bullet \otimes_R R_{f_i}$ が $[a, b]$ に
Tor 振幅を持つならば、$K^\bullet$ は $[a, b]$ に Tor 振幅を持つ。
\end{lemma}

\begin{proof}
これは Lemma \ref{more-algebra-lemma-tor-amplitude-localization} から直ちに従うが、
次のように直接にも分かる。$- \otimes_R R_{f_i}$ は完全関手なので
$$
H^i(K^\bullet)_{f_i} =
H^i(K^\bullet) \otimes_R R_{f_i} = H^i(K^\bullet \otimes_R R_{f_i}).
$$
であることに注意する。同様に、任意の $R$-加群 $M$ に対して
$$
H^i(K^\bullet \otimes_R^{\mathbf{L}} M)_{f_i} =
H^i(K^\bullet \otimes_R^{\mathbf{L}} M) \otimes_R R_{f_i} =
H^i(K^\bullet \otimes_R R_{f_i} \otimes_{R_{f_i}}^{\mathbf{L}} M_{f_i}).
$$
である。従って、$R$-加群 $N$ が零であることと、各 $i$ に対して
$N_{f_i}$ が零であることが同値であるという事実から結果が従う。
Algebra, Lemma \ref{algebra-lemma-cover} を参照せよ。
\end{proof}

\begin{lemma}
\label{more-algebra-lemma-flat-descent-tor-amplitude}
$R$ を環、$a, b \in \mathbf{Z}$ とし、$K^\bullet$ を $R$-加群の複体とする。
$R \to R'$ を忠実平坦な環準同型とする。
複体 $K^\bullet \otimes_R R'$ が $[a, b]$ に Tor 振幅を持つならば、
$K^\bullet$ は $[a, b]$ に Tor 振幅を持つ。
\end{lemma}

\begin{proof}
$M$ を $R$-加群とする。$R \to R'$ は平坦なので
% P02-E0941: frozen source display at lines 17248--17251 has one unmatched opening parenthesis on its right-hand side.
$$
(M \otimes_R^{\mathbf{L}} K^\bullet) \otimes_R R'
=
((M \otimes_R R') \otimes_{R'}^{\mathbf{L}} (K^\bullet \otimes_R R')
$$
であり、コホモロジーを取る操作は $R'$ とのテンソル積と可換である。
従って
$\text{Tor}_i^R(M, K^\bullet) \otimes_R R' =
\text{Tor}_i^{R'}(M \otimes_R R', K^\bullet \otimes_R R')$ である。
$R \to R'$ は忠実平坦なので、$i \not \in [a, b]$ に対する
$\text{Tor}_i^{R'}(M \otimes_R R', K^\bullet \otimes_R R')$ の消滅は、
$\text{Tor}_i^R(M, K^\bullet)$ の消滅を含意する。
\end{proof}

\begin{lemma}
\label{more-algebra-lemma-no-change-tor-amplitude}
環準同型 $R \to A \to B$ で $A \to B$ が忠実平坦であるものと
$K \in D(A)$ を与える。このとき $R$ 上の $K$ の Tor 振幅は、
$R$ 上の $K \otimes_A^\mathbf{L} B$ の Tor 振幅と等しい。
\end{lemma}

\begin{proof}
これは、任意の $R$-加群 $M$ とすべての $i$ に対して
$H^i(K \otimes_R^\mathbf{L} M) \otimes_A B =
H^i((K \otimes_A^\mathbf{L} B) \otimes_R^\mathbf{L} M)$ だからである。
実際、$K$ を $A$-加群の複体 $K^\bullet$ で表し、自由分解
$F^\bullet \to M$ を選ぶ。このとき等式
$$
\text{Tot}(K^\bullet \otimes_A B \otimes_R F^\bullet) =
\text{Tot}(K^\bullet \otimes_R F^\bullet) \otimes_A B
$$
を持つ。左辺のコホモロジー群は
$H^i((K \otimes_A^\mathbf{L} B) \otimes_R^\mathbf{L} M)$ であり、
右辺では $H^i(K \otimes_R^\mathbf{L} M) \otimes_A B$ を得る。
\end{proof}

\begin{lemma}
\label{more-algebra-lemma-finite-gl-dim-tor-dimension}
$R$ を大域次元が有限値 $d$ である環とする。このとき
\begin{enumerate}
\item すべての加群は Tor 次元 $\leq d$ を持つ。
\item $H^i(K^\bullet) \not = 0$ が $i \in [a, b]$ のときに限り得るような
$R$-加群複体 $K^\bullet$ は $[a - d, b]$ に Tor 振幅を持つ。
\item $R$-加群複体 $K^\bullet$ が有限 Tor 次元を持つことと
$K^\bullet \in D^b(R)$ であることは同値である。
\end{enumerate}
\end{lemma}

\begin{proof}
$R$ に関する仮定は、すべての加群が長さ高々 $d$ の有限射影分解を持つことを意味する。
特にすべての加群は Tor 次元 $\leq d$ を持つ。
第二の主張は Lemma \ref{more-algebra-lemma-cohomology-tor-amplitude} と定義から従う。
第三の主張は第二の言い換えである。
\end{proof}

\begin{lemma}
\label{more-algebra-lemma-check-tor-dimension-modulo-nilpotent}
$R' \to R$ を、核が冪零イデアルである全射環準同型とする。
$K' \in D(R')$ とし、$K = K' \otimes_{R'}^\mathbf{L} R$ と置く。
$a, b \in \mathbf{Z}$ とする。このとき $K$ が $[a, b]$ に Tor 振幅を持つことと
$K'$ がそうであることは同値である。
\end{lemma}

\begin{proof}
一方の向きは Lemma \ref{more-algebra-lemma-pull-tor-amplitude} から従う。
逆向きについて、$K$ は $[a, b]$ に Tor 振幅を持つと仮定し、
$M'$ を $R'$-加群とする。$K' \otimes_{R'}^\mathbf{L} M'$ のコホモロジーが
非零となり得る次数は区間 $[a, b]$ に含まれることを示さなければならない。

\medskip\noindent
$I = \Ker(R' \to R)$ と置く。このときある $n$ に対して $I^n = 0$ である。
$IM' = 0$ ならば、$M'$ を $R$-加群とみなして
$$
K' \otimes_{R'}^\mathbf{L} M' =
K' \otimes_{R'}^\mathbf{L} (R \otimes_R^\mathbf{L} M') =
(K' \otimes_{R'}^\mathbf{L} R) \otimes_R^\mathbf{L} M' =
K \otimes_R^\mathbf{L} M'
$$
と論じられる。$K$ に関する仮定により、これは区間 $[a, b]$ の外で
コホモロジーを持たない。$I^{t + 1}M' = 0$ ならば、短完全列
$$
0 \to IM' \to M' \to M'/IM' \to 0
$$
を考える。$t$ に関する帰納法により、
$K' \otimes_{R'}^\mathbf{L} IM'$ と $K' \otimes_{R'}^\mathbf{L} M'/IM'$ の
コホモロジーが非零となり得る次数は区間 $[a, b]$ 内のみである。
そこで完全三角
$$
K' \otimes_{R'}^\mathbf{L} IM' \to
K' \otimes_{R'}^\mathbf{L} M' \to
K' \otimes_{R'}^\mathbf{L} M'/IM' \to
(K' \otimes_{R'}^\mathbf{L} IM')[1]
$$
により、望み通り $K' \otimes_{R'}^\mathbf{L} M'$ についても同じことが分かる。
\end{proof}







\section{Ext のスペクトル列}
\label{more-algebra-section-spectral-sequence-ext}

\noindent
この節では、Ext 関手を考えるときに現れる種々のスペクトル列を集める。
環 $R$ の導来圏 $D(R)$ の任意の二対象 $L$, $K$ に対して、
Derived Categories, Section \ref{derived-section-ext} の一般的な規約に従い
$$
\Ext^n_R(L, K) = \Hom_{D(R)}(L, K[n])
$$
と書く。

\medskip\noindent
$M$ を $R$-加群、$K \in D^+(R)$ とすると、スペクトル列
\begin{equation}
\label{more-algebra-equation-first-ss-ext}
E_2^{i, j} = \Ext_R^i(M, H^j(K)) \Rightarrow \Ext_R^{i + j}(M, K)
\end{equation}
がある。また $K$ が $R$-加群の下に有界な複体 $K^\bullet$ で表されるならば、
スペクトル列
\begin{equation}
\label{more-algebra-equation-second-ss-ext}
E_1^{i, j} = \Ext_R^j(M, K^i) \Rightarrow \Ext_R^{i + j}(M, K)
\end{equation}
がある。これらのスペクトル列は、関手 $\Hom_R(M, -)$ に
Derived Categories, Lemma \ref{derived-lemma-two-ss-complex-functor}
を適用して得られる。




\section{射影次元}
\label{more-algebra-section-projective-dimension}

\noindent
加群の射影次元は Algebra, Definition
\ref{algebra-definition-finite-proj-dim} で定義した。

\begin{definition}
\label{more-algebra-definition-projective-dimension}
$R$ を環、$K$ を $D(R)$ の対象とする。$K$ が射影加群からなる有界複体で
表されるとき、$K$ は{it 有限射影次元を持つ}という。
$K$ が複体
$$
\ldots \to 0 \to P^a \to P^{a + 1} \to \ldots \to
P^{b - 1} \to P^b \to 0 \to \ldots
$$
と擬同型で、すべての $i \in \mathbf{Z}$ に対して $P^i$ が射影
$R$-加群であるとき、$K$ は {it $[a, b]$ に射影振幅を持つ}という。
\end{definition}

\noindent
$K$ が有限射影次元を持つことと、ある $a, b \in \mathbf{Z}$ に対して
$K$ が $[a, b]$ に射影振幅を持つことは明らかに同値である。
さらに、$K$ が有限射影次元を持つならば $K$ は有界である。
$D(R)$ のこのような対象を検出する補題を述べる。

\begin{lemma}
\label{more-algebra-lemma-projective-amplitude}
$R$ を環、$K$ を $D(R)$ の対象、$a, b \in \mathbf{Z}$ とする。
次は同値である。
\begin{enumerate}
\item $K$ は $[a, b]$ に射影振幅を持つ。
\item すべての $R$-加群 $N$ とすべての $i \not \in [-b, -a]$ に対して
$\Ext^i_R(K, N) = 0$ である。
\item $n > b$ に対して $H^n(K) = 0$ であり、すべての $R$-加群 $N$ と
すべての $i > -a$ に対して $\Ext^i_R(K, N) = 0$ である。
\item $n \not \in [a - 1, b]$ に対して $H^n(K) = 0$ であり、
すべての $R$-加群 $N$ に対して $\Ext^{-a + 1}_R(K, N) = 0$ である。
\end{enumerate}
\end{lemma}

\begin{proof}
(1) を仮定する。$K$ は複体
$$
\ldots \to 0 \to P^a \to P^{a + 1} \to \ldots \to
P^{b - 1} \to P^b \to 0 \to \ldots
$$
であり、すべての $i \in \mathbf{Z}$ に対して $P^i$ は射影 $R$-加群であると
仮定してよい。この場合、Ext 群は複体
$$
\ldots \to 0 \to \Hom_R(P^b, N) \to \ldots \to
\Hom_R(P^a, N) \to 0 \to \ldots
$$
によって計算でき、(2) を得る。

\medskip\noindent
(2) が成り立つと仮定する。$I$ が入射 $R$-加群であるような単射
$H^n(K) \to I$ を選ぶ。$\Hom_R(-, I)$ は完全関手なので、
$\Ext^{-n}(K, I) = \Hom_R(H^n(K), I)$ である。
特に $n > b$ に対して $H^n(K)$ は零であると結論できる。
従って (2) は (3) を含意する。

\medskip\noindent
% P02-E0942: frozen source line 17450 is grammatically malformed: "By the same argument as in (2) implies (3) gives that (3) implies (4)."
(2) が (3) を含意することを示したのと同じ議論により、(3) は (4) を含意する。

\medskip\noindent
(4) を仮定する。(2) が (3) を含意することを示したのと同じ議論により
$H^{a - 1}(K) = 0$、すなわち $i \in [a, b]$ でない限り
$H^i(K) = 0$ である。特に $K$ は上に有界であり、
% P02-E0943: frozen source line 17456 repeats the article in "choose a a complex".
$K$ を表す複体 $P^\bullet$ で、すべての $i \in \mathbf{Z}$ に対して
$P^i$ が射影（例えば自由）であり、$i > b$ に対して $P^i = 0$ となるものを選べる。
Derived Categories, Lemma
\ref{derived-lemma-subcategory-left-resolution} を参照せよ。
$Q = \Coker(P^{a - 1} \to P^a)$ と置く。$i < a$ に対して $H^i(K) = 0$ なので、
$K$ は複体
$$
\ldots \to 0 \to Q \to P^{a + 1} \to \ldots \to P^b \to 0 \to \ldots
$$
と擬同型である。
% P02-E0944: frozen source line 17466 uses the malformed denote construction "Denote K' = (...) the corresponding object".
対応する $D(R)$ の対象を $K' = (P^{a + 1} \to \ldots \to P^b)$ と書く。
完全三角
$$
K' \to K \to Q[-a] \to K'[1]
$$
を $D(R)$ で得る。
% P02-E0945: frozen source line 17471 says "Thus ... an exact sequence" without a finite verb.
従って任意の $R$-加群 $N$ に対して完全列
$$
\Ext^{-a}(K', N) \to \Ext^1(Q, N) \to \Ext^{1 - a}(K, N)
$$
を得る。仮定により右端の項は消える。(1) $\Rightarrow$ (2) により左端の項も消える。
従って Algebra, Lemma \ref{algebra-lemma-characterize-projective} により、
$Q$ は射影 $R$-加群である。ゆえに (1) が成り立ち、証明は完了する。
\end{proof}

\begin{example}
\label{more-algebra-example-ext-not-bounded}
$k$ を体とし、$R$ を $k$ 上の双対数環、すなわち $R = k[x]/(x^2)$ とする。
% P02-E0946: frozen source line 17485 uses the malformed denote construction "Denote epsilon in R the class of x".
$x$ の類を $\epsilon \in R$ と書く。$M = R/(\epsilon)$ とする。
このとき $M$ は複体
$$
R \xrightarrow{\epsilon} R \xrightarrow{\epsilon} R \to \ldots
$$
と擬同型だが、Algebra, Definition
\ref{algebra-definition-finite-proj-dim} で定義した意味で
$M$ は有限射影次元を持たない。これは、$D(R)$ の対象の有限射影次元の定義で、
射影加群からなる両方向に有界な複体を考える理由を説明する。
\end{example}






\section{入射次元}
\label{more-algebra-section-injective-dimension}

\noindent
この節は射影次元の節の双対である。

\begin{definition}
\label{more-algebra-definition-injective-dimension}
$R$ を環、$K$ を $D(R)$ の対象とする。
$K$ が入射 $R$-加群からなる有限複体で表されるとき、
$K$ は{it 有限入射次元を持つ}という。
$K$ が複体
$$
\ldots \to 0 \to I^a \to I^{a + 1} \to \ldots \to
I^{b - 1} \to I^b \to 0 \to \ldots
$$
と同型で、すべての $i \in \mathbf{Z}$ に対して $I^i$ が入射
$R$-加群であるとき、$K$ は {it $[a, b]$ に入射振幅を持つ}という。
\end{definition}

\noindent
% P02-E0947: frozen source lines 17524--17527 and 17611--17613 switch from the defined term "finite injective dimension" to "bounded injective dimension".
$K$ が有界入射次元を持つことと、ある $a, b \in \mathbf{Z}$ に対して
$K$ が $[a, b]$ に入射振幅を持つことは明らかに同値である。
さらに、$K$ が有界入射次元を持つならば $K$ は有界である。
必須の補題を述べる。

\begin{lemma}
\label{more-algebra-lemma-injective-amplitude}
$R$ を環、$K$ を $D(R)$ の対象、$a, b \in \mathbf{Z}$ とする。
次は同値である。
\begin{enumerate}
\item $K$ は $[a, b]$ に入射振幅を持つ。
\item すべての $R$-加群 $N$ とすべての $i \not \in [a, b]$ に対して
$\Ext^i_R(N, K) = 0$ である。
% P02-E0948: frozen source line 17537 omits the base-ring subscript on Ext, unlike the adjacent criterion and the proof.
\item すべてのイデアル $I \subset R$ とすべての $i \not \in [a, b]$ に対して
$\Ext^i(R/I, K) = 0$ である。
\end{enumerate}
\end{lemma}

\begin{proof}
(1) を仮定する。$K$ は複体
$$
\ldots \to 0 \to I^a \to I^{a + 1} \to \ldots \to
I^{b - 1} \to I^b \to 0 \to \ldots
$$
であり、すべての $i \in \mathbf{Z}$ に対して $I^i$ は入射 $R$-加群であると
仮定してよい。この場合、Ext 群は複体
$$
\ldots \to 0 \to \Hom_R(N, I^a) \to \ldots \to
\Hom_R(N, I^b) \to 0 \to \ldots
$$
によって計算でき、(2) を得る。(2) が (3) を含意することは明らかである。

\medskip\noindent
(3) が成り立つと仮定する。非零写像 $R \to H^n(K)$ を選ぶ。
$\Hom_R(R, -)$ は完全関手なので、
$\Ext^n_R(R, K) = \Hom_R(R, H^n(K)) = H^n(K)$ である。
$n \not \in [a, b]$ に対して $H^n(K)$ は零であると結論できる。
特に $K$ は下に有界であり、すべての $i \in \mathbf{Z}$ に対して $I^i$ が入射、
$i < a$ に対して $I^i = 0$ となる擬同型
$$
K \to I^\bullet
$$
を選べる。Derived Categories, Lemma
\ref{derived-lemma-subcategory-right-resolution} を参照せよ。
$J = \Ker(I^b \to I^{b + 1})$ と置く。このとき $K$ は複体
$$
\ldots \to 0 \to I^a \to \ldots \to I^{b - 1} \to J \to 0 \to \ldots
$$
と擬同型である。
% P02-E0955: frozen source line 17573 repeats the malformed "Denote K' = (...) the corresponding object" construction.
対応する $D(R)$ の対象を $K' = (I^a \to \ldots \to I^{b - 1})$ と書く。
完全三角
$$
J[-b] \to K \to K' \to J[1 - b]
$$
を $D(R)$ で得る。
% P02-E0956: frozen source line 17578 says "Thus ... an exact sequence" without a finite verb.
従って任意のイデアル $I \subset R$ に対して完全列
$$
\Ext^b(R/I, K') \to \Ext^1(R/I, J) \to \Ext^{1 + b}(R/I, K)
$$
を得る。仮定により右端の項は消える。(1) $\Rightarrow$ (2) により左端の項も消える。
従って Lemma \ref{more-algebra-lemma-characterize-injective-bis} により、
$J$ は入射 $R$-加群である。
\end{proof}

\begin{example}
\label{more-algebra-example-finite-injective-finite-global-dimension}
$R$ を Dedekind 整域とする。このとき任意の非零イデアル $I$ は有限射影加群である。
Lemma \ref{more-algebra-lemma-dedekind-torsion-free-flat} を参照せよ。
従って $R/I$ は射影次元 $1$ を持つ。ゆえに Lemma
\ref{more-algebra-lemma-injective-amplitude} により、任意の $R$-加群 $M$ は入射次元 $\leq 1$ を持つ。
従って任意の $R$-加群の対 $M, N$ と $i \geq 2$ に対して
$\Ext^i_R(M, N) = 0$ である。これから、$D^b(R)$ の任意の対象 $K$ は
そのコホモロジーの直和と同型である：
$K \cong \bigoplus H^i(K)[-i]$。Derived Categories,
Lemma \ref{derived-lemma-ext-2-zero} を参照せよ。
\end{example}

\begin{example}
\label{more-algebra-example-ext-not-bounded-reversed}
$k$ を体とし、$R$ を $k$ 上の双対数環、すなわち $R = k[x]/(x^2)$ とする。
% P02-E0957: frozen source line 17603 repeats the malformed "Denote epsilon in R the class of x" construction.
$x$ の類を $\epsilon \in R$ と書く。$M = R/(\epsilon)$ とする。
このとき $M$ は複体
$$
\ldots \to R \xrightarrow{\epsilon} R \xrightarrow{\epsilon} R
$$
と擬同型であり、$R$ は入射 $R$-加群である。
しかし通常、この状況で $M$ が有限入射次元を持つとは考えない。
これは、$D(R)$ の対象の有界入射次元の定義で、入射加群からなる
両方向に有界な複体を考える理由を説明する。
\end{example}

\begin{lemma}
\label{more-algebra-lemma-finite-injective-dimension}
$R$ を環、$K \in D(R)$ とする。
\begin{enumerate}
\item $K$ が $D^b(R)$ に属し、すべての $i$ に対して $H^i(K)$ が
有限入射次元を持つならば、$K$ は有限入射次元を持つ。
\item $K^\bullet$ が $K$ を表す有界な $R$-加群複体であり、すべての $i$ に対して
$K^i$ が有限入射次元を持つならば、$K$ は有限入射次元を持つ。
\end{enumerate}
\end{lemma}

\begin{proof}
省略する。ヒント：関手 $F = \Hom_R(N, -)$ に Derived Categories,
Lemma \ref{derived-lemma-two-ss-complex-functor} のスペクトル列を適用して
% P02-E0949: frozen source line 17632 writes Ext_A although the section and lemma are over R and A is undefined.
$\Ext^i_A(N, K)$ を計算し、Lemma \ref{more-algebra-lemma-injective-amplitude} の判定法を用いよ。
\end{proof}

\begin{lemma}
\label{more-algebra-lemma-finite-injective-dimension-Noetherian-radical}
$R$ を Noether 環、$I \subset R$ を $R$ の Jacobson 根基に含まれるイデアルとする。
$K \in D^+(R)$ は有限なコホモロジー加群を持つとする。次は同値である。
\begin{enumerate}
\item $K$ は有限入射次元を持つ。
\item ある $b$ が存在し、任意のイデアル $J \supset I$ と $i > b$ に対して
$\Ext^i_R(R/J, K) = 0$ である。
\end{enumerate}
\end{lemma}

\begin{proof}
(1) $\Rightarrow$ (2) は直ちに従う。(2) を仮定する。
$i < a$ に対して $H^i(K) = 0$ であるとしよう。
するとすべての $R$-加群 $M$ と $i < a$ に対して $\Ext^i(M, K) = 0$ である。
% P02-E0950: frozen source line 17652 omits the conjunction/preposition in "for i > b any finite R-module M".
従って Lemma \ref{more-algebra-lemma-injective-amplitude} により、任意の有限 $R$-加群 $M$ と
$i > b$ に対して $\Ext^i(M, K) = 0$ であることを示せば十分である。
Algebra, Lemma \ref{algebra-lemma-filter-Noetherian-module} により、
$M$ は、逐次商が素イデアル $\mathfrak p$ に対する $R/\mathfrak p$ の形である
有限フィルトレーションを持つ。$0 \to M_1 \to M \to M_2 \to 0$ が短完全列で、
$j = 1, 2$ と $i > b$ に対して $\Ext^i(M_j, K) = 0$ ならば、
$i > b$ に対して $\Ext^i(M, K) = 0$ である。
従って $M = R/\mathfrak p$ と仮定してよい。
$I \subset \mathfrak p$ ならば、消滅は仮定から従う。
そうでなければ $f \in I$, $f \not \in \mathfrak p$ を選び、短完全列
$$
0 \to R/\mathfrak p \xrightarrow{f} R/\mathfrak p \to R/(\mathfrak p, f) \to 0
$$
を考える。$R$-加群 $R/(\mathfrak p, f)$ は、逐次商が
$(\mathfrak p, f) \subset \mathfrak q$ を満たす $R/\mathfrak q$ である
フィルトレーションを持つ。従って Noether 帰納法と上の議論により、
$R/(\mathfrak p, f)$ に対して消滅が成り立つと仮定してよい。
一方、加群 $E^i = \Ext^i(R/\mathfrak p, K)$ は、$K$ に関する仮定
（有限コホモロジー加群を持ち下に有界）、スペクトル列
(\ref{more-algebra-equation-first-ss-ext})、および Algebra, Lemma
\ref{algebra-lemma-ext-noetherian} により有限である。
従って $i > b$ に対して $E^i$ は $E^i/fE^i = 0$ を満たす有限 $R$-加群である。
Nakayama の補題（Algebra, Lemma \ref{algebra-lemma-NAK}）により、$E^i$ は零である。
\end{proof}

\begin{lemma}
\label{more-algebra-lemma-finite-injective-dimension-Noetherian-local}
$(R, \mathfrak m, \kappa)$ を Noether 局所環とし、
$K \in D^+(R)$ は有限なコホモロジー加群を持つとする。次は同値である。
\begin{enumerate}
\item $K$ は有限入射次元を持つ。
\item $i \gg 0$ に対して $\Ext^i_R(\kappa, K) = 0$ である。
\end{enumerate}
\end{lemma}

\begin{proof}
これは Lemma \ref{more-algebra-lemma-finite-injective-dimension-Noetherian-radical}
の特別な場合である。
\end{proof}






\section{射影に近い加群}
\label{more-algebra-section-near-projective}

\noindent
% P02-E0951: frozen source line 17707 reads "many different of definitions".
文献には「ほとんど射影的な加群」の異なる定義が多数あるように見える。
この節では多くの可能性のうち一つだけを論じる。

\begin{lemma}
\label{more-algebra-lemma-ideal-factor-through-projective}
$R$ を環とし、$M$, $N$ を $R$-加群とする。
\begin{enumerate}
\item $R$-加群写像 $\varphi : M \to N$ を与えると、次は同値である：
(a) $\varphi$ は射影 $R$-加群を経由して分解する、
(b) $\varphi$ は自由 $R$-加群を経由して分解する。
\item (1) の同値な条件を満たす $\varphi : M \to N$ の集合は、
$\Hom_R(M, N)$ の $R$-部分加群である。
\item 写像 $\psi : M' \to M$ と $\xi : N \to N'$ を与える。
$\varphi : M \to N$ が (1) の同値な条件を満たすならば、
$\xi \circ \varphi \circ \psi : M' \to N'$ もそうである。
\end{enumerate}
\end{lemma}

\begin{proof}
(1)(a) と (1)(b) の同値性は Algebra, Lemma
\ref{algebra-lemma-characterize-projective} から従う。
$\varphi : M \to N$ と $\varphi' : M \to N$ がそれぞれ加群 $P$ と $P'$ を
経由して分解するならば、$\varphi + \varphi'$ は $P \oplus P'$ を経由し、
すべての $\lambda \in R$ に対して $\lambda \varphi$ は $P$ を経由して分解する。
これで (2) が示された。$\varphi : M \to N$ が加群 $P$ を経由して分解し、
$\psi$ と $\xi$ が (3) の通りならば、$\xi \circ \varphi \circ \psi$ は
$P$ を経由して分解する。これで (3) が示された。
\end{proof}

\begin{lemma}
\label{more-algebra-lemma-factor-through-finite-projective}
$R$ を環、$\varphi : M \to N$ を $R$-加群写像とする。
$\varphi$ が射影加群を経由して分解し、$M$ が有限 $R$-加群ならば、
$\varphi$ は有限射影加群を経由して分解する。
\end{lemma}

\begin{proof}
Lemma \ref{more-algebra-lemma-ideal-factor-through-projective} により、ある集合 $I$ に対して
$\sigma$ の終域が $\bigoplus_{i \in I} R$ となるように
$\varphi = \tau \circ \sigma$ と分解できる。
$M$ の生成元 $x_1, \ldots, x_n$ を選び、
$\sigma(x_j) = (a_{ji})_{i \in I}$ と書く。
% P02-E0952: frozen source line 17750 switches the previously defined coefficient order a_{ji} to a_{ij}.
各 $j$ に対して、非零な $a_{ij}$ は有限個しかない。
従って $\sigma$ の像は有限自由 $R$-加群に含まれ、結論を得る。
\end{proof}

\noindent
$R$ を環とする。
% P02-E0953: frozen source lines 17755--17757 characterize an arbitrary R-module using the identity on R; the identity should be on that module.
$R$-加群が射影的であることと、$R$ 上の恒等写像が射影加群を経由して分解することは
同値であることに注意する。

\begin{lemma}
\label{more-algebra-lemma-near-projective}
$R$ を環、$I \subset R$ をイデアル、$M$ を $R$-加群とする。
次の条件は同値である。
\begin{enumerate}
\item すべての $a \in I$ に対して、写像 $a : M \to M$ は射影 $R$-加群を経由して分解する。
\item すべての $a \in I$ に対して、写像 $a : M \to M$ は自由 $R$-加群を経由して分解する。
\item すべての $R$-加群 $N$ に対して、$\Ext^1_R(M, N)$ は $I$ で零化される。
\end{enumerate}
\end{lemma}

\begin{proof}
(1) と (2) の同値性は Lemma \ref{more-algebra-lemma-ideal-factor-through-projective} から従う。
% P02-E0954: frozen source lines 17775--17776 say "because Ext^1_R(P,N) for any..." without the missing vanishing predicate.
(1) が成り立つならば、任意の $N$ と任意の射影加群 $P$ に対する
$\Ext^1_R(P, N)$ により (3) が成り立つ。
逆に (3) が成り立つと仮定する。$P$ が射影的（自由でもよい）である短完全列
$0 \to N \to P \to M \to 0$ を選ぶ。
仮定により、対応する $\Ext^1_R(M, N)$ の元は $I$ で零化される。
従ってすべての $a \in I$ に対して写像 $a : M \to M$ は全射
$P \to M$ を経由して分解でき、(1) が成り立つ。
\end{proof}

\noindent
Lemma \ref{more-algebra-lemma-near-projective} の同値な条件を満たす加群について
簡便に語るため、この性質に名前を与える。

\begin{definition}
\label{more-algebra-definition-near-projective}
$R$ を環、$I \subset R$ をイデアル、$M$ を $R$-加群とする。
Lemma \ref{more-algebra-lemma-near-projective} の同値な条件が成り立つとき、
$M$ は {\it $I$-射影的}\footnote{これは非標準的な記法である。}であるという。
\end{definition}

\noindent
$I$ で零化される加群は $I$-射影的である。

\begin{lemma}
\label{more-algebra-lemma-torsion-near-projective}
$R$ を環、$I \subset R$ をイデアル、$M$ を $R$-加群とする。
$M$ が $I$ で零化されるならば、$M$ は $I$-射影的である。
\end{lemma}

\begin{proof}
定義と、零加群が射影的であることから直ちに従う。
\end{proof}

\begin{lemma}
\label{more-algebra-lemma-ses-near-projective}
$R$ を環、$I \subset R$ をイデアルとし、$R$-加群の短完全列
$$
0 \to K \to P \to M \to 0
$$
を考える。$M$ が $I$-射影的で $P$ が射影的ならば、$K$ は $I$-射影的である。
\end{lemma}

\begin{proof}
元 $\text{id}_K \in \Hom_R(K, K)$ は、与えられた拡大の
$\Ext^1_R(M, K)$ における類へ写る。
仮定によりこの類は任意の $a \in I$ で零化されるので、
$a : K \to K$ は $K \to P$ を経由して分解し、結論を得る。
\end{proof}

\begin{lemma}
\label{more-algebra-lemma-dual-near-projective}
$R$ を環、$I \subset R$ をイデアルとする。
$M$ が有限で $I$-射影的な $R$-加群ならば、
$M^\vee = \Hom_R(M, R)$ は $I$-射影的である。
\end{lemma}

\begin{proof}
$M$ は有限で $I$-射影的であると仮定する。
短完全列 $0 \to K \to R^{\oplus r} \to M \to 0$ を選ぶ。
これは単射 $M^\vee \to R^{\oplus r} = (R^{\oplus r})^\vee$ を生じる。
この短完全列に対応する $\Ext^1_R(M, K)$ の拡大類は $I$ で零化されるので、
任意の $a \in I$ に対して写像 $M \to R^{\oplus r}$ で、与えられた写像
$R^{\oplus r} \to M$ との合成が $a : M \to M$ に等しくなるものを見つけられる。
双対を取ると、$a : M^\vee \to M^\vee$ は上で与えた写像
$M^\vee \to R^{\oplus r}$ を経由して分解することが分かり、結論を得る。
\end{proof}





\section{Hom 複体}
\label{more-algebra-section-hom-complexes}

\noindent
$R$ を環とし、$L^\bullet$ と $M^\bullet$ を二つの $R$-加群複体とする。
複体 $\Hom^\bullet(L^\bullet, M^\bullet)$ を構成する。
すなわち、各 $n$ に対して
$$
\Hom^n(L^\bullet, M^\bullet) =
\prod\nolimits_{n = p + q} \Hom_R(L^{-q}, M^p)
$$
と置く。$\Hom^n$ を、次数付き加群とみなした $L^\bullet$ から
$M^\bullet$ への次数 $n$ の斉次 $R$-線形写像すべてからなる
$R$-加群と考えるとよい。この用語のもとで、微分を規則
$$
\text{d}(f) = \text{d}_M \circ f - (-1)^n f \circ \text{d}_L
$$
によって定義する。ここで $f \in \Hom^n(L^\bullet, M^\bullet)$ である。
$\text{d}^2 = 0$ の確認は省略する。符号規則については
Section \ref{more-algebra-section-sign-rules} を参照せよ。この構成は
Differential Graded Algebra, Example \ref{dga-example-category-complexes}
の特別な場合である。構成から直ちに、すべての $n \in \mathbf{Z}$ に対して
\begin{equation}
\label{more-algebra-equation-cohomology-hom-complex}
H^n(\Hom^\bullet(L^\bullet, M^\bullet)) = \Hom_{K(R)}(L^\bullet, M^\bullet[n])
\end{equation}
を得る。

\begin{lemma}
\label{more-algebra-lemma-compose}
$R$ を環とする。$R$-加群の複体 $K^\bullet, L^\bullet, M^\bullet$ を与えると、
$R$-加群複体の標準同型
$$
\Hom^\bullet(K^\bullet, \Hom^\bullet(L^\bullet, M^\bullet))
=
\Hom^\bullet(\text{Tot}(K^\bullet \otimes_R L^\bullet), M^\bullet)
$$
がある。
\end{lemma}

\begin{proof}
$\alpha$ を左辺の次数 $n$ の元とする。従って
$$
\alpha = (\alpha^{p, q}) \in
\prod\nolimits_{p + q = n} \Hom_R(K^{-q}, \Hom^p(L^\bullet, M^\bullet))
$$
である。各 $\alpha^{p, q}$ は元
$$
\alpha^{p, q} = (\alpha^{r, s, q}) \in
\prod\nolimits_{r + s + q = n} \Hom_R(K^{-q}, \Hom_R(L^{-s}, M^r))
$$
である。同一視
\begin{equation}
\label{more-algebra-equation-identification}
\Hom_R(K^{-q}, \Hom_R(L^{-s}, M^r)) = \Hom_R(K^{-q} \otimes_R L^{-s}, M^r)
\end{equation}
を行うと、符号規則により
\begin{align*}
\text{d}(\alpha^{r, s, q})
& =
\text{d}_{\Hom^\bullet(L^\bullet, M^\bullet)} \circ \alpha^{r, s, q}
- (-1)^n \alpha^{r, s, q} \circ \text{d}_K \\
& =
\text{d}_M \circ \alpha^{r, s, q}
- (-1)^{r + s} \alpha^{r, s, q} \circ \text{d}_L
- (-1)^{r + s + q} \alpha^{r, s, q} \circ \text{d}_K
\end{align*}
を得る。一方、$\beta$ を右辺の次数 $n$ の元とすると
$$
\beta = (\beta^{r, s, q}) \in \prod\nolimits_{r + s + q = n}
\Hom_R(K^{-q} \otimes_R L^{-s}, M^r)
$$
であり、符号規則（Homology, Definition
\ref{homology-definition-associated-simple-complex}）により
\begin{align*}
\text{d}(\beta^{r, s, q})
& =
\text{d}_M \circ \beta^{r, s, q}
- (-1)^n \beta^{r, s, q} \circ
\text{d}_{\text{Tot}(K^\bullet \otimes L^\bullet)} \\
& =
\text{d}_M \circ \beta^{r, s, q}
- (-1)^{r + s + q} \left(
\beta^{r, s, q} \circ \text{d}_K + (-1)^{-q} \beta^{r, s, q} \circ \text{d}_L
\right)
\end{align*}
を得る。従って同一視 (\ref{more-algebra-equation-identification}) が誘導する写像は
実際に複体の射である。
\end{proof}

\begin{remark}
\label{more-algebra-remark-internal-hom}
$R$ を環とする。$R$-加群複体の圏 $\text{Comp}(R)$ は、
テンソル積を $\text{Tot}(- \otimes_R -)$ で与えた対称モノイド圏である。
Lemma \ref{more-algebra-lemma-symmetric-monoidal-cat-complexes} を参照せよ。
$\text{Comp}(R)$ の $L^\bullet$ と $M^\bullet$ を与えると、元
$f \in \Hom^0(L^\bullet, M^\bullet)$ が複体写像 $f : L^\bullet \to M^\bullet$
を定めることと $\text{d}(f) = 0$ は同値である。
従って Lemma \ref{more-algebra-lemma-compose} は、$\text{Comp}(R)$ の
$K^\bullet, L^\bullet, M^\bullet$ に関して関手的な等式
$$
\Mor_{\text{Comp}(R)}(K^\bullet, \Hom^\bullet(L^\bullet, M^\bullet))
=
\Mor_{\text{Comp}(R)}(\text{Tot}(K^\bullet \otimes_R L^\bullet), M^\bullet)
$$
も与える。これは、Categories, Remark
\ref{categories-remark-internal-hom-monoidal} で論じた意味で、
$\Hom^\bullet( - , -)$ が対称モノイド圏 $\text{Comp}(R)$ の内部 Hom であることを意味する。
\end{remark}

\begin{lemma}
\label{more-algebra-lemma-composition}
$R$ を環とする。$R$-加群の複体 $K^\bullet, L^\bullet, M^\bullet$ を与えると、
$R$-加群複体の標準射
$$
\text{Tot}\left(
\Hom^\bullet(L^\bullet, M^\bullet) \otimes_R \Hom^\bullet(K^\bullet, L^\bullet)
\right)
\longrightarrow
\Hom^\bullet(K^\bullet, M^\bullet)
$$
がある。
\end{lemma}

\begin{proof}
Remark \ref{more-algebra-remark-internal-hom} の議論を通せば、このような標準写像の存在は
Categories, Remark \ref{categories-remark-internal-hom-monoidal} から従う。
直接の構成も与える。

\medskip\noindent
$\alpha$ を左辺の次数 $n$ の元とすると
$$
\alpha = (\alpha^{p, q}) \in \bigoplus\nolimits_{p + q = n}
\Hom^p(L^\bullet, M^\bullet) \otimes_R \Hom^q(K^\bullet, L^\bullet)
$$
である。元 $\alpha^{p, q}$ は有限和
$\alpha^{p, q} = \sum \beta^p_i \otimes \gamma^q_i$ であり、ここで
$$
\beta^p_i = (\beta^{r, s}_i)
\in \prod\nolimits_{r + s = p} \Hom_R(L^{-s}, M^r)
$$
および
$$
\gamma^q_i = (\gamma^{u, v}_i)
\in \prod\nolimits_{u + v = q} \Hom_R(K^{-v}, L^u)
$$
である。写像は $\alpha$ を $\delta = (\delta^{r, v})$ へ送り、ここで
$$
\delta^{r, v} =
\sum\nolimits_{i, s}  \beta^{r, s}_i
\circ \gamma^{-s, v}_i
\in \Hom_R(K^{-v}, M^r)
$$
とすることで与えられる。$r + v = n$ を固定すると、非零な
$\alpha^{p, q}$ は有限個しかなく、従って非零な $\beta^p_i$ と $\gamma^q_i$ も
有限個しかないので、この和は有限である。符号規則により
\begin{align*}
\text{d}(\alpha^{p, q})
& =
\text{d}_{\Hom^\bullet(L^\bullet, M^\bullet)}(\alpha^{p, q})
+ (-1)^p \text{d}_{\Hom^\bullet(K^\bullet, L^\bullet)}(\alpha^{p, q}) \\
& =
\sum \Big( \text{d}_M \circ \beta^p_i \circ \gamma^q_i
- (-1)^p \beta^p_i \circ \text{d}_L \circ \gamma^q_i \Big) \\
&
\quad + (-1)^p \sum \Big( \beta^p_i \circ \text{d}_L \circ \gamma^q_i
- (-1)^q \beta^p_i \circ \gamma^q_i \circ \text{d}_K \Big) \\
& =
\sum \Big( \text{d}_M \circ \beta^p_i \circ \gamma^q_i 
-(-1)^n \beta^p_i \circ \gamma^q_i \circ \text{d}_K \Big)
\end{align*}
を得る。
% P02-E0958: frozen source line 18031 says "the rules ... is compatible" although a single assignment/rule is meant.
従って規則 $\alpha \mapsto \delta$ は微分と可換であり、補題が示された。
\end{proof}

\begin{lemma}
\label{more-algebra-lemma-diagonal-better}
$R$ を環とする。$R$-加群の複体 $K^\bullet, L^\bullet, M^\bullet$ を与えると、
三つの複体すべてに関して関手的な $R$-加群複体の標準射
$$
\text{Tot}(K^\bullet \otimes_R \Hom^\bullet(M^\bullet, L^\bullet))
\longrightarrow
\Hom^\bullet(M^\bullet, \text{Tot}(K^\bullet \otimes_R L^\bullet))
$$
がある。
\end{lemma}

\begin{proof}
Remark \ref{more-algebra-remark-internal-hom} の議論を通せば、このような標準写像の存在は
Categories, Remark \ref{categories-remark-internal-hom-monoidal} から従う。
直接の構成も与える。

\medskip\noindent
$\alpha$ を右辺の次数 $n$ の元とする。従って
$$
\alpha = (\alpha^{p, q}) \in \prod\nolimits_{p + q = n}
\Hom_R(M^{-q}, \text{Tot}^p(K^\bullet \otimes_R L^\bullet))
$$
である。各 $\alpha^{p, q}$ は元
$$
\alpha^{p, q} = (\alpha^{r, s, q}) \in
\Hom_R(M^{-q}, \bigoplus\nolimits_{r + s + q = n} K^r \otimes_R L^s)
$$
である。ここで $\alpha^{r, s, q}$ を、すべての $x \in M^{-q}$ に対して
非零な $\alpha^{r, s, q}(x)$ が有限個しかないような写像の族と考える。
符号規則により
\begin{align*}
\text{d}(\alpha^{r, s, q})
& =
\text{d}_{\text{Tot}(K^\bullet \otimes_R L^\bullet)} \circ \alpha^{r, s, q}
- (-1)^n \alpha^{r, s, q} \circ \text{d}_M \\
& =
\text{d}_K \circ \alpha^{r, s, q} + (-1)^r \text{d}_L \circ \alpha^{r, s, q}
- (-1)^n \alpha^{r, s, q} \circ \text{d}_M
\end{align*}
を得る。一方、$\beta$ を左辺の次数 $n$ の元とすると
$$
\beta = (\beta^{p, q}) \in
\bigoplus\nolimits_{p + q = n} K^p \otimes_R \Hom^q(M^\bullet, L^\bullet)
$$
であり、$\beta^{p, q} = \sum \gamma_i^p \otimes \delta_i^q$ と書ける。
ここで $\gamma_i^p \in K^p$ かつ
$$
\delta_i^q = (\delta_i^{r, s}) \in
\prod\nolimits_{r + s = q} \Hom_R(M^{-s}, L^r)
$$
である。符号規則により
\begin{align*}
\text{d}(\beta^{p, q})
& =
\text{d}_K(\beta^{p, q}) +
(-1)^p \text{d}_{\Hom^\bullet(M^\bullet, L^\bullet)}(\beta^{p, q}) \\
& =
\sum \text{d}_K(\gamma_i^p) \otimes \delta_i^q +
(-1)^p \sum \gamma_i^p \otimes
(\text{d}_L \circ \delta_i^q - (-1)^q \delta_i^q \circ \text{d}_M)
\end{align*}
を得る。元 $\beta$ を $\alpha$ へ送り、
$$
\alpha^{r, s, q} = c^{r, s, q}(\sum \gamma_i^r \otimes \delta_i^{s, q})
$$
とする。ここで $c^{r, s, q} : K^r \otimes_R \Hom_R(M^{-q}, L^s) \to
\Hom_R(M^{-q}, K^r \otimes_R L^s)$ は標準写像である。
与えられた $\beta$ と $r$ に対して非零な $\gamma_i^r$ は有限個しかないので、
% P02-E0959: frozen source lines 18106--18108 redundantly say "only finitely many nonzero ... are nonzero".
非零な $\alpha^{r, s, q}$ も（$r$ を固定すれば）有限個しかない。
従ってこの写像族は上の条件を満たし、写像は良定義である。
符号を比較すると、これは微分と可換であることが分かる。
\end{proof}

\begin{lemma}
\label{more-algebra-lemma-diagonal}
$R$ を環とする。$R$-加群の複体 $K^\bullet, L^\bullet$ が与えられると、
両方の複体に関して関手的な $R$-加群の複体の標準射
$$
K^\bullet
\longrightarrow
\Hom^\bullet(L^\bullet, \text{Tot}(K^\bullet \otimes_R L^\bullet))
$$
がある。
\end{lemma}

\begin{proof}
Remark \ref{more-algebra-remark-internal-hom} の議論を通せば、このような標準写像の存在は
Categories, Remark \ref{categories-remark-internal-hom-monoidal} から従う。
直接の構成も与える。

\medskip\noindent
$\alpha$ を右辺の次数 $n$ の元とする。従って
$$
\alpha = (\alpha^{p, q}) \in \prod\nolimits_{p + q = n}
\Hom_R(L^{-q}, \text{Tot}^p(K^\bullet \otimes_R L^\bullet))
$$
である。各 $\alpha^{p, q}$ は元
$$
\alpha^{p, q} = (\alpha^{r, s, q}) \in
\Hom_R(L^{-q}, \bigoplus\nolimits_{r + s + q = n} K^r \otimes_R L^s)
$$
である。ここで $\alpha^{r, s, q}$ を、すべての $x \in L^{-q}$ に対して
非零な $\alpha^{r, s, q}(x)$ が有限個しかないような写像の族と考える。
符号規則により
\begin{align*}
\text{d}(\alpha^{r, s, q})
& =
\text{d}_{\text{Tot}(K^\bullet \otimes_R L^\bullet)} \circ \alpha^{r, s, q}
- (-1)^n \alpha^{r, s, q} \circ \text{d}_L \\
& =
\text{d}_K \circ \alpha^{r, s, q} + (-1)^r \text{d}_L \circ \alpha^{r, s, q}
- (-1)^n \alpha^{r, s, q} \circ \text{d}_L
\end{align*}
を得る。ここで元 $\beta \in K^n$ を、
$\alpha^{n, -q, q} = \beta \otimes \text{id}_{L^{-q}}$ かつ
$r \not = n$ ならば $\alpha^{r, s, q} = 0$ となる $\alpha$ へ送る。
これは実際に上のような元である。なぜなら、$q$ を固定すれば非零な
$\alpha^{r, s, q}$ は一つしかないからである。
微分の記述から、これは微分と可換であることが分かる。
\end{proof}

\begin{lemma}
\label{more-algebra-lemma-evaluate-and-more}
$R$ を環とする。$R$-加群の複体 $K^\bullet, L^\bullet, M^\bullet$ が
与えられると、三つの複体すべてに関して関手的な $R$-加群の複体の標準射
$$
\text{Tot}(\Hom^\bullet(L^\bullet, M^\bullet) \otimes_R K^\bullet)
\longrightarrow
\Hom^\bullet(\Hom^\bullet(K^\bullet, L^\bullet), M^\bullet)
$$
がある。
\end{lemma}

\begin{proof}
Remark \ref{more-algebra-remark-internal-hom} の議論を通せば、このような標準写像の存在は
Categories, Remark \ref{categories-remark-internal-hom-monoidal} から従う。
直接の構成も与える。

\medskip\noindent
$\beta$ を右辺の次数 $n$ の元とする。このとき
$$
\beta = (\beta^{p, s}) \in
\prod\nolimits_{p + s = n} \Hom_R(\Hom^{-s}(K^\bullet, L^\bullet), M^p)
$$
である。符号規則から
\begin{align*}
\text{d}(\beta^{p, s})
& =
\text{d}_M \circ \beta^{p, s}
- (-1)^n \beta^{p, s} \circ
\text{d}_{\Hom^\bullet(K^\bullet, L^\bullet)}
\end{align*}
が分かる。最後の項は次のように記述できる：
$$
(\beta^{p, s} \circ \text{d}_{\Hom^\bullet(K^\bullet, L^\bullet)})(f) =
\beta^{p, s}(\text{d}_L \circ f - (-1)^{s + 1} f \circ \text{d}_K)
$$
ただし $f \in \Hom^{-s - 1}(K^\bullet, L^\bullet)$ である。
従って、ある明示しない意味で $\text{d}(\beta^{p, s})$ は
次の符号を持つ三項の和である：
\begin{equation}
\label{more-algebra-equation-beta}
\text{d}(\beta^{p, s}) =
\text{d}_M(\beta^{p, s})
-(-1)^n\text{d}_L(\beta^{p, s}) +
(-1)^{p + 1}\text{d}_K(\beta^{p, s})
\end{equation}

\noindent
次に、左辺の次数 $n$ の元 $\alpha$ を考える。これは
$$
\alpha = (\alpha^{t, r}) \in
\bigoplus\nolimits_{t + r = n} \Hom^t(L^\bullet, M^\bullet) \otimes K^r
$$
と書ける。各 $\alpha^{t, r}$ は元
$$
\alpha^{t, r} \mapsto (\alpha^{p, q, r}) \in
\prod\nolimits_{p + q = t} \Hom_R(L^{-q}, M^p) \otimes_R K^r
$$
へ写る。符号規則から
\begin{align*}
\text{d}(\alpha^{p, q, r})
& =
\text{d}_{\Hom^\bullet(L^\bullet, M^\bullet)}(\alpha^{p, q, r})
+ (-1)^{p + q} \text{d}_K(\alpha^{p, q, r})
\end{align*}
が分かる。さらに $\alpha^{p, q, r} = \sum g_i^{p, q} \otimes k_i^r$
と書けば
$$
\text{d}_{\Hom^\bullet(L^\bullet, M^\bullet)}(\alpha^{p, q, r}) =
\sum (\text{d}_M \circ g_i^{p, q}) \otimes k_i^r
- (-1)^{p + q} \sum (g_i^{p, q} \circ \text{d}_L) \otimes k_i^r
$$
である。従って、ある明示しない意味で $\text{d}(\alpha^{p, q, r})$ は
次の符号を持つ三項の和である：
\begin{equation}
\label{more-algebra-equation-alpha}
\text{d}(\alpha^{p, q, r}) =
\text{d}_M(\alpha^{p, q, r})
-(-1)^{p + q}\text{d}_L(\alpha^{p, q, r}) +
(-1)^{p + q}\text{d}_K(\alpha^{p, q, r})
\end{equation}

\noindent
写像を定義するため、標準写像
$$
c_{p, q, r} : 
\Hom_R(L^{-q}, M^p) \otimes_R K^r
\longrightarrow
\Hom_R(\Hom_R(K^r, L^{-q}), M^p)
$$
を用いる。これは $\varphi \otimes k$ を写像
$\psi \mapsto \varphi(\psi(k))$ へ送る。この写像は三変数すべてについて関手的である。
$s = q + r$ とすると、射影
$\Hom^{-s}(K^\bullet, L^\bullet) \to \Hom_R(K^r, L^{-q})$ から生じる包含
$$
\Hom_R(\Hom_R(K^r, L^{-q}), M^p) \subset
\Hom_R(\Hom^{-s}(K^\bullet, L^\bullet), M^p)
$$
がある。$\alpha^{p, q, r}$ は有限個の $r$ についてしか非零でないので、
与えられた $s$ に対して $q + r = s$ を満たす $q, r$ は有限個しかない。
従って $\alpha$ を、
$$
\beta^{p, s} =
\sum\nolimits_{q + r = s} \epsilon_{p, q, r} c_{p, q, r}(\alpha^{p, q, r})
$$
を満たす元 $\beta$ へ送れる。
% P02-E0960: frozen source line 18273 duplicates "where" in "where where the sum uses".
ここで、和は上で与えた包含を用い、ここで
$\epsilon_{p, q, r} \in \{\pm 1\}$ とする。
式 (\ref{more-algebra-equation-beta}) と (\ref{more-algebra-equation-alpha}) の符号を比較すると、
\begin{enumerate}
\item $\epsilon_{p, q, r} = \epsilon_{p + 1, q, r}$
\item $-(-1)^n\epsilon_{p, q, r} = -(-1)^{p + q}\epsilon_{p, q - 1, r}$
または同値に $\epsilon_{p, q, r} = (-1)^r\epsilon_{p, q - 1, r}$
\item $(-1)^{p + 1}\epsilon_{p, q, r} = (-1)^{p + q}\epsilon_{p, q, r + 1}$
または同値に
$(-1)^{q + 1}\epsilon_{p, q, r} = \epsilon_{p, q, r + 1}$。
\end{enumerate}
よい解は
% P02-E0961: frozen source line 18287 writes epsilon_{p,r,s}, inconsistent with epsilon_{p,q,r} throughout the sign equations.
$$
\epsilon_{p, r, s} = (-1)^{r + qr}
$$
と取ることである。この符号を選ぶ理由は、証明に続く注意で説明する。
\end{proof}

\begin{remark}
\label{more-algebra-remark-sign-explanation}
Lemma \ref{more-algebra-lemma-evaluate-and-more} の証明の直接構成で用いた符号が、
Remark \ref{more-algebra-remark-internal-hom} と Categories, Remark
\ref{categories-remark-internal-hom-monoidal} の議論を用いる構成から得られる符号と
一致する理由を説明しよう。
$- \otimes - = \text{Tot}(- \otimes_R -)$ および
$hom(-, -) = \Hom^\bullet(-, -)$ と書く。
モノイド圏の言葉を用いる構成によれば、$\text{Comp}(R)$ において射
$$
hom(L^\bullet, M^\bullet) \otimes K^\bullet
\longrightarrow
hom(hom(K^\bullet, L^\bullet), M^\bullet)
$$
に対応する射
$$
hom(L^\bullet, M^\bullet) \otimes K^\bullet \otimes hom(K^\bullet, L^\bullet)
\longrightarrow
M^\bullet
$$
を用いる。これは最後の二つのテンソル積の順序を入れ替え、次いで評価写像
$hom(K^\bullet, L^\bullet) \otimes K^\bullet \to L^\bullet$ および
% P02-E0962: frozen source line 18316 has hom(L,K) in the second evaluation map, inconsistent with its codomain M and the surrounding hom(L,M) construction.
$hom(L^\bullet, K^\bullet) \otimes L^\bullet \to M^\bullet$
を用いて得られる。符号が介在するのは入れ替えにおいてのみである。すなわち同型
$$
K^\bullet \otimes hom(K^\bullet, L^\bullet)
\to
hom(K^\bullet, L^\bullet) \otimes K^\bullet
$$
では、$K^r \otimes_R \Hom_R(K^{-r'}, L^q)$ 上に符号
$(-1)^{r(q + r')}$ がある。Section \ref{more-algebra-section-sign-rules} の項
(\ref{more-algebra-item-flip-tensor-product}) を見よ。
内部 Hom への写像を記述するために用いている対応のため、この符号が問題になるのは
$r = r'$ の場合だけであり、その場合には直接の証明と同じ
$(-1)^{r(q + r)} = (-1)^{r + qr}$ を得ることを読者は確認できる。
\end{remark}

\section{符号規則}
\label{more-algebra-section-sign-rules}

\noindent
この節では、これまで用いてきた符号規則を復習し、そのいくつかの帰結を論じる。
興味深い現象の大部分は既にこの場合に現れるので、これらの問題を環上の加群の複体の
圏という設定で論じることも適切であろう。
読者がこの節のより秘教的な側面を用いずに済むことを心から願う。

\medskip\noindent
この節の残りでは環 $R$ を固定し、微分
$d^n_M : M^n \to M^{n + 1}$ を持つ $R$-加群の複体を
% P02-E0963: frozen source lines 18350--18352 use the malformed construction "we denote M^bullet a complex".
$M^\bullet$ と書く。
\begin{enumerate}
\item $k$ 回シフトした複体 $M^\bullet[k]$ の項は
$(M^\bullet[k])^n = M^{n + k}$ であり、微分は
$d_{M[k]}^n = (-1)^kd^{n + k}_M$ である。Homology, Definition
\ref{homology-definition-shift-cochain} を見よ。
\item 複体の写像 $f : M^\bullet \to N^\bullet$ が与えられたとき、
$f[k] : M^\bullet[k] \to N^\bullet[k]$ を符号を介さずに定義する。
Homology, Definition \ref{homology-definition-shift-cochain} を見よ。
\item $H^n(M^\bullet[k])$ と $H^{n + k}(M^\bullet)$ を符号を介さずに同一視する。
Homology, Definition \ref{homology-definition-cohomology-shift} を見よ。
\item 複体の短完全列の境界写像は、符号を介さず蛇の補題と同じように定義する。
Homology, Lemma \ref{homology-lemma-long-exact-sequence-cochain} を見よ。
\item 複体の各項で分裂する短完全列
$0 \to K^\bullet \to L^\bullet \to M^\bullet \to 0$ に付随する完全三角は
$$
K^\bullet \to L^\bullet \to M^\bullet \to K^\bullet[1]
$$
で与えられる。ここで $s$ と $\pi$ が両立する各項ごとの分裂ならば、
$M^n \to K^{n + 1}$ は写像 $\pi^{n + 1} \circ d^n_L \circ s^n$ である。
言い換えれば符号を介さない。Derived Categories, Definitions
\ref{derived-definition-distinguished-triangle} および
\ref{derived-definition-split-ses} を見よ。
\item 全複体 $\text{Tot}(M^\bullet \otimes_R N^\bullet)$ の微分 $d$ は
Leibniz 則
$d(x \otimes y) = d(x) \otimes y + (-1)^{\deg(x)}x \otimes d(y)$
を満たす。Homology, Example
\ref{homology-example-double-complex-as-tensor-product-of} および
Homology, Definition \ref{homology-definition-associated-simple-complex} を見よ。
\item
\label{more-algebra-item-shift-tensor}
標準同型
$$
\text{Tot}(M^\bullet \otimes_R N^\bullet)[a + b] \to
\text{Tot}(M^\bullet[a] \otimes_R N^\bullet[b])
$$
があり、これは直和因子 $M^p \otimes_R N^q$ 上で符号 $(-1)^{pb}$ を用いる。
Homology, Remark \ref{homology-remark-shift-double-complex} を見よ。
対応するシフトした写像
$\text{Tot}(M^\bullet \otimes_R N^\bullet) \to
\text{Tot}(M^\bullet[a] \otimes_R N^\bullet[b])[-a - b]$
を考える方が便利なことが多い。
\item 複体の標準同型
$$
\text{Tot}(
\text{Tot}(K^\bullet \otimes_R L^\bullet) \otimes_R M^\bullet) \to
\text{Tot}(K^\bullet \otimes_R \text{Tot}(L^\bullet \otimes_R M^\bullet))
$$
があり、符号を介さずに定義される。Section
\ref{more-algebra-section-symmetric-monoidal} を見よ。
\item
\label{more-algebra-item-flip-tensor-product}
標準同型
$$
\text{Tot}(L^\bullet \otimes_R M^\bullet) \to
\text{Tot}(M^\bullet \otimes_R L^\bullet)
$$
があり、これは直和因子 $L^p \otimes_R M^q$ 上で符号 $(-1)^{pq}$ を用いる。
Section \ref{more-algebra-section-symmetric-monoidal} を見よ。
\end{enumerate}
Hom 複体に関する符号規約の議論へ進む前に、上の規約に従って複体の双対を構成する。

\begin{lemma}
\label{more-algebra-lemma-left-dual-module}
$R$ を環とする。$M$ を $R$-加群とする。$N, \eta, \epsilon$ を
$R$-加群のモノイド圏における $M$ の左双対とする。Categories, Definition
\ref{categories-definition-dual} を見よ。このとき
\begin{enumerate}
\item $M$ と $N$ は有限射影 $R$-加群である。
\item 写像
$e : \Hom_R(M, R) \to N$, $\lambda \mapsto (\lambda \otimes 1)(\eta)$
は同型である。
\item $n \in N$ および $m \in M$ に対して
$\epsilon(n, m) = e^{-1}(n)(m)$ である。
\end{enumerate}
\end{lemma}

\begin{proof}
仮定は
$$
M \xrightarrow{\eta \otimes 1} M \otimes_R N \otimes_R M
\xrightarrow{1 \otimes \epsilon} M
\quad\text{and}\quad
N \xrightarrow{1 \otimes \eta} N \otimes_R M \otimes_R N
\xrightarrow{\epsilon \otimes 1} N
$$
が恒等写像であることを意味する。有限自由加群 $F$、$R$-加群写像
$F \to M$、および $\eta$ の持ち上げ
$\tilde \eta : R \to F \otimes_R N$ を選べる。可換図式
$$
\xymatrix{
M \ar[rr]_-{\eta \otimes 1} \ar[rrd]_-{\tilde \eta \otimes 1} & &
M \otimes_R N \otimes_R M \ar[r]_-{1 \otimes \epsilon} &
M \\
& &
F \otimes_R N \otimes_R M \ar[u] \ar[r]^-{1 \otimes \epsilon} &
F \ar[u]
}
$$
を得る。これは $M$ 上の恒等写像が有限自由加群を経由して分解することを示すので、
$M$ は有限射影である。対称性により $N$ も有限射影である。
これで (1) が示された。(2) は Categories, Lemma
\ref{categories-lemma-left-dual} とその証明から従う。
(3) は証明の最初の等式から従う。
\end{proof}

\begin{lemma}
\label{more-algebra-lemma-left-dual-complex}
$R$ を環とする。$M^\bullet$ を $R$-加群の複体とする。
$N^\bullet, \eta, \epsilon$ を $R$-加群の複体のモノイド圏における
$M^\bullet$ の左双対とする。このとき
\begin{enumerate}
\item $M^\bullet$ と $N^\bullet$ は有界である。
\item $M^n$ と $N^n$ は有限射影 $R$-加群である。
\item $\epsilon_n : N^{-n} \otimes_R M^n \to R$ として
$\epsilon = \sum \epsilon_n$ と書き、
$\eta_n : R \to  M^n \otimes_R N^{-n}$ として
$\eta = \sum \eta_n$ と書くと、$(N^{-n}, \eta_n, \epsilon_n)$ は
Lemma \ref{more-algebra-lemma-left-dual-module} における $M^n$ の左双対である。
\item 微分 $d_N^n : N^n \to N^{n + 1}$ は写像
$$
N^n = \Hom_R(M^{-n}, R)
\xrightarrow{d_M^{-n - 1}}
\Hom_R(M^{-n - 1}, R) = N^{n + 1}
$$
の $-(-1)^n$ 倍に等しい。ここで等号は Lemma
\ref{more-algebra-lemma-left-dual-module} の (2) による同一視である。
\end{enumerate}
逆に、有限射影 $R$-加群からなる有界複体 $M^\bullet$ が与えられたとする。
$N^n = \Hom_R(M^{-n}, R)$ と置いて上の微分を入れ、
評価により与えられる
$\epsilon_n : N^{-n} \otimes_R M^n \to R$ を用いて
$\epsilon = \sum \epsilon_n$ と置き、さらに
$\eta_n : R \to M^n \otimes_R N^{-n}$ が $1$ を
% P02-E0964: frozen source line 18489 maps 1 to id_{M_n}, switching from the consistently superscripted M^n.
$\text{id}_{M_n}$ へ送るものとして $\eta = \sum \eta_n$ と置けば、
$R$-加群の複体のモノイド圏における $M^\bullet$ の左双対を得る。
\end{lemma}

\begin{proof}
Categories, Definition \ref{categories-definition-dual} により
$(1 \otimes \epsilon) \circ (\eta \otimes 1) = \text{id}_{M^\bullet}$
および
$(\epsilon \otimes 1) \circ (1 \otimes \eta) = \text{id}_{N^\bullet}$
なので、直ちに
$(1 \otimes \epsilon_n) \circ (\eta_n \otimes 1) = \text{id}_{M^n}$
および
$(\epsilon_n \otimes 1) \circ (1 \otimes \eta_n) = \text{id}_{N^{-n}}$
を得る。これは (3) を示す。Lemma \ref{more-algebra-lemma-left-dual-module} により
(2) が成り立つ。和 $\eta = \sum \eta_n$ は有限なので (1) を得る。
$\eta = \sum \eta_n$ は複体の写像
$R \to \text{Tot}(M^\bullet \otimes_R N^\bullet)$ なので、
全複体 $\text{Tot}(M^\bullet \otimes_R N^\bullet)$ の微分について選んだ符号により
$$
(d_M^{-n - 1} \otimes 1) \circ \eta_{-n - 1}  +
(-1)^n (1 \otimes d_N^{-n}) \circ \eta_{-n} = 0
$$
である。定義を展開すれば (4) が示される。
補題の最後の主張を見るには、上の議論を逆向きに読めばよい。
\end{proof}

\noindent
Lemma \ref{more-algebra-lemma-left-dual-complex} における複体の左双対の記述を、
$\Hom$ 複体の符号規則の動機として用いる。
% P02-E0965: frozen source line 18517 misspells "Namely" as "Namelly".
すなわち、(\ref{more-algebra-item-compatible}) が成り立つように符号を選ぶ。
上と同様に、種々の符号規則の議論を続ける。
\begin{enumerate}
\setcounter{enumi}{9}
\item 複体 $K^\bullet$, $M^\bullet$ が与えられたとき、
$\Hom^\bullet(M^\bullet, K^\bullet)$ を、項
$$
\Hom^n(M^\bullet, K^\bullet) = \prod\nolimits_{n = p + q} \Hom_R(M^{-q}, K^p)
$$
および規則
$$
d(f) = d_K \circ f - (-1)^n f \circ d_M
$$
で与えられる微分を持つ複体とする。
\item
\label{more-algebra-item-compatible}
上の選択は、$M^\bullet$ が Lemma \ref{more-algebra-lemma-left-dual-complex} のような
左双対 $N^\bullet$ を持つとき、標準同型
$$
\text{Tot}(K^\bullet \otimes_R N^\bullet)
\longrightarrow
\Hom^\bullet(M^\bullet, K^\bullet)
$$
が符号を介さずに定義されるようなものである。この同型は直和因子
$K^p \otimes_R N^q$ を、$N^q = \Hom_R(M^{-q}, R)$ と標準写像
$K^p \otimes_R \Hom_R(M^{-q}, R) \to \Hom_R(M^{-q}, K^p)$
により直和因子 $\Hom_R(M^{-q}, K^p)$ へ送る。
\item 合成
$$
\text{Tot}(
\Hom^\bullet(L^\bullet, K^\bullet) \otimes_R
\Hom^\bullet(M^\bullet, L^\bullet))
\longrightarrow
\Hom^\bullet(M^\bullet, K^\bullet)
$$
が符号を介さずに定義される。Lemma \ref{more-algebra-lemma-composition} を見よ。
\item 標準同型
$$
\Hom^\bullet(K^\bullet, \Hom^\bullet(L^\bullet, M^\bullet))
=
\Hom^\bullet(\text{Tot}(K^\bullet \otimes_R L^\bullet), M^\bullet)
$$
が符号を介さずに定義される。Lemma \ref{more-algebra-lemma-compose} を見よ。
\item 標準写像
$$
\text{Tot}(K^\bullet \otimes_R \Hom^\bullet(M^\bullet, L^\bullet))
\longrightarrow
\Hom^\bullet(M^\bullet, \text{Tot}(K^\bullet \otimes_R L^\bullet))
$$
が符号を介さずに定義される。Lemma \ref{more-algebra-lemma-diagonal-better} を見よ。
\item 標準写像
$$
K^\bullet
\longrightarrow
\Hom^\bullet(L^\bullet, \text{Tot}(K^\bullet \otimes_R L^\bullet))
$$
が符号を介さずに定義される。Lemma \ref{more-algebra-lemma-diagonal} を見よ。
\item
% P02-E0966: frozen source line 18580 reads "By Lemma ... is a canonical map", omitting "there".
Lemma \ref{more-algebra-lemma-evaluate-and-more} により標準写像
$$
\text{Tot}(\Hom^\bullet(L^\bullet, M^\bullet) \otimes_R K^\bullet)
\longrightarrow
\Hom^\bullet(\Hom^\bullet(K^\bullet, L^\bullet), M^\bullet)
$$
があり、これは加群 $\Hom_R(L^{-q}, M^p) \otimes_R K^r$ 上で
符号 $(-1)^{r + qr}$ を用いる。その理由は Remark
\ref{more-algebra-remark-sign-explanation} で説明される。
\item
\label{more-algebra-item-evaluation}
$L^\bullet = M^\bullet$ とし、$R \to \Hom^\bullet(M^\bullet, M^\bullet)$
を用いると、前項の写像は評価写像
$$
ev : K^\bullet \longrightarrow
\Hom^\bullet(\Hom^\bullet(K^\bullet, M^\bullet), M^\bullet)
$$
となる。これは $x \in K^n$ を、
$f \in \Hom^m(K^\bullet, M^\bullet)$ を $(-1)^{nm}f(x)$ へ送る写像に送る。
\item
\label{more-algebra-item-shift-hom}
符号を用いる標準同一視
$$
\Hom^\bullet(M^\bullet, K^\bullet)[a - b] \to
\Hom^\bullet(M^\bullet[b], K^\bullet[a])
$$
がある。これは、対応するシフトした写像
$$
\Hom^\bullet(M^\bullet, K^\bullet) \to
\Hom^\bullet(M^\bullet[b], K^\bullet[a])[b - a]
$$
が、$p + q = n$ を満たす加群 $\Hom_R(M^{-q}, K^p)$ 上で
符号 $(-1)^{nb}$ を用いる写像となるように定義される。
実際、$f \in \Hom^n(M^\bullet, K^\bullet)$ ならば、始域では
$$
d(f) = d_K \circ f - (-1)^n f \circ d_M
$$
である。一方、終域では $f$ は
$$
\left(\Hom^\bullet(M^\bullet[b], K^\bullet[a])[b - a]\right)^n =
\Hom^{n + b -a}(M^\bullet[b], K^\bullet[a])
$$
に属し、従って
\begin{align*}
d(f) & =
(-1)^{b - a}
\left(d_{K[a]} \circ f - (-1)^{n + b - a} f \circ d_{M[b]}\right) \\
& =
(-1)^{b - a}
\left((-1)^a d_K \circ f - (-1)^{n + b - a} f \circ (-1)^b d_M \right) \\
& =
(-1)^b d_K \circ f - (-1)^{n + b} f \circ d_M
\end{align*}
を得る。次数 $n$ で選んだ符号 $(-1)^{nb}$ が、これらの微分に関する
複体の写像を与えることが分かる。
\end{enumerate}

\section{導来 Hom}
\label{more-algebra-section-RHom}

\noindent
$R$ を環とする。この節で定義する導来 Hom は関手
$$
D(R)^{opp} \times D(R) \longrightarrow D(R),\quad
(K, L) \longmapsto R\Hom_R(K, L)
$$
である。これは、$D(R)$ の対象 $K, L, M$ に対する式
\begin{equation}
\label{more-algebra-equation-internal-hom}
\Hom_{D(R)}(K, R\Hom_R(L, M))
=
\Hom_{D(R)}(K \otimes_R^\mathbf{L} L, M)
\end{equation}
によって特徴付けられるという意味で、$R$-加群の導来圏における内部 Hom である。
Yoneda の補題により、この式は対象を一意な同型を除いて特徴付けることに注意する。
構成は次のように与えられる。$M$ を表す $R$-加群の K-入射複体
$I^\bullet$ と、$L$ を表す複体 $L^\bullet$ を選び、
$$
R\Hom_R(L, M) = \Hom^\bullet(L^\bullet, I^\bullet)
$$
と置く。記法は Section \ref{more-algebra-section-hom-complexes} のものを用いる。
この構成の一般化は Differential Graded Algebra, Section
\ref{dga-section-restriction} で論じられる。
% P02-E0967: frozen source lines 18670--18672 say "From ... that we have", omitting a finite verb.
(\ref{more-algebra-equation-cohomology-hom-complex}) と Derived Categories, Lemma
\ref{derived-lemma-K-injective} から、すべての $n \in \mathbf{Z}$ に対して
\begin{equation}
\label{more-algebra-equation-h0-RHom}
H^n(R\Hom_R(L, M)) = \Hom_{D(R)}(L, M[n])
\end{equation}
を得る。特に、$D(R)$ の対象 $R\Hom_R(L, M)$ は良定義である。
すなわち K-入射複体 $I^\bullet$ の選択によらない。

\begin{lemma}
\label{more-algebra-lemma-internal-hom}
$R$ を環とする。$K, L, M$ を $D(R)$ の対象とする。
$K, L, M$ に関して関手的な $D(R)$ における標準同型
$$
R\Hom_R(K, R\Hom_R(L, M)) = R\Hom_R(K \otimes_R^\mathbf{L} L, M)
$$
があり、$H^0$ を取ると (\ref{more-algebra-equation-internal-hom}) を回復する。
\end{lemma}

\begin{proof}
$M$ を表す K-入射複体 $I^\bullet$ と、$L$ を表す
$R$-加群の K-平坦複体 $L^\bullet$ を選ぶ。
任意の $R$-加群の複体 $K^\bullet$ に対して、Lemma
\ref{more-algebra-lemma-compose} により
$$
\Hom^\bullet(K^\bullet, \Hom^\bullet(L^\bullet, I^\bullet)) =
\Hom^\bullet(\text{Tot}(K^\bullet \otimes_R L^\bullet), I^\bullet)
$$
である。補題は $R\Hom$ の定義と、
$\text{Tot}(K^\bullet \otimes_R L^\bullet)$ が導来テンソル積を表すことから従う。
\end{proof}

\begin{lemma}
\label{more-algebra-lemma-RHom-out-of-projective}
$R$ を環とする。$P^\bullet$ を射影 $R$-加群からなる上に有界な複体とし、
$L^\bullet$ を $R$-加群の複体とする。このとき
$R\Hom_R(P^\bullet, L^\bullet)$ は複体
$\Hom^\bullet(P^\bullet, L^\bullet)$ によって表される。
\end{lemma}

\begin{proof}
(\ref{more-algebra-equation-cohomology-hom-complex}) と Derived Categories, Lemma
\ref{derived-lemma-morphisms-from-projective-complex} により、
この複体のコホモロジー群は「正しい」。
従って、$I^\bullet$ が $R$-加群の K-入射複体であるような擬同型
$L^\bullet \to I^\bullet$ を選べば、誘導写像
$$
\Hom^\bullet(P^\bullet, L^\bullet)
\longrightarrow
\Hom^\bullet(P^\bullet, I^\bullet)
$$
は擬同型である。右辺は $R\Hom_R(P^\bullet, L^\bullet)$ の定義なので結論を得る。
\end{proof}

\begin{lemma}
\label{more-algebra-lemma-internal-hom-evaluate}
$R$ を環とする。$K, L, M$ を $D(R)$ の対象とする。
$K, L, M$ に関して関手的な $D(R)$ における標準射
$$
R\Hom_R(L, M) \otimes_R^\mathbf{L} K
\longrightarrow
R\Hom_R(R\Hom_R(K, L), M)
$$
がある。
\end{lemma}

\begin{proof}
$M$ を表す K-入射複体 $I^\bullet$、$L$ を表す K-入射複体
$J^\bullet$、および $K$ を表す K-平坦複体 $K^\bullet$ を選ぶ。
写像は、Lemma \ref{more-algebra-lemma-evaluate-and-more} の写像
$$
\text{Tot}(\Hom^\bullet(J^\bullet, I^\bullet) \otimes_R K^\bullet)
\longrightarrow
\Hom^\bullet(\Hom^\bullet(K^\bullet, J^\bullet), I^\bullet)
$$
を用いて定義する。これが $D(R)$ の三つの対象すべてに関して関手的であることの
証明は省略する。
\end{proof}

\begin{lemma}
\label{more-algebra-lemma-internal-hom-composition}
$R$ を環とする。$D(R)$ の $K, L, M$ が与えられると、
$K, L, M$ に関して関手的な $D(R)$ における標準射
$$
R\Hom_R(L, M) \otimes_R^\mathbf{L} R\Hom_R(K, L) \longrightarrow R\Hom_R(K, M)
$$
がある。
\end{lemma}

\begin{proof}
$M$ を表す K-入射複体 $I^\bullet$、$L$ を表す K-入射複体
$J^\bullet$、および $K$ を表す任意の $R$-加群の複体 $K^\bullet$ を選ぶ。
Lemma \ref{more-algebra-lemma-composition} により複体の写像
$$
\text{Tot}\left(
\Hom^\bullet(J^\bullet, I^\bullet) \otimes_R \Hom^\bullet(K^\bullet, J^\bullet)
\right)
\longrightarrow
\Hom^\bullet(K^\bullet, I^\bullet)
$$
がある。$R$-加群の複体
$\Hom^\bullet(J^\bullet, I^\bullet)$,
$\Hom^\bullet(K^\bullet, J^\bullet)$, および
$\Hom^\bullet(K^\bullet, I^\bullet)$ はそれぞれ
$R\Hom_R(L, M)$, $R\Hom_R(K, L)$, および $R\Hom_R(K, M)$ を表す。
K-平坦複体 $H^\bullet$ と擬同型
$H^\bullet \to \Hom^\bullet(K^\bullet, J^\bullet)$ を選べば、写像
$$
\text{Tot}\left(
\Hom^\bullet(J^\bullet, I^\bullet) \otimes_R H^\bullet
\right)
\longrightarrow
\text{Tot}\left(
\Hom^\bullet(J^\bullet, I^\bullet) \otimes_R \Hom^\bullet(K^\bullet, J^\bullet)
\right)
$$
があり、その始域は
$R\Hom_R(L, M) \otimes_R^\mathbf{L} R\Hom_R(K, L)$ を表す。
二つの表示された矢印を合成して望む写像を得る。
この構成が関手的であることの証明は省略する。
\end{proof}

\begin{lemma}
\label{more-algebra-lemma-internal-hom-diagonal-better}
$R$ を環とする。$D(R)$ の複体 $K, L, M$ が与えられると、
$K$, $L$, $M$ に関して関手的な $D(R)$ における標準射
$$
K \otimes_R^\mathbf{L} R\Hom_R(M, L)
\longrightarrow
R\Hom_R(M, K \otimes_R^\mathbf{L} L)
$$
がある。
\end{lemma}

\begin{proof}
$K$ を表す K-平坦複体 $K^\bullet$、$L$ を表す K-入射複体
$I^\bullet$ を選び、$M$ を表す任意の複体 $M^\bullet$ を選ぶ。
$J^\bullet$ が K-入射となるような擬同型
$\text{Tot}(K^\bullet \otimes_R I^\bullet) \to J^\bullet$ を選ぶ。
このとき写像
$$
\text{Tot}\left(
K^\bullet \otimes_R \Hom^\bullet(M^\bullet, I^\bullet)
\right)
\to
\Hom^\bullet(M^\bullet, \text{Tot}(K^\bullet \otimes_R I^\bullet))
\to
\Hom^\bullet(M^\bullet, J^\bullet)
$$
を用いる。ここで最初の写像は Lemma \ref{more-algebra-lemma-diagonal-better} の写像である。
\end{proof}

\begin{lemma}
\label{more-algebra-lemma-internal-hom-diagonal}
$R$ を環とする。$D(R)$ の複体 $K, L$ が与えられると、
$K$ と $L$ の両方に関して関手的な $D(R)$ における標準射
$$
K \longrightarrow R\Hom_R(L, K \otimes_R^\mathbf{L} L)
$$
がある。
\end{lemma}

\begin{proof}
これは Lemma \ref{more-algebra-lemma-internal-hom-diagonal-better} の特別な場合であるが、
直接にも証明する。$K$ を表す K-平坦複体 $K^\bullet$ と、
$L$ を表す任意の複体 $L^\bullet$ を選ぶ。
$J^\bullet$ が K-入射となるような擬同型
$\text{Tot}(K^\bullet \otimes_R L^\bullet) \to J^\bullet$ を選ぶ。
このとき写像
$$
K^\bullet \to
\Hom^\bullet(L^\bullet, \text{Tot}(K^\bullet \otimes_R L^\bullet))
\to \Hom^\bullet(L^\bullet, J^\bullet)
$$
を用いる。ここで最初の写像は Lemma \ref{more-algebra-lemma-diagonal} の写像である。
\end{proof}

\section{Perfect 複体}
\label{more-algebra-section-perfect}

\noindent
perfect 複体とは、有限 Tor 次元を持つ擬連接複体である。
これを定義として用いることはせず、環上の perfect 複体を次のように直接定義する。

\begin{definition}
\label{more-algebra-definition-perfect}
$R$ を環とする。
% P02-E0968: frozen source lines 18872--18873 use the malformed construction "Denote D(R) the derived category".
$R$-加群の Abel 圏の導来圏を $D(R)$ と書く。
\begin{enumerate}
\item $D(R)$ の対象 $K$ が有限射影 $R$-加群からなる有界複体と
擬同型であるとき、それは {\it perfect} であるという。
\item $M[0]$ が $D(R)$ の perfect 対象であるとき、$R$-加群 $M$ は
{\it perfect} であるという。
\end{enumerate}
\end{definition}

\noindent
例えば Noether 環上では、有限加群が perfect であることと有限射影次元を持つことは
同値である。Lemma \ref{more-algebra-lemma-perfect-module} および Algebra, Definition
\ref{algebra-definition-finite-proj-dim} を見よ。

\begin{lemma}
\label{more-algebra-lemma-perfect}
$K^\bullet$ を $D(R)$ の対象とする。次は同値である：
\begin{enumerate}
\item $K^\bullet$ は perfect である。
\item $K^\bullet$ は擬連接であり、有限 Tor 次元を持つ。
\end{enumerate}
(1) と (2) が成り立ち、$K^\bullet$ が $[a, b]$ に Tor 振幅を持つならば、
$K^\bullet$ は、$i \not \in [a, b]$ に対して $E^i = 0$ となる
有限射影 $R$-加群の複体 $E^\bullet$ と擬同型である。
\end{lemma}

\begin{proof}
(1) が (2) を含意することは明らかである。Lemmas
\ref{more-algebra-lemma-pseudo-coherent} および \ref{more-algebra-lemma-tor-amplitude} を見よ。
(2) が成り立ち、$K^\bullet$ が $[a, b]$ に Tor 振幅を持つと仮定する。
特に $i > b$ に対して $H^i(K^\bullet) = 0$ である。
$i > b$ に対して $F^i = 0$ となる有限自由 $R$-加群の複体
$F^\bullet$ と擬同型 $F^\bullet \to K^\bullet$ を選ぶ
（Lemma \ref{more-algebra-lemma-pseudo-coherent}）。
$E^\bullet = \tau_{\geq a}F^\bullet$ と置く。
$E^a$ は有限表示 $R$-加群であり、それ以外の $E^i$ は有限自由であることに注意する。
Lemma \ref{more-algebra-lemma-last-one-flat} により $E^a$ は平坦である。
従って Algebra, Lemma \ref{algebra-lemma-finite-projective} により
$E^a$ は有限射影である。
\end{proof}

\begin{lemma}
\label{more-algebra-lemma-perfect-module}
$M$ を環 $R$ 上の加群とする。次は同値である：
\begin{enumerate}
\item $M$ は perfect 加群である。
\item 各 $F_i$ が有限射影 $R$-加群であるような分解
$$
0 \to F_d \to \ldots \to F_1 \to F_0 \to M \to 0
$$
が存在する。
\end{enumerate}
\end{lemma}

\begin{proof}
(2) を仮定する。$E^{-i} = F_i$ である複体 $E^\bullet$ は $M[0]$ と擬同型である。
従って $M$ は perfect である。
逆に (1) を仮定する。Lemmas \ref{more-algebra-lemma-perfect} および
\ref{more-algebra-lemma-n-pseudo-module} により、各 $E^{-i}$ が有限自由 $R$-加群であるような
% P02-E0969: frozen source line 18934 omits the article in "find resolution".
分解 $E^\bullet \to M$ を見つけられる。
Lemma \ref{more-algebra-lemma-last-one-flat} により、十分大きいある $d$ に対して
% P02-E0970: frozen source lines 18930--18943 switch from E^{-i} and E^{-d+1} to Coker(E^{d-1}->E^d).
$F_d = \Coker(E^{d - 1} \to E^d)$ が平坦である。
Algebra, Lemma \ref{algebra-lemma-finite-projective} により
$F_d$ は有限射影である。従って
$$
0 \to F_d \to E^{-d+1} \to \ldots \to E^0 \to M \to 0
$$
が望む分解である。
\end{proof}

\begin{lemma}
\label{more-algebra-lemma-two-out-of-three-perfect}
$R$ を環とする。$(K^\bullet, L^\bullet, M^\bullet, f, g, h)$ を
$D(R)$ の完全三角とする。$K^\bullet, L^\bullet, M^\bullet$ のうち二つが
perfect ならば、残る一つも perfect である。
\end{lemma}

\begin{proof}
Lemmas \ref{more-algebra-lemma-perfect}, \ref{more-algebra-lemma-two-out-of-three-pseudo-coherent},
および \ref{more-algebra-lemma-cone-tor-amplitude} を組み合わせよ。
\end{proof}

\begin{lemma}
\label{more-algebra-lemma-summands-perfect}
$R$ を環とする。$K^\bullet \oplus L^\bullet$ が perfect ならば、
$K^\bullet$ と $L^\bullet$ も perfect である。
\end{lemma}

\begin{proof}
Lemmas \ref{more-algebra-lemma-perfect}, \ref{more-algebra-lemma-summands-pseudo-coherent},
および \ref{more-algebra-lemma-summands-tor-amplitude} から従う。
\end{proof}

\begin{lemma}
\label{more-algebra-lemma-complex-perfect-modules}
$R$ を環とする。$K^\bullet$ を perfect $R$-加群からなる有界複体とする。
このとき $K^\bullet$ は perfect 複体である。
\end{lemma}

\begin{proof}
有限複体の長さに関する帰納法から従う。Lemma
\ref{more-algebra-lemma-two-out-of-three-perfect} と愚直切断を用いよ。
\end{proof}

\begin{lemma}
\label{more-algebra-lemma-cohomology-perfect}
$R$ を環とする。$K^\bullet \in D^b(R)$ であり、そのすべてのコホモロジー加群が
perfect ならば、$K^\bullet$ は perfect である。
\end{lemma}

\begin{proof}
有限複体の長さに関する帰納法から従う。Lemma
\ref{more-algebra-lemma-two-out-of-three-perfect} と標準切断を用いよ。
\end{proof}

\begin{lemma}
\label{more-algebra-lemma-perfect-push-perfect}
$A \to B$ を環写像とする。$B$ が $A$-加群として perfect であると仮定する。
$K^\bullet$ を perfect な $B$-加群の複体とする。
このとき $K^\bullet$ は $A$-加群の複体として perfect である。
\end{lemma}

\begin{proof}
Lemma \ref{more-algebra-lemma-perfect} を用いると、これは擬連接加群と有限 Tor 次元を持つ加群に
関する対応する結果へ移される。それらの結果については Lemma
\ref{more-algebra-lemma-finite-tor-dimension-push-tor-amplitude} および Lemma
\ref{more-algebra-lemma-finite-push-pseudo-coherent} を見よ。
\end{proof}

\begin{lemma}
\label{more-algebra-lemma-pull-perfect}
$A \to B$ を環写像とする。$K^\bullet$ を perfect な $A$-加群の複体とする。
このとき $K^\bullet \otimes_A^{\mathbf{L}} B$ は perfect な
$B$-加群の複体である。
\end{lemma}

\begin{proof}
Lemma \ref{more-algebra-lemma-perfect} を用いると、これは擬連接加群と有限 Tor 次元を持つ加群に
関する対応する結果へ移される。それらの結果については
Lemma \ref{more-algebra-lemma-pull-tor-amplitude} および
Lemma \ref{more-algebra-lemma-pull-pseudo-coherent} を見よ。
\end{proof}

\begin{lemma}
\label{more-algebra-lemma-flat-base-change-perfect}
$A \to B$ を平坦な環写像とする。$M$ を perfect $A$-加群とする。
このとき $M \otimes_A B$ は perfect $B$-加群である。
\end{lemma}

\begin{proof}
Lemma \ref{more-algebra-lemma-perfect-module} により、仮定から $M$ は有限射影
% P02-E0971: frozen source line 19045 says finite projective R-modules although the ring in the lemma is A.
$R$-加群による有限分解 $F_\bullet$ を持つ。
$A \to B$ は平坦なので、複体 $F_\bullet \otimes_A B$ は
$B$ 上の有限射影加群による $M \otimes_A B$ の有限長分解である。
従って $M \otimes_A B$ は perfect である。
\end{proof}

\begin{lemma}
\label{more-algebra-lemma-tensor-perfect}
$R$ を環とする。$K$ と $L$ が $D(R)$ の perfect 対象ならば、
$K \otimes_R^\mathbf{L} L$ も perfect 対象である。
\end{lemma}

\begin{proof}
定義を用いて次のように証明できる。$K$, resp.\ $L$ を、有限射影
$R$-加群からなる有界複体 $K^\bullet$, resp.\ $L^\bullet$ で表してよい。
このとき $K \otimes_R^\mathbf{L} L$ は有界複体
$\text{Tot}(K^\bullet \otimes_R L^\bullet)$ によって表される。
この複体の各項は加群
% P02-E0972: frozen source lines 19062--19065 use M^a although the two complexes are K^bullet and L^bullet.
$M^a \otimes_R L^b$ の直和である。$M^a$ と $L^b$ は有限自由
$R$-加群の直和因子なので、$M^a \otimes_R L^b$ もそうである。
従って複体 $\text{Tot}(K^\bullet \otimes_R L^\bullet)$ の各項は
有限射影である。

\medskip\noindent
別の証明は、Lemma \ref{more-algebra-lemma-perfect} における perfect 複体の特徴付けと、
擬連接複体に関する対応する補題（Lemma \ref{more-algebra-lemma-tensor-pseudo-coherent}）および
Tor 振幅に関する補題（$A = B = R$ として用いた Lemma
\ref{more-algebra-lemma-push-tor-amplitude}）を用いて与えられる。
\end{proof}

\begin{lemma}
\label{more-algebra-lemma-glue-perfect}
$R$ を環とする。$f_1, \ldots, f_r \in R$ を単位イデアルを生成する元とする。
$K^\bullet$ を $R$-加群の複体とする。各 $i$ に対して複体
$K^\bullet \otimes_R R_{f_i}$ が perfect ならば、$K^\bullet$ は perfect である。
\end{lemma}

\begin{proof}
Lemma \ref{more-algebra-lemma-perfect} を用いると、これは擬連接加群と有限 Tor 次元を持つ加群に
関する対応する結果へ移される。それらの結果については
Lemma \ref{more-algebra-lemma-glue-tor-amplitude} および
Lemma \ref{more-algebra-lemma-glue-pseudo-coherent} を見よ。
\end{proof}

\begin{lemma}
\label{more-algebra-lemma-flat-descent-perfect}
$R$ を環とする。$K^\bullet$ を $R$-加群の複体とする。
$R \to R'$ を忠実平坦な環写像とする。複体
$K^\bullet \otimes_R R'$ が perfect ならば、$K^\bullet$ は perfect である。
\end{lemma}

\begin{proof}
Lemma \ref{more-algebra-lemma-perfect} を用いると、これは擬連接加群と有限 Tor 次元を持つ加群に
関する対応する結果へ移される。それらの結果については
Lemma \ref{more-algebra-lemma-flat-descent-tor-amplitude} および
Lemma \ref{more-algebra-lemma-flat-descent-pseudo-coherent} を見よ。
\end{proof}

\begin{lemma}
\label{more-algebra-lemma-regular-perfect}
$R$ を正則環とする。このとき
\begin{enumerate}
\item $R$-加群が perfect であることと有限 $R$-加群であることは同値である。
\item $R$-加群の複体 $K^\bullet$ が perfect であることと、
$K^\bullet \in D^b(R)$ であり各 $H^i(K^\bullet)$ が有限 $R$-加群であることは
同値である。
\end{enumerate}
\end{lemma}

\begin{proof}
定義により、任意の perfect $R$-加群は有限である。
逆に $M$ を有限 $R$-加群とする。
$F_i$ が有限自由 $R$-加群であるような分解
$$
\ldots \to F_2 \xrightarrow{d_2} F_1 \xrightarrow{d_1} F_0 \to M \to 0
$$
を選ぶ（Algebra, Lemma
\ref{algebra-lemma-resolution-by-finite-free}）。
$M_i = \Ker(d_i)$ と置く。
% P02-E0973: frozen source lines 19132--19134 use the malformed construction "Denote U_i ... the set".
$M_{i, \mathfrak p}$ が自由となる素数 $\mathfrak p$ の集合を
$U_i \subset \Spec(R)$ と書く。Algebra, Lemma
\ref{algebra-lemma-finitely-presented-localization-free} により $U_i$ は開である。
% P02-E0974: frozen source line 19135 reads "a exact sequence".
完全列 $0 \to M_{i + 1} \to F_{i + 1} \to M_i \to 0$ がある。
$\mathfrak p \in U_i$ ならば
$0 \to M_{i + 1, \mathfrak p} \to F_{i + 1, \mathfrak p} \to
M_{i, \mathfrak p} \to 0$ は分裂する。従って
$M_{i + 1, \mathfrak p}$ は有限射影であり、ゆえに自由である
（Algebra, Lemma \ref{algebra-lemma-finite-projective}）。
これは $U_i \subset U_{i + 1}$ を示す。
$\Spec(R) = \bigcup U_i$ であると主張する。
実際、任意の素イデアル $\mathfrak p$ に対して、正則局所環
$R_\mathfrak p$ は Algebra, Proposition
\ref{algebra-proposition-regular-finite-gl-dim} により有限大域次元を持つ。
従って、例えば Algebra, Lemma
\ref{algebra-lemma-independent-resolution} により、$i \gg 0$ に対して
$M_{i, \mathfrak p}$ は有限射影（従って自由）である。
$R$ のスペクトルは Noether なので
（Algebra, Lemma \ref{algebra-lemma-Noetherian-topology}）、
ある $n$ に対して $U_n = \Spec(R)$ であると結論できる。
すると Algebra, Lemma \ref{algebra-lemma-finite-projective} により
$M_n$ は射影 $R$-加群である。従って
% P02-E0975: frozen source lines 19128 and 19155--19159 omit F_0 from the final displayed resolution.
$$
0 \to M_n \to F_n \to \ldots \to F_1 \to M \to 0
$$
は有限射影加群による有界分解であり、ゆえに $M$ は perfect である。
これで (1) が示された。

\medskip\noindent
$K^\bullet$ を $R$-加群の複体とする。
$K^\bullet$ が perfect ならば、それは $D^b(R)$ に属し、有限射影
$R$-加群からなる有限複体と擬同型である。従って、$R$ は Noether なので、
各 $H^i(K^\bullet)$ は確かに有限 $R$-加群である。
逆に $K^\bullet$ が $D^b(R)$ に属し、各 $H^i(K^\bullet)$ が有限
$R$-加群であると仮定する。(1) により各 $H^i(K^\bullet)$ は perfect
$R$-加群であり、従って Lemma \ref{more-algebra-lemma-cohomology-perfect} により
$K^\bullet$ は perfect である。
\end{proof}

\begin{lemma}
\label{more-algebra-lemma-dual-perfect-complex}
$A$ を環とする。$K \in D(A)$ を perfect とする。このとき
$K^\vee = R\Hom_A(K, A)$ は perfect 複体であり、
$K \cong (K^\vee)^\vee$ である。$L \in D(A)$ に対して関手的な同型
$$
L \otimes_A^\mathbf{L} K^\vee = R\Hom_A(K, L)
\quad\text{and}\quad
H^0(L \otimes_A^\mathbf{L} K^\vee) = \Ext_A^0(K, L)
$$
がある。
\end{lemma}

\begin{proof}
$K$ を有限射影 $A$-加群からなる有界複体 $K^\bullet$ で表せる。
Lemma \ref{more-algebra-lemma-RHom-out-of-projective} により対象 $K^\vee$ は複体
$E^\bullet = \Hom^\bullet(K^\bullet, A)$ で表される。
$E^n = \Hom_A(K^{-n}, A)$ も有限射影であることに注意する。
従って $E^\bullet$ は有限射影加群からなる有界複体であり、
$K^\vee$ は perfect であると結論する。各次数で符号を除けば評価写像を用いる標準写像
$$
K^\bullet = \text{Tot}(\Hom^\bullet(A, A) \otimes_A K^\bullet)
\longrightarrow
\Hom^\bullet(\Hom^\bullet(K^\bullet, A), A)
$$
がある。Lemma \ref{more-algebra-lemma-evaluate-and-more} を見よ
（符号規則については Section \ref{more-algebra-section-sign-rules} を見よ）。
有限射影加群の二重双対はそれ自身なので、この写像は標準同型
$(K^\vee)^\vee \cong K$ を定める。

\medskip\noindent
第二の等式は、Lemma \ref{more-algebra-lemma-internal-hom}、Derived Categories, Lemma
\ref{derived-lemma-morphisms-from-projective-complex}、および Ext 群の定義から
第一の等式より従う。Derived Categories, Section
\ref{derived-section-ext} を見よ。
$L^\bullet$ を $L$ を表す $A$-加群の複体とする。
Section \ref{more-algebra-section-sign-rules} の項 (\ref{more-algebra-item-compatible}) により、
$A$-加群の複体の標準同型
$$
\text{Tot}(L^\bullet \otimes_A E^\bullet)
\longrightarrow
\Hom^\bullet(K^\bullet, L^\bullet)
$$
がある（ここでは Lemma \ref{more-algebra-lemma-left-dual-complex} を用いて、
$E^\bullet$ が複体の圏における $K^\bullet$ の左双対であることを見る）。
これは最初の表示された等式を示し、証明が完了する。
\end{proof}

\begin{remark}
\label{more-algebra-remark-dual-perfect-complex}
$A$ を環とする。$K \in D(A)$ を perfect とする。
$K$ が $[a, b]$ に Tor 振幅を持つならば、その双対 $K^\vee$ は
$[-b, -a]$ に Tor 振幅を持つ。これは Lemma \ref{more-algebra-lemma-perfect} の最後の主張と
Lemma \ref{more-algebra-lemma-dual-perfect-complex} の証明における $K^\vee$ の記述から従う。
\end{remark}

\begin{lemma}
\label{more-algebra-lemma-colim-and-lim-of-duals}
\begin{slogan}
perfect 対象の系に対する自明な双対性。
\end{slogan}
$A$ を環とする。$(K_n)_{n \in \mathbf{N}}$ を $D(A)$ の
perfect 対象の系とする。$K = \text{hocolim} K_n$ を導来余極限とする
（Derived Categories, Definition \ref{derived-definition-derived-colimit}）。
このとき $D(A)$ の任意の対象 $E$ に対して
$$
R\Hom_A(K, E) = R\lim E \otimes^\mathbf{L}_A K_n^\vee
$$
である。ここで $(K_n^\vee)$ は双対 perfect 複体からなる逆系である。
\end{lemma}

\begin{proof}
Lemma \ref{more-algebra-lemma-dual-perfect-complex} により
$R\lim E \otimes^\mathbf{L}_A K_n^\vee =
R\lim R\Hom_A(K_n, E)$
であり、これは完全三角
$$
R\lim R\Hom_A(K_n, E) \to
\prod R\Hom_A(K_n, E) \to
\prod R\Hom_A(K_n, E)
$$
に入る。同様に $K$ は完全三角
$\bigoplus K_n \to \bigoplus K_n \to K$ に入るので、
$\prod R\Hom_A(K_n, E) = R\Hom_A(\bigoplus K_n, E)$ を示せば十分である。
これは (\ref{more-algebra-equation-internal-hom}) と、導来テンソル積が直和と可換であることの
形式的帰結である。
\end{proof}

\begin{lemma}
\label{more-algebra-lemma-colimit-perfect-complexes}
$R = \colim_{i \in I} R_i$ を環のフィルター余極限とする。
\begin{enumerate}
\item $D(R)$ の perfect な $K$ が与えられると、ある $i \in I$ と
$D(R_i)$ の perfect な $K_i$ が存在し、$D(R)$ において
$K \cong K_i \otimes_{R_i}^\mathbf{L} R$ となる。
\item $0 \in I$ および $K_0, L_0 \in D(R_0)$ が与えられ、
$K_0$ が perfect であるとき、
$$
\Hom_{D(R)}(K_0 \otimes_{R_0}^\mathbf{L} R, L_0 \otimes_{R_0}^\mathbf{L} R) =
\colim_{i \geq 0}
\Hom_{D(R_i)}(K_0 \otimes_{R_0}^\mathbf{L} R_i,
L_0 \otimes_{R_0}^\mathbf{L} R_i)
$$
である。
\end{enumerate}
言い換えれば、$R$ 上の perfect 複体の三角圏は、
$R_i$ 上の perfect 複体の三角圏の余極限である。
\end{lemma}

\begin{proof}
Algebra, Lemmas \ref{algebra-lemma-module-map-property-in-colimit} および
\ref{algebra-lemma-colimit-category-fp-modules} の結果を以後断りなく用いる。
特にこれらの補題は、有限表示 $R$-加群の圏が有限表示 $R_i$-加群の圏の
余極限であることを述べる。有限射影加群は有限自由加群の直和因子として特徴付けられる
ので（Algebra, Lemma \ref{algebra-lemma-finite-projective}）、
有限射影加群の圏についても同じことが成り立つ。
$D(R)$ の perfect 対象の定義により、これは (1) を示す。

\medskip\noindent
(2) を示すため、$K_0$ を有限射影 $R_0$-加群からなる有界複体
$K_0^\bullet$ で表してよい。$L_0$ を K-平坦複体
$L_0^\bullet$ で表してよい（Lemma \ref{more-algebra-lemma-K-flat-resolution}）。
このとき Derived Categories, Lemma
\ref{derived-lemma-morphisms-from-projective-complex} により
$$
\Hom_{D(R)}(K_0 \otimes_{R_0}^\mathbf{L} R, L_0 \otimes_{R_0}^\mathbf{L} R) =
\Hom_{K(R)}(K_0^\bullet \otimes_{R_0} R, L_0^\bullet \otimes_{R_0} R)
$$
である。$R$ を $R_i$ で置き換えた $\Hom$ についても同様である。
右辺には有限個の項しか関与せず、証明の冒頭で引用した補題により
$$
\Hom_R(K_0^p \otimes_{R_0} R, L_0^q \otimes_{R_0} R) =
\colim_{i \geq 0}
\Hom_{R_i}(K_0^p \otimes_{R_0} R_i, L_0^q \otimes_{R_0} R_i)
$$
であり、さらにフィルター余極限は完全である
（Algebra, Lemma \ref{algebra-lemma-directed-colimit-exact}）。
従って (2) も成り立つ。
\end{proof}

\section{複体の持ち上げ}
\label{more-algebra-section-lifting-complexes}

\noindent
$R$ を環とし、$I \subset R$ をイデアルとする。
ここで考える持ち上げ問題は次のものである。
$D(R)$ の対象 $K$ と $R/I$-加群の複体 $E^\bullet$ が与えられ、
$E^\bullet$ は $D(R)$ において
$K \otimes_R^\mathbf{L} R/I$ を表すとする。
問：$E^\bullet$ を持ち上げ、$D(R)$ において $K$ を表す
$R$-加群の複体 $P^\bullet$ は存在するか？
一般にはこの問いへの答えは否定的だが、よい場合には何かを行える。
まず非輪状複体の持ち上げを論じる。

\begin{lemma}
\label{more-algebra-lemma-lift-acyclic-complex}
$R$ を環とし、$I \subset R$ をイデアルとする。
$\mathcal{P}$ を $R$-加群のクラスとする。次を仮定する：
\begin{enumerate}
\item 各 $P \in \mathcal{P}$ は射影 $R$-加群である。
\item $P_1 \in \mathcal{P}$ かつ $P_1 \oplus P_2 \in \mathcal{P}$ ならば、
$P_2 \in \mathcal{P}$ である。
\item
% P02-E0976: frozen source lines 19350--19351 use the malformed singular construction "if f..., P_1,P_2 ... is surjective".
$f : P_1 \to P_2$, $P_1, P_2 \in \mathcal{P}$ が $I$ を法として全射ならば、
$f$ は全射である。
\end{enumerate}
このとき、各項が $P \in \mathcal{P}$ に対する $P/IP$ の形である任意の
上に有界な非輪状複体 $E^\bullet$ に対し、各項が $\mathcal{P}$ に属し
$E^\bullet$ を持ち上げる上に有界な非輪状複体 $P^\bullet$ が存在する。
\end{lemma}

\begin{proof}
$i > b$ に対して $E^i = 0$ とする。$n$ と複体の射
$$
\xymatrix{
& & P^n \ar[r] \ar[d] & P^{n + 1} \ar[r] \ar[d] & \ldots \ar[r]
& P^b \ar[r] \ar[d] & 0 \ar[r] \ar[d] & \ldots \\
\ldots \ar[r] & E^{n - 1} \ar[r] &
E^n \ar[r] & E^{n + 1} \ar[r] & \ldots \ar[r] &
E^b \ar[r] & 0 \ar[r] & \ldots
}
$$
が与えられ、$P^i \in \mathcal{P}$、複体
$P^n \to P^{n + 1} \to \ldots \to P^b$ は次数 $\geq n + 1$ で非輪状、
かつ垂直写像は同型 $P^i/IP^i \to E^i$ を誘導すると仮定する。
この状況では、写像 $P^i \to P^{i + 1}$ が
$Z^i \oplus Z^{i + 1} \to Z^{i + 1} \to Z^{i + 1} \oplus Z^{i + 2}$
で与えられるような同型
$P^i = Z^i \oplus Z^{i + 1}$ を帰納的に選べる。
性質 (2) と帰納的議論により $Z^i \in \mathcal{P}$ である。
$P^{n - 1} \in \mathcal{P}$ と同型
$P^{n - 1}/IP^{n - 1} \to E^{n - 1}$ を選ぶ。
$P^{n - 1}$ は射影的であり、かつ
$Z^n/IZ^n = \Im(E^{n - 1} \to E^n)$ なので、写像
$P^{n - 1} \to E^{n - 1} \to E^n$ を写像
$P^{n - 1} \to Z^n$ に持ち上げられる。
性質 (3) により写像 $P^{n - 1} \to Z^n$ は全射である。
従って、ちょうど構成した $P^{n - 1}$ と写像を $P^n$ の左に加えることで、
図式の延長を得る。望む形の図式は $n > b$ に対して存在するので、
$n$ に関する帰納法により結論する。
\end{proof}

\begin{lemma}
\label{more-algebra-lemma-lift-complex}
$R$ を環とし、$I \subset R$ をイデアルとする。
$\mathcal{P}$ を $R$-加群のクラスとする。
$K \in D(R)$ とし、$E^\bullet$ を
$K \otimes_R^\mathbf{L} R/I$ を表す $R/I$-加群の複体とする。次を仮定する：
\begin{enumerate}
\item 各 $P \in \mathcal{P}$ は射影 $R$-加群である。
\item $P_1 \in \mathcal{P}$ かつ $P_1 \oplus P_2 \in \mathcal{P}$ であることと、
$P_1, P_2 \in \mathcal{P}$ であることは同値である。
\item $f : P_1 \to P_2$, $P_1, P_2 \in \mathcal{P}$ が $I$ を法として全射ならば、
$f$ は全射である。
\item $E^\bullet$ は上に有界で、$E^i$ はある $P \in \mathcal{P}$ に対する
$P/IP$ の形である。
\item $K$ は各項が $\mathcal{P}$ に属する上に有界な複体で表せる。
\end{enumerate}
このとき、各項が $\mathcal{P}$ に属する上に有界な複体 $P^\bullet$ で、
% P02-E0977: frozen source lines 19406--19408 omit "such that" before "P^bullet/IP^bullet isomorphic".
$P^\bullet/IP^\bullet$ が $E^\bullet$ と同型であり、
$D(R)$ において $K$ を表すものが存在する。
\end{lemma}

\begin{proof}
仮定 (5) により、$K$ を各項が $\mathcal{P}$ に属する上に有界な複体
$K^\bullet$ で表せる。このとき $K \otimes_R^\mathbf{L} R/I$ は
$K^\bullet/IK^\bullet$ で表される。(4) により $E^\bullet$ は射影
$R/I$-加群からなる上に有界な複体なので、擬同型
$\delta : E^\bullet \to K^\bullet/IK^\bullet$ を選べる
（Derived Categories, Lemma
\ref{derived-lemma-morphisms-from-projective-complex}）。
% P02-E0978: frozen source line 19419 omits the article in "be cone on delta".
$C^\bullet$ を $\delta$ のコーンとする
（Derived Categories, Definition \ref{derived-definition-cone}）。
加群 $C^i$ は直和 $K^i/IK^i \oplus E^{i + 1}$ であり、
従ってある $P \in \mathcal{P}$ に対する $P/IP$ の形である。
実際、(2) は特に $\mathcal{P}$ が直和を取ることで保たれることを述べる。
$C^\bullet$ は非輪状なので、Lemma \ref{more-algebra-lemma-lift-acyclic-complex} を適用して
% P02-E0979: frozen source line 19425 reads "find a acyclic lift".
$C^\bullet$ の非輪状な持ち上げ $A^\bullet$ を見つけられる。
複体 $A^\bullet$ は上に有界で、各項が $\mathcal{P}$ に属する。図式
$$
\xymatrix{
K^\bullet \ar@{..>}[r] \ar[d] & A^\bullet \ar[d] \\
K^\bullet/IK^\bullet \ar[r] & C^\bullet \ar[r] & E^\bullet[1]
}
$$
では、Derived Categories, Lemma
\ref{derived-lemma-morphisms-lift-projective} により、図式を可換にする点線の矢印を
見つけられる。以下で、(1), (2), (3) から、すべての $i$ に対して
$K^i \to A^i$ が直和因子の包含であることが従うと示す。
性質 (2) により $P^i = \Coker(K^i \to A^i)$ は $\mathcal{P}$ に属する。
従って $P^\bullet = \Coker(K^\bullet \to A^\bullet)[-1]$ と取ればよい。

\medskip\noindent
証明を完了するには次を示さなければならない：
$f : P_1 \to P_2$, $P_1, P_2 \in \mathcal{P}$ とし、
$P_1/IP_1 \to P_2/IP_2$ が分裂単射で、その余核がある
$P_3 \in \mathcal{P}$ に対する $P_3/IP_3$ の形ならば、$f$ は分裂単射である。
$E_i = P_i/IP_i$ と書く。このとき $E_2 = E_1 \oplus E_3$ である。
$P_2$ は射影的なので、写像 $E_2 \to E_3$ を持ち上げる写像
$g : P_2 \to P_3$ を選べる。条件 (3) により写像 $g$ は全射であり、
$P_3$ は射影的なので分裂する。$P_1' = \Ker(g)$ と置き、分裂
$P_2 = P'_1 \oplus P_3$ を選ぶ。(2) により $P'_1 \in \mathcal{P}$ である。
$g \circ f = 0$ であるとは分からないが、写像
$$
P_1 \xrightarrow{f} P_2 \xrightarrow{projection} P'_1
$$
を考えられる。$I$ を法とした合成は同型である。$P'_1$ は射影的なので、
$P_1 = T \oplus P'_1$ と分裂できる。$T = 0$ ならば証明は完了する。
実際、そのとき $P_2 \to P'_1$ は $f$ の分裂である。
(2) により $T \in \mathcal{P}$ である。
$I$ を法として計算すると $T/IT = 0$ である。
$0 \in \mathcal{P}$ なので（$\mathcal{P}$ の任意の $P$ の直和因子として）、
写像 $0 \to T$ は全射であり、望む通り $T = 0$ と結論する。
\end{proof}

\begin{lemma}
\label{more-algebra-lemma-lift-complex-projectives}
$R$ を環とし、$I \subset R$ をイデアルとする。
$E^\bullet$ を $R/I$-加群の複体とし、$K$ を $D(R)$ の対象とする。次を仮定する：
\begin{enumerate}
\item $E^\bullet$ は射影 $R/I$-加群からなる上に有界な複体である。
\item $D(R/I)$ において $K \otimes_R^\mathbf{L} R/I$ は
$E^\bullet$ で表される。
\item $I$ は冪零イデアルである。
\end{enumerate}
このとき、$D(R)$ において $K$ を表す射影 $R$-加群からなる上に有界な複体
$P^\bullet$ で、$P^\bullet \otimes_R R/I$ が $E^\bullet$ と同型となるものが
存在する。
\end{lemma}

\begin{proof}
すべての射影 $R$-加群からなるクラス $\mathcal{P}$ を用いて
Lemma \ref{more-algebra-lemma-lift-complex} を適用する。補題の性質 (1) と (2) は明らかである。
性質 (3) は Nakayama の補題
（Algebra, Lemma \ref{algebra-lemma-NAK}）から従う。
性質 (4) は射影 $R/I$-加群を射影 $R$-加群へ持ち上げられることから従う。
Algebra, Lemma \ref{algebra-lemma-lift-projective-module} を見よ。
(5) が成り立つことを見るには、$K$ が $D^{-}(R)$ に属することを示せば十分である。
$E^\bullet$ は上に有界なので、$K \otimes_R^\mathbf{L} R/I$ が
$D^{-}(R/I)$ に属することが与えられている。従ってこれは
Lemma \ref{more-algebra-lemma-check-tor-dimension-modulo-nilpotent} から従う。
\end{proof}

\begin{lemma}
\label{more-algebra-lemma-check-pseudo-coherent-modulo-nilpotent}
$R' \to R$ を、核が冪零イデアルである全射環写像とする。
$K' \in D(R')$ とし、$K = K' \otimes_{R'}^\mathbf{L} R$ と置く。
このとき $K$ が擬連接であることと $K'$ が擬連接であることは同値である。
\end{lemma}

\begin{proof}
一方向は Lemma \ref{more-algebra-lemma-pull-pseudo-coherent} から従う。
逆方向について $K$ が擬連接であると仮定する。
Lemma \ref{more-algebra-lemma-pseudo-coherent} により、$K$ を有限自由 $R$-加群からなる
上に有界な複体 $E^\bullet$ で表せる。
Lemma \ref{more-algebra-lemma-lift-complex-projectives} により、$K'$ を射影
$R'$-加群からなる上に有界な複体 $P^\bullet$ で、
$P^n \otimes_{R'} R = E^n$ を満たすものによって表せる。
Nakayama の補題により $P^n$ は有限自由であり、従って $K'$ も擬連接である。
\end{proof}

\begin{lemma}
\label{more-algebra-lemma-lift-complex-stably-frees}
$R$ を環とし、$I \subset R$ をイデアルとする。
$E^\bullet$ を $R/I$-加群の複体とし、$K$ を $D(R)$ の対象とする。次を仮定する：
\begin{enumerate}
\item $E^\bullet$ は有限生成安定自由 $R/I$-加群からなる上に有界な複体である。
\item $D(R/I)$ において $K \otimes_R^\mathbf{L} R/I$ は
$E^\bullet$ で表される。
\item
% P02-E0980: frozen source line 19524 switches from the object K to undefined K^bullet.
$K^\bullet$ は擬連接である。
\item $1 + I$ のすべての元は可逆である。
\end{enumerate}
このとき、$D(R)$ において $K$ を表す有限生成安定自由 $R$-加群からなる
上に有界な複体 $P^\bullet$ で、$P^\bullet \otimes_R R/I$ が
$E^\bullet$ と同型となるものが存在する。さらに $E^i$ が自由ならば
$P^i$ は自由である。
\end{lemma}

\begin{proof}
すべての有限生成安定自由 $R$-加群からなるクラス $\mathcal{P}$ を用いて
Lemma \ref{more-algebra-lemma-lift-complex} を適用する。補題の性質 (1) は明らかである。
性質 (2) は Lemma \ref{more-algebra-lemma-exact-category-stably-free} から従う。
性質 (3) は Nakayama の補題
（Algebra, Lemma \ref{algebra-lemma-NAK}）から従う。
性質 (4) は有限生成安定自由 $R/I$-加群を有限生成安定自由
$R$-加群へ持ち上げられることから従う。
Lemma \ref{more-algebra-lemma-lift-stably-free} を見よ。
% P02-E0981: frozen source line 19541 says "Part (5)" although the numbered hypotheses are consistently called properties.
部分 (5) は、擬連接複体を有限自由 $R$-加群からなる上に有界な複体で
表せることから成り立つ。補題の最後の主張は
Lemma \ref{more-algebra-lemma-isomorphic-finite-projective-lifts} から従う。
\end{proof}


\begin{lemma}
\label{more-algebra-lemma-lift-pseudo-coherent-from-residue-field}
$(R, \mathfrak m, \kappa)$ を局所環とする。
$K \in D(R)$ を擬連接とする。
$d_i = \dim_\kappa H^i(K \otimes_R^\mathbf{L} \kappa)$ と置く。
$d_i < \infty$ であり、ある $b \in \mathbf{Z}$ に対して
$i > b$ ならば $d_i = 0$ である。
% P02-E0982: frozen source lines 19552--19554 repeat consecutive sentence openings with "Then".
このとき、$D(R)$ において $K$ を表す複体
$$
\ldots \to
R^{\oplus d_{b - 2}} \to
R^{\oplus d_{b - 1}} \to
R^{\oplus d_b} \to 0 \to \ldots
$$
が存在する。さらに、この複体は同型を除いて一意である（！）。
\end{lemma}

\begin{proof}
$K \otimes_R^\mathbf{L} \kappa$ は $D(\kappa)$ の対象として擬連接であることに
注意する。Lemma \ref{more-algebra-lemma-pull-pseudo-coherent} を見よ。
従ってコホモロジー空間は有限次元で、ある上限より上で消える。
$D(\kappa)$ の任意の対象は、$D(\kappa)$ において微分が零である複体
$E^\bullet$ と同型である。特に $E^i \cong \kappa^{\oplus d_i}$ は有限自由である。
Lemma \ref{more-algebra-lemma-lift-complex-stably-frees} を適用して存在性を得る。

\medskip\noindent
% P02-E0983: frozen source lines 19576--19578 make "If we have two complexes ... representing K." a sentence fragment before the following "Then".
$F^\bullet$ と $G^\bullet$ が二つの複体で、$F^i$ と $G^i$ は階数
$d_i$ の自由加群であり、$K$ を表すとする。
このとき、同型 $F^\bullet \cong K \cong G^\bullet$ を表す複体の写像
$\beta : F^\bullet \to G^\bullet$ を選べる。Derived Categories, Lemma
\ref{derived-lemma-morphisms-from-projective-complex} を見よ。
誘導される複体の写像
$\beta \otimes 1 : F^\bullet \otimes_R^\mathbf{L} \kappa \to
G^\bullet \otimes_R^\mathbf{L} \kappa$
は複体の同型でなければならない。実際、
$F^\bullet \otimes_R^\mathbf{L} \kappa$ と
$G^\bullet \otimes_R^\mathbf{L} \kappa$ の微分は零である。
従って $\beta^i : F^i \to G^i$ は有限自由 $R$-加群の写像であり、
$\mathfrak m$ を法とした還元は同型である。
ゆえに $\beta^i$ は同型であり、結論を得る。
\end{proof}

\begin{lemma}
\label{more-algebra-lemma-lift-perfect-from-residue-field}
$R$ を環とし、$\mathfrak p \subset R$ を素イデアルとする。
$K \in D(R)$ を perfect とする。
$d_i =
\dim_{\kappa(\mathfrak p)} H^i(K \otimes_R^\mathbf{L} \kappa(\mathfrak p))$
と置く。$d_i < \infty$ であり、非零なものは有限個しかない。
% P02-E0984: frozen source lines 19599--19600 repeat consecutive sentence openings with "Then".
このとき、$f \in R$, $f \not \in \mathfrak p$ と、$D(R_f)$ において
$K \otimes_R^\mathbf{L} R_f$ を表す複体
$$
\ldots \to 0 \to R_f^{\oplus d_a} \to R_f^{\oplus d_{a + 1}} \to
\ldots \to
R_f^{\oplus d_{b - 1}} \to
R_f^{\oplus d_b} \to 0
\to \ldots
$$
が存在する。
\end{lemma}

\begin{proof}
$K \otimes_R^\mathbf{L} \kappa(\mathfrak p)$ は
$D(\kappa(\mathfrak p))$ の対象として perfect であることに注意する。
Lemma \ref{more-algebra-lemma-pull-perfect} を見よ。従って非零な $d_i$ は有限個しかなく、
すべて有限である。Lemma \ref{more-algebra-lemma-lift-pseudo-coherent-from-residue-field}
を適用すると、局所環 $R_\mathfrak p$ 上で望む形を持ち $K$ を表す複体を得る。
$f \in R$, $f \not \in \mathfrak p$ にわたって
$R_\mathfrak p = \colim R_f$ である
（Algebra, Lemma \ref{algebra-lemma-localization-colimit}）。
Lemma \ref{more-algebra-lemma-colimit-perfect-complexes} により結論する。
いくつかの詳細は省略する。
\end{proof}

\begin{lemma}
\label{more-algebra-lemma-compare-representatives-perfect}
$R$ を環とし、$\mathfrak p \subset R$ を素イデアルとする。
$M^\bullet$ と $N^\bullet$ を、$D(R)$ の同じ対象を表す有限射影
$R$-加群からなる有界複体とする。このとき、ある $f \in R$,
$f \not \in \mathfrak p$ が存在し、複体の同型（！）
$$
M^\bullet_f \oplus P^\bullet \cong N^\bullet_f \oplus Q^\bullet
$$
がある。ここで $P^\bullet$ と $Q^\bullet$ は自明な複体、すなわち
% P02-E0985: frozen source lines 19637--19640 duplicate "the form" in "complexes of the form the form".
$\ldots \to 0 \to R_f \xrightarrow{1} R_f \to 0 \to \ldots$
の形の複体（任意の次数に置く）の有限直和である。
\end{lemma}

\begin{proof}
局所化 $R_\mathfrak p$ 上で記述された型の同型があるならば、
$R_\mathfrak p = \colim R_f$ であること
（Algebra, Lemma \ref{algebra-lemma-localization-colimit}）を用いて、
ある $f$ に対する $R_f$ 上の同型へ降下できる。
従って $R$ は局所環で、$\mathfrak p$ は極大イデアルであると仮定してよい。
この場合、結果は perfect 対象を表す「極小」複体の一意性
（Lemma \ref{more-algebra-lemma-lift-pseudo-coherent-from-residue-field}）と、
任意の複体が自明な複体と極小な複体の直和であること
（Algebra, Lemma \ref{algebra-lemma-add-trivial-complex}）から従う。
\end{proof}

\begin{lemma}
\label{more-algebra-lemma-lift-complex-finite-projectives}
$R$ を環とし、$I \subset R$ をイデアルとする。
$E^\bullet$ を $R/I$-加群の複体とし、$K$ を $D(R)$ の対象とする。次を仮定する：
\begin{enumerate}
\item $E^\bullet$ は有限射影 $R/I$-加群からなる上に有界な複体である。
\item $D(R/I)$ において $K \otimes_R^\mathbf{L} R/I$ は
$E^\bullet$ で表される。
\item $K$ は擬連接である。
\item $(R, I)$ は Hensel 対である。
\end{enumerate}
このとき、$D(R)$ において $K$ を表す有限射影 $R$-加群からなる
上に有界な複体 $P^\bullet$ で、$P^\bullet \otimes_R R/I$ が
$E^\bullet$ と同型となるものが存在する。さらに $E^i$ が自由ならば
$P^i$ は自由である。
\end{lemma}

\begin{proof}
すべての有限射影 $R$-加群からなるクラス $\mathcal{P}$ を用いて
Lemma \ref{more-algebra-lemma-lift-complex} を適用する。
補題の性質 (1) と (2) は明らかである。
性質 (3) は Nakayama の補題
（Algebra, Lemma \ref{algebra-lemma-NAK}）から従う。
性質 (4) は有限射影 $R/I$-加群を有限射影 $R$-加群へ持ち上げられることから
従う。Lemma \ref{more-algebra-lemma-lift-finite-projective-module} を見よ。
性質 (5) は、擬連接複体を有限自由 $R$-加群からなる上に有界な複体で
表せることから成り立つ。従って Lemma \ref{more-algebra-lemma-lift-complex} を適用でき、
望む $P^\bullet$ を得る。補題の最後の主張は
Lemma \ref{more-algebra-lemma-isomorphic-finite-projective-lifts} から従う。
\end{proof}

\section{複体の分裂}
\label{more-algebra-section-splitting}

\noindent
この節では、環の導来圏の対象がその切断の直和となることを含意する条件を論じる。
我々の方法は、適切な仮定の下で次の補題を用いて標準完全三角
$$
\tau_{\leq i}K \to K \to \tau_{\geq i + 1}K \to (\tau_{\leq i}K)[1]
$$
を $D(R)$ において分裂させることである。Derived Categories, Remark
\ref{derived-remark-truncation-distinguished-triangle} を見よ。

\begin{lemma}
\label{more-algebra-lemma-splitting-unique}
$R$ を環とし、$K$ と $L$ を $D(R)$ の対象とする。
$L$ が $[a, b]$ に射影振幅を持つと仮定する。
例えば $L$ が $[a, b]$ に Tor 振幅を持つ perfect 対象ならばよい。
\begin{enumerate}
\item $i \geq a$ に対して $H^i(K) = 0$ ならば
$\Hom_{D(R)}(L, K) = 0$ である。
\item $i \geq a + 1$ に対して $H^i(K) = 0$ ならば、任意の完全三角
$K \to M \to L \to K[1]$ に対し、完全三角の写像と両立する
$D(R)$ における同型 $M \cong K \oplus L$ がある。
\item $i \geq a$ に対して $H^i(K) = 0$ ならば、(2) の同型が存在し一意である。
\end{enumerate}
\end{lemma}

\begin{proof}
$L$ が $[a, b]$ に射影振幅を持つという仮定は、$i \not \in [a, b]$ に対して
$L^i = 0$ となる射影 $R$-加群の複体 $L^\bullet$ で $L$ を表せることを意味する。
Definition \ref{more-algebra-definition-projective-dimension} を見よ。
$L$ が $[a, b]$ に Tor 振幅を持つ perfect 対象ならば、
$i \not \in [a, b]$ に対して $L^i = 0$ となる有限射影 $R$-加群の複体
$L^\bullet$ で $L$ を表せる。Lemma \ref{more-algebra-lemma-perfect} を見よ。
$i \geq a$ に対して $H^i(K) = 0$ ならば、$K$ は
$\tau_{\leq a - 1}K$ と擬同型である。従って $K$ を、
$i \geq a$ に対して $K^i = 0$ となる $R$-加群の複体
$K^\bullet$ で表せる。このとき Derived Categories, Lemma
\ref{derived-lemma-morphisms-from-projective-complex} により
$$
\Hom_{D(R)}(L, K) = \Hom_{K(R)}(L^\bullet, K^\bullet) = 0
$$
を得る。これで (1) が示された。(2) の仮定の下では、(1) により
$\Hom_{D(R)}(L, K[1]) = 0$ である。従って Derived Categories, Lemma
\ref{derived-lemma-split} により完全三角は分裂する。
(3) の一意性は (1) から従う。
\end{proof}

\begin{lemma}
\label{more-algebra-lemma-better-cut-complex-in-two}
$R$ を環とし、$\mathfrak p \subset R$ を素イデアルとする。
$K^\bullet$ を擬連接な $R$-加群の複体とする。
ある $i \in \mathbf{Z}$ に対して写像
$$
H^i(K^\bullet) \otimes_R \kappa(\mathfrak p)
\longrightarrow
H^i(K^\bullet \otimes_R^{\mathbf{L}} \kappa(\mathfrak p))
$$
が全射であると仮定する。このとき、ある $f \in R$,
$f \not \in \mathfrak p$ が存在し、
$\tau_{\geq i + 1}(K^\bullet \otimes_R R_f)$ は
$[i + 1, \infty]$ に Tor 振幅を持つ $D(R_f)$ の perfect 対象であり、
$D(R_f)$ における標準同型
$$
K^\bullet \otimes_R R_f \cong
\tau_{\leq i}(K^\bullet \otimes_R R_f) \oplus
\tau_{\geq i + 1}(K^\bullet \otimes_R R_f)
$$
がある。
\end{lemma}

\begin{proof}
この証明ではすべてのテンソル積を $R$ 上で取り、
$\kappa = \kappa(\mathfrak p)$ と書く。
$K^\bullet$ は有限自由 $R$-加群からなる上に有界な複体であると仮定してよい。
次数 $i$ で何が起こるか調べよう：
$$
\ldots \to K^{i - 1} \xrightarrow{d^{i - 1}} K^i \xrightarrow{d^i}
K^{i + 1} \to \ldots
$$
$0 \subset V \subset W \subset K^i \otimes \kappa$ を式
$$
V = \Im\left(
K^{i - 1} \otimes \kappa \to K^i \otimes \kappa
\right)
\quad\text{and}\quad
W = \Ker\left(
K^i \otimes \kappa \to K^{i + 1} \otimes \kappa
\right)
$$
で定義する。
$\dim(V) = r$, $\dim(W/V) = s$, および
$\dim(K^i \otimes \kappa/W) = t$ と置く。
$d^{i - 1}$ によって $V$ の基底へ写る
$x_1, \ldots, x_r \in K^{i - 1}$ を選べる。
仮定により、$W/V$ の基底へ写る
$y_1, \ldots, y_s \in \Ker(d^i)$ を選べる。
最後に、$K^i \otimes \kappa/W$ の基底へ写る
$z_1, \ldots, z_t \in K^i$ を選ぶ。すると元
$d^i(z_1), \ldots, d^i(z_t) \in K^{i + 1}$ は
$K^{i + 1} \otimes_R \kappa$ において線型独立である。
Algebra, Lemma \ref{algebra-lemma-cokernel-flat} により、
ある $f \in R$, $f \not \in \mathfrak p$ に対して $R$ を $R_f$ で
置き換えた後、次を仮定してよい：
\begin{enumerate}
\item
% P02-E0986: frozen source line 19796 applies d^i to x_a in K^{i-1}; the surrounding differential is d^{i-1}.
$d^i(x_a), y_b, z_c$ は $K^i$ の $R$-基底である。
\item $d^i(z_1), \ldots, d^i(z_t)$ は $K^{i + 1}$ において
$R$-線型独立である。
\item 商 $E^{i + 1} = K^{i + 1}/\sum Rd^i(z_c)$ は有限射影である。
\end{enumerate}
$d^i$ は $d^{i - 1}(x_a)$ と $y_b$ を零化するので、条件 (2) から
$E^{i + 1} = \Coker(d^i : K^i \to K^{i + 1})$ と分かる。従って
$$
\tau_{\geq i + 1}K^\bullet =
(\ldots \to 0 \to E^{i + 1} \to K^{i + 2} \to \ldots)
$$
は、ある $b$ に対して次数 $[i + 1, b]$ に置かれた有限射影加群からなる
有界複体である。従って $\tau_{\geq i + 1}K^\bullet$ は
$[i + 1, b]$ に振幅を持つ perfect 対象である。
$\tau_{\leq i}K^\bullet$ は次数 $> i$ にコホモロジーを持たないので、
完全三角
$$
\tau_{\leq i}K^\bullet \to K^\bullet \to \tau_{\geq i + 1}K^\bullet \to
(\tau_{\leq i}K^\bullet)[1]
$$
に Lemma \ref{more-algebra-lemma-splitting-unique} を適用して結論できる
（Derived Categories, Remark
\ref{derived-remark-truncation-distinguished-triangle}）。
\end{proof}

\begin{lemma}
\label{more-algebra-lemma-isolate-a-cohomology-group}
$R$ を環とし、$\mathfrak p \subset R$ を素イデアルとする。
$K^\bullet$ を擬連接な $R$-加群の複体とする。
ある $i \in \mathbf{Z}$ に対して写像
$$
H^i(K^\bullet) \otimes_R \kappa(\mathfrak p)
\longrightarrow
H^i(K^\bullet \otimes_R^{\mathbf{L}} \kappa(\mathfrak p))
\quad\text{and}\quad
H^{i - 1}(K^\bullet) \otimes_R \kappa(\mathfrak p)
\longrightarrow
H^{i - 1}(K^\bullet \otimes_R^{\mathbf{L}} \kappa(\mathfrak p))
$$
が全射であると仮定する。このとき、ある $f \in R$,
$f \not \in \mathfrak p$ が存在し、次が成り立つ：
\begin{enumerate}
\item $\tau_{\geq i + 1}(K^\bullet \otimes_R R_f)$ は
$[i + 1, \infty]$ に Tor 振幅を持つ $D(R_f)$ の perfect 対象である。
\item $H^i(K^\bullet)_f$ は有限自由 $R_f$-加群である。
\item $D(R_f)$ における標準直和分解
$$
K^\bullet \otimes_R R_f \cong
\tau_{\leq i - 1}(K^\bullet \otimes_R R_f) \oplus
H^i(K^\bullet)_f[-i] \oplus
\tau_{\geq i + 1}(K^\bullet \otimes_R R_f)
$$
がある。
\end{enumerate}
\end{lemma}

\begin{proof}
Lemma \ref{more-algebra-lemma-better-cut-complex-in-two} から (1) と、$D(R_f)$ における分裂
$K^\bullet \otimes_R R_f = \tau_{\leq i}K^\bullet \otimes_R R_f \oplus
\tau_{\geq i + 1}K^\bullet \otimes_R R_f$
を得る。Lemma \ref{more-algebra-lemma-better-cut-complex-in-two} を
$\tau_{\leq i}K^\bullet \otimes_R R_f$ にもう一度適用すると、
（$f$ を適切に選んだ後）$D(R_f)$ における分裂
% P02-E0987: frozen source line 19861 omits the required [-i] shift and line 19862 switches from K^bullet to undefined K.
$\tau_{\leq i}K^\bullet \otimes_R R_f =
\tau_{\leq i - 1}K^\bullet \otimes_R R_f \oplus H^i(K^\bullet)_f$
と、$H^i(K)_f$ が平坦な perfect 加群、すなわち有限射影であるという結論を得る。
% P02-E0988: frozen proof concludes only finite projective although statement (2) claims finite free; a further localization/trivialization step is not stated.
\end{proof}

\begin{lemma}
\label{more-algebra-lemma-cut-complex-in-two}
$R$ を環とし、$\mathfrak p \subset R$ を素イデアルとする。
$i \in \mathbf{Z}$ とする。$K^\bullet$ を擬連接な $R$-加群の複体で
$H^i(K^\bullet \otimes_R^{\mathbf{L}} \kappa(\mathfrak p)) = 0$
を満たすものとする。このとき、ある $f \in R$,
$f \not \in \mathfrak p$ と、$D(R_f)$ における標準直和分解
$$
K^\bullet \otimes_R R_f =
\tau_{\geq i + 1}(K^\bullet \otimes_R R_f) \oplus
\tau_{\leq i - 1}(K^\bullet \otimes_R R_f)
$$
が存在し、$\tau_{\geq i + 1}(K^\bullet \otimes_R R_f)$ は
$[i + 1, \infty]$ に Tor 振幅を持つ perfect 複体である。
\end{lemma}

\begin{proof}
これは頻繁に用いられる Lemma \ref{more-algebra-lemma-better-cut-complex-in-two} の特別な場合である。
直接の証明は次の通りである。
$K^\bullet$ は有限自由 $R$-加群からなる上に有界な複体であると仮定してよい。
次数 $i$ で何が起こるか調べよう：
$$
\ldots \to K^{i - 2} \to R^{\oplus l}
\to R^{\oplus m} \to R^{\oplus n} \to K^{i + 2} \to \ldots
$$
$A$ を $K^{i - 1} \to K^i$ に対応する $m \times l$ 行列とし、
$B$ を $K^i \to K^{i + 1}$ に対応する $n \times m$ 行列とする。
仮定は、$A \bmod \mathfrak p$ の階数が $r$ であり、
$B \bmod \mathfrak p$ の階数が $m - r$ であることを述べる。
言い換えれば、$A$ のある $r \times r$ 小行列式 $a$ が
$\mathfrak p$ に属さず、$B$ のある
$(m - r) \times (m - r)$-小行列式 $b$ が $\mathfrak p$ に属さない。
$f = ab$ と置く。$f$ を可逆化した後、直和分解
$K^{i - 1} = R^{\oplus l - r} \oplus R^{\oplus r}$,
$K^i = R^{\oplus r} \oplus R^{\oplus m - r}$,
$K^{i + 1} = R^{\oplus m - r} \oplus R^{\oplus n - m + r}$
を次のように見つけられる：加群写像 $K^{i - 1} \to K^i$ は
% P02-E0989: frozen source lines 19904 and 19907 twice say "kills of" instead of "kills".
$R^{\oplus l - r}$ を零化し、$R^{\oplus r}$ から $K^i$ の対応する
直和因子への同型を誘導する。また、加群写像 $K^i \to K^{i + 1}$ は
$R^{\oplus r}$ を零化し、$R^{\oplus m - r}$ から $K^{i + 1}$ の
対応する直和因子への同型を誘導する。従って $K^\bullet$ は
$$
\ldots \to K^{i - 2} \to R^{\oplus l - r}
\to 0 \to R^{\oplus n - m + r} \to K^{i + 2} \to \ldots
$$
と擬同型になり、すべて明らかである。
\end{proof}

\begin{lemma}
\label{more-algebra-lemma-split-using-ext-injective}
$R$ を環とし、$K \in D^-(R)$, $a \in \mathbf{Z}$ とする。
任意の単射 $R$-加群写像 $M \to M'$ に対して写像
$\Ext^{-a}_R(K, M) \to \Ext^{-a}_R(K, M')$ が単射であると仮定する。
このとき一意な直和分解
$K \cong \tau_{\leq a}K \oplus \tau_{\geq a + 1}K$
があり、ある $b$ に対して $\tau_{\geq a + 1}K$ は
$[a + 1, b]$ に射影振幅を持つ。
\end{lemma}

\begin{proof}
$D(R)$ における完全三角
$$
\tau_{\leq a}K \to K \to \tau_{\geq a + 1}K \to (\tau_{\leq a}K)[1]
$$
を考える。Derived Categories, Remark
\ref{derived-remark-truncation-distinguished-triangle} を見よ。
Derived Categories, Lemma \ref{derived-lemma-negative-exts} により
$\Ext^{-a}_R(\tau_{\leq a}K, M) = \Hom_R(H^a(K), M)$ および
$\Ext^{-a - 1}_R(\tau_{\leq a}K, M) = 0$ である。
従って $\Ext$ の長完全列は、$R$-加群 $M$ に関して関手的な完全列
$$
0 \to
\Ext^{-a}_R(\tau_{\geq a + 1}K, M) \to
\Ext^{-a}_R(K, M) \to
\Hom_R(H^a(K), M)
$$
を与える。ここで $I$ が入射 $R$-加群ならば、例えば Derived Categories, Lemma
\ref{derived-lemma-compute-ext-resolutions} により
$\Ext^{-a}_R(\tau_{\geq a + 1}K, I) = 0$ である。
すべての加群は入射加群へ単射を持つので、任意の $R$-加群 $M$ に対して
$\Ext^{-a}_R(\tau_{\geq a + 1}K, M) = 0$ と結論する。
Lemma \ref{more-algebra-lemma-projective-amplitude} により、ある $b$ に対して
$\tau_{\geq a + 1}K$ は $[a + 1, b]$ に射影振幅を持つ
（ここで $K$ が上に有界であることを用いる）。
Lemma \ref{more-algebra-lemma-splitting-unique} により分裂を得る。
\end{proof}

\begin{lemma}
\label{more-algebra-lemma-split-using-ext-zero}
$R$ を環とし、$K \in D^-(R)$, $a \in \mathbf{Z}$ とする。
任意の $R$-加群 $M$ に対して $\Ext^{-a}_R(K, M) = 0$ と仮定する。
このとき一意な直和分解
$K \cong \tau_{\leq a - 1}K \oplus \tau_{\geq a + 1}K$
があり、ある $b$ に対して $\tau_{\geq a + 1}K$ は
$[a + 1, b]$ に射影振幅を持つ。
\end{lemma}

\begin{proof}
Lemma \ref{more-algebra-lemma-split-using-ext-injective} により直和分解
$K \cong \tau_{\leq a}K \oplus \tau_{\geq a + 1}K$
があり、ある $b$ に対して $\tau_{\geq a + 1}K$ は
$[a + 1, b]$ に射影振幅を持つ。明らかに $H^a(K) = 0$ でなければならず、
$D(R)$ において $\tau_{\leq a}K = \tau_{\leq a - 1}K$ と結論する。
\end{proof}











\section{perfect 複体の判定}
\label{more-algebra-section-recognizing-perfect}

\noindent
ある複体が perfect であることを示すための補題をいくつか与える。

\begin{lemma}
\label{more-algebra-lemma-lift-bounded-pseudo-coherent-to-perfect}
$R$ を環、$\mathfrak p \subset R$ を素イデアルとする。
$K$ は擬コヒーレントかつ下に有界であるとする。次のように置く：
$d_i = \dim_{\kappa(\mathfrak p)}
H^i(K \otimes_R^\mathbf{L} \kappa(\mathfrak p))$.
ある $a \in \mathbf{Z}$ が存在して $i < a$ に対し $d_i = 0$ ならば、
ある $f \in R$, $f \not \in \mathfrak p$ と複体
$$
\ldots \to 0 \to R_f^{\oplus d_a} \to R_f^{\oplus d_{a + 1}} \to
\ldots \to
R_f^{\oplus d_{b - 1}} \to
R_f^{\oplus d_b} \to 0
\to \ldots
$$
が存在し、$D(R_f)$ において
$K \otimes_R^\mathbf{L} R_f$ を表す。
特に $K \otimes_R^\mathbf{L} R_f$ は perfect である。
\end{lemma}

\begin{proof}
$a$ を小さく取り直し、$i < a$ に対して
$H^i(K^\bullet) = 0$ でもあると仮定してよい。
Lemma \ref{more-algebra-lemma-cut-complex-in-two} により、ある
$f \in R$, $f \not \in \mathfrak p$ について $R$ を $R_f$ で置き換えた後、
$D(R)$ において
$K^\bullet = \tau_{\leq a - 1}K^\bullet \oplus
\tau_{\geq a}K^\bullet$ と書け、ここで
$\tau_{\geq a}K^\bullet$ は perfect である。
$i < a$ に対して $H^i(K^\bullet) = 0$ なので、
$D(R)$ において
$\tau_{\leq a - 1}K^\bullet = 0$ である。
従って $K^\bullet$ は perfect である。後は
Lemma \ref{more-algebra-lemma-lift-perfect-from-residue-field} を用いればよい。
\end{proof}

\begin{lemma}
\label{more-algebra-lemma-check-perfect-pointwise}
$R$ を環とし、$a, b \in \mathbf{Z}$ とする。
$K^\bullet$ を $R$-加群の擬コヒーレント複体とする。
次は同値である：
\begin{enumerate}
\item $K^\bullet$ は $[a, b]$ に Tor 振幅を持つ perfect 複体である、
\item 任意の素イデアル $\mathfrak p$ に対し、すべての
$i \not \in [a, b]$ について
$H^i(K^\bullet \otimes_R^{\mathbf{L}} \kappa(\mathfrak p)) = 0$ である、
\item 任意の極大イデアル $\mathfrak m$ に対し、すべての
$i \not \in [a, b]$ について
$H^i(K^\bullet \otimes_R^{\mathbf{L}} \kappa(\mathfrak m)) = 0$ である。
\end{enumerate}
\end{lemma}

\begin{proof}
(1) $\Rightarrow$ (2) $\Rightarrow$ (3) の証明は省略する。
(3) を仮定する。$i \in \mathbf{Z}$, $i \not \in [a, b]$ とする。
Lemma \ref{more-algebra-lemma-cut-complex-in-two} により、仮定から
$R$ のすべての極大イデアルに対して
$H^i(K^\bullet)_{\mathfrak m} = 0$ となる。
従って $H^i(K^\bullet) = 0$ である；Algebra, Lemma
\ref{algebra-lemma-characterize-zero-local} を見よ。
さらに Lemma \ref{more-algebra-lemma-cut-complex-in-two} は、任意の極大イデアル
$\mathfrak m$ に対して、ある $f \in R$, $f \not \in \mathfrak m$
が存在し、$K^\bullet \otimes_R R_f$ が $[a, b]$ に Tor 振幅を持つ
perfect 複体となることも示す。従って
Lemmas \ref{more-algebra-lemma-glue-perfect} and \ref{more-algebra-lemma-glue-tor-amplitude}
を適用して結論を得る。
\end{proof}

\begin{lemma}
\label{more-algebra-lemma-check-perfect-stalks}
$R$ を環とする。$K^\bullet$ を $R$-加群の擬コヒーレント複体とする。
次の条件を考える：
\begin{enumerate}
\item $K^\bullet$ は perfect である、
\item 任意の素イデアル $\mathfrak p$ に対して複体
$K^\bullet \otimes_R R_{\mathfrak p}$ は perfect である、
\item 任意の極大イデアル $\mathfrak m$ に対して複体
$K^\bullet \otimes_R R_{\mathfrak m}$ は perfect である、
\item 任意の素イデアル $\mathfrak p$ に対し、すべての
$i \ll 0$ について
$H^i(K^\bullet \otimes_R^{\mathbf{L}} \kappa(\mathfrak p)) = 0$ である、
\item 任意の極大イデアル $\mathfrak m$ に対し、すべての
$i \ll 0$ について
$H^i(K^\bullet \otimes_R^{\mathbf{L}} \kappa(\mathfrak m)) = 0$ である。
\end{enumerate}
常に次の含意が成り立つ：
% P02-E0990: frozen source line 20082 duplicates (3) in the implication chain.
$$
(1) \Rightarrow (2) \Leftrightarrow (3) \Leftrightarrow (3)
\Leftrightarrow (4) \Leftrightarrow (5)
$$
$K^\bullet$ が下に有界ならば、すべての条件は同値である。
\end{lemma}

\begin{proof}
Lemma \ref{more-algebra-lemma-pull-perfect} により (1) は (2) を含意する。
(2) $\Rightarrow$ (3) は直ちに分かる。任意の素イデアル
$\mathfrak p$ はある極大イデアル $\mathfrak m$ に含まれるので、
写像 $R_\mathfrak m \to R_\mathfrak p$ に
Lemma \ref{more-algebra-lemma-pull-perfect} を適用すれば、(3) が (2) を含意することが分かる。
剰余写像 $R_\mathfrak p \to \kappa(\mathfrak p)$ および
$R_\mathfrak m \to \kappa(\mathfrak m)$ に
Lemma \ref{more-algebra-lemma-pull-perfect} を適用すれば、(2) が (4) を、
(3) が (5) を含意することが分かる。

\medskip\noindent
$R$ は極大イデアル $\mathfrak m$ と剰余体 $\kappa$ を持つ局所環であると仮定する。
ある $a$ について $i < a$ ならば
$H^i(K^\bullet \otimes^\mathbf{L} \kappa) = 0$ であるとき、
$K$ が perfect であることを示す。これにより (4) が (2) を、
(5) が (3) を含意することが分かり、補題の前半が従う。
まず $i = a - 1$ として Lemma \ref{more-algebra-lemma-cut-complex-in-two} を適用すると、
$D(R)$ において
$K^\bullet = \tau_{\leq a - 1}K^\bullet \oplus \tau_{\geq a}K^\bullet$
となり、$\tau_{\geq a}K^\bullet$ は $[a, \infty]$ に含まれる
Tor 振幅を持つ perfect 複体である。証明を終えるには
$\tau_{\leq a - 1}K$ が零、すなわちそのコホモロジー群が零であることを
示せばよい。そうでないとし、
$M = H^i(\tau_{\leq a - 1}K)$ が零でないような最大の添字を $i$ とする。
$\tau_{\leq a - 1}K^\bullet$ は擬コヒーレントなので、
$M$ は有限 $R$-加群である
（Lemmas \ref{more-algebra-lemma-finite-cohomology} and
\ref{more-algebra-lemma-summands-pseudo-coherent}）。
従って Nakayama の補題（Algebra, Lemma \ref{algebra-lemma-NAK}）により
$M \otimes_R \kappa$ は零でない。これは
$$
H^i((\tau_{\leq a - 1}K^\bullet) \otimes_R^\mathbf{L} \kappa) =
H^i(K^\bullet \otimes_R^\mathbf{L} \kappa)
$$
が零でないことを含意し、矛盾である。

\medskip\noindent
同値な条件 (2) -- (5) が成り立ち、$K^\bullet$ が下に有界であると仮定する。
$i < a$ に対して $H^i(K^\bullet) = 0$ であるとする。
$R$ の極大イデアル $\mathfrak m$ を一つ取る。
ある $f \in R$, $f \not \in \mathfrak m$ が存在して
$K^\bullet \otimes_R^\mathbf{L} R_f$ が perfect であることを示せば十分である
（Lemma \ref{more-algebra-lemma-glue-perfect} and
Algebra, Lemma \ref{algebra-lemma-quasi-compact}）。
これは Lemma \ref{more-algebra-lemma-lift-bounded-pseudo-coherent-to-perfect} から従う。
\end{proof}

\begin{lemma}
\label{more-algebra-lemma-projective-amplitude-pseudo-coherent}
$R$ を環とし、$K$ を $D(R)$ の擬コヒーレント対象とする。
$a, b \in \mathbf{Z}$ とする。次は同値である：
\begin{enumerate}
\item $K$ は $[a, b]$ に射影振幅を持つ、
\item $K$ は $[a, b]$ に Tor 振幅を持つ perfect 複体である、
\item すべての有限表示 $R$-加群 $N$ とすべての
$i \not \in [-b, -a]$ に対して $\Ext^i_R(K, N) = 0$ である、
\item $n > b$ に対して $H^n(K) = 0$ であり、すべての有限表示
$R$-加群 $N$ とすべての $i > -a$ に対して
$\Ext^i_R(K, N) = 0$ である、
\item $n \not \in [a - 1, b]$ に対して $H^n(K) = 0$ であり、
すべての有限表示 $R$-加群 $N$ に対して
$\Ext^{-a + 1}_R(K, N) = 0$ である。
\end{enumerate}
\end{lemma}

\begin{proof}
Lemma \ref{more-algebra-lemma-perfect} の最後の主張により (2) は (1) を含意する。
(1) が成り立つとき、$K$ は $i \not \in [a, b]$ に対して
$P^i = 0$ となる射影加群の複体 $P^i$ で表せる。
射影加群は自由加群の直和因子として平坦なので、Lemma
\ref{more-algebra-lemma-tor-amplitude} により $K$ は $[a, b]$ に Tor 振幅を持つ。
従って Lemma \ref{more-algebra-lemma-perfect} により (2) が成り立つ。

\medskip\noindent
条件 (3), (4), (5) において、有限表示の $M$ に対して
Ext 群 $\Ext^i_R(K, M)$ が消滅するという仮定は、
Lemma \ref{more-algebra-lemma-pseudo-coherence-colimit-ext} および
Algebra, Lemma \ref{algebra-lemma-module-colimit-fp} により、
すべての $R$-加群 $M$ に対する消滅と同値である。
従って (1), (3), (4), (5) の同値性は
Lemma \ref{more-algebra-lemma-projective-amplitude} から従う。
\end{proof}

\noindent
% P02-E0991: frozen source line 20171 omits "is" in "The following lemma useful".
次の補題は多項式環 $B = A[x_1, \ldots, x_d]$ 上で
perfect 複体を見つけるのに有用である。

\begin{lemma}
\label{more-algebra-lemma-perfect-over-polynomial-ring}
$A \to B$ を環写像とする。$a, b \in \mathbf{Z}$, $d \geq 0$ とする。
$K^\bullet$ を $B$-加群の複体とし、次を仮定する：
\begin{enumerate}
\item 環写像 $A \to B$ は平坦である、
\item 任意の素イデアル $\mathfrak p \subset A$ に対し、環
$B \otimes_A \kappa(\mathfrak p)$ の大域次元は有限で $\leq d$ である、
\item $K^\bullet$ は $B$-加群の複体として擬コヒーレントである、
\item $K^\bullet$ は $A$-加群の複体として $[a, b]$ に Tor 振幅を持つ。
\end{enumerate}
このとき $K^\bullet$ は $B$-加群の複体として perfect であり、
$[a - d, b]$ に Tor 振幅を持つ。
\end{lemma}

\begin{proof}
$K^\bullet$ は有限自由 $B$-加群からなる上に有界な複体であると仮定してよい。
特に $K^\bullet$ は $A$-加群の複体として平坦であり、任意の
$A$-加群 $M$ に対して
$K^\bullet \otimes_A M = K^\bullet \otimes_A^{\mathbf{L}} M$ である。
$A$ の任意の素イデアル $\mathfrak p$ に対して、複体
$$
K^\bullet \otimes_A \kappa(\mathfrak p)
$$
は有限自由 $B \otimes_A \kappa(\mathfrak p)$-加群からなる
上に有界な複体であり、その $H^i$ は $i \in [a, b]$ の場合を除いて消滅する。
$B \otimes_A \kappa(\mathfrak p)$ の大域次元は $d$ なので、
Lemma \ref{more-algebra-lemma-finite-gl-dim-tor-dimension} により
$K^\bullet \otimes_A \kappa(\mathfrak p)$ は
$[a - d, b]$ に Tor 振幅を持つ。
$\mathfrak q$ を $\mathfrak p$ の上にある $B$ の素イデアルとする。
$K^\bullet \otimes_A \kappa(\mathfrak p)$ は自由
$B \otimes_A \kappa(\mathfrak p)$-加群からなる上に有界な複体なので、
\begin{align*}
K^\bullet \otimes_B^{\mathbf{L}} \kappa(\mathfrak q)
& = K^\bullet \otimes_B \kappa(\mathfrak q) \\
& = (K^\bullet \otimes_A \kappa(\mathfrak p))
\otimes_{B \otimes_A \kappa(\mathfrak p)} \kappa(\mathfrak q) \\
& = (K^\bullet \otimes_A \kappa(\mathfrak p))
\otimes^{\mathbf{L}}_{B \otimes_A \kappa(\mathfrak p)} \kappa(\mathfrak q)
\end{align*}
となる。従って上の議論から、$i \not \in [a - d, b]$ に対し
$H^i(K^\bullet \otimes_B^{\mathbf{L}} \kappa(\mathfrak q)) = 0$
である。Lemma \ref{more-algebra-lemma-check-perfect-pointwise} により結論を得る。
\end{proof}

\noindent
次の補題は Lemma \ref{more-algebra-lemma-perfect-over-polynomial-ring} の局所版である。
これは正則局所環上の perfect 複体を見つけるのに用いられる。

\begin{lemma}
\label{more-algebra-lemma-perfect-over-regular-local-ring}
$A \to B$ を局所環準同型とする。
$a, b \in \mathbf{Z}$, $d \geq 0$ とする。
$K^\bullet$ を $B$-加群の複体とし、次を仮定する：
\begin{enumerate}
\item 環写像 $A \to B$ は平坦である、
\item 環 $B/\mathfrak m_AB$ は次元 $d$ の正則環である、
\item $K^\bullet$ は $B$-加群の複体として擬コヒーレントである、
\item $K^\bullet$ は $A$-加群の複体として $[a, b]$ に Tor 振幅を持つ。
実際には、
$H^i(K^\bullet \otimes_A^\mathbf{L} \kappa(\mathfrak m_A))$
が零でないのは $i \in [a, b]$ の場合に限るという条件で十分である。
\end{enumerate}
このとき $K^\bullet$ は $B$-加群の複体として perfect であり、
$[a - d, b]$ に Tor 振幅を持つ。
\end{lemma}

\begin{proof}
(3) により、$K^\bullet$ は有限自由 $B$-加群からなる上に有界な複体であると
仮定してよい。次を計算する：
\begin{align*}
K^\bullet \otimes_B^{\mathbf{L}} \kappa(\mathfrak m_B)
& = K^\bullet \otimes_B \kappa(\mathfrak m_B) \\
& = (K^\bullet \otimes_A \kappa(\mathfrak m_A))
\otimes_{B/\mathfrak m_A B} \kappa(\mathfrak m_B) \\
& = (K^\bullet \otimes_A \kappa(\mathfrak m_A))
\otimes^{\mathbf{L}}_{B/\mathfrak m_A B} \kappa(\mathfrak m_B)
\end{align*}
% P02-E0992: frozen source line 20257 has the fragment "The first equality because".
最初の等式は $K^\bullet$ が平坦 $B$-加群からなる上に有界な複体であることによる。
二つ目の等式はテンソル積の基本的性質から従う。
三つ目の等式は、$K^\bullet \otimes_A \kappa(\mathfrak m_A) =
K^\bullet/ \mathfrak m_A K^\bullet$ が平坦
$B/\mathfrak m_A B$-加群からなる上に有界な複体だから成り立つ。
(1) により $K^\bullet$ は平坦 $A$-加群からなる上に有界な複体なので、
仮定 (4) により複体 $K^\bullet \otimes_A \kappa(\mathfrak m_A)$ の
コホモロジー加群 $H^i$ が零でないのは $i \in [a, b]$ の場合に限る。
従って Example \ref{more-algebra-example-cohomology-complex-tensored} のスペクトル系列と、
$B/\mathfrak m_AB$ が有限大域次元 $d$ を持つという事実
（(2) および Algebra, Proposition
\ref{algebra-proposition-regular-finite-gl-dim} による）から、
$H^j(K^\bullet \otimes_B^{\mathbf{L}} \kappa(\mathfrak m_B))$
は $j \not \in [a - d, b]$ に対して零である。
Lemma \ref{more-algebra-lemma-check-perfect-pointwise} により証明が完了する。
\end{proof}






\section{perfect 複体の特徴づけ}
\label{more-algebra-section-perfect-compact}

\noindent
この節では、perfect 複体が環の導来圏のコンパクト対象にほかならないことを証明する。
まず次を示す。

\begin{lemma}
\label{more-algebra-lemma-perfect-ring-classical-generator}
$R$ を環とする。perfect 対象からなる充満部分圏
$D_{perf}(R) \subset D(R)$ は、$R = R[0]$ を含む最小の厳密充満で
飽和した三角部分圏である。言い換えれば
$D_{perf}(R) = \langle R \rangle$ である。
特に、$R$ は $D_{perf}(R)$ の古典的生成対象である。
\end{lemma}

\begin{proof}
主張の意味については、Derived Categories, Definitions
\ref{derived-definition-saturated} and
\ref{derived-definition-generators} を見よ。
Lemmas \ref{more-algebra-lemma-two-out-of-three-perfect} and
\ref{more-algebra-lemma-summands-perfect} において、
$D_{perf}(R) \subset D(R)$ が $D(R)$ の厳密充満で飽和した
三角部分圏であることを示した。もちろん $R \in D_{perf}(R)$ である。

\medskip\noindent
$\langle R \rangle = \bigcup \langle R \rangle_n$ であることを思い出そう。
証明を終えるため、$M \in D_{perf}(R)$ が
$$
\ldots \to 0 \to M^a \to M^{a + 1} \to \ldots \to M^b \to 0 \to \ldots
$$
で表され、$M^i$ が有限射影であるならば、
$M \in \langle R \rangle_{b - a + 1}$ であることを示す。
$b - a$ に関する帰納法を用いる。
定義により $\langle R \rangle_1$ は、任意の次数に置かれた有限射影
$R$-加群をすべて含む；従って帰納法の初期段階
$b - a = 0$ は処理できる。一般には完全三角
% P02-E0993: frozen source line 20318 uses M_b although the complex terms are M^i.
$$
M_b[-b] \to M^\bullet \to \sigma_{\leq b - 1}M^\bullet \to M_b[-b + 1]
$$
を考える。帰納法により、切断複体
$\sigma_{\leq b - 1}M^\bullet$ は $\langle R \rangle_{b - a}$ に属し、
$M_b[-b]$ は $\langle R \rangle_1$ に属する。
従って定義により
$M^\bullet \in \langle R \rangle_{b - a + 1}$ である。
\end{proof}

\noindent
$R$ を環とする。$D(R)$ は複体の直和を取ることで与えられる直和を持つことを
思い出そう；Derived Categories, Lemma
\ref{derived-lemma-direct-sums} を見よ。
この節の補題では、以下これを断りなく用いる。

\begin{lemma}
\label{more-algebra-lemma-commutes-with-countable-sums}
$R$ を環とする。$K \in D(R)$ を、任意の可算個の対象
$E_n \in D(R)$ に対して標準写像
$$
\bigoplus \Hom_{D(R)}(K, E_n) \longrightarrow \Hom_{D(R)}(K, \bigoplus E_n)
$$
が全単射となるような対象とする。このとき、$\mathbf{N}$ 上の任意の複体系
$L_n^\bullet$ に対して
$$
\colim \Hom_{D(R)}(K, L^\bullet_n) \longrightarrow \Hom_{D(R)}(K, L^\bullet)
$$
は全単射である。ここで $L^\bullet$ は項別余極限、すなわちすべての
$m \in \mathbf{Z}$ に対して $L^m = \colim L_n^m$ である。
\end{lemma}

\begin{proof}
複体の短完全列
$$
0 \to \bigoplus L_n^\bullet \to \bigoplus L_n^\bullet \to L^\bullet \to 0
$$
を考える。ここで最初の写像は、遷移写像
$t_n : L_n^\bullet \to L_{n + 1}^\bullet$ を用いて、
% P02-E0994: frozen source line 20352 says "in degree n" although n indexes summands.
次数 $n$ では $1 - t_n$ で与えられる。
Derived Categories, Lemma
\ref{derived-lemma-derived-canonical-delta-functor} により、
これは $D(R)$ における完全三角である。
ホモロジー的関手 $\Hom_{D(R)}(K, -)$ を適用する；Derived Categories,
Lemma \ref{derived-lemma-representable-homological} を見よ。
従って次の長完全コホモロジー列を得る：
% P02-E0995: frozen source line 20359 has "Thus a long exact cohomology sequence" without a verb.
$$
\xymatrix{
& \ldots \ar[r] & \Hom_{D(R)}(K, \colim L^\bullet_n[-1]) \ar[lld] \\
\Hom_{D(R)}(K, \bigoplus L^\bullet_n) \ar[r] &
\Hom_{D(R)}(K, \bigoplus L^\bullet_n) \ar[r] &
\Hom_{D(R)}(K, \colim L^\bullet_n) \ar[lld] \\
\Hom_{D(R)}(K, \bigoplus L^\bullet_n[1]) \ar[r] & \ldots
}
$$
仮定により $\Hom_{D(R)}(K, \bigoplus L^\bullet_n)$ は
$\bigoplus \Hom_{D(R)}(K, L^\bullet_n)$ に等しいので、図の各行の
最初の写像は単射である（この写像を $1 - t_n$ が誘導する写像の和として
具体的に記述すれば分かる）。従って
$\Hom_{D(R)}(K, \colim L^\bullet_n)$ は上の図の中央の行の
最初の写像の余核であり、これは示すべきことである。
\end{proof}

\noindent
次の命題は、perfect 複体を導来圏のコンパクト対象
（Derived Categories, Definition \ref{derived-definition-compact-object}）
として特徴づけるもので、さまざまな箇所に現れる。例えば
\cite[proof of Proposition 6.3]{Rickard}（これは有界な場合を扱う）、
\cite[Theorem 2.4.3]{TT}（主張は厳密には一致しない）、および
\cite[Proposition 6.4]{Bokstedt-Neeman}（ひどい記法規約に注意せよ）
を見よ。

\begin{proposition}
\label{more-algebra-proposition-perfect-is-compact}
$R$ を環とする。$D(R)$ の対象 $K$ に対して次は同値である：
\begin{enumerate}
\item $K$ は perfect である、
\item $K$ は $D(R)$ のコンパクト対象である。
\end{enumerate}
\end{proposition}

\begin{proof}
$K$ は perfect であると仮定する。すなわち、Definition
\ref{more-algebra-definition-perfect} により、$K$ は有限射影加群からなる有界複体
$P^\bullet$ と擬同型である。
$E_i$ が複体 $E_i^\bullet$ で表されるならば、
$\bigoplus E_i$ は次数 $n$ の項が $\bigoplus E_i^n$ である複体で表される。
一方、すべての $n$ に対して $P^n$ は射影なので、任意の
$R$-加群の複体 $K^\bullet$ に対して
$\Hom_{D(R)}(P^\bullet, K^\bullet) = \Hom_{K(R)}(P^\bullet, K^\bullet)$
である；Derived Categories, Lemma
\ref{derived-lemma-morphisms-from-projective-complex} を見よ。
従って $\Hom_{D(R)}(P^\bullet, E^\bullet)$ は複体
$$
\prod \Hom_R(P^n, E^{n - 1}) \to
\prod \Hom_R(P^n, E^n) \to
\prod \Hom_R(P^n, E^{n + 1}).
$$
のコホモロジーである。$P^\bullet$ は有界なので、上の複体の
$\prod$ 記号を $\bigoplus$ 記号で置き換えてよい。
各 $P^n$ は有限 $R$-加群なので、すべての $n, m$ に対して
$\Hom_R(P^n, \bigoplus_i E_i^m) = \bigoplus_i \Hom_R(P^n, E_i^m)$
である。これらを合わせると、Derived Categories, Definition
\ref{derived-definition-compact-object} の写像が全単射であることが分かる。

\medskip\noindent
逆に、$K$ はコンパクトであると仮定する。
$K$ を複体 $K^\bullet$ で表し、標準切断を用いた写像
$$
K^\bullet
\longrightarrow
\bigoplus\nolimits_{n \geq 0} \tau_{\geq n} K^\bullet
$$
を考える；Homology, Section \ref{homology-section-truncations} を見よ。
各次数において右辺の直和は有限なので、この写像は意味を持つ。
仮定により、この写像は有限直和を経由する。
従って少なくとも一つの $n$ に対して
$K \to \tau_{\geq n} K$ は零、すなわち $K$ は $D^{-}(R)$ に属する。

\medskip\noindent
$K \in D^{-}(R)$ であり、すべての $R$-加群は自由加群の商なので、
$K$ を自由 $R$-加群からなる上に有界な複体 $K^\bullet$ で表せる；
Derived Categories, Lemma
\ref{derived-lemma-subcategory-left-resolution} を見よ。
愚直切断を用いると
$$
K^\bullet = \bigcup\nolimits_{n \leq 0} \sigma_{\geq n}K^\bullet
$$
である；Homology, Section \ref{homology-section-truncations} を見よ。
従って Lemma \ref{more-algebra-lemma-commutes-with-countable-sums} により、
$1 : K^\bullet \to K^\bullet$ は $D(R)$ において
$\sigma_{\geq n}K^\bullet \to K^\bullet$ を経由する。
従って、ある有界で各項が自由 $R$-加群である複体
$L^\bullet$ について、$1 : K^\bullet \to K^\bullet$ は $D(R)$ において
$$
K^\bullet \xrightarrow{\varphi} L^\bullet \xrightarrow{\psi} K^\bullet
$$
と分解する。$i \not \in [a, b]$ に対して $L^i = 0$ とする。
以下 $a, b$ を固定する。上のような $1_{K^\bullet}$ の分解で、
% P02-E0996: frozen source line 20460 reads "with L^i is finite free".
$i < c$ に対して $L^i$ が有限自由となるものを見つけられるような
$\leq b + 1$ の最大の整数を $c$ とする。
帰納法によって $c = b + 1$ を示す。実際、
$L^c = \bigoplus_{\lambda \in \Lambda} R$ と書く。
$L^{c - 1}$ は有限自由なので、$L^{c - 1} \to L^c$ が
$\bigoplus_{\lambda \in \Lambda'} R \subset L^c$ を経由するような有限部分集合
$\Lambda' \subset \Lambda$ を見つけられる。複体の写像
$$
\pi :
L^\bullet
\longrightarrow
(\bigoplus\nolimits_{\lambda \in \Lambda \setminus \Lambda'} R)[-c]
$$
で、次数 $c$ において $\Lambda \setminus \Lambda'$ に対応する因子への
射影で与えられるものを考える。
$K$ に関する仮定により、必要なら $\Lambda'$ をより大きな有限部分集合で
置き換えて、$D(R)$ において $\pi \circ \varphi = 0$ と仮定してよい。
$(L')^\bullet \subset L^\bullet$ を $\pi$ の核とする。
$\pi$ は全射なので複体の短完全列が得られ、これは $D(R)$ における
完全三角を与える（Derived Categories, Lemma
\ref{derived-lemma-derived-canonical-delta-functor} を見よ）。
$\Hom_{D(R)}(K, -)$ はホモロジー的であり（Derived Categories, Lemma
\ref{derived-lemma-representable-homological} を見よ）、
$\pi \circ \varphi = 0$ なので、$D(R)$ における射
$\varphi' : K^\bullet \to (L')^\bullet$ で、その
$(L')^\bullet \to L^\bullet$ との合成が $\varphi$ となるものを見つけられる。
$\psi'$ を $\psi$ と $(L')^\bullet \to L^\bullet$ の合成に等しいものと
置けば、新しい分解を得る。$(L')^\bullet$ は次数 $c$ を除いて
$L^\bullet$ と一致し、$(L')^c = \bigoplus_{\lambda \in \Lambda'} R$
なので、帰納段階が示された。

\medskip\noindent
前段落の議論から、$D(R)$ において $1_K : K \to K$ は
$$
K \xrightarrow{\varphi} L \xrightarrow{\psi} K
$$
と分解し、ここで $L$ は有限自由 $R$-加群からなる有限複体で表せる。
特に $L$ は perfect である。
$e = \varphi \circ \psi \in \text{End}_{D(R)}(L)$ は冪等であることに注意する。
Derived Categories, Lemma
\ref{derived-lemma-projectors-have-images-triangulated} により
$L = \Ker(e) \oplus \Ker(1 - e)$ である。
写像 $\varphi : K \to L$ は $D(R)$ における
$\Ker(1 - e)$ との同型を誘導する。
従って Lemma \ref{more-algebra-lemma-summands-perfect} により、ついに
$K$ は perfect であると結論する。
\end{proof}

\begin{lemma}
\label{more-algebra-lemma-perfect-modulo-nilpotent-ideal}
$R$ を環とし、$I \subset R$ をイデアルとする。
$K$ を $D(R)$ の対象とする。次を仮定する：
\begin{enumerate}
\item $K \otimes_R^\mathbf{L} R/I$ は $D(R/I)$ において perfect である、
\item $I$ は冪零イデアルである。
\end{enumerate}
このとき $K$ は $D(R)$ において perfect である。
\end{lemma}

\begin{proof}
$K \otimes_R^\mathbf{L} R/I$ を表す有限射影 $R/I$-加群からなる有限複体
$\overline{P}^\bullet$ を選ぶ；Definition
\ref{more-algebra-definition-perfect} を見よ。
Lemma \ref{more-algebra-lemma-lift-complex-projectives} により、
$K$ を表し $\overline{P}^\bullet = P^\bullet/IP^\bullet$ を満たす
射影 $R$-加群の複体 $P^\bullet$ が存在する。
Nakayama の補題（Algebra, Lemma \ref{algebra-lemma-NAK}）から、
$P^\bullet$ は有限射影 $R$-加群からなる有限複体であることが従う。
\end{proof}

\begin{lemma}
\label{more-algebra-lemma-perfect-modulo-two-ideals}
$R$ を環とし、$I, J \subset R$ をイデアルとする。
$K$ を $D(R)$ の対象とし、次を仮定する：
\begin{enumerate}
\item $K \otimes_R^\mathbf{L} R/I$ は $D(R/I)$ において perfect である、
\item $K \otimes_R^\mathbf{L} R/J$ は $D(R/J)$ において perfect である。
\end{enumerate}
このとき $K \otimes_R^\mathbf{L} R/IJ$ は $D(R/IJ)$ において
perfect である。
\end{lemma}

\begin{proof}
% P02-E0997: frozen source line 20543 reads "we may assume replace".
$R$ を $R/IJ$ で、$K$ を $K \otimes_R^\mathbf{L} R/IJ$ で置き換えてよい。
すると $R \to R/(I \cap J)$ は核の二乗が零である全射である。
従って Lemma \ref{more-algebra-lemma-perfect-modulo-nilpotent-ideal} により、
$K \otimes_R^\mathbf{L} R/(I \cap J)$ が perfect であることを示せば十分である。
従って $I \cap J = 0$ と仮定してよい。

\medskip\noindent
$I \cap J = 0$ の場合に補題を証明する。まず、
Lemma \ref{more-algebra-lemma-K-flat-resolution} により、$K$ を各 $K^n$ が平坦である
K-flat 複体 $K^\bullet$ で表してよい。すると複体の短完全列
$$
0 \to
K^\bullet \to K^\bullet/IK^\bullet \oplus K^\bullet/JK^\bullet \to
K^\bullet/(I + J)K^\bullet \to 0
$$
を得る。導来テンソル積の構成により、
$K^\bullet/IK^\bullet$ は $K \otimes^\mathbf{L}_R R/I$ を表す。
$K^\bullet/JK^\bullet$ および $K^\bullet/(I + J)K^\bullet$ についても同様である。
$K^\bullet/(I + J)K^\bullet$ は $R/(I + J)$-加群の perfect 複体である；
Lemma \ref{more-algebra-lemma-pull-perfect} を見よ。
% P02-E0998: frozen source lines 20566-20567 have a stray comma and repeated "and".
従って複体 $K^\bullet/IK^\bullet$, $K^\bullet/JK^\bullet$ および
$K^\bullet/(I + J)K^\bullet$ の零でないコホモロジー群は有限個である
（perfect 複体は有限 Tor 振幅を持つ；Lemma
\ref{more-algebra-lemma-perfect} を見よ）。
上に表示した複体の短完全列に付随するコホモロジー長完全列により、
$K \in D^b(R)$ と結論する。
% P02-E0999: frozen source line 20572 says "we assume" where the resolution lemma permits "may assume".
特に $K^\bullet$ は自由 $R$-加群からなる上に有界な複体であると仮定する
（Derived Categories, Lemma
\ref{derived-lemma-subcategory-left-resolution} を見よ）。

\medskip\noindent
次に Proposition \ref{more-algebra-proposition-perfect-is-compact} の判定法を用いて
$K$ が perfect であることを示す。そこで、集合 $J$ で添字づけられた
対象の族 $E_j \in D(R)$ を取る。
$E_j$ を表す、各項が平坦な複体 $E_j^\bullet$ を選ぶ；例えば
Lemma \ref{more-algebra-lemma-K-flat-resolution} を見よ。複体の列
$$
0 \to
E_j^\bullet \to
E_j^\bullet/IE_j^\bullet \oplus E_j^\bullet/JE_j^\bullet \to
E_j^\bullet/(I + J)E_j^\bullet \to 0
$$
が短完全であることは明らかである。直和を取ると、同様の短完全列
$$
0 \to
\bigoplus E_j^\bullet \to
\bigoplus E_j^\bullet/IE_j^\bullet \oplus E_j^\bullet/JE_j^\bullet \to
\bigoplus E_j^\bullet/(I + J)E_j^\bullet \to 0
$$
を得る（$- \otimes_R R/I$ は直和と可換であることに注意せよ）。
この短完全列は $D(R)$ における完全三角を定める；Derived Categories,
Lemma \ref{derived-lemma-derived-canonical-delta-functor} を見よ。
ホモロジー的関手 $\Hom_{D(R)}(K, -)$ を適用し
（Derived Categories, Lemma
\ref{derived-lemma-representable-homological} を見よ）、可換図
% P02-E1000: frozen source lines 20605-20609 place [-1] outside Hom rather than shifting the target complex.
$$
\xymatrix{
\bigoplus \Hom_{D(R)}(K^\bullet, E_j^\bullet/(I + J))[-1] \ar[r] \ar[d] &
\Hom_{D(R)}(K^\bullet, \bigoplus E_j^\bullet/(I + J))[-1] \ar[d] \\
\bigoplus \Hom_{D(R)}(K^\bullet, E_j^\bullet/I \oplus E_j^\bullet/J)[-1]
\ar[r] \ar[d] &
\Hom_{D(R)}(K^\bullet, \bigoplus E_j^\bullet/I \oplus E_j^\bullet/J)[-1]
\ar[d] \\
\bigoplus \Hom_{D(R)}(K^\bullet, E_j^\bullet) \ar[r] \ar[d] &
\Hom_{D(R)}(K^\bullet, \bigoplus E_j^\bullet) \ar[d] \\
\bigoplus \Hom_{D(R)}(K^\bullet, E_j^\bullet/I \oplus E_j^\bullet/J)
\ar[r] \ar[d] &
\Hom_{D(R)}(K^\bullet, \bigoplus E_j^\bullet/I \oplus E_j^\bullet/J)
\ar[d] \\
\bigoplus \Hom_{D(R)}(K^\bullet, E_j^\bullet/(I + J)) \ar[r] &
\Hom_{D(R)}(K^\bullet, \bigoplus E_j^\bullet/(I + J))
}
$$
を得る。その列は完全である。任意の $R$-加群の複体 $E^\bullet$ に対し
\begin{align*}
\Hom_{D(R)}(K^\bullet, E^\bullet/I) & =
\Hom_{K(R)}(K^\bullet, E^\bullet/I) \\
& =
\Hom_{K(R/I)}(K^\bullet/IK^\bullet, E^\bullet/I) \\
& =
\Hom_{D(R/I)}(K^\bullet/IK^\bullet, E^\bullet/I)
\end{align*}
である。$J$ または $I + J$ で割る場合も同様である；Derived Categories,
Lemma \ref{derived-lemma-morphisms-from-projective-complex} を見よ。
% P02-E1001: frozen source line 20634 contains the orphan fragment "Derived Categories."
従って中央のものを除くすべての水平射は同型である。実際、複体
$K^\bullet/IK^\bullet$, $K^\bullet/JK^\bullet$,
$K^\bullet/(I + J)K^\bullet$ は、それぞれ $R/I$, $R/J$,
$R/(I + J)$-加群の perfect 複体だからである；Proposition
\ref{more-algebra-proposition-perfect-is-compact} を見よ。
$5$-lemma（Homology, Lemma \ref{homology-lemma-five-lemma}）により
中央の写像は同型であり、Proposition
\ref{more-algebra-proposition-perfect-is-compact} によって補題が従う。
\end{proof}






\section{強生成対象と正則環}
\label{more-algebra-section-strong-generators-regular}

\noindent
$R$ を環とする。
% P02-E1002: frozen source line 20653 neither says "by D(R)_c" nor identifies compact objects.
$D(R)$ の飽和充満三角部分圏を $D(R)_c$ と記す。すでに
$$
\langle R \rangle = D_{perf}(R) = D(R)_c
$$
であることを知っている。Lemma
\ref{more-algebra-lemma-perfect-ring-classical-generator} および Proposition
\ref{more-algebra-proposition-perfect-is-compact} を見よ。
$R$ が正則ならば、$R$ は強生成対象であることが分かる
（Derived Categories, Definition \ref{derived-definition-generators}）。

\begin{lemma}
\label{more-algebra-lemma-ghost-lemma}
\begin{reference}
\cite{Kelly}
\end{reference}
$R$ を環とし、$n \geq 1$ とする。
Derived Categories, Section \ref{derived-section-generators} の記法で
$K \in \langle R \rangle_n$ とする。$D(R)$ における写像
$$
K \xrightarrow{f_1} K_1 \xrightarrow{f_2} K_2
\xrightarrow{f_3} \ldots \xrightarrow{f_n} K_n
$$
を考える。すべての $i, j$ に対して $H^i(f_j) = 0$ ならば、
$f_n \circ \ldots \circ f_1 = 0$ である。
\end{lemma}

\begin{proof}
$n = 1$ ならば、$K$ は $D(R)$ において、各項が有限自由
$R$-加群で微分が零である有界複体 $P^\bullet$ の直和因子である。
従って、すべての $i$ に対して $H^i(f) = 0$ となる $D(R)$ の射
$f : P^\bullet \to K_1$ が零であることを示せば十分である。
$P^\bullet$ は有限直和
$P^\bullet = \bigoplus R[m_j]$ なので、$D(R)$ において
$H^{-m}(g) = 0$ となる任意の射 $g : R[m] \to K_1$ が零であることを
示せば十分である。これは
$\Hom_{D(R)}(R[-m], K) = H^m(K)$ という事実から従う。

\medskip\noindent
$n > 1$ の場合は $n$ に関する帰納法を用いる。
すなわち、$K$ は $D(R)$ において、完全三角
$$
P' \xrightarrow{i} P \xrightarrow{p} P'' \to P'[1]
$$
に入る対象 $P$ の直和因子であり、ここで
$P' \in \langle R \rangle_1$ および
$P'' \in \langle R \rangle_{n - 1}$ である。
上と同様に $K$ を $P$ で置き換え、
$$
P \xrightarrow{f_1} K_1 \xrightarrow{f_2} K_2
\xrightarrow{f_3} \ldots \xrightarrow{f_n} K_n
$$
が $D(R)$ における写像であり、各 $f_j$ がコホモロジー上で零であると
仮定してよい。$n = 1$ の場合により、合成
$f_1 \circ i$ は零である。
従って Derived Categories, Lemma
\ref{derived-lemma-representable-homological} により、
$f_1 = h \circ p$ となる射 $h : P'' \to K_1$ を見つけられる。
$f_2 \circ h$ はコホモロジー上で零であることに注意する。
従って帰納法により
$f_n \circ \ldots \circ f_2 \circ h = 0$ であり、望み通り
$f_n \circ \ldots \circ f_1 = 
f_n \circ \ldots \circ f_2 \circ h \circ p = 0$ が従う。
\end{proof}

\begin{lemma}
\label{more-algebra-lemma-not-regular-not-strong}
$R$ を Noether 環とする。$R$ が $D_{perf}(R)$ の強生成対象ならば、
$R$ は有限次元の正則環である。
\end{lemma}

\begin{proof}
ある $n \geq 1$ について
$D_{perf}(R) = \langle R \rangle_n$ であると仮定する。
任意の有限 $R$-加群 $M$ に対して、有限自由 $R$-加群からなる複体
% P02-E1003: frozen source line 20724 labels this differential d^1 instead of d^{-n+1}.
$$
P = (
P^{-n - 1} \xrightarrow{d^{-n - 1}}
P^{-n} \xrightarrow{d^{-n}}
P^{-n + 1} \xrightarrow{d^1}
\ldots \xrightarrow{d^{-1}} P^0)
$$
で、$i = -n, \ldots, - 1$ に対し $H^i(P) = 0$ かつ
$M \cong \Coker(d^{-1})$ となるものを選べる。
$P$ は $D_{perf}(R)$ に属することに注意する。
% P02-E1004: frozen source line 20730 omits "using" before "the finite free resolution".
任意の $R$-加群 $N$ に対し、$M$ の有限自由分解 $P$ を用いて
$\Ext^n_R(M, N)$ を計算できる；Algebra, Section
\ref{algebra-section-ext} を見よ。また Derived Categories, Section
\ref{derived-section-ext} と比較せよ。
特に、上の列は元
$$
\xi \in \Ext^n_R(\Coker(d^{-1}), \Coker(d^{-n - 1})) =
\Ext^n_R(M, \Coker(d^{-n - 1}))
$$
を定める。また $\Ext^n_R(M, N)$ の任意の元
$\overline{\xi}$ に対して、$\varphi$ が $\xi$ を $\overline{\xi}$ に写す
% P02-E1005: frozen source line 20739 says "a R-module map".
$R$-加群写像 $\varphi : \Coker(d^{-n - 1}) \to N$ が存在する。
$j = 1, \ldots, n - 1$ に対し、複体
$$
K_j = (\Coker(d^{-n - 1}) \to P^{-n + 1} \to \ldots \to P^{-j})
$$
を考える。ここで $\Coker(d^{-n - 1})$ は次数 $-n$ に、
$P^t$ は次数 $t$ に置く。
さらに $K_n = \Coker(d^{-n - 1})[n]$ と置く。すると写像
$$
P \to K_1 \to K_2 \to \ldots \to K_n
$$
が得られ、これらがコホモロジー上で誘導する写像は零である。
Lemma \ref{more-algebra-lemma-ghost-lemma} により、
$P \in D_{perf}(R) = \langle R \rangle_n$ なので、
% P02-E1006: frozen source line 20752 says "composition of this maps".
これらの写像の合成は $D(R)$ において零である。
Derived Categories, Lemma
\ref{derived-lemma-morphisms-from-projective-complex} により
$\Hom_{D(R)}(P, K_n) = \Hom_{K(R)}(P, K_n)$ なので、
$\xi = 0$ と結論する。従って上の議論から、すべての
$R$-加群 $N$ に対して $\Ext^n_R(M, N) = 0$ である。
Algebra, Lemma \ref{algebra-lemma-projective-dimension-ext} により、
$M$ の射影次元は $\leq n - 1$ である。
これはすべての有限 $R$-加群 $M$ について成り立つので、
Algebra, Lemma \ref{algebra-lemma-finite-gl-dim} により
$R$ の大域次元は有限である。
最後に Algebra, Lemma
\ref{algebra-lemma-finite-gl-dim-finite-dim-regular} により結論を得る。
\end{proof}

\begin{lemma}
\label{more-algebra-lemma-ext-regular}
$R$ を次元 $d < \infty$ の Noether 正則環とする。
$K, L \in D^-(R)$ とする。
% P02-E1007: frozen source line 20769 says "an k".
ある $k$ が存在して、$i \leq k$ に対し $H^i(K) = 0$、
$i \geq k - d + 1$ に対し $H^i(L) = 0$ であると仮定する。
このとき $\Hom_{D(R)}(K, L) = 0$ である。
\end{lemma}

\begin{proof}
$K^\bullet$ を $K$ を表す上に有界な複体とし、例えば
$i \geq n + 1$ に対して $K^i = 0$ とする。
$K^\bullet$ を $\tau_{\geq k + 1}K^\bullet$ で置き換えれば、
$i \leq k$ に対して $K^i = 0$ と仮定してよい。
完全三角
$$
K^n[-n] \to K^\bullet \to \sigma_{\leq n - 1}K^\bullet
$$
を用いると、補題を $K^n[-n]$ と
$\sigma_{\leq n - 1}K^\bullet$ について証明すれば十分である。
$n$ に関する帰納法により、ある $R$-加群 $M$ と
$m \geq k + 1$ に対し、$K$ が複体 $M[-m]$ で表される場合に
補題を証明すれば十分である。
Algebra, Lemma
\ref{algebra-lemma-finite-gl-dim-finite-dim-regular} により
$R$ の大域次元は $d$ なので、$M$ は射影分解
$0 \to P_d \to \ldots \to P_0 \to M \to 0$ を持つ。
$P_i$ を次数 $m - i$ に置いた複体 $P^\bullet$ は、
$M[-m]$ を表す射影加群からなる有界複体である。
一方、$L$ を表す複体 $L^\bullet$ で、
$i \geq k - d  + 1$ に対して $L^i = 0$ となるものを選べる。
従って複体の任意の写像 $P^\bullet \to L^\bullet$ は零である。
Derived Categories, Lemma
\ref{derived-lemma-morphisms-from-projective-complex} により補題が従う。
\end{proof}

\begin{lemma}
\label{more-algebra-lemma-split-complex-regular}
$R$ を次元 $1 \leq d < \infty$ の Noether 正則環とする。
$K \in D(R)$ を perfect とし、$k \in \mathbf{Z}$ は
$i = k - d + 2, \ldots, k$ に対して $H^i(K) = 0$ を満たすものとする
（$d = 1$ ならば空の条件である）。
このとき
$K = \tau_{\leq k - d + 1}K \oplus \tau_{\geq k + 1}K$ である。
\end{lemma}

\begin{proof}
コホモロジーの消滅により、完全三角
$$
\tau_{\leq k - d + 1}K \to K \to \tau_{\geq k + 1}K \to
(\tau_{\leq k - d + 1}K)[1]
$$
を得る。Derived Categories, Lemma \ref{derived-lemma-split} により、
三つ目の射が零であることを示せば十分である。従って
$$\Hom_{D(R)}(\tau_{\geq k + 1}K, (\tau_{\leq k - d + 1}K)[1]) = 0$$
を示せば十分であり、これは Lemma \ref{more-algebra-lemma-ext-regular} から従う。
\end{proof}

\begin{lemma}
\label{more-algebra-lemma-regular-strong-generator}
$R$ を有限次元の Noether 正則環とする。
このとき $R$ は、perfect 対象からなる充満部分圏
$D_{perf}(R) \subset D(R)$ の強生成対象である。
\end{lemma}

\begin{proof}
$D(R)$ の対象 $K$ が perfect であることと、$K$ が有界で有限
コホモロジー加群を持つことは同値である；Lemma
\ref{more-algebra-lemma-regular-perfect} を見よ。
三角圏の強生成対象は Derived Categories, Definition
\ref{derived-definition-generators} で定義されている。
$d = \dim(R)$ と置く。

\medskip\noindent
$K \in D_{perf}(R)$ とする。
$K \in \langle R \rangle_{d + 1}$ を示す。
Algebra, Lemma
\ref{algebra-lemma-finite-gl-dim-finite-dim-regular} により、
すべての有限 $R$-加群の射影次元は $\leq d$ である。
$0 \leq i \leq d$ に関する帰納法によって、すべての
$n \in \mathbf{Z}$ に対して $H^n(K)$ の射影次元が $\leq i$ ならば、
$K$ は $\langle R \rangle_{i + 1}$ に属することを示す。

\medskip\noindent
初期段階 $i = 0$。この場合 $H^n(K)$ は射影次元 $0$ の有限
$R$-加群である。言い換えれば、各コホモロジーは射影 $R$-加群である。
従って、すべての $i > 0$ および $m, n \in \mathbf{Z}$ に対して
$\Ext^i_R(H^n(K), H^m(K)) = 0$ である。
Derived Categories, Lemma
\ref{derived-lemma-ext-2-zero-pre} により、$K$ はそのコホモロジー加群の
シフトの直和と同型である。各コホモロジー加群は有限射影
$R$-加群なので、$R$ のコピーの直和の直和因子である。
従って定義により $K$ は $\langle R \rangle_1$ に含まれる。

\medskip\noindent
帰納段階。$i < d$ に対して主張が成り立つと仮定し、
$K \in D_{perf}(R)$ は、すべての $n \in \mathbf{Z}$ に対して
$H^n(K)$ の射影次元が $\leq i + 1$ であるという性質を持つとする。
$n \not \in [a, b]$ に対して $H^n(K)$ が零となる
$a \leq b$ を選ぶ。各 $n \in [a, b]$ に対し、
$F^n$ が有限自由 $R$-加群である全射 $F^n \to H^n(K)$ を選ぶ。
$F^n$ は射影なので、$F^n \to H^n(K)$ を $D(R)$ における写像
$F^n[-n] \to K$ に持ち上げられる（細部は省略する）。
従ってコホモロジー加群上で全射となる射
$\bigoplus_{a \leq n \leq b} F^n[-n] \to K$ を得る。
$D(R)$ における完全三角
$$
K' \to \bigoplus\nolimits_{a \leq n \leq b} F^n[-n] \to K \to K'[1]
$$
を選ぶ。もちろん対象 $K'$ は有界で有限コホモロジー加群を持つ。
選択の仕方により、コホモロジー長完全列は短完全列
$$
0 \to H^n(K') \to F^n \to H^n(K) \to 0
$$
に分かれる。Algebra, Lemma
\ref{algebra-lemma-exact-sequence-projective-dimension} により、
$H^n(K')$ の射影次元は $\leq \max(0, i)$ である。
従って $K' \in \langle R \rangle_{i + 1}$ である。
定義により、望み通り $K$ は
$\langle R \rangle_{i + 1 + 1}$ に属する。
\end{proof}

\begin{proposition}
\label{more-algebra-proposition-regular-strong-generator}
$R$ を Noether 環とする。次は同値である：
\begin{enumerate}
\item $R$ は有限次元の正則環である、
\item $D_{perf}(R)$ は強生成対象を持つ、
\item $R$ は $D_{perf}(R)$ の強生成対象である。
\end{enumerate}
\end{proposition}

\begin{proof}
これは Lemmas \ref{more-algebra-lemma-perfect-ring-classical-generator},
\ref{more-algebra-lemma-not-regular-not-strong}, and
\ref{more-algebra-lemma-regular-strong-generator} および Derived Categories, Lemma
\ref{derived-lemma-classical-generator-strong-generator}
の形式的帰結である。
\end{proof}







\section{相対的有限表示加群}
\label{more-algebra-section-relative-finite-presentation}

\noindent
$R$ を環とする。$A \to B$ を有限型 $R$-代数の有限写像とし、
$M$ を有限 $B$-加群とする。この場合
$$
M\text{ of finite presentation over }B
\Leftrightarrow
M\text{ of finite presentation over }A
$$
は {\bf 成り立たない}。反例は
$R = k[x_1, x_2, x_3, \ldots]$, $A = R$, $B = R/(x_i)$,
および $M = B$ である。これを「修正」するため、有限表示の相対的概念を導入する。

\begin{lemma}
\label{more-algebra-lemma-relatively-finitely-presented}
$R \to A$ を有限型環写像とし、$M$ を $A$-加群とする。
次は同値である：
\begin{enumerate}
\item ある表示 $\alpha : R[x_1, \ldots, x_n] \to A$ に対して
$M$ は有限表示 $R[x_1, \ldots, x_n]$-加群である、
\item すべての表示 $\alpha : R[x_1, \ldots, x_n] \to A$ に対して
$M$ は有限表示 $R[x_1, \ldots, x_n]$-加群である、
\item $A'$ が有限表示 $R$-代数である任意の全射
$A' \to A$ に対して、$M$ は $A'$-加群として有限表示である。
\end{enumerate}
この場合 $M$ は有限表示 $A$-加群である。
\end{lemma}

\begin{proof}
% P02-E1008: frozen source lines 20939-20940 end the conditional clause with a period.
$\alpha : R[x_1, \ldots, x_n] \to A$ および
$\beta : R[y_1, \ldots, y_m] \to A$ を表示とする。
$\alpha(f_j) = \beta(y_j)$ を満たす
$f_j \in R[x_1, \ldots, x_n]$ と、
$\beta(g_i) = \alpha(x_i)$ を満たす
$g_i \in R[y_1, \ldots, y_m]$ を選ぶ。このとき可換図
$$
\xymatrix{
R[x_1, \ldots, x_n, y_1, \ldots, y_m]
\ar[d]^{x_i \mapsto g_i} \ar[rr]_-{y_j \mapsto f_j} & &
R[x_1, \ldots, x_n] \ar[d] \\
R[y_1, \ldots, y_m] \ar[rr] & & A
}
$$
を得る。従って (1) と (2) の同値性は、Algebra, Lemmas
\ref{algebra-lemma-finitely-presented-over-subring} and
\ref{algebra-lemma-finite-finitely-presented-extension}
を適用すれば従う。(2) と (3) の同値性は、表示
$A' = R[x_1, \ldots, x_n]/(f_1, \ldots, f_m)$ を選び、
Algebra, Lemma \ref{algebra-lemma-finite-finitely-presented-extension}
を用いて、$M$ が $A'$-加群として有限表示であることと
$M$ が $R[x_1, \ldots, x_n]$-加群として有限表示であることの同値性を示せば従う。
\end{proof}

\begin{definition}
\label{more-algebra-definition-relatively-finitely-presented}
$R \to A$ を有限型環写像とし、$M$ を $A$-加群とする。
Lemma \ref{more-algebra-lemma-relatively-finitely-presented} の同値な条件が
成り立つとき、$M$ は {\it $R$ に関して有限表示な $A$-加群} であるという。
\end{definition}

\noindent
$R \to A$ が有限表示ならば、$M$ が $R$ に関して有限表示な
$A$-加群であることと、$M$ が有限表示 $A$-加群であることは同値である。
同様に、$A$ が $R$ に関して有限表示な $A$-加群であることと、
$A$ が $R$ 上有限表示であることは同値である。
$R$ が Noether 環ならば、この概念に新しい内容はない。
これで求めていた結果を述べられる。

\begin{lemma}
\label{more-algebra-lemma-finite-extension}
$R$ を環とする。$A \to B$ を有限型 $R$-代数の有限写像とし、
$M$ を $B$-加群とする。このとき、
$M$ が $R$ に関して有限表示な $A$-加群であることと、
$M$ が $R$ に関して有限表示な $B$-加群であることは同値である。
\end{lemma}

\begin{proof}
全射 $R[x_1, \ldots, x_n] \to A$ を選ぶ。
$A$ 上 $B$ を生成する $y_1, \ldots, y_m \in B$ を選ぶ。
$A \to B$ は有限なので、各 $y_i$ は係数が $A$ にあるモニック方程式を満たす。
従って、$B$ において $P_j(y_j) = 0$ となるモニック多項式
$P_j(T) \in R[x_1, \ldots, x_n][T]$ を見つけられる。
すると可換図
$$
\xymatrix{
R[x_1, \ldots, x_n] \ar[d] \ar[r] &
R[x_1, \ldots, x_n, y_1, \ldots, y_m]/(P_j(y_j)) \ar[d] \\
A \ar[r] & B
}
$$
を得る。上辺の写像は有限かつ有限表示なので、
Algebra, Lemma \ref{algebra-lemma-finite-finitely-presented-extension}
と定義により結論を得る。
\end{proof}

\noindent
この結果により、相対的概念は意味を持ち、$R$ 上有限型な環の有限写像に
関して良く振る舞うことが分かる。また局所化と基底変換の下で安定であり、
良く「貼り合わさる」。

\begin{lemma}
\label{more-algebra-lemma-localize-relative-finite-presentation}
% P02-E1009: frozen source lines 21016-21019 omit several required verbs/articles.
$R$ を環、$f \in R$ を元とし、$R_f \to A$ は有限型環写像、
$g \in A$、$M$ は $A$-加群であるとする。
$M$ が $R_f$ に関して有限表示ならば、$M_g$ は
$R$ に関して有限表示な $A_g$-加群である。
\end{lemma}

\begin{proof}
表示 $R_f[x_1, \ldots, x_n] \to A$ を選ぶ。
$R_f = R[x_0]/(fx_0 - 1)$ と書く。
与えられた写像を拡張し、$x_0$ を $1/f$ の像に、
$x_{n + 1}$ を $1/g$ に写す表示
$R[x_0, x_1, \ldots, x_n, x_{n + 1}] \to A_g$ を考える。
$g$ に写る $g' \in R[x_0, x_1, \ldots, x_n]$ を選ぶ
（これは可能である）。行列 $(h_{ij})$ で与えられる $M$ の表示
$$
R_f[x_1, \ldots, x_n]^{\oplus s} \to
R_f[x_1, \ldots, x_n]^{\oplus t} \to M \to 0
$$
があるとする。$h_{ij}$ に写る
$h'_{ij} \in R[x_0, x_1, \ldots, x_n]$ を選ぶ。このとき
% P02-E1010: frozen source line 21040 calls this a presentation of M_f rather than M_g.
$$
R[x_0, x_1, \ldots, x_n, x_{n + 1}]^{\oplus s + 2t} \to
R[x_0, x_1, \ldots, x_n, x_{n + 1}]^{\oplus t} \to M_g \to 0
$$
は $M_f$ の表示である。写像を定める $t \times (s + 2t)$ 行列の
最初の $t \times s$ ブロックは行列 $h'_{ij}$ からなり、
% P02-E1011: frozen source line 21043 has the incomplete factor (x_0f - )I_t.
二つ目の $t \times t$ ブロックは $(x_0f - )I_t$、
三つ目のブロックは $(x_{n + 1}g' - 1)I_t$ である。
\end{proof}

\begin{lemma}
\label{more-algebra-lemma-base-change-relative-finite-presentation}
$R \to A$ を有限型環写像とし、$M$ を $R$ に関して有限表示な
$A$-加群とする。任意の環写像 $R \to R'$ に対し、
$A \otimes_R R'$-加群
% P02-E1012: frozen source line 21053 uses the undefined A' in the base-change tensor product.
$$
M \otimes_A A' = M \otimes_R R'
$$
は $R'$ に関して有限表示である。
\end{lemma}

\begin{proof}
全射 $R[x_1, \ldots, x_n] \to A$ と表示
$$
R[x_1, \ldots, x_n]^{\oplus s} \to
R[x_1, \ldots, x_n]^{\oplus t} \to M \to 0
$$
を選ぶ。このとき
$$
R'[x_1, \ldots, x_n]^{\oplus s} \to
R'[x_1, \ldots, x_n]^{\oplus t} \to M \otimes_R R' \to 0
$$
は基底変換の表示であり、主張が従う。
\end{proof}

\begin{lemma}
\label{more-algebra-lemma-pull-relative-finite-presentation}
$R \to A$ を有限型環写像とする。
$M$ を $R$ に関して有限表示な $A$-加群とする。
$A \to A'$ を有限表示環写像とする。
$A'$-加群 $M \otimes_A A'$ は $R$ に関して有限表示である。
\end{lemma}

\begin{proof}
全射 $R[x_1, \ldots, x_n] \to A$ を選ぶ。
表示 $A' = A[y_1, \ldots, y_m]/(g_1, \ldots, g_l)$ を選ぶ。
$g_i$ に写る
$g'_i \in R[x_1, \ldots, x_n, y_1, \ldots, y_m]$ を選ぶ。
行列 $(h_{ij})$ で与えられる $M$ の表示が
$$
R[x_1, \ldots, x_n]^{\oplus s} \to
R[x_1, \ldots, x_n]^{\oplus t} \to M \to 0
$$
であるとする。このとき
% P02-E1013: frozen source line 21093 introduces an undefined x_0 only in the target free module.
$$
R[x_1, \ldots, x_n, y_1, \ldots, y_m]^{\oplus s + tl} \to
R[x_0, x_1, \ldots, x_n, y_1, \ldots, y_m]^{\oplus t} \to M \otimes_A A' \to 0
$$
は $M \otimes_A A'$ の表示である。
写像を定める $t \times (s + lt)$ 行列では、最初の
$t \times s$ ブロックが行列 $h_{ij}$ からなり、その後に
$g'_iI_t$ である大きさ $t \times t$ のブロックが $l$ 個続く。
\end{proof}

\begin{lemma}
\label{more-algebra-lemma-composition-relative-finite-presentation}
$R \to A \to B$ を有限型環写像とし、$M$ を $B$-加群とする。
$M$ が $A$ に関して有限表示で、$A$ が $R$ 上有限表示ならば、
$M$ は $R$ に関して有限表示である。
\end{lemma}

\begin{proof}
全射 $A[x_1, \ldots, x_n] \to B$ を選ぶ。
行列 $(h_{ij})$ で与えられる表示
$$
A[x_1, \ldots, x_n]^{\oplus s} \to
A[x_1, \ldots, x_n]^{\oplus t} \to M \to 0
$$
を選ぶ。また表示
$$
A = R[y_1, \ldots, y_m]/(g_1, \ldots, g_u).
$$
を選ぶ。$h_{ij}$ に写る
$h'_{ij} \in R[y_1, \ldots, y_m, x_1, \ldots, x_n]$
を選ぶ。このとき表示
$$
R[y_1, \ldots, y_m, x_1, \ldots, x_n]^{\oplus s + tu} \to
R[y_1, \ldots, y_m, x_1, \ldots, x_n]^{\oplus t} \to M \to 0
$$
を得る。ここで $t \times (s + tu)$-行列は、$h'_{ij}$ からなる
最初の $t \times s$ ブロックと、それに続く
$g_iI_t$, $i = 1, \ldots, u$ で与えられる大きさ
$t \times t$ の $u$ 個のブロックからなる。
\end{proof}

\begin{lemma}
\label{more-algebra-lemma-glue-relative-finite-presentation}
$R \to A$ を有限型環写像とし、$M$ を $A$-加群とする。
$f_1, \ldots, f_r \in A$ は単位イデアルを生成するとする。
次は同値である：
\begin{enumerate}
\item 各 $M_{f_i}$ は $R$ に関して有限表示である、
\item $M$ は $R$ に関して有限表示である。
\end{enumerate}
\end{lemma}

\begin{proof}
(2) $\Rightarrow$ (1) は
Lemma \ref{more-algebra-lemma-localize-relative-finite-presentation} に含まれる。
(1) を仮定する。$A$ において $1 = \sum f_ig_i$ と書く。
% P02-E1014: frozen source lines 21145-21147 put a period before the dependent "such that" clause.
全射
$R[x_1, \ldots, x_n, y_1, \ldots, y_r, z_1, \ldots, z_r] \to A$
で、$y_i$ が $f_i$ に、$z_i$ が $g_i$ に写るものを選ぶ。
すると全射
$$
P = R[x_1, \ldots, x_n, y_1, \ldots, y_r, z_1, \ldots, z_r]/(\sum y_iz_i - 1)
\longrightarrow
A.
$$
が存在することが分かる。
Lemma \ref{more-algebra-lemma-relatively-finitely-presented} により、
$M_{f_i}$ は有限表示 $A_{f_i}$-加群である。従って
Algebra, Lemma \ref{algebra-lemma-cover} により、
$M$ は有限表示 $A$-加群である。
従って $M$ は有限 $P$-加群である（$P$ は上のもの）。
全射 $P^{\oplus t} \to M$ を選ぶ。
この写像の核 $K$ が有限 $P$-加群であることを示さなければならない。
$P_{y_i}$ は $A_{f_i}$ に全射なので、
Lemma \ref{more-algebra-lemma-relatively-finitely-presented} および
Algebra, Lemma \ref{algebra-lemma-extension} により、
局所化 $K_{y_i}$ は有限生成 $P_{y_i}$-加群である。
元 $k_{i, j} \in K$, $i = 1, \ldots, r$, $j = 1, \ldots, s_i$
を、その $k_{i, j}$ の像が $K_{y_i}$ を生成するように選ぶ。
$K' \subset K$ を元 $k_{i, j}$ が生成する
$P$-加群とする。このとき $K/K'$ は、すべての $i$ に対して
$y_i$ での局所化が零となる加群である。
$(y_1, \ldots, y_r) = P$ なので、望み通り $K/K' = 0$ である。
\end{proof}

\begin{lemma}
\label{more-algebra-lemma-ses-relatively-finite-presentation}
$R \to A$ を有限型環写像とし、
$0 \to M' \to M \to M'' \to 0$ を $A$-加群の短完全列とする。
\begin{enumerate}
\item $M', M''$ が $R$ に関して有限表示ならば、$M$ もそうである。
\item $M'$ が有限型 $A$-加群で、$M$ が $R$ に関して有限表示ならば、
$M''$ は $R$ に関して有限表示である。
\end{enumerate}
\end{lemma}

\begin{proof}
Algebra, Lemma \ref{algebra-lemma-extension} から直ちに従う。
\end{proof}

\begin{lemma}
\label{more-algebra-lemma-sum-relatively-finite-presentation}
$R \to A$ を有限型環写像とする。
$M, M'$ を $A$-加群とする。$M \oplus M'$ が $R$ に関して有限表示ならば、
$M$ と $M'$ もそうである。
\end{lemma}

\begin{proof}
省略する。
\end{proof}
\section{相対的擬コヒーレント加群}
\label{more-algebra-section-relative-pseudo-coherent}

\noindent
この節は擬コヒーレンスについての Section
\ref{more-algebra-section-relative-finite-presentation} の類似である。

\begin{lemma}
\label{more-algebra-lemma-pull-push}
$R$ を環とし、$K^\bullet$ を $R$-加群の複体とする。
$x$ を零に写す $R$-代数写像 $R[x] \to R$ を考える。
このとき $D(R)$ において
$$
K^\bullet \otimes_{R[x]}^{\mathbf{L}} R \cong K^\bullet \oplus K^\bullet[1]
$$
である。
\end{lemma}

\begin{proof}
すべての $n$ に対し $P^n$ が平坦 $R$-加群であるような、
$R$ 上の K-flat 分解 $P^\bullet \to K^\bullet$ を選ぶ；Lemma
\ref{more-algebra-lemma-K-flat-resolution} を見よ。
すると $P^\bullet \otimes_R R[x]$ は、各項が平坦
$R[x]$-加群である $R[x]$-加群の K-flat 複体である；Lemma
\ref{more-algebra-lemma-base-change-K-flat} および Algebra, Lemma
\ref{algebra-lemma-flat-base-change} を見よ。
特に $x : P^n \otimes_R R[x] \to P^n \otimes_R R[x]$ は単射で、
余核は $P^n$ と同型である。従って
$$
P^\bullet \otimes_R R[x] \xrightarrow{x} P^\bullet \otimes_R R[x]
$$
は $R[x]$-加群の二重複体であり、付随する全複体は
$P^\bullet$、従って $K^\bullet$ と擬同型である。
さらに、この付随全複体は $R[x]$-加群の K-flat 複体である。
これは例えば Lemma \ref{more-algebra-lemma-tensor-product-K-flat} または
Lemma \ref{more-algebra-lemma-K-flat-two-out-of-three} による。
従って
\begin{align*}
K^\bullet \otimes_{R[x]}^{\mathbf{L}} R
& \cong
\text{Tot}(P^\bullet \otimes_R R[x] \xrightarrow{x} P^\bullet \otimes_R R[x])
\otimes_{R[x]} R =
\text{Tot}(P^\bullet \xrightarrow{0} P^\bullet) \\
& = P^\bullet \oplus P^\bullet[1] \cong K^\bullet \oplus K^\bullet[1]
\end{align*}
となり、望み通りである。
\end{proof}

\begin{lemma}
\label{more-algebra-lemma-add-variable-pseudo-coherent}
$R$ を環、$K^\bullet$ を $R$-加群の複体とする。
$m \in \mathbf{Z}$ とする。$x$ を零に写す $R$-代数写像
$R[x] \to R$ を考える。このとき、$K^\bullet$ が $R$-加群の複体として
$m$-擬コヒーレントであることと、$K^\bullet$ が $R[x]$-加群の複体として
$m$-擬コヒーレントであることは同値である。
\end{lemma}

\begin{proof}
これは Lemma \ref{more-algebra-lemma-finite-push-pseudo-coherent} の特別な場合である。
次の別の方法でも証明する。

\medskip\noindent
$0 \to R[x] \to R[x] \to R \to 0$ は完全であることに注意する。
従って $R$ は $R[x]$-加群として擬コヒーレントである。
よって補題の一方の含意は
Lemma \ref{more-algebra-lemma-finite-push-pseudo-coherent} から従う。
他方の含意を示すため、$K^\bullet$ が $R[x]$-加群の複体として
$m$-擬コヒーレントであると仮定する。
Lemma \ref{more-algebra-lemma-pull-pseudo-coherent} により、
$K^\bullet \otimes^{\mathbf{L}}_{R[x]} R$ は
$R$-加群の複体として $m$-擬コヒーレントである。
Lemma \ref{more-algebra-lemma-pull-push} により、
$K^\bullet \oplus K^\bullet[1]$ は $R$-加群の複体として
$m$-擬コヒーレントである。
最後に Lemma \ref{more-algebra-lemma-summands-pseudo-coherent} により、
$K^\bullet$ は $R$-加群の複体として $m$-擬コヒーレントであると結論する。
\end{proof}

\begin{lemma}
\label{more-algebra-lemma-relatively-pseudo-coherent}
$R \to A$ を有限型環写像とする。
$K^\bullet$ を $A$-加群の複体とし、$m \in \mathbf{Z}$ とする。
次は同値である：
\begin{enumerate}
\item ある表示 $\alpha : R[x_1, \ldots, x_n] \to A$ に対して、
複体 $K^\bullet$ は $R[x_1, \ldots, x_n]$-加群の
$m$-擬コヒーレント複体である、
\item すべての表示 $\alpha : R[x_1, \ldots, x_n] \to A$ に対して、
複体 $K^\bullet$ は $R[x_1, \ldots, x_n]$-加群の
$m$-擬コヒーレント複体である。
\end{enumerate}
特に、擬コヒーレンスについても同じ同値性が成り立つ。
\end{lemma}

\begin{proof}
% P02-E1015: frozen source lines 21308-21309 repeat the dangling conditional seen in the prior section.
$\alpha : R[x_1, \ldots, x_n] \to A$ および
$\beta : R[y_1, \ldots, y_m] \to A$ を表示とする。
$\alpha(f_j) = \beta(y_j)$ となる
$f_j \in R[x_1, \ldots, x_n]$ と、$\beta(g_i) = \alpha(x_i)$ となる
$g_i \in R[y_1, \ldots, y_m]$ を選ぶ。このとき可換図
$$
\xymatrix{
R[x_1, \ldots, x_n, y_1, \ldots, y_m]
\ar[d]^{x_i \mapsto g_i} \ar[rr]_-{y_j \mapsto f_j} & &
R[x_1, \ldots, x_n] \ar[d] \\
R[y_1, \ldots, y_m] \ar[rr] & & A
}
$$
を得る。座標変換の後、環準同型
$R[x_1, \ldots, x_n, y_1, \ldots, y_m] \to R[x_1, \ldots, x_n]$
は、各 $y_i$ を零に写す環準同型と同型になる。
図の左の縦写像についても同様である。
従って変数の個数に関する帰納法により、この補題は
Lemma \ref{more-algebra-lemma-add-variable-pseudo-coherent} から従う。
擬コヒーレントの場合は、これと
Lemma \ref{more-algebra-lemma-pseudo-coherent} から従う。
\end{proof}

\begin{definition}
\label{more-algebra-definition-relatively-pseudo-coherent}
$R \to A$ を有限型環写像とする。
$K^\bullet$ を $A$-加群の複体、$M$ を $A$-加群とし、
$m \in \mathbf{Z}$ とする。
\begin{enumerate}
\item Lemma \ref{more-algebra-lemma-relatively-pseudo-coherent} の同値な条件が
成り立つとき、$K^\bullet$ は {\it $R$ に関して $m$-擬コヒーレント}
であるという。
\item すべての $m \in \mathbf{Z}$ に対して $K^\bullet$ が
$R$ に関して $m$-擬コヒーレントであるとき、
$K^\bullet$ は {\it $R$ に関して擬コヒーレント} であるという。
\item $M[0]$ が $R$ に関して $m$-擬コヒーレントであるとき、
$M$ は {\it $R$ に関して $m$-擬コヒーレント} であるという。
\item $M[0]$ が $R$ に関して擬コヒーレントであるとき、
$M$ は {\it $R$ に関して擬コヒーレント} であるという。
\end{enumerate}
\end{definition}

\noindent
Part (2) は、任意の全射
$R[y_1, \ldots, y_m] \to A$ に対して、$K^\bullet$ が
$R[x_1, \ldots, x_n]$-加群の複体として擬コヒーレントであることを意味する；
Lemma \ref{more-algebra-lemma-pseudo-coherent} を見よ。
この定義には次の好ましい性質がある。

\begin{lemma}
\label{more-algebra-lemma-finite-extension-pseudo-coherent}
$R$ を環とする。$A \to B$ を有限型 $R$-代数の有限写像とする。
$m \in \mathbf{Z}$ とし、$K^\bullet$ を $B$-加群の複体とする。
このとき、$K^\bullet$ が $R$ に関して $m$-擬コヒーレント
（resp.\ 擬コヒーレント）であることと、$A$-加群の複体とみなした
$K^\bullet$ が $R$ に関して $m$-擬コヒーレント
（擬コヒーレント）であることは同値である。
\end{lemma}

\begin{proof}
全射 $R[x_1, \ldots, x_n] \to A$ を選ぶ。
$A$ 上 $B$ を生成する $y_1, \ldots, y_m \in B$ を選ぶ。
$A \to B$ は有限なので、各 $y_i$ は係数が $A$ にある
モニック方程式を満たす。従って $B$ において $P_j(y_j) = 0$ となる
モニック多項式 $P_j(T) \in R[x_1, \ldots, x_n][T]$ を見つけられる。
すると可換図
$$
\xymatrix{
& R[x_1, \ldots, x_n, y_1, \ldots, y_m] \ar[d] \\
R[x_1, \ldots, x_n] \ar[d] \ar[r] &
R[x_1, \ldots, x_n, y_1, \ldots, y_m]/(P_j(y_j)) \ar[d] \\
A \ar[r] & B
}
$$
を得る。上の水平射と右上の縦射は
Lemma \ref{more-algebra-lemma-finite-push-pseudo-coherent} の仮定を満たす。
従って $K^\bullet$ が $R[x_1, \ldots, x_n]$-加群の複体として
$m$-擬コヒーレント（resp.\ 擬コヒーレント）であることと、
$K^\bullet$ が $R[x_1, \ldots, x_n, y_1, \ldots, y_m]$-加群の複体として
$m$-擬コヒーレント（resp.\ 擬コヒーレント）であることは同値である。
\end{proof}

\begin{lemma}
\label{more-algebra-lemma-cone-relatively-pseudo-coherent}
$R$ を環とし、$R \to A$ を有限型環写像とする。
$m \in \mathbf{Z}$ とする。
$(K^\bullet, L^\bullet, M^\bullet, f, g, h)$ を $D(A)$ の完全三角とする。
\begin{enumerate}
\item $K^\bullet$ が $R$ に関して $(m + 1)$-擬コヒーレントで、
$L^\bullet$ が $R$ に関して $m$-擬コヒーレントならば、
$M^\bullet$ は $R$ に関して $m$-擬コヒーレントである。
\item $K^\bullet, M^\bullet$ が $R$ に関して $m$-擬コヒーレントならば、
$L^\bullet$ は $R$ に関して $m$-擬コヒーレントである。
\item $L^\bullet$ が $R$ に関して $(m + 1)$-擬コヒーレントで、
$M^\bullet$ が $R$ に関して $m$-擬コヒーレントならば、
$K^\bullet$ は $R$ に関して $(m + 1)$-擬コヒーレントである。
\end{enumerate}
さらに、$K^\bullet, L^\bullet, M^\bullet$ のうち二つが
$R$ に関して擬コヒーレントならば、三つ目もそうである。
% P02-E1016: frozen source line 21408 reads "the so is the third".
\end{lemma}

\begin{proof}
Lemma \ref{more-algebra-lemma-cone-pseudo-coherent} と定義から直ちに従う。
\end{proof}

\begin{lemma}
\label{more-algebra-lemma-rel-n-pseudo-module}
$R \to A$ を有限型環写像とし、$M$ を $A$-加群とする。このとき
\begin{enumerate}
\item $M$ が $R$ に関して $0$-擬コヒーレントであることと、
$M$ が有限型 $A$-加群であることは同値である、
\item $M$ が $R$ に関して $(-1)$-擬コヒーレントであることと、
% P02-E1017: frozen source line 21425 says "is a finitely presented relative".
$M$ が $R$ に関して有限表示であることは同値である、
\item $M$ が $R$ に関して $(-d)$-擬コヒーレントであることと、
任意の全射 $R[x_1, \ldots, x_n] \to A$ に対して長さ $d$ の分解
$$
R[x_1, \ldots, x_n]^{\oplus a_d} \to R[x_1, \ldots, x_n]^{\oplus a_{d - 1}}
\to \ldots \to R[x_1, \ldots, x_n]^{\oplus a_0} \to M \to 0
$$
が存在することは同値である、
\item $M$ が $R$ に関して擬コヒーレントであることと、
任意の表示 $R[x_1, \ldots, x_n] \to A$ に対して、有限自由
$R[x_1, \ldots, x_n]$-加群による無限分解
$$
\ldots \to R[x_1, \ldots, x_n]^{\oplus a_1} \to
R[x_1, \ldots, x_n]^{\oplus a_0} \to M \to 0
$$
が存在することは同値である。
\end{enumerate}
\end{lemma}

\begin{proof}
Lemma \ref{more-algebra-lemma-n-pseudo-module} と定義から直ちに従う。
\end{proof}

\begin{lemma}
\label{more-algebra-lemma-summands-relative-pseudo-coherent}
$R \to A$ を有限型環写像とする。
$m \in \mathbf{Z}$ とし、$K^\bullet, L^\bullet \in D(A)$ とする。
$K^\bullet \oplus L^\bullet$ が $R$ に関して $m$-擬コヒーレント
（resp.\ 擬コヒーレント）ならば、$K^\bullet$ と $L^\bullet$ もそうである。
\end{lemma}

\begin{proof}
Lemma \ref{more-algebra-lemma-summands-pseudo-coherent} と定義から直ちに従う。
\end{proof}

\begin{lemma}
\label{more-algebra-lemma-complex-relative-pseudo-coherent-modules}
$R \to A$ を有限型環写像とする。$m \in \mathbf{Z}$ とする。
$K^\bullet$ を上に有界な $A$-加群の複体とし、すべての $i$ に対して
$K^i$ が $R$ に関して $(m - i)$-擬コヒーレントであるとする。
このとき $K^\bullet$ は $R$ に関して $m$-擬コヒーレントである。
特に、$K^\bullet$ が $R$ に関して擬コヒーレントな $A$-加群からなる
上に有界な複体ならば、$K^\bullet$ は $R$ に関して擬コヒーレントである。
\end{lemma}

\begin{proof}
Lemma \ref{more-algebra-lemma-complex-pseudo-coherent-modules} と定義から直ちに従う。
\end{proof}

\begin{lemma}
\label{more-algebra-lemma-cohomology-relative-pseudo-coherent}
$R \to A$ を有限型環写像とし、$m \in \mathbf{Z}$ とする。
$K^\bullet \in D^{-}(A)$ とし、すべての $i$ に対して
$H^i(K^\bullet)$ が $R$ に関して $(m - i)$-擬コヒーレント
（resp.\ 擬コヒーレント）であるとする。
このとき $K^\bullet$ は $R$ に関して $m$-擬コヒーレント
（resp.\ 擬コヒーレント）である。
\end{lemma}

\begin{proof}
Lemma \ref{more-algebra-lemma-cohomology-pseudo-coherent} と定義から直ちに従う。
\end{proof}

\begin{lemma}
\label{more-algebra-lemma-localize-relative-pseudo-coherent}
% P02-E1018: frozen source lines 21500-21504 have a malformed list of hypotheses.
$R$ を環、$f \in R$ を元とし、$R_f \to A$ は有限型環写像、
$g \in A$、$K^\bullet$ は $A$-加群の複体であるとする。
$K^\bullet$ が $R_f$ に関して $m$-擬コヒーレント
（resp.\ 擬コヒーレント）ならば、$K^\bullet \otimes_A A_g$ は
$R$ に関して $m$-擬コヒーレント（resp.\ 擬コヒーレント）である。
\end{lemma}

\begin{proof}
まず $K^\bullet$ が $R$ に関して $m$-擬コヒーレントであることを示す。
実際、$R_f[x_1, \ldots, x_n] \to A$ が全射であるとする。
$R_f = R[x_0]/(fx_0 - 1)$ と書く。このとき
$R[x_0, x_1, \ldots, x_n] \to A$ は全射であり、
$R_f[x_1, \ldots, x_n]$ は
$R[x_0, \ldots, x_n]$-加群として擬コヒーレントである。従って
Lemma \ref{more-algebra-lemma-finite-push-pseudo-coherent} により、
$K^\bullet$ は $R[x_0, x_1, \ldots, x_n]$-加群の複体として
$m$-擬コヒーレントであることが分かる。

\medskip\noindent
$g \in A$ に写る元
$g' \in R[x_0, x_1, \ldots, x_n]$ を選ぶ。
Lemma \ref{more-algebra-lemma-pull-pseudo-coherent} により
% P02-E1019: frozen source line 21526 ends the localization calculation with A_f, although g/g' is localized.
\begin{align*}
K^\bullet \otimes_{R[x_0, x_1, \ldots, x_n]}^{\mathbf{L}}
R[x_0, x_1, \ldots, x_n, \frac{1}{g'}] & =
K^\bullet \otimes_{R[x_0, x_1, \ldots, x_n]}
R[x_0, x_1, \ldots, x_n, \frac{1}{g'}] \\
& = K^\bullet \otimes_A A_f
\end{align*}
は $R[x_0, x_1, \ldots, x_n, \frac{1}{g'}]$-加群の複体として
$m$-擬コヒーレントであることが分かる。
% P02-E1020: frozen source line 21530 starts a new sentence with lowercase "write".
次のように書く：
$$
R[x_0, x_1, \ldots, x_n, \frac{1}{g'}] =
R[x_0, \ldots, x_n, x_{n + 1}]/(x_{n + 1}g' - 1).
$$
$R[x_0, x_1, \ldots, x_n, \frac{1}{g'}]$ は
$R[x_0, \ldots, x_n, x_{n + 1}]$-加群として擬コヒーレントなので、
Lemma \ref{more-algebra-lemma-finite-push-pseudo-coherent} を見れば、望み通り
$K^\bullet \otimes_A A_g$ は
$R[x_0, \ldots, x_n, x_{n + 1}]$-加群の複体として
$m$-擬コヒーレントであると結論できる。
\end{proof}

\begin{lemma}
\label{more-algebra-lemma-base-change-relative-pseudo-coherent}
$R \to A$ を有限型環写像とする。$m \in \mathbf{Z}$ とする。
$K^\bullet$ を $R$ に関して $m$-擬コヒーレント
（resp.\ 擬コヒーレント）な $A$-加群の複体とする。
$R \to R'$ を環写像で、$A$ と $R'$ が $R$ 上 Tor 独立であるものとする。
$A' = A \otimes_R R'$ とおく。このとき
$K^\bullet \otimes_A^{\mathbf{L}} A'$ は
$R'$ に関して $m$-擬コヒーレント（resp.\ 擬コヒーレント）である。
\end{lemma}

\begin{proof}
全射 $R[x_1, \ldots, x_n] \to A$ を選ぶ。次の等式に注意する：
$$
K^\bullet \otimes_A^{\mathbf{L}} A' =
K^\bullet \otimes_R^{\mathbf{L}} R' =
K^\bullet \otimes_{R[x_1, \ldots, x_n]}^{\mathbf{L}} R'[x_1, \ldots, x_n]
$$
これは Lemma \ref{more-algebra-lemma-base-change-comparison} を二度適用すれば従う。
従って Lemma \ref{more-algebra-lemma-pull-pseudo-coherent} により結論を得る。
\end{proof}

\begin{lemma}
\label{more-algebra-lemma-pull-relative-pseudo-coherent}
$R \to A \to B$ を有限型環写像とする。$m \in \mathbf{Z}$ とする。
$K^\bullet$ を $A$-加群の複体とする。
$B$ が $B$-加群として $A$ に関して擬コヒーレントであると仮定する。
$K^\bullet$ が $R$ に関して $m$-擬コヒーレント
（resp.\ 擬コヒーレント）ならば、
$K^\bullet \otimes_A^{\mathbf{L}} B$ は
$R$ に関して $m$-擬コヒーレント（resp.\ 擬コヒーレント）である。
\end{lemma}

\begin{proof}
% P02-E1021: frozen source line 21579 reuses the fixed coherence bound m as the number of polynomial generators.
全射 $A[y_1, \ldots, y_m] \to B$ を選ぶ。
全射 $R[x_1, \ldots, x_n] \to A$ を選ぶ。
これらを合わせて全射
$R[x_1, \ldots, x_n, y_1, \ldots y_m] \to B$ を得る。
% P02-E1022: frozen source line 21583 changes the polynomial generator list from y_m to y_n.
有限自由 $A[y_1, \ldots, y_n]$-加群の複体による $B$ の分解
$E^\bullet \to B$ を選ぶ（環写像 $A \to B$ に関する仮定により可能である）。
$K^\bullet$ は平坦 $A$-加群からなる上に有界な複体であると仮定してよい。
このとき
\begin{align*}
K^\bullet \otimes_A^{\mathbf{L}} B & =
\text{Tot}(K^\bullet \otimes_A B[0]) \\
& = \text{Tot}(K^\bullet \otimes_A A[y_1, \ldots, y_m]
\otimes_{A[y_1, \ldots, y_m]} B[0]) \\
& \cong
\text{Tot}\left(
(K^\bullet \otimes_A A[y_1, \ldots, y_m])
\otimes_{A[y_1, \ldots, y_m]} E^\bullet
\right) \\
& =
\text{Tot}(K^\bullet \otimes_A E^\bullet)
\end{align*}
が $D(A[y_1, \ldots, y_m])$ において成り立つ。
この擬同型 $\cong$ は Lemma
\ref{more-algebra-lemma-derived-tor-quasi-isomorphism} の適用から得られる。
従って、$\text{Tot}(K^\bullet \otimes_A E^\bullet)$ が
$R[x_1, \ldots, x_n, y_1, \ldots y_m]$-加群の複体として
$m$-擬コヒーレントであることを示せばよい。
$\text{Tot}(K^\bullet \otimes_A E^\bullet)$ は部分複体によるフィルトレーションを持ち、
その逐次商は複体 $K^\bullet \otimes_A E^i[-i]$ であることに注意する。
% P02-E1023: frozen source lines 21607-21609 use degrees <= m; the shift direction needed by the argument is the opposite.
$i \ll 0$ ならば、複体 $K^\bullet \otimes_A E^i[-i]$ は次数
$\leq m$ でコホモロジーが零となり、従って（任意の環上で）
$m$-擬コヒーレントであることに注意する。従って Lemma
\ref{more-algebra-lemma-cone-relatively-pseudo-coherent} と帰納法を適用すれば、
% P02-E1024: frozen source line 21612 asks for full pseudo-coherence although the reduction and conclusion are only m-pseudo-coherent.
すべての $i$ に対して $K^\bullet \otimes_A E^i[-i]$ が
$R$ に関して擬コヒーレントであることを示せば十分である。
$i > 0$ ならば $E^i = 0$ であることに注意する。また $E^i$ は有限自由なので、
$K^\bullet \otimes_A A[y_1, \ldots, y_m]$ が
$R$ に関して $m$-擬コヒーレントであることの証明に帰着するが、
これは例えば Lemma \ref{more-algebra-lemma-base-change-relative-pseudo-coherent} から従う。
\end{proof}

\begin{lemma}
\label{more-algebra-lemma-pull-relative-pseudo-coherent-module}
$R \to A \to B$ を有限型環写像とする。
$m \in \mathbf{Z}$ とする。$M$ を $A$-加群とする。
$B$ は $A$ 上平坦で、$B$ は $B$-加群として
$A$ に関して擬コヒーレントであると仮定する。
$M$ が $R$ に関して $m$-擬コヒーレント
（resp.\ 擬コヒーレント）ならば、$M \otimes_A B$ は
$R$ に関して $m$-擬コヒーレント（resp.\ 擬コヒーレント）である。
\end{lemma}

\begin{proof}
Lemma \ref{more-algebra-lemma-pull-relative-pseudo-coherent} から直ちに従う。
\end{proof}

\begin{lemma}
\label{more-algebra-lemma-composition-relative-pseudo-coherent}
$R$ を環とする。$A \to B$ を有限型 $R$-代数の写像とする。
$m \in \mathbf{Z}$ とする。$K^\bullet$ を $B$-加群の複体とする。
$A$ は $R$ に関して擬コヒーレントであると仮定する。このとき次は同値である：
\begin{enumerate}
\item $K^\bullet$ は $A$ に関して $m$-擬コヒーレント
（resp.\ 擬コヒーレント）である、
\item $K^\bullet$ は $R$ に関して $m$-擬コヒーレント
（resp.\ 擬コヒーレント）である。
\end{enumerate}
\end{lemma}

\begin{proof}
全射 $R[x_1, \ldots, x_n] \to A$ を選ぶ。
全射 $A[y_1, \ldots, y_m] \to B$ を選ぶ。このとき全射
$$
R[x_1, \ldots, x_n, y_1, \ldots, y_m] \to A[y_1, \ldots, y_m]
$$
を得るが、これは $R[x_1, \ldots, x_n] \to A$ の平坦な基底変換である。
仮定により $A$ は $R[x_1, \ldots, x_n]$ 上擬コヒーレントな加群なので、
Lemma \ref{more-algebra-lemma-flat-base-change-pseudo-coherent} により
$A[y_1, \ldots, y_m]$ は
$R[x_1, \ldots, x_n, y_1, \ldots, y_m]$ 上擬コヒーレントである。
従ってこの補題は Lemma \ref{more-algebra-lemma-finite-push-pseudo-coherent} と定義から従う。
\end{proof}

\begin{lemma}
\label{more-algebra-lemma-glue-relative-pseudo-coherent}
$R \to A$ を有限型環写像とする。$K^\bullet$ を $A$-加群の複体とする。
$m \in \mathbf{Z}$ とする。
$f_1, \ldots, f_r \in A$ は単位イデアルを生成するとする。
次は同値である：
\begin{enumerate}
\item 各 $K^\bullet \otimes_A A_{f_i}$ は $R$ に関して
$m$-擬コヒーレントである、
\item $K^\bullet$ は $R$ に関して $m$-擬コヒーレントである。
\end{enumerate}
$R$ に関する擬コヒーレンスについても同じ同値性が成り立つ。
\end{lemma}

\begin{proof}
(2) $\Rightarrow$ (1) は Lemma
\ref{more-algebra-lemma-localize-relative-pseudo-coherent} に含まれる。
(1) を仮定する。$A$ において $1 = \sum f_ig_i$ と書く。
% P02-E1025: frozen source lines 21687-21689 put a period before the dependent "such that" clause.
全射
$R[x_1, \ldots, x_n, y_1, \ldots, y_r, z_1, \ldots, z_r] \to A$
で、$y_i$ が $f_i$ に、$z_i$ が $g_i$ に写るものを選ぶ。
すると全射
$$
P = R[x_1, \ldots, x_n, y_1, \ldots, y_r, z_1, \ldots, z_r]/(\sum y_iz_i - 1)
\longrightarrow
A.
$$
が存在することが分かる。$P$ は
$R[x_1, \ldots, x_n, y_1, \ldots, y_r, z_1, \ldots, z_r]$-加群として
擬コヒーレントであり、$P[1/y_i]$ は
$R[x_1, \ldots, x_n, y_1, \ldots, y_r, z_1, \ldots, z_r, 1/y_i]$-加群として
擬コヒーレントであることに注意する。従って Lemma
\ref{more-algebra-lemma-finite-push-pseudo-coherent} により、各 $i$ に対して
$K^\bullet \otimes_A A_{f_i}$ は $P[1/y_i]$-加群の
$m$-擬コヒーレント複体であることが分かる。従って Lemma
\ref{more-algebra-lemma-glue-pseudo-coherent} により、
% P02-E1026: frozen source lines 21703-21711 start from m-pseudo-coherence but conclude full pseudo-coherence twice.
$K^\bullet$ は $P$-加群の複体として擬コヒーレントであり、Lemma
\ref{more-algebra-lemma-finite-push-pseudo-coherent} により $K^\bullet$ は
$R[x_1, \ldots, x_n, y_1, \ldots, y_r, z_1, \ldots, z_r]$-加群の複体として
擬コヒーレントであることが分かる。
\end{proof}

\begin{lemma}
\label{more-algebra-lemma-Noetherian-relative-pseudo-coherent}
$R$ を Noether 環とする。$R \to A$ を有限型環写像とする。このとき
\begin{enumerate}
\item $A$-加群の複体 $K^\bullet$ が $R$ に関して
$m$-擬コヒーレントであることと、$K^\bullet \in D^{-}(A)$ かつ
$i \geq m$ に対して $H^i(K^\bullet)$ が有限 $A$-加群であることは同値である。
\item $A$-加群の複体 $K^\bullet$ が $R$ に関して擬コヒーレントであることと、
$K^\bullet \in D^{-}(A)$ かつ、すべての $i$ に対して
$H^i(K^\bullet)$ が有限 $A$-加群であることは同値である。
\item $A$-加群が $R$ に関して擬コヒーレントであることと、
それが有限であることは同値である。
\end{enumerate}
\end{lemma}

\begin{proof}
Lemma \ref{more-algebra-lemma-Noetherian-pseudo-coherent} と定義から直ちに従う。
\end{proof}











\section{擬コヒーレント環写像と perfect 環写像}
\label{more-algebra-section-pseudo-coherent-perfect-ring-map}

\noindent
これらの種類の環写像を次のように定義できる。

\begin{definition}
\label{more-algebra-definition-pseudo-coherent-perfect}
$A \to B$ を環写像とする。
\begin{enumerate}
\item $A \to B$ が有限型で、$B$ が $B$-加群として
$A$ に関して擬コヒーレントであるとき、この写像を
{\it 擬コヒーレント環写像}という。
\item $A \to B$ が擬コヒーレント環写像で、$A$-加群としての
$B$ が有限 Tor 次元を持つとき、この写像を
{\it perfect 環写像}という。
\end{enumerate}
\end{definition}

\noindent
この用語法は標準的でないかもしれない。Lemma
\ref{more-algebra-lemma-rel-n-pseudo-module} により、$A \to B$ が
擬コヒーレントであることと、$B = A[x_1, \ldots, x_n]/I$ であり、
$A[x_1, \ldots, x_n]$-加群としての $B$ が有限自由
$A[x_1, \ldots, x_n]$-加群による分解を持つことは同値である。
perfect 環写像の定義の動機は Lemma \ref{more-algebra-lemma-perfect} である。
% P02-E1027: frozen source lines 21770-21771 say "The following lemmas gives" although one lemma follows.
次の補題は perfect 環写像の、より有用で直観的な特徴づけを与える。

\begin{lemma}
\label{more-algebra-lemma-perfect-ring-map}
環写像 $A \to B$ が perfect であることと、
$B = A[x_1, \ldots, x_n]/I$ であり、
$A[x_1, \ldots, x_n]$-加群としての $B$ が有限射影
$A[x_1, \ldots, x_n]$-加群による有限分解を持つことは同値である。
\end{lemma}

\begin{proof}
$A \to B$ が perfect ならば、$B = A[x_1, \ldots, x_n]/I$ であり、
$B$ は $A[x_1, \ldots, x_n]$-加群として擬コヒーレントで、
$A$-加群として有限 Tor 次元を持つ。従って Lemma
\ref{more-algebra-lemma-perfect-over-polynomial-ring} により、$B$ は
% P02-E1028: frozen source line 21785 uses the article "a" before A[x_1,...,x_n]-module.
$A[x_1, \ldots, x_n]$-加群として perfect、すなわち有限射影
$A[x_1, \ldots, x_n]$-加群による有限分解を持つ
（Lemma \ref{more-algebra-lemma-perfect-module}）。逆に、
$B = A[x_1, \ldots, x_n]/I$ であり、
$A[x_1, \ldots, x_n]$-加群としての $B$ が有限射影
$A[x_1, \ldots, x_n]$-加群による有限分解を持つならば、
$B$ は $A[x_1, \ldots, x_n]$-加群として擬コヒーレントであり、
従って $A \to B$ は擬コヒーレントである。さらに、与えられた
$A[x_1, \ldots, x_n]$ 上の分解は平坦 $A$-加群による有限分解なので、
$A$-加群としての $B$ は有限 Tor 次元を持つ。
\end{proof}

\noindent
前節までの多くの結果はこの用語法で言い換えられる。
スキームの射について対応する議論は More on Morphisms, Sections
\ref{more-morphisms-section-pseudo-coherent} and
\ref{more-morphisms-section-perfect} も参照されたい。

\begin{lemma}
\label{more-algebra-lemma-Noetherian-pseudo-coherent-ring-map}
Noether 環の有限型環写像は擬コヒーレントである。
\end{lemma}

\begin{proof}
Lemma \ref{more-algebra-lemma-Noetherian-relative-pseudo-coherent} を見よ。
\end{proof}

\begin{lemma}
\label{more-algebra-lemma-flat-finite-presentation-perfect}
平坦かつ有限表示な環写像は perfect である。
\end{lemma}

\begin{proof}
$A \to B$ を平坦かつ有限表示な環写像とする。
$B$ が有限 Tor 次元を持つことは明らかである。Algebra, Lemma
\ref{algebra-lemma-flat-finite-presentation-limit-flat} により、
有限型 $\mathbf{Z}$-代数 $A_0 \subset A$ と平坦有限型環写像
$A_0 \to B_0$ で $B = B_0 \otimes_{A_0} A$ を満たすものが存在する。
Lemma \ref{more-algebra-lemma-Noetherian-relative-pseudo-coherent} により
$A_0 \to B_0$ は擬コヒーレントである。
$A_0 \to B_0$ は平坦なので、$B_0$ と $A$ は $A_0$ 上 Tor 独立である。
従って Lemma \ref{more-algebra-lemma-base-change-relative-pseudo-coherent} を用いて
$A \to B$ が擬コヒーレントであると結論できる。
\end{proof}

\begin{lemma}
\label{more-algebra-lemma-regular-perfect-ring-map}
$A \to B$ を有限型環写像とし、$A$ は有限次元正則環であるとする。
このとき $A \to B$ は perfect である。
\end{lemma}

\begin{proof}
Algebra, Lemma \ref{algebra-lemma-finite-gl-dim-finite-dim-regular} により、
$A$ に関する仮定は $A$ が有限大域次元を持つことを意味する。
従ってすべての加群は有限 Tor 次元を持つ；Lemma
\ref{more-algebra-lemma-finite-gl-dim-tor-dimension} を見よ。特に $B$ もそうである。
Lemma \ref{more-algebra-lemma-Noetherian-pseudo-coherent-ring-map} により、
この写像は擬コヒーレントである。
\end{proof}

\begin{lemma}
\label{more-algebra-lemma-lci-perfect}
局所完全交叉準同型は perfect である。
\end{lemma}

\begin{proof}
$A \to B$ を局所完全交叉準同型とする。Definition
\ref{more-algebra-definition-local-complete-intersection} によれば、これは
$B = A[x_1, \ldots, x_n]/I$ で、$I$ が
$A[x_1, \ldots, x_n]$ の Koszul イデアルであることを意味する。
Lemmas \ref{more-algebra-lemma-perfect-ring-map} and \ref{more-algebra-lemma-perfect-module} により、
$I$ が $A[x_1, \ldots, x_n]$ 上 perfect な加群であることを示せば十分である。
Lemma \ref{more-algebra-lemma-glue-perfect} により、これは局所的な問題である。
従って $I$ は Koszul 正則列で生成されると仮定してよい
（Definition \ref{more-algebra-definition-regular-ideal} による）。
もちろん、これは $I$ が有限自由分解を持つことを意味し、結論を得る。
\end{proof}

\begin{lemma}
\label{more-algebra-lemma-relative-pseudo-coherent-is-moot}
$R \to A$ を擬コヒーレント環写像とする。$K \in D(A)$ とする。
次は同値である：
\begin{enumerate}
\item $K$ は $R$ に関して $m$-擬コヒーレント
（resp.\ 擬コヒーレント）である、
\item $K$ は $D(A)$ において $m$-擬コヒーレント
（resp.\ 擬コヒーレント）である。
\end{enumerate}
\end{lemma}

\begin{proof}
Lemma \ref{more-algebra-lemma-composition-relative-pseudo-coherent} の特別な場合の言い換えである。
\end{proof}

\begin{lemma}
\label{more-algebra-lemma-more-relative-pseudo-coherent-is-moot}
$R \to B \to A$ を環写像とし、$\varphi : B \to A$ は全射、
$R \to B$ と $R \to A$ は平坦かつ有限表示であるとする。
$K \in D(A)$ に対し、$\varphi_*K \in D(B)$ で制限を表す。
次は同値である：
\begin{enumerate}
\item $K$ は擬コヒーレントである、
\item $K$ は $R$ に関して擬コヒーレントである、
\item $K$ は $A$ に関して擬コヒーレントである、
\item $\varphi_*K$ は擬コヒーレントである、
\item $\varphi_*K$ は $R$ に関して擬コヒーレントである。
\end{enumerate}
% P02-E1029: frozen source line 21897 reads "Similar holds" rather than "Similarly" or "The same holds".
$m$-擬コヒーレンスについても同様である。
\end{lemma}

\begin{proof}
$R \to A$ と $R \to B$ は perfect 環写像
（Lemma \ref{more-algebra-lemma-flat-finite-presentation-perfect}）であり、従って特に
擬コヒーレント環写像である。(1) $\Leftrightarrow$ (2) および
(4) $\Leftrightarrow$ (5) は Lemma
\ref{more-algebra-lemma-relative-pseudo-coherent-is-moot} から従う。

\medskip\noindent
$A$ が $R$ に関して擬コヒーレントであることを用い、Lemma
\ref{more-algebra-lemma-composition-relative-pseudo-coherent} により
(2) $\Leftrightarrow$ (3) が分かる。しかし、
% P02-E1030: frozen source line 21911 reverses the surjective map from B -> A to A -> B.
$A \to B$ は全射なので、Definition
\ref{more-algebra-definition-relatively-pseudo-coherent} から直接、
% P02-E1031: frozen source line 21913 uses "equivalent with" instead of "equivalent to".
(3) は (4) と同値であることが分かる。
\end{proof}

\section{相対的 perfect 加群}
\label{more-algebra-section-relatively-perfect}

\noindent
この節は、導来圏の perfect 対象についての Section
\ref{more-algebra-section-relative-pseudo-coherent} の類似である。
% P02-E1032: frozen source line 21925 starts a new sentence with lowercase "we".
一般の場合に正しい定義が何であるか確信がないため、この概念は
限定された一般性のもとでのみ定義する。議論については
Derived Categories of Schemes, Remark
\ref{perfect-remark-discuss-rel-perfect} を見よ。

\begin{definition}
\label{more-algebra-definition-relatively-perfect}
$R \to A$ を平坦かつ有限表示な環写像とする。
$D(A)$ の対象 $K$ が擬コヒーレント
（Definition \ref{more-algebra-definition-pseudo-coherent}）であり、$R$ 上有限 Tor 次元
（Definition \ref{more-algebra-definition-tor-amplitude}）を持つとき、$K$ は
{\it $R$-perfect}、または {it $R$ に関して perfect} であるという。
\end{definition}

\noindent
Lemma \ref{more-algebra-lemma-more-relative-pseudo-coherent-is-moot} により、
$K$ が $R$ に関して擬コヒーレントであることを要求しても同じであった。
必要な補題をいくつか挙げる。

\begin{lemma}
\label{more-algebra-lemma-cone-relatively-perfect}
$R \to A$ を平坦かつ有限表示な環写像とする。
$D(A)$ の $R$-perfect 対象は、飽和な\footnote{Derived Categories,
Definition \ref{derived-definition-saturated}.} 三角厳密充満部分圏をなす。
\end{lemma}

\begin{proof}
Lemmas \ref{more-algebra-lemma-cone-pseudo-coherent},
\ref{more-algebra-lemma-summands-pseudo-coherent},
\ref{more-algebra-lemma-cone-tor-amplitude}, and
\ref{more-algebra-lemma-summands-tor-amplitude} から従う。
\end{proof}

\begin{lemma}
\label{more-algebra-lemma-perfect-relatively-perfect}
$R \to A$ を平坦かつ有限表示な環写像とする。
$D(A)$ の perfect 対象は $R$-perfect である。$K, M \in D(A)$ とすると、
$K$ が perfect で $M$ が $R$-perfect ならば、
$K \otimes_A^\mathbf{L} M$ は $R$-perfect である。
\end{lemma}

\begin{proof}
第一の主張は、第二の主張で $M = A$ とすれば従う。
第二の主張は Lemmas \ref{more-algebra-lemma-perfect},
\ref{more-algebra-lemma-push-tor-amplitude}, and \ref{more-algebra-lemma-tensor-pseudo-coherent} から従う。
\end{proof}

\begin{lemma}
\label{more-algebra-lemma-structure-relatively-perfect}
$R \to A$ を平坦かつ有限表示な環写像とする。
$K \in D(A)$ とする。次は同値である：
\begin{enumerate}
\item $K$ は $R$-perfect である、
\item $K$ は $R$-平坦な有限表示 $A$-加群からなる有限複体と同型である。
\end{enumerate}
\end{lemma}

\begin{proof}
(2) が (1) を含意することを示すには、Lemma
\ref{more-algebra-lemma-cone-relatively-perfect} により、$R$-平坦な有限表示
$A$-加群 $M$ が $D(A)$ の $R$-perfect 対象を定めることを示せば十分である。
$M$ は $R$ 上有限 Tor 次元を持つので、$M$ が擬コヒーレントであることを
示せば十分である。Algebra, Lemma
\ref{algebra-lemma-flat-finite-presentation-limit-flat} により、
有限型 $\mathbf{Z}$-代数 $R_0 \subset R$、平坦有限型環写像
$R_0 \to A_0$、および $R_0$ 上平坦な有限 $A_0$-加群 $M_0$ で、
$A = A_0 \otimes_{R_0} R$ および $M = M_0 \otimes_{R_0} R$ を
満たすものが存在する。Lemma \ref{more-algebra-lemma-Noetherian-pseudo-coherent} により、
% P02-E1033: frozen source line 22002 omits the article before "pseudo-coherent A_0-module".
$M_0$ は擬コヒーレント $A_0$-加群である。
有限自由 $A_0$-加群 $P_0^n$ による分解
$P_0^\bullet \to M_0$ を選ぶ。$A_0$ は $R_0$ 上平坦なので、
これは平坦分解である。$M_0$ は $R_0$ 上平坦なので、
$P^\bullet = P_0^\bullet \otimes_{R_0} R$ はなお
$M = M_0 \otimes_{R_0} R$ を分解する
（これは Lemma \ref{more-algebra-lemma-base-change-comparison} を用いて確認できる）。
従って $P^\bullet$ は $A$ 上の $M$ の有限自由分解であり、
$M$ は擬コヒーレントであると結論する。

\medskip\noindent
(1) を仮定する。$K$ は有限自由 $A$-加群からなる上に有界な複体
$P^\bullet$ で表せる。$D(R)$ の対象とみなした $K$ が
$[a, b]$ に Tor 振幅を持つと仮定する。Lemma
\ref{more-algebra-lemma-last-one-flat} により、$\tau_{\geq a}P^\bullet$ は
$K$ を表す $R$-平坦な有限表示 $A$-加群の複体である。
\end{proof}

\begin{lemma}
\label{more-algebra-lemma-base-change-relatively-perfect}
$R \to A$ を平坦かつ有限表示な環写像とする。
$R \to R'$ を環写像とし、$A' = A \otimes_R R'$ とおく。
$K \in D(A)$ が $R$-perfect ならば、
$K \otimes_A^\mathbf{L} A'$ は $R'$-perfect である。
\end{lemma}

\begin{proof}
Lemma \ref{more-algebra-lemma-pull-pseudo-coherent} により、
$K \otimes_A^\mathbf{L} A'$ は擬コヒーレントである。Lemma
\ref{more-algebra-lemma-base-change-comparison} により、
$K \otimes_A^\mathbf{L} A'$ は $D(R')$ において
$K \otimes_R^\mathbf{L} R'$ に等しい。そこで Lemma
\ref{more-algebra-lemma-pull-tor-amplitude} を適用すれば、$D(R')$ における
$K \otimes_R^\mathbf{L} R'$ が有限 Tor 次元を持つことが分かる。
\end{proof}

\begin{lemma}
\label{more-algebra-lemma-compute-RHom-relatively-perfect}
$R \to A$ を平坦環写像とする。$K, L \in D(A)$ とし、$K$ は
擬コヒーレント、$L$ は $R$ 上有限 Tor 次元を持つとする。次を選べる：
\begin{enumerate}
\item $K$ を表す有限自由 $A$-加群からなる上に有界な複体 $P^\bullet$、
\item $L$ を表す $R$-平坦 $A$-加群からなる有界複体 $F^\bullet$。
\end{enumerate}
これらを選ぶと、次が成り立つ：
\begin{enumerate}
\item[(a)] $E^\bullet = \Hom^\bullet(P^\bullet, F^\bullet)$ は
$R\Hom_A(K, L)$ を表す $R$-平坦 $A$-加群からなる下に有界な複体である、
\item[(b)] 任意の環写像 $R \to R'$ に対し、
$A' = A \otimes_R R'$ とおくと、複体 $E^\bullet \otimes_R R'$ は
$R\Hom_{A'}(K \otimes_A^\mathbf{L} A', L \otimes_A^\mathbf{L} A')$ を表す。
\end{enumerate}
さらに $R \to A$ が有限表示で $L$ が $R$-perfect ならば、
$F^p$ を有限表示 $A$-加群に選べ、その結果 $E^n$ も有限表示
$A$-加群となる。
\end{lemma}

\begin{proof}
$P^\bullet$ の存在は擬コヒーレント複体の定義である。
まず $L$ を自由 $A$-加群からなる上に有界な複体 $F^\bullet$ で表す
（Tor 次元が有界ならば、特に有界なので、これは可能である）。
次に、$D(R)$ の対象とみなした $L$ が $[a, b]$ に Tor 振幅を持つとする。
すると $F^\bullet$ を $\tau_{\geq a}F^\bullet$ で置き換えることにより、
(2) のような複体を得る。これは Lemma \ref{more-algebra-lemma-last-one-flat} から従う。

\medskip\noindent
(a) の証明。
% P02-E1034: frozen source line 22076 says "bounded an since" rather than "bounded and since".
$F^\bullet$ は有界で $P^\bullet$ は上に有界なので、
$n \ll 0$ ならば $E^n = 0$ であり、$E^n$ は有限（！）直和
$$
E^n = \bigoplus\nolimits_{p + q = n} \Hom_A(P^{-q}, F^p)
$$
である。$P^{-q}$ は有限自由なので、これは確かに $R$-平坦
$A$-加群である。$E^\bullet$ が $R\Hom_A(K, L)$ を表すことは
Lemma \ref{more-algebra-lemma-RHom-out-of-projective} から従う。

\medskip\noindent
(b) の証明。
$R \to R'$ を環写像とし、$A' = A \otimes_R R'$ とする。
Lemma \ref{more-algebra-lemma-base-change-comparison} により、
$L \otimes_A^\mathbf{L} A'$ は $A'$-加群の複体とみなした
$F^\bullet \otimes_R R'$ で表される（$F^p$ が $R$ 上平坦であることによる）。
$P^\bullet \otimes_R R'$ についても同様である。上と同じく
$R\Hom_{A'}(K \otimes_A^\mathbf{L} A', L \otimes_A^\mathbf{L} A')$ は
$$
\Hom^\bullet(P^\bullet \otimes_R R', F^\bullet \otimes_R R') =
E^\bullet \otimes_R R'
$$
で表される。この等式は複体の各項を個別に見て、
$$
\Hom_{A'}(P^{-q} \otimes_R R', F^p \otimes_R R') =
\Hom_A(P^{-q}, F^p) \otimes_R R'
$$
を用いれば従う。
\end{proof}

\begin{lemma}
\label{more-algebra-lemma-colimit-relatively-perfect}
$R = \colim_{i \in I} R_i$ を環のフィルター余極限とする。
$0 \in I$ とし、$R_0 \to A_0$ を平坦かつ有限表示な環写像とする。
$i \geq 0$ に対し $A_i = R_i \otimes_{R_0} A_0$ とおき、
$A = R \otimes_{R_0} A_0$ とおく。
\begin{enumerate}
\item $D(A)$ の $R$-perfect な $K$ が与えられると、
$i \in I$ と $D(A_i)$ の $R_i$-perfect な $K_i$ で、
$D(A)$ において $K \cong K_i \otimes_{A_i}^\mathbf{L} A$ を
満たすものが存在する。
\item $K_0, L_0 \in D(A_0)$ が与えられ、$K_0$ は擬コヒーレントで、
$L_0$ は $R_0$ 上有限 Tor 次元を持つとすると、
$$
\Hom_{D(A)}(K_0 \otimes_{A_0}^\mathbf{L} A, L_0 \otimes_{A_0}^\mathbf{L} A) =
\colim_{i \geq 0}
\Hom_{D(A_i)}(K_0 \otimes_{A_0}^\mathbf{L} A_i,
L_0 \otimes_{A_0}^\mathbf{L} A_i)
$$
が成り立つ。
\end{enumerate}
特に、$A$ 上の $R$-perfect 複体の三角圏は、$A_i$ 上の
$R_i$-perfect 複体の三角圏の余極限である。
\end{lemma}

\begin{proof}
Algebra, Lemma \ref{algebra-lemma-colimit-category-fp-modules} により、
有限表示 $A$-加群の圏は有限表示 $A_i$-加群の圏の余極限である。
これを踏まえると、Algebra, Lemma
\ref{algebra-lemma-flat-finite-presentation-limit-flat} により、
% P02-E1035: frozen source line 22135 omits "the" before "category".
$R$-平坦な有限表示 $A$-加群の圏は、$R_i$-平坦な有限表示
$A_i$-加群の圏の余極限である。従って Lemma
\ref{more-algebra-lemma-structure-relatively-perfect} の特徴づけにより (1) が従う。

\medskip\noindent
(2) を示すため、Lemma \ref{more-algebra-lemma-compute-RHom-relatively-perfect} のように、
$K_0$ を表す $P_0^\bullet$ と $L_0$ を表す $F_0^\bullet$ を選ぶ。
すると $E_0^\bullet = \Hom^\bullet(P_0^\bullet, F_0^\bullet)$ は
$$
H^0(E_0^\bullet \otimes_{R_0} R_i) =
\Hom_{D(A_i)}(K_0 \otimes_{A_0}^\mathbf{L} A_i,
L_0 \otimes_{A_0}^\mathbf{L} A_i)
$$
および
$$
H^0(E_0^\bullet \otimes_{R_0} R) =
\Hom_{D(A)}(K_0 \otimes_{A_0}^\mathbf{L} A, L_0 \otimes_{A_0}^\mathbf{L} A)
$$
を満たす。
% P02-E1036: frozen source line 22158 says "Thus the result because" and omits the verb.
これは補題から従う。従って、テンソル積が余極限と可換で、フィルター余極限が
完全であること（Algebra, Lemma \ref{algebra-lemma-directed-colimit-exact}）から
結果が従う。
\end{proof}

\begin{lemma}
\label{more-algebra-lemma-thickening-relatively-perfect}
$R' \to A'$ を平坦かつ有限表示な環写像とする。
$R' \to R$ を、核が冪零イデアルである全射環写像とする。
$A = A' \otimes_{R'} R$ とおく。$K' \in D(A')$ とし、
$D(A)$ において $K = K' \otimes_{A'}^\mathbf{L} A$ とおく。
$K$ が $R$-perfect ならば、$K'$ は $R'$-perfect である。
\end{lemma}

\begin{proof}
$A' \to A$ は冪零核を持つこと、また $R' \to A'$ の平坦性により
$K = K' \otimes_{R'}^\mathbf{L} R$ であることに注意する
（Section \ref{more-algebra-section-tor-independence} を見よ）。従って、この補題は
Lemmas \ref{more-algebra-lemma-check-pseudo-coherent-modulo-nilpotent} and
\ref{more-algebra-lemma-check-tor-dimension-modulo-nilpotent} を組み合わせれば従う。
\end{proof}

\begin{lemma}
\label{more-algebra-lemma-lift-from-fibre-relatively-perfect}
$R$ を環とする。$A = R[x_1, \ldots, x_d]/I$ は
$R$ 上平坦かつ有限表示であるとする。
$\mathfrak q \subset A$ を $\mathfrak p \subset R$ の上にある素イデアルとする。
$K \in D(A)$ は擬コヒーレントであるとする。$a, b \in \mathbf{Z}$ とする。
$H^i(K_\mathfrak q \otimes_{R_\mathfrak p}^\mathbf{L} \kappa(\mathfrak p))$ が
$i \in [a, b]$ に対してのみ非零ならば、$K_\mathfrak q$ は
$R$ 上 $[a - d, b]$ に Tor 振幅を持つ。
\end{lemma}

\begin{proof}
Lemma \ref{more-algebra-lemma-more-relative-pseudo-coherent-is-moot} により、
$K$ は $R[x_1, \ldots, x_d]$-加群の複体として擬コヒーレントである。
従って $A = R[x_1, \ldots, x_d]$ と仮定してよい。仮定を用い、
$R_\mathfrak p \to A_\mathfrak q$ と複体 $K_\mathfrak q$ に Lemma
\ref{more-algebra-lemma-perfect-over-regular-local-ring} を適用すると、
$K_\mathfrak q$ は $D(A_\mathfrak q)$ において perfect で、
$[a - d, b]$ に Tor 振幅を持つことが分かる。
$R_\mathfrak p \to A_\mathfrak q$ は平坦なので、Lemma
\ref{more-algebra-lemma-flat-push-tor-amplitude} により結論を得る。
\end{proof}

\begin{lemma}
\label{more-algebra-lemma-bounded-on-fibres-relatively-perfect}
$R \to A$ を平坦かつ有限表示な環写像とする。
$K \in D(A)$ は擬コヒーレントであるとする。次は同値である：
\begin{enumerate}
\item $K$ は $R$-perfect である、
\item $K$ は下に有界で、すべての素イデアル $\mathfrak p \subset R$ に対して
$K \otimes_R^\mathbf{L} \kappa(\mathfrak p)$ は下に有界である。
\end{enumerate}
\end{lemma}

\begin{proof}
(1) は (2) を含意することに注意する。実際、$R$-perfect 複体は定義により、
$R$-加群の複体として有界 Tor 次元を持つ。逆の含意を示そう。

\medskip\noindent
$A = R[x_1, \ldots, x_d]/I$ と書く。
% P02-E1037: frozen source line 22221 omits "by" in "Denote L ... the restriction".
$D(R[x_1, \ldots, x_d])$ における $K$ の制限を $L$ と表す。Lemma
\ref{more-algebra-lemma-more-relative-pseudo-coherent-is-moot} により、
$L$ は擬コヒーレントである。$L$ と $K$ は $D(R)$ において同じ像を持つので、
$L$ が $R$-perfect であることと $K$ が $R$-perfect であることは同値である。
また $L \otimes_R^\mathbf{L} \kappa(\mathfrak p)$ と
$K \otimes_R^\mathbf{L} \kappa(\mathfrak p)$ は
$D(\kappa(\mathfrak p))$ の同じ対象である。従って
$A = R[x_1, \ldots, x_d]$ の場合に帰着する。

\medskip\noindent
$A = R[x_1, \ldots, x_d]$ で $K$ が (2) を満たすとする。
$\mathfrak q \subset A$ を素イデアル $\mathfrak p \subset R$ の上にある
素イデアルとする。仮定を用い、$R_\mathfrak p \to A_\mathfrak q$ と
複体 $K_\mathfrak q$ に Lemma
\ref{more-algebra-lemma-perfect-over-regular-local-ring} を適用すると、
$K_\mathfrak q$ は $D(A_\mathfrak q)$ において perfect であることが分かる。
$K$ は下に有界なので、Lemma \ref{more-algebra-lemma-check-perfect-stalks} により
$K$ は $D(A)$ において perfect である。従って Lemma
\ref{more-algebra-lemma-perfect-relatively-perfect} により $K$ は $R$-perfect であり、
証明が完了する。
\end{proof}

\section{二項複体}
\label{more-algebra-section-two-term}

\noindent
この節では、素朴余接複体に関する条件を理解する助けとなる、
加群の二項複体についてのいくつかの結果を証明する。

\begin{lemma}
\label{more-algebra-lemma-ext-1-zero}
$R$ を環とする。$K \in D(R)$ とし、$i \not \in \{-1, 0\}$ に対して
$H^i(K) = 0$ とする。次は同値である：
\begin{enumerate}
\item $H^{-1}(K) = 0$ で、$H^0(K)$ は射影加群である、
\item すべての $R$-加群 $M$ に対して $\Ext^1_R(K, M) = 0$ である。
\end{enumerate}
$R$ が Noether 環で、$i = -1, 0$ に対して $H^i(K)$ が有限
$R$-加群ならば、これらは次とも同値である：
\begin{enumerate}
\item[(3)] すべての有限 $R$-加群 $M$ に対して
$\Ext^1_R(K, M) = 0$ である。
\end{enumerate}
\end{lemma}

\begin{proof}
(1) と (2) の同値性は Lemma \ref{more-algebra-lemma-projective-amplitude} から従う。
$R$ が Noether 環で、$i = -1, 0$ に対して $H^i(K)$ が有限
$R$-加群ならば、Lemma \ref{more-algebra-lemma-Noetherian-pseudo-coherent} により
$K$ は擬コヒーレントである。従って (1) と (3) の同値性は
Lemma \ref{more-algebra-lemma-projective-amplitude-pseudo-coherent} から従う。
\end{proof}

\begin{remark}
\label{more-algebra-remark-smoothness-ext-1-zero}
次の二つの主張は Lemma \ref{more-algebra-lemma-ext-1-zero},
Algebra, Definition \ref{algebra-definition-smooth}, and
Algebra, Proposition \ref{algebra-proposition-characterize-formally-smooth}
から従う。
\begin{enumerate}
\item 環写像 $A \to B$ が滑らかであることと、$A \to B$ が有限表示で、
すべての $B$-加群 $N$ に対して $\Ext^1_B(\NL_{B/A}, N) = 0$ であることは
同値である。
\item 環写像 $A \to B$ が形式的に滑らかであることと、すべての
$B$-加群 $N$ に対して $\Ext^1_B(\NL_{B/A}, N) = 0$ であることは同値である。
\end{enumerate}
\end{remark}

\begin{lemma}
\label{more-algebra-lemma-represent-two-term-complex}
$R$ を環とする。$K$ を $D(R)$ の対象とし、$i \not \in \{-1, 0\}$ に対して
$H^i(K) = 0$ とする。このとき
\begin{enumerate}
\item $K$ は、$K^0$ が自由加群である二項複体 $K^{-1} \to K^0$ で表せる、
\item $R$ が Noether 環で、$i = -1, 0$ に対して $H^i(K)$ が有限
$R$-加群ならば、$K$ は、$K^0$ が有限自由で $K^{-1}$ が有限である
二項複体 $K^{-1} \to K^0$ で表せる。
\end{enumerate}
\end{lemma}

\begin{proof}
(1) の証明。$K$ が加群の複体 $M^\bullet$ で与えられているとする。
まず $M^\bullet$ を $\tau_{\leq 0}M^\bullet$ で置き換えてよい。従って
% P02-E1038: frozen source line 22319 uses a comma before the new sentence "Next".
$i > 0$ に対して $M^i = 0$ と仮定してよい。次に、$i > 0$ に対して
$P^i = 0$ である自由分解 $P^\bullet \to M^\bullet$ を選べる；
Derived Categories, Lemma
\ref{derived-lemma-subcategory-left-resolution} を見よ。
最後に $K^\bullet = \tau_{\geq -1}P^\bullet$ とおけばよい。

\medskip\noindent
(2) の証明。$R$ は Noether 環で、$i = -1, 0$ に対して $H^i(K)$ は有限
$R$-加群であると仮定する。Lemma \ref{more-algebra-lemma-pseudo-coherent} により、
$i > 0$ に対して $F^i = 0$ で、$F^i$ が有限自由である擬同型
$F^\bullet \to M^\bullet$ を選べる。そこで
$K^\bullet = \tau_{\geq -1}F^\bullet$ とおけばよい。
\end{proof}

\noindent
素朴余接複体 $\NL_{B/A}$ または $\NL(\alpha)$ から出る導来圏の写像
（Algebra, Section \ref{algebra-section-netherlander} を見よ）は、
次の補題の結果により容易に理解できる。

\begin{lemma}
\label{more-algebra-lemma-map-out-of-almost-free}
$R$ を環とする。$M^\bullet$ を $R$ 上の加群の複体とし、
$i > 0$ に対して $M^i = 0$、$M^0$ は射影 $R$-加群であるとする。
$K^\bullet$ を別の複体とする。
\begin{enumerate}
\item $i \leq -2$ に対して $K^i = 0$ と仮定する。このとき
$\Hom_{D(R)}(M^\bullet, K^\bullet) = \Hom_{K(R)}(M^\bullet, K^\bullet)$ である。
\item $i \not \in [-1, 0]$ に対して $K^i = 0$ で、$K^0$ は射影
$R$-加群であると仮定する。複体の写像
$a^\bullet : M^\bullet \to K^\bullet$ に対して、次は同値である：
\begin{enumerate}
\item $a^\bullet$ はすべての $R$-加群 $N$ に対して零写像
$\Ext^1_R(K^\bullet, N) \to \Ext^1_R(M^\bullet, N)$ を誘導する、
\item 写像 $h^0 : M^0 \to K^{-1}$ で
% P02-E1039: frozen source line 22352 composes h^0 : M^0 -> K^{-1} with d_K^{-1}, which is type-incompatible; d_M^{-1} is required.
$a^{-1} + h^0 \circ d^{-1}_K = 0$ を満たすものが存在する。
\end{enumerate}
\item $i \leq -3$ に対して $K^i = 0$ と仮定する。
$\alpha \in \Hom_{D(R)}(M^\bullet, K^\bullet)$ とする。
$\alpha$ と $K^\bullet \to K^{-2}[2]$ の合成が、
$a \circ d_M^{-3} = 0$ を満たす $R$-加群写像
$a : M^{-2} \to K^{-2}$ から来るならば、$\alpha$ は
$a^{-2} = a$ を満たす複体の写像
$a^\bullet : M^\bullet \to K^\bullet$ で表せる。
% P02-E1040: frozen source line 22361 says "In (2)", but alpha, a, and degree -2 are introduced only in part (3).
\item (2) において、$a = (a')^{-2}$ を満たし $\alpha$ を表す任意の第二の
複体の写像 $(a')^\bullet : M^\bullet \to K^\bullet$ に対し、
$i = 0, -1$ について写像 $h^i : M^i \to K^{i - 1}$ で
$$
h^{-1} \circ d_M^{-2} = 0, \quad
(a')^{-1} = a^{-1} + d_K^{-2} \circ h^{-1} + h^0 \circ d_M^{-1},\quad
(a')^0 = a^0 + d_K^{-1} \circ h^0
$$
を満たすものが存在する。
\end{enumerate}
\end{lemma}

\begin{proof}
$F^0 = M^0$ とおく。自由 $R$-加群 $F^{-1}$ と全射
$F^{-1} \to M^{-1}$ を選ぶ。自由 $R$-加群 $F^{-2}$ と全射
$F^{-2} \to M^{-2} \times_{M^{-1}} F^{-1}$ を選ぶ。
このように続けると、各 $i$ に対して $F^i$ が射影で、項ごとに全射な擬同型
$p^\bullet : F^\bullet \to M^\bullet$ を得る。

\medskip\noindent
(1) の証明。Derived Categories, Lemma
\ref{derived-lemma-morphisms-from-projective-complex} により
$$
\Hom_{D(R)}(M^\bullet, K^\bullet) = \Hom_{K(R)}(F^\bullet, K^\bullet)
$$
である。$i \leq -2$ に対して $K^i = 0$ ならば、複体の任意の射
$F^\bullet \to K^\bullet$ は $p^\bullet$ を通じて分解する。同様に、
任意のホモトピー $\{h^i : F^i \to K^{i - 1}\}$ は
$p^\bullet$ を通じて分解する。従って (1) が成り立つ。

\medskip\noindent
(2) の証明。(2)(b) が成り立つならば、$a^\bullet$ は次数 $-1$ で零である
複体の写像 $(a')^\bullet : M^\bullet \to K^\bullet$ とホモトピックである。
一方、$N \to I^\bullet$ を入射分解とする。Derived Categories, Lemma
\ref{derived-lemma-morphisms-into-injective-complex} により
$$
\Ext^1_R(K^\bullet, N) = \Hom_{D(R)}(K^\bullet, I^\bullet[1]) =
\Hom_{K(R)}(K^\bullet, I^\bullet[1])
$$
である。$b^\bullet : K^\bullet \to I^\bullet[1]$ を複体の写像とする。
$K^1 = 0$ なので、写像 $b^0 : K^0 \to I^1$ の像は
$I^1 \to I^2$ の核、すなわち $I^0 \to I^1$ の像に含まれる。
$K^0$ は射影なので、$b^0$ を写像 $h : K^0 \to I^0$ に持ち上げられる。
従って $b^\bullet$ は $(b')^0 = 0$ を満たす複体の写像
$(b')^\bullet$ とホモトピックである。$i \not \in [-1, 0]$ に対して
$K^i = 0$ なので、$(b')^\bullet \circ (a')^\bullet = 0$ が
複体の写像として成り立つ。従って写像
$\Ext^1_R(K^\bullet, N) \to \Ext^1_R(M^\bullet, N)$ は零である。
このようにして (2)(b) が (2)(a) を含意することが分かる。
逆に (2)(a) を仮定する。$\Ext^1_R(K^\bullet, K^{-1})$ の標準元は
$\Ext^1_R(M^\bullet, K^{-1})$ において零に写る。(1) を用いると、
(2)(b) のような写像 $h^0$ を直ちに得ることが分かる。

\medskip\noindent
(3) の証明。$\alpha$ を表す $b^\bullet : F^\bullet \to K^\bullet$ を選ぶ。
$\alpha$ と $K^\bullet \to K^{-2}[2]$ の合成は
$b^{-2} : F^{-2} \to K^{-2}$ で表される。これは
$a \circ p^{-2} : F^{-2} \to M^{-2} \to K^{-2}$ とホモトピックなので、
$b^{-2} = a \circ p^{-2} + h \circ d_F^{-2}$ を満たす写像
$h : F^{-1} \to K^{-2}$ が存在する。$F^\bullet$ から
$K^\bullet$ へのホモトピーとみなした $h$ により $b^\bullet$ を調整すると、
$b^{-2} = a \circ p^{-2}$ となる。従って $b^{-2}$ は
$p^{-2}$ を通じて分解する。$F^0 = M^0$ なので、$p^{-2}$ の核は
$p^{-1}$ の核に全射する（例えば $p^\bullet$ の核が非輪状複体であること、
または図式追跡による）。従って $b^{-1}$ も必ず $p^{-1}$ を通じて分解し、
これらの分解と $a^0 = b^0$ に対して (3) が成り立つことが分かる。

\medskip\noindent
(4) の証明は省略する。ヒント：$a^\bullet \circ p^\bullet$ と
$(a')^\bullet \circ p^\bullet$ の間にホモトピーがあり、前と同様に、
このホモトピーが $p^\bullet$ を通じて分解することを示す。
\end{proof}

\noindent
$A \to B$ を有限表示環写像とする。イデアル $I \subset B$ が与えられたとき、
次の条件を考えられる：
\begin{enumerate}
\item[(*)] すべての $B$-加群 $N$ に対して
$\Ext^1_B(\NL_{B/A}, N)$ は $I$ で零化される。
\end{enumerate}
% P02-E1041: frozen source line 22447 uses the open compound "scheme theoretically".
この条件は「$A \to B$ の特異点集合がスキーム論的に $V(I)$ に含まれる」
という概念の、一つの可能な精密な数学的定式化である。
Remark \ref{more-algebra-remark-smoothness-ext-1-zero} および次の補題と比較されたい。

\begin{lemma}
\label{more-algebra-lemma-ext-1-annihilated-definite}
$R$ を環、$I \subset R$ をイデアルとする。$K \in D(R)$ とし、
$i \not \in \{-1, 0\}$ に対して $H^i(K) = 0$ と仮定する。
次は同値である：
\begin{enumerate}
\item すべての $R$-加群 $N$ に対して $\Ext^1_R(K, N)$ は
$I$ で零化される、
\item $K$ は、$K^0$ が自由で、任意の $a \in I$ に対して写像
$a : K^{-1} \to K^{-1}$ が $d_K^{-1} : K^{-1} \to K^0$ を通じて
分解するような複体 $K^{-1} \to K^0$ で表せる、
\item $K$ が $K^0$ を射影とする二項複体 $K^{-1} \to K^0$ で
表されるときは常に、任意の $a \in I$ に対して写像
$a : K^{-1} \to K^{-1}$ は $d_K^{-1} : K^{-1} \to K^0$ を通じて分解する。
\end{enumerate}
$R$ が Noether 環で、$i = -1, 0$ に対して $H^i(K)$ が有限
$R$-加群ならば、これらは次とも同値である：
\begin{enumerate}
\item[(4)] すべての有限 $R$-加群 $N$ に対して $\Ext^1_R(K, N)$ は
$I$ で零化される、
\item[(5)] $K$ は、$K^0$ が有限自由で $K^{-1}$ が有限であり、
任意の $a \in I$ に対して写像 $a : K^{-1} \to K^{-1}$ が
$d_K^{-1} : K^{-1} \to K^0$ を通じて分解するような複体
$K^{-1} \to K^0$ で表せる。
\end{enumerate}
\end{lemma}

\begin{proof}

\medskip\noindent
(1) を仮定し、$K$ を表す二項複体 $K^{-1} \to K^0$ で
$K^0$ が射影であるものを取る。以下では Lemma
\ref{more-algebra-lemma-map-out-of-almost-free} が与える $D(R)$ において
$K^\bullet$ から出る写像の記述を断りなく用いる。
$N = K^{-1}$ と選び、$\text{id}_{K^{-1}} : K^{-1} \to K^{-1}$ が与える
$\Ext^1_R(K, N)$ の元 $\xi$ を考える。
% P02-E1042: frozen source line 22485 says "Since is annihilated" and omits xi as the subject.
$a \in I$ によりこれは零化されるので、次の可換図に入る点線の矢印を得る：
$$
\xymatrix{
K^{-1} \ar[d]_a \ar[r]_{d_K^{-1}} & K^0 \ar@{..>}[ld]^h \\
K^{-1}
}
$$
これで (3) が成り立つことが分かる。Lemma
\ref{more-algebra-lemma-represent-two-term-complex} part (1) により、(3) は (2) を含意する。
$K^\bullet$ は (2) のようなものとし、$N$ を任意の $R$-加群とする。
$\Ext^1_R(K, N)$ の任意の元 $\xi$ は、写像
$\varphi : K^{-1} \to N$ の類で与えられる。そこで $a \in I$ に対し、
仮定により上の図のような写像 $h$ を選べるので、
$a\varphi = \varphi \circ a = \varphi \circ h \circ \text{d}_K^{-1}$ となる。
これは $a \xi$ が $\Ext^1_R(K, N)$ において零であることを示す。
従って (1), (2), (3) は同値である。

\medskip\noindent
$R$ は Noether 環で、$i = -1, 0$ に対して $H^i(K)$ は有限
$R$-加群であると仮定する。Lemma
\ref{more-algebra-lemma-represent-two-term-complex} part (2) により (3) は (5) を含意する。
(5) が (2) を含意することは明らかである。(1) は自明に (4) を含意する。
従って証明を終えるには、(4) が他の条件のいずれかを含意することを示せばよい。
Lemma \ref{more-algebra-lemma-represent-two-term-complex} part (2) のように、$K$ を表す
複体 $K^{-1} \to K^0$ で、$K^0$ が有限自由、$K^{-1}$ が有限であるものを取る。
(2) $\Rightarrow$ (1) の証明で与えた議論により、
$\Ext^1_R(K, K^{-1})$ が $I$ で零化されるならば (1) が成り立つ。
このようにして (4) が (1) を含意することが分かり、証明が完了する。
\end{proof}

\begin{lemma}
\label{more-algebra-lemma-two-term-base-change}
$R$ を環とする。$K$ を $D(R)$ の対象とし、$i \not \in \{-1, 0\}$ に対して
$H^i(K) = 0$ とする。$K$ を表す $R$-加群の二項複体
$K^{-1} \to K^0$ で、$K^0$ が平坦 $R$-加群（例えば射影または自由）である
ものを取る。$R \to R'$ を環写像とする。このとき複体
$K^\bullet \otimes_R R'$ は $\tau_{\geq -1}(K \otimes_R^\mathbf{L} R')$ を表す。
\end{lemma}

\begin{proof}
$D(R)$ において完全三角
$$
K^0 \to K^\bullet \to K^{-1}[1] \to K^0[1]
$$
がある。これは完全三角の写像
$$
\xymatrix{
K^0 \otimes_R^\mathbf{L} R' \ar[d] \ar[r] &
K^\bullet \otimes_R^\mathbf{L} R' \ar[r] \ar[d] &
K^{-1} \otimes_R^\mathbf{L} R'[1] \ar[r] \ar[d] &
K^0 \otimes_R^\mathbf{L} R'[1] \ar[d] \\
K^0 \otimes_R R' \ar[r] &
K^\bullet \otimes_R R' \ar[r] &
K^{-1} \otimes_R R'[1] \ar[r] &
K^0 \otimes_R R'[1]
}
$$
を定める。$K^0$ は平坦なので、左端と右端の縦矢印は同型である。
$K^{-1} \otimes_R^\mathbf{L} R' \to K^{-1} \otimes_R R'$ は
次数 $0$ のコホモロジー上で同型なので、結論を得る。
\end{proof}

\begin{lemma}
\label{more-algebra-lemma-base-change-property-ext-1-annihilated}
$I$ を環 $R$ のイデアルとする。$K$ を $D(R)$ の対象とし、
$i \not \in \{-1, 0\}$ に対して $H^i(K) = 0$ とする。
$R \to R'$ を環写像とする。$K$ が $(R, I)$ に関して Lemma
\ref{more-algebra-lemma-ext-1-annihilated-definite} の同値な条件 (1), (2), (3) を満たすならば、
$\tau_{\geq -1}(K \otimes_R^\mathbf{L} R')$ は $(R', IR')$ に関して Lemma
\ref{more-algebra-lemma-ext-1-annihilated-definite} の同値な条件 (1), (2), (3) を満たす。
% P02-E1043: frozen source line 22560 omits terminal punctuation before the lemma ends.
\end{lemma}

\begin{proof}
$K$ は、$K^0$ が自由で、任意の $a \in I$ に対して写像
$a : K^{-1} \to K^{-1}$ がある写像 $h_a : K^0 \to K^{-1}$ により
$h_a \circ d_K^{-1}$ に等しいような二項複体 $K^{-1} \to K^0$ で
表されると仮定してよい。Lemma \ref{more-algebra-lemma-two-term-base-change} により、
$\tau_{\geq -1}(K \otimes_R^\mathbf{L} R')$ は
$K^\bullet \otimes_R R'$ で表される。このとき当然、すべての $a \in I$ に対し
$a \otimes 1 : K^{-1} \otimes_R R' \to K^{-1} \otimes_R R'$ は
$(h_a \otimes 1) \circ (d_K^{-1} \otimes 1)$ に等しい。
$\text{d}_K^{-1} \otimes 1$ を通じて分解する写像
$K^{-1} \otimes_R R' \to K^{-1} \otimes_R R'$ の集まりは
$R'$-加群をなすので、結論を得る。
\end{proof}

\begin{lemma}
\label{more-algebra-lemma-two-term-surjection-map-zero}
$R$ を環とする。$\alpha : K \to K'$ を $D(R)$ の射とする。次を仮定する：
\begin{enumerate}
\item $i \not \in \{-1, 0\}$ に対して $H^i(K) = H^i(K') = 0$ である、
\item $H^0(\alpha)$ は同型で、$H^{-1}(\alpha)$ は全射である。
\end{enumerate}
任意の $f \in R$ に対し、$f : K \to K$ が $0$ ならば、
$f : K' \to K'$ も $0$ である。
\end{lemma}

\begin{proof}
$M = \Ker(H^{-1}(\alpha))$ とおく。このとき $\alpha$ は完全三角
$$
M[1] \to K \to K' \to M[2]
$$
に入る。仮定により $K \to K' \xrightarrow{f} K'$ は零なので、
% P02-E1044: frozen source line 22594 says the map "factors over" another map rather than "factors through" it.
$f : K' \to K'$ は写像 $M[2] \to K'$ を介して分解する。しかし例えば
Derived Categories, Lemma \ref{derived-lemma-negative-exts} により
$\Hom(M[2], K') = 0$ である。
\end{proof}

\begin{lemma}
\label{more-algebra-lemma-surjection-property-ext-1-annihilated}
$I$ を環 $R$ のイデアルとする。$\alpha : K \to K'$ を
$D(R)$ の射とする。次を仮定する：
\begin{enumerate}
\item $i \not \in \{-1, 0\}$ に対して $H^i(K) = H^i(K') = 0$ である、
\item $H^0(\alpha)$ は同型で、$H^{-1}(\alpha)$ は全射である。
\end{enumerate}
$K$ が Lemma \ref{more-algebra-lemma-ext-1-annihilated-definite} の同値な条件
(1), (2), (3) を満たすならば、$K'$ も満たす。
\end{lemma}

\begin{proof}
$M = \Ker(H^{-1}(\alpha))$ とおく。このとき $\alpha$ は完全三角
$$
M[1] \to K \to K' \to M[2]
$$
に入る。任意の $R$-加群 $N$ に対し、これは完全列
$$
\Ext^0_R(M[1], N) \to
\Ext^1_R(K', N) \to
\Ext^1_R(K, N)
$$
を定める。$\Ext^0_R(M[1], N) = \Ext^{-1}_R(M, N) = 0$ なので、
$\Ext^1_R(K', N)$ は $\Ext^1_R(K, N)$ の部分加群である。
従って $\Ext^1_R(K, N)$ が $I$ で零化されるならば、
$\Ext^1_R(K', N)$ もそうである。
\end{proof}

\begin{lemma}
\label{more-algebra-lemma-ext-1-annihilated}
% P02-E1045: frozen source line 22631 says "Let R be ring" and omits the article.
$R$ を環、$I \subset R$ をイデアルとする。$K \in D(R)$ とし、
$i \not \in \{-1, 0\}$ に対して $H^i(K) = 0$ とする。次は同値である：
\begin{enumerate}
\item $c \geq 0$ で、$K$ とイデアル $I^c$ に対して Lemma
\ref{more-algebra-lemma-ext-1-annihilated-definite} の同値な条件 (1), (2), (3) が
成り立つものが存在する、
\item $c \geq 0$ で、(a) $I^c$ が $H^{-1}(K)$ を零化し、かつ
(b) $H^0(K)$ が $I^c$-射影加群であるものが存在する
（Section \ref{more-algebra-section-near-projective} を見よ）。
\end{enumerate}
$R$ が Noether 環で、$i = -1, 0$ に対して $H^i(K)$ が有限
$R$-加群ならば、これらは次とも同値である：
\begin{enumerate}
\item[(3)] $c \geq 0$ で、$K$ とイデアル $I^c$ に対して Lemma
\ref{more-algebra-lemma-ext-1-annihilated-definite} の同値な条件 (4), (5) が
成り立つものが存在する、
\item[(4)] $H^{-1}(K)$ は $I$-冪捩れであり、
$V(f_1, \ldots, f_s) \subset V(I)$ を満たす
$f_1, \ldots, f_s \in R$ で、局所化 $H^0(K)_{f_i}$ が射影
$R_{f_i}$-加群となるものが存在する、
\item[(5)] $H^{-1}(K)$ は $I$-冪捩れであり、
$V(f_1, \ldots, f_s) = V(I)$ を満たす
$f_1, \ldots, f_s \in I$ で、局所化 $H^0(K)_{f_i}$ が射影
$R_{f_i}$-加群となるものが存在する、
\item[(6)] $H^{-1}(K)$ は $I$-冪捩れであり、
$V(f_1, \ldots, f_s) = V(I)$ を満たす任意の
$f_1, \ldots, f_s \in I$ に対して、局所化 $H^0(K)_{f_i}$ は射影
$R_{f_i}$-加群である。
\end{enumerate}
\end{lemma}

\begin{proof}
完全三角 $H^{-1}(K)[1] \to K \to H^0(K)[0] \to H^{-1}(K)[2]$ は
完全列
$$
0 \to \Ext^1_R(H^0(K), N) \to \Ext^1_R(K, N) \to \Hom_R(H^{-1}(K), N) \to
\Ext^2_R(H^0(K), N)
$$
を定める。従って (2) は、すべての $R$-加群 $N$ に対して
$I^{2c}$ が $\Ext^1_R(K, N)$ を零化することを含意する。
(1) を仮定すると、$H^0(K)$ が $I^c$-射影であることが直ちに分かる。
一方、ある入射 $R$-加群 $N$ への単射 $H^{-1}(K) \to N$ を選べる。
完全列の $\Ext^2$ が消えることにより、この写像は
$\Ext^1_R(K, N)$ の元の像であり、$H^{-1}(K)$ は $I^c$ で
零化されると結論する。

\medskip\noindent
$R$ は Noether 環で、$i = -1, 0$ に対して $H^i(K)$ は有限
$R$-加群であると仮定する。Lemma
\ref{more-algebra-lemma-ext-1-annihilated-definite} により、(3) は (1), (2) と同値である。
また (3) が成り立つならば、
% P02-E1046: frozen source lines 22681-22683 use multiplication by f, although condition (3) only supplies factorization for I^c (hence for a power of f).
$f \in I$ に対する $H^0(K)$ 上の $f$ 倍写像は射影加群を通じて分解する。
従って $H^0(K)_f$ は射影 $R_f$-加群の直和因子であり、それ自身が射影
$R_f$-加群である。このように、同値な条件 (1), (2), (3) は (6) を含意する。
もちろん (6) は (5) を含意し、(5) は自明に (4) を含意する。

\medskip\noindent
(4) を仮定する。$H^{-1}(K)$ は有限 $R$-加群で $I$-冪捩れなので、
ある $c_1 \geq 0$ に対して $I^{c_1}$ は $H^{-1}(K)$ を零化する。
短完全列
$$
0 \to M \to R^{\oplus r} \to H^0(K) \to 0
$$
を選ぶ。これは元 $\xi \in \Ext^1_R(H^0(K), M)$ を定める。
任意の $f \in I$ に対し、Lemma
\ref{more-algebra-lemma-pseudo-coherence-and-base-change-ext} により
$\Ext^1_R(H^0(K), M)_f = \Ext^1_{R_f}(H^0(K)_f, M_f)$ である。
従って $H^0(K)_f$ が射影ならば、$f$ のある冪は $\xi$ を零化する。
よって、ある $c_2 \geq 0$ に対して $(f_1, \ldots, f_s)^{c_2}$ は
$\xi$ を零化する。$V(f_1, \ldots, f_s) \subset V(I)$ なので
% P02-E1047: frozen source line 22701 omits the radical on the right; the stated closed-set inclusion only gives sqrt(I) subset sqrt((f_i)).
$$
\sqrt{I} \subset (f_1, \ldots, f_s)
$$
である（Algebra, Lemma \ref{algebra-lemma-Zariski-topology}）。
$R$ は Noether 環なので、ある $c_3 \geq 0$ に対して
$I^{c_3} \subset (f_1, \ldots, f_s)$ である
（Algebra, Lemma \ref{algebra-lemma-Noetherian-power}）。
% P02-E1048: frozen source line 22705 writes the exponent c2c3 without the established subscripts.
従って $I^{c2c3}$ は $\xi$ を零化する。
% P02-E1049: frozen source lines 22706-22707 switch from I^{c_2c_3} to a in I and have agreement errors in "which annihilate ... factor".
これはさらに、$H^0(K)$ が $I^{c_2c_3}$-射影であることを意味する
（$\xi$ を零化する $a \in I$ による乗法は $R^{\oplus r}$ を通じて分解する）。
従って $c = \max(c_1, c_2c_3)$ と取れば (2) が成り立つことが分かる。
\end{proof}

\begin{lemma}
\label{more-algebra-lemma-zero-in-derived}
$R$ を環とする。$K_j \in D(R)$, $j = 1, 2, 3$ とし、
$i \not \in \{-1, 0\}$ に対して $H^i(K_j) = 0$ とする。
$\varphi : K_1 \to K_2$ および $\psi : K_2 \to K_3$ を $D(R)$ の写像とする。
$H^0(\varphi) = 0$ かつ $H^{-1}(\psi) = 0$ ならば、
% P02-E1050: frozen source lines 22713-22722 reverse the only type-correct composition psi o phi as phi o psi.
$\varphi \circ \psi = 0$ である。
\end{lemma}

\begin{proof}
Derived Categories, Lemma
\ref{derived-lemma-trick-vanishing-composition} を適用すると、
$\varphi \circ \psi$ は $\tau_{\leq -2}K_2 = 0$ を通じて分解することが分かる。
\end{proof}

\begin{lemma}
\label{more-algebra-lemma-silly}
$R$ を環とする。$K \in D(R)$ は $R^{\oplus n} \to R^{\oplus n}$ の形の
二項複体で与えられるとする。微分の行列を
$A \in \text{Mat}(n \times n, R)$ と表す。このとき
% P02-E1051: frozen source line 22730 uses undefined lowercase a after naming the matrix A.
$\det(a) : K \to K$ は $D(R)$ において零である。
\end{lemma}

\begin{proof}
省略する。良い練習問題である。
\end{proof}






\section{素朴余接複体}
\label{more-algebra-section-netherlander}

\noindent
この節では Algebra, Section \ref{algebra-section-netherlander} で始めた議論を
続ける。まず基底変換を議論する。最初の補題は、素朴余接複体と環の拡大との
素朴テンソル積を取ることが、考えるほどには素朴でないことを示す。

\begin{lemma}
\label{more-algebra-lemma-tensor-NL}
$R \to S$ および $S \to S'$ を環写像とする。標準写像
$\NL_{S/R} \otimes_S^\mathbf{L} S' \to \NL_{S/R} \otimes_S S'$ は、
$D(S')$ において同型
$\tau_{\geq -1}(\NL_{S/R} \otimes_S^\mathbf{L} S') \to \NL_{S/R} \otimes_S S'$
を誘導する。同様に、$R$ 上の $S$ の表示 $\alpha$ が与えられたとき、
標準写像
$\NL(\alpha) \otimes_S^\mathbf{L} S' \to \NL(\alpha) \otimes_S S'$ は
$D(S')$ において同型 $\tau_{\geq -1}(\NL(\alpha) \otimes_S^\mathbf{L} S') \to
\NL(\alpha) \otimes_S S'$ を誘導する。
\end{lemma}

\begin{proof}
Lemma \ref{more-algebra-lemma-two-term-base-change} の特別な場合である。
\end{proof}

\begin{lemma}
\label{more-algebra-lemma-base-change-NL}
$R \to S$ および $R \to R'$ を環写像とする。
$\alpha : P \to S$ を $R$ 上の $S$ の表示とする。このとき
$\alpha' : P \otimes_R R' \to S \otimes_R R'$ は、$R'$ 上の
$S' = S \otimes_R R'$ の表示である。標準写像
% P02-E1052: frozen source line 22778 uses NL(alpha) without the canonical backslash command.
$$
NL(\alpha) \otimes_S S' \to \NL(\alpha')
$$
は $H^0$ 上同型で、$H^{-1}$ 上全射である。特に標準写像
$$
\NL_{S/R} \otimes_S S' \to \NL_{S'/R'}
$$
は $H^0$ 上同型で、$H^{-1}$ 上全射である。
\end{lemma}

\begin{proof}
% P02-E1053: frozen source lines 22789-22790 use "Denote I = ..." and "Denote P' = ..." without "by".
$I = \Ker(P \to S)$ と表す。$P' = P \otimes_R R'$ および
$I' = \Ker(P' \to S')$ と表す。$P$ は $j \in J$ に対する
$x_j$ 上の多項式代数であるとする。補題に表示された写像は
$$
\xymatrix{
\bigoplus_{j \in J} S' \text{d}x_j \ar[r] &
\bigoplus_{j \in J} S' \text{d}x_j \\
I/I^2 \otimes_S S' \ar[r] \ar[u] &
I'/(I')^2 \ar[u]
}
$$
となる。ここで左の列は $\NL(\alpha) \otimes_S S'$、右の列は
$\NL(\alpha')$ である。テンソル積の右完全性により、
$I \otimes_R R' \to I'$ は全射であることが分かる。
従って下の矢印は全射である。これで補題の第一の主張が示された。
$\NL_{S/R} \otimes_S S' \to \NL_{S'/R'}$ に関する主張は、これらの複体が
$\NL(\alpha) \otimes_S S'$ および $\NL(\alpha')$ とホモトピックであることから従う。
\end{proof}

\begin{lemma}
\label{more-algebra-lemma-base-change-NL-flat}
環の余デカルト図式
$$
\xymatrix{
B \ar[r] & B' \\
A \ar[r] \ar[u] & A' \ar[u]
}
$$
を考える。$B$ が $A$ 上平坦ならば、標準写像
$\NL_{B/A} \otimes_B B' \to \NL_{B'/A'}$ は擬同型である。
さらに $\NL_{B/A}$ が $[-1, 0]$ に Tor 振幅を持つならば、
$\NL_{B/A} \otimes_B^\mathbf{L} B' \to \NL_{B'/A'}$ も擬同型である。
\end{lemma}

\begin{proof}
Algebra, Section \ref{algebra-section-netherlander} のような表示
$\alpha : P \to B$ を選ぶ。$I = \Ker(\alpha)$ とする。
$P' = P \otimes_A A'$ とおき、$A'$ 上の $B'$ の対応する表示を
$\alpha' : P' \to B'$ と表す。$B$ は $A$ 上平坦なので、
$I' = \Ker(\alpha')$ は $I \otimes_A A'$ に等しい。従って
$$
I'/(I')^2 = \Coker(I^2 \otimes_A A' \to I \otimes_A A') =
I/I^2 \otimes_A A' = I/I^2 \otimes_B B'
$$
である。両辺が同じ基底を持つので、
$\Omega_{P'/A'} = \Omega_{P/A} \otimes_A A'$ である。従って
$\Omega_{P'/A'} \otimes_{P'} B' = \Omega_{P/A} \otimes_P B \otimes_B B'$ である。
これは $\NL(\alpha) \otimes_B B' \to \NL(\alpha')$ が複体の同型であることを
示し、第一の主張が従う。

\medskip\noindent
$B$-加群の複体として
$$
\NL(\alpha) = I/I^2 \longrightarrow \Omega_{P/A} \otimes_P B
$$
があり、$I/I^2$ は次数 $-1$ に置かれる。次数 $0$ の項は自由なので、
Lemma \ref{more-algebra-lemma-last-one-flat} により、この複体が $[-1, 0]$ に
Tor 振幅を持つことと、$I/I^2$ が平坦 $B$-加群であることは同値である。
これが成り立つならば
$\NL(\alpha) \otimes_B^\mathbf{L} B' = \NL(\alpha) \otimes_B B'$ となり、
第二の主張を得る。
\end{proof}

\begin{lemma}
\label{more-algebra-lemma-lci-NL}
$A \to B$ を Definition \ref{more-algebra-definition-local-complete-intersection} の意味での
局所完全交叉とする。このとき $\NL_{B/A}$ は $D(B)$ の perfect 対象で、
$[-1, 0]$ に Tor 振幅を持つ。
\end{lemma}

\begin{proof}
$B = A[x_1, \ldots, x_n]/I$ と書く。このとき $\NL_{B/A}$ は
$B$-加群の複体
$$
I/I^2 \longrightarrow \bigoplus B \text{d}x_i
$$
で表され、$I/I^2$ は次数 $-1$ に置かれる。次数 $0$ の項は有限自由なので、
Lemma \ref{more-algebra-lemma-last-one-flat} により、この複体が $[-1, 0]$ に
Tor 振幅を持つことと、$I/I^2$ が平坦 $B$-加群であることは同値である。
定義により $I$ は Koszul 正則イデアル、従って準正則イデアルである；
Section \ref{more-algebra-section-ideals} を見よ。従って $I/I^2$ は有限射影 $B$-加群である
（Lemma \ref{more-algebra-lemma-quasi-regular-ideal-finite-projective}）。これにより
$\NL_{B/A}$ が perfect であり、$[-1, 0]$ に Tor 振幅を持つことの両方を得る。
\end{proof}

\begin{lemma}
\label{more-algebra-lemma-base-change-lci-bis}
環の余デカルト図式
$$
\xymatrix{
B \ar[r] & B' \\
A \ar[r] \ar[u] & A' \ar[u]
}
$$
を考える。$A \to B$ および $A' \to B'$ が Definition
\ref{more-algebra-definition-local-complete-intersection} の意味での局所完全交叉ならば、
$H^{-1}(\NL_{B/A} \otimes_B B') \to H^{-1}(\NL_{B'/A'})$ の核は
有限射影 $B'$-加群である。
\end{lemma}

\begin{proof}
Lemma \ref{more-algebra-lemma-lci-NL} により、複体 $\NL_{B/A}$ と $\NL_{B'/A'}$ は
$[-1, 0]$ に Tor 振幅を持つ perfect 複体である。Lemmas
\ref{more-algebra-lemma-tensor-NL}, \ref{more-algebra-lemma-pull-perfect}, and
\ref{more-algebra-lemma-pull-tor-amplitude} を組み合わせると、
$\NL_{B/A} \otimes_B B' = \NL_{B/A} \otimes_B^\mathbf{L} B'$ であり、
この複体も $[-1, 0]$ に Tor 振幅を持つ perfect 複体である。
$D(B')$ において完全三角
$$
C \to \NL_{B/A} \otimes_B B'  \to \NL_{B'/A'} \to C[1]
$$
を選ぶ。Lemmas \ref{more-algebra-lemma-two-out-of-three-perfect} and
\ref{more-algebra-lemma-cone-tor-amplitude} により、$C$ は $[-1, 1]$ に Tor 振幅を持つ
perfect 複体である。Lemma \ref{more-algebra-lemma-base-change-NL} により、複体 $C$ は
ただ一つの非零コホモロジー加群、すなわち補題の加群を次数 $-1$ に持つ。
この加群は有限表示であり（Lemma \ref{more-algebra-lemma-n-pseudo-module}）、
平坦である（Lemma \ref{more-algebra-lemma-tor-dimension}）。従って Algebra, Lemma
\ref{algebra-lemma-finite-projective} により有限射影である。
\end{proof}

\section{アーベル群の Rlim}
\label{more-algebra-section-Rlim}

\noindent
アーベル群に対する $R\lim$ を簡単に論じる。本節では、アーベル群の逆系の
アーベル圏を $\textit{Ab}(\mathbf{N})$ と表す。この記法はサイト上の
アーベル群の層に対する記法と整合する。実際、アーベル群の逆系は
圏 $\mathbf{N}$（$i \leq j$ のとき一意な射 $i \to j$ を持つ）上の群の層と
同じものである。Remark \ref{more-algebra-remark-Rlim-cohomology} を見よ。本節の議論の
多くは、サイト上のアーベル群の層に対するコホモロジー機構を構成する際の
議論を繰り返すものである。

\begin{lemma}
\label{more-algebra-lemma-compute-Rlim}
関手 $\lim : \textit{Ab}(\mathbf{N}) \to \textit{Ab}$ は右導来関手
\begin{equation}
\label{more-algebra-equation-Rlim}
R\lim : D(\textit{Ab}(\mathbf{N})) \longrightarrow D(\textit{Ab})
\end{equation}
を持つ。通常どおり $R^p\lim(K) = H^p(R\lim(K))$ と置く。さらに次が成り立つ。
\begin{enumerate}
\item $\textit{Ab}(\mathbf{N})$ の任意の $(A_n)$ に対して、$p > 1$ ならば
$R^p\lim A_n = 0$ である。
\item $D(\textit{Ab})$ の対象 $R\lim A_n$ は、次数 $0$ と $1$ に置かれた複体
$$
\prod A_n \to \prod A_n,\quad (x_n) \mapsto (x_n - f_{n + 1}(x_{n + 1}))
$$
によって表される。
\item $(A_n)$ が ML ならば $R^1\lim A_n = 0$ である。すなわち $(A_n)$ は
$\lim$ に関して右非輪状である。
\item 任意の $K^\bullet \in D(\textit{Ab}(\mathbf{N}))$ は、各項が $\lim$ に
関して右非輪状である複体と擬同型である。
% P02-E1054: frozen source lines 22964-22965 contain the stray phrase "i.e., of".
\item 各 $K^p = (K^p_n)$ が $\lim$ に関して右非輪状、すなわち
$R^1\lim_n K^p_n = 0$ ならば、$R\lim K$ は、次数 $p$ の項が
$\lim_n K_n^p$ である複体によって表される。
\end{enumerate}
\end{lemma}

\begin{proof}
$(A_n)$ を任意の逆系とする。射影によって与えられる遷移写像を持つ逆系
$(B_n)$ を
$$
B_n = A_n \oplus A_{n - 1} \oplus \ldots \oplus A_1
$$
で定める。遷移写像を $f_i : A_i \to A_{i - 1}$ とするとき、
$A_n \to B_n$ を
$(1, f_n, f_{n - 1} \circ f_n, \ldots, f_2 \circ \ldots \circ f_n)$
で与える。このようにして、任意の逆系が ML 系の部分対象であることが分かる
（Homology, Section \ref{homology-section-inverse-systems}）。Derived Categories,
Lemma \ref{derived-lemma-subcategory-right-acyclics} に Homology, Lemma
\ref{homology-lemma-Mittag-Leffler} を用いると、任意の ML 系は
$\lim$ に関して右非輪状であること、すなわち (3) が従う。
% P02-E1055: frozen source line 22985 names the otherwise undefined functor RF.
従って既に $RF$ は $D^+(\textit{Ab}(\mathbf{N}))$ 上で定義される。
Derived Categories, Proposition \ref{derived-proposition-enough-acyclics} を見よ。
$n > 1$ に対して $C_n = A_{n - 1} \oplus \ldots \oplus A_1$ と置き、
$C_1 = 0$ と置く。ここでも遷移写像は射影によって与える。このとき逆系の
短完全列 $0 \to (A_n) \to (B_n) \to (C_n) \to 0$ があり、
$B_n \to C_n$ は $(x_i) \mapsto (x_i - f_{i + 1}(x_{i + 1}))$ で与えられる。
$(C_n)$ も ML であるから、(2) が成り立つことが分かり
% P02-E1056: frozen source lines 22992-22993 retain an editorial placeholder.
（上で参照した命題による）、これは (1) も導く。最後に、Derived Categories,
Lemma \ref{derived-lemma-unbounded-right-derived} により、$R\lim$ は実際
$D(\textit{Ab}(\mathbf{N}))$ 全体上で定義される。実際、Derived Categories,
Lemma \ref{derived-lemma-unbounded-right-derived} の証明は主張 (4) と (5) を
証明することによって進む。
\end{proof}

\begin{lemma}
\label{more-algebra-lemma-six-term-Rlim}
アーベル群の逆系の短完全列
$$
0 \to (A_i) \to (B_i) \to (C_i) \to 0
$$
が与えられたとする。このとき付随する $6$ 項完全列
$0 \to \lim A_i \to \lim B_i \to \lim C_i \to
R^1\lim A_i \to R^1\lim B_i \to R^1\lim C_i \to 0$ がある。
\end{lemma}

\begin{proof}
Lemma \ref{more-algebra-lemma-compute-Rlim} の消滅から従う。
\end{proof}

\noindent
以下は Homology, Lemma \ref{homology-lemma-apply-Mittag-Leffler-again} の
``正しい'' 定式化である。

\begin{lemma}
\label{more-algebra-lemma-apply-Mittag-Leffler-again}
アーベル群の複体の逆系
$$
(A^{-2}_n \to A^{-1}_n \to A^0_n \to A^1_n)
$$
が与えられ、その極限を $A^{-2} \to A^{-1} \to A^0 \to A^1$ と表す。
% P02-E1057: frozen source lines 23028-23031 omit "by" in repeated Denote constructions.
コホモロジーの逆系を $(H_n^{-1})$, $(H_n^0)$ と表し、
$A^{-2} \to A^{-1} \to A^0 \to A^1$ のコホモロジーを $H^{-1}$, $H^0$ と表す。
もし
\begin{enumerate}
\item $(A^{-2}_n)$ と $(A^{-1}_n)$ の $R^1\lim$ が消滅し、
\item $(H^{-1}_n)$ の $R^1\lim$ が消滅する
\end{enumerate}
ならば $H^0 = \lim H_n^0$ である。
\end{lemma}

\begin{proof}
$K \in D(\textit{Ab}(\mathbf{N}))$ を、$n$ 番目の成分が複体
$A^{-2}_n \to A^{-1}_n \to A^0_n \to A^1_n$ である複体の系によって
表される対象とする。Derived Categories, Lemma
\ref{derived-lemma-two-ss-complex-functor} の二つのスペクトル系列を用いて
$H^0(R\lim K)$ を計算する\footnote{これらのスペクトル系列を使うには、
% P02-E1058: frozen source line 23047 says "A inverse system".
$\textit{Ab}(\mathbf{N})$ が十分多くの入射対象を持つことを示す必要がある。
アーベル群の逆系 $(I_n)$ が入射的であることと、各 $I_n$ が入射的アーベル群で
あり遷移写像が分裂全射であることは同値である。任意の系はそのような系に
埋め込まれる。詳細は省略する。}。第一の系列の $E_1$-page は
$$
\begin{matrix}
0 & 0 & R^1\lim A^0_n  & R^1\lim A^1_n \\
A^{-2} & A^{-1} & A^0  & A^1
\end{matrix}
$$
であり、水平微分を持つ。
% P02-E1059: frozen source line 23058 joins two predicates ungrammatically.
すべての高次微分は零である。第二の系列の $E_2$ page は
$$
\begin{matrix}
R^1\lim H^{-2}_n & 0 & R^1\lim H^0_n & R^1 \lim H^1_n \\
\lim H^{-2}_n &
\lim H^{-1}_n &
\lim H^0_n &
\lim H^1_n
\end{matrix}
$$
であり、この時点で退化する。従って結果が従う。
\end{proof}

\begin{lemma}
\label{more-algebra-lemma-pro-isomorphism-iso-Rlim}
$(A_n)$ と $(B_n)$ をアーベル群の逆系とする。プロ系の射
$\varphi : (A_n) \to (B_n)$ は写像 $\lim A_n \to \lim B_n$ と
$R^1\lim A_n \to R^1\lim B_n$ を定める。$\varphi$ がプロ同型ならば、
これらの写像は同型である。
\end{lemma}

\begin{proof}
プロ系の射については Categories, Example
\ref{categories-example-pro-morphism-inverse-systems} を見よ。写像 $\varphi$ は
$1 \leq m_1 < m_2 < m_3 < \ldots$ と、遷移写像と両立する写像
$\varphi_n : A_{m_n} \to B_n$ によって与えられる。
$\varphi', \varphi'' : \prod A_n \to \prod B_n$ を、それぞれ
$\varphi'((a_n)) = (\varphi_n(a_{m_n}))$ および
$$
\varphi''((a_n)) =
(\varphi_n(a_{m_n} + f(a_{m_n + 1}) + \ldots + f(a_{m_{n + 1} - 1})))
$$
に等しいものとして定める。
% P02-E1060: frozen source line 23093 misspells occurrence.
ここで現れる各 $f$ は、逆系 $(A_n)$ の適切な遷移写像を表す。このとき図式
$$
\xymatrix{
\prod A_n \ar[r]_\delta \ar[d]^{\varphi'} &
\prod A_n \ar[d]^{\varphi''} \\
\prod B_n \ar[r]^\delta & \prod B_n
}
$$
は可換である。ここで水平射は Lemma \ref{more-algebra-lemma-compute-Rlim} のものとする。
このようにして所望の写像を得る。この構成は次の意味で関手的である。
逆系 $(C_n)$、整数列 $1 \leq m'_1 < m'_2 < m'_3 < \ldots$、および
% P02-E1061: frozen source line 23106 misspells transition.
遷移写像と両立する写像 $\psi_n : B_{m'_n} \to C_n$ が与えられると、
$(\psi \circ \varphi)' = \psi' \circ \varphi'$ および
$(\psi \circ \varphi)'' = \psi'' \circ \varphi''$ が成り立つ。ここで合成
$\psi \circ \varphi$ とは、整数列
$1 \leq m_{m'_1} < m_{m'_2} < \ldots$ と写像
$\psi_n \circ \varphi_{m'_n} : A_{m_{m'_n}} \to C_n$ によるものをいう。
$B_n = A_n$ であり、$\varphi_n$ が遷移写像である場合、得られる写像が
恒等写像であることを示そう。これにより、この構成が $\varphi$ の代表
$(m_n, \varphi_n)$ の選び方に依存しないことと、補題の最後の主張の両方が
証明される。

\medskip\noindent
従って $B_n = A_n$ とし、$\varphi_n$ を遷移写像とする。
$(a_n) \in \lim A_n$ を元とする。このとき $\varphi'((a_n))$ の次数 $n$ の
成分は $a_{m_n}$ の像であり、これは $a_n$ に等しい。これで $\lim A_n$ に
関する主張が示された。$\xi \in R^1\lim A_n$ を $\prod A_n$ の元 $(a_n)$ の
類とする。元 $(b_n)$ を考える。各添字に対し、
$m_{n - 1} < i < m_n$ ならば
$$
b_i = a_i + f(a_{i + 1}) + \ldots + f(a_{m_n - 1})
$$
とし、ある $n$ に対して $i = m_n$ ならば $0$ とする。
このとき $(a'_n) = (a_n) - \delta((b_n))$ は $R^1\lim A_n$ において
同じ類を定め、計算により、ある $n$ に対して $i = m_n$ でない限り
$a'_i = 0$ であることが分かる。従って、ある $n$ に対して $i = m_n$ でない
限り $a_i = 0$ と仮定してよい。さらに
$$
(0, \ldots, -f(a_{m_j}), 0, \ldots, 0, a_{m_j}, 0, \ldots)
$$
は、位置 $j$ と $m_j$ に非零成分を持ち、
$$
c_j = (0, \ldots, 0, f(a_{m_j}), f(a_{m_j}), \ldots, a_{m_j}, 0, \ldots)
$$
の $\delta$ による像である。ここで後者の非零成分は位置
$j + 1, \ldots, m_j$ にある。和 $c = \sum c_j$ は $\prod A_n$ で意味を持つ。
$\varphi''((a_n)) = (a_{m_1}, a_{m_2}, \ldots)$ であることを思い出すと
$$
(a_n) - \varphi''((a_n)) = \delta(c)
$$
が分かり、証明が完了する。
\end{proof}

\begin{lemma}
\label{more-algebra-lemma-map-into-Rlim}
$\mathcal{D}$ を三角圏とする。$(K_n)$ を $\mathcal{D}$ の対象の逆系とし、
$K$ を系 $(K_n)$ の導来極限とする。このとき $\mathcal{D}$ の任意の $L$ に
対して短完全列
$$
0 \to R^1\lim \Hom_\mathcal{D}(L, K_n[-1]) \to
\Hom_\mathcal{D}(L, K) \to
\lim \Hom_\mathcal{D}(L, K_n) \to 0
$$
がある。
\end{lemma}

\begin{proof}
これは Derived Categories, Definition \ref{derived-definition-derived-limit}、
Lemma \ref{derived-lemma-representable-homological}、および上の
Lemma \ref{more-algebra-lemma-compute-Rlim} における $\lim$ と $R^1\lim$ の記述から従う。
\end{proof}

\begin{lemma}
\label{more-algebra-lemma-pro-isomorphism-Rlim}
$\mathcal{D}$ を三角圏とする。$(K_n)$ と $(M_n)$ を $\mathcal{D}$ の対象の
逆系とし、その導来極限をそれぞれ $K$ と $M$ とする。
$a : (K_n) \to (M_n)$ をプロ対象のプロ同型とする。このとき $a$ を用いて
（非標準的な）同型 $K \to M$ を構成できる。
\end{lemma}

\begin{proof}
Derived Categories, Remark
\ref{derived-remark-functoriality-derived-limit} により、定義用の完全三角の射
$$
\xymatrix{
K \ar[r] \ar[d] & \prod K_n \ar[r] \ar[d]^{a'} &
\prod K_n \ar[d]^{a''} \\
M \ar[r] & \prod M_n \ar[r] & \prod M_n
}
$$
に組み込まれる射 $K \to M$ を得る。従って $\mathcal{D}$ の任意の対象 $L$ に
対して、短完全列の写像
$$
\xymatrix{
0 \ar[r] &
R^1\lim \Hom_\mathcal{D}(L, K_n[-1]) \ar[r] \ar[d] &
\Hom_\mathcal{D}(L, K) \ar[r] \ar[d] &
\lim \Hom_\mathcal{D}(L, K_n) \ar[r] \ar[d] &
0 \\
0 \ar[r] &
R^1\lim \Hom_\mathcal{D}(L, M_n[-1]) \ar[r] &
\Hom_\mathcal{D}(L, M) \ar[r] &
\lim \Hom_\mathcal{D}(L, M_n) \ar[r] &
0
}
$$
を得る。Lemma \ref{more-algebra-lemma-map-into-Rlim} とその証明を見よ。
Lemma \ref{more-algebra-lemma-pro-isomorphism-iso-Rlim} により、左右の垂直射は同型である
\footnote{プロ同型 $a$ が写像 $a_n : K_{m_n} \to M_n$ によって与えられるなら、
% P02-E1062: frozen source line 23209 says "the we obtain".
対応する写像 $\Hom_\mathcal{D}(L, K_{m_n}) \to
\Hom_\mathcal{D}(L, M_n)$ を用いて $\Hom$ 上のプロ同型を得る。
従って Lemma \ref{more-algebra-lemma-pro-isomorphism-iso-Rlim} の証明における $\lim$ と
$R^1\lim$ 上の射の構成は、Derived Categories, Remark
\ref{derived-remark-functoriality-derived-limit} における $a'$ と $a''$ の構成と
一致する。}。従って中央の射も同型である。米田の補題により、写像
$K \to M$ は同型である。
\end{proof}

\begin{lemma}
\label{more-algebra-lemma-map-from-hocolim}
$\mathcal{D}$ を三角圏とする。$(K_n)$ を $\mathcal{D}$ の対象の系とし、
$K$ を系 $(K_n)$ の導来余極限とする。このとき $\mathcal{D}$ の任意の $L$ に
対して短完全列
$$
0 \to R^1\lim \Hom_\mathcal{D}(K_n, L[-1]) \to
\Hom_\mathcal{D}(K, L) \to
\lim \Hom_\mathcal{D}(K_n, L) \to 0
$$
がある。
\end{lemma}

\begin{proof}
これは Derived Categories, Definition \ref{derived-definition-derived-colimit}、
Lemma \ref{derived-lemma-representable-homological}、および上の
Lemma \ref{more-algebra-lemma-compute-Rlim} における $\lim$ と $R^1\lim$ の記述から従う。
\end{proof}

\begin{remark}[コホモロジーとしての Rlim]
\label{more-algebra-remark-Rlim-cohomology}
対象を自然数とし、$j \geq i$ のとき一意な射 $i \to j$ を持つ圏
$\mathbf{N}$ を考える。$\mathbf{N}$ にカオス位相
（Sites, Example \ref{sites-example-indiscrete}）を入れる。このとき層
$\mathcal{F}$ は、$\mathbf{N}$ 上の集合の逆系
$$
\mathcal{F}_1 \leftarrow \mathcal{F}_2 \leftarrow \mathcal{F}_3
\leftarrow \ldots
$$
と同じものである。$\Gamma(\mathbf{N}, \mathcal{F}) = \lim \mathcal{F}_n$
であることに注意する。アーベル群の逆系 $\mathcal{F}_n$ に対して
$$
R^p\lim \mathcal{F}_n = H^p(\mathbf{N}, \mathcal{F})
$$
である。実際、両辺は
$\mathcal{F} \mapsto \lim \mathcal{F}_n = H^0(\mathbf{N}, \mathcal{F})$ の
高次右導来関手である。従って $R\lim$ の存在は Cohomology on Sites,
Sections \ref{sites-cohomology-section-cohomology-sheaves} and
\ref{sites-cohomology-section-unbounded} の一般論からも従う。
\end{remark}

\noindent
次の補題の積は、複体の項ごとの積とも、導来圏 $D(\textit{Ab})$ における積とも
見ることができる。Derived Categories, Lemma \ref{derived-lemma-products} を見よ。

\begin{lemma}
\label{more-algebra-lemma-distinguished-triangle-Rlim}
$K = (K_n^\bullet)$ を $D(\textit{Ab}(\mathbf{N}))$ の対象とする。
$D(\textit{Ab})$ において標準的な完全三角
$$
R\lim K \to \prod\nolimits_n K_n^\bullet \to \prod\nolimits_n K_n^\bullet
\to R\lim K[1]
$$
が存在する。言い換えると、$R\lim K$ は $D(\textit{Ab})$ の逆系
$(K_n^\bullet)$ の導来極限である。Derived Categories, Definition
\ref{derived-definition-derived-limit} を見よ。
\end{lemma}

\begin{proof}
各 $p$ に対して逆系 $(K_n^p)$ が $\lim$ に関して右非輪状であると仮定する。
Lemma \ref{more-algebra-lemma-compute-Rlim} により、各 $p$ に対して短完全列
$$
0 \to \lim_n K^p_n \to \prod\nolimits_n K^p_n \to \prod\nolimits_n K^p_n \to 0
$$
を得る。$\lim_n K^p_n$ からなる複体は Lemma
\ref{more-algebra-lemma-compute-Rlim} により $R\lim K$ を計算するから、この場合に補題が
成り立つことが分かる。

\medskip\noindent
次に $K = (K_n^\bullet)$ を一般とする。Lemma
\ref{more-algebra-lemma-compute-Rlim} により、$D(\textit{Ab}(\mathbf{N}))$ における擬同型
$K \to L$ であって、各 $p$ に対して $(L_n^p)$ が非輪状となるものがある。
$\textit{Ab}$ では積が完全であるから、$\prod K_n^\bullet$ は
$\prod L_n^\bullet$ と擬同型である。従って、上で証明した $L$ に対する結果が
$K$ に対する結果を導く。
\end{proof}

\begin{lemma}
\label{more-algebra-lemma-break-long-exact-sequence}
Lemma \ref{more-algebra-lemma-distinguished-triangle-Rlim} の記法のもとで、完全三角に付随する
長完全コホモロジー列は短完全列
$$
0 \to R^1\lim_n H^{p - 1}(K_n^\bullet) \to
H^p(R\lim K) \to
\lim_n H^p(K_n^\bullet) \to 0
$$
に分解する。
\end{lemma}

\begin{proof}
完全三角の長完全列は
$$
\ldots \to H^p(R\lim K) \to \prod\nolimits_n H^p(K_n^\bullet)
\to \prod\nolimits_n H^p(K_n^\bullet) \to
H^{p + 1}(R\lim K) \to \ldots
$$
である。中央の写像の核は、補題で与えた明示的な記述により
$\lim_n H^p(K_n^\bullet)$ である。この写像の余核は
Lemma \ref{more-algebra-lemma-compute-Rlim} により $R^1\lim_n H^p(K_n^\bullet)$ である。
\end{proof}

\noindent
{\bf 警告。}$D(\textit{Ab}(\mathbf{N}))$ の対象はアーベル群の逆系の複体である。
これは複体の逆系 $(K_n^\bullet)$ と考えることもできる。しかし、これは
$D(\textit{Ab})$ の対象の逆系と同じものでは{\bf ない}。次の補題と注意が
その違いを説明する。

\begin{lemma}
\label{more-algebra-lemma-lift-to-system-complexes-Ab}
$(K_n)$ を $D(\textit{Ab})$ の対象の逆系とする。このとき
$D(\textit{Ab}(\mathbf{N}))$ の対象 $M = (M_n^\bullet)$ と、
$D(\textit{Ab})$ における同型 $M_n^\bullet \to K_n$ であって、図式
$$
\xymatrix{
M_{n + 1}^\bullet \ar[d] \ar[r] &
M_n^\bullet \ar[d] \\
K_{n + 1} \ar[r] & K_n
}
$$
を $D(\textit{Ab})$ において可換にするものが存在する。
\end{lemma}

\begin{proof}
まず $M_1^\bullet$ を $K_1$ を表すアーベル群の複体とする。
$M_e^\bullet \to M_{e - 1}^\bullet \to \ldots \to M_1^\bullet$ と写像
$\psi_i : M_i^\bullet \to K_i$ が構成され、補題の図式がすべての $n < e$ に
対して可換であると仮定する。
% P02-E1063: frozen source lines 23351-23364 switch from induction index e to unbound n.
このとき $D(\textit{Ab})$ における図式
$$
\xymatrix{
& M_n^\bullet \ar[d]^{\psi_n} \\
K_{n + 1} \ar[r] & K_n
}
$$
を考える。$D(\textit{Ab})$ における射の定義により、アーベル群の複体
$M_{n + 1}^\bullet$、$D(\textit{Ab})$ における同型
$M_{n + 1}^\bullet \to K_{n + 1}$、および合成
$$
K_{n + 1} \to K_n \xrightarrow{\psi_n^{-1}} M_n^\bullet
$$
を $D(\textit{Ab})$ において表す複体の射
$M_{n + 1}^\bullet \to M_n^\bullet$ を見いだせる。従って帰納法により補題が
成り立つ。
\end{proof}

\begin{remark}
\label{more-algebra-remark-compare-derived-limit}
$(K_n)$ を $D(\textit{Ab})$ の対象の逆系とする。$K = R\lim K_n$ をこの系の
導来極限とする（Derived Categories, Section
\ref{derived-section-derived-limit} を見よ）。$D(\textit{Ab})$ は可算積を持つ
（Derived Categories, Lemma \ref{derived-lemma-products}）ので、このような
導来極限は存在する。Lemma \ref{more-algebra-lemma-lift-to-system-complexes-Ab} により、
$(K_n)$ を $D(\textit{Ab}(\mathbf{N}))$ の対象 $M$ に持ち上げることもできる。
このとき $K \cong R\lim M$ である。ここで $R\lim$ は関手
(\ref{more-algebra-equation-Rlim}) である。実際、Lemma
\ref{more-algebra-lemma-distinguished-triangle-Rlim} により $R\lim M$ も系 $(K_n)$ の
導来極限である。従って、系 $(K_n)$ の持ち上げ $M$ には多数の同型類があり
得るが、その導来極限の同型型は選び方に依存しない。これは $R\lim M$ が系の
導来極限 $K = R\lim K_n$ と同型だからである。従って、本節で証明した
$R\lim$ に関する結果を導来極限に適用してよい。例えば、任意の
$p \in \mathbf{Z}$ に対して標準的な短完全列
$$
0 \to R^1\lim H^{p - 1}(K_n) \to H^p(K) \to \lim H^p(K_n) \to 0
$$
がある。これは Lemma \ref{more-algebra-lemma-break-long-exact-sequence} を $M$ に適用すれば
よいからである。
% P02-E1064: frozen source line 23398 says "can also been seen".
これは $M$ の存在を用いなくても、Lemma
\ref{more-algebra-lemma-break-long-exact-sequence} の証明の議論を（定義用の）完全三角
$K \to \prod K_n \to \prod K_n \to K[1]$ に適用することによって直接分かる。
\end{remark}

\begin{lemma}
\label{more-algebra-lemma-Rlim-pro-equal}
$E \to D$ を $D(\textit{Ab}(\mathbf{N}))$ の射とする。$E$ に付随する
$D(\textit{Ab})$ の対象の系を $(E_n)$、$D$ に付随するものを $(D_n)$ とする。
$(E_n) \to (D_n)$ がプロ対象の同型ならば、$R\lim E \to R\lim D$ は
$D(\textit{Ab})$ における同型である。
\end{lemma}

\begin{proof}
仮定は特に、プロ対象 $H^p(E_n)$ と $H^p(D_n)$ が同型であることを意味する。
結果は Lemma \ref{more-algebra-lemma-break-long-exact-sequence} の短完全列と
Lemma \ref{more-algebra-lemma-pro-isomorphism-iso-Rlim} から従う。
\end{proof}

\begin{lemma}[Emmanouil]
\label{more-algebra-lemma-emmanouil}
\begin{reference}
\cite{Emmanouil} から取った。
\end{reference}
$(A_n)$ をアーベル群の逆系とする。次は同値である。
\begin{enumerate}
\item $(A_n)$ は Mittag-Leffler である。
\item $R^1\lim A_n = 0$ であり、$\bigoplus_{i \in \mathbf{N}} (A_n)$ に
対しても同じことが成り立つ。
\end{enumerate}
\end{lemma}

\begin{proof}
$B = \bigoplus_{i \in \mathbf{N}} (A_n)$ と置く。従って $B = (B_n)$ であり、
$B_n = \bigoplus_{i \in \mathbf{N}} A_n$ である。$(A_n)$ が ML ならば $B$ も
ML であり、Lemma \ref{more-algebra-lemma-compute-Rlim} により $R^1\lim A_n = 0$ および
$R^1\lim B_n = 0$ である。

\medskip\noindent
逆に、$(A_n)$ が ML でないと仮定する。このとき、ある $m$、整数列
$m < m_1 < m_2 < \ldots$、および元 $x_i \in A_{m_i}$ であって、
そこでの像 $y_i \in A_m$ が $A_{m_i + 1} \to A_m$ の像に属さない
ものを選べる。元 $x_i$ と $y_i$ を用いて $R^1\lim B_n \not = 0$ を二通りに
示す。これで補題の証明が完了する。

\medskip\noindent
第一の証明。$C_n = \prod_{i \in \mathbf{N}} A_n$ として $C = (C_n)$ と置く。
余核を $Q_n$ とする標準的な単射 $B_n \to C_n$ がある。$Q = (Q_n)$ と置く。
$Q_n$ の元 $q_n$ は、$q_{n, i} \in A_n$ を満たす元の列
$q_n = (q_{n, 1}, q_{n, 2}, \ldots)$ を、尾部が零である列で割ったものと考えて
よい（言い換えると、有限個の位置でのみ異なる列を同一視する）。逆系の短完全列
$$
0 \to (B_n) \to (C_n) \to (Q_n) \to 0
$$
がある。この系の元 $q_n \in Q_n$ を
$$
q_{n, i} =
\left\{
\begin{matrix}
\text{image of }x_i &\text{if}& m_i \geq n \\
0 & \text{else}
\end{matrix}
\right.
$$
で定める。このとき $q_{n + 1}$ が $q_n$ に写ることは明らかである。従って
$q = (q_n) \in \lim Q_n$ を得る。一方、$q$ は $\lim C_n \to \lim Q_n$ の像に
属さないと主張する。実際、$c = (c_n)$ が $q$ に写ると仮定する。このとき
$c_n = (c_{n, i})$ と書ける。$n' \geq n$ に対して
$c_{n', i} \mapsto c_{n, i}$ であるから、すべての $n, i, n' \geq n$ に対して
$c_{n, i} \in \Im(C_{n'} \to C_n)$ である。特に $A_m$ における
$c_{m, i}$ の像は $\Im(A_{m_i + 1} \to A_m)$ に属するので、$y_i$ と等しく
なり得ない。従って $c_m$ と $q_m = (y_1, y_2, y_3, \ldots)$ は無限個の位置で
異なり、矛盾である。長完全コホモロジー列
$$
0 \to \lim B_n \to \lim C_n \to \lim Q_n \to R^1\lim B_n
$$
を考えると、所望のとおり最後の群が非零であると結論できる。

\medskip\noindent
第二の証明。$n' \geq n$ に対して
$A_{n, n'} = \Im(A_{n'} \to A_n)$ と書く。このとき $y_i \in A_m$ かつ
$y_i \not \in A_{m, m_i + 1}$ である。$\xi = (\xi_n) \in \prod B_n$ を、
$n = m_i$ でない限り $\xi_n = 0$ であり、
$\xi_{m_i} = (0, \ldots, 0, x_i, 0, \ldots)$ であって $x_i$ が第 $i$ 直和因子に
置かれる元とする。$\xi$ は Lemma \ref{more-algebra-lemma-compute-Rlim} の写像
$\prod B_n \to \prod B_n$ の像に属さないと主張する。これにより
$R^1\lim B_n$ が非零であることが示され、証明が完了する。実際、
$\xi$ が $\eta = (z_1, z_2, \ldots)$ の像であると仮定する。ここで
$z_n = \sum z_{n, i} \in \bigoplus_i A_n$ とする。
$x_i = z_{m_i, i} \bmod A_{m_i, m_i + 1}$ であることに注意する。このとき
$z_{m_i - 1, i}$ は $A_{m_i} \to A_{m_i - 1}$ による $z_{m_i, i}$ の像であり、
以下同様である。従って $z_{m, i}$ は $A_{m_i} \to A_m$ による
$z_{m_i, i}$ の像である。ゆえに $z_{m, i}$ は $A_{m, m_i + 1}$ を法として
$y_i$ と合同である。特に $z_{m, i} \not = 0$ である。しかし
$\sum z_{m, i} \in \bigoplus_i A_m$ であるため、非零となり得る $z_{m, i}$ は
有限個しかなく、これは不可能である。
\end{proof}

\begin{lemma}
\label{more-algebra-lemma-Mittag-Leffler}
アーベル群の逆系の短完全列
$$
0 \to (A_i) \to (B_i) \to (C_i) \to 0
$$
が与えられたとする。$(A_i)$ と $(C_i)$ が ML ならば $(B_i)$ も ML である。
\end{lemma}

\begin{proof}
これは Lemma \ref{more-algebra-lemma-emmanouil}、無限直和を取ることの完全性、および
$R\lim$ に付随する長完全コホモロジー列から従う。
\end{proof}

\begin{lemma}
\label{more-algebra-lemma-Rlim-zero-of-direct-sums}
$(A_n)$ をアーベル群の逆系とする。次は同値である。
\begin{enumerate}
\item $(A_n)$ はプロ対象として零である。
\item $\lim A_n = 0$ かつ $R^1\lim A_n = 0$ であり、
$\bigoplus_{i \in \mathbf{N}} (A_n)$ に対しても同じことが成り立つ。
\end{enumerate}
\end{lemma}

\begin{proof}
Lemma \ref{more-algebra-lemma-pro-isomorphism-iso-Rlim} により (1) は (2) を導く。
(2) を仮定する。このとき Lemma \ref{more-algebra-lemma-emmanouil} により $(A_n)$ は ML で
ある。$m \geq n$ に対して $A_{n, m} = \Im(A_m \to A_n)$ と置くと、
$A_n = A_{n, n} \supset A_{n, n + 1} \supset \ldots$ である。$(A_n)$ が
プロ対象として零であることと、任意の $n$ に対して、ある $m \geq n$ が存在して
$A_{n, m} = 0$ となることは同値である。$(A_n)$ が ML であることと、任意の
$n$ に対して、ある $m_n \geq n$ が存在して
$A_{n, m} = A_{n, m + 1} = \ldots$ となることは同値である。ML の場合、
$\lim A_n = 0$ ならば $A_{n, m_n} = 0$ であることは明らかである。
% P02-E1065: frozen source line 23548 has an unbound target index m in this map.
実際、写像 $A_{n + 1, m_{n + 1}} \to A_{n, m}$ は全射である。これで証明が
完了する。
\end{proof}

\section{加群の Rlim}
\label{more-algebra-section-Rlim-modules}

\noindent
加群に対する $R\lim$ を簡単に論じる。本節の議論の多くは、環付きサイト上の
加群に対するコホモロジー機構を構成する際の議論を繰り返すものである。

\medskip\noindent
$(A_n)$ を環の逆系とする。アーベル群の逆系 $(M_n)$ であって、各 $M_n$ に
% P02-E1066: frozen source line 23567 says "a A_n-module".
$A_n$-加群の構造が与えられ、遷移写像 $M_{n + 1} \to M_n$ が
$A_{n + 1}$-加群写像であるものの圏を
$\textit{Mod}(\mathbf{N}, (A_n))$ と表す。これはアーベル圏である。
$A = \lim A_n$ と置く。$\textit{Mod}(\mathbf{N}, (A_n))$ の対象 $(M_n)$ が
与えられると、その極限 $\lim M_n$ は $A$-加群である。

\begin{lemma}
\label{more-algebra-lemma-compute-Rlim-modules}
% P02-E1067: frozen source line 23575 is the sentence fragment "In the situation above.".
上の状況において、関手
$\lim : \textit{Mod}(\mathbf{N}, (A_n)) \to \text{Mod}_A$ は右導来関手
$$
R\lim :
D(\textit{Mod}(\mathbf{N}, (A_n)))
\longrightarrow
D(A)
$$
を持つ。通常どおり $R^p\lim(K) = H^p(R\lim(K))$ と置く。さらに次が成り立つ。
\begin{enumerate}
\item $\textit{Mod}(\mathbf{N}, (A_n))$ の任意の $(M_n)$ に対して、$p > 1$
ならば $R^p\lim M_n = 0$ である。
\item $D(\text{Mod}_A)$ の対象 $R\lim M_n$ は、次数 $0$ と $1$ に置かれた複体
$$
\prod M_n \to \prod M_n,\quad
(x_n) \mapsto (x_n - f_{n + 1}(x_{n + 1}))
$$
によって表される。
\item $(M_n)$ が ML ならば $R^1\lim M_n = 0$ である。すなわち $(M_n)$ は
$\lim$ に関して右非輪状である。
\item 任意の $K^\bullet \in D(\textit{Mod}(\mathbf{N}, (A_n)))$ は、各項が
$\lim$ に関して右非輪状である複体と擬同型である。
% P02-E1068: frozen source lines 23599-23600 repeat the stray phrase "i.e., of".
\item 各 $K^p = (K^p_n)$ が $\lim$ に関して右非輪状、すなわち
$R^1\lim_n K^p_n = 0$ ならば、$R\lim K$ は、次数 $p$ の項が
$\lim_n K_n^p$ である複体によって表される。
\end{enumerate}
\end{lemma}

\begin{proof}
この証明は Lemma \ref{more-algebra-lemma-compute-Rlim} の証明と逐語的に同じである。
\end{proof}

\begin{remark}
\label{more-algebra-remark-Rlim-cohomology-modules}
これは Remark \ref{more-algebra-remark-Rlim-cohomology} の続きである。$\mathbf{N}$ 上の
環の層は環の逆系 $(A_n)$ にほかならない。$(A_n)$ 上の加群の層は、上で定義した
圏 $\textit{Mod}(\mathbf{N}, (A_n))$ の対象と全く同じものである。
Lemma \ref{more-algebra-lemma-compute-Rlim-modules} の導来関手 $R\lim$ は、加群の導来圏から
構造層の大域切断環上の加群の導来圏への $R\Gamma(\mathbf{N}, -)$ にほかならない。
一般に群としてのコホモロジーと加群としてのコホモロジーは一致する。
Cohomology on Sites, Lemma
\ref{sites-cohomology-lemma-cohomology-modules-abelian-agree} を見よ。
\end{remark}

\noindent
次の補題の積は、複体の項ごとの積とも、導来圏 $D(A)$ における積とも見ることが
できる。Derived Categories, Lemma \ref{derived-lemma-products} を見よ。

\begin{lemma}
\label{more-algebra-lemma-distinguished-triangle-Rlim-modules}
$K = (K_n^\bullet)$ を $D(\textit{Mod}(\mathbf{N}, (A_n)))$ の対象とする。
$D(A)$ において標準的な完全三角
$$
R\lim K \to \prod\nolimits_n K_n^\bullet \to \prod\nolimits_n K_n^\bullet
\to R\lim K[1]
$$
が存在する。言い換えると、$R\lim K$ は $D(A)$ の逆系 $(K_n^\bullet)$ の
導来極限である。Derived Categories, Definition
\ref{derived-definition-derived-limit} を見よ。
\end{lemma}

\begin{proof}
この証明は Lemma \ref{more-algebra-lemma-distinguished-triangle-Rlim} の証明において
% P02-E1069: frozen source line 23648 spells "instead" as two words.
Lemma \ref{more-algebra-lemma-compute-Rlim-modules} を
Lemma \ref{more-algebra-lemma-compute-Rlim} の代わりに用いたものと全く同じである。
\end{proof}

\begin{lemma}
\label{more-algebra-lemma-break-long-exact-sequence-modules}
Lemma \ref{more-algebra-lemma-distinguished-triangle-Rlim-modules} の記法のもとで、完全三角に
付随する長完全コホモロジー列は、$A$-加群の短完全列
$$
0 \to R^1\lim_n H^{p - 1}(K_n^\bullet) \to
H^p(R\lim K) \to
\lim_n H^p(K_n^\bullet) \to 0
$$
に分解する。
\end{lemma}

\begin{proof}
この証明は Lemma \ref{more-algebra-lemma-break-long-exact-sequence} の証明において
% P02-E1070: frozen source line 23668 again spells "instead" as two words.
Lemma \ref{more-algebra-lemma-compute-Rlim-modules} を
Lemma \ref{more-algebra-lemma-compute-Rlim} の代わりに用いたものと全く同じである。
\end{proof}

\noindent
{\bf 警告。} アーベル群の場合と同様に、
$D(\textit{Mod}(\mathbf{N}, (A_n)))$ の対象 $M = (M_n^\bullet)$ は加群の複体の
逆系であり、導来圏の対象の逆系と同じものでは{\bf ない}。次の補題では、
導来圏の対象の逆系が常に $D(\textit{Mod}(\mathbf{N}, (A_n)))$ の対象へ
持ち上がることを示す。

\begin{lemma}
\label{more-algebra-lemma-lift-to-system-complexes}
$(A_n)$ を環の逆系とする。次が与えられていると仮定する。
\begin{enumerate}
\item 各 $n$ に対して $D(A_n)$ の対象 $K_n$。
\item 各 $n$ に対して、$A_{n + 1} \to A_n$ によるスカラー制限によって
$K_n$ を $D(A_{n + 1})$ の対象とみなしたときの、$D(A_{n + 1})$ の写像
$\varphi_n : K_{n + 1} \to K_n$。
\end{enumerate}
このとき対象 $M = (M_n^\bullet) \in D(\textit{Mod}(\mathbf{N}, (A_n)))$ と、
$D(A_n)$ における同型 $\psi_n : M_n^\bullet \to K_n$ であって、図式
$$
\xymatrix{
M_{n + 1}^\bullet \ar[d]_{\psi_{n + 1}} \ar[r] & M_n^\bullet \ar[d]^{\psi_n} \\
K_{n + 1} \ar[r]^{\varphi_n} & K_n
}
$$
を $D(A_{n + 1})$ において可換にするものが存在する。
\end{lemma}

\begin{proof}
証明を詳しく書く。$A_n$-加群 $T$ に対し、同じ加群を $A_{n + 1}$-加群として
% P02-E1071: frozen source line 23705 misspells viewed.
見たものを $T_{A_{n + 1}}$ と書く。$K_n^\bullet$ を $K_n$ を表す
$A_n$-加群の複体とする。このとき $K_{n, A_{n + 1}}^\bullet$ は同じ複体を
$A_{n + 1}$-加群の複体として見たものである。導来圏の構成により、
% P02-E1072: frozen source line 23709 names psi_n although sigma_n and tau_n represent phi_n.
写像 $\psi_n$ は
$$
\psi_n = \tau_n \circ \sigma_n^{-1}
$$
で与えられる。ここで $\sigma_n : L_{n + 1}^\bullet \to K_{n + 1}^\bullet$ は
$A_{n + 1}$-加群の複体の擬同型であり、
$\tau_n : L_{n + 1}^\bullet \to K_{n, A_{n + 1}}^\bullet$ は
$A_{n + 1}$-加群の複体の写像である。

\medskip\noindent
次に複体 $M_n^\bullet$ を帰納的に構成する。基底の場合として
$M_1^\bullet = K_1^\bullet$ と置く。
$M_e^\bullet \to M_{e - 1}^\bullet \to \ldots \to M_1^\bullet$ と複体の写像
$\psi_i : M_i^\bullet \to K_i^\bullet$ が既に構成され、図式
$$
\xymatrix{
M_{n + 1}^\bullet \ar[d]_{\psi_{n + 1}} \ar[rr] & &
M_{n, A_{n + 1}}^\bullet \ar[d]^{\psi_{n, A_{n + 1}}} \\
K_{n + 1}^\bullet & L_{n + 1}^\bullet \ar[l]_{\sigma_n} \ar[r]^{\tau_n} &
K_{n, A_{n + 1}}^\bullet
}
$$
がすべての $n < e$ に対して上の $D(A_{n + 1})$ において可換であると仮定する。
このとき $D(A_{e + 1})$ における図式
$$
\xymatrix{
& & M_{e, A_{e + 1}}^\bullet \ar[d]^{\psi_{e, A_{e + 1}}} \\
K_{e + 1}^\bullet & L_{e + 1}^\bullet \ar[r]^{\tau_e} \ar[l]_{\sigma_e} &
K_{e, A_{e + 1}}^\bullet
}
$$
を考える。$\psi_e$ は擬同型なので、$\psi_{e, A_{e + 1}}$ も擬同型である。
$D(A_{e + 1})$ における射の定義により、$A_{e + 1}$-加群の複体の擬同型
$\psi_{e + 1} : M_{e + 1}^\bullet \to K_{e + 1}^\bullet$ であって、
$D(A_{e + 1})$ における合成
$\psi_{e, A_{e + 1}}^{-1} \circ \tau_e \circ \sigma_e^{-1}$ を表す
$A_{e + 1}$-加群の複体の射
$M_{e + 1}^\bullet \to M_{e, A_{e + 1}}^\bullet$ が存在するものを見いだせる。
従って帰納法により補題が成り立つ。
\end{proof}

\begin{remark}
\label{more-algebra-remark-how-unique}
Lemma \ref{more-algebra-lemma-lift-to-system-complexes} と同じ仮定のもとで考える。補題の系
$(K_n, \varphi_n)$ を生じる $D(\textit{Mod}(\mathbf{N}, (A_n)))$ の対象 $M$ には、
先験的には多数の同型類がある。そのような各 $M$ に対して、
$A = \lim A_n$ と置けば、複体 $R\lim M \in D(A)$ を考えられる。
Lemma \ref{more-algebra-lemma-distinguished-triangle-Rlim-modules} により、$R\lim M$ は
$D(A)$ の逆系 $(K_n)$ の導来極限である。従って $D(A)$ における
$R\lim M$ の同型類は、$M$ の構成で行った選択に依存しない。特に、本節で
証明した $R\lim$ に関する結果を $D(A)$ における逆系の導来極限へ適用してよい。
例えば、任意の $p \in \mathbf{Z}$ に対して標準的な短完全列
$$
0 \to R^1\lim H^{p - 1}(K_n) \to H^p(R\lim K_n) \to \lim H^p(K_n) \to 0
$$
がある。これは Lemma \ref{more-algebra-lemma-distinguished-triangle-Rlim-modules} を $M$ に
適用すればよいからである。
% P02-E1073: frozen source line 23771 repeats "can also been seen".
これは $M$ の存在を用いなくても、Lemma
\ref{more-algebra-lemma-distinguished-triangle-Rlim-modules} の証明の議論を、導来極限の
（定義用の）完全三角
$R\lim K_n \to \prod K_n \to \prod K_n \to (R\lim K_n)[1]$ に適用することに
よって直接分かる。
\end{remark}

\begin{lemma}
\label{more-algebra-lemma-get-ML-system}
$(A_n)$ を環の逆系とする。任意の
$K \in D(\textit{Mod}(\mathbf{N}, (A_n)))$ は、すべての遷移写像
$M_{n + 1}^\bullet \to M_n^\bullet$ が全射である複体の系 $(M_n^\bullet)$ に
よって表すことができる。
\end{lemma}

\begin{proof}
$K$ が系 $(K_n^\bullet)$ によって表されるとする。
$M_1^\bullet = K_1^\bullet$ と置く。複体の全射
$M_n^\bullet \to M_{n - 1}^\bullet \to \ldots \to M_1^\bullet$ と
ホモトピー同値 $\psi_e : K_e^\bullet \to M_e^\bullet$ が構成され、図式
$$
\xymatrix{
K_{e + 1}^\bullet \ar[d] \ar[r] & K_e^\bullet \ar[d] \\
M_{e + 1}^\bullet \ar[r] & M_e^\bullet
}
$$
がすべての $e < n$ に対して可換であると仮定する。このとき図式
$$
\xymatrix{
K_{n + 1}^\bullet \ar[r] & K_n^\bullet \ar[d] \\
& M_n^\bullet
}
$$
を考える。Derived Categories, Lemma \ref{derived-lemma-make-surjective} により、
合成 $K_{n + 1}^\bullet \to M_n^\bullet$ を
$K_{n + 1}^\bullet \to M_{n + 1}^\bullet \to M_n^\bullet$ と分解できる。
ここで第一の射はホモトピー同値であり、第二の射は各項で分裂する全射である。
これと帰納法から補題が従う。
\end{proof}

\begin{lemma}
\label{more-algebra-lemma-get-K-flat-system}
$(A_n)$ を環の逆系とする。任意の
$K \in D(\textit{Mod}(\mathbf{N}, (A_n)))$ は、各 $K_n^\bullet$ が K-flat である
複体の系 $(K_n^\bullet)$ によって表すことができる。
\end{lemma}

\begin{proof}
まず Lemma \ref{more-algebra-lemma-get-ML-system} を用いて、$K$ を、すべての遷移写像
$M_{n + 1}^\bullet \to M_n^\bullet$ が全射である複体の系 $(M_n^\bullet)$ に
よって表す。次に $K_1^\bullet$ を $A_1$-加群の K-flat 複体として、擬同型
$K_1^\bullet \to M_1^\bullet$ を取る
（Lemma \ref{more-algebra-lemma-K-flat-resolution}）。
$K_n^\bullet \to K_{n - 1}^\bullet \to \ldots \to K_1^\bullet$ と複体の写像
$\psi_e : K_e^\bullet \to M_e^\bullet$ が構成され、図式
$$
\xymatrix{
K_{e + 1}^\bullet \ar[d] \ar[r] & K_e^\bullet \ar[d] \\
M_{e + 1}^\bullet \ar[r] & M_e^\bullet
}
$$
がすべての $e < n$ に対して可換であると仮定する。このとき
$D(A_{n + 1})$ における図式
$$
\xymatrix{
C^\bullet \ar@{..>}[d] \ar@{..>}[r] & K_n^\bullet \ar[d]^{\psi_n} \\
M_{n + 1}^\bullet \ar[r]^{\varphi_n} & M_n^\bullet
}
$$
を考える。$M_{n + 1}^\bullet \to M_n^\bullet$ は各項で全射であるから、項が
$$
C^p = M_{n + 1}^p \times_{M_n^p} K_n^p
$$
である左上隅の複体 $C^\bullet$ は $M_{n + 1}^\bullet$ と擬同型である
（詳細は省略する）。$K_{n +1}^\bullet$ が K-flat となる擬同型
$K_{n + 1}^\bullet \to C^\bullet$ を選ぶ。従って帰納法により補題が成り立つ。
\end{proof}

\begin{lemma}
\label{more-algebra-lemma-derived-tensor-product-systems}
$(A_n)$ を環の逆系とする。
$K, L \in D(\textit{Mod}(\mathbf{N}, (A_n)))$ が与えられると、
$D(\mathbf{N}, (A_n))$ に標準的な導来テンソル積
$K \otimes^\mathbf{L} L$ があり、$D(A_n)$ への写像と両立する。この構成は
$K$ と $L$ に関して対称であり、各変数について三角圏の完全関手である。
\end{lemma}

\begin{proof}
$K$ の代表 $(K_n^\bullet)$ であって、各 $K_n^\bullet$ が K-flat 複体である
ものを選ぶ（Lemma \ref{more-algebra-lemma-get-K-flat-system}）。このとき $L$ の任意の代表
$(L_n^\bullet)$ に対して、複体の系
$$
(\text{Tot}(K_n^\bullet \otimes_{A_n} L_n^\bullet))
$$
によって表される対象として $K \otimes^\mathbf{L} L$ を定義できる。
Lemmas \ref{more-algebra-lemma-K-flat-quasi-isomorphism} and
\ref{more-algebra-lemma-derived-tor-quasi-isomorphism-other-side} により、これは両変数に
関して良定義である。$D(A_n)$ への写像との両立性は明らかである。完全性は
Lemma \ref{more-algebra-lemma-derived-tor-exact} と全く同様に従う。
\end{proof}

\begin{remark}
\label{more-algebra-remark-constructing-tensor-with-limits-functorially}
$A$ を環とし、$(E_n)$ を $D(A)$ の対象の逆系とする。上で導来極限
$R\lim E_n$ が存在することを見た。従って $D(A)$ の任意の対象 $K$ に対しても
導来極限 $R\lim( K \otimes_A^\mathbf{L} E_n )$ が存在する。実際、これらの
導来極限を $K$ に関して関手的に構成し、三角圏の完全関手
$$
R\lim(- \otimes_A^\mathbf{L} E_n) : D(A) \longrightarrow D(A)
$$
を得ることができる。まず Lemma \ref{more-algebra-lemma-lift-to-system-complexes} により、
$(E_n)$ を $D(\mathbf{N}, A)$ の対象 $E$ へ持ち上げる
（この関手は持ち上げの選び方に依存する）。次に、複体 $K^\bullet$ を、値
$K^\bullet$ を持つ複体の定値逆系に写す ``対角'' または ``定値'' 関手
$$
\Delta : D(A) \longrightarrow D(\mathbf{N}, A)
$$
があることに注意する。そこで単に
$$
R\lim(K \otimes_A^\mathbf{L} E_n) = R\lim(\Delta(K)\otimes^\mathbf{L} E)
$$
と定義する。右辺では Lemma \ref{more-algebra-lemma-compute-Rlim-modules} の関手 $R\lim$ と
Lemma \ref{more-algebra-lemma-derived-tensor-product-systems} の関手
$- \otimes^\mathbf{L} -$ を用いる。
\end{remark}

\begin{lemma}
\label{more-algebra-lemma-tensor-Rlim-exact}
$A$ を環とする。$E \to D \to F \to E[1]$ を $D(\mathbf{N}, A)$ の完全三角と
する。$E$, $D$, $F$ に付随する $D(A)$ の対象の系をそれぞれ
$(E_n)$, $(D_n)$, $(F_n)$ とする。このとき任意の $K \in D(A)$ に対して、
Remark \ref{more-algebra-remark-constructing-tensor-with-limits-functorially} の記法のもとで、
$D(A)$ に標準的な完全三角
$$
R\lim (K \otimes^\mathbf{L}_A E_n) \to
R\lim (K \otimes^\mathbf{L}_A D_n) \to
R\lim (K \otimes^\mathbf{L}_A F_n) \to
R\lim (K \otimes^\mathbf{L}_A E_n)[1]
$$
がある。
\end{lemma}

\begin{proof}
これは Remark \ref{more-algebra-remark-constructing-tensor-with-limits-functorially} の構成と、
$\Delta : D(A) \to D(\mathbf{N}, A)$、$- \otimes^\mathbf{L} -$、および
$R\lim$ が三角圏の完全関手であることから明らかである。
\end{proof}

\begin{lemma}
\label{more-algebra-lemma-tensor-Rlim-pro-equal}
$A$ を環とする。$E \to D$ を $D(\mathbf{N}, A)$ の射とする。$E$, $D$ に
付随する $D(A)$ の対象の系をそれぞれ $(E_n)$, $(D_n)$ とする。
$(E_n) \to (D_n)$ がプロ対象の同型ならば、任意の $K \in D(A)$ に対して
対応する写像
$$
R\lim (K \otimes^\mathbf{L}_A E_n)
\longrightarrow
R\lim (K \otimes^\mathbf{L}_A D_n)
$$
は $D(A)$ における同型である
（記法は Remark \ref{more-algebra-remark-constructing-tensor-with-limits-functorially} のもの）。
\end{lemma}

\begin{proof}
定義と Lemma \ref{more-algebra-lemma-Rlim-pro-equal} から従う。
\end{proof}

\section{ねじれ加群}
\label{more-algebra-section-torsion}

\noindent
本節で ``ねじれ加群'' とは、アフィンスキーム $\Spec(R)$ の与えられた閉部分集合
$V(I)$ 上に台を持つ加群をいう。これは整域上のねじれ加群の概念
（Definition \ref{more-algebra-definition-torsion}）とは異なるが、それと類似している。

\begin{definition}
\label{more-algebra-definition-f-power-torsion}
$R$ を環とし、$M$ を $R$-加群とする。
\begin{enumerate}
\item $I \subset R$ をイデアルとする。任意の $m \in M$ に対して、ある
$n > 0$ が存在して $I^n m = 0$ となるとき、$M$ を
{\it $I$-冪ねじれ加群}という。
\item $f \in R$ とする。各 $m \in M$ に対して、ある $n > 0$ が存在して
$f^n m = 0$ となるとき、$M$ を {\it $f$-冪ねじれ加群}という。
\end{enumerate}
\end{definition}

\noindent
従って $f$-冪ねじれ加群は、$I = (f)$ に対する $I$-冪ねじれ加群と同じもので
ある。$R$-加群 $M$ に対して記法
$$
M[I^n] = \{m \in M \mid I^nm = 0\}
$$
および
$$
M[I^\infty] = \bigcup M[I^n]
$$
を用いる。従って $M$ が $I$-冪ねじれであることと、$M = M[I^\infty]$ で
あることと、$M = \bigcup M[I^n]$ であることは同値である。

\begin{lemma}
\label{more-algebra-lemma-I-power-torsion-presentation}
$R$ を環、$I$ を $R$ のイデアル、$M$ を $I$-冪ねじれ加群とする。このとき
$M$ は分解
$$
\ldots \to K_2 \to K_1 \to K_0 \to M \to 0
$$
を持ち、各 $K_i$ は $R/I^n$ のコピーの直和である。
% P02-E1074: frozen source line 24010 says "for n variable" rather than "for varying n".
ここで $n$ は変化し得る。
\end{lemma}

\begin{proof}
標準的な全射
$$
\oplus_{m \in M} R/I^{n_m} \to M \to 0
$$
がある。ここで $n_m$ は $I^{n_m} \cdot m = 0$ となる最小の正整数である。
この全射の核も $I$-冪ねじれ加群である。帰納的に進めると、$M$ の所望の分解が
構成される。
\end{proof}

\begin{lemma}
\label{more-algebra-lemma-torsion-free}
$R$ を環、$I$ を $R$ のイデアルとする。任意の $R$-加群 $M$ に対して
$M[I^n] = \{m \in M \mid I^nm = 0\}$ と置く。$I$ が有限生成ならば次は同値で
ある。
\begin{enumerate}
\item $M[I] = 0$。
\item すべての $n \geq 1$ に対して $M[I^n] = 0$。
\item $I = (f_1, \ldots, f_t)$ ならば写像 $M \to \bigoplus M_{f_i}$ は単射で
ある。
\end{enumerate}
\end{lemma}

\begin{proof}
Algebra, Lemma \ref{algebra-lemma-when-injective-covering} から従う。
\end{proof}

\begin{lemma}
\label{more-algebra-lemma-divide-by-torsion}
$R$ を環、$I$ を $R$ の有限生成イデアルとする。
\begin{enumerate}
\item 任意の $R$-加群 $M$ に対して $(M/M[I^\infty])[I] = 0$ である。
\item $I$-冪ねじれ加群の拡大は $I$-冪ねじれである。
\end{enumerate}
\end{lemma}

\begin{proof}
$m \in M$ とする。$m$ が $(M/M[I^\infty])[I]$ の元に写るならば
$Im \subset M[I^\infty]$ である。$I = (f_1, \ldots, f_t)$ と書く。このとき
$f_i m \in M[I^\infty]$、すなわち、ある $n_i > 0$ に対して
$I^{n_i}f_i m = 0$ である。従って $N = \sum n_i + 2$ とすれば $I^Nm = 0$ で
ある。ゆえに $m$ は $(M/M[I^\infty])$ において零に写る。これで補題の第一の
主張が示された。

\medskip\noindent
第二の主張について、$0 \to M' \to M \to M'' \to 0$ を加群の短完全列とし、
$M'$ と $M''$ がともに $I$-冪ねじれ加群であると仮定する。このとき
$M[I^\infty] \supset M'$ であり、従って $M/M[I^\infty]$ は $M''$ の商なので
$I$-冪ねじれである。これと第一の主張および Lemma \ref{more-algebra-lemma-torsion-free} を
合わせると、この商が零であることが従う。
\end{proof}

\begin{lemma}
\label{more-algebra-lemma-I-power-torsion}
$I$ を環 $R$ の有限生成イデアルとする。$I$-冪ねじれ加群はアーベル圏
$\text{Mod}_R$ の Serre 部分圏をなす。Homology, Definition
\ref{homology-definition-serre-subcategory} を見よ。
\end{lemma}

\begin{proof}
% P02-E1075: frozen source lines 24077-24078 use singular "is" for a compound subject.
$I$-冪ねじれ加群の部分加群と商加群が $I$-冪ねじれであることは明らかである。
さらに、二つの $I$-冪ねじれ加群の拡大は Lemma
\ref{more-algebra-lemma-divide-by-torsion} により $I$-冪ねじれである。
% P02-E1076: frozen source lines 24081-24082 omit the verb "follows".
従って補題の主張は Homology, Lemma
\ref{homology-lemma-characterize-serre-subcategory} から従う。
\end{proof}

\begin{lemma}
\label{more-algebra-lemma-local-cohomology-closed}
$R$ を環、$I \subset R$ を有限生成イデアルとする。部分圏
$I^\infty\text{-torsion} \subset \text{Mod}_R$ は閉部分集合
$Z = V(I) \subset \Spec(R)$ のみに依存する。実際、$R$-加群 $M$ が
$I$-冪ねじれであることと、その台が $Z$ に含まれることは同値である。
\end{lemma}

\begin{proof}
$M$ を $R$-加群、$x \in M$ とする。$x \in M[I^\infty]$ ならば、すべての
$f \in I$ に対して $x$ は $M_f$ で零に写る。従って、すべての
$\mathfrak p \not \supset I$ に対して $x$ は $M_\mathfrak p$ で零に写る。
逆に、すべての $\mathfrak p \not \supset I$ に対して $x$ が
$M_\mathfrak p$ で零に写るならば、すべての $f \in I$ に対して $x$ は
$M_f$ で零に写る。従って $I = (f_1, \ldots, f_r)$ ならば、ある
$n_i \geq 1$ に対して $f_i^{n_i}x = 0$ である。ゆえに
$x \in M[I^{\sum n_i}]$ である。従って $M[I^\infty]$ は写像
$M \to \prod_{\mathfrak p \not \in Z} M_\mathfrak p$ の核である。補題の第二の
主張が従い、それは第一の主張を導く。
\end{proof}

\noindent
次の二つの補題は、おそらく別の場所に置くべきである。

\begin{lemma}
\label{more-algebra-lemma-bounded-above-mod-I}
$R$ を環、$I \subset R$ をイデアルとする。$K$ を
$i > 0$ に対して $H^i(K \otimes_R^\mathbf{L} R/I) = 0$ を満たす $D(R)$ の
対象とする。このとき次が成り立つ。
\begin{enumerate}
\item すべての $n \geq 1$ と $i > 0$ に対して
$H^i(K \otimes_R^\mathbf{L} R/I^n) = 0$ である。
\item 任意の $I$-冪ねじれ $R$-加群 $N$ と $i > 0$ に対して
$H^i(K \otimes_R^\mathbf{L} N) = 0$ である。
\item コホモロジー加群 $H^i(M)$ が $I$-冪ねじれであり、$i > 0$ に対して
$0$ である任意の $M \in D^b(R)$ に対して
$H^i(K \otimes_R^\mathbf{L} M) = 0$ が $i > 0$ で成り立つ。
\end{enumerate}
\end{lemma}

\begin{proof}
(2) の証明。$N = \bigcup N[I^n]$ と書ける。テンソル積は余極限と可換なので
（詳細は省略する。ヒント：$K$ を K-flat 複体で表して直接計算せよ）、
$K \otimes_R^\mathbf{L} N = \text{hocolim}_n K \otimes_R^\mathbf{L} N[I^n]$
である。従って $N$ は $I^n$ によって零化されると仮定してよい。
$N$ を平方零イデアルとする $R$-代数 $R' = R/I^n \oplus N$ を考える。
$D(R')$ の対象 $K' = K \otimes_R^\mathbf{L} R'$ が正次数で消滅する
コホモロジーを持つことを示せば十分である。核を $J$ とする $R$-代数の全射
$R' \to R/I$ がある。
% P02-E1077: frozen source line 24135 gives an insufficient nilpotence exponent.
この核は冪零である（核の任意の $n$ 個の元の積は零である）。さらに
$$
K \otimes_R^\mathbf{L} R/I =
(K \otimes_R^\mathbf{L} R') \otimes_{R'}^\mathbf{L} R/I =
K' \otimes_{R'}^\mathbf{L} R/I
$$
である。仮定により複体 $K \otimes_R^\mathbf{L} R/I$ は $[-\infty, 0]$ に
Tor 振幅を持つ。
% P02-E1078: frozen source lines 24143-24144 omit the verb "follows".
従って結論は Lemma \ref{more-algebra-lemma-check-tor-dimension-modulo-nilpotent} から従う。

\medskip\noindent
(1) は (2) から直ちに従う。(3) は (2)、$M$ の非零コホモロジー加群の個数に
関する帰納法、および Derived Categories, Remark
\ref{derived-remark-truncation-distinguished-triangle} の切断の完全三角から従う。
詳細は省略する。
\end{proof}

\begin{lemma}
\label{more-algebra-lemma-derived-vanishing-mod-I}
$R$ を環、$I \subset R$ をイデアルとする。$K$ を
$D(R)$ において $K \otimes_R^\mathbf{L} R/I = 0$ を満たす $D(R)$ の対象と
する。このとき次が成り立つ。
\begin{enumerate}
\item すべての $n \geq 1$ に対して $K \otimes_R^\mathbf{L} R/I^n = 0$ である。
\item 任意の $I$-冪ねじれ $R$-加群 $N$ に対して
$K \otimes_R^\mathbf{L} N = 0$ である。
\item コホモロジー加群が $I$-冪ねじれである任意の $M \in D^b(R)$ に対して
$K \otimes_R^\mathbf{L} M = 0$ である。
\end{enumerate}
\end{lemma}

\begin{proof}
Lemma \ref{more-algebra-lemma-bounded-above-mod-I} の帰結である（いくつかの詳細は省略する）。
\end{proof}

\begin{lemma}
\label{more-algebra-lemma-torsion-completion}
\begin{reference}
\cite[Lemme~1]{Beauville-Laszlo} のわずかな一般化である。
\end{reference}
$R \to R'$ を環写像とする。$I \subset R$ を、$n > 0$ に対して
$R/I^n \to R'/I^nR'$ が同型となるイデアルとする。任意の $I$-冪ねじれ
$R$-加群 $M$ に対して、写像 $M \to M \otimes_R R'$ は同型である。
例えば $I$ が有限生成であり、$R^\wedge$ が $I$ に関する $R$ の完備化ならば、
$M \cong M \otimes_R R^\wedge$ である。
\end{lemma}

\begin{proof}
$M$ が $I^n$ によって零化されるならば
$$
M \otimes_R R' \cong
M \otimes_{R/I^n} R'/I^n R' \cong
M \otimes_{R/I^n} R/I^n \cong M.
$$
$M$ が $I$-冪ねじれならば $M = \bigcup M[I^n]$ である。テンソル積は
順極限と可換するから
（Algebra, Lemma \ref{algebra-lemma-tensor-products-commute-with-limits}）、
所望の同型を得る。最後の主張は Algebra, Lemma
\ref{algebra-lemma-hathat-finitely-generated} により第一の主張の特別な場合で
ある。
\end{proof}

\section{加群圏の形式的貼り合わせ}
\label{more-algebra-section-formal-glueing}

\noindent
Noether スキーム $X$ と、補集合を $U$ とする閉部分スキーム $Z$ を固定する。
目標は、$X$ 上の連接層を、$Z$ に沿う $X$ の形式完備化上の連接層と、$U$ 上の
連接層および重なり上の適切な両立性から（一意に）構成できることを説明する
ことである。まず可換代数のみを用いてこれを行い（本節）、後に代数空間の枠組み
で説明する（Pushouts of Spaces, Section
\ref{spaces-pushouts-section-formal-glueing}）。

\medskip\noindent
本節の内容の一部を扱う参考文献として次がある：
\cite[Section 2]{ArtinII},
\cite[Appendix]{Ferrand-Raynaud},
\cite{Beauville-Laszlo},
\cite{MB}, and
\cite[Section 4.6]{dJ-crystalline}。

\begin{lemma}
\label{more-algebra-lemma-characterize-flatness-on-torsion}
$\varphi : R \to S$ を環写像、$I \subset R$ をイデアルとする。次は同値である。
\begin{enumerate}
\item $\varphi$ は平坦であり、$R/I \to S/IS$ は忠実平坦である。
\item $\varphi$ は平坦であり、写像
$\Spec(S/IS) \to \Spec(R/I)$ は全射である。
\item $\varphi$ は平坦であり、基底変換関手 $M \mapsto M \otimes_R S$ は
$I$ によって零化される加群上で忠実である。
\item $\varphi$ は平坦であり、基底変換関手 $M \mapsto M \otimes_R S$ は
$I$-冪ねじれ加群上で忠実である。
\end{enumerate}
\end{lemma}

\begin{proof}
$R \to S$ が平坦ならば、任意の $n$ に対して $R/I^n \to S/I^nS$ は平坦で
ある。Algebra, Lemma \ref{algebra-lemma-flat-base-change} を見よ。従って (1) と
(2) は Algebra, Lemma \ref{algebra-lemma-ff-rings} により同値である。
$I$ によって零化される $R$-加群を $R/I$-加群と同一視し、
$I$ によって零化される $R$-加群 $M$ に対して
$$
M \otimes_R S = M \otimes_{R/I} S/IS
$$
であることと Algebra, Lemma \ref{algebra-lemma-easy-ff} を用いると、(1) と
(3) の同値性が従う。(4) $\Rightarrow$ (3) は直ちに分かる。(3) を仮定する。
上で $R/I^n \to S/I^nS$ が平坦であることを見た。また仮定により、この写像は
スペクトル上の全射を導く。実際、$\Spec(R/I^n) = \Spec(R/I)$ であり、$S$ に
ついても同様である。従って基底変換関手は $I^n$ によって零化される加群上で
忠実である。任意の $I$-冪ねじれ加群 $M$ は、$M_n$ が $I^n$ によって零化
される合併 $M = \bigcup M_n$ であるから、基底変換関手はすべての
$I$-冪ねじれ加群の圏上で忠実である（テンソル積は余極限と可換する）。
\end{proof}

\begin{lemma}
\label{more-algebra-lemma-neighbourhood-isomorphism}
$(\varphi : R \to S, I)$ が Lemma
\ref{more-algebra-lemma-characterize-flatness-on-torsion} の同値な条件を満たすと仮定する。
次は同値である。
\begin{enumerate}
\item 任意の $I$-冪ねじれ加群 $M$ に対して、自然な写像
$M \to M \otimes_R S$ は同型である。
\item $R/I \to S/IS$ は同型である。
\end{enumerate}
\end{lemma}

\begin{proof}
(1) $\Rightarrow$ (2) は直ちに分かる。(2) を仮定する。まず $M$ が $I$ に
よって零化されると仮定する。この場合 $M$ は $R/I$-加群であり、同型
$$
M \otimes_R S = M \otimes_{R/I} S/IS = M \otimes_{R/I} R/I = M
$$
を得るので主張が示される。次に、$I^n$ によって零化される任意の加群 $M$ に
対して $M \to M \otimes_R S$ が同型であることを帰納法で示す。
% P02-E1079: frozen source line 24283 omits "which/that" before "is annihilated".
$n$ に対して帰納法の仮定が成り立つとし、$M$ が $I^{n + 1}$ によって零化
されると仮定する。このとき短完全列
$$
0 \to I^nM \to M \to M/I^nM \to 0
$$
がある。$R \to S$ は平坦なので、これは短完全列
$$
0 \to I^nM \otimes_R S \to M \otimes_R S \to M/I^nM \otimes_R S \to 0
$$
を生じる。$M' = I^nM$ と $M'' = M/I^nM$ に対して標準写像が同型である
こと（帰納法の仮定による）を用いると、$M$ に対しても同じことが成り立つ。
最後に $M$ を一般の $I$-冪ねじれ加群とする。このとき、$M_n$ が $I^n$ に
よって零化される合併 $M = \bigcup M_n$ であり、テンソル積が余極限と可換する
ことから結論を得る。
\end{proof}

\begin{lemma}
\label{more-algebra-lemma-neighbourhood-equivalence}
$\varphi : R \to S$ を平坦な環写像、$I \subset R$ を
$R/I \to S/IS$ が同型となる有限生成イデアルとする。このとき次が成り立つ。
\begin{enumerate}
\item 任意の $R$-加群 $M$ に対して、写像 $M \to M \otimes_R S$ は
$I$-冪ねじれ部分加群の同型
$M[I^\infty] \to (M \otimes_R S)[(IS)^\infty]$ を導く。
\item $M$ または $N$ のいずれかが $I$-冪ねじれならば、自然な写像
$$
\Hom_R(M, N) \longrightarrow \Hom_S(M \otimes_R S, N \otimes_R S)
$$
は同型である。
\item 基底変換関手 $M \mapsto M \otimes_R S$ は、$I$-冪ねじれ加群と
$IS$-冪ねじれ加群の圏の間の圏同値を定める。
\end{enumerate}
\end{lemma}

\begin{proof}
Lemma \ref{more-algebra-lemma-characterize-flatness-on-torsion} と
Lemma \ref{more-algebra-lemma-neighbourhood-isomorphism} の同値な条件がともに満たされることに
注意する。以後これを断りなく用いる。まず (1) を証明する。$M$ を任意の
$R$-加群とする。$M' = M/M[I^\infty]$ と置き、完全列
$$
0 \to M[I^\infty] \to M \to M' \to 0
$$
を考える。$M[I^\infty] = M[I^\infty] \otimes_R S$ であるから、
$(M' \otimes_R S)[(IS)^\infty] = 0$ を示せば十分である。
$I = (f_1, \ldots, f_t)$ と書く。Lemma \ref{more-algebra-lemma-divide-by-torsion} により
$M'[I^\infty] = 0$ である。従って任意の $n > 0$ に対して写像
$$
M' \longrightarrow \bigoplus\nolimits_{i = 1, \ldots t} M',
\quad
x \longmapsto (f_1^n x, \ldots, f_t^n x)
$$
は単射である。$S$ は $R$ 上平坦なので、対応する写像
$M' \otimes_R S \to \bigoplus_{i = 1, \ldots t} M' \otimes_R S$ も単射で
ある。これは所望のとおり $(M' \otimes_R S)[I^n] = 0$ を意味する。

\medskip\noindent
次に (2) を証明する。$N$ が $I$-冪ねじれならば $N \otimes_R S = N$ であり、
(2) の表示写像は Algebra, Lemma
\ref{algebra-lemma-adjoint-tensor-restrict} により同型である。$M$ が
$I$-冪ねじれならば、任意の写像 $M \to N$ の像は $N[I^\infty]$ を経由し、
任意の写像 $M \otimes_R S \to N \otimes_R S$ の像は
$(N \otimes_R S)[(IS)^\infty]$ を経由する。従ってこの場合、(1) により
$N$ を $N[I^\infty]$ で置き換えてよく、結果は先に論じた $N$ が
$I$-冪ねじれである場合から従う。

\medskip\noindent
次に (3) を証明する。関手は (2) により充満忠実である。本質的全射性については、
任意の $IS$-冪ねじれ $S$-加群 $N$ に対して自然な写像
$N \otimes_R S \to N$ が同型であることに注意すればよい。
\end{proof}

\begin{lemma}
\label{more-algebra-lemma-map-identifies-koszul-and-cech-complexes}
$\varphi : R \to S$ を平坦な環写像、$I \subset R$ を
$R/I \to S/IS$ が同型となる有限生成イデアルとする。
$V(f_1, \ldots, f_r) = V(I)$ を満たす任意の $f_1, \ldots, f_r \in R$ に対し、
次が成り立つ。
\begin{enumerate}
\item Koszul 複体の写像
$K(R, f_1, \ldots, f_r) \to K(S, f_1, \ldots, f_r)$ は擬同型である。
\item 拡張交代 {\v C}ech 複体の写像
$$
\xymatrix{
R \to \prod_{i_0} R_{f_{i_0}} \to \prod_{i_0 < i_1} R_{f_{i_0}f_{i_1}}
\to \ldots \to R_{f_1\ldots f_r} \ar[d] \\
S \to \prod_{i_0} S_{f_{i_0}} \to \prod_{i_0 < i_1} S_{f_{i_0}f_{i_1}}
\to \ldots \to S_{f_1\ldots f_r}
}
$$
は擬同型である。
\end{enumerate}
\end{lemma}

\begin{proof}
どちらの場合も $R$-加群の複体 $K_\bullet$ があり、
% P02-E1080: frozen source line 24384 omits the hyphen in "R-modules".
$K_\bullet \to K_\bullet \otimes_R S$ が擬同型であることを示したい。
Lemma \ref{more-algebra-lemma-neighbourhood-isomorphism} と $R \to S$ の平坦性により、
$K$ のすべてのホモロジー群が $I$-冪ねじれならばこれが成り立つ。Koszul 複体に
対しては Lemma \ref{more-algebra-lemma-homotopy-koszul} により、拡張交代 {\v C}ech 複体に
対しては Lemma \ref{more-algebra-lemma-extended-alternating-torsion} により、これは真である。
\end{proof}

\begin{lemma}
\label{more-algebra-lemma-naive-Koszul-complex}
$R$ を環、$I = (f_1, \ldots, f_n)$ を $R$ の有限生成イデアルとする。
$M$ を、元 $e_1, \ldots, e_n$ と関係式 $f_i e_j - f_j e_i = 0$ によって生成
される $R$-加群とする。このとき短完全列
$$
0 \to K \to M \to I \to 0
$$
であって、$K$ が $I$ によって零化されるものが存在する。
\end{lemma}

\begin{proof}
これは Koszul 複体の切断にほかならない。写像 $M \to I$ は規則
$e_i \mapsto f_i$ によって定まる。$m = \sum a_i e_i$ が $M \to I$ の核に
属する、すなわち $\sum a_i f_i = 0$ ならば、
$f_j m = \sum f_j a_i e_i = (\sum f_i a_i) e_j = 0$ である。
\end{proof}

\begin{lemma}
\label{more-algebra-lemma-explicit-ext}
$R$ を環、$I = (f_1, \ldots, f_n)$ を $R$ の有限生成イデアルとする。任意の
$R$-加群 $N$ に対して
$$
H_1(N, f_\bullet) =
\frac{\{(x_1, \ldots, x_n) \in N^{\oplus n} \mid f_i x_j = f_j x_i \}}
{\{f_1x, \ldots, f_nx) \mid x \in N\}}
$$
と置く。
% P02-E1081: frozen source lines 24417-24420 have mismatched tuple delimiters in the denominator.
任意の $R$-加群 $N$ に対して標準的な短完全列
% P02-E1082: frozen source lines 24421-24424 call a one-sided exact sequence "short exact".
$$
0 \to \Ext_R(R/I, N) \to H_1(N, f_\bullet) \to \Hom_R(K, N)
$$
が存在する。ここで $K$ は Lemma \ref{more-algebra-lemma-naive-Koszul-complex} のものとする。
\end{lemma}

\begin{proof}
上の記法における $\Ext$-群は Homology, Section
\ref{homology-section-extensions} で定義された $\text{Mod}_R$ におけるものを
表す。これらは $\Ext_R(M, N)$ と書かれる。短完全列
$0 \to I \to R \to R/I \to 0$ に付随する Homology, Lemma
\ref{homology-lemma-six-term-sequence-ext} の長完全列と
$\Ext_R(R, N) = 0$ を用いると
$$
\Ext_R(R/I, N) =
\Coker(N \longrightarrow \Hom(I, N))
$$
が分かる。Lemma \ref{more-algebra-lemma-naive-Koszul-complex} の短完全列を用いると、複体
$$
N \to \Hom(M, N) \to \Hom_R(K, N)
$$
を得て、その中央のホモロジーは標準的に $\Ext_R(R/I, N)$ と同型である。
第一の写像の余核は標準的に $H_1(N, f_\bullet)$ と同型なので、補題の証明が
完了する。
\end{proof}

\begin{lemma}
\label{more-algebra-lemma-koszul-homology-annihilated}
$R$ を環、$I = (f_1, \ldots, f_n)$ を $R$ の有限生成イデアルとする。任意の
$R$-加群 $N$ に対して、Lemma \ref{more-algebra-lemma-explicit-ext} で定義した Koszul
ホモロジー群 $H_1(N, f_\bullet)$ は $I$ によって零化される。
\end{lemma}

\begin{proof}
$(x_1, \ldots, x_n) \in N^{\oplus n}$ が $f_i x_j = f_j x_i$ を満たすとする。
% P02-E1083: frozen source line 24464 gives the wrong component pattern; the boundary vector is (f_1 x_i, ..., f_n x_i).
このとき $f_i(x_1, \ldots, x_n) = (f_i x_i, \ldots, f_i x_n)$ である。
言い換えると、$f_i$ は $H_1(N, f_\bullet)$ を零化する。
\end{proof}

\noindent
Lemma \ref{more-algebra-lemma-neighbourhood-equivalence} の充満忠実性を改良し、始域が
$I$-冪ねじれである $\Ext$-群も $S$ への移行で変化しないことを示せる。
次の補題の導来的な版については Dualizing Complexes, Lemma
\ref{dualizing-lemma-neighbourhood-extensions} を見よ。

\begin{lemma}
\label{more-algebra-lemma-neighbourhood-extensions}
$\varphi : R \to S$ を平坦な環写像、$I \subset R$ を
$R/I \to S/IS$ が同型となる有限生成イデアルとする。$M$, $N$ を $R$-加群とし、
$M$ が $I$-冪ねじれであると仮定する。
% P02-E1084: frozen source line 24481 says "an short exact sequence".
短完全列
$$
0 \to N \otimes_R S \to \tilde E \to M \otimes_R S \to 0
$$
が与えられると、行が完全である可換図式
$$
\xymatrix{
0 \ar[r] &
N \ar[r] \ar[d] &
E \ar[r] \ar[d] &
M \ar[r] \ar[d] &
0 \\
0 \ar[r] &
N \otimes_R S \ar[r] &
\tilde E \ar[r] &
M \otimes_R S \ar[r] &
0
}
$$
が存在する。
\end{lemma}

\begin{proof}
$M$ は $I$-冪ねじれなので $M \otimes_R S = M$ である。Lemma
\ref{more-algebra-lemma-neighbourhood-isomorphism} を見よ。以後この同一視を断りなく用いる。
$R \to S$ は平坦なので基底変換関手は完全であり、$\Ext$-群の関手的な写像
$$
\Ext_R(M, N)
\longrightarrow
\Ext_S(M \otimes_R S, N \otimes_R S),
$$
を得る。Homology, Lemma \ref{homology-lemma-exact-functor-ext} を見よ。補題の
主張は、$M$ が $I$-冪ねじれであるとき、この写像が全射であるということである。
実際、これが同型であることを示す。Lemma
\ref{more-algebra-lemma-I-power-torsion-presentation} により、$R/I^n$ の形の加群の直和である
$M'$ からの全射 $M' \to M$ を見いだせる。Homology, Lemma
\ref{homology-lemma-six-term-sequence-ext} の長完全列と Lemma
\ref{more-algebra-lemma-neighbourhood-equivalence} を用いると、$M'$ に対して補題を証明すれば
十分であることが分かる。$\Ext$ と直和との両立性（詳細は省略する）を用いて、
ある $n$ に対する $M = R/I^n$ の場合に帰着する。

\medskip\noindent
$f_1, \ldots, f_t$ を $I^n$ の生成元とする。Lemma
\ref{more-algebra-lemma-explicit-ext} により、行が完全である可換図式
$$
\xymatrix{
0 \ar[r] &
\Ext_R(R/I^n, N) \ar[r] \ar[d] &
H_1(N, f_\bullet) \ar[r] \ar[d] &
\Hom_R(K, N) \ar[d] \\
0 \ar[r] &
\Ext_S(S/I^nS, N \otimes S) \ar[r] &
H_1(N \otimes S, f_\bullet) \ar[r] &
\Hom_S(K \otimes S, N \otimes S)
}
$$
を得る。ここで $K$ は Lemma \ref{more-algebra-lemma-naive-Koszul-complex} のものとする。
従って右側の二つの垂直射が同型であることを示せば十分である。$K$ は $I^n$ に
よって零化されるから、Lemma \ref{more-algebra-lemma-neighbourhood-equivalence} により
$\Hom_R(K, N) = \Hom_S(K \otimes_R S, N \otimes_R S)$ である。
$R \to S$ は平坦なので
$H_1(N, f_\bullet) \otimes_R S = H_1(N \otimes_R S, f_\bullet)$ である。
Lemma \ref{more-algebra-lemma-koszul-homology-annihilated} により
$H_1(N, f_\bullet)$ は $I^n$ によって零化されるから、Lemma
\ref{more-algebra-lemma-neighbourhood-isomorphism} により
$H_1(N, f_\bullet) \otimes_R S = H_1(N, f_\bullet)$ である。
\end{proof}

\noindent
$R \to S$ を環写像とする。$f_1, \ldots, f_t \in R$ とし、
$I = (f_1, \ldots, f_t)$ と置く。このとき任意の $R$-加群 $M$ に対して複体
\begin{equation}
\label{more-algebra-equation-glueing-complex}
0 \to M \xrightarrow{\alpha}
M \otimes_R S \times \prod M_{f_i} \xrightarrow{\beta}
\prod (M \otimes_R S)_{f_i}
\times
\prod M_{f_if_j}
\end{equation}
を定義できる。ここで
$\alpha(m) = (m \otimes 1, m/1, \ldots, m/1)$ および
$$
\beta(m', m_1, \ldots, m_t) =
((m'/1 - m_1 \otimes 1, \ldots, m'/1 - m_t \otimes 1),
(m_1 - m_2, \ldots, m_{t - 1} - m_t).
$$
% P02-E1085: frozen source lines 24572-24575 leave the outer beta tuple unclosed.
この複体がいつ完全であるかを知りたい。

\begin{lemma}
\label{more-algebra-lemma-recover-module-from-glueing-data}
$\varphi : R \to S$ を平坦な環写像とし、
$I = (f_1, \ldots, f_t) \subset R$ を $R/I \to S/IS$ が同型となるイデアルと
する。$M$ を $R$-加群とする。このとき複体
(\ref{more-algebra-equation-glueing-complex}) は完全である。
\end{lemma}

\begin{proof}
第一の証明。Lemma \ref{more-algebra-lemma-map-identifies-koszul-and-cech-complexes} の拡張交代
{\v C}ech 複体の擬同型を
% P02-E1086: frozen source lines 24589-24591 omit "by" in the Denote construction.
$\check{\mathcal{C}}_R \to \check{\mathcal{C}}_S$ と表す。これらの複体は平坦な
項を持つ有界複体なので、
$M \otimes_R \check{\mathcal{C}}_R \to
M \otimes_R \check{\mathcal{C}}_S$ も擬同型である（Lemmas
\ref{more-algebra-lemma-derived-tor-quasi-isomorphism} and
\ref{more-algebra-lemma-derived-tor-quasi-isomorphism-other-side}）。複体
(\ref{more-algebra-equation-glueing-complex}) は写像
$M \otimes_R \check{\mathcal{C}}_R \to
M \otimes_R \check{\mathcal{C}}_S$ の錐の切断なので、結論が従う。

\medskip\noindent
第二の計算的証明。$m \in M$ とする。$\alpha(m) = 0$ ならば
$m \in M[I^\infty]$ である。Lemma \ref{more-algebra-lemma-torsion-free} を見よ。
$I^n m = 0$ となる $n$ を選び、写像 $\varphi : R/I^n \to M$ を考える。
$m \otimes 1 = 0$ ならば $\varphi \otimes 1_S = 0$ であり、従って
Lemma \ref{more-algebra-lemma-neighbourhood-equivalence} により $\varphi = 0$、ゆえに
$m = 0$ である。このようにして $\alpha$ が単射であることが分かる。

\medskip\noindent
$(m', m'_1, \ldots, m'_t) \in \Ker(\beta)$ とする。ある $n > 0$ と
$m_i \in M$ によって $m'_i = m_i/f_i^n$ と書く。必要なら $n$ を大きくして、
$M \otimes_R S$ において $f_i^n m' = m_i \otimes 1$、また $M$ において
$f_j^nm_i - f_i^nm_j = 0$ と仮定してよい。特に
$(m_1, \ldots, m_t)$ は
$H_1(M, (f_1^n, \ldots, f_t^n))$ の元 $\xi$ を定める。
$H_1(M, (f_1^n, \ldots, f_t^n))$ は $I^{tn + 1}$ によって零化される
（Lemma \ref{more-algebra-lemma-koszul-homology-annihilated} を見よ）。また $R \to S$ は
平坦なので、Lemma \ref{more-algebra-lemma-neighbourhood-isomorphism} により
$$
H_1(M, (f_1^n, \ldots, f_t^n)) =
H_1(M, (f_1^n, \ldots, f_t^n)) \otimes_R S =
H_1(M \otimes_R S, (f_1^n, \ldots, f_t^n))
$$
である。$m'$ の存在は $\xi$ が最後の群で零に写ること、すなわち元 $\xi$ が
零であることを意味する。従って $m_i = f_i^n m$ となる $m \in M$ が存在する。
このとき、ある $m'' \in (M \otimes_R S)[(IS)^\infty]$ に対して
$(m', m'_1, \ldots, m'_t) - \alpha(m) = (m'', 0, \ldots, 0)$ である。
Lemma \ref{more-algebra-lemma-neighbourhood-equivalence} により $m'' \in M[I^\infty]$ と結論
でき、証明が完了する。
\end{proof}

\begin{remark}
\label{more-algebra-remark-glueing-data}
この注意では貼り合わせデータの圏を定義する。$R \to S$ を環写像とする。
$f_1, \ldots, f_t \in R$ とし、$I = (f_1, \ldots, f_t)$ と置く。
圏 $\text{Glue}(R \to S, f_1, \ldots, f_t)$ を次の圏として考える。
\begin{enumerate}
\item 対象は系 $(M', M_i, \alpha_i, \alpha_{ij})$ である。ここで $M'$ は
$S$-加群、$M_i$ は $R_{f_i}$-加群、
$\alpha_i : (M')_{f_i} \to M_i \otimes_R S$ は同型、
$\alpha_{ij} : (M_i)_{f_j} \to (M_j)_{f_i}$ は次を満たす同型である。
\begin{enumerate}
\item 写像 $(M')_{f_if_j} \to (M_j)_{f_i}$ として
$\alpha_{ij} \circ \alpha_i = \alpha_j$ である。
\item 写像 $(M_i)_{f_jf_k} \to (M_k)_{f_if_j}$ として
$\alpha_{jk} \circ \alpha_{ij} = \alpha_{ik}$ である（コサイクル条件）。
\end{enumerate}
\item 射
$(M', M_i, \alpha_i, \alpha_{ij}) \to (N', N_i, \beta_i, \beta_{ij})$ は、
与えられた写像 $\alpha_i, \beta_i, \alpha_{ij}, \beta_{ij}$ と両立する写像
$\varphi' : M' \to N'$ および $\varphi_i : M_i \to N_i$ によって与えられる。
\end{enumerate}
標準的な関手
$$
\text{Can} : \text{Mod}_R
\longrightarrow
\text{Glue}(R \to S, f_1, \ldots, f_t),
\quad
M \longmapsto (M \otimes_R S, M_{f_i}, \text{can}_i, \text{can}_{ij})
$$
がある。ここで
$\text{can}_i : (M \otimes_R S)_{f_i} \to M_{f_i} \otimes_R S$ および
$\text{can}_{ij} : (M_{f_i})_{f_j} \to (M_{f_j})_{f_i}$ は標準的な同型である。
圏 $\text{Glue}(R \to S, f_1, \ldots, f_t)$ の任意の対象
$\mathbf{M} = (M', M_i, \alpha_i, \alpha_{ij})$ に対して
$$
H^0(\mathbf{M}) =
\{(m', m_i) \mid \alpha_i(m') = m_i \otimes 1, \alpha_{ij}(m_i) = m_j\}
$$
と定義する。言い換えると、これは
(\ref{more-algebra-equation-glueing-complex}) と同様な完全列
$$
0 \to H^0(\mathbf{M}) \to
M' \times \prod M_i \to
\prod M'_{f_i}
\times
\prod (M_i)_{f_j}
$$
によって定義される。$H^0(\mathbf{M})$ を $R$-加群と考える。従って関手
$$
H^0 :
\text{Glue}(R \to S, f_1, \ldots, f_t)
\longrightarrow
\text{Mod}_R
$$
も得る。次の目標は、関手 $\text{Can}$ と $H^0$ が場合によって互いに
擬逆となることを示すことである。
\end{remark}

\begin{lemma}
\label{more-algebra-lemma-H0-Can-adjoint}
Remark \ref{more-algebra-remark-glueing-data} において、関手
$H^0 : \text{Glue}(R \to S, f_1, \ldots, f_t) \to \text{Mod}_R$ は関手
$\text{Can} : \text{Mod}_R \to \text{Glue}(R \to S, f_1, \ldots, f_t)$ の
右随伴である。
\end{lemma}

\begin{proof}
$\mathbf{M} = (M', M_i, \alpha_i, \alpha_{ij})$ を
$\text{Glue}(R \to S, f_1, \ldots, f_t)$ の対象とする。任意の $R$-加群 $N$ に
対して写像
$$
\Hom_{\text{Glue}(R \to S, f_1, \ldots, f_t)}(\text{Can}(N), \mathbf{M})
\to
\Hom_R(N, H^0(\mathbf{M}))
$$
があり、$\psi$ を、$H^0(\psi)$ と明らかな写像
$N \to H^0(\text{Can}(N))$ との合成に写す。構成により、この表示写像は
$N = R$ に対して同型である（一般には $R \to H^0(\text{Can}(R))$ が同型で
なくてもよい）。圏 $\text{Glue}(R \to S, f_1, \ldots, f_t)$ は直和と余核を
持つ。関手 $\text{Can}$ は直和および余核と可換する。$N$ を自由
$R$-加群間の写像の余核として書くことにより、これらの観察から表示写像が
全単射であることが分かる。詳細は省略する。
\end{proof}

\begin{lemma}
\label{more-algebra-lemma-H0-inverse}
$\varphi : R \to S$ を平坦な環写像とし、
$I = (f_1, \ldots, f_t) \subset R$ を $R/I \to S/IS$ が同型となるイデアルと
する。このとき関手 $H^0$ は Remark \ref{more-algebra-remark-glueing-data} の関手
$\text{Can}$ の左擬逆である。
\end{lemma}

\begin{proof}
Lemma \ref{more-algebra-lemma-recover-module-from-glueing-data} の言い換えである。
\end{proof}

\begin{lemma}
\label{more-algebra-lemma-exact}
$\varphi : R \to S$ を平坦な環写像とし、
$I = (f_1, \ldots, f_t) \subset R$ をイデアルとする。このとき
$\text{Glue}(R \to S, f_1, \ldots, f_t)$ はアーベル圏であり、関手
$\text{Can}$ は完全で、任意の余極限と可換する。
\end{lemma}

\begin{proof}
圏 $\text{Glue}(R \to S, f_1, \ldots, f_t)$ の射
$$
(\varphi', \varphi_i) :
(M', M_i, \alpha_i, \alpha_{ij})
\to
(N', N_i, \beta_i, \beta_{ij})
$$
が与えられると、その核は存在し、対象
$(\Ker(\varphi'), \Ker(\varphi_i), \alpha_i, \alpha_{ij})$ に等しい。またその
余核も存在し、対象
$(\Coker(\varphi'), \Coker(\varphi_i), \beta_i, \beta_{ij})$ に等しい。
$R \to S$ が平坦であり、従って核と余核を取ることが
$- \otimes_R S$ と可換するため、これは機能する。詳細は省略する。完全性は
$R_{f_i}$ と $S$ が $R$-平坦であることから従い、余極限との可換性はテンソル積が
余極限と可換することから従う。
\end{proof}

\begin{lemma}
\label{more-algebra-lemma-equivalence-I-unit}
$\varphi : R \to S$ を平坦な環写像とし、$(f_1, \ldots, f_t) = R$ とする。
このとき $\text{Can}$ と $H^0$ は圏の擬逆同値
$$
\text{Mod}_R = \text{Glue}(R \to S, f_1, \ldots, f_t)
$$
を与える。
\end{lemma}

\begin{proof}
$\mathbf{M} = (M', M_i, \alpha_i, \alpha_{ij})$ を
$\text{Glue}(R \to S, f_1, \ldots, f_t)$ の対象とする。Algebra, Lemma
\ref{algebra-lemma-glue-modules} により、一意な加群 $M$ と、貼り合わせデータ
$\alpha_{ij}$ を復元する同型 $M_{f_i} \to M_i$ が存在する。このとき
$M'$ と $M \otimes_R S$ はともに、$f_i$ で局所化すると加群
$M_i \otimes_R S$ を復元する $S$-加群である。従って標準的な同型
$M \otimes_R S \to M'$ がある。これは $\mathbf{M}$ が $\text{Can}$ の本質像に
属することを示す。これと Lemma \ref{more-algebra-lemma-H0-inverse} から補題が従う。
\end{proof}

\begin{lemma}
\label{more-algebra-lemma-base-change-glue}
$\varphi : R \to S$ を平坦な環写像とし、$I = (f_1, \ldots, f_t)$ を
% P02-E1087: frozen source lines 24793-24794 say "and ideal" instead of "an ideal".
イデアルとする。$R \to R'$ を平坦な環写像とし、$S' = S \otimes_R R'$ と置く。
このとき圏と関手の可換図式
$$
\xymatrix{
\text{Mod}_R \ar[r]_-{\text{Can}} \ar[d]_{-\otimes_R R'} &
\text{Glue}(R \to S, f_1, \ldots, f_t) \ar[r]_-{H^0} \ar[d]^{-\otimes_R R'} &
\text{Mod}_R \ar[d]^{-\otimes_R R'} \\
\text{Mod}_{R'} \ar[r]^-{\text{Can}} &
\text{Glue}(R' \to S', f_1, \ldots, f_t) \ar[r]^-{H^0} &
\text{Mod}_{R'}
}
$$
を得る。
\end{lemma}

\begin{proof}
省略する。
\end{proof}

\begin{proposition}
\label{more-algebra-proposition-equivalence}
$\varphi : R \to S$ を平坦な環写像とし、
$I = (f_1, \ldots, f_t) \subset R$ を $R/I \to S/IS$ が同型となるイデアルと
する。このとき $\text{Can}$ と $H^0$ は圏の擬逆同値
$$
\text{Mod}_R = \text{Glue}(R \to S, f_1, \ldots, f_t)
$$
を与える。
\end{proposition}

\begin{proof}
$H^0 \circ \text{Can}$ が恒等関手と同型であることは既に見た。Lemma
\ref{more-algebra-lemma-H0-inverse} を見よ。
$\mathbf{M} = (M', M_i, \alpha_i, \alpha_{ij})$ を
$\text{Glue}(R \to S, f_1, \ldots, f_t)$ の対象とする。自然な射
$$
\Psi :
(H^0(\mathbf{M}) \otimes_R S, H^0(\mathbf{M})_{f_i},
\text{can}_i, \text{can}_{ij})
\longrightarrow
(M', M_i, \alpha_i, \alpha_{ij}).
$$
を得る。実際、定義により $H^0(\mathbf{M})$ には両立する $R$-加群写像
$H^0(\mathbf{M}) \to M'$ および $H^0(\mathbf{M}) \to M_i$ が備わっている。
この写像が同型であることを示す必要がある。

\medskip\noindent
添字 $i$ を一つ選び、$R' = R_{f_i}$ と置く。Lemmas
\ref{more-algebra-lemma-base-change-glue} and \ref{more-algebra-lemma-equivalence-I-unit} を組み合わせると、
$\Psi \otimes_R R'$ が同型であることが分かる。従って $\Psi$ の核と余核は、
それぞれ $(K, 0, 0, 0)$ および $(Q, 0, 0, 0)$ の形の系である。
$H^0((K, 0, 0, 0)) = K$ であること、$H^0$ が左完全であること、および構成に
より $H^0(\Psi)$ が全単射であることに注意する。従って $K = 0$、すなわち
$\Psi$ の核は零である。

\medskip\noindent
以上から、短完全列
$$
0 \to H^0(\mathbf{M}) \otimes_R S \to M' \to Q \to 0
$$
を得て、$M_i = H^0(\mathbf{M})_{f_i}$ である。$Q$ は
$Q = Q \otimes_R S$ を満たす $I$-冪ねじれ $R$-加群と考えてよい。Lemma
\ref{more-algebra-lemma-neighbourhood-extensions} により、行が完全である可換図式
$$
\xymatrix{
0 \ar[r] &
H^0(\mathbf{M}) \ar[r] \ar[d] &
E \ar[r] \ar[d] &
Q \ar[r] \ar[d] &
0 \\
0 \ar[r] &
H^0(\mathbf{M}) \otimes_R S \ar[r] &
M' \ar[r] &
Q \ar[r] &
0
}
$$
が存在する。これは圏 $\text{Glue}(R \to S, f_1, \ldots, f_t)$ における同型
$\text{Can}(E) \to (M', M_i, \alpha_i, \alpha_{ij})$ を明らかに定め、証明が
完了する（もちろん、事後的には $Q = 0$ である）。
\end{proof}

\begin{lemma}
\label{more-algebra-lemma-application-formal-glueing}
$\varphi : R \to S$ を平坦な環写像とし、$I \subset R$ を
$R/I \to S/IS$ が同型となる有限生成イデアルとする。
\begin{enumerate}
\item $R$-加群 $N$、$S$-加群 $M'$、および核と余核が $I$-冪ねじれである
$S$-加群写像 $\varphi : M' \to N \otimes_R S$ が与えられると、
$R$-加群写像 $\psi : M \to N$ と、$\varphi$ および $\psi$ と両立する同型
$M \otimes_R S = M'$ が存在する。
\item $R$-加群 $M$、$S$-加群 $N'$、および核と余核が $I$-冪ねじれである
$S$-加群写像 $\varphi : M \otimes_R S \to N'$ が与えられると、
$R$-加群写像 $\psi : M \to N$ と、$\varphi$ および $\psi$ と両立する同型
$N \otimes_R S = N'$ が存在する。
\end{enumerate}
どちらの場合も $\Ker(\varphi) \cong \Ker(\psi)$ および
$\Coker(\varphi) \cong \Coker(\psi)$ である。
\end{lemma}

\begin{proof}
(1) の証明。$I = (f_1, \ldots, f_t)$ とする。局所化 $\varphi_{f_i}$ が同型で
あることは明らかである。従って
$(M', N_{f_i}, \varphi_{f_i}, can_{ij})$ は
$\text{Glue}(R \to S, f_1, \ldots, f_t)$ の対象である。Remark
\ref{more-algebra-remark-glueing-data} を見よ。Proposition \ref{more-algebra-proposition-equivalence} により、
同型 $\varphi_{f_i}$ および $can_{ij}$ と両立して
$M' = M \otimes_R S$ および $N_{f_i} = M_{f_i}$ を満たす $R$-加群 $M$ が
存在する。第一成分で $\varphi$ を用いる
$\text{Glue}(R \to S, f_1, \ldots, f_t)$ の射
$$
(M \otimes_R S, M_{f_i}, can_i, can_{ij}) =
(M', N_{f_i}, \varphi_{f_i}, can_{ij})
\to
(N \otimes_R S, N_{f_i}, can_i, can_{ij})
$$
がある。これは Proposition \ref{more-algebra-proposition-equivalence} の圏同値により
$R$-加群写像 $\psi : M \to N$ に対応する。$M \to N$ の基底変換と同型
$M' \cong M \otimes_R S$ との合成は $\varphi$ である。言い換えると
$M \to N$ は $\varphi$ と両立する。

\medskip\noindent
(2) の証明。これは上の議論の双対にほかならない。実際、局所化
$\varphi_{f_i}$ は同型である。従って
$(N', M_{f_i}, \varphi_{f_i}^{-1}, can_{ij})$ は
$\text{Glue}(R \to S, f_1, \ldots, f_t)$ の対象である。Remark
\ref{more-algebra-remark-glueing-data} を見よ。Proposition \ref{more-algebra-proposition-equivalence} により、
同型 $\varphi_{f_i}^{-1}$ および $can_{ij}$ と両立して
$N' = N \otimes_R S$ および $N_{f_i} = M_{f_i}$ を満たす $R$-加群 $N$ が
存在する。第一成分で $\varphi$ を用いる
$\text{Glue}(R \to S, f_1, \ldots, f_t)$ の射
$$
(M \otimes_R S, M_{f_i}, can_i, can_{ij}) \to
(N', M_{f_i}, \varphi_{f_i}, can_{ij}) =
(N \otimes_R S, N_{f_i}, can_i, can_{ij})
$$
がある。
% P02-E1088: frozen source line 24939 uses phi_{f_i} where the glue datum just defined is its inverse.
これは Proposition \ref{more-algebra-proposition-equivalence} の圏同値により $R$-加群写像
$\psi : M \to N$ に対応する。$M \to N$ の基底変換と同型
$N' \cong N \otimes_R S$ との合成は $\varphi$ である。言い換えると
$M \to N$ は $\varphi$ と両立する。

\medskip\noindent
最後の主張は、例えば Lemma \ref{more-algebra-lemma-neighbourhood-equivalence} から従う。
\end{proof}

\noindent
次に Proposition \ref{more-algebra-proposition-equivalence} を特殊化して、より使いやすい
結果を得る。すなわち $I = (f)$ が単項イデアルならば、
$\text{Glue}(R \to S, f)$ の対象は単に三つ組 $(M', M_1, \alpha_1)$ であり、
確認すべきコサイクル条件は{\it ない}。

\begin{theorem}
\label{more-algebra-theorem-formal-glueing}
$R$ を環、$f \in R$ とする。$\varphi : R \to S$ を同型
$R/fR \to S/fS$ を導く平坦な環写像とする。このとき関手
$$
\text{Mod}_R
\longrightarrow
\text{Mod}_S \times_{\text{Mod}_{S_f}} \text{Mod}_{R_f},
\quad
M
\longmapsto
(M \otimes_R S, M_f, \text{can})
$$
は圏同値である。
\end{theorem}

\begin{proof}
矢印の右辺に現れる圏は、$M'$ が $S$-加群、$M_1$ が $R_f$-加群、
$\alpha_1 : M'_f \to M_1 \otimes_R S$ が $S_f$-同型である三つ組
$(M', M_1, \alpha_1)$ の圏である。Categories, Example
\ref{categories-example-2-fibre-product-categories} を見よ。従ってこの定理は
Proposition \ref{more-algebra-proposition-equivalence} の特別な場合である。
\end{proof}

\noindent
Theorem \ref{more-algebra-theorem-formal-glueing} の有用な特別な場合は、$R$ が Noether で、
$S$ が元 $f$ における $R$ の完備化である場合である。完備化
$R \to S$ は平坦であり、$M$ が有限生成ならば、関手
$M \mapsto M \otimes_R S$ を $f$-進完備化関手と同一視できる。これをより正確に
述べるため、有限生成 $R$-加群の圏を $\text{Mod}^{fg}_R$ と表す。

\begin{proposition}
\label{more-algebra-proposition-formal-glueing}
$R$ を Noether 環、$f \in R$ を元とする。$R^\wedge$ を $R$ の $f$-進完備化と
する。このとき関手 $M \mapsto (M^\wedge, M_f, \text{can})$ は圏同値
$$
\text{Mod}^{fg}_R
\longrightarrow
\text{Mod}^{fg}_{R^\wedge}
\times_{\text{Mod}^{fg}_{(R^\wedge)_f}}
\text{Mod}^{fg}_{R_f}
$$
を定める。
\end{proposition}

\begin{proof}
環写像 $R \to R^\wedge$ は Algebra, Lemma
\ref{algebra-lemma-completion-flat} により平坦である。
$R/fR = R^\wedge/fR^\wedge$ であることは明らかである。Algebra, Lemma
\ref{algebra-lemma-completion-tensor} により、有限 $R$-加群 $M$ の完備化は
$M \otimes_R R^\wedge$ に等しい。従って命題の表示関手は Theorem
\ref{more-algebra-theorem-formal-glueing} に現れる関手に等しい。特にこれは充満忠実である。
$(M_1, M_2, \psi)$ を右辺の対象とする。Theorem
\ref{more-algebra-theorem-formal-glueing} により、$M_1 = M \otimes_R R^\wedge$ および
$M_2 = M_f$ を満たす $R$-加群 $M$ が存在する。
$R \to R^\wedge \times R_f$ は忠実平坦なので、Algebra, Lemma
\ref{algebra-lemma-cover} により $M$ が有限生成、すなわち
$M \in \text{Mod}^{fg}_R$ であると結論できる。これで命題が証明された。
\end{proof}

\begin{remark}
\label{more-algebra-remark-formal-glueing-algebras}
Proposition \ref{more-algebra-proposition-equivalence}、Theorem
\ref{more-algebra-theorem-formal-glueing}、および Proposition
\ref{more-algebra-proposition-formal-glueing} の同値は加群の性質を保つ。例えば $M$ が
$\mathbf{M} = (M', M_i, \alpha_i, \alpha_{ij})$ に対応するならば、$M$ が $R$ 上
有限、有限表示、平坦、または射影的であることと、$M'$ および $M_i$ がそれぞれ
$S$ および $R_{f_i}$ 上で対応する性質を持つことは同値である。これは
$R \to S \times \prod R_{f_i}$ が忠実平坦であることと、忠実平坦写像に沿う
% P02-E1090: frozen source lines 25045-25049 have an ungrammatical "and descend and ascent" construction.
これらの性質の降下および上昇から従う。Algebra, Lemma
\ref{algebra-lemma-descend-properties-modules} and
Theorem \ref{algebra-theorem-ffdescent-projectivity} を見よ。これらの関手は両辺の
$\otimes$-構造も保つ。
% P02-E1089: frozen source line 25051 uses singular "it" after plural "These functors".
従って、代数の圏など、組 $(\text{Mod}_R, \otimes)$ から構成される種々の圏の
同値を定める。
\end{remark}

\begin{remark}
\label{more-algebra-remark-topological-analogue}
微分可能多様体 $X$ と、補集合を $U$ とするコンパクトな閉部分多様体 $Z$ が
与えられたとする。$X$ 上の層を指定することは、$U$ 上の層、$X$ における
$Z$ の特定されていない管状近傍 $T$ 上の層、および $T \cap U$ に沿って得られる
二つの層の間の同型を指定することと同じである。代数幾何において管状近傍は
そのままの形では存在しないが、Proposition \ref{more-algebra-proposition-equivalence}、
Theorem \ref{more-algebra-theorem-formal-glueing}、および Proposition
\ref{more-algebra-proposition-formal-glueing} のような結果により、代わりに形式近傍を用いる
ことができる。
\end{remark}

\section{Beauville--Laszlo の定理}
\label{more-algebra-section-beauville-laszlo}

\noindent
$R$ を環、$f$ を $R$ の元とする。$R$ の $f$-進完備化を
% P02-E1091: frozen source lines 25078-25079 omit "by" in the Denote construction.
$R^\wedge = \lim R/f^n R$ と表す。本節では Beauville と Laszlo の定理を論じ、
わずかに一般化する。\cite{Beauville-Laszlo} を見よ。この定理は、適切な条件の
もとで $R$ 上の加群を、$R^\wedge$ 上の加群と $R_f$ 上の加群を、それらの
$(R^\wedge)_f$ への基底拡大の間の同型に沿って ``貼り合わせる'' ことにより
構成できると主張する。

\medskip\noindent
\cite{Beauville-Laszlo} では、$f$ が $R$ と $M$ の両方で非零因子であることを
仮定する。実際には、
$$
R[f^\infty] \longrightarrow R^\wedge[f^\infty]
$$
が全単射であり、
$$
M[f^\infty] \longrightarrow M \otimes_R R^\wedge
$$
が単射であることのみを仮定すればよい。この最適化は、非アルキメデス解析空間の
理論で用いるために \cite[\S 1.3]{Kedlaya-Liu-I} で導入された貼り合わせへの別の
アプローチから一部着想を得たものである。

\medskip\noindent
実際には、環写像
$$
R \longrightarrow R'
$$
であって、任意の $n > 0$ に対して同型 $R/f^nR \to R'/f^nR'$ を、また同型
$R[f^\infty] \to R'[f^\infty]$ を導くものという、より一般の設定で
Beauville--Laszlo の定理を確立する。この設定は大域化により適しており、
$R'$ が $R$ の完備化である場合から形式的には従わない。例えば、
$R[f^\infty] \to R'[f^\infty]$ が全単射であるという条件は
$R[f^\infty] \to R^\wedge[f^\infty]$ が全単射であることを意味しない。

\medskip\noindent
本節で証明する Beauville と Laszlo の定理は Theorem
\ref{more-algebra-theorem-formal-glueing} の非平坦版と見ることができ、$R' = R^\wedge$ の
場合には Proposition \ref{more-algebra-proposition-formal-glueing} の非 Noether 版と見る
ことができる。平坦降下との比較については Remark \ref{more-algebra-remark-not-descent} を
見よ。

\medskip\noindent
さらに強い結果（例えば $M$ に制限を課さないもの）も確立できるが、そのためには
導来圏の水準で作業しなければならない。詳しくは
\cite[\S 5]{Bhatt-Algebraize} を見よ。

\begin{lemma}
\label{more-algebra-lemma-same-quotients}
$R$ を環、$f \in R$ とする。任意の正整数 $n$ に対して写像
$R/f^nR \to R^\wedge/f^n R^\wedge$ は同型である。
\end{lemma}

\begin{proof}
Algebra, Lemma \ref{algebra-lemma-hathat-finitely-generated} の特別な場合である。
\end{proof}

\noindent
Section \ref{more-algebra-section-torsion} で導入した記法を用いる。従って $R$-加群 $M$ に
対して、$f^n$ によって零化される $M$ の部分加群を
% P02-E1092: frozen source lines 25140-25142 omit "by" in the Denote construction.
$M[f^n]$ と表し、
$$
M[f^\infty] = \bigcup\nolimits_{n = 1}^\infty M[f^n] = \Ker(M \to M_f).
$$
と置く。$M = M[f^\infty]$ ならば、$M$ を $f$-冪ねじれ加群という。

\begin{lemma}
\label{more-algebra-lemma-BL-faithful}
$R$ を環、$f \in R$ とし、$R \to R'$ を、$n > 0$ に対して同型
$R/f^nR \to R'/f^nR'$ を導く環写像とする。$R$-加群 $R' \oplus R_f$ は
忠実である。すなわち任意の非零 $R$-加群 $M$ に対して、加群
$M \otimes_R (R' \oplus R_f)$ も非零である。例えば $M$ が非零ならば
$M \otimes_R (R^\wedge \oplus R_f)$ は非零である。
\end{lemma}

\noindent
しかし写像 $M \to M \otimes_R (R' \oplus R_f)$ は単射とは限らない。
Example \ref{more-algebra-example-not-glueing-pair} を見よ。

\begin{proof}
$M \neq 0$ だが $M \otimes_R R_f = 0$ ならば、$M$ は $f$-冪ねじれである。
Lemma \ref{more-algebra-lemma-torsion-completion} により
$M \otimes_R R' \cong M \neq 0$ である。最後の主張は Lemma
\ref{more-algebra-lemma-same-quotients} により第一の主張の特別な場合である。
\end{proof}

\begin{lemma}
\label{more-algebra-lemma-cover-spec}
$R$ を環、$f \in R$ とし、$R \to R'$ を同型 $R/fR \to R'/fR'$ を導く環写像と
する。写像 $\Spec(R') \amalg \Spec(R_f) \to \Spec(R)$ は全射である。例えば
写像 $\Spec(R^\wedge) \amalg \Spec(R_f) \to \Spec(R)$ は全射である。
\end{lemma}

\begin{proof}
$\Spec(R) = V(f) \amalg D(f)$ であり、$V(f) = \Spec(R/fR)$ および
$D(f) = \Spec(R_f)$ であることを思い出す。Algebra, Section
\ref{algebra-section-spectrum-ring}、特に Lemmas
\ref{algebra-lemma-spec-closed} and \ref{algebra-lemma-standard-open} を見よ。
従って、写像 $R \to R/fR$ が $R'$ を経由することから補題が従う。最後の主張は
Lemma \ref{more-algebra-lemma-same-quotients} により第一の主張の特別な場合である。
\end{proof}

\begin{lemma}
\label{more-algebra-lemma-faithful-descent}
\begin{reference}
\cite[Lemme~2(a)]{Beauville-Laszlo} のわずかな一般化である。
\end{reference}
$R$ を環、$f \in R$ とし、$R \to R'$ を、$n > 0$ に対して同型
$R/f^nR \to R'/f^nR'$ を導く環写像とする。$R$-加群 $M$ が有限生成であることと、
($R' \oplus R_f$)-加群 $M \otimes_R (R' \oplus R_f)$ が有限生成であることは
同値である。例えば $M \otimes_R (R^\wedge \oplus R_f)$ が
$R^\wedge \oplus R_f$ 上有限生成ならば、$M$ は有限生成 $R$-加群である。
\end{lemma}

\begin{proof}
``のみならば'' は明らかなので、$M \otimes_R (R' \oplus R_f)$ が有限生成で
あると仮定する。この場合、各生成元を単純テンソルの和として書くことにより、
$M \otimes_R (R' \oplus R_f)$ は $M$ の元からなる有限生成系を持つ。
すなわち、有限自由 $R$-加群から $M$ への射であって、その余核が
$R' \oplus R_f$ とのテンソル積によって零化されるものが存在する。従ってこの
余核に Lemma \ref{more-algebra-lemma-BL-faithful} を適用すると、
% P02-E1093: frozen source line 25211 says "finite generated".
$M$ が有限生成であると結論できる。最後の主張は Lemma
\ref{more-algebra-lemma-same-quotients} により第一の主張の特別な場合である。
\end{proof}

\begin{remark}
\label{more-algebra-remark-not-descent}
$R \to R_f$ は常に平坦であるが、$R$ が Noether でない限り
$R \to R^\wedge$ は一般には平坦でない（Algebra, Lemma
\ref{algebra-lemma-completion-flat} および Examples, Section
\ref{examples-section-nonflat} の議論を見よ）。従って一般には、Descent,
Section \ref{descent-section-descent-modules} で論じた忠実平坦降下を射
$R \to R^\wedge \oplus R_f$ に適用できない。さらに Noether の場合でさえ、
この射に対する降下データの通常の定義は環
$R^\wedge \otimes_R R^\wedge$ に言及するが、本節ではこれを考えることを
避ける。
\end{remark}

\noindent
{\bf 貼り合わせ対。}$R \to R'$ を、$n > 0$ に対して同型
$R/f^nR \to R'/f^nR'$ を導く環写像とする。列
\begin{equation}
\label{more-algebra-equation-BL-cech-re}
0 \to R \to R' \oplus R_f \to R'_f \to 0,
\end{equation}
を考える。ここで右側の写像は二つの標準的な準同型の差である。この列が完全なら
ば、$(R \to R', f)$ を \emph{貼り合わせ対}という。また
$(R \to R^\wedge, f)$ が貼り合わせ対であるとき、$(R, f)$ を
\emph{貼り合わせ対}という。これは Lemma \ref{more-algebra-lemma-same-quotients} により
意味を持つ。従って $(R, f)$ が貼り合わせ対であることと、列
\begin{equation}
\label{more-algebra-equation-BL-cech}
0 \to R \to R^\wedge \oplus R_f \to (R^\wedge)_f \to 0,
\end{equation}
が完全であることは同値である。

\begin{lemma}
\label{more-algebra-lemma-same-f-torsion}
$R$ を環、$f \in R$ とし、$R \to R'$ を、$n > 0$ に対して同型
$R/f^nR \to R'/f^nR'$ を導く環写像とする。列
(\ref{more-algebra-equation-BL-cech-re}) は次を満たす。
\begin{enumerate}
\item 右で完全である。
\item 左で完全であることと $R[f^\infty] \to R'[f^\infty]$ が単射であることは
同値である。
\item 中央で完全であることと $R[f^\infty] \to R'[f^\infty]$ が全射である
ことは同値である。
\end{enumerate}
特に、$(R \to R', f)$ が貼り合わせ対であることと
$R[f^\infty] \to R'[f^\infty]$ が全単射であることは同値である。例えば
$(R, f)$ が貼り合わせ対であることと
$R[f^\infty] \to R^\wedge[f^\infty]$ が全単射であることは同値である。
\end{lemma}

\begin{proof}
$x \in R'_f$ とする。$x' \in R'$ によって $x = x'/f^n$ と書く。
$x'' \in R$ および $y \in R'$ によって $x' = x'' + f^n y$ と書く。このとき
$(y, -x''/f^n)$ が $x$ に写ることが分かる。従って (1) が成り立つ。

\medskip\noindent
(2) は $\Ker(R \to R_f) = R[f^\infty]$ であることから従う。

\medskip\noindent
列が中央で完全ならば、$x \in R'[f^\infty]$ に対して $(x, 0)$ の形の元は
第一の矢印の像に属する。これは $R[f^\infty] \to R'[f^\infty]$ が全射である
ことを意味する。逆に、$R[f^\infty] \to R'[f^\infty]$ が全射であると仮定する。右で零に写る中央の元
$(x, y)$ を取る。ある $y' \in R$ によって $y = y'/f^n$ と書く。このとき
$R'$ において $f^n x - y'$ は $f$ のある冪によって零化される。仮定により、
ある $z \in R[f^\infty]$ に対して $f^nx - y' = z$ と書ける。このとき
$y = y''/f^n$ であり、ここで $y'' = y' + z$ は
$R \to R/f^nR$ の核に属する。従って、ある $y''' \in R$ によって
$y$ を $y'''/1$ と表せる。すると $x - y'''$ は $R'[f^\infty]$ に属する。
従って $x - y''' = z' \in R[f^\infty]$ である。このとき所望のとおり
$(x, y'''/1) = (y''' + z', (y''' + z')/1)$ である。

\medskip\noindent
補題の最後の主張は Lemma \ref{more-algebra-lemma-same-quotients} により最後から二番目の主張の
特別な場合である。
\end{proof}

\begin{remark}
\label{more-algebra-remark-BL-special-case}
$f$ が非零因子であると仮定する。このとき Algebra, Lemma
\ref{algebra-lemma-completion-differ-by-torsion} により、$f$ は $R^\wedge$ に
おいても非零因子である。従って $(R, f)$ は貼り合わせ対である。
\end{remark}

\begin{remark}
\label{more-algebra-remark-noetherian-case}
$R \to R^\wedge$ が平坦ならば、各正整数 $n$ に対して列
$0 \to R[f^n] \to R \to R$ と $R^\wedge$ とのテンソル積を取ると列
$0 \to R[f^n] \otimes_R R^\wedge \to R^\wedge \to R^\wedge$ を得る。
Lemma \ref{more-algebra-lemma-torsion-completion} と合わせると、
$R[f^n] \to R^\wedge[f^n]$ が同型であると結論できる。従って $(R, f)$ は
貼り合わせ対である。特に $R$ が Noether ならばこれは成り立つ。Algebra,
Lemma \ref{algebra-lemma-completion-flat} を見よ。
\end{remark}

\begin{example}
\label{more-algebra-example-not-glueing-pair}
$k$ を体とし、
$$
R = k[f, T_1, T_2, \ldots]/(fT_1, fT_2 - T_1, fT_3 - T_2, \ldots).
$$
と置く。このとき $R^\wedge$ における $T_1$ の像は $f$-可除なので、写像
$R[f^\infty] \to R^\wedge[f^\infty]$ は単射でなく、$(R, f)$ は
貼り合わせ対でない。また
$$
R = k[f, T_1, T_2, \ldots]/(fT_1, f^2T_2, \ldots),
$$
に対しては、元 $T_1 + fT_2 + f^2 T_3 + \ldots$ が像に属さないため、写像
$R[f^\infty] \to R^\wedge[f^\infty]$ は全射でない。特に Remark
\ref{more-algebra-remark-noetherian-case} により、どちらも $R \to R^\wedge$ が平坦でない
例である。
\end{example}

\noindent
{\bf 貼り合わせ可能加群。}
$R \to R'$ を、$n > 0$ に対して同型
$R/f^nR \to R'/f^nR'$ を導く環写像とする。任意の $R$-加群 $M$ に対して、
(\ref{more-algebra-equation-BL-cech-re}) と $M$ とのテンソル積を取り、列
\begin{equation}
\label{more-algebra-equation-BL-cech-mod-re}
0 \to M \to (M \otimes_R R') \oplus (M \otimes_R R_f) \to
M \otimes_R R'_f \to 0
\end{equation}
を得る。$M \otimes_R R_f = M_f$ および
$M \otimes_R R'_f = (M \otimes_R R')_f$ であることに注意する。
この列が完全ならば、$M$ は
\emph{$(R \to R', f)$ に対して貼り合わせ可能}であるという。
$R$ が環で $f \in R$ であるとき、$R$-加群が\emph{貼り合わせ可能}であるとは、
$M$ が $(R \to R^\wedge, f)$ に対して貼り合わせ可能であることをいう。
従って $M$ が貼り合わせ可能であることと、列
\begin{equation}
\label{more-algebra-equation-BL-cech-mod}
0 \to M \to (M \otimes_R R^\wedge) \oplus (M \otimes_R R_f) \to
M \otimes_R (R^\wedge)_f \to 0
\end{equation}
が完全であることは同値である。

\begin{lemma}
\label{more-algebra-lemma-same-f-torsion-module}
$R$ を環、$f \in R$ とし、$R \to R'$ を、$n > 0$ に対して同型
$R/f^nR \to R'/f^nR'$ を導く環写像とする。列
(\ref{more-algebra-equation-BL-cech-mod-re}) は次を満たす。
\begin{enumerate}
\item 右で完全である。
\item 左で完全であることと
$M[f^\infty] \to (M \otimes_R R')[f^\infty]$
が単射であることは同値である。
\item 中央で完全であることと
$M[f^\infty] \to (M \otimes_R R')[f^\infty]$
が全射であることは同値である。
\end{enumerate}
従って $M$ が $(R \to R', f)$ に対して貼り合わせ可能であることと
$M[f^\infty] \to (M \otimes_R R')[f^\infty]$ が全単射であることは同値である。
$(R \to R', f)$ が貼り合わせ対ならば、$M$ が $(R \to R', f)$ に対して
貼り合わせ可能であることと
$M[f^\infty] \to (M \otimes_R R')[f^\infty]$ が単射であることは同値である。
例えば $(R, f)$ が貼り合わせ対ならば、$M$ が貼り合わせ可能であることと
$M[f^\infty] \to (M \otimes_R R^\wedge)[f^\infty]$ が単射であることは
同値である。
\end{lemma}

\begin{proof}
Lemma \ref{more-algebra-lemma-same-f-torsion} の結果を断りなく用いる。関手
$M \otimes_R -$ は右完全なので（Algebra, Lemma
\ref{algebra-lemma-tensor-product-exact}）、(1) を得る。

\medskip\noindent
$M \to M \otimes_R R_f = M_f$ の核は $M[f^\infty]$ である。
従って (2) が従う。

\medskip\noindent
列が中央で完全ならば、$x \in (M \otimes_R R')[f^\infty]$ に対して
$(x, 0)$ の形の元は第一の矢印の像に属する。これは
$M[f^\infty] \to (M \otimes_R R')[f^\infty]$
が全射であることを意味する。逆に、
$M[f^\infty] \to (M \otimes_R R')[f^\infty]$
が全射であると仮定する。右で零に写る中央の元 $(x, y)$ を取る。
ある $y' \in M$ によって $y = y'/f^n$ と書く。このとき
$M \otimes_R R'$ において $f^n x - y'$ は $f$ のある冪によって零化される。
仮定により、ある $z \in M[f^\infty]$ に対して $f^nx - y' = z$ と書ける。
このとき $y = y''/f^n$ であり、ここで $y'' = y' + z$ は
$M \to M/f^nM$ の核に属する。従って、ある $y''' \in M$ によって
$y$ を $y'''/1$ と表せる。すると $x - y'''$ は
$(M \otimes_R R')[f^\infty]$ に属する。従って
$x - y''' = z' \in M[f^\infty]$ である。このとき所望のとおり
$(x, y'''/1) = (y''' + z', (y''' + z')/1)$ である。

\medskip\noindent
$(R \to R', f)$ が貼り合わせ対ならば、
(\ref{more-algebra-equation-BL-cech-mod-re}) は、Algebra, Lemma
\ref{algebra-lemma-tensor-product-exact} により任意の $M$ に対して中央で完全である。
これより補題の最後から二番目の主張を得る。補題の最後の主張は、これと、
$(R, f)$ が貼り合わせ対であることと $(R \to R^\wedge, f)$ が貼り合わせ対で
あることが同値であるという事実から従う。
\end{proof}

\begin{remark}
\label{more-algebra-remark-glueable}
$(R \to R', f)$ を貼り合わせ対とし、$M$ を $R$-加群とする。
以下は、$M$ が $(R \to R', f)$ に対して貼り合わせ可能かどうかを判定するために
用いることのできるいくつかの所見である。
\begin{enumerate}
\item
% P02-E1094: frozen source lines 25427-25435 assume only (R -> R', f) is a glueing pair but assert the injectivity criterion for the unrelated completion pair.
Lemma \ref{more-algebra-lemma-same-f-torsion-module} により、$M$ が
$(R \to R^\wedge, f)$ に対して貼り合わせ可能であることと
$M[f^\infty] \to M \otimes_R R^\wedge$ が単射であることは同値である。
これは $M[f] \to M^\wedge$ が単射であるとき、すなわち
$M[f] \cap \bigcap_{n = 1}^\infty f^n M = 0$ のときに成り立つ。
\item $\text{Tor}_1^R(M, R'_f) = 0$ ならば、$M$ は
$(R \to R', f)$ に対して貼り合わせ可能である
（Algebra, Lemma \ref{algebra-lemma-long-exact-sequence-tor} を用いよ）。
これは $\text{Tor}_1^R(M, R')$ が $f$-冪ねじれであることと同値である。
特に任意の平坦 $R$-加群は $(R \to R', f)$ に対して貼り合わせ可能である。
\item $R \to R'$ が平坦ならば、任意の $R$-加群に対して
$\text{Tor}_1^R(M, R') = 0$ なので、すべての $R$-加群が
$(R \to R', f)$ に対して貼り合わせ可能である。特に $R$ が Noether で
$R' = R^\wedge$ のとき、これは成り立つ。Algebra, Lemma
\ref{algebra-lemma-completion-flat} を見よ。
\end{enumerate}
\end{remark}

\begin{example}[貼り合わせ可能でない加群]
\label{more-algebra-example-not-glueable-module}
\begin{reference}
\cite[\S 4, Remarques]{Beauville-Laszlo}
\end{reference}
$R$ を $\mathbf{R}$ 上の $0$ における $C^\infty$ 関数の芽の環とする。
$f \in R$ を関数 $f(x) = x$ とする。このとき $f$ は $R$ の非零因子なので
$(R, f)$ は貼り合わせ対であり、$R^\wedge \cong \mathbf{R}[[x]]$ である。
% P02-E1095: frozen source line 25458 gives exp(-1/x^2) as a function germ at 0 without specifying its smooth extension value at x = 0.
$\varphi \in R$ を関数 $\varphi(x) = \text{exp}(-1/x^2)$ とする。
すると $\varphi$ の Taylor 級数は零なので
$\varphi \in \Ker(R \to R^\wedge)$ である。$x \neq 0$ に対して
$\varphi(x) \neq 0$ なので、$\varphi$ は $R$ の非零因子である。
関数 $\varphi/f$ も Taylor 級数が零なので、
$M = R/\varphi R$ におけるその像は $M[f]$ の非零元であって、
$M \otimes_R R^\wedge = R^\wedge/\varphi R^\wedge = R^\wedge$
において零に写る。従って $M$ は貼り合わせ可能でない。
\end{example}

\noindent
次に Tor 群をいくつか計算する。

\begin{lemma}
\label{more-algebra-lemma-first-tor}
$(R \to R', f)$ を貼り合わせ対とする。このとき各 $n > 0$ に対して
$\text{Tor}^R_1(R', f^n R) = 0$ である。
\end{lemma}

\begin{proof}
完全列 $0 \to R[f^n] \to R \to f^n R \to 0$ から、
$R[f^n] \otimes_R R' \to R'$ が単射であることを確認すれば十分である。
Lemma \ref{more-algebra-lemma-torsion-completion} により
$R[f^n] \otimes_R R' = R[f^n]$ であり、$(R \to R', f)$ は貼り合わせ対なので、
Lemma \ref{more-algebra-lemma-same-f-torsion} により $R[f^n] \to R'$ は単射である。
\end{proof}

\begin{lemma}
\label{more-algebra-lemma-first-tor-total}
$(R \to R',f)$ を貼り合わせ対とする。
このとき $\text{Tor}^R_1(R', R/R[f^\infty]) = 0$ である。
\end{lemma}

\begin{proof}
$R/R[f^\infty] = \colim R/R[f^n] = \colim f^nR$ である。
Tor 群の形成はフィルター付き余極限と可換なので
（Algebra, Lemma \ref{algebra-lemma-tor-commutes-filtered-colimits}）、
Lemma \ref{more-algebra-lemma-first-tor} を適用できる。
\end{proof}

\begin{lemma}
\label{more-algebra-lemma-BL3}
\begin{reference}
\cite[Lemme 3(a)]{Beauville-Laszlo} のわずかな一般化である。
\end{reference}
$(R \to R', f)$ を貼り合わせ対とする。任意の $R$-加群 $M$ に対して
$\text{Tor}^R_1(R', \Coker(M \to M_f)) = 0$ である。
\end{lemma}

\begin{proof}
$\overline{M} = M/M[f^\infty]$ と置く。このとき
$\Coker(M \to M_f) \cong \Coker(\overline{M} \to \overline{M}_f)$ なので、
$f$ が $M$ 上の非零因子であると仮定してよいし、そう仮定する。
この場合 $M \subset M_f$ であり、$M_f/M = \colim M/f^nM$ である。ここで
遷移写像は $f$ による乗法で与えられる。Tor 群の形成は余極限と可換するので
（Algebra, Lemma \ref{algebra-lemma-tor-commutes-filtered-colimits}）、
$\text{Tor}^R_1(R', M/f^n M) = 0$ を示せば十分である。

\medskip\noindent
まず $M = R/R[f^\infty]$ の場合を扱う。
Lemma \ref{more-algebra-lemma-same-f-torsion} により
$M \otimes_R R' = R'/R'[f^\infty]$ である。
短完全列 $0 \to M \to M \to M/f^nM \to 0$ から、
Algebra, Lemma \ref{algebra-lemma-long-exact-sequence-tor} により完全列
$$
\xymatrix{
\text{Tor}_1^R(R', R/R[f^\infty]) \ar[r] &
\text{Tor}_1^R(R', M/f^n M) \ar[r] &
R'/R'[f^\infty] \ar[dll]_{f^n} \\
R'/R'[f^\infty] \ar[r] &
(R'/R'[f^\infty])/(f^n
(R'/R'[f^\infty])) \ar[r] & 0
}
$$
を得る。ここで斜めの矢印は単射である。また第一の群
$\text{Tor}_1^R(R', R/R[f^\infty])$ は
Lemma \ref{more-algebra-lemma-first-tor-total} により零なので、所望のとおり
$\text{Tor}_1^R(R', M/f^nM) = 0$ と結論する。

\medskip\noindent
一般の場合を扱うため、$F$ が自由 $R/R[f^\infty]$-加群であるような全射
$F \to M$ を選び、完全列
$$
0 \to N \to F/f^n F \to M/f^n M \to 0.
$$
を作る。Lemma \ref{more-algebra-lemma-torsion-completion} により、この列は $R'$ と
テンソル積を取っても変わらず、従って完全なままである。前段落により
$\text{Tor}^R_1(R', F/f^n F) = 0$ なので、所望のとおり
$\text{Tor}^R_1(R', M/f^n M) = 0$ と結論する。
\end{proof}

\noindent
$(R \to R', f)$ を貼り合わせ対とする。これは、$n > 0$ に対して
$R/f^nR \to R'/f^nR'$ が同型であり、列
$$
0 \to R \to R' \oplus R_f \to R_f' \to 0
$$
が完全であることを意味する。Remark \ref{more-algebra-remark-glueing-data} で導入した圏
$\text{Glue}(R \to R', f)$ を考える。
$\text{Glue}(R \to R', f)$ の対象 $(M', M_1, \alpha_1)$ を
\emph{貼り合わせデータ}と呼ぶ。これは $R'$-加群 $M'$、$R_f$-加群 $M_1$、
および同型 $\alpha_1 : (M')_f \to M_1 \otimes_R R'$ からなる。
明らかな関手
$$
\text{Can} : \text{Mod}_R \longrightarrow \text{Glue}(R \to R', f),\quad
M \longmapsto (M \otimes_R R', M_f, \text{can}),
$$
があり、逆方向には関手
$$
H^0 : \text{Glue}(R \to R', f) \longrightarrow \text{Mod}_R,\quad
(M', M_1, \alpha_1) \longmapsto \Ker(M' \oplus M_1 \to (M')_f)
$$
がある。正確な定義については Remark \ref{more-algebra-remark-glueing-data} を見よ。

\begin{theorem}
\label{more-algebra-theorem-BL-glueing}
\begin{reference}
\cite{Beauville-Laszlo} の主定理のわずかな一般化である。
\end{reference}
$(R \to R',f)$ を貼り合わせ対とする。関手
$\text{Can} : \text{Mod}_R \longrightarrow \text{Glue}(R \to R', f)$
は、$(R \to R', f)$ に対して貼り合わせ可能な $R$-加群の圏と、
貼り合わせデータの圏 $\text{Glue}(R \to R', f)$ との圏同値を定める。
\end{theorem}

\begin{proof}
$(M', M_1, \alpha_1)$ を貼り合わせデータとする。
$M = H^0((M', M_1, \alpha_1))$ が $(R \to R', f)$ に対して貼り合わせ可能な加群
であり、$(M', M_1, \alpha_1) \cong \text{Can}(M)$ であることを示す。

\medskip\noindent
まず、関手 $H^0$ の定義で用いた写像
$\text{d} : M' \oplus M_1 \to (M')_f$ が全射であることを確認する。
$(x, y) \in M' \oplus M_1$ は $(M')_f$ において
$\text{d}(x, y) = x/1 - \alpha_1^{-1}(y \otimes 1)$ に写ることに注意する。
$z \in (M')_f$ ならば、$g_i \in R'$ および $y_i \in M_1$ によって
$\alpha_1(z) = \sum y_i \otimes g_i$ と書ける。ある $y'_i \in M'$ と
$n \geq 0$ によって
$\alpha_1^{-1}(y_i \otimes 1) = y_i'/f^n$ と書く（すべての $i$ に同じ $n$ を
選べる）。$a_i \in R$ および $b_i \in R'$ によって
$g_i = a_i + f^n b_i$ と書く。このとき
$y = \sum a_i y_i \in M_1$ および $x = \sum b_i y'_i \in M'$ とすれば、
所望のとおり $\text{d}(x, -y) = z$ である。

\medskip\noindent
$M = H^0((M', M_1, \alpha_1)) = \Ker(\text{d})$ なので、$R$-加群の完全列
\begin{equation}
\label{more-algebra-equation-define-M}
0 \to M \to M' \oplus M_1 \to (M')_f \to 0.
\end{equation}
を得る。写像 $M \to M'$ と $M \to M_1$ が同型
$M \otimes_R R' \to M'$ および $M \otimes_R R_f \to M_1$ を導くことを示す。
これより $M$ は $(R \to R', f)$ に対して貼り合わせ可能であり、
所望のとおり $\text{Can}(M) \cong (M', M_1, \alpha_1)$ であることが従う。

\medskip\noindent
$f$ は $M_1$ 上の非零因子なので、
$M[f^\infty] \cong M'[f^\infty]$ である。これより完全列
\begin{equation}
\label{more-algebra-equation-exact-mod-torsion}
0 \to M/M[f^\infty] \to M_1 \to (M')_f/M' \to 0.
\end{equation}
を得る。$R \to R_f$ は平坦なので、この完全列と $R_f$ とのテンソル積を取ると、
$M \otimes_R R_f = (M/M[f^\infty]) \otimes_R R_f \to M_1$
が同型であることが従う。

\medskip\noindent
Lemma \ref{more-algebra-lemma-BL3} により
$\text{Tor}_1^R(R', \Coker(M' \to (M')_f)) = 0$ である。
従って列 (\ref{more-algebra-equation-exact-mod-torsion}) は $R$ 上で $R'$ とテンソル積を取っても
完全なままである。$\alpha_1$ と Lemma
\ref{more-algebra-lemma-torsion-completion} を用いると、得られた完全列は
\begin{equation}
\label{more-algebra-equation-mod-torsion-sequence}
0 \to (M/M[f^\infty]) \otimes_R R' \to
(M')_f \to (M')_f/M' \to 0
\end{equation}
と書ける。これより同型
$(M/M[f^\infty]) \otimes_R R' \cong M'/M'[f^\infty]$ を得る。
従って図式
$$
\xymatrix{
& M[f^\infty] \otimes_R R' \ar[r] \ar[d] &
M \otimes_R R'  \ar[r] \ar[d] &
(M/M[f^\infty]) \otimes_R R' \ar[r] \ar[d] & 0 \\
0 \ar[r] &
M'[f^\infty] \ar[r] &
M' \ar[r] &
M'/M'[f^\infty] \ar[r] & 0,
}
$$
における第三の垂直矢印は同型である。行は完全であり、第一の垂直矢印は
Lemma \ref{more-algebra-lemma-torsion-completion} および
$M[f^\infty] = M'[f^\infty]$ により同型なので、五項補題から
$M \otimes_R R' \to M'$ は同型である。

\medskip\noindent
以上により $\text{Can}$ は本質的全射であり、関手 $H^0$ は貼り合わせ可能加群の
圏に値を取る。貼り合わせ可能加群に対する
(\ref{more-algebra-equation-BL-cech-mod-re}) の完全性により、貼り合わせ可能加群の圏上で
$H^0 \circ \text{Can} = \text{id}$ である。Lemma
\ref{more-algebra-lemma-H0-Can-adjoint} と Categories, Lemma
\ref{categories-lemma-adjoint-fully-faithful} を合わせると、
これより $\text{Can}$ は充満忠実である。以上で証明が完了する。
\end{proof}

\begin{remark}
\label{more-algebra-remark-what-you-get-for-general-modules}
$(R \to R', f)$ を貼り合わせ対とする。$M$ を、$(R \to R', f)$ に対して
貼り合わせ可能とは限らない $R$-加群とする。
$M' = M \otimes_R R'$ および $M_1 = M_f$ と置くと、貼り合わせデータ
$\text{Can}(M) = (M', M_1, \text{can})$ を得る。このとき
$\tilde M = H^0(M', M_1, \text{can})$ は $(R \to R', f)$ に対して
貼り合わせ可能な $R$-加群であり、標準写像 $M \to \tilde M$ は同型
$M \otimes_R R' \to \tilde M \otimes_R R'$ および
$M_f \to \tilde M_f$ を導く。Theorem \ref{more-algebra-theorem-BL-glueing} を見よ。
列
$$
M \to (M \otimes_R R' )\oplus M_f \to M \otimes_R (R')_f  \to 0
$$
および
$$
0 \to \tilde M \to (\tilde M \otimes_R R') \oplus \tilde M_f \to
\tilde M \otimes_R  (R')_f \to 0
$$
の完全性から、写像 $M \to \tilde M$ は全射であると結論する。
\end{remark}

\noindent
貼り合わせ対 $(R \to R', f)$ 上の平坦 $R$-加群は貼り合わせ可能であることを
思い出す（Remark \ref{more-algebra-remark-glueable}）。従って次の補題により、Theorem
\ref{more-algebra-theorem-BL-glueing} は平坦 $R$-加群の圏と、$M'$ および $M_1$ がそれぞれ
$R'$ および $R_f$ 上平坦である貼り合わせデータ $(M', M_1, \alpha_1)$ の圏との
圏同値を定める。

\begin{lemma}
\label{more-algebra-lemma-BL-flat}
$(R \to R', f)$ を貼り合わせ対とする。$M$ を、$(R \to R', f)$ に対して
貼り合わせ可能とは限らない $R$-加群とする。このとき $M$ が $R$ 上平坦である
ことと、$M \otimes_R R'$ が $R'$ 上平坦かつ $M_f$ が $R_f$ 上平坦であることは
同値である。
\end{lemma}

\begin{proof}
一方向は Algebra, Lemma \ref{algebra-lemma-flat-base-change} から従う。
逆方向について、$M \otimes_R R'$ が $R'$ 上平坦で、$M_f$ が $R_f$ 上平坦で
あると仮定する。$\tilde M$ を Remark
\ref{more-algebra-remark-what-you-get-for-general-modules} のものとする。
$\tilde M$ が $R$ 上平坦ならば、短完全列
$0 \to \Ker(M \to \tilde M) \to M \to \tilde M \to 0$
に Algebra, Lemma \ref{algebra-lemma-flat-tor-zero} を適用すると、
$\Ker(M \to \tilde M) \otimes_R (R' \oplus R_f)$ が零であることが分かる。
従って Lemma \ref{more-algebra-lemma-BL-faithful} により $M = \tilde M$ であり、結論を得る。
換言すれば $M$ を $\tilde M$ で置き換え、$M$ は $(R \to R', f)$ に対して
貼り合わせ可能であると仮定してよい。$N$ をもう一つの $R$-加群とする。
Algebra, Lemma \ref{algebra-lemma-characterize-flat} により、
$\text{Tor}_1^R(M, N) = 0$ を示せば十分である。

\medskip\noindent
短完全列 $0 \to R \to R' \oplus R_f \to (R')_f \to 0$ と $N$ に付随する
Tor の長完全列は完全列
$$
0 \to \text{Tor}_1^R(R', N) \to \text{Tor}_1^R((R')_f, N)
$$
および $i \geq 2$ に対する同型
$\text{Tor}_i^R(R', N) = \text{Tor}_i^R((R')_f, N)$
を与える。
$\text{Tor}_i^R((R')_f, N) = \text{Tor}_i^R(R', N)_f$ なので、
$f$ は $\text{Tor}_1^R(R', N)$ 上の非零因子であり、$i \geq 2$ に対して
$\text{Tor}_i^R(R', N)$ 上可逆であると結論する。
$M \otimes_R R'$ は $R'$ 上平坦なので、Example
\ref{more-algebra-example-tor-change-rings} のスペクトル系列により
$$
\text{Tor}_i^R(M \otimes_R R', N) =
(M \otimes_R R') \otimes_{R'} \text{Tor}_i^R(R', N)
$$
である。$M \otimes_R R'$ を有限自由 $R'$-加群のフィルター付き余極限として
書くと（Algebra, Theorem \ref{algebra-theorem-lazard}）、
$f$ は $\text{Tor}_1^R(M \otimes_R R', N)$ 上の非零因子であり、
% P02-E1096: frozen source lines 25753-25755 omit the already required qualifier i >= 2 from the second Tor conclusion.
$\text{Tor}_i^R(M \otimes_R R', N)$ 上可逆であると結論する。
次に、$M$ が貼り合わせ可能であることから来る完全列
$0 \to M \to M \otimes_R R' \oplus M_f \to M \otimes_R (R')_f \to 0$
と、それに付随する $\text{Tor}$ の長完全列を考える。関連する部分は
$$
\xymatrix{
\text{Tor}_1^R(M, N) \ar[r] &
\text{Tor}_1^R(M \otimes_R R', N) \ar[r] &
\text{Tor}_1^R(M \otimes_R (R')_f, N) \\
& \text{Tor}_2^R(M \otimes_R R', N) \ar[r] &
\text{Tor}_2^R(M \otimes_R (R')_f, N) \ar[llu]
}
$$
である。上で述べた、$\text{Tor}_i^R(M \otimes_R R', N)$ 上の $f$ の作用により、
% P02-E1097: frozen source lines 25769-25770 say “the action on f on” instead of “the action of f on”.
$\text{Tor}_1^R(M, N) = 0$ と結論する。
\end{proof}

\noindent
Lemma \ref{more-algebra-lemma-faithful-descent} では、次の補題の結果を「有限生成」について
既に見たことに注意する。

\begin{lemma}
\label{more-algebra-lemma-BL-properties}
$(R \to R', f)$ を貼り合わせ対とする。$M$ を、$(R \to R', f)$ に対して
貼り合わせ可能とは限らない $R$-加群とする。このとき $M$ が有限射影
$R$-加群であることと、$M \otimes_R R'$ が $R'$ 上有限射影であり、
$M_f$ が $R_f$ 上有限射影であることは同値である。
\end{lemma}

\begin{proof}
$M \otimes_R R'$ が $R'$ 上有限射影加群であり、$M_f$ が $R_f$ 上有限射影加群
であると仮定する。$M$ が $R$ 上有限射影であることを示すのが目標である。
Algebra, Lemma \ref{algebra-lemma-finite-projective} を断りなく用いる。
Lemma \ref{more-algebra-lemma-BL-flat} により $M$ は平坦である。
Lemma \ref{more-algebra-lemma-faithful-descent} により $M$ は有限生成である。
短完全列 $0 \to K \to R^{\oplus n} \to M \to 0$ を選ぶ。
有限射影加群は有限表示であり、$R'$（Algebra, Lemma
\ref{algebra-lemma-flat-tor-zero} による）および $R_f$ とテンソル積を取った
後も列は完全なままなので、$K \otimes_R R'$ と $K_f$ は有限生成加群である。
上の補題を用いて $K$ は有限生成であると結論する。従って $M$ は有限表示であり、
従って有限射影である。
\end{proof}

\begin{remark}
\label{more-algebra-remark-compare-BL}
\cite{Beauville-Laszlo} では $f$ が $R$ の非零因子であり
$R' = R^\wedge$ であることを仮定する。これは Lemma
\ref{more-algebra-lemma-same-f-torsion} により貼り合わせ対を与える。
この設定でさえ Theorem \ref{more-algebra-theorem-BL-glueing} は新しいことを述べる。
\cite{Beauville-Laszlo} の結果は、$f$ が非零因子として作用する加群にしか
適用されない（従ってそのような加群は我々の意味で貼り合わせ可能である。
Lemma \ref{more-algebra-lemma-same-f-torsion-module} を見よ）。
Lemma \ref{more-algebra-lemma-BL-properties} は \cite{Beauville-Laszlo} の結果をさらに
わずかに拡張する。すなわち、$M$ が非零の $f$-冪ねじれを持つことを許せるだけで
なく、それが貼り合わせ可能であることさえ要求しない。
\end{remark}

\section{導来完備化}
\label{more-algebra-section-derived-completion}

\noindent
本節の内容の参考文献として
\cite{Dwyer-Greenlees}、\cite{Greenlees-May}、\cite{PSY}、\cite{dag12}
（特に Chapter 4）がある。本節の解説は \cite{BS} に従う。
ねじれ加群に対する本節の類似物（または「双対」）は Dualizing Complexes,
Section \ref{dualizing-section-local-cohomology} である。
ねじれコホモロジーを持つ複体の導来圏と導来完備複体との関係については
Dualizing Complexes, Section
\ref{dualizing-section-torsion-and-complete} を見よ。

\medskip\noindent
$K \in D(A)$ とし、$f \in A$ とする。系
% P02-E1098: frozen source line 25842 omits "by" in the Denote construction.
$$
\ldots \to K \xrightarrow{f} K \xrightarrow{f} K
$$
の $D(A)$ における導来極限を $T(K, f)$ と表す。

\begin{lemma}
\label{more-algebra-lemma-hom-from-Af}
$A$ を環、$f \in A$、$K \in D(A)$ とする。次は同値である。
\begin{enumerate}
\item すべての $n$ に対して $\Ext^n_A(A_f, K) = 0$ である。
\item $D(A_f)$ のすべての $E$ に対して $\Hom_{D(A)}(E, K) = 0$ である。
\item $T(K, f) = 0$ である。
\item すべての $p \in \mathbf{Z}$ に対して $T(H^p(K), f) = 0$ である。
\item すべての $p \in \mathbf{Z}$ に対して
$\Hom_A(A_f, H^p(K)) = 0$ かつ $\Ext^1_A(A_f, H^p(K)) = 0$ である。
\item $R\Hom_A(A_f, K) = 0$ である。
\item 写像 $\prod_{n \geq 0} K \to \prod_{n \geq 0} K$,
$(x_0, x_1, \ldots) \mapsto (x_0 - fx_1, x_1 - fx_2, \ldots)$
は $D(A)$ における同型である。
\item
% P02-E1099: frozen source line 25864 contains the unfinished placeholder "add more here."
ここにさらに加える。
\end{enumerate}
\end{lemma}

\begin{proof}
(2) が (1) を導くこと、および (1) と (6) が同値であることは明らかである。
(1) を仮定する。$I^\bullet$ を、$K$ を表す $A$-加群の K-単射的複体とする。
条件 (1) は $\Hom_A(A_f, I^\bullet)$ が非輪状であることを意味する。
$M^\bullet$ を、$E$ を表す $A_f$-加群の複体とする。このとき Algebra, Lemma
\ref{algebra-lemma-adjoint-hom-restrict} により
$$
\Hom_{D(A)}(E, K) =
\Hom_{K(A)}(M^\bullet, I^\bullet) =
\Hom_{K(A_f)}(M^\bullet, \Hom_A(A_f, I^\bullet))
$$
である。Lemma \ref{more-algebra-lemma-hom-K-injective} により
$\Hom_A(A_f, I^\bullet)$ は $A_f$-加群の K-単射的複体なので、非輪状であること
から零とホモトピー同値であることが従う（Derived Categories, Lemma
\ref{derived-lemma-K-injective}）。従って (2) を得る。

\medskip\noindent
$A$-加群 $A_f$ の自由分解は
$$
0 \to \bigoplus\nolimits_{n \in \mathbf{N}} A \to
\bigoplus\nolimits_{n \in \mathbf{N}} A
\to A_f \to 0
$$
で与えられる。ここで第一の写像は
$(a_0, a_1, a_2, \ldots)$ を
$(a_0, a_1 - fa_0, a_2 - fa_1, \ldots)$ に写し、第二の写像は
$(a_0, a_1, a_2, \ldots)$ を $a_0 + a_1/f + a_2/f^2 + \ldots$ に写す。
$\Hom_A(-, I^\bullet)$ を適用すると
$$
0 \to \Hom_A(A_f, I^\bullet) \to \prod I^\bullet \to \prod I^\bullet \to 0
$$
を得る。$\prod I^\bullet$ は $\prod_{n \geq 0} K$ を表すので、これにより
(1) と (7) の同値性が証明される。一方、Derived Categories, Section
\ref{derived-section-derived-limit} における導来極限の構成により、上の完全列は
$T(K, f)$ が $D(A)$ における $R\Hom_A(A_f, K)$ の代表であることを示す。
% P02-E1100: frozen source lines 25904-25906 say "Thus the equivalence..." without a finite verb.
従って (1) と (3) は同値である。

\medskip\noindent
スペクトル系列
$$
E_2^{p, q} = \Ext^p_A(A_f, H^q(K)) \Rightarrow
\Ext^{p + q}_A(A_f, K)
$$
がある。Equation (\ref{more-algebra-equation-first-ss-ext})\footnote{$K$ が下に有界でない
ならば、示された参照を用いることはできない。その場合は、$K$ を表す複体
% P02-E1101: frozen source line 25916 says "we chooses".
$K^\bullet$ と自由二項分解
$0 \to F^{-1} \to F^0 \to A_f \to 0$ を選ぶ（上を見よ）。
すると項 $\Hom_A(F^{-q}, K^p)$ を持つ二重複体を得て、その全複体は
$R\Hom(A_f, K)$ を計算する。この二重複体に付随する第二スペクトル系列
（Homology, Section \ref{homology-section-double-complex}）が所望のものを与える。
詳細は省略する。} を見よ。このスペクトル系列は $E_2$ で退化する。実際、
$A_f$ は射影 $A$-加群による長さ $1$ の分解を持つので（上を見よ）、
$E_2$-頁は非零の列を二つしか持たない。従って短完全列
$$
0 \to \Ext^1_A(A_f, H^{p - 1}(K)) \to
\Ext^p_A(A_f, K) \to
\Hom_A(A_f, H^p(K)) \to 0
$$
を得る。これにより (4) と (5) が (1) と同値であることが証明される。
\end{proof}

\begin{lemma}
\label{more-algebra-lemma-ideal-of-elements-complete-wrt}
$A$ を環、$K \in D(A)$ とする。$T(K, f) = 0$ となる $f \in A$ の集合 $I$ は
$A$ の根基イデアルである。
\end{lemma}

\begin{proof}
Lemma \ref{more-algebra-lemma-hom-from-Af} の結果を断りなく用いる。
$f \in I$ かつ $g \in A$ ならば、$A_{gf}$ は $A_f$-加群なので、
すべての $n$ に対して $\Ext^n_A(A_{gf}, K) = 0$ であり、従って $gf \in I$ で
ある。$f, g \in I$ と仮定する。このとき $f, g$ は $A_{f + g}$ において
単位イデアルを生成するので、短完全列
$$
0 \to A_{f + g} \to A_{f(f + g)} \oplus A_{g(f + g)} \to A_{gf(f + g)} \to 0
$$
がある。これは Algebra, Lemma \ref{algebra-lemma-standard-covering} と、
最後の矢印が全射であるという容易な事実から従う。$\Ext$ の長完全列と、
すべての $n$ に対する
$\Ext^n_A(A_{f(f + g)}, K)$、
$\Ext^n_A(A_{g(f + g)}, K)$、および
$\Ext^n_A(A_{gf(f + g)}, K)$ の消滅から、すべての $n$ に対して
$\Ext^n_A(A_{f + g}, K)$ が消滅することを導く。最後に、ある $n > 0$ に対して
$f^n \in I$ ならば、$T(K, f) = T(K, f^n)$ または
$A_f \cong A_{f^n}$ であることから $f \in I$ である。
\end{proof}

\begin{lemma}
\label{more-algebra-lemma-complete-derived-complete}
$A$ を環、$I \subset A$ をイデアル、$M$ を $A$-加群とする。
\begin{enumerate}
\item $M$ が $I$-進完備ならば、すべての $f \in I$ に対して
$T(M, f) = 0$ である。
\item 逆に、すべての $f \in I$ に対して $T(M, f) = 0$ であり、$I$ が有限生成
ならば、$M \to \lim M/I^nM$ は全射である。
\end{enumerate}
\end{lemma}

\begin{proof}
(1) の証明。$M$ が $I$-進完備であると仮定する。
Lemma \ref{more-algebra-lemma-hom-from-Af} により
$\Ext^1_A(A_f, M) = 0$ および $\Hom_A(A_f, M) = 0$ を示せば十分である。
$M = \lim M/I^nM$ かつ $\Hom_A(A_f, M/I^nM) = 0$ なので、
$\Hom_A(A_f, M) = 0$ が従う。拡大
$$
0 \to M \to E \to A_f \to 0
$$
が与えられたとする。$n \geq 0$ に対して、$1/f^n$ に写る $e_n \in E$ を選ぶ。
$n \geq 0$ に対して $\delta_n = fe_{n + 1} - e_n \in M$ と置く。
$e_n$ を
$$
e'_n = e_n + \delta_n + f\delta_{n + 1} + f^2 \delta_{n + 2} + \ldots
$$
で置き換える。$M$ は $I$ に関して完備で $f \in I$ なので、この無限和は存在する。
簡単な計算により $fe'_{n + 1} = e'_n$ である。従って $1/f^n$ を $e'_n$ に写す
ことにより、拡大の分裂を得る。

\medskip\noindent
(2) の証明。$I = (f_1, \ldots, f_r)$ かつ、$i = 1, \ldots, r$ に対して
$T(M, f_i) = 0$ であると仮定する。Algebra, Lemma
\ref{algebra-lemma-when-surjective-to-completion} により、
$I = (f)$ かつ $T(M, f) = 0$ と仮定してよい。
$n \geq 0$ に対して $x_n \in M$ とする。
$$
0 \to M \to E \to A_f \to 0
$$
を、
$$
E = M \oplus \bigoplus Ae_n\Big/\langle x_n - fe_{n + 1} + e_n\rangle
$$
で与えられ、$e_n$ を $A_f$ の $1/f^n$ に写す拡大とする（上を見よ）。
仮定と Lemma \ref{more-algebra-lemma-hom-from-Af} によりこの拡大は分裂する。従って
$E$ において $A_f$ のコピーを生成する元 $x + e_0$ を得る。このとき
$$
x + e_0 = x - x_0 + fe_1 = x - x_0 - f x_1 + f^2 e_2 = \ldots
$$
である。蛇の補題により $M/f^nM = E/f^nE$ なので、
$f^nM$ を法として
$x = x_0 + fx_1 + \ldots + f^{n - 1}x_{n - 1}$ であることが分かる。
換言すれば、所望のとおり写像 $M \to \lim M/f^nM$ は全射である。
\end{proof}

\noindent
以上の結果に動機づけられ、次の定義を置く。

\begin{definition}
\label{more-algebra-definition-derived-complete}
$A$ を環、$K \in D(A)$、$I \subset A$ をイデアルとする。すべての $f \in I$ に
対して $T(K, f) = 0$ であるとき、$K$ は
{\it $I$ に関して導来完備}であるという。
$M$ が $A$-加群であるとき、$M[0] \in D(A)$ が $I$ に関して導来完備ならば、
$M$ は {\it $I$ に関して導来完備}であるという。
\end{definition}

\noindent
導来完備な対象からなる充満部分圏
$D_{comp}(A) = D_{comp}(A, I) \subset D(A)$ は、真に充満な飽和三角部分圏である。
Derived Categories, Definitions
\ref{derived-definition-triangulated-subcategory} および
\ref{derived-definition-saturated} を見よ。
Lemma \ref{more-algebra-lemma-ideal-of-elements-complete-wrt} により、部分圏
$D_{comp}(A, I)$ は $I$ の根基 $\sqrt{I}$ のみに依存する。換言すれば、
$\Spec(A)$ の閉部分集合 $Z = V(I)$ のみに依存する。
部分圏 $D_{comp}(A, I)$ は $D(A)$ における積とホモトピー極限で保たれる。
しかし一般には可算直和で保たれない。イデアル $I$ の選択が文脈から明らかなときは、
$M$ を単に導来完備加群と呼ぶことが多い。

\begin{proposition}
\label{more-algebra-proposition-derived-complete-modules}
$I \subset A$ を環 $A$ の有限生成イデアルとし、$M$ を $A$-加群とする。
次は同値である。
\begin{enumerate}
\item $M$ は $I$-進完備である。
\item $M$ は $I$ に関して導来完備であり、$\bigcap I^nM = 0$ である。
\end{enumerate}
\end{proposition}

\begin{proof}
Lemma \ref{more-algebra-lemma-complete-derived-complete} の結果から明らかである。
\end{proof}

\noindent
次の補題は、導来完備加群の圏 $\mathcal{C}$ が Abel 圏であることを示す。
$\mathcal{C}$ は Grothendieck Abel 圏ではない。Examples, Section
\ref{examples-section-derived-complete-modules} を見よ。

\begin{lemma}
\label{more-algebra-lemma-serre-subcategory}
$I$ を環 $A$ のイデアルとする。
\begin{enumerate}
\item 導来完備 $A$-加群は $\text{Mod}_A$ の弱 Serre 部分圏
$\mathcal{C}$ をなす。
\item $D_\mathcal{C}(A) \subset D(A)$ は導来完備対象の充満部分圏である。
\end{enumerate}
\end{lemma}

\begin{proof}
(2) は Lemma \ref{more-algebra-lemma-hom-from-Af} と定義から直ちに従う。
(1) について、$M \to N$ を導来完備加群の写像とする。
% P02-E1102: frozen source lines 26080-26082 omit "by" in the Denote construction.
対応する $D(A)$ の対象を $K = (M \to N)$ と表す。
$f \in I$ を選ぶ。すべての $n$ に対して
$\Ext_A^n(A_f, M)$ と $\Ext_A^n(A_f, N)$ が零なので、
$\Ext_A^n(A_f, K)$ もすべての $n$ に対して零である。従って $K$ は導来完備で
ある。(2) により $\Ker(M \to N)$ と $\Coker(M \to N)$ は
$\mathcal{C}$ の対象である。最後に
$0 \to M_1 \to M_2 \to M_3 \to 0$
を $A$-加群の短完全列とし、$M_1$, $M_3$ が導来完備であると仮定する。
$\Ext$ の長完全列から $M_2$ が導来完備であることが従う。
従って Homology, Lemma
\ref{homology-lemma-characterize-weak-serre-subcategory} により
$\mathcal{C}$ は弱 Serre 部分圏である。
\end{proof}

\noindent
次の補題を Lemma \ref{more-algebra-lemma-derived-complete-zero-bis} で一般化する。

\begin{lemma}
\label{more-algebra-lemma-derived-complete-zero}
$I$ を環 $A$ の有限生成イデアルとし、$M$ を導来完備 $A$-加群とする。
$M/IM = 0$ ならば $M = 0$ である。
\end{lemma}

\begin{proof}
$M/IM$ が零であると仮定する。$I = (f_1, \ldots, f_r)$ とする。
$N = M/(f_1, \ldots, f_i)M$ が非零となる最大の整数 $i < r$ を取る。
そのような $i$ が存在しなければ $M = 0$ であり、これが示すべきことである。
$N$ は導来完備加群の間の写像の余核なので、Lemma
\ref{more-algebra-lemma-serre-subcategory} により導来完備である。$i$ の選び方から
$f_{i + 1} : N \to N$ は全射である。従って
$$
\lim (\ldots \to N \xrightarrow{f_{i + 1}} N \xrightarrow{f_{i + 1}} N)
$$
は非零であり、$N$ の導来完備性に矛盾する。
\end{proof}

\noindent
環が $I$-進完備ならば、導来完備複体が豊富に得られる。

\begin{lemma}
\label{more-algebra-lemma-pseudo-coherent-is-derived-complete}
$A$ を環、$I \subset A$ をイデアルとする。$A$ が導来完備ならば
（例えば $I$-進完備ならば）、$D(A)$ の任意の擬コヒーレント対象は導来完備である。
\end{lemma}

\begin{proof}
（Lemma \ref{more-algebra-lemma-complete-derived-complete} は補題の括弧内の主張を説明する。）
$K$ を $D(A)$ の擬コヒーレント対象とする。定義により、これは $K$ が有限自由
$A$-加群の上に有界な複体 $K^\bullet$ によって表されることを意味する。
$A$ は導来完備なので、Lemma \ref{more-algebra-lemma-serre-subcategory} の (1) により、
すべての $n$ に対して $H^n(K)$ は導来完備である。さらに同じ補題の (2) により、
$K$ は導来完備である。
\end{proof}

\begin{lemma}
\label{more-algebra-lemma-double-localize}
$A$ を環、$f, g \in A$ とする。このとき $K \in D(A)$ に対して
$$
R\Hom_A(A_f, R\Hom_A(A_g, K)) = R\Hom_A(A_{fg}, K)
$$
である。
\end{lemma}

\begin{proof}
Lemma \ref{more-algebra-lemma-internal-hom} から従う。
\end{proof}

\begin{lemma}
\label{more-algebra-lemma-derived-completion}
\begin{slogan}
有限生成イデアルに沿う導来完備化は存在し、{\v C}ech 手続きによって計算できる。
\end{slogan}
$I$ を環 $A$ の有限生成イデアルとする。包含関手
$D_{comp}(A, I) \to D(A)$ は左随伴を持つ。すなわち、$D(A)$ の任意の対象 $K$ に
対して、$K$ を $D(A)$ の導来完備対象へ写す写像 $K \to K^\wedge$ であって、
$E$ が $D(A)$ の導来完備対象であるとき常に写像
$$
\Hom_{D(A)}(K^\wedge, E) \longrightarrow \Hom_{D(A)}(K, E)
$$
が全単射となるものが存在する。実際、$I$ が $f_1, \ldots, f_r \in A$ によって
生成されるならば、$K$ に関して関手的に
$$
K^\wedge = R\Hom\left((A \to \prod\nolimits_{i_0} A_{f_{i_0}} \to
\prod\nolimits_{i_0 < i_1} A_{f_{i_0}f_{i_1}}
\to \ldots \to A_{f_1\ldots f_r}), K\right)
$$
である。
\end{lemma}

\begin{proof}
$K^\wedge$ を補題の最後の表示式によって定義する。複体の写像
$$
(A \to \prod\nolimits_{i_0} A_{f_{i_0}} \to
\prod\nolimits_{i_0 < i_1} A_{f_{i_0}f_{i_1}} \to
\ldots \to A_{f_1\ldots f_r}) \longrightarrow A
$$
があり、これは写像 $K \to K^\wedge$ を導く。$K^\wedge$ が導来完備であり、
$K$ が導来完備ならば $K \to K^\wedge$ が同型であることを示せば十分である
\footnote{実際、$E$ が導来完備で $a : K \to E$ が写像ならば、可換図式
$$
\xymatrix{
K \ar[r]_a \ar[d] & E \ar[d]^{\cong} \\
K^\wedge \ar[r]^{a^\wedge} & E^\wedge
}
$$
は $E$ への任意の写像が $K^\wedge$ を経由することを示す。
区別三角 $K \to K^\wedge \to C$ を選び、${}^\wedge$ を適用する。
$C^\wedge = 0$ を得る。上で述べたことにより、これは導来完備な $E$ に対して
$\Hom(C, E) = 0$ であることを意味し、所望のとおり
$\Hom(K, E) = \Hom(K^\wedge, E)$ が証明される。}。

\medskip\noindent
$f \in A$ とする。Lemma \ref{more-algebra-lemma-double-localize} により、対象
$R\Hom_A(A_f, K^\wedge)$ は
$$
R\Hom\left((A_f \to \prod\nolimits_{i_0} A_{ff_{i_0}} \to
\prod\nolimits_{i_0 < i_1} A_{ff_{i_0}f_{i_1}} \to
\ldots \to A_{ff_1\ldots f_r}), K\right)
$$
に等しい。$f \in I$ ならば $f_1, \ldots, f_r$ は $A_f$ において単位イデアルを
生成するので、増広交代 {\v C}ech 複体
$$
A_f \to \prod\nolimits_{i_0} A_{ff_{i_0}} \to
\prod\nolimits_{i_0 < i_1} A_{ff_{i_0}f_{i_1}} \to
\ldots \to A_{ff_1\ldots f_r}
$$
は Lemma \ref{more-algebra-lemma-extended-alternating-torsion} により $D(A)$ において零である。
（実際、ある $i$ に対して $f = f_i$ ならば、Lemma
\ref{more-algebra-lemma-extended-alternating-homotopy-zero} によりこの複体は零とホモトピー
同値である。必要なのはこの場合だけである。）従って
$R\Hom_A(A_f, K^\wedge) = 0$ であり、Lemma
\ref{more-algebra-lemma-hom-from-Af} により $K^\wedge$ は導来完備であると結論する。

\medskip\noindent
逆に $K$ が導来完備ならば、$f = f_{i_0} \ldots f_{i_p}$,
$p \geq 0$ であるものすべてに対して $R\Hom_A(A_f, K)$ は零である。
従って $K \to K^\wedge$ は $D(A)$ における同型である。
\end{proof}

\begin{remark}
\label{more-algebra-remark-derived-completion}
$A$ を環、$I \subset A$ を有限生成イデアルとする。Lemma
\ref{more-algebra-lemma-derived-completion} により存在する包含関手
$D_{comp}(A, I) \to D(A)$ の左随伴を {\it 導来完備化}という。
これを示すため「$K^\wedge$ を $K$ の導来完備化とする」という。
随伴の単位が関手的な写像 $K \to K^\wedge$ であることに留意せよ。
\end{remark}

\begin{lemma}
\label{more-algebra-lemma-derived-completion-vanishes}
$A$ を環、$I \subset A$ を有限生成イデアルとする。$K^\bullet$ を $A$-加群の複体
とし、ある $f \in I$ に対して
$f : K^\bullet \to K^\bullet$ が同型、すなわち $K^\bullet$ が
$A_f$-加群の複体であるとする。このとき $K^\bullet$ の導来完備化は零である。
\end{lemma}

\begin{proof}
実際、この場合、任意の導来完備複体 $L$ に対して
% P02-E1103: frozen source line 26253 has the ungrammatical article in "the RHom_A(K, L)".
$R\Hom_A(K, L)$ は零である。Lemma \ref{more-algebra-lemma-hom-from-Af} を見よ。
従って Lemma \ref{more-algebra-lemma-derived-completion} の普遍性により $K^\wedge$ は零である。
\end{proof}

\begin{lemma}
\label{more-algebra-lemma-completion-RHom}
$A$ を環、$I \subset A$ を有限生成イデアルとし、$K, L \in D(A)$ とする。
このとき
$$
R\Hom_A(K, L)^\wedge = R\Hom_A(K, L^\wedge) = R\Hom_A(K^\wedge, L^\wedge)
$$
である。
\end{lemma}

\begin{proof}
Lemma \ref{more-algebra-lemma-derived-completion} により、ある $C \in D(A)$ に対して導来完備化が
$R\Hom_A(C, -)$ で与えられることを知っている。このとき Lemma
\ref{more-algebra-lemma-internal-hom} により
\begin{align*}
R\Hom_A(C, R\Hom_A(K, L))
& =
R\Hom_A(C \otimes_A^\mathbf{L} K, L) \\
& =
R\Hom_A(K, R\Hom_A(C, L))
\end{align*}
である。これにより第一の等式が証明される。写像 $K \to K^\wedge$ は写像
$$
R\Hom_A(K^\wedge, L^\wedge) \to R\Hom_A(K, L^\wedge)
$$
を導く。包含関手の左随伴としての導来完備化の定義により、これは
$D(A)$ における同型である。
\end{proof}

\begin{lemma}
\label{more-algebra-lemma-naive-derived-completion}
$A$ を環、$I \subset A$ をイデアルとする。$(K_n)$ を $D(A)$ の対象の逆系とし、
すべての $f \in I$ と $n$ に対して、$K_n$ 上で $f^e$ が零となるような
$e = e(n, f)$ が存在するとする。このとき $K \in D(A)$ に対して、対象
$K' = R\lim (K \otimes_A^\mathbf{L} K_n)$ は $I$ に関して導来完備である。
\end{lemma}

\begin{proof}
導来完備対象の圏は $R\lim$ で保たれるので、各
$K \otimes_A^\mathbf{L} K_n$ が導来完備であることを示せば十分である。
仮定により、すべての $f \in I$ に対して、$K \otimes_A^\mathbf{L} K_n$ 上で
$f^e$ が零となるような $e$ がある。これはもちろん
$T(K \otimes_A^\mathbf{L} K_n, f) = 0$ を意味し、結論を得る。
\end{proof}

\begin{situation}
\label{more-algebra-situation-koszul}
$A$ を環、$I = (f_1, \ldots, f_r) \subset A$ とする。
$K_n^\bullet = K_\bullet(A, f_1^n, \ldots, f_r^n)$ を
$f_1^n, \ldots, f_r^n$ 上の Koszul 複体とし、次数
$-r, -r + 1, \ldots, 0$ の余鎖複体と見る。Lemma
\ref{more-algebra-lemma-functorial} の関手性を用いると、逆系
$$
\ldots \to K_3^\bullet \to K_2^\bullet \to K_1^\bullet
$$
を得る。これは逆系
$H^0(K_n^\bullet) = A/(f_1^n, \ldots, f_r^n)$ および写像
$A \to K_n^\bullet$ と両立する。
\end{situation}

\noindent
以下の議論の要点として、$m > n$ のとき写像
$$
K_m^{-p} = \wedge^p(A^{\oplus r}) \to \wedge^p(A^{\oplus r}) = K_n^{-p}
$$
は基底元 $e_{i_1} \wedge \ldots \wedge e_{i_p}$ 上で
$f_{i_1}^{m - n} \ldots f_{i_p}^{m - n}$ による乗法で与えられる。

\begin{lemma}
\label{more-algebra-lemma-koszul-derived-completion-complete}
Situation \ref{more-algebra-situation-koszul} において、$K \in D(A)$ に対する対象
$K' = R\lim (K \otimes_A^\mathbf{L} K_n^\bullet)$ は
$I$ に関して導来完備である。
\end{lemma}

\begin{proof}
これは Lemma \ref{more-algebra-lemma-naive-derived-completion} の特別な場合である。実際、
$f_i^n$ は $K_n^\bullet$ 上、Lemma \ref{more-algebra-lemma-homotopy-koszul} により零と
ホモトピー同値な自己準同型として作用する。
\end{proof}

\begin{lemma}
\label{more-algebra-lemma-characterize-derived-complete-Koszul}
Situation \ref{more-algebra-situation-koszul} において、$K \in D(A)$ とする。次は同値である。
\begin{enumerate}
\item $K$ は $I$ に関して導来完備である。
\item 標準写像 $K \to R\lim (K \otimes_A^\mathbf{L} K_n^\bullet)$ は
$D(A)$ の同型である。
\end{enumerate}
\end{lemma}

\begin{proof}
(2) が成り立つならば、Lemma
\ref{more-algebra-lemma-koszul-derived-completion-complete} により
$K$ は $I$ に関して導来完備である。逆に、$K$ は $I$ に関して導来完備であると
仮定する。愚直な切断によるフィルトレーション
$$
K_n^\bullet \supset
\sigma_{\geq -r + 1}K_n^\bullet \supset
\sigma_{\geq -r + 2}K_n^\bullet \supset \ldots \supset
\sigma_{\geq -1}K_n^\bullet \supset
\sigma_{\geq 0}K_n^\bullet = A
$$
を考える（Homology, Section \ref{homology-section-truncations}）。
構成 $R\lim(K \otimes E)$ は第二変数について完全なので
（Lemma \ref{more-algebra-lemma-tensor-Rlim-exact}）、$p < 0$ に対して
$$
R\lim \left(
K \otimes_A^\mathbf{L}
(\sigma_{\geq p}K_n^\bullet/ \sigma_{\geq p + 1}K_n^\bullet)
\right) = 0
$$
を示せば十分である。上の Koszul 複体の明示的記述により
$$
R\lim \left(
K \otimes_A^\mathbf{L}
(\sigma_{\geq p}K_n^\bullet/ \sigma_{\geq p + 1}K_n^\bullet)
\right) =
\bigoplus\nolimits_{i_1, \ldots, i_{-p}}
T(K, f_{i_1}\ldots f_{i_{-p}})
$$
であり、$K$ に対する仮定により $p < 0$ ならば零である。
\end{proof}

\begin{lemma}
\label{more-algebra-lemma-derived-completion-koszul}
Situation \ref{more-algebra-situation-koszul} において、$K \in D(A)$ を導来極限
$K' = R\lim( K \otimes_A^\mathbf{L} K_n^\bullet )$ に送る関手は、
Lemma \ref{more-algebra-lemma-derived-completion} で構成した包含関手
$D_{comp}(A) \to D(A)$ の左随伴である。
\end{lemma}

\begin{proof}[第一の証明]
対応 $K \leadsto K'$ は関手であり、Lemma
\ref{more-algebra-lemma-koszul-derived-completion-complete} により $K'$ は $I$ に関して
導来完備である。形式的な議論（省略）により、$K$ が $I$ に関して導来完備ならば
$K \to K'$ が同型であることを示せば十分である。これは Lemma
\ref{more-algebra-lemma-characterize-derived-complete-Koszul} である。
\end{proof}

\begin{proof}[第二の証明]
% P02-E1104: frozen source line 26407 omits "by" in the Denote construction.
Lemma \ref{more-algebra-lemma-derived-completion} で構成した随伴を
$K \mapsto K^\wedge$ と表す。同補題により
$$
K^\wedge = R\Hom\left((A \to \prod\nolimits_{i_0} A_{f_{i_0}} \to
\prod\nolimits_{i_0 < i_1} A_{f_{i_0}f_{i_1}}
\to \ldots \to A_{f_1\ldots f_r}), K\right)
$$
である。Lemma \ref{more-algebra-lemma-extended-alternating-Cech-is-colimit-koszul} では、
増広交代 {\v C}ech 複体
$$
A \to \prod\nolimits_{i_0} A_{f_{i_0}} \to
\prod\nolimits_{i_0 < i_1} A_{f_{i_0}f_{i_1}}
\to \ldots \to A_{f_1\ldots f_r}
$$
が、次数 $0, \ldots, r$ に置かれた Koszul 複体
$K^n = K(A, f_1^n, \ldots, f_r^n)$ の余極限であることを見た。
$K^n$ は有限自由 $A$-加群の有限鎖複体であり、その双対は
（Lemma \ref{more-algebra-lemma-dual-perfect-complex} の意味で）
$R\Hom_A(K^n, A) = K_n$ であることに注意する。ここで $K_n$ は通常どおり次数
$-r, \ldots, 0$ に置かれた Koszul 余鎖複体である。
従って
$$
R\Hom_A(\text{hocolim} K^n, K) =  R\lim (K \otimes_A^\mathbf{L} K_n)
$$
を示せば十分である。これは Lemma
\ref{more-algebra-lemma-colim-and-lim-of-duals} から従う。
\end{proof}

\begin{lemma}
\label{more-algebra-lemma-derived-complete-bounded}
$I = (f_1, \ldots, f_r)$ を環 $A$ の有限生成イデアルとする。
$K$ を $D(A)$ の導来完備対象とする。次は同値である。
\begin{enumerate}
\item $i > 0$ に対して $H^i(K) = 0$ である。
\item $i > 0$ に対して $H^i(K \otimes_A^\mathbf{L} A/I) = 0$ である。
\item $i > 0$ に対して
$H^i(K \otimes_A^\mathbf{L} K_1^\bullet) = 0$ である。ここで
$K_1^\bullet$ は Situation \ref{more-algebra-situation-koszul} のものである。
\end{enumerate}
\end{lemma}

\begin{proof}
(1) $\Rightarrow$ (2) は常に成り立つ。(2) $\Rightarrow$ (3) は Lemma
\ref{more-algebra-lemma-bounded-above-mod-I} から従う。(3) を仮定する。
$s = 0, \ldots, r$ に対して複体
$$
K(s) =  K \otimes_A^\mathbf{L} K_\bullet(A, f_1, \ldots, f_s)
$$
を考える。ここで $K_\bullet(A, f_1, \ldots, f_s)$ はコホモロジー次数
$-s, \ldots, 0$ に置かれた Koszul 複体である。Lemma
\ref{more-algebra-lemma-cone-koszul} により
% P02-E1105: frozen source line 26457 says "We have a distinguished triangles".
$$
K(s - 1) \xrightarrow{f_s} K(s - 1) \to K(s) \to K(s - 1)[1]
$$
という区別三角がある。$K(0) = K$ は導来完備なので、昇納法により各
$K(s)$ は $D(A)$ の導来完備対象である。降納法により、$K(s)$ は次数
$> 0$ に非零コホモロジーを持たないことを示す。実際、$s = r$ については仮定
(3) により成り立つ。$s > 0$ で $K(s)$ について成り立つならば、$i > 0$ に
対して $f_s : H^i(K(s - 1)) \to H^i(K(s - 1))$ は全射であり、
Lemmas \ref{more-algebra-lemma-serre-subcategory} and
\ref{more-algebra-lemma-derived-complete-zero} により $H^i(K(s - 1))$ は消滅すると結論する。
\end{proof}

\begin{lemma}
\label{more-algebra-lemma-derived-complete-zero-bis}
\begin{slogan}
導来 Nakayama
\end{slogan}
\begin{reference}
関連する結果は \cite[Proposition 6.5]{Dwyer-Greenlees} である。
導来 Nakayama の補題は、例えば Columbia University における 2018 年秋の
Bhatt による Prismatic cohomology 第 3 講の Section 2 property (2) に見られる。
Leonid Positselski は
\url{https://mathoverflow.net/a/331501} で証明を提案した。しかしここでは、
コメント中で Anonymous が提案した証明に従う。
\end{reference}
$I$ を環 $A$ の有限生成イデアルとし、$K$ を $D(A)$ の導来完備対象とする。
$K \otimes_A^\mathbf{L} A/I = 0$ ならば $K = 0$ である。
\end{lemma}

\begin{proof}
これは Lemma \ref{more-algebra-lemma-derived-complete-bounded} の特別な場合であるが、直接の
証明も与える。$I$ の生成元 $f_1, \ldots, f_r$ を選ぶ。
$K_n$ を $A$ 上の $f_1^n, \ldots, f_r^n$ に関する Koszul 複体と表す。
$K_n$ は有界であり、$K_n$ のコホモロジー加群は
$f_1^n, \ldots, f_r^n$ によって、従って $I^{nr}$ によって零化されることを
思い出す。Lemma \ref{more-algebra-lemma-derived-vanishing-mod-I} により
$K \otimes_A^\mathbf{L} K_n = 0$ である。$K$ は導来完備なので、Lemma
\ref{more-algebra-lemma-derived-completion-koszul} により
$K = R\lim K \otimes_A^\mathbf{L} K_n = 0$ であり、所望の結論を得る。
\end{proof}

\noindent
Koszul 複体との関係の応用として、導来完備化が有限コホモロジー次元を持つことを
得る。

\begin{lemma}
\label{more-algebra-lemma-derived-completion-finite-cohomological-dimension}
$A$ を環とし、$I \subset A$ を $r$ 個の元で生成できるイデアルとする。
このとき導来完備化は有限コホモロジー次元を持つ。
\begin{enumerate}
\item $K \to L$ を $D(A)$ の射とし、$i \geq 1$ に対して
$H^i(K) \to H^i(L)$ が同型で、$i = 0$ に対して全射であるとする。このとき
$i \geq 1$ に対して $H^i(K^\wedge) \to H^i(L^\wedge)$ は同型であり、
$i = 0$ に対して全射である。
\item $K \to L$ を $D(A)$ の射とし、$i \leq -1$ に対して
$H^i(K) \to H^i(L)$ が同型で、$i = 0$ に対して単射であるとする。このとき
$i \leq -r - 1$ に対して $H^i(K^\wedge) \to H^i(L^\wedge)$ は同型であり、
$i = -r$ に対して単射である。
\end{enumerate}
\end{lemma}

\begin{proof}
$I$ が $f_1, \ldots, f_r$ によって生成されるとする。任意の $K \in D(A)$ に
対して、Lemma \ref{more-algebra-lemma-derived-completion-koszul} により
$K^\wedge = R\lim K \otimes_A^\mathbf{L} K_n$ である。ここで $K_n$ は
$f_1^n, \ldots, f_r^n$ に関する Koszul 複体である。従って Lemma
\ref{more-algebra-lemma-break-long-exact-sequence-modules} により短完全列
$$
0 \to R^1\lim H^{i - 1}(K \otimes_A^\mathbf{L} K_n)
\to H^i(K^\wedge) \to \lim H^i(K \otimes_A^\mathbf{L} K_n) \to 0
$$
を得る。

\medskip\noindent
(1) の証明。区別三角 $K \to L \to C \to K[1]$ を選ぶ。
このとき $i \geq 0$ に対して $H^i(C) = 0$ である。
$K_n$ は次数 $\leq 0$ に置かれているので、
$i \geq 0$ に対して $H^i(C \otimes_A^\mathbf{L} K_n) = 0$ であり、
$H^{-1}(C \otimes_A^\mathbf{L} K_n) =
H^{-1}(C) \otimes_A A/(f_1^n, \ldots, f_r^n)$ は全射な遷移写像を持つ系である。
上の表示式は $i \geq 0$ に対して $H^i(C^\wedge) = 0$ であることを示す。
区別三角 $K^\wedge \to L^\wedge \to C^\wedge \to K^\wedge[1]$ を適用して
(1) を得る。

\medskip\noindent
(2) の証明。区別三角 $K \to L \to C \to K[1]$ を選ぶ。
このとき $i < 0$ に対して $H^i(C) = 0$ である。
$K_n$ は次数 $-r, \ldots, 0$ に置かれているので、$i < -r$ に対して
$H^i(C \otimes_A^\mathbf{L} K_n) = 0$ である。上の表示式は
% P02-E1106: frozen source line 26551 says i < r, but the preceding bound and conclusion require i < -r.
$i < r$ に対して $H^i(C^\wedge) = 0$ であることを示す。
区別三角 $K^\wedge \to L^\wedge \to C^\wedge \to K^\wedge[1]$ を適用して
(2) を得る。
\end{proof}

\begin{lemma}
\label{more-algebra-lemma-derived-completion-spectral-sequence}
$A$ を環、$I \subset A$ を有限生成イデアルとする。
$K^\bullet$ を $A$-加群のフィルター付き複体とする。二重次数付き導来完備
$A$-加群からなる標準スペクトル系列
$(E_r, \text{d}_r)_{r \geq 1}$ が存在する。ここで $d_r$ の双次数は
$(r, -r + 1)$ であり、
$$
E_1^{p, q} = H^{p + q}((\text{gr}^pK^\bullet)^\wedge)
$$
である。各 $K^n$ 上のフィルトレーションが有限ならば、このスペクトル系列は有界で
$H^*((K^\bullet)^\wedge)$ に収束する。
\end{lemma}

\begin{proof}
Lemma \ref{more-algebra-lemma-derived-completion} により、ある $C \in D^b(A)$ に対して導来完備化
が $R\Hom_A(C, -)$ で与えられることを知っている。Lemmas
\ref{more-algebra-lemma-derived-completion-finite-cohomological-dimension} and
\ref{more-algebra-lemma-projective-amplitude} により、$C$ は有限射影次元を持つ。
従って $C$ を表す射影加群の有界複体 $P^\bullet$ を選べる。このとき
$$
M^\bullet = \Hom^\bullet(P^\bullet, K^\bullet)
$$
は $(K^\bullet)^\wedge$ を表す $A$-加群の複体である。これは
$F^pM^\bullet = \Hom^\bullet(P^\bullet, F^pK^\bullet)$
で与えられるフィルトレーションを持つ。
$F^pM^\bullet$ は $(F^pK^\bullet)^\wedge$ を表し、従って
% P02-E1107: frozen source line 26583 drops the filtration index p from gr^p K in the proof.
$\text{gr}^pM^\bullet$ は $(\text{gr}K^\bullet)^\wedge$ を表す。
従ってフィルター付き複体 $M^\bullet$ のスペクトル系列を取ることにより、
所望のスペクトル系列を得る。Homology, Section
\ref{homology-section-filtered-complex} を見よ。各 $K^n$ 上のフィルトレーションが
有限ならば、$P^\bullet$ は有界複体なので各 $M^n$ 上のフィルトレーションも
有限である。従って最後の主張は Homology, Lemma
\ref{homology-lemma-biregular-ss-converges} から従う。
\end{proof}

\begin{example}
\label{more-algebra-example-derived-completion-spectral-sequence}
$A$ を環、$I \subset A$ を有限生成イデアルとし、$K^\bullet$ を $A$-加群の複体
とする。$F^pK^\bullet = \tau_{\leq -p}K^\bullet$ として Lemma
\ref{more-algebra-lemma-derived-completion-spectral-sequence} を適用できる。このとき
$$
E_1^{p, q} = H^{p + q}(H^{-p}(K^\bullet)^\wedge[p]) =
H^{2p + q}(H^{-p}(K^\bullet)^\wedge)
$$
を持ち、$H^{p + q}((K^\bullet)^\wedge)$ に収束する有界スペクトル系列を得る。
$p = -j$ および $q = i + 2j$ と番号を付け直すと、任意の $K \in D(A)$ に対して、
双次数付き導来完備加群からなる有界スペクトル系列
$(E'_r, d'_r)_{r \geq 2}$ がある。ここで $d'_r$ の双次数は
$(r, -r + 1)$ であり、
$$
(E'_2)^{i, j} = H^i(H^j(K)^\wedge)
$$
を持ち、$H^{i + j}(K^\wedge)$ に収束する。
\end{example}

\begin{lemma}
\label{more-algebra-lemma-restriction-derived-complete}
$A \to B$ を環写像、$I \subset A$ をイデアルとする。制限関手
$D(B) \to D(A)$ による $D_{comp}(A, I)$ の逆像は
$D_{comp}(B, IB)$ である。
\end{lemma}

\begin{proof}
Lemma \ref{more-algebra-lemma-ideal-of-elements-complete-wrt} を用いると、
$L \in D(B)$ が $D_{comp}(B, IB)$ に属することと、すべての局所切断
% P02-E1108: frozen source line 26625 calls an element f in the ring ideal I a "local section".
$f \in I$ に対して $T(L, f)$ が零であることは同値である。
$T(L, f)$ のコホモロジーは Abel 群の圏で計算されるので、$f$ を $A$ の元と
考えるか、$f$ の $B$ における像を取るかは問題にならない。
補題はこれと導来完備対象の定義から直ちに従う。
\end{proof}

\begin{lemma}
\label{more-algebra-lemma-restriction-derived-complete-equivalence}
$A \to B$ を環写像、$I \subset A$ を有限生成イデアルとする。
$A \to B$ が平坦で $A/I \cong B/IB$ ならば、制限関手
$D(B) \to D(A)$ は圏同値
$D_{comp}(B, IB) \to D_{comp}(A, I)$ を導く。
\end{lemma}

\begin{proof}
$I$ の生成元 $f_1, \ldots, f_r$ を選ぶ。
% P02-E1109: frozen source lines 26643-26646 omit "by" in the Denote construction.
Lemma \ref{more-algebra-lemma-map-identifies-koszul-and-cech-complexes} の増広交代
{\v C}ech 複体の擬同型を
$\check{\mathcal{C}}^\bullet_A \to \check{\mathcal{C}}^\bullet_B$
と表す。$K \in D_{comp}(A, I)$ とする。$I^\bullet$ を $K$ を表す
$A$-加群の K-単射的複体とする。すべての $f \in I$ および
$n \in \mathbf{Z}$ に対して $\Ext^n_A(A_f, K)$ と
$\Ext^n_A(B_f, K)$ は零なので（Lemma \ref{more-algebra-lemma-hom-from-Af}）、
$\check{\mathcal{C}}^\bullet_A \to A$ および
$\check{\mathcal{C}}^\bullet_B \to B$ は準同型
$$
I^\bullet = \Hom_A(A, I^\bullet) \longrightarrow
\text{Tot}(\Hom_A(\check{\mathcal{C}}^\bullet_A, I^\bullet))
$$
および
$$
\Hom_A(B, I^\bullet) \longrightarrow
\text{Tot}(\Hom_A(\check{\mathcal{C}}^\bullet_B, I^\bullet))
$$
を導き、これらは擬同型である。いくつかの詳細は省略する。
$\check{\mathcal{C}}^\bullet_A \to \check{\mathcal{C}}^\bullet_B$ は擬同型であり、
$I^\bullet$ は K-単射的なので、
$\Hom_A(B, I^\bullet) \to I^\bullet$ は擬同型であると結論する。
複体 $\Hom_A(B, I^\bullet)$ は $B$-加群の複体なので、$K$ は制限写像の像に属する。
すなわち関手は本質的全射である。
% P02-E1110: frozen source line 26668 lacks terminal punctuation.

\medskip\noindent
実際、この議論は
$F : D_{comp}(A, I) \to D_{comp}(B, IB)$,
$K \mapsto \Hom_A(B, I^\bullet)$
が制限の左逆であることを示す。最後に $L \in D_{comp}(B, IB)$ とする。
$L$ を $B$-加群の K-単射的複体 $J^\bullet$ で表す。
Lemma \ref{more-algebra-lemma-K-injective-flat} により、$J^\bullet$ は $A$-加群の複体としても
K-単射的なので、
$F(\text{restriction of }L) = \Hom_A(B, J^\bullet)$ である。
$B$-加群の複体の写像 $J^\bullet \to \Hom_A(B, J^\bullet)$ があり、
これと $\Hom_A(B, J^\bullet) \to J^\bullet$ との合成は恒等写像である。
従って $F$ は制限の右逆でもあり、証明が完了する。
\end{proof}

\section{導来完備加群の圏}
\label{more-algebra-section-derived-complete-modules}

\noindent
$A$ を環、$I$ をイデアルとする。導来完備加群の圏を
% P02-E1111: frozen source lines 26696-26698 omit "by" in the Denote construction.
$\mathcal{C}$ と表す。Definition \ref{more-algebra-definition-derived-complete} を見よ。
本節ではこの圏のいくつかの性質を論じる。Examples, Section
\ref{examples-section-derived-complete-modules} では、一般に
$\mathcal{C}$ が Grothendieck Abel 圏でないことを示す。

\medskip\noindent
Lemma \ref{more-algebra-lemma-serre-subcategory} により、圏 $\mathcal{C}$ は Abel 圏であり、
包含関手 $\mathcal{C} \to \text{Mod}_A$ は完全である。

\medskip\noindent
$D_{comp}(A) \subset D(A)$ は積で閉じており
（Definition \ref{more-algebra-definition-derived-complete} に続く議論を見よ）、
$D(A)$ における積は複体の水準で計算されるので、$\mathcal{C}$ は
$\text{Mod}_A$ における積と一致する積を持つ。従って $\mathcal{C}$ は実際に
任意の極限を持ち、包含関手 $\mathcal{C} \to \text{Mod}_A$ はそれらと可換する。
Categories, Lemma \ref{categories-lemma-limits-products-equalizers} を見よ。

\medskip\noindent
$I$ が有限生成であると仮定する。${}^\wedge : D(A) \to D(A)$ を Lemma
\ref{more-algebra-lemma-derived-completion} の導来完備化関手とする。関手
$$
\text{Mod}_A \longrightarrow \mathcal{C},\quad
M \longmapsto H^0(M^\wedge)
$$
が包含関手 $\mathcal{C} \to \text{Mod}_A$ の左随伴であることを示そう。
例えば Lemma \ref{more-algebra-lemma-derived-completion-finite-cohomological-dimension} により
$i > 0$ に対して $H^i(M^\wedge) = 0$ であることに注意する。
従って $N$ が導来完備 $A$-加群ならば
\begin{align*}
\Hom_\mathcal{C}(H^0(M^\wedge), N)
& =
\Hom_{D_{comp}(A)}(M^\wedge, N)\\
& =
\Hom_{D(A)}(M, N) \\
& =
\Hom_A(M, N)
\end{align*}
であり、所望の結論を得る。

\medskip\noindent
$T$ を前順序集合とし、$t \mapsto M_t$ を導来完備 $A$-加群の系、すなわち
$\mathcal{C}$ における $T$ 上の系とする。Categories, Section
\ref{categories-section-posets-limits} を見よ。$\text{Mod}_A$ におけるこの系の
余極限を
% P02-E1112: frozen source lines 26741-26744 omit "by" in the Denote construction.
$\colim_{t \in T} M_t$ と表す。上の議論から形式的に
$$
H^0((\colim_{t \in T} M_t)^\wedge)
$$
が $\mathcal{C}$ におけるこの系の余極限であることが従う。このようにして
$\mathcal{C}$ がすべての余極限を持つことが分かる。一般に包含関手
$\mathcal{C} \to \text{Mod}_A$ は余極限と可換しない。Examples, Section
\ref{examples-section-derived-complete-modules} を見よ。

\begin{lemma}
\label{more-algebra-lemma-derived-complete-modules}
$A$ を環、$I \subset A$ をイデアルとする。導来完備加群の圏
$\mathcal{C}$ は Abel 圏であって任意の極限を持ち、包含関手
$F : \mathcal{C} \to \text{Mod}_A$ は完全で極限と可換する。
$I$ が有限生成ならば、$\mathcal{C}$ は任意の余極限を持ち、$F$ は左随伴を持つ。
% P02-E1113: frozen source line 26760 ends the lemma sentence without terminal punctuation.
\end{lemma}

\begin{proof}
以上の議論を要約したものである。
\end{proof}

\section{単項イデアルに対する導来完備化}
\label{more-algebra-section-derived-completion-principal}

\noindent
本節では、イデアルが一つの元によって生成されるとき導来完備化に何が起こるかを
論じる。

\begin{lemma}
\label{more-algebra-lemma-lift-universally}
$A$ を環、$f \in A$ とする。
$A[f^c] = A[f^{c + 1}] = A[f^{c + 2}] = \ldots$ となる整数 $c \geq 1$ が
存在するならば（例えば $A$ が Noether ならば）、すべての $n \geq 1$ に対して
$D(A)$ に写像
$$
(A \xrightarrow{f^n} A) \longrightarrow A/(f^n),
\quad\text{and}\quad
A/(f^{n + c}) \longrightarrow (A \xrightarrow{f^n} A)
$$
が存在し、$D(A)$ におけるプロ対象 $\{A/(f^n)\}$ と
$\{(f^n : A \to A)\}$ の同型を導く。
\end{lemma}

\begin{proof}
第一の表示矢印は明らかである。補題の第二の矢印は図式
$$
\xymatrix{
A/A[f^c] \ar[r]_-{f^{n + c}} \ar[d]_{f^c} & A \ar[d]^1 \\
A \ar[r]^{f^n} & A
}
$$
によって定義できる。上の水平矢印は単射なので、上段の複体は
$A/f^{n + c}A$ と擬同型である。プロ対象についての主張を示すために必要な合成の
計算は省略する。
\end{proof}

\begin{lemma}
\label{more-algebra-lemma-when-does-it-work}
$A$ を環、$f \in A$ とし、$I = (f)$ と置く。この状況では、素朴な導来完備化
$K \mapsto K' = R\lim (K \otimes_A^\mathbf{L} A/f^nA)$ と、Lemma
\ref{more-algebra-lemma-derived-completion-koszul} の導来完備化
$$
K \mapsto K^\wedge = R\lim (K \otimes_A^\mathbf{L} (A \xrightarrow{f^n} A))
$$
がある。関手の自然変換 $K^\wedge \to K'$ が同型であることと、$A$ の
$f$-冪ねじれが有界であることは同値である。
\end{lemma}

\begin{proof}
$f$-冪ねじれが有界ならば、Lemma \ref{more-algebra-lemma-lift-universally} によりプロ対象
$\{(f^n : A \to A)\}$ と $\{A/f^nA\}$ は同型である。従って Lemma
\ref{more-algebra-lemma-Rlim-pro-equal} により関手は同型である。逆に、Lemma
\ref{more-algebra-lemma-tensor-Rlim-exact} から、この条件は、すべての $K \in D(A)$ に対して
$$
R\lim (K \otimes_A^\mathbf{L} A[f^n])
$$
が零であることにほかならない。ここで系 $(A[f^n])$ の写像は $f$ による乗法で
与えられる。$K = A$ および $K = \bigoplus_{i \in \mathbf{N}} A$ とすると、
Lemma \ref{more-algebra-lemma-Rlim-zero-of-direct-sums} により、これは $(A[f^n])$ がプロ対象
として零、すなわちある $n$ に対して $f^{n - 1}A[f^n] = 0$、すなわち
$A[f^{n - 1}] = A[f^n]$、すなわち $f$-冪ねじれが有界であることを意味する。
\end{proof}

\begin{example}
\label{more-algebra-example-derived-complete-modules}
$A$ を環、$f \in A$ を非零因子とする。念頭に置くべき例は
$A = \mathbf{Z}_p$ および $f = p$ である。$M$ を $A$-加群とする。
主張：$M$ が $f$ に関して導来完備であることと、短完全列
$$
0 \to K \to L \to M \to 0
$$
であって、$K, L$ が $f$-進完備かつ $f$-ねじれが零であるものが存在することは
同値である。実際、
% P02-E1117: frozen source line 26853 says "there is a such a short exact sequence".
そのような短完全列が存在するならば、$f$ は $K$ と $L$ 上の非零因子なので
$$
M \otimes_A^\mathbf{L} (A \xrightarrow{f^n} A) = (K/f^nK \to L/f^nL)
$$
であり、$R\lim (M \otimes_A^\mathbf{L} (A \xrightarrow{f^n} A))$ は
$K \to L$、すなわち $M$ と擬同型であると結論する。これは Lemma
\ref{more-algebra-lemma-characterize-derived-complete-Koszul} により $M$ が導来完備である
ことを示す。逆に $M$ が導来完備であると仮定する。$F$ が自由 $A$-加群である
ような全射 $F \to M$ を選ぶ。$f$ は $F$ 上の非零因子なので、$F$ の導来完備化
は $L = \lim F/f^nF$ である。$L$ は $f$-ねじれなしであることに注意する。
% P02-E1118: frozen source line 26864 uses the compound predicate "f-torsion free" without the second hyphen.
実際、$x_n \in F$ を持つ $(x_n)$ が $L$ の元 $\xi$ を表し $f\xi = 0$ ならば、
ある $z_n, y_n \in F$ に対して
$x_n = x_{n + 1} + f^nz_n$ および $fx_n = f^ny_n$ である。このとき
$f^n y_n = fx_n =
fx_{n + 1} + f^{n + 1}z_n = f^{n + 1}y_{n + 1} + f^{n + 1}z_n$ であり、
$f$ は $F$ 上の非零因子なので $y_n \in fF$ であることが分かる。これは
$x_n \in f^nF$、すなわち $\xi = 0$ を意味する。$L$ は導来完備化なので、
普遍性は $F \to M$ を経由させる写像 $L \to M$ を与える。
$K = \Ker(L \to M)$ を核とする。再び $K$ は $f$-ねじれなしなので、$K$ の
導来完備化は $\lim K/f^nK$ である。一方、$M$ と $L$ はともに導来完備なので、
Lemma \ref{more-algebra-lemma-serre-subcategory} により $K$ も導来完備である。
従って $K = \lim K/f^nK$ であり、主張が証明された。
\end{example}

\begin{example}
\label{more-algebra-example-derived-complete-not-complete}
$p$ を素数とする。$x$ を $py$ に写す多項式環の写像
$\mathbf{Z}_p[x] \to \mathbf{Z}_p[y]$ を考える。誘導される通常の
$p$-進完備化の写像の余核
$M = \Coker(\mathbf{Z}_p[x]^\wedge \to \mathbf{Z}_p[y]^\wedge)$
を考える。このとき Proposition \ref{more-algebra-proposition-derived-complete-modules} と
Lemma \ref{more-algebra-lemma-serre-subcategory} により、$M$ は導来完備
$\mathbf{Z}_p$-加群である。Example \ref{more-algebra-example-derived-complete-modules} の
議論も見よ。しかし $1 + py + p^2 y^2 + \ldots$ は $M$ の非零元に写り、
これは $\bigcap p^nM$ に含まれるので、$M$ は $p$-進完備でない。
\end{example}

\begin{example}
\label{more-algebra-example-spectral-sequence-principal}
$A$ を環、$f \in A$ とする。
% P02-E1114: frozen source line 26897 omits "by" in the Denote construction.
$(f)$ に関する導来完備化を $K \mapsto K^\wedge$ と表す。
$M$ を $A$-加群とする。
Lemma \ref{more-algebra-lemma-derived-completion-koszul} による
$$
M^\wedge = R\lim (M \xrightarrow{f^n} M)
$$
および Lemma \ref{more-algebra-lemma-break-long-exact-sequence-modules} を用いると
$$
H^{-1}(M^\wedge) = \lim M[f^n] = T_f(M)
$$
を得る。これは {\it $M$ の $f$-進 Tate 加群}である。ここで写像
$M[f^n] \to M[f^{n - 1}]$ は $f$ による乗法で与えられる。さらに
$H^0(M^\wedge)$ を記述する短完全列
$$
0 \to R^1\lim M[f^n] \to H^0(M^\wedge) \to \lim M/f^n M \to 0
$$
がある。遷移写像は全射なので（Lemma
\ref{more-algebra-lemma-compute-Rlim-modules}）、$H^1(M^\wedge) = R^1\lim M/f^nM = 0$
である。$M^\wedge$ の他のすべてのコホモロジーは自明な理由で零である。
最後に、$K \in D(A)$ および $p \in \mathbf{Z}$ に対して短完全列
$$
0 \to H^0(H^p(K)^\wedge) \to H^p(K^\wedge) \to T_f(H^{p + 1}(K)) \to 0
$$
がある。これは Example \ref{more-algebra-example-derived-completion-spectral-sequence} の
スペクトル系列から従う。このスペクトル系列は $E_2$ で退化する
（$i = -1, 0$ だけが非零項を与えるため）。次の補題はさらに多くの情報を与える。
\end{example}

\begin{lemma}
\label{more-algebra-lemma-compare-completion-derived-completion}
$A$ を環、$f \in A$ とし、$K$ を $D(A)$ の対象とする。
% P02-E1115: frozen source line 26930 omits "by" in the Denote construction.
$K_n = K \otimes_A^\mathbf{L} (A \xrightarrow{f^n} A)$ と表す。
すべての $p \in \mathbf{Z}$ に対して可換図式
$$
\xymatrix{
& 0 & 0 \\
0 \ar[r] &
\widehat{H^p(K)} \ar[r] \ar[u] &
\lim H^p(K_n) \ar[r] \ar[u] &
T_f(H^{p + 1}(K)) \ar[r] &
0 \\
0 \ar[r] &
H^0(H^p(K)^\wedge) \ar[r] \ar[u] &
H^p(K^\wedge) \ar[r] \ar[u] &
T_f(H^{p + 1}(K)) \ar[r] \ar@{=}[u] &
0 \\
&
R^1\lim H^p(K)[f^n] \ar[u] \ar[r]^\cong &
R^1\lim H^{p - 1}(K_n) \ar[u] \\
& 0 \ar[u] & 0 \ar[u]
}
$$
があり、その行と列は完全である。ここで
$\widehat{H^p(K)} = \lim H^p(K)/f^nH^p(K)$ は通常の $f$-進完備化である。
左の垂直短完全列と中央の水平短完全列は Example
\ref{more-algebra-example-spectral-sequence-principal} から取ったものである。
% P02-E1116: frozen source line 26954 lacks terminal punctuation before the next sentence.
中央の垂直短完全列は Lemma
\ref{more-algebra-lemma-break-long-exact-sequence-modules} のものである。
\end{lemma}

\begin{proof}
上の水平短完全列を構成するため、$K_n$ が
$f^n : K \to K$ の錐のシフトとして構成されることから来る逆系の短完全列
$$
0 \to H^p(K)/f^nH^p(K) \to H^p(K_n) \to H^{p + 1}(K)[f^n] \to 0
$$
があることに注意する。これらの逆極限を取ると、上の水平短完全列を得る。
Homology, Lemma \ref{homology-lemma-Mittag-Leffler} を見よ。

\medskip\noindent
補題のような可換図式があることを証明しよう。写像
$L = \tau_{\leq p}K \to K$ を考える。
$L_n = L \otimes_A^\mathbf{L} (A \xrightarrow{f^n} A)$
と置くと、逆系の写像 $(L_n) \to (K_n)$ が得られ、これは短完全列の写像
$$
\xymatrix{
0 & 0 \\
\lim H^p(L_n) \ar[r] \ar[u] &
\lim H^p(K_n) \ar[u] \\
H^p(L^\wedge) \ar[r] \ar[u] &
H^p(K^\wedge) \ar[u] \\
R^1\lim H^{p - 1}(L_n) \ar[r] \ar[u] &
R^1\lim H^{p - 1}(K_n) \ar[u] \\
0 \ar[u] & 0 \ar[u]
}
$$
を導く。$i > p$ に対して $H^i(L) = 0$ であり、$H^p(L) = H^p(K)$ なので、
補題の主張中の参照を用いた計算により
$H^p(L^\wedge) = H^0(H^p(K)^\wedge)$ および
$H^p(L_n) = H^p(K)/f^nH^p(K)$ であることが分かる。一方、
$H^{p - 1}(L_n) = H^{p - 1}(K_n)$ である。さらに
$H^0(H^p(K)^\wedge) \to \widehat{H^p(K)}$ の核が
$R^1\lim H^p(K)[f^n]$ に等しいことを既に知っているので、補題の主張に示した同型を
得る。上段を証明の第一段落の構成によって定義したとき、図式の最右の正方形が可換で
あることの確認は省略する。
\end{proof}

\begin{remark}
\label{more-algebra-remark-ML-condition}
Lemma \ref{more-algebra-lemma-compare-completion-derived-completion} の記法のもとで、逆系
$H^p(K_n)$ が ML を持つことと、逆系 $H^{p + 1}(K)[f^n]$ が ML を持つことは
同値であることも分かる。これは逆系の短完全列
$0 \to H^p(K)/f^nH^p(K) \to H^p(K_n) \to H^{p + 1}(K)[f^n] \to 0$
（補題の証明を見よ）と Homology, Lemma
\ref{homology-lemma-Mittag-Leffler} および Lemma
\ref{more-algebra-lemma-Mittag-Leffler} から従う。
\end{remark}

\begin{lemma}[Bhatt]
\label{more-algebra-lemma-torsion-and-derived-complete}
$I$ を環 $A$ の有限生成イデアルとし、$M$ を導来完備 $A$-加群とする。
$M$ が $I$-冪ねじれ加群ならば、ある $n$ に対して $I^nM = 0$ である。
\end{lemma}

\begin{proof}
$I = (f_1, \ldots, f_r)$ とする。各 $i$ に対して、ある $n_i$ が存在して
$f_i^{n_i}M = 0$ となることを示せば十分である。従って
$I = (f)$ が単項イデアルであると仮定してよい。
$B = \mathbf{Z}[x] \to A$ を、$x$ を $f$ に写す環写像とする。
Lemma \ref{more-algebra-lemma-restriction-derived-complete} により、$M$ はイデアル $(x)$ に
関して $B$-加群として導来完備である。$A$ を $B$ で置き換えると、
$f$ が $A$ の非零因子であると仮定してよい。

\medskip\noindent
$I = (f)$ で $f \in A$ が非零因子であると仮定する。
Example \ref{more-algebra-example-derived-complete-modules} によれば、短完全列
$$
0 \to K \xrightarrow{u} L \to M \to 0
$$
であって、$K$ と $L$ が $I$-進完備 $A$-加群で、その $f$-ねじれが零であるものが
存在する\footnote{証明には、$K$ と $L$ が $I$-進完備 $A$-加群であるような列
$K \xrightarrow{u} L \to M \to 0$ が存在することを示せば十分である。
これは、$F_i$ が自由であるような表示 $F_1 \to F_0 \to M \to 0$ を選び、
$K$ と $L$ をそれぞれ $F_1$ と $F_0$ の $f$-進完備化とすることで示せる。
実際、$f$ は非零因子なので、これらの完備化は導来完備化であり、列は完全なままで
ある。}。$K$ と $L$ を $I$-進位相を持つ位相加群と考える。このとき $u$ は
連続である。
$$
L_n = \{x \in L \mid f^n x \in u(K)\}
$$
と置く。$M$ は $f$-冪ねじれなので $L = \bigcup L_n$ である。
$N_n$ を $L$ における $L_n$ の閉包とする。Lemma
\ref{more-algebra-lemma-consequence-baire-complete-module} により、ある $n$ に対して
$N_n$ は $L$ で開である。そのような $n$ を固定する。
$f^{n + m} : L \to L$ は連続開写像であり、かつ
$f^{n + m} L_n \subset u(f^m K)$ なので、すべての $m \geq 1$ に対して
$u(f^mK)$ の閉包は開であると結論する。従って Lemma
\ref{more-algebra-lemma-open-mapping} により $u$ は開である。従って、ある $t$ に対して
$f^tL \subset \Im(u)$ であり、所望のとおり $f^t$ は $M$ を零化すると結論する。
\end{proof}

\begin{lemma}
\label{more-algebra-lemma-kernel-to-completion-square-zero}
% P02-E1119: frozen source line 27068 says "Let f in A be an element of a ring" instead of introducing the ring as A.
$f \in A$ を環の元とし、$J = \bigcap f^nA$ と置く。
$M$ を $f$ に関して導来完備な $A$-加群とする。このとき
$M' = \Ker(M \to \lim M/f^nM)$ と置けば $JM' = 0$ である。
特に $A$ が導来完備ならば、$J$ は平方零イデアルである。
\end{lemma}

\begin{proof}
$x \in M'$ および $g \in J$ を取る。すべての $n \geq 1$ に対して
$x = f^n x_n$ と書ける。$g$ は $f^nA$ に属するので、$M'$ の元
$y_n = gx_n$ は $x_n$ の選び方によらない。特に
$x_n = fx_{n + 1}$ と取ることができ、$y_n = fy_{n + 1}$ を得る。
従って $1/f^n$ を $y_n$ に写す写像 $A_f \to M$ を得る。
$M$ は導来完備なので、この写像は零でなければならない
（Lemma \ref{more-algebra-lemma-hom-from-Af}）。従ってすべての $n$ に対して $y_n = 0$ である。
$gx = gfx_1 = fy_1$ なので、これで証明が完了する。
\end{proof}

\begin{lemma}
\label{more-algebra-lemma-derived-complete-henselian}
$A$ をイデアル $I$ に関して導来完備な環とする。このとき $(A, I)$ は Hensel 対
である。
\end{lemma}

\begin{proof}
$f \in I$ とする。Lemma \ref{more-algebra-lemma-largest-ideal-henselian} により
$(A, fA)$ が Hensel 対であることを示せば十分である。
$A$ は $fA$ に関して導来完備であることに注意する
（Definition \ref{more-algebra-definition-derived-complete} から直ちに従う）。
Lemma \ref{more-algebra-lemma-complete-derived-complete} により、$A$ から $A$ の $f$-進完備化
$A'$ への写像は全射である。Lemma \ref{more-algebra-lemma-complete-henselian} により
対 $(A', fA')$ は Hensel 的である。従って Lemma
\ref{more-algebra-lemma-henselian-henselian-pair} により、
$(A, \bigcap f^nA)$ が Hensel 対であることを示せば十分である。これは Lemmas
\ref{more-algebra-lemma-kernel-to-completion-square-zero} and
\ref{more-algebra-lemma-locally-nilpotent-henselian} から従う。
\end{proof}

\begin{lemma}
\label{more-algebra-lemma-kernel-to-completion-nilpotent}
$A$ をイデアル $I$ に関して導来完備な環とし、$J = \bigcap I^n$ と置く。
$I$ が $r$ 個の元で生成できるならば、$N = 2^r$ として $J^N = 0$ である。
\end{lemma}

\begin{proof}
$r = 1$ のとき、これは Lemma
\ref{more-algebra-lemma-kernel-to-completion-square-zero} である。
$r > 1$ として $I = (f_1, \ldots, f_r)$ とする。Lemma
\ref{more-algebra-lemma-serre-subcategory} により、環 $A_t = A/f_r^tA$ は $I$ に関して
導来完備であり、従ってなおさら
$I_t = (f_1, \ldots, f_{r - 1})A_t$ に関して導来完備である。
$A \to A_t$ は $J$ を $J_t = \bigcap I_t^n$ に写すことに注意する。
帰納法により、$N = 2^r$ として $J_t^{N/2} = 0$ である。
イデアル $\bigcap \Ker(A \to A_t) = \bigcap f_r^t A$ は
$r = 1$ の場合により平方零である。これで証明が完了する。
\end{proof}

\begin{lemma}
\label{more-algebra-lemma-reduced-derived-complete-complete}
$A$ を有限生成イデアル $I$ に関して導来完備な被約環とする。
このとき $A$ は $I$-進完備である。
\end{lemma}

\begin{proof}
% P02-E1120: frozen source line 27132 is the sentence fragment "Follows from ...".
Lemma \ref{more-algebra-lemma-kernel-to-completion-nilpotent} と Proposition
\ref{more-algebra-proposition-derived-complete-modules} から従う。
\end{proof}


\section{Noether 環に対する導来完備化}
\label{more-algebra-section-derived-completion-noetherian}

\noindent
$A$ を環、$I \subset A$ をイデアルとする。任意の
$K \in D(A)$ に対して導来逆極限
$$
K' = R\lim (K \otimes_A^\mathbf{L} A/I^n)
$$
を考えることができる。これは $K$ に関する関手である。Remark
\ref{more-algebra-remark-constructing-tensor-with-limits-functorially} を参照せよ。
写像の系 $A \to A/I^n$ は写像 $K \to K'$ を誘導し、
$K'$ は $I$ に関して導来完備である
（Lemma \ref{more-algebra-lemma-naive-derived-completion}）。
この「素朴な」導来完備化の構成は、一般には Lemma
\ref{more-algebra-lemma-derived-completion} の随伴とは一致しない。
例えば、第2成分を平方零イデアルとして
$A = \mathbf{Z}_p \oplus \mathbf{Q}_p/\mathbf{Z}_p$、
$K = A[0]$、$I = (p)$ とすると、素朴な導来完備化は
$\mathbf{Z}_p[0]$ を与えるが、Lemma \ref{more-algebra-lemma-derived-completion} の構成は
$K^\wedge \cong \mathbf{Z}_p[1] \oplus \mathbf{Z}_p[0]$ を与える
（計算は省略する）。Lemma \ref{more-algebra-lemma-when-does-it-work} は、
$I$ が一つの元で生成される場合に二つの関手が一致する条件を特徴づける。

\medskip\noindent
% P02-E1121: frozen source line 27167 says "is the show" instead of "is to show".
この節の主な目標は、$A$ が Noether 環ならば素朴な導来完備化が
導来完備化に等しいことを示すことである。

\begin{lemma}
\label{more-algebra-lemma-sequence-Koszul-complexes}
% P02-E1123: frozen source line 27173 begins with the fragment "In Situation ... .".
Situation \ref{more-algebra-situation-koszul} のもとで、$A$ が Noether 環ならば、
$D(A)$ の pro-対象 $\{K_n^\bullet\}$ と
$\{A/(f_1^n, \ldots, f_r^n)\}$ は同型である\footnote{特に、任意の
$n$ に対してある $m \geq n$ が存在し、$K_m^\bullet \to K_n^\bullet$
は写像 $K_m^\bullet \to A/(f_1^m, \ldots, f_r^m)$ を経由する。}。
\end{lemma}

\begin{proof}
完全三角形の逆系
% P02-E1122: the frozen source display uses m in the quotient although every surrounding object is indexed by n.
$$
\tau_{\leq -1}K_n^\bullet \to K_n^\bullet \to A/(f_1^m, \ldots, f_r^m) \to
(\tau_{\leq -1}K_n^\bullet)[1]
$$
がある。Derived Categories, Remark
\ref{derived-remark-truncation-distinguished-triangle} を参照せよ。
Derived Categories, Lemma \ref{derived-lemma-pro-isomorphism} により、
逆系 $\tau_{\leq -1}K_n^\bullet$ が pro-零であることを示せば十分である。
$K_n^\bullet$ は $-r \leq i \leq 0$ を満たす次数 $i$ にしか
非零項をもたないことを思い出そう。従って Derived Categories, Lemma
\ref{derived-lemma-essentially-constant-cohomology} により、
$p \leq -1$ に対して $H^p(K_n^\bullet)$ が pro-零であることを示せばよい。
言い換えると、任意の $n \in \mathbf{N}$ に対し、ある $m \geq n$ が存在して
$H^p(K_m^\bullet) \to H^p(K_n^\bullet)$ が零となることを示す必要がある。
$A$ は Noether 環なので、
$$
H^p(K_n^\bullet) =
\frac{\Ker(K_n^p \to K_n^{p + 1})}{\Im(K_n^{p - 1} \to K_n^p)}
$$
は有限 $A$-加群である。さらに、$p < 0$ のとき写像
$K_m^p \to K_n^p$ は、成分がイデアル
$(f_1^{m - n}, \ldots, f_r^{m - n})$ に属する対角行列で与えられる。
Lemma \ref{more-algebra-lemma-homotopy-koszul} により、$H^p(K_n^\bullet)$ は
$J = (f_1^n, \ldots, f_r^n)$ により消去されることに注意する。
ここで $m - n \geq tn$ ならば
$(f_1^{m - n}, \ldots, f_r^{m - n}) \subset J^t$ である。
従って、イデアル $J$、加群 $M = K_n^p$、部分加群
$N = \Ker(K_n^p \to K_n^{p + 1})$ に Algebra, Lemma
\ref{algebra-lemma-Artin-Rees}（Artin--Rees）を適用すると、
$m$ が十分大きいとき、$K_m^p \to K_n^p$ の像と
$\Ker(K_n^p \to K_n^{p + 1})$ との共通部分は
$J \Ker(K_n^p \to K_n^{p + 1})$ に含まれる。そのような $m$ に対して
零写像を得る。
\end{proof}

\begin{proposition}
\label{more-algebra-proposition-noetherian-naive-completion-is-completion}
$A$ を Noether 環、$I \subset A$ をイデアルとする。
$K \in D(A)$ を導来逆極限
$K' = R\lim( K \otimes_A^\mathbf{L} A/I^n )$ に送る関手は、
Lemma \ref{more-algebra-lemma-derived-completion} で構成した包含関手
$D_{comp}(A) \to D(A)$ の左随伴である。
\end{proposition}

\begin{proof}
$(f_1, \ldots, f_r) = I$ とし、$K_n^\bullet$ を
$f_1^n, \ldots, f_r^n$ に関する Koszul 複体とする。Lemma
\ref{more-algebra-lemma-derived-completion-koszul} により、
$$
R\lim (K \otimes_A^\mathbf{L} K_n^\bullet) =
R\lim (K \otimes_A^\mathbf{L} A/(f_1^n, \ldots, f_r^n) ) =
R\lim (K \otimes_A^\mathbf{L} A/I^n ).
$$
を示せば十分である。Lemma \ref{more-algebra-lemma-sequence-Koszul-complexes} により、
$D(A)$ の pro-対象 $\{K_n^\bullet\}$ と
$\{A/(f_1^n, \ldots, f_r^n)\}$ は同型である。pro-対象
$\{A/(f_1^n, \ldots, f_r^n)\}$ と $\{A/I^n\}$ が同型であることは明らかである。
従って Lemma \ref{more-algebra-lemma-tensor-Rlim-pro-equal} により、
左から右への写像は同型である。
\end{proof}

\begin{lemma}
\label{more-algebra-lemma-noetherian-calculate}
$I$ を Noether 環 $A$ のイデアルとする。$M$ を $A$-加群、
$M^\wedge$ をその導来完備化とする。このとき短完全列
$$
0 \to R^1\lim \text{Tor}_{i + 1}^A(M, A/I^n) \to
H^{-i}(M^\wedge) \to \lim \text{Tor}_i^A(M, A/I^n) \to 0
$$
がある。$M \in D^-(A)$ に対しても同様の結果が成り立つ。
\end{lemma}

\begin{proof}
Proposition \ref{more-algebra-proposition-noetherian-naive-completion-is-completion} と
Lemma \ref{more-algebra-lemma-break-long-exact-sequence-modules} の直ちに得られる帰結である。
\end{proof}

\noindent
上の命題の応用として、Noether の場合に擬コヒーレント複体の
導来完備化を同定する。

\begin{lemma}
\label{more-algebra-lemma-derived-completion-pseudo-coherent}
$A$ を Noether 環、$I \subset A$ をイデアルとする。$K$ を $D(A)$ の対象とし、
% P02-E1124: frozen source line 27268 omits "is" after H^n(K).
すべての $n \in \mathbf{Z}$ に対して $H^n(K)$ が有限 $A$-加群であるとする。
このとき導来完備化のコホモロジー加群 $H^n(K^\wedge)$ は、
コホモロジー加群 $H^n(K)$ の $I$-進完備化である。
\end{lemma}

\begin{proof}
Lemma \ref{more-algebra-lemma-Noetherian-pseudo-coherent} により、
複体 $\tau_{\leq m}K$ はすべての $m$ に対して擬コヒーレントである。
従って $\tau_{\leq m}K$ は有限自由 $A$-加群からなる上に有界な複体
$P^\bullet$ で表される。このとき
$\tau_{\leq m}K \otimes_A^\mathbf{L} A/I^n = P^\bullet/I^nP^\bullet$ である。
従って
$(\tau_{\leq m}K)^\wedge = R\lim P^\bullet/I^nP^\bullet$
（Proposition \ref{more-algebra-proposition-noetherian-naive-completion-is-completion}）である。
$R\lim$ は項ごとの $\lim$ で与えられ
（Lemma \ref{more-algebra-lemma-compute-Rlim-modules}）、また $I$-進完備化は
有限 $A$-加群上の完全関手である
（Algebra, Lemma \ref{algebra-lemma-completion-flat}）から、
結果は $\tau_{\leq m}K$ に対して成り立つ。従って、導来完備化は
有限コホモロジー次元をもつので、結果は $K$ に対して成り立つ。
Lemma \ref{more-algebra-lemma-derived-completion-finite-cohomological-dimension} を参照せよ。
\end{proof}

\begin{lemma}
\label{more-algebra-lemma-derived-complete-finite}
$I$ を Noether 環 $A$ のイデアルとする。
$M$ を導来完備 $A$-加群とする。$M/IM$ が有限 $A/I$-加群ならば、
$M = \lim M/I^nM$ であり、$M$ は有限 $A^\wedge$-加群である。
\end{lemma}

\begin{proof}
$M/IM$ が有限であると仮定する。$M/IM$ の生成元に写る
$x_1, \ldots, x_t \in M$ を選ぶ。第 $i$ 基底ベクトルを $x_i$ に写す写像
$A^{\oplus t} \to M$ を得る。Proposition
\ref{more-algebra-proposition-noetherian-naive-completion-is-completion} により、
$A$ の導来完備化は $A^\wedge = \lim A/I^n$ である。$M$ は導来完備なので、
この写像は写像 $q : (A^\wedge)^{\oplus t} \to M$ を経由する。
Lemma \ref{more-algebra-lemma-derived-complete-zero} により加群 $\Coker(q)$ は零である。
従って $M$ は有限 $A^\wedge$-加群である。
$A^\wedge$ は Noether 環であり $IA^\wedge$ に関して完備なので、
$M$ は $I$-進完備である（Algebra, Lemmas
\ref{algebra-lemma-completion-Noetherian},
\ref{algebra-lemma-when-finite-module-complete-over-complete-ring}, and
\ref{algebra-lemma-Artin-Rees} を用いる）。
\end{proof}

\begin{lemma}
\label{more-algebra-lemma-when-derived-completion-is-completion}
$I$ を Noether 環 $A$ のイデアルとする。
\begin{enumerate}
\item $M$ が有限 $A$-加群で $N$ が平坦 $A$-加群ならば、
$M \otimes_A N$ の導来 $I$-進完備化は $M \otimes_A N$ の
通常の $I$-進完備化である。
\item $M$ が有限 $A$-加群で $f \in A$ ならば、
$M_f$ の導来 $I$-進完備化は $M_f$ の通常の $I$-進完備化である。
\end{enumerate}
\end{lemma}

\begin{proof}
% P02-E1125: frozen source line 27331 omits "by" in the Denote construction.
$A$-加群 $M$ に対し、その導来完備化を $M^\wedge$、
通常の完備化を $\lim M/I^nM$ と表す。$M$ が有限であると仮定する。
$i > 0$ に対して系 $\text{Tor}^A_i(M, A/I^n)$ は pro-零である。
Lemma \ref{more-algebra-lemma-tor-strictly-pro-zero} を参照せよ。
$N$ は平坦なので
$\text{Tor}_i^A(M \otimes_A N, A/I^n) =
\text{Tor}_i^A(M, A/I^n) \otimes_A N$ であり、従って系
$\text{Tor}^A_i(M \otimes_A N, A/I^n)$ についても同じことが成り立つ。
Lemma \ref{more-algebra-lemma-noetherian-calculate} により、
% P02-E1127: frozen source lines 27340-27342 leave the Rlim argument and usual-completion quotient ungrouped.
$R\lim (M \otimes_A N) \otimes_A^\mathbf{L} A/I^n$
は次数 $0$ にしかコホモロジーをもたず、それは通常の完備化
$\lim M \otimes_A N/ I^n(M \otimes_A N)$ で与えられる。これで (1) が従う。
(2) は (1) と $M_f = M \otimes_A A_f$ であることから従う。
\end{proof}

\begin{lemma}
\label{more-algebra-lemma-derived-completion-tensor-finite}
$I$ を Noether 環 $A$ のイデアルとする。
${}^\wedge$ は $I$ に関する導来完備化を表すものとする。
$K \in D^-(A)$ とする。
\begin{enumerate}
\item $M$ が有限 $A$-加群ならば、
$(K \otimes_A^\mathbf{L} M)^\wedge = K^\wedge \otimes_A^\mathbf{L} M$ である。
\item $L \in D(A)$ が擬コヒーレントならば、
$(K \otimes_A^\mathbf{L} L)^\wedge = K^\wedge \otimes_A^\mathbf{L} L$ である。
\end{enumerate}
\end{lemma}

\begin{proof}
$L$ を (2) のとおりとする。$K$ を自由 $A$-加群からなる
上に有界な複体 $P^\bullet$ で表すことができる。$L$ を有限自由
$A$-加群からなる上に有界な複体 $F^\bullet$ で表すことができる。
$\text{Tot}(P^\bullet \otimes_A F^\bullet)$ は
$K \otimes_A^\mathbf{L} L$ を表すので、
$(K \otimes_A^\mathbf{L} L)^\wedge$ は
$$
\text{Tot}((P^\bullet)^\wedge \otimes_A F^\bullet)
$$
で表される。ここで $(P^\bullet)^\wedge$ は、各項が通常の
% P02-E1126: frozen source line 27371 contains a stray standalone equality between "usual" and "derived".
$=$ 導来完備化 $(P^n)^\wedge$ である複体である。例えば Proposition
\ref{more-algebra-proposition-noetherian-naive-completion-is-completion} および Lemma
\ref{more-algebra-lemma-when-derived-completion-is-completion} を参照せよ。
これで (2) が従う。(1) は (2) の特別な場合である。
\end{proof}


\section{Berthelot と Ogus が導入した作用素}
\label{more-algebra-section-eta}

\noindent
この節では、\cite[Section 8]{Berthelot-Ogus} で導入され、
\cite[Section 6]{BMS} で一般化された構成を論じる。
この概念を論じる原論文を読者に強く推奨する。

\medskip\noindent
$A$ を環、$f \in A$ を非零因子とする。$M$ が $A$-加群ならば、
% P02-E1128: frozen source line 27394 omits the article in "following are equivalent".
Lemma \ref{more-algebra-lemma-torsion-free} により次は同値である。
\begin{enumerate}
\item $f$ は $M$ 上の非零因子である、
\item $M[f] = 0$、
\item すべての $n \geq 1$ に対して $M[f^n] = 0$ である、および
\item 写像 $M \to M_f$ は単射である。
\end{enumerate}
これらの同値な条件が成り立つとき、この節では
{\it $M$ は $f$-ねじれをもたない}ということにする。
このとき、$f^ix$（$x \in M$）の形の元からなる部分加群を
$f^iM \subset M_f$ と表す。もちろん $f^iM$ は $A$-加群として $M$ と同型である。
% P02-E1129: the frozen section repeatedly uses unhyphenated "f-torsion free" attributively.
$M^\bullet$ を $f$-ねじれをもたない $A$-加群の複体とし、
微分を $d^i : M^i \to M^{i + 1}$ とする。このとき
$\eta_fM^\bullet$ を、各項が
$$
(\eta_fM)^i = \{x \in f^iM^i \mid d^i(x) \in f^{i + 1}M^{i + 1}\}
$$
で微分が $d^i$ から誘導される複体として定義する。
$\eta_fM^\bullet$ もまた $f$-ねじれをもたない $A$-加群の複体である。
$a^\bullet : M^\bullet \to N^\bullet$ が $f$-ねじれをもたない
$A$-加群の複体の写像ならば、写像 $f^iM^i \to f^iN^i$ から誘導される複体の写像
$$
\eta_fa^\bullet : \eta_fM^\bullet \longrightarrow \eta_fN^\bullet
$$
を得る。これにより $f$-ねじれをもたない $A$-加群の複体の圏上の自己関手が
得られることを読者は確認できる。
$a^\bullet, b^\bullet : M^\bullet \to N^\bullet$ を
$f$-ねじれをもたない $A$-加群の複体の二つの写像とし、
$h = \{h^i : M^i \to N^{i - 1}\}$ を $a^\bullet$ と $b^\bullet$ の間の
ホモトピーとする。このとき $\eta_fh$ を、$x$ を $h^i(x)$ に写す写像族
$(\eta_fh)^i : (\eta_fM)^i \to (\eta_fN)^{i - 1}$ として定義する。
これは、$x \in f^iM^i$ ならば $h^i(x) \in f^iN^{i - 1}$ であり、
これは確かに $(\eta_fN)^{i - 1}$ に含まれるので意味をもつ。
$\eta_fh$ が $\eta_fa^\bullet$ と $\eta_fb^\bullet$ の間のホモトピーであることを
読者は確認できる。以上により、$f$-ねじれをもたない $A$-加群の加法圏の
ホモトピー圏（Derived Categories, Section
\ref{derived-section-homotopy}）上の関手
$$
\eta_f :
K(f\text{-torsion free }A\text{-modules})
\longrightarrow
K(f\text{-torsion free }A\text{-modules})
$$
を得る。$\eta_f$ が三角圏の完全関手となるような意味は存在しない。
Example \ref{more-algebra-example-eta-not-distinguished} を参照せよ。

\begin{example}
\label{more-algebra-example-eta-not-distinguished}
$A$ を環、$f \in A$ を非零因子とする。関手
$\eta_f : K(f\text{-torsion free }A\text{-modules})
\to K(f\text{-torsion free }A\text{-modules})$ を考える。
$M^\bullet$ を $f$-ねじれをもたない $A$-加群の複体とする。
$f$ 倍は同型 $\eta_f(M^\bullet[1]) \to (\eta_fM^\bullet)[1]$ を定めるので、
この意味で $\eta_f$ はシフトと両立する。しかし、図式
$$
\xymatrix{
A \ar[r]_f &
A \ar[r]_1 &
A \ar[r] &
0 \\
0 \ar[r] \ar[u] &
0 \ar[r] \ar[u] &
A \ar[r]^{-1} \ar[u]^f &
A \ar[u]
}
$$
を考える。各列を $f$-ねじれをもたない $A$-加群の複体とみなし、
上の加群を次数 $1$、その下の加群を次数 $0$ に置く。
するとこの図式は、Derived Categories, Section
\ref{derived-section-homotopy-triangulated} で与えられる三角構造をもつ
$K(f\text{-torsion free }A\text{-modules})$ における完全三角形を与える。
実際、第3の複体は最初の二つの複体の間の写像の錐である。
しかし各列に $\eta_f$ を適用すると
$$
\xymatrix{
fA \ar[r]_f &
fA \ar[r]_1 &
fA \ar[r] &
0 \\
0 \ar[r] \ar[u] &
0 \ar[r] \ar[u] &
A \ar[r]^{-1} \ar[u]^f &
A \ar[u]
}
$$
を得る。ところが第3の複体は非輪体であり、さらに零にホモトピックである。
従って、これが完全三角形ならば最初の矢印はホモトピー圏で同型でなければ
ならないが、$f$ が単元でない限りこれは成り立たない。
\end{example}

\begin{lemma}
\label{more-algebra-lemma-eta-first-property}
$A$ を環、$f \in A$ を非零因子とする。$M^\bullet$ を
$f$-ねじれをもたない $A$-加群の複体とする。$f^i$ 倍で与えられる標準同型
$$
f^i : H^i(M^\bullet)/H^i(M^\bullet)[f] \longrightarrow H^i(\eta_fM^\bullet)
$$
がある。
\end{lemma}

\begin{proof}
$\Ker(d^i : (\eta_fM)^i \to (\eta_fM)^{i + 1})$ は
$\Ker(d^i : f^iM^i \to f^iM^{i + 1}) =
f^i\Ker(d^i : M^i \to M^{i + 1})$ に等しいことに注意する。
% P02-E1130: frozen source line 27507 says "This we get" instead of "Thus we get".
従って、$z \in \Ker(d^i : M^i \to M^{i + 1})$ の類を
$f^iz$ の類に写すことにより、全射
$f^i : H^i(M^\bullet) \to H^i(\eta_fM^\bullet)$ を得る。
$H^i(\eta_fM^\bullet)$ で零類を得たとすると、ある $y \in M^{i - 1}$ に対して
$f^i z = d^{i - 1}(f^{i - 1}y)$ である。関与するすべての加群上で
$f$ は非零因子なので、これは $f z = d^{i - 1}(y)$ を意味する。
これはまさに $z$ の類が $f$-ねじれであることを意味し、望む結論を得る。
\end{proof}

\begin{lemma}
\label{more-algebra-lemma-eta-qis}
$A$ を環、$f \in A$ を非零因子とする。
$M^\bullet \to N^\bullet$ が $f$-ねじれをもたない $A$-加群の複体の
擬同型ならば、誘導される写像
$\eta_fM^\bullet \to \eta_fN^\bullet$ も擬同型である。
\end{lemma}

\begin{proof}
Lemma \ref{more-algebra-lemma-eta-first-property} の同型が複体の写像と両立するからである。
\end{proof}

\begin{lemma}
\label{more-algebra-lemma-Leta}
$A$ を環、$f \in A$ を非零因子とする。加法関手\footnote{この関手は
完全ではない、すなわち完全三角形を完全三角形に写さないことに注意せよ。
Example \ref{more-algebra-example-eta-not-distinguished} を参照せよ。}
$L\eta_f : D(A) \to D(A)$ が存在し、$M \in D(A)$ が
$f$-ねじれをもたない $A$-加群の複体 $M^\bullet$ で表されるならば
$L\eta_fM = \eta_fM^\bullet$ であり、射についても同様である。
\end{lemma}

\begin{proof}
% P02-E1131: frozen source line 27543 omits "by" in the Denote construction.
$f$-ねじれをもたない $A$-加群の充満部分圏を
$\mathcal{T} \subset \text{Mod}_A$ と表す。$K(A)$ の充満三角部分圏として
$K(\mathcal{T})$ の対応する包含
$$
K(\mathcal{T}) \quad\subset\quad K(\text{Mod}_A) = K(A)
$$
がある。$S \subset \text{Arrows}(K(\mathcal{T}))$ を擬同型全体とする。
Derived Categories, Lemma \ref{derived-lemma-localization-subcategory} を適用し、
写像
$$
S^{-1}K(\mathcal{T}) \longrightarrow D(A)
$$
が三角圏の同値であることを示す。この補題により、次を証明すれば十分である：
$A$-加群の複体 $M^\bullet$ が与えられたとき、$K^\bullet$ が
$f$-ねじれをもたない加群の複体であるような擬同型
$K^\bullet \to M^\bullet$ が存在する。Lemma
\ref{more-algebra-lemma-K-flat-resolution} により、複体 $K^\bullet$ が K-平坦であり
（これは用いない）、平坦 $A$-加群 $K^i$ からなるような擬同型
$K^\bullet \to M^\bullet$ を見つけられる。特に、望みどおりすべての $i$ に対して
$f$ は $K^i$ 上の非零因子である。

\medskip\noindent
以上の準備により $L\eta_f$ を定義できる。実際、この節の冒頭の議論により、
既に良定義な関手
$$
K(\mathcal{T}) \xrightarrow{\eta_f} K(\mathcal{T}) \to K(A) \to D(A)
$$
があり、Lemma \ref{more-algebra-lemma-eta-qis} によればこれは擬同型を擬同型に写す。
従って Categories, Lemma
\ref{categories-lemma-properties-left-localization} により、この関手は
$S^{-1}K(\mathcal{T}) = D(A)$ を経由する。
\end{proof}

\begin{remark}
\label{more-algebra-remark-eta-BZ}
$A$ を環、$f \in A$ を非零因子とする。$M^\bullet$ を
$f$-ねじれをもたない $A$-加群の複体とする。
すべての $i$ に対して $\overline{M}^i = M^i/fM^i$ と置く。
% P02-E1132: frozen source line 27582 omits "by" in the Denote construction.
複体 $\overline{M}^\bullet = M^\bullet \otimes_A A/fA$ の微分に関する
境界とコサイクルを $B^i \subset Z^i \subset \overline{M}^i$ と表す。
完全な行をもつ可換図式
$$
\xymatrix{
0 \ar[r] &
B^{i + 1} \ar[r] \ar@{=}[d] &
B^{i + 1} \oplus B^i \ar[r] \ar[d]^{s, s'} &
B^i \ar[r] \ar[d] & 0 \\
0 \ar[r] &
B^{i + 1} \ar[r]^-s &
(\eta_fM)^i /f(\eta_fM)^i \ar[r]^-t &
Z^i \ar[r] & 0
}
$$
が存在すると主張する。写像の構成は次のとおりである。
\begin{enumerate}
\item $x \in (\eta_fM)^i$ ならば、$x = f^ix'$ かつ
$d^i(x') = 0$ は $\overline{M}^{i + 1}$ において成り立つ。従って $x$ を $x'$ の類に
写すことにより写像 $t$ を定義できる。
\item $y \in M^{i + 1}$ の類 $\overline{y}$ が
$B^{i + 1} \subset \overline{M}^{i + 1}$ に属するならば、
ある $y' \in M^{i + 1}$ および $x \in M^i$ に対して
$y = fy' + d^i(x)$ と書ける。従って $\overline{y}$ を
$(\eta_fM)^i /f(\eta_fM)^i$ における $f^{i + 1}x$ の類に写す写像
$s$ を定義できる。これが良定義であることの確認は省略する。
\item $x \in M^i$ の類 $\overline{x}$ が
$B^i \subset \overline{M}^i$ に属するならば、ある $x' \in M^i$ および
$z \in M^{i - 1}$ に対して $x = fx' + d^{i - 1}(z)$ と書ける。
$\overline{x}$ を $(\eta_fM)^i/f(\eta_fM)^i$ における
$f^i d^{i - 1}(z)$ の類に写すことで写像 $s'$ を定義する。
これは良定義である。実際、$fx' + d^{i - 1}(z) = 0$ ならば
$f^ix'$ は $(\eta_fM)^i$ に属し、従って
$f^id^{i - 1}(z)$ は $f(\eta_fM)^i$ に属する。
\end{enumerate}
上の図式の下段が加群の短完全列であることの確認は省略する。
これらの構成から直ちに、可換図式
$$
\xymatrix{
B^{i + 1} \oplus B^i \ar[d]^{s, s'} \ar[r] &
B^{i + 2} \oplus B^{i + 1} \ar[d]^{s, s'} \\
(\eta_fM)^i /f(\eta_fM)^i \ar[r] &
(\eta_fM)^{i + 1} /f(\eta_fM)^{i + 1}
}
$$
を得る。ここで上の水平矢印は、始域と終域の直和因子
$B^{i + 1}$ の同一視で与えられる。言い換えると、
$\eta_fM^\bullet / f(\eta_fM^\bullet) =
\eta_fM^\bullet \otimes_A A/fA$ の非輪部分複体を見つけ、その部分複体による
商は、各項 $Z^i/B^i$ が複体
$\overline{M}^\bullet = M^\bullet \otimes_A A/fA$ の
コホモロジー加群である複体となる。
\end{remark}


\noindent
Remark \ref{more-algebra-remark-eta-BZ} で観察した現象をより標準的な形で説明するため、
Bockstein 作用素を構成する。$A$ を環、$f \in A$ を非零因子とする。
$M^\bullet$ を $f$-ねじれをもたない $A$-加群の複体とする。
すべての $i \in \mathbf{Z}$ に対して、可換図式
（テンソル積は $A$ 上）
$$
\xymatrix{
0 \ar[r] &
M^\bullet \otimes f^{i + 1}A \ar[r] \ar[d] &
M^\bullet \otimes f^iA \ar[r] \ar[d] &
M^\bullet \otimes f^iA/f^{i + 1}A \ar[r] \ar@{=}[d] &
0 \\
0 \ar[r] &
M^\bullet \otimes f^{i + 1}A/f^{i + 2}A \ar[r] &
M^\bullet \otimes f^iA/f^{i + 2}A \ar[r] &
M^\bullet \otimes f^iA/f^{i + 1}A \ar[r] &
0
}
$$
があり、その各行は複体の短完全列である。もちろん、異なる $i$ に対する
これらの短完全列は、$f$ の適切な冪を掛けることにより互いに同型となる。
特に、下段の列の長完全コホモロジー列は、すべての
$i \in \mathbf{Z}$ に対して {\it Bockstein 作用素}
$$
\beta = \beta^i : H^i(M^\bullet \otimes f^iA/f^{i + 1}A) \to
H^{i + 1}(M^\bullet \otimes f^{i + 1}A/f^{i + 2}A)
$$
を定める。後で用いるため、上の可換図式により Bockstein 作用素が
\begin{equation}
\label{more-algebra-equation-factorization-bockstein}
\vcenter{
\xymatrix{
H^i(M^\bullet \otimes f^iA/f^{i + 1}A)
\ar[r]_\delta \ar[rd]_\beta &
H^{i + 1}(M^\bullet \otimes f^{i + 1}A) \ar[d] \\
&
H^{i + 1}(M^\bullet \otimes f^{i + 1}A/f^{i + 2}A)
}
}
\end{equation}
と分解することを記しておく。ここで $\delta$ は上の可換図式の上段から生じる
境界作用素である。複体
\begin{equation}
\label{more-algebra-equation-complex-bocksteins}
H^\bullet(M^\bullet/f) =
\left[
\begin{matrix}
\ldots \\
\downarrow \\
H^{i - 1}(M^\bullet \otimes f^{i - 1}A/f^iA) \\
\downarrow \beta \\
H^i(M^\bullet \otimes f^iA/f^{i + 1}A) \\
\downarrow \beta \\
H^{i + 1}(M^\bullet \otimes f^{i + 1}A/f^{i + 2}A) \\
\downarrow \\
\ldots
\end{matrix}
\right]
\end{equation}
が得られること、すなわち
$\beta \circ \beta = 0$ であることを示そう\footnote{別の方法として、
$\beta$ は部分複体 $f^iM^\bullet$ によるフィルターを入れた複体
$(M^\bullet)_f$ のスペクトル系列の微分として現れる、と論じてもよい。
さらに別の、より強い結果を証明する方法は、まず $M^\bullet = A$ の場合を
考えることである。この場合、短完全列
$0 \to f^{i + 1}A/f^{i + 2}A \to f^iA/f^{i + 2}A \to f^iA/f^{i + 1}A \to 0$
は $D(A)$ における写像
$\beta^i : f^iA/f^{i + 1}A \to f^{i + 1}A/f^{i + 2}A[1]$ を定める。
このとき（本文と同様に論じて）合成
$f^iA/f^{i + 1}A \to f^{i + 1}A/f^{i + 2}A[1] \to f^{i + 2}A/f^{i + 3}A[2]$
が $D(A)$ で零であることを計算できる。$M^\bullet$ の各項が
$f$-ねじれをもたないという仮定により
$M^\bullet \otimes f^iA/f^{i + 1}A =
M^\bullet \otimes^\mathbf{L} f^iA/f^{i + 1}A$ なので、合成
% P02-E1133: frozen source line 27715 misspells "composition" as "compostion".
$$
(M^\bullet \otimes f^iA/f^{i + 1}A)
\to
(M^\bullet \otimes f^{i + 1}A/f^{i + 2}A)[1]
\to
(M^\bullet \otimes f^{i + 2}A/f^{i + 3}A)[2]
$$
も $D(A)$ で零である。}。実際、分解
(\ref{more-algebra-equation-factorization-bockstein}) を用いると、
$$
H^{i + 1}(M^\bullet \otimes f^{i + 1}A)
\to
H^{i + 1}(M^\bullet \otimes f^{i + 1}A/f^{i + 2}A)
\xrightarrow{\beta^{i + 1}}
H^{i + 2}(M^\bullet \otimes f^{i + 2}A/f^{i + 3}A)
$$
が零であることを示せば十分である。これは、$\beta^{i + 1}$ の核が
$H^{i + 1}(M^\bullet \otimes f^{i + 1}A/f^{i + 3}A)$ に持ち上げられる
コホモロジー類からなり、最初の写像の像に属する類は確かに持ち上げられるからである。

\begin{lemma}
\label{more-algebra-lemma-eta-second-property}
$A$ を環、$f \in A$ を非零因子とする。$M^\bullet$ を
$f$-ねじれをもたない $A$-加群の複体とする。標準的な複体の写像
$$
\eta_fM^\bullet \otimes_A A/fA
\longrightarrow
H^\bullet(M^\bullet/f)
$$
があり、これは擬同型である。ここで右辺は複体
(\ref{more-algebra-equation-complex-bocksteins}) である。
\end{lemma}

\begin{proof}
$x \in (\eta_fM)^i$ とする。このとき
$x = f^ix' \in f^iM$ かつ
$d^i(x) = f^{i + 1}y \in f^{i + 1}M^{i + 1}$ である。
従って $d^i$ は $x' \otimes f^i$ を
$M^{i + 1} \otimes f^iA/f^{i + 1}A$ で零に写す。この証明では
すべてのテンソル積を $A$ 上でとる。
そこで $x$ を $H^i(M^\bullet \otimes f^iA/f^{i + 1}A)$ における
$x' \otimes f^i$ の類に写すことができる。この規則が $A/fA$-加群の写像
$$
(\eta_fM)^i \otimes A/fA
\longrightarrow
H^i(M^\bullet \otimes f^iA/f^{i + 1}A)
$$
を定めることは明らかである。上の状況で、
$x' \otimes f^i$ を $M^i \otimes f^iA/f^{i + 2}A$ の元とみなすと、
その微分は $d^i(x' \otimes f^i) = y \otimes f^{i + 1}$ である。
上の $\beta$ の構成により
$\beta(x' \otimes f^i) = y \otimes f^{i + 1}$ となるので、
これらの写像は微分と両立する、すなわち複体の写像をなす。

\medskip\noindent
証明を終えるため、前段落で与えた構成が Remark \ref{more-algebra-remark-eta-BZ} で論じた写像
$(\eta_fM)^i \otimes A/fA \to Z^i/B^i$ と一致することに注意する。
これらの写像の核が $\eta_fM^\bullet \otimes A/fA$ の非輪部分複体であることを
既に見たので、補題が従う。
\end{proof}

\begin{lemma}
\label{more-algebra-lemma-vanishing-beta}
$A$ を環、$f \in A$ を非零因子とする。$M^\bullet$ を
$f$-ねじれをもたない $A$-加群の複体とする。
$i \in \mathbf{Z}$ に対して次は同値である。
\begin{enumerate}
\item $\Ker(d^i \bmod f^2)$ は $\Ker(d^i \bmod f)$ の上に全射する、
\item $\beta : H^i(M^\bullet \otimes_A f^iA/f^{i + 1}A) \to
H^{i + 1}(M^\bullet \otimes_A f^{i + 1}A/f^{i + 2}A)$ は零である。
\end{enumerate}
条件 $H^{i + 1}(M^\bullet)[f] = 0$ はこれらの同値な条件を含意する。
\end{lemma}

\begin{proof}
(1) と (2) の同値性は、$\beta$ が、複体の短完全列
$0 \to fM^\bullet/f^2M^\bullet \to M^\bullet/f^2M^\bullet \to
M^\bullet/fM^\bullet \to 0$ と同型な複体の短完全列の
コホモロジー上の境界写像として定義されることから従う。
$\beta \not = 0$ ならば、分解
(\ref{more-algebra-equation-factorization-bockstein}) により
$H^{i + 1}(M^\bullet)[f] \not = 0$ である。
\end{proof}

\begin{lemma}
\label{more-algebra-lemma-eta-vanishing-beta}
$A$ を環、$f \in A$ を非零因子とする。$M^\bullet$ を
$f$-ねじれをもたない $A$-加群の複体とする。
$\Ker(d^i \bmod f^2)$ が $\Ker(d^i \bmod f)$ の上に全射するならば、
標準写像
$$
(1, d^i) :
(\eta_fM)^i / f(\eta_fM)^i \longrightarrow
f^iM^i/f^{i + 1}M^i \oplus f^{i + 1}M^{i + 1}/f^{i + 2}M^{i + 1}
$$
は左辺を右辺の部分加群の直和と同一視する。
\end{lemma}

\begin{proof}
Remark \ref{more-algebra-remark-eta-BZ} の記法を用いて、写像
$t^{-1} : Z^i \to (\eta_fM)^i / f(\eta_fM)^i$ を定義する。
すなわち、$d^i(x) = f^2y$ を満たす $x \in M^i$ に対し、
$Z^i$ における $x$ の類を
$(\eta_fM)^i / f(\eta_fM)^i$ における $f^ix$ の類に写す。
これが良定義であることの確認は省略する。この演算の定義域が
$Z^i$ の全体であることを補題の仮定がちょうど意味している。
このとき $t \circ t^{-1} = \text{id}_{Z^i}$ である。
従って $t^{-1}$ は Remark \ref{more-algebra-remark-eta-BZ} の短完全列の分裂を定め、
その結果得られる直和分解
$$
(\eta_fM)^i / f(\eta_fM)^i = Z^i \oplus B^{i + 1}
$$
は補題に表示した写像と両立する。
\end{proof}

\begin{lemma}
\label{more-algebra-lemma-eta-third-property}
$A$ を環、$f, g \in A$ を非零因子とする。$M^\bullet$ を
すべての $M^i$ 上で $fg$ が非零因子となる $A$-加群の複体とする。
このとき $\eta_f\eta_gM^\bullet = \eta_{fg}M^\bullet$ である。
\end{lemma}

\begin{proof}
この主張は、次数 $i$ において局所化
$M^i_{fg} = (M^i_g)_f$ の同じ部分加群を得ることを意味する。
詳細は省略する。
\end{proof}

\begin{lemma}
\label{more-algebra-lemma-eta-flat-base-change}
$A$ を環、$f \in A$ を非零因子とする。$A \to B$ を平坦な環写像とし、
% P02-E1134: frozen source line 27854 omits "be" before "the image of f".
$g \in B$ を $f$ の像とする。$M^\bullet$ を
$f$-ねじれをもたない $A$-加群の複体とする。
このとき $g$ は非零因子であり、$M^\bullet \otimes_A B$ は
$g$-ねじれをもたない加群の複体であり、
$\eta_fM^\bullet \otimes_A B = \eta_g(M^\bullet \otimes_A B)$ である。
\end{lemma}

\begin{proof}
省略する。
\end{proof}

\section{完全複体と eta 作用素}
\label{more-algebra-section-perfect-eta}

\noindent
この節では、More on Flatness, Section
\ref{flat-section-blowup-complexes-III} を参照して、Macpherson のグラフ構成の
我々の版を準備するための代数を行う。Section \ref{more-algebra-section-eta} で導入した
$\eta_f$ 作用素を用いる。

\medskip\noindent
$A$ を環、$f \in A$ を非零因子とする。$M^\bullet$ を有限自由
$A$-加群からなる有界複体とする。各 $i$ に対して $r_i$ を $M^i$ の階数とし、
$$
I_i(M^\bullet, f) = \text{ideal generated by the }
r_i \times r_i\text{-minors of }
(f, d^i) : M^i \to M^i \oplus M^{i + 1}
$$
と置く。$f^{r_i} \in I_i(M^\bullet, f)$ であることに注意する。

\begin{lemma}
\label{more-algebra-lemma-ideal-well-defined}
$A$ を環、$f \in A$ を非零因子とする。$M^\bullet$ と $N^\bullet$ を、
$D(A)$ の同じ対象を表す有限自由 $A$-加群の二つの有界複体とする。
整数 $n, m \geq 0$ が
$$
m + \sum\nolimits_{j \geq i} (-1)^{j - i}rk(M^j) =
n + \sum\nolimits_{j \geq i} (-1)^{j - i}rk(N^j)
$$
を満たすならば、$A$ のイデアルとして
$$
f^m I_i(M^\bullet, f) = f^n I_i(N^\bullet, f)
$$
である。
\end{lemma}

\begin{proof}
$A$ のすべての素イデアルで局所化した後に等式を証明すれば十分である。
% P02-E1135: frozen source line 27911 omits "by" before "an induction argument".
従って Lemma \ref{more-algebra-lemma-compare-representatives-perfect} と、
詳細を省略する帰納法により、ある自明な複体 $Q^\bullet$ に対して
$N^\bullet = M^\bullet \oplus Q^\bullet$ と仮定してよい。すなわち
$$
Q^\bullet = \ldots \to 0 \to A \xrightarrow{1} A \to 0 \to \ldots
$$
である。ここで
% P02-E1136: frozen source line 27917 says singular "A is placed" although two copies occupy degrees j and j+1.
$A$ は次数 $j$ と $j + 1$ に置かれている。
$j \not = i - 1, i, i + 1$ ならば、明らかに
$I_i(M^\bullet, f) = I_i(N^\bullet, f)$ かつ $m = n$ であり、望む等式を得る。
$j = i + 1$ ならば、写像
$$
(f, d^i) : M^i \to M^i \oplus M^{i + 1}
\quad\text{and}\quad
(f, d^i, 0) : M^i \to M^i \oplus M^{i + 1} \oplus A
$$
は同じ非零小行列式をもつ。従ってこの場合にも
$I_i(M^\bullet, f) = I_i(N^\bullet, f)$ かつ $m = n$ である。
$j = i$ ならば、$I_i(M^\bullet, f)$ は
$$
(f, d^i) : M^i \to M^i \oplus M^{i + 1}
$$
の $r_i \times r_i$ 小行列式で生成され、$I_i(N^\bullet, f)$ は
$$
(f \oplus f, d^i \oplus 1) :
(M^i \oplus A) \to (M^i \oplus A) \oplus (M^{i + 1} \oplus A)
$$
の $(r_i + 1) \times (r_i + 1)$ 小行列式で生成される。
適切に座標を選ぶと、第2の写像の行列はブロック形
$$
T =
\left(
\begin{matrix}
T_1 & 0 \\
0 & T_2
\end{matrix}
\right),
\quad
T_1 = \text{matrix of first map},
\quad
T_2 =
\left(
\begin{matrix}
f \\
1
\end{matrix}
\right)
$$
である。Lemma \ref{more-algebra-lemma-ideals-generated-by-minors} の記法を用いると、
$I_0(T_2) = A$、$I_1(T_2) = A$、$p \geq 2$ に対して $I_p(T_2) = 0$ なので、
$I_{r_i + 1}(T) = I_{r_i + 1}(T_1) + I_{r_i}(T_1) =
I_{r_i}(T_1)$ である。これは
$I_i(M^\bullet, f) = I_i(N^\bullet, f)$ を意味する。
また $m = n$ なので $j = i$ の場合が完了する。
最後に $j = i - 1$ とする。このとき $m = n + 1$ なので、
$fI_i(M^\bullet, f) = I_i(N^\bullet, f)$ を示す必要がある。
この場合 $I_i(M^\bullet, f)$ は
$$
(f, d^i) : M^i \to M^i \oplus M^{i + 1}
$$
の $r_i \times r_i$ 小行列式で生成され、$I_i(N^\bullet, f)$ は
$$
(f \oplus f, d^i) : (M^i \oplus A) \to (M^i \oplus A) \oplus M^{i + 1}
$$
の $(r_i + 1) \times (r_i + 1)$ 小行列式で生成される。
適切に座標を選ぶと、第2の写像の行列はブロック形
$$
T =
\left(
\begin{matrix}
T_1 & 0 \\
0 & T_2
\end{matrix}
\right),
\quad
T_1 = \text{matrix of first map},
\quad
T_2 =
\left(
\begin{matrix}
f
\end{matrix}
\right)
$$
である。上と同様に論じると、実際
$fI_i(M^\bullet, f) = I_i(N^\bullet, f)$ を得る。
\end{proof}

\begin{lemma}
\label{more-algebra-lemma-eta-change-unit}
$f \in A$ を環 $A$ の非零因子、$u \in A$ を単元とする。
$M^\bullet$ を有限自由 $A$-加群の有界複体とする。このとき
$I_i(M^\bullet, f) = I_i(M^\bullet, uf)$ である。
\end{lemma}

\begin{proof}
省略する。
\end{proof}

\begin{lemma}
\label{more-algebra-lemma-eta-base-change-pre}
$A \to B$ を環写像、$f \in A$ を非零因子とする。
$M^\bullet$ を有限自由 $A$-加群の有界複体とする。
$f$ が $B$ の非零因子 $g$ に写ると仮定する。このとき
$I_i(M^\bullet, f)B = I_i(M^\bullet \otimes_A B, g)$ である。
\end{lemma}

\begin{proof}
$(f, d^i) : M^i \to M^i \oplus M^{i + 1}$ の小行列式は、対応する
$(g, d^i) : M^i \otimes_A B \to M^i \otimes_A B \oplus M^{i + 1} \otimes_A B$
の小行列式に写る。
\end{proof}

\begin{lemma}
\label{more-algebra-lemma-eta-cohomology-free}
$A$ を環、$\mathfrak p \subset A$ を素イデアル、
$f \in A$ を非零因子とする。$M^\bullet$ を有限自由 $A$-加群の有界複体とする。
すべての $i$ に対して $H^i(M^\bullet)_\mathfrak p$ が自由ならば、
すべての $i$ に対して $I_i(M^\bullet, f)_\mathfrak p$ は単項イデアルであり、
実際 $f$ の冪で生成される。
\end{lemma}

\begin{proof}
Lemma \ref{more-algebra-lemma-eta-base-change-pre} により、$A$ は極大イデアル
$\mathfrak p$ をもつ局所環であると仮定してよい。Lemma
\ref{more-algebra-lemma-ideal-well-defined} により、$M^\bullet$ を擬同型な複体で
置き換えてもよい。コホモロジー加群が自由であるという仮定から、
$M^\bullet$ は、次数 $i$ の項が $H^i(M^\bullet)$ で微分が零である複体と
擬同型である。例えば Derived Categories, Lemma
\ref{derived-lemma-ext-2-zero-pre} を参照せよ。
言い換えると、複体 $M^\bullet$ のすべての微分が零であると仮定してよい。
この場合 $I_i(M^\bullet, f) = (f^{r_i})$ が単項であることは明らかである。
\end{proof}

\begin{lemma}
\label{more-algebra-lemma-eta-locally-free}
$A$ を環、$f \in A$ を非零因子とする。$M^\bullet$ を有限自由
$A$-加群の有界複体とする。$I_i(M^\bullet, f)$ が単項イデアルであると仮定する。
このとき $(\eta_fM)^i$ は階数 $r_i$ の局所自由加群であり、写像
$(1, d^i) : (\eta_fM)^i \to f^iM^i \oplus f^{i + 1}M^{i + 1}$
は直和因子の包含である。
\end{lemma}

\begin{proof}
$I_i(M^\bullet, f)$ の生成元 $g$ を選ぶ。
$f^{r_i} \in I_i(M^\bullet, f)$ なので、$g$ は $f$ のある冪を割る。
特に $g$ は $A$ の非零因子である。写像
$(f, d^i) : M^i \to M^i \oplus M^{i + 1}$ の
$r_i \times r_i$ 小行列式はイデアル $I_i(M^\bullet, f)$ を生成し、
$(f, d^i)$ の $(r_i + 1) \times (r_i + 1)$ 小行列式は零である。
これは $f$ で局所化した後に確認でき、その場合写像の階数は $r_i$ に等しい。
全射
$$
M^i \oplus M^{i + 1}
\longrightarrow
Q = \Coker(f, d^i)/g\text{-torsion}
$$
を考える。Lemma \ref{more-algebra-lemma-fitting-ideals-and-pd1} により、
加群 $Q$ は階数 $r_{i + 1}$ の有限局所自由加群である。
従って $Q$ は $f$-ねじれをもたず、$(f, d^i)$ の余核を
$f$-冪ねじれで割ったものも $Q$ である。

\medskip\noindent
有限自由 $A$-加群の複体
$$
0 \to f^{i + 1}M^i \xrightarrow{1, d^i} f^iM^i \oplus f^{i + 1}M^{i + 1}
\xrightarrow{d^i, -1} f^iM^{i + 1} \to 0
$$
を考える。これは $f$ で局所化した後に分裂完全となる。
写像 $(1, d^i) : f^{i + 1}M^i \to f^iM^i \oplus f^{i + 1}M^{i + 1}$ は、
上で調べた写像 $(f, d^i) : M^i \to M^i \oplus M^{i + 1}$ と同型である。
従って像
$$
Q' = \Im(f^iM^i \oplus f^{i + 1}M^{i + 1} \xrightarrow{d^i, -1} f^iM^{i + 1})
$$
は $Q$ と同型であり、特に射影的である。他方、Section \ref{more-algebra-section-eta} の
$\eta_f$ の構成により、単射
$(1, d^i) : (\eta_fM)^i \to f^iM^i \oplus f^{i + 1}M^{i + 1}$ の像は
$(d^i, -1)$ の核である。従って同型
$(\eta_fM)^i \oplus Q' = f^iM^i \oplus f^{i + 1}M^{i + 1}$ を得る。
% P02-E1137: frozen source line 28099 writes eta_fM^i instead of the degree-i term (eta_fM)^i.
これにより確かに $\eta_fM^i$ は階数 $r_i$ の有限局所自由加群であり、
$(1, d^i)$ は直和因子の包含である。
\end{proof}

\begin{lemma}
\label{more-algebra-lemma-eta-base-change}
$A \to B$ を環写像、$f \in A$ を非零因子とする。
$M^\bullet$ を有限自由 $A$-加群の有界複体とする。$f$ は $B$ の非零因子
$g$ に写り、すべての $i \in \mathbf{Z}$ に対して
$I_i(M^\bullet, f)$ が単項イデアルであると仮定する。
このとき標準同型
$\eta_fM^\bullet \otimes_A B = \eta_g(M^\bullet \otimes_A B)$ がある。
\end{lemma}

\begin{proof}
$N^i = M^i \otimes_A B$ と置く。$(N^i)_g$ の部分加群として
$f^iM^i \otimes_A B = g^iN^i$ であることに注意する。写像
% P02-E1138: frozen source lines 28118-28120 use tensor signs where the surrounding direct-sum inclusions require op⊕.
$$
(\eta_fM)^i \otimes_A B \to g^iN^i \otimes g^{i + 1}N^{i + 1}
\quad\text{and}\quad
(\eta_gN)^i \to g^iN^i \otimes g^{i + 1}N^{i + 1}
$$
は Lemma \ref{more-algebra-lemma-eta-locally-free} により直和因子の包含である。
それらの像は $g$ で局所化した後に一致するので、結論を得る。
\end{proof}


\begin{lemma}
\label{more-algebra-lemma-ideal-direct-summand}
$A$ を環とし、$M$, $N_1$, $N_2$ を有限射影 $A$-加群とする。
$s : M \to N_1 \oplus N_2$ を分裂単射とする。次の性質をもつ有限生成イデアル
$J \subset A$ が存在する：環写像 $A \to B$ が $A/J$ を経由することと、
$s \otimes \text{id}_B$ が $M \otimes_A B$ を
$N_1 \otimes_A B \oplus N_2 \otimes_A B$ の部分加群の直和と同一視することは
同値である。
\end{lemma}

\begin{proof}
$s$ の分裂 $\pi : N_1 \oplus N_2 \to M$ を選ぶ。
% P02-E1139: frozen source line 28138 omits "by" in the Denote construction.
$N_i$ への射影を $q_i : N_1 \oplus N_2 \to N_1 \oplus N_2$ と表し、
$p_i = \pi \circ q_i \circ s$ と置く。
$p_1 + p_2 = \text{id}_M$ であることに注意する。
$M$ が $N_1 \oplus N_2$ の部分加群の直和であることと、
$p_1$ と $p_2$ が直交射影子であることは同値であると主張する。
従って $J$ は、
$p_1 \circ p_1 - p_1$, $p_2 \circ p_2 - p_2$, $p_1 \circ p_2$, および
$p_2 \circ p_1$ が $J \otimes_A \text{End}_A(M)$ に含まれるような
$A$ の最小のイデアルである。いくつかの詳細は省略する。
\end{proof}

\noindent
$A$ を環、$f \in A$ を非零因子とする。$M^\bullet$ を有限自由
$A$-加群の有界複体とする。すべての $i \in \mathbf{Z}$ に対して
イデアル $I_i(M^\bullet, f)$ が単項であると仮定する。このとき Lemma
\ref{more-algebra-lemma-eta-locally-free} により写像
$$
(1, d^i) : (\eta_fM)^i / f(\eta_fM)^i \longrightarrow
f^iM^i/f^{i + 1}M^i \oplus f^{i + 1}M^{i + 1}/f^{i + 2}M^{i + 1}
$$
は分裂単射である。
% P02-E1140: frozen source line 28159 omits "by" in the Denote construction.
上に表示した分裂単射 $(1, d^i)$ に対応する Lemma
\ref{more-algebra-lemma-ideal-direct-summand} の有限生成イデアルを
$J_i(M^\bullet, f) \subset A/fA$ と表す。

\begin{lemma}
\label{more-algebra-lemma-eta-vanishing-beta-plus-pre}
$A$ を環、$f \in A$ を非零因子とする。$M^\bullet$ と $N^\bullet$ を
$D(A)$ の同じ対象を表す有限自由 $A$-加群の二つの有界複体とする。
すべての $i \in \mathbf{Z}$ に対して $I_i(M^\bullet, f)$ が
単項イデアルであると仮定する。このとき $A/fA$ のイデアルとして
$J_i(M^\bullet, f) = J_i(N^\bullet, f)$ である。
\end{lemma}

\begin{proof}
% P02-E1141: frozen source lines 28174-28176 repeats M where N is needed to make J_i(N,f) defined.
$I_i(M^\bullet, f)$ が単項イデアルであることから、Lemma
\ref{more-algebra-lemma-ideal-well-defined} により $I_i(M^\bullet, f)$ が単項イデアルであり、
従って主張が意味をもつことに注意する。
Lemma \ref{more-algebra-lemma-ideal-well-defined} の証明と同様、ある自明な複体
$Q^\bullet$ に対して $N^\bullet = M^\bullet \oplus Q^\bullet$ と仮定してよい。
すなわち
$$
Q^\bullet = \ldots \to 0 \to A \xrightarrow{1} A \to 0 \to \ldots
$$
であり、$A$ は次数 $j$ と $j + 1$ に置かれている。
$\eta_f$ は直和と両立するので、写像
$$
(1, d^i) : (\eta_fN)^i / f(\eta_fN)^i \longrightarrow
f^iN^i/f^{i + 1}N^i \oplus f^{i + 1}N^{i + 1}/f^{i + 2}N^{i + 1}
$$
は $M^\bullet$ に対応する写像と $Q^\bullet$ に対応する写像の直和である。
問題のイデアルを定める普遍性により
$J_i(N^\bullet, f) = J_i(M^\bullet, f) + J_i(Q^\bullet, f)$ を得る。
従って、すべての $i$ に対して $J_i(Q^\bullet, f) = 0$ を示せば十分である。
これは省略する計算である。
\end{proof}

\begin{lemma}
\label{more-algebra-lemma-eta-vanishing-beta-plus}
$A$ を環、$f \in A$ を非零因子とする。$M^\bullet$ を有限自由
$A$-加群の有界複体とする。すべての $i \in \mathbf{Z}$ に対して
$I_i(M^\bullet, f)$ が単項イデアルであると仮定する。
$A/fA$ のイデアル
$J(M^\bullet, f) = \sum_i J_i(M^\bullet, f)$ を考える。
素イデアルの集合
\begin{align*}
E
& =
\{f \in \mathfrak p \subset A \mid
\Ker(d^i \bmod f^2)_\mathfrak p \text{ surjects onto }
\Ker(d^i \bmod f)_\mathfrak p \text{ for all }i \in \mathbf{Z}\} \\
& =
\{f \in \mathfrak p \subset A \mid
\text{the localizations }\beta_\mathfrak p
\text{ of the Bockstein operators are zero}\}
\end{align*}
を考える。このとき次が成り立つ。
\begin{enumerate}
\item $J(M^\bullet, f)$ は有限生成である、
\item $A/fA \to C = (A/fA)/J(M^\bullet, f)$ は有限表示の全射である、
\item $\mathfrak p \in E$ に対して $J(M^\bullet, f)_\mathfrak p = 0$ である、
\item $f \in \mathfrak p$ かつすべての $i \in \mathbf{Z}$ に対して
$H^i(M^\bullet)_\mathfrak p$ が自由ならば $\mathfrak p \in E$ である、および
\item $\eta_f M^\bullet \otimes_A C$ のコホモロジー加群は
有限局所自由 $C$-加群である。
\end{enumerate}
\end{lemma}

\begin{proof}
$E$ の定義における等式は Lemma \ref{more-algebra-lemma-vanishing-beta} から従い、
さらに同補題の最後の主張は (4) を含意する。

\medskip\noindent
(1) は、イデアル $J_i(M^\bullet, f)$ が有限生成であり、$M^\bullet$ が有界なので
ほとんどすべての $i$ に対して $J_i(M^\bullet, f)$ が零であることから従う。
(2) は (1) の言い換えにすぎない。

\medskip\noindent
(3) の証明。Lemma \ref{more-algebra-lemma-eta-locally-free} により、
すべての $i$ に対して $(\eta_fM)^i$ は階数 $r_i$ の有限局所自由加群である。
写像
$$
(1, d^i) :
(\eta_fM)^i / f(\eta_fM)^i
\longrightarrow
f^iM^i/f^{i + 1}M^i \oplus f^{i + 1}M^{i + 1}/f^{i + 2}M^{i + 1}
$$
を考える。$\mathfrak p \in E$ を選ぶ。Lemma
\ref{more-algebra-lemma-eta-vanishing-beta} と加群 $(\eta_fM)^i$ の局所自由性により、
上の矢印 $(1, d^i)$ と両立する形で
$$
\left((\eta_fM)^i / f(\eta_fM)^i\right)_\mathfrak p =
(A/fA)_\mathfrak p^{\oplus m_i} \oplus (A/fA)_\mathfrak p^{\oplus n_i}
$$
と書ける。イデアル $J_i(M^\bullet, f)$ の普遍性により
$J_i(M^\bullet, f)_\mathfrak p = 0$ を得る。
% P02-E1142: frozen source line 28259 concludes with an undefined ideal I.
従って $\mathfrak p \in E$ に対して
$I_\mathfrak p = fA_\mathfrak p$ である。

\medskip\noindent



\medskip\noindent
(5) の証明。$\eta_fM^\bullet$ 上の微分は可換図式に入る。
% P02-E1143: frozen source lines 28278 and 28297 use M^i where degree consistency requires M^{i+1}.
$$
\xymatrix{
(\eta_fM)^i \ar[d] \ar[r] &
f^iM^i \oplus f^{i + 1}M^{i + 1}
\ar[d]^{\left(
\begin{matrix}
0 & 1 \\
0 & 0
\end{matrix}
\right)} \\
(\eta_fM)^{i + 1} \ar[r] &
f^{i + 1}M^i \oplus f^{i + 2}M^{i + 2}
}
$$
構成により、$C$ とテンソル積をとった後、左の加群は右の直和因子の
直和となる。次の図を描ける。
$$
\xymatrix{
(\eta_fM)^i \otimes_A C \ar[d] \ar@{=}[r] &
K^i \oplus L^i \ar[r] \ar[d] &
f^iM^i \otimes_A C \oplus f^{i + 1}M^{i + 1} \otimes_A C
\ar[d]^{\left(
\begin{matrix}
0 & 1 \\
0 & 0
\end{matrix}
\right)} \\
(\eta_fM)^{i + 1} \otimes_A C \ar@{=}[r] &
K^{i + 1} \oplus L^{i + 1} \ar[r] &
f^{i + 1}M^i \otimes_A C \oplus f^{i + 2}M^{i + 2} \otimes_A C
}
$$
ここで水平矢印は直和分解および直和因子の包含と両立する。
従って微分は $L^i$ を $K^{i + 1}$ の直和因子と同一視する。
% P02-E1144: frozen source lines 28303-28305 shifts the claimed degree-i cohomology quotient by one.
よって $\eta_fM^\bullet \otimes_A C$ の次数 $i$ のコホモロジーは
加群 $K^{i + 1}/L^i$ であり、これは望みどおり有限射影的である。
\end{proof}


\section{複体の極限をとること}
\label{more-algebra-section-limits}

\noindent
この節では、複体の「形式変形」とその極限をとるときに何が起こるかを論じる。
次の二つの場合を考える。
\begin{enumerate}
\item 遷移写像が局所冪零な核をもつ全射である環の逆系について
極限 $A = \lim A_n$ があり、対象 $K_n \in D(A_n)$ が
$K_n = K_{n + 1} \otimes_{A_{n + 1}}^\mathbf{L} A_n$ という意味で整合する場合、または
\item 環 $A$、イデアル $I$、対象 $K_n \in D(A/I^n)$ があり、
$K_n = K_{n + 1} \otimes_{A/I^{n + 1}}^\mathbf{L} A/I^n$
という意味で整合する場合。
\end{enumerate}
追加の仮定のもとで、$K = R\lim K_n$ が
$K_n = K \otimes_A^\mathbf{L} A_n$ または
$K_n = K \otimes_A^\mathbf{L} A/I^n$ という意味で系を再現することを示せる。

\begin{lemma}
\label{more-algebra-lemma-Rlim-pseudo-coherent-gives-pseudo-coherent}
$A = \lim A_n$ を環の逆系 $(A_n)$ の極限とする。
$K_n \in D(A_n)$ と $D(A_{n + 1})$ における写像
$K_{n + 1} \to K_n$ が与えられているとする。次を仮定する。
\begin{enumerate}
\item 遷移写像 $A_{n + 1} \to A_n$ は局所冪零な核をもつ全射である、
\item すべての $K_n$ が擬コヒーレントであるか、またはある整数 $n_0$ が存在して
$K_{n_0}$ が擬コヒーレントであり、$n \geq n_0$ に対して
$A_{n + 1} \to A_n$ の核が冪零イデアルである、
\item 写像は同型
$K_{n + 1} \otimes_{A_{n + 1}}^\mathbf{L} A_n \to K_n$ を誘導する。
\end{enumerate}
このとき $K = R\lim K_n$ は $D(A)$ の擬コヒーレント対象であり、
すべての $n$ に対して $K \otimes_A^\mathbf{L} A_n \to K_n$ は同型である。
\end{lemma}

\begin{proof}
仮定により、$K_1$ を表す有限自由 $A_1$-加群の上に有界な複体
$P_1^\bullet$ を見つけられる。Definition
\ref{more-algebra-definition-pseudo-coherent} を参照せよ。Lemma
\ref{more-algebra-lemma-check-pseudo-coherent-modulo-nilpotent} により、
すべての $n$ に対して $K_n$ は擬コヒーレントである。
次に Lemma \ref{more-algebra-lemma-lift-complex-stably-frees} により、
$n > 1$ に関する帰納法で、$K_n$ を表す有限自由 $A_n$-加群の複体
$P_n^\bullet$ と、写像 $K_n \to K_{n - 1}$ を表し、複体の同型 (!)
$P_n^\bullet \otimes_{A_n} A_{n - 1} \to P_{n - 1}^\bullet$ を誘導する写像
$P_n^\bullet \to P_{n - 1}^\bullet$ を見つけられる。
従って Lemma \ref{more-algebra-lemma-compute-Rlim-modules} および Remark
\ref{more-algebra-remark-how-unique} により、$K = R\lim K_n$ は
$P^\bullet = \lim P_n^\bullet$ で表される。
各 $n$ に対して $P_n^i$ は有限自由 $A_n$-加群であり、
$A = \lim A_n$ なので、各 $i$ に対して $P^i$ は $P_1^i$ と同じ階数の
有限自由加群である。これは $K$ が擬コヒーレントであることを意味する。
また $K \otimes_A^\mathbf{L} A_n$ は
$P^\bullet \otimes_A A_n = P_n^\bullet$ で表され、最後の主張が従う。
\end{proof}

\begin{lemma}
\label{more-algebra-lemma-Rlim-pseudo-coherent-gives-complete-pseudo-coherent}
$A$ を環、$I \subset A$ をイデアルとする。
$K_n \in D(A/I^n)$ と $D(A/I^{n + 1})$ における写像
$K_{n + 1} \to K_n$ が与えられているとする。次を仮定する。
\begin{enumerate}
\item $A$ は $I$-進完備である、
\item $K_1$ は擬コヒーレントである、および
\item 写像は同型
$K_{n + 1} \otimes_{A/I^{n + 1}}^\mathbf{L} A/I^n \to K_n$ を誘導する。
\end{enumerate}
このとき $K = R\lim K_n$ は $D(A)$ の擬コヒーレントな導来完備対象であり、
すべての $n$ に対して $K \otimes_A^\mathbf{L} A/I^n \to K_n$ は同型である。
\end{lemma}

\begin{proof}
Lemma \ref{more-algebra-lemma-Rlim-pseudo-coherent-gives-pseudo-coherent} により、
$K$ が擬コヒーレントであり、すべての $n$ に対して
$K \otimes_A^\mathbf{L} A/I^n \to K_n$ が同型であることは既に分かっている。
最後に Lemma \ref{more-algebra-lemma-naive-derived-completion} により $K$ は導来完備である。
\end{proof}

\begin{lemma}
\label{more-algebra-lemma-Rlim-perfect-gives-perfect}
\begin{reference}
\cite[Lemma 4.2]{Bhatt-Algebraize}
\end{reference}
$A = \lim A_n$ を環の逆系 $(A_n)$ の極限とする。
$K_n \in D(A_n)$ と $D(A_{n + 1})$ における写像
$K_{n + 1} \to K_n$ が与えられているとする。次を仮定する。
\begin{enumerate}
\item 遷移写像 $A_{n + 1} \to A_n$ は局所冪零な核をもつ全射である、
\item すべての $K_n$ が完全であるか、またはある整数 $n_0$ が存在し、
$K_{n_0}$ が完全で、$n \geq n_0$ に対して
$A_{n + 1} \to A_n$ の核が冪零である、および
\item 写像は同型
$K_{n + 1} \otimes_{A_{n + 1}}^\mathbf{L} A_n \to K_n$ を誘導する。
\end{enumerate}
このとき $K = R\lim K_n$ は $D(A)$ の完全対象であり、
すべての $n$ に対して $K \otimes_A^\mathbf{L} A_n \to K_n$ は同型である。
\end{lemma}

\begin{proof}
Lemma \ref{more-algebra-lemma-Rlim-pseudo-coherent-gives-pseudo-coherent} により、
$K$ が擬コヒーレントであり、すべての $n$ に対して
$K \otimes_A^\mathbf{L} A_n \to K_n$ が同型であることは既に分かっている。
核が極大イデアル $\mathfrak m$ である全射 $A \to \kappa$ を考える。
$A_1$ で単元に写る $A$ の任意の元は、Algebra, Lemma
\ref{algebra-lemma-locally-nilpotent-unit} により $A$ の単元である。
従って Algebra, Lemma \ref{algebra-lemma-contained-in-radical} により
$\Ker(A \to A_1)$ は $A$ の Jacobson 根基に含まれる。
ゆえに $A \to \kappa$ は $A \to A_1 \to \kappa$ と分解する。従って
$$
K \otimes_A^\mathbf{L} \kappa =
K \otimes_A^\mathbf{L} A_1 \otimes_{A_1}^\mathbf{L} \kappa =
K_1 \otimes_{A_1}^\mathbf{L} \kappa
$$
である。仮定 (2) により $K_1$ は完全であることに注意する。従って Lemma
\ref{more-algebra-lemma-perfect} により $K_1$ は有限 Tor 次元をもつ。
従ってある $a, b \in \mathbf{Z}$ が存在し、すべての
$i \not \in [a, b]$ に対して
$H^i(K \otimes_A^\mathbf{L} \kappa) = 0$ である。
Lemma \ref{more-algebra-lemma-check-perfect-pointwise} により $K$ は完全である。
\end{proof}

\begin{lemma}
\label{more-algebra-lemma-Rlim-perfect-gives-complete}
$A$ を環、$I \subset A$ をイデアルとする。
$K_n \in D(A/I^n)$ と $D(A/I^{n + 1})$ における写像
$K_{n + 1} \to K_n$ が与えられているとする。次を仮定する。
\begin{enumerate}
\item $A$ は $I$-進完備である、
\item $K_1$ は完全対象である、および
\item 写像は同型
$K_{n + 1} \otimes_{A/I^{n + 1}}^\mathbf{L} A/I^n \to K_n$ を誘導する。
\end{enumerate}
このとき $K = R\lim K_n$ は $D(A)$ の完全な導来完備対象であり、
すべての $n$ に対して $K \otimes_A^\mathbf{L} A/I^n \to K_n$ は同型である。
\end{lemma}

\begin{proof}
Lemmas \ref{more-algebra-lemma-Rlim-perfect-gives-perfect} と
\ref{more-algebra-lemma-Rlim-pseudo-coherent-gives-complete-pseudo-coherent}
（導来完備性を得るため）を組み合わせる。
\end{proof}

\noindent
次の補題が非有界複体に対して成り立つかどうかは分からない。

\begin{lemma}
\label{more-algebra-lemma-Rlim-gives-complete}
$A$ を環、$I \subset A$ をイデアルとする。
$K_n \in D(A/I^n)$ と $D(A/I^{n + 1})$ における写像
$K_{n + 1} \to K_n$ が与えられているとする。
\begin{enumerate}
\item $A$ は Noether 環である、
\item $K_1$ は上に有界である、および
\item 写像は同型
$K_{n + 1} \otimes_{A/I^{n + 1}}^\mathbf{L} A/I^n \to K_n$ を誘導する
\end{enumerate}
ならば、$K = R\lim K_n$ は $D^-(A)$ の導来完備対象であり、
すべての $n$ に対して $K \otimes_A^\mathbf{L} A/I^n \to K_n$ は同型である。
\end{lemma}

\begin{proof}
$D(A)$ の対象 $K$ は Lemma \ref{more-algebra-lemma-naive-derived-completion} により
導来完備である。

\medskip\noindent
$i > b$ に対して $H^i(K_1) = 0$ であると仮定する。
$i > b$ に対して $P_1^i = 0$ であり、$K_1$ を表す自由
$A/I$-加群の複体 $P_1^\bullet$ を見つけられる。Lemma
\ref{more-algebra-lemma-lift-complex-projectives} により、$n > 1$ に関する帰納法で、
$K_n$ を表す自由 $A/I^n$-加群の複体 $P_n^\bullet$ と、
写像 $K_n \to K_{n - 1}$ を表し、複体の同型 (!)
$P_n^\bullet/I^{n - 1}P_n^\bullet \to P_{n - 1}^\bullet$ を誘導する写像
$P_n^\bullet \to P_{n - 1}^\bullet$ を見つけられる。

\medskip\noindent
従って Lemma \ref{more-algebra-lemma-compute-Rlim-modules} および Remark
\ref{more-algebra-remark-how-unique} により、$R\lim K_n$ が
$P^\bullet = \lim P_n^\bullet$ で表される状況に到達した。
複体 $P_n^\bullet$ は平坦 $A/I^n$-加群の一様に上に有界な複体であり、
遷移写像は項ごとに全射である。Lemma \ref{more-algebra-lemma-limit-flat} により、
$P^\bullet$ は平坦 $A$-加群の上に有界な複体である。
従って $K \otimes_A^\mathbf{L} A/I^t$ は
$P^\bullet \otimes_A A/I^t$ で表される。Lemma
\ref{more-algebra-lemma-limit-flat} により項ごとに
$P^\bullet \otimes_A A/I^t = \lim P_n^\bullet \otimes_A A/I^t$ である。
$P_n^\bullet$ の選び方により、$n \geq t$ に対して遷移写像
$P_{n + 1}^\bullet \otimes_A A/I^t \to P_n^\bullet \otimes_A A/I^t$
は同型なので、
$\lim P_n^\bullet \otimes_A A/I^t = P_t^\bullet \otimes_A A/I^t = P_t^\bullet$
である。$P_t^\bullet$ は $K_t$ を表すので、
$K \otimes_A^\mathbf{L} A/I^t \to K_t$ は同型である。
\end{proof}

\noindent
別の種類の結果を挙げる。

\begin{lemma}[Koll\'ar--Kov\'acs]
\label{more-algebra-lemma-kollar-kovacs}
\begin{reference}
Email from Kovacs of 23/02/2018.
\end{reference}
$I$ を Noether 環 $A$ のイデアル、$K \in D(A)$ とする。
$K_n = K \otimes_A^\mathbf{L} A/I^n$ と置く。
すべての $i \in \mathbf{Z}$ に対して次を仮定する。
\begin{enumerate}
\item $H^i(K)$ は有限 $A$-加群である、および
\item 系 $H^i(K_n)$ は Mittag--Leffler 条件を満たす。
\end{enumerate}
このとき、すべての $i \in \mathbf{Z}$ に対して
$\lim H^i(K)/I^nH^i(K)$ は $\lim H^i(K_n)$ に等しい。
\end{lemma}

\begin{proof}
Proposition \ref{more-algebra-proposition-noetherian-naive-completion-is-completion} により、
$K^\wedge = R\lim K_n$ は $K$ の導来完備化であることを思い出そう。
Lemma \ref{more-algebra-lemma-derived-completion-pseudo-coherent} により
$H^i(K^\wedge) = \lim H^i(K)/I^nH^i(K)$ である。
Lemma \ref{more-algebra-lemma-break-long-exact-sequence-modules} により短完全列
$$
0 \to R^1\lim H^{i - 1}(K_n) \to H^i(K^\wedge) \to \lim H^i(K_n) \to 0
$$
を得る。Mittag--Leffler 条件により左の項は零である
（Lemma \ref{more-algebra-lemma-compute-Rlim-modules}）ので、補題が従う。
\end{proof}


\section{いくつかの評価写像}
\label{more-algebra-section-evaluation}

\noindent
この節では、適切な種類の複体に対して、ある標準的な $R\Hom$ の写像が
同型であることを証明する。

\begin{lemma}
\label{more-algebra-lemma-internal-hom-evaluate-isomorphism}
$R$ を環とし、$K, L, M$ を $D(R)$ の対象とする。
% P02-E1145: frozen source line 28564 begins a new sentence with lowercase "the".
Lemma \ref{more-algebra-lemma-internal-hom-evaluate} の写像
$$
R\Hom_R(L, M) \otimes_R^\mathbf{L} K \longrightarrow R\Hom_R(R\Hom_R(K, L), M)
$$
は次の二つの場合に同型である。
\begin{enumerate}
\item $K$ が完全である場合、または
\item $K$ が擬コヒーレントで、$L \in D^+(R)$、かつ $M$ が有限単射次元をもつ場合。
\end{enumerate}
\end{lemma}

\begin{proof}
$M$ を表す K-単射複体 $I^\bullet$、$L$ を表す K-単射複体 $J^\bullet$、
および $K$ を表す有限射影加群の上に有界な複体 $K^\bullet$ を選ぶ。
Lemma \ref{more-algebra-lemma-evaluate-and-more} の複体の写像
$$
\text{Tot}(\Hom^\bullet(J^\bullet, I^\bullet) \otimes_R K^\bullet)
\longrightarrow
\Hom^\bullet(\Hom^\bullet(K^\bullet, J^\bullet), I^\bullet)
$$
を考える。$K^s$ は有限射影的なので
$$
\left(\prod\nolimits_{p + r = t} \Hom_R(J^{-r}, I^p)\right) \otimes_R K^s =
\prod\nolimits_{p + r = t} \Hom_R(J^{-r}, I^p) \otimes_R K^s
$$
であることに注意する。この写像は、$K^s$ が有限射影的なので同型である写像
$$
c_{p, r, s} :
\Hom_R(J^{-r}, I^p) \otimes_R K^s
\longrightarrow
\Hom_R(\Hom_R(K^s, J^{-r}), I^p)
$$
で与えられる。左辺の次数 $n$ の任意の元
$\alpha = (\alpha^{p, r, s})$ に対して、ある
$n = p + r + s$ を満たす $p, r$ について $\alpha^{p, r, s}$ が非零となる
$s$ の値は有限個しかない。従って、右辺の元
$\beta = (\beta^{p, r, s})$ にも同じ消滅条件が強制されるならば、この写像は同型である。
$K^\bullet$ が有限射影加群の有界複体ならばこれは明らかである。
他方、$I^\bullet$ を有界、$J^\bullet$ を下に有界に選べるならば、
$\beta^{p, r, s}$ は $p$ がある固定範囲の外にあるとき、$s \gg 0$ のとき、
および $r \gg 0$ のとき零である。従って
$\beta^{p, r, s}$ が非零である $n = p + r + s$ の解のうち、
現れる $s$ の値は有限個しかない。
\end{proof}

\begin{lemma}
\label{more-algebra-lemma-internal-hom-evaluate-isomorphism-technical}
$R$ を環とし、$K, L, M$ を $D(R)$ の対象とする。
% P02-E1146: frozen source line 28619 begins a new sentence with lowercase "the".
Lemma \ref{more-algebra-lemma-internal-hom-evaluate} の写像
$$
R\Hom_R(L, M) \otimes_R^\mathbf{L} K \longrightarrow R\Hom_R(R\Hom_R(K, L), M)
$$
は次の三条件が満たされるならば同型である。
\begin{enumerate}
\item $L, M$ は有限単射次元をもつ、
\item $R\Hom_R(L, M)$ は有限 Tor 次元をもつ、
\item すべての $n \in \mathbf{Z}$ に対して切断 $\tau_{\leq n}K$ は擬コヒーレントである。
\end{enumerate}
\end{lemma}

\begin{proof}
整数 $n$ を選び、完全三角形
$$
\tau_{\leq n}K \to K \to \tau_{\geq n + 1}K \to \tau_{\leq n}K[1]
$$
を考える。Derived Categories, Remark
\ref{derived-remark-truncation-distinguished-triangle} を参照せよ。
仮定 (3) と Lemma \ref{more-algebra-lemma-internal-hom-evaluate-isomorphism} により、
写像は $\tau_{\leq n}K$ に対して同型である。従って、ある $c$ に対して
$$
R\Hom_R(L, M) \otimes_R^\mathbf{L} \tau_{\geq n + 1}K
\quad\text{and}\quad
R\Hom_R(R\Hom_R(\tau_{\geq n + 1}K, L), M)
$$
の双方が次数 $\leq n - c$ で消滅するコホモロジーをもつことを示せば十分である。
これは仮定 (2) と (1) から直ちに従う。
\end{proof}

\begin{lemma}
\label{more-algebra-lemma-internal-hom-evaluate-tensor-isomorphism}
$R$ を環とし、$K, L, M$ を $D(R)$ の対象とする。Lemma
\ref{more-algebra-lemma-internal-hom-diagonal-better} の写像
$$
K \otimes_R^\mathbf{L} R\Hom_R(M, L) \longrightarrow
R\Hom_R(M, K \otimes_R^\mathbf{L} L)
$$
は次の場合に同型である。
\begin{enumerate}
\item $M$ が完全である場合、または
\item $K$ が完全である場合、または
\item $M$ が擬コヒーレントで、$L \in D^+(R)$、かつ $K$ が
% P02-E1147: frozen source line 28665 closes an interval at infinity with a square bracket.
$[a, \infty]$ に Tor 振幅をもつ場合。
\end{enumerate}
\end{lemma}

\begin{proof}
$M$ が完全である場合の証明。矢印の両辺は $M$ の完全三角形を完全三角形に写し、
直和と可換であることに注意する。従って $M = R[n]$ の場合に成り立つことを
確認すれば十分である。Derived Categories, Remark
\ref{derived-remark-check-on-generator} および Lemma
\ref{more-algebra-lemma-perfect-ring-classical-generator} を参照せよ。この場合、結果は明らかである。

\medskip\noindent
$K$ が完全である場合の証明。前の場合と同じ議論である。

\medskip\noindent
(3) の場合の証明。$K$ と $L$ を $R$-加群の下に有界な複体
$K^\bullet$ と $L^\bullet$ で表すことができる。Lemma
\ref{more-algebra-lemma-bounded-below-tor-amplitude} により、$K^\bullet$ は平坦
$R$-加群からなる K-平坦複体であると仮定してよい。
Definition \ref{more-algebra-definition-pseudo-coherent} により、$M$ を有限自由
$R$-加群の上に有界な複体 $M^\bullet$ で表すことができる。
このとき左辺の対象は
$$
\text{Tot}(K^\bullet \otimes_R \Hom^\bullet(M^\bullet, L^\bullet))
$$
で表され、右辺の対象は
$$
\Hom^\bullet(M^\bullet, \text{Tot}(K^\bullet \otimes_R L^\bullet))
$$
で表される。ここで Lemma \ref{more-algebra-lemma-RHom-out-of-projective} を用いた。
両複体の次数 $n$ の加群は
$$
\bigoplus\nolimits_{p + q + r = n} K^p \otimes \Hom_R(M^{-r}, L^q) =
\bigoplus\nolimits_{p + q + r = n} \Hom_R(M^{-r}, K^p \otimes_R L^q)
$$
である。これは $M^{-r}$ が有限自由であるからである
% P02-E1148: frozen source line 28703 has the awkward parenthetical "as well these are finite direct sums".
（また、これらは有限直和である）。
Lemma \ref{more-algebra-lemma-internal-hom-diagonal-better} で定義した写像は、
これらの加群の間の標準同型を用いる Lemma
\ref{more-algebra-lemma-diagonal-better} で定義した複体の写像から生じる。
\end{proof}

\begin{lemma}
\label{more-algebra-lemma-hom-complex-K-flat}
$R$ を環とする。$P^\bullet$ を射影 $R$-加群の上に有界な複体、
$K^\bullet$ を $R$-加群の K-平坦複体とする。$P^\bullet$ が $D(R)$ の
完全対象ならば、$\Hom^\bullet(P^\bullet, K^\bullet)$ は K-平坦であり、
$R\Hom_R(P^\bullet, K^\bullet)$ を表す。
\end{lemma}

\begin{proof}
最後の主張は Lemma \ref{more-algebra-lemma-RHom-out-of-projective} である。
$P^\bullet$ は完全対象を表すので、有限射影 $R$-加群の有限複体
$F^\bullet$ が存在し、$P^\bullet$ と $F^\bullet$ は $D(R)$ で同型である。
Definition \ref{more-algebra-definition-perfect} を参照せよ。
Derived Categories, Lemma
\ref{derived-lemma-morphisms-from-projective-complex} により、
$P^\bullet$ と $F^\bullet$ はホモトピー同値である。
従って $\Hom^\bullet(P^\bullet, K^\bullet)$ と
$\Hom^\bullet(F^\bullet, K^\bullet)$ はホモトピー同値である。
ゆえに前者が K-平坦であることと後者が K-平坦であることは同値である
（Definition \ref{more-algebra-definition-K-flat} および Lemma
\ref{more-algebra-lemma-derived-tor-homotopy} から従う）。
明らかに
$$
\Hom^\bullet(F^\bullet, K^\bullet) =
\text{Tot}(E^\bullet \otimes_R K^\bullet)
$$
である。ここで $E^\bullet$ は $F^\bullet$ の双対複体で、
各項は $E^n = \Hom_R(F^{-n}, R)$ である。Lemma
\ref{more-algebra-lemma-dual-perfect-complex} とその証明を参照せよ。
$E^\bullet$ は射影加群の有界複体なので、Lemma
\ref{more-algebra-lemma-derived-tor-quasi-isomorphism} により K-平坦である。
従って Lemma \ref{more-algebra-lemma-tensor-product-K-flat} から結論が従う。
\end{proof}


\section{導来 hom の底変換}
\label{more-algebra-section-base-change-RHom}

\noindent
これについて論じる内容の一部は既に Lemma
\ref{more-algebra-lemma-pseudo-coherence-and-base-change-ext} および
Algebra, Section \ref{algebra-section-functoriality-ext} で見た。

\begin{lemma}
\label{more-algebra-lemma-upgrade-adjoint-tensor-RHom}
$R \to R'$ を環写像とする。$K \in D(R)$ および $M \in D(R')$ に対して
標準同型
$$
R\Hom_R(K, M) = R\Hom_{R'}(K \otimes_R^\mathbf{L} R', M)
$$
がある。
\end{lemma}

\begin{proof}
$M$ を表す K-単射な $R'$-加群の複体 $J^\bullet$ を選ぶ。
$I^\bullet$ が K-単射な $R$-加群の複体であるような擬同型
$J^\bullet \to I^\bullet$ を選ぶ。$K$ を表す K-平坦な
$R$-加群の複体 $K^\bullet$ を選ぶ。写像
$$
\Hom^\bullet(K^\bullet \otimes_R R', J^\bullet)
\longrightarrow
\Hom^\bullet(K^\bullet, I^\bullet)
$$
を考える。次数 $n$ の項上の写像は、
$K^{-q} \to K^{-q} \otimes_R R'$ との前合成および
$J^p \to I^p$ との後合成から生じる写像
$$
\prod\nolimits_{n = p + q} \Hom_{R'}(K^{-q} \otimes_R R', J^p)
\longrightarrow
\prod\nolimits_{n = p + q} \Hom_R(K^{-q}, I^p)
$$
で与えられる。証明を終えるには、コホモロジー群上で同型
$$
\Hom_{D(R)}(K, M) = \Hom_{D(R')}(K \otimes_R^\mathbf{L} R', M)
$$
を得ることを示せば十分である。これは Lemma
\ref{more-algebra-lemma-tensor-hom-adjoint} により、底変換
$- \otimes_R^\mathbf{L} R' : D(R) \to D(R')$ が制限関手
$D(R') \to D(R)$ の左随伴だからである。
\end{proof}

\noindent
$R \to R'$ を環写像とする。$K, M \in D(R)$ に関して関手的な
$D(R')$ の底変換写像
\begin{equation}
\label{more-algebra-equation-base-change-RHom}
R\Hom_R(K, M) \otimes_R^\mathbf{L} R'
\longrightarrow
R\Hom_{R'}(K \otimes_R^\mathbf{L} R', M \otimes_R^\mathbf{L} R')
\end{equation}
がある。実際、$- \otimes_R^\mathbf{L} R' : D(R) \to D(R')$ と制限関手
$D(R') \to D(R)$ の随伴性により、これは写像
$$
R\Hom_R(K, M)
\longrightarrow
R\Hom_{R'}(K \otimes_R^\mathbf{L} R', M \otimes_R^\mathbf{L} R') =
R\Hom_R(K, M \otimes_R^\mathbf{L} R')
$$
（等式は Lemma \ref{more-algebra-lemma-upgrade-adjoint-tensor-RHom} による）と同じものであり、
これには標準写像 $M \to M \otimes_R^\mathbf{L} R'$
（随伴の単位）を用いることができる。

\begin{lemma}
\label{more-algebra-lemma-base-change-RHom}
$R \to R'$ を環写像とし、$K, M \in D(R)$ とする。写像
(\ref{more-algebra-equation-base-change-RHom})
$$
R\Hom_R(K, M) \otimes_R^\mathbf{L} R'
\longrightarrow
R\Hom_{R'}(K \otimes_R^\mathbf{L} R', M \otimes_R^\mathbf{L} R')
$$
は次の場合に $D(R')$ で同型である。
\begin{enumerate}
\item $K$ が完全である、
\item $R'$ が $R$-加群として完全である、
\item $R \to R'$ が平坦で、$K$ が擬コヒーレント、かつ
$M \in D^{+}(R)$ である、または
\item $R'$ が $R$-加群として有限 Tor 次元をもち、
$K$ が擬コヒーレント、かつ $M \in D^{+}(R)$ である。
\end{enumerate}
\end{lemma}

\begin{proof}
制限関手 $D(R') \to D(R)$ を適用した後に写像が同型であることを
確認してよい。この関手を適用すると、写像は
% P02-E1149: frozen source lines 28843-28845 use undefined L where the lemma's object is M.
$$
R\Hom_R(K, L) \otimes_R^\mathbf{L} R'
\longrightarrow
R\Hom_R(K, L \otimes_R^\mathbf{L} R')
$$
という Lemma \ref{more-algebra-lemma-internal-hom-diagonal-better} の写像になる。
左右を対応させるには補題の前の議論を参照せよ。特に Lemma
\ref{more-algebra-lemma-upgrade-adjoint-tensor-RHom} を用いる。
従って Lemma \ref{more-algebra-lemma-internal-hom-evaluate-tensor-isomorphism} から結論が従う。
\end{proof}


\section{加群の系}
\label{more-algebra-section-systems}

\noindent
$I$ を Noether 環 $A$ のイデアルとする。この節では、有限
$A$-加群と有限 $A/I^n$-加群の系との関係についての知見を補う。

\begin{lemma}
\label{more-algebra-lemma-consequence-Artin-Rees}
$I$ を Noether 環 $A$ のイデアルとする。
$$
K \xrightarrow{\alpha} L \xrightarrow{\beta} M
$$
を有限 $A$-加群の複体とする。$H = \Ker(\beta)/\Im(\alpha)$ とおく。
$n \geq 0$ に対し
$$
K/I^nK \xrightarrow{\alpha_n} L/I^nL \xrightarrow{\beta_n} M/I^nM
$$
を誘導された複体とし、$H_n = \Ker(\beta_n)/\Im(\alpha_n)$ とおく。
このとき可換図式
$$
\xymatrix{
& & & H \ar[lld] \ar[ld] \ar[d] \\
\ldots \ar[r] & H_3 \ar[r] & H_2 \ar[r] & H_1
}
$$
を与える標準的な $A$-加群写像がある。さらに、ある $c > 0$ と、
$n \geq c$ に対する標準的な $A$-加群写像
$H_n \to H/I^{n - c}H$ が存在し、合成
$$
H/I^n H \to H_n \to  H/I^{n - c}H
\quad\text{and}\quad
H_n \to H/I^{n - c}H \to H_{n - c}
$$
はいずれも標準的な写像である。さらに次が成り立つ。
\begin{enumerate}
\item $(H_n)$ と $(H/I^nH)$ は $\text{Mod}_A$ の pro-対象として同型である。
\item $\lim H_n = \lim H/I^n H$ である。
\item 逆系 $(H_n)$ は Mittag--Leffler 条件を満たす。
\item $H_{n + c} \to H_n$ の像は $H \to H_n$ の像に等しい。
\item 合成
$I^cH_n \to H_n \to H/I^{n - c}H \to H_n/I^{n - c}H_n$
は包含 $I^cH_n \to H_n$ に商写像
$H_n \to H_n/I^{n - c}H_n$ を続けたものである。
% P02-E1150: source line 28901 has singular agreement for the joint kernel-and-cokernel subject.
\item $H/I^nH \to H_n$ の核と余核は $I^c$ で零化される。
\end{enumerate}
\end{lemma}

\begin{proof}
$H_n = \beta^{-1}(I^nM)/\Im(\alpha) + I^nL$ であることに注意する。
$n \geq 2$ に対して
$\beta^{-1}(I^nM) \subset \beta^{-1}(I^{n - 1}M)$ および
$\Im(\alpha) + I^nL \subset \Im(\alpha) + I^{n - 1}L$ が成り立つ。
従って標準写像 $H_n \to H_{n - 1}$ を得る。同様に、
$\Ker(\beta) \subset \beta^{-1}(I^nM)$ および
$\Im(\alpha) \subset \Im(\alpha) + I^nL$ から標準写像
$H \to H_n$ を得る。図式が可換であることの確認は省略する。

\medskip\noindent
Artin--Rees により、$n \geq c_1$ に対して
$\beta^{-1}(I^nM) \subset \Ker(\beta) + I^{n - c_1}L$、また
$n \geq c_2$ に対して
$\Ker(\beta) \cap I^nL \subset I^{n - c_2}\Ker(\beta)$
となる $c_1, c_2 \geq 0$ を選べる。Algebra, Lemmas
\ref{algebra-lemma-map-AR} および
\ref{algebra-lemma-Artin-Rees} を参照せよ。$c = c_1 + c_2$ とおく。

\medskip\noindent
$n \geq c$ とする。$\psi_n : H_n \to H/I^{n - c}H$ を次のように定める。
$x \in H_n$ とし、$x$ を表す
$y \in \beta^{-1}(I^nM)$ を選ぶ。
$y = z + w$ と書く。ここで $z \in \Ker(\beta)$、
$w \in I^{n - c_1}L$ である（$c_1$ の選び方により可能である）。
$\psi_n(x)$ を $H/I^{n - c}H$ における $z$ の類と定める。
これが well-defined であることを見るため、同じ $x$ に対して上と同様の
別の選択 $y', z', w'$ をしたとする。記号の意味は明らかである。
このとき $y' - y \in \Im(\alpha) + I^nL$ なので、
ある $v \in K$ と $u \in I^nL$ によって
$y' - y = \alpha(v) + u$ と書ける。従って
$$
y' = z' + w' = \alpha(v) + u + z + w
\Rightarrow
z' = z + \alpha(v) + u + w - w'
$$
である。$\beta(z' - z - \alpha(v)) = 0$ だから
$u + w - w' \in \Ker(\beta) \cap I^{n - c_1}L$ であり、
$c_2$ の選び方によりこれは
$I^{n - c_1 - c_2}\Ker(\beta) = I^{n - c}\Ker(\beta)$
に含まれる。ゆえに $z'$ と $z$ は $H/I^{n - c}H$ で同じ像をもち、
望む結論を得る。

\medskip\noindent
合成 $H/I^n H \to H_n \to  H/I^{n - c}H$ は標準写像である。
実際、$z \in \Ker(\beta)$ が
$H/I^nH = \Ker(\beta)/\Im(\alpha) + I^n\Ker(\beta)$
における元 $x$ を表すなら、上の構成による写像のもとで
$x$ が $H/I^{n - c}H$ における $z$ の類へ写ることは明らかである。

\medskip\noindent
合成 $H_n \to H/I^{n - c}H \to H_{n - c}$ を考える。
上の $\psi_n$ の構成に現れる $x, y, z, w$ を取る。
% P02-E1151: source line 28953 misspells “class” as “cass”.
$x$ は $H_{n - c}$ における $z$ の類へ写る。一方、
証明の第1段落で構成した標準写像 $H_n \to H_{n - c}$ は
$x$ を $y$ の類へ送る。従って
$y - z \in \Im(\alpha) + I^{n - c}L$ を示せばよいが、
$y - z = w \in I^{n - c_1}L \subset I^{n - c}L$ なのでこれは成り立つ。

\medskip\noindent
(1)--(4) は今示したことの形式的帰結である。実際、(1) は写像の存在と
Categories, Remark \ref{categories-remark-pro-category}
における pro-対象の射の定義から従う。
(2) は、同型な pro-対象の極限が同型だから成り立つ。
(3) は (4) から直ちに従う。
(4) は標準写像 $H_{n + c} \to H_n$ の分解
$H_{n + c} \to H/I^nH \to H_n$ から従う。

\medskip\noindent
(5) の証明。$x \in I^cH_n$ とする。
$x_i \in H_n$ および $f_i \in I^c$ によって
$x = \sum f_i x_i$ と書く。各 $x_i$ について
$\psi_n$ の構成に現れる $y_i, z_i, w_i$ を選ぶ。
すると $x$ の $\psi_n$ を計算する際に
$y = \sum f_iy_i$、$z = \sum f_i z_i$、$w = \sum f_i w_i$
と選べるので、$\psi_n(x)$ は $z$ の類で与えられる。
$w = \sum f_i w_i$ は $I^nL$ に属するから、
この元の $H_n/I^{n - c}H_n$ における像は $y$ の類に等しい。
これで (5) が証明された。

\medskip\noindent
(6) の証明。$y \in \Ker(\beta)$ とし、$H$ におけるその類を $x$ とする。
$x$ が $H_n$ で零へ写るなら
$y \in I^nL + \Im(\alpha)$ である。従って、ある $v \in K$ に対し
$y - \alpha(v) \in \Ker(\beta) \cap I^nL$ となる。
このとき $y - \alpha(v) \in I^{n - c_2}\Ker(\beta)$ なので、
$H/I^nH$ における $y$ の類は $I^{c_2}$ で零化される。
最後に $x \in H_n$ を
$y \in \beta^{-1}(I^nM)$ の類とする。上と同様に、
$z \in \Ker(\beta)$ および $w \in I^{n - c_1}L$ によって
$y = z + w$ と書く。明らかに $f \in I^{c_1}$ なら
$fx$ は
$fy + fw \equiv fy$ modulo $\Im(\alpha) + I^nL$
の類であり、従って $fx$ は望む通り
$H$ における $fy$ の類の像である。
\end{proof}

\begin{lemma}
\label{more-algebra-lemma-kollar-kovacs-pseudo-coherent}
\begin{reference}
2018年2月23日付 Kovacs の電子メール。
\end{reference}
$I$ を Noether 環 $A$ のイデアルとする。
$K \in D(A)$ を擬コヒーレントとし、
$K_n = K \otimes_A^\mathbf{L} A/I^n$ とおく。
このとき任意の $i \in \mathbf{Z}$ に対して系 $H^i(K_n)$ は
Mittag--Leffler 条件を満たし、
$\lim H^i(K)/I^nH^i(K)$ は $\lim H^i(K_n)$ に等しい。
\end{lemma}

\begin{proof}
$K$ を有限自由 $A$-加群の上に有界な複体 $P^\bullet$ で表してよい。
このとき $K_n$ は $P^\bullet/I^nP^\bullet$ で表される。
% P02-E1152: source lines 29008-29009 leave the Mittag–Leffler conclusion as a sentence fragment.
従って Mittag--Leffler 性は
Lemma \ref{more-algebra-lemma-consequence-Artin-Rees} から従う。
最後の主張はさらに Lemma \ref{more-algebra-lemma-kollar-kovacs} から従う。
\end{proof}

\begin{lemma}
\label{more-algebra-lemma-derived-completion-plain-completion}
$A$ を Noether 環、$I \subset A$ をイデアルとする。
$M^\bullet$ を有限 $A$-加群の有界複体とする。写像の逆系
$$
M^\bullet \otimes_A^\mathbf{L} A/I^n \longrightarrow M^\bullet/I^nM^\bullet
$$
は $D(A)$ の pro-対象の同型を定める。
\end{lemma}

\begin{proof}
$I = (f_1, \ldots, f_r)$ とする。
$K_n \in D(A)$ を $f_1^n, \ldots, f_r^n$ 上の Koszul 複体で表される
対象とする。Lemma \ref{more-algebra-lemma-sequence-Koszul-complexes}
を参照すると、pro-同型な逆系を誘導する写像
$K_n \to A/I^n$ があることを思い出す。
従って
$$
M^\bullet \otimes_A^\mathbf{L} K_n \longrightarrow M^\bullet/I^nM^\bullet
$$
が $D(A)$ の pro-対象の同型を定めることを示せば十分である。
$K_n$ は次数 $-r, \ldots, 0$ にある有限自由 $A$-加群の複体で
表されるので、上の矢印の始域と終域のコホモロジーが、すべての
$n$ に対して $[a, b]$ の外で消えるような
$a, b \in \mathbf{Z}$ が存在する。
従って Derived Categories, Lemma
\ref{derived-lemma-pro-isomorphism-bis} を適用でき、任意の
$i \in \mathbf{Z}$ に対して写像
$$
H^i(M^\bullet \otimes_A^\mathbf{L} A/I^n) \to H^i(M^\bullet/I^nM^\bullet)
$$
が $A$-加群の pro-系の同型を定めることを示せば十分だと分かる。
これを見るため、$P^\bullet$ が有限自由 $A$-加群の上に有界な複体である
ような擬同型 $P^\bullet \to M^\bullet$ を選ぶ。上の矢印は写像
$$
H^i(P^\bullet/I^nP^\bullet) \to H^i(M^\bullet/I^nM^\bullet)
$$
で与えられる。これらは Lemma
\ref{more-algebra-lemma-consequence-Artin-Rees} により pro-系の同型を定める。
実際、この補題は両者が
$H^i = H^i(M^\bullet) = H^i(P^\bullet)$ としたときの pro-系
$H^i/I^nH^i$ と同型であることを示す。
\end{proof}

\begin{lemma}
\label{more-algebra-lemma-hom-systems-ML}
$A$ を Noether 環、$I \subset A$ をイデアルとする。
$M$, $N$ を有限 $A$-加群とする。
$M_n = M/I^nM$ および $N_n = N/I^nN$ とおく。このとき
\begin{enumerate}
\item 系 $(\Hom_A(M_n, N_n))$ と $(\text{Isom}_A(M_n, N_n))$ は
Mittag--Leffler 条件を満たす。
\item ある $c \geq 0$ が存在し、すべての $n$ に対して
$$
\Hom_A(M, N)/I^n\Hom_A(M, N) \to \Hom_A(M_n, N_n)
$$
の核と余核は $I^c$ で零化される。
\item
$\lim \Hom_A(M_n, N_n) =\Hom_A(M, N)^\wedge =
\Hom_{A^\wedge}(M^\wedge, N^\wedge)$
である。
\item $\lim \text{Isom}_A(M_n, N_n) =
\text{Isom}_{A^\wedge}(M^\wedge, N^\wedge)$ である。
\end{enumerate}
ここで ${}^\wedge$ は通常の $I$-進完備化を表す。
\end{lemma}

\begin{proof}
$\Hom_A(M_n, N_n) = \Hom_A(M, N_n)$ であることに注意する。
表示
$$
A^{\oplus t} \to A^{\oplus s} \to M \to 0
$$
を選ぶ。
% P02-E1153: source line 29085 incorrectly calls the contravariant Hom functor right exact.
右完全関手 $\Hom_A(-, N)$ を適用すると複体
$$
0 \xrightarrow{\alpha} N^{\oplus s} \xrightarrow{\beta} N^{\oplus t}
$$
を得る。その中間におけるコホモロジーは $\Hom_A(M, N)$ であり、
$n \geq 0$ に対して複体
$$
0 \xrightarrow{\alpha_n} N_n^{\oplus s} \xrightarrow{\beta_n} N_n^{\oplus t}
$$
のコホモロジーは $\Hom_A(M_n, N_n)$ である。
この $A$, $I$, $\alpha$, $\beta$ に対して
Lemma \ref{more-algebra-lemma-consequence-Artin-Rees}
に現れる $c \geq 0$ を取る。この補題の (3) から
$(\Hom_A(M_n, N_n))$ の Mittag--Leffler 性が従う。
% P02-E1154: source line 29101 contains the stray token “[art”.
写像
$\Hom_A(M, N)/I^n\Hom_A(M, N) \to \Hom_A(M_n, N_n)$
の核と余核は補題の (6) により $I^c$ で零化される。
補題の (2) から
$\lim \Hom_A(M_n, N_n) = \Hom_A(M, N)^\wedge$ を得る。等式
$$
\Hom_{A^\wedge}(M^\wedge, N^\wedge) = \lim \Hom_A(M_n, N_n)
$$
は、$M^\wedge = \lim M_n$ および
$M_n = M^\wedge/I^nM^\wedge$ と、$N$ に関する対応する事実から
形式的に従う。Algebra, Lemma
\ref{algebra-lemma-completion-complete} を参照せよ。

\medskip\noindent
同型に関する結果は、$(\Hom(M_n, N_n))$ と
$(\Hom(N_n, M_n))$ の双方に準同型の場合を適用し、次の事実を
用いれば従う。$n > m > 0$ とし、写像
$\alpha : M_n \to N_n$ および $\beta : N_n \to M_n$ が、
% P02-E1155: source line 29116 has the ungrammatical phrase “induce an isomorphisms”.
$M_m \to N_m$ および $N_m \to M_m$ 上に同型を誘導するとする。
このとき $\alpha$ と $\beta$ は同型である。
% P02-E1156: source proves only the alpha-after-beta composite; the symmetric composite is also needed to conclude both maps are isomorphisms.
実際、Nakayama の補題（Algebra, Lemma
\ref{algebra-lemma-NAK}）により $\alpha \circ \beta$ は全射であり、
従って $\alpha \circ \beta$ は
Algebra, Lemma \ref{algebra-lemma-fun} により同型である。
\end{proof}

\begin{lemma}
\label{more-algebra-lemma-isomorphic-completions}
$A$ を Noether 環、$I \subset A$ をイデアルとする。
$M$, $N$ を有限 $A$-加群とする。
$M_n = M/I^nM$ および $N_n = N/I^nN$ とおく。
すべての $n$ に対して $M_n \cong N_n$ なら、
$A^\wedge$-加群として $M^\wedge \cong N^\wedge$ である。
\end{lemma}

\begin{proof}
Lemma \ref{more-algebra-lemma-hom-systems-ML} により
系 $(\text{Isom}_A(M_n, N_n))$ は Mittag--Leffler 条件を満たす。
仮定により集合 $\text{Isom}_A(M_n, N_n)$ はすべて空でない。
従って $\lim \text{Isom}_A(M_n, N_n)$ は空でない。
$$
\lim \text{Isom}_A(M_n, N_n) = \text{Isom}_{A^\wedge}(M^\wedge, N^\wedge)
$$
だから同型を得る。
\end{proof}

\begin{remark}
\label{more-algebra-remark-weird-systems}
$I$ を Noether 環 $A$ のイデアルとする。
$n \geq 1$ に対して $A_n = A/I^n$ とおく。次の圏を考える。
\begin{enumerate}
\item 対象は列 $\{E_n\}_{n \geq 1}$ であり、
$E_n$ は有限 $A_n$-加群である。
\item 射 $\{E_n\} \to \{E'_n\}$ は写像
$$
\varphi_n : I^cE_n \longrightarrow E'_n/E'_n[I^c]
\quad\text{for }n \geq c
$$
で与えられる。ここで $E'_n[I^c]$ はねじれ部分加群
（Section \ref{more-algebra-section-torsion}）であり、次の同値関係で割る。
$(c, \varphi_n)$ と $(c + 1, \overline{\varphi}_n)$ を同じものとする。
ここで
$\overline{\varphi}_n : I^{c + 1}E_n \longrightarrow E'_n/E'_n[I^{c + 1}]$
は誘導された写像である。
\end{enumerate}
$(c, \varphi_n) : \{E_n\} \to \{E'_n\}$ と
$(c', \varphi'_n) : \{E'_n\} \to \{E''_n\}$ の合成は、
$n \geq c + c'$ に対する明らかな合成
$$
I^{c + c'}E_n \to I^{c'}E'_n/E'_n[I^{c}] \to E''_n/E''_n[I^{c + c'}]
$$
で定める。これが圏であることの確認は省略する。
\end{remark}

\begin{lemma}
\label{more-algebra-lemma-iso}
Remark \ref{more-algebra-remark-weird-systems} の圏の射
$(c, \varphi_n)$ が同型であるための必要十分条件は、
ある $c' \geq 0$ が存在し、すべての $n \gg 0$ に対して
$\Ker(\varphi_n)$ と $\Coker(\varphi_n)$ が
$I^{c'}$-ねじれとなることである。
\end{lemma}

\begin{proof}
$c' \geq c$ であり、すべての $n$ に対して
$\Ker(\varphi_n)$ と $\Coker(\varphi_n)$ が
$I^{c'}$-ねじれであると仮定してよく、そうする。
$n \geq c'$ とし $x \in I^{c'}E'_n$ とする。
$\Coker(\varphi_n)$ は $I^{c'}$ で零化されるので、
$x = \varphi_n(y) \bmod E'_n[I^c]$ となる
$y \in I^cE_n$ を選べる。$\psi_n(x)$ を
$E_n/E_n[I^{c'}]$ における $y$ の類と定める。
別の $y' \in I^cE_n$ が
$x = \varphi_n(y') \bmod E'_n[I^c]$ を満たすなら、
差 $y - y'$ は $E'_n/E'_n[I^c]$ で零へ写り、従って
$I^cE_n$ において $I^{c'}$ で零化される。
ゆえに写像
$\psi_n : I^{c'}E'_n \to E_n/E_n[I^{c'}]$
は well-defined である。
$(c', \psi_n)$ がこの圏における $(c, \varphi_n)$ の逆であることの
確認は省略する。
\end{proof}

\begin{lemma}
\label{more-algebra-lemma-dejong-kollar-kovacs}
\begin{reference}
2018年2月23日付の Janos Kollar、Sandor Kovacs、Johan de Jong
の間の電子メールによる書簡。
\end{reference}
$I$ を Noether 環 $A$ のイデアルとする。
$M$ と $N$ を有限 $A$-加群とする。
$A_n = A/I^n$、$M_n = M/I^nM$、$N_n = N/I^nN$ と書く。
任意の $i \geq 0$ に対して対象
$$
\{\Ext^i_A(M, N)/I^n\Ext^i_A(M, N)\}_{n \geq 1}
\quad\text{and}\quad
\{\Ext^i_{A_n}(M_n, N_n)\}_{n \geq 1}
$$
は Remark \ref{more-algebra-remark-weird-systems} の圏 $\mathcal{C}$ で同型である。
\end{lemma}

\begin{proof}
短完全列
$$
0 \to K \to A^{\oplus r} \to M \to 0
$$
を選び、$K_n = K/I^nK$ とおく。
$n \geq 1$ に対して
$K(n) = \Ker(A_n^{\oplus r} \to M_n)$ と定めると、完全列
$$
0 \to K(n) \to A_n^{\oplus r} \to M_n \to 0
$$
と全射 $K_n \to K(n)$ がある。実際、Lemma
\ref{more-algebra-lemma-consequence-Artin-Rees} により、ある
$c \geq 0$ と「ほとんど逆」である写像
$K(n) \to K_n/I^{n - c}K_n$ が存在する。
$I^{n - c}K_n \subset K_n[I^c]$ だから、
% P02-E1157: source lines 29227-29229 omit the finite verb before “which witness”.
これらの写像は、系 $\{K(n)\}_{n \geq 1}$ と
$\{K_n\}_{n \geq 1}$ が $\mathcal{C}$ で同型であることを示す。

\medskip\noindent
系
$$
\{\Ext^i_{A_n}(K(n), N_n)\}_{n \geq 1}
\quad\text{and}\quad
\{\Ext^i_{A_n}(K_n, N_n)\}_{n \geq 1}
$$
は圏 $\mathcal{C}$ で同型であると主張する。
実際、全射 $K_n \to K(n)$ の核は $I^c$ で零化されるので、
これらは写像
$$
\Ext^i_{A_n}(K(n), N_n) \to \Ext^i_{A_n}(K_n, N_n)
$$
を定め、その核と余核は $I^c$ で零化される。
従って Lemma \ref{more-algebra-lemma-iso} により主張が従う。

\medskip\noindent
$i \geq 2$ に対して同型
$$
\Ext^{i - 1}_A(K, N) = \Ext^i_A(M, N)
\quad\text{and}\quad
\Ext^{i - 1}_{A_n}(K(n), N_n) = \Ext^i_{A_n}(M_n, N_n)
$$
がある。このようにして、$i = 0, 1$ の場合に補題を証明すれば十分だと分かる。

\medskip\noindent
$i = 0, 1$ に対して可換図式
$$
\xymatrix{
0 \ar[r] &
\Hom(M, N) \ar[r] \ar[dd] &
N^{\oplus r} \ar[r]_-\varphi \ar[dd] &
\Hom(K, N) \ar[r] \ar[d] &
\Ext^1(M, N) \ar[r] &
0 \\
& & &
\Hom(K_n, N_n)
\\
0 \ar[r] &
\Hom(M_n, N_n) \ar[r] &
N_n^{\oplus r} \ar[r] &
\Hom(K(n), N_n) \ar[r] \ar[u] &
\Ext^1(M_n, N_n) \ar[r] &
0
}
$$
を考える。Lemma \ref{more-algebra-lemma-hom-systems-ML} により、
% P02-E1158: source lines 29277-29280 leave an extra conjunction before the predicate.
$\Hom(M, N)/I^n \Hom(M, N) \to \Hom(M_n, N_n)$ と\\
$\Hom(K, N)/I^n \Hom(K, N) \to \Hom(K_n, N_n)$
の核と余核は、$n$ に依存しないある $c \geq 0$ に対して
$I^c$-ねじれである。
% P02-E1159: source lines 29281-29282 use a singular subject with the plural predicate “are”.
上で、単射
$\Hom(K(n), N_n) \to \Hom(K_n, N_n)$ の余核は、
必要なら $c$ を大きくすることにより $I^c$ で零化されることを見た。
このような $c$ に対して、図式に収まる写像
$\delta_n : I^c\Hom(K, N)/I^n\Hom(K, N) \to \Hom(K(n), N_n)$
を得る（正確な定式化は省略する）。
$\Hom(K, N)/I^n \Hom(K, N) \to \Hom(K_n, N_n)$ と
$\Hom(K(n), N_n) \to \Hom(K_n, N_n)$
について同じことが成り立つと分かっているので、
必要なら $c$ を大きくすることにより、
$\delta_n$ の核と余核は $I^c$ で零化される。
そこで実線からなる図式
{\scriptsize
$$
\xymatrix{
\varphi^{-1}(I^c\Hom(K, N)) \ar[r]_-\varphi \ar[d] &
I^c\Hom(K, N)/I^n\Hom(K, N) \ar[r] \ar[d]^{\delta_n} &
I^c\Ext^1(M, N)/I^n\Ext^1(M, N) \ar[r] \ar@{..>}[d] & 0 \\
N_n^{\oplus r} \ar[r] &
\Hom(K(n), N_n) \ar[r] &
\Ext^1(M_n, N_n) \ar[r] & 0
}
$$
}
の可換性を用いて点線の矢印を定められる。
省略した素直な図式追跡により、最後にもう一度
$c$ を大きくすれば、点線の矢印の核と余核は
% P02-E1160: source line 29303 misspells “by” as “buy”.
$I^c$ で零化されることが分かる。
\end{proof}

\begin{remark}
\label{more-algebra-remark-awkward}
Lemma \ref{more-algebra-lemma-dejong-kollar-kovacs} の主張が扱いにくいのは、
一つには $n$ を動かしたときの加群
$\Ext^i_{A_n}(M_n, N_n)$ の間に明らかな写像がないためである。
補題から結論できるのは、ある $c \geq 0$ が存在し、
$m \gg n \gg 0$ に対して（標準的な）写像
% P02-E1161: source lines 29316-29317 form Ext over A_n with M_m and N_m, which are not generally A_n-modules when m is larger than n.
$$
I^c\Ext^i_{A_n}(M_m, N_m)/I^n\Ext^i_{A_n}(M_m, N_m) \to
\Ext^i_{A_n}(M_n, N_n)/\Ext^i_{A_n}(M_n, N_n)[I^c]
$$
があり、その核と余核が $I^c$ で零化されるということである。
これは加群の系が得られるという（弱い）意味である。
\end{remark}

\begin{example}
\label{more-algebra-example-has-to-be-awkward}
$k$ を体とする。$A = k[[x, y]]/(xy)$ とおく。
記号の濫用により、$A$ における $x$ と $y$ の像も
$x$ と $y$ で表す。$I = (x)$、$M = A/(y)$ とおく。自由分解
$$
\ldots \to
A \xrightarrow{y}
A \xrightarrow{x}
A \xrightarrow{y} A \to M \to 0
$$
がある。従って
$$
\Ext^2_A(M, N) = N[y]/xN
$$
である。ここで $N[y] = \Ker(y : N \to N)$ である。
$A_n = A/I^n$、$M_n = M/I^nM$、$N_n = N/I^nN$ と書く。
各 $n$ に対して自由分解
$$
\ldots \to
A_n^{\oplus 2} \xrightarrow{y, x^{n - 1}}
A_n \xrightarrow{x}
A_n \xrightarrow{y} A_n \to M_n \to 0
$$
がある。従って
$$
\Ext^2_{A_n}(M_n, N_n) = (N_n[y] \cap N_n[x^{n - 1}])/xN_n
$$
である。ここで $N_n[y] = \Ker(y : N_n \to N_n)$ であり、
% P02-E1162: source line 29352 uses N instead of N_n in the x-power torsion notation.
$N[x^{n - 1}] = \Ker(x^{n - 1} : N_n \to N_n)$ である。
$N = A/(y)$ と取る。このとき
$$
\Ext^2_A(M, N) = N[y]/xN = N/xN \cong k
$$
であるが、
$$
\Ext^2_{A_n}(M_n, N_n) = (N_n[y] \cap N_n[x^{n - 1}])/xN_n =
N_n[x^{n - 1}]/xN_n = 0
$$
である。
% P02-E1163: source line 29362 says “for all r” although the system is indexed by n.
これはすべての $r$ に対して成り立つ。実際
$N_n = k[x]/(x^n)$ であり、列
$$
N_n \xrightarrow{x} N_n \xrightarrow{x^{n - 1}} N_n
$$
は完全である。従って Lemma
\ref{more-algebra-lemma-dejong-kollar-kovacs} のような結果を得るには、
何らかの種類の $I$-べきねじれを無視する必要がある。
\end{example}

\begin{lemma}
\label{more-algebra-lemma-not-awkward}
\begin{reference}
2018年2月23日付の Janos Kollar、Sandor Kovacs、Johan de Jong
の間の電子メールによる書簡。
\end{reference}
$A \to B$ を Noether 環の平坦な準同型とする。
$I \subset A$ をイデアルとし、$M, N$ を $A$-加群とする。
$B_n = B/I^nB$、$M_n = M/I^nM$、$N_n = N/I^nN$ とおく。
$M$ が $A$ 上平坦なら、すべての $i \in \mathbf{Z}$ に対して
$$
\lim \Ext^i_B(M, N)/I^n \Ext^i_B(M, N) =
\lim \Ext^i_{B_n}(M_n, N_n)
$$
が成り立つ。
% P02-E1164: the source states that M and N are A-modules, but Ext over B and the finite-free B-resolution require at least finite B-module structures.
\end{lemma}

\begin{proof}
分解
$$
\ldots \to P_2 \to P_1 \to P_0 \to M \to 0
$$
を有限自由 $B$-加群 $P_i$ により選ぶ。
% P02-E1165: source line 29392 misspells “modules” as “modues”.
$P_{i, n} = P_i/I^nP_i$ とおく。
$M$ と $B$ は $A$ 上平坦なので、列
$$
\ldots \to P_{2, n} \to P_{1, n} \to P_{0, n} \to M_n \to 0
$$
は完全である。一方では複体
$$
\Hom_B(P_0, N) \to \Hom_B(P_1, N) \to \Hom_B(P_2, N) \to \ldots
$$
が加群 $\Ext^i_B(M, N)$ を計算し、他方では複体
$$
\Hom_{B_n}(P_{0, n}, N_n) \to \Hom_{B_n}(P_{1, n}, N_n) \to
\Hom_{B_n}(P_{2, n}, N_n) \to \ldots
$$
が加群 $\Ext^i_{B_n}(M_n, N_n)$ を計算することが分かる。
$$
\Hom_{B_n}(P_{i, n}, N_n) = \Hom_B(P_i, N)/I^n \Hom_B(P_i, N)
$$
だから、Lemma \ref{more-algebra-lemma-consequence-Artin-Rees} の (2) から
結論を得る。
\end{proof}








\section{加群の系、再論}
\label{more-algebra-section-systems-bis}

\noindent
$I$ を Noether 環 $A$ のイデアルとする。
Section \ref{more-algebra-section-systems} では、有限 $A$-加群 $M$ に対する
$M/I^nM$ の形の系を考えると何が起こるかを検討した。
この節では代わりに系 $I^nM$ を考える。

\begin{lemma}
\label{more-algebra-lemma-consequence-Artin-Rees-bis}
$I$ を Noether 環 $A$ のイデアルとする。
$$
K \xrightarrow{\alpha} L \xrightarrow{\beta} M
$$
を有限 $A$-加群の複体とする。$H = \Ker(\beta)/\Im(\alpha)$ とおく。
$n \geq 0$ に対して
$$
I^nK \xrightarrow{\alpha_n} I^nL \xrightarrow{\beta_n} I^nM
$$
を誘導された複体とし、$H_n = \Ker(\beta_n)/\Im(\alpha_n)$ とおく。
このとき標準的な $A$-加群写像
$$
\ldots \to H_3 \to H_2 \to H_1 \to H
$$
がある。ある $c > 0$ が存在し、$n \geq c$ に対して
$H_n \to H$ の像は $I^{n - c}H$ に含まれ、標準的な
$A$-加群写像 $I^nH \to H_{n - c}$ があって、合成
$$
I^n H \to H_{n - c} \to  I^{n - 2c}H
\quad\text{and}\quad
H_n \to I^{n - c}H \to H_{n - 2c}
$$
はいずれも標準的な写像である。特に逆系 $(H_n)$ と
$(I^nH)$ は $\text{Mod}_A$ の pro-対象として同型である。
\end{lemma}

\begin{proof}
$H_n = \Ker(\beta) \cap I^nL/\alpha(I^nK)$ である。
$\Ker(\beta) \cap I^nL \subset \Ker(\beta) \cap I^{n - 1}L$
かつ $\alpha(I^nK) \subset \alpha(I^{n - 1}K)$ なので、
写像 $H_n \to H_{n - 1}$ を得る。$H_1 \to H$ についても同様である。

\medskip\noindent
Artin--Rees により、$n \geq c_1$ に対して
$\Im(\alpha) \cap I^nL \subset \alpha(I^{n - c_1}K)$、また
$n \geq c_2$ に対して
$\Ker(\beta) \cap I^nL \subset I^{n - c_2}\Ker(\beta)$
となる $c_1, c_2 \geq 0$ を選べる。Algebra, Lemmas
\ref{algebra-lemma-map-AR} および
\ref{algebra-lemma-Artin-Rees} を参照せよ。$c = c_1 + c_2$ とおく。

\medskip\noindent
$c \geq c_2$ の選び方から、$n \geq c$ に対して
$H_n \to H$ の像が $I^{n - c}H$ に含まれることは直ちに従う。

\medskip\noindent
$n \geq c$ とする。$\psi_n : I^nH \to H_{n - c}$ を次のように定める。
$x \in I^nH$ とし、$x$ を表す
$y \in I^n\Ker(\beta)$ を選ぶ。$\psi_n(x)$ を
$H_{n - c}$ における $y$ の類と定める。
これが well-defined であることを見るため、同じ $x$ に対する
別の選択 $y'$ を取る。このとき $y' - y \in \Im(\alpha)$ である。
$c \geq c_1$ の選び方から
$y' - y \in \alpha(I^{n - c}K)$ と結論できるので、
$y$ と $y'$ は $H_{n - c}$ の同じ元を表す。
従って $\psi_n$ は well-defined である。

\medskip\noindent
合成 $I^n H \to H_{n - c} \to  I^{n - 2c}H$ および
$H_n \to I^{n - c}H \to H_{n - 2c}$ に関する主張は、
定義から直ちに従う。
\end{proof}

\begin{lemma}
\label{more-algebra-lemma-ext-factors}
$A$ を Noether 環、$I \subset A$ をイデアルとする。
$M$, $N$ を $A$-加群とし、$M$ は有限であるとする。
各 $p > 0$ に対し、ある $c \geq 0$ が存在し、$n \geq c$ なら
写像 $\Ext_A^p(M, N) \to \Ext_A^p(I^nM, N)$ は
$\Ext^p_A(I^nM, I^{n - c}N) \to \Ext_A^p(I^nM, N)$
を経由して分解する。
\end{lemma}

\begin{proof}
$p = 0$ の場合、$\varphi : M \to N$ を $A$-線型写像とすると、
$f_i \in A$、$m_i \in M$ に対して
$\varphi(\sum f_i m_i) = \sum f_i \varphi(m_i)$ である。
従って $\varphi$ はすべての $n$ に対して写像
$I^nM \to I^nN$ を誘導し、$c = 0$ として結果が成り立つ。

\medskip\noindent
短完全列 $0 \to K \to A^{\oplus t} \to M \to 0$ を選ぶ。
各 $n$ に対して短完全列
$0 \to L_n \to A^{\oplus s_n} \to I^nM \to 0$ を選ぶ。
$A^{\oplus s_n} \to A^{\oplus t}$ の像が
$(I^n)^{\oplus t}$ に含まれるような短完全列の写像
$$
\xymatrix{
0 \ar[r] &
L_n \ar[r] \ar[d] &
A^{\oplus s_n} \ar[r] \ar[d] &
I^nM \ar[r] \ar[d] & 0 \\
0 \ar[r] &
K \ar[r] &
A^{\oplus t} \ar[r] &
M \ar[r] & 0
}
$$
を構成できることは明らかである。
Artin--Rees（Algebra, Lemma
\ref{algebra-lemma-Artin-Rees}）により、
ある $c \geq 0$ が存在し、$n \geq c$ なら
$L_n \to K$ は $I^{n - c}K$ を経由して分解する。

\medskip\noindent
$p = 1$ に対して、上の選択は実線からなる可換図式
{\footnotesize
$$
\xymatrix{
\Hom_A(A^{\oplus s_n}, N) \ar[r] &
\Hom_A(L_n, N) \ar[r] &
\Ext_A^1(I^nM, N) \ar[r] & 0 \\
\Hom_A((I^n)^{\oplus t}, I^{n - c}N) \ar[r] \ar[u] &
\Hom_A(K \cap (I^n)^{\oplus t}, I^{n - c}N) \ar[r] \ar[u] &
\Ext_A^1(I^nM, I^{n - c}N) \ar[u] \\
\Hom_A(A^{\oplus t}, N) \ar[r] \ar[u] &
\Hom_A(K, N) \ar[r] \ar[u] &
\Ext_A^1(M, N) \ar@{..>}[u] \ar[r] & 0
}
$$
}
を誘導し、その横向きの矢印は完全である。
下段中央の縦矢印が存在するのは
$K \cap (I^n)^{\oplus t} \subset I^{n - c}K$ であり、
% P02-E1166: source lines 29545-29546 names the whole ambient power instead of its intersection with K as the induced map's domain.
従って第1段落の議論により、任意の $A$-線型写像 $K \to N$ が
$A$-線型写像 $(I^n)^{\oplus t} \to I^{n - c}N$
を誘導するからである。これにより望む点線の矢印を得る。

\medskip\noindent
$p > 1$ に対して可換図式
$$
\xymatrix{
\Ext^{p - 1}_A(I^{n - c}K, N) \ar[r] &
\Ext^{p - 1}_A(L_n, N) \ar[r] &
\Ext_A^p(I^nM, N) \\
\Ext^{p - 1}_A(K, N) \ar[rr] \ar[u] & &
\Ext_A^p(M, N) \ar[u]
}
$$
を得る。下段の横向きの矢印は同型である。
$p$ に関する帰納法により、ある $c' \geq 0$ と
すべての $n \geq c + c'$ に対して、左の縦写像は
$\Ext^{p - 1}_A(I^{n - c}K, I^{n - c - c'}N)$ を経由して分解する。
合成
$\Ext^{p - 1}_A(I^{n - c}K, I^{n - c - c'}N) \to
\Ext^{p - 1}_A(L_n, I^{n - c - c'}N) \to \Ext^p_A(I^nM, I^{n - c - c'}N)$
を用いると、望む分解を得る（$n \geq c + c'$ とし、
$c$ を $c + c'$ で置き換える）。
\end{proof}

\begin{lemma}
\label{more-algebra-lemma-ext-annihilated}
$A$ を Noether 環、$I \subset A$ をイデアルとする。
$M$, $N$ を $A$-加群とし、$M$ は有限、$N$ は $I$ のあるべきで
零化されるとする。各 $p > 0$ に対し、写像
$\Ext_A^p(M, N) \to \Ext_A^p(I^nM, N)$ が零となるような
$n$ が存在する。
\end{lemma}

\begin{proof}
Lemma \ref{more-algebra-lemma-ext-factors} と、ある $m > 0$ に対して
$I^mN = 0$ であることから直ちに従う。
\end{proof}

\begin{lemma}
\label{more-algebra-lemma-ext-induced-toplogy}
$A$ を Noether 環、$I \subset A$ をイデアルとする。
$K \in D(A)$ を擬コヒーレントとし、$M$ を有限 $A$-加群とする。
% P02-E1167: source line 29587 uses “an c” instead of “a c”.
各 $p \in \mathbf{Z}$ に対し、ある $c$ が存在し、$n \geq c$ なら
$\Ext_A^p(K, I^nM) \to \Ext_A^p(K, M)$ の像は
$I^{n - c}\Ext_A^p(K, M)$ に含まれる。
\end{lemma}

\begin{proof}
$K$ を表す有限自由 $A$-加群の上に有界な複体
$P^\bullet$ を選ぶ。このとき $\Ext_A^p(K, M)$ は
% P02-E1168: source lines 29593-29598 choose P but then use the undefined complex F in all three Hom terms.
$$
\Hom_A(F^{-p + 1}, M) \xrightarrow{a}
\Hom_A(F^{-p}, M) \xrightarrow{b}
\Hom_A(F^{-p - 1}, M)
$$
のコホモロジーであり、$\Ext_A^p(K, I^nM)$ はこれらの
有限 $A$-加群を各々その $I^n$ 倍で置き換えて計算される。
% P02-E1169: source lines 29601-29603 leave “Thus the result by Lemma ...” without a finite verb.
従って結果は Lemma
\ref{more-algebra-lemma-consequence-Artin-Rees-bis} から従う
（しかも、はるかに多くのことが成り立つ）。
\end{proof}

\noindent
Situation \ref{more-algebra-situation-koszul} において、複体 $I_n^\bullet$ を
三角圏 $K(A)$、すなわちホモトピーを除いた $A$-加群の複体の圏に
おける識別三角形
$$
I_n^\bullet \to A \to K_n^\bullet \to I_n^\bullet[1]
$$
が存在するように定める。具体的には
$I_n^\bullet = \sigma_{\leq -1}K_n^\bullet[-1]$ とおく。
複体の各次数で分裂する短完全列
$$
0 \to A \to K_n^\bullet \to I_n^\bullet[1] \to 0
$$
があり、これは $K(A)$ の三角圏構造の定義により識別三角形を定める。
それらを回転させると上の望む識別三角形を得る。
$I_n^0 = A^{\oplus r} \to A$ は第 $i$ 成分上で
$f_i^n$ を掛ける写像によって与えられることに注意する。
従って $I_n^\bullet \to A$ は
$$
I_n^\bullet \to (f_1^n, \ldots, f_r^n) \to A
$$
と分解する。実際、短完全列
$$
0 \to H^{-1}(K_n^\bullet) \to H^0(I_n^\bullet) \to (f_1^n, \ldots, f_r^n) \to 0
$$
があり、すべての $i < 0$ に対して
$H^i(I_n^\bullet) = H^{i - 1}(K_n^\bullet)$ である。
写像 $K_{n + 1}^\bullet \to K_n^\bullet$ は
$I_{n + 1}^\bullet \to I_n^\bullet$ を誘導し、$K(A)$ における可換図式
$$
\xymatrix{
\ldots \ar[r] &
I_3^\bullet \ar[d] \ar[r] &
I_2^\bullet \ar[d] \ar[r] &
I_1^\bullet \ar[d] \\
\ldots \ar[r] &
(f_1^3, \ldots, f_r^3) \ar[r] &
(f_1^2, \ldots, f_r^2) \ar[r] &
(f_1, \ldots, f_r)
}
$$
を得る。

\begin{lemma}
\label{more-algebra-lemma-sequence-powers-pro-bounded}
Situation \ref{more-algebra-situation-koszul} において $A$ が Noether 環であるとする。
上の記号のもとで、逆系 $(I^n)$ は $D(A)$ において逆系
$(I_n^\bullet)$ と pro-同型である。
\end{lemma}

\begin{proof}
逆系 $I^n$ が $A$-加群の圏において逆系
$(f_1^n, \ldots, f_r^n)$ と pro-同型であることは初等的に示せる。
識別三角形の逆系
$$
I_n^\bullet \to (f_1^n, \ldots, f_r^n) \to C_n^\bullet \to I_n^\bullet[1]
$$
を考える。ここで $C_n^\bullet$ は最初の矢印の錐である。
Derived Categories, Lemma
\ref{derived-lemma-pro-isomorphism} により、
逆系 $C_n^\bullet$ が pro-零であることを示せば十分である。
複体 $I_n^\bullet$ の非零な項は
$-r + 1 \leq i \leq 0$ を満たす次数 $i$ にしかないので、
$C_n^\bullet$ も同様に有界である。従って Derived Categories, Lemma
\ref{derived-lemma-essentially-constant-cohomology} により、
$H^p(C_n^\bullet)$ が pro-零であることを示せば十分である。
上の議論から、$p \leq -1$ に対して
$H^p(C_n^\bullet) = H^p(K_n^\bullet)$ であり、
$p \geq 0$ に対して $H^p(C_n^\bullet) = 0$ である。
逆系 $H^p(K_n^\bullet)$ が pro-零であることは
Lemma \ref{more-algebra-lemma-sequence-Koszul-complexes} の証明で示した
（ここで $A$ が Noether 環であるという仮定を用いる）。
\end{proof}

\begin{lemma}
\label{more-algebra-lemma-tensoring-Deligne-system}
$A$ を Noether 環、$I \subset A$ をイデアルとする。
$M^\bullet$ を有限 $A$-加群の有界複体とする。写像の逆系
$$
I^n \otimes_A^\mathbf{L} M^\bullet \longrightarrow I^nM^\bullet
$$
は $D(A)$ の pro-対象の同型を定める。
\end{lemma}

\begin{proof}
$I$ の生成元 $f_1, \ldots, f_r \in I$ を選ぶ。
逆系 $I^n$ は $A$-加群の圏において逆系
$(f_1^n, \ldots, f_r^n)$ と pro-同型である。
Lemma \ref{more-algebra-lemma-sequence-powers-pro-bounded} の記号のもとで、
写像の逆系
$$
I_n^\bullet \otimes_A^\mathbf{L} M^\bullet
\longrightarrow
(f_1^n, \ldots, f_r^n)M^\bullet
$$
が $D(A)$ の pro-対象の同型を定めることを証明すれば十分だと分かる。
$i \not \in [a, b]$ なら $M^i = 0$ となるような
$a \leq b$ があるとする。このとき上の矢印の始域と終域の
コホモロジーは次数 $[-r + a, b]$ にしか存在しない。
従って任意の $p \in \mathbf{Z}$ に対して写像の逆系
$$
H^p(I_n^\bullet \otimes_A^\mathbf{L} M^\bullet)
\longrightarrow
H^p((f_1^n, \ldots, f_r^n)M^\bullet)
$$
が $A$-加群の pro-対象の同型を定めることを示せば十分である。
Derived Categories, Lemma
\ref{derived-lemma-pro-isomorphism-bis} を参照せよ。
$I_n^\bullet \otimes_A^\mathbf{L} M^\bullet$ と
$I^n \otimes_A^\mathbf{L} M^\bullet$ の間の pro-同型、および
$(f_1^n, \ldots, f_r^n)M^\bullet$ と $I^nM^\bullet$
の間の pro-同型を用いると、これは写像の逆系
$$
H^p(I^n \otimes_A^\mathbf{L} M^\bullet)
\longrightarrow
H^p(I^nM^\bullet)
$$
が $A$-加群の pro-対象の同型を定めることと同値である。
% P02-E1170: source line 29722 omits the sentence-ending period before “Choose”.
有限自由 $A$-加群の上に有界な複体 $P^\bullet$ と擬同型
$P^\bullet \to M^\bullet$ を選ぶ。このとき、写像の逆系
$$
H^p(I^nP^\bullet)
\longrightarrow
H^p(I^nM^\bullet)
$$
が pro-同型であることを示せば十分である。
$H^p(P^\bullet) = H^p(M^\bullet)$ なので、これは Lemma
\ref{more-algebra-lemma-consequence-Artin-Rees-bis} から従う。
\end{proof}

\begin{lemma}
\label{more-algebra-lemma-factor-through-derived-tensor-product}
$A$ を Noether 環、$I \subset A$ をイデアルとする。
$M$ を有限 $A$-加群とする。$D(A)$ において
$I^nM \to M$ が写像 $I \otimes_A^\mathbf{L} M \to M$
を経由して分解するような整数 $n > 0$ が存在する。
\end{lemma}

\begin{proof}
これは Lemma \ref{more-algebra-lemma-tensoring-Deligne-system} から従う。
% P02-E1171: source line 29746 uses the malformed phrase “can also been seen”.
次のように直接見ることもできる。識別三角形
$$
I \otimes_A^\mathbf{L} M \to M \to A/I \otimes_A^\mathbf{L} M \to
I \otimes_A^\mathbf{L} M[1]
$$
を考える。三角圏の公理により、ある $n$ に対して
$I^nM \to A/I \otimes_A^\mathbf{L} M$ が $D(A)$ で零であることを
証明すれば十分である。$I$ の生成元 $f_1, \ldots, f_r$ を選び、
$K = K_\bullet(A, f_1, \ldots, f_r)$ を Koszul 複体とする。
商写像の分解 $A \to K \to A/I$ を考える。
すると、ある $n > 0$ に対して
$I^nM \to K \otimes_A M$ が $D(A)$ で零であることを示せば
十分だと分かる。ある $n > 0$ が見つかり、
$I^nM \to K \otimes_A M$ が $D(A)$ で
$\tau_{\geq t}(K \otimes_A M)$ を経由して分解すると仮定する。
$\tau_{\geq t + 1}(K \otimes_A M)$ を経由して分解することへの障害は
$\Ext^t(I^nM, H_t(K \otimes_A M))$ の元である。
有限 $A$-加群 $H_t(K \otimes_A M)$ は $I$ で零化される。
従って Lemma \ref{more-algebra-lemma-ext-annihilated} により、
$n$ を大きくした後、この障害元が零であると仮定できる。
これを有限回繰り返すと、
$I^nM \to K \otimes_A M$ が $D(A)$ で
$0 = \tau_{\geq r + 1}(K \otimes_A M)$ を経由して分解するような
$n$ を得て、結論が従う。
\end{proof}






\section{雑録}
\label{more-algebra-section-misc}

\noindent
他のどこにも収まらないいくつかの結果を述べる。

\begin{lemma}
\label{more-algebra-lemma-ext-annihilated-into}
$A$ を Noether 環、$I \subset A$ をイデアルとする。
$K \in D(A)$ を擬コヒーレントとし、$a \in \mathbf{Z}$ とする。
すべての有限 $A$-加群 $M$ に対して、$i \geq a$ なら
$\Ext^i_A(K, M)$ が $I$-べきねじれ加群であると仮定する。
このとき $i \geq a$ かつ $M$ が有限なら、系
$\Ext^i_A(K, M/I^nM)$ は
$$
\Ext^i_A(K, M) = \lim \Ext^i_A(K, M/I^nM)
$$
を値として本質的に定値である。
\end{lemma}

\begin{proof}
$M$ を有限 $A$-加群とする。$K$ は擬コヒーレントだから
$\Ext^i_A(K, M)$ は有限 $A$-加群である。
従って $i \geq a$ なら、ある $t \geq 0$ に対して
$I^t$ で零化される。Lemma
\ref{more-algebra-lemma-ext-induced-toplogy} により、ある $n > 0$ に対して
$\Ext^i_A(K, I^nM) \to \Ext^i_A(K, M)$ の像は零である。
短完全列 $0 \to I^nM \to M \to M/I^n M \to 0$ から長完全列
$$
\Ext^i_A(K, I^nM) \to \Ext^i_A(K, M) \to
\Ext^i_A(K, M/I^nM) \to \Ext^{i + 1}_A(K, I^nM)
$$
を得る。先ほど述べたことを有限 $A$-加群 $I^mM$ に適用すると、
系 $\Ext^i_A(K, I^nM)$ および
$\Ext^{i + 1}_A(K, I^nM)$ は値 $0$ で本質的に定値である。
図式追跡により、$\Ext^i_A(K, M/I^nM)$ は
$\Ext^i_A(K, M)$ を値として本質的に定値である。
\end{proof}

\begin{lemma}
\label{more-algebra-lemma-tor-annihilated}
$A$ を Noether 環、$I \subset A$ をイデアルとする。
$M$ を有限 $A$-加群、$N$ を $I$ で零化される $A$-加群とする。
すべての $p \geq 0$ に対して
$\text{Tor}^A_p(I^nM, N) \to \text{Tor}^A_p(M, N)$ が零となるような
整数 $n > 0$ が存在する。
\end{lemma}

\begin{proof}
Lemma \ref{more-algebra-lemma-factor-through-derived-tensor-product} により、
$I^nM \to M$ を
$I^nM \to M \otimes_A^\mathbf{L} I \to M$ と分解できる。
合成
$$
I^nM \otimes_A^\mathbf{L} N \to
(M \otimes_A^\mathbf{L} I) \otimes_A^\mathbf{L} N
\to M \otimes_A^\mathbf{L} N
$$
は零であると主張する。実際、図式
$$
\xymatrix{
(M \otimes_A^\mathbf{L} I) \otimes_A^\mathbf{L} N \ar[rr] \ar[rd] & &
M \otimes_A^\mathbf{L} (I \otimes_A^\mathbf{L} N) \ar[ld] \\
& M \otimes_A^\mathbf{L} N
}
$$
は可換であり（詳細は省略する）、$N$ は $I$ で零化されるので
写像 $I \otimes_A^\mathbf{L} N \to N$ は零である。
\end{proof}

\begin{lemma}
\label{more-algebra-lemma-pseudo-coherent-tensor-limit}
$R$ を環、$K \in D(R)$ を擬コヒーレントとする。
$(M_n)$ を $R$-加群の逆系とする。このとき
$R\lim K \otimes_R^\mathbf{L} M_n = K \otimes_R^\mathbf{L} R\lim M_n$
である。
\end{lemma}

\begin{proof}
定義する識別三角形
$$
R\lim M_n \to \prod M_n \to \prod M_n \to R\lim M_n[1]
$$
を考え、Lemma \ref{more-algebra-lemma-pseudo-coherent-tensor} を適用する。
\end{proof}

\begin{lemma}
\label{more-algebra-lemma-additivity-of-pd}
$R$ を Noether 局所環、$I \subset R$ をイデアルとし、
$E$ を非零な有限 $R/I$-加群とする。
$R/I$ が有限射影次元をもち、$E$ が $R/I$ 上有限射影次元をもつなら、
$E$ は $R$ 上有限射影次元をもち、
$$
\text{pd}_R(E) = \text{pd}_R(R/I) + \text{pd}_{R/I}(E)
$$
である。
\end{lemma}

\begin{proof}
有限加群について、$R$ 上、resp.\ $R/I$ 上で有限射影次元をもつことは
完全加群であることと同値であるという事実を用いる。
Definition \ref{more-algebra-definition-perfect} に続く議論を参照せよ。
Lemma \ref{more-algebra-lemma-cohomology-perfect} により、
$E$ は $R$ 上有限射影次元をもつ。
従って Auslander--Buchsbaum（Algebra, Proposition
\ref{algebra-proposition-Auslander-Buchsbaum}）を適用して
$$
\text{pd}_R(E) + \text{depth}(E) = \text{depth}(R),\quad
\text{pd}_{R/I}(E) + \text{depth}(E) = \text{depth}(R/I),
$$
および
$$
\text{pd}_R(R/I) + \text{depth}(R/I) = \text{depth}(R)
$$
を得る。最初の等式では $E$ の深さを $R$-加群として取り、
第2の等式では $R/I$-加群として取ることに注意する。
しかし、これらの深さは等しい
（これは自明であるが、Algebra, Lemma
\ref{algebra-lemma-depth-goes-down-finite} からも従う）。
これで証明が終わる。
\end{proof}

\begin{lemma}
\label{more-algebra-lemma-enlarge}
$A \to B$ を環写像とする。次の性質をもつ基数
$\kappa = \kappa(A \to B)$ が存在する。
$M^\bullet$, resp.\ $N^\bullet$ を
$A$-加群, resp.\ $B$-加群の複体とする。
$a : M^\bullet \to N^\bullet$ を $A$-加群の複体の写像とし、
$D(B)$ において同型
$M^\bullet \otimes_A^\mathbf{L} B \to N^\bullet$ を誘導するとする。
$M_1^\bullet \subset M^\bullet$, resp.\
$N_1^\bullet \subset N^\bullet$ を
$A$-加群, resp.\ $B$-加群の部分複体とし、
$a(M_1^\bullet) \subset N_1^\bullet$ とする。
このとき部分複体
$$
M_1^\bullet \subset M_2^\bullet \subset M^\bullet
\quad\text{and}\quad
N_1^\bullet \subset N_2^\bullet \subset N^\bullet
$$
で $a(M_2^\bullet) \subset N_2^\bullet$ を満たし、
次の性質をもつものが存在する。
\begin{enumerate}
\item $\Ker(H^i(M_1^\bullet \otimes_A^\mathbf{L} B) \to H^i(N_1^\bullet))$
は $H^i(M_2^\bullet \otimes_A^\mathbf{L} B)$ で零へ写る。
\item $\Im(H^i(N_1^\bullet) \to H^i(N_2^\bullet))$ は
% P02-E1172: source line 29917 uses H^2 in the target of a statement quantified in degree i.
$\Im(H^i(M_2^\bullet \otimes_A^\mathbf{L} B) \to H^2(N_2^\bullet))$
に含まれる。
\item $|\bigcup M_2^i \cup \bigcup N_2^i| \leq
\max(\kappa, |\bigcup M_1^i \cup \bigcup N_1^i|)$ である。
\end{enumerate}
\end{lemma}

\begin{proof}
$\kappa = \max(|A|, |B|, \aleph_0)$ とおく。
$|M^\bullet| = |\bigcup M^i|$ とおき、他の複体についても同様に定める。
この記法のもとで、補題の主張で用いた量について
% P02-E1173: source lines 29927-29929 replaces N_1 by the not-yet-constructed M_2 in the cardinal identity.
$$
\max(\kappa, |\bigcup M_1^i \cup \bigcup N_1^i|) =
\max(\kappa, |M_1^\bullet|, |M_2^\bullet|)
$$
が成り立つ。以下ではこれと、基数算術から従う他の観察を
断りなく用いる。

\medskip\noindent
まず「小さい」部分複体が十分多く存在することを示す。
$E^i \subset M^i$、$i \in \mathbf{Z}$ である有限部分集合の族
$E = \{E^i\}$ と、$F^i \subset N^i$、$i \in \mathbf{Z}$ である
有限部分集合の族 $F = \{F^i\}$ の各組に対して、
$$
M_1^\bullet \subset M_1(E, F)^\bullet \subset M^\bullet
\quad\text{and}\quad
N_1^\bullet \subset N_1(E, F)^\bullet \subset N^\bullet
$$
を、$a(M_1(E, F)^\bullet) \subset N_1(E, F)^\bullet$、
$E^i \subset M_1(E, F)^i$、および
% P02-E1174: source line 29946 uses the undefined M_2(E,F) where the F-subsets belong to the N-side subcomplex.
$F^i \subset M_2(E, F)^i$ を満たす最小の
$A$-加群および $B$-加群の部分複体とする。
このとき
% P02-E1175: source line 29951 again uses undefined M_2(E,F) and M_2 instead of the N_1-side quantities.
$$
|M_1(E, F)^\bullet| \leq \max(\kappa, |M_1^\bullet|)
\quad\text{and}\quad
|M_2(E, F)^\bullet| \leq \max(\kappa, |M_2^\bullet|)
$$
であることは容易に分かる。詳細は省略する。明らかに
$$
M^\bullet = \colim_{(E, F)} M_1(E, F)^\bullet
\quad\text{and}\quad
N^\bullet = \colim_{(E, F)} N_1(E, F)^\bullet
$$
であり、これらの余極限は各次数でフィルター余極限である。

\medskip\noindent
自由 $A$-加群 $F_i$ による分解
% P02-E1176: source lines 29962-29963 writes the resolution with superscripted F but then names the free modules F_i with subscripts.
$\ldots \to F^{-1} \to F^0 \to B$ で
$|F_i| \leq \kappa$ となるものが存在する（詳細は省略する）。
$M_1^\bullet \otimes_A^\mathbf{L} B$ のコホモロジー加群は
$\text{Tot}(M_1^\bullet \otimes_A F^\bullet)$ によって計算される。
従って
$|H^i(M_1^\bullet \otimes_A^\mathbf{L} B)| \leq
\max(\kappa, |M_1^\bullet|)$ である。

\medskip\noindent
$i \in \mathbf{Z}$ とし、
$\xi \in H^i(M_1^\bullet \otimes_A^\mathbf{L} B)$ を
$H^i(N_1^\bullet)$ で零へ写る元とする。
このとき $\xi$ は $H^i(N^\bullet)$ で零へ写り、従って
$\xi$ は $H^i(M^\bullet \otimes_A^\mathbf{L} B)$ でも零へ写る。
導来テンソル積はフィルター余極限と可換だから、上のような有限族
$E_\xi$ と $F_\xi$ で、$\xi$ が
$H^i(M_1(E_\xi, F_\xi)^\bullet \otimes_A^\mathbf{L} B)$
で零へ写るものを見つけられる。

\medskip\noindent
$i \in \mathbf{Z}$ とし、$\eta \in H^i(N_1^\bullet)$ とする。
$H^i(N^\bullet)$ における $\eta$ の像は
$H^i(M^\bullet \otimes_A^\mathbf{L} B) \to H^i(N^\bullet)$
の像に属する。従って先ほどと同様に、上のような有限族
$E_\eta$ と $F_\eta$ で、$\eta$ が
% P02-E1177: source lines 29985-29988 omit the closing parenthesis and complex superscript in both N_1(E_eta,F_eta) cohomology terms.
$H^i(N_1(E_\eta, F_\eta)$ の元へ写り、その元が写像
$H^i(M_1(E_\eta, F_\eta)^\bullet \otimes_A^\mathbf{L} B) \to
H^i(N_1(E_\eta, F_\eta)$
の像に属するものを見つけられる。

\medskip\noindent
ここで単に
$$
M_2^\bullet =
\sum\nolimits_\xi M_1(E_\xi, F_\xi)^\bullet +
\sum\nolimits_\eta M_1(E_\eta, F_\eta)^\bullet
$$
と定める。和は前の2段落に現れる $\xi$ と $\eta$ にわたり、
$M^\bullet$ の内部で取る。同様に
$$
N_2^\bullet =
\sum\nolimits_\xi N_1(E_\xi, F_\xi)^\bullet +
\sum\nolimits_\eta N_1(E_\eta, F_\eta)^\bullet
$$
とおき、和は $N^\bullet$ の内部で取る。
構成により、これらの選択は性質 (1) と (2) を満たす。
$\xi$ と $\eta$ の集合の基数は高々
$\max(\kappa, |M_1^\bullet|, |N_1^\bullet|)$ だから、
上界 (3) も従う。
\end{proof}

\begin{lemma}
\label{more-algebra-lemma-f-bounded}
$R$ を環、$f \in R$ とする。$M \in D(R)$ とし、
$C$ を $f : M \to M$ の錐とする。
$i < 0$ に対して $H^i(M)_f = 0$ であり、
$i < -1$ に対して $H^i(C) = 0$ なら、
$i < 0$ に対して $H^i(M) = 0$ である。
\end{lemma}

\begin{proof}
$M_f = M \otimes_R^\mathbf{L} R_f$ と書き、識別三角形
$F \to M \to M_f$ を選ぶ。仮定により $i < 0$ に対して
$H^i(M_f) = H^i(M)_f = 0$ である。
従って $i < 0$ に対して $H^i(F) = 0$ であることを示せば十分である。
すべての $i \in \mathbf{Z}$ に対して $H^i(F)$ は
$f$-べきねじれであることに注意する。
一方、$f : M_f \to M_f$ は同型だから、
$C$ は $f : F \to F$ の錐と同型である
（Derived Categories, Proposition
\ref{derived-proposition-9} を用いる）。
さて $H^i(F) \not = 0$ なら
$f : H^i(F) \to H^i(F)$ の核は非零であり、
従って $H^{i - 1}(C)$ は非零である。
仮定により $i - 1 < -1$ のときこれは起こりえず、証明が終わる。
\end{proof}

\begin{lemma}
\label{more-algebra-lemma-f-flat}
$R$ を環、$f \in R$ とし、$M \in D(R)$ とする。次を仮定する。
\begin{enumerate}
\item $i > 0$ に対して $H^i(M) = 0$ である。
\item $M \otimes_R^\mathbf{L} R_f$ は次数 $0$ に置かれた
平坦 $R_f$-加群と同型である。
\item $M \otimes_R^\mathbf{L} R/fR$ は次数 $0$ に置かれた
平坦 $R/fR$-加群と同型である。
\end{enumerate}
このとき $M$ は次数 $0$ に置かれた平坦 $R$-加群と同型である。
\end{lemma}

\begin{proof}
$K$ を $R$-加群とする。Section \ref{more-algebra-section-tor} を参照すると、
$N = M \otimes_R^\mathbf{L} K$ が次数 $0$ に置かれた
$R$-加群と同型であることを示せば十分である。
仮定 (1) により、$N$ のコホモロジーは次数 $\leq 0$ にしか存在しない。
従って Lemma \ref{more-algebra-lemma-f-bounded} により、
$i < 0$ に対して $H^i(N)_f = 0$ であり、また $C$ を
$f : N \to N$ の錐としたとき $i < -1$ に対して
$H^i(C) = 0$ であることを示せば十分である。
前者について、
$$
H^i(N)_f = H^i(N \otimes_R^\mathbf{L} R_f) =
H^i(M \otimes_R^\mathbf{L} K \otimes_R^\mathbf{L} R_f ) =
H^i((M \otimes_R^\mathbf{L} R_f) \otimes_{R_f}^\mathbf{L}
(K \otimes_R^\mathbf{L} R_f))
$$
であることに注意する。仮定 (2) により、これは $i = 0$ の場合を除いて
零であり、その場合には $H^0(M)_f \otimes_R K$ を得る。
後者について、$C'$ を $f : K \to K$ の錐とする。このとき
$$
C = M \otimes_R^\mathbf{L} C'
$$
である。$C'$ の非零なコホモロジーは次数 $0$ と $-1$ にしかなく、
それぞれ $K/fK$ と $K[f]$ に等しい。言い換えると、識別三角形
$$
(K[f])[1] \to C' \to K/fK
$$
がある。一方、$f$ で零化される任意の $R$-加群 $K'$ に対して
$$
M \otimes_R^\mathbf{L} K' =
M \otimes_R^\mathbf{L} R/fR \otimes_{R/fR}^\mathbf{L} K'
$$
である。仮定 (3) により、これは次数 $0$ に置かれた加群
$H^0(M \otimes_R^\mathbf{L} R/fR) \otimes_{R/fR} K'$ に等しい。
以上を合わせると、$C$ の非零なコホモロジーは次数 $0$ と $-1$
にしかないことが分かり、証明が完了する。
\end{proof}







\section{二重複体に関する技巧}
\label{more-algebra-section-tricks}

\noindent
この節では Homology, Section
\ref{homology-section-double-complexes-abelian-groups}
の議論を続ける。

\begin{lemma}
\label{more-algebra-lemma-vanishing-coh-prod-totalization}
$A_0^\bullet \to A_1^\bullet \to A_2^\bullet \to \ldots$
を Abel 群の複体からなる複体とする。
すべての $p \geq 0$ に対して $H^{-p}(A_p^\bullet) = 0$ と仮定する。
$A^{p, q} = A_p^q$ とおき、$A^{\bullet, \bullet}$ を二重複体とみなす。
このとき $H^0(\text{Tot}_\pi(A^{\bullet, \bullet})) = 0$ である。
\end{lemma}

\begin{proof}
% P02-E1178: source line 30105 omits “by” in “Denote f_p ... the given maps”.
$f_p : A_p^\bullet \to A_{p + 1}^\bullet$ で与えられた複体の写像を表す。
$\text{Tot}_\pi(A^{\bullet, \bullet})$ の次数 $n$ の元上の微分は
$$
\prod\nolimits_{p + q = n} (f^q_p + (-1)^p\text{d}^q_{A_p^\bullet})
$$
で与えられることを思い出す。
$\xi \in H^0(\text{Tot}_\pi(A^{\bullet, \bullet}))$ を
コホモロジー類とする。$\xi$ が零であることを示す。
% P02-E1179: source line 30113 uses “an cocycle” instead of “a cocycle”.
$\xi$ をコサイクル
$x = (x_p) \in \prod A^{p, -p}$ の類として表す。
$\text{d}(x) = 0$ だから
$\text{d}_{A_0^\bullet}(x_0) = 0$ である。
$H^0(A_0^\bullet) = 0$ だから
$\text{d}_{A_0^\bullet}(y_{-1}) = x_0$ となる
$y_{-1} \in A^{0, -1}$ が存在する。
すると
$\text{d}_{A_1^\bullet}(x_1 + f_0(y_{-1})) = 0$ である。
$H^{-1}(A_1^\bullet) = 0$ だから
$-\text{d}_{A_1^\bullet}(y_{-2}) = x_1 + f_0(y_{-1})$
となる $y_{-2} \in A^{1, -2}$ を見つけられる。
帰納法により
$$
(-1)^p\text{d}_{A_p^\bullet}(y_{-p - 1}) = x_p + f_{p - 1}(y_{-p})
$$
となる $y_{-p - 1} \in A^{p, -p - 1}$ を見つけられる。
これは $y = (y_{-p - 1})$ としたとき
$\text{d}(y) = x$ であることを意味する。
\end{proof}

\begin{lemma}
\label{more-algebra-lemma-prod-qis-gives-qis}
$$
(A_0^\bullet \to A_1^\bullet \to A_2^\bullet \to \ldots)
\longrightarrow
(B_0^\bullet \to B_1^\bullet \to B_2^\bullet \to \ldots)
$$
を Abel 群の複体からなる2つの複体の間の写像とする。
$A^{p, q} = A_p^q$、$B^{p, q} = B_p^q$ とおいて二重複体を得る。
$\text{Tot}_\pi(A^{\bullet, \bullet})$ および
$\text{Tot}_\pi(B^{\bullet, \bullet})$ を、それらの二重複体に付随する
積全複体とする。各 $A_p^\bullet \to B_p^\bullet$ が擬同型なら
$\text{Tot}_\pi(A^{\bullet, \bullet}) \to
\text{Tot}_\pi(B^{\bullet, \bullet})$ は擬同型である。
\end{lemma}

\begin{proof}
% P02-E1180: source line 30151 uses p+1=n instead of p+q=n in the second product index.
$\text{Tot}_\pi(A^{\bullet, \bullet})$ の次数 $n$ の項は
$\prod_{p + q = n} A^{p, q} = \prod_{p + 1 = n} A^q_p$ で与えられる
ことを思い出す。$C_p^\bullet$ を写像
$A_p^\bullet \to B_p^\bullet$ の錐とする。
Derived Categories, Section \ref{derived-section-cones} を参照せよ。
錐の構成の関手性により、複体からなる複体
$$
C_0^\bullet \to C_1^\bullet \to C_2^\bullet \to \ldots
$$
を得る。このとき $\text{Tot}_\pi(C^{\bullet, \bullet})$
の次数 $n$ の項は
$$
\prod_{p + q = n} C^{p, q} = \prod_{p + q = n} C^q_p =
\prod_{p + q = n} (B^q_p \oplus A^{q + 1}_p) =
\prod_{p + q = n} B^q_p \oplus \prod_{p + q = n} A^{q + 1}_p
$$
で与えられる。従って
$\text{Tot}_\pi(C^{\bullet, \bullet})$ は写像
$\text{Tot}_\pi(A^{\bullet, \bullet}) \to
\text{Tot}_\pi(B^{\bullet, \bullet})$
の錐である（微分が一致することの確認は省略する）。
従って、各 $A_p^\bullet$ が非輪状なら
$\text{Tot}_\pi(A^{\bullet, \bullet})$ が非輪状であることを
示せば十分である。これは Lemma
\ref{more-algebra-lemma-vanishing-coh-prod-totalization} から従う。
\end{proof}






\section{弱 \'etale 環写像}
\label{more-algebra-section-weakly-etale}

\noindent
この節の結果の大部分は Olivier の論文 \cite{Olivier-AF} による。
関連する論文 \cite{Ferrand-epi} も参照せよ。

\begin{definition}
\label{more-algebra-definition-weakly-etale}
すべての $A$-加群が $A$ 上平坦であるとき、環 $A$ を
{\it 絶対平坦}という。環写像 $A \to B$ について、
$A \to B$ と $B \otimes_A B \to B$ がともに平坦であるとき、
これを{\it 弱 \'etale} または{\it 絶対平坦}という。
\end{definition}

\noindent
絶対平坦環は（しばしば非可換環の文脈で）
von Neumann 正則環と呼ばれることもある。
局所化は弱 \'etale 環写像である。
\'etale 環写像は弱 \'etale である。
次が単純だが鍵となる性質である。

\begin{lemma}
\label{more-algebra-lemma-key}
$A \to B$ を、$B \otimes_A B \to B$ が平坦である環写像とする。
$N$ を $B$-加群とする。$N$ が $A$-加群として平坦なら、
$N$ は $B$-加群として平坦である。
\end{lemma}

\begin{proof}
% P02-E1181: source line 30209 uses the malformed article phrase “is a flat”.
$N$ が $A$-加群として平坦であると仮定する。このとき関手
$$
\text{Mod}_B \longrightarrow \text{Mod}_{B \otimes_A B},\quad
N' \mapsto N \otimes_A N'
$$
は完全である。$B \otimes_A B \to B$ は平坦だから、関手
$$
\text{Mod}_B \longrightarrow \text{Mod}_B,\quad
N' \mapsto (N \otimes_A N') \otimes_{B \otimes_A B} B = N \otimes_B N'
$$
は完全であると結論できる。従って $N$ は $B$ 上平坦である。
\end{proof}

\begin{definition}
\label{more-algebra-definition-weak-dimension}
$A$ を環、$d \geq 0$ を整数とする。
すべての $A$-加群が Tor 次元 $\leq d$ をもつとき、
$A$ の{\it 弱次元は $\leq d$} であるという。
\end{definition}

\begin{lemma}
\label{more-algebra-lemma-weak-dimension-goes-up}
$A \to B$ を弱 \'etale 環写像とする。
$A$ の弱次元が高々 $d$ なら、$B$ についても同様である。
\end{lemma}

\begin{proof}
$N$ を $B$-加群とする。$d = 0$ なら $N$ は $A$-加群として平坦であり、
従って Lemma \ref{more-algebra-lemma-key} により $B$-加群として平坦である。
$d > 0$ と仮定する。自由 $B$-加群による分解
$F_\bullet \to N$ を選ぶ。仮定により
$K = \Im(F_d \to F_{d - 1})$ は $A$-平坦である。
Lemma \ref{more-algebra-lemma-last-one-flat} を参照せよ。
従って Lemma \ref{more-algebra-lemma-key} により $B$-平坦でもある。
ゆえに
$0 \to K \to F_{d - 1} \to \ldots \to F_0 \to N \to 0$
は長さ $d$ の平坦分解であり、$N$ の Tor 次元は高々 $d$ である。
\end{proof}

\begin{lemma}
\label{more-algebra-lemma-absolutely-flat}
$A$ を環とする。次は同値である。
\begin{enumerate}
\item $A$ の弱次元は $\leq 0$ である。
\item $A$ は絶対平坦である。
\item $A$ は被約であり、すべての素イデアルが極大である。
\item $A$ のすべての局所環は体である。
\end{enumerate}
\end{lemma}

\begin{proof}
(1) と (2) の同値性は直ちに分かる。

\medskip\noindent
(1) $\Rightarrow$ (3) の証明。
$A$ が絶対平坦であると仮定する。
このとき $A$ のすべてのイデアルは純粋である。
Algebra, Definition \ref{algebra-definition-pure-ideal} を参照せよ。
従って Algebra, Lemma
\ref{algebra-lemma-finitely-generated-pure-ideal} により、
すべての有限生成イデアルは冪等元で生成される。
$f \in A$ なら、ある冪等元 $e \in A$ に対して
$(f) = (e)$ であり、$D(f) = D(e)$ は開かつ閉である
（Algebra, Lemma \ref{algebra-lemma-idempotent-spec}）。
Algebra, Lemma
\ref{algebra-lemma-ring-with-only-minimal-primes} により、
これは $A$ のすべての素イデアルが極大であることを意味する。
さらに $f$ が冪零なら $e = 0$、従って $f = 0$ である。
ゆえに $A$ は被約である。

\medskip\noindent
(3) $\Rightarrow$ (4) の証明。
$A$ が被約であり、すべての素イデアルが極大なら、
Algebra, Lemma
\ref{algebra-lemma-minimal-prime-reduced-ring} により
すべての局所環は体である。

\medskip\noindent
(4) $\Rightarrow$ (2) の証明。
$A$ のすべての局所環が体なら、Algebra, Lemma
\ref{algebra-lemma-flat-localization} により
すべての $A$-加群は平坦である。
\end{proof}

\begin{lemma}
\label{more-algebra-lemma-product-fields-absolutely-flat}
体の積は絶対平坦環である。
\end{lemma}

\begin{proof}
$K_i$ を体の族とする。$f = (f_i) \in \prod K_i$ とすると、
$f$ が生成するイデアルは、$f_i$ が $0$ であるか否かに応じて
$e_i = 0, 1$ と定めた冪等元 $e = (e_i)$ が生成するイデアルに等しい。
従って $D(f) = D(e)$ は開かつ閉であり、
Lemma \ref{more-algebra-lemma-absolutely-flat} および Algebra, Lemma
\ref{algebra-lemma-ring-with-only-minimal-primes} から結論を得る。
\end{proof}

\begin{lemma}
\label{more-algebra-lemma-base-change-weakly-etale}
$A \to B$ および $A \to A'$ を環写像とする。
$B' = B \otimes_A A'$ を $B$ の底変換とする。
\begin{enumerate}
\item $B \otimes_A B \to B$ が平坦なら、
$B' \otimes_{A'} B' \to B'$ は平坦である。
\item $A \to B$ が弱 \'etale なら、
$A' \to B'$ は弱 \'etale である。
\end{enumerate}
\end{lemma}

\begin{proof}
$B \otimes_A B \to B$ が平坦であると仮定する。
環写像 $B' \otimes_{A'} B' \to B'$ は
$A \to A'$ による $B \otimes_A B \to B$ の底変換である。
従って Algebra, Lemma \ref{algebra-lemma-flat-base-change}
により平坦である。これで (1) が示された。
(2) は (1) と、今用いた平坦環写像の底変換が平坦であるという事実から従う。
\end{proof}

\begin{lemma}
\label{more-algebra-lemma-absolutely-flat-over-absolutely-flat}
$A \to B$ を、$B \otimes_A B \to B$ が平坦である環写像とする。
\begin{enumerate}
\item $A$ が絶対平坦環なら $B$ も絶対平坦環である。
\item $A$ が被約で $A \to B$ が弱 \'etale なら $B$ は被約である。
\end{enumerate}
\end{lemma}

\begin{proof}
(1) は Lemma \ref{more-algebra-lemma-key} と定義から直ちに従う。
$A$ が被約なら、$A$ から絶対平坦環への単射
$$
A \to A' = \prod_{\mathfrak p \subset A\text{ minimal}} A_\mathfrak p
$$
が存在する（Algebra, Lemma
\ref{algebra-lemma-reduced-ring-sub-product-fields} および
Lemma \ref{more-algebra-lemma-product-fields-absolutely-flat}）。
$A \to B$ が平坦なら、誘導写像
$B \to B' = B \otimes_A A'$ も単射である。
Lemma \ref{more-algebra-lemma-base-change-weakly-etale} により
環写像 $A' \to B'$ は弱 \'etale である。
(1) により $B'$ は絶対平坦である。
Lemma \ref{more-algebra-lemma-absolutely-flat} により環 $B'$ は被約である。
従って $B$ は被約である。
\end{proof}

\begin{lemma}
\label{more-algebra-lemma-composition-weakly-etale}
$A \to B$ および $B \to C$ を環写像とする。
\begin{enumerate}
\item $B \otimes_A B \to B$ と $C \otimes_B C \to C$ が平坦なら、
$C \otimes_A C \to C$ は平坦である。
\item $A \to B$ と $B \to C$ が弱 \'etale なら、
$A \to C$ は弱 \'etale である。
\end{enumerate}
\end{lemma}

\begin{proof}
(1) は乗法写像の分解
$$
C \otimes_A C \longrightarrow C \otimes_B C \longrightarrow C
$$
と、等式
$$
C \otimes_B C = (C \otimes_A C) \otimes_{B \otimes_A B} B,
$$
および平坦写像の底変換が平坦であることと、平坦環写像の合成が
平坦であることから従う。Algebra, Lemmas
\ref{algebra-lemma-flat-base-change} および
\ref{algebra-lemma-composition-flat} を参照せよ。
(2) は (1) と、今用いた平坦環写像の合成が平坦であるという事実から従う。
\end{proof}

\begin{lemma}
\label{more-algebra-lemma-go-down}
$A \to B \to C$ を環写像とする。
\begin{enumerate}
\item $B \to C$ が忠実平坦で $C \otimes_A C \to C$ が平坦なら、
$B \otimes_A B \to B$ は平坦である。
\item $B \to C$ が忠実平坦で $A \to C$ が弱 \'etale なら、
$A \to B$ は弱 \'etale である。
\end{enumerate}
\end{lemma}

\begin{proof}
$B \to C$ が忠実平坦で $C \otimes_A C \to C$ が平坦であると仮定する。
可換図式
$$
\xymatrix{
C \otimes_A C \ar[r] & C \\
B \otimes_A B \ar[r] \ar[u] & B \ar[u]
}
$$
を考える。縦向きの矢印は平坦であり、上の横向きの矢印も平坦である。
従って $C$ は $B \otimes_A B$-加群として平坦である。
写像 $B \to C$ は忠実平坦であり、$C = B \otimes_B C$ である。
従って Algebra, Lemma
\ref{algebra-lemma-flatness-descends-more-general} により、
$B$ は $B \otimes_A B$-加群として平坦である。これで (1) が示された。
(2) は (1) と、$A \to C$ が平坦で $B \to C$ が忠実平坦なら
$A \to B$ が平坦であるという事実から従う
（Algebra, Lemma
\ref{algebra-lemma-flatness-descends-more-general}）。
\end{proof}

\begin{lemma}
\label{more-algebra-lemma-weakly-etale-permanence}
$A$ を環とする。$B \to C$ を、弱 \'etale $A$-代数の間の
$A$-代数写像とする。このとき $B \to C$ は弱 \'etale である。
\end{lemma}

\begin{proof}
Lemma \ref{more-algebra-lemma-key} により環写像 $B \to C$ は平坦である。
環写像 $C \otimes_A C \to C \otimes_B C$ は全射、従ってエピ射である。
% P02-E1182: source lines 30413-30414 place a spurious comma between “implies” and “that”.
従って Lemma \ref{more-algebra-lemma-key} から、$C$ は $C \otimes_A C$ 上平坦なので
$C$ は $C \otimes_B C$ 上でも平坦であることが従う。
\end{proof}

\begin{lemma}
\label{more-algebra-lemma-formally-unramified}
$A \to B$ を、$B \otimes_A B \to B$ が平坦である環写像とする。
このとき $\Omega_{B/A} = 0$、すなわち $B$ は $A$ 上形式的非分岐である。
\end{lemma}

\begin{proof}
$I \subset B \otimes_A B$ を平坦な全射
$B \otimes_A B \to B$ の核とする。
このとき $I$ は純粋イデアルである
（Algebra, Definition \ref{algebra-definition-pure-ideal}）から、
$I^2 = I$ である
（Algebra, Lemma \ref{algebra-lemma-pure}）。
$\Omega_{B/A} = I/I^2$ だから
（Algebra, Lemma \ref{algebra-lemma-differentials-diagonal}）、
消滅を得る。Algebra, Lemma
\ref{algebra-lemma-characterize-formally-unramified} により、
これは $B$ が $A$ 上形式的非分岐であることを意味する。
\end{proof}

\begin{lemma}
\label{more-algebra-lemma-weakly-etale-finite-type}
$A \to B$ を、$B \otimes_A B \to B$ が平坦である環写像とする。
\begin{enumerate}
\item $A \to B$ が有限型なら $A \to B$ は非分岐である。
\item $A \to B$ が有限表示かつ平坦なら $A \to B$ は \'etale である。
\end{enumerate}
特に、有限表示な弱 \'etale 環写像は \'etale である。
\end{lemma}

\begin{proof}
(1) は Lemma \ref{more-algebra-lemma-formally-unramified} と Algebra, Definition
\ref{algebra-definition-unramified} から従う。
(2) は (1) と Algebra, Lemma
\ref{algebra-lemma-etale-flat-unramified-finite-presentation} から従う。
\end{proof}

\begin{lemma}
\label{more-algebra-lemma-when-weakly-etale}
$A \to B$ を環写像とする。次の各場合に $A \to B$ は弱 \'etale である。
\begin{enumerate}
\item $B = S^{-1}A$ は $A$ の局所化である。
\item $A \to B$ は \'etale である。
\item $B$ は弱 \'etale $A$-代数のフィルター余極限である。
\end{enumerate}
\end{lemma}

\begin{proof}
\'etale 環写像は平坦であり、写像 $B \otimes_A B \to B$ も
\'etale $A$-代数の間の写像として \'etale である
（Algebra, Lemma \ref{algebra-lemma-map-between-etale}）。
これで (2) が示された。

\medskip\noindent
$B_i$ を弱 \'etale $A$-代数の有向系とする。
このとき Algebra, Lemma \ref{algebra-lemma-colimit-flat} により
$B = \colim B_i$ は $A$ 上平坦である。
Lemma \ref{more-algebra-lemma-weakly-etale-permanence} により、
推移写像 $B_i \to B_{i'}$ は平坦であることに注意する。
従って各 $i$ に対して $B$ は $B_i$ 上平坦であり、
Algebra, Lemma \ref{algebra-lemma-composition-flat} により
$B$ は $B_i \otimes_A B_i$ 上平坦である。
ゆえに Algebra, Lemma
\ref{algebra-lemma-colimit-rings-flat} により、
$B$ は $B \otimes_A B = \colim B_i \otimes_A B_i$ 上平坦である。

\medskip\noindent
(1) は直接証明することもできるが、(2) と (3) を組み合わせても従う。
\end{proof}

\begin{lemma}
\label{more-algebra-lemma-absolutely-flat-fields}
$L/K$ を体の拡大とする。$L \otimes_K L \to L$ が平坦なら、
$L$ は $K$ の代数的分離拡大である。
\end{lemma}

\begin{proof}
% P02-E1183: source lines 30494-30495 omit “for” before the quantified subfield.
Lemma \ref{more-algebra-lemma-go-down} により、任意の中間体
$K \subset L' \subset L$ に対して写像
$L' \otimes_K L' \to L'$ は平坦である。
従って $L$ が $K$ の有限生成体拡大であると仮定してよい。
この場合、$L/K$ が形式的非分岐であること
（Lemma \ref{more-algebra-lemma-formally-unramified}）から、
$L/K$ は有限分離拡大であることが従う。Algebra, Lemma
\ref{algebra-lemma-characterize-separable-algebraic-field-extensions}
を参照せよ。
\end{proof}

\begin{lemma}
\label{more-algebra-lemma-absolutely-flat-over-field}
$B$ を体 $K$ 上の代数とする。次は同値である。
\begin{enumerate}
\item $B \otimes_K B \to B$ は平坦である。
\item $K \to B$ は弱 \'etale である。
\item $B$ は \'etale $K$-代数のフィルター余極限である。
\end{enumerate}
さらに、$B$ のすべての有限生成 $K$-部分代数は $K$ 上 \'etale である。
\end{lemma}

\begin{proof}
すべての $K$-代数は $K$ 上平坦だから、(1) と (2) は同値である。
% P02-E1184: source line 30518 omits the terminal period after the cited lemma.
(3) なら Lemma \ref{more-algebra-lemma-when-weakly-etale} により (1) と (2) が成り立つ。

\medskip\noindent
(1) と (2) が成り立つと仮定し、(3) と補題の有限生成に関する主張を示す。
% P02-E1185: source lines 30521-30523 omit “an” twice in “is absolutely flat ring” and “is a absolutely flat ring”.
体は絶対平坦環だから、Lemma
\ref{more-algebra-lemma-absolutely-flat-over-absolutely-flat} により $B$ は絶対平坦環である。
従って Lemma \ref{more-algebra-lemma-absolutely-flat} により
$B$ は被約であり、すべての局所環は体である。

\medskip\noindent
$\mathfrak q \subset B$ を素イデアルとする。環写像
$B \to B_\mathfrak q$ は弱 \'etale であり、従って
$B_\mathfrak q$ は $K$ 上弱 \'etale である
（Lemma \ref{more-algebra-lemma-composition-weakly-etale}）。
ゆえに Lemma \ref{more-algebra-lemma-absolutely-flat-fields} により
$B_\mathfrak q$ は $K$ の分離代数拡大である。

\medskip\noindent
% P02-E1186: source line 30535 writes the compound noun “sub algebra” as two words.
$K \subset A \subset B$ を有限生成 $K$-部分代数とする。
$A$ が $K$ 上 \'etale であることを示せば補題の証明が終わる。
このときすべての極小素イデアル
$\mathfrak p \subset A$ は $B$ のある素イデアル $\mathfrak q$ の像である。
Algebra, Lemma
\ref{algebra-lemma-injective-minimal-primes-in-image} を参照せよ。
従って $\kappa(\mathfrak p)$ は
$B_\mathfrak q = \kappa(\mathfrak q)$ の部分体として
$K$ 上分離代数的である。ゆえに $\Spec(A)$ のすべての生成点は閉である
（Algebra, Lemma
\ref{algebra-lemma-finite-residue-extension-closed}）。
従って $\dim(A) = 0$ である。このとき、例えば Algebra, Proposition
\ref{algebra-proposition-dimension-zero-ring} により、
$A$ はその局所環の積である。さらに $A$ は被約だから、
すべての局所環はその剰余体に等しく、それらは $K$ の有限分離拡大である。
Algebra, Lemma \ref{algebra-lemma-etale-over-field} により、
これは $A$ が $K$ 上 \'etale であることを意味し、証明が終わる。
\end{proof}

\begin{lemma}
\label{more-algebra-lemma-weakly-etale-residue-field-extensions}
$A \to B$ を環写像とする。$A \to B$ が弱 \'etale なら、
$A \to B$ は分離代数的な剰余体拡大を誘導する。
\end{lemma}

\begin{proof}
$\mathfrak p$ を $A$ の素イデアルとする。
Lemma \ref{more-algebra-lemma-base-change-weakly-etale} により
$\kappa(\mathfrak p) \to B \otimes_A \kappa(\mathfrak p)$ は弱 \'etale である。
従って Lemma \ref{more-algebra-lemma-absolutely-flat-over-field} により
$B \otimes_A \kappa(\mathfrak p)$ は
\'etale $\kappa(\mathfrak p)$-代数のフィルター余極限である。
ゆえに $\mathfrak p$ の上にある $\mathfrak q \subset B$ に対して、
拡大 $\kappa(\mathfrak q)/\kappa(\mathfrak p)$ は有限分離拡大の
フィルター余極限である。Algebra, Lemma
\ref{algebra-lemma-etale-over-field} を参照せよ。
\end{proof}

\begin{lemma}
\label{more-algebra-lemma-weak-dimension-at-most-1}
$A$ を環とする。次は同値である。
\begin{enumerate}
\item $A$ の弱次元は $\leq 1$ である。
\item $A$ のすべてのイデアルは平坦である。
\item $A$ のすべての有限生成イデアルは平坦である。
\item 平坦 $A$-加群のすべての部分加群は平坦である。
\item $A$ のすべての局所環は付値環である。
\end{enumerate}
\end{lemma}

\begin{proof}
$A$ の弱次元が $\leq 1$ なら、分解
$0 \to I \to A \to A/I \to 0$ と Lemma
\ref{more-algebra-lemma-last-one-flat} により、すべてのイデアル $I$ は平坦である。
従って (1) $\Rightarrow$ (2) である。

\medskip\noindent
(4) を仮定する。$M$ を $A$-加群とする。
$F$ が自由 $A$-加群であるような全射 $F \to M$ を選ぶ。
仮定により $\Ker(F \to M)$ は平坦であり、Lemma
\ref{more-algebra-lemma-tor-dimension} によって $M$ の Tor 次元は
$\leq 1$ である。従って (4) $\Rightarrow$ (1) である。

\medskip\noindent
すべてのイデアルは、その中に含まれる有限生成イデアルの合併である。
従って Algebra, Lemma \ref{algebra-lemma-colimit-flat} により
(3) から (2) が従う。ゆえに (3) $\Leftrightarrow$ (2) である。

\medskip\noindent
(2) を仮定する。$M$ を平坦 $A$-加群とし、$N \subset M$ とする。
$N$ が平坦であることを示す。
Algebra, Theorem \ref{algebra-theorem-lazard} により、
各 $M_i$ が有限自由であるように $M = \colim M_i$ と書ける。
$N_i \subset M_i$ を $N$ の逆像とおくと $N = \colim N_i$ である。
% P02-E1187: source lines 30611-30613 insert a period after the citation and leave the following “it suffices” sentence malformed.
Algebra, Lemma \ref{algebra-lemma-colimit-flat} により、
$N_i$ が平坦であることを示せば十分であり、
% P02-E1188: source line 30614 switches from the ambient ring A to the undefined ring R.
$M = R^{\oplus n}$ の場合に帰着する。この場合、加群 $N$ は
部分加群 $R^{\oplus j} \cap N$ による有限フィルトレーションをもち、
その部分商はイデアルである。(2) によりこれらのイデアルは平坦であり、
従って Algebra, Lemma \ref{algebra-lemma-flat-ses} により
$N$ は平坦である。ゆえに (2) $\Rightarrow$ (4) である。

\medskip\noindent
$A$ が (1) を満たすと仮定し、$\mathfrak p \subset A$ を素イデアルとする。
Lemmas \ref{more-algebra-lemma-when-weakly-etale} および
\ref{more-algebra-lemma-weak-dimension-goes-up} により
$A_\mathfrak p$ は (1) を満たす。
$A$ が (3) を満たす局所環なら $A$ は付値環であることを示す。
$f \in \mathfrak m$ を非零元とする。
このとき $(f)$ は1元で生成される非零平坦加群である。
従って Algebra, Lemma
\ref{algebra-lemma-finite-flat-local} により自由 $A$-加群である。
ゆえに $f$ は非零因子であり、$A$ は整域である。
$I \subset A$ を有限生成イデアルとすると、同様に
$I$ は有限自由 $A$-加群である。
% P02-E1189: source lines 30630-30632 declares every finitely generated ideal free of rank 1, overlooking the zero ideal.
従って（階数を考えると）階数 $1$ の自由加群であり、
$I$ は単項イデアルである。ゆえに Algebra, Lemma
\ref{algebra-lemma-characterize-valuation-ring} により
$A$ は付値環である。従って (1) $\Rightarrow$ (5) である。

\medskip\noindent
(5) を仮定する。$I \subset A$ を有限生成イデアルとする。
$I_\mathfrak p \subset A_\mathfrak p$ は付値環の有限生成イデアルだから、
単項であり
（Algebra, Lemma \ref{algebra-lemma-characterize-valuation-ring}）、
従って平坦である。ゆえに Algebra, Lemma
\ref{algebra-lemma-flat-localization} により $I$ は平坦である。
従って (5) $\Rightarrow$ (3) であり、補題の証明が終わる。
\end{proof}

\begin{lemma}
\label{more-algebra-lemma-product-weak-dimension-at-most-1}
$J$ を集合とする。各 $j \in J$ に対し、
$A_j$ を分数体 $K_j$ をもつ付値環とする。
$A = \prod A_j$ および $K = \prod K_j$ とおく。
このとき $A$ の弱次元は高々 $1$ であり、
$A \to K$ は局所化である。
\end{lemma}

\begin{proof}
$I \subset A$ を有限生成イデアルとする。
Lemma \ref{more-algebra-lemma-weak-dimension-at-most-1} により、
$I$ が平坦 $A$-加群であることを示せば十分である。
$I_j \subset A_j$ を $I$ の像とする。
$I_j = I \otimes_A A_j$ であることに注意する。
従って Algebra, Proposition
\ref{algebra-proposition-fg-tensor} により
$I \to \prod I_j$ は全射である。ゆえに $I = \prod I_j$ である。
$A_j$ は付値環だから、イデアル $I_j$ は1元で生成される
（Algebra, Lemma
\ref{algebra-lemma-characterize-valuation-ring}）。
$I_j = (f_j)$ とする。このとき $I$ は元 $f = (f_j)$ で生成される。
$e \in A$ を、$f_j$ が $0$ であるか否かに応じて
$A_j$ における成分が $0$ または $1$ となる冪等元とする。
するとある非零因子 $g \in A$ に対して $f = g e$ である。
実際、$f_j = 0$ なら $g_j = 1$、そうでなければ $g_j = f_j$
として $g = (g_j)$ と取ればよい。
従って加群として $I \cong (e)$ である。
$(e)$ は $A$ の直和因子だから $I$ は平坦であると結論する。
最後の主張は、
$S = \prod (A_j \setminus \{0\})$ としたとき
$K = S^{-1}A$ だから成り立つ。
\end{proof}

\begin{lemma}
\label{more-algebra-lemma-product-found-valuation-rings}
$A$ を分数体 $K$ をもつ正規整域とする。
環の Cartesian 図式
$$
\xymatrix{
A \ar[d] \ar[r] & K \ar[d] \\
V \ar[r] & L
}
$$
で、$V$ の弱次元が高々 $1$ であり、
$V \to L$ が平坦かつ単射な環のエピ射であるものが存在する。
\end{lemma}

\begin{proof}
各 $x \in K$, $x \not \in A$ に対して、
Algebra, Lemma \ref{algebra-lemma-find-valuation-rings}
に現れる $V_x \subset K$ を選ぶ。
$V = \prod_{x \in K \setminus A} V_x$ および
$L = \prod_{x \in K \setminus A} K$ とおく。
Lemma \ref{more-algebra-lemma-product-weak-dimension-at-most-1} により
環 $V$ の弱次元は高々 $1$ であり、同じ補題から
$V \to L$ は局所化であることも分かる。
局所化は平坦かつエピ射である。Algebra, Lemmas
\ref{algebra-lemma-flat-localization} および
\ref{algebra-lemma-epimorphism-local} を参照せよ。
\end{proof}

\begin{lemma}
\label{more-algebra-lemma-weak-dimension-at-most-1-integrally-closed}
$A$ を弱次元が高々 $1$ の環とする。
$A \to B$ が平坦かつ単射な環のエピ射なら、
$A$ は $B$ の中で整閉である。
\end{lemma}

\begin{proof}
$x \in B$ を $A$ 上整な元とし、$A' = A[x] \subset B$ とおく。
Algebra, Lemma
\ref{algebra-lemma-characterize-finite-in-terms-of-integral} により
$A'$ は $A$ の有限環拡大である。
$A = A'$ を示すには、Algebra, Lemma
\ref{algebra-lemma-finite-epimorphism-surjective} により
$A \to A'$ がエピ射であることを示せば十分である。
$A$ に関する仮定と $B$ が $A$ 上平坦であることから、
$A'$ は $A$ 上平坦であることに注意する
（Lemma \ref{more-algebra-lemma-weak-dimension-at-most-1}）。
従って合成
$$
A' \otimes_A A' \to B \otimes_A A' \to B \otimes_A B \to B
$$
は単射である。すなわち $A' \otimes_A A' \cong A'$ であり、
補題が証明された。
\end{proof}

\begin{lemma}
\label{more-algebra-lemma-normality-goes-up}
$A$ を分数体 $K$ をもつ正規整域とする。
$A \to B$ を弱 \'etale とする。このとき
$B$ は $B \otimes_A K$ の中で整閉である。
\end{lemma}

\begin{proof}
Lemma \ref{more-algebra-lemma-product-found-valuation-rings} に現れる図式を選ぶ。
$A \to B$ は平坦だから、底変換により環の Cartesian 図式
$$
\xymatrix{
B \ar[d] \ar[r] & B \otimes_A K \ar[d] \\
B \otimes_A V \ar[r] & B \otimes_A L
}
$$
を得る。$V \to B \otimes_A V$ は弱 \'etale であることに注意する
（Lemma \ref{more-algebra-lemma-base-change-weakly-etale}）。
従って Lemma \ref{more-algebra-lemma-weak-dimension-goes-up} により
$B \otimes_A V$ の弱次元は高々 $1$ である。
また $B \otimes_A V \to B \otimes_A L$ は、
そのような写像の平坦な底変換として、平坦かつ単射な環のエピ射である
（Algebra, Lemmas \ref{algebra-lemma-flat-base-change} および
\ref{algebra-lemma-base-change-epimorphism}）。
Lemma \ref{more-algebra-lemma-weak-dimension-at-most-1-integrally-closed} により
$B \otimes_A V$ は $B \otimes_A L$ の中で整閉である。
図式の Cartesian 性から、
$B$ は $B \otimes_A K$ の中で整閉であることが従う。
\end{proof}

\begin{lemma}
\label{more-algebra-lemma-integral-over-henselian}
$A \to B$ を環準同型とする。次を仮定する。
\begin{enumerate}
\item $A$ は Hensel 局所環である。
\item $A \to B$ は整である。
\item $B$ は整域である。
\end{enumerate}
このとき $B$ は Hensel 局所環であり、$A \to B$ は局所準同型である。
$A$ が強 Hensel なら、$B$ は強 Hensel 局所環であり、
剰余体の拡大
$\kappa(\mathfrak m_B)/\kappa(\mathfrak m_A)$ は純非分離である。
\end{lemma}

\begin{proof}
% P02-E1190: source line 30772 writes the compound noun “subalgebras” as “sub algebras”.
$B$ を有限 $A$-部分代数のフィルター余極限
$B = \colim B_i$ と書く。各 $B_i$ について結果を証明すれば、
$B$ に対する結果が従う。Algebra, Lemma
\ref{algebra-lemma-colimit-henselian} を参照せよ。
$A \to B$ が有限なら、Algebra, Lemma
\ref{algebra-lemma-finite-over-henselian} により
$B$ は Hensel 局所環の積である。
$B$ は整域だから $B$ は局所環である。
$A \to B$ に対する going up により、$B$ の極大イデアルは
$A$ の極大イデアルの上にある
（Algebra, Lemma \ref{algebra-lemma-integral-going-up}）。
$A$ が強 Hensel なら、
% P02-E1191: source lines 30780-30782 use a dangling “being algebraic” construction separated from its subject by a comma.
体拡大 $\kappa(\mathfrak m_B)/\kappa(\mathfrak m_A)$ は代数的なので、
純非分離でなければならない。もちろん、このとき
$\kappa(\mathfrak m_B)$ は分離代数閉であり、
$B$ は強 Hensel である。
\end{proof}

\begin{theorem}[Olivier]
\label{more-algebra-theorem-olivier}
$A \to B$ を局所環の局所準同型とする。
$A$ が強 Hensel で $A \to B$ が弱 \'etale なら $A = B$ である。
\end{theorem}

\begin{proof}
すべての $\mathfrak p \subset A$ に対して、
$\mathfrak p$ の上にある素イデアル $\mathfrak q \subset B$ は一意であり、
$\kappa(\mathfrak p) = \kappa(\mathfrak q)$ であることを示す。
これは $B \otimes_A B \to B$ が Spec 上全単射であると同時に
全射かつ平坦であることを意味する。
従って、例えば Algebra, Lemma
\ref{algebra-lemma-pure-open-closed-specializations}
における純粋イデアルの記述により、これは同型である。
従って $A \to B$ は忠実平坦な環のエピ射である。
Algebra, Lemma
\ref{algebra-lemma-faithfully-flat-epimorphism} により $A = B$ を得る。

\medskip\noindent
Lemmas \ref{more-algebra-lemma-base-change-weakly-etale} および
\ref{more-algebra-lemma-absolutely-flat-over-field} により、
ファイバー環 $B \otimes_A \kappa(\mathfrak p)$ は
$\kappa(\mathfrak p)$ の \'etale 拡大の余極限であることに注意する。
従って $\mathfrak p$ の上に2つ以上の素イデアルがあるか、
ある $\mathfrak q$ に対して
$\kappa(\mathfrak p) \not = \kappa(\mathfrak q)$ なら、
ある（分離）代数体拡大 $L/\kappa(\mathfrak p)$ に対して
$B \otimes_A L$ は非自明な冪等元をもつ。

\medskip\noindent
$L/\kappa(\mathfrak p)$ を代数体拡大とする。
$A' \subset L$ を $L$ における $A/\mathfrak p$ の整閉包とする。
Lemma \ref{more-algebra-lemma-integral-over-henselian} により、
$A'$ は強 Hensel 局所環であり、その剰余体は $A$ の剰余体の
純非分離拡大である。従って Algebra, Lemma
\ref{algebra-lemma-local-tensor-with-integral} により
$B \otimes_A A'$ は局所環である。
一方、Lemma \ref{more-algebra-lemma-normality-goes-up} により
$B \otimes_A A'$ は $B \otimes_A L$ の中で整閉である。
$B \otimes_A A'$ は局所環だから、
$B \otimes_A L$ は非自明な冪等元をもたない。
これが示すべきことであった。
\end{proof}







\section{体上の弱 \'etale 代数}
\label{more-algebra-section-weakly-etale-over-field}

\noindent
$K$ が体なら、代数 $B$ が $K$ 上弱 \'etale であることと、
\'etale $K$-代数のフィルター余極限であることは同値である。
これは Lemma \ref{more-algebra-lemma-absolutely-flat-over-field} である。

\begin{lemma}
\label{more-algebra-lemma-class-weakly-etale-over-field}
$K$ を体とする。$B$ が $K$ 上弱 \'etale なら、次が成り立つ。
\begin{enumerate}
\item $B$ は被約である。
\item $B$ は $K$ 上整である。
\item $B$ の任意の有限生成 $K$-部分代数は、
$K$ の有限分離拡大の有限積である。
\item $B$ が体であることと $B$ が非自明な冪等元をもたないことは同値であり、
この場合それは $K$ の分離代数拡大である。
\item $B$ の任意の部分 $K$-代数または商代数は
$K$ 上弱 \'etale である。
\item $B'$ が $K$ 上弱 \'etale なら、
$B \otimes_K B'$ は $K$ 上弱 \'etale である。
\end{enumerate}
\end{lemma}

\begin{proof}
(1) は Lemma \ref{more-algebra-lemma-absolutely-flat-over-absolutely-flat} から従うが、
もちろん (3) からも従う。
(3) は Lemma \ref{more-algebra-lemma-absolutely-flat-over-field} と、
\'etale $K$-代数が $K$ の有限分離拡大の有限積であることから従う。
Algebra, Lemma \ref{algebra-lemma-etale-over-field} を参照せよ。
(3) は (2) を含意する。
体の積が体であることと非自明な冪等元をもたないことは同値なので、
(4) は (3) から従う。

\medskip\noindent
$S \subset B$ が部分代数なら、これはその有限生成部分代数の
フィルター余極限である。上で見たようにそれらはすべて
$K$ 上 \'etale だから、Lemma \ref{more-algebra-lemma-absolutely-flat-over-field}
により $S$ は $K$ 上弱 \'etale である。
$B \to Q$ が商代数なら、$Q$ は有限分離拡大 $L_i/K$ の有限積
$\prod_{i \in I} L_i$ の $K$-代数としての商のフィルター余極限である。
このような商はある部分集合 $J \subset I$ に対して
$\prod_{i \in J} L_i$ の形をしているので、同じ議論により
商についても結果が成り立つ。

\medskip\noindent
テンソル積に関する主張も同様に従う。または Lemmas
\ref{more-algebra-lemma-base-change-weakly-etale} および
\ref{more-algebra-lemma-composition-weakly-etale} を組み合わせてもよい。
\end{proof}

\begin{lemma}
\label{more-algebra-lemma-max-weakly-etale-subalgebra}
$K$ を体とし、$A$ を $K$-代数とする。
最大の弱 \'etale $K$-部分代数 $B_{max} \subset A$ が存在する。
\end{lemma}

\begin{proof}
$B_1, B_2 \subset A$ を弱 \'etale $K$-部分代数とする。
このとき $B_1 \otimes_K B_2$ は $K$ 上弱 \'etale であり、
$B_1 \otimes_K B_2 \to A$ の像もそうである
（Lemma \ref{more-algebra-lemma-class-weakly-etale-over-field}）。
従って、弱 \'etale $K$-部分代数 $B \subset A$ 全体の集合
$\mathcal{B}$ は有向であり、その余極限
$B_{max} = \colim_{B \in \mathcal{B}} B$ は
Lemma \ref{more-algebra-lemma-when-weakly-etale} により弱 \'etale $K$-代数である。
ゆえに $B_{max} \to A$ の像は $K$ 上弱 \'etale である
（前の補題を参照）。
従ってこの像は $\mathcal{B}$ に属し、$\mathcal{B}$ は最大元をもつ
（そしてこの像は $B_{max}$ と同じである）。
\end{proof}

\begin{lemma}
\label{more-algebra-lemma-properties-of-max-weakly-etale-subalgebra}
$K$ を体とする。$K$-代数 $A$ に対し、Lemma
\ref{more-algebra-lemma-max-weakly-etale-subalgebra} における $A$ の最大弱
\'etale $K$-部分代数を $B_{max}(A)$ と書く。このとき
% P02-E1192: source lines 30914-30916 omit “by” in the notation clause.
\begin{enumerate}
\item 任意の $K$-代数準同型 $A' \to A$ は $K$-代数準同型
$B_{max}(A') \to B_{max}(A)$ を誘導する。
\item $A' \subset A$ なら
$B_{max}(A') = B_{max}(A) \cap A'$ である。
\item $A = \colim A_i$ がフィルター余極限なら
$B_{max}(A) = \colim B_{max}(A_i)$ である。
\item 写像 $B_{max}(A) \to B_{max}(A_{red})$ は同型である。
\item $B_{max}(A_1 \times \ldots \times A_n) =
B_{max}(A_1) \times \ldots \times B_{max}(A_n)$ である。
\item $A$ が非自明な冪等元をもたないなら、$B_{max}(A)$ は体であり、
$K$ の分離代数拡大である。
% P02-E1193: source line 30928 contains the unfinished placeholder “add more here.”
\item ここにさらに追加する。
\end{enumerate}
\end{lemma}

\begin{proof}
(1) の証明。Lemma \ref{more-algebra-lemma-class-weakly-etale-over-field} により
$B_{max}(A') \to A$ の像は $K$ 上弱 \'etale だからである。

\medskip\noindent
(2) の証明。(1) により
$B_{max}(A') \subset B_{max}(A)$ である。
逆に、Lemma \ref{more-algebra-lemma-class-weakly-etale-over-field} により
$B_{max}(A) \cap A'$ は弱 \'etale $K$-代数であり、
従って $B_{max}(A')$ に含まれる。

\medskip\noindent
(3) の証明。(1) により写像
$\colim B_{max}(A_i) \to A$ がある。
系はフィルター系で $B_{max}(A_i) \subset A_i$ だから、この写像は単射である。
余極限 $\colim B_{max}(A_i)$ は Lemma
\ref{more-algebra-lemma-when-weakly-etale} により $K$ 上弱 \'etale である。
従って単射
$\colim B_{max}(A_i) \to B_{max}(A)$ を得る。
$a \in B_{max}(A)$ とする。このとき $a$ は有限表示
$K$-部分代数 $B \subset B_{max}(A)$ を生成する。
Algebra, Lemma \ref{algebra-lemma-characterize-finite-presentation}
により、ある $i$ と、与えられた写像 $B \to A$ を持ち上げる
$K$-代数準同型 $f : B \to A_i$ が存在する。
Lemma \ref{more-algebra-lemma-class-weakly-etale-over-field} により
$B$ は弱 \'etale だから $f(B) \subset B_{max}(A_i)$ であり、
$a$ は $\colim B_{max}(A_i) \to B_{max}(A)$ の像に属すると結論できる。

\medskip\noindent
(4) の証明。Lemma \ref{more-algebra-lemma-absolutely-flat-over-field} により
$B_{max}(A_{red}) = \colim B_i$ を \'etale $K$-代数の
フィルター余極限として書く。
Algebra, Lemma \ref{algebra-lemma-smooth-strong-lift} により、
各 $i$ に対し、与えられた写像 $B_i \to A_{red}$ を持ち上げる
$K$-代数準同型 $f_i : B_i \to A$ が存在する。
% P02-E1194: source line 30966 reverses the only canonical map; the canonical map is from B_max(A) to B_max(A_red).
従って標準写像 $B_{max}(A_{red}) \to B_{max}(A)$ は全射である。
核は冪零元からなるが、Lemma
\ref{more-algebra-lemma-class-weakly-etale-over-field} により
$B_{max}(A_{red})$ は被約なので、核は零である。

\medskip\noindent
% P02-E1195: source line 30972 leaves the proof of assertion (5) as “Omitted.”
(5) の証明。省略する。

\medskip\noindent
(6) の証明。Lemma \ref{more-algebra-lemma-class-weakly-etale-over-field}
の (4) から従う。
\end{proof}

\begin{lemma}
\label{more-algebra-lemma-change-fields-max-weakly-etale-subalgebra}
$L/K$ を体拡大とし、$A$ を $K$-代数とする。
Lemma \ref{more-algebra-lemma-max-weakly-etale-subalgebra} における
$A$ の最大弱 \'etale $K$-部分代数を $B \subset A$ とする。
このとき $B \otimes_K L$ は $A \otimes_K L$ の
最大弱 \'etale $L$-部分代数である。
\end{lemma}

\begin{proof}
$K$ 上の代数 $A$ に対して、$A$ の最大弱 \'etale $K$-部分代数を
$B_{max}(A/K)$ と書く。同様に、$A'$ が $L$-代数なら、
$A'$ の最大弱 \'etale $L$-部分代数を $B_{max}(A'/L)$ と書く。
Lemma \ref{more-algebra-lemma-base-change-weakly-etale} により
$B_{max}(A/K) \otimes_K L$ は $L$ 上弱 \'etale であり、
$B_{max}(A/K) \otimes_K L \subset A \otimes_K L$ だから、
標準的な単射
$$
B_{max}(A/K) \otimes_K L \to B_{max}((A \otimes_K L)/L)
$$
を得る。補題はこの写像が同型であることを主張している。

\medskip\noindent
$L$ といまの $K$-代数 $A$ に対して補題を証明するには、
$L$ の任意の体拡大 $L'$ に対して補題を証明すれば十分である。
実際、分解
$$
B_{max}(A/K) \otimes_K L' \to
B_{max}((A \otimes_K L)/L) \otimes_L L' \to
B_{max}((A \otimes_K L')/L')
$$
があるので、合成が全射であるためには
$B_{max}(A/K) \otimes_K L \to B_{max}((A \otimes_K L)/L)$
が全射でなければならない。従って $L$ は代数閉であると仮定してよい。

\medskip\noindent
有限型 $K$-代数の場合への帰着。$A$ はその有限型 $K$-部分代数の
フィルター余極限として書ける。
% P02-E1196: source lines 31016-31017 say “write A is” instead of “write A as”.
Lemma \ref{more-algebra-lemma-properties-of-max-weakly-etale-subalgebra} を用いると、
有限型 $K$-代数について補題を証明すれば十分であることが分かる。

\medskip\noindent
$A$ を有限型 $K$-代数と仮定する。
$A \to A_{red}$ の核は冪零なので、
$A \otimes_K L \to A_{red} \otimes_K L$ の核もそうである。
すると
$$
B_{max}((A \otimes_K L)/L) \to B_{max}((A_{red} \otimes_K L)/L)
$$
は単射である。実際、その核は冪零であり、弱 \'etale $L$-代数
$B_{max}((A \otimes_K L)/L)$ は被約である
（Lemma \ref{more-algebra-lemma-class-weakly-etale-over-field}）。
Lemma \ref{more-algebra-lemma-properties-of-max-weakly-etale-subalgebra} により
$B_{max}(A/K) = B_{max}(A_{red}/K)$ だから、
$A_{red}$ に対して補題を証明すれば十分である。

\medskip\noindent
$A$ を被約な有限型 $K$-代数と仮定する。
$Q = Q(A)$ を $A$ の全商環とする。このとき
% P02-E1197: source line 31040 tensors Q with L over A, although L has only the stated K-algebra structure.
$A \subset Q$ および $A \otimes_K L \subset Q \otimes_A L$ であり、
従って Lemma \ref{more-algebra-lemma-properties-of-max-weakly-etale-subalgebra}
により
$$
B_{max}(A/K) = A \cap B_{max}(Q/K)
$$
および
$$
B_{max}((A \otimes_K L)/L) =
(A \otimes_K L) \cap B_{max}((Q \otimes_K L)/L)
$$
である。$-\otimes_K L$ は完全関手だから、$Q$ に対して結果を証明すれば
$A$ に対する結果が従う。
$Q$ は体の有限積であり（Algebra, Lemmas
\ref{algebra-lemma-total-ring-fractions-no-embedded-points},
\ref{algebra-lemma-minimal-prime-reduced-ring},
\ref{algebra-lemma-Noetherian-irreducible-components}, および
\ref{algebra-lemma-Noetherian-permanence}）、また
$B_{max}$ は積と可換なので
（Lemma \ref{more-algebra-lemma-properties-of-max-weakly-etale-subalgebra}）、
$A$ が体の場合に補題を証明すれば十分である。

\medskip\noindent
$A$ を体と仮定する。この証明の第3段落の議論により、
$A$ が $K$ 上有限生成である場合に帰着する。
（実際、体の場合への帰着の仕方により、$K$ の有限生成体拡大が得られる。）

\medskip\noindent
$A$ を $K$ の有限生成体拡大と仮定する。
Lemma \ref{more-algebra-lemma-properties-of-max-weakly-etale-subalgebra} の (6) により
$K' = B_{max}(A/K)$ は $K$ 上分離代数的な体である。
従って $K'$ は $K$ の有限分離体拡大であり、Algebra, Lemma
\ref{algebra-lemma-make-geometrically-irreducible} により
$A$ は $K'$ 上幾何学的既約である。
$L$ は代数閉で $K'/K$ は有限分離だから、
$$
K' \otimes_K L \to \prod\nolimits_{\sigma \in \Hom_K(K', L)} L,\quad
\alpha \otimes \beta \mapsto (\sigma(\alpha)\beta)_\sigma
$$
は同型である
（Fields, Lemma \ref{fields-lemma-finite-separable-tensor-alg-closed}）。
従って
$$
A \otimes_K L = A \otimes_{K'} (K' \otimes_K L) =
\prod\nolimits_{\sigma \in \Hom_K(K', L)} A \otimes_{K', \sigma} L
$$
である。$A$ は $K'$ 上幾何学的既約だから、
$A \otimes_{K', \sigma} L$ は唯一の極小素イデアルをもつ。
$L$ は代数閉なので、
Lemma \ref{more-algebra-lemma-properties-of-max-weakly-etale-subalgebra} の (6) により
この $L$-代数は $L$ 上代数的な体であり、従って
$B_{max}((A \otimes_{K', \sigma} L)/L) = L$ である。
上のようにこれらの部分代数を積に分解できるので、
$A \otimes_K L$ の最大弱 \'etale
$K' \otimes_K L$-部分代数は $K' \otimes_K L$ である。
従って包含
$K' \otimes_K L \subset B_{max}((A \otimes_K L)/L)$ は等号である。
実際、Lemma \ref{more-algebra-lemma-weakly-etale-permanence} により環準同型
$K' \otimes_K L \to B_{max}((A \otimes_K L)/L)$ は弱 \'etale である。
\end{proof}













\section{局所既約性}
\label{more-algebra-section-unibranch}

\noindent
次の定義が一般に受け入れられているものと思われる。
これを理解するため、$A \subset B$ が局所整域の整拡大なら、
going up により $A \to B$ は局所環準同型であることに注意する
（Algebra, Lemma \ref{algebra-lemma-integral-going-up}）。

\begin{definition}
\label{more-algebra-definition-unibranch}
\begin{reference}
\cite[Chapter 0 (23.2.1)]{EGA4}
\end{reference}
$A$ を局所環とする。被約化 $A_{red}$ が整域であり、
$A_{red}$ の分数体における整閉包 $A'$ が局所環であるとき、
$A$ は {it 単枝} であるという。
$A$ が単枝であり、さらに $A'$ の剰余体が $A$ の剰余体上
純非分離であるとき、$A$ は {it 幾何学的単枝} であるという。
\end{definition}

\noindent
$A$ を局所環とする。同値な定式化は次の通りである。
\begin{enumerate}
\item $A$ が唯一の極小素イデアル $\mathfrak p$ をもち、
$A/\mathfrak p$ の分数体における整閉包が局所環であるとき、
$A$ は単枝である。
\item $A$ が唯一の極小素イデアル $\mathfrak p$ をもち、
$A/\mathfrak p$ の分数体における整閉包が局所環であって、
その剰余体が $A$ の剰余体上純非分離であるとき、
$A$ は幾何学的単枝である。
\end{enumerate}
正規局所環は幾何学的単枝である
（Definition \ref{more-algebra-definition-unibranch} および Algebra, Definition
\ref{algebra-definition-ring-normal} から従う）。
Lemmas \ref{more-algebra-lemma-unibranch} および
\ref{more-algebra-lemma-geometrically-unibranch} は、（幾何学的）単枝性が
考察するに値する性質であることを示唆する。

\begin{lemma}
\label{more-algebra-lemma-branches}
$A$ を局所環とし、$A$ は有限個の極小素イデアルをもつと仮定する。
$A'$ を、$A_{red}$ の全商環における $A$ の整閉包とする。
$A^h$ を $A$ の Hensel 化とする。写像
$$
\Spec(A') \leftarrow \Spec((A')^h) \rightarrow \Spec(A^h)
$$
を考える。ここで $(A')^h = A' \otimes_A A^h$ である。このとき
\begin{enumerate}
\item 左の矢印は極大イデアル上全単射である。
\item 右の矢印は極小素イデアル上全単射である。
\item $(A')^h$ の各極小素イデアルは唯一の極大イデアルに含まれ、
各極大イデアルはちょうど1つの極小素イデアルを含む。
\end{enumerate}
\end{lemma}

\begin{proof}
$I \subset A$ を冪零元全体のイデアルとする。
Algebra, Lemma \ref{algebra-lemma-quotient-henselization} により
$(A/I)^h = A^h/IA^h$ である。
$A$, $A^h$, $A'$, $(A')^h$ のスペクトルは、それぞれ
$A/I$, $A^h/IA^h$, $A'$,
$(A')^h = A' \otimes_{A/I} A^h/IA^h$ のスペクトルと同じである。
従って $A$ を $A_{red} = A/I$ で置き換え、
$A$ は被約であると仮定してよい。
このとき $A \subset A'$ であり、以下では断りなくこれを用いる。

\medskip\noindent
(1) の証明。$A'$ は $A$ 上整なので、$(A')^h$ は $A^h$ 上整である。
going up（Algebra, Lemma
\ref{algebra-lemma-integral-going-up}）により、
$A'$ および $(A')^h$ の各極大イデアルは、それぞれ
$A$ および $A^h$ の極大イデアル
$\mathfrak m$ および $\mathfrak m^h$ の上にある。
従って (1) は同型
$$
(A')^h \otimes_{A^h} \kappa^h =
A' \otimes_A A^h \otimes_{A^h} \kappa^h = A' \otimes_A \kappa
$$
から従う。実際、$A \to A^h$ が誘導する剰余体拡大
$\kappa^h/\kappa$ は自明である。
% P02-E1198: source lines 31196-31197 omit the article in “hence spectrum of this ring”.
以下では、表示された環は体上整であり、従ってそのスペクトルが
副有限空間であることも用いる。Algebra, Lemmas
\ref{algebra-lemma-integral-over-field} および
\ref{algebra-lemma-ring-with-only-minimal-primes} を参照せよ。

\medskip\noindent
(3) の証明。$A'$ は正規環であり、実際、正規整域の有限積である。
Algebra, Lemma
\ref{algebra-lemma-characterize-reduced-ring-normal} を参照せよ。
$A^h$ は \'etale $A$-代数のフィルター余極限だから、
% P02-E1199: source line 31206 omits “a” before “filtered colimit”.
$(A')^h$ は \'etale $A'$-代数のフィルター余極限であり、
従って Algebra, Lemmas \ref{algebra-lemma-normal-goes-up} および
\ref{algebra-lemma-colimit-normal-ring} により $(A')^h$ は正規環である。
従って $(A')^h$ の各局所環は正規整域であり、
各極大イデアルは唯一の極小素イデアルを含む。
Lemma \ref{more-algebra-lemma-integral-over-henselian-pair} を
$A^h \to (A')^h$ に適用すると、
$((A')^h, \mathfrak m(A')^h)$ は Hensel 対であることが分かる。
$\mathfrak q \subset (A')^h$ が極小素イデアル（または任意の素イデアル）
なら、Lemma \ref{more-algebra-lemma-irreducible-henselian-pair-connected} により
$V(\mathfrak q)$ と $V(\mathfrak m (A')^h)$ の共通部分は連結である。
% P02-E1200: source line 31217 lacks terminal punctuation.
% P02-E1201: source lines 31218-31220 contain a dangling “by” and omit the intended citation or connective.
$V(\mathfrak m (A')^h) = \Spec((A')^h \otimes \kappa^h)$ は
副有限空間なので、$\mathfrak q$ を含む極大イデアルは唯一である。

\medskip\noindent
(2) の証明。$A'$ の極小素イデアルは、ちょうど $A$ の
極小素イデアルの上にある素イデアルである（構成による）。
$A' \to (A')^h$ は平坦なので、going down
（Algebra, Lemma \ref{algebra-lemma-flat-going-down}）により
$(A')^h$ の各極小素イデアルは $A'$ の極小素イデアルの上にある。
逆に、$A'$ の極小素イデアルの上にある $(A')^h$ の任意の
素イデアルは極小である。実際、$(A')^h$ は $A'$ 上 \'etale、
従って準有限な代数のフィルター余極限である
% P02-E1202: source line 31230 explicitly leaves this implication as “small detail omitted”.
（細部は省略する）。
従って $(A')^h$ の極小素イデアルは、ちょうど $A$ の
極小素イデアルの上にある素イデアルである。
同様に、$A^h$ の極小素イデアルは、ちょうど $A$ の
極小素イデアルの上にある素イデアルである。
構成により $A' \otimes_A Q(A) = Q(A)$ である。ここで
$Q(A)$ は被約局所環 $A$ の全分数環である。
もちろん $Q(A)$ は $A$ の極小素イデアルの剰余体の有限積である。
従って
$$
(A')^h \otimes_A Q(A) = A^h \otimes_A A' \otimes_A Q(A) = A^h \otimes_A Q(A)
$$
である。上の議論により、左辺の環のスペクトルは
$(A')^h$ の極小素イデアル全体であり、右辺の環のスペクトルは
% P02-E1203: source lines 31244-31245 duplicate “is the”.
$A^h$ の極小素イデアル全体である。これで証明が完了する。
\end{proof}

\begin{lemma}
\label{more-algebra-lemma-unibranch}
\begin{reference}
\cite[Chapter IV Proposition 18.6.12]{EGA4}
\end{reference}
$A$ を局所環とし、$A^h$ を $A$ の Hensel 化とする。
% P02-E1204: source line 31254 lacks punctuation after “The following are equivalent”.
次は同値である。
\begin{enumerate}
\item $A$ は単枝である。
\item $A^h$ は唯一の極小素イデアルをもつ。
\end{enumerate}
\end{lemma}

\begin{proof}
これは Lemma \ref{more-algebra-lemma-branches} から従うが、直接証明も与える。
% P02-E1205: source line 31263 omits “by” in “Denote m the maximal ideal”.
環 $A$ の極大イデアルを $\mathfrak m$ と書く。
剰余体 $\kappa = A/\mathfrak m$ は $A^h$ の剰余体と同じであることを
思い出そう。

\medskip\noindent
(2) を仮定する。$\mathfrak p^h$ を $A^h$ の唯一の極小素イデアルとする。
$A \to A^h$ の平坦性により
$\mathfrak p = A \cap \mathfrak p^h$ は $A$ の唯一の
極小素イデアルである（going down による。Algebra, Lemma
\ref{algebra-lemma-flat-going-down} を参照せよ）。
また、Algebra, Lemma
\ref{algebra-lemma-quotient-henselization} を参照すると
$A^h/\mathfrak pA^h = (A/\mathfrak p)^h$ であり、これは
Lemma \ref{more-algebra-lemma-henselization-reduced} により被約なので、
$\mathfrak p^h = \mathfrak pA^h$ である。
$A'$ を $A/\mathfrak p$ の分数体における整閉包とする。
$A'$ が局所環であることを示さなければならない。
$A \to A'$ は整だから、$A'$ のすべての極大イデアルは
$\mathfrak m$ の上にある（整環準同型に対する going up による。
Algebra, Lemma \ref{algebra-lemma-integral-going-up} を参照せよ）。
$A'$ が局所環でなければ、相異なる極大イデアル
$\mathfrak m_1$, $\mathfrak m_2$ をとれる。
$f_i \in \mathfrak m_i$ かつ
$f_i \not \in \mathfrak m_{3 - i}$ となる元
$f_1, f_2 \in A'$ を選ぶ。
有限部分代数
$B = A/\mathfrak p[f_1, f_2] \subset A'$ は、
相異なる極大イデアル $B \cap \mathfrak m_i$, $i = 1, 2$ をもつ。
包含
$$
A/\mathfrak p \subset B \subset \kappa(\mathfrak p)
$$
を平坦な環準同型 $A \to A^h$ でテンソルすると、包含
$$
A^h/\mathfrak p^h \subset
B \otimes_A A^h \subset
\kappa(\mathfrak p) \otimes_A A^h \subset
\kappa(\mathfrak p^h)
$$
を得る。最後の包含が成り立つのは
$$
\kappa(\mathfrak p) \otimes_A A^h =
\kappa(\mathfrak p) \otimes_{A/\mathfrak p} A^h/\mathfrak p^h
$$
が整域 $A^h/\mathfrak p^h$ の局所化だからである。
$B/\mathfrak mB$ は2つの極大イデアルをもつので、
$B \otimes_A \kappa$ は少なくとも2つの極大イデアルをもつことに注意する。
従って $A^h$ は Hensel だから、Algebra, Lemma
\ref{algebra-lemma-mop-up} により
$B \otimes_A A^h$ は $\geq 2$ 個の局所環の積である。
しかし $B \otimes_A A^h$ は整域の部分環であることを既に見たので、
矛盾を得る。

\medskip\noindent
(1) を仮定する。$\mathfrak p \subset A$ を唯一の極小素イデアルとし、
$A'$ を $A/\mathfrak p$ の分数体における整閉包とする。
$A \to B$ を、剰余体の同型を誘導し、かつ \'etale $A$-代数の
局所化である局所環の局所準同型とする。特に
$\mathfrak m_B$ は $\mathfrak m B$ を含む唯一の素イデアルである。
このとき $B' = A' \otimes_A B$ は $B$ 上整であり、
$A \to A'$ が局所的であるという仮定から $B'$ は局所環である
（Algebra, Lemma \ref{algebra-lemma-local-tensor-with-integral}）。
一方、$A' \to B'$ は \'etale 環準同型の局所化なので、
Algebra, Lemma \ref{algebra-lemma-normal-goes-up} により
$B'$ は正規である。従って $B'$ は（局所）正規整域である。
最後に
$$
B/\mathfrak pB \subset B \otimes_A \kappa(\mathfrak p)
= B' \otimes_{A'} (\text{fraction field of }A') \subset
\text{fraction field of }B'
$$
である。従って $B/\mathfrak pB$ は整域であり、これは $B$ が
唯一の極小素イデアルをもつことを含意する
（$A \to B$ の平坦性により、極小素イデアルはすべて
$\mathfrak p$ の上になければならない）。
$A^h$ は局所環 $B$ のフィルター余極限だから、
$A^h$ は唯一の極小素イデアルをもつ。
実際、ある非冪零元 $f, g$ に対して $A^h$ 内で $fg = 0$ なら、
$f$ と $g$ の両方を含む上のような $B$ をとることができ、矛盾に至る。
\end{proof}

\begin{lemma}
\label{more-algebra-lemma-geometric-branches}
$(A, \mathfrak m, \kappa)$ を局所環とし、
$A$ は有限個の極小素イデアルをもつと仮定する。
$A'$ を $A_{red}$ の全商環における $A$ の整閉包とする。
$\kappa$ の代数閉包 $\overline{\kappa}$ を選び、
% P02-E1206: source lines 31344-31345 omit “by” in the notation clause.
$\kappa$ の分離代数閉包を
$\kappa^{sep} \subset \overline{\kappa}$ と書く。
$A^{sh}$ を $\kappa^{sep}$ に関する $A$ の強 Hensel 化とする。
写像
$$
\Spec(A') \xleftarrow{c} \Spec((A')^{sh}) \xrightarrow{e} \Spec(A^{sh})
$$
を考える。ここで $(A')^{sh} = A' \otimes_A A^{sh}$ である。このとき
\begin{enumerate}
\item 極大イデアル $\mathfrak m' \subset A'$ に対し、
剰余体 $\kappa'$ は $\kappa$ 上代数的であり、
$\mathfrak m'$ 上の $c$ のファイバーは標準的に
$\Hom_\kappa(\kappa', \overline{\kappa})$ と同一視できる。
\item 右の矢印は極小素イデアル上全単射である。
\item $(A')^{sh}$ の各極小素イデアルは唯一の極大イデアルに含まれ、
各極大イデアルは唯一の極小素イデアルを含む。
\end{enumerate}
\end{lemma}

\begin{proof}
証明は Lemma \ref{more-algebra-lemma-branches} の証明とほとんど同じである。
$I \subset A$ を冪零元全体のイデアルとする。
Algebra, Lemma \ref{algebra-lemma-quotient-henselization} により
$(A/I)^{sh} = A^{sh}/IA^{sh}$ である。
% P02-E1207: source line 31369 writes (A')^h although the strict-henselization construction requires (A')^{sh}.
$A$, $A^{sh}$, $A'$, $(A')^h$ のスペクトルは、それぞれ
$A/I$, $A^{sh}/IA^{sh}$, $A'$,
$(A')^{sh} = A' \otimes_{A/I} A^{sh}/IA^{sh}$ のスペクトルと同じである。
従って $A$ を $A_{red} = A/I$ で置き換え、
$A$ は被約であると仮定してよい。
このとき $A \subset A'$ であり、以下では断りなくこれを用いる。

\medskip\noindent
(1) の証明。$A'$ は $A$ 上整なので、体拡大
$\kappa'/\kappa$ は代数的である。
$A'$ は $A$ 上整だから、$(A')^{sh}$ は $A^{sh}$ 上整である。
going up（Algebra, Lemma
\ref{algebra-lemma-integral-going-up}）により、
$A'$ および $(A')^{sh}$ の各極大イデアルは、それぞれ
$A$ および $A^h$ の極大イデアル
% P02-E1208: source lines 31381-31386 switch from the strict henselization A^{sh} to A^h in the base ring and tensor display.
$\mathfrak m$ および $\mathfrak m^{sh}$ の上にある。さらに
$$
(A')^{sh} \otimes_{A^{sh}} \kappa^{sep} =
A' \otimes_A A^h \otimes_{A^h} \kappa^{sep} =
(A' \otimes_A \kappa) \otimes_{\kappa} \kappa^{sep}
$$
である。実際、$A^{sh}$ の剰余体は $\kappa^{sep}$ である。
従って $\mathfrak m'$ 上の $c$ のファイバーは
$\kappa' \otimes_\kappa \kappa^{sep}$ のスペクトルである。
全単射
$$
\Hom_\kappa(\kappa', \overline{\kappa}) \to
\Spec(\kappa' \otimes_\kappa \kappa^{sep}),\quad
\sigma \mapsto \Ker(
\sigma \otimes 1 : \kappa' \otimes_\kappa \kappa^{sep} \to \overline{\kappa}
)
$$
があるので (1) が成り立つ。
% P02-E1209: source lines 31399-31400 omit the article in “hence spectrum of this ring”.
以下では、表示された環が体上整であり、従ってそのスペクトルが
副有限空間であることも用いる。Algebra, Lemmas
\ref{algebra-lemma-integral-over-field} および
\ref{algebra-lemma-ring-with-only-minimal-primes} を参照せよ。

\medskip\noindent
(3) の証明。$A'$ は正規環であり、実際、正規整域の有限積である。
Algebra, Lemma
\ref{algebra-lemma-characterize-reduced-ring-normal} を参照せよ。
$A^{sh}$ は \'etale $A$-代数のフィルター余極限だから、
% P02-E1210: source line 31409 omits “a” before “filtered colimit”.
$(A')^{sh}$ は \'etale $A'$-代数のフィルター余極限であり、
従って Algebra, Lemmas \ref{algebra-lemma-normal-goes-up} および
\ref{algebra-lemma-colimit-normal-ring} により $(A')^{sh}$ は正規環である。
従って $(A')^{sh}$ の各局所環は正規整域であり、
各極大イデアルは唯一の極小素イデアルを含む。
% P02-E1211: source lines 31415-31417 say “By Lemma ... applied ... to see”, with a duplicated infinitive.
Lemma \ref{more-algebra-lemma-integral-over-henselian-pair} を
$A^{sh} \to (A')^{sh}$ に適用すると、
$((A')^{sh}, \mathfrak m(A')^{sh})$ は Hensel 対であることが分かる。
$\mathfrak q \subset (A')^{sh}$ が極小素イデアル
（または任意の素イデアル）なら、Lemma
\ref{more-algebra-lemma-irreducible-henselian-pair-connected} により
$V(\mathfrak q)$ と $V(\mathfrak m (A')^{sh})$ の共通部分は連結である。
% P02-E1212: source line 31420 lacks terminal punctuation.
% P02-E1213: source lines 31421-31423 use undefined kappa^{sh} instead of the chosen kappa^{sep} and contain a dangling “by”.
$V(\mathfrak m (A')^{sh}) = \Spec((A')^{sh} \otimes \kappa^{sh})$
は副有限空間なので、$\mathfrak q$ を含む極大イデアルは唯一である。

\medskip\noindent
(2) の証明。$A'$ の極小素イデアルは、ちょうど $A$ の
極小素イデアルの上にある素イデアルである（構成による）。
$A' \to (A')^{sh}$ は平坦なので、going down
（Algebra, Lemma \ref{algebra-lemma-flat-going-down}）により
$(A')^{sh}$ の各極小素イデアルは $A'$ の極小素イデアルの上にある。
逆に、$A'$ の極小素イデアルの上にある $(A')^{sh}$ の任意の
素イデアルは極小である。実際、$(A')^{sh}$ は $A'$ 上 \'etale、
従って準有限な代数のフィルター余極限である
% P02-E1214: source line 31433 explicitly leaves this implication as “small detail omitted”.
（細部は省略する）。
従って $(A')^{sh}$ の極小素イデアルは、ちょうど $A$ の
極小素イデアルの上にある素イデアルである。
同様に、$A^{sh}$ の極小素イデアルは、ちょうど $A$ の
極小素イデアルの上にある素イデアルである。
構成により $A' \otimes_A Q(A) = Q(A)$ である。ここで
$Q(A)$ は被約局所環 $A$ の全分数環である。
もちろん $Q(A)$ は $A$ の極小素イデアルの剰余体の有限積である。
従って
$$
(A')^{sh} \otimes_A Q(A) =
A^{sh} \otimes_A A' \otimes_A Q(A) = A^{sh} \otimes_A Q(A)
$$
である。上の議論により、左辺の環のスペクトルは
$(A')^{sh}$ の極小素イデアル全体であり、右辺の環のスペクトルは
% P02-E1215: source lines 31448-31449 duplicate “is the”.
$A^{sh}$ の極小素イデアル全体である。これで証明が完了する。
\end{proof}

\begin{lemma}
\label{more-algebra-lemma-geometrically-unibranch}
\begin{reference}
\cite[Lemma 2.2]{Etale-coverings} および
\cite[Chapter IV Proposition 18.8.15]{EGA4}
\end{reference}
$A$ を局所環とし、$A^{sh}$ を $A$ の強 Hensel 化とする。
% P02-E1216: source line 31459 lacks punctuation after “The following are equivalent”.
次は同値である。
\begin{enumerate}
\item $A$ は幾何学的単枝である。
\item $A^{sh}$ は唯一の極小素イデアルをもつ。
\end{enumerate}
\end{lemma}

\begin{proof}
これは Lemma \ref{more-algebra-lemma-geometric-branches} から従うが、直接証明も与える。
この直接証明は Lemma \ref{more-algebra-lemma-unibranch} の直接証明とほとんど同じである。
% P02-E1217: source lines 31471-31472 omit “by” in both notation clauses and omit “respectively”.
環 $A$ の極大イデアルを $\mathfrak m$ と書く。
$A$, $A^{sh}$ の剰余体をそれぞれ $\kappa$, $\kappa^{sh}$ と書く。

\medskip\noindent
(2) を仮定する。$\mathfrak p^{sh}$ を $A^{sh}$ の唯一の
極小素イデアルとする。$A \to A^{sh}$ の平坦性により
$\mathfrak p = A \cap \mathfrak p^{sh}$ は $A$ の唯一の
極小素イデアルである（going down による。Algebra, Lemma
\ref{algebra-lemma-flat-going-down} を参照せよ）。
また、Algebra, Lemma
\ref{algebra-lemma-quotient-strict-henselization} を参照すると
$A^{sh}/\mathfrak pA^{sh} = (A/\mathfrak p)^{sh}$ であり、これは
Lemma \ref{more-algebra-lemma-henselization-reduced} により被約なので、
$\mathfrak p^{sh} = \mathfrak pA^{sh}$ である。
$A'$ を $A/\mathfrak p$ の分数体における整閉包とする。
$A'$ が局所環であり、その剰余体が $\kappa$ 上純非分離であることを
示さなければならない。
$A \to A'$ は整だから、$A'$ のすべての極大イデアルは
$\mathfrak m$ の上にある（整環準同型に対する going up による。
Algebra, Lemma \ref{algebra-lemma-integral-going-up} を参照せよ）。
$A'$ が局所環でなければ、相異なる極大イデアル
$\mathfrak m_1$, $\mathfrak m_2$ をとれる。
$f_i \in \mathfrak m_i, f_i \not \in \mathfrak m_{3 - i}$ を満たす
元 $f_1, f_2 \in A'$ を選ぶと、有限部分代数
$B = A[f_1, f_2] \subset A'$ は、相異なる極大イデアル
$B \cap \mathfrak m_i$, $i = 1, 2$ をもつ。
$A'$ が極大イデアル $\mathfrak m'$ をもつ局所環であるが
$A/\mathfrak m \subset A'/\mathfrak m'$ が純非分離でなければ、
$A'/\mathfrak m'$ における像が $A/\mathfrak m$ の有限な
非純非分離拡大を生成する $f \in A'$ をとれる。
従って有限局所部分代数 $B = A[f] \subset A'$ で、
その剰余体が $A/\mathfrak m$ の純非分離拡大でないものをとれる。
包含
$$
A/\mathfrak p \subset B \subset \kappa(\mathfrak p)
$$
を平坦な環準同型 $A \to A^{sh}$ でテンソルすると、包含
$$
A^{sh}/\mathfrak p^{sh} \subset
B \otimes_A A^{sh} \subset
\kappa(\mathfrak p) \otimes_A A^{sh} \subset
\kappa(\mathfrak p^{sh})
$$
を得る。最後の包含が成り立つのは
$$
\kappa(\mathfrak p) \otimes_A A^{sh} =
\kappa(\mathfrak p) \otimes_{A/\mathfrak p} A^{sh}/\mathfrak p^{sh}
$$
が整域 $A^{sh}/\mathfrak p^{sh}$ の局所化だからである。
$B/\mathfrak mB$ は2つの極大イデアルをもつか、またはその唯一の
剰余体が $\kappa$ 上純非分離でなく、かつ $\kappa^{sh}$ は
分離代数閉なので、$B \otimes_A \kappa^{sh}$ は少なくとも
2つの極大イデアルをもつことに注意する。
従って $A^{sh}$ は強 Hensel だから、Algebra, Lemma
\ref{algebra-lemma-mop-up-strictly-henselian} により
$B \otimes_A A^{sh}$ は $\geq 2$ 個の局所環の積である。
しかし $B \otimes_A A^{sh}$ は整域の部分環であることを既に見たので、
矛盾を得る。

\medskip\noindent
(1) を仮定する。$\mathfrak p \subset A$ を唯一の極小素イデアルとし、
$A'$ を $A/\mathfrak p$ の分数体における整閉包とする。
$A \to B$ を、\'etale $A$-代数の局所化である局所環の
局所準同型とする。特に $\mathfrak m_B$ は
$\mathfrak m_AB$ を含む唯一の素イデアルである。
このとき $B' = A' \otimes_A B$ は $B$ 上整であり、
$A \to A'$ が純非分離な剰余体拡大をもつ局所準同型であるという仮定から、
$B'$ は局所環である
（Algebra, Lemma \ref{algebra-lemma-local-tensor-with-integral}）。
一方、$A' \to B'$ は \'etale 環準同型の局所化なので、
Algebra, Lemma \ref{algebra-lemma-normal-goes-up} により
$B'$ は正規である。従って $B'$ は（局所）正規整域である。
最後に
$$
B/\mathfrak pB \subset B \otimes_A \kappa(\mathfrak p)
= B' \otimes_{A'} (\text{fraction field of }A') \subset
\text{fraction field of }B'
$$
である。従って $B/\mathfrak pB$ は整域であり、これは $B$ が
唯一の極小素イデアルをもつことを含意する
（$A \to B$ の平坦性により、極小素イデアルはすべて
$\mathfrak p$ の上になければならない）。
$A^{sh}$ は局所環 $B$ のフィルター余極限だから、
$A^{sh}$ は唯一の極小素イデアルをもつ。
実際、ある非冪零元 $f, g$ に対して $A^{sh}$ 内で $fg = 0$ なら、
$f$ と $g$ の両方を含む上のような $B$ をとることができ、矛盾に至る。
\end{proof}

\begin{definition}
\label{more-algebra-definition-number-of-branches}
$A$ を、Hensel 化 $A^h$ と強 Hensel 化 $A^{sh}$ をもつ局所環とする。
有限ならば $A^h$ の極小素イデアルの個数を、そうでなければ
$\infty$ を、{it $A$ の枝数}という。
有限ならば $A^{sh}$ の極小素イデアルの個数を、そうでなければ
$\infty$ を、{it $A$ の幾何学的枝数}という。
\end{definition}

\noindent
Definition \ref{more-algebra-definition-unibranch} との関係を詳しく述べる。

\begin{lemma}
\label{more-algebra-lemma-number-of-branches-1}
$(A, \mathfrak m, \kappa)$ を局所環とする。
\begin{enumerate}
\item $A$ が無限個の極小素イデアルをもつなら、
$A$ の（幾何学的）枝数は $\infty$ である。
\item $A$ の枝数が $1$ であることと $A$ が単枝であることは同値である。
\item $A$ の幾何学的枝数が $1$ であることと
$A$ が幾何学的単枝であることは同値である。
\end{enumerate}
$A$ は有限個の極小素イデアルをもち、$A'$ は
$A_{red}$ の全商環における $A$ の整閉包であると仮定する。このとき
\begin{enumerate}
\item[(4)] $A$ の枝数は $A'$ の極大イデアル
$\mathfrak m'$ の個数である。
\item[(5)] $A$ の幾何学的枝数を得るには、$A'$ の各極大イデアル
$\mathfrak m'$ を、$\kappa(\mathfrak m')/\kappa$ の
分離次数を重複度として数えなければならない。
\end{enumerate}
\end{lemma}

\begin{proof}
この補題は定義、Lemma \ref{more-algebra-lemma-branches}、
Lemma \ref{more-algebra-lemma-geometric-branches}、および Fields, Lemma
\ref{fields-lemma-separable-degree} から直ちに従う。
\end{proof}

\begin{lemma}
\label{more-algebra-lemma-invariance-number-branches-smooth}
$A \to B$ を、滑らかな環準同型の局所化である局所環の
局所準同型とする。
\begin{enumerate}
\item $A$ の幾何学的枝数は $B$ の幾何学的枝数に等しい。
\item $A \to B$ が剰余体の純非分離拡大を誘導するなら、
$A$ の枝数は $B$ の枝数に等しい。
\end{enumerate}
\end{lemma}

\begin{proof}
以下を用いる。滑らかな環準同型は平坦であり
（Algebra, Lemma \ref{algebra-lemma-smooth-syntomic}）、
局所化は平坦であり（Algebra, Lemma
\ref{algebra-lemma-flat-localization}）、
平坦な環準同型の合成は平坦であり
（Algebra, Lemma \ref{algebra-lemma-composition-flat}）、
平坦な環準同型の底変換は平坦であり
（Algebra, Lemma \ref{algebra-lemma-flat-base-change}）、
平坦な局所準同型は忠実平坦であり
（Algebra, Lemma \ref{algebra-lemma-local-flat-ff}）、
（強）Hensel 化は平坦であり
（Lemma \ref{more-algebra-lemma-dumb-properties-henselization}）、そして
平坦な環準同型に対して going down が成り立つ
（Algebra, Lemma \ref{algebra-lemma-flat-going-down}）。

\medskip\noindent
(2) の証明。$A^h$, $B^h$ を $A$, $B$ の Hensel 化とする。
このとき $B^h$ は、$\mathfrak m_B$ の上にある唯一の極大イデアルにおける
$A^h \otimes_A B$ の Hensel 化である。Algebra, Lemma
\ref{algebra-lemma-henselian-functorial-improve} を参照せよ。
従って $A$ は Hensel であると仮定してよく、そう仮定する。
$A \to B \to B^h$ は平坦だから、$B^h$ の各極小素イデアルは
$A$ の極小素イデアルの上にある。また $A \to B^h$ は忠実平坦だから、
$A$ の各極小素イデアルの上には $B^h$ の極小素イデアルがある。
いずれの場合も平坦環準同型に対する going down を用いる。
従って、与えられた極小素イデアル $\mathfrak p \subset A$ に対し、
$\mathfrak p$ の上にある $B^h$ の極小素イデアルが
高々1つであることを示せば十分である。
$A$ を $A/\mathfrak p$ で、$B$ を $B/\mathfrak p B$ で置き換えると、
$A$ は整域であると仮定してよい。
% P02-E1218: source line 31637 says “the A is still henselian”.
Algebra, Lemma \ref{algebra-lemma-quotient-henselization} により、
$A$ はなお Hensel である。
Lemma \ref{more-algebra-lemma-unibranch} により、$A$ の分数体における
整閉包 $A'$ は局所整域である。もちろん $A'$ は正規整域である。
Algebra, Lemma \ref{algebra-lemma-normal-goes-up} により
$A' \otimes_A B^h$ は正規環である
（この補題が直接与えるのは $A' \otimes_A B$ に対する主張だけである。
$A' \otimes_A B^h$ まで進むには、$B^h$ が \'etale
$B$-代数の余極限であることと Algebra, Lemma
\ref{algebra-lemma-colimit-normal-ring} を用いる）。
Algebra, Lemma \ref{algebra-lemma-local-tensor-with-integral} により
$A' \otimes_A B^h$ は局所環である
（ここで $A$ と $B$ の剰余体に関する仮定を用いる）。
従って $A' \otimes_A B^h$ は局所正規環、従って局所整域である。
$A \to B^h$ の平坦性により
$B^h \subset A' \otimes_A B^h$ だから、
望む通り $B^h$ は整域である。

\medskip\noindent
(1) の証明。$A^{sh}$, $B^{sh}$ を $A$, $B$ の強 Hensel 化とする。
% P02-E1219: source lines 31655-31659 use A^h and m_{A^h} in the strict-henselization argument, where A^{sh} and m_{A^{sh}} are required.
このとき $B^{sh}$ は、$\mathfrak m_B$ および
$\mathfrak m_{A^h}$ の上にある極大イデアルにおける
$A^h \otimes_A B$ の強 Hensel 化である。Algebra, Lemma
\ref{algebra-lemma-strictly-henselian-functorial-improve} を参照せよ。
従って $A$ は強 Hensel であると仮定してよく、そう仮定する。
$A \to B \to B^{sh}$ は平坦だから、$B^{sh}$ の各極小素イデアルは
$A$ の極小素イデアルの上にある。また $A \to B^{sh}$ は忠実平坦だから、
$A$ の各極小素イデアルの上には $B^{sh}$ の極小素イデアルがある。
いずれの場合も平坦環準同型に対する going down を用いる。
従って、与えられた極小素イデアル $\mathfrak p \subset A$ に対し、
$\mathfrak p$ の上にある $B^{sh}$ の極小素イデアルが
高々1つであることを示せば十分である。
$A$ を $A/\mathfrak p$ で、$B$ を $B/\mathfrak p B$ で置き換えると、
$A$ は整域であると仮定してよい。このとき Algebra, Lemma
\ref{algebra-lemma-quotient-strict-henselization} により
$A$ はなお強 Hensel である。
Lemma \ref{more-algebra-lemma-geometrically-unibranch} により、$A$ の分数体における
整閉包 $A'$ は局所整域で、その剰余体は $A$ の剰余体の
純非分離拡大である。もちろん $A'$ は正規整域である。
Algebra, Lemma \ref{algebra-lemma-normal-goes-up} により
$A' \otimes_A B^{sh}$ は正規環である
（この補題が直接与えるのは $A' \otimes_A B$ に対する主張だけである。
$A' \otimes_A B^{sh}$ まで進むには、$B^{sh}$ が \'etale
$B$-代数の余極限であることと Algebra, Lemma
\ref{algebra-lemma-colimit-normal-ring} を用いる）。
Algebra, Lemma \ref{algebra-lemma-local-tensor-with-integral} により
$A' \otimes_A B^{sh}$ は局所環である
（$A \subset A'$ が純非分離な剰余体拡大を誘導するためである）。
従って $A' \otimes_A B^{sh}$ は局所正規環、従って局所整域である。
$A \to B^{sh}$ の平坦性により
$B^{sh} \subset A' \otimes_A B^{sh}$ だから、
望む通り $B^{sh}$ は整域である。
\end{proof}







\section{枝に関する雑録}
\label{more-algebra-section-branches}

\noindent
Section \ref{more-algebra-section-unibranch} で定義した局所環の枝に関する
いくつかの結果を述べる。

\begin{lemma}
\label{more-algebra-lemma-ideal-inverse-nonzero}
$A$ と $B$ を整域とし、$A \to B$ を環準同型とする。
$A \to B$ はさらに次の性質の少なくとも1つをもつと仮定する。
\begin{enumerate}
\item \'etale 環準同型の局所化である。
\item 平坦で、不分岐環準同型の局所化である。
\item 平坦で、準有限環準同型の局所化である。
\item 平坦で、整環準同型の局所化である。
\item 平坦であり、$\Spec(B) \to \Spec(A)$ のファイバーの点の間に
非自明な特殊化がない。
\item $\Spec(B) \to \Spec(A)$ は生成点を生成点に写し、
ファイバーの点の間に非自明な特殊化がない。
\item $\Spec(B)$ のちょうど1点が $\Spec(A)$ の生成点に写る。
\end{enumerate}
このとき $B$ の任意の非零イデアル $J$ に対して
$A \cap J$ は非零である。
\end{lemma}

\begin{proof}
(7) の場合の証明。
$A$, $B$ の分数体をそれぞれ $K$, $L$ とする。
Algebra, Lemma
\ref{algebra-lemma-minimal-prime-image-minimal-prime} により、
生成点 $(0) \in \Spec(A)$ に写る
$\Spec(B)$ の唯一の点は $(0) \in \Spec(B)$ である。
従って $B \otimes_A K$ は唯一の素イデアルをもつ環であり、
その剰余体は $L$ である
（実際、この環は $L$ に等しいが、それは必要ない）。
非零元 $b \in J$ を選ぶ。このとき $b$ は $L$ の単元に写る。
従って Algebra, Lemma
\ref{algebra-lemma-surjective-on-spec-units} により
$b$ は $B \otimes_A K$ の単元に写る。
$$
B \otimes_A K = \colim_{f \in A \setminus \{0\}} B_f
$$
だから、ある非零 $f \in A$ に対して $b$ は $B_f$ の単元に写る。
これはある $b' \in B$ と $n \geq 1$ に対して
$b b' = f^n$ であることを意味する。
従って望む通り $f^n \in A \cap J$ である。

\medskip\noindent
証明の残りでは、他の各仮定が (7) を含意することを示す。
仮定 (1) -- (5) のもとで環準同型 $A \to B$ は平坦であり、
従って $A \to B$ は単射である。
% P02-E1220: source lines 31740-31742 invokes faithful flatness of flat local maps although A -> B was not assumed local.
（平坦な局所準同型は Algebra, Lemma
\ref{algebra-lemma-local-flat-ff} により忠実平坦だからである。）
従って $\Spec(B)$ の生成点は $\Spec(A)$ の生成点に写る。
$\Spec(B) \to \Spec(A)$ のファイバーの点の間に非自明な特殊化がなければ、
もちろん $\Spec(B)$ のこの生成点は $\Spec(A)$ の生成点に写る
唯一の点でなければならない。従って (6) は (7) を含意する。
最後に、(1) -- (5) の場合に
$\Spec(B) \to \Spec(A)$ のファイバーの点の間に非自明な特殊化が
ないことを示せば証明が完了する。
整の場合は Algebra, Lemma
\ref{algebra-lemma-integral-no-inclusion} を、
準有限の場合は Algebra, Definition
\ref{algebra-definition-quasi-finite} を参照し、
不分岐および \'etale 環準同型が準有限であること
（Algebra, Lemmas \ref{algebra-lemma-unramified-quasi-finite} および
\ref{algebra-lemma-etale-quasi-finite}）を用いればよい。
\end{proof}

\begin{lemma}
\label{more-algebra-lemma-local-unramified-extension-unibranch-domain-is-etale}
$A \to B$ を環準同型とする。$\mathfrak q \subset B$ を
素イデアル $\mathfrak p \subset A$ の上にある素イデアルとする。
次を仮定する。
\begin{enumerate}
\item $A$ は整域である。
\item $A_\mathfrak p$ は幾何学的単枝である。
\item $A \to B$ は $\mathfrak q$ において不分岐である。
\item $A_\mathfrak p \to B_\mathfrak q$ は単射である。
\end{enumerate}
このとき $g \in B$, $g \not \in \mathfrak q$ で、
$B_g$ が $A$ 上 \'etale となるものが存在する。
\end{lemma}

\begin{proof}
Algebra, Proposition
\ref{algebra-proposition-unramified-locally-standard} により、
$B$ を主局所化で置き換えた後、標準 \'etale 環準同型
$A \to B'$ と全射 $B' \to B$ をとれる。
$\mathfrak q$ の逆像を $\mathfrak q' \subset B'$ と書く。
主局所化で $B'$ をさらに置き換えた後、
$B' \to B$ が単射であることを示す。

\medskip\noindent
この段落では $B'$ が整域である場合に帰着する。
$A$ は整域だから、Algebra, Lemma
\ref{more-algebra-lemma-reduced-goes-up} により $B'$ は被約である。
$K$ を $A$ の分数体とする。このとき $B' \otimes_A K$ は
体上 \'etale であり、従って体の有限積である。
Algebra, Lemma \ref{algebra-lemma-etale-over-field} を参照せよ。
$A \to B'$ は \'etale（従って平坦）なので、
% P02-E1221: source line 31792 says “the minimal primes ... are lie over”.
$B'$ の極小素イデアルは $(0) \subset A$ の上にある
（平坦環準同型に対する going down による）。
従って $B'$ は有限個の極小素イデアル
$\mathfrak r_1, \ldots, \mathfrak r_r \subset B'$ をもつ。
$A_\mathfrak p$ は幾何学的単枝で $A \to B'$ は \'etale だから、
Lemmas \ref{more-algebra-lemma-invariance-number-branches-smooth} および
\ref{more-algebra-lemma-number-of-branches-1} により
$B'_{\mathfrak q'}$ は整域である。
従って $\mathfrak q' \supset \mathfrak r_i$ となるのは、
ちょうど1つの $i = i_0$ に対してである。
Algebra, Lemma \ref{algebra-lemma-silly} により
$g' \in B'$, $g' \not \in \mathfrak r_{i_0}$ かつ
$i \not = i_0$ に対して $g' \in \mathfrak r_i$ となるものを選ぶ。
$B'$ と $B$ を $B'_{g'}$ と $B_{g'}$ で置き換えた後、
$B'$ は整域となる。

\medskip\noindent
$B'$ は整域であると仮定する。特に
$B' \subset B'_{\mathfrak q'}$ である。
$B' \to B$ が単射でなければ、
$J = \Ker(B'_{\mathfrak q'} \to B_\mathfrak q)$ は非零である。
Lemma \ref{more-algebra-lemma-ideal-inverse-nonzero} を
$A_\mathfrak p \to B'_{\mathfrak q'}$ に適用すると、
$B_\mathfrak q$ で零に写る非零元 $a \in A_\mathfrak p$ が得られる。
これは仮定 (4) に反する。これで証明が完了する。
\end{proof}

\begin{lemma}
\label{more-algebra-lemma-unramified-extension-unibranch-domain-is-etale}
\begin{reference}
\cite[Expose I, Theorem 9.5 part (ii)]{SGA1} の一般化
\end{reference}
$(A, \mathfrak m)$ を幾何学的単枝な局所整域とする。
$A \to B$ を、本質的有限型な、局所環の単射局所準同型とする。
$\mathfrak m B$ が $B$ の極大イデアルで、誘導される剰余体拡大が
分離的なら、$A \to B$ は \'etale 環準同型の局所化である。
\end{lemma}

\begin{proof}
$B = C_\mathfrak q$ と書ける。ここで $A \to C$ は有限型環準同型で、
$\mathfrak q \subset C$ は $\mathfrak m$ の上にある素イデアルである。
Algebra, Lemma \ref{algebra-lemma-characterize-unramified} により
環準同型 $A \to C$ は $\mathfrak q$ において不分岐である。
Algebra, Proposition
\ref{algebra-proposition-unramified-locally-standard} により、
$C$ を主局所化で置き換えた後、標準 \'etale 環準同型
$A \to C'$ と全射 $C' \to C$ をとれる。
% P02-E1222: source line 31839 omits “by” in “Denote q' ... the inverse image”.
$\mathfrak q$ の逆像を $\mathfrak q' \subset C'$ と書き、
$B' = C'_{\mathfrak q'}$ とおく。このとき $B' \to B$ は全射である。
$B' \to B$ も単射であることを示せば十分である。

\medskip\noindent
$A$ は整域だから、Algebra, Lemma
\ref{more-algebra-lemma-reduced-goes-up} により環 $C'$ と $B'$ は被約である。
$A$ は幾何学的単枝だから、Lemmas
\ref{more-algebra-lemma-invariance-number-branches-smooth} および
\ref{more-algebra-lemma-number-of-branches-1} により $B'$ は整域である。
% P02-E1223: source lines 31847-31849 say “see by Lemmas”, with a redundant connective.
$B' \to B$ が単射でなければ、Lemma
\ref{more-algebra-lemma-ideal-inverse-nonzero} により
$A \cap \Ker(B' \to B)$ は非零である。
これは $A \to B$ が単射であるという仮定に反する。
\end{proof}

\begin{lemma}
\label{more-algebra-lemma-minimal-primes-tensor-strictly-henselian}
$k$ を代数閉体とする。$A$, $B$ を、剰余体が $k$ に等しい
強 Hensel 局所 $k$-代数とする。
$C$ を、極大イデアル
$\mathfrak m_A \otimes_k B + A \otimes_k \mathfrak m_B$
における $A \otimes_k B$ の強 Hensel 化とする。
% P02-E1224: source lines 31861-31862 omit “are in” in the one-to-one correspondence clause.
このとき $C$ の極小素イデアルは、$A$ と $B$ の極小素イデアルの対と
$1$ 対 $1$ に対応する。
\end{lemma}

\begin{proof}
まず、$C$ の極小素イデアル $\mathfrak r$ は $A$ の極小素イデアル
$\mathfrak p$ に写り、$B$ の極小素イデアル $\mathfrak q$ に写る。
実際、環準同型 $A \to C$ と $B \to C$ は平坦である
（平坦環準同型に対する going down、Algebra, Lemma
\ref{algebra-lemma-flat-going-down} による）。
従って
$$
(A/\mathfrak p \otimes_k B/\mathfrak q)_{
\mathfrak m_A \otimes_k B + A \otimes_k \mathfrak m_B}
$$
の強 Hensel 化が唯一の極小素イデアルをもつことを示せば十分である。
Algebra, Lemma
\ref{algebra-lemma-quotient-strict-henselization} により、
環 $A/\mathfrak p$, $B/\mathfrak q$ は強 Hensel である。
従って $A$ と $B$ は強 Hensel 局所整域であると仮定してよく、
$C$ が唯一の極小素イデアルをもつことを示すのが目標となる。
Lemma \ref{more-algebra-lemma-geometrically-unibranch} により、
$A$ の分数体における整閉包 $A'$ は剰余体 $k$ をもつ
正規局所整域である。$B$ の分数体における整閉包 $B'$ についても同様である。
% P02-E1225: source line 31881 says “integral closure ... into its fraction field”.
Algebra, Lemma
\ref{algebra-lemma-geometrically-normal-tensor-normal} により
$A' \otimes_k B'$ は正規環である。従ってその局所化
$$
R = (A' \otimes_k B')_{
\mathfrak m_{A'} \otimes_k B' + A' \otimes_k \mathfrak m_{B'}}
$$
は正規局所整域である。
$A \otimes_k B \to A' \otimes_k B'$ は整であること
% P02-E1226: source line 31889 misspells “going up” as “gong up”.
（従って going up が成り立つ。Algebra, Lemma
\ref{algebra-lemma-integral-going-up}）、
および
$\mathfrak m_{A'} \otimes_k B' + A' \otimes_k \mathfrak m_{B'}$
が $\mathfrak m_A \otimes_k B + A \otimes_k \mathfrak m_B$
の上にある $A' \otimes_k B'$ の唯一の極大イデアルであることに注意する。
従って Algebra, Lemma
\ref{algebra-lemma-unique-prime-over-localize-below} により
$$
R = (A' \otimes_k B')_{
\mathfrak m_A \otimes_k B + A \otimes_k \mathfrak m_B}
$$
である。従って
$$
(A \otimes_k B)_{
\mathfrak m_A \otimes_k B + A \otimes_k \mathfrak m_B}
\longrightarrow
R
$$
は整である。従って $R$ は
$$
(A \otimes_k B)_{
\mathfrak m_A \otimes_k B + A \otimes_k \mathfrak m_B}
$$
の分数体における整閉包である。
Lemma \ref{more-algebra-lemma-geometrically-unibranch} をもう一度用いると、
% P02-E1227: source line 31913 concludes that C has a unique prime ideal, although the proved goal is a unique minimal prime.
$C$ は唯一の素イデアルをもつと結論できる。
\end{proof}







\section{完備化の枝}
\label{more-algebra-section-branches-completion}

\noindent
$(A, \mathfrak m)$ を Noether 局所環とする。写像
$A \to A^h \to A^\wedge$ を考える。
一般に写像 $A^h \to A^\wedge$ は極小素イデアル上の全単射を
誘導するとは限らない。Examples, Section
\ref{examples-section-bad} を参照せよ。言い換えると、
Definition \ref{more-algebra-definition-number-of-branches} で定義した
$A$ の枝数は $A^\wedge$ の枝数と異なることがある。
しかし、ある条件のもとでは枝数は同じである。
例えば $A$ の次元が $1$ ならそうである。

\begin{lemma}
\label{more-algebra-lemma-nr-branches-completion}
$(A, \mathfrak m)$ を Noether 局所環とする。
\begin{enumerate}
\item 写像 $A^h \to A^\wedge$ は、$A^\wedge$ の極小素イデアルから
$A^h$ の極小素イデアルへの全射を定める。
\item $A$ の枝数は $A^\wedge$ の枝数以下である。
\item $A$ の幾何学的枝数は $A^\wedge$ の幾何学的枝数以下である。
\end{enumerate}
\end{lemma}

\begin{proof}
Lemma \ref{more-algebra-lemma-henselization-noetherian} により写像
$A^h \to A^\wedge$ は平坦かつ単射である。
going down（Algebra, Lemma
\ref{algebra-lemma-flat-going-down}）と Algebra, Lemma
\ref{algebra-lemma-injective-minimal-primes-in-image} を組み合わせると、
(1) が従う。(2) はこれと Definition
\ref{more-algebra-definition-number-of-branches}、および $A^\wedge$ が
Hensel であること（Algebra, Lemma
\ref{algebra-lemma-complete-henselian}）から従う。
Lemma \ref{more-algebra-lemma-henselization-noetherian} により
$(A^\wedge)^{sh} = A^{sh} \otimes_{A^h} A^\wedge$ である。
従って平坦単射
$A^{sh} \to (A^\wedge)^{sh}$ を用いて上の議論を繰り返せば、
(3) が証明される。
\end{proof}

\begin{lemma}
\label{more-algebra-lemma-equal-nr-branches-completion}
$(A, \mathfrak m)$ を Noether 局所環とする。
$A$ の枝数が $A^\wedge$ の枝数と同じであることと、
Hensel 化の各極小素イデアル $\mathfrak q \subset A^h$ に対して
$\sqrt{\mathfrak qA^\wedge}$ が素イデアルであることは同値である。
\end{lemma}

\begin{proof}
Lemma \ref{more-algebra-lemma-nr-branches-completion} と、
$A$ および $A^\wedge$ の枝数がともに有限であることから従う。
後者は Algebra, Lemma
\ref{algebra-lemma-Noetherian-irreducible-components} と、
$A^h$ および $A^\wedge$ が Noether であること
（Lemma \ref{more-algebra-lemma-henselization-noetherian}）による。
\end{proof}

\noindent
簡単な貼り合わせ補題を述べる。

\begin{lemma}
\label{more-algebra-lemma-glueing-sum-components-open}
$A$ を環とし、$I$ を有限生成イデアルとする。
$A \to C$ を、すべての $f \in I$ に対して環準同型
$A_f \to C_f$ が冪等元における局所化となる環準同型とする。
このとき全射 $A \to C'$ で、すべての $f \in I$ に対して
$A_f \to (C \times C')_f$ が同型となるものが存在する。
\end{lemma}

\begin{proof}
$I$ の生成元 $f_1, \ldots, f_r$ を選ぶ。ある冪等元
$e_i \in A_{f_i}$ に対して
$$
C_{f_i} = (A_{f_i})_{e_i}
$$
と書く。ある $a_i \in A$ と $n \geq 0$ に対して
$e_i = a_i/f_i^n$ と書く。すべての $i = 1, \ldots, r$ に対して
同じ $n$ を用いてよい。適当な $m \gg 0$ によって
$a_i$ を $f_i^ma_i$ で、$n$ を $n + m$ で置き換えれば、
すべての $i$ に対して $a_i^2 = f_i^n a_i$ と仮定してよい。
$e_i$ は
$C_{f_if_j} = (A_{f_if_j})_{e_j} = A_{f_if_ja_j}$ において
$1$ に写るので、ある $N$ に対して
$$
(f_if_ja_j)^N(f_j^n a_i  - f_i^na_j) = 0
$$
である（すべての対 $i, j$ に対して同じ $N$ を選べる）。
$a_j^2 = f_j^na_j$ を用いると
$$
f_i^{N + n} f_j^{N + nN} a_j = f_i^N f_j^{N + n} a_ia_j^N
$$
を得る。$n$ を $n + N + nN$ に増やし、$a_i$ を
$f_i^{N + nN}a_i$ で置き換えると、すべての対 $i, j$ に対し
$f_i^n a_j$ は $a_i$ の生成するイデアルに属する。
$C' = A/(a_1, \ldots, a_r)$ とおく。このとき
$$
C'_{f_i} = A_{f_i}/(a_i) = A_{f_i}/(e_i)
$$
である。実際、$f_i$ を可逆にした後、$a_j$ は $a_i$ の生成する
イデアルに属する。
環 $B$ の冪等元 $e$ に対して $B = B_e \times B/(e)$ だから、
補題の結論は $f$ が $f_1, \ldots, f_r$ のいずれかに等しい場合に成り立つ。
Algebra, Lemma \ref{algebra-lemma-cover} の形の関数の貼り合わせを用いると、
結果はすべての $f \in I$ に対して成り立つと結論できる。
実際、$f \in I$ に対し元 $f_1, \ldots, f_r$ は $A_f$ の
単位イデアルを生成するので、$A_f \to (C \times C')_f$ が
同型であることと、$f_1, \ldots, f_r$ で局所化した後に
そうであることは同値である。
\end{proof}

\noindent
Lemma \ref{more-algebra-lemma-quotient-by-idempotent} を用いると、
形式完備化の与えられた有限型拡大から有限型拡大を構成できる。
この補題は Algebraization of Formal Spaces, Lemma
\ref{restricted-lemma-approximate-by-etale-over-complement} で一般化する。

\begin{lemma}
\label{more-algebra-lemma-quotient-by-idempotent}
$A$ を Noether 環、$I$ をイデアルとする。
$B$ を有限型 $A$-代数とする。
$B^\wedge$ を $I$-進完備化とし、$B^\wedge \to C$ を
核 $J$ をもつ全射環準同型とする。
ある $c \geq 0$ に対して $J/J^2$ が $I^c$ で零化されるなら、
$C$ は有限型 $A$-代数の完備化と同型である。
\end{lemma}

\begin{proof}
$f \in I$ とする。$B^\wedge$ は Noether だから
（Algebra, Lemma
\ref{algebra-lemma-completion-Noetherian-Noetherian}）、
$J$ は有限生成イデアルである。従って Algebra, Lemma
\ref{algebra-lemma-ideal-is-squared-union-connected} から、
ある冪等元 $e \in (B^\wedge)_f$ に対して
$$
C_f = ((B^\wedge)_f)_e
$$
であることが従う。Lemma \ref{more-algebra-lemma-glueing-sum-components-open}
により全射 $B^\wedge \to C'$ で、
任意の $f \in I$ を可逆にした後に
$B^\wedge \to C \times C'$ が同型となるものをとれる。
$C \times C'$ は有限 $B^\wedge$-代数であることに注意する。

\medskip\noindent
生成元 $f_1, \ldots, f_r \in I$ を選ぶ。
% P02-E1228: source lines 32063-32067 omit “by” in both notation clauses for alpha_i and alpha_ij.
上で得た $(B^\wedge)_{f_i}$-代数の同型の逆を
$\alpha_i : (C \times C')_{f_i} \to B_{f_i} \otimes_B B^\wedge$
と書く。明らかな $B$-代数同型を
$\alpha_{ij} : (B_{f_i})_{f_j} \to (B_{f_j})_{f_i}$ と書く。
Remark \ref{more-algebra-remark-glueing-data} で導入した圏
$\text{Glue}(B \to B^\wedge, f_1, \ldots, f_r)$ の対象
$$
(C \times C', B_{f_i}, \alpha_i, \alpha_{ij})
$$
を考える。
% P02-E1229: source lines 32072-32074 explicitly omit verification of the gluing conditions (1)(a) and (1)(b).
条件 (1)(a) および (1)(b) の検証は省略する。
$B \to B^\wedge$ は平坦であり
（Algebra, Lemma \ref{algebra-lemma-completion-flat}）、
同型 $B/IB \to B^\wedge/IB^\wedge$ を誘導するので、
Proposition \ref{more-algebra-proposition-equivalence} および Remark
\ref{more-algebra-remark-formal-glueing-algebras} を適用できる。
従って、ある有限 $B$-代数 $D$ に対して
$C \times C'$ は $D \otimes_B B^\wedge$ と同型である。
すると $D/ID \cong C/IC \times C'/IC'$ である。
因子 $C/IC$ に対応する冪等元を
$\overline{e} \in D/ID$ とする。
Lemma \ref{more-algebra-lemma-lift-idempotent-upstairs} により、
同型 $B/IB \to B'/IB'$ を誘導する \'etale 環準同型
$B \to B'$ で、$D' = D \otimes_B B'$ が
$\overline{e}$ を持ち上げる冪等元 $e$ を含むものが存在する。
$C \times C'$ は $I$-進完備だから、対
$(C \times C', IC \times IC')$ は Hensel である
（Lemma \ref{more-algebra-lemma-complete-henselian}）。
従って写像 $B \to C \times C'$ を $B'$ を経由して分解できる
（Lemma \ref{more-algebra-lemma-characterize-henselian-pair} を参照せよ）。
このようにして $B$ を $B'$ で、$D$ を $D'$ で置き換えてよい
% P02-E1230: source line 32093 misspells “completions” as “completetions”.
（これによって $I$-進完備化は変わらない）。
すると
$D = D_e \times D_{1 - e} = D/(1 - e) \times D/(e)$ は
有限型 $A$-代数の積であり、第1因子の完備化は $C$、
第2因子の完備化は $C'$ である。
\end{proof}

\begin{lemma}
\label{more-algebra-lemma-one-dimensional-formal-branch}
$(A, \mathfrak m)$ を Hensel 化 $A^h$ をもつ Noether 局所環とする。
$\mathfrak q \subset A^\wedge$ を
$\dim(A^\wedge/\mathfrak q) = 1$ を満たす極小素イデアルとする。
このとき $A^h$ の極小素イデアル $\mathfrak q^h$ で
$\mathfrak q = \sqrt{\mathfrak q^hA^\wedge}$ となるものが存在する。
\end{lemma}

\begin{proof}
$A$ と $A^h$ の完備化は同じなので、Lemma
\ref{more-algebra-lemma-henselization-noetherian} により
$A$ は Hensel であると仮定してよい。
$$
J = \Ker(A^\wedge \to (A^\wedge)_{\mathfrak q})
$$
とおき、$A^\wedge \to A^\wedge/J$ に
Lemma \ref{more-algebra-lemma-quotient-by-idempotent} を適用する。
$\dim((A^\wedge)_\mathfrak q) = 0$ だから、ある $n$ に対して
$\mathfrak q^n \subset J$ である。従って $J/J^2$ は
$\mathfrak q^n$ で零化される。
一方、$J_\mathfrak q = 0$ なので
$(J/J^2)_\mathfrak q = 0$ である。
従って $\mathfrak m$ は $J/J^2$ の唯一の随伴素イデアルであり、
$\mathfrak m$ のある冪が $J/J^2$ を零化する。
ゆえに補題を適用でき、ある有限型 $A$-代数 $C$ に対して
$A^\wedge/J = C^\wedge$ を得る。

\medskip\noindent
$A^\wedge/J$ が同じ性質をもつので
$C/\mathfrak m C = A/\mathfrak m$ である。
従って $\mathfrak m_C = \mathfrak m C$ は極大イデアルであり、
$A \to C$ は $\mathfrak m_C$ において不分岐である
（Algebra, Lemma \ref{algebra-lemma-characterize-unramified}）。
$C$ を主局所化で置き換えた後、$C$ は \'etale $A$-代数 $B$ の
商であると仮定してよい。Algebra, Proposition
\ref{algebra-proposition-unramified-locally-standard} を参照せよ。
しかし $A \to C_{\mathfrak m_C}$ の剰余体拡大は自明で
$A$ は Hensel だから、もう一度局所化した後 $B = A$ と結論できる。
従ってあるイデアル $I \subset A$ に対して $C = A/I$ であり、
$J = IA^\wedge$ が従う
（この状況では Algebra, Lemma
\ref{algebra-lemma-completion-flat} により完備化が完全だからである）。
また $A \to A^\wedge$ の平坦性により $I = J \cap A$ である。
$\mathfrak q^n \subset J \subset \mathfrak q$ だから、
$\mathfrak p = \mathfrak q \cap A$ は
$\mathfrak p^n \subset I \subset \mathfrak p$ を満たす。
すると $\sqrt{\mathfrak p A^\wedge} = \mathfrak q$ であり、
証明が完了する。
\end{proof}

\begin{lemma}
\label{more-algebra-lemma-completion-disconnected}
$(A, \mathfrak m)$ を Noether 局所環とする。
$A^\wedge$ の穴あきスペクトルが非連結であることと、
$A^h$ の穴あきスペクトルが非連結であることは同値である。
\end{lemma}

\begin{proof}
$A$ と $A^h$ の完備化は同じなので、Lemma
\ref{more-algebra-lemma-henselization-noetherian} により
$A$ は Hensel であると仮定してよい。

\medskip\noindent
$A \to A^\wedge$ は忠実平坦だから（直前の参照を見よ）、
$A^\wedge$ の穴あきスペクトルから $A$ の穴あきスペクトルへの写像は
全射である（Algebra, Lemma \ref{algebra-lemma-ff-rings} を参照せよ）。
従って $A$ の穴あきスペクトルが非連結なら、
$A^\wedge$ についても同じことが成り立つ。

\medskip\noindent
$A^\wedge$ の穴あきスペクトルは非連結であると仮定する。
これは、閉集合 $Z$ と $Z'$ によって
$$
\Spec(A^\wedge) \setminus \{\mathfrak m^\wedge\} = Z \amalg Z'
$$
と書けることを意味する。
$\overline{Z}, \overline{Z}' \subset \Spec(A^\wedge)$ を閉包とする。
あるイデアル $J, J' \subset A^\wedge$ に対して
$\overline{Z} = V(J)$, $\overline{Z}' = V(J')$ と書く。
このとき $V(J + J') = \{\mathfrak m^\wedge\}$ および
$V(JJ') = \Spec(A^\wedge)$ である。
第1の等式は $\mathfrak m^\wedge = \sqrt{J + J'}$ を意味し、
従ってある $e \geq 1$ に対して
$(\mathfrak m^\wedge)^e \subset J + J'$ である。
第2の等式により $JJ'$ の各元は冪零なので、ある $n \geq 1$ に対して
$(JJ')^n = 0$ である。これらを合わせると $J^n/J^{2n}$ は
$J^n$ および $(J')^n$ で、従って
$(\mathfrak m^\wedge)^{2en}$ で零化される。
よって Lemma \ref{more-algebra-lemma-quotient-by-idempotent} を適用すると、
有限型 $A$-代数 $C$ と同型 $A^\wedge/J^n = C^\wedge$ が得られる。

\medskip\noindent
証明の残りは Lemma
\ref{more-algebra-lemma-one-dimensional-formal-branch} の証明の後半とまったく同じである。
もちろん、あの補題はこの補題の特別な場合である。
$A^\wedge/J^n$ が同じ性質をもつので
$C/\mathfrak m C = A/\mathfrak m$ である。
従って $\mathfrak m_C = \mathfrak m C$ は極大イデアルであり、
$A \to C$ は $\mathfrak m_C$ において不分岐である
（Algebra, Lemma \ref{algebra-lemma-characterize-unramified}）。
$C$ を主局所化で置き換えた後、$C$ は \'etale $A$-代数 $B$ の
商であると仮定してよい。Algebra, Proposition
\ref{algebra-proposition-unramified-locally-standard} を参照せよ。
しかし $A \to C_{\mathfrak m_C}$ の剰余体拡大は自明で
$A$ は Hensel だから、もう一度局所化した後 $B = A$ と結論できる。
従ってあるイデアル $I \subset A$ に対して $C = A/I$ であり、
$J^n = IA^\wedge$ が従う
（この状況では Algebra, Lemma
\ref{algebra-lemma-completion-flat} により完備化が完全だからである）。
また $A \to A^\wedge$ の平坦性により $I = J^n \cap A$ である。
対称性により $I' = (J')^n \cap A$ は
$(J')^n = I'A^\wedge$ を満たす。
% P02-E1231: source line 32203 retains the original exponent e although passing from J,J' to J^n,(J')^n requires a larger power.
このとき $\mathfrak m^e \subset I + I'$ および $II' = 0$ である。
従って $V(I)$ と $V(I')$ は、$A$ の穴あきスペクトルの
求める非交和分解を与える閉部分スキームである。
\end{proof}

\begin{lemma}
\label{more-algebra-lemma-one-dimensional-number-of-branches}
$(A, \mathfrak m)$ を次元 $1$ の Noether 局所環とする。
$A$ と $A^\wedge$ の（幾何学的）枝数は同じである。
\end{lemma}

\begin{proof}
枝数についてこれを見るには、Lemmas
\ref{more-algebra-lemma-nr-branches-completion},
\ref{more-algebra-lemma-equal-nr-branches-completion}, および
\ref{more-algebra-lemma-one-dimensional-formal-branch} を組み合わせ、
Lemma \ref{more-algebra-lemma-completion-dimension} により
$A^\wedge$ の次元が1であることを用いる。
幾何学的枝数についてこれを見るには、枝数に対する結果、
強 Hensel 化で次元が変わらないこと
（Lemma \ref{more-algebra-lemma-henselization-dimension}）、および Lemma
\ref{more-algebra-lemma-henselization-noetherian} による
$(A^{sh})^\wedge = ((A^\wedge)^{sh})^\wedge$ を用いる。
\end{proof}







\begin{lemma}
\label{more-algebra-lemma-geometrically-normal-formal-fibres-number-of-branches}
\begin{reference}
\cite[Theorem 2.3]{Beddani}
\end{reference}
$(A, \mathfrak m)$ を Noether 局所環とする。
$A$ の形式ファイバーが幾何学的正規なら
（例えば $A$ が excellent または quasi-excellent なら）、
$A$ は Nagata であり、$A$ と $A^\wedge$ の（幾何学的）枝数は同じである。
\end{lemma}

\begin{proof}
正規環は被約なので、Lemma \ref{more-algebra-lemma-Nagata-local-ring} により
$A$ は Nagata である。証明の残りでは Lemma
\ref{more-algebra-lemma-formal-fibres-normal}、Proposition
\ref{more-algebra-proposition-finite-type-over-P-ring}、および Lemma
\ref{more-algebra-lemma-check-P-ring-maximal-ideals} を用いる。
これらにより、$A$ は
$P(k \to R) = $``$R$ は $k$ 上幾何学的正規である''
という P-ring であり、任意の（本質的）有限型 $A$-代数についても
同じことが成り立つ。

\medskip\noindent
$\mathfrak q \subset A$ を極小素イデアルとする。このとき
$A^\wedge/\mathfrak q A^\wedge = (A/\mathfrak q)^\wedge$
および
$A^h/\mathfrak qA^h = (A/\mathfrak q)^h$
である（Algebra, Lemma
\ref{algebra-lemma-quotient-henselization}）。
従って $A$ の枝数は環 $A/\mathfrak q$ の枝数の和であり、
$A^\wedge$ についても同様である。
このようにして $A$ が整域である場合に帰着する。

\medskip\noindent
$A$ を整域と仮定する。$A'$ を $A$ の分数体 $K$ における
$A$ の整閉包とする。$A$ は Nagata だから $A \to A'$ は有限である。
$A$ の枝数は $A'$ の極大イデアル $\mathfrak m'$ の個数であることを
思い出そう（Lemma \ref{more-algebra-lemma-branches}）。
また Algebra, Lemma
\ref{algebra-lemma-completion-finite-extension} により
$$
(A')^\wedge = A' \otimes_A A^\wedge =
\prod\nolimits_{\mathfrak m' \subset A'} (A'_{\mathfrak m'})^\wedge
$$
であることも思い出そう。
$A'_{\mathfrak m'}$ は形式ファイバーが幾何学的正規な局所環だから、
Lemma \ref{more-algebra-lemma-completion-normal-local-ring} により
$(A'_{\mathfrak m'})^\wedge$ は正規である。
従って $A' \otimes_A A^\wedge$ の極小素イデアルは、
上の分解の因子と $1$ 対 $1$ に対応する。
$A \to A^\wedge$ の平坦性により
$$
A^\wedge \subset A' \otimes_A A^\wedge \subset K \otimes_A A^\wedge
$$
である。左端と右端の環は同じ極小素イデアル集合をもつので、
中間の環についても同じことが成り立つ
% P02-E1232: source lines 32282-32284 explicitly leave the minimal-prime comparison as “small detail omitted”.
（細部は省略する）。これで証明が完了する。

\medskip\noindent
幾何学的枝数についてこれを見るには、枝数に対する結果、
$A^{sh}$ の形式ファイバーが幾何学的正規であること
（Lemmas \ref{more-algebra-lemma-formal-fibres-normal} および
\ref{more-algebra-lemma-henselization-P-ring}）、および Lemma
\ref{more-algebra-lemma-henselization-noetherian} による
$(A^{sh})^\wedge = ((A^\wedge)^{sh})^\wedge$ を用いる。
\end{proof}







\section{形式鎖状環}
\label{more-algebra-section-formally-catenary}

\noindent
この節では、Noether 局所環が普遍鎖状であることと形式鎖状であることが
同値であるという Ratliff の定理 \cite{Ratliff} を証明する。

\begin{definition}
\label{more-algebra-definition-formally-catenary}
Noether 局所環 $A$ のすべての極小素イデアル
$\mathfrak p \subset A$ に対して
$A^\wedge/\mathfrak p A^\wedge$ のスペクトルが等次元であるとき、
この環を {\it 形式鎖状} であるという。
\end{definition}

\noindent
$A$ を形式鎖状 Noether 局所環とする。
Ratliff の結果（Proposition \ref{more-algebra-proposition-ratliff}）により
$A$ の任意の商も形式鎖状である
（普遍鎖状環の類は商で安定だからである）。
従って $A$ のすべての素イデアル $\mathfrak p$ に対して
$A^\wedge/\mathfrak p A^\wedge$ のスペクトルは等次元である。

\begin{lemma}
\label{more-algebra-lemma-not-formally-catenary}
$(A, \mathfrak m)$ を、形式鎖状でない Noether 局所環とする。
このとき $A$ は普遍鎖状でない。
\end{lemma}

\begin{proof}
仮定により、極小素イデアル $\mathfrak p \subset A$ で
$A^\wedge /\mathfrak p A^\wedge$ のスペクトルが
等次元でないものが存在する。
$A$ を $A/\mathfrak p$ で置き換えると、$A$ は整域であり、
$A^\wedge$ のスペクトルは等次元でないと仮定してよい。
$A^\wedge$ の極小素イデアル $\mathfrak q$ で
$d = \dim(A^\wedge/\mathfrak q)$ が最小となるものをとる。
このとき $0 < d < \dim(A)$ である。$d$ に関する帰納法で補題を証明する。

\medskip\noindent
$d = 1$ の場合。この場合
$\dim(A^\wedge_\mathfrak q) = 0$ である。
従って $A^\wedge_\mathfrak q$ は Artin 局所環であり、
ある $n > 0$ に対してイデアル $J = \mathfrak q^n$ は
$A^\wedge_\mathfrak q$ で零に写る。
従って $\mathfrak m$ は $J/J^2$ の唯一の随伴素イデアルであり、
ある $m > 0$ に対して $\mathfrak m^m$ は $J/J^2$ を零化する。
従って Lemma \ref{more-algebra-lemma-quotient-by-idempotent} を用いると、
$B^\wedge \cong A^\wedge/J$ となる有限型 $A$-代数
$A \to B$ が得られる。
すると $\mathfrak m_B = \sqrt{\mathfrak mB}$ は
$\mathfrak m$ と同じ剰余体をもつ極大イデアルであり、
$B^\wedge$ は $\mathfrak m_B$-進完備化である
（Algebra, Lemma \ref{algebra-lemma-finite-after-completion}）。
このとき
$$
\dim(B_{\mathfrak m_B}) = \dim(B^\wedge) = 1 = d.
$$
分解 $A \to B \to A^\wedge/J$ があるので、
$\mathfrak q/J$ の逆像は $(0)$ の上にある素イデアル
$\mathfrak q' \subset \mathfrak m_B$ である。
従って $A$ が普遍鎖状なら、次元公式
（Algebra, Lemma \ref{algebra-lemma-dimension-formula}）により
\begin{align*}
\dim(B_{\mathfrak m_B})
& \geq
\dim((B/\mathfrak q')_{\mathfrak m_B}) \\
& =
\dim(A) + \text{trdeg}_A(B/\mathfrak q') -
\text{trdeg}_{\kappa(\mathfrak m)}(\kappa(\mathfrak m_B)) \\
& =
\dim(A) + \text{trdeg}_A(B/\mathfrak q')
\end{align*}
となる。この矛盾により $d = 1$ の場合の議論が完了する。

\medskip\noindent
$d > 1$ と仮定する。$Z \subset \Spec(A^\wedge)$ を
$V(\mathfrak q)$ と異なる既約成分の和とする。
$\mathfrak r_1, \ldots, \mathfrak r_m \subset A^\wedge$ を、
$V(\mathfrak q) \cap Z$ の次元が $> 0$ の既約成分に対応する素イデアルとする。
素回避（Algebra, Lemma \ref{algebra-lemma-silly}）を用いて
$f \in \mathfrak m$, $f \not \in A \cap \mathfrak r_j$ を選ぶ。
このとき $\dim(A/fA) = \dim(A) - 1$ であり、
$V(\mathfrak q, f)$ は次元 $d - 1$ の既約成分をもつ。
従って $A/fA$ は形式鎖状でなく、不変量 $d$ は小さくなった。
帰納法により $A/fA$ は普遍鎖状でなく、従って $A$ も普遍鎖状でない。
\end{proof}

\begin{lemma}
\label{more-algebra-lemma-flat-under-catenary-equidimensional}
$A \to B$ を Noether 局所環の平坦な局所環準同型とする。
% P02-E1233: source lines 32390--32391 contain “and is Spec(B) equidimensional”, with a superfluous “is”.
$B$ が鎖状で、$\Spec(B)$ が等次元であると仮定する。このとき
\begin{enumerate}
\item すべての $\mathfrak p \subset A$ に対して
$\Spec(B/\mathfrak p B)$ は等次元であり、
\item $A$ は鎖状で、$\Spec(A)$ は等次元である。
\end{enumerate}
\end{lemma}

\begin{proof}
$\mathfrak p \subset A$ を素イデアルとし、$\mathfrak q \subset B$ を
$\mathfrak pB$ 上の極小素イデアルとする。$A \to B$ に対する下降定理
（Algebra, Lemma \ref{algebra-lemma-flat-going-down}）により
$\mathfrak q \cap A = \mathfrak p$ である。
従って $A_\mathfrak p \to B_\mathfrak q$ は、次元 $0$ の
特殊ファイバーをもつ平坦な局所環準同型であり、
$$
\dim(A_\mathfrak p) = \dim(B_\mathfrak q) = \dim(B) - \dim(B/\mathfrak q)
$$
を得る（Algebra, Lemma
\ref{algebra-lemma-dimension-base-fibre-equals-total}）。
% P02-E1234: source line 32410 is a sentence fragment (“The second equality because ...”).
第 2 の等号は、$\Spec(B)$ が等次元であり $B$ が鎖状であることによる。
従って $\dim(B/\mathfrak q)$ は $\mathfrak q$ の選び方によらず、
$\Spec(B/\mathfrak p B)$ は次元
$\dim(B) - \dim(A_\mathfrak p)$ の等次元スキームである。一方、平坦性により
（上で引用した補題を参照）、
$\dim(B/\mathfrak p B) = \dim(A/\mathfrak p) + \dim(B/\mathfrak m_A B)$
および
$\dim(B) = \dim(A) + \dim(B/\mathfrak m_A B)$
が成り立つので、
$$
\dim(A_\mathfrak p) = \dim(A) - \dim(A/\mathfrak p)
$$
を得る。これは $A$ のすべての $\mathfrak p$ に対して成り立つ。
$A$ のすべての極小素イデアルに適用すれば、$A$ が等次元であることが分かる。
$\mathfrak p \subset \mathfrak p'$ が、その間に素イデアルをもたない
真の包含なら、$A/\mathfrak p \to B/\mathfrak p B$ は平坦であり
$\Spec(B/\mathfrak p B)$ は等次元だから、上の議論を
$A/\mathfrak p$ の素イデアル $\mathfrak p'/\mathfrak p$ に適用して
$$
1 = \dim((A/\mathfrak p)_{\mathfrak p'}) =
\dim(A/\mathfrak p) - \dim(A/\mathfrak p')
$$
を得る。従って $\mathfrak p \mapsto \dim(A/\mathfrak p)$ は次元関数であり、
$A$ は鎖状である。
\end{proof}

\begin{lemma}
\label{more-algebra-lemma-formally-catenary}
$A$ を形式鎖状 Noether 局所環とする。
このとき $A$ は普遍鎖状である。
\end{lemma}

\begin{proof}
$\mathfrak p$ を $A$ の極小素イデアルとして、$A$ を
$A/\mathfrak p$ で置き換えてよい（Algebra, Lemma
\ref{algebra-lemma-catenary-check-irreducible}）。
従って $A^\wedge$ のスペクトルが等次元であると仮定してよい。
$A$ 上本質的有限型なすべての局所環が鎖状であることを示せば十分である
（例えば Algebra, Lemma \ref{algebra-lemma-catenary-check-local} を参照）。
従って、$\mathfrak m \subset A[x_1, \ldots, x_n]$ を
$\mathfrak m_A$ 上にある極大イデアルとして、
$A[x_1, \ldots, x_n]_\mathfrak m$ が鎖状であることを示せば十分である。
ここでは Algebra, Lemma
\ref{algebra-lemma-localization-at-closed-point-special-fibre}
（および Algebra, Lemmas \ref{algebra-lemma-quotient-catenary} と
\ref{algebra-lemma-localization-catenary}）を用いた。
$\mathfrak m' \subset A^\wedge[x_1, \ldots, x_n]$ を
$\mathfrak m$ 上にある唯一の極大イデアルとする。このとき
$$
A[x_1, \ldots, x_n]_\mathfrak m \to A^\wedge[x_1, \ldots, x_n]_{\mathfrak m'}
$$
は局所的かつ平坦である（Algebra, Lemma
\ref{algebra-lemma-completion-flat}）。
% P02-E1235: source lines 32462--32463 omit “is” before “catenary”.
従って、Lemma \ref{more-algebra-lemma-flat-under-catenary-equidimensional} により、
右辺の環が鎖状で等次元スペクトルをもつことを示せば十分である。
完備局所環は普遍鎖状だから、この環は鎖状である
（Algebra, Remark
\ref{algebra-remark-Noetherian-complete-local-ring-universally-catenary}）。
$A^\wedge[x_1, \ldots, x_n]_{\mathfrak m'}$ の任意の極小素イデアル
$\mathfrak q$ をとる。このとき $A^\wedge$ のある極小素イデアル
$\mathfrak p$ に対して
$\mathfrak q = \mathfrak p A^\wedge[x_1, \ldots, x_n]_{\mathfrak m'}$
である
% P02-E1236: source lines 32468--32471 explicitly omit the minimal-prime extension detail.
（細部は省略する）。従って
$$
\dim(A^\wedge[x_1, \ldots, x_n]_{\mathfrak m'}/\mathfrak q) =
\dim(A^\wedge/\mathfrak p) + n = \dim(A^\wedge) + n
$$
である。第 1 の等号は Algebra, Lemma
\ref{algebra-lemma-dimension-base-fibre-equals-total} により、
第 2 の等号は $A^\wedge$ のスペクトルが等次元であることによる。
これで証明が完了する。
\end{proof}

\begin{proposition}[Ratliff]
\label{more-algebra-proposition-ratliff}
\begin{reference}
\cite{Ratliff}
\end{reference}
Noether 局所環が普遍鎖状であるための必要十分条件は、
それが形式鎖状であることである。
\end{proposition}

\begin{proof}
Lemmas \ref{more-algebra-lemma-not-formally-catenary} と
\ref{more-algebra-lemma-formally-catenary} を組み合わせればよい。
\end{proof}

\begin{lemma}
\label{more-algebra-lemma-geometrically-normal-fibres-universally-catenary}
\begin{reference}
\cite[Corollary 2.3]{Heinzer-Rotthaus-Wiegand}
\end{reference}
$(A, \mathfrak m)$ を、形式ファイバーが幾何学的正規な
Noether 局所環とする。このとき
\begin{enumerate}
\item $A^h$ は普遍鎖状であり、
\item $A$ が単枝（例えば正規）なら、$A$ は普遍鎖状である。
\end{enumerate}
\end{lemma}

\begin{proof}
Lemma \ref{more-algebra-lemma-geometrically-normal-formal-fibres-number-of-branches}
により $A$ と $A^\wedge$ の枝数は等しいので、Lemma
\ref{more-algebra-lemma-equal-nr-branches-completion} を適用できる。
従って任意の極小素イデアル $\mathfrak q \subset A^h$ に対して、
$A^\wedge/\mathfrak q A^\wedge$ はただ一つの極小素イデアルをもつ。
従って $A^h$ は（定義により）形式鎖状であり、Proposition
\ref{more-algebra-proposition-ratliff} により普遍鎖状である。
$A$ が単枝なら、$A^h$ はただ一つの極小素イデアルをもち、
従って $A^\wedge$ もただ一つの極小素イデアルをもち、
従って $A$ は形式鎖状である。同様にして結論を得る。
\end{proof}






\section{群作用と整閉包}
\label{more-algebra-section-group-actions-integral}

\noindent
この節は、ある意味で Algebra, Section
\ref{algebra-section-going-down-integral-over-normal} の続きである。
同種の内容は Fundamental Groups, Section
\ref{pione-section-group-actions-integral} にも見られる
% P02-E1237: source line 32537 lacks terminal punctuation.

\begin{lemma}
\label{more-algebra-lemma-pol-lifting}
$\varphi : A \to B$ を環の全射とする。位数 $n$ の有限群 $G$ が
$\varphi : A \to B$ に作用しているとする。$b \in B^G$ なら、
$B^G[T]$ において $(T - b)^n$ に写るモニック多項式
$P \in A^G[T]$ が存在する。
\end{lemma}

\begin{proof}
$b$ の持ち上げ $a \in A$ を選び、
$P = \prod_{\sigma \in G} (T - \sigma(a))$ とおけばよい。
\end{proof}

\begin{lemma}
\label{more-algebra-lemma-invariants-modulo}
$R$ を環とし、有限群 $G$ が $R$ に作用しているとする。
$I \subset R$ を、すべての $\sigma \in G$ に対して
$\sigma(I) \subset I$ となるイデアルとする。このとき
$R^G/I^G \subset (R/I)^G$ は整環拡大であり、スペクトル上の
同相写像と剰余体の純非分離拡大を誘導する。
\end{lemma}

\begin{proof}
$I^G = R^G \cap I$ なので、写像が単射であることは明らかである。
Lemma \ref{more-algebra-lemma-pol-lifting} により Algebra, Lemma
\ref{algebra-lemma-universally-bijective} を適用できる。
\end{proof}

\begin{lemma}
\label{more-algebra-lemma-pol-in-ideal}
位数 $n$ の有限群 $G$ が環 $R$ に作用しているとする。
$J \subset R^G$ をイデアルとする。$x \in JR$ に対して
$\prod_{\sigma \in G} (T - \sigma(x)) =
T^n + a_1 T^{n - 1} + \ldots + a_n$、ただし $a_i \in J$、が成り立つ。
\end{lemma}

\begin{proof}
この多項式が確かにモニックで、その係数が $R^G$ に属することに注意する。
$x = f_1 b_1 + \ldots + f_m b_m$、ただし $f_j \in J$ および $b_j \in R$、
と書ける。従って $m$ に関する帰納法により、$f \in J$、$b \in R$、
および $y \in JR$（この $y$ に対して結論が成り立つ）を用いて
$x = y - fb$ と書ける場合を考えればよい。このとき
$$
\prod\nolimits_{\sigma \in G} (T - \sigma(x)) =
\prod\nolimits_{\sigma \in G} (T - \sigma(y) + f\sigma(b)) =
\prod\nolimits_{\sigma \in G} (T - \sigma(y)) +
\sum_{i = 1, \ldots, n} f^i a_i
$$
であり、ここで
$$
a_i =
\sum\nolimits_{S \subset G,\ |S| = i}
\prod\nolimits_{\sigma \in S} \sigma(b)
\prod\nolimits_{\sigma \not \in S} (T - \sigma(y))
$$
である。
% P02-E1238: source line 32595 explicitly omits the invariance computation.
省略した計算により $a_i \in R^G$ であることが分かる
（ヒント：上の式は対称的である）。従って補題の $x$ に対する多項式は、
$J$ を法として $y$ に対する多項式と合同である。これで帰納段階が示された。
\end{proof}

\begin{lemma}
\label{more-algebra-lemma-invariants-modulo-bis}
$R$ を環とし、位数 $n$ の有限群 $G$ が $R$ に作用しているとする。
$J \subset R^G$ をイデアルとする。このとき
% P02-E1239: source line 32605 omits the article in “is ring map”.
$R^G/J \to (R/JR)^G$ は環準同型であり、次を満たす。
\begin{enumerate}
\item $b \in (R/JR)^G$ に対して、$(R/JR)^G[T]$ における像が
$(T - b)^n$ であるモニック多項式 $P \in R^G/J[T]$ が存在する。
\item $a \in \Ker(R^G/J \to (R/JR)^G)$ に対して、
$R^G/J[T]$ において $(T - a)^n = T^n$ である。
\end{enumerate}
特に $R^G/J \to (R/JR)^G$ は整環準同型であり、スペクトル上の
同相写像と剰余体の純非分離拡大を誘導する。
\end{lemma}

\begin{proof}
% P02-E1240: source line 32619 has subject–verb disagreement in “Part (1) follow”.
(1) は $I = JR$ として Lemma \ref{more-algebra-lemma-pol-lifting} を適用すれば従う。
(2) の $a$ をとると、$a$ はある $x \in R^G \cap JR$ の像である。
従って Lemma \ref{more-algebra-lemma-pol-in-ideal} により、
$(T - x)^n = \prod_{\sigma \in G} (T - \sigma(x))$ は
$J$ を法として $T^n$ と合同である。これで (2) が示された。
最後の主張には Algebra, Lemma
\ref{algebra-lemma-universally-bijective} を適用すればよい。
\end{proof}

\begin{remark}
\label{more-algebra-remark-when-iso}
Lemma \ref{more-algebra-lemma-invariants-modulo-bis} において、$n$ が $R$ で可逆なら
$R^G/J \to (R/JR)^G$ は同型である。
\end{remark}

\begin{lemma}
\label{more-algebra-lemma-functor-invariants-tensor}
$R$ を環とし、位数 $n$ の有限群 $G$ が $R$ に作用しているとし、
$A$ を $R^G$-代数とする。
\begin{enumerate}
\item $b \in (A \otimes_{R^G} R)^G$ に対して、
$(A \otimes_{R^G} R)^G[T]$ における像が $(T - b)^n$ である
モニック多項式 $P \in A[T]$ が存在する。
\item $a \in \Ker(A \to (A \otimes_{R^G} R)^G)$ に対して、
$A[T]$ において $(T - a)^n = T^n$ である。
\end{enumerate}
\end{lemma}

\begin{proof}
$E$ を $R^G$ 上の多項式代数として、全射 $E \to A$ を選ぶ。
$E$ は $R^G$-加群として自由なので、
$(E \otimes_{R^G} R)^G = E$ である。
$J = \Ker(E \to A)$ とおく。テンソル積は右完全だから、
$A \otimes_{R^G} R$ は $E \otimes_{R^G} R$ を $J$ の生成する
イデアルで割った商である。従って本補題は Lemma
\ref{more-algebra-lemma-invariants-modulo-bis} の特別な場合である。
\end{proof}

\begin{lemma}
\label{more-algebra-lemma-base-change-invariants}
$R$ を環とし、有限群 $G$ が $R$ に作用しているとする。
$R^G \to A$ を環準同型とする。写像
$$
A \to (A \otimes_{R^G} R)^G
$$
は $R^G \to A$ が平坦なら同型である。一般には、この写像は整であり、
スペクトル上の同相写像と剰余体の純非分離拡大を誘導する。
\end{lemma}

\begin{proof}
第 1 の主張を示すため、完全列
$0 \to R^G \to R \to \bigoplus_{\sigma \in G} R$ を考える。
ここで第 2 の写像は $x$ を $(\sigma(x) - x)_{\sigma \in G}$ に送る。
$R^G \to A$ が平坦なら、$A$ とテンソルをとった列も完全である。
% P02-E1241: source line 32672 takes G-invariants inside an expression already marked invariant.
従って $A$ は $(A \otimes_{R^G} R)^G$ における $G$-不変部分である。

\medskip\noindent
第 2 の主張は Lemma \ref{more-algebra-lemma-functor-invariants-tensor} と
Algebra, Lemma \ref{algebra-lemma-universally-bijective} から従う。
\end{proof}

\begin{lemma}
\label{more-algebra-lemma-one-orbit}
有限群 $G$ が環 $R$ に作用しているとする。$R^G$ の同じ素イデアル上にある
任意の二つの素イデアル $\mathfrak q, \mathfrak q' \subset R$ に対して、
$\sigma(\mathfrak q) = \mathfrak q'$ となる $\sigma \in G$ が存在する。
\end{lemma}

\begin{proof}
すべての $x \in R$ は $R^G[T]$ のモニック多項式
$\prod_{\sigma \in G}(T - \sigma(x))$ の根なので、
$R^G \subset R$ は整環拡大である。従って、与えられた素イデアル
$\mathfrak p$ 上にある素イデアルの間には包含関係がない
（Algebra, Lemma \ref{algebra-lemma-integral-no-inclusion}）。
補題が誤りなら、Algebra, Lemma \ref{algebra-lemma-silly} により、
すべての $\sigma \in G$ に対して
$x \in \mathfrak q'$、$x \not \in \sigma(\mathfrak q)$ となる元を選べる。
このとき $y = \prod_{\sigma \in G} \sigma(x)$ は $R^G$ に属し、
$\mathfrak p = R^G \cap \mathfrak q'$ にも属する。一方、すべての $\sigma$ に対する
$x \not \in \sigma(\mathfrak q)$ は、すべての $\sigma$ に対する
$\sigma(x) \not \in \mathfrak q$ を意味する。$\mathfrak q$ は素イデアルだから
$y \not \in \mathfrak q$ である。しかし $y \in \mathfrak p \subset \mathfrak q$
なので、これは不可能である。
\end{proof}

\begin{lemma}
\label{more-algebra-lemma-one-orbit-geometric}
有限群 $G$ が環 $R$ に作用しているとする。
$\mathfrak q \subset R$ を $\mathfrak p \subset R^G$ 上にある
素イデアルとする。このとき $\kappa(\mathfrak q)/\kappa(\mathfrak p)$ は
代数的正規拡大であり、写像
$$
D = \{\sigma \in G \mid \sigma(\mathfrak q) = \mathfrak q\}
\longrightarrow
\text{Aut}(\kappa(\mathfrak q)/\kappa(\mathfrak p))
$$
は全射である\footnote{記号 $\text{Gal}$ は Galois 拡大の場合にのみ
用いるという約束を思い出されたい。}。
\end{lemma}

\begin{proof}
$A = (R^G)_\mathfrak p$、$B = A \otimes_{R^G} R$ とおくと、
局所化は平坦なので Lemma \ref{more-algebra-lemma-base-change-invariants} により
$A = B^G$ である。$\mathfrak pA$ と $\mathfrak qB$ は素イデアルであり、
$D$ は $\mathfrak qB$ の安定化群で、
$\kappa(\mathfrak p) = \kappa(\mathfrak pA)$ および
$\kappa(\mathfrak q) = \kappa(\mathfrak qB)$ であることに注意する。
従って $R$ を $B$ で置き換え、$\mathfrak p$ が極大イデアルであると
仮定してよい。$R^G \subset R$ は整環拡大だから、$R$ の極大イデアルは
$\mathfrak p$ 上にある素イデアルとちょうど一致する
（Algebra, Lemmas \ref{algebra-lemma-integral-no-inclusion} および
\ref{algebra-lemma-integral-going-up} から従う）。
Lemma \ref{more-algebra-lemma-one-orbit} により、それらは有限個
$\mathfrak q = \mathfrak q_1, \mathfrak q_2, \ldots, \mathfrak q_m$ であり、
$G$ の一つの軌道をなす。中国剰余定理（Algebra, Lemma
\ref{algebra-lemma-chinese-remainder}）により、写像
$R \to \prod_{j = 1, \ldots, m} R/\mathfrak q_j$ は全射である。

\medskip\noindent
まず拡大が正規であることを示す。元
$\alpha \in \kappa(\mathfrak q)$ をとる。$\alpha$ の
$\kappa(\mathfrak p)$ 上の最小多項式 $P$ が完全分解することを
示さなければならない。上の全射により、$\kappa(\mathfrak q)$ で
$\alpha$ に写る元 $a \in \mathfrak q_2 \cap \ldots \cap \mathfrak q_m$
を選べる。$R^G[T]$ の多項式
$Q = \prod_{\sigma \in G} (T - \sigma(a))$ を考える。
$Q$ の $R[T]$ における像は一次因子に完全分解するので、その
$\kappa(\mathfrak q)[T]$ における像も同様である。
$P$ は $Q$ の $\kappa(\mathfrak p)[T]$ における像を割り切るから、
$P$ は望みどおり $\kappa(\mathfrak q)$ 上で一次因子に完全分解する。

\medskip\noindent
$\kappa(\mathfrak q)/\kappa(\mathfrak p)$ は正規なので、Fields, Lemma
\ref{fields-lemma-normal-case} により
$\kappa(\mathfrak q) = \kappa_1 \otimes_{\kappa(\mathfrak p)} \kappa_2$、
ただし $\kappa_1/\kappa(\mathfrak p)$ は純非分離、
$\kappa_2/\kappa(\mathfrak p)$ は Galois、であると仮定してよい。
それが有限なら $\kappa(\mathfrak p)$ 上 $\kappa_2$ を生成する元
$\alpha \in \kappa_2$ を選び、無限なら矛盾を得るため次数 $> |G|$ の
部分体を選ぶ。これは Fields, Lemma
\ref{fields-lemma-primitive-element} により可能である。
前段落と同様に $a$、$P$、$Q$ を選ぶ。
$\alpha' \in \kappa_2$ を $\alpha$ の $\kappa(\mathfrak p)$ 上の
Galois 共役とする。
% P02-E1242: source lines 32765--32767 say P divides the image of P; the preceding argument requires Q.
$P$ が $\kappa(\mathfrak p)[T]$ における $P$ の像を割り切ることから、
$\sigma(a)$ が $\alpha'$ に写るような $\sigma \in G$ が存在する。
$a$ は他の極大イデアルで消えるように選んだので、これは
$\sigma \in D$ を意味し、$\sigma$ の
$\text{Aut}(\kappa(\mathfrak q)/\kappa(\mathfrak p))$ における像は
$\alpha$ を $\alpha'$ に写す。
% P02-E1243: source lines 32771--32773 form an incomplete concluding sentence.
従って全射性が得られ、$\alpha$ の $\kappa(\mathfrak p)$ 上の次数が
$> |G|$ の場合には、求めていた矛盾が得られる。
\end{proof}

\begin{lemma}
\label{more-algebra-lemma-one-orbit-geometric-galois}
$A$ を分数体 $K$ をもつ正規整域とする。
$L/K$ を（無限でもよい）Galois 拡大とする。
$G = \text{Gal}(L/K)$ とし、$B$ を $L$ における $A$ の整閉包とする。
\begin{enumerate}
\item $A$ の同じ素イデアル上にある任意の二つの素イデアル
$\mathfrak q, \mathfrak q' \subset B$ に対して、
$\sigma(\mathfrak q) = \mathfrak q'$ となる $\sigma \in G$ が存在する。
\item $\mathfrak q \subset B$ を $\mathfrak p \subset A$ 上にある
素イデアルとする。このとき $\kappa(\mathfrak q)/\kappa(\mathfrak p)$ は
代数的正規拡大であり、写像
$$
D = \{\sigma \in G \mid \sigma(\mathfrak q) = \mathfrak q\}
\longrightarrow
\text{Aut}(\kappa(\mathfrak q)/\kappa(\mathfrak p))
$$
は全射である。
\end{enumerate}
\end{lemma}

\begin{proof}
(1) の証明。$K \subset M \subset L$ が $M/K$ を Galois 拡大とする
部分体で、$\sigma \in \text{Gal}(M/K)$ が
$\sigma(\mathfrak q \cap M) = \mathfrak q' \cap M$ を満たすような
組 $(M, \sigma)$ を考える。
$M \subset M'$ かつ $\sigma'|_M = \sigma$ であることを、
$(M', \sigma') \geq (M, \sigma)$ と定義する。
$A$ は正規整域なので $A = K \cap B$ であり、
$(K, \text{id}_K)$ はこのような組である。
これらの組の集合は Zorn の補題の仮定を満たすので、極大な組
$(M, \sigma)$ が存在する。$M \not = L$ なら、
$M'/M$ が非自明な有限拡大で $M'/K$ が Galois となるような
$M \subset M' \subset L$ を選べる（Fields, Lemma
\ref{fields-lemma-normal-closure-inside-normal}）。
$M$ への制限が $\sigma$ である $\sigma' \in \text{Gal}(M'/K)$ を選ぶ
（Fields, Lemma \ref{fields-lemma-galois-infinite}）。
このとき素イデアル $\sigma'(\mathfrak q \cap M')$ と
$\mathfrak q' \cap M'$ は $B \cap M$ の同じ素イデアルに制限される。
$B \cap M = (B \cap M')^{\text{Gal}(M'/M)}$ なので、Lemma
\ref{more-algebra-lemma-one-orbit} により
$\tau(\sigma'(\mathfrak q \cap M')) = \mathfrak q' \cap M'$ となる
$\tau \in \text{Gal}(M'/M)$ を選べる。従って
$(M', \tau \circ \sigma') > (M, \sigma)$ となり、
$(M, \sigma)$ の極大性に矛盾する。

\medskip\noindent
(2) も (1) とまったく同様に証明されるが、詳細を記す。
$\overline{\sigma} \in
\text{Aut}(\kappa(\mathfrak q)/\kappa(\mathfrak p))$ をとる。
$K \subset M \subset L$ が $M/K$ を Galois 拡大とする部分体で、
$\sigma \in \text{Gal}(M/K)$ が
$\sigma(\mathfrak q \cap M) = \mathfrak q \cap M$ を満たし、図式
$$
\xymatrix{
\kappa(\mathfrak q \cap M) \ar[r] \ar[d]_\sigma &
\kappa(\mathfrak q) \ar[d]_{\overline{\sigma}} \\
\kappa(\mathfrak q \cap M) \ar[r] & \kappa(\mathfrak q)
}
$$
が可換となる組 $(M, \sigma)$ を考える。
$M \subset M'$ かつ $\sigma'|_M = \sigma$ であることを、
$(M', \sigma') \geq (M, \sigma)$ と定義する。
上と同様に $(K, \text{id}_K)$ はこのような組である。
これらの組の集合は Zorn の補題の仮定を満たすので、極大な組
$(M, \sigma)$ が存在する。$M \not = L$ なら、$M'/M$ が有限で
$M'/K$ が Galois となるような $M \subset M' \subset L$ を選べる
（Fields, Lemma \ref{fields-lemma-normal-closure-inside-normal}）。
$M$ への制限が $\sigma$ である $\sigma' \in \text{Gal}(M'/K)$ を選ぶ
（Fields, Lemma \ref{fields-lemma-galois-infinite}）。
このとき素イデアル $\sigma'(\mathfrak q \cap M')$ と
$\mathfrak q \cap M'$ は $B \cap M$ の同じ素イデアルに制限される。
第 1 段落と同様に $\sigma'$ の選び方を調整し、
$\sigma'(\mathfrak q \cap M') = \mathfrak q \cap M'$ と仮定してよい。
すると $\sigma'$ と $\overline{\sigma}$ は写像
$\kappa(\mathfrak q \cap M') \to \kappa(\mathfrak q)$ を定め、
$\kappa(\mathfrak q \cap M)$ 上で一致する。
$B \cap M = (B \cap M')^{\text{Gal}(M'/M)}$ なので、Lemma
\ref{more-algebra-lemma-one-orbit-geometric} により、
$\tau(\mathfrak q \cap M') = \mathfrak q \cap M'$ を満たし、
% P02-E1244: source line 32856 composes tau with sigma, although sigma prime is the map on M prime used here and in line 32864.
$\tau \circ \sigma$ と $\overline{\sigma}$ が
$\kappa(\mathfrak q \cap M')$ 上で同じ写像を誘導するような
$\tau \in \text{Gal}(M'/M)$ を選べる。ここには少し細部がある。
まず補題により
$\kappa(\mathfrak q \cap M')/\kappa(\mathfrak q \cap M)$ が
正規であることが保証される。
従って二つの写像の差はこの拡大の自己同型である
（Fields, Lemma \ref{fields-lemma-normal-embeddings-differ-by-aut}）。
そこでこの自己同型に補題を適用して $\tau$ を得る。
従って $(M', \tau \circ \sigma') > (M, \sigma)$ となり、
$(M, \sigma)$ の極大性に矛盾する。
\end{proof}

\begin{lemma}
\label{more-algebra-lemma-one-orbit-geometric-galois-compare}
$A$ を分数体 $K$ をもつ正規整域とする。
$M/L/K$ を $K$ の（無限でもよい）Galois 拡大の塔とする。
$H = \text{Gal}(M/K)$、$G = \text{Gal}(L/K)$ とし、
% P02-E1245: source line 32873 uses singular “the integral closure” for the two closures C and B.
$C$ と $B$ をそれぞれ $M$ と $L$ における $A$ の整閉包とする。
$\mathfrak r \subset C$ および $\mathfrak q = B \cap \mathfrak r$ とする。
さらに
$D_\mathfrak r = \{\tau \in H \mid \tau(\mathfrak r) = \mathfrak r\}$
および
$I_\mathfrak r = \{\tau \in D_\mathfrak r \mid
\tau \bmod \mathfrak r = \text{id}_{\kappa(\mathfrak r)}\}$
とおき、$D_\mathfrak q$ と $I_\mathfrak q$ も同様に定める。
写像 $H \to G$ の下で、誘導される写像
$D_\mathfrak r \to D_\mathfrak q$ および
$I_\mathfrak r \to I_\mathfrak q$ は全射である。
\end{lemma}

\begin{proof}
$\sigma \in D_\mathfrak q$ とする。$\sigma$ に写る $\tau \in H$ を選ぶ。
これは Fields, Lemma \ref{fields-lemma-galois-infinite} により可能である。
$\tau(\mathfrak r)$ と $\mathfrak r$ はともに $\mathfrak q$ 上にある。
従って Lemma \ref{more-algebra-lemma-one-orbit-geometric-galois} により
$\sigma'(\tau(\mathfrak r)) = \mathfrak r$ となる
$\sigma' \in \text{Gal}(M/L)$ が存在する。従って
$\sigma'\tau \in D_\mathfrak r$ は $\sigma$ に写る。
惰性群の場合も、自己同型群への全射性を用いてまったく同様に証明される。
\end{proof}








\section{離散付値環の拡大}
\label{more-algebra-section-discrete-valuation-rings}

\noindent
この節と続く数節では、以下の定義を用いる。

\begin{definition}
\label{more-algebra-definition-extension-discrete-valuation-rings}
$A$ と $B$ が離散付値環で、$A \to B$ が単射かつ局所的であるとき、
$A \to B$ または $A \subset B$ を {\it 離散付値環の拡大} という。
特に $\pi_A$ と $\pi_B$ がそれぞれ $A$ と $B$ の素元なら、
ある $e \geq 1$ と $B$ の単元 $u$ に対して
$\pi_A = u \pi_B^e$ である。整数 $e$ は素元の選び方によらない。
実際、これは
$$
\mathfrak m_A B = \mathfrak m_B^e
$$
を満たす唯一の整数 $\geq 1$ でもある。整数 $e$ を $A$ 上の $B$ の
{\it 分岐指数} という。$e = 1$ なら、$B$ は $A$ 上
{\it 弱不分岐} であるという。剰余体の拡大
$\kappa_A = A/\mathfrak m_A \subset \kappa_B = B/\mathfrak m_B$
が有限なら、$f = [\kappa_B : \kappa_A]$ とおき、これを拡大
$A \subset B$ の {\it 剰余次数} または {\it 残留次数} という。
\end{definition}

\noindent
分数体の拡大が有限であることは仮定していない点に注意する。

\begin{lemma}
\label{more-algebra-lemma-inequality}
$A \subset B$ を、分数体 $K \subset L$ をもつ離散付値環の拡大とする。
拡大 $L/K$ が有限なら剰余体拡大も有限であり、
$ef \leq [L : K]$ である。
\end{lemma}

\begin{proof}
剰余体拡大の有限性は Algebra, Lemma
\ref{algebra-lemma-finite-extension-residue-fields-dimension-1} である。
不等式は Algebra, Lemmas \ref{algebra-lemma-finite-length} および
\ref{algebra-lemma-pushdown-module} から従う。
\end{proof}

\begin{lemma}
\label{more-algebra-lemma-multiplicative-e-f}
$A \subset B \subset C$ を離散付値環の拡大とする。
$B/A$ と $C/B$ の分岐指数の積は $C/A$ の分岐指数である。式では
$e_{C/A} = e_{B/A} e_{C/B}$ となる。剰余次数が有限な場合には、
それについても同様である。
\end{lemma}

\begin{proof}
定義と Fields, Lemma \ref{fields-lemma-multiplicativity-degrees} から直ちに従う。
\end{proof}

\begin{lemma}
\label{more-algebra-lemma-ramification-index-a-power-of-p}
$A \subset B$ を、体の拡大 $K \subset L$ を誘導する
離散付値環の拡大とする。$K$ の標数が $p > 0$ で、$L$ が $K$ 上
純非分離なら、分岐指数 $e$ は $p$ の冪である。
\end{lemma}

\begin{proof}
ある $u \in B^*$ に対して $\pi_A = u \pi_B^e$ と書く。
一方、ある $p$-冪 $q$ に対して $\pi_B^q \in K$ である。
ある $v \in A^*$ と $k \in \mathbf{Z}$ によって
$\pi_B^q = v \pi_A^k$ と書く。このとき
$\pi_A^q = u^q \pi_B^{qe} = u^q v^e \pi_A^{ke}$ である。
$B$ で付値をとると $ke = q$ を得る。
\end{proof}

\noindent
次の補題では、離散付値環の拡大 $A \subset B$ が「不分岐」、すなわち
分岐指数が $1$ で剰余体拡大が分離的（代数的でなくてもよい）であるとは
何を意味するかを論じる。ただし、不分岐環準同型という概念が既にあるので、
「不分岐」という語そのものは使えない。Algebra, Section
\ref{algebra-section-unramified} を参照。

\begin{lemma}
\label{more-algebra-lemma-extension-dvrs-formally-smooth}
$A \subset B$ を離散付値環の拡大とする。次は同値である。
\begin{enumerate}
\item $A \to B$ は $\mathfrak m_B$-進位相について形式的滑らかである。
\item $A \to B$ は弱不分岐であり、$\kappa_B/\kappa_A$ は分離拡大である。
\end{enumerate}
\end{lemma}

\begin{proof}
Proposition \ref{more-algebra-proposition-fs-flat-fibre-fs} と Algebra, Proposition
\ref{algebra-proposition-characterize-separable-field-extensions} から従う。
\end{proof}

\begin{remark}
\label{more-algebra-remark-finite-separable-extension}
$A$ を分数体 $K$ をもつ離散付値環とし、$L/K$ を有限分離体拡大とする。
$B \subset L$ を $L$ における $A$ の整閉包とする。図式は
$$
\xymatrix{
B \ar[r] & L \\
A \ar[u] \ar[r] & K \ar[u]
}
$$
である。Algebra, Lemma
\ref{algebra-lemma-Noetherian-normal-domain-finite-separable-extension}
により環拡大 $A \subset B$ は有限なので、$B$ は Noether 環である。
Algebra, Lemma \ref{algebra-lemma-integral-sub-dim-equal} により
$B$ の次元は $1$ なので、Algebra, Lemma
\ref{algebra-lemma-characterize-Dedekind} により $B$ は Dedekind 整域である。
$\mathfrak m_1, \ldots, \mathfrak m_n$ を $B$ の極大イデアル
（すなわち $\mathfrak m_A$ 上にある素イデアル）とする。
離散付値環の拡大
$$
A \subset B_{\mathfrak m_i}
$$
が得られ、従って分岐指数 $e_i$ と剰余次数 $f_i$ が得られる。
$A$ の素元に Algebra, Lemma
\ref{algebra-lemma-finite-extension-dim-1} を適用すると
$$
[L : K] = \sum\nolimits_{i = 1, \ldots, n} e_i f_i
$$
を得る。$A$ が Hensel 環なら $n = 1$ である
（Algebra, Lemma \ref{algebra-lemma-finite-over-henselian}）。
例えば $A$ が完備ならそうである。
\end{remark}

\begin{definition}
\label{more-algebra-definition-types-of-extensions}
$A$ を分数体 $K$ をもつ離散付値環とし、$L/K$ を有限分離拡大とする。
$B$ および $\mathfrak m_i$、$i = 1, \ldots, n$ を Remark
\ref{more-algebra-remark-finite-separable-extension} のとおりとする。拡大 $L/K$ が
\begin{enumerate}
\item すべての $i$ に対して $e_i = 1$ であり、拡大
$\kappa(\mathfrak m_i)/\kappa_A$ が分離的であるとき、
$A$ に関して {\it 不分岐} であるという。
% P02-E1246: source lines 33050--33054 have ambiguous coordination between the characteristic-zero and positive-characteristic tame cases.
\item $\kappa_A$ の標数が $0$ であるか、または $\kappa_A$ の標数が
$p > 0$ であって体拡大 $\kappa(\mathfrak m_i)/\kappa_A$ が分離的かつ
分岐指数 $e_i$ が $p$ と互いに素であるとき、$A$ に関して
{\it 順分岐} であるという。
\item $n = 1$ で剰余体拡大 $\kappa(\mathfrak m_1)/\kappa_A$ が
自明であるとき、$A$ に関して {\it 完全分岐} であるという。
\end{enumerate}
離散付値環 $A$ が文脈から明らかな場合、略して $L/K$ は不分岐、
完全分岐、または順分岐であるともいう。
\end{definition}

\noindent
不分岐拡大について、次の基本補題がある。

\begin{lemma}
\label{more-algebra-lemma-permanence-unramified}
$A$ を分数体 $K$ をもつ離散付値環とする。
\begin{enumerate}
\item $M/L/K$ が有限分離拡大で、$M$ が $A$ に関して不分岐なら、
$L$ も $A$ に関して不分岐である。
\item $L/K$ が $A$ に関して不分岐な有限分離拡大なら、$L$ を含み
$A$ に関して不分岐な Galois 拡大 $M/K$ が存在する。
% P02-E1247: source line 33077 duplicates the article in “a a finite”.
\item $L_1/K$ と $L_2/K$ が $A$ に関して不分岐な有限分離拡大なら、
$L_1$ と $L_2$ を含み $A$ に関して不分岐な有限分離拡大 $L/K$ が存在する。
\end{enumerate}
\end{lemma}

\begin{proof}
以下では Remark \ref{more-algebra-remark-finite-separable-extension} の議論の結果を
断りなく用いる。

\medskip\noindent
(1) の証明。$C/B/A$ を $M/L/K$ における $A$ の整閉包とする。
$C$ は $B$ の有限環拡大なので、$\Spec(C) \to \Spec(B)$ は全射である。
% P02-E1248: source line 33091 has “for ever maximal ideal” instead of “every”.
従って任意の極大イデアル $\mathfrak m \subset B$ に対して、$\mathfrak m$ 上にある
極大イデアル $\mathfrak m' \subset C$ が存在する。分岐指数の乗法性
（Lemma \ref{more-algebra-lemma-multiplicative-e-f}）と仮定により、
$A$ 上の $B_\mathfrak m$ の分岐指数は $1$ である。
$\kappa(\mathfrak m')/\kappa_A$ は有限分離拡大なので、
$\kappa(\mathfrak m)/\kappa_A$ も同様である。

\medskip\noindent
(2) の証明。$M$ を $K$ 上の $L$ の正規閉包とする
（Fields, Definition \ref{fields-definition-normal-closure}）。
Fields, Lemma \ref{fields-lemma-normal-closure-galois} により
$M/K$ は Galois である。一方、Fields, Lemma
\ref{fields-lemma-normal-closure-tensor-product} により、$K$-代数の全射
$$
L \otimes_K \ldots \otimes_K L \longrightarrow M
$$
がある。Remark \ref{more-algebra-remark-finite-separable-extension} のとおり、
$B$ を $L$ における $A$ の整閉包とする。$L$ が $A$ に関して
不分岐であるという条件は、$A \to B$ が \'etale 環準同型であることと
ちょうど同値である（Algebra, Lemma
\ref{algebra-lemma-characterize-etale}）。\'etale 環準同型の安定性により
$$
B \otimes_A \ldots \otimes_A B
$$
は $A$ 上 \'etale である（Algebra, Lemma
\ref{algebra-lemma-etale}）。従って Lemma
\ref{more-algebra-lemma-Dedekind-etale-extension} により、表示した環は
Dedekind 整域の積である。ゆえに $M$ は $A$ 上有限 \'etale な
Dedekind 整域の分数体である。これは $M$ が $A$ に関して
不分岐であることを意味し、望む結論を得る。

\medskip\noindent
(3) の証明。
% P02-E1249: source line 33130 omits “in L_i” and uses a singular closure for the indexed family.
$B_i \subset L_i$ を、それぞれ対応する体における $A$ の整閉包とする。
上と同様に論じて、$B_1 \otimes_A B_2$ が $A$ 上有限 \'etale であることを示す。
% P02-E1250: source lines 33130--33132 explicitly omit the proof details for part (3).
詳細は省略する。
\end{proof}

\begin{lemma}
\label{more-algebra-lemma-composition-unramified}
$A$ を分数体 $K$ をもつ離散付値環とし、$M/L/K$ を有限分離拡大とする。
$B$ を $L$ における $A$ の整閉包とする。$L/K$ が $A$ に関して
不分岐であり、$B$ のすべての極大イデアル $\mathfrak m$ に対して
$M/L$ が $B_\mathfrak m$ に関して不分岐なら、$M/K$ は $A$ に関して
不分岐である。
\end{lemma}

\begin{proof}
$C$ を $M$ における $A$ の整閉包とする。$C$ の任意の極大イデアル
$\mathfrak m'$ は $B$ のある極大イデアル $\mathfrak m$ 上にある。
このとき補題は、分岐指数の乗法性（Lemma
\ref{more-algebra-lemma-multiplicative-e-f}）と、有限体拡大の塔
$\kappa(\mathfrak m')/\kappa(\mathfrak m)/\kappa_A$ があることから従う。
\end{proof}















\section{Galois 拡大と分岐}
\label{more-algebra-section-ramification}

\noindent
Galois 拡大の場合には、Section
\ref{more-algebra-section-discrete-valuation-rings} の議論をさらに詳しくできる。

\begin{lemma}
\label{more-algebra-lemma-galois}
$A$ を分数体 $K$ をもつ離散付値環とし、$L/K$ を Galois 群 $G$ をもつ
有限 Galois 拡大とする。このとき $G$ は Remark
\ref{more-algebra-remark-finite-separable-extension} の環 $B$ に作用し、
$B$ の極大イデアルの集合に推移的に作用する。
\end{lemma}

\begin{proof}
$A$ は $K$ の中で整閉であり $K = L^G$ なので $A = B^G$ である。
従って本補題は Lemma \ref{more-algebra-lemma-one-orbit} の特別な場合である。
\end{proof}

\begin{lemma}
\label{more-algebra-lemma-galois-conclusion}
$A$ を分数体 $K$ をもつ離散付値環とし、$L/K$ を有限 Galois 拡大とする。
すべての $i$ に対して $e_i = e$ および $f_i = f$ となる
$e \geq 1$ と $f \geq 1$ が存在する
（記号は Remark \ref{more-algebra-remark-finite-separable-extension} のとおり）。
特に $[L : K] = n e f$ である。
\end{lemma}

\begin{proof}
Lemma \ref{more-algebra-lemma-galois} と定義から直ちに従う。
\end{proof}

\begin{definition}
\label{more-algebra-definition-decomposition-inertia}
$A$ を分数体 $K$ をもつ離散付値環とし、$L/K$ を Galois 群 $G$ をもつ
有限 Galois 拡大とする。$B$ を $L$ における $A$ の整閉包とし、
$\mathfrak m \subset B$ を極大イデアルとする。
\begin{enumerate}
\item $\mathfrak m$ の {\it 分解群} とは、部分群
$D = \{\sigma \in G \mid \sigma(\mathfrak m) = \mathfrak m\}$ である。
\item $\mathfrak m$ の {\it 惰性群} とは、写像
$D \to \text{Aut}(\kappa(\mathfrak m)/\kappa_A)$ の核 $I$ である。
\end{enumerate}
\end{definition}

\noindent
体 $\kappa(\mathfrak m)$ は $\kappa_A$ 上非分離であることもある。
特に体拡大 $\kappa(\mathfrak m)/\kappa_A$ は Galois とは限らない。
$\kappa_A$ が完全なら Galois である。

\begin{lemma}
\label{more-algebra-lemma-galois-galois}
$A$ を分数体 $K$ と剰余体 $\kappa$ をもつ離散付値環とする。
$L/K$ を Galois 群 $G$ をもつ有限 Galois 拡大とする。
$B$ を $L$ における $A$ の整閉包とし、$\mathfrak m$ を $B$ の
極大イデアルとする。このとき
\begin{enumerate}
\item 体拡大 $\kappa(\mathfrak m)/\kappa$ は正規であり、
\item $D \to \text{Aut}(\kappa(\mathfrak m)/\kappa)$ は全射である。
\end{enumerate}
ある（同値に、すべての）極大イデアル $\mathfrak m \subset B$ に対して
体拡大 $\kappa(\mathfrak m)/\kappa$ が分離的なら、
\begin{enumerate}
\item[(3)] $\kappa(\mathfrak m)/\kappa$ は Galois であり、
\item[(4)] $D \to \text{Gal}(\kappa(\mathfrak m)/\kappa)$ は全射である。
\end{enumerate}
ここで $D \subset G$ は $\mathfrak m$ の分解群である。
\end{lemma}

\begin{proof}
$A$ は $K$ の中で整閉であり $K = L^G$ なので $A = B^G$ である。
従って (1) と (2) は Lemma \ref{more-algebra-lemma-one-orbit-geometric} から従う。
補題の「同値に、すべての」という部分は Lemma \ref{more-algebra-lemma-galois} から従う。
$\kappa(\mathfrak m)/\kappa$ が分離的であると仮定する。
すると (3) と (4) は (1) と (2) から直ちに従う。
\end{proof}

\begin{lemma}
\label{more-algebra-lemma-galois-inertia}
$A$ を分数体 $K$ をもつ離散付値環とし、$L/K$ を Galois 群 $G$ をもつ
有限 Galois 拡大とする。$B$ を $L$ における $A$ の整閉包とし、
$\mathfrak m \subset B$ を極大イデアルとする。$\mathfrak m$ の惰性群 $I$ は
標準完全列
$$
1 \to P \to I \to I_t \to 1
$$
に入り、次を満たす。
\begin{enumerate}
\item $D$ を分解群とすると、
$$
P = \{\sigma \in I \mid \sigma|_{\mathfrak m/\mathfrak m^2} =
\text{id}_{\mathfrak m/\mathfrak m^2}\}
=
\{\sigma \in D \mid \sigma \text{ acts trivially on }\text{Gr}_\mathfrak m(B)\}
$$
である。
\item $P$ は $D$ の正規部分群である。
\item $\kappa_A$ の標数が $p > 0$ なら $P$ は $p$-群であり、
$\kappa_A$ の標数が零なら $P = \{1\}$ である。
\item $I_t$ は、整数 $e$ の $p$ と互いに素な部分を位数とする巡回群である。
\item 標準同型 $\theta : I_t \to \mu_e(\kappa(\mathfrak m))$ が存在する。
\item
$\kappa(m)$ が $A$ の剰余体上分離的なら、
% P02-E1251: source line 33278 writes kappa(m) instead of kappa(mathfrak m).
$P =
\{\sigma \in D \mid \sigma|_{B/\mathfrak m^2} = \text{id}_{B/\mathfrak m^2}\}$
である。
\end{enumerate}
ここで $e$ は Lemma \ref{more-algebra-lemma-galois-conclusion} の整数である。
\end{lemma}

\begin{proof}
Lemma \ref{more-algebra-lemma-galois-conclusion} により
$|G| = [L : K] = nef$ であることを思い出す。
$G$ は $B$ の極大イデアルの集合
$\{\mathfrak m_1, \ldots, \mathfrak m_n\}$ に推移的に作用し
（Lemma \ref{more-algebra-lemma-galois}）、$D$ は一つの元の安定化群なので
$|D| = ef$ である。Lemma \ref{more-algebra-lemma-galois-galois} により
$$
ef = |D| = |I| \cdot |\text{Aut}(\kappa(\mathfrak m)/\kappa)|
$$
である。ここで $\kappa$ は $A$ の剰余体である。
$\kappa(\mathfrak m)$ は $\kappa$ 上正規なので、
$\text{Aut}(\kappa(\mathfrak m)/\kappa)$ の位数と $f$ との差は
$p$ の冪である（Fields, Lemma
\ref{fields-lemma-normal-and-automorphisms} と Fields, Definition
\ref{fields-definition-insep-degree} に続く議論を参照）。
従って $|I|$ の $p$ と互いに素な部分は、$e$ の $p$ と互いに素な部分に等しい。

\medskip\noindent
$C = B_\mathfrak m$ とおく。このとき $I$ は $A$ 上 $C$ に作用し、
$C$ の剰余体には自明に作用する。$\pi_A \in A$ と $\pi_C \in C$ を
素元とし、$C$ のある単元 $u$ によって $\pi_A = u \pi_C^e$ と書く。
$\sigma \in I$ に対して、$C$ のある単元 $\theta_\sigma$ によって
$\sigma(\pi_C) = \theta_\sigma \pi_C$ と書く。このとき
$$
\pi_A = \sigma(\pi_A) = \sigma(u) (\theta_\sigma \pi_C)^e
= \sigma(u) \theta_\sigma^e \pi_C^e = \frac{\sigma(u)}{u} \theta_\sigma^e \pi_A
$$
である。$\sigma \in I$ なので
$\sigma(u) \equiv u \bmod \mathfrak m_C$ である。従って
$\theta_\sigma$ の $\kappa_C = \kappa(\mathfrak m)$ における像
$\overline{\theta}_\sigma$ は $e$ 次の 1 の根である。写像
\begin{equation}
\label{more-algebra-equation-inertia-character}
\theta : I \longrightarrow \mu_e(\kappa(\mathfrak m)),\quad
\sigma \mapsto \overline{\theta}_\sigma
\end{equation}
を得る。$\theta$ は群準同型であり、素元 $\pi_C$ の選び方によらないと主張する。
実際、$\tau$ を $I$ の別の元とすると、
$\tau(\sigma(\pi_C)) = \tau(\theta_\sigma \pi_C) =
\tau(\theta_\sigma) \theta_\tau \pi_C$ であるから、
$\theta_{\tau \sigma} = \tau(\theta_\sigma) \theta_\tau$ である。
$\tau \in I$ なので
$\overline{\theta}_{\tau \sigma} =
\overline{\theta}_\sigma \overline{\theta}_\tau$ を得る。
$\pi'_C$ を別の素元とすると、$C$ のある単元 $w$ によって
$\pi'_C = w \pi_C$ と書け、
$\sigma(\pi'_C) = w^{-1}\sigma(w)\theta_\sigma \pi'_C$ である。
従って $\theta'_\sigma = w^{-1}\sigma(w)\theta_\sigma$ となり、
$\theta'_\sigma$ と $\theta_\sigma$ は先ほどと同じく剰余体の同じ元に写る。

\medskip\noindent
$\kappa(\mathfrak m)$ の標数は $p$ なので、群
$\mu_e(\kappa(\mathfrak m))$ は、$e$ の $p$ と互いに素な部分以下の
位数をもつ巡回群である（Fields, Section
\ref{fields-section-roots-of-1} を参照）。

\medskip\noindent
$P = \Ker(\theta)$ とおく。構成により $P$ の元は、
$\mathfrak m/\mathfrak m^2 = \text{Gr}_\mathfrak m^1(B)$ に自明に作用する
$I$ の元とちょうど一致する。すなわち
$P = \{\sigma \in I \mid \sigma|_{\mathfrak m/\mathfrak m^2} =
\text{id}_{\mathfrak m/\mathfrak m^2}\}$ である。また $I$ は
$\kappa(\mathfrak m) = B/\mathfrak m$ に自明に作用する $D$ の元からなる。
次数付き環 $\text{Gr}_\mathfrak m(B)$ は
$\text{Gr}^0_\mathfrak m(B) = B/\mathfrak m$ 上
$\text{Gr}^1_\mathfrak m(B) = \mathfrak m/\mathfrak m^2$ で生成されるので、
$P = \{\sigma \in D \mid \sigma \text{ acts trivially on }
\text{Gr}_\mathfrak m(B)\}$ である。従って (1) が成り立つ。
$P$ は準同型 $D \to \text{Aut}(\text{Gr}_\mathfrak m(B))$ の核だから、
これから (2) が従う。(3) を証明できれば、$I_t$ は
$\mu_e(\kappa(\mathfrak m))$ と同型になり、(4) と (5) も従う。
実際、上の議論により
$|I_t| \geq |\mu_e(\kappa(\mathfrak m))|$ である。
% P02-E1252: source line 33361 reverses the subgroup-order inequality; the image I_t of theta has order at most that of mu_e.

\medskip\noindent
従って $\theta$ の核 $P$ が $p$-群であることを示せば十分である。
核の非自明な元 $\sigma$ をとる。するとすべての $i$ に対して
$\sigma - \text{id}$ は $\mathfrak m_C^i$ を $\mathfrak m_C^{i + 1}$ に写す。
$m$ を $\sigma$ の位数とする。$\sigma(c) \not = c$ となる $c \in C$ を選ぶ。
ある $i$ に対して $\sigma(c) - c \in \mathfrak m_C^i$ かつ
$\sigma(c) - c \not \in \mathfrak m_C^{i + 1}$ であり、
\begin{align*}
0
& =
\sigma^m(c) - c \\
& =
\sigma^m(c) - \sigma^{m - 1}(c) + \ldots + \sigma(c) - c \\
& =
\sum\nolimits_{j = 0, \ldots, m - 1} \sigma^j(\sigma(c) - c) \\
& \equiv
m(\sigma(c) - c) \bmod \mathfrak m_C^{i + 1}
\end{align*}
である。従って $p | m$ である（または $p = 1$ なら $m = 0$ である）。
% P02-E1253: source line 33383 invokes the impossible order m = 0 under the unexplained convention p = 1 for characteristic zero.
従って $\theta$ の核のすべての元の位数は $p$ で割り切れ、すなわち
$\Ker(\theta)$ は $p$-群である。

\medskip\noindent
(6) の証明。$\kappa(\mathfrak m)/\kappa$ が分離的であると仮定する。
この場合 Lemma \ref{more-algebra-lemma-galois-galois} により $f = |D/I|$ であり、
$I$ の位数は $e$ である。$e = 1$ なら
$P = I = \{\text{id}\}$ なので結論は成り立つ。$e > 1$ なら、
$B/\mathfrak m^2 = C/\pi_C^2C$ は $\kappa$-代数であり、
$\kappa(\mathfrak m)$ の小拡大である。$\kappa(\mathfrak m)$ は
$\kappa$ 上形式的 \'etale なので（Algebra, Lemma
\ref{algebra-lemma-characterize-separable-algebraic-field-extensions}）、
$B/\mathfrak m^2 \to \kappa(\mathfrak m)$ の右逆となる $\kappa$-代数準同型
$\kappa(\mathfrak m) \to B/\mathfrak m^2$ が一意に存在する。
言い換えると、$\kappa$-代数の $D$-同変同型
$B/\mathfrak m^2 = \mathfrak m/\mathfrak m^2 \oplus \kappa(\mathfrak m)$
がある。これから直ちに、この場合
$P =
\{\sigma \in D \mid \sigma|_{B/\mathfrak m^2} = \text{id}_{B/\mathfrak m^2}\}$
であることが分かる。
\end{proof}

\begin{example}
\label{more-algebra-example-counter-description-P}
Lemma \ref{more-algebra-lemma-galois-inertia} の (6) の等式は、
$\kappa(\mathfrak m)$ に関する仮定なしには成り立たない。
反例は $A = Z_2[t]_{(2)}$ の拡大 $B = A[x]/(x^2 - t)$ である。
実際、このとき $\sigma(x) = -x = x - 2x$ であり、
$2x$ は $\mathfrak m^2$ に属さない。
\end{example}

\begin{definition}
\label{more-algebra-definition-wild-inertia}
% P02-E1254: source line 33419 is a sentence fragment (“With assumptions and notation as in ...”).
仮定と記号は Lemma \ref{more-algebra-lemma-galois-inertia} と同じとする。
\begin{enumerate}
\item $\mathfrak m$ の {\it 野性惰性群} とは部分群 $P$ である。
% P02-E1255: source lines 33422--33423 define the tame inertia group as the quotient map rather than the quotient group I_t.
\item $\mathfrak m$ の {\it 順惰性群} とは商 $I \to I_t$ である。
\end{enumerate}
% P02-E1256: source line 33425 omits “by” in “We denote theta ... the surjective map”.
(\ref{more-algebra-equation-inertia-character}) の全射で、核が $P$ であり同型
$I_t \to \mu_e(\kappa(\mathfrak m))$ を誘導するものを
$\theta : I \to \mu_e(\kappa(\mathfrak m))$ と表す。
\end{definition}

\begin{lemma}
\label{more-algebra-lemma-inertia-character}
仮定と記号は Lemma \ref{more-algebra-lemma-galois-inertia} と同じとする。
惰性指標 $\theta : I \to \mu_e(\kappa(\mathfrak m))$ は
$\tau \in D$ と $\sigma \in I$ に対して
$$
\theta(\tau \sigma \tau^{-1}) = \tau(\theta(\sigma))
$$
を満たす。
\end{lemma}

\begin{proof}
$I$ は $D$ の正規部分群であり、$\tau$ は、例えば Lemma
\ref{more-algebra-lemma-galois-galois} で論じた写像
$D \to \text{Aut}(\kappa(\mathfrak m))$ を通して
$\kappa(\mathfrak m)$ に作用するので、この式には意味がある。
$\theta$ の構成を思い出す。$B_\mathfrak m$ の素元 $\pi$ を選び、
$\sigma \in I$ に対して $\sigma(\pi) = \theta_\sigma \pi$ と書く。
このとき $\theta(\sigma)$ は $\theta_\sigma$ の剰余体における像
$\overline{\theta}_\sigma$ である。任意の $\tau \in D$ に対して、
ある単元 $\theta_\tau$ によって $\tau(\pi) = \theta_\tau \pi$ と書ける。
このとき $\theta_{\tau^{-1}} = \tau^{-1}(\theta_\tau^{-1})$ である。
計算すると
\begin{align*}
\theta_{\tau \sigma \tau^{-1}}
& =
\tau(\sigma(\tau^{-1}(\pi)))/\pi \\
& =
\tau(\sigma(\tau^{-1}(\theta_\tau^{-1}) \pi))/\pi \\
& =
\tau(\sigma(\tau^{-1}(\theta_\tau^{-1})) \theta_\sigma \pi)/\pi \\
& =
\tau(\sigma(\tau^{-1}(\theta_\tau^{-1}))) \tau(\theta_\sigma) \theta_\tau
\end{align*}
となる。しかし $\sigma$ は $\pi$ を法として自明に作用するので、積
$\tau(\sigma(\tau^{-1}(\theta_\tau^{-1}))) \theta_\tau$ は剰余体で
$1$ に写る。これで補題が示された。
\end{proof}

\noindent
次の補題は Fundamental Groups, Lemma
\ref{pione-lemma-inertial-invariants-unramified} で一般化する。

\begin{lemma}
\label{more-algebra-lemma-inertial-invariants-unramified}
$A$ を分数体 $K$ をもつ離散付値環とする。$L/K$ を有限 Galois 拡大とし、
$\mathfrak m \subset B$ を $L$ における $A$ の整閉包の極大イデアルとする。
$I \subset G$ を $\mathfrak m$ の惰性群とする。このとき $B^I$ は
$L^I$ における $A$ の整閉包であり、
$A \to (B^I)_{B^I \cap \mathfrak m}$ は \'etale である。
\end{lemma}

\begin{proof}
$B' = B^I$ と書く。定義から $B' = B^I$ は $L^I$ における $A$ の
整閉包である。さらに
$\mathfrak m' = B^I \cap \mathfrak m = B' \cap \mathfrak m \subset B'$
と書く。Lemma \ref{more-algebra-lemma-one-orbit} により、$\mathfrak m$ は
$\mathfrak m'$ 上にある $B$ の唯一の素イデアルである。
$I$ は $\kappa(\mathfrak m)$ に自明に作用するので、Lemma
\ref{more-algebra-lemma-invariants-modulo} により拡大
$\kappa(\mathfrak m)/\kappa(\mathfrak m')$ は純非分離である
（より容易な別証として Lemma
\ref{more-algebra-lemma-one-orbit-geometric} の結果を適用してもよい）。
$D/I$ は $\kappa(\mathfrak m')$ に忠実に作用するので、
$D/I$ は $\kappa(\mathfrak m)$ にも忠実に作用する。もちろん $A$ の剰余体
$\kappa$ の元はこの作用で固定される。Galois 理論により
$[\kappa(\mathfrak m') : \kappa] \geq |D/I|$ である
（Fields, Lemma \ref{fields-lemma-galois-over-fixed-field} を参照）。

\medskip\noindent
$\pi$ を $A$ の素元とする。
$\text{Norm}_{L/K}(\pi) = \pi^{[L : K]}$ なので、Algebra, Lemma
\ref{algebra-lemma-finite-extension-dim-1} から
$$
|G| = [L : K] = [L : K]\ \text{ord}_A(\pi) =
|G/D|\ [\kappa(\mathfrak m) : \kappa]\ \text{ord}_{B_\mathfrak m}(\pi)
$$
を得る。実際、$B$ の極大イデアルは $n = |G/D|$ 個で、すべて $G$ の下で
共役である（Remark \ref{more-algebra-remark-finite-separable-extension} および
Lemma \ref{more-algebra-lemma-galois} を参照）。同じ議論を次数 $|I|$ の有限拡大
% P02-E1257: source lines 33515--33516 duplicate “the finite extension”.
$L/L^I$ に適用すると
$$
|I|\ \text{ord}_{B'_{\mathfrak m'}}(\pi) =
[\kappa(\mathfrak m) : \kappa(\mathfrak m')]\ \text{ord}_{B_\mathfrak m}(\pi)
$$
を得る。従って
$$
\text{ord}_{B'_{\mathfrak m'}}(\pi) =
\frac{|D/I|}{[\kappa(\mathfrak m') : \kappa]}
$$
である。左辺は正の整数であり、上で見たように右辺は $\leq 1$ なので、
等号が成り立ち、$\text{ord}_{B'_{\mathfrak m'}}(\pi) = 1$ かつ
$\kappa(\mathfrak m')/\kappa$ の次数は $|D/I|$ である。
従って
% P02-E1258: source line 33530 uses B localized at m prime on the right, where B prime localized at m prime is required.
$\pi B'_{\mathfrak m'} = \mathfrak m' B_\mathfrak m'$ であり、
$\kappa(\mathfrak m')$ は $\kappa$ 上 Galois で、その Galois 群は
$D/I$ である。特にこれは分離的である（Fields, Lemma
\ref{fields-lemma-finite-Galois} を参照）。Algebra, Lemma
\ref{algebra-lemma-characterize-etale} により、望みどおり
$A \to B'_{\mathfrak m'}$ は \'etale である。
\end{proof}

\begin{remark}
\label{more-algebra-remark-tower-of-rings}
$A$ を分数体 $K$ をもつ離散付値環とする。$L/K$ を有限 Galois 拡大とし、
$\mathfrak m \subset B$ を $L$ における $A$ の整閉包の極大イデアルとする。
% P02-E1259: source line 33548 uses a singular predicate for the three named groups.
$\mathfrak m$ の野性惰性群、惰性群、分解群を
$$
P \subset I \subset D \subset G
$$
とする。図式
$$
\xymatrix{
\mathfrak m \ar@{-}[d] \ar@{-}[r] &
\mathfrak m^P \ar@{-}[d] \ar@{-}[r] &
\mathfrak m^I \ar@{-}[d] \ar@{-}[r] &
\mathfrak m^D \ar@{-}[d] \ar@{-}[r] &
A \cap \mathfrak m \ar@{-}[d] \\
B & B^P \ar[l] & B^I \ar[l] & B^D \ar[l] & A \ar[l]
}
$$
を考える。$B^P, B^I, B^D$ はそれぞれ体 $L^P$, $L^I$, $L^D$ における
$A$ の整閉包である。従って $B^P$ は $L^P$ における $B^I$ の整閉包でもあり、
以下同様である。また $\mathfrak m^P = \mathfrak m \cap B^P$、
$\mathfrak m^I = \mathfrak m \cap B^I$、および
$\mathfrak m^D = \mathfrak m \cap B^D$ である。従って図式の上段は、
スペクトルの誘導写像による $\mathfrak m \in \Spec(B)$ の像に対応する。
以上を踏まえると、次が成り立つ。
\begin{enumerate}
\item 拡大 $L^I/L^D$ は群 $D/I$ をもつ Galois 拡大である。
\item 拡大 $L^P/L^I$ は群 $I_t = I/P$ をもつ Galois 拡大である。
\item 拡大 $L^P/L^D$ は群 $D/P$ をもつ Galois 拡大である。
\item $\mathfrak m^I$ は $\mathfrak m^D$ 上にある $B^I$ の唯一の素イデアルである。
\item $\mathfrak m^P$ は $\mathfrak m^I$ 上にある $B^P$ の唯一の素イデアルである。
\item $\mathfrak m$ は $\mathfrak m^P$ 上にある $B$ の唯一の素イデアルである。
\item $\mathfrak m^P$ は $\mathfrak m^D$ 上にある $B^P$ の唯一の素イデアルである。
\item $\mathfrak m$ は $\mathfrak m^I$ 上にある $B$ の唯一の素イデアルである。
\item $\mathfrak m$ は $\mathfrak m^D$ 上にある $B$ の唯一の素イデアルである。
\item $A \to B^D_{\mathfrak m^D}$ は \'etale で、自明な剰余体拡大を誘導する。
\item $B^D_{\mathfrak m^D} \to B^I_{\mathfrak m^I}$ は \'etale で、
Galois 群 $D/I$ をもつ剰余体の Galois 拡大を誘導する。
\item $A \to B^I_{\mathfrak m^I}$ は \'etale である。
\item $B^I_{\mathfrak m^I} \to B^P_{\mathfrak m^P}$ は、
$p$ と互いに素な分岐指数 $|I/P|$ をもち、自明な剰余体拡大を誘導する。
\item $B^D_{\mathfrak m^D} \to B^P_{\mathfrak m^P}$ は、
$p$ と互いに素な分岐指数 $|I/P|$ をもち、分離剰余体拡大を誘導する。
\item $A \to B^P_{\mathfrak m^P}$ は、$p$ と互いに素な分岐指数
$|I/P|$ をもち、分離剰余体拡大を誘導する。
\end{enumerate}
(1)、(2)、(3) は Galois 理論（Fields, Section
\ref{fields-section-galois-theory}）と Lemma
\ref{more-algebra-lemma-galois-inertia} から直ちに従う。(4)--(9) は Lemma
\ref{more-algebra-lemma-galois} から明らかである。(12) は Lemma
\ref{more-algebra-lemma-inertial-invariants-unramified} である。分解
$A \to B^D_{\mathfrak m^D} \to B^I_{\mathfrak m^I}$ があるので、
(10) と (11) の \'etale 性も得られる。(10) の剰余体拡大は分離的であり、
Lemma \ref{more-algebra-lemma-galois-galois} で示したように $D/I$ は
$\text{Aut}(\kappa(\mathfrak m)/\kappa_A)$ に全射するので、自明でなければならない。
同じ議論で (11) の剰余体に関する主張も示される。
(13)、(14)、(15) を示すには (13) を示せば十分である。
上で見たように、拡大 $L^P/L^I$ は $p$ と互いに素な位数の巡回 Galois 群をもち、
素イデアル $\mathfrak m^P$ は $\mathfrak m^I$ 上にある唯一の素イデアルで、
$I/P$ の剰余体への作用は自明である。従ってこの拡大と離散付値環
$B^I_{\mathfrak m^I}$ に Lemma \ref{more-algebra-lemma-galois-inertia} を適用すると、
(13) が成り立つことが分かる。
\end{remark}

\begin{lemma}
\label{more-algebra-lemma-compare-inertia}
$A$ を分数体 $K$ をもつ離散付値環とする。$M/L/K$ を、$M/K$ と $L/K$ が
有限 Galois である塔とする。
% P02-E1260: source line 33623 uses singular “the integral closure” for C and B in M and L.
$C$、$B$ をそれぞれ $M$、$L$ における $A$ の整閉包とする。
$\mathfrak m' \subset C$ を極大イデアルとし、
$\mathfrak m = \mathfrak m' \cap B$ とおく。さらに
$$
P \subset I \subset D \subset \text{Gal}(L/K)
\quad\text{and}\quad
P' \subset I' \subset D' \subset \text{Gal}(M/K)
$$
を $\mathfrak m$ と $\mathfrak m'$ の野性惰性群、惰性群、分解群とする。
このとき標準全射 $\text{Gal}(M/K) \to \text{Gal}(L/K)$ は全射
$P' \to P$、$I' \to I$、$D' \to D$ を誘導する。さらに、これらは可換図式
$$
\vcenter{
\xymatrix{
D' \ar[r] \ar[d] &
\text{Aut}(\kappa(\mathfrak m')/\kappa_A) \ar[d] \\
D \ar[r] &
\text{Aut}(\kappa(\mathfrak m)/\kappa_A)
}
}
\quad\text{and}\quad
\vcenter{
\xymatrix{
I' \ar[r]_-{\theta'} \ar[d] &
\mu_{e'}(\kappa(\mathfrak m')) \ar[d]^{(-)^{e'/e}} \\
I \ar[r]^-\theta &
\mu_e(\kappa(\mathfrak m))
}
}
$$
に入る。ここで $e'$ と $e$ は、それぞれ
$A \to C_{\mathfrak m'}$ と $A \to B_\mathfrak m$ の分岐指数である。
\end{lemma}

\begin{proof}
写像 $\text{Gal}(M/K) \to \text{Gal}(L/K)$ の下で群
$P', I', D'$ がそれぞれ $P, I, D$ に写ることは、これらの群の定義から
直ちに分かる。第 1 の図式の可換性も明らかである。実際、
$\kappa(\mathfrak m)/\kappa_A$ は正規なので、$\kappa_A$ 上の
$\kappa(\mathfrak m')$ の任意の自己同型は $\kappa_A$ 上の
$\kappa(\mathfrak m)$ の自己同型を誘導し、従って第 1 図式の右の縦矢印が得られる
（Lemma \ref{more-algebra-lemma-galois-galois} および Fields, Lemma
\ref{fields-lemma-lift-maps} を参照）。

\medskip\noindent
写像 $I' \to I$ と $D' \to D$ は Lemma
\ref{more-algebra-lemma-one-orbit-geometric-galois-compare} により全射である。
$P'$ と $P$ は $I'$ と $I$ の p-Sylow 部分群なので、
$P' \to P$ の全射性も従う。

\medskip\noindent
第 2 の図式の可換性を見るため、$C_{\mathfrak m'}$ の素元 $\pi'$ と
$B_\mathfrak m$ の素元 $\pi$ を選ぶ。このとき $C_{\mathfrak m'}$ の
ある単元 $c'$ によって $\pi = c' (\pi')^{e'/e}$ と書ける。
$\sigma' \in I'$ に対して、その像 $\sigma \in I$ は単に $\sigma'$ の
$L$ への制限である。ある単元 $c \in C_{\mathfrak m'}$ によって
$\sigma'(\pi') = c \pi'$ と書き、$B_\mathfrak m$ のある単元 $b$ によって
$\sigma(\pi) = b \pi$ と書く。このとき $\sigma'(\pi) = b \pi$ であり、
$$
b \pi = \sigma'(\pi) = \sigma'(c' (\pi')^{e'/e}) =
\sigma'(c') c^{e'/e} (\pi')^{e'/e} =
\frac{\sigma'(c')}{c'} c^{e'/e} \pi
$$
を得る。$\sigma' \in I'$ なので、$b$ と $c^{e'/e}$ は剰余体で同じ像をもつ。
これが示すべきことであった。
\end{proof}

\begin{remark}
\label{more-algebra-remark-canonical-inertia-character}
無限 Galois 拡大に惰性指標
$\theta : I \to \mu_e(\kappa(\mathfrak m))$ を用いるには、
これをスケールしておくと便利である。
$A, K, L, B, \mathfrak m, G, P, I, D, e, \theta$ を Lemma
\ref{more-algebra-lemma-galois-inertia} と Definition
\ref{more-algebra-definition-wild-inertia} のとおりとする。
% P02-E1261: source line 33702 says “with q is a power” instead of “where q is a power”.
このとき $e = q |I_t|$ であり、$q$ は $\kappa(\mathfrak m)$ の標数 $p$ が
正ならその冪、標数が零なら $1$ である。
$\kappa(\mathfrak m)$ の標数は $p$ なので
$\mu_e(\kappa(\mathfrak m)) = \mu_{|I_t|}(\kappa(\mathfrak m))$ である。
写像
$$
\theta_{can} = q\theta : I \longrightarrow \mu_{|I_t|}(\kappa(\mathfrak m))
$$
を考える。この写像は同型
$\theta_{can} : I_t \to \mu_{|I_t|}(\kappa(\mathfrak m))$ を誘導する。
Lemma \ref{more-algebra-lemma-inertia-character} により、$\tau \in D$ と $\sigma \in I$ に対して
$\theta_{can}(\tau \sigma \tau^{-1}) = \tau(\theta_{can}(\sigma))$ である。
最後に、$M/L$ が $M/K$ を Galois とする拡大で、$\mathfrak m'$ が
$M$ における $A$ の整閉包の $\mathfrak m$ 上にある素イデアルなら、Lemma
\ref{more-algebra-lemma-compare-inertia} により可換図式
$$
\xymatrix{
I' \ar[r]_-{\theta'_{can}} \ar[d] &
\mu_{|I'_t|}(\kappa(\mathfrak m')) \ar[d]^{(-)^{|I'_t|/|I_t|}} \\
I \ar[r]^-{\theta_{can}} &
\mu_{|I_t|}(\kappa(\mathfrak m))
}
$$
を得る。
\end{remark}













+\section{Krasner の補題}
\label{more-algebra-section-krasner}

\noindent
離散付値体の場合の Krasner の補題は次のとおりである。

\begin{lemma}[Krasner の補題]
\label{more-algebra-lemma-krasner}
$A$ を次元 $1$ の完備局所整域とする。$P(t) \in A[t]$ を
$A$ に係数をもつ多項式とする。$\alpha \in A$ を $P$ の根で、
導多項式 $P' = \text{d}P/\text{d}t$ の根ではないものとする。
すべての $c \geq 0$ に対して、次を満たす整数 $n$ が存在する。
係数が $\mathfrak m_A^n$ に属する任意の $Q \in A[t]$ に対して、
多項式 $P + Q$ は $\beta - \alpha \in \mathfrak m_A^c$ となる根
$\beta \in A$ をもつ。
\end{lemma}

\begin{proof}
非零元 $\pi \in \mathfrak m$ を選ぶ。$A$ の次元は $1$ なので
$\mathfrak m = \sqrt{(\pi)}$ である。仮定により、ある $m \geq 0$ と
$a \in A$ に対して $P'(\alpha)^{-1} = \pi^{-m} a$ と書ける。
$c \geq m + 1$ と仮定してよいし、そう仮定する。
$\mathfrak m_A^n \subset (\pi^{c + m})$ となる $n$ を選び、
主張にある任意の $Q$ をとる。後で用いるため、ある
$R(x, y) \in A[x, y]$ に対して
$$
P(x + y) = P(x) + P'(x)y + R(x, y)y^2
$$
と書けることに注意する。帰納法により、次を満たす列
$\alpha_m, \alpha_{m + 1}, \alpha_{m + 2}, \ldots$ を構成する。
\begin{enumerate}
\item $\alpha_k \equiv \alpha \bmod \pi^c$。
\item $\alpha_{k + 1} - \alpha_k \in (\pi^k)$。
\item $(P + Q)(\alpha_k) \in (\pi^{m + k})$。
\end{enumerate}
$\beta = \lim \alpha_k$ とおけば証明が完了する。

\medskip\noindent
初期段階。$Q$ の係数は $(\pi^{c + m})$ に属するので
$(P + Q)(\alpha) \in (\pi^{c + m})$ である。従って
$\alpha_m = \alpha$ とおけばよい。この選択により、すべての
$k \geq m$ に対して $\alpha_k \equiv \alpha \bmod \pi^c$ となる。

\medskip\noindent
帰納段階。$\alpha_k$ が与えられたとし、ある $\delta \in (\pi^k)$ によって
$\alpha_{k + 1} = \alpha_k + \delta$ と書く。このとき
$$
(P + Q)(\alpha_{k + 1}) =
P(\alpha_k + \delta) + Q(\alpha_k + \delta)
$$
である。$Q$ の係数は $(\pi^{c + m})$ に属するので、
$Q(\alpha_k + \delta) \equiv Q(\alpha_k) \bmod \pi^{c + m + k}$ である。
一方、
$$
P(\alpha_k + \delta) =
P(\alpha_k) + P'(\alpha_k)\delta + R(\alpha_k, \delta)\delta^2
$$
である。$\alpha_k \equiv \alpha \bmod \pi^{m + 1}$ なので
$P'(\alpha_k) \equiv P'(\alpha) \bmod (\pi^{m + 1})$ である。
従って
$$
P(\alpha_k + \delta) \equiv P(\alpha_k) + P'(\alpha) \delta
\bmod \pi^{k + m + 1}
$$
を得る。二項を合わせると
$$
(P + Q)(\alpha_{k + 1}) \equiv (P + Q)(\alpha_k) + P'(\alpha) \delta
\bmod \pi^{k + m + 1}
$$
となる。従って
$\delta = -P'(\alpha)^{-1} (P + Q)(\alpha_k) =  - \pi^{-m} a (P + Q)(\alpha_k)$
ととればよい。帰納法の仮定により、これは $(\pi^k)$ に属する。
\end{proof}

\begin{lemma}
\label{more-algebra-lemma-approximate-separable-extension}
$A$ を分数体 $K$ をもつ離散付値環とし、$A^\wedge$ を $A$ の完備化、
$K^\wedge$ をその分数体とする。$M/K^\wedge$ が有限分離拡大なら、
$M = K^\wedge \otimes_K L$ となる有限分離拡大 $L/K$ が存在する。
\end{lemma}

\begin{proof}
Lemmas \ref{more-algebra-lemma-completion-regular} と
\ref{more-algebra-lemma-completion-dimension} により、$A^\wedge$ も離散付値環である。
特に $A^\wedge$ は整域である。証明はより一般に、$A^\wedge$ が次元 $1$ の
局所整域であるような Noether 局所環 $A$ に対して機能する。

\medskip\noindent
$M$ を $K^\wedge$ 上生成する元 $\theta \in M$ をとる
（原始元定理）。$P(t) \in K^\wedge[t]$ を $\theta$ の
$K^\wedge$ 上の最小多項式とし、非零元 $\pi \in \mathfrak m_A$ をとる。
$\theta$ を $\pi^n\theta$ で置き換えることにより、$P(t)$ の係数が
$A^\wedge$ に属すると仮定してよい。
$B = A^\wedge[\theta] = A^\wedge[t]/(P(t))$ とおく。
% P02-E1262: source line 33837 says B is finite over A, but the displayed construction makes it finite over A wedge.
$B$ は $A$ 上有限で $M$ に含まれるので、次元 $1$ の完備局所整域である。
% P02-E1263: source line 33838 invokes separability of M over K instead of the stated extension M/K wedge.
$M$ は $K$ 上分離的なので、$\theta$ は $P$ の導多項式の根ではない。
任意の整数 $n$ に対して、$P - P_1$ の係数が
$\pi^nA^\wedge[t]$ に属するようなモニック多項式 $P_1 \in A[t]$ を選べる。
Krasner の補題（Lemma \ref{more-algebra-lemma-krasner}）により、$n$ が十分大きければ
$P_1$ は $B$ に根 $\beta$ をもつ。さらに $n$ を十分大きく選べば
$\theta - \beta \in \pi B$ と仮定してよい。$t$ を $\beta$ に写す
$A^\wedge$-代数準同型 $\Phi : A^\wedge[t]/(P_1) \to B$ を考える。
$B = \pi B + \sum_{i < \deg(P)} A^\wedge \theta^i$ なので、Nakayama の補題により
$\Phi$ は全射である。$\deg(P_1) = \deg(P)$ だから、$\Phi$ は同型である。
従って環拡大 $L = K[t]/(P_1(t))$ は
$K^\wedge \otimes_K L \cong M$ を満たす。これから $L$ は体であることが従い、
証明が完了する。
\end{proof}

\begin{definition}
\label{more-algebra-definition-mixed}
$A$ を離散付値環とする。$A$ の剰余体の標数が $p > 0$ で、
$A$ の分数体の標数が $0$ であるとき、$A$ は {\it 混標数} をもつという。
この場合、離散付値環の拡大 $\mathbf{Z}_{(p)} \subset A$ が得られ、
$A$ の {\it 絶対分岐指数} とはこの拡大の分岐指数である。
\end{definition}









\section{Abhyankar の補題と順分岐}
\label{more-algebra-section-abhyankar-tame}

\noindent
この節では、離散付値環に対する Abhyankar の補題のうち、我々が最も一般的と
考える形を証明する。その後、これを応用して、離散付値環の分数体の
順分岐拡大に関するいくつかの結果を証明する。

\begin{remark}
\label{more-algebra-remark-construction}
$A \to B$ を、分数体 $K \subset L$ をもつ離散付値環の拡大とする。
$K_1/K$ を体の有限拡大とする。$A_1 \subset K_1$ を
$K_1$ における $A$ の整閉包とする。一方、
$L_1 = (L \otimes_K K_1)_{red}$ とおく。このとき $L_1$ は、
$L$ の有限次体拡大の空でない有限直積である。$B_1$ を
$L_1$ における $B$ の整閉包とする。整合的な可換図式
$$
\vcenter{
\xymatrix{
L \ar[r] & L_1 \\
K \ar[u] \ar[r] & K_1 \ar[u]
}
}
\quad\text{and}\quad
\vcenter{
\xymatrix{
B \ar[r] & B_1 \\
A \ar[u] \ar[r] & A_1 \ar[u]
}
}
$$
を得る。この状況では次が成り立つ。
\begin{enumerate}
\item Algebra, Lemma \ref{algebra-lemma-integral-closure-Dedekind}
により、環 $A_1$ は Dedekind 整域であり、$B_1$ は
Dedekind 整域の有限直積である。
\item $\pi \in A$ を一意化元とすると
$L \otimes_K K_1 = (B \otimes_A A_1)_\pi$ であり、
$\pi$ は $B \otimes_A A_1$ 上の非零因子であることに注意する。
従って環準同型 $B \otimes_A A_1 \to B_1$ は整であり、その核は
冪零元からなる。ゆえに
$\Spec(B_1) \to \Spec(B \otimes_A A_1)$ はスペクトル上全射である
（Algebra, Lemma \ref{algebra-lemma-integral-overring-surjective}）。
$A_1/\mathfrak m_A A_1$ が Artin 環であり、環準同型
$A_1/\mathfrak m_A A_1 \to
B/\mathfrak m_AB \otimes_{\kappa_A} A_1/\mathfrak m_A A_1$
が単射なので、
$\Spec(B \otimes_A A_1) \to \Spec(A_1)$ は全射である。
従って $\Spec(B_1) \to \Spec(A_1)$ は全射である。
\item $\mathfrak m_i$, $i = 1, \ldots n$ を、$n \geq 1$ 個ある
$A_1$ の極大イデアルとする。各 $i = 1, \ldots, n$ に対し、
$\mathfrak m_{ij}$, $j = 1, \ldots, m_i$ を、$\mathfrak m_i$ の上にある
$B_1$ の $m_i \geq 1$ 個の極大イデアルとする。このとき
$$
\xymatrix{
B \ar[r] & (B_1)_{\mathfrak m_{ij}} \\
A \ar[u] \ar[r] & (A_1)_{\mathfrak m_i} \ar[u]
}
$$
という離散付値環の拡大の図式を得る。
\item $A$ が Hensel 環（例えば完備）ならば、$A_1$ は
離散付値環である、すなわち $n = 1$ である。
実際、$A_1$ は整域である $A$ の有限拡大の合併であり、従って
Algebra, Lemma \ref{algebra-lemma-finite-over-henselian} により局所環である。
\item $B$ が Hensel 環（例えば完備）ならば、$B_1$ は
離散付値環の直積である、すなわち
$i = 1, \ldots, n$ に対して $m_i = 1$ である。
\item $K \subset K_1$ が純非分離ならば、$A_1$ と $B_1$ はともに
離散付値環である、すなわち $n = 1$ かつ $m_1 = 1$ である。
% P02-E1264: source line 33941 repeats “power” in “a p-power power”.
これは、任意の $b \in B_1$ に対して $b$ のある $p$ 冪が
$B$ に属し、従って $B_1$ は極大イデアルを一つしかもち得ないからである。
\item $K \subset K_1$ が有限分離拡大ならば、
$L_1 = L \otimes_K K_1$ であり、これは有限分離拡大の有限直積でもある。
従って Algebra, Lemma
\ref{algebra-lemma-Noetherian-normal-domain-finite-separable-extension}
により、$A \subset A_1$ と $B \subset B_1$ は有限である。
\item $A$ が Nagata 環ならば、$A \subset A_1$ は有限である。
\item $B$ が Nagata 環ならば、$B \subset B_1$ は有限である。
\end{enumerate}
\end{remark}

\begin{lemma}
\label{more-algebra-lemma-pull-root-uniformizer}
$A$ を一意化元 $\pi$ をもつ離散付値環とし、$n \geq 2$ とする。
$K_1 = K[\pi^{1/n}]$ とおく。このとき
\begin{enumerate}
\item $[K_1 : K] = n$,
\item $K_1$ における $A$ の整閉包 $A_1$ は環 $A[\pi^{1/n}]$ である、
\item $A_1$ は離散付値環である、
\item $A$ 上の $A_1$ の分岐指数は $n$ である、
\item $K_1$ は $A$ に関して完全分岐している、さらに
\item $n$ が $A$ の剰余標数と互いに素ならば、$K_1/K$ は順分岐し、
$K_1/K$ の任意の中間拡大は、$n$ のある約数 $d$ に対する
$\pi^{1/d}$ によって生成される。
\end{enumerate}
\end{lemma}

\begin{proof}
環 $A' = A[x]/(x^n - \pi)$ を考える。
% P02-E1265: source line 33970 omits “by” in “denote pi-prime the image”.
$\pi'$ を $A_1$ における $x$ の像とする。
$A'$ は $A$-加群として階数 $n$ の有限自由加群なので、
元 $\pi$ は $A'$ の非零因子であり、従って $\pi'$ もそうである。
さらに $A'/\pi'A'$ は体 $A/\pi A$ と同型である。
従って $A'$ は、一意化元 $\pi'$ をもつ離散付値環である。
$A'$ の分数体 $K'$ は次数 $n$ の拡大であることが従う。
% P02-E1266: source line 33976 uses “by adjoint” instead of “by adjoining”.
明らかに $K'$ は、$\pi$ の $n$ 乗根
$\pi' = \pi^{1/n}$ を $K$ に添加して得られる、すなわち
$K' \cong K_1$ である。$A'$ は正規なので、$A'$ は
$K'$ における $A$ の整閉包、すなわち $A' = A_1$ である。
これで (1) -- (5) が証明された。

\medskip\noindent
$n$ が剰余体で可逆ならば $K_1/K$ が順分岐するという主張は、
Definition \ref{more-algebra-definition-types-of-extensions} と (3), (4) から直ちに従う。
最後の主張を証明するには、Fields, Lemma
\ref{fields-lemma-subfields-kummer} により、
任意の $n$ 乗根 $\zeta \in K'$ が $K$ に含まれることを示せば十分である。
$\zeta \in (A')^*$ は明らかである。
$a \in A$ および
$b \in A \pi' \oplus \ldots \oplus A (\pi')^{n - 1}$
を用いて $\zeta = a + b$ と書ける。
% P02-E1267: source line 33991 calls b, rather than u, a unit of A'.
$b$ が非零ならば、$t \geq 1$ かつ $b$ が $A'$ の単元となるように
$b = u(\pi')^t$ と書ける。このとき
$$
1 = \zeta^n \equiv a^n + n a^{n - 1} u (\pi')^t \bmod (\pi')^{t + 1}
$$
$n$ と $a$ は $A$ の単元なので、これは
$b = u(\pi')^t$ が $(\pi')^{t + 1}$ を法として
$(1 - a^n)/(n a^{n - 1}) \in A$ と合同であることを意味する。
直和分解
$A' = A \oplus \ldots \oplus A(\pi')^{n - 1}$
から、これは不可能である。実際、$(\pi')^{t + 1}$ による乗法の像は
対応する直和であり、その $A(\pi')^t$ との共通部分は
$(\pi A)(\pi')^t$ である。従って $b = 0$ かつ $\zeta \in A$ となり、
望む結論を得る。
\end{proof}

\begin{lemma}
\label{more-algebra-lemma-formally-smooth-goes-up}
$A \to B$ を、分数体 $K \subset L$ をもつ離散付値環の拡大とする。
$A \to B$ が $\mathfrak m_B$-進位相に関して形式的に滑らかであると仮定する。
このとき任意の有限拡大 $K_1/K$ に対して
$L_1 = L \otimes_K K_1$, $B_1 = B \otimes_A A_1$ であり、
各拡大
$(A_1)_{\mathfrak m_i} \subset (B_1)_{\mathfrak m_{ij}}$
（Remark \ref{more-algebra-remark-construction} 参照）は
$\mathfrak m_{ij}$-進位相に関して形式的に滑らかである。
\end{lemma}

\begin{proof}
以下、Lemma \ref{more-algebra-lemma-extension-dvrs-formally-smooth} の同値条件を
断りなく用いる。$\pi \in A$ および
$\pi_i \in (A_1)_{\mathfrak m_i}$ を一意化元とする。
$\kappa_A \subset \kappa_B$ は分離的なので、環
$$
(B \otimes_A (A_1)_{\mathfrak m_i})/\pi_i (B \otimes_A (A_1)_{\mathfrak m_i}) =
B/\pi B \otimes_{A/\pi A} (A_1)_{\mathfrak m_i}/\pi_i (A_1)_{\mathfrak m_i}
$$
は体の直積であり、各因子は $\kappa_{\mathfrak m_i}$ 上分離的である。
従って $B \otimes_A (A_1)_{\mathfrak m_i}$ において
$\pi_i$ は非零因子であり、それによる商は体の直積である。
% P02-E1268: source line 34027 calls B tensor_A A_1 a Dedekind domain although its generic fibre may be a product of fields.
$B \otimes_A A_1$ は特に被約な Dedekind 整域であることが従う。
従って $B \otimes_A A_1 \subset B_1$ は等号である。
\end{proof}

\noindent
次の補題が、離散付値環に対する我々の Abhyankar の補題である。
$\kappa_B/\kappa_A$ は体の代数拡大であるとは仮定されていないことに注意する。

\begin{lemma}[Abhyankar の補題]
\label{more-algebra-lemma-abhyankar}
$A \subset B$ を離散付値環の拡大とする。
$A$ の剰余標数が $0$ であるか、またはそれが $p$ であり、
分岐指数 $e$ が $p$ と互いに素で、
$\kappa_B/\kappa_A$ が体の分離拡大であると仮定する。
$K_1/K$ を有限拡大とする。Remark \ref{more-algebra-remark-construction} の記法を用い、
ある $i$ に対して $e$ が $A \subset (A_1)_{\mathfrak m_i}$ の
分岐指数を割り切ると仮定する。このとき、すべての
$j = 1, \ldots, m_i$ に対し、
$(A_1)_{\mathfrak m_i} \subset (B_1)_{\mathfrak m_{ij}}$
は $\mathfrak m_{ij}$-進位相に関して形式的に滑らかである。
\end{lemma}

\begin{proof}
$\pi \in A$ を一意化元とする。
$\pi_1$ を $(A_1)_{\mathfrak m_i}$ の一意化元とする。
$(A_1)_{\mathfrak m_i}$ の単元 $u$ と
$A \subset (A_1)_{\mathfrak m_i}$ の分岐指数 $e_1$ を用いて、
$\pi = u \pi_1^{e_1}$ と書く。

\medskip\noindent
主張：$u$ が $K_1$ の $e$ 乗元であると仮定してよい。
実際、$x^e = u$ の根を添加して $K_1$ から得られる拡大を $K_2$ とする。
従って $K_2$ は $K_1[x]/(x^e - u)$ の一因子である。
$K_2/K_1$ は（$e$ に関する仮定により）有限分離拡大であり、
従って $A_1 \subset A_2$ は有限である。
$(A_1)_{\mathfrak m_i} \to
(A_1)_{\mathfrak m_i}[x]/(x^e - u)$ は
（$e$ が剰余標数と互いに素で $u$ が単元なので）有限 \'etale であるから、
% P02-E1269: source line 34066 localizes A_2 at the ideal m_i of A_1 without choosing an ideal of A_2 over it.
$(A_2)_{\mathfrak m_i}$ は $(A_1)_{\mathfrak m_i}$ の
有限 \'etale 拡大の一因子であり、それ自身
$(A_1)_{\mathfrak m_i}$ 上有限 \'etale である。
同じ議論により、$B_1 \subset B_2$ は有限 \'etale 拡大
$(B_1)_{\mathfrak m_{ij}} \subset (B_2)_{\mathfrak m_{ij}}$
を誘導する。$\mathfrak m_{ij} \subset B_1$ の上にある
極大イデアル $\mathfrak m'_{ij} \subset B_2$ を一つ選び
（もちろん複数あり得る）、次を考える。
$$
\xymatrix{
(B_1)_{\mathfrak m_{ij}} \ar[r] & (B_2)_{\mathfrak m'_{ij}} \\
(A_1)_{\mathfrak m_i} \ar[u] \ar[r] &
(A_2)_{\mathfrak m'_i} \ar[u]
}
$$
ここで $\mathfrak m'_i \subset A_2$ はその像である。
いま水平射の分岐指数は $1$ で、有限分離な剰余体拡大を誘導する。
従って Lemma \ref{more-algebra-lemma-extension-dvrs-formally-smooth} の同値条件により、
右側の垂直射が $\mathfrak m'_{ij}$-進位相に関して形式的に滑らかであることを
示せば十分である。$u$ は $K_2$ に $e$ 乗根をもつので、主張を得る。

\medskip\noindent
$u$ が $K_1$ に $e$ 乗根をもつと仮定する。
$e | e_1$ であり、$u$ は $K_1$ に $e$ 乗根をもつので、
ある $\theta \in K_1$ に対して $\pi = \theta^e$ である。
$K'_1 = K[\theta] \subset K_1$ を $\theta$ が生成する部分体とする。
Lemma \ref{more-algebra-lemma-pull-root-uniformizer} により、$K[\theta]$ における
$A$ の整閉包 $A'_1$ は離散付値環 $A'_1 = A[\theta]$ であり、
$A$ 上の分岐指数は $e$ である。
拡大 $K'_1/K$ に対して補題を証明できれば、すべての
$j = 1, \ldots, m_i$ について Lemma
\ref{more-algebra-lemma-formally-smooth-goes-up} を図式
$$
\xymatrix{
(B'_1)_{B'_1 \cap \mathfrak m_{ij}} \ar[r] & (B_1)_{\mathfrak m_{ij}} \\
A'_1 \ar[u] \ar[r] &
(A_1)_{\mathfrak m_i} \ar[u]
}
$$
に適用して結論を得る。これにより、次の段落で扱う場合に帰着する。

\medskip\noindent
$K_1 = K[\pi^{1/e}]$ と仮定し、$\theta = \pi^{1/e}$ とおく。
$\pi_B$ を $B$ の一意化元とし、$B$ のある単元 $w$ を用いて
$\pi = w \pi_B^e$ と書く。このとき $L_1 = L \otimes_K K_1$ は、
単元 $w$ の $e$ 乗根である $\theta / \pi_B$ を添加して得られる。
従って $B \subset B_1$ は有限 \'etale である。
そこで任意の極大イデアル $\mathfrak m \subset B_1$ に対し、
可換図式
$$
\xymatrix{
B \ar[r]_1 & (B_1)_{\mathfrak m} \\
A \ar[u]^e \ar[r]^e & A_1 \ar[u]_{e_\mathfrak m}
}
$$
を考える。矢印に沿う数は分岐指数である。
分岐指数の乗法性（Lemma \ref{more-algebra-lemma-multiplicative-e-f}）により
$e_\mathfrak m = 1$ を得る。剰余体拡大を見ると、
$\kappa(\mathfrak m)$ は $\kappa_B$ の有限分離拡大であり、
後者は $\kappa_A$ 上分離的である。
従って $\kappa(\mathfrak m)$ は、$A_1$ の剰余体に等しい
$\kappa_A$ 上分離的である。
Lemma \ref{more-algebra-lemma-extension-dvrs-formally-smooth} により結論を得る。
\end{proof}

\begin{lemma}
\label{more-algebra-lemma-composition-tame}
$A$ を分数体 $K$ をもつ離散付値環とする。
$M/L/K$ を有限分離拡大とする。
$B$ を $L$ における $A$ の整閉包とする。
$L/K$ が $A$ に関して順分岐し、
かつ $B$ のすべての極大イデアル $\mathfrak m$ に対して
$M/L$ が $B_\mathfrak m$ に関して順分岐するならば、
$M/K$ は $A$ に関して順分岐する。
\end{lemma}

\begin{proof}
$C$ を $M$ における $A$ の整閉包とする。
$C$ の任意の極大イデアル $\mathfrak m'$ は、
$B$ のある極大イデアル $\mathfrak m$ の上にある。
従って補題は、分岐指数の乗法性
（Lemma \ref{more-algebra-lemma-multiplicative-e-f}）と、
有限体拡大の塔
$\kappa(\mathfrak m')/\kappa(\mathfrak m)/\kappa_A$
があることから従う。
\end{proof}

\begin{lemma}
\label{more-algebra-lemma-subextension-tame}
$A$ を分数体 $K$ をもつ離散付値環とする。
$M/L/K$ が有限分離拡大で、$M$ が $A$ に関して順分岐するならば、
$L$ も $A$ に関して順分岐する。
\end{lemma}

\begin{proof}
以下、Remark \ref{more-algebra-remark-finite-separable-extension} の議論の結果を
断りなく用いる。$C/B/A$ を、$M/L/K$ における $A$ の整閉包とする。
$C$ は $B$ の有限環拡大なので、$\Spec(C) \to \Spec(B)$ は全射である。
% P02-E1270: source line 34173 says “for ever maximal” instead of “for every maximal”.
従って任意の極大イデアル $\mathfrak m \subset B$ に対し、
$\mathfrak m$ の上にある極大イデアル $\mathfrak m' \subset C$ が存在する。
分岐指数の乗法性（Lemma \ref{more-algebra-lemma-multiplicative-e-f}）と仮定により、
$A$ 上の $B_\mathfrak m$ の分岐指数は剰余標数と互いに素である。
$\kappa(\mathfrak m')/\kappa_A$ は有限分離拡大なので、
$\kappa(\mathfrak m)/\kappa_A$ もそうである。
\end{proof}

\begin{lemma}
\label{more-algebra-lemma-characterize-tame}
$A$ を分数体 $K$ と一意化元 $\pi \in A$ をもつ離散付値環とする。
$L/K$ を有限分離拡大とする。次は同値である。
\begin{enumerate}
\item $L$ は $A$ に関して順分岐する、
\item $\kappa_A$ で可逆な $e \geq 1$ と、
$A' = A[\pi^{1/e}]$ に関して不分岐である拡大
$L'/K' = K[\pi^{1/e}]$ が存在し、$L$ は $L'$ に含まれる、さらに
\item $\kappa_A$ で可逆な $e_0 \geq 1$ が存在し、
$\kappa_A$ で可逆なすべての $d \geq 1$ に対して、
$e = de_0$ とすれば (2) が成り立つ。
\end{enumerate}
\end{lemma}

\begin{proof}
Lemma \ref{more-algebra-lemma-pull-root-uniformizer} により、$A'$ は
分数体 $K'$ をもつ離散付値環である。
もちろん $A$ 上の $A'$ の分岐指数は $e$ である。
従って (2) が成り立てば、Lemma \ref{more-algebra-lemma-composition-tame} により
$L'$ は $A$ に関して順分岐する。
ゆえに Lemma \ref{more-algebra-lemma-subextension-tame} により、
$L$ は $A$ に関して順分岐する。

\medskip\noindent
(3) $\Rightarrow$ (2) は直ちに従う。

\medskip\noindent
(1) を仮定する。$B$ を $L$ における $A$ の整閉包とし、
$\mathfrak m_1, \ldots, \mathfrak m_n$ をその極大イデアルとする。
$e_i$ を $A \to B_{\mathfrak m_i}$ の分岐指数とする。
% P02-E1271: source line 34216 takes the lcm through e_r although only e_1,...,e_n were defined.
$e_0$ を $e_1, \ldots, e_r$ の最小公倍数とする。
仮定 (1) により、これは $\kappa_A$ で可逆である。
(3) のように $e = de_0$ とし、$A' = A[\pi^{1/e}]$ とおく。
Lemma \ref{more-algebra-lemma-pull-root-uniformizer} により、
$A \to A'$ は分数体 $K' = K[\pi^{1/e}]$ をもつ離散付値環の拡大である。
直積分解
$$
L \otimes_K K' = \prod L'_j
$$
を選ぶ。ここで $L'_j$ は体である。
$B'_j$ を $L'_j$ における $A$ の整閉包とする。
$\mathfrak m_{ijk}$ を、$\mathfrak m_i$ の上にある
$B'_j$ の極大イデアルとする。
% P02-E1272: source line 34229 localizes B'_j at m_i rather than the just-defined m_{ijk}.
$(B'_j)_{\mathfrak m_i}$ は $L'_j$ における
$B_{\mathfrak m_i}$ の整閉包であることに注意する。
Abhyankar の補題（Lemma \ref{more-algebra-lemma-abhyankar}）を
$A \subset B_{\mathfrak m_i}$ と拡大 $K'/K$ に適用すると、
$A' \to (B'_j)_{\mathfrak m_{ijk}}$ は
$\mathfrak m_{ijk}$-進位相に関して形式的に滑らかである。
これは、分岐指数が $1$ で剰余体拡大が分離的であることを意味する
（Lemma \ref{more-algebra-lemma-extension-dvrs-formally-smooth}）。
このようにして $L'_j$ は $A'$ に関して不分岐であることが分かる。
ある $j$ に対して $L' = L'_j$ とすれば、証明は完了する。
\end{proof}

\begin{lemma}
\label{more-algebra-lemma-permanence-tame}
$A$ を分数体 $K$ をもつ離散付値環とする。
\begin{enumerate}
\item $L/K$ が $A$ に関して順分岐する有限分離拡大ならば、
$L$ を含み $A$ に関して順分岐する Galois 拡大 $M/K$ が存在する。
\item $L_1/K$, $L_2/K$ が $A$ に関して順分岐する有限分離拡大ならば、
$L_1$ と $L_2$ を含み $A$ に関して順分岐する有限分離拡大 $L/K$ が存在する。
\end{enumerate}
\end{lemma}

\begin{proof}
(2) の証明。$\pi \in A$ を一意化元として選ぶ。
Lemma \ref{more-algebra-lemma-characterize-tame} により、
$\kappa_A$ で可逆な整数 $e$ と、
$A' = A[\pi^{1/e}]$ に関して不分岐である拡大
$L_i'/K' = K[\pi^{1/e}]$ を、$K$ の拡大として $L'_i/L_i$ となるように
選べる。Lemma \ref{more-algebra-lemma-permanence-unramified} により、
$A'$ に関して不分岐である拡大 $L'/K'$ であって、
$i = 1, 2$ に対して $L'_i/K$ が $L'/K'$ の中間拡大と同型になるものを
見つけられる。
% P02-E1273: source line 34268 calls this the proof of nonexistent item (3), rather than item (2).
上と同じ補題を用いると $L'/K$ は順分岐するので、これで (3) の証明が完了する。

\medskip\noindent
(1) の証明。まず $L$ をより大きな拡大で置き換え、
$L$ は $\kappa_A$ で可逆な $e$ に対する $K' = K[\pi^{1/e}]$ の拡大であり、
$A' = A[\pi^{1/e}]$ に関して不分岐であると仮定してよい
（Lemma \ref{more-algebra-lemma-characterize-tame}）。
$M$ を $K$ 上の $L$ の正規閉包とする
（Fields, Definition \ref{fields-definition-normal-closure}）。
Fields, Lemma \ref{fields-lemma-normal-closure-galois} により
$M/K$ は Galois 拡大である。一方、Fields, Lemma
\ref{fields-lemma-normal-closure-tensor-product} により、
$K$-代数の全射
$$
L \otimes_K \ldots \otimes_K L \longrightarrow M
$$
がある。Remark \ref{more-algebra-remark-finite-separable-extension} のように、
$B$ を $L$ における $A$ の整閉包とする。
$L$ が $A' = A[\pi^{1/e}]$ に関して不分岐であるという条件は、
Algebra, Lemma \ref{algebra-lemma-characterize-etale} により、
環準同型 $A' \to B$ が \'etale であることとちょうど同値である。
主張：
$$
K' \otimes_K \ldots \otimes_K K' = \prod K'_i
$$
は、$A$ に関して順分岐する体拡大 $K'_i/K$ の直積である。
このとき $A'_i$ を $K'_i$ における $A$ の整閉包とすれば、
$$
\prod A'_i \otimes_{(A' \otimes_A \ldots \otimes_A A')}
(B \otimes_A \ldots \otimes_A B)
$$
は $\prod A'_i$ 上有限 \'etale であり、従って Dedekind 整域の直積である
（Lemma \ref{more-algebra-lemma-Dedekind-etale-extension}）。
$M$ はこれらの Dedekind 整域の一つの分数体であり、その整域は
ある $i$ に対して $A'_i$ 上有限 \'etale であると結論する。
従って $M/K'_i$ は $A'_i$ のすべての極大イデアルに関して不分岐であり、
Lemma \ref{more-algebra-lemma-composition-tame} により $M/K$ は順分岐する。

\medskip\noindent
主張を証明すればよい。そのため
$A' = A[x]/(x^e - \pi)$ と書くと、
$$
A' \otimes_A \ldots \otimes_A A' =
A'[x_1, \ldots, x_r]/(x_1^e - \pi, \ldots, x_r^e - \pi)
$$
である。この環の正規化は、$i = 2, \ldots, r$ に対する
元 $y_i = x_i/x_1$ を確かに含み、それらは関係式
$y_i^e - 1 = 0$ を満たす。従って
$$
A[x_1, y_2, \ldots, y_r]/(x_1^e - \pi, y_2^e - 1, \ldots, y_r^e - 1) =
A'[y_2, \ldots, y_r]/(y_2^e - 1, \ldots, y_r^e - 1)
$$
を得る。$e$ は $A'$ で可逆なので、この環は $A'$ 上有限 \'etale である。
従って望む通り、これはそれぞれ $A'$ 上不分岐な Dedekind 整域の直積である
（不明な点があれば上記の参照文献を参照）。
\end{proof}

\begin{lemma}
\label{more-algebra-lemma-tame-goes-up}
$A \subset B$ を離散付値環の拡大とし、対応する分数体の拡大を
$L/K$ と記す。$K'/K$ を有限分離拡大とする。このとき
$$
K' \otimes_K L = \prod L'_i
$$
は体の有限直積であり、次が成り立つ。
\begin{enumerate}
\item $K'$ が $A$ に関して不分岐ならば、各 $L'_i$ は
$B$ に関して不分岐である。
\item $K'$ が $A$ に関して順分岐するならば、各 $L'_i$ は
$B$ に関して順分岐する。
\end{enumerate}
\end{lemma}

\begin{proof}
代数 $K' \otimes_K L$ は $L$ 上有限 \'etale なので、体の有限直積である。
$A'$ を $K'$ における $A$ の整閉包とする。

\medskip\noindent
(1) の場合、環準同型 $A \to A'$ は有限 \'etale である。
従って $B' = B \otimes_A A'$ は $B$ 上有限 \'etale であり、
Dedekind 整域の有限直積である
（Lemma \ref{more-algebra-lemma-Dedekind-etale-extension}）。
ゆえに $B'$ は $K' \otimes_K L$ における $B$ の整閉包である。
各 $L'_i$ が $B$ に関して不分岐であることが直ちに従う。

\medskip\noindent
$\pi \in A$ を一意化元として選ぶ。(2) を証明するため、
$K'$ を $A$ に関して順分岐するより大きな拡大で置き換えてよい
（詳細は省略する。ヒント：Lemma \ref{more-algebra-lemma-subextension-tame} を用いる）。
% P02-E1274: source line 34363 uses the ungrammatical “tame ramified”.
従って Lemma \ref{more-algebra-lemma-characterize-tame} により、
$\kappa_A$ で可逆なある $e \geq 1$ が存在し、
$K'$ は $K[\pi^{1/e}]$ を含み、$K'$ は $A[\pi^{1/e}]$ に関して
不分岐であると仮定してよい。直積分解
$$
K[\pi^{1/e}] \otimes_K L = \prod L_{e, j}
$$
を選ぶ。各 $i$ に対し、$L'_i/L_{e, j_i}$ が有限分離拡大となる
$j_i$ が存在する。$B_{e, j}$ を $L_{e, j}$ における
$B$ の整閉包とする。(1) を $K'/K[\pi^{1/e}]$ と
$A[\pi^{1/e}] \subset (B_{e, j_i})_\mathfrak m$ に適用すると、
任意の極大イデアル $\mathfrak m \subset B_{e, j_i}$ に対して、
$L'_i$ は $(B_{e, j_i})_\mathfrak m$ に関して不分岐である。
従って Lemma \ref{more-algebra-lemma-composition-tame} により、証明を完了するには
$L_{e, j}$ が $B$ に関して順分岐することを示せば十分である。

\medskip\noindent
$\theta$ を $B$ の一意化元として選ぶ。
$B$ の単元 $u$ と $t \geq 1$ を用いて $\pi = u \theta^t$ と書く。
このとき
$$
A[\pi^{1/e}] \otimes_A B = B[x]/(x^e - u \theta^t)
\subset B[y, z]/(y^{e'} - \theta, z^e - u)
$$
である。ここで $e' = e/\gcd(e, t)$ である。この写像は $x$ を
$z y^{t/\gcd(e, t)}$ に写す。右辺は Dedekind 整域の直積であり、
各因子は $B$ 上順分岐するので、証明は完了する（詳細は省略する）。
\end{proof}












\section{分岐の除去}
\label{more-algebra-section-eliminating-ramification}

\noindent
この節では Helmut Epp の結果を論じる。文献 \cite{Epp} を参照されたい。
読者には原論文を読むことを強く勧める。我々の方法は、混標数の場合と
等標数の場合を同じ方法で扱おうとする点で少し異なる。
関連する結果については、\cite{Ponomarev}, \cite{Ponomarev-Abhyankar},
\cite{Kuhlmann}, \cite{ZK} も参照されたい。

\medskip\noindent
$A \subset B$ を、分数体 $K \subset L$ をもつ離散付値環の拡大とする。
この節での目標は、Remark \ref{more-algebra-remark-construction} のように
$$
\vcenter{
\xymatrix{
L \ar[r] & L_1 \\
K \ar[u] \ar[r] & K_1 \ar[u]
}
}
\quad\text{and}\quad
\vcenter{
\xymatrix{
B \ar[r] & B_1 \ar[r] & (B_1)_{\mathfrak m_{ij}} \\
A \ar[u] \ar[r] & A_1 \ar[r] \ar[u] & (A_1)_{\mathfrak m_i} \ar[u]
}
}
$$
としたとき、拡大
$(A_1)_{\mathfrak m_i} \subset (B_1)_{\mathfrak m_{ij}}$
がすべて弱不分岐となるか、さらには該当する進位相に関して
形式的に滑らかとなるような有限拡大 $K_1/K$ を見つけることである。
この最も簡単な（ただし自明ではない）例は Abhyankar の補題、
すなわち Lemma \ref{more-algebra-lemma-abhyankar} である。

\begin{definition}
\label{more-algebra-definition-solution}
$A \to B$ を、分数体 $K \subset L$ をもつ離散付値環の拡大とする。
\begin{enumerate}
\item Remark \ref{more-algebra-remark-construction} の拡大
$(A_1)_{\mathfrak m_i} \subset (B_1)_{\mathfrak m_{ij}}$ が
すべて弱不分岐であるとき、有限体拡大 $K_1/K$ を
$A \subset B$ に対する {\it 弱解} という。
\item Remark \ref{more-algebra-remark-construction} の各拡大
$(A_1)_{\mathfrak m_i} \subset (B_1)_{\mathfrak m_{ij}}$ が
$\mathfrak m_{ij}$-進位相に関して形式的に滑らかであるとき、
有限体拡大 $K_1/K$ を $A \subset B$ に対する {\it 解} という。
\end{enumerate}
$K_1/K$ が分離的である解 $K_1/K$ を {\it 分離解} という。
\end{definition}

\noindent
一般には（弱）解は存在しない。\cite{Epp} に例がある。
剰余体拡大に関する緩やかな仮定の下で、\cite{Epp} に従い
Theorem \ref{more-algebra-theorem-epp} で弱解の存在を証明する。
次節では、幾何学的に意味のある場合に解が、また時には分離解が存在することを
導く。Proposition \ref{more-algebra-proposition-epp-essentially-finite-type} と
Lemma \ref{more-algebra-lemma-epp-essentially-finite-type-separable} を参照されたい。
しかし次の例は、一般には弱解を得るためにさえ非分離拡大が必要であることを示す。

\begin{example}
\label{more-algebra-example-inseparable-necessary}
$k$ を標数 $p > 0$ の完全体とする。$A = k[[x]]$ および
$K = k((x))$ とする。$B = A[x^{1/p}]$ とおく。
$A \to B$ に対する任意の弱解 $K_1/K$ は非分離的である
（そして $K$ の任意の有限非分離拡大は解である）。証明は省略する。
\end{example}

\noindent
解はさらなる拡大の下で安定である。Lemma \ref{more-algebra-lemma-solution-goes-up}
を参照されたい。弱解についてはこれは成り立たないことがある。
弱解は、ある意味で完全分岐拡大の下で安定である。
Lemma \ref{more-algebra-lemma-weakly-unramified-goes-up-along-totally-ramified}
を参照されたい。

\begin{lemma}
\label{more-algebra-lemma-weakly-unramified-goes-up-along-totally-ramified}
$A \to B$ を、分数体 $K \subset L$ をもつ離散付値環の拡大とし、
$A \to B$ は弱不分岐であると仮定する。
$A$ に関して完全分岐する任意の有限分離拡大 $K_1/K$ に対して、
$L_1 = L \otimes_K K_1$ は体であり、$A_1$ と
$B_1 = B \otimes_A A_1$ は離散付値環であり、
拡大 $A_1 \subset B_1$（Remark \ref{more-algebra-remark-construction} 参照）は
弱不分岐である。
\end{lemma}

\begin{proof}
$\pi \in A$ と $\pi_1 \in A_1$ を一意化元とする。
$K_1/K$ は $A$ に関して完全分岐しているので、
$A_1$ のある単元 $u_1$ に対して $\pi_1^e = u_1 \pi$ である。
従って $A_1$ は $A$ 上 $\pi_1$ によって生成され、
$K$ 上の $\pi_1$ の最小多項式 $P(t)$ は
$$
P(t) = t^e + a_{e - 1} t^{e - 1} + \ldots + a_0
$$
という形をもつ。ここで $a_i \in (\pi)$ であり、
$A$ のある単元 $u$ に対して $a_0 = u\pi$ である。
さらに $e = [K_1 : K]$ であることにも注意する。
$A \to B$ は弱不分岐なので、$\pi$ は $B$ の一意化元である。
従って $B_1 = B[t]/(P(t))$ は $t$ の類を一意化元とする
離散付値環である。これで補題は明らかである。
\end{proof}

\begin{lemma}
\label{more-algebra-lemma-solutions-go-down}
$A \to B \to C$ を、分数体 $K \subset L \subset M$ をもつ
離散付値環の拡大とする。$K_1/K$ を有限拡大とする。
\begin{enumerate}
\item $K_1$ が $A \to C$ に対する（弱）解ならば、
$K_1$ は $A \to B$ に対する（弱）解である。
\item $K_1$ が $A \to B$ に対する（弱）解であり、
$L_1 = (L \otimes_K K_1)_{red}$ が、いずれも
$B \to C$ に対する（弱）解である体の直積ならば、
$K_1$ は $A \to C$ に対する（弱）解である。
\end{enumerate}
\end{lemma}

\begin{proof}
$L_1 = (L \otimes_K K_1)_{red}$ および
$M_1 = (M \otimes_K K_1)_{red}$ とし、
$B_1 \subset L_1$ と $C_1 \subset M_1$ をそれぞれ
$B$ と $C$ の整閉包とする。
$M_1 = (M \otimes_L L_1)_{red}$ であり、$L_1$ は
$L$ の有限拡大の（空でない）有限直積であることに注意する。
従って環準同型 $B_1 \to C_1$ は、
Remark \ref{more-algebra-remark-construction} で論じた形の環準同型の有限直積である。
特に $\Spec(C_1) \to \Spec(B_1)$ は全射である。
極大イデアル $\mathfrak m \subset C_1$ を選び、
離散付値環の拡大
$$
(A_1)_{A_1 \cap \mathfrak m} \to
(B_1)_{B_1 \cap \mathfrak m} \to
(C_1)_\mathfrak m
$$
を考える。合成が弱不分岐ならば、
$(A_1)_{A_1 \cap \mathfrak m} \to
(B_1)_{B_1 \cap \mathfrak m}$ も弱不分岐である。
剰余体拡大
$\kappa_{A_1 \cap \mathfrak m} \to \kappa_\mathfrak m$
が分離的ならば、その中間拡大
$\kappa_{A_1 \cap \mathfrak m} \to
\kappa_{B_1 \cap \mathfrak m}$ も分離的である。
Lemma \ref{more-algebra-lemma-extension-dvrs-formally-smooth} を考慮すると、
これで (1) が証明される。(2) も同様に証明できる。
\end{proof}

\begin{lemma}
\label{more-algebra-lemma-solution-after-strict-henselization}
$A \to B$ を離散付値環の拡大とする。次の可換図式
$$
\xymatrix{
B \ar[r] & B' \\
A \ar[r] \ar[u] & A' \ar[u]
}
$$
をなす離散付値環の拡大で、以下を満たすものが存在する。
\begin{enumerate}
\item 分数体の拡大 $K'/K$ と $L'/L$ は分離代数拡大である、
\item $A'$ と $B'$ の剰余体はそれぞれ $A$ と $B$ の剰余体の
分離代数閉包である、さらに
\item $A' \to B'$ に対して解、弱解、または分離解が存在するならば、
$A \to B$ に対してそれぞれ解、弱解、または分離解が存在する。
\end{enumerate}
\end{lemma}

\begin{proof}
Algebra, Lemma \ref{algebra-lemma-colimit-finite-etale-given-residue-field}
により、有限 \'etale 拡大のフィルター余極限であり、
$A'$ の剰余体が $A$ の剰余体の分離代数閉包となるような
拡大 $A \subset A'$ が存在する。このとき $A \subset A'$ は
離散付値環の拡大であり、誘導される分数体拡大 $K'/K$ は
分離代数拡大である。

\medskip\noindent
$B \subset B'$ を $B$ の強 Hensel 化とする。
$B \subset B'$ は離散付値環の拡大であり、その分数体拡大は
分離代数拡大である。Algebra, Lemma
\ref{algebra-lemma-strictly-henselian-functorial-prepare} により、
補題の主張にある可換図式が存在する。補題の (1), (2) は明らかである。

\medskip\noindent
$K'_1/K'$ を $A' \to B'$ に対する（弱）解とする。
$A'$ は余極限なので、有限 \'etale 拡大 $A \subset A_1'$ と、
$A_1'$ の分数体 $F$ の有限拡大 $K_1$ で
$K'_1 = K' \otimes_F K_1$ を満たすものを見つけられる。
$A \subset A_1'$ は有限 \'etale で $B'$ は強 Hensel 環なので、
$B' \otimes_A A_1'$ は $B'$ と同型な環の有限直積である。従って
$$
L' \otimes_K K_1 = L' \otimes_K F \otimes_F K_1
$$
は $L' \otimes_{K'} K'_1$ と同型な環の有限直積である。
これより $K_1/K$ は $A \to B'$ に対する（弱）解である。
従って Lemma \ref{more-algebra-lemma-solutions-go-down} により、
$A \to B$ に対する（弱）解でもある。
\end{proof}

\begin{lemma}
\label{more-algebra-lemma-galois-relative}
$A \to B$ を、分数体 $K \subset L$ をもつ離散付値環の拡大とする。
$K_1/K$ を正規拡大とし、$G = \text{Aut}(K_1/K)$ とおく。
このとき $G$ は Remark \ref{more-algebra-remark-construction} の環
$K_1$, $L_1$, $A_1$, $B_1$ に作用し、
$B_1$ の極大イデアル全体に推移的に作用する。
\end{lemma}

\begin{proof}
最後の主張以外はすべて明らかである。作用の軌道が二つ以上あるならば、
一つの軌道のすべての極大イデアルで消え、別の軌道のすべての極大イデアルで
剰余類 $1$ をもつ元 $b \in B_1$ を見つけられる。このとき
% P02-E1275: the source assumes only “normal” but uses a finite product over G and K = (K_1)^G.
$b' = \prod_{\sigma \in G} \sigma(b)$ は、
$B_1 \subset L_1 = (L \otimes_K K_1)_{red}$ の $G$-不変元であり、
$B_1$ のいくつかの極大イデアルには属するが、
$B_1$ のすべての極大イデアルに属するわけではない。
これを $L \otimes_K K_1$ の元に持ち上げ、高い冪をとることにより、
ある $N > 0$ に対して $(b')^N$ に写る
$L \otimes_K K_1$ の $G$-不変元 $b''$ を得る。
実際、これが必要なのは標数が $p > 0$ の場合だけであり、この場合、
十分大きな $p$ 冪 $q$ をとる写像が標準的な写像
$(L \otimes_K K_1)_{red} \to L \otimes_K K_1$ を定める。
$K = (K_1)^G$ なので $b'' \in L$ と結論する。
$b''$ は $B_1$ の元に写るので、$B$ は正規であることから $b'' \in B$ である。
一方では、$b'$ が $B_1$ のある極大イデアルに属するので
$b'' \in \mathfrak m_B$ でなければならず、他方では、
$b'$ が $B_1$ のすべての極大イデアルに属するわけではないので
$b'' \not \in \mathfrak m_B$ でなければならない。
この矛盾により補題の証明は完了する。
\end{proof}

\begin{lemma}
\label{more-algebra-lemma-make-degree-q-extension}
$A$ を一意化元 $\pi$ をもつ離散付値環とする。
$A$ の剰余標数が $p > 0$ ならば、任意の $n > 1$ と $p$ 冪 $q$ に対し、
$A$ に関して完全分岐する次数 $q$ の分離拡大 $L/K$ であって、
$L$ における $A$ の整閉包 $B$ の分岐指数が $q$ であり、
ある $b, b' \in B$ に対して
$\pi_B^q = \pi + \pi^n b$ および
$\pi_B^q = \pi + (\pi_B)^{nq}b'$ を満たす一意化元 $\pi_B$ を
もつものが存在する。
\end{lemma}

\begin{proof}
$K$ の標数が零ならば、Lemma \ref{more-algebra-lemma-pull-root-uniformizer} により、
$\pi_B^q = \pi$ で与えられる拡大をとればよい。
$K$ の標数が $p > 0$ ならば、
$z^q - \pi^n z = \pi^{1 - q}$ で与えられる $K$ の拡大をとればよい。
実際、$y = \pi z$ とすれば
$y^q - \pi^{n + q - 1} y = \pi$ である。
$\pi_B = y$ ととれば望む結果を得る。
\end{proof}

\begin{lemma}
\label{more-algebra-lemma-pre-purely-inseparable-case}
% P02-E1276: source line 34664 says “reside field” instead of “residue field”.
$A$ を離散付値環とする。剰余体 $\kappa_A$ の標数が $p > 0$ であり、
$a \in A$ の $\kappa_A$ における剰余類が $p$ 乗元でないと仮定する。
このとき $a$ は $K$ の $p$ 乗元ではなく、
$K[a^{1/p}]$ における $A$ の整閉包は環 $A[a^{1/p}]$ である。
これは $A$ 上弱不分岐な離散付値環である。
\end{lemma}

\begin{proof}
この補題は自明である。
\end{proof}

\begin{lemma}
\label{more-algebra-lemma-purely-inseparable-case}
$A \subset B \subset C$ を、分数体 $K \subset L \subset M$ をもつ
離散付値環の拡大とする。$\pi \in A$ を一意化元とし、
\begin{enumerate}
\item $B$ は Nagata 環である、
\item $A \subset B$ は弱不分岐である、
\item $M$ は $L$ の次数 $p$ の純非分離拡大である
\end{enumerate}
と仮定する。このとき次のいずれかが成り立つ。
\begin{enumerate}
\item $A \to C$ は弱不分岐である、
\item $C = B[\pi^{1/p}]$ である、または
\item $A$ に関して完全分岐する次数 $p$ の分離拡大 $K_1/K$ が存在し、
$L_1 = L \otimes_K K_1$ と $M_1 = M \otimes_K K_1$ は体であり、
整閉包の写像 $A_1 \to B_1 \to C_1$ は弱不分岐な
離散付値環の拡大である。
\end{enumerate}
\end{lemma}

\begin{proof}
$e$ を $B$ 上の $C$ の分岐指数とする。$e = 1$ ならば証明は終わる。
そうでなければ、Lemmas \ref{more-algebra-lemma-inequality} と
\ref{more-algebra-lemma-ramification-index-a-power-of-p} により $e = p$ である。
これはさらに、$B$ と $C$ の剰余体が一致することを意味する。
$C$ の一意化元 $\pi_C$ を選ぶ。$C$ のある単元 $u$ に対し
$\pi_C^p = u \pi$ と書く。
$\pi_C^p \in L$ なので $u \in B^*$ である。また $M = L[\pi_C]$ である。

\medskip\noindent
ある整数 $m \geq 0$ が存在し、
$$
u = \sum\nolimits_{0 \leq i < m} b_i^p \pi^i + b \pi^m
$$
と書けると仮定する。ここで $b_i \in B$ であり、
$b \in B$ の $\kappa_B$ における像は $p$ 乗元ではない。
Lemma \ref{more-algebra-lemma-make-degree-q-extension} のような拡大 $K_1/K$ を
$n = m + 2$ として選び、$K_1$ における $A$ の整閉包 $A_1$ の
% P02-E1277: source line 34714 omits “by” in “denote pi-prime the uniformizer”.
一意化元を $\pi'$ と記す。ある $a \in A_1$ に対し
$\pi = (\pi')^p + (\pi')^{np} a$ である。
$B_1$ を $L \otimes_K K_1$ における $B$ の整閉包とする。
Lemma \ref{more-algebra-lemma-weakly-unramified-goes-up-along-totally-ramified}
により $A_1 \to B_1$ は弱不分岐であることに注意する。
$B_1$ において
$$
u \pi =
\left(\sum\nolimits_{0 \leq i < m} b_i (\pi')^{i + 1}\right)^p +
b (\pi')^{(m + 1)p} + (\pi')^{np} b_1
$$
である。ここで $b_1 \in B_1$ である（計算は省略する）。
$M_1$ は
$$
b + (\pi')^{n - m - 1} b_1
$$
の $p$ 乗根を $L_1$ に添加して得られると結論する。
$B_1$ の剰余体は $B$ の剰余体に等しいので、
Lemma \ref{more-algebra-lemma-pre-purely-inseparable-case} により
$M_1/L_1$ の次数は $p$ であり、$B_1$ の整閉包 $C_1$ は
$B_1$ 上弱不分岐である。従ってこの場合の結論を得る。

\medskip\noindent
前段落のような整数 $m$ が存在しないならば、
% P02-E1278: source lines 34740-34742 refer to B_1, which is defined only in the preceding branch; the argument needs B.
$u$ は $B_1$ の $\pi$-進完備化における $p$ 乗元である。
$B$ は Nagata 環なので、Algebra, Lemma
\ref{algebra-lemma-nagata-pth-roots} により、これは
$u$ が $B_1$ の $p$ 乗元であることを意味する。
従って補題の主張の第二の場合が成り立つ。
\end{proof}

\begin{lemma}
\label{more-algebra-lemma-cohen}
$A$ を素数 $p$ で零化される局所環とし、その極大イデアルは
冪零であるとする。剰余写像 $A \to \kappa_A$ の切断となる環準同型
$\sigma : \kappa_A \to A$ が存在する。
$A \to A'$ が局所環の局所準同型ならば、
拡大 $\kappa_{A'}/\kappa_A$ が分離的である限り、
$\sigma$ と整合する同様の環準同型
$\sigma' : \kappa_{A'} \to A'$ を選べる。
\end{lemma}

\begin{proof}
Algebra, Proposition
\ref{algebra-proposition-characterize-separable-field-extensions}
により、分離拡大は形式的に滑らかである。
従って $\sigma$ の存在は $\mathbf{F}_p \to \kappa_A$ が
分離的であることから従う。$\sigma$ と整合する $\sigma'$ の存在も同様である。
\end{proof}

\begin{lemma}
\label{more-algebra-lemma-pre-characteristic-p-case}
$A$ を、標数 $p > 0$ の分数体 $K$ をもつ離散付値環とする。
$\xi \in K$ とする。$L$ を $z^p - z = \xi$ の根を添加して得られる
$K$ の拡大とする。このとき $L/K$ は Galois 拡大であり、
次のいずれかが成り立つ。
\begin{enumerate}
\item $L = K$,
\item $L/K$ は $A$ に関して不分岐な次数 $p$ の拡大である、
\item $L/K$ は $A$ に関して完全分岐し、分岐指数は $p$ である、または
\item $L$ における $A$ の整閉包 $B$ は離散付値環であり、
$A \subset B$ は弱不分岐で、$A \to B$ は次数 $p$ の
純非分離な剰余体拡大を誘導する。
\end{enumerate}
$\pi$ を $A$ の一意化元とする。次の含意が成り立つ。
\begin{enumerate}
\item[(A)] $\xi \in A$ ならば、場合 (1) または (2) である。
\item[(B)] $n > 0$ は $p$ で割り切れず、$a$ は $A$ の単元で、
$\xi = \pi^{-n}a$ ならば、場合 (3) である。
\item[(C)] $n > 0$ は $p$ で割り切れ、$a$ の $\kappa_A$ における像は
$p$ 乗元でなく、$\xi = \pi^{-n} a$ ならば、場合 (4) である。
\end{enumerate}
\end{lemma}

\begin{proof}
Fields, Section \ref{fields-section-Artin-Schreier} の議論により、
この拡大は位数が $p$ を割り切る Galois 拡大である。
Section \ref{more-algebra-section-ramification} の議論から、補題に列挙した
場合 (1) -- (4) のいずれかであることが直ちに従う。

\medskip\noindent
場合 (A)。ここでは $A \to A[x]/(x^p - x - \xi)$ は
有限 \'etale 環拡大である。従って場合 (1) または (2) である。

\medskip\noindent
場合 (B)。$p$ が $n$ を割り切らないように
$\xi = \pi^{-n}a$ と書く。$B \subset L$ を
$L$ における $A$ の整閉包とする。
ある極大イデアル $\mathfrak m$ に対して
$C = B_\mathfrak m$ ならば、
$p \text{ord}_C(z) = -n \text{ord}_C(\pi)$ は明らかである。
特に $A \subset C$ の分岐指数は $p$ で割り切れる。
従ってそれは $p$ であり、$B = C$ である。

\medskip\noindent
場合 (C)。$k = n/p$ とおく。このとき方程式を
$$
(\pi^kz)^p - \pi^{n - k} (\pi^kz) = a
$$
と書き直せる。
$A[y]/(y^p - \pi^{n - k}y - a)$ は
$A$ 上弱不分岐な離散付値環なので、補題が従う。
\end{proof}

\begin{lemma}
\label{more-algebra-lemma-characteristic-p-case}
$A \subset B \subset C$ を、分数体 $K \subset L \subset M$ をもつ
離散付値環の拡大とする。次を仮定する。
\begin{enumerate}
\item $A \subset B$ は弱不分岐である、
\item $K$ の標数は $p$ である、
\item $M$ は $L$ の次数 $p$ の Galois 拡大である、さらに
\item $\kappa_A = \bigcap_{n \geq 1} \kappa_B^{p^n}$.
\end{enumerate}
このとき、$A$ に関して完全分岐し、
$A \to C$ に対する弱解である有限 Galois 拡大 $K_1/K$ が存在する。
\end{lemma}

\begin{proof}
$L$ の標数は $p$ なので、$M$ は $L$ の Artin--Schreier 拡大である
（Fields, Lemma \ref{fields-lemma-Artin-Schreier}）。
従って $z \in M$, $z \not \in L$ で
$\xi = z^p - z \in L$ となるものを選べる。
$\pi^n\xi \in B$ となる $n \geq 0$ を選ぶ。
この $n$ が最小となるように $z$ を選ぶ。$n = 0$ ならば、
Lemma \ref{more-algebra-lemma-pre-characteristic-p-case} により
$M/L$ は $B$ に関して不分岐であり、証明は終わる。
従って $n > 0$ である。

\medskip\noindent
仮定 (4) は $\kappa_A$ が完全であることを意味する。
Lemma \ref{more-algebra-lemma-cohen} のように、整合する環準同型
$\overline{\sigma} : \kappa_A \to A/\pi^n A$ と
$\overline{\sigma} : \kappa_B \to B/\pi^n B$ を選べる。
二番目の写像を集合の写像
$\sigma : \kappa_B \to B$\footnote{$B$ が完備ならば、
$\sigma$ を環準同型として選べる。$A$ も完備で $\sigma$ が環準同型ならば、
$\sigma$ は $\kappa_A$ を $A$ に写す。}
に持ち上げる。このとき
$$
\xi = \sum\nolimits_{i = n, \ldots, 1} \sigma(\lambda_i) \pi^{-i} + b
$$
と書ける。ここで $\lambda_i \in \kappa_B$, $b \in B$ である。
$$
I = \{i \in \{n, \ldots, 1\} \mid \lambda_i \in \kappa_A\}
$$
および
% P02-E1279: source line 34859 uses lambda_i in the defining condition for the set indexed by j.
$$
J = \{j \in \{n, \ldots, 1\} \mid \lambda_i \not \in \kappa_A\}
$$
とおく。有限集合 $J$ の大きさに関する帰納法で論じる。

\medskip\noindent
$J = \emptyset$ の場合。このときすべての
$i \in \{n, \ldots, 1\}$ に対し、$\sigma$ の選び方により、
ある $a_i \in A$, $b_i \in B$ に対して
$\sigma(\lambda_i) = a_i + \pi^n b_i$ である。
従ってある $a \in A$, $b \in B$ に対して
$\xi = \pi^{-n} a + b$ である。
$p | n$ ならば、$A$ の剰余体は完全なので、
ある $a_0, a_1 \in A$ に対して $a = a_0^p + \pi a_1$ と書ける。
ある $b' \in B$ に対し
$$
(z - \pi^{-n/p}a_0)^p - (z - \pi^{-n/p}a_0) =
\pi^{-(n - 1)}(a_1 + \pi^{n - 1 - n/p}a_0) + b'
$$
と計算できる。これは $n$ の最小性に反する。
従って $p$ は $n$ を割り切らない。
$w^p - w = \pi^{-n}a$ で与えられる $K$ の次数 $p$ の拡大を
$K_1$ とする。Lemma \ref{more-algebra-lemma-pre-characteristic-p-case} により、
この拡大は Galois であり、$A$ に関して完全分岐する。
従って $L_1 = L \otimes_K K_1$ は体で、
$A_1 \subset B_1$ は弱不分岐である
（Lemma \ref{more-algebra-lemma-weakly-unramified-goes-up-along-totally-ramified}）。
Lemma \ref{more-algebra-lemma-pre-characteristic-p-case} により、
環 $M_1 = M \otimes_K K_1$ は $L_1$ の $p$ 個のコピーの直積
（この場合は証明終了）であるか、$L_1$ の次数 $p$ の体拡大である。
さらに後者の場合、$C_1$ が $B_1$ 上弱不分岐
（この場合も証明終了）であるか、$M_1/L_1$ が次数 $p$ の Galois 拡大で、
$B_1$ に関して完全分岐する。この最後の場合、拡大 $M_1/L_1$ は
元 $z - w$ によって生成され、
$$
(z - w)^p - (z - w) = z^p - z - (w^p - w) = b
$$
である。ここで $b \in B$ である（上記参照）。
もう一度 Lemma \ref{more-algebra-lemma-pre-characteristic-p-case} を用いると、
拡大 $M_1/L_1$ は $B_1$ に関して不分岐である。
従って $K_1$ は $A \to C$ に対する弱解である。
以下では $J \not = \emptyset$ と仮定する。

\medskip\noindent
ある $r > 0$ に対して $j' = p^r j$ となる
$j', j \in J$ が存在すると仮定する。このとき $z$ の選択を
$$
z' = z -
(\sigma(\lambda_j) \pi^{-j} + \sigma(\lambda_j^p) \pi^{-pj} + \ldots +
\sigma(\lambda_j^{p^{r - 1}}) \pi^{-p^{r - 1}j})
$$
に変更する。すると $\xi$ は次のように
$\xi' = (z')^p - (z')$ に変わる。
$$
\xi' =
\xi - \sigma(\lambda_j) \pi^{-j} + \sigma(\lambda_j^{p^r}) \pi^{-j'}
+ \text{something in }B
$$
前と同様に
$\xi' = \sum\nolimits_{i = n, \ldots, 1}
\sigma(\lambda'_i) \pi^{-i} + b'$
と書けば、$i \not = j, j'$ に対して
$\lambda'_i = \lambda_i$ であり、$\lambda'_j = 0$ である。
従って集合 $J$ は小さくなった。
$J$ の大きさに関する帰納法により、このような組 $j, j'$ は
存在しないと仮定してよい
（この手続きの途中で、上で扱った $J = \emptyset$ の場合に
戻ることがあり得ることに注意されたい）。

\medskip\noindent
$j \in J$ に対し、$r_j \geq 0$ と、$p$ 乗元でない
$\mu_j \in \kappa_B$ を用いて
$\lambda_j = \mu_j^{p^{r_j}}$ と書く。仮定 (4) によりこれは可能である。
$j p^{-r_j}$ が最大となる唯一の添字 $j \in J$ をとる
（前段落の結果により添字は一意である）。
すべての $j \in J$ に対して $j p^{r - r_j} > n$ となるように、
$r > \max(r_j + 1)$ を選ぶ。
$A$ に関して完全分岐する次数 $p^r$ の分離拡大 $K_1/K$ で、
対応する離散付値環 $A_1 \subset K_1$ が一意化元 $\pi'$ をもち、
ある $a \in A_1$ に対して
$(\pi')^{p^r} = \pi + \pi^{n + 1}a$ となるものを選ぶ
（Lemma \ref{more-algebra-lemma-make-degree-q-extension}）。
$L_1 = L \otimes_K K_1$ は体であり、$L_1/L$ は $B$ に関して
完全分岐することに注意する
（Lemma \ref{more-algebra-lemma-weakly-unramified-goes-up-along-totally-ramified}）。
整閉包 $B_1$ で計算すると、
$$
\xi = \sum\nolimits_{i \in I} \sigma(\lambda_i) (\pi')^{-i p^r} +
\sum\nolimits_{j \in J} \sigma(\mu_j)^{p^{r_j}} (\pi')^{-j p^r} + b_1
$$
を得る。ここで $b_1 \in B_1$ である。
% P02-E1280: source line 34940 calls this a qth power although q is undefined; the chosen exponent is p^r.
$i \in I$ に対する $\sigma(\lambda_i)$ は $\pi^n$ を法として、
すなわち $(\pi')^{n p^r}$ を法として $q$ 乗元であることに注意する。
従って上式を
$$
\xi = \sum\nolimits_{i \in I} x_i^{p^r} (\pi')^{-i p^r} +
\sum\nolimits_{j \in J} \sigma(\mu_j)^{p^{r_j}} (\pi')^{- j p^r}
+ b_1
$$
と書き直せる。前段落と同様に $z$ の選択を
\begin{align*}
z' & = z \\
& -
\sum\nolimits_{i \in I}
\left(x_i (\pi')^{-i} + \ldots + x_i^{p^{r - 1}} (\pi')^{-i p^{r - 1}}\right)
\\
& -
\sum\nolimits_{j \in J}
\left(
\sigma(\mu_j) (\pi')^{- j p^{r - r_j}}
+ \ldots +
% P02-E1281: the displayed terminal exponent p^(r_j-1) is invalid when the allowed value r_j = 0.
\sigma(\mu_j)^{p^{r_j - 1}} (\pi')^{- j p^{r - 1}}
\right)
\end{align*}
に変更すると、
$$
(z')^p - z' =
\sum\nolimits_{i \in I} x_i (\pi')^{-i} +
\sum\nolimits_{j \in J} \sigma(\mu_j) (\pi')^{- j p^{r - r_j}} + b_1'
$$
を得る。ここで $b'_1 \in B_1$ である。
$j p^{r - r_j}$ が最大となる $j$ は一意であり、
$j p^{r - r_j}$ は $i \in I$ より大きく $p$ で割り切れるので、
$M_1 / L_1$ は Lemma \ref{more-algebra-lemma-pre-characteristic-p-case} の
場合 (C) に属する。これで証明は完了する。
\end{proof}

\begin{lemma}
\label{more-algebra-lemma-prepare}
$A$ を原始 $p$ 乗根 $\zeta$ を含む環とし、$w = 1 - \zeta$ とおく。
このとき
$$
P(z) = \frac{(1 + wz)^p - 1}{w^p} =
z^p - z + \sum\nolimits_{0 < i < p} a_i z^i
$$
は $A[z]$ の元であり、実際 $a_i \in (w)$ である。さらに多項式環
$A[z_1, z_2]$ において
$$
P(z_1 + z_2 + w z_1 z_2) = P(z_1) + P(z_2) + w^p P(z_1) P(z_2)
$$
が成り立つ。
\end{lemma}

\begin{proof}
円分体の整数環
$$
A = \mathbf{Z}[\zeta] =
\mathbf{Z}[x]/(x^{p - 1} + \ldots + x + 1)
$$
の場合を証明すれば十分である。多項式恒等式
$t^p - 1 = (t - 1)(t - \zeta) \ldots (t - \zeta^{p - 1})$
（両辺の根を見ることで証明できる）は
$t^{p - 1} + \ldots + t + 1 =
(t - \zeta) \ldots (t - \zeta^{p - 1})$
を示す。$t = 1$ を代入すると
$p = (1 - \zeta)(1 - \zeta^2) \ldots (1 - \zeta^{p - 1})$
を得る。極大イデアル $(p, w) = (w)$ は $p$ の上にある
$A$ の唯一の素イデアルである
（標数 $p$ の体は非自明な $1$ の $p$ 乗根をもたない）。
従ってある単元 $u$ に対して $p = u w^{p - 1}$ である。
これは、$p > i > 1$ に対して
$$
a_i = \frac{1}{p} {p \choose i} u w^{i - 1}
$$
であり、$- 1 + a_1 = pw/w^p = u$ であることを意味する。
$P(-1) = 0$ なので、$(w)$ を法として $0 = (-1)^p - u$ である。
従って $a_1 \in (w)$ であり、
% P02-E1282: source line 35012 says “the proof if the first part” instead of “of”.
第一部分の証明は終わる。第二部分は、省略する直接計算から従う。
\end{proof}

\begin{lemma}
\label{more-algebra-lemma-extension-defined-by-nice-polynial}
$A$ を原始 $1$ の $p$ 乗根を含む混標数 $(0, p)$ の離散付値環とする。
$P(t) \in A[t]$ を Lemma \ref{more-algebra-lemma-prepare} の多項式とする。
$\xi \in K$ とし、$L$ を $P(z) = \xi$ の根を添加して得られる
$K$ の拡大とする。このとき $L/K$ は Galois 拡大であり、
次のいずれかが成り立つ。
\begin{enumerate}
\item $L = K$,
\item $L/K$ は $A$ に関して不分岐な次数 $p$ の拡大である、
\item $L/K$ は $A$ に関して完全分岐し、分岐指数は $p$ である、または
\item $L$ における $A$ の整閉包 $B$ は離散付値環であり、
$A \subset B$ は弱不分岐で、$A \to B$ は次数 $p$ の
純非分離な剰余体拡大を誘導する。
\end{enumerate}
$\pi$ を $A$ の一意化元とする。次の含意が成り立つ。
\begin{enumerate}
\item[(A)] $\xi \in A$ ならば、場合 (1) または (2) である。
\item[(B)] $n > 0$ は $p$ で割り切れず、$a$ は $A$ の単元で、
$\xi = \pi^{-n}a$ ならば、場合 (3) である。
\item[(C)] $n > 0$ は $p$ で割り切れ、$a$ の $\kappa_A$ における像は
$p$ 乗元でなく、$\xi = \pi^{-n} a$ ならば、場合 (4) である。
\end{enumerate}
\end{lemma}

\begin{proof}
$P(z) = \xi$ の根を添加することは、
$y^p = w^p(1 + \xi)$ の根を添加することと同じである。
$K$ は原始 $1$ の $p$ 乗根を含むので、Fields, Section
\ref{fields-section-Kummer} の議論により、この拡大は位数が
$p$ を割り切る Galois 拡大である。
Section \ref{more-algebra-section-ramification} の議論から、補題に列挙した
場合 (1) -- (4) のいずれかであることが直ちに従う。

\medskip\noindent
場合 (A)。ここでは $A \to A[x]/(P(x) - \xi)$ は
有限 \'etale 環拡大である。従って場合 (1) または (2) である。

\medskip\noindent
場合 (B)。$p$ が $n$ を割り切らないように
$\xi = \pi^{-n}a$ と書く。$B \subset L$ を
$L$ における $A$ の整閉包とする。
ある極大イデアル $\mathfrak m$ に対して
$C = B_\mathfrak m$ ならば、
$p \text{ord}_C(z) = -n \text{ord}_C(\pi)$ は明らかである。
特に $A \subset C$ の分岐指数は $p$ で割り切れる。
従ってそれは $p$ であり、$B = C$ である。

\medskip\noindent
場合 (C)。$k = n/p$ とおく。このとき方程式を
$$
(\pi^kz)^p - \pi^{n - k} (\pi^kz) +
\sum a_i \pi^{n - ik} (\pi^kz)^i = a
$$
と書き直せる。
% P02-E1283: source line 35070 changes the displayed plus sign before the a_i-sum to a minus in the quotient polynomial.
$A[y]/(y^p - \pi^{n - k}y - \sum  a_i \pi^{n - ik} y^i - a)$
は $A$ 上弱不分岐な離散付値環なので、補題が従う。
\end{proof}

\noindent
$A$ を原始 $1$ の $p$ 乗根を含む混標数 $(0, p)$ の
離散付値環とする。$w \in A$ と $P(t) \in A[t]$ を
Lemma \ref{more-algebra-lemma-prepare} のようにとる。$L$ を $K$ の有限拡大とする。
$L$ が、$w^p \xi \in \mathfrak m_A$ を満たす元 $\xi \in K$ に対する
方程式 $P(z) = \xi$ の根を添加して得られる $K$ の次数 $p$ の拡大ならば、
$L/K$ を {\it 有限レベルの次数 $p$ の拡大} という。

\medskip\noindent
次の直接的な補題により、この定義はこの節の議論に関係する。

\begin{lemma}
\label{more-algebra-lemma-make-finite-level}
$A \subset B \subset C$ を、分数体 $K \subset L \subset M$ をもつ
離散付値環の拡大とする。次を仮定する。
\begin{enumerate}
\item $A$ は混標数 $(0, p)$ をもつ、
\item $A \subset B$ は弱不分岐である、
\item $B$ は原始 $1$ の $p$ 乗根を含む、さらに
\item $M/L$ は次数 $p$ の Galois 拡大である。
\end{enumerate}
このとき、$A$ に関して完全分岐する有限 Galois 拡大 $K_1/K$ で、
$A \to C$ に対する弱解であるか、または $M_1/L_1$ が
有限レベルの次数 $p$ の拡大となるものが存在する。
\end{lemma}

\begin{proof}
$\pi \in A$ を一意化元とする。Kummer 理論
（Fields, Lemma \ref{fields-lemma-Kummer}）により、
$M$ はある $b \in L$ に対して $y^p = b$ の根を
$L$ に添加して得られる。

\medskip\noindent
$\text{ord}_B(b)$ が $p$ と互いに素ならば、
$A$ に関して完全分岐する次数 $p$ の分離拡大 $K_1/K$ を選ぶ
（例えば Lemma \ref{more-algebra-lemma-make-degree-q-extension} を用いる）。
$A_1$ を $K_1$ における $A$ の整閉包とする。
Lemma \ref{more-algebra-lemma-weakly-unramified-goes-up-along-totally-ramified}
により、$L_1 = L \otimes_K K_1$ における $B$ の整閉包 $B_1$ は、
$A_1$ 上弱不分岐な離散付値環である。
$K_1/K$ が $A \to C$ に対する弱解でなければ、
$M_1 = M \otimes_K K_1$ における $C$ の整閉包 $C_1$ は
離散付値環であり、$B_1 \to C_1$ の分岐指数は $p$ である。
この場合、体 $M_1$ は $b$ の $p$ 乗根を $L_1$ に添加して得られ、
$\text{ord}_{B_1}(b)$ は $p$ で割り切れる。
$A$ を $A_1$ で置き換えるなどして、$u \in B$ が単元で、
$n$ が $p$ で割り切れ、$b = \pi^n u$ であると仮定してよい。
もちろんこの場合、拡大 $M$ は単元の $p$ 乗根を $L$ に
添加して得られる。

\medskip\noindent
$M$ が $B$ のある単元 $u$ に対する $y^p = u$ の根を
$L$ に添加して得られると仮定する。
$\kappa_B$ における $u$ の剰余類が $p$ 乗元でなければ、
Lemma \ref{more-algebra-lemma-pre-purely-inseparable-case} により
$B \subset C$ は弱不分岐であり、証明は終わる。
そうでなければ、$v^p$ と $u$ が $\kappa_B$ で同じ像をもつように選び、
$y$ の選択を $y/v$ に置き換えられる。
このように置き換えた後、
$$
y^p = 1 + \pi b
$$
と書ける。ここで $b \in B$ である。このとき
$z = (y - 1)/w$ とすれば $P(z) = \pi b/ w^p$ である。
従ってこの拡大は $\xi = \pi b / w^p$ をもつ
有限レベルの次数 $p$ の拡大である。
\end{proof}

\noindent
$A$ を原始 $1$ の $p$ 乗根を含む混標数 $(0, p)$ の
離散付値環とする。$w \in A$ と $P(t) \in A[t]$ を
Lemma \ref{more-algebra-lemma-prepare} のようにとる。
$L$ を $K$ の有限レベルの次数 $p$ の拡大とする。
$L$ を $K$ 上生成し、$\xi = P(z) \in K$ を満たす $z \in L$ を選ぶ。
$A$ の一意化元 $\pi$ を選び、整数
$e_1 = \text{ord}_A(w)$ と単元 $u \in A$ を用いて
$w = u \pi^{e_1}$ と書く。最後に
$$
\pi^n \xi \in A
$$
となる $n \geq 0$ をとる。
% P02-E1284: source lines 35153-35154 say “taking over” instead of “taken over”.
$L/K$ の {\it レベル} とは、$\xi = P(z) \in K$ を満たし
$L/K$ を生成するすべての $z$ にわたって取った量 $n/e_1$ の最小値である。

\medskip\noindent
いくつか注意する。この拡大は有限レベルなので、$n < pe_1$ となるように
$z$ を選べる。従ってレベルは $[0, p)$ に属する有理数である。
レベルが零ならば、Lemma \ref{more-algebra-lemma-extension-defined-by-nice-polynial}
により $L/K$ は $A$ に関して不分岐である。次の目標はレベルを下げることである。

\begin{lemma}
\label{more-algebra-lemma-lowering-the-level}
$A \subset B \subset C$ を、分数体 $K \subset L \subset M$ をもつ
離散付値環の拡大とする。次を仮定する。
\begin{enumerate}
\item $A$ は混標数 $(0, p)$ をもつ、
\item $A \subset B$ は弱不分岐である、
\item $B$ は原始 $1$ の $p$ 乗根を含む、
\item $M/L$ は有限レベル $l > 0$ の次数 $p$ の拡大である、
\item $\kappa_A = \bigcap_{n \geq 1} \kappa_B^{p^n}$.
\end{enumerate}
このとき、$A$ に関して完全分岐する $K$ の有限分離拡大 $K_1$ で、
$K_1$ が $A \to C$ に対する弱解であるか、または拡大 $M_1/L_1$ が
レベル $\leq \max(0, l - 1, 2l - p)$ の有限レベルの
次数 $p$ の拡大となるものが存在する。
\end{lemma}

\begin{proof}
$\pi \in A$ を一意化元とする。
$w \in B$ と $P \in B[t]$ を Lemma \ref{more-algebra-lemma-prepare} のように
（$B$ に対して）とる。$e_1 = \text{ord}_B(w)$ とおく。
従って $w$ と $\pi^{e_1}$ は $B$ で同伴である。
$M$ を $L$ 上生成し、$\xi = P(z) \in K$ を満たす $z \in M$ と、
$M$ の $L$ 上のレベルの定義のように $\pi^n\xi \in B$ を満たす $n$ を選ぶ。
すなわち $l = n/e_1$ である。

\medskip\noindent
この補題の証明は Lemma \ref{more-algebra-lemma-characteristic-p-case} の証明と
完全に同様である。何が起きているかを説明するため、$z \in L$ が
$\pi^{-n} P(z) \in B$ を満たすならば
（$z$ の付値が最低でも $-n/p$ であることと多項式 $P$ の形を用いて）、
\begin{equation}
\label{more-algebra-equation-first-congruence}
P(z) \equiv z^p - z \bmod \pi^{-n + e_1}B
\end{equation}
であることに注意する。さらに、$\xi_1, \xi_2 \in \pi^{-n}B$ に対して
\begin{equation}
\label{more-algebra-equation-second-congruence}
\xi_1 + \xi_2 + w^p \xi_1 \xi_2 \equiv
\xi_1 + \xi_2 \bmod \pi^{-2n + pe_1}B
\end{equation}
である。最後に
% P02-E1285: source line 35205 divides by e_1 where l = n/e_1 requires multiplication.
$n - e_1 = (l - 1)/e_1$ および
$-2n + pe_1 = -(2l - p)e_1$ であることに注意する。
$m = n - e_1 \max(0, l - 1, 2l - p)$ とおく。
上の式は、$\pi^{-n}B / \pi^{-n + m}B$ で計算すると、
多項式 $P$ が多項式 $z^p -  z$ とまったく同じように振る舞うことを示す。
これは補題が成り立つ理由を説明するが、以下で詳細も与える。

\medskip\noindent
% P02-E1286: source line 35212 cites Assumption (4), but perfection follows from item (5).
仮定 (4) は $\kappa_A$ が完全であることを意味する。
$m \leq e_1$ であり、従って $A/\pi^m$ は $w$ で、ゆえに $p$ で
零化されることに注意する。そこで Lemma \ref{more-algebra-lemma-cohen} のように、
整合する環準同型
$\overline{\sigma} : \kappa_A \to A/\pi^mA$ と
$\overline{\sigma} : \kappa_B \to B/\pi^mB$ を選べる。
二番目の写像を集合の写像 $\sigma : \kappa_B \to B$ に持ち上げる。
このとき
% P02-E1287: source line 35222 has an unmatched closing parenthesis in the exponent -n + m).
$$
\xi =
\sum\nolimits_{i = n, \ldots, n - m + 1} \sigma(\lambda_i) \pi^{-i} +
\pi^{-n + m)} b
$$
と書ける。ここで $\lambda_i \in \kappa_B$, $b \in B$ である。
$$
I = \{i \in \{n, \ldots, n - m + 1\} \mid \lambda_i \in \kappa_A\}
$$
および
% P02-E1288: source line 35230 uses lambda_i in the defining condition for the set indexed by j.
$$
J = \{j \in \{n, \ldots, n - m + 1\} \mid \lambda_i \not \in \kappa_A\}
$$
とおく。有限集合 $J$ の大きさに関する帰納法で論じる。

\medskip\noindent
$J = \emptyset$ の場合。このときすべての
$i \in \{n, \ldots, n - m + 1\}$ に対し、
$\overline{\sigma}$ の選び方により、ある $a_i \in A$, $b_i \in B$ に対して
$\sigma(\lambda_i) = a_i + \pi^{n - m}b_i$ である。
従ってある $a \in A$, $b \in B$ に対して
$\xi = \pi^{-n} a + \pi^{-n + m} b$ である。
$p | n$ ならば、$A$ の剰余体は完全なので、ある
$a_0, a_1 \in A$ に対して $a = a_0^p + \pi a_1$ と書ける。
$z_1 = - \pi^{-n/p} a_0$ とおく。
$P(z_1) \in \pi^{-n}B$ であり、
$z + z_1 + w z z_1$ は $M$ を $L$ 上生成する元であることに注意する
（$n < pe_1$ なので $wz_1 \not = -1$ である）。
さらに Lemma \ref{more-algebra-lemma-prepare} により
$$
P(z + z_1 + w z z_1) = P(z) + P(z_1) + w^p P(z) P(z_1) \in K
$$
であり、式 (\ref{more-algebra-equation-first-congruence}) と
(\ref{more-algebra-equation-second-congruence}) により
$$
P(z) + P(z_1) + w^p P(z) P(z_1)
\equiv
\xi + z_1^p - z_1 \bmod \pi^{-n + m}B
$$
である。
% P02-E1289: source line 35254 appends “for some b' in B” although no b' occurs in the congruence.
ここである $b' \in B$ をとる。
% P02-E1290: source line 35254 says “This contradict” instead of “contradicts”.
これは $n$ の最小性に反する！従って $p$ は $n$ を割り切らない。
$P(y) = -\pi^{-n}a$ で与えられる $K$ の次数 $p$ の拡大を $K_1$ とする。
Lemma \ref{more-algebra-lemma-extension-defined-by-nice-polynial} により、
この拡大は分離的であり、$A$ に関して完全分岐する。
従って $L_1 = L \otimes_K K_1$ は体で、$A_1 \subset B_1$ は
弱不分岐である
（Lemma \ref{more-algebra-lemma-weakly-unramified-goes-up-along-totally-ramified}）。
Lemma \ref{more-algebra-lemma-extension-defined-by-nice-polynial} により、
環 $M_1 = M \otimes_K K_1$ は $L_1$ の $p$ 個のコピーの直積
（この場合は証明終了）であるか、$L_1$ の次数 $p$ の体拡大である。
さらに後者の場合、$C_1$ が $B_1$ 上弱不分岐
（この場合も証明終了）であるか、$M_1/L_1$ が次数 $p$ の Galois 拡大で、
$B_1$ に関して完全分岐する。この最後の場合、拡大 $M_1/L_1$ は
元 $z + y + wzy$ によって生成され、$P(z + y + wzy) \in L_1$ であり、
\begin{align*}
P(z + y + wzy)
& = P(z) + P(y) + w^p P(z) P(y) \\
& \equiv
\xi - \pi^{-n}a \bmod \pi^{-n + m}B_1 \\
& \equiv
0 \bmod \pi^{-n + m}B_1
\end{align*}
であることが、上とまったく同じ方法で分かる。
$m$ の選び方により、これはちょうど $M_1/L_1$ のレベルが
高々 $\max(0, l - 1, 2l - p)$ であることを意味する。
以下では $J \not = \emptyset$ と仮定する。

\medskip\noindent
ある $r > 0$ に対して $j' = p^r j$ となる
$j', j \in J$ が存在すると仮定する。このとき
$$
z_1 = - \sigma(\lambda_j) \pi^{-j} - \sigma(\lambda_j^p) \pi^{-pj} -
\ldots - \sigma(\lambda_j^{p^{r - 1}}) \pi^{-p^{r - 1}j}
$$
とおき、$z$ を $z' = z + z_1 + wzz_1$ に変更する。
$z' \in M$ は $M$ を $L$ 上生成し、
% P02-E1291: source line 35291 uses w instead of the w^p factor in Lemma prepare's group law.
$\xi' = P(z') = P(z) + P(z_1) + wP(z)P(z_1) \in L$ であり、
上と同様に式 (\ref{more-algebra-equation-first-congruence}) と
(\ref{more-algebra-equation-second-congruence}) を用いると
$$
\xi' \equiv
\xi - \sigma(\lambda_j) \pi^{-j} + \sigma(\lambda_j^{p^r}) \pi^{-j'}
\bmod \pi^{-n + m}B
$$
となることに注意する。前と同様に
$$
\xi' =
\sum\nolimits_{i = n, \ldots, n - m + 1} \sigma(\lambda'_i) \pi^{-i}
+ \pi^{-n + m}b'
$$
と書けば、$i \not = j, j'$ に対して
$\lambda'_i = \lambda_i$ であり、$\lambda'_j = 0$ である。
従って集合 $J$ は小さくなった。
$J$ の大きさに関する帰納法により、$J$ の元の組 $j, j'$ で
$j'/j$ が $p$ の冪となるものは存在しないと仮定してよい
（この手続きの途中で、上で扱った $J = \emptyset$ の場合に
戻ることがあり得ることに注意されたい）。

\medskip\noindent
$j \in J$ に対し、$r_j \geq 0$ と、$p$ 乗元でない
$\mu_j \in \kappa_B$ を用いて
$\lambda_j = \mu_j^{p^{r_j}}$ と書く。仮定 (4) によりこれは可能である。
$j p^{-r_j}$ が最大となる唯一の添字 $j \in J$ をとる
（前段落の結果により添字は一意である）。
すべての $j \in J$ に対して $j p^{r - r_j} > n$ となるように、
$r > \max(r_j + 1)$ を選ぶ。
$(\pi')^{p^r} = \pi$ で定義される、$A$ に関して完全分岐する
次数 $p^r$ の拡大を $K_1/K$ とする。
$\pi'$ は、対応する離散付値環 $A_1 \subset K_1$ の
一意化元であることに注意する。
$L_1 = L \otimes_K K_1$ は体であり、$L_1/L$ は $B$ に関して
完全分岐することに注意する
（Lemma \ref{more-algebra-lemma-weakly-unramified-goes-up-along-totally-ramified}）。
整閉包 $B_1$ で計算すると、
$$
\xi = \sum\nolimits_{i \in I} \sigma(\lambda_i) (\pi')^{-i p^r} +
\sum\nolimits_{j \in J} \sigma(\mu_j)^{p^{r_j}} (\pi')^{-j p^r} +
\pi^{-n + m} b_1
$$
を得る。ここで $b_1 \in B_1$ である。
% P02-E1292: source line 35331 calls this a qth power although q is undefined; the chosen exponent is p^r.
$i \in I$ に対する $\sigma(\lambda_i)$ は $\pi^m$ を法として、
すなわち $(\pi')^{m p^r}$ を法として $q$ 乗元であることに注意する。
従って上式を
$$
\xi = \sum\nolimits_{i \in I} x_i^{p^r} (\pi')^{-i p^r} +
\sum\nolimits_{j \in J} \sigma(\mu_j)^{p^{r_j}} (\pi')^{- j p^r}
+ \pi^{-n + m}b_1
$$
と書き直せる。前段落の選択と同様に
% P02-E1293: source line 35340 has “z_1 & -” and omits the equality sign in the align environment.
\begin{align*}
z_1 & - \sum\nolimits_{i \in I}
\left(x_i (\pi')^{-i} + \ldots + x_i^{p^{r - 1}} (\pi')^{-i p^{r - 1}}\right)
\\
& - \sum\nolimits_{j \in J}
\left(
\sigma(\mu_j) (\pi')^{- j p^{r - r_j}}
+ \ldots +
% P02-E1294: the displayed terminal exponent p^(r_j-1) is invalid when the allowed value r_j = 0.
\sigma(\mu_j)^{p^{r_j - 1}} (\pi')^{- j p^{r - 1}}
\right)
\end{align*}
とおき、$z$ の選択を $z' = z + z_1 + wzz_1$ に変更する。
すると $z'$ は $M_1$ を $L_1$ 上生成し、
$\xi' = P(z') = P(z) + P(z_1) + w^p P(z) P(z_1) \in L_1$
であり、計算により
$$
\xi' \equiv
\sum\nolimits_{i \in I} x_i (\pi')^{-i} +
\sum\nolimits_{j \in J} \sigma(\mu_j) (\pi')^{- j p^{r - r_j}} +
(\pi')^{(-n + m)p^r}b'_1
$$
であることが分かる。ここで $b'_1 \in B_1$ である。
$j p^{r - r_j}$ が最大となる $j$ は一意であり、
$j p^{r - r_j}$ は $i \in I$ より大きい。
$j p^{r - r_j} \leq (n - m)p^r$ ならば、拡大 $M_1/L_1$ のレベルは
$\max(0, l - 1, 2l - p)$ より小さい。
そうでなければ $p$ が $j p^{r - r_j}$ を割り切るので、
$M_1 / L_1$ は
Lemma \ref{more-algebra-lemma-extension-defined-by-nice-polynial} の場合 (C) に属する。
これで証明は完了する。
\end{proof}

\begin{lemma}
\label{more-algebra-lemma-special-case}
$A \subset B \subset C$ を、分数体 $K \subset L \subset M$ をもつ
離散付値環の拡大とする。次を仮定する。
\begin{enumerate}
\item $A$ の剰余体 $k$ は標数 $p > 0$ の代数閉体である、
\item $A$ と $B$ は完備である、
\item $A \to B$ は弱不分岐である、
\item $M$ は $L$ の有限拡大である、
\item $k = \bigcap\nolimits_{n \geq 1} \kappa_B^{p^n}$.
\end{enumerate}
このとき、$A \to C$ に対する弱解である有限拡大 $K_1/K$ が存在する。
\end{lemma}

\begin{proof}
$M'$ を $L$ の任意の有限拡大とし、$C'$ を
$M'$ における $B$ の整閉包とする。
$B$ は Algebra, Lemma
\ref{algebra-lemma-Noetherian-complete-local-Nagata} により
Nagata 環なので、$C'$ は $B$ 上有限である。
さらに Remark \ref{more-algebra-remark-construction} の議論により、
$C'$ は離散付値環である。また $C'$ は $B$-加群として完備であり、
従って離散付値環として完備である
（Algebra, Section \ref{algebra-section-completion}）。
特に $C$ は $M$ における $B$ の整閉包であることが従う
（優越関係について付値環が極大であるという定義による）。

\medskip\noindent
$M \subset M'$ を有限拡大とし、$C' \subset M'$ を上のような
$B$ の整閉包とする。Lemma \ref{more-algebra-lemma-solutions-go-down} により、
$A \to B \to C'$ に対して結果を証明すれば十分である。
従って Fields, Lemma \ref{fields-lemma-normal-closure} により、
$M/L$ は正規拡大であると仮定してよい。

\medskip\noindent
$M / L$ が正規ならば、有限拡大の鎖
$$
L = L^0 \subset L^1 \subset L^2 \subset \ldots \subset L^r = M
$$
で、各拡大 $L^{j + 1}/L^j$ が次のいずれかとなるものを見つけられる。
\begin{enumerate}
\item[(a)] 次数 $p$ の純非分離拡大、
\item[(b)] $B^j$ に関して完全分岐する次数 $p$ の Galois 拡大、
\item[(c)] $B^j$ に関して完全分岐し、位数が $p$ と互いに素な
Galois 巡回拡大、
\item[(d)] $B^j$ に関して不分岐な Galois 拡大。
\end{enumerate}
ここで $B^j$ は $L^j$ における $B$ の整閉包である。
実際、$M/L$ は正規なので、Galois 拡大と純非分離拡大の合成として書ける
（Fields, Lemma \ref{fields-lemma-normal-case}）。
純非分離拡大についてフィルトレーションの存在は明らかである。
Galois の場合、$G$ は「その」分解群であることに注意し、
$I \subset G$ を惰性群とする。一方では Lemma
\ref{more-algebra-lemma-galois-inertia} により $I$ は可解であり、他方では Lemma
\ref{more-algebra-lemma-inertial-invariants-unramified} により拡大 $M^I/L$ は
$B$ に関して不分岐である。これで、主張したフィルトレーションが
存在することが証明された。

\medskip\noindent
整数 $r$ に関する帰納法で論じる。
$B^1$ を $L^1$ における $B$ の整閉包として、
$A \to B^1$ に対する弱解である有限拡大 $K_1/K$ を見つけられたとする。
$K'_1$ を $K_1/K$ の正規閉包とする
（Fields, Lemma \ref{fields-lemma-normal-closure}）。
$A$ は完備で、$A$ の剰余体は代数閉体なので、
$K'_1/K_1$ は $A_1$ に関して完全分岐する分離拡大である
（詳細の一部は省略する）。
従って Lemma
\ref{more-algebra-lemma-weakly-unramified-goes-up-along-totally-ramified} により、
$K'_1/K$ も $A \to B^1$ に対する弱解である。
言い換えれば、$K_1$ は $K$ の正規拡大であると仮定してよく、そうする。
その上で列
$$
L^0_1 = (L^0 \otimes_K K_1)_{red} \subset
L^1_1 = (L^1 \otimes_K K_1)_{red} \subset \ldots \subset
L^r_1 = (L^r \otimes_K K_1)_{red}
$$
と、対応する整閉包 $B^i_1$ を考える。
$C_1 = B^r_1$ は離散付値環の直積であり、
Lemma \ref{more-algebra-lemma-galois-relative} により
$G = \text{Aut}(K_1/K)$ がその因子を推移的に置換することに注意する。
特に離散付値環の拡大 $A_1 \to (C_1)_\mathfrak m$ はすべて同型であり、
一つに対する弱解はすべてに対する弱解となる。
列
$$
A_1 \to (B^1_1)_{B^1_1 \cap \mathfrak m} \to
(B^2_1)_{B^2_1 \cap \mathfrak m} \to
\ldots \to
(B^r_1)_{B^r_1 \cap \mathfrak m} =
(C_1)_\mathfrak m
$$
に帰納法の仮定を適用すると、
$A_1 \to (C_1)_\mathfrak m$ に対する弱解 $K_2/K_1$ を得る。
そのとき、前述のことから拡大 $K_2/K$ は $A \to C$ に対する弱解である。
帰納法の仮定を適用できることに注意する。
環準同型 $A_1 \to (B^1_1)_{B^1_1 \cap \mathfrak m}$ は
$K_1$ の選び方により弱不分岐であり、
% P02-E1295: source lines 35467-35468 use singular “sequence ... each still have” instead of “extensions ... each still have”.
分数体拡大の列は、それぞれ上に列挙した性質 (a), (b), (c), (d) の
いずれかをなおもつ。さらに、任意の有限拡大
$\kappa_B \subset \kappa$ に対して
$k = \bigcap \kappa^{p^n}$ がなお成り立つことに注意する。

\medskip\noindent
従って、すべては体拡大 $M/L$ が性質 (a), (b), (c), (d) の
いずれかを満たす場合に $A \subset C$ に対する弱解を
見つけることに帰着する。

\medskip\noindent
場合 (d)。この場合 $B \to C$ は既に不分岐なので自明である。

\medskip\noindent
場合 (c)。$M/L$ は $p$ と互いに素な位数 $n$ の巡回拡大であるとする。
$M/L$ は $B$ に関して完全分岐するので、$B \subset C$ の
分岐指数は $n$ であり、従って $A \subset C$ の分岐指数も $n$ である。
一意化元 $\pi \in A$ を選び、$K_1 = K[\pi^{1/n}]$ とおく。
Abhyankar の補題（Lemma \ref{more-algebra-lemma-abhyankar}）により、
$K_1/K$ は $A \subset C$ に対する解である。

\medskip\noindent
場合 (b)。この場合を混標数の場合と等標数の場合に分ける。
等標数の場合は Lemma \ref{more-algebra-lemma-characteristic-p-case} である。
混標数の場合、まず Lemma \ref{more-algebra-lemma-make-finite-level} を用い、
$K$ を有限拡大で置き換えて、$M/L$ が有限レベルの次数 $p$ の拡大である
状況にする。このときレベルは、Lemma
\ref{more-algebra-lemma-lowering-the-level} の前の議論のように、
有理数 $l \in [0, p)$ である。レベルが $0$ ならば
$B \to C$ は弱不分岐であり、証明は終わる。
そうでなければ、Lemma \ref{more-algebra-lemma-lowering-the-level} により
% P02-E1296: source line 35498 says “can replacing” instead of “can replace”.
体 $K$ を有限拡大で置き換えて、レベル
$l' \leq \max(0, l - 1, 2l - p)$ の新しい状況を得られる。
$\epsilon < 1$ に対して $l = p - \epsilon$ ならば、
$l' \leq p - 2\epsilon$ である。従って有限回の置換の後、
レベルが $\leq p - 1$ の場合を得る。
さらに高々 $p - 1$ 回の同様の置換の後、レベルが零の状況に達する。

\medskip\noindent
場合 (a) は Lemma \ref{more-algebra-lemma-purely-inseparable-case} である。
これは $K$ の純非分離拡大が必要となり得る唯一の場合である。
すなわち補題の主張の (2) の場合には、
% P02-E1297: source lines 35509-35510 says to adjoin a pth power of pi; the argument needs a pth root.
元 $\pi$ の $p$ 乗を添加することで結論を得る。これで補題の証明は完了する。
\end{proof}

\noindent
これで、この節の主結果を証明するために必要な補題がすべて揃った。

\begin{theorem}[Epp]
\label{more-algebra-theorem-epp}
$A \subset B$ を、分数体 $K \subset L$ をもつ離散付値環の拡大とする。
$\kappa_A$ の標数が $p > 0$ ならば、
$$
\bigcap\nolimits_{n \geq 1} \kappa_B^{p^n}
$$
のすべての元が $\kappa_A$ 上分離代数的であると仮定する。
このとき Definition \ref{more-algebra-definition-solution} の意味で
$A \to B$ に対する弱解となる有限拡大 $K_1/K$ が存在する。
\end{theorem}

\begin{proof}
$\kappa_A$ の標数が零であるか、または剰余標数が $p$ で、
分岐指数が $p$ と互いに素で、剰余体拡大が分離的ならば、
これは Abhyankar の補題（Lemma \ref{more-algebra-lemma-abhyankar}）から従う。
実際、分岐指数を $e$ とする。一意化元 $\pi \in A$ を選び、
$\pi$ の $e$ 乗根を添加して得られる拡大を $K_1/K$ とする。
Lemma \ref{more-algebra-lemma-pull-root-uniformizer} により、
$K_1$ における $A$ の整閉包 $A_1$ は
% P02-E1298: source lines 35538-35539 says only “with ramification index over A” and omits the value e.
$A$ 上の分岐指数をもつ離散付値環である。
従って Lemma \ref{more-algebra-lemma-abhyankar} により、$B_1$ のすべての
極大イデアル $\mathfrak m$ に対し
$A_1 \to (B_1)_\mathfrak m$ は $\mathfrak m$-進位相に関して
形式的に滑らかであり、特にこれらは弱不分岐な離散付値環の拡大である。

\medskip\noindent
以下、$p$ を素数とし、$\kappa_A$ の標数が $p$ であると仮定する。
まず Lemma \ref{more-algebra-lemma-solution-after-strict-henselization} を適用し、
$A$ と $B$ の剰余体が分離閉体である場合に帰着する。
この手続きで $\kappa_A$ と $\kappa_B$ はそれぞれの分離代数閉包で
置き換えられるので、定理の条件から
$$
\kappa_A \supset \bigcap\nolimits_{n \geq 1} \kappa_B^{p^n}
$$
を得る。

\medskip\noindent
$\pi \in A$ を一意化元とする。$A^\wedge$, $B^\wedge$ をそれぞれ
$A$, $B$ の完備化とする。離散付値環の拡大の可換図式
$$
\xymatrix{
B \ar[r] & B^\wedge \\
A \ar[u] \ar[r] & A^\wedge \ar[u]
}
$$
がある。$K^\wedge$ を $A^\wedge$ の分数体とする。
有限拡大 $M/K^\wedge$ で、(a) $A^\wedge \to B^\wedge$ に対する弱解であり、
(b) 分離拡大と $\pi$ の $p$ 冪根を添加して得られる拡大との合成であるものを
見つけられたとする。すると
Lemma \ref{more-algebra-lemma-approximate-separable-extension} により、
$K^\wedge \otimes_K K_1 = M$ となる有限拡大 $K_1/K$ を見つけられる。
$A_1$, resp.\ $A_1^\wedge$ を、それぞれ
$K_1$, resp.\ $M$ における $A$, resp.\ $A^\wedge$ の整閉包とする。
$A \to A^\wedge$ は $\mathfrak m^\wedge$-進位相に関して
形式的に滑らかなので
（Lemma \ref{more-algebra-lemma-extension-dvrs-formally-smooth}）、
$A_1 \to A_1^\wedge$ は $\mathfrak m_1^\wedge$-進位相に関して
形式的に滑らかである
（Lemma \ref{more-algebra-lemma-formally-smooth-goes-up}、また
Remark \ref{more-algebra-remark-construction} の議論により
$A_1$ と $A_1^\wedge$ は離散付値環である）。
Lemma \ref{more-algebra-lemma-solutions-go-down} の (2) により、
$K_1/K$ は $A \to B^\wedge$ に対する弱解であると結論する。
Lemma \ref{more-algebra-lemma-solutions-go-down} の (1) を適用すると、
$K_1/K$ は $A \to B$ に対する弱解である。

\medskip\noindent
従って、$A$ と $B$ は標数 $p$ の分離閉な剰余体をもつ
完備離散付値環であり、
$\kappa_A \supset \bigcap\nolimits_{n \geq 1} \kappa_B^{p^n}$
であると仮定してよい。また一意化元 $\pi \in A$ が与えられ、
$A \to B$ に対する弱解で、分離拡大と $\pi$ の $p$ 冪根を
とって得られる体との合成であるものを見つけなければならない。
$A$ が混標数をもつならば、第二の条件は自動的であることに注意する。

\medskip\noindent
$k = \bigcap\nolimits_{n \geq 1} \kappa_B^{p^n}$ とおく。
$k$ は標数 $p$ の代数閉体であることに注意する。
$A$ が混標数をもつならば $\Lambda$ を $k$ の Cohen 環とし、
等標数の場合は $\Lambda = k[[t]]$ とおく。
等標数の場合に $t$ を $\pi$ に写す環準同型
$\Lambda \to A$ を選べる。等標数の場合、これは Cohen 構造定理
（Algebra, Theorem \ref{algebra-theorem-cohen-structure-theorem}）から従い、
混標数の場合、これは $\mathbf{Z}_p \to \Lambda$ が進位相に関して
形式的に滑らかであることから従う
（Lemmas \ref{more-algebra-lemma-extension-dvrs-formally-smooth} と
\ref{more-algebra-lemma-lift-continuous}）。
Lemma \ref{more-algebra-lemma-solutions-go-down} を適用すると、
$\Lambda \to B$ に対する弱解の存在を証明すれば十分であり、
等標数 $p$ の場合、その弱解は分離拡大と $t$ の $p$ 冪根を
とって得られる体との合成でなければならない。
しかし等標数の場合は $\Lambda = k[[t]]$ であり、
$k((t))$ の任意の拡大はそのような合成なので、
この要件を捨てられる。

\medskip\noindent
従って次の状況に到達する。$A$ と $B$ は完備であり、
$A$ の剰余体 $k$ は標数 $p > 0$ の代数閉体で、
$k = \bigcap \kappa_B^{p^n}$ である。混標数の場合には
$p$ は $A$ の一意化元である（すなわち $A$ は $k$ の Cohen 環である）。
$A$ が混標数をもつならば $\Lambda$ を $\kappa_B$ の Cohen 環とし、
等標数の場合は $\Lambda = \kappa_B[[t]]$ とおく。
上と同様に論じると、$k \to \kappa_B$ を持ち上げ、
一意化元を一意化元に写す環準同型 $A \to \Lambda$ を選べる。
$k \subset \kappa_B$ は分離的なので、環準同型 $A \to \Lambda$ は
進位相に関して形式的に滑らかである
（Lemma \ref{more-algebra-lemma-extension-dvrs-formally-smooth}）。
従って合成 $A \to \Lambda \to B$ が与えられた環準同型
$A \to B$ となるような環準同型 $\Lambda \to B$ を見つけられる
（Lemma \ref{more-algebra-lemma-lift-continuous} 参照）。
$\Lambda$ と $B$ は同じ剰余体をもつ完備離散付値環なので、
$B$ は $\Lambda$ 上有限である
（Algebra, Lemma \ref{algebra-lemma-finite-over-complete-ring}）。
これにより Lemma \ref{more-algebra-lemma-special-case} で論じた特別な場合に帰着する。
\end{proof}







\section{分岐の除去 II}
\label{more-algebra-section-eliminating-ramification-bis}

\noindent
この節では Section \ref{more-algebra-section-eliminating-ramification} の結果を用いて、
いくつかの場合に（分離）解を得る。

\begin{lemma}
\label{more-algebra-lemma-solution-goes-up}
$A \to B$ を、分数体 $K \subset L$ をもつ離散付値環の拡大とする。
$K_1/K$ が $A \subset B$ に対する解ならば、任意の有限拡大
$K_2/K_1$ に対して拡大 $K_2/K$ は $A \subset B$ に対する解である。
\end{lemma}

\begin{proof}
これは Lemma \ref{more-algebra-lemma-formally-smooth-goes-up} から従う。詳細は省略する。
\end{proof}

\begin{lemma}
\label{more-algebra-lemma-nagata-goes-down}
$A \subset B$ を離散付値環の拡大とする。$B$ が Nagata 環で、
分数体の拡大 $L/K$ が分離的ならば、$A$ は Nagata 環である。
\end{lemma}

\begin{proof}
離散付値環が Nagata であることと N-2 であることは同値である。
$K_1/K$ を有限純非分離体拡大とする。
$K_1$ における $A$ の整閉包 $A_1$ が $A$ 上有限であることを
示さなければならない。Algebra, Lemma
\ref{algebra-lemma-domain-char-p-N-1-2} を参照されたい。
$L/K$ は分離的で $K_1/K$ は純非分離なので、Algebra, Lemmas
\ref{algebra-lemma-separable-extension-preserves-reducedness} と
\ref{algebra-lemma-radicial-integral-bijective} により、
代数 $L \otimes_K K_1$ は体である。
$B_1$ を $L \otimes_K K_1$ における $B$ の整閉包とする。
$B$ は Nagata 環なので $B_1$ は $B$ 上有限である。
$B \otimes_A A_1 \subset B_1$ で $B$ は Noether 環なので、
$B \otimes_A A_1$ は $B$ 上有限である。
$A \to B$ は忠実平坦なので、Algebra, Lemma
\ref{algebra-lemma-descend-properties-modules} により、
これは $A_1$ が $A$ 上有限であることを意味する。
\end{proof}

\begin{lemma}
\label{more-algebra-lemma-construct-extension}
$A' \subset A$ を環の拡大とし、$f \in A'$ とする。
(a) $A$ は $A'$ 上有限である、(b) $f$ は $A$ 上の非零因子である、
(c) $A'_f = A_f$ と仮定する。このとき整数 $n_0 > 0$ が存在して、
すべての $n \geq n_0$ に対して次が成り立つ。
環 $B'$、非零因子 $g \in B'$、および
$\varphi'(f) \equiv g$ を満たす同型
$\varphi' : A'/f^n A' \to B'/g^n B'$ が与えられると、
有限拡大 $B' \subset B$ と、$\varphi'$ と整合する同型
$\varphi : A/fA \to B/gB$ が存在する。
\end{lemma}

\begin{proof}
$A$ は $A'$ 上有限で $A'_f = A_f$ なので、
% P02-E1299: source line 35704 begins “Cchoose” instead of “Choose”.
$f^t A \subset A'$ となる $t > 0$ を Cchoose できる。
$n_0 = 2t$ とおく。補題の主張のような $n, B', g, \varphi'$ が
与えられたとする。
% P02-E1300: source line 35706 omits “by” in “denote N ... the set”.
$b \bmod g^nB' \in \varphi'(f^tA)$ を満たす元 $b \in B'$ の集合を
$N \subset B'$ と記す。$B = g^{-t}N$ とおく。
$f^tA' \subset f^tA$ で $\varphi'$ は $f$ を $g$ に写すので、
$g^tB' \subset N$ であり、従って $B' \subset B$ である。
$f^tA \cdot f^tA \subset f^t \cdot f^tA$ で
$\varphi'$ は $f$ を $g$ に写すので、$N \cdot N \subset g^t N$ である。
従って $B'$ の乗法を拡張する乗法を $B$ 上に得る。
右辺の加群構造を $\varphi'$ を用いて定めると、
$A'/f^nA'$-加群の同型
$$
A/f^tA' \xrightarrow{f^t} f^tA/f^nA' \xrightarrow{\varphi'}
g^tB/g^nB' \xrightarrow{g^{-t}} B/g^tB'
$$
を得る。$A/f^tA'$ は有限 $A'$-加群なので、$B/g^tB'$ は
有限 $B'$-加群であると結論する。従って $B' \to B$ は有限である。
最後に、表示した加群同型が $fA$ を $gB$ に写し、
$\varphi'$ と整合する環の同型
$\varphi : A/fA \to B/gB$ を誘導することは読者に任せる
（実際、同型 $A/f^tA \to B/g^tB$ も誘導するが、これは必要ない）。
\end{proof}

\begin{remark}
\label{more-algebra-remark-functoriality-construct-extension}
Lemma \ref{more-algebra-lemma-construct-extension} の構成は次の「関手性」を満たす。
単射な水平射をもつ可換図式
$$
\xymatrix{
A'_2 \ar[r] & A_2 \\
A'_1 \ar[r] \ar[u] & A_1 \ar[u]
}
$$
があるとする。$(A'_1 \subset A_1, f)$ と
$(A'_2 \subset A_2, f)$ が Lemma \ref{more-algebra-lemma-construct-extension} の
性質 (a), (b), (c) を満たすような元 $f \in A'_1$ が与えられたとする。
この二つの状況に対して補題で得られた整数を
$n_{0, 1}$ と $n_{0, 2}$ とする。
最後に、環準同型 $B'_1 \to B'_2$ と元 $g \in B'_1$ が与えられ、
% P02-E1301: source line 35743 says g is a nonzerodivisor on B_1 and B_2 before those rings are constructed; the given rings are B'_1 and B'_2.
この元は $B_1$ と $B_2$ 上の非零因子であり、
$n \geq \max(n_{0, 1}, n_{0, 2})$ とする。また、水平射が同型で
$\varphi'_1(f) \equiv g$ を満たす可換図式
% P02-E1302: source line 35748 repeats B'_2 in the lower-right corner; functoriality requires B'_1.
$$
\xymatrix{
A'_2/f^nA'_2 \ar[r]_{\varphi'_2} & B'_2/g^nB'_2 \\
A'_1/f^nA'_1 \ar[r]^{\varphi'_1} \ar[u] & B'_2/g^nB'_2 \ar[u]
}
$$
が与えられたとする。このとき可換図式
% P02-E1303: source line 35764 repeats B_2 in the lower-right corner; the lower row requires B_1.
$$
\vcenter{
\xymatrix{
B'_2 \ar[r] & B_2 \\
B'_1 \ar[r] \ar[u] & B_1 \ar[u]
}
}
\quad\text{and}\quad
\vcenter{
\xymatrix{
A_2/fA_2 \ar[r]_{\varphi_2} & B_2/gB_2 \\
A_1/fA_1 \ar[r]^{\varphi_1} \ar[u] & B_2/gB_2 \ar[u]
}
}
$$
を得る。ここで $(B'_1 \subset B_1, \varphi_1)$ と
$(B'_2 \subset B_2, \varphi_2)$ は
Lemma \ref{more-algebra-lemma-construct-extension} の証明のように構成される。
詳細な検証は省略する。
\end{remark}


\begin{lemma}
\label{more-algebra-lemma-approximate-solution}
$p$ を素数とする。$A \subset B$ を、分数体拡大 $L/K$ をもつ
離散付値環の拡大とする。$K_2/K_1/K$ を有限体拡大の塔とする。
次を仮定する。
\begin{enumerate}
\item $K$ の標数は $p$ である、
\item $L/K$ は分離的である、
\item $B$ は Nagata 環である、
\item $K_2$ は $A \subset B$ に対する解である、
\item $K_2/K_1$ は次数 $p$ の純非分離拡大である。
\end{enumerate}
このとき、$A \subset B$ に対する解である分離拡大 $K_3/K_1$ が存在する。
\end{lemma}

\begin{proof}
Remark \ref{more-algebra-remark-construction} の記法を用い、そこで述べたすべての
観察を用いる。$L/K$ は分離的なので、代数
$L_1 = L \otimes_K K_1$ は被約である
（Algebra, Lemma
\ref{algebra-lemma-separable-extension-preserves-reducedness}）。
$B$ は Nagata 環なので、$L_1$ における $B$ の整閉包を $B_1$ とすれば、
環拡大 $B \subset B_1$ は有限で、$B_1$ は Nagata 環である。
同様に Lemma \ref{more-algebra-lemma-nagata-goes-down} により $A$ は Nagata 環なので、
$A \subset A_1$ は有限で、$A_1$ も Nagata 環である。
さらに $K_2$ に対しても同じ主張が成り立つ。
すなわち $L_2 = L \otimes_K K_2$ は被約であり、
環拡大 $A_1 \subset A_2$, $B_1 \subset B_2$ は有限である。
ここで $A_2$, resp.\ $B_2$ は
$K_2$, resp.\ $L_2$ における $A$, resp.\ $B$ の整閉包である。

\medskip\noindent
$\pi \in A$ を一意化元とする。
$\pi$ は $K_1$, $K_2$, $A_1$, $A_2$, $L_1$, $L_2$, $B_1$, $B_2$
上の非零因子であり、
$K_1 = (A_1)_\pi$, $K_2 = (A_2)_\pi$,
$L_1 = (B_1)_\pi$, $L_2 = (B_2)_\pi$ であることに注意する。
Fields, Lemma \ref{fields-lemma-finite-purely-inseparable} により、
$\alpha^p = a_1 \in K_1$ を満たす元を用いて
$K_2 = K_1(\alpha)$ と書ける。
$\alpha$ に $\pi$ の冪を掛けることにより、
$a_1 \in A_1$ と仮定してよく、そうする。
証明の残りでは
$K_2 = K_1[x]/(x^p - a_1)$ および
$L_2 = L_1[x]/(x^p - a_1)$ と書くと便利である。
環の拡大
$$
A'_2 = A_1[x]/(x^p - a_1) \subset A_2
\quad\text{and}\quad
B'_2 = B_1[x]/(x^p - a_1) \subset B_2
$$
を考える。Lemma \ref{more-algebra-lemma-construct-extension} を
$A'_2 \subset A_2$ と $f = \pi^2$ に、また
$B'_2 \subset B_2$ と $f = \pi^2$ に適用できる。
両方に使える十分大きな整数 $n$ を選ぶ。

\medskip\noindent
代数
$$
K_3 = K_1[x]/(x^p - \pi^{2n} x - a_1)
\quad\text{and}\quad
L_3 = L_1[x]/(x^p - \pi^{2n} x - a_1)
$$
を考える。$K_3/K_1$ と $L_3/L_1$ は次数 $p$ の
有限 \'etale 代数拡大であることに注意する。
% P02-E1304: source lines 35834-35838 use pi^n in A'_3 and B'_3, inconsistent with K_3, L_3 and the modulo-pi^{2n} identity.
部分環
$$
A'_3 = A_1[x]/(x^p - \pi^n x - a_1)
\quad\text{and}\quad
B'_3 = B_1[x]/(x^p - \pi^n x - a_1)
$$
を
% P02-E1305: source line 35838 identifies K_3 and L_3 with localizations of A'_2 and B'_3 inconsistently; K_3 should be the localization of the correctly defined A'_3.
$K_3 = (A'_2)_\pi$ と $L_3 = (B'_3)_\pi$ の部分環として考える。
可換図式
$$
\xymatrix{
B'_2/\pi^{2n} B'_2 \ar[r]_{\psi'} & B'_3/\pi^{2n} B'_3 \\
A'_2/\pi^{2n} A'_2 \ar[r]^{\varphi'} \ar[u] &
A'_3/\pi^{2n} A'_3 \ar[u]
}
$$
を構成する。すなわち $\varphi'$ は、$x$ の類を $x$ の類に写す
一意な $A_1$-代数同型である。同様に $\psi'$ は、$x$ の類を
$x$ の類に写す一意な $B_1$-代数同型である。
$n$ の選び方により、Lemma \ref{more-algebra-lemma-construct-extension} と
Remark \ref{more-algebra-remark-functoriality-construct-extension} を通じて、
有限環拡大 $A'_3 \subset A_3$, $B'_3 \subset B_3$ で、
$A'_3 \to B'_3$ が環準同型 $A_3 \to B_3$ に拡張され、
可換図式
$$
\xymatrix{
B_2/\pi^2 B_2 \ar[r]_\psi & B_3/\pi^2B_3 \\
A_2/\pi^2 A_2 \ar[r]^\varphi \ar[u] & A_3/\pi^2A_3 \ar[u]
}
$$
が上記の参照箇所で主張されたすべての性質を満たすものを得る
（特に $\varphi$ と $\psi$ は同型である）。

\medskip\noindent
これらのデータを用いて証明を終える。
まず $A_3$ と $B_3$ は Dedekind 整域の有限直積であり、
$\pi$ はそのすべての極大イデアルに含まれることに注意する。
実際、$\mathfrak p \subset A_3$ が極大イデアルならば、
$A \to A_3$ は有限なので $\pi \in \mathfrak p$ である。
すると $\mathfrak p/\pi^2 A_3$ は $\varphi$ を通じて
$A_2 / \pi^2A_2$ の極大イデアルに対応し、$A_2$ は
Dedekind 整域の有限直積なので、それは単項である。
従って $\mathfrak p/\pi^2 A_3$ は単項であり、Nakayama の補題により
$\mathfrak p (A_3)_\mathfrak p$ は単項である。
$B_3$ に対しても同じ議論が成り立つ。
従って $A_3$ は $K_3$ における $A$ の整閉包であり、
$B_3$ は $L_3$ における $B$ の整閉包である。
$\mathfrak p \subset A_3$ の上にある極大イデアル
$\mathfrak q \subset B_3$ をとる。
証明を終えるには
$(A_3)_\mathfrak p \to (B_3)_\mathfrak q$ が
$\mathfrak q$-進位相に関して形式的に滑らかであることを示さなければならない。
Lemma \ref{more-algebra-lemma-extension-dvrs-formally-smooth} の判定法により、
$\mathfrak p (B_3)_\mathfrak q =
\mathfrak q (B_3)_\mathfrak q$ であり、体拡大
$\kappa(\mathfrak q)/\kappa(\mathfrak p)$ が分離的であることを
示せば十分である。両主張は環準同型
$A_3/\pi^2 A_3 \to B_3/\pi^2 B_3$ を見れば検査でき、
これは対応する主張が成り立つ環準同型
$A_2/\pi^2 A_2 \to B_2/\pi^2 B_2$ と同型である。
ここで対応する主張は、$K_2$ が $A \subset B$ に対する解であるという
仮定により成り立つ。詳細の一部は省略する。
\end{proof}

\begin{lemma}
\label{more-algebra-lemma-separable-solution-separable-solution}
$A \subset B$ を離散付値環の拡大とする。次を仮定する。
\begin{enumerate}
\item 分数体の拡大 $L/K$ は分離的である、
\item $B$ は Nagata 環である、さらに
\item $A \subset B$ に対する解が存在する。
\end{enumerate}
このとき $A \subset B$ に対する分離解が存在する。
\end{lemma}

\begin{proof}
$K$ の標数が零ならば補題は自明である。従って $K$ の標数は
$p > 0$ であると仮定してよく、そうする。

\medskip\noindent
$K_2/K$ を $A \to B$ に対する解とする。
$K_2/K$ の非分離次数 $[K_2 : K]_i$
（Fields, Definition \ref{fields-definition-insep-degree}）
に関する帰納法を用いる。
$[K_2 : K]_i = 1$ ならば $K_2$ は $K$ 上分離的であり、証明は終わる。
そうでなければ、$K_2/K_1$ が次数 $p$ の純非分離拡大となる
部分体の塔 $K_2/K_1/K$ が存在する
（Fields, Lemmas \ref{fields-lemma-separable-first} と
\ref{fields-lemma-finite-purely-inseparable}）。
Lemma \ref{more-algebra-lemma-approximate-solution} により、
$A \subset B$ に対する解となる分離拡大 $K_3/K_1$ が存在する。
このとき
$[K_3 : K]_i = [K_1 : K]_i = [K_2 : K]_i/p$
（Fields, Lemma \ref{fields-lemma-multiplicativity-all-degrees}）
はより小さいので、帰納法により結論を得る。
\end{proof}

\begin{lemma}
\label{more-algebra-lemma-big-extension-is-ok}
$A \to B$ を、分数体 $K \subset L$ をもつ離散付値環の拡大とする。
$B$ は $A$ 上本質的有限型であると仮定する。
$K'/K$ を体の代数拡大とし、$K'$ における $A$ の整閉包 $A'$ は
Noether 環であるとする。このとき
$L' = (L \otimes_K K')_{red}$ における $B$ の整閉包 $B'$ も
Noether 環である。さらに写像 $\Spec(B') \to \Spec(A')$ は全射であり、
対応する剰余体拡大は有限生成体拡大である。
\end{lemma}

\begin{proof}
$A \to C$ を有限型環準同型で、$B$ が $C$ の素イデアル
$\mathfrak p$ における局所化となるものとする。
$C' = C \otimes_A A'$ は有限型 $A'$-代数なので、特に Noether 環である。
$A \to A'$ は整なので $C \to C'$ も整である。
従って $B = C_\mathfrak p \subset C'_\mathfrak p$ も整である。
Algebra, Lemma \ref{algebra-lemma-integral-sub-dim-equal} により、
$C'_\mathfrak p$ の次元は $1$ である。もちろん
$C'_\mathfrak p$ は Noether 環である。
$\mathfrak q_1, \ldots, \mathfrak q_n$ を
$C'_\mathfrak p$ の極小素イデアルとする。
$B'_i$ を $B = C_\mathfrak p$ の整閉包、または上記により同値に、
% P02-E1306: source line 35948 uses the undefined localization subscript p-prime.
$C'_{\mathfrak p'}/\mathfrak q_i$ の分数体における
$C'_\mathfrak p$ の整閉包とする。
Krull--Akizuki、すなわち $C'_\mathfrak p$ の有限個の極大イデアルでの
局所化に適用した Algebra, Lemma
\ref{algebra-lemma-krull-akizuki} により、各 $B'_i$ は Noether 環である。
さらに Algebra, Lemma
\ref{algebra-lemma-finite-extension-residue-fields-dimension-1}
により、$C'_\mathfrak p \to B'_i$ の剰余体拡大は有限である。
最後に、$B' = \prod B'_i$ が
$L' = (L \otimes_K K')_{red}$ における $B$ の整閉包であることに注意する。
\end{proof}

\begin{proposition}
\label{more-algebra-proposition-epp-essentially-finite-type}
\begin{reference}
特別な場合については \cite[Lemma 2.13]{alterations} を参照されたい。
\end{reference}
$A \to B$ を、分数体 $K \subset L$ をもつ離散付値環の拡大とする。
$B$ が $A$ 上本質的有限型ならば、Definition \ref{more-algebra-definition-solution}
の意味で $A \to B$ に対する解となる有限拡大 $K_1/K$ が存在する。
\end{proposition}

\begin{proof}
$A$ の剰余体が完全ならば弱解は解であることに注意する。
Lemma \ref{more-algebra-lemma-extension-dvrs-formally-smooth} を参照されたい。
従って $A$ の剰余標数が $0$ ならば、命題は Theorem
\ref{more-algebra-theorem-epp} から直ちに従う
（実際、$A \to B$ が本質的有限型であるという仮定は必要ない）。
$A$ の剰余標数が $p > 0$ ならば、この場合も Epp の定理から導く。

\medskip\noindent
$x_i \in A$, $i \in I$ を、$A$ の剰余体 $\kappa$ の
$p$-基底に写る元の集合とする。
$$
A' = \bigcup\nolimits_{n \geq 1} A[t_{i, n}]/(t_{i, n}^{p^n} - x_i)
$$
とおく。ここで遷移写像は $t_{i, n + 1}$ を $t_{i, n}^p$ に写す。
$A'$ は $A$ 上の弱不分岐な有限離散付値環拡大のフィルター余極限である
ことに注意する。従って $A'$ は離散付値環であり、
$A \to A'$ は弱不分岐である。構成により剰余体
$\kappa' = A'/\mathfrak m_A A'$ は $\kappa$ の完全化である。

\medskip\noindent
$K'$ を $A'$ の分数体とする。
拡大 $K'/K$ に Lemma \ref{more-algebra-lemma-big-extension-is-ok} を適用できる。
従って $B'$ は Dedekind 整域の有限直積である。
$\mathfrak m_1, \ldots, \mathfrak m_n$ を $B'$ の極大イデアルとする。
Epp の定理（Theorem \ref{more-algebra-theorem-epp}）を用いると、各拡大
$A' \subset B'_{\mathfrak m_i}$ に対する弱解 $K'_i/K'$ が見つかる。
$A'$ の剰余体は完全なので、これらは実際には解である。
各 $K'_i$ を含む有限拡大を $K'_1/K'$ とする。
Lemma \ref{more-algebra-lemma-solution-goes-up} により、$K'_1/K'$ は
各 $A' \subset B'_{\mathfrak m_i}$ に対してなお解である。

\medskip\noindent
$A'_1$ を $K'_1$ における $A$ の整閉包とする。
Remark \ref{more-algebra-remark-construction} の $K' \subset K'_1$ への適用の議論により、
$A'_1$ は Dedekind 整域であることに注意する。
従って Lemma \ref{more-algebra-lemma-big-extension-is-ok} は $K'_1/K$ に適用できる。
ゆえに $L'_1 = (L \otimes_K K'_1)_{red}$ における $B$ の整閉包 $B'_1$ は
% P02-E1307: source line 36012 says B'_1 is a Dedekind domain although the reduced tensor product may have several field factors.
Dedekind 整域である。また $K'_1/K'$ は各
$A' \subset B'_{\mathfrak m_i}$ に対する解なので、
各極大イデアル $\mathfrak m \subset B'_1$ に対して
$(A'_1)_{A'_1 \cap \mathfrak m} \to (B'_1)_{\mathfrak m}$ は
$\mathfrak m$-進位相に関して形式的に滑らかである。

\medskip\noindent
構成により、体 $K'_1$ は $K$ の有限拡大のフィルター余極限である。
$K'_1 = \colim_{i \in I} K_i$ と書く。
各 $i$ に対し、$A_i$, resp.\ $B_i$ を
$K_i$, resp.\ $L_i = (L \otimes_K K_i)_{red}$ における
$A$, resp.\ $B$ の整閉包とする。このとき明らかに
$$
A'_1 =  \colim A_i\quad\text{and}\quad B'_1 = \colim B_i
$$
である。環準同型 $A_i \to A'_1$, $B_i \to B'_1$ は単射整準同型であり、
$A'_1$, $B'_1$ のスペクトルは有限なので、十分大きいすべての $i$ に対して
環準同型 $A_i \to A'_1$, $B_i \to B'_1$ はスペクトル上全単射となる。
これが成り立てば、さらに十分大きいすべての $i$ に対して
$A_i \to A'_1$, $B_i \to B'_1$ は弱不分岐となる
（一意化元が像に入った後）。
分岐指数の乗法性により、そのような $i$ に対して
$A_i \to B_i$ は $B_i$ のすべての極大イデアルでの局所化上、
弱不分岐な写像を誘導する。$i$ をもう少し大きくすると、
$$
B_i \otimes_{A_i} A'_1 \longrightarrow B'_1
$$
は剰余体上全射な写像を誘導する
（Lemma \ref{more-algebra-lemma-big-extension-is-ok} により、$B'_1$ の剰余体は
$A'_1$ の剰余体上有限生成だからである）。
$B'_1$ の選んだ極大イデアルの下にある極大イデアルでの
剰余体の図は
$$
\xymatrix{
\kappa_{B_i} \ar[r] &
\kappa_{B_{i'}} \ar[r] &
 & \ldots &
\kappa_{B'_1} \\
\kappa_{A_i} \ar[r] \ar[u] &
\kappa_{A_{i'}} \ar[r] \ar[u] &
 & \ldots &
\kappa_{A'_1} \ar[u]
}
$$
である。従って $\kappa_{B_i}$ は $\kappa_{A_i}$ の有限生成拡大であり、
$\kappa_{B'_1}$ における $\kappa_{B_i}$ と $\kappa_{A'_1}$ の合成は
$\kappa_{A'_1}$ 上分離的である。するとこれは既に有限段階で成り立つ。
例えば $\kappa_{B'_1}$ が
$\kappa_{A'_1}(x_1, \ldots, x_n)$ 上有限分離的であるとする。
$x_1, \ldots, x_n$ が $\kappa_{B_i}$ に入り、すべての生成元が
$\kappa_{A_i}(x_1, \ldots, x_n)$ 上分離多項式方程式を満たすように、
$i$ を大きくすればよい。これは、$B_i$ のすべての極大イデアル
$\mathfrak m$ に対して $A_i \to (B_i)_\mathfrak m$ が
$\mathfrak m$-進位相に関して形式的に滑らかであることを意味する。
これで証明は完了する。
\end{proof}

\begin{lemma}
\label{more-algebra-lemma-epp-essentially-finite-type-separable}
$A \to B$ を、分数体 $K \subset L$ をもつ離散付値環の拡大とする。
次を仮定する。
\begin{enumerate}
\item $B$ は $A$ 上本質的有限型である、
\item $A$ または $B$ は Nagata 環である、さらに
\item $L/K$ は分離的である。
\end{enumerate}
このとき $A \to B$ に対する分離解が存在する
（Definition \ref{more-algebra-definition-solution}）。
\end{lemma}

\begin{proof}
$A$ が Nagata 環ならば $B$ も Nagata 環であることに注意する
（Algebra, Lemma \ref{algebra-lemma-nagata-localize} と
Proposition \ref{algebra-proposition-nagata-universally-japanese}）。
従って Proposition \ref{more-algebra-proposition-epp-essentially-finite-type} と
Lemma \ref{more-algebra-lemma-separable-solution-separable-solution} を組み合わせれば
補題が従う。
\end{proof}







\section{環の Picard 群}
\label{more-algebra-section-picard}

\noindent
まず可逆加群を次のように定義する。

\begin{definition}
\label{more-algebra-definition-invertible}
$R$ を環とする。$R$-加群 $M$ が {\it 可逆} であるとは、関手
$$
\text{Mod}_R \longrightarrow \text{Mod}_R,\quad
N \longmapsto M \otimes_R N
$$
が圏同値であることをいう。可逆 $R$-加群が {\it 自明} であるとは、
$R$-加群として $R$ と同型であることをいう。
\end{definition}

\begin{lemma}
\label{more-algebra-lemma-invertible}
$R$ を環とし、$M$ を $R$-加群とする。次は同値である。
\begin{enumerate}
\item $M$ は階数 $1$ の有限局所自由加群である、
\item $M$ は可逆である、さらに
\item $M \otimes_R N \cong R$ となる $R$-加群 $N$ が存在する。
\end{enumerate}
さらに、この場合 (3) の加群 $N$ は
$\Hom_R(M, R)$ と同型である。
\end{lemma}

\begin{proof}
(1) を仮定する。加群 $N = \Hom_R(M, R)$ と評価写像
$M \otimes_R N = M \otimes_R \Hom_R(M, R) \to R$ を考える。
$M_f \cong R_f$ となる $f \in R$ を取れば、この $f$ で局所化すると、
評価写像は同型となる（詳細は省略する）。従って Algebra,
Lemma \ref{algebra-lemma-cover} により評価写像は同型である。
よって (1) $\Rightarrow$ (3) である。

\medskip\noindent
(3) を仮定する。このとき関手 $K \mapsto K \otimes_R N$ は
関手 $K \mapsto K \otimes_R M$ の擬逆である。従って
(3) $\Rightarrow$ (2) である。逆に (2) が成り立つならば、
$K \mapsto K \otimes_R M$ は本質的全射なので (3) が従う。

\medskip\noindent
同値な条件 (2), (3) が成り立つと仮定する。(3) の同型を
$\psi : M \otimes_R N \to R$ と書く。
$\psi(\xi) = 1$ となる元
$\xi = \sum_{i = 1, \ldots, n} x_i \otimes y_i$ を選ぶ。同型
$$
M \to M \otimes_R M \otimes_R N \to M
$$
を考える。ここで第一の矢印は $x$ を
$\sum x_i \otimes x \otimes y_i$ に送り、第二の矢印は
$x \otimes x' \otimes y$ を $\psi(x' \otimes y)x$ に送る。
従って $x \mapsto \sum \psi(x \otimes y_i)x_i$ は
$M$ の自己同型である。この自己同型は
$$
M \to R^{\oplus n} \to M
$$
と分解する。第一の矢印は
$x \mapsto (\psi(x \otimes y_1), \ldots, \psi(x \otimes y_n))$
で与えられ、第二の矢印は
$(a_1, \ldots, a_n) \mapsto \sum a_i x_i$ で与えられる。
このようにして $M$ は有限自由 $R$-加群の直和因子であることが分かる。
これは $M$ が有限局所自由であることを意味する
（Algebra, Lemma \ref{algebra-lemma-finite-projective}）。
対称性により $N$ についても同様であり、かつ
$M \otimes_R N \cong R$ なので、
$M$ と $N$ はともに階数 $1$ でなければならない。
\end{proof}

\noindent
これらの加群の同型類の集合は、しばしば $R$ の
{\it 類群} または {\it Picard 群} と呼ばれる。
群構造は、可逆加群 $L$, $L'$ の同型類に
$L \otimes_R L'$ の同型類を対応させることで定まる。
可逆加群 $L$ の逆元は加群
$$
L^{\otimes -1} = \Hom_R(L, R),
$$
である。実際、Lemma \ref{more-algebra-lemma-invertible} の証明で見たように、
評価写像 $L \otimes_R L^{\otimes -1} \to R$ は同型である。
$R$ の Picard 群を $\Pic(R)$ と書くことにする。

\begin{lemma}
\label{more-algebra-lemma-UFD-Pic-trivial}
$R$ を UFD とする。このとき $\Pic(R)$ は自明である。
\end{lemma}

\begin{proof}
$L$ を可逆 $R$-加群とする。Lemma \ref{more-algebra-lemma-invertible} により、
$L$ は有限局所自由 $R$-加群である。特に $L$ はねじれなしで
$R$ 上有限である。双対可逆加群の零でない元
$\varphi \in \Hom_R(L, R)$ を取る。
このとき $I = \varphi(L) \subset R$ は、可逆加群でもあるイデアルである。
零でない $f \in I$ を取り、
$$
f = u p_1^{e_1} \ldots p_r^{e_r}
$$
を、$p_i$ が相異なる素元となる素元分解とする。
$L$ は有限局所自由なので、$a_i \in R$, $a_i \not \in (p_i)$ であって、
ある $g_i \in R_{a_i}$ に対し $I_{a_i} = (g_i)$ となるものが存在する。
すると $p_i$ は UFD $R_{a_i}$ においても素元であり、
$p_i$ で割り切れないある $g'_i \in R_{a_i}$ に対して
$g_i = p_i^{c_i} g'_i$ と書ける。$f \in I_{a_i}$ なので
$e_i \geq c_i$ である。$I$ は
$h = p_1^{c_1} \ldots p_r^{c_r}$ で生成されると主張する。
これが証明を完了する。

\medskip\noindent
主張を証明するには、$I_a$ が主イデアルとなる任意の $a \in R$ に対し、
$I_a$ が $h$ で生成されることを示せば十分である
（Algebra, Lemma \ref{algebra-lemma-cover}）。$I_a = (g)$ とする。
$J \subset \{1, \ldots, r\}$ を、$p_i$ が $R_a$ で非単元
（従って素元）となるような $i$ の集合とする。
$f \in I_a = (g)$ なので、$v$ を単元、$b_j \leq e_j$ として
% P02-E1308: source line 36214 binds i in the product but uses p_j and b_j.
$g = v \prod_{i \in J} p_j^{b_j}$
という素元分解を得る。各 $j \in J$ に対して
$I_{aa_j} = g R_{aa_j} = g_j R_{aa_j}$ である。言い換えれば、
$g$ と $g_j$ の $R_{aa_j}$ における像は同伴である。
分解の一意性により $b_j = c_j$ であり、証明は完了する。
\end{proof}







\section{行列式}
\label{more-algebra-section-determinants}

\noindent
$R$ を環とし、$M$ を有限射影 $R$-加群とする。
積分解 $R = R_0 \times \ldots \times R_t$ であって、対応する
$M$ の分解 $M = M_0 \times \ldots \times M_t$ において、
$M_i$ が $R_i$ 上階数 $i$ の有限局所自由加群となるものが存在する。
これは Algebra, Lemma \ref{algebra-lemma-finite-projective}
（階数が局所定数であることを見るため）と、Algebra, Lemmas
\ref{algebra-lemma-disjoint-decomposition} および
\ref{algebra-lemma-disjoint-implies-product}
（$R$ を積に分解するため）から従う。この状況で、$R$-加群として
$$
\det(M) = \wedge^0_{R_0}(M_0) \times \ldots \times \wedge^t_{R_t}(M_t)
$$
と定義する。各項が階数 $1$ の有限局所自由加群なので、
これは階数 $1$ の有限局所自由加群である。
$\varphi : M \to N$ が有限射影 $R$-加群の同型ならば、
階数 $1$ の局所自由加群の標準的同型
$$
\det(\varphi) : \det(M) \longrightarrow \det(N)
$$
を得る。より一般に、$R$ のすべての素イデアル $\mathfrak p$ に対し、
自由加群 $M_\mathfrak p$ と $N_\mathfrak p$ の階数が等しいならば、
任意の $R$-加群準同型 $\varphi : M \to N$ は $R$-加群準同型
$\det(\varphi) : \det(M) \to \det(N)$ を誘導する。
最後に $M = N$ ならば、
$\det(\varphi) : \det(M) \to \det(M)$ は可逆 $R$-加群の自己準同型である。
可逆 $R$-加群について $R = \Hom_R(L, L)$ なので、
$\det(\varphi)$ を $R$ の元と見なすことができ、またそうする。
このようにして {\it 行列式}
$$
\det : \Hom_R(M, M) \longrightarrow R
$$
を得る。これは乗法的写像である。

\begin{remark}
\label{more-algebra-remark-determinant-as-socle}
$R$ を環とし、$M$ を有限射影 $R$-加群とする。
このとき $R$ 上 $M$ の次数付き可換 $R$-代数である外積代数
$\wedge^*_R(M)$ を考えられる。$\det(M)$ の一つの公式は、
$\det(M) \subset \wedge^*_R(M)$ が
$M \subset \wedge^*_R(M)$ の零化域であるというものである。
この記述は $R$ の積分解を参照しないので、ときに有用である。
もちろん、これが所望の性質を満たすことを証明するには、
（上のように）$R$ を積に分解するか、または
$R$ の素イデアルにおける局所化を調べなければならない。
\end{remark}

\noindent
% P02-E1309: source line 36279 says “determinant give a short exact sequence”; “given” is intended.
次に、有限射影加群の短完全列が与えられたとき行列式がどうなるかを考える。

\begin{lemma}
\label{more-algebra-lemma-det-ses}
$R$ を環とする。
$$
0 \to M' \to M \to M'' \to 0
$$
を有限射影 $R$-加群の短完全列とする。このとき標準的同型
$$
\gamma : \det(M') \otimes \det(M'') \longrightarrow \det(M)
$$
が存在する。
\end{lemma}

\begin{proof}[第一の証明]
% P02-E1310: source repeats “First proof.” immediately after the proof heading.
$R$ を環 $R_{ij}$ の積に分解し、
$M' = \prod M'_{ij}$, $M'' = \prod M''_{ij}$ で、
$M'_{ij}$ の階数が $i$、$M''_{ij}$ の階数が $j$ となるようにする。
もちろんこのとき $M = \prod M_{ij}$ であり、$M_{ij}$ の階数は $i + j$ である。
従って $M'$ と $M''$ の階数がそれぞれ定数 $i$, $j$ である場合に帰着する。
この場合、標準的写像
$$
\wedge^i(M') \otimes \wedge^j(M'') \longrightarrow \wedge^{i + j}(M)
$$
を構成すればよい。そのため $M'$ の元
$m'_1, \ldots, m'_i$ と $M''$ の元
$m''_1, \ldots, m''_j$ を選ぶ。
$m'_1, \ldots, m'_i$ の像を $m_1, \ldots, m_i \in M$ と書き、
$M''$ において $m''_1, \ldots, m''_j$ に写る元を
$m_{i + 1}, \ldots , m_{i + j} \in M$ と書く。対応を
$$
m'_1 \wedge \ldots \wedge m'_i \otimes
m''_1 \wedge \ldots \wedge m''_j
\longmapsto
m_1 \wedge \ldots \wedge m_{i + j}
$$
と定める。これが矛盾なく定まり、同型であることの詳細な証明は省略する。
\end{proof}

\begin{proof}[第二の証明]
Remark \ref{more-algebra-remark-determinant-as-socle} にある
$\det(M)$, $\det(M')$, $\det(M'')$ の記述を用いる。
$R$-代数準同型 $\wedge^*_R(M') \to \wedge^*_R(M)$ および
$\wedge^*_R(M) \to \wedge^*_R(M'')$ を考える。前者は単射であり、
後者は全射である。
$x' \in \det(M') \subset \wedge^*_R(M')$ と
$x'' \in \det(M'') \subset \wedge^*_R(M'')$ を取る。
$x''$ に写る元 $y'' \in \wedge^*(M)$ を選び、
$$
\gamma(x' \otimes x'') = x' \wedge y'' \in \det(M) \subset \wedge^*_R(M)
$$
と置く。
% P02-E1311: source line 36332 omits “is” in “that this well defined”.
素イデアルでの局所化を調べれば、これが矛盾なく定まり同型であることは
容易に確かめられる。さらに、この構成は第一の証明の構成と同じ写像を与える。
\end{proof}

\begin{lemma}
\label{more-algebra-lemma-det-ses-functorial}
$R$ を環とする。
$$
\xymatrix{
0 \ar[r] &
M' \ar[r] \ar[d]^u &
M \ar[r] \ar[d]^v &
M'' \ar[r] \ar[d]^w &
0 \\
0 \ar[r] &
K' \ar[r] &
K \ar[r] &
K'' \ar[r] &
0
}
$$
を有限射影 $R$-加群の可換図式で、垂直矢印が同型であるものとする。
このとき同型の可換図式
$$
\xymatrix{
\det(M') \otimes \det(M'') \ar[r]_-\gamma \ar[d]_{\det(u) \otimes \det(w)} &
\det(M) \ar[d]^{\det(v)} \\
\det(K') \otimes \det(K'') \ar[r]^-\gamma & \det(K)
}
$$
を得る。ここで水平矢印は Lemma \ref{more-algebra-lemma-det-ses} で構成したものである。
\end{lemma}

\begin{proof}
省略する。ヒント：Lemma \ref{more-algebra-lemma-det-ses} における写像 $\gamma$ の
第二の構成を用いる。
\end{proof}

\begin{lemma}
\label{more-algebra-lemma-det-filtration}
$R$ を環とする。
$$
K \subset L \subset M
$$
を $R$-加群で、$K$, $L/K$, $M/L$ が有限射影 $R$-加群であるものとする。
このとき図式
$$
\xymatrix{
\det(K) \otimes \det(L/K) \otimes \det(M/L) \ar[r] \ar[d] &
\det(L) \otimes \det(M/L) \ar[d] \\
\det(K) \otimes \det(M/K) \ar[r] &
\det(M)
}
$$
は可換である。ここで写像は Lemma \ref{more-algebra-lemma-det-ses} のものである。
\end{lemma}

\begin{proof}
省略する。ヒント：$R$ の素イデアルで局所化した後、
$K \subset L \subset M$ は
$R^{\oplus a} \subset R^{\oplus a + b}  \subset R^{\oplus a + b + c}$
と同型であると仮定でき、この場合は明らかな計算である。
\end{proof}

\begin{lemma}
\label{more-algebra-lemma-det-direct-sum}
$R$ を環とし、$M'$, $M''$ を二つの有限射影 $R$-加群とする。
このとき図式
$$
\xymatrix{
\det(M') \otimes \det(M'') \ar[r]
\ar[d]_{\epsilon \cdot (\text{switch tensors})} &
\det(M' \oplus M'') \ar[d]^{\det(\text{switch summands})} \\
\det(M'') \otimes \det(M') \ar[r] &
\det(M'' \oplus M')
}
$$
は可換である。ここで
$\epsilon = \det( -\text{id}_{M' \otimes M''}) \in R^*$ であり、
水平矢印は Lemma \ref{more-algebra-lemma-det-ses} のものである。
\end{lemma}

\begin{proof}
省略する。
\end{proof}

\begin{lemma}
\label{more-algebra-lemma-det-switch}
$R$ を環とし、$M$, $N$ を有限射影 $R$-加群とする。
$a : M \to N$, $b : N \to M$ を $R$-線形写像とする。このとき
$$
\det(\text{id} + a \circ b) = \det(\text{id} + b \circ a)
$$
が $R$ の元として成り立つ。
\end{lemma}

\begin{proof}
素イデアルにおける局所化で $R$ を置き換えた後に主張を証明すれば十分である。
従って $R$ は局所環で、$M$, $N$ は有限自由であると仮定してよい。
この場合、$A$ のサイズが $n \times m$、$B$ のサイズが
$m \times n$ である行列の通常の行列式について
$$
\det(I_n + AB) = \det(I_m + BA)
$$
を示せばよい。これは
$R = \mathbf{Z}[a_{ij}, b_{ji}; 1 \leq i \leq n, 1 \leq j \leq m]$
で、$a_{ij}$ と $b_{ij}$ が変数かつ行列 $A$, $B$ の成分である場合に帰着する。
% P02-E1312: source names the second variable family b_ji and immediately calls it b_ij.
分数体を取れば標数零の体の場合に帰着する。
標数零では、サイズ $\leq N$ の行列の行列式を
その行列の冪の跡で表す普遍多項式がある。従ってすべての $k \geq 1$ に対し
% P02-E1313: the displayed trace identity has unequal constant traces when n != m and cannot prove the determinant identity as stated.
$$
\text{Trace}((I_n + AB)^k) = \text{Trace}((I_m + BA)^k)
$$
を証明すれば十分である。展開すると、すべての $k \geq 0$ に対し
$\text{Trace}((AB)^k) = \text{Trace}((BA)^k)$ を証明すれば十分である。
$k = 1$ については、これは周知の事実
$\text{Trace}(AB) = \text{Trace}(BA)$ である。$k > 1$ については、
$(AB)^k = A(BA)^{k - 1}B$ および
$(BA)^k = (BA)^{k - 1} A B$ と書けばこれから従う。
\end{proof}

\noindent
Algebra, Section \ref{algebra-section-K-groups} で、
有限射影 $R$-加群の同型類上の自由群を関係式
$[M'] + [M''] = [M' \oplus M'']$ で割った群 $K_0(R)$ を
定義したことを思い出そう。

\begin{lemma}
\label{more-algebra-lemma-det}
$R$ を環とする。写像
$$
\det : K_0(R) \longrightarrow \Pic(R)
$$
であって、$M$ が階数 $n$ の有限局所自由加群ならば、
$[M]$ を可逆加群 $\wedge^n(M)$ の類に送るものが存在する。
\end{lemma}

\begin{proof}
これは上の構成から直ちに従う。特に関係式が $0$ に写ることを見るには
Lemma \ref{more-algebra-lemma-det-ses} を用いる。
\end{proof}






\section{完全複体と K 群}
\label{more-algebra-section-perfect-K-group}

\noindent
環 $R$ の完全複体の導来圏の第零 K 群が、
Algebra, Section \ref{algebra-section-K-groups} で定義した
$K_0(R)$ と同じであることを手短に示す。

\begin{lemma}
\label{more-algebra-lemma-perfect-to-K-group}
$R$ を環とする。写像
$$
c : \text{perfect complexes over }R \longrightarrow K_0(R)
$$
で次の性質をもつものが存在する。
\begin{enumerate}
\item 完全複体 $K$ に対して $c(K[n]) = (-1)^nc(K)$、
\item 完全複体の識別三角形 $K \to L \to M \to K[1]$ に対して
$c(L) = c(K) + c(M)$、さらに
\item $K$ が有限射影加群からなる有限複体 $M^\bullet$ で表されるならば、
% P02-E1314: source line 36503 uses M_i although the complex is cohomologically indexed M^bullet.
$c(K) = \sum (-1)^i[M_i]$。
\end{enumerate}
\end{lemma}

\begin{proof}
$K$ を $D(R)$ の完全対象とする。定義により、$K$ は有限射影
$R$-加群の有限複体 $M^\bullet$ で表せる。$K_0(R)$ において
$$
c(K) = \sum (-1)^n[M^n]
$$
と置いて $c$ を定義する。もちろんこれが矛盾なく定まることを
示さなければならないが、それが示されれば (1), (3) は直ちに従う。
当面は写像 $c$ が有限射影 $R$-加群の複体上で定義されていると見る。

\medskip\noindent
$L^\bullet \to M^\bullet$ を、有限射影 $R$-加群の有限複体の
全射な写像とする。その核を $K^\bullet$ とする。このとき
$R$-加群の短完全列
$$
0 \to K^n \to L^n \to M^n \to 0
$$
を得る。$M^n$ は射影的なので、これらは分裂する。
従って $K^\bullet$ も有限射影 $R$-加群の有限複体であり、
$K_0(R)$ において
$c(L^\bullet) = c(K^\bullet) + c(M^\bullet)$ である。

\medskip\noindent
$M^\bullet$ を有限射影 $R$-加群の非輪状有限複体とする。
$n \not \in [a, b]$ に対して $M^n = 0$ とする。
このとき $M^\bullet$ を短完全列
% P02-E1315: source line 36538 jumps from N^{a+1} to N^{a+3}; the cokernel should be N^{a+2}.
$$
\begin{matrix}
0 \to M^a \to M^{a + 1} \to N^{a + 1} \to 0, \\
0 \to N^{a + 1} \to M^{a + 2} \to N^{a + 3} \to 0, \\
\ldots \\
0 \to N^{b - 3} \to M^{b - 2} \to N^{b - 2} \to 0, \\
0 \to N^{b - 2} \to M^{b - 1} \to M^b \to 0
\end{matrix}
$$
に分けられる。降下帰納法により
$N^{b - 2}, \ldots, N^{a + 1}$ は有限射影 $R$-加群であり、
これらの列は分裂完全で、かつ
% P02-E1316: source line 36547 omits the exponent n from the first alternating-sign sum.
$$
c(M^\bullet) = \sum (-1)[M^n] = \sum (-1)^n([N^{n - 1}] + [N^n]) = 0
$$
である。従って我々の構成は非輪状複体上で零となる。

\medskip\noindent
直前の二段落の結果から形式的に、$c$ が矛盾なく定まり (2) を満たすことが従う。
実際、有限射影 $R$-加群の有限複体 $M^\bullet$, $L^\bullet$ が
$D(R)$ の同じ対象を表すとする。この同型は複体の写像
$f : M^\bullet \to L^\bullet$ で表せる。Derived Categories,
Lemma \ref{derived-lemma-morphisms-from-projective-complex} を見よ。
複体の短完全列
% P02-E1317: source lines 36554-36570 introduce only M and L but use an undefined K in the cone sequence and calculation.
$$
0 \to L^\bullet \to C(f)^\bullet \to K^\bullet[1] \to 0
$$
を得る。Derived Categories, Definition \ref{derived-definition-cone} を見よ。
$f$ は擬同型なので、錐 $C(f)^\bullet$ は非輪状である
（これは例えば Derived Categories, Section
\ref{derived-section-canonical-delta-functor} の議論から従う）。
従って所望の通り
$$
0 = c(C(f)^\bullet) = c(L^\bullet) + c(K^\bullet[1]) =
c(L^\bullet) - c(K^\bullet)
$$
である。同様な (2) の証明は省略する。
\end{proof}

\noindent
次の補題は $K_0(R)$ が $K_0(D_{perf}(R))$ に等しいことを示す。

\begin{lemma}
\label{more-algebra-lemma-perfect-to-K-group-universal}
$R$ を環とする。$D_{perf}(R)$ を完全対象の導来圏とする。
Lemma \ref{more-algebra-lemma-perfect-ring-classical-generator} を見よ。
Lemma \ref{more-algebra-lemma-perfect-to-K-group} の写像 $c$ は同型
$K_0(D_{perf}(R)) = K_0(R)$ を与える。
\end{lemma}

\begin{proof}
$K_0(D_{perf}(R))$ の定義
（Derived Categories, Definition \ref{derived-definition-K-zero}）
から、$c$ は準同型
$K_0(D_{perf}(R)) \to K_0(R)$ を誘導する。

\medskip\noindent
$R$ 上の有限射影加群 $M$ が与えられたとき、
次数 $0$ に $M$ をもち、他の次数で零である $R$ 上の完全複体を
% P02-E1318: source line 36593 says “denote M[0] the perfect complex” and omits “by”.
$M[0]$ と書く。有限射影加群の短完全列
$0 \to M \to M' \to M'' \to 0$ から識別三角形
$M[0] \to M'[0] \to M''[0] \to M[1]$ を得る。
Derived Categories, Section \ref{derived-section-canonical-delta-functor} を見よ。
従って、ひどい記法で恐縮だが、$[M]$ を $[M[0]]$ に送ることで
写像 $K_0(R) \to K_0(D_{perf}(R))$ を得る。

\medskip\noindent
$K_0(R) \to K_0(D_{perf}(R)) \to K_0(R)$ が恒等写像であることは明らかである。
一方、$M^\bullet$ が有限射影 $R$-加群の有界複体ならば、
% P02-E1319: source line 36605 repeats “the” in “then the the existence”.
「愚直切断」の識別三角形
（Homology, Section \ref{homology-section-truncations} を見よ）
$$
\sigma_{\geq n}M^\bullet \to \sigma_{\geq n - 1}M^\bullet \to
M^{n - 1}[-n + 1] \to (\sigma_{\geq n}M^\bullet)[1]
$$
の存在と帰納法により、$K_0(D_{perf}(R))$ において
$$
[M^\bullet] = \sum (-1)^i[M^i[0]]
$$
である（再びこの記法については恐縮である）。
従って写像 $K_0(R) \to K_0(D_{perf}(R))$ は全射であり、証明が完了する。
\end{proof}







\section{有限長加群の自己準同型の行列式}
\label{more-algebra-section-determinants-finite-length}

\noindent
$(R, \mathfrak m, \kappa)$ を局所環とする。
有限長 $R$-加群と自己準同型 $\varphi : M \to M$ からなる対
$(M, \varphi)$ の圏を考える。この圏は Abel 圏であり、
すべての対象は Artin 的かつ Noether 的である。
定義については Homology, Section
\ref{homology-section-jordan-holder} を見よ。

\medskip\noindent
$(M, \varphi)$ がこの圏の単純対象ならば、$M$ は
$\mathfrak m$ で零化される。そうでなければ
$(\mathfrak m M, \varphi|_{\mathfrak m M})$ が非自明な部分対象となるからである。
% P02-E1320: source line 36642 misspells “subobject” as “suboject”.
また $\dim_\kappa(M) = \text{length}_R(M)$ は有限である。
従って線形代数により、行列式と跡
$$
\det\nolimits_\kappa(\varphi),\quad \text{Trace}_\kappa(\varphi)
$$
を $\kappa$ の元として定義できる。この場合の $\varphi$ の
特性多項式についても同様である。

\medskip\noindent
Homology, Lemma \ref{homology-lemma-finite-length} により、
この圏の任意の対象 $(M, \varphi)$ に対して
$$
0 \subset M_1 \subset \ldots \subset M_n = M
$$
という $\varphi$ の下で安定な部分加群による有限フィルトレーションであって、
$(M_i/M_{i - 1}, \varphi_i)$ がこの圏の単純対象となるものが存在する。
ここで $\varphi_i : M_i/M_{i - 1} \to M_i/M_{i - 1}$ は誘導写像である。
$(M, \varphi)$ の $\kappa$ 上の {\it 行列式} を
$$
\det\nolimits_\kappa(\varphi) = \prod \det\nolimits_\kappa(\varphi_i)
$$
と定義する。ここで $\det_\kappa(\varphi_i)$ は前段落で定義したものである。
$(M, \varphi)$ の $\kappa$ 上の {\it 跡} を
$$
\text{Trace}_\kappa(\varphi) = \sum \text{Trace}_\kappa(\varphi_i)
$$
と定義する。ここで $\text{Trace}_\kappa(\varphi_i)$ は
前段落で定義したものである。同様に $\kappa$ 上の $\varphi$ の
特性多項式を、前段落で定義した $\varphi_i$ の特性多項式の積として
定義できる。Jordan--H\"older
（Homology, Lemma \ref{homology-lemma-jordan-holder}）により、
これは矛盾なく定まる。

\begin{lemma}
\label{more-algebra-lemma-ses}
$(R, \mathfrak m, \kappa)$ を局所環とする。
上で述べた圏の短完全列
$0 \to (M, \varphi) \to (M', \varphi') \to (M'', \varphi'') \to 0$
を取る。このとき
$$
\det\nolimits_\kappa(\varphi') =
\det\nolimits_\kappa(\varphi)\det\nolimits_\kappa(\varphi''),\quad
\text{Trace}_\kappa(\varphi') = \text{Trace}_\kappa(\varphi) + 
\text{Trace}_\kappa(\varphi'')
$$
である。また $\kappa$ 上の $\varphi'$ の特性多項式は、
$\varphi$ と $\varphi''$ の特性多項式の積である。
\end{lemma}

\begin{proof}
演習として残す。
\end{proof}

\begin{lemma}
\label{more-algebra-lemma-multiplication}
$(R, \mathfrak m, \kappa) \to (R', \mathfrak m', \kappa')$
を局所環の局所準同型とし、$\kappa'/\kappa$ は有限拡大であると仮定する。
$u \in R'$ とする。このとき任意の有限長 $R'$-加群 $M'$ に対して
$$
\det\nolimits_\kappa(u : M' \to M') =
\text{Norm}_{\kappa'/\kappa}(u \bmod \mathfrak m')^m
$$
である。ここで $m = \text{length}_{R'}(M')$ である。
\end{lemma}

\begin{proof}
$\text{length}_R(M') = \text{length}_{R'}(M') [\kappa' : \kappa]$
なので、主張は意味をもつ。$M' = \kappa'$ のとき、
等式は $u$ による乗法という線形作用素の行列式としてのノルムの定義から従う。
一般の場合は適切なフィルトレーションを選び、
Lemma \ref{more-algebra-lemma-ses} の乗法性を用いてこの場合に帰着する。
いくつかの詳細は省略する。
\end{proof}

\begin{lemma}
\label{more-algebra-lemma-flat-base-change-det}
$$
(R, \mathfrak m, \kappa) \to (R', \mathfrak m', \kappa')
$$
を局所環の平坦局所準同型で
$$
m = \text{length}_{R'}(R'/\mathfrak mR') < \infty
$$
となるものとする。上の任意の $(M, \varphi)$ に対し、元
$\det_\kappa(\varphi)^m$ は $\kappa'$ において
$$
\det_{\kappa'}(\varphi \otimes 1 : M \otimes_R R' \to M \otimes_R R')
$$
に写る。
\end{lemma}

\begin{proof}
$R \to R'$ の平坦性により、Lemma \ref{more-algebra-lemma-ses} のような短完全列は
$R'$ への基底変換後も短完全列となる。従って
Lemma \ref{more-algebra-lemma-ses} の乗法性により、
$(M, \varphi)$ がこの圏の単純対象である場合を仮定してよい
（この節の導入を見よ）。単純な場合 $M$ は $\mathfrak m$ で零化される。
逐次商が $R'$-加群として $\kappa'$ と同型となるフィルトレーション
$$
0 \subset I_1 \subset I_2 \subset \ldots \subset I_{m - 1} \subset
R'/\mathfrak mR'
$$
を選ぶ。このとき逐次商が $M \otimes_\kappa \kappa'$ と同型となる
フィルトレーション
$$
0 \subset
M \otimes_\kappa I_1 \subset
M \otimes_\kappa I_2 \subset
\ldots \subset
M \otimes_\kappa I_{m - 1} \subset
M \otimes_\kappa R'/\mathfrak mR' = M \otimes_R R'
$$
を得る。また、これらの部分加群は $\varphi \otimes 1$ の下で不変である。
Lemma \ref{more-algebra-lemma-ses} により
$$
\det\nolimits_{\kappa'}(\varphi \otimes 1 : M \otimes_R R' \to M \otimes_R R')
=
\det\nolimits_{\kappa'}(\varphi \otimes 1 :
M \otimes_\kappa \kappa' \to M \otimes_\kappa \kappa')^m =
\det\nolimits_\kappa(\varphi)^m
$$
を得る。最後の等式は線形写像の行列式と体拡大との両立性による。
これで補題が証明された。
\end{proof}







\section{正則局所環は UFD である}
\label{more-algebra-section-regular-local-UFD}

\noindent
節の表題で述べた結果を証明する。

\begin{lemma}
\label{more-algebra-lemma-regular-local-Pic-zero}
$R$ を正則局所環とし、$f \in R$ とする。このとき
$\Pic(R_f) = 0$ である。
\end{lemma}

\begin{proof}
$L$ を可逆 $R_f$-加群とする。特に $L$ は有限 $R_f$-加群である。
$M_f \cong L$ となる有限 $R$-加群 $M$ が存在する。
Algebra, Lemma \ref{algebra-lemma-construct-fp-module} を見よ。
Algebra, Proposition \ref{algebra-proposition-regular-finite-gl-dim} により、
$M$ は $R$ 上有限自由分解 $F_\bullet$ をもつ。
従って $L$ は {\it 自由} $R_f$-加群の有限複体と擬同型である。
よって Lemma \ref{more-algebra-lemma-perfect-to-K-group} により、
% P02-E1321: source line 36795 writes [L_f], although L is already only defined as an R_f-module.
$K_0(R_f)$ において、ある $n \in \mathbf{Z}$ に対し
$[L_f] = n[R_f]$ である。Lemma \ref{more-algebra-lemma-det} の写像を適用すると、
$L$ が自明であることが分かる。
\end{proof}

\begin{lemma}
\label{more-algebra-lemma-regular-local-UFD}
正則局所環は UFD である。
\end{lemma}

\begin{proof}
正則局所環が整域であることを思い出そう。
Algebra, Lemma \ref{algebra-lemma-regular-domain} を見よ。
正則局所環 $R$ の次元に関する帰納法により一意分解性を証明する。
$\dim(R) = 0$ ならば $R$ は体であり、特に UFD である。
$\dim(R) > 0$ と仮定する。
$x \in \mathfrak m$, $x \not \in \mathfrak m^2$ を取る。
Algebra, Lemma \ref{algebra-lemma-regular-ring-CM} により
$R/(x)$ は正則であり、従って Algebra, Lemma
\ref{algebra-lemma-regular-domain} により整域であるから、
$x$ は素元である。$\mathfrak p \subset R$ を高さ $1$ の素イデアルとする。
Algebra, Lemma \ref{algebra-lemma-characterize-UFD-height-1} により、
$\mathfrak p$ が単項であることを示せばよい。
$x \in \mathfrak p$ ならば $\mathfrak p = (x)$ で終わるので、
$x \not \in \mathfrak p$ と仮定してよい。
任意の非極大素イデアル $\mathfrak q \subset R$ に対し、
局所環 $R_\mathfrak q$ は正則局所環である。
Algebra, Lemma
\ref{algebra-lemma-localization-of-regular-local-is-regular} を見よ。
帰納法により $\mathfrak pR_\mathfrak q$ は単項である。
特に $R_x$-加群
$\mathfrak p_x = \mathfrak pR_x \subset R_x$ は有限表示 $R_x$-加群であり、
任意の素イデアルにおける局所化は階数 $1$ の自由加群である。
Algebra, Lemma \ref{algebra-lemma-finite-projective} により、
$\mathfrak p_x$ は可逆 $R_x$-加群である。
Lemma \ref{more-algebra-lemma-regular-local-Pic-zero} により、
ある $y \in R_x$ に対し $\mathfrak p_x = (y)$ である。
ある $f \in \mathfrak p$, $e \in \mathbf{Z}$ に対し
$y = x^e f$ と書ける。
$f = a_1 \ldots a_r$ を $R$ の既約元への分解とする
（Algebra, Lemma \ref{algebra-lemma-factorization-exists}）。
$\mathfrak p$ は素なので、ある $i$ に対して
$a_i \in \mathfrak p$ である。
$\mathfrak p_x = (y)$ は素であり、$R_x$ において $a_i | y$ なので、
$\mathfrak p_x$ は $R_x$ において $a_i$ で生成される。
すなわち $R_x$ における $a_i$ の像は素元である。
$x$ は素元なので、Algebra, Lemma
\ref{algebra-lemma-invert-prime-elements} により
$a_i$ は $R$ において素元である。
$(a_i) \subset \mathfrak p$ で $\mathfrak p$ の高さは $1$ だから、
所望の通り $(a_i) = \mathfrak p$ である。
\end{proof}

\begin{lemma}
\label{more-algebra-lemma-picard-group-generic-fibre-regular}
$R$ を分数体 $K$、剰余体 $\kappa$ をもつ付値環とする。
$R \to A$ を次を満たす環準同型とする。
\begin{enumerate}
\item $A$ は局所環であり、$R \to A$ は局所的である、
\item $A$ は $R$ 上平坦かつ本質的有限型である、
% P02-E1322: source line 36853 omits “is” in the condition “A tensor_R kappa regular”.
\item $A \otimes_R \kappa$ は正則である。
\end{enumerate}
このとき $\Pic(A \otimes_R K) = 0$ である。
\end{lemma}

\begin{proof}
$L$ を可逆 $A \otimes_R K$-加群とする。特に $L$ は有限加群である。
$M \otimes_R K \cong L$ となる有限 $A$-加群 $M$ が存在する。
Algebra, Lemma \ref{algebra-lemma-construct-fp-module} を見よ。
$M$ は $R$-加群としてねじれなしであると仮定してよい。
従って $M$ は $R$-加群として平坦である
（Lemma \ref{more-algebra-lemma-valuation-ring-torsion-free-flat}）。
Lemma \ref{more-algebra-lemma-flat-finite-type-valuation-ring-finite-presentation} から、
$M$ は $A$-加群として有限表示であり、
$A$ は $R$-代数として本質的有限表示であると分かる。
Lemma \ref{more-algebra-lemma-structure-relatively-perfect} により、
$M$ は $R$ に関して相対的完全であり、特に $M$ は $A$-加群として擬連接である。
Lemma \ref{more-algebra-lemma-perfect-over-regular-local-ring} により
$M$ は完全であり、従って $M$ は $A$ 上有限自由分解 $F_\bullet$ をもつ。
すると $L$ は {\it 自由} $A \otimes_R K$-加群の有限複体と擬同型である。
よって Lemma \ref{more-algebra-lemma-perfect-to-K-group} により、
ある $n \in \mathbf{Z}$ に対し $K_0(A \otimes_R K)$ において
$[L] = n[A \otimes_R K]$ である。
Lemma \ref{more-algebra-lemma-det} の写像を適用すると $L$ は自明である。
\end{proof}







\section{複体の行列式}
\label{more-algebra-section-determinants-complexes}

\noindent
Section \ref{more-algebra-section-perfect-K-group} では、環 $R$ 上の完全複体 $K$ に
可逆 $R$-加群の同型類、すなわち $\Pic(R)$ の元を対応させる方法を見た。
実は Section \ref{more-algebra-section-determinants} と同様に、関手
$$
\det :
\left\{
\begin{matrix}
\text{category of perfect complexes} \\
\text{morphisms are isomorphisms}
\end{matrix}
\right\}
\longrightarrow
\left\{
\begin{matrix}
\text{category of invertible modules} \\
\text{morphisms are isomorphisms}
\end{matrix}
\right\}
$$
が存在することが分かる。さらに、$R$ のフィルター付き導来圏 $DF(R)$ の対象
$(L, F)$ で、フィルトレーションが有限かつ次数付き部分が完全複体であるものには、
標準的同型 $\det(\text{gr}L) \to \det(L)$ が存在する。
原典の説明については \cite{determinant} を見よ。
% P02-E1323: source line 36918 retains an unresolved editorial placeholder “insert future reference”.
この内容は後で追加する（将来の参照を挿入）。

\medskip\noindent
当面は、$D(R)$ における Tor 振幅 $[-1, 0]$ の完全対象 $L$ の場合に、
個別的な構成を与える。このような対象は
$$
L^\bullet = \ldots \to 0 \to L^{-1} \to L^0 \to 0 \to \ldots
$$
という複体で表せる。ここで $L^{-1}$ と $L^0$ は有限射影 $R$-加群である。
Lemma \ref{more-algebra-lemma-perfect} を見よ。この場合
$$
\det(L^\bullet) = \det(L^0) \otimes_R \det(L^{-1})^{\otimes -1} =
\Hom_R(\det(L^{-1}), \det(L^0))
$$
と置く。この形の複体が {\it 階数 $0$} であるとは、
$R$ のすべての素イデアルに対して
$L^{-1}_\mathfrak p$ と $L^0_\mathfrak p$ の階数が等しいことをいう。
$L^\bullet$ の階数が $0$ ならば、Section \ref{more-algebra-section-determinants} で見たように、
標準元
$$
\delta(L^\bullet) \in \det(L^\bullet)
$$
が存在する。これは単に $d : L^{-1} \to L^0$ の行列式である。
% P02-E1324: source line 36941 misspells “determinant” as “determininant”.
$\delta(L^\bullet)$ が $\det(L^\bullet)$ の自明化であることと、
$L^\bullet$ が非輪状であることは同値であることに注意する。

\medskip\noindent
複体の写像 $a^\bullet : K^\bullet \to L^\bullet$ で次を満たすものを考える。
\begin{enumerate}
\item $a^\bullet$ は擬同型である、
\item すべての $n$ に対して $a^n : K^n \to L^n$ は全射である、
\item $K^n$, $L^n$ は有限射影 $R$-加群であり、
$n \in \{-1, 0\}$ 以外では零である。
\end{enumerate}
この状況で同型
$$
\det(a^\bullet) : \det(K^\bullet) \longrightarrow \det(L^\bullet)
$$
を構成する。完全列
$0 \to \Ker(a^i) \to K^i \to L^i \to 0$ を用い、
Lemma \ref{more-algebra-lemma-det-ses} により $i = -1, 0$ に対して同型
$$
\gamma^i : \det(\Ker(a^i)) \otimes \det(L^i) \to \det(K^i)
$$
を得る。$a^\bullet$ は擬同型なので複体 $\Ker(a^\bullet)$ は
非輪状かつ階数 $0$ である。従って標準元
$\delta(\Ker(a^\bullet))$ は可逆 $R$-加群
$\det(\Ker(a^\bullet))$ の自明化である。上を見よ。
$\det(a^\bullet) : \det(K^\bullet) \to \det(L^\bullet)$ を、
図式
$$
\xymatrix{
\det(K^\bullet)
\ar[rr]_{\det(a^\bullet)}
& &
\det(L^\bullet)
\ar[ld]^{\delta(\Ker(a^\bullet))}
\\
& \det(L^\bullet) \otimes \det(\Ker(a^\bullet))
\ar[lu]^{\gamma^0 \otimes (\gamma^{-1})^{\otimes -1}}
}
$$
が可換となる一意な同型として定義する。

\begin{lemma}
\label{more-algebra-lemma-canonical-element-well-defined}
$R$ を環とし、$a^\bullet : K^\bullet \to L^\bullet$ を
上の (1), (2), (3) を満たす $R$-加群の複体の写像とする。
$L^\bullet$ の階数が $0$ ならば、$\det(a^\bullet)$ は
標準元 $\delta(K^\bullet)$ を $\delta(L^\bullet)$ に写す。
\end{lemma}

\begin{proof}
$M^i = \Ker(a^i)$ と書く。従って短完全列の写像
$$
\xymatrix{
0 \ar[r] &
M^{-1} \ar[r] \ar[d]_{d_M} &
K^{-1} \ar[r] \ar[d]_{d_K} &
L^{-1} \ar[r] \ar[d]_{d_L} &
0 \\
0 \ar[r] &
M^0 \ar[r] &
K^0 \ar[r] &
L^0 \ar[r] &
0
}
$$
を得る。Lemma \ref{more-algebra-lemma-det-ses-functorial} により、写像として
$\det(d_K)$ は $\det(d_M) \otimes \det(d_L)$ に対応する。
定義を書き下せば、これが必要な等式を与える。
\end{proof}

\begin{lemma}
\label{more-algebra-lemma-homotopic-surjections}
$R$ を環とし、$a^\bullet : K^\bullet \to L^\bullet$ を
上の (1), (2), (3) を満たす $R$-加群の複体の写像とする。
$h : K^0 \to L^{-1}$ を、
$b^0 = a^0 + d \circ h$ と
$b^{-1} = a^{-1} + h \circ d$ が全射となる写像とする。このとき
$\det(K^\bullet) \to \det(L^\bullet)$ という写像として
$\det(a^\bullet) = \det(b^\bullet)$ である。
\end{lemma}

\begin{proof}
$h = a^{-1} \circ \tilde h$ となる写像
$\tilde h : K^0 \to K^{-1}$ であって、
% P02-E1325: source line 37028 calls the degree-minus-one endomorphism k^1, while every later occurrence is k^{-1}.
$k^0 = \text{id} + d \circ \tilde h : K^0 \to K^0$ と
$k^1 = \text{id} + \tilde h \circ d : K^{-1} \to K^{-1}$ が
同型となるものが存在すると仮定する。このとき複体の可換図式
$$
\xymatrix{
0 \ar[r] &
\Ker(b^\bullet) \ar[r] \ar[d]_{c^\bullet} &
K^\bullet \ar[r]_{b^\bullet}
\ar[d]_{k^\bullet} &
L^\bullet \ar[r] \ar[d]^{\text{id}} &
0 \\
0 \ar[r] &
\Ker(a^\bullet) \ar[r] &
K^\bullet \ar[r]^{a^\bullet} &
L^\bullet \ar[r] &
0
}
$$
を得る。ここで $c^\bullet$ は核の間の誘導同型である。
Lemma \ref{more-algebra-lemma-det-ses-functorial} を用いると、図式
$$
\xymatrix{
\det(\Ker(b^i)) \otimes \det(L^i) \ar[r] \ar[d]_{\det(c^i) \otimes 1} &
\det(K^i) \ar[d]^{\det(k^i)} \\
\det(\Ker(a^i)) \otimes \det(L^i) \ar[r] &
\det(K^i)
}
$$
が可換であることが分かる。
% P02-E1326: source lines 37056-37058 reverse the typed direction of det(c^bullet), omit det around Ker(b), and misspell “trivialization”.
$\det(c^\bullet)$ は $\det(\Ker(a^\bullet))$ の標準的自明化を
$\Ker(b^\bullet)$ の標準的自明化へ写すので
（Lemma \ref{more-algebra-lemma-canonical-element-well-defined}）、
結論を得るための必要十分条件は、$R$ の元として
$$
\det(k^0) = \det(k^{-1})
$$
である。これは Lemma \ref{more-algebra-lemma-det-switch} から従う。

\medskip\noindent
直和因子 $U \subset K^{-1}$ で、
$a^{-1}|_U : U \to L^{-1}$ と $b^{-1}|_U : U \to L^{-1}$ が
ともに同型となるものが存在すると仮定する。
$\tilde h$ を $h$ と $a^{-1}|_U$ の逆写像との合成として定義する。
$\tilde h$ は証明の第一段落のような写像であると主張する。
実際、構成により $h = a^{-1} \circ \tilde h$ である。
$k^{-1} : K^{-1} \to K^{-1}$ が同型であることを示すには、
それが全射であることを示せば十分である
（Algebra, Lemma \ref{algebra-lemma-fun}）。
$u \in U$ とする。$b^{-1}(u') = a^{-1}(u)$ となる
$u' \in U$ を選べる。このとき $u = k^{-1}(u')$ である。
実際 $u$ と $k^{-1}(u')$ はともに $U$ に属し、計算
\footnote{$a^{-1}(k^{-1}(u')) =
a^{-1}(u') + a^{-1}(\tilde h(d(u'))) =
a^{-1}(u') + h(d(u')) = b^{-1}(u') = a^{-1}(u)$}
により $a^{-1}(u) = a^{-1}(k^{-1}(u'))$ である。
$a^{-1}|_U$ は同型なので等式を得る。従って
$U \subset \Im(k^{-1})$ である。一方 $x \in \Ker(a^{-1})$ ならば
$x = k^{-1}(x) \bmod U$ である。
$K^{-1} = \Ker(a^{-1}) + U$ なので $k^{-1}$ は全射である。
最後に $k^0 : K^0 \to K^0$ が全射であることを示す。
まず $a^0 \circ k^0 = b^0$ なので $a^0 \circ k^0$ は全射である。
$x \in \Ker(a^0)$ ならば、ある $y \in \Ker(a^{-1})$ に対して
$x = d(y)$ である。上で示したことにより、ある $z \in K^{-1}$ に対し
$y = k^{-1}(z)$ と書ける。このとき $x = k^0(d(z))$ であり、結論を得る。

\medskip\noindent
証明の最後の段階である。直前の段落のような $U$ を見つければ十分だが、
これは常に可能とは限らない。しかし二つの $R$-加群の写像の等しさを示すには、
$R$ の素イデアルでの局所化後に示せば十分である。従って
$R$ は局所環であると仮定してよい。すると次の問題を得る。
$$
\alpha, \beta : R^{\oplus n} \longrightarrow R^{\oplus m}
$$
を二つの全射 $R$-線形写像とする。
$\alpha|_U$ と $\beta|_U$ がともに同型となる直和因子
$U \subset R^{\oplus n}$ を見つけよ。
$R$ が体ならば線形代数により可能である。一般には剰余体上の解を取り、
これを局所環 $R$ 上の解へ持ち上げる。いくつかの詳細は省略する。
\end{proof}

\begin{lemma}
\label{more-algebra-lemma-compose-surjections}
$R$ を環とする。$a^\bullet : K^\bullet \to L^\bullet$ と
$b^\bullet : L^\bullet \to M^\bullet$ を、上の (1), (2), (3) を
満たす $R$-加群の複体の写像とする。このとき
$\det(b^\bullet) \circ \det(a^\bullet) = 
\det(b^\bullet \circ a^\bullet)$ である。
% P02-E1327: source line 37114 reverses the source and target of the composed determinant map.
これは写像 $\det(M^\bullet) \to \det(K^\bullet)$ としての等式である。
\end{lemma}

\begin{proof}
省略する。ヒント：Lemmas
\ref{more-algebra-lemma-det-ses}, \ref{more-algebra-lemma-det-ses-functorial},
\ref{more-algebra-lemma-det-filtration} から直ちに従う。
\end{proof}

\begin{lemma}
\label{more-algebra-lemma-determinant-two-term-complexes}
$R$ を環とする。上の構成は関手
$$
\det :
\left\{
\begin{matrix}
\text{category of perfect complexes} \\
\text{with tor amplitude in }[-1, 0] \\
\text{morphisms are isomorphisms}
\end{matrix}
\right\}
\longrightarrow
\left\{
\begin{matrix}
\text{category of invertible modules} \\
\text{morphisms are isomorphisms}
\end{matrix}
\right\}
$$
を定める。さらに $D(R)$ の階数 $0$ の完全対象 $L$ で
Tor 振幅が $[-1, 0]$ であるものには標準元
$\delta(L) \in \det(L)$ が存在し、$D(R)$ における任意の同型
$a : L \to K$ に対し $\det(a)(\delta(L)) = \delta(K)$ である。
\end{lemma}

\begin{proof}
Lemma \ref{more-algebra-lemma-perfect} により、始域の圏のすべての対象は複体
$$
L^\bullet = \ldots \to 0 \to L^{-1} \to L^0 \to 0 \to \ldots
$$
で表せる。ここで $L^{-1}$ と $L^0$ は有限射影 $R$-加群である。
この型の複体を一時的に良い複体と呼ぶことにする。
Derived Categories, Lemma
\ref{derived-lemma-morphisms-from-projective-complex} により、
導来圏における良い複体間の射は複体写像のホモトピー類である。
従って良い複体を用いて議論でき、上で調べた行列式
$\det(L^\bullet) = \det(L^0) \otimes \det(L^{-1})^{\otimes -1}$
を用いられる。

\medskip\noindent
$a^\bullet : L^\bullet \to K^\bullet$ を、
$D(R)$ における同型、すなわち擬同型である良い複体の射とする。
図式
$$
\xymatrix{
L^\bullet \ar[rr]_{a^\bullet} & & K^\bullet \\
& M^\bullet \ar[lu]^{b^\bullet} \ar[ru]_{c^\bullet}
}
$$
がホモトピーまで可換であり、$b^\bullet$, $c^\bullet$ が
上の条件 (1), (2), (3) を満たすとき、これを良い図式と呼ぶ。
このような図式があれば
$$
\det(a^\bullet) = \det(c^\bullet) \circ \det(b^\bullet)^{-1}
$$
と定義できる。ここで $\det(c^\bullet)$ と $\det(b^\bullet)$ は
上で構成した同型である。良い図式が常に存在し、得られる写像
$\det(a^\bullet)$ が良い図式の選択に依存しないことを示す。

\medskip\noindent
良い複体の擬同型 $a^\bullet : L^\bullet \to K^\bullet$ に対する
良い図式の存在を示す。全射 $p : R^{\oplus n} \to K^{-1}$ を選ぶ。
新たな良い複体
$$
M^\bullet = \ldots \to 0 \to
L^{-1} \oplus R^{\oplus n} \xrightarrow{d \oplus 1}
L^0 \oplus R^{\oplus n} \to 0 \to \ldots
$$
を考えられる。射影写像 $b^\bullet : M^\bullet \to L^\bullet$ と、
次数 $-1$ では $a^{-1} \oplus p$、次数 $0$ では
$a^0 \oplus d \circ p$ を用いる写像
$c^\bullet : M^\bullet \to K^\bullet$ を取る。
写像 $b^\bullet : M^\bullet \to L^\bullet$ と
$c^\bullet : M^\bullet \to K^\bullet$ は
上の条件 (1), (2), (3) を満たすので、良い図式を得る。

\medskip\noindent
良い図式
$$
\xymatrix{
L^\bullet \ar[rr]_{\text{id}^\bullet} & & L^\bullet \\
& M^\bullet \ar[lu]^{b^\bullet} \ar[ru]_{c^\bullet}
}
$$
が与えられたとする。Lemma \ref{more-algebra-lemma-homotopic-surjections} により
$\det(c^\bullet) = \det(b^\bullet)$ である。従って
$\det(\text{id}^\bullet) = \text{id}$ は良い図式の選択に依存しない。

\medskip\noindent
一般の場合の独立性を証明する前に、合成について考える。
良い複体の擬同型 $L_1^\bullet \to L_2^\bullet$ と
$L_2^\bullet \to L_3^\bullet$ および良い図式
$$
\vcenter{
\xymatrix{
L_1^\bullet \ar[rr] & &
L_2^\bullet \\
& M_{12}^\bullet \ar[lu] \ar[ru]
}
}
\quad\text{and}\quad
\vcenter{
\xymatrix{
L_2^\bullet \ar[rr] & &
L_3^\bullet \\
& M_{23}^\bullet \ar[lu] \ar[ru]
}
}
$$
があるとする。これを図式
$$
\xymatrix{
L_1^\bullet \ar[rr] & &
L_2^\bullet \ar[rr] & &
L_3^\bullet \\
& M_{12}^\bullet \ar[lu] \ar[ru]
& & M_{23}^\bullet \ar[lu] \ar[ru] \\
& &
M_{123}^\bullet \ar[lu] \ar[ru]
}
$$
へ拡張できる。ここで $M_{123}^\bullet \to M_{12}^\bullet$ と
$M_{123}^\bullet \to M_{23}^\bullet$ は性質 (1), (2), (3) をもち、
図式の正方形は可換である。単に
$M_{123}^n = M_{12}^n \times_{L_2^n} M_{23}^n$ と取ればよい。
すると Lemma \ref{more-algebra-lemma-compose-surjections} により図式
$$
\xymatrix{
\det(L_2^\bullet) &
\det(M_{23}^\bullet) \ar[l] \\
\det(M_{12}^\bullet) \ar[u] &
\det(M_{123}^\bullet) \ar[l] \ar[u]
}
$$
は可換である。図式追跡により、
$M_{12}^\bullet$, $M_{23}^\bullet$ を用いた二つの良い図式に
付随する写像の合成
$\det(L_1^\bullet) \to \det(L_2^\bullet) \to \det(L_3^\bullet)$
は、良い図式
$$
\xymatrix{
L_1^\bullet \ar[rr] & & L_3^\bullet \\
& M_{123}^\bullet \ar[lu] \ar[ru]
}
$$
に付随する写像と等しい。従って、これらの写像が選択に依存しないことを
示せれば合成則も満たされ、所望の関手を得る。

\medskip\noindent
独立性を示す。良い複体の擬同型
$a^\bullet : L^\bullet \to K^\bullet$ が与えられたとする。
逆擬同型 $b^\bullet : K^\bullet \to L^\bullet$ を選ぶ。
% P02-E1328: source lines 37280-37282 omit the bullet on the first L and omit “we” before “may fix”.
$L_1^\bullet = L$, $L_2^\bullet = K^\bullet$,
$L_3^\bullet = L^\bullet$ と置くと、$b^\bullet$ に対する良い図式の
選択を固定し、$a^\bullet$ に対する良い図式を動かしてよい。
前段落の結果によれば、どのような選択をしても合成は常に
$\det(L^\bullet)$ 上の恒等写像に等しい。
これは明らかに、それらの選択からの独立性を証明する。

\medskip\noindent
標準元についての主張は Lemma
\ref{more-algebra-lemma-canonical-element-well-defined} と我々の構成から直ちに従う。
\end{proof}







\section{付値環の拡大}
\label{more-algebra-section-valuation-rings}

\noindent
この節は一般の付値環に対する
Section \ref{more-algebra-section-discrete-valuation-rings} の類似である。

\begin{definition}
\label{more-algebra-definition-extension-valuation-rings}
$A$, $B$ が付値環で、$A \to B$ が単射かつ局所的であるとき、
$A \to B$ または $A \subset B$ を
{\it 付値環の拡大} という。このような拡大は可換図式
$$
\xymatrix{
A \setminus \{0\} \ar[r] \ar[d]_v & B \setminus \{0\} \ar[d]^v \\
\Gamma_A \ar[r] & \Gamma_B
}
$$
を誘導する。ここで $\Gamma_A$, $\Gamma_B$ は付値群である。
下の水平矢印が全単射であるとき、$B$ は $A$ 上
{\it 弱不分岐} であるという。
剰余体の拡大
$\kappa_A = A/\mathfrak m_A \subset \kappa_B = B/\mathfrak m_B$
が有限ならば $f = [\kappa_B : \kappa_A]$ と置き、
これを拡大 $A \subset B$ の {\it 剰余次数} または
{\it 剰余体次数} と呼ぶ。
\end{definition}

\noindent
% P02-E1329: source line 37329 says the units of A are “the inverse” of the units of B; the inverse image is intended.
$A \to B$ の下で $A$ の単元は $B$ の単元の逆なので、
$\Gamma_A \to \Gamma_B$ は単射であることに注意する。
また、分数体の拡大が有限であることは仮定しない。

\begin{lemma}
\label{more-algebra-lemma-inequality-general}
$A \subset B$ を、分数体 $K \subset L$ をもつ付値環の拡大とする。
拡大 $L/K$ が有限ならば剰余体の拡大は有限であり、
$\Gamma_B$ における $\Gamma_A$ の指数も有限で、かつ
$$
[\Gamma_B : \Gamma_A] [\kappa_B : \kappa_A] \leq [L : K].
$$
\end{lemma}

\begin{proof}
$b_1, \ldots, b_n \in B$ を、その $\kappa_B$ における像が
$\kappa_A$ 上線形独立となる単元とする。
$c_1, \ldots, c_m \in B$ を、その $\Gamma_B/\Gamma_A$ における像が
相異なる零でない元とする。$b_i c_j$ は $L$ において
$K$-線形独立であると主張する。実際、すべてが零ではない
$a_{ij} \in K$ を係数とする和
$$
\sum a_{ij} b_i c_j
$$
は零にならないと主張する。
$v(a_{i_0j_0}b_{i_0}c_{j_0})$ が最小となる $(i_0, j_0)$ を選ぶ。
$a_{ij}$ を $a_{ij}/a_{i_0j_0}$ で置き換え、
$a_{i_0 j_0} = 1$ とする。
$$
P = \{(i, j) \mid
v(a_{ij}b_ic_j) = v(a_{i_0j_0}b_{i_0}c_{j_0}) \}
$$
と置く。$c_1, \ldots, c_m$ の選び方により、
$(i, j) \in P$ ならば $j = j_0$ である。
従って $(i, j) \in P$ ならば
$v(a_{ij}) = v(a_{i_0j_0}) = 0$、すなわち $a_{ij}$ は単元である。
$b_1, \ldots, b_n$ の選び方により
$$
\sum\nolimits_{(i, j) \in P} a_{ij}b_i
$$
は $B$ の単元である。従って
$\sum\nolimits_{(i, j) \in P} a_{ij}b_ic_j$ の付値は
$v(c_{j_0}) = v(a_{i_0j_0}b_{i_0}c_{j_0})$ である。
最初に表示した和のうち $(i, j) \not \in P$ の項の付値は
これより真に大きいので、その和は零ではあり得ない。
これで補題が証明された。
\end{proof}

\begin{lemma}
\label{more-algebra-lemma-valuation-ring-purely-inseparable}
$A$ を、標数 $p > 0$ の分数体 $K$ をもつ付値環とする。
$L/K$ を純非分離拡大とする。このとき $L$ における $A$ の整閉包 $B$ は
分数体 $L$ をもつ付値環であり、
$A \subset B$ は付値環の拡大である。
\end{lemma}

\begin{proof}
省略する。ヒント：例えば Algebra, Lemmas
\ref{algebra-lemma-x-or-x-inverse-valuation-ring} および
\ref{algebra-lemma-integral-overring-surjective} を用いる。
\end{proof}

\begin{lemma}
\label{more-algebra-lemma-extension-normal-domains-and-roots}
$A \to B$ を Noether 局所正規整域の平坦局所準同型とする。
$f \in A$, $h \in B$ で、ある $n > 1$ と $B$ の単元 $w$ に対し
$f = w h^n$ となるものを取る。高さ $1$ のすべての素イデアル
$\mathfrak p \subset A$ に対し、その上にある高さ $1$ の素イデアル
$\mathfrak q \subset B$ で、$\mathfrak p$ の上にあり、拡大
$A_\mathfrak p \subset B_\mathfrak q$ が弱不分岐となるものが
存在すると仮定する。このとき、ある $g \in A$ と $A$ の単元 $u$ に対し
$f = u g^n$ である。
\end{lemma}

\begin{proof}
$A$, $B$ の高さ $1$ の素イデアルにおける局所環は離散付値環である
（Algebra, Lemma \ref{algebra-lemma-characterize-dvr}）。
従って仮定は Definition
\ref{more-algebra-definition-extension-discrete-valuation-rings} を通じて意味をもつ。
$\mathfrak p_1, \ldots, \mathfrak p_r$ を $f$ 上極小な $A$ の素イデアルとする。
Algebra, Lemma \ref{algebra-lemma-minimal-over-1} により高さは $1$ である。
各 $i$ に対し、$\mathfrak p_i$ 上にある高さ $1$ の $B$ の素イデアルを
$\mathfrak q_{i, j} \subset B$, $j = 1, \ldots, r_i$ とする。
$A_{\mathfrak p_i} \to B_{\mathfrak q_{i, 1}}$ が弱不分岐となるように
番号を付けるとする。$f$ は $B_{\mathfrak q_{i, 1}}$ において
$n$ 乗と単元の積に写るので、$A_{\mathfrak p_i}$ における $f$ の付値
$v_i$ は $n$ で割り切れる。ある $w_i \geq 0$ に対し
$v_i = n w_i$ と書く。イデアル $I$ を定める完全列
$$
0 \to I \to A \to
\prod\nolimits_{i = 1, \ldots, r}
A_{\mathfrak p_i}/\mathfrak p_i^{w_i}A_{\mathfrak p_i}
$$
を考える。完全関手 $- \otimes_A B$ を適用して完全列
$$
0 \to I \otimes_A B \to B \to
\prod\nolimits_{i = 1, \ldots, r}
(A_{\mathfrak p_i}/\mathfrak p_i^{w_i}A_{\mathfrak p_i}) \otimes_A B
$$
を得る。$i$ を固定する。標準写像
$$
(A_{\mathfrak p_i}/\mathfrak p_i^{w_i}A_{\mathfrak p_i}) \otimes_A B
\to
\prod\nolimits_{j = 1, \ldots, r_i}
B_{\mathfrak q_{i, j}}/\mathfrak q_{i, j}^{e_{i, j}w_i}B_{\mathfrak q_{i, j}}
$$
は単射であると主張する。ここで $e_{i, j}$ は
$A_{\mathfrak p_i} \to B_{\mathfrak q_{i, j}}$ の分岐指数である。
主張は、$\mathfrak p_i^{w_i}B_{\mathfrak p_i}$ が
$B_{\mathfrak p_i}$ の元 $b$ で、$\mathfrak q_{i, j}$ における付値が
$\geq e_{i, j}w_i$ となるもの全体に等しいというものである。
単項イデアル $\mathfrak p_i^{w_i}$ の生成元
$a \in A_{\mathfrak p_i}$ を選ぶ。このとき
$\mathfrak q_{i, j}$ における $a$ の付値は $e_{i, j}w_i$ に等しい。
$B_{\mathfrak p_i}$ は高さ一の素イデアルが
$\mathfrak q_{i, j}$, $j = 1, \ldots, r_i$ である正規整域なので、
上のような $b$ に対し Algebra, Lemma
\ref{algebra-lemma-normal-domain-intersection-localizations-height-1} により
$b/a \in B_{\mathfrak p_i}$ である。これで主張が従う。

\medskip\noindent
この主張と上の第二の完全列を組み合わせると、完全列
$$
0 \to I \otimes_A B \to B \to
\prod\nolimits_{i = 1, \ldots, r}
\prod\nolimits_{j = 1, \ldots, r_i}
B_{\mathfrak q_{i, j}}/\mathfrak q_{i, j}^{e_{i, j}w_i}B_{\mathfrak q_{i, j}}
$$
が定まる。従って $I \otimes_A B$ は、
$\mathfrak q_{i, j}$ における付値が $\geq e_{i, j}w_i$ となる
$B$ の元 $h'$ 全体である。$B$ において $f = wh^n$ なので、
$\mathfrak q_{i, j}$ における $h$ の付値は $e_{i, j}w_i$ である。
従って Algebra, Lemma
\ref{algebra-lemma-normal-domain-intersection-localizations-height-1} により
$h'/h \in B$ である。よって $I \otimes_A B$ は
（$h$ で生成される）階数 $1$ の自由 $B$-加群である。
従って $I$ は階数 $1$ の自由 $A$-加群である。
Algebra, Lemma \ref{algebra-lemma-finite-projective-descends} を見よ。
$g \in I$ を生成元とする。このとき $g$ と $h$ は
$B$ の単元だけ異なる。議論を逆向きにたどると、
$A_{\mathfrak p_i}$ における $g$ の付値は
$w_i = v_i/n$ である。従って所望の通り、Algebra, Lemma
\ref{algebra-lemma-normal-domain-intersection-localizations-height-1} により
$g^n$ と $f$ は $A$ の単元だけ異なる。
\end{proof}

\begin{lemma}
\label{more-algebra-lemma-etale-extension-valuation-ring}
$A$ を付値環とする。$A \to B$ をエタール環準同型とし、
$\mathfrak m \subset B$ を $A$ の極大イデアル上にある素イデアルとする。
このとき $A \subset B_\mathfrak m$ は弱不分岐な付値環の拡大である。
\end{lemma}

\begin{proof}
Lemma \ref{more-algebra-lemma-weak-dimension-at-most-1} により
$A$ の弱次元は $\leq 1$ である。Lemmas
\ref{more-algebra-lemma-weak-dimension-goes-up} と \ref{more-algebra-lemma-when-weakly-etale} により
$B$ の弱次元も $\leq 1$ である。
% P02-E1330: source line 37491 begins a new sentence with lowercase “hence”.
従って Lemma \ref{more-algebra-lemma-weak-dimension-at-most-1} により
局所環 $B_\mathfrak m$ は付値環である。
拡大 $A \subset B_\mathfrak m$ は分数体の有限拡大を誘導するので、
% P02-E1331: source line 37495 has the ungrammatical article in “the Gamma_A”.
$\Gamma_A$ は $B_{\mathfrak m}$ の付値群において有限指数をもつ。
従って任意の $h \in B_\mathfrak m$ に対して、
ある $n > 0$、元 $f \in A$、単元 $w \in B_\mathfrak m$ が存在し、
$B_\mathfrak m$ において $f = w h^n$ である。
これがある $g \in A$ と単元 $u \in A$ に対し $f = ug^n$ を導くことを示す。
これにより、補題で主張したように $A$ と $B_\mathfrak m$ の付値群が一致する。

\medskip\noindent
$\mathbf{Z}$ 上本質的有限型な局所部分環の余極限
$A = \colim A_i$ と書く。$A$ は正規整域なので
（Algebra, Lemma \ref{algebra-lemma-valuation-ring-normal}）、
各 $A_i$ は正規であると仮定してよい
（ここでは正規化を取った後も局所環が $\mathbf{Z}$ 上本質的有限型で
あり続けることを Algebra, Proposition
\ref{algebra-proposition-ubiquity-nagata} により用いる）。
ある $i$ に対し、$B = A \otimes_{A_i} B_i$ となるエタール拡大
$A_i \to B_i$ を見つけられる。Algebra, Lemma
\ref{algebra-lemma-etale} を見よ。
$\mathfrak m_i$ を $B_i$ と $\mathfrak m$ の共通部分とする。
このとき環準同型 $A_i \to (B_i)_{\mathfrak m_i}$ に
Lemma \ref{more-algebra-lemma-extension-normal-domains-and-roots} を適用して結論を得る。
補題の仮定は次により満たされる。
\begin{enumerate}
\item $A_i$ と $(B_i)_{\mathfrak m_i}$ は $\mathbf{Z}$ 上
本質的有限型なので Noether 環である、
\item $A_i \to B_i$ はエタールなので
$A_i \to (B_i)_{\mathfrak m_i}$ は平坦である、
\item $A_i \to B_i$ はエタールなので $B_i$ は正規である。
Algebra, Lemma \ref{algebra-lemma-normal-goes-up} を見よ、
\item $A_i$ の高さ $1$ のすべての素イデアルに対し、
その上にある $(B_i)_{\mathfrak m_i}$ の高さ $1$ の素イデアルが存在する。
これは Algebra, Lemma \ref{algebra-lemma-finite-in-codim-1} と、
$\Spec((B_i)_{\mathfrak m_i}) \to \Spec(A_i)$ が全射であることによる、
\item $A_i \to B_i$ はエタールなので、素イデアル $\mathfrak p$ 上にある
すべての素イデアル $\mathfrak q$ に対して誘導拡大
$(A_i)_\mathfrak p \to (B_i)_\mathfrak q$ は不分岐である。
\end{enumerate}
これで補題の証明が完了する。
\end{proof}

\begin{lemma}
\label{more-algebra-lemma-henselization-valuation-ring}
$A$ を付値環とし、$A^h$, resp.\ $A^{sh}$ をその
Hensel 化, resp.\ 強 Hensel 化とする。このとき
$$
A \subset A^h \subset A^{sh}
$$
は付値群上全単射を誘導する、すなわち弱不分岐な付値環の拡大である。
\end{lemma}

\begin{proof}
$A^h = \colim (B_i)_{\mathfrak q_i}$ と書く。ここで
$A \to B_i$ はエタールで、$\mathfrak q_i \subset B_i$ は
$\mathfrak m_A$ 上にある素イデアルである。
Algebra, Lemma \ref{algebra-lemma-henselization-different} を見よ。
Lemma \ref{more-algebra-lemma-etale-extension-valuation-ring} により
$(B_i)_{\mathfrak q_i}$ は付値環であり、誘導写像
$$
(A \setminus \{0\})/A^* \longrightarrow
((B_i)_{\mathfrak q_i} \setminus \{0\}) / (B_i)_{\mathfrak q_i}^*
$$
は全単射である。Algebra, Lemma
\ref{algebra-lemma-colimit-valuation-rings} により
$A^h$ は付値環であると結論する。また
$(A \setminus \{0\})/A^* \to (A^h \setminus \{0\})/(A^h)^*$
が全単射であることも従う。これで包含 $A \subset A^h$ に対して補題が示された。
$A \subset A^{sh}$ については、Algebra, Lemma
\ref{algebra-lemma-henselization-different} を Algebra, Lemma
\ref{algebra-lemma-strict-henselization-different} に置き換え、
全く同じ議論を用いればよい。$A^{sh} = (A^h)^{sh}$ なので、
これにより包含 $A^h \subset A^{sh}$ に対する補題の主張も示される。
\end{proof}







\section{PID 上の加群の構造}
\label{more-algebra-section-structure-modules}

\noindent
証明が付値環上でも機能するよう、Warfield の論文
\cite{Warfield-Purity}, \cite{Warfield-Decomposition} に従って、
少し一般的に議論する。

\begin{lemma}
\label{more-algebra-lemma-characterize-PD-modules}
\begin{reference}
\cite[Corollary 1]{Warfield-Purity}
\end{reference}
$P$ を環 $R$ 上の加群とする。次は同値である。
\begin{enumerate}
\item $P$ は、$f \in R$ を動かした形 $R/fR$ の加群の直和の直和因子である。
\item すべての $f \in R$ に対して $fA = A \cap fB$ を満たす
$R$-加群の任意の短完全列 $0 \to A \to B \to C \to 0$ に対し、
写像 $\Hom_R(P, B) \to \Hom_R(P, C)$ は全射である。
\end{enumerate}
\end{lemma}

\begin{proof}
(2) のような完全列 $0 \to A \to B \to C \to 0$ を取る。
(1) が (2) を導くことを証明するには、すべての $f \in R$ に対し
$\Hom_R(R/fR, B) \to \Hom_R(R/fR, C)$ が全射であることを示せば十分である。
写像 $\psi : R/fR \to C$ を取る。$\psi(1)$ は
ある $b \in B$ の像であるとする。このとき $fb \in A$ である。
従って $fa = fb$ となる $a \in A$ が存在する。
すると $f(b - a) = 0$ なので、$1$ を $b - a$ に写す射
$\varphi : R/fR \to B$ を得る。これは $\psi$ を持ち上げる。

\medskip\noindent
逆に (2) が成り立つと仮定する。$I$ を、
$f \in R$ と $\varphi : R/fR \to P$ の対 $(f, \varphi)$ 全体の集合とする。
$i \in I$ に対し、対応する対を $(f_i, \varphi_i)$ と書く。
直和因子 $R/f_iR$ の元 $r$ を $P$ の元 $\varphi_i(r)$ に送る写像
$$
B = \bigoplus\nolimits_{i \in I} R/f_iR \longrightarrow P
$$
を考える。$A = \Ker(B \to P)$ とする。列
$$
0 \to A \to B \to P \to 0
$$
が (2) のような完全列ならば (1) が成り立つ。
これを見るため $f \in R$ とし、$a \in A$ が $B$ において
$f b$ に写るとする。ほとんどすべての $r_i = 0$ であるように
$b = (r_i)_{i \in I}$ と書く。このとき $P$ において
$$
f\sum \varphi_i(r_i) = 0
$$
である。従って $f_{i_0} = f$ かつ
$\varphi_{i_0}(1) = \sum \varphi_i(r_i)$ となる $i_0 \in I$ が存在する。
$x_{i_0} \in R/f_{i_0}R$ を $1$ の類とする。このとき
$$
a' = (r_i)_{i \in I} - (0, \ldots, 0, x_{i_0}, 0, \ldots )
$$
は $A$ の元であり、所望の通り $fa' = a$ である。
\end{proof}

\begin{lemma}[一般化付値環]
\label{more-algebra-lemma-generalized-valuation-ring}
\begin{reference}
\cite{Warfield-Decomposition}
\end{reference}
$R$ を零でない環とする。次は同値である。
\begin{enumerate}
\item $a, b \in R$ に対して、$a$ が $b$ を割るか $b$ が $a$ を割る。
\item すべての有限生成イデアルは単項であり、$R$ は局所環である。
\item $R$ のイデアル全体は包含関係により全順序付けられる。
\end{enumerate}
特に $R$ が付値環ならば、これらは成り立つ。
\end{lemma}

\begin{proof}
(2) を仮定し、$a, b \in R$ とする。このとき $(a, b) = (c)$ である。
$c = 0$ ならば $a = b = 0$ であり、$a$ は $b$ を割る。
$c \not = 0$ と仮定する。
$c = ua + vb$, $a = wc$, $b = zc$ と書く。
すると $c(1 - uw - vz) = 0$ である。$R$ は局所環なので、
これは $1 - uw - vz \in \mathfrak m$ を意味する。
従って $w$, $z$ のいずれかは単元であり、$a$ が $b$ を割るか
$b$ が $a$ を割る。よって (2) は (1) を導く。

\medskip\noindent
(1) を仮定する。$R$ が二つの極大イデアル $\mathfrak m_i$ をもつならば、
$a \in \mathfrak m_1$, $a \not \in \mathfrak m_2$ と
$b \in \mathfrak m_2$, $b \not \in \mathfrak m_1$ を選べる。
このとき $a$ は $b$ を割らず、$b$ も $a$ を割らない。
従って $R$ は一意な極大イデアルをもち、局所環である。
条件 (1) と帰納法から、すべての有限生成イデアルが単項であることが容易に従う。
よって (1) は (2) を導く。

\medskip\noindent
(1) と (3) が同値であることは直ちに示せる。
最後の主張は Algebra, Lemma
\ref{algebra-lemma-valuation-ring-x-or-x-inverse} である。
\end{proof}

\begin{lemma}
\label{more-algebra-lemma-generalized-valuation-ring-modules}
\begin{reference}
\cite[Theorem 1]{Warfield-Decomposition}
\end{reference}
$R$ を Lemma \ref{more-algebra-lemma-generalized-valuation-ring} の
同値な条件を満たす環とする。このときすべての有限表示
$R$-加群は、$R/fR$ の形の加群の有限直和と同型である。
\end{lemma}

\begin{proof}
$M$ を有限表示 $R$-加群とする。以後断りなく
Lemma \ref{more-algebra-lemma-generalized-valuation-ring} にある $R$ の
同値な性質をすべて用いる。極大イデアルを
$\mathfrak m \subset R$、剰余体を
$\kappa = R/\mathfrak m$ と書く。
$I \subset R$ を $M$ の零化イデアルとする。
有限次元 $\kappa$-ベクトル空間 $M/\mathfrak m M$ の基底
$y_1, \ldots, y_n$ を選ぶ。$n$ に関する帰納法で議論する。

\medskip\noindent
Nakayama の補題により、$M/\mathfrak m M$ において
$y_1, \ldots, y_n$ を持ち上げる任意の元
$x_1, \ldots, x_n \in M$ は $M$ を生成する。
Algebra, Lemma \ref{algebra-lemma-NAK} を見よ。
これは帰納法の基底の場合 $n = 0$ を直ちに証明する。

\medskip\noindent
任意の $y_i$ に写る $x_i \in M$ の零化イデアルが $I$ となるような
添字 $i$ が存在すると主張する。この $x_i$ の零化イデアルについて、
実際、そうでなければ
$i$ のすべてに対して
$I_i = \text{Ann}(x_i) \not = I$ となる
$x_1, \ldots, x_n$ を選べる。しかしすべての $i$ に対し
$I \subset I_i$ であり、イデアルが全順序付けられているので、
$i = 1, \ldots, n$ に対し $I_i$ は $I$ より真に大きい。
再び全順序性により
$\text{Ann}(M) = I_1 \cap \ldots \cap I_n$ は $I$ より大きくなり、
矛盾である。番号を付け直して、$y_1$ が次の性質をもつと仮定してよい。
$y_1$ を持ち上げる任意の $x_1 \in M$ について、
この $x_1$ の零化イデアルは $I$ である。

\medskip\noindent
$A = Rx_1 \subset M$ と置く。完全列
$0 \to A \to M \to M/A \to 0$ を考える。
$A$ は有限なので、$M/A$ は生成元がより少ない有限表示 $R$-加群である
（Algebra, Lemma \ref{algebra-lemma-extension}）。
帰納法により
$M/A \cong \bigoplus_{j = 1, \ldots, m} R/f_jR$ である。
一方、$A \to M$ は次の性質を満たすと主張する。
$f \in R$ ならば $fA = A \cap fM$ である。
包含 $fA \subset A \cap fM$ は明らかである。
逆に $x \in A \cap fM$ ならば、ある $g \in R$, $y \in M$ に対し
$x = gx_1 = f y$ である。$f$ が $g$ を割るならば所望の通り
$x \in fA$ である。そうでなければ、ある $h \in \mathfrak m$ に対し
$f = hg$ と書ける。前段落により、元
$x'_1 = x_1 - hy$ の零化イデアルは $I$ である。
従って $g \in I$ であり、所望の通り $x = 0$ である。
この主張と Lemma \ref{more-algebra-lemma-characterize-PD-modules} により、
列 $0 \to A \to M \to M/A \to 0$ は分裂し、
$M \cong A \oplus \bigoplus_{j = 1, \ldots, m} R/f_jR$ を得る。
すると $A = R/I$ は（$M$ の直和因子として）有限表示なので、
$I$ は有限生成、従って単項である。これで証明が完了する。
\end{proof}

\begin{lemma}
\label{more-algebra-lemma-warfield}
\begin{reference}
\cite[Theorem 3]{Warfield-Decomposition}
\end{reference}
$R$ を、各極大イデアルにおける $R$ のすべての局所環が
Lemma \ref{more-algebra-lemma-generalized-valuation-ring} の同値な条件を満たす環とする。
このときすべての有限表示 $R$-加群は、$f$ を $R$ の中で動かした
$R/fR$ の形の加群の有限直和の直和因子である。
\end{lemma}

\begin{proof}
$M$ を有限表示 $R$-加群とする。まず $M$ が $R/fR$ の形の加群の
直和の直和因子であることを示し、最後に直和を有限に取れることを論じる。
すべての $f \in R$ に対して $fA = A \cap fB$ を満たす
$R$-加群の短完全列
$$
0 \to A \to B \to C \to 0
$$
を取る。Lemma \ref{more-algebra-lemma-characterize-PD-modules} により、
$\Hom_R(M, B) \to \Hom_R(M, C)$ が全射であることを示せばよい。
これは極大イデアル $\mathfrak m$ における局所化後に証明すれば十分である。
Algebra, Lemma \ref{algebra-lemma-characterize-zero-local} を見よ。
局所化した列
$0 \to A_\mathfrak m \to B_\mathfrak m \to C_\mathfrak m \to 0$
は、すべての $f \in R_\mathfrak m$ に対し
$fA_\mathfrak m = A_\mathfrak m \cap fB_\mathfrak m$ という条件を満たす。
実際 $f = uf'$ と書ける。ここで
$u \in R_\mathfrak m$ は単元、$f' \in R$ であり、局所化は完全だからである。
$M$ は有限表示なので、Algebra, Lemma
\ref{algebra-lemma-hom-from-finitely-presented} により
$$
\Hom_R(M, B)_\mathfrak m = \Hom_{R_\mathfrak m}(M_\mathfrak m, B_\mathfrak m)
\quad\text{and}\quad
\Hom_R(M, C)_\mathfrak m = \Hom_{R_\mathfrak m}(M_\mathfrak m, C_\mathfrak m)
$$
である。加群 $M_\mathfrak m$ は有限表示 $R_\mathfrak m$-加群である。
Lemma \ref{more-algebra-lemma-generalized-valuation-ring-modules} により
$M_\mathfrak m$ は $R_\mathfrak m/fR_\mathfrak m$ の形の
加群の直和である。従って Lemma \ref{more-algebra-lemma-characterize-PD-modules} により、
局所化上の写像は全射である。

\medskip\noindent
この時点で $M$ は $\bigoplus_{i \in I} R/f_i R$ の直和因子である。
写像 $M \to \bigoplus_{i \in I} R/f_i R$ を考える。
$M$ は有限 $R$-加群なので、ある有限部分集合 $I' \subset I$ に対し
像は $\bigoplus_{i \in I'} R/f_i R$ に含まれる。
これで証明が完了する。
\end{proof}

\begin{definition}
\label{more-algebra-definition-bezout}
$R$ を整域とする。
\begin{enumerate}
\item $R$ のすべての有限生成イデアルが単項であるとき、
$R$ を {\it B\'ezout 整域} という。
\item すべての $n , m \geq 1$ と任意の $n \times m$ 行列 $A$ に対し、
サイズ $n \times n, m \times m$ の可逆行列 $U, V$ で
$$
U A V =
\left(
\begin{matrix}
f_1 & 0 & 0 & \ldots \\
0 & f_2 & 0 & \ldots \\
0 & 0 & f_3 & \ldots \\
\ldots & \ldots & \ldots & \ldots
\end{matrix}
\right)
$$
かつ $f_1, \ldots, f_{\min(n, m)} \in R$,
$f_1 | f_2 | \ldots$ となるものが存在するとき、
$R$ を {\it 単因子整域} という。
\end{enumerate}
\end{definition}

\noindent
すべての B\'ezout 整域 $R$ が単因子整域であるか否かは、
どうやら依然として未解決問題である。
これは $R$ 上のすべての有限表示加群が巡回加群の直和であるか
という問題と同値である。逆の含意は真である。

\begin{lemma}
\label{more-algebra-lemma-elementary-divisor-is-bezout}
単因子整域は B\'ezout 整域である。
\end{lemma}

\begin{proof}
$a, b \in R$ を零でない元とする。$1 \times 2$ 行列 $A = (a\ b)$ を考える。
このとき $u \in R$ は単元、$V = (g_{ij})$ は可逆
$2 \times 2$ 行列として $u(a\ b)V = (f\ 0)$ となる。
従って
$f = u a g_{11} + u b g_{2 1}$ かつ $(g_{11}, g_{2 1}) = R$ である。
よって $(a, b) = (f)$ である。帰納法（省略する）により、
$R$ の任意の有限生成イデアルは一つの元で生成される。
\end{proof}

\begin{lemma}
\label{more-algebra-lemma-localize-bezout}
B\'ezout 整域の局所化は B\'ezout 整域である。
B\'ezout 整域のすべての局所環は付値環である。
局所整域が B\'ezout 整域であることと付値環であることは同値である。
\end{lemma}

\begin{proof}
局所化についての主張の証明は省略する。最後の主張は
Algebra, Lemma \ref{algebra-lemma-characterize-valuation-ring} である。
第二の主張は他の二つから従う。
\end{proof}

\begin{lemma}
\label{more-algebra-lemma-split-off-free-part}
$R$ を B\'ezout 整域とする。
\begin{enumerate}
\item 自由加群のすべての有限部分加群は有限自由である。
\item すべての有限表示 $R$-加群 $M$ は、有限自由加群と
ねじれ加群 $M_{tors}$ の直和である。ここでこのねじれ加群は、
零でない $f_1, \ldots, f_n \in R$ に対する形
$\bigoplus_{i = 1, \ldots, n} R/f_iR$ の加群の直和因子である。
\end{enumerate}
\end{lemma}

\begin{proof}
(1) を証明する。$M \subset F$ を自由加群 $F$ の有限部分加群とする。
$M$ は有限なので、$F$ は有限自由加群であると仮定してよい
（詳細は省略する）。$F = R^{\oplus n}$ とし、$n$ に関する帰納法で論じる。
$n = 1$ ならば $M$ は有限生成イデアルなので、
$R$ が B\'ezout であるという仮定により単項である。
$n > 1$ ならば、最後の直和因子への射影
$R^{\oplus n} \to R$ の下での $M$ の像 $I$ を考える。
$I = (0)$ ならば $M \subset R^{\oplus n - 1}$ であり、
帰納法により終わる。$I \not = 0$ ならば $I = (f) \cong R$ である。
従って $M \cong R \oplus \Ker(M \to I)$ であり、
やはり帰納法により終わる。

\medskip\noindent
$M$ を有限表示 $R$-加群とする。
% P02-E1332: source line 37872 says “localizations of R are maximal ideals”; “at maximal ideals” is intended.
$R$ の極大イデアルにおける局所化は付値環なので
（Lemma \ref{more-algebra-lemma-localize-bezout}）、Lemma \ref{more-algebra-lemma-warfield} を適用できる。
従って $M$ は、$f_i \not = 0$ として
$R^{\oplus r} \oplus \bigoplus_{i = 1, \ldots, n} R/f_iR$
の形の加群の直和因子である。
ねじれ部分加群を取ることは関手なので、
$M_{tors}$ は加群 $\bigoplus_{i = 1, \ldots, n} R/f_iR$ の直和因子であり、
$M/M_{tors}$ は $R^{\oplus r}$ の直和因子である。
証明の第一部により $M/M_{tors}$ は有限自由である。
従って所望の通り $M \cong M_{tors} \oplus M/M_{tors}$ である。
\end{proof}

\begin{lemma}
\label{more-algebra-lemma-modules-PID}
$R$ を PID とする。すべての有限 $R$-加群 $M$ は、
% P02-E1333: source line 37888 says “is of isomorphic”.
ある $r, n \geq 0$ と零でない $f_1, \ldots, f_n \in R$ に対する形
$$
R^{\oplus r} \oplus \bigoplus\nolimits_{i = 1, \ldots, n} R/f_iR
$$
の加群と同型である。
\end{lemma}

\begin{proof}
PID は Noether B\'ezout 環である。
Lemma \ref{more-algebra-lemma-split-off-free-part} により、
$M$ がねじれ加群である場合に結果を証明すれば十分である。
$M$ は有限なので、これは $M$ の零化イデアルが零でないことを意味する。
ある零でない $f \in R$ に対して $fM = 0$ とする。
このとき $M$ を $R/fR$ 上の加群と見なせる。
$R/fR$ は次元 $0$ の Noether 環なので（小さな詳細は省略する）、
$R/fR = \prod R_j$ は Artin 局所環 $R_i$ の有限積である
% P02-E1334: source lines 37902-37903 switch from factor index j to i.
（Algebra, Proposition \ref{algebra-proposition-dimension-zero-ring}）。
各 $R_i$ は局所環かつ PID の商なので、
Lemma \ref{more-algebra-lemma-generalized-valuation-ring} の意味で一般化付値環である
（小さな詳細は省略する）。
$R_j$ に対応する冪等元を $e_j \in R/fR$ として、
$M_j = e_j M$ となるように $M = \prod M_j$ と書く。
Lemma \ref{more-algebra-lemma-generalized-valuation-ring-modules} により、
ある $\overline{f}_{ji} \in R_j$ に対して
$M_j = \bigoplus_{i = 1, \ldots, n_j} R_j/\overline{f}_{ji}R_j$ である。
持ち上げ $f_{ji} \in R$ を選び、
% P02-E1335: source line 37913 uses undefined f_j when defining g_{ji}.
$(g_{ji}) = (f_j, f_{ji})$ となる $g_{ji} \in R$ を選ぶ。
すると
$$
M \cong \bigoplus R/g_{ji}R
$$
であると結論できる。これは $R$-加群としての同型であり、証明を完了する。
\end{proof}

\noindent
% P02-E1336: source line 37922 says “a elementary divisor domain” instead of “an”.
PID が単因子整域であることも、次のような補題を証明することで示せる
% P02-E1337: source lines 37922-37923 retain the unresolved placeholder “insert future reference here”.
（将来の参照をここに挿入）。

\begin{lemma}
\label{more-algebra-lemma-unimodular-vector}
$R$ を B\'ezout 整域とする。$n \geq 1$ とし、
$f_1, \ldots, f_n \in R$ は単位イデアルを生成するとする。
第一行が $f_1 \ldots f_n$ である $R$ 上の可逆
$n \times n$ 行列が存在する。
\end{lemma}

\begin{proof}
これは Lemma \ref{more-algebra-lemma-split-off-free-part} から従うが、
次のように直接証明することもできる。$n$ に関する帰納法を用いる。
$n = 1$ の場合は成り立つ。$n > 1$ と仮定する。
番号を付け直して $f_1 \not = 0$ と仮定してよい。
$(f) = (f_1, \ldots, f_{n - 1})$ となる $f \in R$ を選ぶ。
% P02-E1338: source lines 37938-37939 do not say that A is invertible, although induction supplies and the proof needs an invertible matrix.
第一行が $f_1/f, \ldots, f_{n - 1}/f$ である
$(n - 1) \times (n - 1)$ 行列 $A$ を取る。
$1 \in (f_1, \ldots, f_n) = (f, f_n)$ なので、
$af - bf_n = 1$ となる $a, b \in R$ を選べる。このとき解は行列
$$
\left(
\begin{matrix}
f & 0 & \ldots & 0 & f_n \\
0 & 1 & \ldots & 0 & 0 \\
  &   & \ldots \\
0 & 0 & \ldots & 1 & 0 \\
b & 0 & \ldots & 0 & a
\end{matrix}
\right)
\left(
\begin{matrix}
  & & & 0 \\
  & A \\
  & & & 0 \\
0 & \ldots & 0 & 1
\end{matrix}
\right)
$$
である。左の行列は行列式が $1$ なので可逆であることに注意する。
\end{proof}







\section{単項根基イデアル}
\label{more-algebra-section-principal-radical-ideals}

\noindent
% P02-E1339: source lines 37974-37975 omit “in” in “that a ... domain there exists”.
この節では、鎖状 Noether 局所正規整域には非自明な単項根基イデアルが
存在することを証明する。この結果は \cite{Artin-Lipman} にある。

\begin{lemma}
\label{more-algebra-lemma-polypoly}
$(R,\mathfrak m)$ を次元一の Noether 局所環とし、
$x\in\mathfrak m$ を $R$ のいずれの極小素イデアルにも含まれない元とする。
このとき
\begin{enumerate}
\item 関数 $P : n \mapsto \text{length}_R(R/x^n R)$ は
$n \geq 0$ に対して $P(n) \leq n P(1)$ を満たす、
\item $x$ が非零因子ならば、$n \geq 0$ に対して
$P(n) = nP(1)$ である。
\end{enumerate}
\end{lemma}

\begin{proof}
$\dim(R) = 1$ なので $\dim(R/x^n R) = 0$ であり、
各 $n$ に対して $\text{length}_R(R/x^n R)$ は有限である
（Algebra, Lemma \ref{algebra-lemma-support-point}）。
$n$ に関する帰納法で補題を示す。$x^0 R = R$ なので
$P(0) = \text{length}_R(R/x^0R) = \text{length}_R 0 = 0$ である。
$n = 1$ に対しても主張は成り立つ。
$n \geq 2$ とし、$n - 1$ に対して主張が成り立つと仮定する。
列
$$
R/x^{n-1}R \xrightarrow{x} R/x^nR \to R/xR \to 0
$$
は完全である。ここで $x$ は $x$ による乗法写像を表す。
長さは加法的なので
（Algebra, Lemma \ref{algebra-lemma-length-additive}）、
$P(n) \leq P(n - 1) + P(1)$ である。
帰納法により $P(n - 1) \leq (n - 1)P(1)$ だから
$P(n) \leq nP(1)$ である。これで帰納段階が示された。

\medskip\noindent
$x$ が非零因子ならば、上の表示された完全列は左でも完全である。
従ってすべての $n \geq 1$ に対し
$P(n) = P(n - 1) + P(1)$ を得る。
\end{proof}

\begin{lemma}
\label{more-algebra-lemma-minprimespoly}
$(R, \mathfrak m)$ を次元 $1$ の Noether 局所環とする。
$x \in \mathfrak m$ を $R$ のいずれの極小素イデアルにも含まれない元とし、
$t$ を $R$ の極小素イデアルの個数とする。このとき
$t \leq \text{length}_R(R/xR)$ である。
\end{lemma}

\begin{proof}
$\mathfrak p_1, \ldots, \mathfrak p_t$ を $R$ の極小素イデアルとする。
$R' = R/\sqrt{0} = R/(\bigcap_{i = 1}^t \mathfrak p_i)$ と置く。
$R'$ に対して補題を証明すれば十分であると主張する。
実際 $R'$ も $t$ 個の極小素イデアルをもち、
$\text{length}_{R'}(R'/xR') = \text{length}_R(R'/xR')$ は
全射 $R/xR \to R'/xR'$ があるため
$\text{length}_R(R/xR)$ より小さい。従って $R$ は被約と仮定してよい。

\medskip\noindent
$R$ は被約で、極小素イデアルが
$\mathfrak p_1, \ldots, \mathfrak p_t$ であると仮定する。
これは完全列
$$
0 \to R \to
\prod\nolimits_{i = 1}^t R/\mathfrak p_i \to Q \to 0
$$
があることを意味する。ここで $Q$ は第一の写像の余核である。
$M = \prod_{i = 1}^t R/\mathfrak p_i$ と書く。
$\mathfrak p_j$ で局所化すると
$$
R_{\mathfrak p_j} \to M_{\mathfrak p_j} =
\left(\prod\nolimits_{i=1}^t R/\mathfrak p_i\right)_{\mathfrak p_j} =
(R/\mathfrak p_j)_{\mathfrak p_j}
$$
は全射である。従ってすべての $j$ に対し
$Q_{\mathfrak p_j} = 0$ である。
$\mathfrak m$ は $\mathfrak p_i$ と異なる $R$ の唯一の素イデアルなので、
$\text{Supp}(Q) = \{\mathfrak m\}$ と結論する。
従って $Q$ は有限長である
（Algebra, Lemma \ref{algebra-lemma-support-point}）。
$\text{Supp}(Q) = \{\mathfrak m\}$ なので、
$x^n$ が $Q$ 上 $0$ として作用するような $n \gg 0$ を選べる
（Algebra, Lemma
\ref{algebra-lemma-Noetherian-power-ideal-kills-module}）。
図式
$$
\xymatrix{
0 \ar[r] & R \ar[r] \ar[d]^-{x^n} & M
\ar[r] \ar[d]^-{x^n} & Q \ar[r] \ar[d]^-{x^n} & 0 \\
0 \ar[r] & R \ar[r] & M \ar[r] & Q \ar[r] & 0
}
$$
を考える。ここで垂直写像は $x^n$ による乗法である。
$x$ はいずれの $\mathfrak p_i$ にも含まれないので、
これは $R$ 上および $M$ 上で単射である。
蛇の補題（Algebra, Lemma \ref{algebra-lemma-snake}）により列
$$
0 \to Q \to R/x^nR \to M/x^nM \to Q \to 0
$$
は完全である。従って十分大きい $n$ に対して
$\text{length}_R(R/x^nR) = \text{length}_R(M/x^nM)$ である。
$R_i = R/\mathfrak p_i$ と書けば
$\text{length}(M/x^nM) =
\sum_{i = 1}^t \text{length}_R(R_i/x^nR_i)$ である。
Lemma \ref{more-algebra-lemma-polypoly} と、$x$ が $R$ および $R_i$ 上
非零因子であることを用いると
$$
n \text{length}_R(R/xR) =
\sum\nolimits_{i = 1}^t n \text{length}_{R_i}(R_i/x R_i)
$$
を得る。$\text{length}_{R_i}(R_i/x R_i) \geq 1$ なので補題が従う。
\end{proof}

\begin{lemma}
\label{more-algebra-lemma-sopexists}
$(R,\mathfrak m)$ を次元 $d > 1$ の Noether 局所環とし、
$f \in \mathfrak m$ を $R$ のいずれの極小素イデアルにも
含まれない元、$k\in\mathbf{N}$ とする。このとき
$f, g_1, \ldots, g_{d - 1}$ がパラメータ系となるような元
$g_1, \ldots, g_{d - 1} \in \mathfrak m^k$ が存在する。
\end{lemma}

\begin{proof}
Algebra, Lemma \ref{algebra-lemma-one-equation} により
$\dim(R/fR) = d - 1$ である。
$R/fR$ のパラメータ系
$\overline{g}_1, \ldots, \overline{g}_{d - 1}$ を選び
（Algebra, Proposition \ref{algebra-proposition-dimension}）、
$R$ における持ち上げ $g_1, \ldots, g_{d - 1}$ を取る。
$f, g_1, \ldots, g_{d - 1}$ が $R$ のパラメータ系であることは
直ちに分かる。すると
$f, g_1^k, \ldots, g_{d - 1}^k$ もパラメータ系であり、証明が完了する。
\end{proof}

\begin{lemma}
\label{more-algebra-lemma-syspar}
$(R,\mathfrak m)$ を次元二の Noether 局所環とし、
$f \in \mathfrak m$ を $R$ のいずれの極小素イデアルにも
含まれない元とする。このとき $g \in \mathfrak m$ と
$N \in \mathbf{N}$ で次を満たすものが存在する。
\begin{enumerate}
\item[(a)] $f,g$ は $R$ のパラメータ系をなす。
\item[(b)] $h \in \mathfrak m^N$ ならば
$f + h, g$ はパラメータ系であり、
$\text{length}_R (R/(f, g)) = \text{length}_R(R/(f + h, g))$ である。
\end{enumerate}
\end{lemma}

\begin{proof}
Lemma \ref{more-algebra-lemma-sopexists} により、$f, g$ が $R$ のパラメータ系となる
$g \in \mathfrak m$ が存在する。このとき
$\mathfrak m = \sqrt{(f, g)}$ である。
従って Algebra, Lemma \ref{algebra-lemma-Noetherian-power} により、
$\mathfrak m^n \subset (f, g)$ となる $n$ が存在する。
$N = n + 1$ でよいと主張する。
$h \in \mathfrak m^N$ とする。$N$ の選び方により
$a, b \in \mathfrak m$ を用いて $h = af + bg$ と書ける。
従って
$$
(f + h, g) = (f + af + bg, g) = ((1 + a)f, g) = (f, g)
$$
である。$1 + a$ は $R$ の単元だからである。
これは長さの等式と $f + h, g$ がパラメータ系であることを証明する。
\end{proof}

\begin{lemma}
\label{more-algebra-lemma-radical-element}
$R$ を次元 $2$ の Noether 局所正規整域とする。
$\mathfrak p_1, \ldots, \mathfrak p_r$ を相異なる高さ $1$ の
素イデアルとする。このとき
$f \in \mathfrak p_1 \cap \ldots \cap \mathfrak p_r$ で零でなく、
$R/fR$ が被約となるものが存在する。
\end{lemma}

\begin{proof}
零でない元
$f \in \mathfrak p_1 \cap \ldots \cap \mathfrak p_r$ を取る。
$f$ を少し修正して根基イデアルを生成する元を得る。
高さ一の各素イデアル $\mathfrak p$ における $R$ の局所化
$R_\mathfrak p$ は離散付値環である。Algebra, Lemma
\ref{algebra-lemma-characterize-dvr} または Algebra, Lemma
\ref{algebra-lemma-criterion-normal} を見よ。
$R_{\mathfrak p}$ における $f$ の対応する付値を
$\text{ord}_\mathfrak p(f)$ と書く。
$f$ を含む相異なる高さ一の素イデアルを
$\mathfrak q_1, \ldots, \mathfrak q_s$ とする。
各 $j$ に対し $\text{ord}_{\mathfrak q_j}(f) = m_j \geq 1$ と書く。
そして整数係数をもつ高さ一の素イデアルの形式的線形結合として
$\text{div}(f) = \sum_{j = 1}^s m_j\mathfrak q_j$ を定義する。
後で用いるため、各 $\mathfrak p_i$ は $\mathfrak q_j$ のいずれかとして
現れることに注意する。$R/fR$ が被約であることと、
$j = 1, \ldots, s$ に対し $m_j = 1$ であることは同値である。
実際 $m_j$ が $1$ ならば
% P02-E1340: source line 38165 omits the localization subscript in (R/fR) q_j.
$(R/fR)\mathfrak q_j$ は被約であり、
% P02-E1341: source line 38167 ends the associated-prime list at q_j instead of q_s.
$\mathfrak q_1, \ldots, \mathfrak q_j$ は $R/fR$ の随伴素イデアルなので
$R/fR \subset \prod (R/fR)_{\mathfrak q_j}$ である。
Algebra, Lemmas \ref{algebra-lemma-zero-at-ass-zero} および
\ref{algebra-lemma-normal-domain-intersection-localizations-height-1} を見よ。

\medskip\noindent
Lemma \ref{more-algebra-lemma-syspar} のような $g$, $N$ を選んで固定する。
零でない $y \in R$ に対し、$y$ 上極小な素イデアルの個数を
$t(y)$ と書く。$R$ は正規整域なので、これらの素イデアルは高さ一であり、
$R/yR$ の極小素イデアルと $1$ 対 $1$ に対応する
（Algebra, Lemmas \ref{algebra-lemma-minimal-over-1} および
\ref{algebra-lemma-normal-domain-intersection-localizations-height-1}）。
例えば $t(f) = s$ は $\text{div}(f)$ に現れる
素イデアル $\mathfrak q_j$ の個数である。
$h \in \mathfrak m^N$ とする。Lemma \ref{more-algebra-lemma-minprimespoly} により
\begin{align*}
t(f + h) & \leq \text{length}_{R/(f + h)}(R/(f + h, g)) \\
& = \text{length}_R(R/(f + h, g)) \\
& = \text{length}_R(R/(f, g))
\end{align*}
である。第一の等式については Algebra, Lemma
\ref{algebra-lemma-length-independent} を見よ。
従って $t(f + h)$ は $h \in \mathfrak m^N$ によらず有界である。

\medskip\noindent
上で示した有界性により、
$h \in \mathfrak m^N \cap \mathfrak p_1 \cap \ldots \cap \mathfrak p_r$
で、そのような $h$ の中で $t(f + h)$ が最大となるものを選べる。
$f' = f + h$ と置く。
$h' \in \mathfrak m^N \cap \mathfrak p_1 \cap \ldots \cap \mathfrak p_r$
に対し $t(f' + h') \leq t(f + h)$ である。
従って $f$ を $f'$ で置き換え、すべての
$h \in \mathfrak m^N \cap \mathfrak p_1 \cap \ldots \cap \mathfrak p_r$
に対し $t(f + h) \leq s$ と仮定してよい。

\medskip\noindent
次に、各 $j$ に対し
$\text{ord}_{\mathfrak q_j}(h) \geq 1$ かつ
$\text{ord}_{\mathfrak q_j}(h) = 1 \Leftrightarrow m_j > 1$
となる元 $h \in \mathfrak m^N$ を見つけられると仮定する。
$h \in \mathfrak m^N \cap \mathfrak p_1 \cap \ldots \cap \mathfrak p_r$
であることに注意する。付値の基本的性質により、すべての $j$ に対し
$\text{ord}_{\mathfrak q_j}(f + h) = 1$ である。従って
$$
\text{div}(f + h) = \sum\nolimits_{j = 1}^s \mathfrak q_j +
\sum\nolimits_{k = 1}^v e_k \mathfrak r_k
$$
である。ここで $\mathfrak r_1, \ldots, \mathfrak r_v$ は
相異なる高さ一の素イデアル、$e_k \geq 1$ である。
しかし $s = t(f) \geq t(f + h)$ なので $v = 0$ であり、
所望の元が見つかった。

\medskip\noindent
上の条件を満たす $h$ を選ぶ。
素イデアル回避（Algebra, Lemma \ref{algebra-lemma-silly}）により、
各 $1 \leq j \leq s$ に対し
$a_j \in \mathfrak q_j$ で、
$j' \not = j$ に対して $a_j \not \in \mathfrak q_{j'}$ かつ
$a_j \not \in \mathfrak q_j^{(2)}$ となるものを見つけられる。
ここで
$\mathfrak q_j^{(2)} = \{x \in R \mid \text{ord}_{\mathfrak q_j}(x) \geq 2\}$
は $\mathfrak q_j$ の第二記号的冪である。そこで
$$
h = \prod\nolimits_{m_j = 1} a_j^2 \times
\prod\nolimits_{m_j > 1} a_j
$$
と取る。$h$ が上で課した付値の条件を満たすことは明らかである。
$h \not \in \mathfrak m^N$ ならば、すべての $j$ に対して
$\mathfrak q_j$ に含まれない $\mathfrak m^N$ の元を掛ける。
\end{proof}

\begin{lemma}
\label{more-algebra-lemma-divides-radical}
$(A, \mathfrak m, \kappa)$ を次元 $2$ の Noether 局所正規整域とする。
$a \in \mathfrak m$ が零でないならば、$A/cA$ が被約で、
ある $n$ に対し $a$ が $c^n$ を割るような元 $c \in A$ が存在する。
\end{lemma}

\begin{proof}
Lemma \ref{more-algebra-lemma-radical-element} の証明の記法を用いて
$\text{div}(a) = \sum_{i = 1, \ldots, r} n_i \mathfrak p_i$ とする。
Lemma \ref{more-algebra-lemma-radical-element} により、
$A/cA$ が被約となる
$c \in \mathfrak p_1 \cap \ldots \cap \mathfrak p_r$ を選ぶ。
$n \geq \max(n_i)$ に対して
$-\text{div}(a) + \text{div}(c^n)$ は有効因子
（すべての係数が非負）である。
従って Algebra, Lemma
\ref{algebra-lemma-normal-domain-intersection-localizations-height-1} により
$c^n/a \in A$ である。
\end{proof}

\noindent
この節の残りでは、次元 $> 2$ の結果を証明する。

\begin{lemma}
\label{more-algebra-lemma-multiplicity}
$(R, \mathfrak m)$ を次元 $d$ の Noether 局所環とし、
$g_1, \ldots, g_d$ をパラメータ系、
$I = (g_1, \ldots, g_d)$ とする。数値多項式
$n \mapsto \text{length}_R(R/I^{n+1})$ の最高次係数が
$e_I/d!$ ならば、$e_I \leq \text{length}_R(R/I)$ である。
\end{lemma}

\begin{proof}
この関数は Algebra, Proposition
\ref{algebra-proposition-hilbert-function-polynomial} により
数値多項式であり、Algebra, Proposition
\ref{algebra-proposition-dimension} により次数 $d$ である。
$d = 0$ ならば結果は自明であり、$d = 1$ ならば
Lemma \ref{more-algebra-lemma-polypoly} が結果を与える。
一般の場合を示すため、全射
$$
\bigoplus\nolimits_{i_1, \ldots, i_d \geq 0,\ \sum i_j = n} R/I
\longrightarrow
I^n/I^{n + 1}
$$
があることに注意する。これは $i_1, \ldots, i_d$ に対応する基底元を
$I^n/I^{n + 1}$ における $g_1^{i_1} \ldots g_d^{i_d}$ の類へ送る写像である。
従って
$$
\text{length}_R(R/I^{n + 1}) - \text{length}_R(R/I^n)
\leq  \text{length}_R(R/I) {n + d - 1 \choose d - 1}
$$
を得る。$d \geq 2$ なので、左辺の数値多項式は次数 $d - 1$、
最高次係数 $e_I / (d - 1)!$ をもつ。右辺の多項式も次数 $d - 1$ で、
その最高次係数は $\text{length}_R(R/I)/ (d - 1)!$ である。
これで補題が示された。
\end{proof}

\begin{lemma}
\label{more-algebra-lemma-minprimespolyhigher}
$(R, \mathfrak m)$ を次元 $d$ の Noether 局所環とし、
$t$ を次元 $d$ の $R$ の極小素イデアルの個数、
$(g_1,\ldots,g_d)$ をパラメータ系とする。このとき
% P02-E1342: source line 38301 uses g_n although the parameter list ends in g_d.
$t \leq \text{length}_R(R/(g_1,\ldots,g_n))$ である。
\end{lemma}

\begin{proof}
$d = 0$ ならば補題は自明である。$d = 1$ ならば補題は
Lemma \ref{more-algebra-lemma-minprimespoly} である。従って $d > 1$ と仮定してよい。
$\mathfrak p_1, \ldots, \mathfrak p_s$ を $R$ の極小素イデアルとし、
最初の $t$ 個が次元 $d$ をもつとする。また
% P02-E1343: source line 38309 defines I with g_n although the parameter list ends in g_d.
$I = (g_1, \ldots, g_n)$ と書く。
Lemma \ref{more-algebra-lemma-minprimespoly} の証明と全く同じ議論により、
$R$ は被約と仮定してよい。

\medskip\noindent
% P02-E1344: source lines 38313-38319 retain only the first t minimal primes after merely reducing R; when lower-dimensional minimal primes remain, the displayed map into the truncated product need not be injective.
$R$ は被約で、極小素イデアルが
$\mathfrak p_1, \ldots, \mathfrak p_t$ であると仮定する。
これは完全列
$$
0 \to R \to
\prod\nolimits_{i = 1}^t R/\mathfrak p_i \to Q \to 0
$$
があることを意味する。ここで $Q$ は第一の写像の余核である。
$M = \prod_{i = 1}^t R/\mathfrak p_i$ と書く。
$\mathfrak p_j$ で局所化すると
$$
R_{\mathfrak p_j} \to M_{\mathfrak p_j} =
\left(\prod\nolimits_{i=1}^t R/\mathfrak p_i\right)_{\mathfrak p_j} =
(R/\mathfrak p_j)_{\mathfrak p_j}
$$
は全射である。従ってすべての $j$ に対し $Q_{\mathfrak p_j} = 0$ である。
従って高さ $0$ の $R$ の素イデアルは $Q$ の台に含まれない。
ゆえに数値多項式
$n \mapsto \text{length}_R(Q/I^nQ)$ の次数は
$\dim(\text{Supp}(Q)) < d$ に等しい。Algebra, Lemma
\ref{algebra-lemma-support-dimension-d} を見よ。
Algebra, Lemma \ref{algebra-lemma-hilbert-ses-chi}
（$R$ は有限長でないため適用できる）により、多項式
$$
n \longmapsto
\text{length}_R(M/I^nM) - \text{length}_R(R/I^n) - \text{length}_R(Q/I^nQ)
$$
の次数は $< d$ である。$M = \prod R/\mathfrak p_i$ であり、また
% P02-E1345: source line 38340 uses n \to instead of the function notation n \mapsto.
$n \to \text{length}_R(R/\mathfrak p_i + I^n)$ は
$i = 1, \ldots, t$ に対し次数がまさに $d$ の数値多項式である
（Algebra, Lemma \ref{algebra-lemma-support-dimension-d}）から、
$n \mapsto \text{length}_R(M/I^nM)$ の最高次係数は少なくとも
$t/d!$ である。従って Lemma \ref{more-algebra-lemma-multiplicity} により結論を得る。
\end{proof}

\begin{lemma}
\label{more-algebra-lemma-sysparhigher}
$(R, \mathfrak m)$ を次元 $d$ の Noether 局所環とし、
$f \in \mathfrak m$ を $R$ のいずれの極小素イデアルにも
含まれない元とする。このとき元
$g_1, \ldots, g_{d - 1} \in \mathfrak m$ と $N \in \mathbf{N}$ で
次を満たすものが存在する。
\begin{enumerate}
\item $f, g_1, \ldots, g_{d - 1}$ は $R$ のパラメータ系をなす。
\item $h \in \mathfrak m^N$ ならば
$f + h, g_1, \ldots, g_{d - 1}$ はパラメータ系であり、
$$
\text{length}_R R/(f, g_1, \ldots, g_{d-1}) =
\text{length}_R R/(f + h, g_1, \ldots, g_{d-1})
$$
である。
\end{enumerate}
\end{lemma}

\begin{proof}
Lemma \ref{more-algebra-lemma-sopexists} により、
$f, g_1, \ldots, g_{d - 1}$ が $R$ のパラメータ系となる元
$g_1, \ldots, g_{d - 1} \in \mathfrak m$ が存在する。
このとき
$\mathfrak m = \sqrt{(f, g_1, \ldots, g_{d - 1})}$ である。
従って Algebra, Lemma \ref{algebra-lemma-Noetherian-power} により、
% P02-E1346: source line 38367 uses (f,g) although only g_1,...,g_{d-1} are defined here.
$\mathfrak m^n \subset (f, g)$ となる $n$ が存在する。
$N = n + 1$ でよいと主張する。
$h \in \mathfrak m^N$ とする。$N$ の選び方により、
$a, b_i \in \mathfrak m$ を用いて $h = af + \sum b_ig_i$ と書ける。
従って
\begin{align*}
(f + h, g_1, \ldots, g_{d - 1})
& =
(f + af + \sum b_ig_i, g_1, \ldots, g_{d - 1}) \\
& =
((1 + a)f, g_1, \ldots, g_{d - 1}) \\
& =
(f, g_1, \ldots, g_{d - 1})
\end{align*}
である。$1 + a$ は $R$ の単元だからである。
これは長さの等式と $f + h, g_1, \ldots, g_{d - 1}$ が
パラメータ系であることを証明する。
\end{proof}

\begin{proposition}
\label{more-algebra-proposition-propdimd}
\begin{reference}
\cite[Lemma 3.14]{Artin-Lipman} には、環が鎖状であるという
仮定なしにこの結果が述べられている。
\end{reference}
$R$ を鎖状 Noether 局所正規整域とする。
% P02-E1347: source proposition lines 38392-38395 omit J != 0; the stated case J = 0 is a counterexample and the proof immediately assumes a nonzero f in J.
$J \subset R$ を根基イデアルとする。このとき $R/fR$ が被約となる
零でない元 $f \in J$ が存在する。
\end{proposition}

\begin{proof}
証明は Lemma \ref{more-algebra-lemma-radical-element} の証明と同じであり、
Lemma \ref{more-algebra-lemma-minprimespolyhigher} を
Lemma \ref{more-algebra-lemma-minprimespoly} の代わりに、
Lemma \ref{more-algebra-lemma-sysparhigher} を
Lemma \ref{more-algebra-lemma-syspar} の代わりに用いる。
$R$ は鎖状整域なので $R$ の高さ一の各素イデアルは次元 $d - 1$ をもち、
従って $R/(f + h)$ のスペクトルは等次元である。このため
Lemma \ref{more-algebra-lemma-minprimespolyhigher} を用いることができる。
読者の便宜のため詳細を書き下す。

\medskip\noindent
零でない元 $f \in J$ を取る。
$f$ を少し修正して根基イデアルを生成する元を得る。
高さ一の各素イデアル $\mathfrak p$ における $R$ の局所化
$R_\mathfrak p$ は離散付値環である。Algebra, Lemma
\ref{algebra-lemma-characterize-dvr} または Algebra, Lemma
\ref{algebra-lemma-criterion-normal} を見よ。
$R_{\mathfrak p}$ における $f$ の対応する付値を
$\text{ord}_\mathfrak p(f)$ と書く。
$f$ を含む相異なる高さ一の素イデアルを
$\mathfrak q_1, \ldots, \mathfrak q_s$ とする。
各 $j$ に対し $\text{ord}_{\mathfrak q_j}(f) = m_j \geq 1$ と書く。
そして整数係数をもつ高さ一の素イデアルの形式的線形結合として
$\text{div}(f) = \sum_{j = 1}^s m_j\mathfrak q_j$ を定義する。
$R/fR$ が被約であることと、$j = 1, \ldots, s$ に対し
$m_j = 1$ であることは同値である。実際 $m_j$ が $1$ ならば
% P02-E1348: source line 38429 omits the localization subscript in (R/fR) q_j.
$(R/fR)\mathfrak q_j$ は被約であり、
% P02-E1349: source line 38431 ends the associated-prime list at q_j instead of q_s.
$\mathfrak q_1, \ldots, \mathfrak q_j$ は $R/fR$ の随伴素イデアルなので
$R/fR \subset \prod (R/fR)_{\mathfrak q_j}$ である。
Algebra, Lemmas \ref{algebra-lemma-zero-at-ass-zero} および
\ref{algebra-lemma-normal-domain-intersection-localizations-height-1} を見よ。

\medskip\noindent
% P02-E1350: source line 38437 starts the system at g_2 although Lemma sysparhigher supplies g_1,...,g_{d-1} and the formulas below use g_1.
Lemma \ref{more-algebra-lemma-sysparhigher} のような
$g_2, \ldots, g_{d - 1}$ と $N$ を選んで固定する。
零でない $y \in R$ に対し、$y$ 上極小な素イデアルの個数を
$t(y)$ と書く。$R$ は正規整域なので、これらの素イデアルは高さ一であり、
$R/yR$ の極小素イデアルと $1$ 対 $1$ に対応する
（Algebra, Lemmas \ref{algebra-lemma-minimal-over-1} および
\ref{algebra-lemma-normal-domain-intersection-localizations-height-1}）。
例えば $t(f) = s$ は $\text{div}(f)$ に現れる
素イデアル $\mathfrak q_j$ の個数である。
$h \in \mathfrak m^N$ とする。$R$ は鎖状なので、
$R$ の高さ一の各素イデアル $\mathfrak p$ に対し
% P02-E1351: source line 38448 says dim(R/p)=d for a height-one prime, contradicting d-1 stated at line 38407.
$\dim(R/\mathfrak p) = d$ である。従って
Lemma \ref{more-algebra-lemma-minprimespolyhigher} により
\begin{align*}
t(f + h) & \leq \text{length}_{R/(f + h)}(R/(f + h, g_1, \ldots, g_{d - 1})) \\
& = \text{length}_R(R/(f + h, g_1, \ldots, g_{d - 1})) \\
& = \text{length}_R(R/(f, g_1, \ldots, g_{d - 1}))
\end{align*}
である。第一の等式については Algebra, Lemma
\ref{algebra-lemma-length-independent} を見よ。
従って $t(f + h)$ は $h \in \mathfrak m^N$ によらず有界である。

\medskip\noindent
上で示した有界性により、$h \in \mathfrak m^N \cap J$ で、
そのような $h$ の中で $t(f + h)$ が最大となるものを選べる。
$f' = f + h$ と置く。$h' \in \mathfrak m^N \cap J$ に対し
$t(f' + h') \leq t(f + h)$ である。
従って $f$ を $f'$ で置き換え、すべての
$h \in \mathfrak m^N \cap J$ に対し
$t(f + h) \leq s$ と仮定してよい。

\medskip\noindent
次に、各 $j$ に対し
$\text{ord}_{\mathfrak q_j}(h) \geq 1$ かつ
$\text{ord}_{\mathfrak q_j}(h) = 1 \Leftrightarrow m_j > 1$
となる元 $h \in \mathfrak m^N \cap J$ を見つけられると仮定する。
付値の基本的性質により、すべての $j$ に対し
$\text{ord}_{\mathfrak q_j}(f + h) = 1$ である。従って
$$
\text{div}(f + h) = \sum\nolimits_{j = 1}^s \mathfrak q_j +
\sum\nolimits_{k = 1}^v e_k \mathfrak r_k
$$
である。ここで $\mathfrak r_1, \ldots, \mathfrak r_v$ は
相異なる高さ一の素イデアル、$e_k \geq 1$ である。
しかし $s = t(f) \geq t(f + h)$ なので $v = 0$ であり、
所望の元が見つかった。

\medskip\noindent
上の条件を満たす $h$ を選ぶ。
素イデアル回避（Algebra, Lemma \ref{algebra-lemma-silly}）により、
各 $1 \leq j \leq s$ に対し、
$a_j \in \mathfrak q_j \cap J$ で、$j' \not = j$ ならば
$a_j \not \in \mathfrak q_{j'}$ となるものを見つけられる。
次に
% P02-E1352: source lines 38491 and 38496 use bare q_s rather than \mathfrak q_s.
$b_j \in J \cap \mathfrak q_1 \cap \ldots \cap q_s$ で
$b_j \not \in \mathfrak q_j^{(2)}$ となるものを選べる。ここで
$\mathfrak q_j^{(2)} = \{x \in R \mid \text{ord}_{\mathfrak q_j}(x) \geq 2\}$
は $\mathfrak q_j$ の第二記号的冪である。素イデアル回避を適用できるのは、
イデアル $J' = J \cap \mathfrak q_1 \cap \ldots \cap q_s$ が根基的で、
従って $R/J'$ が被約、従って $(R/J')_{\mathfrak q_j}$ が被約であり、
従って $J'$ は $\text{ord}_{\mathfrak q_j}(x) = 1$ となる元 $x$ を含み、
従って $J' \not \subset \mathfrak q_j^{(2)}$ だからである。すると元
$$
c = \sum\nolimits_{j = 1, \ldots, s}
b_j \times \prod\nolimits_{j' \not = j} a_{j'}
$$
は $J$ の元であり、付値の基本的性質により
すべての $j = 1, \ldots, s$ に対し
$\text{ord}_{\mathfrak q_j}(c) = 1$ を満たす。最後に
$$
h = c \times \prod\nolimits_{m_j = 1} a_j \times y
$$
と置く。ここで $y \in \mathfrak m^N$ は、すべての $j$ に対し
$\mathfrak q_j$ に含まれない元である。
\end{proof}








\section{導来圏における可逆対象}
\label{more-algebra-section-invertible-D-or-R}

\noindent
環の導来圏における可逆対象を特徴づける。

\begin{lemma}
\label{more-algebra-lemma-symmetric-monoidal-derived}
$R$ を環とする。$R$ の導来圏 $D(R)$ は、導来テンソル積を
テンソル積とし、Section \ref{more-algebra-section-sign-rules} のような結合制約と
可換制約をもつ対称モノイド圏である。
\end{lemma}

\begin{proof}
省略する。ヒント：結合制約は Lemma \ref{more-algebra-lemma-triple-tensor-product} の
同型であり、可換制約は Lemma
\ref{more-algebra-lemma-flip-douoble-tensor-product} の同型である。
こう述べれば、種々の図式の可換性は $R$-加群の複体の圏に対する
対応する結果から従う。Section \ref{more-algebra-section-symmetric-monoidal} を見よ。
\end{proof}

\noindent
従って、$D(R)$ の対象が（左）双対をもつ、または可逆であるとは
何を意味するかが分かった。その具体的な内容を求める前に、
簡単な補題が必要である。

\begin{lemma}
\label{more-algebra-lemma-complex-bounded-above-free-colim-bounded-finite-free}
$R$ を環とし、$F^\bullet$ を自由 $R$-加群からなる上に有界な複体とする。
$n_i \in \mathbf{Z}$ および $f_i \in F^{n_i}$ を満たす対
$(n_i, f_i)$, $i = 1, \ldots, N$ が与えられたとする。
このとき、すべての $f_i$ を含み、有界で有限自由 $R$-加群からなる
部分複体 $G^\bullet \subset F^\bullet$ が存在する。
\end{lemma}

\begin{proof}
$a = \min(n_i; i = 1, \ldots, N)$ に関する下降帰納法を用いる。
$n \geq a$ に対して $F^n = 0$ ならば、$G^\bullet$ を零複体として
結果が成り立つ。一般の場合、番号を付け直して、
% P02-E1353: source line 38565 says “an 1 <= r” instead of “a 1 <= r”.
$n_1 = \ldots = n_r = a$ かつ $i > r$ に対し $n_i > a$ となる
$1 \leq r \leq N$ が存在すると仮定してよい。
$F^a$ の基底を $b_j, j \in J$ とする。
$i = 1, \ldots, r$ に対し
$f_i \in \bigoplus_{j \in J'} Rb_j$ となる有限部分集合
$J' \subset J$ を選べる。$F^{a + 1}$ の基底を $c_k, k \in K$ とする。
$j \in J'$ に対し
$\text{d}_F^a(b_j) \in \bigoplus_{k \in K'} Rc_k$ となる有限部分集合
$K' \subset K$ を選べる。そこで帰納法の仮定を適用し、
$k \in K'$ に対する $c_k \in F^{a + 1}$ と、$i > r$ に対する
$f_i \in F^{n_i}$ とを含む部分複体
$H^\bullet \subset F^\bullet$ を見つけられる。
次数 $> a$ では $H^\bullet$ に等しく、次数 $a$ では
$\bigoplus_{j \in J'} Rb_j$ に等しいものを $G^\bullet$ とする。
\end{proof}

\begin{lemma}
\label{more-algebra-lemma-have-dual-derived}
$R$ を環とし、$M$ を $D(R)$ の対象とする。次は同値である。
\begin{enumerate}
\item Categories, Definition \ref{categories-definition-dual} の意味で
$M$ は $D(R)$ における左双対をもつ。
\item $M$ は $D(R)$ の完全対象である。
\end{enumerate}
さらにこの場合、$M$ の左双対は Lemma
\ref{more-algebra-lemma-dual-perfect-complex} の対象 $M^\vee$ である。
\end{lemma}

\begin{proof}
$M$ が完全ならば、有限射影 $R$-加群からなる有界複体
$M^\bullet$ で $M$ を表せる。この場合、Lemma
\ref{more-algebra-lemma-left-dual-complex} により $M^\bullet$ は複体の圏で左双対をもち、
まして $D(R)$ でも左双対をもつ。

\medskip\noindent
(1) を仮定する。Categories, Definition \ref{categories-definition-dual} の
意味で $N$, $\eta : R \to M \otimes_R^\mathbf{L} N$, および
% P02-E1354: source lines 38601-38603 types epsilon as M tensor N -> R, but the displayed triangle identity and diagram require epsilon: N tensor M -> R.
$\epsilon : M \otimes_R^\mathbf{L} N \to R$ が左双対を与えるとする。
$M$ を表す複体 $M^\bullet$ を選ぶ。
% P02-E1355: source line 38605 has the number-agreement error “a K-flat complexes”.
平坦な項をもつ K-平坦複体 $N^\bullet$ で $N$ を表すものを選ぶ。
Lemma \ref{more-algebra-lemma-K-flat-resolution} を見よ。このとき $\eta$ は複体の写像
$$
\eta : R \longrightarrow \text{Tot}(M^\bullet \otimes_R N^\bullet)
$$
で与えられる。$1$ の像を有限和
$$
\eta(1) = \sum\nolimits_n \sum\nolimits_i m_{n, i} \otimes n_{-n, i}
$$
と書ける。ここで $m_{n, i} \in M^n$ かつ
$n_{-n, i} \in N^{-n}$ である。すべての元 $m_{n, i}$ と
$\text{d}(m_{n, i})$ で生成される部分複体を
$K^\bullet \subset M^\bullet$ とする。$N^\bullet$ の選び方により
$\text{Tot}(K^\bullet \otimes_R N^\bullet) \subset
\text{Tot}(M^\bullet \otimes_R N^\bullet)$ であり、上の選び方により
$\eta(1)$ はこの部分複体に含まれる。
$K^\bullet$ が表す $D(R)$ の対象を $K$ と書く。
すると $\eta$ は写像
$\tilde \eta : R \longrightarrow K \otimes_R^\mathbf{L} N$ を経由する。
$$
(1 \otimes \epsilon) \circ (\eta \otimes 1) = \text{id}_M
$$
なので、次の可換図式により $M$ 上の恒等写像は $K$ を経由する。
$$
\xymatrix{
M \ar[rr]_-{\eta \otimes 1} \ar[rrd]_{\tilde \eta \otimes 1} & &
M \otimes_R^\mathbf{L} N \otimes_R^\mathbf{L} M
\ar[r]_-{1 \otimes \epsilon} &
M \\
& &
K \otimes_R^\mathbf{L} N \otimes_R^\mathbf{L} M
\ar[u] \ar[r]^-{1 \otimes \epsilon} &
K \ar[u]
}
$$
$K$ は上に有界なので $M \in D^-(R)$ が従う。
従って、例えば Derived Categories, Lemma
\ref{derived-lemma-subcategory-left-resolution} により、自由 $R$-加群からなる
上に有界な複体 $M^\bullet$ で $M$ を表せる。
前と同様に
$\eta(1) = \sum\nolimits_n \sum\nolimits_i m_{n, i} \otimes n_{-n, i}$
と書く。Lemma
\ref{more-algebra-lemma-complex-bounded-above-free-colim-bounded-finite-free} により、
すべての元 $m_{n, i}$ を含み、有界で有限自由 $R$-加群からなる部分複体
$K^\bullet \subset M^\bullet$ を見つけられる。
上と同様に $M$ 上の恒等写像は $K$ を経由する。
$K$ は完全なので、Lemma \ref{more-algebra-lemma-summands-perfect} により
$M$ も完全である。
\end{proof}

\begin{lemma}
\label{more-algebra-lemma-invertible-derived}
$R$ を環とし、$M$ を $D(R)$ の対象とする。次は同値である。
\begin{enumerate}
\item Categories, Definition \ref{categories-definition-invertible} の意味で
$M$ は $D(R)$ において可逆である。
\item すべての素イデアル $\mathfrak p \subset R$ に対し、
$f \not \in \mathfrak p$ で、ある $n \in \mathbf{Z}$ に対して
$M_f \cong R_f[-n]$ となる $f \in R$ が存在する。
\end{enumerate}
さらにこの場合、次が成り立つ。
\begin{enumerate}
\item[(a)] $M$ は $D(R)$ の完全対象である。
\item[(b)] $D(R)$ において $M = \bigoplus H^n(M)[-n]$ である。
\item[(c)] 各 $H^n(M)$ は有限射影 $R$-加群である。
\item[(d)] $H^n(M)$ が可逆 $R_n$-加群に対応するように
$R = \prod_{a \leq n \leq b} R_n$ と書ける。
\end{enumerate}
\end{lemma}

\begin{proof}
(2) を仮定する。対象 $R\Hom_R(M, R)$ と合成写像
$$
R\Hom(M, R) \otimes_R^\mathbf{L} M \to R
$$
を考える。局所的に調べると、これは同型であることが分かる。
詳細は省略する。$D(R)$ は対称モノイド圏なので、$M$ は可逆である。

\medskip\noindent
(1) を仮定する。モノイド圏の可逆対象は、その逆対象を左双対として
もつことに注意する。従って Lemma \ref{more-algebra-lemma-have-dual-derived} により
$M$ は完全である。素イデアル $\mathfrak p \subset R$ を取り、
その剰余体を $\kappa$ とする。このとき
$M \otimes_R^\mathbf{L} \kappa$ は $D(\kappa)$ の可逆対象である。
これは明らかに、$\dim H^i(M \otimes_R^\mathbf{L} \kappa)$ が
ちょうど一つの $i$ に対して零でなく、その場合には $1$ に等しいことを
意味する。Lemma \ref{more-algebra-lemma-lift-perfect-from-residue-field} により (2) を得る。

\medskip\noindent
上の証明で (a) が成り立つことを見た。
$M_f \cong R_f[-n]$ となる形の開集合 $D(f)$ の和集合を
$U_n \subset \Spec(R)$ とする。$n \not = n'$ ならば、明らかに
$U_n \cap U_{n'} = \emptyset$ である。$M$ が $[a, b]$ に Tor 振幅を
もつならば、$n \not \in [a, b]$ に対し $U_n = \emptyset$ である。
従って (d) のような積分解
$R = \prod_{a \leq n \leq b} R_n$ があり、$U_n$ は
$\Spec(R_n)$ に対応する。Algebra, Lemma
\ref{algebra-lemma-disjoint-implies-product} を見よ。
$D(R) = \prod_{a \leq n \leq b} D(R_n)$ であり、加群の圏についても
同様なので、(b), (c), (d) は直ちに従う。
\end{proof}





\section{自由加群を分離すること}
\label{more-algebra-section-serre-splitting}

\noindent
この節の議論は Serre による。\cite{Serre-projective} を見よ。

\begin{situation}
\label{more-algebra-situation-splitting}
ここで $R$ は環、$M$ は有限表示 $R$-加群である。
閉点全体の集合に誘導位相を入れたものを
$\Omega \subset \Spec(R)$ と書く。$x \in \Omega$ に対し、
$x$ における $M$ のファイバーを $M(x) = M/xM$ と書く。
これは $x$ における剰余体 $\kappa(x)$ 上の有限次元ベクトル空間である。
$s \in M$ が与えられたとき、$M(x)$ における $s$ の像を $s(x)$ と書く。
\end{situation}

\begin{lemma}
\label{more-algebra-lemma-which-elements-split}
Situation \ref{more-algebra-situation-splitting} において $x \in \Omega$ とする。
$\kappa(x)$-ベクトル空間の標準的な短完全列
$$
0 \to B(x) \to M(x) \to V(x) \to 0
$$
で次の性質をもつものが存在する。
% P02-E1356: source line 38737 omits “has” in “which the following property”.
$s_1, \ldots, s_r \in M$ に対し、次は同値である。
\begin{enumerate}
\item $f \not \in x$ となる $f \in R$ で、$f$ を逆にすると写像
$s_1, \ldots, s_r : R^{\oplus r} \to M$ が直和因子の包含になるものが存在する。
\item $s_1(x), \ldots, s_r(x)$ は $V(x)$ の線形独立な元へ写る。
\end{enumerate}
\end{lemma}

\begin{proof}
写像
$$
\Hom_R(M, R) \to \Hom_{\kappa(x)}(M(x), \kappa(x))
$$
の像の直交補空間を $B(x) \subset M(x)$ と定義し、
$V(x) = M(x)/B(x)$ と置く。このとき任意の $R$-線形写像
$\varphi : M \to R$ は写像
$\overline{\varphi} : V(x) \to \kappa(x)$ を誘導し、逆に任意の
$\kappa(x)$-線形写像 $\lambda : V(x) \to \kappa(x)$ は、ある
$\varphi$ に対して $\overline{\varphi}$ に等しい。
$s_1, \ldots, s_r \in M$ とする。

\medskip\noindent
$s_1, \ldots, s_r$ が $V(x)$ の線形独立な元へ写ると仮定する。
$\varphi_i(s_j)$ が $\kappa(x)$ における
$\delta_{ij}$\footnote{Kronecker のデルタ。} へ写るような
$\varphi_1, \ldots, \varphi_r \in \Hom_R(M, R)$ を見つけられる。
従って合成
$$
R^{\oplus r}
\xrightarrow{s_1, \ldots, s_r}
M
\xrightarrow{\varphi_1, \ldots, \varphi_r}
R^{\oplus r}
$$
の行列は、$\kappa(x)$ で $1$ へ写る行列式 $f \in R$ をもつ。
% P02-E1357: source line 38770 lacks terminal punctuation after kappa(x).
これは明らかに、$f$ を逆にすると
$s_1, \ldots, s_r : R^{\oplus r} \to M$ が直和因子の包含になることを
意味する。

\medskip\noindent
逆に、$f \not \in x$ となる $f \in R$ で、$f$ を逆にすると
$s_1, \ldots, s_r : R^{\oplus r} \to M$ が直和因子の包含になるものが
あると仮定する。従って
$\varphi_i(s_j) = \delta_{ij} \in R_f$ となる $R_f$-線形写像
$\varphi_i : M_f \to R_f$ を見つけられる。
Algebra, Lemma \ref{algebra-lemma-hom-from-finitely-presented} により
$\Hom_R(M, R)_f = \Hom_{R_f}(M_f, R_f)$ なので、
$\varphi'_i(s_j) = f^n\delta_{ij} \in R$ となる $n \geq 0$ と
$\varphi'_i \in \Hom_R(M, R)$ を見つけられる。
$\overline{\varphi}'_i(s_j) = f^n\delta_{ij}$ だから、
$s_1, \ldots, s_r$ は $V(x)$ の線形独立な元へ写る。
\end{proof}

\noindent
Situation \ref{more-algebra-situation-splitting} において
$s_1, \ldots, s_r \in M$ が与えられたとき、
$s_1(x), \ldots, s_r(x)$ が $V(x)$ の線形従属な元へ写るような
$x \in \Omega$ の集合を $Z(s_1, \ldots, s_r) \subset \Omega$ と書く。
補題により、これは $\Omega$ の閉部分集合である。

\begin{lemma}
\label{more-algebra-lemma-choose-values}
Situation \ref{more-algebra-situation-splitting} において、
$x_1, \ldots, x_n \in \Omega$ は相異なるものとし、
$v_i \in V(x_i)$ とする。このとき $i = 1, \ldots, n$ に対し
$s(x_i)$ が $v_i$ へ写るような $s \in M$ が存在する。
\end{lemma}

\begin{proof}
$x_i$ は $R$ の極大イデアルなので、Algebra, Lemma
\ref{algebra-lemma-chinese-remainder} により
$M(x_1) \oplus \ldots \oplus M(x_n)$ は $M$ の商である。
\end{proof}

\begin{proposition}
\label{more-algebra-proposition-splitting}
\begin{reference}
\cite[Theorem 2]{Serre-projective}
\end{reference}
Situation \ref{more-algebra-situation-splitting} において、$\Omega$ は Noether 位相空間と
仮定する。$s_1, \ldots, s_h \in M$ とする。
$Z(s_1, \ldots, s_h) \subset F \subset \Omega$ は閉であるとし、
$x_1, \ldots, x_n \in F$ は相異なるものとする。
$v_i  \in V(x_i)$ とする。また
$$
(*)\quad h + k \leq \dim_{\kappa(x)} V(x)\text{ for all }x \in \Omega
$$
を満たす整数 $k \geq 0$ を取る。このとき、$s \in M$ と閉集合
$F' \subset \Omega$ で次を満たすものが存在する。
\begin{enumerate}
\item[(a)] $s(x_i)$ は $v_i$ へ写る。
\item[(b)] $Z(s_1, \ldots, s_h, s) \subset F \cup F'$ である。
\item[(c)] $F'$ の各既約成分は $\Omega$ において余次元 $\geq k$ をもつ。
\end{enumerate}
\end{proposition}

\begin{proof}
余次元は Topology, Section \ref{topology-section-catenary-spaces} で
定義されたこと、および Topology, Section
\ref{topology-section-noetherian} に含まれる Noether 位相空間に関する
いくつかの結果を用いることに注意する。

\medskip\noindent
$k$ に関する帰納法で証明する。$k = 0$ ならば、Lemma
\ref{more-algebra-lemma-choose-values} のような $s \in M$ を選び、
$F' = \Omega$ と取る。

\medskip\noindent
$k > 0$ と仮定する。帰納法の仮定により、
$s_1, \ldots, s_h$, $F$, $x_1, \ldots, x_n$,
$v_1, \ldots, v_n$, $k - 1$ に対して (a), (b), (c) を満たす
$u \in M$ と閉集合 $G \subset \Omega$ を選べる。

\medskip\noindent
$G = G_1 \cup \ldots \cup G_m$ を $G$ の既約成分への分解とする。
$G_j \subset F$ ならば、それをリストから除ける。
従って $j = 1, \ldots, m$ に対し $G_j$ は $F$ に含まれないと
仮定してよい。$j = 1, \ldots, m$ に対し
$y_j \not \in F$ かつ $j' \not = j$ に対し
$y_j \not \in G_{j'}$ となる $y_j \in G_j$ を選ぶ。
$G$ の既約成分間には包含関係がないので、これは可能である。
$s_1(y_j), \ldots, s_h(y_j)$ の像が張る部分空間に含まれない
$w_j \in V(y_j)$ を選ぶ。$h + k \leq \dim V(y_j)$ かつ
$k > 0$ なので、これは可能である。

\medskip\noindent
$h + 1$ 個の切断 $s_1, \ldots, s_h, u$、閉集合 $F \cup G$、
点 $x_1, \ldots, x_n, y_1, \ldots, y_m \in F \cup G$、
元 $0 \in V(x_i)$ および $w_j \in V(y_j)$、整数 $k - 1$ に対して
帰納法の仮定を適用する。$h$ を $1$ 増やし $k$ を $1$ 減らしたので、
命題の仮定 $(*)$ は依然として成り立つことに注意する。
これにより $t \in M$ と閉集合 $H \subset \Omega$ が得られ、
これらは $s_1, \ldots, s_h, u$, $F \cup G$,
$x_1, \ldots, x_n, y_1, \ldots, y_m$,
$0, \ldots, 0, w_1, \ldots, w_m$, $k - 1$ に対し
(a), (b), (c) を満たす。

\medskip\noindent
$F \cup G$ に含まれない $H$ の既約成分を
$H_1, \ldots, H_p \subset H$ とする。前と同様に、
$z_l \not \in F \cup G$ かつ $l' \not = l$ に対し
$z_l \not \in H_{l'}$ となる $z_l \in H_l$ を選ぶ。
Algebra, Lemma \ref{algebra-lemma-chinese-remainder} により
$f(y_j) = 1$, $j = 1, \ldots, m$ かつ
$f(z_l) = 0$, $l = 1, \ldots, p$ となる $f \in R$ を選べる。
元 $s = u + f t$ が条件を満たすと主張する。

\medskip\noindent
まず、
% P02-E1358: source line 38884 asserts t(x_i)=0, while induction only says its image in V(x_i) is 0; that weaker statement is sufficient.
$t(x_i) = 0$ なので $s(x_i)$ の値は $u(x_i)$ と一致し、従って
$s(x_i)$ は $v_i$ へ写る。これで (a) が示された。
証明を終えるには、$F$ に含まれない
$Z(s_1, \ldots, s_h, s)$ の各既約成分 $Z$ が $\Omega$ において
余次元 $\geq k$ をもつことを示せば十分である。
実際、そのような成分の和集合を $F'$ とすれば (b), (c) を得る。
余次元 $\leq k - 1$ の $Z(s_1, \ldots, s_h, s)$ の既約成分 $Z$ が
存在しないことを次のように示せる。
\begin{enumerate}
\item $s = u + ft$ なので
% P02-E1359: source line 38894 incorrectly states equality with F union H and omits G; induction gives only containment in F union G union H, as the later G_j case itself confirms.
$Z(s_1, \ldots, s_h, s) \subset Z(s_1, \ldots, s_h, u, t) = F \cup H$
であることに注意する。従って $Z \subset H$ である。
\item $H$ の既約成分は余次元 $\geq k - 1$ をもつ。
従って $Z$ は $H$ の既約成分に等しい。実際、$Z$ は
余次元 $\leq k - 1$ をもつからである。
ゆえに、ある $l \in \{1, \ldots, p\}$ に対して $Z = H_l$、または
% P02-E1360: source line 38900 omits the comma in the range {1, ..., m}.
ある $j \in \{1, \ldots m\}$ に対して $Z = G_j$ である。
\item しかし $Z = G_j$ は不可能である。実際
$s_1(y_j), \ldots, s_h(y_j)$ は $V(y_j)$ の線形独立な元へ写り、
$s(y_j) = u(y_j) + f(y_j) t(y_j) = u(y_j) + t(y_j)$ は形
$$
\text{linear combination images of }s_i(y_j) + w_j
$$
の元へ写る。これは $w_j$ の選び方により、$V(y_j)$ における
$s_1(y_j), \ldots, s_h(y_j)$ の像と線形独立である。
% P02-E1361: source line 38906 has the ungrammatical displayed phrase “linear combination images of ...”; it is missing “of the”.
\item $Z = H_l$ も不可能である。実際、再び
$s_1(z_l), \ldots, s_h(z_l)$ は $V(z_l)$ の線形独立な元へ写り、
$s(z_l) = u(z_l) + f(z_l) t(z_l) = u(z_l)$ は
$z_l \not \in F \cup G$ なので、それらと線形独立な $V(z_l)$ の元へ写る。
\end{enumerate}
これで証明が完了する。
\end{proof}

\begin{theorem}
\label{more-algebra-theorem-splitting}
\begin{reference}
\cite[Theorem 1]{Serre-projective}
\end{reference}
極大スペクトル $\Omega \subset \Spec(R)$ が次元 $d < \infty$ の
Noether 位相空間であるような環を $R$ とする。
$M$ を有限表示 $R$-加群とし、すべての $\mathfrak m \in \Omega$ に対し
$R_\mathfrak m$-加群 $M_\mathfrak m$ が階数 $> d$ の自由直和因子を
もつとする。このとき $M \cong R \oplus M'$ である。
\end{theorem}

\begin{proof}
$\mathfrak m \in \Omega$ に対し、$R_\mathfrak m^{\oplus r}$ が
$M_\mathfrak m$ の直和因子であると仮定する。Algebra, Lemmas
\ref{algebra-lemma-localization-colimit} および
\ref{algebra-lemma-colimit-category-fp-modules} により、
ある $f \not \in \mathfrak m$ となる $f \in R$ に対して
$R_f^{\oplus r}$ は $M_f$ の直和因子である。
従ってこの仮定は、Lemma \ref{more-algebra-lemma-which-elements-split} のような
$V(x)$ に対し、すべての $x \in \Omega$ について
$\dim V(x) > d$ が成り立つことを意味する。
Proposition \ref{more-algebra-proposition-splitting} を $F = \emptyset$,
$h = 0$ かつ $s_i$ なし、$n = 0$ かつ $x_i, v_i$ なし、
$k = d + 1$ として適用すると、$s \in M$ と
$F' \subset \Omega$ で、$F'$ の各既約成分が余次元 $\geq d + 1$ をもち、
$Z(s) \subset F'$ となるものを得る。$d = \dim(\Omega)$ なので
$F' = \emptyset$ が強制される。従って $s : R \to M$ は
すべての極大イデアルにおいて直和因子の包含である。
Algebra, Lemma
\ref{algebra-lemma-universally-injective-check-stalks} により、
$s$ は普遍単射である。すると Algebra, Lemma
\ref{algebra-lemma-universally-exact-split} により $s$ は分裂単射である。
\end{proof}








\section{大きな射影加群は自由である}
\label{more-algebra-section-big-projective-free}

\noindent
この節では \cite{Bass} の結果の一つを論じる。読者には原論文を
参照することを勧める。ここでの議論には、論文 \cite{Akasaki} と
\cite{Hinohara} で与えられた、やや簡略化された証明を用いる。

\begin{lemma}[Eilenberg の補題]
\label{more-algebra-lemma-eilenberg-swindle}
\begin{slogan}
Eilenberg のトリック
\end{slogan}
\begin{history}
\cite{Bass} には次の記述がある。「これは、数年前に Eilenberg が見いだした
優美で小さなトリックであり、Barry Mazur の額から生まれ出てもおかしくない。」
\end{history}
\begin{reference}
\cite[Eilenberg's lemma]{Bass}
\end{reference}
$P \oplus Q \cong F$ で $F$ が有限生成でない自由加群ならば、
$P \oplus F \cong F$ である。
\end{lemma}

\begin{proof}
$$
F \cong F \oplus F \oplus \ldots \cong
P \oplus Q \oplus P \oplus Q \oplus \ldots \cong
P \oplus F \oplus F \oplus \ldots \cong P \oplus F
$$
\end{proof}

\begin{lemma}
\label{more-algebra-lemma-projective-plus-free-is-free}
$R$ を環、$P$ を射影加群とする。
$P \oplus F$ が自由となる自由加群 $F$ が存在する。
\end{lemma}

\begin{proof}
$P$ は射影的なので、ある加群 $Q$ に対し
$F_0 = P \oplus Q$ は自由加群である。
$F = \bigoplus_{n \geq 1} F_0$ と置く。Lemma
\ref{more-algebra-lemma-eilenberg-swindle} により $P \oplus F \cong F$ である。
\end{proof}

\begin{lemma}
\label{more-algebra-lemma-element-projective}
$R$ を環、$P$ を射影加群、$s \in P$ とする。
有限自由加群 $F$ と、$(0, s) \in K$ を満たす有限自由直和因子
$K \subset F \oplus P$ が存在する。
\end{lemma}

\begin{proof}
Lemma \ref{more-algebra-lemma-projective-plus-free-is-free} により、
$F \oplus P$ が自由となる（無限でもよい）自由加群 $F$ を見つけられる。
するともちろん $(0, s)$ は有限自由直和因子
$K \subset F \oplus P$ に含まれる。さらに $K$ は、有限自由直和因子
$F' \subset F$ に対する $F' \oplus P$ に含まれる。
\end{proof}

\begin{lemma}
\label{more-algebra-lemma-trick-to-find-good-element}
$R$ を、Jacobson 根基 $J$ をもち $R/J$ が Noether である環とする。
$P$ を、$R$ のすべての極大イデアル $\mathfrak m$ に対し
$P_\mathfrak m$ が無限階数をもつ射影 $R$-加群とする。
$s \in P$ および $Rs + M = P$ を満たす $M \subset P$ を取る。
このとき $R(s + m)$ が $P$ の自由直和因子となるような
$m \in M$ を見つけられる。
\end{lemma}

\begin{proof}
Algebra, Theorem
\ref{algebra-theorem-projective-free-over-local-ring} により
$P_\mathfrak m$ は自由なので、この主張には意味がある。

\medskip\noindent
$M$ の像を $M' \subset P/JP$、
$s$ の像を $s' \in P/JP$ と書く。
% P02-E1362: source lines 39041 and 39043 write R/J s' and R/J(s'+m') without parentheses around the scalar ring R/J, making the cyclic-submodule notation parse as a quotient.
$R/J s' + M' = P/JP$ であることに注意する。
$R/J(s' + m')$ が $M'$ の自由直和因子となるような
% P02-E1363: source line 39043 says the cyclic module generated by s'+m' is a direct summand of M'; it need not lie in M', and the following splitting map shows the intended ambient module is P/JP.
$m' \in M'$ を見つけられると仮定する。
分裂を与える $\varphi' : P/JP \to R/J$、すなわち
$R/J$ において $\varphi'(s' + m') = 1$ となるものを選ぶ。
$P$ は射影 $R$-加群なので、$\varphi'$ の持ち上げ
$\varphi : P \to R$ を見つけられる。$m'$ へ写る $m \in M$ を選ぶ。
すると $\varphi(s + m) \in R$ は $J$ を法として $1$ と合同であり、
従って $R$ の単元である（Algebra, Lemma
\ref{algebra-lemma-contained-in-radical}）。ゆえに $R(s + m)$ は
$P$ の自由直和因子である。これにより次の段落で論じる場合へ帰着される。

\medskip\noindent
$R$ は Noether であると仮定する。元 $m \in M$ と $R$-線形写像
$\varphi_1, \ldots, \varphi_n : P \to R$ を取る。
$\varphi_1(s + m), \ldots, \varphi_n(s + m) \in R$ の消滅集合を
$$
Z(s + m, \varphi_1, \ldots, \varphi_n) \subset \Spec(R)
$$
と書く。

\medskip\noindent
$\mathfrak m$ を $R$ の極大イデアルとし、
$\mathfrak m \in Z(s + m, \varphi_1, \ldots, \varphi_n)$ と仮定する。
$K = M \cap \bigcap \Ker(\varphi_i)$ と置く。写像
$$
K/\mathfrak mK \to P/\mathfrak m P
$$
の像は無限次元であると主張する。実際、$P/K$ は
$P/M \oplus R^{\oplus n}$ の部分加群に同型なので有限 $R$-加群である。
従って、表示された矢印の核は
$\text{Tor}_1^R(P/K, \kappa(\mathfrak m))$ の商であり、これは
Algebra, Lemma \ref{algebra-lemma-tor-noetherian} により有限である。
$P/\mathfrak mP$ が無限次元であることと合わせて主張を得る。
従って、$P/\mathfrak mP$ において $s + m$ の像
$\overline{s + m}$ と線形独立な元 $\overline{t}$ へ写る
$t \in K$ を見つけられる。線形代数により、
$\overline{\varphi}(\overline{t}) = 1$ かつ
$\overline{\varphi}(\overline{s + m}) = 0$ となる $R$-線形写像
$\overline{\varphi} : P \to \kappa(\mathfrak m)$ が存在する。
$P$ は射影的なので、$\overline{\varphi}$ を持ち上げる $R$-線形写像
$\varphi : P \to R$ を見つけられる。すると消滅集合
$Z(s + m + t, \varphi_1, \ldots, \varphi_n, \varphi)$ は
$Z(s + m, \varphi_1, \ldots, \varphi_n)$ に含まれるが
$\mathfrak m$ を含まず、すなわち
$Z(s + m, \varphi_1, \ldots, \varphi_n)$ より真に小さい。

\medskip\noindent
$\Spec(R)$ は Noether 位相空間なので、上の議論から、閉部分集合
$Z(s + m, \varphi_1, \ldots, \varphi_n)$ が $\Spec(R)$ の閉点を
一つも含まないような $m \in M$ と
$\varphi_1, \ldots, \varphi_n : P \to R$ を見つけられる。
従って $Z(s + m, \varphi_1, \ldots, \varphi_n) = \emptyset$ である。
従って $\sum r_i\varphi_i(s + m) = 1$ となる
$r_1, \ldots, r_n \in R$ を見つけられる。ゆえに
$$
R \xrightarrow{s + m} P \xrightarrow{\sum r_i \varphi_i} R
$$
が所望の分裂である。
\end{proof}

\begin{lemma}
\label{more-algebra-lemma-element-in-free-summand}
$R$ を、Jacobson 根基 $J$ をもち $R/J$ が Noether である環とする。
$P$ を、$R$ のすべての極大イデアル $\mathfrak m$ に対し
$P_\mathfrak m$ が無限階数をもつ射影 $R$-加群とする。
$s \in P$ とする。このとき $s \in M$ となる有限安定自由直和因子
$M \subset P$ を見つけられる。
\end{lemma}

\begin{proof}
Lemma \ref{more-algebra-lemma-element-projective} により、有限自由加群 $F$ と、
$(0, s) \in K$ となる有限自由直和因子 $K \subset F \oplus P$ を
見つけられる。$F$ の階数に関する帰納法により、次の段落で論じる場合へ
帰着される。

\medskip\noindent
$(0, s) \in K$ となる有限安定自由直和因子
$K \subset R \oplus P$ が存在すると仮定する。
$K$ の補加群 $K'$、すなわち $R \oplus P = K \oplus K'$ となるものを選ぶ。
射影 $\pi : R \oplus P \to K'$ は全射なので、
% P02-E1364: source line 39126 misspells “surjective” as “surjectve”.
Lemma \ref{more-algebra-lemma-trick-to-find-good-element} により、
$\pi(1, p) \in K'$ が自由直和因子を生成するような $p \in P$ を
見つけられる。従って $K' = R\pi(1, p) \oplus K''$ と書く。このとき
$$
R \oplus P = K \oplus K' = K \oplus R\pi(1, p) \oplus K''
$$
である。射影 $\pi' : P \to K''$ は全射である\footnote{実際、
$k'' \in K''$ ならば、$K'$ の元と見た $k''$ は、ある
$\lambda \in R$ と $q \in P$ によって
$k'' = \lambda \pi(1, 0) + \pi(0, q)$ と書ける。
これは $k'' = \lambda \pi(1, p) + \pi(0, q - \lambda p)$ を意味する。
さらに、$K' \to K''$ は $\pi(1, p)$ を零にするので、
$q - \lambda p$ は合成
$P \to R \oplus P \xrightarrow{\pi} K' \to K''$ によって
$k''$ へ写ることを意味する。}。
% P02-E1365: source line 39142 has the agreement error “and hence split” rather than “and hence splits”.
$K''$ は射影的なので、この全射は分裂する。従って
$\Ker(\pi') \subset P$ は $s$ を含む直和因子である。
最後に、構成により同型
$$
R \oplus \Ker(\pi') \cong K \oplus R\pi(1, p)
$$
がある。従って $K$ が有限安定自由なので、$\Ker(\pi')$ もそうである。
\end{proof}

\begin{theorem}
\label{more-algebra-theorem-countable-free}
\begin{reference}
\cite[Theorem 3.1]{Bass} の可換な場合
\end{reference}
$R$ を、Jacobson 根基 $J$ をもち $R/J$ が Noether である環とする。
$P$ を可算生成射影 $R$-加群とし、$R$ のすべての極大イデアル
$\mathfrak m$ に対し $P_\mathfrak m$ は無限階数をもつとする。
このとき $P$ は自由である。
\end{theorem}

\begin{proof}
まず、$P$ が有限安定自由加群の可算直和であることを証明する。
$x_1, x_2, \ldots$ を $P$ の可算生成系とする。
すべての $n$ に対し $F_1 \oplus \ldots \oplus F_n$ が
$x_1, \ldots, x_n$ を含む $P$ の直和因子となるような、
$P$ の有限安定自由直和因子 $F_1, F_2, \ldots$ を帰納的に構成する。
実際、所望の性質をもつ $F_1, \ldots, F_n$ が与えられたら
$$
P = F_1 \oplus \ldots \oplus F_n \oplus P'
$$
と書き、その成分における $x_{n + 1}$ の像を $s \in P'$ とする。
Lemma \ref{more-algebra-lemma-element-in-free-summand} により、$s$ を含む有限安定自由
直和因子 $F_{n + 1} \subset P'$ を見つけられる。
すると $P = \bigoplus_{i = 1}^{\infty} F_i$ である。

\medskip\noindent
$P$ は零でない有限安定自由加群の無限直和
$P = \bigoplus_{i = 1}^{\infty} F_i$ であると仮定する。
加群 $F_i$ の安定自由性を次のように用いる。各 $F_i$ の階数は
一定で（かつ正で）ある。従って $R$ の各極大イデアル
$\mathfrak m$ に対し、$P_\mathfrak m$ は可算無限階数の自由加群である。
Lemma \ref{more-algebra-lemma-trick-to-find-good-element} を $s = 0$, $M = P$ として
適用すると、$Rt_1$ が $P$ の自由直和因子となる $t_1 \in P$ を
見つけられる。このとき、ある $n_1 > n_0 = 0$ に対して
$t_1$ は $F_1 \oplus \ldots \oplus F_{n_1}$ に含まれる。
$\bigoplus_{n > n_1} F_n$ に同じ議論を適用すると、
$n_1 < n_2$ と、自由直和因子を生成する
$t_2 \in F_{n_1 + 1} \oplus \ldots \oplus F_{n_2}$ が得られる。
このように続けると、無限階数の自由直和因子
$$
\bigoplus\nolimits_{i \geq 1} t_i :
\bigoplus\nolimits_{i \geq 1} R
\longrightarrow
\bigoplus\nolimits_{i \geq 1}
\bigoplus\nolimits_{n_i \geq n > n_{i - 1}} F_n = P
$$
を得る。従って、可算階数のある自由 $R$-加群 $F$ とある補加群に対し
$P \cong Q \oplus F$ である。$Q$ は可算生成なので、ある加群 $Q'$ に対し
$Q \oplus Q' \cong F$ となる。すると Eilenberg のトリック
（Lemma \ref{more-algebra-lemma-eilenberg-swindle}）により
$Q \oplus F \cong F$ であり、$P$ は自由である。
\end{proof}
